


\section*{CONTENTS}
Foreword ..... iii\\
Preface to the second edition ..... v\\
Preface to the first edition ..... vi\\
PART ONE\\
PHYSICAL FUNDAMENTALS OF MECHANICS\\
1.1 Kinematics ..... 1-34\\
1.2 The Fundamental Equation of Dynamics ..... 35-65\\
1.3 Laws of Conservation of Energy, Momemtum, and Angular Momentum ..... 66-101\\
1.4 Universal Gravitation ..... 102-117\\
1.5 Dynamics of a Solid Body ..... 118-143\\
1.6 Elastic Deformations of a Solid Body ..... 144-155\\
1.7 Hydrodynamics ..... 156-167\\
1.8 Relativistic Mechanics ..... 168-183\\
PART TWO\\
THERMODYNAMICS AND MOLECULAR PHYSICS\\
2.1 Equation of the Gas State. Processes ..... 184-195\\
2.2 The first Law of Thermodynamics. Heat Capacity ..... 196-212\\
2.3 Kinetic theory of Gases. Boltzmann's Law and Maxwell's Distribution ..... 213-226\\
2.4 The Second Law of Thermodynamics. Entropy ..... 227-241\\
2.5 Liquids. Capillary Effects ..... 242-247\\
2.6 Phase Transformations ..... 248-256\\
2.7 Transport Phenomena ..... 257-266\\
PART THREE ELECTRODYNAMICS\\
3.1 Constant Electric Field in Vacuum ..... 267-288\\
3.2 Conductors and Dielectrics in an Electric Field ..... 289-305\\
3.3 Electric Capacitance. Energy of an Electric Field ..... 306-324\\
3.4 Electric Current ..... 325-353\\
3.5 Constant Magnetic Field. Magnetics ..... 354-379\\
3.6 Electromagnetic Induction. Maxwell's Equations ..... 380-407\\
3.7 Motion of Charged Particles in Electric and Magnetic Fields ..... 408-424

\section*{PART ONE}
\section*{PHYSICAL FUNDAMENTALS OF MECHANICS}
\subsection*{1.1 KINEMATICS}
1.1 Let \(v_{0}\) be the stream velocity and \(v^{\prime}\) the velocity of motorboat with respect to water. The motorboat reached point \(B\) while going downstream with velocity ( \(v_{0}+v^{\prime}\) ) and then returned with velocity ( \(v^{\prime}-v_{0}\) ) and passed the raft at point \(C\). Let \(t\) be the time for the raft (which flows with stream with velocity \(v_{\mathrm{o}}\) ) to move from point \(A\) to \(C\), during which the motorboat moves from \(A\) to \(B\) and then from \(B\) to \(C\).\\
Therefore\\
\(\frac{l}{v_{0}}=\tau+\frac{\left(v_{0}+v^{\prime}\right) \tau-l}{\left(v^{\prime}-v_{0}\right)}\)\\
On solving we get \(\nu_{0}=\frac{l}{2 \tau}\)\\
\includegraphics[max width=\textwidth, alt={}, center]{9f3c5a8e-ac5d-42f4-8eb6-ea2af623f2d8-008_179_690_1012_622}\\
1.2 Let \(s\) be the total distance traversed by the point and \(t_{1}\) the time taken to cover half the distance. Further let \(2 t\) be the time to cover the rest half of the distance.

Therefore


\begin{equation*}
\frac{s}{2}=v_{0} t_{1} \quad \text { or } \quad t_{1}=\frac{s}{2 v_{0}} \tag{1}
\end{equation*}


and


\begin{equation*}
\frac{s}{2}=\left(v_{1}+v_{2}\right) t \quad \text { or } \quad 2 t=\frac{s}{v_{1}+v_{2}} \tag{2}
\end{equation*}


Hence the sought average velocity

\[
<v>=\frac{s}{t_{1}+2 t}=\frac{s}{\left[s / 2 v_{0}\right]+\left[s /\left(v_{1}+v_{2}\right)\right]}=\frac{2 v_{0}\left(v_{1}+v_{2}\right)}{v_{1}+v_{2}+2 v_{0}}
\]

1.3 As the car starts from rest and finally comes to a stop, and the rate of acceleration and deceleration are equal, the distances as well as the times taken are same in these phases of motion.\\
Let \(\Delta t\) be the time for which the car moves uniformly. Then the acceleration / deceleration time is \(\frac{\tau-\Delta t}{2}\) each. So,

\[
<v>\tau=2\left\{\frac{1}{2} w \frac{(\tau-\Delta t) \cdot}{4}\right\}+w, \frac{(\tau-\Delta t)}{2} \Delta t
\]

or

\[
\Delta t^{2}=\tau^{2}-\frac{4<v>\tau}{w}
\]

Hence

\[
\Delta t=\tau \sqrt{1-\frac{4<v>}{w \tau}}=15 \mathrm{~s} .
\]

1.4 (a) Sought average velocity\\
\(\langle\nu\rangle=\frac{s}{t}=\frac{200 \mathrm{~cm}}{20 \mathrm{~s}}=10 \mathrm{~cm} / \mathrm{s}\)\\
(b) For the maximum velocity, \(\frac{d s}{d t}\) should be maximum. From the figure \(\frac{d s}{d t}\) is maximum for all points on the line ac, thus the sought maximum velocity becomes average velocity for the line \(a c\) and is equal to :\\
\includegraphics[max width=\textwidth, alt={}, center]{9f3c5a8e-ac5d-42f4-8eb6-ea2af623f2d8-009_443_505_421_851}

\[
\frac{b c}{a b}=\frac{100 \mathrm{~cm}}{4 \mathrm{~s}}=25 \mathrm{~cm} / \mathrm{s}
\]

(c) Time \(t_{0}\) should be such that corresponding to it the slope \(\frac{d s}{d t}\) should pass through the point \(O\) (origin), to satisfy the relationship \(\frac{d s}{d t}=\frac{s}{t_{0}}\). From figure the tangent at point \(d\) passes through the origin and thus corresponding time \(t=t_{0}=16 \mathrm{~s}\).\\
1.5 .Let the particles collide at the point \(A\) (Fig.), whose position vector is \(\overrightarrow{r_{3}}\) (say). If \(t\) be the time taken by each particle to reach at point \(A\), from triangle law of vector addition :


\begin{align*}
& \begin{aligned}
& \overrightarrow{r_{3}}=\overrightarrow{r_{1}}+\overrightarrow{v_{1}} t=\overrightarrow{r_{2}}+\overrightarrow{v_{2}} \\
& \text { so, } \quad \overrightarrow{r_{1}}-\overrightarrow{r_{2}}=\left(\overrightarrow{v_{2}}-\overrightarrow{v_{1}}\right) t
\end{aligned} \\
& \text { therefore, } \quad t=\frac{\left|\overrightarrow{r_{1}}-\overrightarrow{r_{2}}\right|}{\left|\overrightarrow{v_{2}}-\overrightarrow{v_{1}}\right|} \tag{1}
\end{align*}


From Eqs. (1) and (2)

\[
\overrightarrow{r_{1}}=\overrightarrow{r_{2}}-\left(\overrightarrow{v_{2}}-\overrightarrow{v_{1}}\right) \frac{\left|\overrightarrow{r_{1}}-\overrightarrow{r_{2}}\right|}{\left|\overrightarrow{v_{2}}-\overrightarrow{v_{1}}\right|}
\]

\includegraphics[max width=\textwidth, alt={}, center]{9f3c5a8e-ac5d-42f4-8eb6-ea2af623f2d8-009_424_458_1285_893}\\
or, \(\frac{\overrightarrow{r_{1}}-\overrightarrow{r_{2}}}{\left|\overrightarrow{r_{1}}-\overrightarrow{r_{2}}\right|}=\frac{\overrightarrow{v_{2}}-\overrightarrow{v_{1}}}{\left|\overrightarrow{v_{2}}-\overrightarrow{v_{1}}\right|}\), which is the sought relationship.\\
1.6 We have


\begin{equation*}
\vec{v}^{\prime}=\vec{v}-\vec{v}_{0} \tag{1}
\end{equation*}


From the vector diagram [of Eq. (1)] and using properties of triangle


\begin{equation*}
v^{\prime}=\sqrt{v_{0}^{2}+v^{2}+2 v_{0} v \cos \varphi}=39.7 \mathrm{~km} / \mathrm{hr} \tag{2}
\end{equation*}


and

\[
\frac{\nu^{\prime}}{\sin (\pi-\varphi)}=\frac{\nu}{\sin \theta} \quad \text { or, } \quad \sin \theta=\frac{\nu \sin \varphi}{v^{\prime}}
\]

or

\[
\theta=\sin ^{-1}\left(\frac{\nu \sin \varphi}{\nu^{\prime}}\right)
\]

\begin{center}
\includegraphics[max width=\textwidth, alt={}]{9f3c5a8e-ac5d-42f4-8eb6-ea2af623f2d8-010_320_465_142_886}
\end{center}

Using (2) and putting the values of \(v\) and \(d\)

\[
\theta=19.1^{\circ}
\]

1.7 Let one of the swimmer (say 1 ) cross the river along \(A B\), which is obviously the shortest path. Time taken to cross the river by the swimmer 1.\\
\(t_{1}=\frac{d}{\sqrt{v^{\prime 2}-v_{0}^{2}}},(\) where \(A B=d\) is the width of the river \()\)

For the other swimmer (say 2), which follows the quickest path, the time taken to cross the river.


\begin{equation*}
t_{2}=\frac{d}{v^{\prime}} \tag{2}
\end{equation*}


\includegraphics[max width=\textwidth, alt={}, center]{9f3c5a8e-ac5d-42f4-8eb6-ea2af623f2d8-010_415_531_869_93}\\
\includegraphics[max width=\textwidth, alt={}, center]{9f3c5a8e-ac5d-42f4-8eb6-ea2af623f2d8-010_417_587_868_739}

In the time \(t_{2}\), drifting of the swimmer 2 , becomes


\begin{equation*}
\left.x=v_{0} t_{2}=\frac{v_{0}}{v^{\prime}} d, \text { (using Eq. } 2\right) \tag{3}
\end{equation*}


If \(t_{3}\) be the time for swimmer 2 to walk the distance \(x\) to come from \(C\) to \(B\) (Fig.), then


\begin{equation*}
t_{3}=\frac{x}{u}=\frac{v_{0} d}{v^{\prime} u} \quad \text { (using Eq. 3) } \tag{4}
\end{equation*}


According to the problem \(t_{1}=t_{2}+t_{3}\)\\
or,

\[
\frac{d}{\sqrt{v^{\prime 2}-v_{0}^{2}}}=\frac{d}{v^{\prime}}+\frac{v_{0} d}{v^{\prime} u}
\]

On solving we get

\[
u=\frac{v_{0}}{\left(\frac{1-v_{0}^{2}}{v^{\prime 2}}\right)^{-\frac{1}{2}}-1}=3 \mathrm{~km} / \mathrm{hr} .
\]

1.8 Let \(l\) be the distance covered by the boat \(A\) along the river as well as by the boat \(B\) acre the river. Let \(v_{0}\) be the stream velocity and \(v^{\prime}\) the velocity of each boat with respect water. Therefore time taken by the boat \(A\) in its journey

\[
\begin{aligned}
\text { en by the boat } A \text { in its journey } & \\
t_{A} & =\frac{l}{v^{\prime}+v_{0}}+\frac{l}{v^{\prime}-v_{0}} \\
t_{B} & =\frac{l}{\sqrt{v^{\prime 2}-v_{0}^{2}}}+\frac{l}{\sqrt{v^{\prime 2}-v_{0}^{2}}}=\frac{2 l}{\sqrt{v^{\prime 2}-v_{0}^{2}}} \\
\frac{t_{A}}{t_{B}} & =\frac{v^{\prime}}{\sqrt{v^{\prime 2}-v_{0}^{2}}}=\frac{\eta}{\sqrt{\eta^{2}-1}}\left(\text { where } \eta=\frac{v^{\prime}}{v}\right)
\end{aligned}
\]

and for the boat \(B\)

Hence,

\[
t_{A} / t_{B}=1.8
\]

On substitution\\
1.9 Let \(v_{0}\) be the stream velocity and \(v^{\prime}\) the velocity of boat with respect to water. A \(\frac{\nu_{0}}{\nu^{\prime}}=\eta=2>0\), some drifting of boat is inevitable.\\
Let \(\vec{v}^{\prime}\) make an angle \(\theta\) with flow direction. (Fig.), then the time taken to cross the rive

\[
t=\frac{d}{v^{\prime} \sin \theta} \text { (where } d \text { is the width of the river) }
\]

In this time interval, the drifting of the boat

\[
\begin{aligned}
& \qquad x=\left(v^{\prime} \cos \theta+v_{0}\right) t \\
& =\left(v^{\prime} \cos \theta+v_{0}\right) \frac{d}{v^{\prime} \sin \theta}=(\cot \theta+\eta \operatorname{cosec} \theta) d \rightarrow \\
& \text { For } x_{\text {min }}(\text { minimum drifting }) \\
& \frac{d}{d \theta}(\cot \theta+\eta \operatorname{cosec} \theta)=0 \text {, which yields } \\
& \cos \theta=-\frac{1}{\eta}=-\frac{1}{2}
\end{aligned}
\]

1.10 The solution of this problem becomes simple in the frame attached with one of the bodies. Let the body thrown straight up be 1 and the other body be 2 , then for the body 1 in the frame of 2 from the kinematic equation for constant acceleration :

\[
\overrightarrow{r_{12}}=\overrightarrow{r_{0(12)}}+\overrightarrow{v_{0(12)}} t+\frac{1}{2} \vec{w}_{12} t^{2}
\]

So, \(\quad \overrightarrow{r_{12}}=\overrightarrow{v_{0(12)}} t,\left(\right.\) because \(\vec{w}_{12}=0\) and \(\left.\vec{r}_{0(12)}=0\right)\)\\
or, \(\quad\left|\overrightarrow{r_{12}}\right|=\left|\overrightarrow{v_{0(12)}}\right| t\)

But \(\quad\left|\vec{v}_{01}\right|=\left|\vec{v}_{02}\right|=v_{0}\)\\
So, from properties of triangle

\[
v_{\alpha(12)}=\sqrt{v_{0}^{2}+v_{0}^{2}-2 v_{0} v_{0} \cos \left(\pi / 2-\theta_{0}\right)}
\]

Hence, the sought distance

\[
\left|\overrightarrow{r_{12}}\right|=v_{0} \sqrt{2(1-\sin \theta)} t=22 \mathrm{~m}
\]

1.11 Let the velocities of the paricles (say \(\overrightarrow{v_{1}}{ }^{\prime}\) and \(\overrightarrow{v_{2}}\) ) becomes mutually perpendicular after time \(t\). Then their velocitis become

\[
\overrightarrow{v_{1}^{\prime}}=\overrightarrow{v_{1}}+\overrightarrow{g t} ; \overrightarrow{v_{2}^{\prime}}=\overrightarrow{v_{2}}+\overrightarrow{g t}
\]

As \(\overrightarrow{v_{1}} \perp \overrightarrow{v_{2}}\) so, \(\overrightarrow{v_{1}} \cdot \overrightarrow{v_{2}}=0\)\\
or, \(\left(\overrightarrow{v_{1}}+\overrightarrow{g t}\right) \cdot\left(\overrightarrow{v_{2}}+\overrightarrow{g t}\right)=0\)\\
or \(\quad-v_{1} v_{2}+g^{2} t^{2}=0\)\\
\includegraphics[max width=\textwidth, alt={}, center]{9f3c5a8e-ac5d-42f4-8eb6-ea2af623f2d8-012_314_682_185_704}

Hence, \(\quad t=\frac{\sqrt{v_{1} v_{2}}}{g}\)\\
Now form the Eq. \(\quad \vec{r}_{12}=\vec{r}_{0(12)}+\vec{v}_{0(12)} t+\frac{1}{2} \vec{w}_{12} t^{2}\)\\
\(\left|\overrightarrow{r_{12}}\right|=\left|\overrightarrow{v_{0(12)}}\right| t\), (because here \(\vec{w}_{12}=0\) and \(\vec{r}_{0(12)}=0\) )

Hence the sought distance

\[
\left|\overrightarrow{r_{12}}\right|=\frac{v_{1}+v_{2}}{g} \sqrt{v_{1} v_{2}} \quad\left(\text { as }\left|\vec{v}_{0(12)}\right|=v_{1}+v_{2}\right)
\]

1.12 From the symmetry of the problem all the three points are always located at the vertices of equilateral triangles of varying side length and finally meet at the centriod of the initial equilateral triangle whose side length is \(a\), in the sought time interval (say \(t\) ).\\
\includegraphics[max width=\textwidth, alt={}, center]{9f3c5a8e-ac5d-42f4-8eb6-ea2af623f2d8-012_509_506_1059_156}\\
\includegraphics[max width=\textwidth, alt={}, center]{9f3c5a8e-ac5d-42f4-8eb6-ea2af623f2d8-012_372_423_1095_937}

Let us consider an arbitrary equilateral triangie of edge length \(l\) (say).\\
Then the rate by which 1 approaches 2,2 approches 3 , and 3 approches 1 , becomes :

On integrating :

\[
\begin{gathered}
\frac{-d l}{d t}=v-v \cos \left(\frac{2 \pi}{3}\right) \\
-\int_{a}^{0} d l=\frac{3 v}{2} \int_{0}^{1} d t \\
a=\frac{3}{2} v t \quad \text { so } t=\frac{2 a}{3 v}
\end{gathered}
\]

1.13 Let us locate the points \(A\) and \(B\) at an arbitrary instant of time (Fig.).

If \(A\) and \(B\) are separated by the distance \(s\) at this moment, then the points converge or point \(A\) approaches \(B\) with velocity \(\frac{-d s}{d t}=v-u \cos \alpha\) where angle \(\alpha\) varies with time.

On intergating,

\[
-\int_{l}^{0} d s=\int_{0}^{T}(v-u \cos \alpha) d t
\]

(where \(T\) is the sought time.)\\
or \(\quad l=\int_{0}^{T}(v-u \cos \alpha) d t\)

As both \(A\) and \(B\) cover the same distance in \(x\)-direction during the sought time interval, so the other condition which is required, can be obtained by the equation

So,\\
\includegraphics[max width=\textwidth, alt={}, center]{9f3c5a8e-ac5d-42f4-8eb6-ea2af623f2d8-013_435_690_254_697}


\begin{gather*}
\Delta x=\int v_{x} d t \\
u T=\int_{0}^{T} v \cos \alpha d t \tag{2}
\end{gather*}


Solving (1) and (2), we get \(T=\frac{u l}{v^{2}-u^{2}}\)\\
One can see that if \(u=v\), or \(u<v\), point \(A\) cannot catch \(B\).\\
1.14 In the reference frame fixed to the train, the distance between the two events is obviously equal to \(l\). Suppose the train starts moving at time \(t=0\) in the positive \(x\) direction and take the origin \((x=0)\) at the head-light of the train at \(t=0\). Then the coordinate of first event in the earth's frame is

\[
x_{1}=\frac{1}{2} w t^{2}
\]

and similarly the coordinate of the second event is

\[
x_{2}=\frac{1}{2} w(t+\tau)^{2}-l
\]

The distance between the two events is obviously.

\[
x_{1}-x_{2}=l-w \tau(t+\tau / 2)=0.242 \mathrm{~km}
\]

in the reference frame fixed on the earth..\\
For the two events to occur at the same point in the reference frame \(K\), moving with constant velocity \(V\) relative to the earth, the distance travelled by the frame in the time interval \(T\) must be equal to the above distance.\\
Thus

\[
V \tau=l-w \tau(t+\tau / 2)
\]

So,

\[
V=\frac{l}{\tau}-w(t+\tau / 2)=4.03 \mathrm{~m} / \mathrm{s}
\]

The frame \(K\) must clearly be moving in a direction opposite to the train so that if (for example) the origin of the frame coincides with the point \(x_{1}\) on the earth at time \(t\), it coincides with the point \(x_{2}\) at time \(t+\tau\).\\
1.15 (a) One good way to solve the problem is to work in the elevator's frame having the observer at its bottom (Fig.).\\
Let us denote the separation between floor and celing by \(h=2.7 \mathrm{~m}\). and the acceleration of the elevator by \(w=1.2 \mathrm{~m} / \mathrm{s}^{2}\)\\
From the kinematical formula


\begin{equation*}
y=y_{0}+v_{0 y} t+\frac{1}{2} w_{y} t^{2} \tag{1}
\end{equation*}


Here \(y=0, y_{0}=+h, v_{0 y}=0\)\\
and

\[
\begin{aligned}
w_{y} & =w_{\text {bolt }(y)}-w_{\text {cle }(y)} \\
& =(-g)-(w)=-(g+w)
\end{aligned}
\]

So, \(\quad 0=h+\frac{1}{2}\{-(g+w)\} t^{2}\)\\
\includegraphics[max width=\textwidth, alt={}, center]{9f3c5a8e-ac5d-42f4-8eb6-ea2af623f2d8-014_374_480_353_869}\\
or, \(\quad t=\sqrt{\frac{2 h}{g+w}}=0.7 \mathrm{~s}\).\\
(b) At the moment the bolt loses contact with the elevator, it has already aquired the velocity equal to elevator,given by :

\[
v_{0}=(1.2)(2)=2.4 \mathrm{~m} / \mathrm{s}
\]

In the reference frame attached with the elevator shaft (ground) and pointing the \(y\)-axis upward, we have for the displacement of the bolt,

\[
\begin{aligned}
\Delta y & =v_{0 y} t+\frac{1}{2} w_{y} t^{2} \\
& =v_{0} t+\frac{1}{2}(-g) t^{2} \\
\text { or, } \quad \Delta y & =(2.4)(0.7)+\frac{1}{2}(-9.8)(0.7)^{2}=-0.7 \mathrm{~m} .
\end{aligned}
\]

\begin{center}
\includegraphics[max width=\textwidth, alt={}]{9f3c5a8e-ac5d-42f4-8eb6-ea2af623f2d8-014_331_190_1005_1127}
\end{center}

Hence the bolt comes down or displaces downward relative to the point, when it loses contact with the elevator by the amount 0.7 m (Fig.).\\
Obviously the total distance covered by the bolt during its free fall time

\[
s=|\Delta y|+2\left(\frac{v_{0}^{2}}{2 g}\right)=0.7 \mathrm{~m}+\frac{(2 \cdot 4)^{2}}{(9 \cdot 8)} \mathrm{m}=1.3 \mathrm{~m} .
\]

1.16 Let the particle 1 and 2 be at points \(B\) and \(A\) at \(t=0\) at the distances \(l_{1}\) and \(l_{2}\) from intersection point \(O\).\\
Let us fix the inertial frame with the particle 2 . Now the particle 1 moves in relative to this reference frame with a relative velocity \(\overrightarrow{v_{12}}=\overrightarrow{v_{1}}-\overrightarrow{v_{2}}\) and its trajectory is the straight line \(B P\). Obviously, the minimum distance between the particles is equal to the length of the perpendicular \(A P\) dropped from point \(A\) on to the straight line \(B P\) (Fig.).\\
\includegraphics[max width=\textwidth, alt={}, center]{9f3c5a8e-ac5d-42f4-8eb6-ea2af623f2d8-015_482_665_112_138}\\
\includegraphics[max width=\textwidth, alt={}, center]{9f3c5a8e-ac5d-42f4-8eb6-ea2af623f2d8-015_349_300_178_984}

From Fig. (b),


\begin{equation*}
v_{12}=\sqrt{v_{1}^{2}+v_{2}^{2}}, \text { and } \tan \theta=\frac{v_{1}}{v_{2}} \tag{1}
\end{equation*}


The shortest distance\\
or

\[
A P=A M \sin \theta=(O A-O M) \sin \theta=\left(l_{2}-l_{1} \cot \theta\right) \sin \theta
\]

\[
\left.A P=\left(l_{2}-l_{1} \frac{v_{2}}{v_{1}}\right) \frac{v_{1}}{\sqrt{v_{1}^{2}+v_{2}^{2}}}=\frac{v_{1} l_{2}-v_{2} l_{1}}{\sqrt{v_{1}^{2}+v_{2}^{2}}} \quad \text { (using } 1\right)
\]

The sought time can be obtained directly from the condition that \(\left(l_{1}-v_{1} t\right)^{2}+\left(l_{2}-v_{2} t\right)^{2}\) is minimum. This gives \(t=\frac{l_{1} v_{1}+l_{2} v_{2}}{v_{1}^{2}+v_{2}^{2}}\).\\
1.17 Let the car turn off the highway at a distance \(x\) from the point \(D\).

So, \(C D=x\), and if the speed of the car in the field is \(v\), then the time taken by the car to cover the distance \(A C=A D-x\) on the highway


\begin{equation*}
t_{1}=\frac{A D-x}{\eta v} \tag{1}
\end{equation*}


and the time taken to travel the distance \(C B\) in the field


\begin{equation*}
t_{2}=\frac{\sqrt{l^{2}+x^{2}}}{v} \tag{2}
\end{equation*}


So, the total time elapsed to move the car from point \(A\) to \(B\)

\[
t=t_{1}+t_{2}=\frac{A D-x}{\eta v}+\frac{\sqrt{l^{2}+x^{2}}}{v}
\]

For \(t\) to be minimum\\
\includegraphics[max width=\textwidth, alt={}, center]{9f3c5a8e-ac5d-42f4-8eb6-ea2af623f2d8-015_483_470_1365_871}

\[
\frac{d t}{d x}=0 \text { or } \frac{1}{v}\left[-\frac{1}{\eta}+\frac{x}{\sqrt{l^{2}+x^{2}}}\right]=0
\]

or

\[
\eta^{2} x^{2}=l^{2}+x^{2} \text { or } x=\frac{l}{\sqrt{\eta^{2}-1}}
\]

1.18 To plot \(x(t), s(t)\) and \(w_{x}(t)\) let us partion the given plot \(v_{x}(t)\) into five segments (for detailed analysis) as shown in the figure.\\
For the part \(o a: w_{x}=1\) and \(v_{x}=t=v\)\\
Thus, \(\Delta x_{1}(t)=\int v_{x} d t=\int_{0}^{t} d t=\frac{t^{2}}{2}=s_{1}(t)\)\\
Putting \(t=1\), we get, \(\Delta x_{1}=s=\frac{1}{2}\) unit\\
For the part ab :\\
\includegraphics[max width=\textwidth, alt={}, center]{9f3c5a8e-ac5d-42f4-8eb6-ea2af623f2d8-016_267_445_224_783}\\
\(w_{x}=0\) and \(v_{x}=v=\) constant \(=1\)

Thus

\[
\Delta x_{2}(t)=\int v_{x} d t=\int_{1}^{t} d t=(t-1)=s_{2}(t)
\]

Putting

\[
t=3, \Delta x_{2}=s_{2}=2 \text { unit }
\]

For the part \(b 4: \quad w_{x}=1\) and \(\left.v_{x}=1-(t-3)=4-t\right)=v\)

Thus

\[
\Delta x_{3}(t)=\int_{3}^{t}(4-t) d t=4 t-\frac{t^{2}}{2}-\frac{15}{2}=s_{3}(t)
\]

Putting

\[
t=4, \Delta x_{3}=x_{3}=\frac{1}{2} \text { unit }
\]

For the part 4d:

\[
v_{x}=-1 \text { and } v_{x}=-(1-4)=4-1
\]

So,

\[
v=\left|v_{x}\right|=t-4 \text { for } t>4
\]

Thus

\[
\Delta x_{4}(t)=\int_{4}^{t}(1-t) d t=4 t-\frac{t^{2}}{2}-8
\]

Putting

\[
t=6, \Delta x_{4}=-1
\]

Similarly

\[
s_{4}(t)=\int\left|v_{x}\right| d t=\int_{4}^{t}(t-4) d t=\frac{t^{2}}{2}-4 t+8
\]

Putting

\[
t=6, s_{4}=2 \text { unit }
\]

For the part \(d 7\) :

\[
\begin{aligned}
w_{x} & =2 \text { and } v_{x}=-2+2(t-6)=2(t-7) \\
v & =\left|v_{x}\right|=2(7-t) \text { for } t \leftarrow 7
\end{aligned}
\]

Now,

\[
\Delta x(t)=\int_{t} 2(t-7) d t=t^{2}-14 t+48
\]

Putting

\[
t=4, \Delta x_{5}=-1
\]

Similarly

\[
s_{5}(t)=\int_{t}^{6} 2(7-t) d t=14 t-t^{2}-48
\]

Putting

\[
t=7, s_{5}=1
\]

On the basis of these obtained expressions \(w_{x}(t), x(t)\) and \(s(t)\) plots can be easily plotted as shown in the figure of answersheet.\\
1.19 (a) Mean velocity


\begin{align*}
<v> & =\frac{\text { Total distance covered }}{\text { Time elapsed }} \\
& =\frac{s}{t}=\frac{\pi R}{\tau}=50 \mathrm{~cm} / \mathrm{s} \tag{1}
\end{align*}


(b) Modulus of mean velocity vector\\
\includegraphics[max width=\textwidth, alt={}, center]{9f3c5a8e-ac5d-42f4-8eb6-ea2af623f2d8-017_261_437_120_767}


\begin{equation*}
|<\vec{v}>|=\frac{|\Delta \vec{r}|}{\Delta t}=\frac{2 R}{\tau}=32 \mathrm{~cm} / \mathrm{s} \tag{2}
\end{equation*}


(c) Let the point moves from \(i\) to \(f\) along the half circle (Fig.) and \(v_{0}\) and \(\nu\) be the spe at the points respectively.

We have

\[
\frac{d v}{d t}=w_{t}
\]

or, \(\quad v=v_{0}+w_{t} t\) (as \(w_{t}\) is constant, according to the problem)

So,

\[
<v>=\frac{\int_{0}^{t}\left(v_{0}+w_{t} t\right) d t}{\int_{0}^{t} d t}=\frac{v_{0}+\left(v_{0}+w_{t} t\right)}{2}=\frac{v_{0}+v}{2}
\]

So, from (1) and (3)


\begin{equation*}
\frac{v_{0}+v}{2}=\frac{\pi R}{\tau} \tag{(3}
\end{equation*}


Now the modulus of the mean vector of total acceleration


\begin{equation*}
|<\vec{w}>|=\frac{|\Delta \vec{v}|}{\Delta t}=\frac{\left|\vec{v}-\overrightarrow{v_{0}}\right|}{\tau}=\frac{v_{0}+v}{\tau} \text { (see Fig.) } \tag{(5}
\end{equation*}


Using (4) in (5), we get :

\[
|<\vec{w}>|=\frac{2 \pi R}{\tau^{2}}
\]

1.20 (a) we have

\[
\vec{r}=\vec{a} t(1-\alpha t)
\]

So,

\[
\vec{v}=\frac{d \vec{r}}{d t}=\vec{a}(1-2 \alpha t)
\]

and

\[
\vec{w}=\frac{d \vec{v}}{d t}=-2 \alpha \vec{a}
\]

(b) From the equation

\[
\vec{r}=\overrightarrow{a t}(1-\alpha t),
\]

\[
\vec{r}=0, \text { at } t=0 \text { and also at } t=\Delta t=\frac{1}{\alpha}
\]

So, the sought time

\[
\begin{aligned}
\Delta t & =\frac{1}{\alpha} \\
\vec{v} & =\vec{a}(1-2 \alpha t)
\end{aligned}
\]

As

\[
\left.v=|\vec{v}|=\begin{array}{l}
a(1-2 \alpha t) \\
a(2 \alpha t-1)
\end{array}\right\} \begin{aligned}
& \text { for } t \leq \frac{1}{2 \alpha} \\
& \text { for } t>\frac{1}{2 \alpha}
\end{aligned}
\]

So,

Hence, the sought distance

\[
s=\int v d t=\int_{0}^{1 / 2 \alpha} a(1-2 \alpha t) d t+\int_{1 / 2 \alpha}^{1 / \alpha} a(2 \alpha t-1) d t
\]

Simplifying, we get, \(s=\frac{a}{2 \alpha}\)\\
1.21 (a) As the particle leaves the origin at \(t=0\)

So,


\begin{equation*}
\Delta x=x=\int v_{x} d t \tag{1}
\end{equation*}


As

\[
\vec{v}=\vec{v}_{0}\left(1-\frac{t}{\tau}\right),
\]

where \(\overrightarrow{v_{0}}\) is directed towards the \(+\mathrm{ve} x\)-axis\\
So,


\begin{equation*}
v_{x}=v_{0}\left(1-\frac{t}{\tau}\right) \tag{2}
\end{equation*}


From (1) and (2),


\begin{equation*}
x=\int_{0}^{t} v_{0}\left(1-\frac{t}{\tau}\right) d t=v_{0} t\left(1-\frac{t}{2 \tau}\right) \tag{3}
\end{equation*}


Hence \(x\) coordinate of the particle at \(t=6 \mathrm{~s}\).

\[
x=10 \times 6\left(1-\frac{6}{2 \times 5}\right)=24 \mathrm{~cm}=0.24 \mathrm{~m}
\]

Similarly at

\[
t=10 \mathrm{~s}
\]

\[
x=10 \times 10\left(1-\frac{10}{2 \times 5}\right)=0
\]

and at

\[
t=20 \mathrm{~s}
\]

\[
x=10 \times 20\left(1-\frac{20}{2 \times 5}\right)=-200 \mathrm{~cm}=-2 \mathrm{~m}
\]

(b) At the moments the particle is at a distance of 10 cm from the origin, \(x= \pm 10 \mathrm{~cm}\).

Putting

\[
x=+10 \text { in Eq. (3) }
\]

\[
10=10 t\left(1-\frac{t}{10}\right) \text { or, } t^{2}-10 t+10=0
\]

So,

\[
t=t=\frac{10 \pm \sqrt{100-40}}{2}=5 \pm \sqrt{15} \mathrm{~s}
\]

Now putting

\[
x=-10 \text { in Eqn }(3)
\]

\[
-10=10\left(1-\frac{t}{10}\right)
\]

On solving,

\[
t=5 \pm \sqrt{35} \mathrm{~s}
\]

As \(t\) cannot be negative, so,

\[
t=(5+\sqrt{35}) \mathrm{s}
\]

Hence the particle is at a distance of 10 cm from the origin at three moments of time :

\[
t=5 \pm \sqrt{15} \mathrm{~s}, 5+\sqrt{35} \mathrm{~s}
\]

(c) We have

\[
\vec{v}=\overrightarrow{v_{0}}\left(1-\frac{t}{\tau}\right)
\]

So,

\[
\left.v=|\vec{v}|=\begin{array}{l}
v_{0}\left(1-\frac{t}{\tau}\right) \\
v_{0}\left(\frac{t}{\tau}-1\right)
\end{array}\right\} \begin{aligned}
& \text { for } t \leq \tau \\
& \text { for } t>\tau
\end{aligned}
\]

So

\[
s=\int_{0}^{t} v_{0}\left(1-\frac{t}{\tau}\right) d t \text { for } t \leq \tau=v_{0} t(1-t / 2 \tau)
\]

and


\begin{align*}
& s=\int_{0} v_{0}\left(1-\frac{t}{2}\right) d t+\int_{\tau} v_{0}\left(\frac{t}{\tau}-1\right) d t \text { for } t>\tau \\
= & v_{0} \tau\left[1+(1-t / \tau)^{2}\right] / 2 \text { for } t>\tau \tag{A}
\end{align*}


\[
s=\int_{0}^{4} v_{0}\left(1-\frac{t}{\tau}\right) d t=\int_{0}^{4} 10\left(1-\frac{t}{5}\right) d t=24 \mathrm{~cm}
\]

And for \(t=8 \mathrm{~s}\)

\[
s=\int_{0}^{5} 10\left(1-\frac{t}{5}\right) d t+\int_{5}^{8} 10\left(\frac{t}{5}-1\right) d t
\]

On integrating and simplifying, we get

\[
s=34 \mathrm{~cm} .
\]

On the basis of Eqs. (3) and (4), \(x(t)\) and \(s(t)\) plots can be drawn as shown in the answer sheet.\\
1.22 As particle is in unidirectional motion it is directed along the \(x\)-axis all the time. As at \(t=0, x=0\)

So,

\[
\begin{aligned}
\Delta x & =x=s, \text { and } \frac{d v}{d t}=w \\
v & =\alpha \sqrt{x}=\alpha \sqrt{s}
\end{aligned}
\]

Therefore,\\
or,


\begin{align*}
w & =\frac{d v}{d t}=\frac{\alpha}{2 \sqrt{s}} \frac{d s}{d t}=\frac{\alpha}{2 \sqrt{s}} \\
& =\frac{\alpha v}{2 \sqrt{s}}=\frac{\alpha \alpha \sqrt{s}}{2 \sqrt{s}}=\frac{\alpha^{2}}{2} \tag{1}
\end{align*}


As,

\[
w=\frac{d v}{d t}=\frac{\alpha^{2}}{2}
\]

On integrating,


\begin{equation*}
\int_{0}^{v} d v=\int_{0}^{t} \frac{\alpha^{2}}{2} d t \quad \text { or, } \quad v=\frac{\alpha^{2}}{2} t \tag{2}
\end{equation*}


(b) Let \(s\) be the time to cover first \(s m\) of the path. From the Eq.


\begin{align*}
s & =\int v d t \\
s & =\int_{0}^{t} \frac{\alpha^{2}}{2} d t=\frac{\alpha^{2}}{2} \frac{t^{2}}{2}
\end{align*}


or


\begin{equation*}
t=\frac{2}{\alpha} \sqrt{s} \tag{3}
\end{equation*}


The mean velocity of particle

\[
<v>=\frac{\int v(t) d t}{\int d t}=\frac{\int_{0}^{2 \sqrt{s} / \alpha} \frac{\alpha^{2}}{2} t d t}{2 \sqrt{s} / \alpha}=\frac{\alpha \sqrt{s}}{2}
\]

1.23 According to the problem

\[
-\frac{v d v}{d s}=a \sqrt{v}(\text { as } v \text { decreases with time })
\]

or,

\[
-\int_{v_{0}}^{0} \sqrt{v} d v=a \int_{0}^{s} d s
\]

On integrating we get \(s=\frac{2}{3 a} v_{0}^{3 / 2}\)\\
Again according to the problem\\
\(-\frac{d v}{d t}=a \sqrt{v}\) or \(-\frac{d v}{\sqrt{v}}=a d t\)\\
or,

\[
\int_{v_{0}}^{0} \frac{d v}{\sqrt{v}}=a \int_{0}^{t} d t
\]

Thus

\[
t=\frac{2^{\sqrt{v_{0}}}}{a}
\]

1.24 (a) As

So,\\
and therefore

\[
\vec{r}=a t \vec{i}-b t^{2} \vec{j}
\]

\[
y=\frac{-b x^{2}}{a^{2}}
\]

\includegraphics[max width=\textwidth, alt={}, center]{9f3c5a8e-ac5d-42f4-8eb6-ea2af623f2d8-020_332_361_1710_838}\\
which is Eq. of a parabola, whose graph is shown in the Fig.\\
(b) As

\[
\vec{r}=a t \vec{i}-b t^{2} \vec{j}
\]


\begin{equation*}
\vec{v}=\frac{d \vec{r}}{d t}=a \vec{i}-2 b t \vec{j} \tag{1}
\end{equation*}


So,

\[
v=\sqrt{a^{2}(-2 b t)^{2}}=\sqrt{a^{2}+4 b^{2} t^{2}}
\]

Diff. Eq. (1) w.r.t. time, we get

\[
\vec{w}=\frac{d \vec{v}}{d t}=-2 b \vec{j}
\]

So,

\[
|\vec{w}|=w=2 b
\]

(c)

\[
\cos \alpha=\frac{\vec{v} \cdot \vec{w}}{v w}=\frac{(a \vec{i}-2 b t \vec{j}) \cdot(-2 b \vec{j})}{\left(\sqrt{a^{2}+4 b^{2} t^{2}}\right) 2 b}
\]

or,

\[
\cos \alpha=\frac{2 b t}{\sqrt{a^{2}+4 b^{2} t^{2}}},
\]

so,

\[
\tan \alpha=\frac{a}{2 b t}
\]

or, \(\alpha=\tan ^{-1}\left(\frac{a}{2 b t}\right)\)\\
(d) The mean velocity vector

\[
<\vec{v}>=\frac{\int \overrightarrow{v d t}}{\int d t}=\frac{\int_{0}^{t}(a \vec{i}-2 b t \vec{j}) d t}{t}=a \vec{i}-b t \vec{j}
\]

Hence,

\[
|<\vec{v}>|=\sqrt{a^{2}+(-b t)^{2}}=\sqrt{a^{2}+b^{2} t^{2}}
\]

1.25 (a) We have


\begin{equation*}
x=a t \text { and } y=a t(1-\alpha t) \tag{1}
\end{equation*}


Hence, \(y(x)\) becomes,

\[
y=\frac{a x}{a}\left(1-\frac{\alpha x}{a}\right)=x-\frac{\alpha}{a} x^{2} \text { (parabola) }
\]

(b) Differentiating Eq. (1) we get


\begin{equation*}
v_{x}=a \text { and } v_{y}=a(1-2 \alpha t) \tag{2}
\end{equation*}


So,

\[
v=\sqrt{v_{x}^{2}+v_{y}^{2}}=a \sqrt{1+(1-2 \alpha t)^{2}}
\]

Diff. Eq. (2) with respect to time

So,

\[
\begin{aligned}
w_{x} & =0 \text { and } \quad w_{y}=-2 a \alpha \\
w & =\sqrt{w_{x}^{2}+w_{y}^{2}}=2 a \alpha
\end{aligned}
\]

(c) From Eqs.\\
(2) and\\
(3)

We have

\[
\vec{v}=a \vec{i}+a(1-2 \alpha t) \vec{j} \text { and } \vec{w}=2 a \alpha \vec{j}
\]

So,

\[
\cos \frac{\pi}{4}=\frac{1}{\sqrt{2}}=\frac{\vec{v} \cdot \vec{w}}{v w}=\frac{-a\left(1-2 \alpha t_{0}\right) 2 a \alpha}{a \sqrt{1+\left(1-2 \alpha t_{0}\right)^{2}} 2 a \alpha}
\]

On simplifying.

\[
1-2 \alpha t_{0}= \pm 1
\]

As,

\[
t_{0} \approx 0, t_{0}=\frac{1}{\alpha}
\]

1.26 Differentiating motion law : \(x=a \sin \omega t, y=a(1-\cos \omega t)\), with respect to time, \(v_{x}=a \omega \cos \omega t, v_{y}=a \omega \sin \omega t\)\\
So,


\begin{equation*}
\vec{v}=a \omega \cos \omega t \vec{i}+a \omega \sin \omega t \vec{j} \tag{1}
\end{equation*}


and


\begin{equation*}
v=a \omega=\text { Const. } \tag{2}
\end{equation*}


Differentiating Eq. (1) with respect to time


\begin{equation*}
\vec{w}=\frac{d \vec{v}}{d t}=-a \omega^{2} \sin \omega t \vec{i}+a \omega^{2} \cos \omega t \vec{j} \tag{3}
\end{equation*}


(a) The distance \(s\) traversed by the point during the time \(\tau\) is given by

\[
s=\int_{0}^{\tau} v d t=\int_{0}^{\tau} a \omega d t=a \omega \tau \quad(\text { using } 2)
\]

(b) Taking inner product of \(\overrightarrow{\boldsymbol{v}}\) and \(\overrightarrow{\boldsymbol{w}}\)

We get, \(\quad \vec{v} \cdot \vec{w}=(a \omega \cos \omega t \vec{i}+a \omega \sin \omega t \vec{j}) \cdot\left(a \omega^{2} \sin \omega t(-i)+a \omega^{2} \cos \omega t-\vec{j}\right)\)\\
So,

\[
\vec{v} \cdot \vec{w}=-a^{2} \omega^{2} \sin \omega t \cos \omega t+a^{2} \omega^{3} \sin \omega t \cos \omega t=0
\]

Thus, \(\vec{v} \perp \vec{w}\), i.e., the angle between velocity vector and acceleration vector equals \(\frac{\pi}{2}\).\\
1.27 Accordiing to the problem

\[
\vec{w}=w(-\vec{j})
\]

So,


\begin{equation*}
w_{x}=\frac{d v_{x}}{d t}=0 \text { and } w_{y}=\frac{d v_{y}}{d t}=-w \tag{1}
\end{equation*}


Differentiating Eq. of trajectory, \(y=a x-b x^{2}\), with respect to time


\begin{equation*}
\frac{d y}{d t}=\frac{a d x}{d t}-2 b x \frac{d x}{d t} \tag{2}
\end{equation*}


So,

\[
\left.\frac{d y}{d t}\right|_{x=0}=\left.a \frac{d x}{d t}\right|_{x=0}
\]

Again differentiating with respect to time

\[
\frac{d^{2} y}{d t^{2}}=\frac{a d^{2} x}{d t^{2}}-2 b\left(\frac{d x}{d t}\right)^{2}-2 b x \frac{d^{2} x}{d t^{2}}
\]

or,

\[
-w=a(0)-2 b\left(\frac{d x}{d t}\right)^{2}-2 b x(0)(\text { using } 1)
\]

or,


\begin{equation*}
\frac{d x}{d t}=\sqrt{\frac{w}{2 b}}(\text { using } 1) \tag{∴}
\end{equation*}


Using (3) in (2)


\begin{equation*}
\left.\frac{d y}{d t}\right|_{x=0}=a \sqrt{\frac{w}{2 b}} \tag{\((<\)}
\end{equation*}


Hence, the velocity of the particle at the origin\\
\(v=\sqrt{\left(\frac{d x}{d t}\right)_{x=0}^{2}+\left(\frac{d y}{d t}\right)_{x=0}^{2}}=\sqrt{\frac{w}{2 b}+a^{2} \frac{w}{2 b}}\) (using Eqns (3) and (4))

Hence,

\[
v=\sqrt{\frac{w}{2 b}\left(1+a^{2}\right)}
\]

1.28 As the body is under gravity of constant accelration \(\vec{g}\), it's velocity vector and displacemen vectors are:\\
and


\begin{gather*}
\vec{v}=\overrightarrow{v_{0}}+\vec{g} t  \tag{1}\\
\Delta \vec{r}=\vec{r}=\overrightarrow{v_{0}} t+\frac{1}{2} g t^{2}(\vec{r}=0 \text { at } t=0) \tag{2}
\end{gather*}


So, \(<\overrightarrow{v>}\) over the first \(t\) seconds


\begin{equation*}
\left\langle\overrightarrow{v\rangle}=\frac{\Delta \vec{r}}{\Delta t}=\frac{\vec{r}}{t}=\overrightarrow{v_{0}}+\frac{\overrightarrow{g t}}{2}\right. \tag{3}
\end{equation*}


Hence from Eq. (3), << >> over the first \(t\) seconds


\begin{equation*}
\left\langle\overrightarrow{v\rangle}=\overrightarrow{v_{0}}+\frac{\vec{g}}{2} \tau\right. \tag{4}
\end{equation*}


For evaluating \(t\), take

\[
\vec{v} \cdot \vec{v}=\left(\overrightarrow{v_{0}}+\vec{g} t\right) \cdot\left(\overrightarrow{v_{0}}+\vec{g} t\right)=v_{0}^{2}+2\left(\overrightarrow{v_{0}} \cdot \vec{g}\right) t+g^{2} t^{2}
\]

or, \(v^{2}=v_{0}^{2}+\left(\overrightarrow{v_{0}} \cdot \vec{g}\right) t+g^{2} t^{2}\)\\
\includegraphics[max width=\textwidth, alt={}, center]{9f3c5a8e-ac5d-42f4-8eb6-ea2af623f2d8-023_475_489_1491_904}

But we have \(v=v_{0}\) at \(t=0\) and

Also at \(t=\tau\) (Fig.) (also from energy conservation)

Hence using this propety in Eq. (5)

\[
v_{0}^{2}=v_{0}^{2}+2\left(\overrightarrow{v_{0}} \cdot \vec{g}\right) \tau+g^{2} \tau^{2}
\]

As

\[
\tau \neq 0, \text { so, } \tau=-\frac{2\left(\vec{v}_{o} \cdot \vec{g}\right)}{g^{2}}
\]

Putting this value of \(\tau\) in Eq. (4), the average velocity over the time of flight

\[
<\overrightarrow{v>}=\overrightarrow{v_{0}}-\vec{g} \frac{\left(\overrightarrow{v_{0}} \cdot \vec{g}\right)}{g^{2}}
\]

1.29 The body thrown in air with velocity \(v_{0}\) at an angle \(\alpha\) from the horizontal lands at point \(P\) on the Earth's surface at same horizontal level (Fig.). The point of projection is taken as origin, so, \(\Delta x=x\) and \(\Delta y=y\)\\
(a) From the Eq. \(\Delta y=v_{0} t+\frac{1}{2} w_{y} t^{2}\)

\[
0=v_{0} \sin \alpha \tau-\frac{1}{2} g \tau^{2}
\]

As \(\tau \neq 0\), so, time of motion \(\tau=\frac{2 v_{0} \sin \alpha}{g}\)\\
(b) At the maximum height of ascent, \(\nu_{y}=0\) so, from the Eq. \(\quad v_{y}^{2}=v_{0 y}^{2}+2 w_{y} \Delta y\)

\[
0=\left(v_{0} \sin \alpha\right)^{2}-2 g H
\]

Hence maximum height \(H=\frac{v_{0}^{2} \sin ^{2} \alpha}{2 g}\)\\
\includegraphics[max width=\textwidth, alt={}, center]{9f3c5a8e-ac5d-42f4-8eb6-ea2af623f2d8-024_480_476_766_851}

During the time of motion the net horizontal displacement or horizontal range, will be obtained by the equation

\[
\Delta x=v_{0 x} t+\frac{1}{2} w_{x} \tau^{2}
\]

or,

\[
R=v_{0} \cos \alpha \tau-\frac{1}{2}(0) \tau^{2}=v_{0} \cos \alpha \tau=\frac{v_{0}^{2} \sin 2 \alpha}{g}
\]

when

\[
R=H
\]

\[
\frac{v_{0}^{2} \sin ^{2} \alpha}{g}=\frac{v_{0}^{2} \sin ^{2} \alpha}{2 g} \text { or } \tan \alpha=4, \text { so, } \alpha=\tan ^{-1} 4
\]

(c) For the body, \(x(t)\) and \(y(t)\) are


\begin{equation*}
x=v_{0} \cos \alpha t \tag{1}
\end{equation*}


and


\begin{equation*}
y=v_{0} \sin \alpha t-\frac{1}{2} g t^{2} \tag{2}
\end{equation*}


Hence putting the value of \(\boldsymbol{t}\) from (1) into (2) we get,

\[
y=v_{0} \sin \alpha\left(\frac{x}{v_{0} \cos \alpha}\right)-\frac{1}{2} g\left(\frac{x}{v_{0} \cos \alpha}\right)^{2}=x \tan \alpha-\frac{g x^{2}}{2 v_{0}^{2} \cos ^{2} \alpha}
\]

Which is the sought equation of trajectory i.e. \(y(x)\)\\
(d) As the body thrown in air follows a curve, it has some normal acceleration at all the moments of time during it's motion in air.\\
At the initial point \((x=0, y=0)\), from the equation :

\[
\begin{aligned}
w_{n} & =\frac{v^{2}}{R},(\text { where } R \text { is the radius of curvature }) \\
g \cos \alpha & =\frac{v_{0}^{2}}{R_{0}} \text { (see Fig.) or } R_{0}=\frac{v_{0}^{2}}{g \cos \alpha}
\end{aligned}
\]

At the peak point \(v_{y}=0, v=v_{x}=v_{0} \cos \alpha\) and the angential acceleration is zero.\\
Now from the Eq.

\[
\begin{aligned}
w_{n} & =\frac{v^{2}}{R} \\
g & =\frac{v_{0}^{2} \cos ^{2} \alpha}{R}, \text { or } R=\frac{v_{0}^{2} \cos ^{2} \alpha}{g}
\end{aligned}
\]

Note : We may use the formula of curvature radius of a trajectory \(y(x)\), to solve part (d),

\[
R=\frac{\left[1+(d y / d x)^{2}\right]^{\frac{3}{2}}}{\left|d^{2} y / d x^{2}\right|}
\]

1.30 We have, \(v_{x}=v_{0} \cos \alpha, v_{y}=v_{0} \sin \alpha-g t\)

As

\[
\vec{v} \uparrow \uparrow \hat{u}_{t} \text { all the moments of time. }
\]

Thus

\[
v^{2}=v_{t}^{2}-2 g t v_{0} \sin \alpha+g^{2} t^{2}
\]

Now,

\[
\begin{aligned}
w_{t} & =\frac{d v_{t}}{d t}=\frac{1}{2 v_{t}} \frac{d}{d t}\left(v_{t}^{2}\right)=\frac{1}{v_{t}}\left(g^{2} t-g v_{0} \sin \alpha\right) \\
& =-\frac{g}{v_{t}}\left(v_{0} \sin \alpha-g t\right)=-g \frac{v_{y}}{v_{t}}
\end{aligned}
\]

Hence

\[
\left|w_{t}\right|=g \frac{\left|v_{y}\right|}{v}
\]

Now

\[
w_{n}=\sqrt{w^{2}-w_{t}^{2}}=\sqrt{g^{2}-g^{2} \frac{v_{y}^{2}}{v_{t}^{2}}}
\]

or

\[
w_{n}=g \frac{v_{x}}{v_{t}}\left(\text { where } v_{x}=\sqrt{v_{t}^{2}-v_{y}^{2}}\right)
\]

As \(\vec{v} \uparrow \uparrow \hat{v}_{t}\), during time of motion

\[
w_{v}=w_{t}=-g \frac{v_{y}}{v}
\]

On the basis of obtained expressions or facts the sought plots can be drawn as shown in the figure of answer sheet.\\
1.31 The ball strikes the inclined plane \((O x)\) at point \(O(\) origin \()\) with velocity \(v_{0}=\sqrt{2 g h}\) As the ball elastically rebounds, it recalls with same velocity \(v_{0}\), at the same angle \(\alpha\) from the normal or \(y\) axis (Fig.). Let the ball strikes the incline second time at \(P\), which is at a distance \(l\) (say) from the point \(O\), along the incline. From the equation

\[
\begin{aligned}
& y=v_{0 y} t+\frac{1}{2} w_{y} t^{2} \\
& 0=v_{0} \cos \alpha \tau-\frac{1}{2} g \cos \alpha t^{2}
\end{aligned}
\]

where \(\tau\) is the time of motion of ball in air while moving from \(O\) to \(P\).


\begin{equation*}
\text { As } \quad \tau \neq 0, \text { so, } \tau=\frac{2 v_{0}}{g} \tag{2}
\end{equation*}


Now from the equation.

\[
\begin{aligned}
& x=v_{0 x} t+\frac{1}{2} w_{x} t^{2} \\
& l=v_{0} \sin \alpha \tau+\frac{1}{2} g \sin \alpha \tau^{2}
\end{aligned}
\]

\includegraphics[max width=\textwidth, alt={}, center]{9f3c5a8e-ac5d-42f4-8eb6-ea2af623f2d8-026_576_570_678_795}\\
so,

\[
\begin{gathered}
l=v_{0} \sin \alpha\left(\frac{2 v_{0}}{g}\right)+\frac{1}{2} g \sin \alpha\left(\frac{2 v_{0}}{g}\right) \\
=\frac{4 v_{0}^{2} \sin \alpha}{g}(\text { using } 2)
\end{gathered}
\]

Hence the sought distance, \(l=\frac{4(2 g h) \sin \alpha}{g}=8 h \sin \alpha\) (Using Eq. 1)\\
1.32 Total time of motion


\begin{equation*}
\tau=\frac{2 v_{0} \sin \alpha}{g} \text { or } \sin \alpha=\frac{\tau g}{2 v_{0}}=\frac{9.8 \tau}{2 \times 240} \tag{1}
\end{equation*}


and horizontal range


\begin{equation*}
R=v_{0} \cos \alpha \tau \text { or } \cos \alpha=\frac{R}{v_{0} \tau}=\frac{5100}{240 \tau}=\frac{85}{4 \tau} \tag{2}
\end{equation*}


From Eqs. (1) and (2)

\[
\frac{(9 \cdot 8)^{2} \tau^{2}}{(480)^{2}}+\frac{(85)^{2}}{\left(4 \tau^{2}\right)^{2}}=1
\]

On simplifying

\[
\tau^{4}-2400 \tau^{2}+1083750=0
\]

Solving for \(\boldsymbol{\tau}^{\mathbf{2}}\) we get :

\[
\tau^{2}=\frac{2400 \pm \sqrt{1425000}}{2}=\frac{2400 \pm 1194}{2}
\]

Thus

\[
\begin{aligned}
& \tau=42.39 \mathrm{~s}=0.71 \mathrm{~min} \text { and } \\
& \tau=24.55 \mathrm{~s}=0.41 \mathrm{~min} \text { depending on the angle } \alpha .
\end{aligned}
\]

1.33 Let the shells collide at the point \(P(x, y)\). If the first shell takes \(t \mathrm{~s}\) to collide with second and \(\Delta t\) be the time interval between the firings, then


\begin{equation*}
x=v_{0} \cos \theta_{1} t=v_{0} \cos \theta_{2}(t-\Delta t) \tag{1}
\end{equation*}


and


\begin{align*}
y & =v_{0} \sin \theta_{1} t-\frac{1}{2} g t^{2} \\
& =v_{0} \sin \theta_{2}(t-\Delta t)-\frac{1}{2} g(t-\Delta t)^{2} \tag{2}
\end{align*}


From Eq. (1) \(t=\frac{\Delta t \cos \theta_{2}}{\cos \theta_{2}-\cos \theta_{1}}\)

From Eqs. (2) and (3)

\[
\Delta t=\frac{2 v_{0} \sin \left(\theta_{1}-\theta_{2}\right)}{g\left(\cos \theta_{2}+\cos \theta_{1}\right)} \text { as } \Delta t \neq 0
\]

\includegraphics[max width=\textwidth, alt={}, center]{9f3c5a8e-ac5d-42f4-8eb6-ea2af623f2d8-027_396_567_571_804}\\
1.34 According to the problem\\
(a) \(\frac{d y}{d t}=v_{0}\) or \(d y=v_{0} d t\)

Integrating \(\quad \int_{0}^{y} d y=v_{0} \int_{0}^{t} d t\) or \(y=v_{0} t\)

And also we have \(\frac{d x}{d t}=a y\) or \(d x=a y d t=a v_{0} t d t \quad\) (using 1)\\
So, \(\quad \int_{0}^{x} d x=a v_{0} \int_{0}^{t} t d t\), or, \(x=\frac{1}{2} a v_{0} t^{2}=\frac{1}{2} \frac{a y^{2}}{v_{0}}\) (using 1)\\
(b) According to the problem


\begin{equation*}
v_{y}=v_{0} \text { and } v_{x}=a y \tag{2}
\end{equation*}


So,

\[
v=\sqrt{v_{x}^{2}+v_{y}^{2}}=\sqrt{v_{0}^{2}+a^{2} y^{2}}
\]

Therefore

\[
w_{t}=\frac{d v}{d t}=\frac{a^{2} y}{\sqrt{v_{0}^{2}+a y^{2}}} \frac{d y}{d t}=\frac{a^{2} y}{\sqrt{1+\left(a y / v_{0}\right)^{2}}}
\]

Diff. Eq. (2) with respect to time.

\[
\frac{d v_{y}}{d t}=w_{y}=0 \text { and } \frac{d v_{x}}{d t}=w_{x}=a \frac{d y}{d t}=a v_{0}
\]

So,

\[
w=\left|w_{x}\right|=a v_{0}
\]

Hence \(w_{n}=\sqrt{w^{2}-w_{t}^{2}}=\sqrt{a^{2} v_{0}^{2}-\frac{a^{4} y^{2}}{1+\left(a y / v_{0}\right)^{2}}}=\frac{a v_{0}}{\sqrt{1+\left(a y / v_{0}\right)^{2}}}\)\\
1.35 (a) The velocity vector of the particle

\[
\vec{v}=a \vec{i}+b x \vec{j}
\]

So,


\begin{equation*}
\frac{d x}{d t}=a: \frac{d y}{d t}=b x \tag{1}
\end{equation*}


From (1)


\begin{equation*}
\int_{0}^{x} d x=a \int_{0}^{t} d t \text { or, } x=a t \tag{2}
\end{equation*}


And

\[
d y=b x d t=b a t d t
\]

Integrating


\begin{equation*}
\int_{0}^{y} d y=a b \int_{0}^{t} t d t \text { or, } y=\frac{1}{2} a b t^{2} \tag{3}
\end{equation*}


From Eqs. (2) and (3), we get, \(y=\frac{b}{2 a} x^{2}\)\\
(b) The curvature radius of trajectory \(y(x)\) is :


\begin{equation*}
R=\frac{\left[1+(d y / d x)^{2}\right]^{\frac{3}{2}}}{\left|d^{2} y / d x^{2}\right|} \tag{5}
\end{equation*}


Let us differentiate the path Eq. \(y=\frac{b}{2 a} x^{2}\) with respect to \(x\),


\begin{equation*}
\frac{d y}{d x}=\frac{b}{a} x \text { and } \frac{d^{2} y}{d x^{2}}=\frac{b}{a} \tag{6}
\end{equation*}


From Eqs. (5) and (6), the sought curvature radius :

\[
R=\frac{a}{b}\left[1+\left(\frac{b}{a} x\right)^{2}\right]^{\frac{3}{2}}
\]

1.36 In accordance with the problem

\[
w_{t}=\vec{a} \cdot \vec{\tau}
\]

But

\[
w_{t}=\frac{v d v}{d s} \text { or } v d v=w_{t} d s
\]

So,

\[
v d v=(\vec{a} \cdot \vec{\tau}) d s=\vec{a} \cdot d \vec{r}
\]

or, \(\quad v d v=a \vec{i} \cdot d \vec{r}=a d x\) (because \(\vec{a}\) is directed towards the \(x\)-axis)

So,

\[
\int_{0}^{v} v d v=a \int_{0}^{x} d x
\]

Hence

\[
v^{2}=2 a x \text { or, } v=\sqrt{2 a x}
\]

1.37 The velocity of the particle \(v=a t\)

So,


\begin{equation*}
\frac{d v}{d t}=w_{t}=a \tag{1}
\end{equation*}


And


\begin{equation*}
w_{n}=\frac{v^{2}}{R}=\frac{a^{2} t^{2}}{R}(\text { using } v=a t) \tag{2}
\end{equation*}


From

\[
s=\int v d t
\]

\[
2 \pi R \eta=a \int_{0}^{t} v d t=\frac{1}{2} a t^{2}
\]

So,


\begin{equation*}
\frac{4 \pi \eta}{a}=\frac{t^{2}}{R} \tag{3}
\end{equation*}


From Eqs. (2) and (3) \(w_{n}=4 \pi a \eta\)\\
Hence

\[
\begin{aligned}
w & =\sqrt{w_{t}^{2}+w_{n}^{2}} \\
& =\sqrt{a^{2}+(4 \pi a \eta)^{2}}=a \sqrt{1+16 \pi^{2} \eta^{2}}=0.8 \mathrm{~m} / \mathrm{s}^{2}
\end{aligned}
\]

1.38 According to the problem

For \(v(t)\),

\[
\begin{aligned}
& \left|w_{t}\right|=\left|w_{n}\right| \\
& \frac{-d v}{d t}=\frac{v^{2}}{R}
\end{aligned}
\]

Integrating this equation from \(v_{0} \leq v \leq v\) and \(0 \leq t \leq t\)

\[
-\int_{v_{0}}^{v} \frac{d v}{v^{2}}=\frac{1}{R} \int_{0}^{t} d t \text { or, } v=\frac{v_{0}}{\left(1+\frac{v_{0} t}{R}\right)}
\]

N Jw for \(v(s),-\frac{v d v}{d s}=\frac{v^{2}}{R}\), Integrating this equation from \(v_{0} \leq v \leq v\) and \(0 \leq s \leq s\)

S0,

\[
\int_{v_{0}}^{v} \frac{d v}{v}=-\frac{1}{R} \int_{0}^{s} d s \text { or, } \ln \frac{v}{v_{0}}=-\frac{s}{R}
\]

Hence


\begin{equation*}
v=v_{0} e^{-s / R} \tag{2}
\end{equation*}


(b) The normal acceleration of the point

\[
\left.w_{n}=\frac{v^{2}}{R}=\frac{v^{2} e^{-2 s / R}}{R} \text { (using } 2\right)
\]

And as accordance with the problem

\[
\left|w_{t}\right|=\left|w_{n}\right| \text { and } w_{t} \hat{u}_{t} \perp w_{n} \hat{u}_{n}
\]

so,

\[
w=\sqrt{2} w_{n}=\sqrt{2} \frac{v_{0}^{2}}{R} e^{-2 s / R}=\sqrt{2} \frac{v^{2}}{R}
\]

1.39 From the equation

\[
\begin{aligned}
v & =a \sqrt{s} \\
w_{t} & =\frac{d v}{d t}=\frac{a}{2 \sqrt{s}} \frac{d s}{d t}=\frac{a}{2 \sqrt{s}} a \sqrt{s}=\frac{a^{2}}{2}, \text { and } \\
w_{n} & =\frac{v^{2}}{R}=\frac{a^{2} s}{R}
\end{aligned}
\]

As \(w_{t}\) is a positive constant, the speed of the particle increases with time, and the tangential acceleration vector and velocity vector coincides in direction.\\
Hence the angle between \(\vec{v}\) and \(\vec{w}\) is equal to between \(w_{t} \hat{w}_{t}\) an \(\vec{w}\), and \(\alpha\) can be found by means of the formula \(: \tan \alpha=\frac{\left|w_{n}\right|}{\left|w_{t}\right|}=\frac{a^{2} s / R}{a^{2} / 2}=\frac{2 s}{R}\)\\
1.40 From the equation


\begin{align*}
l & =a \sin \omega t \\
\frac{d l}{d t} & =v=a \omega \cos \omega t \\
\text { So, } \quad w_{t} & =\frac{d v}{d t}=-a \omega^{2} \sin \omega t, \text { and }  \tag{1}\\
w_{n} & =\frac{v^{2}}{R}=\frac{a^{2} \omega^{2} \cos ^{2} \omega t}{R} \tag{2}
\end{align*}


(a) At the point \(l=0, \sin \omega t=0\) and \(\cos \omega t= \pm 1\) so, \(\omega t=0, \pi\) etc.

Hence

\[
w=w_{n}=\frac{a^{2} \omega^{2}}{R}
\]

Similarly at \(l= \pm a, \sin \omega t= \pm 1\) and \(\cos \omega t=0\), so, \(w_{n}=0\)\\
Hence

\[
w=\left|w_{t}\right|=a \omega^{2}
\]

1.41 As \(w_{t}=a\) and at \(t=0\), the point is at rest

So,


\begin{equation*}
v(t) \text { and } s(t) \text { are, } v=a t \text { and } s=\frac{1}{2} a t^{2} \tag{1}
\end{equation*}


Let \(R\) be the curvature radius, then

\[
\left.w_{n}=\frac{v^{2}}{R}=\frac{a^{2} t^{2}}{R}=\frac{2 a s}{R} \text { (using } 1\right)
\]

But according to the problem

\[
w_{n}=b t^{4}
\]

So, \(b t^{4}=\frac{a^{2} t^{2}}{R}\) or, \(R=\frac{a^{2}}{b t^{2}}=\frac{a^{2}}{2 b s}\) (using 1)

Therefore \(\quad w=\sqrt{w_{t}^{2}+w_{n}^{2}}=\sqrt{a^{2}+(2 a s / R)^{2}}=\sqrt{a^{2}+\left(4 b s^{2} / a^{2}\right)^{2}}\) (using 2)\\
Hence

\[
w=a \sqrt{1+\left(4 b s^{2} / a^{3}\right)^{2}}
\]

1.42 (a) Let us differentiate twice the path equation \(y(x)\) with respect to time.

\[
\frac{d y}{d t}=2 a x \frac{d x}{d t} ; \frac{d^{2} y}{d t^{2}}=2 a\left[\left(\frac{d x}{d t}\right)^{2}+x \frac{d^{2} x}{d t^{2}}\right]
\]

Since the particle moves uniformly, its acceleration at all points of the path is normal and at the point \(x=0\) it coincides with the direction of derivative \(d^{2} y / d t^{2}\). Keeping in mind that at the point \(x=0,\left|\frac{d x}{d t}\right|=v\),

We get

\[
w=\left|\frac{d^{2} y}{d t^{2}}\right|_{x=0}=2 a v^{2}=w_{n}
\]

So,

\[
w_{n}=2 a v^{2}=\frac{v^{2}}{R}, \text { or } R=\frac{1}{2 a}
\]

Note that we can also calculate it from the formula of problem (1.35 b)\\
(b) Differentiating the equation of the trajectory with respect to time we see that


\begin{equation*}
b^{2} x \frac{d x}{d t}+a^{2} y \frac{d y}{d t}=0 \tag{1}
\end{equation*}


which implies that the vector \(\left(b^{2} x \vec{i}+a^{2} y \vec{j}\right)\) is normal to the velocity vector \(\vec{v}=\frac{d x}{d t} \vec{i}+\frac{d y}{d t} \vec{j}\) which, of course, is along the tangent. Thus the former vactor is along the normal and the normal component of acceleration is clearly

\[
w_{n}=\frac{b^{2} x \frac{d^{2} x}{d t^{2}}+a^{2} y \frac{d^{2} y}{d t^{2}}}{\left(b^{4} x^{2}+a^{4} y^{2}\right)^{1 / 2}}
\]

on using \(\quad w_{n}=\vec{w} \cdot \vec{n} /|\vec{n}|\). At \(x=0, y= \pm b\) and so at \(x=0\)

\[
w_{n}= \pm\left.\frac{d^{2} y}{d t^{2}}\right|_{x=0}
\]

Differentiating (1)

\[
b^{2}\left(\frac{d x}{d t}\right)^{2}+b^{2} x\left(\frac{d^{2} x}{d t^{2}}\right)+a^{2}\left(\frac{d y}{d t}\right)^{2}+a^{2} y\left(\frac{d^{2} y}{d t^{2}}\right)=0
\]

Also from (1)

\[
\frac{d y}{d t}=0 \text { at } x=0
\]

So \(\quad\left(\frac{d x}{d t}\right)= \pm v(\) since tangential velocity is constant \(=v)\)

\[
\text { Thus }\left(\frac{d^{2} y}{d t^{2}}\right)= \pm \frac{b}{a^{2}} v^{2}
\]

and

\[
\left|w_{n}\right|=\frac{b v^{2}}{a^{2}}=\frac{v^{2}}{R}
\]

This gives \(R=a^{2} / b\).\\
1.43 Let us fix the co-ordinate system at the point \(O\) as shown in the figure, such that the radius vector \(\vec{r}\) of point \(A\) makes an angle \(\theta\) with \(x\) axis at the moment shown.\\
Note that the radius vector of the particle \(A\) rotates clockwise and we here take line \(o x\) as reference line, so in this case obviously the angular velocity \(\omega=\left(-\frac{d \theta}{d t}\right)\) taking anticlockwise sense of angular displacement as positive.\\
Also from the geometry of the triangle \(O A C \frac{R}{\sin \theta}=\frac{r}{\sin (\pi-2 \theta)}\) or , \(r=2 R \cos \theta\).\\
Let us write,\\
\includegraphics[max width=\textwidth, alt={}, center]{9f3c5a8e-ac5d-42f4-8eb6-ea2af623f2d8-032_443_578_248_772}\\
\(\vec{r}=r \cos \theta_{\vec{i}}+r \sin \theta_{\vec{j}}=2 R \cos ^{2} \theta_{\vec{i}+} R \sin 2 \theta_{\vec{j}}\)\\
Differentiating with respect to time.\\
\(\frac{\overrightarrow{d r}}{d t}\) or \(\vec{v}=2 R 2 \cos \theta(-\sin \theta) \frac{d \theta}{d t} \vec{i}+2 R \cos 2 \theta \frac{d \theta}{d t} \vec{j}\)\\
or, \(\vec{v}=2 R\left(\frac{-d \theta}{d t}\right)[\sin 2 \theta \vec{i}-\cos 2 \theta \vec{j}]\)\\
or, \(\vec{v}=2 R \omega\left(\sin 2 \theta \vec{i}-\cos ^{2} \theta \vec{j}\right)\)\\
So,

\[
|\vec{v}| \text { or } v=2 \omega R=0.4 \mathrm{~m} / \mathrm{s} .
\]

As \(\omega\) is constant, \(v\) is also constant and \(w_{t}=\frac{d v}{d t}=0\),\\
So, \(\quad w=w_{n}=\frac{v^{2}}{R}=\frac{(2 \omega R)^{2}}{R}=4 \omega^{2} R=0.32 \mathrm{~m} / \mathrm{s}^{2}\)\\
Alternate : From the Fig. the angular velocity of the point \(A\), with respect to centre of the circle \(\boldsymbol{C}\) becomes

\[
-\frac{d(2 \theta)}{d t}=2\left(\frac{-d \theta}{d t}\right)=2 \omega=\text { constant }
\]

Thus we have the problem of finding the velocity and acceleration of a particle moving along a circle of radius \(R\) with constant angular velocity \(2 \omega\).

Hence

\[
\begin{gathered}
v=2 \omega R \text { and } \\
w=w_{n}=\frac{v^{2}}{R}=\frac{(2 \omega R)^{2}}{R}=4 \omega^{2} R
\end{gathered}
\]

1.44 Differentiating \(\varphi(t)\) with respect to time


\begin{equation*}
\frac{d \varphi}{d t}=\omega_{z}=2 a t \tag{1}
\end{equation*}


For fixed axis rotation, the speed of the point A :


\begin{equation*}
v=\omega R=2 a t R \text { or } R=\frac{v}{2 a t} \tag{2}
\end{equation*}


Differentiating with respect to time

But

\[
\begin{aligned}
w_{t} & =\frac{d v}{d t}=2 a R=\frac{v}{t},(\text { using } 1) \\
w_{n} & =\frac{v^{2}}{R}=\frac{v^{2}}{v / 2 a t}=2 a t v \quad(\text { using } 2) \\
w & =\sqrt{w_{t}^{2}+w_{n}^{2}}=\sqrt{(v / t)^{2}+(2 a t v)^{2}} \\
& =\frac{v}{t} \sqrt{1+4 a^{2} t^{4}}
\end{aligned}
\]

1.45 The shell acquires a constant angular acceleration at the same time as it accelerates linearly. The two are related by (assuming both are constant)

\[
\frac{w}{l}=\frac{\beta}{2 \pi n}
\]

Where \(\boldsymbol{w}=\) linear acceleration and \(\boldsymbol{\beta}=\) angular acceleration\\
Then, \(\omega=\sqrt{2 \cdot \beta 2 \pi n}=\sqrt{2 \cdot \frac{w}{l}(2 \pi n)^{2}}\)\\
But \(v^{2}=2 w l\), hence finally

\[
\omega=\frac{2 \pi n v}{l}
\]

1.46 Let us take the rotation axis as \(z\)-axis whose positive direction is associated with the positive direction of the cordinate \(\varphi\), the rotation angle, in accordance with the right-hand screw rule (Fig.)\\
(a) Defferentiating \(\varphi(t)\) with respect to time.


\begin{align*}
\frac{d \varphi}{d t} & =a-3 b t^{2}=\omega_{z}  \tag{1}\\
\frac{d^{2} \varphi}{d t^{2}} & =\frac{d \omega_{z}}{d t}=\beta_{z}=-6 b t \tag{2}
\end{align*}


From (1) the solid comes to stop at \(\Delta t=t=\sqrt{\frac{a}{3 b}}\)\\
The angular velocity \(\omega=a-3 b t^{2}\), for \(0 \leq t \leq \sqrt{a / 3 b}\)

So, \(\langle\omega\rangle=\frac{\int \omega d t}{\int d t}=\frac{\int_{0}^{\sqrt{a / 3 b}}\left(a-3 b t^{2}\right) d t}{\sqrt{a / 3 b}}=\left[a t-b t^{3}\right]_{0}^{\sqrt{a / 3 b}} / \sqrt{a / 3 b}=2 a / 3\)\\
Similarly \(\beta=\left|\beta_{z}\right|=6 b t\) for all values of \(t\).

So,

\[
<\beta>=\frac{\int \beta d t}{\int d t}=\frac{\int_{0}^{\sqrt{a / 3 b}} 6 b t d t}{\sqrt{a / 3 b}}=\sqrt{3 a b}
\]

(b) From Eq. (2) \(\beta_{z}=-6 b t\)

So,

\[
\left(\beta_{z}\right)_{t}=\sqrt{a / 3 b}=-6 b \sqrt{\frac{a}{3 b}}=-2 \sqrt{a b}
\]

Hence

\[
\beta=\left|\left(\beta_{z}\right)_{t-\sqrt{a / 3 b}}\right|=2 \sqrt{3 a b}
\]

1.47 Angle \(\alpha\) is related with \(\left|w_{t}\right|\) and \(w_{n}\) by means of the fomula :


\begin{equation*}
\tan \alpha=\frac{w_{n}}{\left|w_{t}\right|}, \text { where } w_{n}=\omega^{2} R \text { and }\left|w_{t}\right|=\beta R \tag{1}
\end{equation*}


where \(R\) is the radius of the circle which an arbitrary point of the body circumscribes.\\
From the given equation \(\beta=\frac{d \omega}{d t}=a t\) (here \(\beta=\frac{d \omega}{d t}\), as \(\beta\) is positive for all values of \(t\) )\\
Integrating within the limit \(\int_{0}^{\infty} d \omega=a \int_{0}^{t} t d t\) or , \(\omega=\frac{1}{2} a t^{2}\)

So,

\[
w_{n}=\omega^{2} R=\left(\frac{a t^{2}}{2}\right)^{2} R=\frac{a^{2} t^{4}}{4} R
\]

and

\[
\left|w_{t}\right|=\beta R=a t R
\]

Putting the values of \(\left|w_{t}\right|\) and \(w_{n}\) in Eq. (1), we get,

\[
\tan \alpha=\frac{a^{2} t^{4} R / 4}{a t R}=\frac{a t^{3}}{4} \text { or , } t=\left[\left(\frac{4}{a}\right) \tan \alpha\right]^{1 / 3}
\]

1.48 In accordance with the problem, \(\beta_{z}<0\)

Thus \(-\frac{d \omega}{d t}=k \sqrt{\omega}\), where \(k\) is proportionality constant\\
or,


\begin{equation*}
-\int_{0_{0}}^{\infty} \frac{d \omega}{\sqrt{\omega}}=k \int_{0}^{t} d t \text { or, } \sqrt{\omega}=\sqrt{\omega_{0}}-\frac{k t}{2} \tag{1}
\end{equation*}


When \(\omega=0\), total time of rotation \(t=\tau=\frac{2 \sqrt{\omega_{0}}}{k}\)

\[
2 \sqrt{\omega_{0}} / k
\]

Average angular velocity \(<\omega>=\frac{\int \omega d t}{\int d t}=\frac{\int_{0}\left(\omega_{0}+\frac{k^{2} t^{2}}{4}-k t \sqrt{\omega_{0}}\right) d t}{2 \sqrt{\omega_{0}} / k}\)\\
Hence \(<\omega>=\left[\omega_{0} t+\frac{k^{2} t^{3}}{12}-\frac{k}{2} \sqrt{\omega_{0}} t^{2}\right]_{0}^{2 \sqrt{\omega_{0}} / k} / 2 \frac{\sqrt{\omega_{0}}}{k}=\omega_{0} / 3\)\\
1.49 We have \(\omega=\omega_{0}-a \varphi=\frac{d \varphi}{d t}\)

Integratin this Eq. within its limit for \((\varphi) t\)

\[
\int_{0}^{\varphi} \frac{d \varphi}{\omega_{0}-k \varphi}=\int_{0}^{t} d t \text { or, } \ln \frac{\omega_{0}-k \varphi}{\omega_{0}}=-k t
\]

Hence


\begin{equation*}
\varphi=\frac{\omega_{0}}{k}\left(1-e^{-k t}\right) \tag{1}
\end{equation*}


(b) From the Eq., \(\omega=\omega_{0}-k \varphi\) and Eq. (1) or by differentiating Eq. (1)

\[
\omega=\omega_{0} e^{-k t}
\]

1.50 Let us choose the positive direction of \(z\)-axis (stationary rotation axis) along the vector \(\beta_{0}\). In accordance with the equation

\[
\frac{d \omega_{z}}{d t}=\beta_{z} \quad \text { or } \quad \omega_{z} \frac{d \omega_{z}}{d \varphi}=\beta_{z}
\]

or, \(\omega_{z} d \omega_{z}=\beta_{z} d \varphi=\beta \cos \varphi d \varphi\),\\
Integrating this Eq. within its limit for \(\omega_{z}(\varphi)\)\\
or, \(\quad \int_{0}^{\omega_{z}} d \omega_{z}=\beta_{0} \int_{0}^{\varphi} \cos \varphi d \varphi\)\\
\includegraphics[max width=\textwidth, alt={}, center]{9f3c5a8e-ac5d-42f4-8eb6-ea2af623f2d8-035_439_610_1257_795}\\
or, \(\quad \frac{\omega_{z}^{2}}{2}=\beta_{0} \sin \varphi\)\\
Hence \(\quad \omega_{z}= \pm \sqrt{2 \beta_{0} \sin \varphi}\)\\
The plot \(\omega_{z}(\varphi)\) is shown in the Fig. It can be seen that as the angle \(\varphi\) grows, the vector \(\vec{\omega}\) first increases, coinciding with the direction of the vector \(\overrightarrow{\beta_{0}}\left(\omega_{z}>0\right)\), reaches the maximum at \(\varphi=\varphi / 2\), then starts decreasing and finally turns into zero at \(\varphi=\pi\). After that the body starts rotating in the opposite direction in a similar fashion ( \(\omega_{z}<0\) ). As a result, the body will oscillate about the position \(\varphi=\varphi / 2\) with an amplitude equal to \(\pi / 2\).\\
1.51 Rotating disc moves along the \(x\)-axis, in plane motion in \(x-y\) plane. Plane motion of a solid can be imagined to be in pure rotation about a point (say \(I\) ) at a certain instant known as instantaneous centre of rotation. The instantaneous axis whose positive sense is directed along \(\vec{\omega}\) of the solid and which passes through the point \(I\), is known as instantaneous axis of rotation.\\
Therefore the velocity vector of an arbitrary point \((P)\) of the solid can be represented as :


\begin{equation*}
\overrightarrow{v_{P}}=\vec{\omega} \times \overrightarrow{r_{P I}} \tag{1}
\end{equation*}


On the basis of Eq. (1) for the C. M. (C) of the disc


\begin{equation*}
\overrightarrow{v_{c}}=\vec{\omega} \times \overrightarrow{r_{c I}} \tag{2}
\end{equation*}


According to the problem \(\vec{v}_{c} \uparrow \uparrow \vec{i}\) and \(\vec{\omega} \uparrow \uparrow \vec{k}\) i.e. \(\vec{\omega} \perp x-y\) plane, so to satisy the Eqn. (2) \(\overrightarrow{r_{C I}}\) is directed along \((-\vec{j})\). Hence point \(I\) is at a distance \(r_{C I}=y\), above the centre of the disc along \(y\)-axis. Using all these facts in Eq. (2), we get\\
\includegraphics[max width=\textwidth, alt={}, center]{9f3c5a8e-ac5d-42f4-8eb6-ea2af623f2d8-036_456_558_513_793}


\begin{equation*}
v_{C}=\omega y \text { or } y=\frac{v_{c}}{\omega} \tag{3}
\end{equation*}


(a) From the angular kinematical equation


\begin{gather*}
\omega_{z}=\omega_{0 z}+\beta_{z} t  \tag{4}\\
\omega=\beta t .
\end{gather*}


On the other hand \(x=v t\), (where \(x\) is the \(x\) coordinate of the C.M.)\\
or,


\begin{equation*}
t=\frac{x}{v} \tag{5}
\end{equation*}


From Eqs. (4) and (5), \(\omega=\frac{\beta x}{v}\)\\
Using this value of \(\omega\) in Eq. (3) we get \(\dot{y}=\frac{\nu_{c}}{\omega}=\frac{\nu}{\beta x / v}=\frac{\nu^{2}}{\beta x}\) (hyperbola)\\
(b) As centre \(C\) moves with constant acceleration \(\boldsymbol{w}\), with zero initial velocity

So,

\[
x=\frac{1}{2} w t^{2} \text { and } v_{c}=w t
\]

Therefore,

\[
v_{c}=w \sqrt{\frac{2 x}{w}}=\sqrt{2 x w}
\]

Hence

\[
y=\frac{v_{c}}{\omega}=\frac{\sqrt{2 w x}}{\omega} \text { (parabola) }
\]

1.52 The plane motion of a solid can be imagined as the combination of translation of the C.M. and rotation about C.M.\\
So, we may write \(\overrightarrow{v_{A}}=\overrightarrow{v_{C}}+\overrightarrow{v_{A}} C\)


\begin{align*}
& =\overrightarrow{v_{C}}+\vec{\omega} \times \overrightarrow{r_{A C}}  \tag{1}\\
\overrightarrow{w_{A}} & =\overrightarrow{w_{C}}+\overrightarrow{w_{A C}} \\
& =\overrightarrow{w_{C}}+\omega^{2}\left(-\overrightarrow{r_{A C}}\right)+\left(\vec{\beta} \times \overrightarrow{r_{A C}}\right) \tag{2}
\end{align*}


\(\overrightarrow{r_{A C}}\) is the position of vector of \(A\) with respect to \(C\).\\
\includegraphics[max width=\textwidth, alt={}, center]{9f3c5a8e-ac5d-42f4-8eb6-ea2af623f2d8-037_310_458_193_885}

In the problem \(v_{C}=v=\) constant, and the rolling is without slipping i.e., \(v_{C}=v=\omega R\),\\
So, \(w_{C}=0\) and \(\beta=0\). Using these conditions in Eq. (2)

\[
\overrightarrow{w_{A}}=\omega^{2}\left(-\vec{r}_{A C}\right)=\omega^{2} R\left(-\hat{u}_{A C}\right)=\frac{v^{2}}{R}\left(-\hat{u}_{A C}\right)
\]

Here, \(\hat{u}_{A C}\) is the unit vector directed along \(\vec{r}_{A C}\).\\
Hence \(w_{A}=\frac{v^{2}}{R}\) and \(\vec{w}_{A}\) is directed along \(\left(-\hat{u}_{A C}\right)\) or directed toward the centre of the wheel.\\
(b) Let the centre of the wheel move toward right (positive \(\boldsymbol{x}\)-axis) then for pure tolling on the rigid horizontal surface, wheel will have to rotate in clockwise sense. If \(\omega\) be the angular velocity of the wheel then \(\omega=\frac{\nu_{C}}{R}=\frac{\nu}{R}\).\\
Let the point \(A\) touches the horizontal surface at \(t=0\), further let us locate the point \(A\) at \(t=t\),\\
When it makes \(\theta=\omega t\) at the centre of the wheel.\\
From Eqn. (1)

\[
\begin{aligned}
\overrightarrow{v_{A}} & =\overrightarrow{v_{C}}+\vec{\omega} \times \overrightarrow{r_{A C}} \\
& =v \vec{i}+\omega(-\vec{k}) \times[R \cos \theta(-\vec{j})+R \sin \theta(-\vec{i})]
\end{aligned}
\]

or,

\[
\begin{aligned}
\overrightarrow{v_{A}} & =v \vec{i}+\omega R[\cos \omega t(-\vec{i})+\sin \omega t \vec{j}] \\
& =(v-\cos \omega t) \vec{i}+v \sin \omega t \vec{j}(\text { as } v=\omega R)
\end{aligned}
\]

So,

\[
\begin{aligned}
v_{A} & =\sqrt{(v-v \cos \omega t)^{2}+(v \sin \omega t)^{2}} \\
& =v \sqrt{2(1-\cos \omega t)}=2 v \sin (\omega t / 2)
\end{aligned}
\]

Hence distance covered by the point \(A\) during \(T=2 \pi / \omega\)

\[
s=\int v_{A} d t=\int_{0}^{2 \pi / \omega} 2 v \sin (\omega t / 2) d t=\frac{8 v}{\omega}=8 R
\]

1.53 Let us fix the co-ordinate axis \(x y z\) as shown in the fig. As the ball rolls without slipping along the rigid surface so, on the basis of the solution of problem 1.52 :

Thus

\[
\left.\begin{array}{c}
\overrightarrow{v_{0}}=\overrightarrow{v_{c}}+\vec{\omega} \times \overrightarrow{r_{o c}}=0 \\
v_{c}=\omega R \text { and } \vec{\omega} \uparrow \uparrow(-\vec{k}) \text { as } \overrightarrow{v_{c}} \uparrow \uparrow \vec{i} \tag{1}
\end{array}\right\}
\]

and

\[
\left.\begin{array}{c}
\vec{\omega}_{c}+\vec{\beta} \times \vec{r}_{o c}=0 \\
w_{c}=\beta R \text { and } \vec{\beta} \uparrow \uparrow(-\vec{k}) \text { as } \vec{w}_{c} \uparrow \uparrow \vec{i}
\end{array}\right\}
\]

At the position corresponding to that of Fig., in accordance with the problem,\\
and

\[
\begin{aligned}
& w_{c}=w, \text { so } v_{c}=w t \\
& \omega=\frac{v_{c}}{R}=\frac{w t}{R} \text { and } \beta=\frac{w}{R}(\text { using 1) }
\end{aligned}
\]

(a) Let us fix the co-ordinate system with the frame\\
\includegraphics[max width=\textwidth, alt={}, center]{9f3c5a8e-ac5d-42f4-8eb6-ea2af623f2d8-038_363_514_147_801}\\
attached with the rigid surface as shown in the Fig.\\
As point \(O\) is the instantaneous centre of rotation of the ball at the moment shown in Fig.\\
so,

\[
\overrightarrow{v_{0}}=0,
\]

Now,

\[
\begin{aligned}
\overrightarrow{v_{A}} & =\overrightarrow{v_{C}}+\vec{\omega} \times \overrightarrow{r_{A C}} \\
& =v_{C} \vec{i}+\omega(-\vec{k}) \times R(\vec{j})=\left(v_{C}+\omega R\right) \vec{i}
\end{aligned}
\]

So,

\[
\overrightarrow{v_{A}}=2 v_{C} \vec{i}=2 w t \vec{i} \text { (using 1) }
\]

Similalry \(\overrightarrow{v_{B}}=\overrightarrow{v_{C}}+\vec{\omega} \times \overrightarrow{r_{B C}}=v_{C} \vec{i}+\omega(-\vec{k}) \times R(\vec{i})\)

\[
=v_{C} \vec{i}+\omega R(-\vec{j})=v_{C} \vec{i}+v_{C}(-\vec{j})
\]

So, \(v_{B}=\sqrt{2} v_{c}=\sqrt{2} w t\) and \(\overrightarrow{v_{B}}\) is at an angle \(45^{\circ}\) from both \(\vec{i}\) and \(\vec{j}\) (Fig.)\\
(b) \(\vec{w}_{0}=\vec{w}_{C}+\omega^{2}\left(-\overrightarrow{r_{o c}}\right)+\vec{\beta} \times \overrightarrow{r_{O C}}\)


\begin{equation*}
=\omega^{2}\left(-\vec{r}_{O C}\right)=\frac{v_{C}^{2}}{R}\left(-\hat{u}_{O C}\right) \tag{using1}
\end{equation*}


where \(\hat{u}_{O C}\) is the unit vector along \(\vec{r}_{O C}\)\\
so, \(w_{0}=\frac{v_{0}^{2}}{R}=\frac{w^{2} t^{2}}{R}\) (using 2) and \(\vec{w}_{0}\) is directed towards the centre of the ball\\
Now \(\vec{w}_{A}=\vec{w}_{C}+\omega^{2}\left(-\vec{r}_{A C}\right)+\vec{\beta} \times \overrightarrow{r_{A C}}\)\\
\includegraphics[max width=\textwidth, alt={}, center]{9f3c5a8e-ac5d-42f4-8eb6-ea2af623f2d8-038_374_451_1107_732}\\
\(=w \vec{i}+\omega^{2} R(-\vec{j})+\beta(-\vec{k}) \times R \vec{j}\)\\
\(=(w+\beta R) \vec{i}+\frac{v_{c}^{2}}{R}(-\vec{j})\) (using 1) \(=2 w \vec{i}+\frac{w^{2} t^{2}}{R}(-\vec{j})\)

So,

\[
w_{A}=\sqrt{4 w^{2}+\frac{w^{4} t^{4}}{R^{2}}}=2 w \sqrt{1+\left(\frac{w t^{2}}{2 R}\right)^{2}}
\]

Similarly

\[
\begin{aligned}
& \vec{w}_{B}=\vec{w}_{c}+\omega^{2}\left(-\overrightarrow{r_{B C}}\right)+\vec{\beta} \times \overrightarrow{r_{B C}} \\
= & w \vec{i}+\omega^{2} R(-\vec{i})+\beta(-\vec{k}) \times R(\vec{i}) \\
= & \left(v-\frac{v_{C}^{2}}{R}\right) \vec{i}+\beta R(-\vec{j}) \text { (using 1) }
\end{aligned}
\]

\[
=\left(w-\frac{w^{2} t^{2}}{R}\right) \vec{i}+w(-\vec{j}) \text { (using 2) }
\]

So,

\[
w_{B}=\sqrt{\left(w-\frac{w^{2} t^{2}}{R}\right)^{2}+w^{2}}
\]

1.54 Let us draw the kinematical diagram of the rolling cylinder on the basis of the solution of problem 1.53.\\
\includegraphics[max width=\textwidth, alt={}, center]{9f3c5a8e-ac5d-42f4-8eb6-ea2af623f2d8-039_439_504_434_198}\\
\includegraphics[max width=\textwidth, alt={}, center]{9f3c5a8e-ac5d-42f4-8eb6-ea2af623f2d8-039_447_507_425_831}

As, an arbitrary point of the cylinder follows a curve, its normal acceleration and radius of curvature are related by the well known equation

\[
w_{n}=\frac{v^{2}}{R}
\]

so, for point \(\boldsymbol{A}\),

\[
w_{A(n)}=\frac{v_{A}^{2}}{R_{A}}
\]

or,

\[
R_{A}=\frac{4 v_{c}^{2}}{\omega_{r}^{2}}=4 r \text { (because } v_{c}=\omega r, \text { for pure rolling) }
\]

Similarly for point \(B\),\\
or,

\[
\begin{aligned}
w_{B(n)} & =\frac{v_{B}^{2}}{R_{B}} \\
\omega^{2} r \cos 45^{\circ} & =\frac{\left(\sqrt{2} v_{c}\right)^{2}}{R_{B}}, \\
R_{B} & =2 \sqrt{2} \frac{v_{C}^{2}}{\omega^{2} r}=2 \sqrt{2} r
\end{aligned}
\]

1.55 The angular velocity is a vector as infinitesimal rotation commute. Then the relative angular velocity of the body 1 with respect to the body 2 is clearly.

\[
\vec{\omega}_{12}=\vec{\omega}_{1}-\vec{\omega}_{2}
\]

as for relative linear velocity. The relative acceleration of 1 w.r.t. 2 is

\[
\left(\frac{d \overrightarrow{\omega_{1}}}{d t}\right)_{s^{\prime}}
\]

where \(S^{\prime}\) is a frame corotating with the second body and \(S\) is a space fixed frame with origin coinciding with the point of intersection of the two axes,\\
but

\[
\left(\frac{d \vec{\omega}_{1}}{d t}\right)_{s}=\left(\frac{d \vec{\omega}_{1}}{d t}\right)_{s^{\prime}}+\vec{\omega}_{2} \times \vec{\omega}_{1}
\]

Since \(S^{\prime}\) rotates with angular velocity \(\vec{w}_{2}\). However \(\left(\frac{d \vec{\omega}_{1}}{d t}\right)_{S}=0\) as the first body rotates with constant angular velocity in space, thus

\[
\overrightarrow{\beta_{12}}=\overrightarrow{\omega_{1}} \times \overrightarrow{\omega_{2}} .
\]

Note that for any vector \(\vec{b}\), the relation in space forced frame \((k)\) and a frame \(\left(k^{\prime}\right)\) rotating with angular velocity \(\vec{\omega}\) is


\begin{equation*}
\left.\frac{d \vec{b}}{d t}\right|_{K}=\left.\frac{d \vec{b}}{d t}\right|_{K^{\prime}}+\vec{\omega} \times \vec{b} \tag{1}
\end{equation*}


1.56 We have \(\vec{\omega}=a t \vec{i}+b t^{2} \vec{j}\)

So, \(\quad \omega=\sqrt{(a t)^{2}+\left(b t^{2}\right)^{2}}\), thus, \(\left.\omega\right|_{t-10 \mathrm{~s}}=7.81 \mathrm{rad} / \mathrm{s}\)\\
Differentiating Eq. (1) with respect to time


\begin{equation*}
\vec{\beta}=\frac{d \vec{\omega}}{d t}=a \vec{i}+2 b t \vec{j} \tag{2}
\end{equation*}


So,

\[
\beta=\sqrt{a^{2}+(2 b t)^{2}}
\]

and

\[
\left.\beta\right|_{t-10 \mathrm{~s}}=1.3 \mathrm{rad} / \mathrm{s}^{2}
\]

(b)

\[
\cos \alpha=\frac{\vec{\omega} \cdot \vec{\beta}}{\omega \beta}=\frac{\left(a t \vec{i}+b t^{2} \vec{j}\right) \cdot(a \vec{i}+2 b t \vec{j})}{\sqrt{(a t)^{2}+\left(b t^{2}\right)^{2}} \sqrt{a^{2}+(2 b t)^{2}}}
\]

Putting the values of (a) and (b) and taking \(t=10 \mathrm{~s}\), we get

\[
\alpha \cong 17^{\circ}
\]

1.57 (a) Let the axis of the cone ( \(O C\) ) rotates in anticlockwise sense with constant angular velocity \(\vec{\omega}^{\prime}\) and the cone itself about it's own axis \((O C)\) in clockwise sense with angular velocity \(\vec{\omega}_{0}\) (Fig.). Then the resultant angular velocity of the cone.


\begin{equation*}
\vec{\omega}=\vec{\omega}+\overrightarrow{\omega_{0}} \tag{1}
\end{equation*}


As the rolling is pure the magnitudes of the vectors \(\vec{\omega}^{\prime}\) and \(\overrightarrow{\omega_{0}}\) can be easily found from Fig.


\begin{equation*}
\omega^{\prime}=\frac{v}{R \cot \alpha}, \omega_{0}=v / R \tag{2}
\end{equation*}


As \(\vec{\omega}^{\prime} \perp \vec{\omega}_{0}\), from Eq. (1) and (2)\\
\includegraphics[max width=\textwidth, alt={}, center]{9f3c5a8e-ac5d-42f4-8eb6-ea2af623f2d8-040_481_536_1621_801}

\[
\begin{aligned}
\omega= & \sqrt{\omega^{\prime 2}+\omega_{0}^{2}} \\
& \sqrt{\left(\frac{v}{R \cot \alpha}\right)^{2}+\left(\frac{v}{R}\right)^{2}}=\frac{v}{R \cos \alpha}=2.3 \mathrm{rad} / \mathrm{s}
\end{aligned}
\]

(b) Vector of angular acceleration

\[
\vec{\beta}=\frac{d \vec{\omega}}{d t}=\frac{d\left(\vec{\omega}^{\prime}+\vec{\omega}_{0}\right)}{d t}\left(\text { as } \vec{\omega}^{\prime}=\text { constant. }\right)
\]

The vector \(\vec{\omega}_{0}\) which rotates about the \(O O^{\prime}\) axis with the angular velocity \(\vec{\omega}^{\prime}\), retains i magnitude. This increment in the time interval dt is equal to


\begin{equation*}
\left|d \vec{\omega}_{0}\right|=\omega_{0} \cdot \omega^{\prime} d t \text { or in vector form } d \vec{\omega}_{0}=\left(\vec{\omega}^{\prime} \times \vec{\omega}_{0}\right) d t . \tag{\((3\)}
\end{equation*}


Thus \(\quad \vec{\beta}=\vec{\omega}^{\prime} \times \vec{\omega}_{0}\)

The magnitude of the vector \(\vec{\beta}\) is equal to

\[
\beta=\omega^{\prime} \omega_{0}\left(\text { as } \vec{\omega}^{\prime} \perp \vec{\omega}_{0}\right)
\]

So,

\[
\beta=\frac{v}{R \cot \alpha} \frac{v}{R}=\frac{v^{2}}{R^{2}} \tan \alpha=2.3 \mathrm{rad} / \mathrm{s}
\]

1.58 The axis \(A B\) acquired the angular velocity


\begin{equation*}
\vec{\omega}=\vec{\beta}_{o} t \tag{1}
\end{equation*}


Using the facts of the solution of 1.57 , the angular velocity of the body

\[
\begin{aligned}
\omega & =\sqrt{\omega_{o}^{2}+\omega^{\prime 2}} \\
& =\sqrt{\omega_{0}^{2}+\beta_{0}^{2} t^{2}}=0.6 \mathrm{rad} / \mathrm{s}
\end{aligned}
\]

\begin{center}
\includegraphics[max width=\textwidth, alt={}]{9f3c5a8e-ac5d-42f4-8eb6-ea2af623f2d8-041_319_529_1178_806}
\end{center}

And the angular acceleration.

\[
\vec{\beta}=\frac{d \vec{\omega}}{d t}=\frac{d\left(\vec{\omega}+\overrightarrow{\omega_{0}}\right)}{d t}=\frac{d \vec{\omega}}{d t}+\frac{d \overrightarrow{\omega_{0}}}{d t}
\]

But

\[
\frac{d \overrightarrow{\omega_{0}}}{d t}=\overrightarrow{\omega^{\prime}} \times \overrightarrow{\omega_{0}}, \text { and } \frac{d \vec{\omega}}{d t}=\overrightarrow{\beta_{0}} t
\]

So,

\[
\vec{\beta}=\left(\overrightarrow{\beta_{0}} t \times \overrightarrow{\omega_{0}}\right)+\overrightarrow{\beta_{0}}
\]

As, \(\quad \overrightarrow{\beta_{0}} \perp \overrightarrow{\omega_{0}}\) so, \(\beta=\sqrt{\left(\omega_{0} \beta_{0} t\right)^{2}+\beta_{0}^{2}}=\beta_{0} \sqrt{1+\left(\omega_{0} t\right)^{2}}=0.2 \mathrm{rad} / \mathrm{s}^{2}\)

\subsection*{1.2 THE FUNDAMENTAL EQUATION OF DYNAMICS}
1.59 Let \(R\) be the constant upward thurst on the aerostat of mass \(m\), coming down with a constant acceleration w. Applying Newton's second law of motion for the aerostat in projection form


\begin{align*}
F_{y} & =m w_{y} \\
m g-R & =m w \tag{1}
\end{align*}


Now, if \(\Delta m\) be the mass, to be dumped, then using the Eq. \(F_{y}=m w_{y}\)


\begin{equation*}
R-(m-\Delta m) g=(m-\Delta m) w, \tag{2}
\end{equation*}


From Eqs. (1) and (2), we get, \(\quad \Delta m=\frac{2 m w}{g+w}\)\\
1.60 Let us write the fundamental equation of dynamics for all the three blocks in terms of projections, having taken the positive direction of \(x\) and \(y\) axes as shown in Fig; and using the fact that kinematical relation between the accelerations is such that the blocks move with same value of acceleration (say \(w\) )


\begin{align*}
m_{0} g-T_{1} & =m_{0} w  \tag{1}\\
T_{1}-T_{2}-k m_{1} g & =m_{1} w \tag{2}
\end{align*}


and


\begin{equation*}
T_{2}-k m_{2} g=m_{2} w \tag{3}
\end{equation*}


The simultaneous solution of Eqs. (1), (2) and (3) yields,

\[
\begin{aligned}
w & =g \frac{\left[m_{0}-k\left(m_{1}+m_{2}\right)\right]}{m_{0}+m_{1}+m_{2}} \\
T_{2} & =\frac{(1+k) m_{0}}{m_{0}+m_{1}+m_{2}} m_{2} g
\end{aligned}
\]

\includegraphics[max width=\textwidth, alt={}, center]{9f3c5a8e-ac5d-42f4-8eb6-ea2af623f2d8-042_499_619_791_752}\\
and\\
As the block \(m_{0}\) moves down with acceleration \(w\), so in vector form

\[
\vec{w}=\frac{\left[m_{0}-k\left(m_{1}+m_{2}\right)\right] \vec{g}}{m_{0}+m_{1}+m_{2}}
\]

1.61 Let us indicate the positive direction of \(x\)-axis along the incline (Fig.). Figures show the force diagram for the blocks.\\
Let, \(R\) be the force of interaction between the bars and they are obviously sliding down with the same constant acceleration \(w\).\\
\includegraphics[max width=\textwidth, alt={}, center]{9f3c5a8e-ac5d-42f4-8eb6-ea2af623f2d8-042_405_451_1662_158}\\
\includegraphics[max width=\textwidth, alt={}, center]{9f3c5a8e-ac5d-42f4-8eb6-ea2af623f2d8-042_448_516_1616_696}

Newton's second law of motion in projection form along \(x\)-axis for the blocks gives :


\begin{align*}
& m_{1} g \sin \alpha-k_{1} m_{1} g \cos \alpha+R=m_{1} w  \tag{1}\\
& m_{2} g \sin \alpha-R-k_{2} m_{2} g \cos \alpha=m_{2} w \tag{2}
\end{align*}


Solving Eqs. (1) and (2) simultaneously, we get


\begin{gather*}
w=g \sin \alpha-g \cos \alpha \frac{k_{1} m_{1}+k_{2} m_{2}}{m_{1}+m_{2}} \text { and } \\
R=\frac{m_{1} m_{2}\left(k_{1}-k_{2}\right) g \cos \alpha}{m_{1}+m_{2}} \tag{3}
\end{gather*}


(b) when the blocks just slide down the plane, \(\boldsymbol{w} \boldsymbol{=} \mathbf{0}\), so from Eqn. (3)

\[
\begin{gathered}
g \sin \alpha-g \cos \alpha \frac{k_{1} m_{1}+k_{2} m_{2}}{m_{1}+m_{2}}=0 \\
\text { or, }\left(m_{1}+m_{2}\right) \sin \alpha=\left(k_{1} m_{1}+k_{2} m_{2}\right) \cos \alpha \\
\tan \alpha=\frac{\left(k_{1} m_{1}+k_{2} m_{2}\right)}{m_{1}+m_{2}}
\end{gathered}
\]

Hence\\
1.62 Case 1. When the body is launched up :

Let \(k\) be the coefficeint of friction, \(u\) the velocity of projection and \(l\) the distance traversed along the incline. Retarding force on the block \(=m g \sin \alpha+k m g \cos \alpha\) and hence the retardation \(=g \sin \alpha+k g \cos \alpha\).\\
Using the equation of particle kinematics along the incline,

\[
0=u^{2}-2(g \sin \alpha+k g \cos \alpha) l
\]

or,


\begin{equation*}
l=\frac{u^{2}}{2(g \sin \alpha+k g \cos \alpha)} \tag{1}
\end{equation*}


and

\[
0=u-(g \sin \alpha+k g \cos \alpha) t
\]

or,


\begin{equation*}
u=(g \sin \alpha+k g \cos \alpha) t \tag{2}
\end{equation*}


Using (2) in (1) \(l=\frac{1}{2}(g \sin \alpha+k g \cos \alpha) t^{2}\)

Case (2). When the block comes downward, the net force on the body\\
\(=m g \sin \alpha-k m g \cos \alpha\) and hence its acceleration \(=g \sin \alpha-k g \cos \alpha\)\\
Let, \(t\) be the time required then,


\begin{equation*}
l=\frac{1}{2}(g \sin \alpha-k g \cos \alpha) t^{\prime 2} \tag{4}
\end{equation*}


From Eqs. (3) and (4)

\[
\begin{aligned}
& \frac{t^{2}}{t^{\prime 2}}=\frac{\sin \alpha-k \cos \alpha}{\sin \alpha+k \cos \alpha} \\
& \text { But } \frac{t}{t^{\prime}}=\frac{1}{\eta}
\end{aligned}
\]

(according to the question),\\
Hence on solving we get

\[
k=\frac{\left(\eta^{2}-1\right)}{\left(\eta^{2}+1\right)} \tan \alpha=0.16
\]

1.63 At the initial moment, obviously the tension in the thread connecting \(m_{1}\) and \(m_{2}\) equals the weight of \(m_{2}\).\\
(a) For the block \(m_{2}\) to come down or the block \(m_{1}\) to go up, the conditions is

\[
m_{2} g-T \geq 0 \text { and } T-m_{1} g \sin \alpha-f r \geq 0
\]

where \(T\) is tension and \(f_{r}\) is friction which in the limiting case equals \(k m_{1} g \cos \alpha\). Then or \(\quad m_{2} g-m_{1} \sin \alpha>k m_{1} g \cos \alpha\)\\
or \(\quad \frac{m_{2}}{m_{1}}>(k \cos \alpha+\sin \alpha)\)\\
(b) Similarly in the case

\[
m_{1} g \sin \alpha-m_{2} g>f r_{\lim }
\]

or, \(m_{1} g \sin \alpha-m_{2} g>k m_{1} g \cos \alpha\)\\
\includegraphics[max width=\textwidth, alt={}, center]{9f3c5a8e-ac5d-42f4-8eb6-ea2af623f2d8-044_295_485_455_705}\\
or, \(\quad \frac{m_{2}}{m_{1}}<(\sin \alpha-k \cos \alpha)\)\\
(c) For this case, neither kind of motion is possible, and fr need not be limiting.

Hence, \(\quad(k \cos \alpha+\sin \alpha)>\frac{m_{2}}{m_{1}}>(\sin \alpha-k \cos \alpha)\)\\
1.64 From the conditions, obtained in the previous problem, first we will check whether the mass \(m_{2}\) goes up or down.\\
Here, \(m_{2} / m_{1}=\eta>\sin \alpha+k \cos \alpha\), (substituting the values). Hence the mass \(m_{2}\) will come down with an acceleration (say \(w\) ). From the free body diagram of previous problem,


\begin{equation*}
m_{2}-g-T=m_{2} w \tag{1}
\end{equation*}


and


\begin{equation*}
T-m_{1} g \sin \alpha-k m_{1} g \cos \alpha=m_{1} w \tag{2}
\end{equation*}


Adding (1) and (2), we get,

\[
\begin{gathered}
m_{2} g-m_{1} g \sin \alpha-k m_{1} g \cos \alpha=\left(m_{1}+m_{2}\right) w \\
w=\frac{\left(m_{2} / m_{1}-\sin \alpha-k \cos \alpha\right) g}{\left(1+m_{2} / m_{1}\right)}=\frac{(\eta-\sin \alpha-k \cos \alpha) g}{1+\eta}
\end{gathered}
\]

Substituting all the values, \(w=0.048 \mathrm{~g} \sim 0.05 \mathrm{~g}\)\\
As \(m_{2}\) moves down with acceleration of magnitude \(w=0.05 g>0\), thus in vector form acceleration of \(m_{2}\) :

\[
\vec{w}_{2}=\frac{(\eta-\sin \alpha-k \cos \alpha) \vec{g}}{1+\eta}=0.05 \vec{g} .
\]

1.65 Let us write the Newton's second law in projection form along positive \(x\)-axis for the plank and the bar


\begin{equation*}
f r=m_{1} w_{1}, f r=m_{2} w_{2} \tag{1}
\end{equation*}


At the initial moment, fr represents the static friction, and as the force \(F\) grows so does the friction force \(\boldsymbol{f}\), but up to it's limiting value i.e. \(f r=f r_{s(\max )}=k N=k m_{2} g\).\\
Unless this value is reached, both bodies moves as a single body with equal acceleration. But as soon as the force fr reaches the limit, the\\
\includegraphics[max width=\textwidth, alt={}, center]{9f3c5a8e-ac5d-42f4-8eb6-ea2af623f2d8-045_185_441_210_830}\\
bar starts sliding over the plank i.e. \(w_{2} z w_{1}\).\\
Substituting here the values of \(w_{1}\) and \(w_{2}\) taken from Eq. (1) and taking into account that \(f_{r}=k m_{2} g\), we obtain, \(\left(a t-k m_{2} g\right) / m_{2} \geq \frac{k m_{2}}{m_{1}} g\), were the sign " \(=\) " corresponds to the moment \(t=t_{0}\) (say)

Hence,

\[
t_{0}=\frac{k g m_{2}\left(m_{1}+m_{2}\right)}{a m_{1}}
\]

If

\[
\begin{gathered}
t \leq t_{0}, \text { then } w_{1}=\frac{k m_{2} g}{m_{1}}(\text { constant }) . \text { and } \\
w_{2}=\left(a t-k m_{2} g\right) / m_{2}
\end{gathered}
\]

On this basis \(w_{1}(t)\) and \(w_{2}(t)\), plots are as shown in the figure of answersheet.\\
1.66 Let us designate the \(x\)-axis (Fig.) and apply \(F_{x}=m w_{x}\) for body \(A\) :

\[
m g \sin \alpha-k m g \cos \alpha=m w
\]

or,

\[
w=g \sin \alpha-k g \cos \alpha
\]

Now, from kinematical equation :

\[
l \sec \alpha=0+(1 / 2) w t^{2}
\]

or,

\[
\begin{aligned}
t & =\sqrt{2 l \sec \alpha /(\sin \alpha-k \cos \alpha) g} \\
& =\sqrt{2 l /\left(\sin 2 \alpha / 2-k \cos ^{2} \alpha\right) g}
\end{aligned}
\]

(using Eq. (1)).\\
for \(t_{\text {min }}\),

\[
\frac{d\left(\frac{\sin 2 \alpha}{2}-k \cos ^{2} \alpha\right)}{d \alpha}=0
\]

i.e.

\[
\frac{2 \cos 2 \alpha}{2}+2 k \cos \alpha \sin \alpha=0
\]

\includegraphics[max width=\textwidth, alt={}, center]{9f3c5a8e-ac5d-42f4-8eb6-ea2af623f2d8-045_438_215_1294_1079}\\
or,

\[
\tan 2 \alpha=-\frac{1}{k} \Rightarrow \alpha=49^{\circ}
\]

and putting the values of \(\alpha, k\) and \(l\) in Eq. (2) we get \(t_{\text {min }}=1 \mathrm{~s}\).\\
1.67 Let us fix the \(x\)-y co-ordinate system to the wedge, taking the \(x\)-axis up, along the incline and the \(y\)-axis perpendicular to it (Fig.).

Now, we draw the free body diagram for the bar.\\
Let us apply Newton's second law in projection form along \(x\) and \(y\) axis for the bar :


\begin{equation*}
T \cos \beta-m g \sin \alpha-f r=0 \tag{1}
\end{equation*}



\begin{equation*}
T \sin \beta+N-m g \cos \alpha=0 \tag{2}
\end{equation*}


or, \(N=m g \cos \alpha-T \sin \beta\)\\
\includegraphics[max width=\textwidth, alt={}, center]{9f3c5a8e-ac5d-42f4-8eb6-ea2af623f2d8-046_386_377_63_787}

But \(f_{r}=k N\) and using (2) in (1), we get


\begin{equation*}
T=m g \sin \alpha+k m g \cos \alpha /(\cos \beta+k \sin \beta) \tag{3}
\end{equation*}


For \(T_{\text {min }}\) the value of \((\cos \beta+k \sin \beta)\) should be maximum

So,

\[
\frac{d(\cos \beta+k \sin \beta)}{d \beta}=0 \text { or } \tan \beta=k
\]

Putting this value of \(\beta\) in Eq. (3) we get,

\[
T_{\min }=\frac{m g(\sin \alpha+k \cos \alpha)}{1 / \sqrt{1+k^{2}}+k^{2} / \sqrt{1+k^{2}}}=\frac{m g(\sin \alpha+k \cos \alpha)}{\sqrt{1+k^{2}}}
\]

1.68 First of all let us draw the free body diagram for the small body of mass \(m\) and indicate \(x\)-axis along the horizontal plane and \(y\)-axis, perpendicular to it, as shown in the figure.\\
Let the block breaks off the plane at \(t=t_{0}\) i.e. \(N=0\)\\
So, \(\quad N=m g-a t_{0} \sin \alpha=0\)\\
or, \(\quad t_{0}=\frac{m g}{a \sin \alpha}\)

From \(F_{x}=m w_{x}\), for the body under investigation :\\
\(m d \psi / d t=a t \cos \alpha ;\) Integrating within the limits for \(v(t)\)\\
\includegraphics[max width=\textwidth, alt={}, center]{9f3c5a8e-ac5d-42f4-8eb6-ea2af623f2d8-046_352_403_1093_831}

\[
m \int_{0}^{v} d v_{x}=a \cos \alpha \int_{0} t d t \quad \text { (using Eq. 1) }
\]

So,


\begin{equation*}
v=\frac{d s}{d t}=\frac{a \cos \alpha}{2 m} t^{2} \tag{2}
\end{equation*}


Integrating, Eqn. (2) for \(s(t)\)


\begin{equation*}
s=\frac{a \cos \alpha}{2 m} \frac{t^{3}}{3} \tag{3}
\end{equation*}


Using the value of \(t=t_{0}\) from Eq. (1), into Eqs. (2) and (3)

\[
v=\frac{m g^{2} \cos \alpha}{2 a \sin ^{2} \alpha} \text { and } s=\frac{m^{2} g^{3} \cos \alpha}{6 a^{2} \sin ^{3} \alpha}
\]

1.69 Newton's second law of motion in projection form, along horizontal or \(x\)-axis i.e. \(F_{x}=m w_{x}\) gives.

\[
F \cos (a s)=m v \frac{d v}{d s} \quad(a s \alpha=a s)
\]

or, \(\quad F \cos (a s) d s=m v d v\)\\
Integrating, over the limits for \(v(s)\)\\
\includegraphics[max width=\textwidth, alt={}, center]{9f3c5a8e-ac5d-42f4-8eb6-ea2af623f2d8-047_357_463_225_743}

\[
\begin{aligned}
& \quad \frac{F}{m} \int_{0}^{\infty} \cos (a s) d s=\frac{v^{2}}{2} \\
& \text { or } \quad v=\sqrt{\frac{2 F \sin \alpha}{m a}} \\
& =\sqrt{2 g \sin \alpha / 3 a}\left(\text { using } F=\frac{m g}{3}\right)
\end{aligned}
\]

which is the sought relationship.\\
1.70 From the Newton's second law in projection from :

For the bar,


\begin{equation*}
T-2 k m g=(2 m) w \tag{1}
\end{equation*}


For the motor,


\begin{equation*}
T-k m g=m w^{\prime} \tag{2}
\end{equation*}


Now, from the equation of kinematics in the frame of bar or motor :


\begin{equation*}
l=\frac{1}{2}\left(w+w^{\prime}\right) t^{2} \tag{3}
\end{equation*}


From (1), (2) and (3) we get on eliminating \(T\) and \(w^{\prime}\)\\
\includegraphics[max width=\textwidth, alt={}, center]{9f3c5a8e-ac5d-42f4-8eb6-ea2af623f2d8-047_293_885_1112_213}\\
1.71 Let us write Newton's second law in vector from \(\vec{F}=m \vec{w}\), for both the blocks (in the frame of ground).


\begin{align*}
& \vec{T}+m_{1} \vec{g}=m_{1} \overrightarrow{w_{1}}  \tag{1}\\
& \vec{T}+m_{2} \vec{g}=m_{2} \overrightarrow{w_{2}} \tag{2}
\end{align*}


These two equations contain three unknown quantities \(\vec{w}_{1}, \vec{w}_{2}\) and \(\vec{T}\). The third equation is provided by the kinematic relationship between the accelerations :


\begin{equation*}
\overrightarrow{w_{1}}=\overrightarrow{w_{0}}+\overrightarrow{w^{\prime}}, \overrightarrow{w_{2}}=\overrightarrow{w_{0}}-\overrightarrow{w^{\prime}} \tag{3}
\end{equation*}


where \(\vec{w}^{\prime}\) is th acceleration of the mass \(m_{1}\) with respect to the pulley or elevator car. Summing up termwise the left hand and the right-hand sides of these kinematical equations, we get\\
\includegraphics[max width=\textwidth, alt={}, center]{9f3c5a8e-ac5d-42f4-8eb6-ea2af623f2d8-047_571_392_1479_948}


\begin{equation*}
\overrightarrow{w_{1}}+\overrightarrow{w_{2}}=2 \overrightarrow{w_{0}} \tag{4}
\end{equation*}


The simultaneous solution of Eqs•(1), (2) and (4) yields

\[
\overrightarrow{w_{1}}=\frac{\left(m_{1}-m_{2}\right) \vec{g}+2 m_{2} \overrightarrow{w_{0}}}{m_{1+} m_{2}}
\]

Using this result in Eq. (3), we get,

\[
\vec{w}=\frac{m_{1}-m_{2}}{m_{1}+m_{2}}\left(\vec{g}-\vec{w}_{0}\right) \quad \text { and } \quad \vec{T}=\frac{2 m_{1} m_{2}}{m_{1}+m_{2}}\left(\overrightarrow{w_{0}}-\vec{g}\right)
\]

Using the results in Eq. (3) we get \(\vec{w}^{\prime}=\frac{m_{1}-m_{2}}{m_{1}+m_{2}}\left(\vec{g}-\vec{w}_{0}\right)\)\\
(b) obviously the force exerted by the pulley on the celing of the car

\[
\vec{F}=-2 \vec{T}=\frac{4 m_{1} m_{2}}{m_{1}+m_{2}}\left(\vec{g}-\vec{w}_{0}\right)
\]

Note : one could also solve this problem in the frame of elevator car.\\
1.72 Let us write Newton's second law for both, bar 1 and body 2 in terms of projection having taken the positive direction of \(x_{1}\) and \(x_{2}\) as shown in the figure and assuming that body 2 starts sliding, say, upward along the incline


\begin{align*}
T_{1}-m_{1} g \sin \alpha & =m_{1} w_{1}  \tag{1}\\
m_{2} g-T_{2} & =m_{2} w \tag{2}
\end{align*}


For the pulley, moving in vertical direction from the equation \(F_{x}=m w_{x}\)

\[
2 T_{2}-T_{1}=\left(m_{p}\right) w_{1}=0
\]

(as mass of the pulley \(m_{p}=0\) )\\
or


\begin{equation*}
T_{1}=2 T_{2} \tag{3}
\end{equation*}


As the length of the threads are constant, the kinematical relationship of accelerations\\
\includegraphics[max width=\textwidth, alt={}, center]{9f3c5a8e-ac5d-42f4-8eb6-ea2af623f2d8-048_474_611_1007_758}\\
becomes


\begin{equation*}
w=2 w_{1} \tag{4}
\end{equation*}


Simultaneous solutions of all these equations yields :

\[
w=\frac{2 g\left(2 \frac{m_{2}}{m_{1}}-\sin \alpha\right)}{\left(4 \frac{m_{2}}{m_{1}}+1\right)}=\frac{2 g(2 \eta-\sin \alpha)}{(4 \eta+1)}
\]

As \(\eta>1, w\) is directed vertically downward, and hence in vector form

\[
\vec{w}=\frac{2 \vec{g}(2 \eta-\sin \alpha)}{4 \eta+1}
\]

1.73 Let us write Newton's second law for masses \(m_{1}\) and \(m_{2}\) and moving pully in vertical direction along positive \(x\)-axis (Fig.) :


\begin{gather*}
m_{1} g-T=m_{1} w_{1 x}  \tag{1}\\
m_{2} g-T=m_{2} w_{2 x}  \tag{2}\\
T_{1}-2 T=0(\text { as } m=0)
\end{gather*}


or


\begin{equation*}
T_{1}=2 T \tag{3}
\end{equation*}


Again using Newton's second law in projection form for mass \(m_{0}\) along positive \(x_{1}\) direction (Fig.), we get


\begin{equation*}
T_{1}=m_{0} w_{0} \tag{4}
\end{equation*}


The kinematical relationship between the accelerations of masses gives in terms of projection on the \(x\)-axis\\
\includegraphics[max width=\textwidth, alt={}, center]{9f3c5a8e-ac5d-42f4-8eb6-ea2af623f2d8-049_563_584_219_794}


\begin{equation*}
w_{1 x}+w_{2 x}=2 w_{0} \tag{5}
\end{equation*}


Simultaneous solution of the obtained five equations yields :

\[
w_{1}=\frac{\left[4 m_{1} m_{2}+m_{0}\left(m_{1}-m_{2}\right)\right] g}{4 m_{1} m_{2}+m_{0}\left(m_{1}+m_{2}\right)}
\]

In vector form

\[
\overrightarrow{w_{1}}=\frac{\left[4 m_{1} m_{2}+m_{0}\left(m_{1}-m_{2}\right)\right] \vec{g}}{4 m_{1} m_{2}+m_{0}\left(m_{1}+m_{2}\right)}
\]

1.74 As the thread is not tied with \(m\), so if there were no friction between the thread and the ball \(m\), the tension in the thread would be zero and as a result both bodies will have free fall motion. Obviously in the given problem it is the friction force exerted by the ball on the thread, which becomes the tension in the thread. From the condition or language of the problem \(w_{M}>w_{m}\) and as both are directed downward so, relative acceleration of \(M=w_{M}-w_{m}\) and is directed downward. Kinematical equation for the ball in the frame of rod in projection form along upward direction gives :


\begin{equation*}
l=\frac{1}{2}\left(w_{M}-w_{m}\right) t^{2} \tag{1}
\end{equation*}


Newton's second law in projection form along vertically down direction for both, rod and ball gives,


\begin{align*}
& M g-f r=M w_{M}  \tag{2}\\
& m g-f r=m w_{m} \tag{3}
\end{align*}


Multiplying Eq. (2) by \(m\) and Eq. (3) by \(M\) and then subtracting Eq. (3) from (2) and after using Eq. (1) we get

\[
f r=\frac{2 l M m}{(M-m) t^{2}}
\]

\includegraphics[max width=\textwidth, alt={}, center]{9f3c5a8e-ac5d-42f4-8eb6-ea2af623f2d8-049_582_399_1453_931}\\
1.75 Suppose, the ball goes up with accleration \(w_{1}\) and the rod comes down with the acceleration \(w_{2}\). As the length of the thread is constant, \(2 w_{1}=w_{2}\)

From Newton's second law in projection form along vertically upward for the ball and vertically downward for the rod respectively gives,


\begin{equation*}
T-m g=m w_{1} \tag{2}
\end{equation*}


and


\begin{equation*}
M g-T^{\prime}=M w_{2} \tag{3}
\end{equation*}


but


\begin{equation*}
T=2 T \quad \text { (because pulley is massless) } \tag{4}
\end{equation*}


From Eqs. (1), (2), (3) and (4)

\[
w_{1}=\frac{(2 M-m) g}{m+4 M}=\frac{(2-\eta) g}{\eta+4}(\text { in upward direction })
\]

and \(w_{2}=\frac{2(2-\eta) g}{(\eta+4)}\) (downwards)\\
From kinematical equation in projection form, we get

\[
l=\frac{1}{2}\left(w_{1}+w_{2}\right) t^{2}
\]

as, \(w_{1}\) and \(w_{2}\) are in the opposite direction.\\
Putting the values of \(w_{1}\) and \(w_{2}\), the sought time becomes\\
\includegraphics[max width=\textwidth, alt={}, center]{9f3c5a8e-ac5d-42f4-8eb6-ea2af623f2d8-050_535_201_411_1045}

\[
t=\sqrt{2 l(\eta+4) / 3(2-\eta) g}=1.4 \mathrm{~s}
\]

1.76 Using Newton's second law in projection form along \(x\)-axis for the body 1 and along negative \(x\)-axis for the body 2 respectively, we get


\begin{align*}
& m_{1} g-T_{1}=m_{1} w_{1}  \tag{1}\\
& T_{2}-m_{2} g=m_{2} w_{2} \tag{2}
\end{align*}


For the pulley lowering in downward direction from Newton's law along \(x\) axis,

\[
T_{1}-2 T_{2}=0 \text { (as pulley is mass less) }
\]

or,


\begin{equation*}
T_{1}=2 T_{2} \tag{3}
\end{equation*}


As the length of the thread is constant so,


\begin{equation*}
w_{2}=2 w_{1} \tag{4}
\end{equation*}


The simultaneous solution of above equations yields,\\
\includegraphics[max width=\textwidth, alt={}, center]{9f3c5a8e-ac5d-42f4-8eb6-ea2af623f2d8-050_544_355_1089_965}


\begin{equation*}
w_{2}=\frac{2\left(m_{1}-2 m_{2}\right) g}{4 m_{2}+m_{1}}=\frac{2(\eta-2)}{\eta+4}\left(\text { as } \frac{m_{1}}{m_{2}}=\eta\right) \tag{5}
\end{equation*}


Obviously during the time interval in which the body 1 comes to the horizontal floor covering the distance \(h\), the body 2 moves upward the distance \(2 h\). At the moment when the body 2 is at the height \(2 h\) from the floor its velocity is given by the expression :

\[
v_{2}^{2}=2 w_{2}(2 h)=2\left[\frac{2(\eta-2) g}{\eta+4}\right] 2 h=\frac{8 h(\eta-2) g}{\eta+4}
\]

After the body \(m_{1}\) touches the floor the thread becomes slack or the tension in the thread zero, thus as a result body 2 is only under gravity for it's subsequent motion.

Owing to the velocity \(v_{2}\) at that moment or at the height \(2 h\) from the floor, the body 2 further goes up under gravity by the distance,

\[
h^{\prime}=\frac{\nu_{2}^{2}}{2 g}=\frac{4 h(\eta-2)}{\eta+4}
\]

Thus the sought maximum height attained by the body 2 :

\[
H=2 h+h^{\prime}=2 h+\frac{4 h(\eta-2)}{(\eta+4)}=\frac{6 \eta h}{\eta+4}
\]

1.77 Let us draw free body diagram of each body, i.e. of \(\operatorname{rod} A\) and of wedge \(B\) and also draw the kinemetical diagram for accelerations, after analysing the directions of motion of \(A\) and \(B\). Kinematical relationship of accelarations is :


\begin{equation*}
\tan \alpha=\frac{w_{A}}{w_{B}} \tag{1}
\end{equation*}


Let us write Newton's second law for both bodies in terms of projections having taken positive directions of \(y\) and \(x\) axes as shown in the figure.


\begin{equation*}
m_{A} g-N \cos \alpha=m_{A} w_{A} \tag{2}
\end{equation*}


and


\begin{equation*}
N \sin \alpha=m_{B} w_{B} \tag{3}
\end{equation*}


Simultaneous solution of (1), (2) and (3) yields :

\[
\begin{gathered}
w_{A}=\frac{m_{A} g \sin \alpha}{m_{A} \sin \alpha+m_{B} \cot \alpha \cos \alpha}=\frac{g}{\left(1+\eta \cot ^{2} \alpha\right)} \text { and } \\
w_{B}=\frac{w_{A}}{\tan \alpha}=\frac{g}{(\tan \alpha+\eta \cot \alpha)}
\end{gathered}
\]

\includegraphics[max width=\textwidth, alt={}, center]{9f3c5a8e-ac5d-42f4-8eb6-ea2af623f2d8-051_361_437_1168_270}\\
\includegraphics[max width=\textwidth, alt={}, center]{9f3c5a8e-ac5d-42f4-8eb6-ea2af623f2d8-051_354_374_1208_892}

Note : We may also solve this problem using conservation of mechanical energy instead of Newton's second law.\\
1.78 Let us draw free body diagram of each body and fix the coordinate system, as shown in the figure. After analysing the motion of \(M\) and \(m\) on the basis of force diagrams, let us draw the kinematical diagram for accelerations (Fig.).\\
As the length of threads are constant so, \(d s_{m M}=d s_{M}\) and as \(\vec{v}_{m M}\) and \(\vec{v}_{M}\) do not change their directions that why

\[
\begin{aligned}
& \left|\vec{w}_{m M}\right|=\left|\vec{w}_{M}\right|=w \text { (say) and } \\
& \vec{w}_{m M} \uparrow \uparrow \vec{v}_{M} \text { and } \vec{w}_{M} \uparrow \uparrow \vec{v}_{M}
\end{aligned}
\]

\begin{center}
\includegraphics[max width=\textwidth, alt={}]{9f3c5a8e-ac5d-42f4-8eb6-ea2af623f2d8-052_437_826_139_298}
\end{center}

As \(\vec{w}_{m}=\vec{w}_{m M}+\vec{w}_{M}\)\\
so, from the triangle law of vector addition


\begin{equation*}
w_{m}=\sqrt{2} w \tag{1}
\end{equation*}


From the Eq. \(F_{x}=m w_{x}\), for the wedge and block:


\begin{align*}
T-N & =M w,  \tag{2}\\
N & =m w \tag{3}
\end{align*}


and

Now, from the Eq. \(F_{y}=m w_{y}\), for the block


\begin{equation*}
m g-T-k N=m w \tag{4}
\end{equation*}


Simultaneous solution of Eqs. (2), (3) and (4) yields :

\[
w=\frac{m g}{(k m+2 m+M)}=\frac{g}{(k+2+M / m)}
\]

Hence using Eq. (1)

\[
w_{m}=\frac{g \sqrt{2}}{(2+k+M / m)}
\]

1.79 Bodies 1 and 2 will remain at rest with repect to bar \(A\) for \(w_{\min } \leq w \leq w_{\max }\), where \(w_{\min }\) is the sought minimum acceleration of the bar. Beyond these limits there will be a relative motion between bar and the bodies. For \(0 \leq w \leq w_{\text {min }}\), the tendency of body 1 in relation to the bar \(A\) is to move towards right and is in the opposite sense for \(w \geq w_{\max }\). On the basis of above argument the static friction on 2 by \(A\) is directed upward and on 1 by \(A\) is directed towards left for the purpose of calculating \(\boldsymbol{w}_{\text {min }}\).\\
Let us write Newton's second law for bodies 1 and 2 in terms of projection along positive \(x\)-axis (Fig.).


\begin{align*}
T-f r_{1}=m w \text { or, }, f r_{1} & =T-m w  \tag{1}\\
N_{2} & =m w \tag{2}
\end{align*}


As body 2 has no acceleration in vertical direction, so


\begin{equation*}
f r_{2}=m g-T \tag{3}
\end{equation*}


From (1) and (3)


\begin{equation*}
\left(f r_{1}+f r_{2}\right)=m(g-w) \tag{4}
\end{equation*}


But

\[
f r_{1}+f r_{2} \leq k\left(N_{1}+N_{2}\right)
\]

or


\begin{equation*}
f r_{1}+f r_{2} \leq k(m g+m w) \tag{5}
\end{equation*}


\begin{center}
\includegraphics[max width=\textwidth, alt={}]{9f3c5a8e-ac5d-42f4-8eb6-ea2af623f2d8-052_484_511_1558_734}
\end{center}

From (4) and (5)

Hence

\[
\begin{gathered}
m(g-w) \leq \underset{c}{m k}(g+w), \text { or } w \geq \frac{g(1-k)}{(1+k)} \\
w_{\min }=\frac{g(1-k)}{(1+k)}
\end{gathered}
\]

1.80 On the basis of the initial argument of the solution of 1.79 , the tendency of bar 2 with respect to 1 will be to move up along the plane.\\
Let us fix \((x-y)\) coordinate system in the frame of ground as shown in the figure.\\
From second law of motion in projection form along \(y\) and \(x\) axes :\\
\(m g \cos \alpha-N=m w \sin \alpha\)\\
or,


\begin{align*}
& N=m(g \cos \alpha-w \sin \alpha)  \tag{1}\\
& m g \sin \alpha+f r=m w \cos \alpha \\
& f r=m(w \cos \alpha-g \sin \alpha) \tag{2}
\end{align*}


but \(f r \leq k N\), so from (1) and (2)

\[
(w \cos \alpha-g \sin \alpha) \leq k(g \cos \alpha+w \sin \alpha)
\]

or, \(w(\cos \alpha-k \sin \alpha) \leq g(k \cos \alpha+\sin \alpha)\)\\
or, \(\quad w \leq g \frac{(\cos \alpha+\sin \alpha)}{\cos \alpha-k \sin \alpha}\),\\
So, the sought maximum acceleration of the wedge :\\
\includegraphics[max width=\textwidth, alt={}, center]{9f3c5a8e-ac5d-42f4-8eb6-ea2af623f2d8-053_440_404_611_878}\\
\(w_{\text {max }}=\frac{(k \cos \alpha+\sin \alpha) g}{\cos \alpha-k \sin \alpha}=\frac{(k \cot \alpha+1) g}{\cot \alpha-k}\) where \(\cot \alpha>k\)\\
1.81 Let us draw the force diagram of each body, and on this basis we observe that the prism moves towards right say with an acceleration \(w_{1}\) and the bar 2 of mass \(m_{2}\) moves down the plane with respect to 1 , say with acceleration \(\vec{w}_{21}\), then , \(\overrightarrow{w_{2}}=\overrightarrow{w_{21}}+\overrightarrow{w_{1}}\) (Fig.)\\
Let us write Newton's second law for both bodies in projection form along positive \(y_{2}\) and \(x_{1}\) axes as shown in the Fig.\\
\(m_{2} g \cos \alpha-N=m_{2} w_{2\left(y_{2}\right)}=m_{2}\left[w_{21\left(y_{2}\right)}+w_{1\left(y_{2}\right)}\right]=m_{2}\left[0+w_{1} \sin \alpha\right]\)\\
or,


\begin{equation*}
m_{2} g \cos \alpha-N=m_{2} w_{1} \sin \alpha \tag{1}
\end{equation*}


and


\begin{equation*}
N \sin \alpha=m_{1} w_{1} \tag{2}
\end{equation*}


Solving (1) and (2), we get

\[
w_{1}=\frac{m_{2} g \sin \alpha \cos \alpha}{m_{1}+m_{2} \sin ^{2} \alpha}=\frac{g \sin \alpha \cos \alpha}{\left(m_{1} / m_{2}\right)+\sin ^{2} \alpha}
\]

\includegraphics[max width=\textwidth, alt={}, center]{9f3c5a8e-ac5d-42f4-8eb6-ea2af623f2d8-053_298_373_1727_319}\\
\includegraphics[max width=\textwidth, alt={}, center]{9f3c5a8e-ac5d-42f4-8eb6-ea2af623f2d8-053_264_350_1780_708}\\
1.82 To analyse the kinematic relations between the bodies, sketch the force diagram of each body as shown in the figure.\\
On the basis of force diagram, it is obvious that the wedge \(M\) will move towards right and the block will move down along the wedge. As the length of the thread is constant, the distance travelled by the block on the wedge must be equal to the distance travelled by the wedge on the floor. Hence \(d s_{m M}=d s_{M}\) As \(\vec{v}_{m M}\) and \(\vec{v}_{M}\) do not change their directions and acceleration that's why \(\vec{w}_{m} \uparrow \uparrow \vec{v}_{m M}\) and \(\vec{w}_{M} \uparrow \uparrow \vec{v}_{M}\) and \(w_{m M}=w_{M}=w\) (say) and accordingly the diagram of kinematical dependence is shown in figure.\\
\includegraphics[max width=\textwidth, alt={}, center]{9f3c5a8e-ac5d-42f4-8eb6-ea2af623f2d8-054_372_672_471_182}\\
\includegraphics[max width=\textwidth, alt={}, center]{9f3c5a8e-ac5d-42f4-8eb6-ea2af623f2d8-054_384_342_489_973}

As \(\vec{w}_{m}=\vec{w}_{m}+\vec{w}_{M}\), so from triangle law of vector addition.


\begin{equation*}
w_{m}=\sqrt{w_{M}^{2}+w_{m M}^{2}-2 w_{m M} w_{M} \cos \alpha}=w \sqrt{2(1-\cos \alpha)} \tag{1}
\end{equation*}


From \(F_{x}=m w_{x}\), (for the wedge),


\begin{equation*}
T=T \cos \alpha+N \sin \alpha=M w \tag{2}
\end{equation*}


For the bar \(m\) let us fix ( \(x-y\) ) coordinate system in the frame of ground Newton's law in projection form along \(x\) and \(y\) axes (Fig.) gives


\begin{gather*}
m g \sin \alpha-T=m w_{m(x)}=m\left[w_{m M(x)}+w_{M(x)}\right] \\
=m\left[w_{m M}+w_{M} \cos (\pi-\alpha)\right]=m w(1-\cos \alpha)  \tag{3}\\
m g \cos \alpha-N=m w_{m(y)}=m\left[w_{m M(y)}+w_{M(y)}\right]=m[0+w \sin \alpha] \tag{4}
\end{gather*}


Solving the above Eqs. simultaneously, we get

\[
w=\frac{m g \sin \alpha}{M+2 m(1-\cos \alpha)}
\]

Note : We can study the motion of the block \(m\) in the frame of wedge also, alternately we may solve this problem using conservation of mechanical energy.\\
1.83 Let us sketch the diagram for the motion of the particle of mass \(m\) along the circle of radius \(R\) and indicate \(x\) and \(y\) axis, as shown in the figure.\\
(a) For the particle, change in momentum \(\Delta \vec{p}=m v(-\vec{i})-m v(\vec{j})\)\\
so, \(\quad|\Delta \vec{p}|=\sqrt{2} m v\)\\
and time taken in describing quarter of the circle,

\[
\Delta t=\frac{\pi R}{2 v}
\]

Hence, \(\langle\vec{F}\rangle=\frac{|\Delta \vec{p}|}{\Delta t}=\frac{\sqrt{2} m v}{\pi R / 2 v}=\frac{2 \sqrt{2} m v^{2}}{\pi R}\)\\
(b) In this case

\[
\overrightarrow{p_{i}}=0 \text { and } \overrightarrow{p_{f}}=m w_{t} t(-\vec{i}),
\]

so \(\quad|\Delta \vec{p}|=m w_{t} t\)\\
Hence, \(|\langle\vec{F}\rangle|=\frac{|\Delta \vec{p}|}{t}=m w_{t}\)\\
\includegraphics[max width=\textwidth, alt={}, center]{9f3c5a8e-ac5d-42f4-8eb6-ea2af623f2d8-055_361_434_188_893}\\
1.84 While moving in a loop, normal reaction exerted by the flyer on the loop at different points and uncompensated weight if any contribute to the weight of flyer at those points.\\
(a) When the aircraft is at the lowermost point, Newton's second law of motion in projection form \(F_{n}=m w_{n}\) gives

\[
N-m g=\frac{m v^{2}}{R}
\]

or , \(N=m g+\frac{m v^{2}}{R}=2.09 \mathrm{kN}\)\\
(b) When it is at the upper most point, again from \(F_{n}=m w_{n}\) we get

\[
N^{\prime \prime}+m g=\frac{m v^{2}}{R}
\]

\begin{figure}[h]
\begin{center}
  \includegraphics[alt={},max width=\textwidth]{9f3c5a8e-ac5d-42f4-8eb6-ea2af623f2d8-055_406_384_786_900}
\captionsetup{labelformat=empty}
\caption{\(m g\)}
\end{center}
\end{figure}

\[
N^{\prime \prime}=\frac{m v^{2}}{R}-m g=0.7 \mathrm{kN}
\]

(c) When the aircraft is at the middle point of the loop, again from \(F_{n}=m w_{n}\)

\[
N^{\prime}=\frac{m v^{2}}{R}=1.4 \mathrm{kN}
\]

The uncompensated weight is \(m g\). Thus effective weight \(=\sqrt{N^{\prime 2}+m^{2} g^{2}}=1.56 \mathrm{kN}\) acts obliquely.\\
1.85 Let us depict the forces acting on the small sphere \(m\), (at an arbitrary position when the thread makes an angle \(\theta\) from the vertical) and write equation \(\vec{F}=m \vec{w}\) via projection on the unit vectors \(\hat{u}_{t}\) and \(\hat{u}_{n}\). From \(F_{t}=m w_{t}\), we have

\[
\begin{gathered}
m g \sin \theta=m \frac{d v}{d t} \\
=m \frac{v d v}{d s}=m \frac{v d v}{l(-d \theta)}
\end{gathered}
\]

(as vertical is refrence line of angular position)\\
or

\[
v d v=-g l \sin \theta d \theta
\]

Integrating both the sides :

\[
\int_{0}^{v} v d v=-g l \int_{\pi / 2}^{\theta} \sin \theta d \theta
\]

or,

\[
\frac{v^{2}}{2}=g l \cos \theta
\]

Hence \(\frac{v^{2}}{l}=2 g \cos \theta=w_{n}\)\\
(Eq. (1) can be easily obtained by the conservation of mechanical energy).\\
From

\[
F_{n}=m w_{n}
\]

\[
T-m g \cos \theta=\frac{\dot{m} v^{2}}{l}
\]

Using (1) we have\\
\includegraphics[max width=\textwidth, alt={}, center]{9f3c5a8e-ac5d-42f4-8eb6-ea2af623f2d8-056_495_297_385_794}


\begin{equation*}
T=3 m g \cos \theta \tag{2}
\end{equation*}


Again from the Eq. \(F_{t}=m w_{t}\) :\\
\(m g \sin \theta=m w_{t}\) or \(w_{t}=g \sin \theta\)

Hence \(w=\sqrt{w_{t}^{2}+w_{n}^{2}}=\sqrt{(g \sin \theta)^{2}+(2 g \cos \theta)^{2}}\) ( using 1 and 3\()\)

\[
=g \sqrt{1+3 \cos ^{2} \theta}
\]

(b) Vertical component of velocity, \(\boldsymbol{v}_{\boldsymbol{y}}=\boldsymbol{v} \sin \theta\)

So,

\[
v_{y}^{2}=v^{2} \sin ^{2} \theta=2 g l \cos \theta \sin ^{2} \theta(\text { using } 1)
\]

For maximum

\[
v_{y} \text { or } v_{y}^{2}, \frac{d\left(\cos \theta \sin ^{2} \theta\right)}{d \theta}=0
\]

which yields

\[
\cos \theta=\frac{1}{\sqrt{3}}
\]

Therefore from (2) \(T=3 m g \frac{1}{\sqrt{3}}=\sqrt{3} m g\)\\
(c) We have \(\vec{w}=w_{t} \hat{u}_{t}+w_{n} \hat{u}_{n}\) thus \(w_{y}=w_{t(y)}+w_{n(y)}\)

But in accordance with the problem \(w_{y}=0\)\\
So,

\[
w_{t(y)}+w_{n(y)}=0
\]

or,

\[
g \sin \theta \sin \theta+2 g \cos ^{2} \theta(-\cos \theta)=0
\]

or,

\[
\cos \theta=\frac{1}{\sqrt{3}} \quad \text { or, } \quad \theta=54.7^{\circ}
\]

1.86 The ball has only normal acceleration at the lowest position and only tangential acceleration at any of the extreme position. Let \(v\) be the speed of the ball at its lowest position and \(l\) be the length of the thread, then according to the problem


\begin{equation*}
\frac{v^{2}}{l}=g \sin \alpha \tag{1}
\end{equation*}


where \(\alpha\) is the maximum deflection angle\\
From Newton's law in projection form : \(\boldsymbol{F}_{\boldsymbol{t}}=\boldsymbol{m} \boldsymbol{w}_{\boldsymbol{t}}\)

\[
-m g \sin \theta=m v \frac{d v}{l d \theta}
\]

or, \(\quad-g l \sin \theta d \theta=v d v\)\\
On integrating both the sides within their limits.

\[
\begin{array}{r}
-g l \int_{0}^{\alpha} \sin \theta d \theta=\int_{v}^{0} v d v \\
\text { or, } \quad v^{2}=2 g l(1-\cos \alpha) \tag{2}
\end{array}
\]

\begin{center}
\includegraphics[max width=\textwidth, alt={}]{9f3c5a8e-ac5d-42f4-8eb6-ea2af623f2d8-057_476_303_288_971}
\end{center}

Note : Eq. (2) can easily be obtained by the conservation of mechanical energy of the ball in the uniform field of gravity.\\
From Eqs. (1) and (2) with \(\theta=\alpha\)\\
ball in the uniform field of gravity.\\
From Eqs. (1) and (2) with \(\theta=\alpha\)

\[
\begin{array}{ll} 
& 2 g l(1-\cos \alpha)=\lg \cos \alpha \\
\text { or, } & \cos \alpha=\frac{2}{3} \text { so, } \alpha=53^{\circ}
\end{array}
\]

1.87 Let us depict the forces acting on the body \(\underset{\rightarrow}{A}\) (which are the force of gravity \(m \vec{g}\) and the normal reaction \(\vec{N}\) ) and write equation \(\vec{F}=m \vec{w}\) via projection on the unit vectors \(\hat{u}_{t}\) and \(\hat{u}_{n}\) (Fig.)\\
From \(F_{t}=m w_{t}\)

\[
\begin{aligned}
m g \sin \theta & =m \frac{d v}{d t} \\
& =m \frac{v d v}{d s}=m \frac{v d v}{R d \theta}
\end{aligned}
\]

or, \(\quad g R \sin \theta d \theta=v d v\)\\
Integrating both side for obtaining \(v(\theta)\)

\[
\int_{0}^{\theta} g R \sin \theta d \theta=\int_{0}^{v} v d v
\]

or,\\
From \(F_{n}=m w_{n}\)


\begin{equation*}
m g \cos \theta-N=m \frac{v^{2}}{R}, \quad R\left(\frac{\theta}{\hat{u}_{n}}\right)^{N}{\hat{u}}^{-} \tag{1}
\end{equation*}


At the moment the body loses contact with the surface, \(N=0\) and therefore the Eq. (2) becomes


\begin{equation*}
v^{2}=g R \cos \theta \tag{3}
\end{equation*}


where \(v\) and \(\theta\) correspond to the moment when the body loses contact with the surface. Solving Eqs. (1) and (3) we obtain \(\cos \theta=\frac{2}{3}\) or, \(\theta=\cos ^{-1}(2 / 3)\) and \(v=\sqrt{2 g R / 3}\).\\
1.88 At first draw the free body diagram of the device as, shown. The forces, acting on the sleeve are it's weight, acting vertically downward, spring force, along the length of the spring and normal reaction by the rod, perpendicular to its length.\\
Let \(F\) be the spring force, and \(\Delta l\) be the elongation.\\
From, \(F_{n}=m w_{n}\) :


\begin{equation*}
N \sin \theta+F \cos \theta=m \omega^{2} r \tag{1}
\end{equation*}


where \(r \cos \theta=\left(l_{0}+\Delta l\right)\).\\
Similarly from \(F_{t}=m w_{t}\)


\begin{equation*}
N \cos \theta-F \sin \theta=0 \text { or, } N=F \sin \theta / \cos \theta \tag{2}
\end{equation*}


From (1) and (2)

\[
\begin{gathered}
F(\sin \theta / \cos \theta) \cdot \sin \theta+F \cos \theta=m \omega^{2} r \\
=m \omega^{2}\left(l_{0}+\Delta l\right) / \cos \theta
\end{gathered}
\]

On putting \(F=\kappa \Delta l\),\\
\(\kappa \Delta l \sin ^{2} \theta+\kappa \Delta l \cos ^{2} \theta=m \omega^{2}\left(l_{0}+\Delta l\right)\)\\
on solving, we get,\\
\(\Delta l=m \omega^{2} \frac{l_{0}}{\kappa-m \omega^{2}}=\frac{l_{0}}{\left(\kappa / m \omega^{2}-1\right)}\)\\
and it is independent of the direction of rotation.\\
\includegraphics[max width=\textwidth, alt={}, center]{9f3c5a8e-ac5d-42f4-8eb6-ea2af623f2d8-058_410_435_703_808}\\
1.89 According to the question, the cyclist moves along the circular path and the centripetal force is provided by the frictional force. Thus from the equation \(F_{n}=m w_{n}\)

\[
f r=\frac{m v^{2}}{r} \text { or } k m g=\frac{m v^{2}}{r}
\]

or


\begin{equation*}
k_{0}\left(1-\frac{r}{R}\right) g=\frac{v^{2}}{r} \text { or } v^{2}=k_{0}\left(r-r^{2} / R\right) g \tag{1}
\end{equation*}


For \(v_{\max }\), we should have \(\frac{d\left(r-\frac{r^{2}}{R}\right)}{d r}=0\)\\
or, \(\quad 1-\frac{2 r}{R}=0\), so \(r=R / 2\)\\
Hence \(v_{\max }=\frac{1}{2} \sqrt{k_{0} g R}\)\\
1.90 As initial velocity is zero thus


\begin{equation*}
v^{2}=2 w_{t} s \tag{1}
\end{equation*}


As \(w_{t}>0\) the speed of the car increases with time or distance. Till the moment, sliding starts, the static friction provides the required centripetal acceleration to the car.\\
Thus

\[
f r=m w, \text { but } f r \leq k m g
\]

So,

\[
w^{2} \leq k^{2} g^{2} \text { or, } w_{t}^{2}+\frac{v^{2}}{R} \leq k^{2} g^{2}
\]

or,\\
Hence

\[
\begin{gathered}
v^{2} \leq\left(k^{2} g^{2}-w_{t}^{2}\right) R \\
v_{\max }=\sqrt{\left(k^{2} g^{2}-w_{t}^{2}\right) R}
\end{gathered}
\]

so, from Eqn. (1), the sought distance \(s=\frac{v_{\text {max }}^{2}}{2 w_{t}}=\frac{1}{2} \sqrt{\left(\frac{k g}{w_{t}}\right)^{2}-1}=60 \mathrm{~m}\).\\
1.91 Since the car follows a curve, so the maximum velocity at which it can ride without sliding at the point of minimum radius of curvature is the sought velocity and obviously in this case the static friction between the car and the road is limiting.\\
Hence from the equation \(F_{n}=m w\)\\
so


\begin{gather*}
k m g \geq \frac{m v^{2}}{R} \text { or } v \leq \sqrt{k R g} \\
v_{\max }=\sqrt{k R_{\min } g} . \tag{1}
\end{gather*}


We know that, radius of curvature for a curve at any point \((x, y)\) is given as,


\begin{equation*}
R=\left|\frac{\left[1+(d y / d x)^{2}\right]^{3 / 2}}{\left(d^{2} y\right) / d x^{2}}\right| \tag{2}
\end{equation*}


For the given curve,

\[
\frac{d y}{d x}=\frac{a}{\alpha} \cos \left(\frac{x}{\alpha}\right) \text { and } \frac{d^{2} y}{d x^{2}}=\frac{-a}{\alpha^{2}} \sin \frac{x}{\alpha}
\]

Substituting this value in (2) we get,

\[
R=\frac{\left[1+\left(a^{2} / \alpha^{2}\right) \cos ^{2}(x / \alpha)\right]^{3 / 2}}{\left(a / \alpha^{2}\right) \sin (x / \alpha)}
\]

For the minimum \(R, \frac{x}{\alpha}=\frac{\pi}{2}\)\\
and therefore, corresponding radius of curvature


\begin{equation*}
R_{\min }=\frac{\alpha^{2}}{a} \tag{3}
\end{equation*}


Hence from (1) and (2)

\[
v_{\max }=\alpha \sqrt{\mathrm{kg} / \mathrm{a}}
\]

1.92 The sought tensile stress acts on each element of the chain. Hence divide the chain into small, similar elements so that each element may be assumed as a particle. We consider one such element of mass \(d m\), which subtends angle \(d \alpha\) at the centre. The chain moves along a circle of known radius \(R\) with a known angular speed \(\omega\) and certain forces act on it. We have to find one of these forces.\\
From Newton's second law in projection form, \(\boldsymbol{F}_{\boldsymbol{x}} \boldsymbol{=} \boldsymbol{m} \boldsymbol{w}_{\boldsymbol{x}}\) we get\\
\(2 T \sin (d \alpha / 2)-d N \cos \theta=d m \omega^{2} R\)\\
and from \(F_{Z}=m w_{Z}\) we get

\[
d N \sin \theta=g d m
\]

Then putting \(d m=m d \alpha / 2 \pi\) and \(\sin (d \alpha / 2)=d \alpha / 2\) and solving, we get,

\[
T=\frac{m\left(\omega^{2} R+g \cot \theta\right)}{2 \pi}
\]

\includegraphics[max width=\textwidth, alt={}, center]{9f3c5a8e-ac5d-42f4-8eb6-ea2af623f2d8-060_714_1230_47_53}\\
1.93 Let, us consider a small element of the thread and draw free body diagram for this element.\\
(a) Applying Newton's second law of motion in projection form, \(F_{n}=m w_{n}\) for this element, \((T+d T) \sin (d \theta / 2)+T \sin (d \theta / 2)-d N=d m \omega^{2} R=0\)\\
or, \(\quad 2 T \sin (d \theta / 2)=d N\), [negelecting the term \((d T \sin d \theta / 2)\) ]\\
or,


\begin{equation*}
T d \theta=d N, \text { as } \sin \frac{d \theta}{2}=\frac{d \theta}{2} \tag{1}
\end{equation*}


Also,


\begin{equation*}
d f r=k d N=(T+d T)-T=d T \tag{2}
\end{equation*}


From Eqs. (1) and (2),\\
\(k T d \theta=d T\) or \(\frac{d T}{T}=k d \theta\)\\
In this case \(Q=\pi\) so,\\
or, or, \(\ln \frac{T_{2}}{T_{1}}=k \pi\)

So, \(\quad k=\frac{1}{\pi} \ln \frac{T_{2}}{T_{1}}=\frac{1}{\pi} \ln \eta_{0}\)\\
as \(\quad \frac{T_{2}}{T_{1}}=\frac{m_{2} g}{m_{1} g}=\frac{m_{2}}{m_{1}}=\eta_{0}\)\\
\includegraphics[max width=\textwidth, alt={}, center]{9f3c5a8e-ac5d-42f4-8eb6-ea2af623f2d8-060_571_519_1092_786}\\
(b) When \(\frac{m_{2}}{m_{1}}=\eta\), which is greater than \(\eta_{0}\), the blocks will move with same value of acceleration. (say \(w\) ) and clearly \(m_{2}\) moves downward. From Newton's second law in projection form (downward for \(m_{2}\) and upward for \(m_{1}\) ) we get :


\begin{equation*}
m_{2} g-T_{2}=m_{2} w \tag{4}
\end{equation*}


and


\begin{equation*}
T_{1}-m_{1} g=m_{1} w \tag{5}
\end{equation*}


Also


\begin{equation*}
\frac{T_{2}}{T_{1}}=\eta_{0} \tag{6}
\end{equation*}


Simultaneous solution of Eqs. (4), (5) and (6) yields :\\
\(w=\frac{\left(m_{2}-\eta_{0} m_{1}\right) g}{\left(m_{2}+\eta_{0} m_{1}\right)}=\frac{\left(\eta-\eta_{0}\right)}{\left(\eta+\eta_{0}\right)} g\left(\right.\) as \(\left.\frac{m_{2}}{m_{1}}=\eta\right)\)\\
1.94 The force with which the cylinder wall acts on the particle will provide centripetal force necessary for the motion of the particle, and since there is no acceleration acting in the horizontal direction, horizontal component of the velocity will remain constant througout the motion.\\
So

\[
v_{x}=v_{0} \cos \alpha
\]

Using, \(F_{n}=m w_{n}\), for the particle of mass \(m\),\\
\(N=\frac{m v_{x}^{2}}{R}=\frac{m v_{0}^{2} \cos ^{2} \alpha}{R}\),\\
\includegraphics[max width=\textwidth, alt={}, center]{9f3c5a8e-ac5d-42f4-8eb6-ea2af623f2d8-061_534_277_305_937}\\
which is the required normal force.\\
1.95 Obviously the radius vector describing the position of the particle relative to the origin of coordinate is

\[
\vec{r}=x \overrightarrow{i+} y \vec{j}=a \sin \omega t \overrightarrow{i+} b \cos \omega t \vec{j}
\]

Differentiating twice with respect the time :


\begin{equation*}
\vec{w}=\frac{d^{2} \vec{r}}{d t^{2}}=-\omega^{2}(a \sin \omega t \vec{i}+b \cos \omega t \vec{j})=-\omega^{2} \vec{r} \tag{1}
\end{equation*}


Thus

\[
\vec{F}=m \vec{w}=-m \omega^{2} \vec{r}
\]

1.96 (a) We have \(\Delta \vec{p}=\int \vec{F} d t\)


\begin{equation*}
=\int_{0}^{t} m \vec{g} d t=m \overrightarrow{g t} \tag{1}
\end{equation*}


(b) Using the solution of problem 1.28 (b), the total time of motion, \(\tau=-\frac{2\left(\overrightarrow{v_{0}} \cdot \vec{g}\right)}{g^{2}}\) Hence using

\[
\begin{gathered}
t=\tau \text { in (1) } \\
|\Delta \vec{p}|=m g \tau \\
=-2 m\left(\overrightarrow{v_{0} \cdot g}\right) / g\left(\overrightarrow{v_{0}} \cdot \vec{g} \text { is }-\mathrm{ve}\right)
\end{gathered}
\]

1.97 From the equation of the given time dependence force \(\vec{F}=\vec{a} t(\tau-t)\) at \(t=\tau\), the force vanishes,\\
(a) Thus

\[
\overrightarrow{\Delta p}=\vec{p}=\int_{0}^{\tau} \vec{F} d t
\]

or,

\[
\vec{p}=\int_{0}^{\tau} \vec{a} t(\tau-t) d t \frac{\overrightarrow{a \tau^{3}}}{6}
\]

but

\[
\vec{p}=m \vec{v} \text { so } \vec{v}=\frac{\vec{a}^{3}}{6 m}
\]

(b) Again from the equation \(\vec{F}=m \vec{w}\)\\
or,

\[
\begin{gathered}
\overrightarrow{a t}(\tau-t)=m \frac{d \vec{v}}{d t} \\
\vec{a}\left(t \tau-t^{2}\right) d t=m d \vec{v}
\end{gathered}
\]

Integrating within the limits for \(\vec{v}(t)\),\\
or,

Thus

\[
\begin{gathered}
\int_{0}^{t} \vec{a}\left(t \tau-t^{2}\right) d t=m \int_{0}^{\vec{v}} d \vec{v} \\
\vec{v}=\frac{\vec{a}}{m}\left(\frac{\tau t^{2}}{2}-\frac{t^{3}}{3}\right)=\frac{\vec{a}^{2}}{m}\left(\frac{\tau}{2}-\frac{t}{3}\right) \\
v=\frac{a t^{2}}{m}\left(\frac{\tau}{2}-\frac{t}{3}\right) \text { for } t \leq \tau
\end{gathered}
\]

Hence distance covered during the time interval \(t=\tau\),

\[
\begin{gathered}
s=\int_{0}^{\tau} v d t \\
=\int_{0}^{\tau} \frac{a t^{2}}{m}\left(\frac{\tau}{2}-\frac{t}{3}\right) d t=\frac{a}{m} \frac{\tau^{4}}{12}
\end{gathered}
\]

1.98 We have \(F=F_{0} \sin \omega t\)\\
or

\[
m \frac{d \vec{v}}{d t}=\vec{F}_{0} \sin \omega t \text { or } m d \vec{v}=\vec{F}_{0} \sin \omega t d t
\]

On integrating,

\[
\overrightarrow{m v}=\frac{-\vec{F}_{0}}{\omega} \cos \omega t+C,(\text { where } C \text { is integration constant })
\]

When

\[
t=0, v=0, \text { so } C=\frac{\vec{F}_{0}}{m \omega}
\]

Hence, \(\vec{v}=\frac{-\vec{F}_{0}}{m \omega} \cos \omega t+\frac{\vec{F}_{0}}{m \omega}\)\\
As \(\mid \cos \omega t \leq 1\) so, \(v=\frac{F_{0}}{m \omega}(1-\cos \omega t)\)

Thus

\[
\begin{gathered}
s=\int_{0}^{t} v d t \\
=\frac{F_{0} t}{m \omega}-\frac{F_{0} \sin \omega t}{m \omega^{2}}=\frac{F_{0}}{m \omega^{2}}(\omega t-\sin \omega t)
\end{gathered}
\]

(Figure in the answer sheet).\\
1.99 According to the problem, the force acting on the particle of mass \(m\) is, \(\vec{F}=\overrightarrow{F_{0}} \cos \omega t\) So, \(\quad m \frac{d \vec{v}}{d t}=\vec{F}_{0} \cos \omega t\) or \(d \vec{v}=\frac{\vec{F}_{0}}{m} \cos \omega t d t\)

Integrating, within the limits.

\[
\int_{0}^{\vec{v}} d \vec{v}=\frac{\vec{F}_{0}}{m} \int_{0}^{t} \cos \omega t d t \text { or } \vec{v}=\frac{\vec{F}_{0}}{m \omega} \sin \omega t
\]

It is clear from equation (1), that after starting at \(t=0\), the particle comes to rest fro the first time at \(t=\frac{\pi}{\omega}\).\\
From Eq. (1), \(v=|\vec{v}|=\frac{F_{0}}{m \omega} \sin \omega t\) for \(t \leq \frac{\pi}{\omega}\)

Thus during the time interval \(t=\pi / \omega\), the sought distance

\[
s=\frac{F_{0}}{m w} \int_{0}^{\pi / \omega} \sin \omega t d t=\frac{2 F}{m \omega^{2}}
\]

From Eq. (1)

\[
v_{\max }=\frac{F_{0}}{m \omega} \text { as }|\sin \omega t| \leq 1
\]

1.100 (a) From the problem

\[
\vec{F}=-\overrightarrow{r v} \text { so } m \frac{d \vec{v}}{d t}=-\overrightarrow{r v}
\]

Thus

\[
m \frac{d v}{d t}=-r v[\text { as } d \vec{v} \uparrow \downarrow \vec{v}]
\]

or,

\[
\frac{d v}{v}=-\frac{r}{m} d t
\]

On integrating

\[
\ln v=-\frac{r}{m} t+C
\]

But at

\[
t=0, v=v_{0}, \text { so, } C=\ln v_{0}
\]

or,

\[
\ln \frac{v}{v_{0}}=-\frac{r}{m} t \text { or, } v=v_{0} e^{-\frac{r}{m} t}
\]

Thus for

\[
t \rightarrow \infty, v=0
\]

(b) We have \(m \frac{d v}{d t}=-r v\) so \(d v=\frac{-r}{m} d s\)

Integrating within the given limits to obtain \(\nu(s)\) :\\
or,


\begin{equation*}
\int_{v_{0}}^{v} d v=-\frac{r}{m} \int_{0}^{s} d s \text { or } v=v_{0}-\frac{r s}{m} \tag{1}
\end{equation*}


Thus for

\[
v=0, s=s_{\text {total }}=\frac{m v_{0}}{r}
\]

(c) Let we have

\[
\frac{m d v}{v}=-r v \text { or } \frac{d v}{v}=\frac{-r}{m} d t
\]

or,

\[
\int_{0}^{v_{0} / \eta} \frac{d v}{v}=\frac{-r}{m} \int_{0}^{t} d t, \text { or, } \ln \frac{v_{0}}{\eta v_{0}}=-\frac{r}{m} t
\]

So

\[
t=\frac{-m \ln (1 / \eta)}{r}=\frac{m \ln \eta}{r}
\]

Now, average velocity over this time interval,

\[
<v>=\frac{\int v d t}{\int d t}=\frac{\int_{0}^{\frac{m}{r} \ln \eta} v_{0} e^{-\frac{r t}{m}} d t}{\frac{m}{r} \ln \eta}=\frac{v_{0}(\eta-1)}{\eta \ln \eta}
\]

1.101 According to the problem

\[
m \frac{d v}{d t}=-k v^{2} \text { or, } m \frac{d v}{v^{2}}=-k d t
\]

Integrating, withing the limits,


\begin{equation*}
\int_{v_{0}}^{v} \frac{d v}{v^{2}}=-\frac{k}{m} \int_{0}^{t} d t \text { or, } t=\frac{m}{k} \frac{\left(v_{0}-v\right)}{v_{0} v} \tag{1}
\end{equation*}


To finc the value of \(k\), rewrite

\[
m v \frac{d v}{d s}=-k v^{2} \text { or, } \frac{d v}{v}=-\frac{k}{m} d s
\]

On integrating

So,


\begin{gather*}
\int_{v_{0}}^{v} \frac{d v}{v}=-\frac{k}{m} \int_{0}^{h} d s \\
k=\frac{m}{h} \ln \frac{v_{0}}{v} \tag{2}
\end{gather*}


Putting the value of \(k\) from (2) in (1), we get

\[
t=\frac{h\left(v_{0}-v\right)}{v_{0} v \ln \frac{v_{0}}{v}}
\]

1.102 From Newton's second law for the bar in projection from, \(F_{x}=m w_{x}\) along \(x\) direction we get\\
or,

\[
v \frac{d v}{d x}=g \sin \alpha-a x g \cos \alpha,(\text { as } k=a x),
\]

or,

\[
v d v=(g \sin \alpha-a x g \cos \alpha) d x
\]

or,


\begin{equation*}
\int_{0}^{v} v d v=g \int_{0}^{x}(\sin \alpha-x \cos \alpha) d x \tag{1}
\end{equation*}


So, \(\frac{v^{2}}{2}=g\left(\sin \alpha x-\frac{x^{2}}{2} a \cos \alpha\right)\)

From (1) \(v=0\) at either

\[
x=0, \text { or } x=\frac{2}{a} \tan \alpha
\]

As the motion of the bar is unidirectional it stops after going through a distance of \(\frac{2}{a} \tan \alpha\).\\
\includegraphics[max width=\textwidth, alt={}, center]{9f3c5a8e-ac5d-42f4-8eb6-ea2af623f2d8-065_315_349_511_926}

From (1), for \(v_{\text {max }}\),

\[
\frac{d}{d x}\left(\sin \alpha x-\frac{x^{2}}{2} a \cos \alpha\right)=0, \text { which yields } x=\frac{1}{a} \tan \alpha
\]

Hence, the maximum velocity will be at the distance, \(x=\tan \alpha / a\) Putting this value of \(x\) in (1) the maximum velocity,

\[
v_{\max }=\sqrt{\frac{g \sin \alpha \tan \alpha}{a}}
\]

1.103 Since, the applied force is proportional to the time and the frictional force also exists, the motion does not start just after applying the force. The body starts its motion when \(F\) equals the limiting friction.\\
Let the motion start after time \(t_{0}\), then\\
\(F=a t_{0}=k m g\) or, \(t_{0}=\frac{k m g}{a}\)\\
So, for \(t=s t_{0}\), the body remains at rest and for \(t>t_{0}\) obviously\\
\(\frac{m d v}{d t}=a\left(t-t_{0}\right)\) or, \(m d v=a\left(t-t_{0}\right) d t\)\\
Integrating, and noting \(v=0\) at \(t=t_{0}\), we have for \(t>t_{0}\)

\[
\int_{0}^{v} m d v=a \int_{t_{0}}^{t}\left(t-t_{0}\right) d t \text { or } v=\frac{a}{2 m}\left(t-t_{0}\right)^{2}
\]

Thus

\[
s=\int v d t=\frac{a}{2 m} \int_{t_{0}}^{t}\left(t-t_{0}\right)^{2} d t=\frac{a}{6 m}\left(t-t_{0}\right)^{3}
\]

1.104 While going upward; from Newton's second law in vertical direction :

\[
m \frac{v d v}{d s}=-\left(m g+k v^{2}\right) \text { or } \frac{v d v}{\left(g+\frac{k v^{2}}{m}\right)}=-d s
\]

At the maximum height \(h\), the speed \(v=0\), so

\[
\int_{v_{0}}^{0} \frac{v d v}{g+\left(k v^{2} / m\right)}=-\int_{0}^{h} d s
\]

Integrating and solving, we get,


\begin{equation*}
h=\frac{m}{2 k} \ln \left(1+\frac{k v_{o}^{2}}{m g}\right) \tag{1}
\end{equation*}


When the body falls downward, the net force acting on the body in downward direction equals ( \(m g-k v^{2}\) ),\\
Hence net acceleration, in downward direction, according to second law of motion

Thus

\[
\begin{gathered}
\frac{v d v}{d s}=g-\frac{k v^{2}}{m} \text { or, } \frac{v d v}{g-\frac{k v^{2}}{m}}=d s \\
\int_{0}^{v^{\prime}} \frac{v d v}{g-k v^{2} / m}=\int_{0}^{h} d s
\end{gathered}
\]

Integrating and putting the value of \(h\) from (1), we get,

\[
v^{\prime}=v_{0} / \sqrt{1+k v_{0}^{2} / m g} .
\]

1.105 Let us fix \(x-y\) co-ordinate system to the given plane, taking \(x\)-axis in the direction along which the force vector was oriented at the moment \(t=0\), then the fundamental equation of dynamics expressed via the projection on \(x\) and \(y\)-axes gives,


\begin{equation*}
F \cos \omega t=m \frac{d v_{x}}{d t} \tag{1}
\end{equation*}


and


\begin{equation*}
F \sin \omega t=m \frac{d v_{y}}{d t} \tag{2}
\end{equation*}


(a) Using the condition \(v(0)=0\), we obtain \(v_{x}=\frac{F}{m \omega} \sin \omega t\)\\
and


\begin{equation*}
v_{y}=\frac{F}{m \omega}(1-\cos \omega t) \tag{4}
\end{equation*}


Hence,

\[
v=\sqrt{v_{x}^{2}+v_{y}^{2}}=\left(\frac{2 F}{m \omega}\right)\left|\sin \left(\frac{\omega t}{2}\right)\right|
\]

(b) It is seen from this that the velocity \(v\) turns into zero after the time interval \(\Delta t\), which can be found from the relation, \(\omega \frac{\Delta t}{2}=\pi\). Consequentely,\\
the sought distance, is

\[
s=\int_{0}^{\Delta t} v d t=\frac{8 F}{m \omega^{2}}
\]

Average velocity, \(\langle v\rangle=\frac{\int v d t}{\int d t}\)\\
\includegraphics[max width=\textwidth, alt={}, center]{9f3c5a8e-ac5d-42f4-8eb6-ea2af623f2d8-067_393_514_176_855}

So, \(\langle v\rangle=\int_{0}^{2 \pi / \omega} \frac{2 F}{m \omega} \sin \left(\frac{\omega t}{2}\right) d t /(2 \pi \omega)=\frac{4 F}{\pi m \omega}\)\\
1.106 The acceleration of the disc along the plane is determined by the projection of the force of gravity on this plane \(F_{x}=m g \sin \alpha\) and the friction force \(f r=k m g \cos \alpha\). In our case \(k=\tan \alpha\) and therefore

\[
f r=F_{x}=m g \sin \alpha
\]

Let us find the projection of the acceleration on the derection of the tangent to the trajectory and on the \(x\)-axis :\\
\(m w_{t}=F_{x} \cos \varphi-f r=m g \sin \alpha(\cos \varphi-1)\)\\
\(m w_{x}=F_{x}-f r \cos \varphi=m g \sin \alpha(1-\cos \varphi)\)\\
It is seen fromthis that \(w_{t}=-w_{x}\), which means that the velocity \(v\) and its projection \(v_{x}\) differ only by a constant value \(C\) which does not change with time, i.e.\\
\includegraphics[max width=\textwidth, alt={}, center]{9f3c5a8e-ac5d-42f4-8eb6-ea2af623f2d8-067_459_532_949_847}

\[
v=v_{x}+C,
\]

where \(v_{x}=v \cos \varphi\). The constant \(C\) is found from the initial condition \(v=v_{0}\), whence \(C=v_{0}\) since \(\varphi=\frac{\pi}{2}\) initially. Finally we obtain

\[
v=v_{0} /(1+\cos \varphi) .
\]

In the cource of time \(\varphi \rightarrow 0\) and \(v \rightarrow v_{0} / 2\). (Motion then is unaccelerated.)\\
1.107 Let us consider an element of length \(d s\) at an angle \(\varphi\) from the vertical diameter. As the speed of this element is zero at initial instant of time, it's centripetal acceleration is zero, and hence, \(d N-\lambda d s \cos \varphi=0\), where \(\lambda\) is the linear mass density of the chain Let \(T\) and \(T+d T\) be the tension at the upper and the lower ends of \(d s\). we have from, \(F_{t}=m w_{t}\)\\
or,

\[
\begin{gathered}
(T+d T)+\lambda d s g \sin \varphi-T=\lambda d s w_{t} \\
d T+\lambda R d \varphi g \sin \varphi=\lambda d s w_{t}
\end{gathered}
\]

If we sum the above equation for all elements, the term \(\int d T=0\) because there is no tension at the free ends, so

Hence

\[
\lambda g R \int_{0}^{l / R} \sin \varphi d \varphi=\lambda w_{t} \int d s=\lambda l w_{t}
\]

Hence \(w_{t}=\frac{g R}{l}\left(1-\cos \frac{l}{R}\right)\)\\
As \(\boldsymbol{w}_{n}=a\) at initial moment\\
So, \(w=\left|w_{t}\right|=\frac{g R}{l}\left(1-\cos \frac{l}{R}\right)\)\\
\includegraphics[max width=\textwidth, alt={}, center]{9f3c5a8e-ac5d-42f4-8eb6-ea2af623f2d8-068_386_411_222_870}\\
1.108 In the problem, we require the velocity of the body, realtive to the sphere, which itself moves with an acceleration \(w_{0}\) in horizontal direction (say towards left). Hence it is advisible to solve the problem in the frame of sphere (non-inertial frame).\\
At an arbitary moment, when the body is at an angle \(\theta\) with the vertical, we sketch the force diagram for the body and write the second law of motion in projection form \(F_{n}=m w_{n}\)\\
or, \(m g \cos \theta-N-m w_{0} \sin \theta=\frac{m v^{2}}{R}\)

At the break off point, \(N=0, \theta=\theta_{0}\) and let\\
\(v=v_{0}\), so the Eq. (1) becomes,


\begin{equation*}
\frac{v_{0}^{2}}{R}=g \cos \theta_{0}-w_{0} \sin \theta_{0} \tag{2}
\end{equation*}


From, \(F_{t}=m w_{t}\)\\
\(m g \sin \theta-m w_{0} \cos \theta=m \frac{v d v}{d s}=m \frac{v d v}{R d \theta}\)\\
or, \(v d v=R\left(g \sin \theta+w_{0} \cos \theta\right) d \theta\)\\
Integrating, \(\int_{0}^{\nu_{0}} v d v=\int_{0}^{\theta_{0}} R\left(g \sin \theta+w_{0} \cos \theta\right) d \theta\)\\
\includegraphics[max width=\textwidth, alt={}, center]{9f3c5a8e-ac5d-42f4-8eb6-ea2af623f2d8-068_352_454_1133_908}


\begin{equation*}
\frac{v_{0}^{2}}{2 R}=g\left(1-\cos \theta_{0}\right)+w_{0} \sin \theta_{0} \tag{3}
\end{equation*}


Note that the Eq. (3) can also be obtained by the work-energy theorem \(A=\Delta T\) (in the frame of sphere)\\
therefore, \(\quad m g R\left(1-\cos \theta_{0}\right)+m w_{0} R \sin \theta_{0}=\frac{1}{2} m v_{0}^{2}\)\\[0pt]
[here \(m w_{0} R \sin \theta_{0}\) is the work done by the pseudoforce \(\left(-m \vec{w}_{o}\right)\) ]\\
or,

\[
\frac{v_{0}^{2}}{2 R}=g\left(1-\cos \theta_{0}\right)+w_{0} \sin \theta_{0}
\]

Solving Eqs. (2) and (3) we get,\\
\(v_{0}=\sqrt{2 g R / 3}\) and \(\theta_{0}=\cos ^{-1}\left[\frac{2+k \sqrt{5+9 k^{2}}}{3\left(1+k^{2}\right)}\right]\), where \(k=\frac{w_{0}}{g}\)\\
Hence

\[
\left.\theta_{0}\right|_{w_{0}=g}=17^{\circ}
\]

1.109 This is not central force problem unless the path is a circle about the said point. Rather here \(F_{t}\) (tangential force) vanishes. Thus equation of motion becomes,\\
and,

\[
\begin{gathered}
v_{t}=v_{0}=\text { constant } \\
\frac{m v_{0}^{2}}{r}=\frac{A}{r^{2}} \text { for } r=r_{0}
\end{gathered}
\]

We can consider the latter equation as the equilibrium under two forces. When the motion is perturbed, we write \(r=r_{0}+x\) and the net force acting on the particle is,

\[
-\frac{A}{\left(r_{0}+x\right)^{n}}+\frac{m v_{0}^{2}}{r_{0}+x}=\frac{-A}{r_{0}^{n}}\left(1-\frac{n x}{r_{0}}\right)+\frac{m v_{0}^{2}}{r_{0}}\left(1-\frac{x}{r_{0}}\right)=-\frac{m v_{0}^{2}}{r_{0}^{2}}(1-n) x
\]

This is opposite to the displacement \(x\), if \(n<1 \cdot\left(\frac{m v_{0}^{2}}{r}\right.\) is an outward directed centrifugul force while \(\frac{-A}{r^{n}}\) is the inward directed external force).\\
1.110 There are two forces on the sleeve, the weight \(F_{1}\) and the centrifugal force \(F_{2}\). We resolve both forces into tangential and normal component then the net downward tangential force on the sleeve is,

\[
m g \sin \theta\left(1-\frac{\omega^{2} R}{g} \cos \theta\right)
\]

This vanishes for \(\theta=0\) and for \(\theta=\theta_{0}=\cos ^{-1}\left(\frac{g}{\omega^{2} R}\right)\), which is real if \(\omega^{2} R>g\). If \(\omega^{2} R<g\), then \(1-\frac{\omega^{2} R}{g} \cos \theta\) is always positive for small values of \(\theta\) and hence the net tangential force near \(\theta=0\) opposes any displacement away from it. \(\theta=0\) is then stable.\\
\includegraphics[max width=\textwidth, alt={}, center]{9f3c5a8e-ac5d-42f4-8eb6-ea2af623f2d8-069_469_620_1238_783}

If \(\omega^{2} R>g, 1-\frac{\omega^{2} R}{g} \cos \theta\) is negative for small\\
\(\theta\) near \(\theta=0\) and \(\theta=0\) is then unstable.\\
However \(\theta=\theta_{0}\) is stable because the force tends to bring the sleeve near the equilibrium position \(\theta=\theta_{0}\).\\
If \(\omega^{2} R=g\), the two positions coincide and becomes a stable equilibrium point.\\
1.111 Define the axes as shown with \(z\) along the local vertical, \(x\) due east and \(y\) due north. (We assume we are in the northern hemisphere). Then the Coriolis force has the components.

\[
\vec{F}_{\text {cor }}=-2 m(\vec{\omega} \times \vec{v})
\]

\(\left.=2 m \omega\left[v_{y} \cos \theta-v_{z} \sin \theta\right) \vec{i}-v_{x} \cos \theta \vec{j}+v_{x} \cos \theta \vec{k}\right]=2 m \omega\left(v_{y} \cos \theta-v_{z} \sin \theta\right) \vec{i}\)\\
since \(v_{x}\) is small when the direction in which the gun is fired is due north. Thus the equation of motion (neglecting centrifugal forces) are\\
\(\ddot{x}=2 m \omega\left(v_{y} \sin \varphi-v_{z} \cos \varphi\right), \ddot{y}=0\) and \(\ddot{z}=-g\) Integrating we get \(\dot{y}=v\) (constant), \(\dot{z}=-g t\) and \(\dot{x}=2 \omega v \sin \varphi t+\omega g t^{2} \cos \varphi\)\\
Finally,\\
\(x=\omega \nu t^{2} \sin \varphi+\frac{1}{3} \omega g t^{3} \cos \varphi\)\\
Now \(v \gg g t\) in the present case. so,

\[
\begin{aligned}
x & \approx \omega v \sin \varphi\left(\frac{s}{v}\right)^{2}=\omega \sin \varphi \frac{s^{2}}{v} \\
& \approx 7 \mathrm{~cm} \text { (to the east) }
\end{aligned}
\]

\includegraphics[max width=\textwidth, alt={}, center]{9f3c5a8e-ac5d-42f4-8eb6-ea2af623f2d8-070_503_567_416_804}\\
1.112 The disc exerts three forces which are mutually perpendicular. They are the reaction of the weight, \(m g\), vertically upward, the Coriolis force \(2 m v^{\prime} \omega\) perpendicular to the plane of the vertical and along the diameter, and \(m \omega^{2} r\) outward along the diameter. The resultant force is,

\[
F=m \sqrt{g^{2}+\omega^{4} r^{2}+\left(2 v^{\prime} \omega\right)^{2}}
\]

1.113 The sleeve is free to slide along the \(\operatorname{rod} A B\). Thus only the centrifugal force acts on it. The equation is,

\[
m \dot{v}=m \omega^{2} r \text { where } v=\frac{d r}{d t} .
\]

\[
\text { But } \dot{v}=v \frac{d v}{d r}=\frac{d}{d r}\left(\frac{1}{2} v^{2}\right)
\]

so,

\[
\frac{1}{2} v^{2}=\frac{1}{2} \omega^{2} r^{2}+\text { constant }
\]

or,

\[
v^{2}=v_{0}^{2}+\omega^{2} r^{2}
\]

\(v_{0}\) being the initial velocity when \(r=0\). The Coriolis force is then,

\[
\begin{gathered}
2 m \omega \sqrt{v_{0}^{2}+\omega^{2} r^{2}}=2 m \omega^{2} r \sqrt{1+v_{0}^{2} / \omega^{2} r^{2}} \\
=2.83 \mathrm{~N} \text { on putting the values. }
\end{gathered}
\]

1.114 The disc \(O B A C\) is rotating with angular velocity \(\omega\) about the axis \(O O^{\prime}\) passing through the edge point \(O\). The equation of motion in rotating frame is,\\
\(m \vec{w}=\vec{F}+m \omega^{2} \vec{R}+2 m \vec{v} \times \vec{\omega}=\vec{F}+\vec{F}_{i n}\)\\
where \(\vec{F}_{\text {in }}\) is the resultant inertial forc (pseudo force) which is the vector sum of centrifugal and Coriolis forces.\\
(a) At \(A, F_{i n}\) vanishes. Thus \(0=-2 m \omega^{2} R \hat{n}+2 m v^{\prime} \omega \hat{n}\)\\
\includegraphics[max width=\textwidth, alt={}, center]{9f3c5a8e-ac5d-42f4-8eb6-ea2af623f2d8-071_455_456_105_942}\\
where \(\hat{n}\) is the inward drawn unit vector to the centre from the point in question, here \(A\).\\
Thus,

\[
v^{\prime}=\omega R
\]

so,

\[
w=\frac{v^{\prime 2}}{\rho}=\frac{v^{\prime 2}}{R}=\omega^{2} R .
\]

(b) At B

\[
\vec{F}_{i n}=m \omega^{2} \overrightarrow{O C}+m \omega^{2} \overrightarrow{B C}
\]

its magnitude is \(m \omega^{2} \sqrt{4 R^{2}-r^{2}}\), where \(r=O B\).\\
1.115 The equation of motion in the rotating coordinate system is,

\[
m \vec{w}=\vec{F}+m \omega^{2} \vec{R}+2 m(\vec{v} \times \vec{\omega})
\]

Now, \(\quad \vec{v}=R \theta \overrightarrow{e_{\theta}}+R \sin \theta \dot{\varphi} \vec{e}_{\varphi}\)\\
and

\[
\vec{w}=w^{\prime} \cos \theta \overrightarrow{e_{r}}-w^{\prime} \sin \theta \overrightarrow{e_{\theta}}
\]

\[
\frac{1}{2 m} \vec{F}_{\operatorname{cor}}=\left|\begin{array}{ccc}
\vec{e}_{r} & \vec{e}_{\theta} & \vec{e}_{\varphi} \\
0 & R \theta & R \sin \theta \dot{\varphi} \\
\omega \cos \theta & -\omega \sin \theta & 0
\end{array}\right|
\]

\(=\overrightarrow{e_{r}}\left(\omega R \sin ^{2} \theta \dot{\varphi}\right)+\omega R \sin \theta \cos \theta \dot{\varphi} \overrightarrow{e_{\theta}}-\omega R \theta \cos \theta \overrightarrow{e_{\theta}}\)\\
Now on the sphere,

\[
\begin{gathered}
\vec{v}=\left(-R \dot{\theta}^{2}-R \sin ^{2} \theta \varphi^{2}\right) \vec{e}_{r} \\
+\left(R \dot{\theta}^{\cdot}-R \sin \theta \cos \theta \dot{\varphi}^{2}\right) \overrightarrow{e_{\theta}} \\
+(R \sin \theta \ddot{\varphi}+2 R \cos \theta \dot{\theta} \dot{\varphi}) \overrightarrow{e_{\varphi}}
\end{gathered}
\]

\begin{center}
\includegraphics[max width=\textwidth, alt={}]{9f3c5a8e-ac5d-42f4-8eb6-ea2af623f2d8-071_461_489_929_916}
\end{center}

Thus the equation of motion are,\\
\(m\left(-R \dot{\theta}^{2}-R \sin ^{2} \theta \dot{\varphi}^{2}\right)=N-m g \cos \theta+m \omega^{2} R \sin ^{2} \theta+2 m \omega R \sin ^{2} \theta \dot{\varphi}\)\\
\(m\left(R \dot{\theta}^{\cdot}-R \sin \theta \cos \theta \dot{\varphi}^{2}\right)=m g \sin \theta+m \omega^{2} R \sin \theta \cos \theta+2 m \omega R \sin \theta \cos \theta \dot{\varphi}\)\\
\(m(R \sin \theta \ddot{\varphi}+2 R \cos \theta \dot{\theta} \dot{\varphi})=-2 m \omega R \dot{\theta} \cos \theta\)\\
From the third equation, we get, \(\dot{\varphi}=-\omega\)\\
A result that is easy to understant by considering the motion in non-rotating frame. The eliminating \(\dot{\varphi}\) we get,

\[
\begin{aligned}
& m R \dot{\theta}^{2}=m g \cos \theta-N \\
& m R \dot{\theta}^{3}=m g \sin \theta
\end{aligned}
\]

Integrating the last equation,

\[
\frac{1}{2} m R \dot{\theta}^{2}=m g(1-\cos \theta)
\]

Hence

\[
N=(3-2 \cos \theta) m g
\]

So the body must fly off for \(\theta=\theta_{0}=\cos ^{-1} \frac{2}{3}\), exactly as if the sphere were nonrotating.\\
Now, at this point \(F_{c f}=\) centrifugal force \(=m \omega^{2} R \sin \theta_{0}=\sqrt{\frac{5}{9}} m \omega^{2} R\)

\[
\begin{aligned}
F_{c o r} & =\sqrt{\omega^{2} R^{2} \theta^{2} \cos ^{2} \theta+\left(\omega^{2} R^{2}\right)^{2} \sin ^{2} \theta} \times 2 m \\
& =\sqrt{\frac{5}{9}\left(\omega^{2} R\right)^{2}+\omega^{2} R^{2} \times \frac{4}{9} \times \frac{2 g}{3 R}} \times 2 m=\frac{2}{3} m \omega^{2} R \sqrt{5+\frac{8 g}{3 \omega^{2} R}}
\end{aligned}
\]

1.116 (a) When the train is moving along a meridian only the Coriolis force has a lateral component and its magnitude (see the previous problem) is,

\[
2 m \omega v \cos \theta=2 m \omega \sin \lambda
\]

(Here we have put \(R \dot{\theta} \rightarrow v\) )\\
So,

\[
\begin{aligned}
F_{\text {lateral }} & =2 \times 2000 \times 10^{3} \times \frac{2 \pi}{86400} \times \frac{54000}{3600} \times \frac{\sqrt{3}}{2} \\
& =3.77 \mathrm{kN},(\text { we write } \lambda \text { for the latitude })
\end{aligned}
\]

(b) The resultant of the inertial forces acting on the train is,

\[
\vec{F}_{i n}=-2 m \omega R \hat{\theta} \cos \theta \vec{e}_{\varphi}
\]

\(+\left(m \omega^{2} R \sin \theta \cos \theta+2 m \omega R \sin \theta \cos \theta \dot{\varphi}\right) \overrightarrow{e_{\theta}}\)\\
\(+\left(m \omega^{2} R \sin ^{2} \theta+2 m \omega R \sin ^{2} \theta \dot{\varphi}\right) \vec{e}_{r}\)\\
This vanishes if \(\dot{\theta}=0, \dot{\varphi}=-\frac{1}{2} \omega\)\\
\includegraphics[max width=\textwidth, alt={}, center]{9f3c5a8e-ac5d-42f4-8eb6-ea2af623f2d8-072_396_542_764_817}

Thus

\[
\vec{v}=v_{\varphi} \vec{e}_{\varphi}, v_{\varphi}=-\frac{1}{2} \omega R \sin \theta=-\frac{1}{2} \omega R \cos \lambda
\]

(We write \(\dot{\lambda}\) for the latitude here)\\
Thus the train must move from the east to west along the \(60^{\text {th }}\) parallel with a speed, \(\frac{1}{2} \omega R \cos \lambda=\frac{1}{4} \times \frac{2 \pi}{8.64} \times 10^{-4} \times 6.37 \times 10^{6}=115.8 \mathrm{~m} / \mathrm{s} \approx 417 \mathrm{~km} / \mathrm{hr}\)\\
1.117 We go to the equation given in 1.111 . Here \(v_{y}=0\) so we can take \(y=0\), thus we get for the motion in the \(x \dot{z}\) plane.\\
and

\[
\begin{gathered}
\ddot{x}=-2 \omega v_{z} \cos \theta \\
\ddot{z}=-g
\end{gathered}
\]

Integrating,

\[
z=-\frac{1}{2} g t^{2}
\]

\[
\dot{x}=\omega g \cos \varphi t^{2}
\]

So

\[
\begin{gathered}
x=\frac{1}{3} \omega g \cos \varphi t^{3}=\frac{1}{3} \omega g \cos \varphi\left(\frac{2 h}{g}\right)^{3 / 2} \\
=\frac{2 \omega h}{3} \cos \varphi \sqrt{\frac{2 h}{g}}
\end{gathered}
\]

There is thus a displacement to the east of

\[
\frac{2}{3} \times \frac{2 \pi}{8} 64 \times 500 \times 1 \times \sqrt{\frac{400}{9 \cdot 8}} \approx 26 \mathrm{~cm} .
\]

\subsection*{1.3 Laws of Conservation of Energy, Momentum and Angular Momentum.}
1.118 As \(\vec{F}\) is constant so the sought work done

\[
A=\vec{F} \cdot \Delta \vec{r}=\vec{F} \cdot\left(\overrightarrow{r_{2}}-\overrightarrow{r_{1}}\right)
\]

or, \(A=(3 \vec{i}+4 \vec{j}) \cdot[(2 \vec{i}-3 \vec{j})-(\vec{i}+2 \vec{j})]=(3 \vec{i}+4 \vec{j}) \cdot(\vec{i}-5 \vec{j})=17 \mathrm{~J}\)\\
1.119 Differentating \(v(s)\) with respect to time

\[
\frac{d v}{d t}=\frac{a}{2 \sqrt{s}} \frac{d s}{d t}=\frac{a}{2 \sqrt{s}} a \sqrt{s}=\frac{a^{2}}{2}=w
\]

(As locomotive is in unidrectional motion)\\
Hence force acting on the locomotive \(F=m w=\frac{m a^{2}}{2}\)\\
Let, at \(v=0\) at \(t=0\) then the distance covered during the first \(t\) seconds

\[
s=\frac{1}{2} w t^{2}=\frac{1}{2} \frac{a^{2}}{2} t^{2}=\frac{a^{2}}{4} t^{2}
\]

Hence the sought work, \(A=F s=\frac{m a^{2}}{2} \frac{\left(a^{2} t^{2}\right)}{4}=\frac{m a^{4} t^{2}}{8}\)\\
1.120 We have


\begin{equation*}
T=\frac{1}{2} m v^{2}=a s^{2} \text { or, } v^{2}=\frac{2 a s^{2}}{m} \tag{1}
\end{equation*}


Differentating Eq. (1) with respect to time


\begin{equation*}
2 v w_{t}=\frac{4 a s}{m} v \text { or, } w_{t}=\frac{2 a s}{m} \tag{2}
\end{equation*}


Hence net acceleration of the particle

\[
w=\sqrt{w_{t}^{2}+w_{n}^{2}}=\sqrt{\left(\frac{2 a s}{m}\right)^{2}+\left(\frac{2 a s^{2}}{m R}\right)^{2}}=\frac{2 a s}{m} \sqrt{1+(s / R)^{2}}
\]

Hence the sought force, \(F=m w=2 a s \sqrt{1+(s / R)^{2}}\)\\
1.121 Let \(\vec{F}\) makes an angle \(\theta\) with the horizontal at any instant of time (Fig.). Newton's second law in projection form along the direction of the force, gives :\\
\(F=k m g \cos \theta+m g \sin \theta\) (because there is no acceleration of the body.)\\
As \(\vec{F} \uparrow \uparrow d \vec{r}\) the differential work done by the force \(\vec{F}\),

\[
\begin{aligned}
d A & =\vec{F} \cdot d \vec{r}=F d s,(\text { where } d s=|d \vec{r}|) \\
& =k m g d s(\cos \theta)+m g d s \sin \theta \\
& =k m g d x+m g d y .
\end{aligned}
\]

Hence, \(A=k m g \int_{0}^{l} d x+m g \int_{0}^{h} d y\)

\[
=k m g l+m g h=m g(k l+h) .
\]

\includegraphics[max width=\textwidth, alt={}, center]{9f3c5a8e-ac5d-42f4-8eb6-ea2af623f2d8-073_460_605_1621_799}\\
1.122 Let \(s\) be the distance covered by the disc along the incline, from the Eq. of increment of M.E. of the disc in the field of gravity : \(\Delta T+\Delta U=A_{f r}\)


\begin{equation*}
0+(-m g s \sin \alpha)=-k m g \cos \alpha s-k m g l \tag{1}
\end{equation*}


or, \(\quad s=\frac{k l}{\sin \alpha-k \cos \alpha}\)

Hence the sought work

\[
A_{f r}=-k m g[s \cos \alpha+l]
\]

\(A_{f r}=-\frac{k l m g}{1-k \cot \alpha}\) [Using the Eqn. (1)]\\
On puting the values \(A_{f r}=-0.05 \mathrm{~J}\)\\
\includegraphics[max width=\textwidth, alt={}, center]{9f3c5a8e-ac5d-42f4-8eb6-ea2af623f2d8-074_476_562_135_780}\\
1.123 Let \(x\) be the compression in the spring when the bar \(m_{2}\) is about to shift. Therefore at this moment spring force on \(m_{2}\) is equal to the limiting friction between the bar \(m_{2}\) and horizontal floor. Hence


\begin{equation*}
\kappa x=k m_{2} g \quad[\text { where } \kappa \text { is the spring constant (say)] } \tag{1}
\end{equation*}


For the block \(m_{1}\) from work-energy theorem : \(A=\Delta T=0\) for minimum force. ( \(A\) here indudes the work done in stretching the spring.)\\
so,


\begin{equation*}
F x-\frac{1}{2} \kappa x^{2}-k m g x=0 \text { or } \kappa \frac{x}{2}=F-k m_{1} g \tag{2}
\end{equation*}


From (1) and (2),

\[
F=k g\left(m_{1}+\frac{m_{2}}{2}\right) .
\]

1.124 From the initial condition of the problem the limiting fricition between the chain lying on the horizontal table equals the weight of the over hanging part of the chain, i.e.\\
\(\lambda \eta \lg =k \lambda(1-\eta) \lg\) (where \(\lambda\) is the linear mass density of the chain)

So,


\begin{equation*}
k=\frac{\eta}{1-\eta} \tag{1}
\end{equation*}


Let (at an arbitrary moment of time) the length of the chain on the table is \(x\). So the net friction force between the chain and the table, at this moment :


\begin{equation*}
f_{r}=k N=k \lambda x g \tag{2}
\end{equation*}


The differential work done by the friction forces :\\
\includegraphics[max width=\textwidth, alt={}, center]{9f3c5a8e-ac5d-42f4-8eb6-ea2af623f2d8-074_449_625_1216_758}


\begin{equation*}
d A=\overrightarrow{f_{r}} \cdot d \vec{r}=-f_{r} d s=-k \lambda x g(-d x)=\lambda g\left(\frac{\eta}{1-\eta}\right) x d x \tag{3}
\end{equation*}


(Note that here we have written \(d s=-d x\)., because \(d s\) is essentially a positive term and as the length of the chain decreases with time, \(d x\) is negative)\\
Hence, the sought work done

\[
A=\int_{(1-\eta) l}^{0} \lambda g \frac{\eta}{1-\eta} x d x=-(1-\eta) \eta \frac{m g l}{2}=-1.3 \mathrm{~J}
\]

1.125 The velocity of the body, \(t\) seconds after the begining of the motion becomes \(\vec{v}=\overrightarrow{v_{0}}+\overrightarrow{g t}\). The power developed by the gravity ( \(m \vec{g}\) ) at that moment, is


\begin{equation*}
P=m \vec{g} \cdot \vec{v}=m\left(\vec{g} \cdot \overrightarrow{v_{0}}+g^{2} t\right)=m g\left(g t-v_{0} \sin \alpha\right) \tag{1}
\end{equation*}


As \(m \vec{g}\) is a constant force, so the average power

\[
<P>=\frac{A}{\tau}=\frac{m \vec{g} \cdot \Delta \vec{r}}{\tau}
\]

where \(\Delta \vec{r}\) is the net displacement of the body during time of flight.\\
As,

\[
m \vec{g} \perp \Delta \vec{r} \text { so }\langle P\rangle=0
\]

1.126 We have

\[
w_{n}=\frac{v^{2}}{R}=a t^{2}, \quad \text { or, } \quad v=\sqrt{a R} t,
\]

\(t\) is defined to start from the begining of motion from rest.\\
So,

\[
w_{t}=\frac{d v}{d t}=\sqrt{a R}
\]

Instantaneous power, \(P=\vec{F} \cdot \vec{v}=m\left(w_{t} \hat{u}_{t}+w_{n} \hat{u}_{t}\right) \cdot\left(\sqrt{a R} t \hat{u}_{t}\right)\),\\
(where \(\hat{u}_{t}\) and \(\hat{u}_{t}\) are unit vectors along the direction of tangent (velocity) and normal respectively)\\
So,

\[
P=m w_{t} \sqrt{a R} t=m a R t
\]

Hence the sought average power

Hence

\[
\begin{aligned}
& <P>=\frac{\int_{0}^{t} P d t}{\int_{0}^{t} d t}=\frac{\int_{0}^{t} m a R t d t}{t} \\
& <P>=\frac{m a R t^{2}}{2 t}=\frac{m a R t}{2}
\end{aligned}
\]

1.127 Let the body \(m\) acquire the horizontal velocity \(v_{0}\) along positive \(x\)-axis at the point \(O\).\\
(a) Velocity of the body \(t\) seconds after the begining of the motion,


\begin{equation*}
\vec{v}=\overrightarrow{v_{0}}+\vec{w} t=\left(v_{0}-k g t\right) \vec{i} \tag{1}
\end{equation*}


Instantaneous power \(P=\vec{F} \cdot \vec{v}=(-k m g \vec{i}) \cdot\left(v_{0}-k g t\right) \vec{i}=-k m g\left(v_{0}-k g t\right)\)\\
From Eq. (1), the time of motion \(\tau=v_{0} / k g\)\\
Hence sought average power during the time of motion

\[
<P>=\frac{\int_{0}^{\tau}-k m g\left(v_{0}-k g t\right) d t}{\tau}=-\frac{k m g v_{0}}{2}=-2 \mathrm{~W} \text { (On substitution) }
\]

From \(F_{x}=m w_{x}\)\\
or,

\[
\begin{gathered}
-k m g=m w_{x}=m v_{x} \frac{d v_{x}}{d x} \\
v_{x} d v_{x}=-k g d x \doteq-\alpha g x d x
\end{gathered}
\]

To find \(v(x)\), let us integrate the above equation


\begin{equation*}
\int_{v_{0}}^{v} v_{x} d v_{x}=-\alpha g \int_{0}^{x} x d x \text { or, } v^{2}=v_{0}^{2}-\alpha g x^{2} \tag{1}
\end{equation*}


Now,


\begin{equation*}
\vec{P}=\vec{F} \cdot \vec{v}=-m \alpha x g \sqrt{v_{0}^{2}-\alpha g x^{2}} \tag{2}
\end{equation*}


For maximum power, \(\frac{d}{d t}\left(\sqrt{v_{0}^{2} x^{2}-\lambda g x^{4}}\right)=0\) which yields \(x=\frac{v_{0}}{\sqrt{2 \alpha g}}\) Putting this value of \(x\), in Eq. (2) we get,

\[
P_{\max }=-\frac{1}{2} m v_{0}^{2} \sqrt{\alpha g}
\]

Centrifugal force of inertia is directed outward along radial line, thus the sought work

\[
A=\int_{r_{1}}^{r_{2}} m \omega^{2} r d r=\frac{1}{2} m \omega^{2}\left(r_{2}^{2}-r_{1}^{2}\right)=0.20 \mathrm{~T} \text { (On substitution) }
\]

1.129 Since the springs are connected in series, the combination may be treated as a single spring of spring constant.

\[
\kappa=\frac{\kappa_{1} \kappa_{2}}{\kappa_{1}+\kappa_{2}}
\]

From the equation of increment of M.E., \(\quad \Delta T+\Delta U=A_{\text {ext }}\)

\[
0+\frac{1}{2} \kappa \Delta l^{2}=A, \text { or, } A=\frac{1}{2}\left(\frac{\kappa \kappa_{2}}{\kappa_{1}+\kappa_{2}}\right) \Delta l^{2}
\]

1.130 First, let us find the total height of ascent. At the beginning and the end of the path of velocity of the body is equal to zero, and therefore the increment of the kinetic energy of the body is also equal to zero. On the other hand, in according with work-energy theorem \(\Delta T\) is equal to the algebraic sum of the works \(A\) performed by all the forces, i.e. by the force \(F\) and gravity, over this path. However, since \(\Delta T=0\) then \(A=0\). Taking into account that the upward direction is assumed to coincide with the positive direction of the \(y\)-axis, we can write

\[
\begin{aligned}
& A=\int_{0}^{h}(\vec{F}+m \vec{g}) \cdot d \vec{r}=\int_{0}^{h}\left(F_{y}-m g\right) d y \\
&=m g \int_{0}^{h}(1-2 a y) d y=m g h(1-a h)=0 \\
& \text { whence } h=1 / a
\end{aligned}
\]

The work performed by the force \(F\) over the first half of the ascent is

\[
A_{F}=\int_{0}^{h / 2} F_{y} d y=2 m g \int_{0}^{h / 2}(1-a y) d y=3 m g / 4 a
\]

The corresponding increment of the potential energy is

\[
\Delta U=m g h / 2=m g / 2 a .
\]

1.131 From the equation \(F_{r}=-\frac{d U}{d r}\) we get \(F_{r}=\left[-\frac{2 a}{r^{3}}+\frac{b}{r^{2}}\right]\)\\
(a) we have at \(r=r_{0}\), the particle is in equilibrium position. i.e. \(F_{r}=0\) so, \(r_{0}=\frac{2 a}{b}\)

To check, whether the position is steady (the position of stable equilibrium), we have to satisfy

\[
\frac{d^{2} U}{d r^{2}}>0
\]

We have

\[
\frac{d^{2} U}{d r^{2}}=\left[\frac{6 a}{r^{4}}-\frac{2 b}{r^{3}}\right]
\]

Putting the value of \(r=r_{0}=\frac{2 a}{b}\), we get\\
\(\frac{d^{2} U}{d r^{2}}=\frac{b^{4}}{8 a^{3}},(\) as \(a\) and \(b\) are positive constant)\\
So,

\[
\frac{d^{2} U}{d r^{2}}=\frac{b^{2}}{8 a^{3}}>0
\]

which indicates that the potential energy of the system is minimum, hence this position is steady.\\
(b) We have

\[
F_{r}=-\frac{d U}{d r}=\left[-\frac{2 a}{r^{3}}+\frac{b}{r^{2}}\right]
\]

For \(F\), to be maximum, \(\frac{d F_{r}}{d r}=0\)\\
So, \(r=\frac{3 a}{b}\) and then \(F_{r(\max )}=\frac{-b^{3}}{27 a^{2}}\),\\
As \(F_{r}\) is negative, the force is attractive.\\
1.132 (a) We have

\[
F_{x}=-\frac{\partial U}{d x}=-2 \alpha x \text { and } F_{y}=\frac{-\partial U}{d y}=-2 \beta y
\]

So,


\begin{equation*}
\vec{F}=2 \alpha x \vec{i}-2 \beta y \vec{i} \text { and, } F=2 \sqrt{\alpha^{2} x^{2}+\beta^{2} y^{2}} \tag{1}
\end{equation*}


For a central force, \(\overrightarrow{r \times} \vec{F}=0\)\\
Here,

\[
\begin{aligned}
\vec{r} \times \vec{F} & =(x \vec{i}+y \vec{j}) \times(-2 \alpha x \vec{i}-2 \beta y \vec{j}) \\
& =-2 \beta x y \vec{k}-2 \alpha x y(\vec{k})=0
\end{aligned}
\]

Hence the force is not a central force.\\
(b) As \(U=\alpha x^{2}+\beta y^{2}\)

So,

\[
F_{x}=\frac{\partial U}{\partial x}=-2 \alpha x \text { and } F_{y}=\frac{-\partial U}{\partial y}=-2 \beta y .
\]

So,

\[
F=\sqrt{F_{x}^{2}+F_{y}^{2}}=\sqrt{4 \alpha^{2} x^{2}+4 \beta^{2} y^{2}}
\]

According to the problem

\[
F=2 \sqrt{\alpha^{2} x^{2}+\beta^{2} y^{2}}=C \text { (constant) }
\]

or,

\[
\alpha^{2} x^{2}+\beta^{2} y^{2}=\frac{C^{2}}{2}
\]

or,


\begin{equation*}
\frac{x^{2}}{\beta^{2}}+\frac{y^{2}}{\alpha^{2}}=\frac{C^{2}}{2 \alpha^{2} \beta^{2}}=k(\mathrm{say}) \tag{2}
\end{equation*}


Therefore the surfaces for which \(F\) is constant is an ellipse.\\
For an equipotential surface \(U\) is constant.\\
So,

\[
\alpha x^{2}+\beta y^{2}=C_{0} \text { (constant) }
\]

or,

\[
\frac{x^{2}}{\sqrt{\beta^{2}}}+\frac{y^{2}}{\sqrt{\alpha^{2}}}=\frac{C_{0}}{\alpha \beta}=K_{0}(\text { constant })
\]

Hence the equipotential surface is also an ellipse.

\begin{enumerate}
  \item 133 Let us calculate the work performed by the forces of each field over the path from a certain point \(1\left(x_{1}, y_{1}\right)\) to another certain point \(2\left(x_{2}, y_{2}\right)\)\\
(i) \(d A=\vec{F} \cdot d \vec{r}=a y \vec{i} \cdot d \vec{r}=a y d x\) or, \(A=a \int_{x_{1}}^{x_{2}} y d x\).\\
(ii) \(d A=\vec{F} \cdot d \vec{r}=(\overrightarrow{a x i}+b y \vec{i}) \cdot d \vec{r}=a x d x+b y d y\)
\end{enumerate}

Hence

\[
A=\int_{x_{1}}^{x_{2}} a x d x+\int_{y_{1}}^{y_{2}} b y d y
\]

In the first case, the integral depends on the function of type \(y(x)\), i.e. on the shape of the path. Consequently, the first field of force is not potential. In the second case, both the integrals do not depend on the shape of the path. They are defined only by the coordinate of the initial and final points of the path, therefore the second field of force is potential.\\
1.134 Let \(s\) be the sought distance, then from the equation of increment of M.E. \(\Delta T+\Delta U=A_{f r}\)\\
or,

\[
\left(0-\frac{1}{2} m v_{0}^{2}\right)+(+m g s \sin \alpha)=-k m g \cos \alpha s
\]

,

\[
s=\frac{v_{0}^{2}}{2 g} /(\sin \alpha+k \cos \alpha)
\]

Hence

\[
A_{f r}=-k m g \cos \alpha s=\frac{-k m v_{0}^{2}}{2(k+\tan \alpha)}
\]

1.135 Velocity of the body at height \(h, v_{h}=\sqrt{2 g(H-h)}\), horizontally (from the figure given in the problem). Time taken in falling through the distance \(h\).\\
\(t=\sqrt{\frac{2 h}{g}}\) (as initial vertical component of the velocity is zero.)\\
Now

\[
s=v_{h} t=\sqrt{2 g(H+h)} \times \sqrt{\frac{2 h}{g}}=\sqrt{4\left(H h-h^{2}\right)}
\]

For \(s_{\max }, \frac{d}{d s}\left(H h-h^{2}\right)=0\), which yields \(h=\frac{H}{2}\)\\
Putting this value of \(h\) in the expression obtained for \(s\), we get,

\[
S_{\max }=H
\]

1.136 To complete a smooth vertical track of radius \(R\), the minimum height at which a particle starts, must be equal to \(\frac{5}{2} R\) (one can proved it from energy conservation). Thus in our problem body could not reach the upper most point of the vertical track of radius \(R / 2\). Let the particle \(A\) leave the track at some point \(O\) with speed \(v\) (Fig). Now from energy conservation for the body \(A\) in the field of gravity :

\[
m g\left[h-\frac{h}{2}(1+\sin \theta)\right]=\frac{1}{2} m v^{2}
\]

or,


\begin{equation*}
v^{2}=g h(1-\sin \theta) \tag{1}
\end{equation*}


From Newton's second law for the particle at the point \(O ; F_{n}=m w_{n}\),

\[
N+m g \sin \theta=\frac{m v^{2}}{(h / 2)}
\]

But, at the point \(O\) the normal reaction \(N=0\)\\
So, \(\quad v^{2}=\frac{g h}{2} \sin \theta\)\\
\includegraphics[max width=\textwidth, alt={}, center]{9f3c5a8e-ac5d-42f4-8eb6-ea2af623f2d8-079_350_621_793_797}

From (3) and (4), \(\sin \theta=\frac{2}{3}\) and \(v=\sqrt{\frac{g h}{3}}\)\\
After leaving the track at \(O\), the particle \(A\) comes in air and further goes up and at maximum height of it's trajectory in air, it's velocity (say \(v^{\prime}\) ) becomes horizontal (Fig.). Hence, the sought velocity of \(A\) at this point.

\[
v^{\prime}=v \cos (90-\theta)=v \sin \theta=\frac{2}{3} \sqrt{\frac{g h}{3}}
\]

1.137 Let, the point of suspension be shifted with velocity \(v_{A}\) in the horizontal direction towards left then in the rest frame of point of suspension the ball starts with same velocity horizontally towards right. Let us work in this, frame. From Newton's second law in projection form towards the point of suspension at the upper most point (say \(B\) ) :


\begin{equation*}
m g+T=\frac{m v_{B}^{2}}{l} \text { or, } \quad T=\frac{m v_{B}^{2}}{l}-m g \tag{1}
\end{equation*}


Condition required, to complete the vertical circle is that \(T \geq 0\). But


\begin{equation*}
\frac{1}{2} m v_{A}^{2}=m g(2 l)+\frac{1}{2} m v_{B}^{2} \text { So, } v_{B}^{2}=v_{A}^{2}-4 g l \tag{2}
\end{equation*}


From (1), (2) and (3)\\
\(T=\frac{m\left(v_{A}^{2}-4 g l\right)}{l}-m g \geq 0\) or, \(v_{A} \geq \sqrt{5 g l}\)\\
Thus \(\quad v_{A(\min )}=\sqrt{5 g l}\)\\
From the equation \(F_{n}=m w_{n}\) at point \(C\)


\begin{equation*}
T^{\prime}=\frac{m v_{c}^{2}}{l} \tag{4}
\end{equation*}


\begin{center}
\includegraphics[max width=\textwidth, alt={}]{9f3c5a8e-ac5d-42f4-8eb6-ea2af623f2d8-080_517_450_131_787}
\end{center}

From (4) and (5)


\begin{equation*}
\frac{1}{2} m v_{A}^{2}=\frac{1}{2} m v_{c}^{2}+m g l \tag{5}
\end{equation*}


Again from energy conservation

\[
T=3 \mathrm{mg}
\]

1.138 Since the tension is always perpendicular to the velocity vector, the work done by the tension force will be zero. Hence, according to the work energy theorem, the kinetic energy or velocity of the disc will remain constant during it's motion. Hence, the sought time \(t=\frac{s}{v_{0}}\), where \(s\) is the total distance traversed by the small disc during it's motion.\\
Now, at an arbitary position (Fig.)\\
\(d s=\left(l_{0}-R \theta\right) d \theta\),\\
so, \(s=\int_{0}^{L / R}\left(l_{0}-R \theta\right) d \theta\)\\
or,

\[
s=\frac{l_{0}^{2}}{R}-\frac{R l_{0}^{2}}{2 R^{2}}=\frac{l_{0}^{2}}{2 R}
\]

Hence, the required time, \(\quad t=\frac{l_{0}^{2}}{2 R v_{0}}\)\\
\includegraphics[max width=\textwidth, alt={}, center]{9f3c5a8e-ac5d-42f4-8eb6-ea2af623f2d8-080_344_592_1101_742}

It should be clearly understood that the only uncompensated force acting on the disc A in this case is the tension T, of the thread. It is easy to see that there is no point here, relative to which the moment of force \(T\) is invarible in the process of motion. Hence conservation of angular momentum is not applicable here.

Suppose that \(\Delta l\) is the elongation of the rubbler cord. Then from energy conservation,

\[
\Delta U_{g r}+\Delta U_{e l}=0(\text { as } \Delta T=0)
\]

or, \(\quad-m g(l+\Delta l)+\frac{1}{2} \kappa \Delta l^{2}=0\)\\
or, \(\quad \frac{1}{2} \kappa \Delta l^{2}-m g \Delta l-m g l=0\)\\
or, \(\quad \Delta l=\frac{m g \pm \sqrt{(m g)^{2}+4 \times \frac{\kappa}{2} m g l}}{2 \times \frac{\kappa}{2}} \times \frac{\kappa}{2}=\frac{m g}{\kappa}\left[1+\sqrt{1 \pm \frac{2 \kappa l}{m g}}\right]\)\\
Since the value of \(\sqrt{1+\frac{2 \kappa l}{m g}}\) is certainly greater than 1 , hence negative sign is avoided.

So,

\[
\Delta l=\frac{m g}{\kappa}\left(1+\sqrt{1+\frac{2 \kappa l}{m g}}\right)
\]

1.140 When the thread PA is burnt, obviously the speed of the bars will be equal at any instant of time until it breaks off. Let \(v\) be the speed of each block and \(\theta\) be the angle, which the elongated spring makes with the vertical at the moment, when the bar \(A\) breaks off the plane. At this stage the elongation in the spring.


\begin{equation*}
\Delta l=l_{0} \sec \theta-l_{0}=l_{0}(\sec \theta-1) \tag{1}
\end{equation*}


Since the problem is concerned with position and there are no forces other than conservative forces, the mechanical energy of the system (both bars + spring) in the field of gravity is conserved, i.e. \(\Delta T+\Delta U=0\)

So,


\begin{equation*}
2\left(\frac{1}{2} \dot{m} v^{2}\right)+\frac{1}{2} \kappa l_{0}^{2}(\sec \theta-1)^{2}-m g l_{0} \tan \theta=0 \tag{2}
\end{equation*}


From Newton's second law in projection form along vertical direction :\\
\(m g=N+\kappa l_{0}(\sec \theta-1) \cos \theta\)\\
But, at the moment of break off, \(N=0\).\\
Hence, \(\kappa l_{0}(\sec \theta-1) \cos \theta=m g\)\\
or, \(\quad \cos \theta=\frac{\kappa l_{0}-m g}{\kappa l_{0}}\)

Taking \(\kappa=\frac{5 m g}{l_{0}}\), simultaneous solution of (2) and (3) yields :\\
\includegraphics[max width=\textwidth, alt={}, center]{9f3c5a8e-ac5d-42f4-8eb6-ea2af623f2d8-081_400_440_928_915}

\[
v=\sqrt{\frac{19 g l_{0}}{32}}=1.7 \mathrm{~m} / \mathrm{s} .
\]

1.141 Obviously the elongation in the cord, \(\Delta l=l_{0}(\sec \theta-1)\), at the moment the sliding first starts and at the moment horizontal projection of spring force equals the limiting friction.\\
So,


\begin{equation*}
\kappa_{1} \Delta l \sin \theta=k N \tag{1}
\end{equation*}


(where \(\kappa_{1}\) is the elastic constant).\\
From Newton's law in projection form along vertical direction :

\[
\kappa_{1} \Delta l \cos \theta+N=m g .
\]

or,


\begin{equation*}
N=m g-\kappa_{1} \Delta l \cos \theta \tag{2}
\end{equation*}


From (1) and (2),

\[
\kappa_{1} \Delta l \sin \theta=k\left(m g-\kappa_{1} \Delta l \cos \theta\right)
\]

\includegraphics[max width=\textwidth, alt={}, center]{9f3c5a8e-ac5d-42f4-8eb6-ea2af623f2d8-081_398_415_1605_949}\\
or, \(\quad \kappa_{1}=\frac{k m g}{\Delta l \sin \theta+k \Delta l \cos \theta}\)\\
From the equation of the increment of mechanical energy : \(\Delta U+\Delta T=A_{f r}\)\\
or,

\[
\left(\frac{1}{2} \kappa_{1} \Delta l^{2}\right)=A_{f r}
\]

or, \(\frac{k m g \Delta l^{2}}{2 \Delta l(\sin \theta+k \cos \theta)}=A_{f r}\)\\
Thus \(A_{f r}=\frac{k m g l_{0}(\sec \theta-1)}{2(\sin \theta-k \cos \theta)}=0.09 \mathrm{~J}\) (on substitution)\\
1.142 Let the deformation in the spring be \(\Delta l\), when the \(\operatorname{rod} A B\) has attained the angular velocity \(\omega\). From the second law of motion in projection form \(F_{n}=m w_{n}\).\\
\(\kappa \Delta l=m \omega^{2}\left(l_{0}+\Delta l\right)\) or, \(\Delta l=\frac{m \omega^{2} l_{0}}{\kappa-m \omega^{2}}\)\\
From the energy equation, \(A_{e x t}=\frac{1}{2} m v^{2}+\frac{1}{2} \kappa \Delta l^{2}\)

\[
\begin{aligned}
& =\frac{1}{2} m \omega^{2}\left(l_{0}+\Delta l\right)^{2}+\frac{1}{2} \kappa \Delta l^{2} \\
& =\frac{1}{2} m \omega^{2}\left(l_{0}+\frac{m \omega^{2} l_{0}}{\kappa-m \omega^{2}}\right)^{2}+\frac{1}{2} \kappa\left(\frac{m \omega^{2} l_{0}^{2}}{\kappa-m \omega^{2}}\right)^{2}
\end{aligned}
\]

On solving \(\quad A_{e x t}=\frac{\kappa}{2} \frac{l_{0}^{2} \eta\left(1+r_{c}\right)}{(1-\eta)^{2}}\), where \(\eta=\frac{m \omega^{2}}{\kappa}\)\\
1.143 We know that acceleration of centre of mass of the system is given by the expression.

\[
\vec{w}_{C}=\frac{m_{1} \vec{w}_{1}+m_{2} \vec{w}_{2}}{m_{1}+m_{2}}
\]

Since


\begin{gather*}
\vec{w}_{1}=-\vec{w}_{2} \\
\vec{w}_{C}=\frac{\left(m_{1}-m_{2}\right) \vec{w}_{1}}{m_{1}+m_{2}} \tag{1}
\end{gather*}


Now from Newton's second law \(\vec{F}=m \vec{w}\), for the bodies \(m_{1}\) and \(m_{2}\) respectively.


\begin{equation*}
\vec{T}+m_{1} \vec{g}=m_{1} \overrightarrow{w_{1}} \tag{2}
\end{equation*}


and


\begin{equation*}
\vec{T}+m_{2} \vec{g}=m_{2} \overrightarrow{w_{2}}=-m_{2} \overrightarrow{w_{1}} \tag{3}
\end{equation*}


Solving (2) and (3)


\begin{equation*}
\overrightarrow{u_{1}}=\frac{\left(m_{1}-m_{2}\right) \vec{g}}{m_{1}+m_{2}} \tag{4}
\end{equation*}


\begin{center}
\includegraphics[max width=\textwidth, alt={}]{9f3c5a8e-ac5d-42f4-8eb6-ea2af623f2d8-082_620_358_1358_863}
\end{center}

Thus from (1), (2) and (4),

\[
\vec{w}_{c}=\frac{\left(m_{1}-m_{2}\right)^{2} \vec{g}}{\left(m_{1}+m_{2}\right)^{2}}
\]

1.144 As the closed system consisting two particles \(m_{1}\) and of \(m_{2}\) is initially at rest the C.M. of the system will remain at rest. Further as \(m_{2}=m_{1} / 2\), the C.M. of the system divides the line joining \(m_{1}\) and \(m_{2}\) at all the moments of time in the ratio \(1: 2\). In addition to it the total linear momentum of the system at all the times is zero. So, \(\overrightarrow{p_{1}}=-\overrightarrow{p_{2}}\) and therefore the velocities of \(m_{1}\) and \(m_{2}\) are also directed in opposite sense. Bearing in mind all these thing, the sought trajectory is as shown in the figure.\\
\includegraphics[max width=\textwidth, alt={}, center]{9f3c5a8e-ac5d-42f4-8eb6-ea2af623f2d8-083_487_350_391_937}\\
1.145 First of all, it is clear that the chain does not move in the vertical direction during the uniform rotation. This means that the vertical component of the tension \(T\) balances gravity. As for the horizontal component of the tension \(T\), it is constant in magnitude and permanently directed toward the rotation axis. It follows from this that the C.M. of the chain, the point \(C\), travels along horizontal circle of radius \(\rho\) (say). Therefore we have,

\[
T \cos \theta=m g \text { and } T \sin \theta=m \omega^{2} \rho
\]

Thus \(\rho=\frac{g \tan \theta}{\omega^{2}}=0.8 \mathrm{~cm}\)\\
\includegraphics[max width=\textwidth, alt={}, center]{9f3c5a8e-ac5d-42f4-8eb6-ea2af623f2d8-083_474_319_943_929}\\
and \(\quad T=\frac{m g}{\cos \theta}=5 \mathrm{~N}\)\\
1.146 (a) Let us draw free body diagram and write Newton's second law in terms of projection along vertical and horizontal direction respectively.


\begin{gather*}
N \cos \alpha-m g+f r \sin \alpha=0  \tag{1}\\
f r \cos \alpha-N \sin \alpha=m \omega^{2} l \tag{2}
\end{gather*}


From (1) and (2)\\
\(f r \cos \alpha-\frac{\sin \alpha}{\cos \alpha}(-f r \sin \alpha+m g)=m \omega^{2} l\)\\
\includegraphics[max width=\textwidth, alt={}, center]{9f3c5a8e-ac5d-42f4-8eb6-ea2af623f2d8-083_472_490_1576_905}

So,


\begin{equation*}
f r=m g\left(\sin \alpha+\frac{\omega^{2} l}{g} \cos \alpha\right)=6 \mathrm{~N} \tag{3}
\end{equation*}


(b) For rolling, without sliding,

\[
f r \leq k \mathrm{~N}
\]

but, \(N=m g \cos \alpha-m \omega^{2} l \sin \alpha\)

\[
m g\left(\sin \alpha+\frac{\omega^{2} l}{g} \cos \alpha\right) \leq k\left(m g \cos \alpha-m \omega^{2} l \sin \alpha\right) \text { [Using (3)] }
\]

Rearranging, we get,

Thus

\[
\begin{gathered}
m \omega^{2} l(\cos \alpha+k \sin \alpha) \leq(k m g \cos \alpha-m g \sin \alpha) \\
\omega \leq \sqrt{g(k-\tan \alpha) /(1+k \tan \alpha) l}=2 \mathrm{rad} / \mathrm{s}
\end{gathered}
\]

1.147 (a) Total kinetic energy in frame \(K^{\prime}\) is

\[
T=\frac{1}{2} m_{1}\left(\overrightarrow{v_{1}}-\vec{V}\right)^{2}+\frac{1}{2} m_{2}\left(\overrightarrow{v_{2}}-\vec{V}\right)^{2}
\]

This is minimum with respect to variation in \(\vec{V}\), when\\
or

\[
\begin{gathered}
\frac{\delta T^{\prime}}{\delta \vec{V}}=0, \text { i.e. } m_{1}\left(\overrightarrow{v_{1}}-\vec{V}\right)^{2}+m_{2}\left(\overrightarrow{v_{2}}-\vec{V}\right)=0 \\
\vec{V}=\frac{m_{1} \overrightarrow{v_{1}}+m_{2} \overrightarrow{V_{2}}}{m_{1}+m_{2}}=\overrightarrow{v_{c}}
\end{gathered}
\]

Hence, it is the frame of C.M. in which kinetic energy of a system is minimum.\\
(b) Linear momentum of the particle 1 in the \(K^{\prime}\) or \(C\) frame

\[
\tilde{\overrightarrow{p_{1}}}=m_{1}\left(\overrightarrow{v_{1}}-\overrightarrow{v_{c}}\right)=\frac{m_{1} m_{2}}{m_{1}+m_{2}}\left(\overrightarrow{v_{1}}-\overrightarrow{v_{2}}\right)
\]

or,

\[
\tilde{\overrightarrow{p_{1}}}=\mu\left(\overrightarrow{v_{1}}-\overrightarrow{v_{2}}\right), \text { where, } \mu=\frac{m_{1} m_{2}}{m_{1}+m_{2}}=\text { reduced mass }
\]

Similarly, \(\quad \overrightarrow{\overrightarrow{p_{2}}}=\mu\left(\overrightarrow{v_{2}}-\overrightarrow{v_{1}}\right)\)\\
So,


\begin{equation*}
\left|\widetilde{\overrightarrow{p_{1}}}\right|=\left|\widetilde{\overrightarrow{p_{2}}}\right|=\tilde{p}=\mu v_{\text {rel }} \text { where, } v_{\text {rel }}=\left|\overrightarrow{v_{1}}-\overrightarrow{v_{2}}\right| \tag{3}
\end{equation*}


Now the total kinetic energy of the system in the \(C\) frame is

Hence

\[
\begin{aligned}
& \tilde{T}=\tilde{T}_{1}+\tilde{T}_{2}=\frac{\tilde{p}^{2}}{2 m_{1}}+\frac{\tilde{p}^{2}}{2 m_{2}}=\frac{\tilde{p}^{2}}{2 \mu} \\
& \tilde{T}=\frac{1}{2} \mu v_{r e l}^{2}=\frac{1}{2} \mu\left|\overrightarrow{v_{1}}-\overrightarrow{v_{2}}\right|^{2}
\end{aligned}
\]

1.148 To find the relationship between the values of the mechanical energy of a system in the \(K\) and \(C\) reference frames, let us begin with the kinetic energy \(T\) of the system. The velocity of the \(i\)-th particle in the \(K\) frame may be represented as \(\overrightarrow{v_{i}}=\overrightarrow{v_{i}}+\overrightarrow{v_{c}}\). Now we can write

\[
\begin{aligned}
T & =\sum \frac{1}{2} m_{i} v_{i}^{2}=\sum \frac{1}{2} m_{i}\left(\tilde{v_{1}}+\overrightarrow{v_{C}}\right) \cdot\left(\tilde{v_{i}}+\overrightarrow{v_{C}}\right) \\
& =\sum \frac{1}{2} m_{i} \tilde{v}_{i}^{2}+\overrightarrow{v_{C}} \sum m_{i} \tilde{v_{1}}+\sum \frac{1}{2} m_{i} v_{C}^{2}
\end{aligned}
\]

Since in the \(C\) frame \(\sum m_{i} \overrightarrow{v_{i}}=0\), the previous expression takes the form


\begin{equation*}
\left.T=\tilde{T}+\frac{1}{2} m v_{C}^{2}=\tilde{T}+\frac{1}{2} m V^{2} \text { (since according to the problem } v_{C}=V\right) \tag{1}
\end{equation*}


Since the internal potential energy \(U\) of a system depends only on its configuration, the magnitude \(U\) is the same in all refrence frames. Adding \(U\) to the left and right hand sides of Eq. (1), we obtain the sought relationship

\[
E=\tilde{E}+\frac{1}{2} m V^{2}
\]

1.149 As initially \(U=\tilde{U}=0\), so, \(\tilde{E}=\tilde{T}\)

From the solution of 1.147 (b)

\[
\tilde{T}=\frac{1}{2} \mu\left|\overrightarrow{v_{1}}-\overrightarrow{v_{2}}\right|
\]

As

\[
\overrightarrow{v_{1}} \perp \overrightarrow{v_{2}}
\]

Thus

\[
\tilde{T}=\frac{1}{2} \frac{m_{1} m_{2}}{m_{1}+m_{2}}\left(v_{1}^{2}+v_{2}^{2}\right)
\]

1.150 Velocity of masses \(m_{1}\) and \(m_{2}\), after \(t\) seconds are respectively.

\[
\overrightarrow{v_{1}}=\overrightarrow{v_{1}}+\overrightarrow{g t} \text { and } \overrightarrow{v_{2}}=\overrightarrow{v_{2}}+\overrightarrow{g t}
\]

Hence the final momentum of the system,

\[
\begin{aligned}
\vec{p} & =m_{1} \overrightarrow{v_{1}}+m_{2} \overrightarrow{v_{2}}=m_{1} \overrightarrow{v_{1}}+m_{2} \overrightarrow{v_{2}}+\left(m_{1}+m_{2}\right) \vec{g} t \\
& =\overrightarrow{p_{0}}+m \overrightarrow{g t},\left(\text { where, } \overrightarrow{p_{0}}=m_{2} \overrightarrow{v_{1}}+m_{2} \overrightarrow{v_{2}} \text { and } m=m_{1}+m_{2}\right)
\end{aligned}
\]

And radius vector,

\[
\overrightarrow{r_{c}}=\overrightarrow{v_{c}} t+\frac{1}{2} \overrightarrow{w_{c}} t^{2}
\]

\[
\begin{gathered}
\frac{\left(m_{1} \overrightarrow{v_{1}}+m_{2} \overrightarrow{v_{2}}\right) t}{\left(m_{1}+m_{2}\right)}+\frac{1}{2} \vec{g}^{2} \\
=\overrightarrow{v_{0}} t+\frac{1}{2} \vec{g} t^{2}, \text { where } \overrightarrow{v_{0}}=\frac{m_{1} \overrightarrow{v_{1}}+m_{2} \overrightarrow{v_{2}}}{m_{1}+m_{2}}
\end{gathered}
\]

1.151 After releasing the bar 2 acquires the velocity \(v_{2}\), obtained by the energy, conservation :


\begin{equation*}
\frac{1}{2} m_{2} v_{2}^{2}=\frac{1}{2} \kappa x^{2} \text { or, } v_{2}=x \sqrt{\frac{\kappa}{m_{2}}} \tag{1}
\end{equation*}


Thus the sought velocity of C.M.

\[
v_{c m}=\frac{0+m_{2} x \sqrt{\frac{\kappa}{m_{2}}}}{m_{1}+m_{2}}=\frac{x \sqrt{m_{2} \kappa}}{\left(m_{1}+m_{2}\right)}
\]

1.152 Let us consider both blocks and spring as the physical system. The centre of mass of the system moves with acceleration \(a=\frac{F}{m_{1}+m_{2}}\) towards right. Let us work in the frame of centre of mass. As this frame is a non-inertial frame (accelerated with respect to the ground) we have to apply a pseudo force \(m_{1} a\) towards left on the block \(m_{1}\) and \(m_{2} a\) towards left on the block \(m_{2}\)\\
As the center of mass is at rest in this frame, the blocks move in opposite directions and come to instantaneous rest at some instant. The elongation of the spring will be maximum or minimum at this instant. Assume that the block \(m_{1}\) is displaced by the distance \(x_{1}\) and the block\\
\includegraphics[max width=\textwidth, alt={}, center]{9f3c5a8e-ac5d-42f4-8eb6-ea2af623f2d8-086_164_545_853_812}\\
\(m_{2}\) through a distance \(x_{2}\) from the initial positions.\\
From the energy equation in the frame of C.M.

\[
\Delta \tilde{T}+U=A_{e x t},
\]

(where \(A_{\text {ext }}\) also includes the work done by the pseudo forces)\\
Here,\\
\(\Delta \tilde{T}=0, U=\frac{1}{2} k\left(x_{1}+x_{2}\right)^{2}\) and\\
\(W_{\text {ext }}=\left(\frac{F-m_{2} F}{m_{1}+m_{2}}\right) x_{2}+\frac{m_{1} F}{m_{1}+m_{2}} x_{1}=\frac{m_{1} F\left(x_{1}+x_{2}\right)}{m_{1}+m_{2}}\),\\
or,

\[
\frac{1}{2} k\left(x_{1}+x_{2}\right)^{2}=\frac{m_{1}\left(x_{1}+x_{2}\right) F}{m_{1}+m_{2}}
\]

So,

\[
x_{1}+x_{2}=0 \text { or, } x_{1}+x_{2}=\frac{2 m_{1} F}{k\left(m_{1}+m_{2}\right)}
\]

Hence the maximum separation between the blocks equals : \(l_{0}+\frac{2 m_{1} F}{k\left(m_{1}+m_{2}\right)}\)\\
Obviously the minimum sepation corresponds to zero elongation and is equal to \(l_{0}\)\\
1.153 (a) The initial compression in the spring \(\Delta l\) must be such that after burning of the thread, the upper cube rises to a height that produces a tension in the spring that is atleast equal to the weight of the lower cube. Actually, the spring will first go from its compressed\\
state to its natural length and then get elongated beyond this natural length. Let \(l\) be the maximum elongation produced under these circumstances.\\
Then


\begin{equation*}
\kappa l=m g \tag{1}
\end{equation*}


Now, from energy conservation,


\begin{equation*}
\frac{1}{2} \kappa \Delta l^{2}=m g(\Delta l+l)+\frac{1}{2} \kappa l^{2} \tag{2}
\end{equation*}


(Because at maximum elongation of the spring, the speed of upper cube becomes zero) From (1) and (2),

\[
\Delta l^{2}-\frac{2 m g \Delta l}{\kappa}-\frac{3 m^{2} g^{2}}{\kappa^{2}}=0 \text { or, } \Delta l=\frac{3 m g}{\kappa}, \frac{-m g}{\kappa}
\]

Therefore, acceptable solution of \(\Delta l\) equals \(\frac{3 m g}{\kappa}\)\\
(b) Let \(v\) the velocity of upper cube at the position (say, at \(C\) ) when the lower block breaks off the floor, then from energy conservation.


\begin{gather*}
\frac{1}{2} m v^{2}=\frac{1}{2} \kappa\left(\Delta l^{2}-l^{2}\right)-m g(l+\Delta l) \\
\left(\text { where } l=m g / \kappa \text { and } \Delta l=7 \frac{m g}{\kappa}\right) \\
v^{2}=32 \frac{m g^{2}}{\kappa} \tag{2}
\end{gather*}


At the position \(C\), the velocity of \(C . M ; v_{C}=\frac{m v+0}{2 m}=\frac{v}{2}-L e t\), the C.M. of the system (spring+ two cubes) further rises up to \(\Delta y_{C 2}\).\\
Now, from energy conservation,

\[
\frac{1}{2}(2 m) v_{C}^{2}=(2 m) g \Delta r_{C 2}
\]

or, \(\Delta y_{C 2}=\frac{v_{C}^{2}}{2 g}=\frac{v^{2}}{8 g}=\frac{4 m g}{\kappa}\)\\
But, uptil position C, the C.M. of the system has already elevated by,

\[
\Delta y_{C 1}=\frac{(\Delta l+l) m+0}{2 m}=\frac{4 m g}{\kappa}
\]

Hence, the net displacement of the C.M. of the system, in upward direction

\[
\Delta y_{C}=\Delta y_{C 1}+\Delta y_{C 2}=\frac{8 m g}{\kappa}
\]

\includegraphics[max width=\textwidth, alt={}, center]{9f3c5a8e-ac5d-42f4-8eb6-ea2af623f2d8-087_667_415_1170_980}\\
1.154 Due to ejection of mass from a moving system (which moves due to inertia) in a direction perpendicular to it, the velocity of moving system does not change. The momentum change being adjusted by the forces on the rails. Hence in our problem velocities of buggies change only due to the entrance of the man coming from the other buggy. From the

Solving (1) and (2), we get

\[
v_{1}=\frac{m v}{M-m} \text { and } v_{2}=\frac{M v}{M-m}
\]

As

\[
\overrightarrow{v_{1}} \uparrow \downarrow \vec{v} \text { and } \overrightarrow{v_{2}} \uparrow \uparrow \vec{v}
\]

So,

\[
\overrightarrow{v_{1}}=\frac{-m \vec{v}}{(M-m)} \text { and } \overrightarrow{v_{2}}=\frac{M \vec{v}}{(M-m)}
\]

1.155 From momentum conservation, for the system "rear buggy with man"


\begin{equation*}
(M+m) \overrightarrow{v_{0}}=m\left(\vec{u}+\overrightarrow{v_{R}}\right)+M \overrightarrow{v_{R}} \tag{1}
\end{equation*}


From momentum conservation, for the system (front buggy + man coming from rear buggy)

\[
M \overrightarrow{v_{0}}+m\left(\vec{u}+\overrightarrow{v_{R}}\right)=(M+m) \overrightarrow{v_{F}}
\]

So,

\[
\overrightarrow{v_{F}}=\frac{M \overrightarrow{v_{0}}}{M+m}+\frac{m}{M+m}\left(\vec{u}+\overrightarrow{v_{R}}\right)
\]

Putting the value of \(\overrightarrow{v_{R}}\) from (1), we get

\[
\overrightarrow{v_{F}}=\overrightarrow{v_{0}}+\frac{m M}{(M+m)^{2}} \vec{u}
\]

1.156 (i) Let \(\overrightarrow{v_{1}}\) be the velocity of the buggy after both man jump off simultaneously. For the closed system (two men + buggy), from the conservation of linear momentum,\\
or,


\begin{gather*}
M \overrightarrow{v_{1}}+2 m\left(\vec{u}+\overrightarrow{v_{1}}\right)=0 \\
\overrightarrow{v_{1}}=\frac{-2 m \vec{u}}{M+2 m} \tag{1}
\end{gather*}


(ii) Let \(\vec{v}^{\prime}\) be the velocity of buggy with man, when one man jump off the buggy. For the closed system (buggy with one man + other man) from the conservation of linear momentum :


\begin{equation*}
0=(M+m) \vec{v}^{\prime}+m\left(\vec{u}+\overrightarrow{v^{\prime}}\right) \tag{2}
\end{equation*}


Let \(\overrightarrow{v_{2}}\) be the sought velocity of the buggy when the second man jump off the buggy; then from conservation of linear momentum of the system (buggy + one man) :


\begin{equation*}
(M+m) \overrightarrow{v^{\prime}}=M \overrightarrow{v_{2}}+m\left(\vec{u}+\overrightarrow{v_{2}}\right) \tag{3}
\end{equation*}


Solving equations (2) and (3) we get


\begin{equation*}
\overrightarrow{v_{2}}=\frac{m(2 M+3 m) \vec{u}}{(M+m)(M+2 m)} \tag{4}
\end{equation*}


From (1) and (4)

\[
\frac{v_{2}}{v_{1}}=1+\frac{m}{2(M+m)}>1
\]

Hence \(\boldsymbol{v}_{\mathbf{2}}>\boldsymbol{v}_{\mathbf{1}}\)\\
1.157 The descending part of the chain is in free fall, it has speed \(v=\sqrt{2 g h}\) at the instant, all its points have descended a distance \(y\). The length of the chain which lands on the floor during the differential time interval \(d t\) following this instant is \(v d t\).

For the incoming chain element on the floor :\\
From \(d p_{y}=F_{v} d t\) (where \(y\) - axis is directed down)\\
or

\[
\begin{gathered}
0-(\lambda v d t) v=F_{y} d t \\
F_{y}=-\lambda v^{2}=-2 \lambda g y
\end{gathered}
\]

Hence, the force exerted on the falling chain equals \(\lambda \nu^{2}\) and is directed upward. Therefore from third law the force exerted by the falling chain on the table at the same instant of time becomes \(\lambda \nu^{2}\) and is directed downward.\\
\includegraphics[max width=\textwidth, alt={}, center]{9f3c5a8e-ac5d-42f4-8eb6-ea2af623f2d8-089_418_274_111_1011}

Since a length of chain of weight ( \(\lambda y g\) ) already lies on the table the total force on the floor is \((2 \lambda y g)+(\lambda y g)=(3 \lambda y g)\) or the weight of a length \(3 y\) of chain.\\
1.158 Velocity of the ball, with which it hits the slab, \(v=\sqrt{2 g h}\)

After first impact, \(v^{\prime}=e \nu\) (upward) but according to the problem \(v^{\prime}=\frac{\nu}{\eta}\), so \(e=\frac{1}{\eta}\)\\
and momentum, imparted to the slab,

\[
=m v-\left(-m v^{\prime}\right)=m v(1+e)
\]

Similarly, velocity of the ball after second impact,

\[
v^{\prime \prime}=e v^{\prime}=e^{2} v
\]

And momentum imparted \(=m\left(v^{\prime}+v^{\prime \prime}\right)=m(1+e) e v\)\\
Again, momentum imparted during third impact,

\[
=m(1+e) e^{2} v, \text { and so on, }
\]

Hence, net momentum, imparted \(=m v(1+e)+m v e(1+e)+m v e^{2}(1+e)+\ldots\)

\[
\begin{gathered}
=m v(1+e)\left(1+e+e^{2}+\ldots\right) \\
=m v \frac{(1+e)}{(1-e)},(\text { from summation of G.P. }) \\
=\sqrt{2 g h} \frac{\left(1+\frac{1}{\eta}\right)}{\left(1-\frac{1}{\eta}\right)}=m \sqrt{2 g h} /(\eta+1) /(\eta-1)(\text { Using Eq. } 1) \\
=0.2 \mathrm{~kg} \mathrm{~m} / \mathrm{s} . \text { (On substitution) }
\end{gathered}
\]

1.159 (a) Since the resistance of water is negligibly small, the resultant of all external forces acting on the system "a man and a raft" is equal to zero. This means that the position of the C.M. of the given system does not change in the process of motion.\\
i.e. \(\overrightarrow{r_{C}}=\) constant or, \(\quad \Delta \overrightarrow{r_{C}}=0\) i.e. \(\sum m_{i} \Delta \overrightarrow{r_{i}}=0\)\\
or,

\[
m\left(\Delta \vec{r}_{m M}+\Delta \vec{r}_{M}\right)+M \Delta \vec{r}_{M}=0
\]

Thus,

\[
m\left(\vec{l}^{\prime}+\vec{l}\right)+M \vec{l}=0, \text { or, } \vec{l}=-\frac{m \vec{l}^{\prime}}{m+M}
\]

(b) As net external force on "man-raft" system is equal to zero, therefore the momentum of this system does not change,\\
So,

\[
0=m\left[\overrightarrow{v^{\prime}}(t)+\overrightarrow{v_{2}}(t)\right]+M \overrightarrow{v_{2}}(t)
\]

1.159 (a) Since the resistance of water is negligibly small, the resultant of all external forces acting on the system "a man and a raft" is equal to zero. This means that the position of the C.M. of the given system does not change in the process of motion.\\
i.e. \(\overrightarrow{r_{C}}=\) constant or, \(\quad \Delta \overrightarrow{r_{C}}=0\) i.e. \(\sum m_{i} \Delta \overrightarrow{r_{i}}=0\)\\
or,

\[
m\left(\Delta \vec{r}_{m M}+\Delta \overrightarrow{r_{M}}\right)+M \Delta \overrightarrow{r_{M}}=0
\]

Thus,

\[
m\left(\vec{l}^{\prime}+\vec{l}\right)+M \vec{l}=0, \text { or, } \vec{l}=-\frac{m \vec{l}^{\prime}}{m+M}
\]

(b) As net external force on "man-raft" system is equal to zero, therefore the momentum of this system does not change,

So,

\[
0=m\left[\overrightarrow{v^{\prime}}(t)+\overrightarrow{v_{2}}(t)\right]+M \overrightarrow{v_{2}}(t)
\]

or,


\begin{equation*}
\vec{v}_{2}(t)=-\frac{m \vec{v}(t)}{m+M} \tag{1}
\end{equation*}


As \(\vec{v}^{\prime}(t)\) or \(\overrightarrow{v_{2}}(t)\) is along horizontal direction, thus the sought force on the raft

\[
=\frac{M d \vec{v}_{2}(t)}{d t}=-\frac{M m}{m+M} \frac{d \vec{v}^{\prime}(t)}{d t}
\]

Note : we may get the result of part (a), if we integrate Eq. (1) over the time of motion of man or raft.\\
1.160 In the refrence frame fixed to the pulley axis the location of C.M. of the given system is described by the radius vector

\[
\Delta \vec{r}_{C}=\frac{M \Delta \vec{r}_{M}+(M-m) \Delta \vec{r}_{(M-m)}+m \Delta \vec{r}_{m}}{2 M}
\]

But \(\quad \Delta \overrightarrow{r_{M}}=-\Delta \overrightarrow{r_{(M-m)}}\)\\
and \(\quad \Delta \overrightarrow{r_{m}}=\Delta \overrightarrow{r_{m}(M-m)}+\Delta \overrightarrow{r_{(M-m)}}\)\\
Thus \(\quad \Delta \overrightarrow{r_{c}}=\frac{m \vec{l}}{2 M}\)\\
\includegraphics[max width=\textwidth, alt={}, center]{9f3c5a8e-ac5d-42f4-8eb6-ea2af623f2d8-090_431_348_1014_930}

Note : one may also solve this problem using momentum conservation.\\
1.161 Velocity of cannon as well as that of shell equals \(\sqrt{2 g l \sin \alpha}\) down the inclined plane taken as the positive \(x\)-axis. From the linear impulse momentum theorem in projection form along \(x\)-axis for the system (connon + shell) i.e. \(\Delta p_{x}=F_{x} \Delta t\) :\\
\(p \cos \alpha-M \sqrt{2 g l \sin \alpha}=M g \sin \alpha \Delta t\) (as mass of the shell is neligible)\\
or,

\[
\Delta t=\frac{p \cos \alpha-M \sqrt{2 g l \sin \alpha}}{M g \sin \alpha}
\]

1.162 From conservation of momentum, for the system (bullet + body) along the initial direction of bullet

\[
m v_{0}=(m+M) v, \text { or, } v=\frac{m v_{0}}{m+M}
\]

1.163 When the disc breaks off the body \(M\), its velocity towards right (along \(x\)-axis) equals the velocity of the body \(M\), and let the disc's velocity' in upward direction (along \(y\)-axis) at that moment be \(v_{y}^{\prime}\)\\
From conservation of momentum, along \(x\)-axis for the system (disc + body)


\begin{equation*}
m v=(m+M) v_{x}^{\prime} \text { or } v_{x}^{\prime}=\frac{m v}{m+M} \tag{1}
\end{equation*}


And from energy conservation, for the same system in the field of gravity :

\[
\frac{1}{2} m v^{2}=\frac{1}{2}(m+M) v_{x}^{\prime 2}+\frac{1}{2} m v_{y}^{\prime 2}+m g h^{\prime},
\]

where \(h^{\prime}\) is the height of break off point from initial level. So,

\[
\begin{array}{r}
\frac{1}{2} m v^{2}=\frac{1}{2}(m+M) \frac{m^{2} v^{2}}{(M+m)}+\frac{1}{2} m v_{y}^{\prime 2}+m g h^{\prime}  \tag{1}\\
\text { or, } \quad v_{y}^{\prime 2}=v^{2}-\frac{m v^{2}}{(m+M)}-2 g h^{\prime}
\end{array}
\]

Also, if \(h^{\prime \prime}\) is the height of the disc, from the break-off point,\\
then,

\[
v_{y}^{\prime 2}=2 g h^{\prime \prime}
\]

So,

\[
2 g\left(h^{\prime \prime}+h^{\prime}\right)=v^{2}-\frac{m v^{2}}{(M+m)}
\]

Hence, the total height, raised from the initial level

\[
=h^{\prime}+h^{\prime \prime}=\frac{M v^{2}}{2 g(M+m)}
\]

1.164 (a) When the disc slides and comes to a plank, it has a velocity equal to \(v=\sqrt{2 g h}\). Due to friction between the disc and the plank the disc slows down and after some time the disc moves in one piece with the plank with velocity \(v^{\prime}\) (say).\\
From the momentum conservation for the system (disc + plank) along horizontal towards right :

\[
m v=(m+M) v^{\prime} \text { or } v^{\prime}=\frac{m v}{m+M}
\]

Now from the equation of the increment of total mechanical energy of a system :\\
or,\\
so,

Hence,

\[
\begin{gathered}
\frac{1}{2}(M+m) v^{\prime 2}-\frac{1}{2} m v^{2}=A_{f r} \\
\frac{1}{2}(M+m) \frac{m^{2} v^{2}}{(m+M)^{2}}-\frac{1}{2} m v^{2}=A_{f r} \\
\frac{1}{2} v^{2}\left[\frac{m^{2}}{M+m}-m\right]=A_{f r} \\
A_{f r}=-\left(\frac{m M}{m+M}\right) g h=-\mu g h \\
\left(\text { where } \mu=\frac{m M}{m+M}=\text { reduced mass }\right)
\end{gathered}
\]

(b) We look at the problem from a frame in which the hill is moving (together with the disc on it) to the right with speed \(u\). Then in this frame the speed of the disc when it just gets onto the plank is, by the law of addition of velocities, \(\bar{v}=u+\sqrt{2 g h}\). Similarly the common speed of the plank and the disc when they move together is

\[
\vec{v}=u+\frac{m}{m+M} \sqrt{2 g h}
\]

Then as above

\[
\bar{A}_{f r}=\frac{1}{2}(m+M) \bar{v}^{2}-\frac{1}{2} m \bar{v}^{2}-\frac{1}{2} M u^{2}
\]

\[
=\frac{1}{2}(m+M)\left\{u^{2}+\frac{2 m}{m+M} u \sqrt{2 g h}+\frac{m^{2}}{(m+M)^{2}} 2 g h\right\}-\frac{1}{2}(m+M) u^{2}-\frac{1}{2} m 2 u \sqrt{2 g h}-m g h
\]

We see that \(\bar{A}_{f r}\) is independent of \(u\) and is in fact just - \(\mu g h\) as in (a). Thus the result obtained does not depend on the choice of reference frame.\\
Do note however that it will be in correct to apply "conservation of enegy" formula in the frame in which the hill is moving. The energy carried by the hill is not negligible in this frame. See also the next problem.\\
1.165 In a frame moving relative to the earth, one has to include the kinetic energy of the earth as well as earth's acceleration to be able to apply conservation of energy to the problem. In a reference frame falling to the earth with velocity \(v_{\mathrm{o}}\), the stone is initially going up with velocity \(v_{o}\) and so is the earth. The final velocity of the stone is \(0=v_{o}-\mathrm{gt}\) and that of the earth is \(v_{o}+\frac{m}{M} g t\) ( \(M\) is the mass of the earth), from Newton's third law, where \(t=\) time of fall. From conservation of energy

Hence

\[
\begin{aligned}
\frac{1}{2} m v_{0}^{2}+\frac{1}{2} M v_{0}^{2}+m g h & =\frac{1}{2} M\left(v_{0}+\frac{m}{M} v_{0}\right)^{2} \\
\frac{1}{2} v_{0}^{2}\left(m+\frac{m^{2}}{M}\right) & =m g h
\end{aligned}
\]

Negecting \(\frac{m}{M}\) in comparison with 1 , we get

\[
v_{0}^{2}=2 g h \text { or } v_{o}=\sqrt{2 g h}
\]

The point is this in earth's rest frame the effect of earth's accleration is of order \(\frac{m}{M}\) and can be neglected but in a frame moving with respect to the earth the effect of earth's acceleration must be kept because it is of order one (i.e. large).\\
1.166 From conservation of momentum, for the closed system "both colliding particles"

\[
m_{1} \overrightarrow{v_{1}}+m_{2} \overrightarrow{v_{2}}=\left(m_{1}+m_{2}\right) \vec{v}
\]

or,

\[
\vec{v}=\frac{m_{1} \overrightarrow{v_{1}}+m_{2} \overrightarrow{v_{2}}}{m_{1}+m_{2}}=\frac{1(3 \vec{i}-2 \vec{j})+2(4 \vec{j}-6 \vec{k})}{3}=\vec{i}+2 \vec{j}-4 \vec{k}
\]

Hence

\[
|\vec{v}|=\sqrt{1+4+16} \mathrm{~m} / \mathrm{s}=4.6 \mathrm{~m} / \mathrm{s}
\]

1.167 For perfectly inelastic collision, in the C.M. frame, final kinetic energy of the colliding system (both spheres) becomes zero. Hence initial kinetic energy of the system in C.M. frame completely turns into the internal energy ( \(Q\) ) of the formed body. Hence

\[
Q=\tilde{T}_{i}=\frac{1}{2} \mu\left|\overrightarrow{v_{1}}-\overrightarrow{v_{2}}\right|^{2}
\]

Now from energy conservation \(\Delta T=-Q=-\frac{1}{2} \mu\left|\overrightarrow{v_{1}}-\overrightarrow{v_{2}}\right|^{2}\),\\
In lab frame the same result is obtained as

\[
\begin{gathered}
\Delta T=\frac{1}{2} \frac{\left(m_{1} \overrightarrow{v_{1}}+m_{2} \overrightarrow{v_{2}}\right)^{2}}{m_{1}+m_{2}}-\frac{1}{2} m_{1}\left|\overrightarrow{v_{1}}\right|^{2}+m_{2}\left|\overrightarrow{v_{2}}\right|^{2} \\
=-\frac{1}{2} \mu\left|\overrightarrow{v_{1}}-\overrightarrow{v_{2}}\right|^{2}
\end{gathered}
\]

1.168 (a) Let the initial and final velocities of \(m_{1}\) and \(m_{2}\) are \(\vec{u}_{1}, \overrightarrow{u_{2}}\) and \(\vec{v}, \overrightarrow{v_{2}}\) respectively.

Then from conservation of momentum along horizontal and vertical directions, we get :


\begin{equation*}
m_{1} u_{1}=m_{2} v_{2} \cos \theta \tag{1}
\end{equation*}


and


\begin{equation*}
m_{1} v_{1}=m_{2} v_{2} \sin \theta \tag{2}
\end{equation*}


Squaring (1) and (2) and then adding them,

\[
m_{2}^{2} v_{2}^{2}=m_{1}^{2}\left(u_{1}^{2}+v_{1}^{2}\right)
\]

\begin{center}
\includegraphics[max width=\textwidth, alt={}]{9f3c5a8e-ac5d-42f4-8eb6-ea2af623f2d8-093_390_626_816_772}
\end{center}

Now, from kinetic energy conservation,


\begin{equation*}
\frac{1}{2} m_{1} u_{1}^{2}=\frac{1}{2} m_{2} v_{2}^{2}+\frac{1}{2} m_{1} v_{1}^{2} \tag{3}
\end{equation*}


or, \(m\left(u_{1}^{2}-v_{1}^{2}\right)=m_{2} v_{2}^{2}=m_{2} \frac{m_{1}^{2}}{m_{2}^{2}}\left(u_{1}^{2}+v_{1}^{2}\right)\) [Using (3)]\\
or,

\[
u_{1}^{2}\left(1-\frac{m_{1}}{m_{2}}\right)=v_{1}^{2}\left(1+\frac{m_{1}}{m_{2}}\right)
\]

or,


\begin{equation*}
\left(\frac{v_{1}}{u_{1}}\right)^{2}=\frac{m_{2}-m_{1}}{m_{1}+m_{2}} \tag{4}
\end{equation*}


So, fraction of kinetic energy lost by the particle 1 ,


\begin{gather*}
=\frac{\frac{1}{2} m_{1} u_{1}^{2}-\frac{1}{2} m_{1} v_{1}^{2}}{\frac{1}{2} m_{1} u_{1}^{2}}=1-\frac{v_{1}^{2}}{u_{1}^{2}} \\
=1-\frac{m_{2}-m_{1}}{m_{1}+m_{2}}=\frac{2 m_{1}}{m_{1}+m_{2}}[\text { Using (4) }] \tag{5}
\end{gather*}


(b) When the collision occurs head on,


\begin{equation*}
m_{1} u_{1}=m_{1} v_{1}+m_{2} v_{2} \tag{1}
\end{equation*}


and from conservation of kinetic energy,

\[
\begin{array}{r}
\frac{1}{2} m_{1} u_{1}^{2}=\frac{1}{2} m_{1} v_{1}^{2}+\frac{1}{2} m_{2} v_{2}^{2} \\
=\frac{1}{2} m_{1} v_{1}^{2}+\frac{1}{2} m_{2}\left[\frac{m_{1}\left(u_{1}-v_{1}\right)^{2}}{m_{2}}\right]^{2}[\mathrm{U}
\end{array}
\]

or,

\[
v_{1}\left(1+\frac{m_{1}}{m_{2}}\right)=u_{1}\left(\frac{m_{1}}{m_{2}}-1\right)
\]

or,


\begin{equation*}
\frac{v_{1}}{u_{1}}=\frac{\left(m_{1} / m_{2}-1\right)}{\left(1+m_{1} / m_{2}\right)} \tag{6}
\end{equation*}


Fraction of kinetic energy, lost

\[
=1-\frac{v_{1}^{2}}{u_{1}^{2}}=1-\left(\frac{m_{1}-m_{2}}{m_{1}+m_{2}}\right)^{2}=\frac{4 m_{1} m_{2}}{\left(m_{1}+m_{2}\right)^{2}}[\text { Using (6)] }
\]

1.169 (a) When the particles fly apart in opposite direction with equal velocities (say \(v\) ), then from conservatin of momentum,


\begin{equation*}
m_{1} u+0=\left(m_{2}-m_{1}\right) v \tag{1}
\end{equation*}


and from conservation of kinetic energy,\\
or,


\begin{align*}
\frac{1}{2} m_{1} u^{2} & =\frac{1}{2} m_{1} v^{2}+\frac{1}{2} m_{2} v^{2} \\
m_{1} u^{2} & =\left(m_{1}+m_{2}\right) v^{2} \tag{2}
\end{align*}


From Eq. (1) and (2),\\
\includegraphics[max width=\textwidth, alt={}, center]{9f3c5a8e-ac5d-42f4-8eb6-ea2af623f2d8-094_353_593_1040_640}

Hence \(\frac{m_{1}}{m_{2}}=\frac{1}{3}\) as \(m_{2}=0\)

\[
m_{1} u^{2}=\left(m_{1}+m_{2}\right) \frac{m_{1}^{2} u^{2}}{\left(m_{2}-m_{1}\right)^{2}}
\]

\[
m_{2}^{2}-3 m_{1} m_{2}=0
\]

(b) When they fly apart symmetrically relative to the initial motion direction with the angle of divergence \(\theta=60^{\circ}\),\\
From conservation of momentum, along horizontal and vertical direction,


\begin{equation*}
m_{1} u_{1}=m_{1} v_{1} \cos (\theta / 2)+m_{2} v_{2} \cos (\theta / 2) \tag{1}
\end{equation*}


and

\[
m_{1} v_{1} \sin (\theta / 2)=m_{2} v_{2} \sin (\theta / 2)
\]

or,


\begin{equation*}
m_{1} v_{1}=m_{2} v_{2} \tag{2}
\end{equation*}


Now, from conservation of kinetic energy,


\begin{equation*}
\frac{1}{2} m_{1} u_{1}^{2}+0=\frac{1}{2} m_{1} v_{1}^{2}+\frac{1}{2} m_{2} v_{2}^{2} \tag{3}
\end{equation*}


From (1) and (2),\\
\(m_{1} u_{1}=\cos (\theta / 2)\left(m_{1} v_{1}+\frac{m_{1} v_{1}}{m_{2}} m_{2}\right)=2 m_{1} v_{1} \cos (\theta / 2)\)

So,


\begin{equation*}
u_{1}=2 v_{1} \cos (\theta / 2) \tag{4}
\end{equation*}


From (2), (3), and (4)

\[
\begin{gathered}
4 m_{1} \cos ^{2}(\theta / 2) v_{1}^{2}=m_{1} v_{1}^{2}+\frac{m_{2} m_{1}^{2} v_{1}^{2}}{m_{2}^{2}} \\
\text { or, } 4 \cos ^{2}(\theta / 2)=1+\frac{m_{1}}{m_{2}}
\end{gathered}
\]

or, \(\quad \frac{m_{1}}{m_{2}}=4 \cos ^{2} \frac{\theta}{2}-1\)\\
and putting the value of \(\theta\), we get, \(\frac{m_{1}}{m_{2}}=2\)\\
1.170 If \(\left(v_{1 x}, v_{1 y}\right)\) are the instantaneous velocity components of the incident ball and ( \(v_{2 x}, v_{2 y}\) ) are the velocity components of the struck ball at the same moment, then since there are no external impulsive forces (i.e. other than the mutual interaction of the balls) We have

\[
\begin{gathered}
u \sin \alpha=v_{1 y}, v_{2 y}=0 \\
m u \cos \alpha=m v_{1 x}+m v_{2 x}
\end{gathered}
\]

The impulsive force of mutual interaction satisfies

\[
\frac{d}{d t}\left(v_{1 x}\right)=\frac{F}{m}=-\frac{d}{d t}\left(v_{2 x}\right)
\]

( \(F\) is along the \(x\) axis as the balls are smooth. Thus \(Y\) component of momentum is not transferred.) Since loss of K.E. is stored as deformation energy \(D\), we have

\[
\begin{aligned}
D & =\frac{1}{2} m u^{2}-\frac{1}{2} m v_{1}^{2}-\frac{1}{2} m v_{2}^{2} \\
& =\frac{1}{2} m u^{2} \cos ^{2} \alpha-\frac{1}{2} m v_{1 x}^{2}-\frac{1}{2} m v_{2 x}^{2} \\
& =\frac{1}{2 m}\left[m^{2} u^{2} \cos ^{2} \alpha-m^{2} v_{1 x}^{2}-\left(m u \cos \alpha-m v_{1 x}\right)^{2}\right] \\
& =\frac{1}{2 m}\left[2 m^{2} u \cos \alpha v_{1 x}-2 m^{2} v_{1 x}^{2}\right]=m\left(v_{1 x} u \cos \alpha-v_{1 x}^{2}\right) \\
& =m\left[\frac{u^{2} \cos ^{2} \alpha}{4}-\left(\frac{u \cos \alpha}{2}-v_{1 x}\right)^{2}\right]
\end{aligned}
\]

We see that \(D\) is maximum when

\[
\frac{u \cos \alpha}{2}=v_{1 x}
\]

and

\[
D_{\max }=\frac{m u^{2} \cos ^{2} \alpha}{4}
\]

Then \(\eta=\frac{D_{\text {max }}}{\frac{1}{2} m u^{2}}=\frac{1}{2} \cos ^{2} \alpha=\frac{1}{4}\)\\
On substiuting \(\alpha=45^{\circ}\)\\
\includegraphics[max width=\textwidth, alt={}, center]{9f3c5a8e-ac5d-42f4-8eb6-ea2af623f2d8-095_302_467_1729_824}\\
1.171 From the conservation of linear momentum of the shell just before and after its fragmentation


\begin{equation*}
3 \vec{v}=\overrightarrow{v_{1}}+\overrightarrow{v_{2}}+\overrightarrow{v_{3}} \tag{1}
\end{equation*}


where \(\overrightarrow{v_{1}}, \overrightarrow{v_{2}}\) and \(\overrightarrow{v_{3}}\) are the velocities of its fragments.\\
From the energy conservation \(3 \eta v^{2}=v_{1}^{2}+v_{2}^{2}+v_{3}^{2}\)

Now


\begin{equation*}
\tilde{\overrightarrow{v_{i}}} \text { or } \overrightarrow{v_{i C}}=\overrightarrow{v_{i}}-\overrightarrow{v_{C}}=\overrightarrow{v_{i}}-\vec{v} \tag{2}
\end{equation*}


where \(\vec{v}_{C}=\vec{v}=\) velocity of the C.M. of the fragments the velocity of the shell. Obviously in the C.M. frame the linear momentum of a system is equal to zero, so


\begin{equation*}
\tilde{\overrightarrow{v_{1}}}+\tilde{\overrightarrow{v_{2}}}+\tilde{\overrightarrow{v_{3}}}=0 \tag{4}
\end{equation*}


Using (3) and (4) in (2), we get


\begin{gather*}
3 \eta v^{2}=\left(\vec{v}+\tilde{\overrightarrow{v_{1}}}\right)^{2}+\left(\vec{v}+\tilde{\overrightarrow{v_{2}}}\right)^{2}+\left(\vec{v}-\tilde{\overrightarrow{v_{1}}}-\tilde{\overrightarrow{v_{2}}}\right)^{2}=3 v^{2}+2 \tilde{v}_{1}^{2}+2 \tilde{v}_{2}^{2}+2 \tilde{\overrightarrow{v_{1}}} \cdot \tilde{\overrightarrow{v_{2}}} \\
\text { or, } \quad 2 \tilde{v}_{1}^{2}+2 \tilde{v}_{1} \tilde{v}_{2} \cos \theta+2 \tilde{v}_{2}^{2}+3(1-\eta) v^{2}=0
\end{gather*}


If we have had used \(\tilde{\overrightarrow{v_{2}}}=-\tilde{\overrightarrow{v_{1}}}-\tilde{\overrightarrow{v_{3}}}\), then Eq. 5 were contain \(\tilde{v}_{3}\) instead of \(\tilde{v}_{2}\) and so on.\\
The problem being symmetrical we can look for the maximum of any one. Obviously it will be the same for each.\\
For \(\tilde{v}_{1}\) to be real in Eq. (5)

\[
4 \tilde{v}_{2}^{2} \cos ^{2} \theta \geq 8\left(2 \tilde{v}_{2}^{2}+3(1-\eta) v^{2}\right) \text { or } 6(\eta-1) v^{2} \geq\left(4-\cos ^{2} \theta\right) \tilde{v}_{2}^{2}
\]

So,

\[
\tilde{v}_{2} \leq v \sqrt{\frac{6(\eta-1)}{4-\cos ^{2} \theta}} \text { or } \tilde{v}_{2(\max )}=\sqrt{2(\eta-1)} v
\]

Hence

\[
v_{2(\max )}=\left|\vec{v}+\tilde{v}_{2}\right|_{\max }=v+\sqrt{2(\eta-1)} v=v(1+\sqrt{2(\eta-1)})=1 \mathrm{~km} / \mathrm{s}
\]

Thus owing to the symmetry

\[
v_{1(\max )}=v_{2(\max )}=v_{3(\max )}=v(1+\sqrt{2(\eta-1)})=1 \mathrm{~km} / \mathrm{s}
\]

1.172 Since, the collision is head on, the particle 1 will continue moving along the same line as before the collision, but there will be a change in the magnitude of it's velocity vector. Let it starts moving with velocity \(v_{1}\) and particle 2 with \(v_{2}\) after collision, then from the conservation of momentum


\begin{equation*}
m u=m v_{1}+m v_{2} \text { or, } u=v_{1}+v_{2} \tag{1}
\end{equation*}


And from the condition, given,\\
or,


\begin{gather*}
\eta=\frac{\frac{1}{2} m u^{2}-\left(\frac{1}{2} m v_{1}^{2}+\frac{1}{2} m v_{2}^{2}\right)}{\frac{1}{2} m u^{2}}=1-\frac{v_{1}^{2}+v_{2}^{2}}{u^{2}} \\
v_{1}^{2}+v_{2}^{2}=(1-\eta) u^{2} \tag{2}
\end{gather*}


From (1) and (2),\\
or,

\[
\begin{gathered}
v_{1}^{2}+\left(u-v_{1}\right)^{2}=(1-\eta) u^{2} \\
v_{1}^{2}+u^{2}-2 u v_{1}+v_{1}^{2}=(1-\eta) u^{2}
\end{gathered}
\]

or,

\[
2 v_{1}^{2}-2 v_{1} u+\eta u^{2}=0
\]

So,

\[
\begin{gathered}
v_{1}=2 u \pm \frac{\sqrt{4 u^{2}-8 \eta u^{2}}}{4} \\
=\frac{1}{2}\left[u \pm \sqrt{u^{2}-2 \eta u^{2}}\right]=\frac{1}{2} u(1 \pm \sqrt{1-2 \eta})
\end{gathered}
\]

Positive sign gives the velocity of the 2 nd particle which lies ahead. The negative sign is correct for \(v_{1}\).\\
So, \(v_{1}=\frac{1}{2} u(1-\sqrt{1-2 \eta})=5 \mathrm{~m} / \mathrm{s}\) will continue moving in the same direction.\\
Note that \(\nu_{1}=0\) if \(\eta=0\) as it must.\\
1.173 Since, no external impulsive force is effective on the system " \(M+m\) ", its total momentum along any direction will remain conserved.\\
So from \(p_{x}=\) const.


\begin{equation*}
m u=M v_{1} \cos \theta \text { or, } v_{1}=\frac{m}{M} \frac{u}{\cos \theta} \tag{1}
\end{equation*}


and from \(p_{y}=\) const\\
\(m v_{2}=M v_{1} \sin \theta\) or, \(v_{2}=\frac{M}{m} v_{1} \sin \theta=u \tan \theta\), [using (1)]\\
Final kinetic energy of the system

\[
T_{f}=\frac{1}{2} m v_{2}^{2}+\frac{1}{2} M v_{1}^{2}
\]

And initial kinetic energy of the system \(=\frac{1}{2} m u^{2}\)\\
So, \(\quad \%\) change \(=\frac{T_{f}-T_{i}}{T_{i}} \times 100\)

\[
\begin{aligned}
& =\frac{\frac{1}{2} m u^{2} \tan ^{2} \theta+\frac{1}{2} M \frac{m^{2}}{M^{2}} \frac{u^{2}}{\cos ^{2} \theta}-\frac{1}{2} m u^{2}}{\frac{1}{12} m u^{2}} \times 100 \\
& =\frac{\frac{1}{2} u^{2} \tan ^{2} \theta+\frac{1}{2} \frac{m}{M} u^{2} \sec ^{2} \theta-\frac{1}{2} u^{2}}{\frac{1}{2} u^{2}} \times 100 \\
& =\left(\tan ^{2} \theta+\frac{m}{M} \sec ^{2} \theta-1\right) \times 100
\end{aligned}
\]

\includegraphics[max width=\textwidth, alt={}, center]{9f3c5a8e-ac5d-42f4-8eb6-ea2af623f2d8-097_381_599_1404_776}\\
and putting the values of \(\theta\) and \(\frac{m}{M}\), we get \(\%\) of change in kinetic energy \(=-40 \%\) 1.174 (a) Let the particles \(m_{1}\) and \(m_{2}\) move with velocities \(\overrightarrow{v_{1}}\) and \(\overrightarrow{v_{2}}\) respectively. On the basis of solution of problem 1.147 (b)

\[
\tilde{p}=\mu v_{r e l}=\mu\left|\overrightarrow{v_{1}}-\overrightarrow{v_{2}}\right|
\]

As \(\quad \overrightarrow{v_{1}} \perp \overrightarrow{v_{2}}\)\\
So, \(\quad \tilde{p}=\mu \sqrt{v_{1}^{2}+v_{2}^{2}}\) where \(\mu=\frac{m_{1} m_{2}}{m_{1}+m_{2}}\)\\
(b) Again from 1.147 (b)

\[
\tilde{T}=\frac{1}{2} \mu v_{\text {rel }}^{2}=\frac{1}{2} \mu\left|\overrightarrow{v_{1}}-\overrightarrow{v_{2}}\right|^{2}
\]

So,

\[
\tilde{T}=\frac{1}{2} \mu\left(v_{1}^{2}+v_{2}^{2}\right)
\]

1.175 From conservation of momentum

\[
\overrightarrow{p_{1}}=\overrightarrow{p_{1}^{\prime}}+\overrightarrow{p_{2}^{\prime}}
\]

so

\[
\left(\overrightarrow{p_{1}}-\overrightarrow{p_{1}^{\prime}}\right)^{2}=p_{1}^{2}-2 p_{1} p_{1}^{\prime} \cos \theta_{1}+p_{1}^{\prime 2}=p_{2}^{\prime 2}
\]

From conservation of energy

\[
\frac{p_{1}^{2}}{2 m_{1}}=\frac{p_{1}^{\prime 2}}{2 m_{1}}+\frac{p_{2}^{\prime 2}}{2 m_{2}}
\]

Eliminating \(p_{2}{ }^{\prime}\) we get

\[
0=p_{1}^{\prime 2}\left(1+\frac{m_{2}}{m_{1}}\right)-2 p_{1}^{\prime} p_{1} \cos \theta_{1}+p_{1}^{2}\left(1-\frac{m_{2}}{m_{1}}\right)
\]

This quadratic equation for \(p_{1}^{\prime}\) has a real solution in terms of \(p_{1}\) and \(\cos \theta_{1}\) only if

\[
4 \cos ^{2} \theta_{1} \geq 4\left(1-\frac{m_{2}^{2}}{m_{1}^{2}}\right)
\]

or

\[
\sin ^{2} \theta_{1} \leq \frac{m_{2}^{2}}{m_{1}^{2}}
\]

\begin{center}
\includegraphics[max width=\textwidth, alt={}]{9f3c5a8e-ac5d-42f4-8eb6-ea2af623f2d8-098_365_454_1096_869}
\end{center}

This clearly implies (since only + sign makes sense) that

\[
\sin \theta_{1 \max }=\frac{m_{2}}{m_{1}} .
\]

1.176 From the symmetry of the problem, the velocity of the disc \(A\) will be directed either in the initial direction or opposite to it just after the impact. Let the velocity of the disc \(A\) after the collision be \(v^{\prime}\) and be directed towards right after the collision. It is also clear from the symmetry of problem that the discs \(B\) and \(C\) have equal speed (say \(v^{\prime \prime}\) ) in the directions, shown. From the condition of the problem,


\begin{equation*}
\cos \theta=\frac{\eta \frac{d}{2}}{d}=\frac{\eta}{2} \text { so, } \sin \theta=\sqrt{4-\eta^{2}} / 2 \tag{1}
\end{equation*}


For the three discs, system, from the conservation of linear momentum in the symmetry direction (towards right)


\begin{equation*}
m v=2 m v^{\prime \prime} \sin \theta+m v^{\prime} \text { or, } v=2 v^{\prime \prime} \sin \theta+v^{\prime} \tag{2}
\end{equation*}


From the definition of the coefficeint of restitution, we have for the discs \(A\) and \(B\) (or \(C\) )\\
\(e=\frac{\nu^{\prime \prime}-\nu^{\prime} \sin \theta}{\nu \sin \theta-0}\)\\
But \(e=1\), for perfectly elastic collision,

\[
\begin{aligned}
& \text { So, } v \sin \theta=v^{\prime \prime}-v^{\prime} \sin \theta \\
& \text { From (2) and (3), } \\
& \qquad \begin{aligned}
v^{\prime} & =\frac{v\left(1-2 \sin ^{2} \theta\right)}{\left(1+2 \sin ^{2} \theta\right)} \\
& =\frac{v\left(\eta^{2}-2\right)}{6-\eta^{2}} \quad\{\text { using }(1)\}
\end{aligned}
\end{aligned}
\]

Hence we have,\\
\includegraphics[max width=\textwidth, alt={}, center]{9f3c5a8e-ac5d-42f4-8eb6-ea2af623f2d8-099_434_736_362_653}

\[
v^{\prime}=\frac{v\left(\eta^{2}-2\right)}{6-\eta^{2}}
\]

Therefore, the disc \(A\) will recoil if \(\eta<\sqrt{2}\) and stop if \(\eta=\sqrt{2}\).\\
Note : One can write the equations of momentum conservation along the direction perpendicular to the initial direction of disc \(A\) and the consevation of kinetic energy instead of the equation of restitution.\\
1.177 (a) Let a molecule comes with velocity \(\overrightarrow{v_{1}}\) to strike another stationary molecule and just after collision their velocities become \(\vec{v}_{1}^{\prime}\) and \(\vec{v}_{2}^{\prime}\) respectively. As the mass of the each molecule is same, conservation of linear momentum and conservation of kinetic energy for the system (both molecules) respectively gives :\\
and

\[
\begin{aligned}
& \overrightarrow{v_{1}}=\vec{v}_{1}^{\prime}+\vec{v}_{2}^{\prime} \\
& v_{1}^{2}=v_{1}^{\prime 2}+v_{2}^{\prime 2}
\end{aligned}
\]

From the property of vector addition it is obvious from the obtained Eqs. that

\[
\vec{v}_{1}^{\prime} \perp \vec{v}_{2}^{\prime} \text { or } \vec{v}_{1}^{\prime} \cdot \vec{v}_{2}^{\prime}=0
\]

(b) Due to the loss of kinetic energy in inelastic collision \(v_{1}^{2}>v_{1}^{\prime 2}+v_{2}^{\prime 2}\)\\
so, \(\vec{v}_{1}^{\prime} \cdot \vec{v}_{2}^{\prime}>0\) and therefore angle of divergence \(<90^{\circ}\).\\
1.178 Suppose that at time \(t\), the rocket has the mass \(m\) and the velocity \(\vec{v}\), relative to the reference frame, employed. Now consider the inertial frame moving with the velocity that the rocket has at the given moment. In this reference frame, the momentum increament that the rocket \& ejected gas system acquires during time \(d t\) is,

\[
d \vec{p}=\dot{m} d \vec{v}+\mu d t \vec{u}=\vec{F} d t
\]

or,

\[
m \frac{d \vec{v}}{d t}=\vec{F}-\mu \vec{u}
\]

or,

\[
m \vec{w}=\vec{F}-\mu \vec{u}
\]

1.179 According to the question, \(\vec{F}=0\) and \(\mu=-d m / d t\) so the equation for this system becomes,

\[
m \frac{d \vec{v}}{d t}=\frac{d m}{d t} \vec{u}
\]

As \(d \vec{v} \uparrow \downarrow \vec{u} \quad\) so, \(m d v=-u d m\).\\
Integrating within the limits :

\[
\frac{1}{u} \int_{0}^{v} d v=-\int_{m_{0}}^{m} \frac{d m}{m} \text { or } \frac{v}{u}=\ln \frac{m_{0}}{m}
\]

Thus, \(v=u \ln \frac{m_{0}}{m}\)\\
As \(d \vec{v} \uparrow \downarrow \vec{u}\), so in vector form \(\vec{v}=-\vec{u} \ln \frac{m_{0}}{m}\)\\
1.180 According to the question, \(\vec{F}\) (external force) \(=0\)

So,

\[
m \frac{d \vec{v}}{d t}=\frac{d m}{d t} \vec{u}
\]

As

\[
d \vec{v} \uparrow \downarrow \vec{u},
\]

so, in scalar form,\\
or,

\[
m d v=-u d m
\]

\[
\frac{w d t}{u}=-\frac{d m}{m}
\]

Integrating within the limits for \(m(t)\)

\[
\frac{w t}{u}=-\int_{m_{0}}^{m} \frac{d m}{m} \text { or, } \frac{v}{u}=-\ln \frac{m}{m_{0}}
\]

Hence,

\[
m=m_{0} e^{-(w t / u)}
\]

1.181 As \(\vec{F}=0\), from the equation of dynamics of a body with variable mass;


\begin{equation*}
m \frac{d \vec{v}}{d t}=\vec{u} \frac{d m}{d t} \quad \text { or, } \quad d \vec{v}=\vec{u} \frac{d m}{m} \tag{1}
\end{equation*}


Now \(d \vec{v} \uparrow \downarrow \vec{u}\) and since \(\vec{u} \perp \vec{v}\), we must have \(|d \vec{v}|=v_{0} d \alpha\) (because \(v_{0}\) is constant) where \(d \alpha\) is the angle by which the spaceship turns in time \(d t\).

So,

\[
-u \frac{d m}{m}=v_{0} d \alpha \text { or, } d \alpha=-\frac{u}{v_{0}} \frac{d m}{m}
\]

or,

\[
\alpha=-\frac{u}{v_{0}} \int_{m_{0}}^{m} \frac{d m}{m}=\frac{u}{v_{0}} \ln \left(\frac{m_{0}}{m}\right)
\]

1.182 We have \(\frac{d m}{d t}=-\mu\) or, \(d m=-\mu d t\)

Integrating \(\quad \int_{m_{0}}^{m} d m=-\mu \int_{0}^{t} d t\) or, \(m=m_{0}-\mu t\)\\
As \(\vec{u}=0\) so, from the equation of variable mass system :

\[
\begin{gathered}
\left(m_{0}-\mu t\right) \frac{d \vec{v}}{d t}=\vec{F} \text { or, } \frac{d \vec{v}}{d t}=\vec{w}=\vec{F} /\left(m_{0}-\mu t\right) \\
\int_{0} d \vec{v}=\vec{F} \int_{0}^{t} \frac{d t}{\left(m_{0}-\mu t\right)} \\
\vec{v}=\frac{\vec{F}}{\mu} \ln \left(\frac{m_{0}}{m_{0}-\mu t}\right)
\end{gathered}
\]

or,\\
1.183 Let the car be moving in a reference frame to which the hopper is fixed and at any instant of time, let its mass be \(m\) and velocity \(\overrightarrow{v .}\)\\
Then from the general equation, for variable mass system.

\[
m \frac{d \vec{v}}{d t}=\vec{F}+\vec{u} \frac{d m}{d t}
\]

We write the equation, for our system as,


\begin{equation*}
m \frac{d \vec{v}}{d t}=\vec{F}-\vec{v} \frac{d m}{d t} \text { as, } \vec{u}=-\vec{v} \tag{1}
\end{equation*}


So

\[
\frac{d}{d t}(\vec{m} v)=\vec{F}
\]

and

\[
\vec{v}=\frac{\overrightarrow{F t}}{m} \text { on integration. }
\]

But

\[
\begin{gathered}
m=m_{0}+\mu t \\
\vec{v}=\frac{\vec{F}_{t}}{m_{0}\left(1+\frac{\mu t}{m_{0}}\right)}
\end{gathered}
\]

Thus the sought acceleration, \(\vec{w}=\frac{d \vec{v}}{d t}=\frac{\vec{F}}{m_{0}\left(1+\frac{\mu t}{m_{0}}\right)^{2}}\)\\
1.184 Let the length of the chain inside the smooth horizontal tube at an arbitrary instant is \(x\). From the equation,

\[
m \vec{w}=\vec{F}+\vec{u} \frac{d m}{d t}
\]

as \(\vec{u}=0, \vec{F} \uparrow \uparrow \vec{w}\), for the chain inside the tube


\begin{equation*}
\lambda x w=T \text { where } \lambda=\frac{m}{l} \tag{1}
\end{equation*}


Similarly for the overhanging part,

\[
\vec{u}=0
\]

Thus \(\quad m w=F\)\\
or \(\quad \lambda h w=\lambda h g-T\)\\
From (1) and (2),\\
\(\lambda(x+h) w=\lambda h g\) or, \((x+h) v \frac{d v}{d s}=h g\)\\
\includegraphics[max width=\textwidth, alt={}, center]{9f3c5a8e-ac5d-42f4-8eb6-ea2af623f2d8-102_383_584_269_715}\\
or,

\[
(x+h) v \frac{d v}{(-d x)}=g h,
\]

[As the length of the chain inside the tube decreases with time, \(d s=-d x\).]\\
or,

\[
v d v=-g h \frac{d x}{x+h}
\]

Integrating,

\[
\int_{0}^{v} v d v=-g h \int_{(l-h)}^{0} \frac{d x}{x+h}
\]

or, \(\quad \frac{v^{2}}{2}=g h \ln \left(\frac{l}{h}\right)\) or \(v=\sqrt{2 g h \ln \left(\frac{l}{h}\right)}\)\\
1.185 Force moment relative to point \(O\);

\[
\vec{N}=\frac{d \vec{M}}{d t}=2 \overrightarrow{b t}
\]

Let the angle between \(\vec{M}\) and \(\vec{N}\),\\
\(\alpha=45^{\circ}\) at \(t=t_{0}\),\\
Then \(\frac{1}{\sqrt{2}}=\frac{\vec{M} \cdot \vec{N}}{|\vec{M}||\vec{N}|}=\frac{\left(\vec{a}+\vec{b} t_{0}^{2}\right) \cdot\left(2 \overrightarrow{b t_{0}}\right)}{\sqrt{a^{2}+b^{2} t_{0}^{4}} 2 b t_{0}}\)

\[
=\frac{2 b^{2} t_{0}^{3}}{\sqrt{a^{2}+b^{2} t_{0}^{4}} 2 b t_{0}}=\frac{b t_{0}^{2}}{\sqrt{a^{2}+b^{2} t_{0}^{4}}}
\]

\begin{center}
\includegraphics[max width=\textwidth, alt={}]{9f3c5a8e-ac5d-42f4-8eb6-ea2af623f2d8-102_431_451_1323_838}
\end{center}

So, \(2 b^{2} t_{0}^{4}=a^{2}+b^{2} t_{0}^{4}\) or, \(t_{0}=\sqrt{\frac{a}{b}}\) (as \(t_{0}\) cannot be negative)\\
It is also obvious from the figure that the angle \(\alpha\) is equal to \(45^{\circ}\) at the moment \(t_{0}\), when \(a=b t_{0}^{2}\), i.e. \(t_{0}=\sqrt{a / b}\) and \(\vec{N}=2 \sqrt{\frac{a}{b}} \vec{b}\).\\
1.186

\[
\begin{aligned}
\vec{M}(t) & =\vec{r} \times \vec{p}=\left(\vec{v}_{0} t+\frac{1}{2} \overrightarrow{g t}^{2}\right) \times m\left(\overrightarrow{v_{0}}+\overrightarrow{g t}\right) \\
& =m v_{0} g t^{2} \sin \left(\frac{\pi}{2}+\alpha\right)(-\vec{k})+\frac{1}{2} m v_{0} g t^{2} \sin \left(\frac{\pi}{2}+\alpha\right)(\vec{k}) \\
& =\frac{1}{2} m v_{0} g t^{2} \cos \alpha(-\vec{k}):
\end{aligned}
\]

Thus \(M(t)=\frac{m v_{0} g t^{2} \cos \alpha}{2}\)\\
Thus angular momentum at maximum height i.e. at \(t=\frac{\tau}{2}=\frac{v_{0} \sin \alpha}{g}\),\\
\(M\left(\frac{\tau}{2}\right)=\left(\frac{m v_{0}^{3}}{2 g}\right) \sin ^{2} \alpha \cos \alpha=37 \mathrm{~kg}-\mathrm{m}^{2} / \mathrm{s}\) Alternate :\\
\includegraphics[max width=\textwidth, alt={}, center]{9f3c5a8e-ac5d-42f4-8eb6-ea2af623f2d8-103_301_527_476_810}

\[
\begin{aligned}
& \vec{M}(0)=0 \text { so, } \vec{M}(t)=\int_{0}^{t} \vec{N} d t=\int_{0}^{t}(\vec{r} \times m \vec{g}) \\
& =\int_{0}^{t}\left[\left(\overrightarrow{v_{0}} t+\frac{1}{2} \vec{g}^{2}\right) \times m \vec{g}\right] d t=\left(\overrightarrow{v_{0}} \times m \vec{g}\right) \frac{t^{2}}{2}
\end{aligned}
\]

1.187 (a) The disc experiences gravity, the force of reaction of the horizontal surface, and the force \(\vec{R}\) of reaction of the wall at the moment of the impact against it. The first two forces counter-balance each other, leaving only the force \(\vec{R}\). It's moment relative to any point of the line along which the vector \(\vec{R}\) acts or along normal to the wall is equal to zero and therefore the angular momentum of the disc relative to any of these points does not change in the given process.\\
(b) During the course of collision with wall the position of disc is same and is equal to\\
\includegraphics[max width=\textwidth, alt={}, center]{9f3c5a8e-ac5d-42f4-8eb6-ea2af623f2d8-103_447_475_1259_834}\\
\(\overrightarrow{r_{o o}}\) Obviously the increment in linear momentum of the ball \(\Delta \vec{p}=2 m v \cos \alpha \hat{n}\)\\
Here, \(\Delta \vec{M}=\vec{r}_{o o^{\prime}} \times \Delta \vec{p}=2 m v l \cos \alpha \hat{n}\) and directed normally emerging from the plane of figure\\
Thus \(|\Delta \vec{M}|=2 m v l \cos \alpha\)\\
1.188 (a) The ball is under the influence of forces \(\vec{T}\) and \(m \vec{g}\) at all the moments of time, while moving along a horizontal circle. Obviously the vertical component of \(\vec{T}\) balance \(m \vec{g}\) and\\
so the net moment of these two about any point becoems zero. The horizontal component of \(\vec{T}\), which provides the centripetal acceleration to ball is already directed toward the centre ( \(C\) ) of the horizontal circle, thus its moment about the point \(C\) equals zero at all the moments of time. Hence the net moment of the force acting on the ball about point \(C\) equals zero and that's why the angular mommetum of the ball is conserved about the horizontal circle.\\
(b) Let \(\alpha\) be the angle which the thread forms with the vertical.\\
Now from equation of particle dynamics :\\
\(T \cos \alpha=m g\) and \(T \sin \alpha=m \omega^{2} l \sin \alpha\)\\
Hence on solving \(\cos \alpha=\frac{g}{\omega^{2} l}\)

As \(|\vec{M}|\) is constant in magnitude so from figure.\\
\(|\Delta \vec{M}|=2 M \cos \alpha\) where

\[
\begin{aligned}
& M=\left|\vec{M}_{i}\right|=\left|\vec{M}_{f}\right| \\
& =\left|\vec{r}_{b o} \times m \vec{v}\right|=m v l\left(\text { as } \vec{r}_{b o} \perp \vec{v}\right)
\end{aligned}
\]

Thus \(|\Delta \vec{M}|=2 m v l \cos \alpha=2 m \omega l^{2} \sin \alpha \cos \alpha\)\\
\includegraphics[max width=\textwidth, alt={}, center]{9f3c5a8e-ac5d-42f4-8eb6-ea2af623f2d8-104_570_405_520_898}

\[
=\frac{2 m g l}{\omega} \sqrt{1-\left(\frac{g}{\omega^{2} t}\right)^{2}}(\text { using } 1)
\]

1.189 During the free fall time \(t=\tau=\sqrt{\frac{2 h}{g}}\), the reference point \(O\) moves in hoizontal direction (say towards right) by the distance \(V \tau\). In the translating frame as \(\vec{M}(O)=0\), so

\[
\begin{aligned}
\Delta \vec{M} & =\vec{M}_{f}=\vec{r} \\
& =(-V \tau \vec{i}+h \vec{j}) \times m[g \tau \vec{j}-V \vec{i}] \quad \vec{i} \vec{j}(\vec{g}) \\
& =-m V g \tau^{2} \vec{h}+m V h(+\vec{k}) \\
& =-m V g\left(\frac{2 h}{g}\right) \vec{k}+m V h(+\vec{k})=-m V h \vec{k}
\end{aligned} \quad{ }^{\rightarrow}
\]

Hence

\[
|\Delta \vec{M}|=m V h
\]

1.190 The Coriolis force is. \(\left(2 m \overrightarrow{v^{\prime}} \times \vec{\omega}\right)\).

Here \(\vec{\omega}\) is along the \(z\)-axis (vertical). The moving disc is moving with velocity \(v_{0}\) which is constant. The motion is along the \(x\)-axis say. Then the Coriolis force is along \(y\)-axis and has the magnitude \(2 m v_{0} \omega\). At time \(t\), the distance of the centre of moving disc from \(O\) is \(v_{0} t\) (along \(x\)-axis). Thus the torque \(N\) due to the coriolis force is\\
\(N=2 m v_{0} \omega \cdot v_{0} t\) along the \(z\)-axis.

Hence equating this to \(\frac{d M}{d t}\)

\[
\frac{d M}{d t}=2 m v_{0}^{2} \omega t \text { or } M=m v_{0}^{2} \omega t^{2}+\text { constant } .
\]

The constant is irrelevant and may be put equal to zero if the disc is originally set in motion from the point \(O\).\\
This discussion is approximate. The Coriolis force will cause the disc to swerve from straight line motion and thus cause deviation from the above formula which will be substantial for large \(t\).\\
1.191 If \(\dot{r}=\) radial velocity of the particle then the total energy of the particle at any instant is


\begin{equation*}
\frac{1}{2} m \dot{r}^{2}+\frac{M^{2}}{2 m r^{2}}+k r^{2}=E \tag{1}
\end{equation*}


where the second term is the kinetic energy of angular motion about the centre \(O\). Then the extreme values of \(r\) are determined by \(r=0\) and solving the resulting quadratic equation

\[
k\left(r^{2}\right)^{2}-E r^{2}+\frac{M^{2}}{2 m}=0
\]

we get

\[
r^{2}=\frac{E \pm \sqrt{E^{2}-\frac{2 M^{2} k}{m}}}{2 k}
\]

From this we see that


\begin{equation*}
E=k\left(r_{1}^{2}+r_{2}^{2}\right) \tag{2}
\end{equation*}


where \(r_{1}\) is the minimum distance from O and \(r_{2}\) is the maximum distance. Then

\[
\begin{gathered}
\frac{1}{2} m v_{2}^{2}+2 k r_{2}^{2}=k\left(r_{1}^{2}+r_{2}^{2}\right) \\
m=\frac{2 k r^{2}}{v_{2}^{2}}
\end{gathered}
\]

-Note : Eq. (1) can be derived from the standard expression for kinetic energy and angular momentum in plane poler coordinates :

\[
T=\frac{1}{2} m \dot{r}^{2}+\frac{1}{2} m r^{2} \dot{\theta}^{2}
\]

\(M=\) angular momentum \(=m r^{2} \dot{\theta}\)\\
1.192 The swinging sphere experiences two forces : The gravitational force and the tension of the thread. Now, it is clear from the condition, given in the problem, that the moment of these forces about the vertical axis, passing through the point of suspension \(N_{z}=0\). Consequently, the angular momentum \(M_{z}\) of the sphere relative to the given axis ( \(z\) ) is constant.\\
Thus


\begin{equation*}
m v_{0}(l \sin \theta)=m v l \tag{1}
\end{equation*}


where \(m\) is the mass of the sphere and \(v\) is it \(s\) velocity in the position, when the thread forms an angle \(\frac{\pi}{2}\) with the vertical. Mechanical energy is also conserved, as the sphere is\\
under the influence if only one other force, i.e. tension, which does not perform any work, as it is always perpendicular to the velocity.

So,


\begin{equation*}
\frac{1}{2} m v_{0}^{2}+m g l \cos \theta=\frac{1}{2} m v^{2} \tag{2}
\end{equation*}


From (1) and (2), we get,

\[
v_{0}=\sqrt{2 g l / \cos \theta}
\]

1.193 Forces, acting on the mass \(m\) are shown in the figure. As \(\vec{N}=m \vec{g}\), the net torque of these two forces about any fixed point must be equal to zero. Tension \(T\), acting on the mass \(m\) is a central force, which is always directed towards the centre \(\boldsymbol{O}\). Hence the moment of force \(T\) is also zero about the point \(O\) and therefore the angular momentum of the particle \(m\) is conserved about \(O\).\\
Let, the angular velocity of the particle be \(\omega\), when the separation between hole and particle \(m\) is \(r\), then from the conservation of momentum about the point \(O\) : :

\[
m\left(\omega_{0} r_{0}\right) r_{0}=m(\omega r) r,
\]

or

\[
\omega=\frac{\omega_{0} r_{0}^{2}}{r^{2}}
\]

Now, from the second law of motion for \(m\),

\[
T=F=m \omega^{2} r
\]

Hence the sought tension;

\[
F=\frac{m \omega_{0}^{2} r_{0}^{4} r}{r^{4}}=\frac{m \omega_{0}^{2} r_{0}^{4}}{r^{3}}
\]

\includegraphics[max width=\textwidth, alt={}, center]{9f3c5a8e-ac5d-42f4-8eb6-ea2af623f2d8-106_431_642_680_758}\\
1.194 On the given system the weight of the body \(m\) is the only force whose moment is effective about the axis of pulley. Let us take the sense of \(\vec{\omega}\) of the pulley at an arbitrary instant as the positive sense of axis of rotation (z-axis)

As

\[
M_{z}(0)=0, \text { so, } \Delta M_{z}=M_{z}(t)=\int N_{z} d t
\]

So,

\[
M_{z}(t)=\int_{0}^{t} m g R d t=m g R t
\]

1.195 Let the point of contact of sphere at initial moment ( \(t=0\) ) be at \(O\). At an arbitrary moment, the forces acting on the sphere are shown in the figure. We have normal reaction \(N_{r}=m g \sin \alpha\) and both pass through same line and the force of static friction passes through the point \(O\), thus the moment about point \(O\) becomes zero. Hence \(m g \sin \alpha\) is the only force which has effective torque about point \(O\), and is given by \(|\vec{N}|=m g R \sin \alpha\) normally emerging from the plane of figure.\\
As \(\vec{M}(t=0)=0\), so, \(\Delta \vec{M}=\vec{M}(t)=\int \vec{N} d t\)\\
\includegraphics[max width=\textwidth, alt={}, center]{9f3c5a8e-ac5d-42f4-8eb6-ea2af623f2d8-106_392_531_1601_862}

Hence, \(\quad M(t)=N t=m g R \sin \alpha t\)\\
1.196 Let position vectors of the particles of the system be \(\vec{r}_{i}\) and \(\vec{r}_{i}^{\prime}\) with respect to the points \(O\) and \(O^{\prime}\) respectively. Then we have,


\begin{equation*}
\overrightarrow{r_{i}}=\overrightarrow{r_{i}}+\overrightarrow{r_{0}} \tag{1}
\end{equation*}


where \(\overrightarrow{r_{O}}\) is the radius vector of \(O^{\prime}\) with respect to \(O\).\\
Now, the angular momentum of the system relative to the point \(O\) can be written as follows;


\begin{equation*}
\vec{M}=\sum\left(\overrightarrow{r_{i}} \times \overrightarrow{p_{i}}\right)=\sum\left(\vec{r}_{i}^{\prime} \times \overrightarrow{p_{i}}\right)+\sum\left(\overrightarrow{r_{0}} \times \overrightarrow{p_{i}}\right) \tag{1}
\end{equation*}


or,


\begin{equation*}
\vec{M}=\vec{M}^{\prime}+\left(\overrightarrow{r_{0}} \times \vec{p}\right), \text { where, } \vec{p}=\sum \overrightarrow{p_{i}} \tag{2}
\end{equation*}


From (2), if the total linear momentum of the system, \(\vec{p}=0\), then its angular momentum does not depend on the choice of the point \(O\).\\
Note that in the C.M. frame, the system of particles, as a whole is at rest.\\
1.197 On the basis of solution of problem 1.196, we have concluded that; "in the C.M. frame, the angular momentum of system of particles is independent of the choice of the point, relative to which it is determined" and in accordance with the problem, this is denoted by \(\vec{M}\).

We denote the angular momentum of the system of particles, relative to the point \(O\), by \(\vec{M}\). Since the internal and proper angular momentum \(\widetilde{M}\), in the C.M. frame, does not depend on the choice of the point \(O^{\prime}\), this point may be taken coincident with the point \(O\) of the \(K\)-frame, at a given moment of time. Then at that moment, the radius vectors of all the particles, in both reference frames, are equal ( \(\vec{r}_{i}^{\prime}=\overrightarrow{r_{i}}\) ) and the velocities are related by the equation,


\begin{equation*}
\overrightarrow{v_{i}}=\overrightarrow{v_{i}}+\overrightarrow{v_{c}}, \tag{1}
\end{equation*}


where \(\vec{v}_{c}\) is the velocity of C.M. frame, relative to the \(K\)-frame. Consequently, we may write,

\[
\vec{M}=\sum m_{i}\left(\overrightarrow{r_{i}} \times \overrightarrow{v_{i}}\right)=\sum m_{i}\left(\overrightarrow{r_{i}^{\prime}} \times \overrightarrow{v_{i}}\right)+\sum m_{i}\left(\overrightarrow{r_{i}} \times \overrightarrow{v_{c}}\right)
\]

or, \(\quad \vec{M}=\vec{M}+m\left(\overrightarrow{r_{c}} \times \overrightarrow{v_{c}}\right)\), as \(\sum m_{i} \overrightarrow{r_{i}}=m \overrightarrow{r_{c}}\), where \(m=\sum m_{i}\).\\
or,

\[
\vec{M}=\tilde{M}+\left(\overrightarrow{r_{c}} \times m \overrightarrow{v_{c}}\right)=\tilde{M}+\left(\overrightarrow{r_{c}} \times \vec{p}\right)
\]

1.198 From conservation of linear momentum along the direction of incident ball for the system consists with colliding ball and phhere


\begin{equation*}
m v_{0}=m v^{\prime}+\frac{m}{2} v_{1} \tag{1}
\end{equation*}


where \(v^{\prime}\) and \(v_{1}\) are the velocities of ball and sphere 1 respectively after collision. (Remember that the collision is head on).\\
As the collision is perfectly elastic, from the definition of co-efficeint of restitution,


\begin{equation*}
1=\frac{v^{\prime}-v_{1}}{0-v_{0}} \text { or, } v^{\prime}-v_{1}=-v_{0} \tag{2}
\end{equation*}


Solving (1) and (2), we get,\\
\(v_{1}=\frac{4 v_{0}}{3}\), directed towards right.\\
In the C.M. frame of spheres 1 and 2 (Fig.)\\
\(\tilde{\overrightarrow{p_{1}}}=-\tilde{\overrightarrow{p_{2}}}\) and \(\left|\tilde{\overrightarrow{p_{1}}}\right|=\left|\tilde{\overrightarrow{p_{2}}}\right|=\mu\left|\overrightarrow{v_{1}}-\overrightarrow{v_{2}}\right|\)\\
Also, \(\overrightarrow{r_{1 C}}=-\overrightarrow{r_{2 C}}\), thus \(\tilde{M}=2\left[\overrightarrow{r_{1 C}} \times \tilde{\vec{p}_{1}}\right]\)\\
As \(\overrightarrow{r_{1 C}} \perp \overrightarrow{p_{1}}, \quad\) so, \(\tilde{M}=2\left[\frac{l}{2} \frac{m / 2}{2} \frac{4 v_{0}}{3} \hat{n}\right]\)\\
(where \(\hat{n}\) is the unit vector in the sense of \(\overrightarrow{r_{1 C}} \times \overrightarrow{p_{1}}\) )\\
\includegraphics[max width=\textwidth, alt={}, center]{9f3c5a8e-ac5d-42f4-8eb6-ea2af623f2d8-108_499_362_144_885}

Hence \(\tilde{M}=\frac{m v_{0} l}{3}\)\\
1.199 In the C.M. frame of the system (both the discs + spring), the linear momentum of the discs are related by the relation, \(\tilde{p_{1}}=-\tilde{p_{2}}\), at all the moments of time.\\
where,

\[
\tilde{p}_{1}=\tilde{p}_{2}=\tilde{p}=\mu v_{r e l}
\]

And the total kinetic energy of the system,

\[
T=\frac{1}{2} \mu v_{\text {rel }}^{2} \text { [See solution of } 1.147 \text { (b)] }
\]

Bearing in mind that at the moment of maximum deformation of the spring, the projection of \(\vec{v}_{\text {rel }}\) along the length of the spring becomes zero, i.e. \(v_{\text {rel }(x)}=0\).\\
The conservation of mechanical energy of the considered system in the C.M. frame gives.


\begin{equation*}
\frac{1}{2}\left(\frac{m}{2}\right) v_{0}^{2}=\frac{1}{2} \kappa x^{2}+\frac{1}{2}\left(\frac{m}{2}\right) v_{r d(y)}^{2} . \tag{1}
\end{equation*}


Now from the conservation of angular momentum of the system about the C.M.,\\
or,


\begin{gather*}
\frac{1}{2}\left(\frac{l_{0}}{2}\right)\left(\frac{m}{2} v_{0}\right)=2\left(\frac{l_{0}+x}{2}\right) \frac{m}{2} v_{r e l(y)} \\
v_{r e l(y)}=\frac{v_{0} l_{0}}{\left(l_{0}+x\right)}=v_{0}\left(1+\frac{x}{l_{0}}\right)^{-1} \sim v_{0}\left(1-\frac{x}{l_{0}}\right), \text { as } x \ll l_{0} \tag{2}
\end{gather*}


Using (2) in (1),

\[
\frac{1}{2} m v_{0}^{2}\left[1-\left(1-\frac{x}{l_{0}}\right)^{2}\right]=\kappa x^{2}
\]

or,

\[
\frac{1}{2} m v_{0}^{2}\left[1-\left(1-\frac{2 x}{l_{0}}+\frac{x^{2}}{l_{0}^{2}}\right)^{2}\right]=\kappa x^{2}
\]

or,

\[
\frac{m v_{0}^{2} x}{l_{0}} \approx \kappa x^{2},\left[\text { neglecting } x^{2} / l_{0}^{2}\right]
\]

As

\[
x \neq 0, \text { thus } x=\frac{m v_{0}^{2}}{\kappa l_{0}}
\]

\subsection*{1.4 UNIVERSAL GRAVITATION}
1.200 We have

Thus

\[
\begin{aligned}
\frac{M v^{2}}{r} & =\frac{\gamma M m_{s}}{r^{2}} \text { or } r=\frac{\gamma m_{s}}{v^{2}} \\
\omega & =\frac{v}{r}=\frac{v}{\gamma m_{s} / v^{2}}=\frac{v^{3}}{\gamma m_{s}}
\end{aligned}
\]

(Here \(m_{s}\) is the mass of the Sun.)\\
So \(\quad T=\frac{2 \pi \gamma m_{s}}{v^{3}}=\frac{2 \pi \times 6.67 \times 10^{-11} \times 1.97 \times 10^{30}}{\left(34.9 \times 10^{3}\right)^{3}}=1.94 \times 10^{7} \mathrm{sec}=225\) days.\\
(The answer is incorrectly written in terms of the planetary mass \(M\) )\\
1.201 For any planet

\[
M R \omega^{2}=\frac{\gamma M m_{s}}{R^{2}} \text { or } \omega=\sqrt{\frac{\gamma m_{s}}{R^{3}}}
\]

So,

\[
T=\frac{2 \pi}{\omega}=2 \pi R^{3 / 2} / \sqrt{\gamma m_{s}}
\]

(a) Thus

\[
\frac{T_{J}}{T_{E}}=\left(\frac{R_{J}}{R_{E}}\right)^{3 / 2}
\]

So

\[
\frac{R_{J}}{R_{E}}=\left(T_{J} / T_{E}\right)^{2 / 3}=(12)^{2 / 3}=5 \cdot 24 .
\]

(b)

\[
V_{J}^{2}=\frac{\gamma m_{s}}{R_{J}}, \text { and } R_{J}=\left(T \frac{\sqrt{\gamma m_{s}}}{2 \pi}\right)^{2 / 3}
\]

So

\[
V_{J}^{2}=\frac{\left(\gamma m_{s}\right)^{2 / 3}(2 \pi)^{2 / 3}}{T^{2 / 3}} \text { or, } V_{J}=\left(\frac{2 \pi \gamma m_{s}}{T}\right)^{2 / 3}
\]

where \(T=12\) years. \(m_{s}=\) mass of ths Sun. Putting the values we get \(V_{J}=12.97 \mathrm{~km} / \mathrm{s}\)

\[
\begin{aligned}
\text { Acceleration } & =\frac{v_{J}^{2}}{R_{J}}=\left(\frac{2 \pi \gamma m_{s}}{T}\right)^{2 / 3} \times\left(\frac{2 \pi}{T \sqrt{\gamma m_{s}}}\right)^{2 / 3} \\
& =\left(\frac{2 \pi}{T}\right)^{4 / 3}\left(\gamma m_{s}\right)^{1 / 3} \\
& =2.15 \times 10^{-4} \mathrm{~km} / \mathrm{s}^{2}
\end{aligned}
\]

1.202 Semi-major axis \(=(r+R) / 2\)

It is sufficient to consider the motion be along a circle of semi-major axis \(\frac{r+R}{2}\) for \(T\) does not depend on eccentricity.

Hence

\[
T=\frac{2 \pi\left(\frac{r+R}{2}\right)^{3 / 2}}{\sqrt{\gamma m_{s}}}=\pi \sqrt{(r+R)^{3} / 2 \gamma m_{s}}
\]

(again \(m_{s}\) is the mass of the Sun)\\
1.203 We can think of the body as moving in a very elongated orbit of maximum distance \(R\) and minimum distance 0 so semi major axis \(=R / 2\). Hence if \(\tau\) is the time of fall then\\
or

\[
\begin{aligned}
& \left(\frac{2 \tau}{T}\right)^{2}=\left(\frac{R / 2}{R}\right)^{3} \text { or } \tau^{2}=T^{2} / 32 \\
& \tau=T / 4 \sqrt{2}=365 / 4 \sqrt{2}=64.5 \text { days. }
\end{aligned}
\]

\(1.204 T=2 \pi R^{3 / 2} / \sqrt{\gamma m_{s}}\)\\
If the distances are scaled down, \(R^{3 / 2}\) decreases by a factor \(\eta^{3 / 2}\) and so does \(m_{s}\). Hence \(T\) does not change.\\
1.205 The double star can be replaced by a single star of mass \(\frac{m_{1} m_{2}}{m_{1}+m_{2}}\) moving about the centre of mass subjected to the force \(\gamma m_{1} m_{2} / r^{2}\). Then

So

\[
T=\frac{2 \pi r^{3 / 2}}{\sqrt{\gamma m_{1} m_{2} / \frac{m_{1} m_{2}}{m_{1}+m_{2}}}}=\frac{2 \pi r^{3 / 2}}{\sqrt{\gamma M}}
\]

or,

\[
r=\left(\frac{T}{2 \pi}\right)^{2 / 3}(\gamma M)^{1 / 3}=\sqrt[3]{\gamma M(T / 2 \pi)^{2}}
\]

1.206 (a) The gravitational potential due to \(m_{1}\) at the point of location of \(m_{2}\) :

\[
V_{2}=\int_{r}^{\infty} \vec{G} \cdot d \vec{r}=\int_{r}^{\infty}-\frac{\gamma m_{1}}{x^{2}} d x=-\frac{\gamma m_{1}}{r}
\]

So,

\[
U_{21}=m_{2} V_{2}=-\frac{\gamma m_{1} m_{2}}{r}
\]

Similarly

\[
U_{12}=-\frac{\gamma m_{1} m_{2}}{r}
\]

Hence

\[
U_{12}=U_{21}=U=-\frac{\gamma m_{1} m_{2}}{r}
\]

\includegraphics[max width=\textwidth, alt={}, center]{9f3c5a8e-ac5d-42f4-8eb6-ea2af623f2d8-111_168_697_136_698}\\
(b) Choose the location of the point mass as the origin. Then the potential energy \(d U\) of an element of mass \(d M=\frac{M}{l} d x\) of the rod in the field of the point mass is

\[
d U=-\gamma m \frac{M}{l} d x \frac{1}{x}
\]

where \(x\) is the distance between the element and the point. (Note that the rod and the point mass are on a straight line.) If then \(a\) is the distance of the nearer end of the rod from the point mass.\\
\includegraphics[max width=\textwidth, alt={}, center]{9f3c5a8e-ac5d-42f4-8eb6-ea2af623f2d8-111_170_942_724_340}

\[
U=-\gamma \frac{m M}{l} \int_{a}^{a+l} \frac{d x}{x}=-\gamma m \frac{M}{l} \ln \left(1+\frac{l}{a}\right)
\]

The force of interaction is

\[
\begin{aligned}
F & =-\frac{\partial U}{\partial a} \\
& =\gamma \frac{m M}{l} \times \frac{1}{1+\frac{l}{a}}\left(-\frac{l}{a^{2}}\right)=-\frac{\gamma m M}{a(a+l)}
\end{aligned}
\]

Minus sign means attraction.\\
1.207 As the planet is under central force (gravitational interaction), its angular momentum is conserved about the Sun (which is situated at one of the focii of the ellipse)\\
So, \(\quad m v_{1} r_{1}=m v_{2} r_{2} \quad\) or, \(\quad v_{1}^{2}=\frac{v_{2}^{2} r_{2}^{2}}{r_{1}^{2}}\)

From the conservation of mechanical energy of the system (Sun + planet),

\[
-\frac{\gamma m_{s} m}{r_{1}}+\frac{1}{2} m v_{1}^{2}=-\frac{\gamma m_{s} m}{r_{2}}+\frac{1}{2} m v_{2}^{2}
\]

or,

\[
-\frac{\gamma m_{s}}{r_{1}}+\frac{1}{2} v_{2}^{2} \frac{r_{2}^{2}}{r_{1}^{2}}=-\left(\frac{\gamma m_{s}}{r_{2}}\right)+\frac{1}{2} v_{2}^{2}[\text { Using (1) }]
\]

Thus,


\begin{equation*}
v_{2}=\sqrt{2 \gamma m_{s} r_{1} / r_{2}\left(r_{1}+r_{2}\right)} \tag{2}
\end{equation*}


Hence

\[
M=m v_{2} r_{2}=m \sqrt{2 \gamma m_{s} r_{1} r_{2} /\left(r_{1}+r_{2}\right)}
\]

1.208 From the previous problem, if \(r_{1}, r_{2}\) are the maximum and minimum distances from the sun to the planet and \(v_{1}, v_{2}\) are the corresponding velocities, then, say,

\[
E=\frac{1}{2} m v_{2}^{2}-\frac{\gamma m m_{s}}{r_{2}}
\]

\[
=\frac{\gamma m m_{s}}{r_{1}+r_{2}} \cdot \frac{r_{1}}{r_{2}}-\frac{\gamma m m_{s}}{r_{2}}=-\frac{\gamma m m_{s}}{r_{1}+r_{2}}=-\frac{\gamma m m_{s}}{2 a} \text { [Using Eq. (2) of 1.207] }
\]

where \(2 a=\) major axis \(=r_{1}+r_{2}\). The same result can also be obtained directly by writing an equation analogous to Eq (1) of problem 1.191.

\[
E=\frac{1}{2} m \dot{r}^{2}+\frac{M^{2}}{2 m r^{2}}-\frac{\gamma m m_{s}}{r}
\]

(Here \(M\) is angular momentum of the planet and \(m\) is its mass). For extreme position \(\dot{r}=0\) and we get the quadratic

\[
E r^{2}+\gamma m m_{s} r-\frac{M^{2}}{2 m}=0
\]

The sum of the two roots of this equation are

\[
r_{1}+r_{2}=-\frac{\gamma m m_{s}}{E}=2 a
\]

Thus

\[
E=-\frac{\gamma m m_{s}}{2 a}=\text { constant }
\]

1.209 From the conservtion of angular momentum about the Sun.


\begin{equation*}
m v_{0} r_{0} \sin \alpha=m v_{1} r_{1}=m v_{2} r_{2} \text { or, } v_{1} r_{1}=v_{2} r_{2}=v_{0} r_{2} \sin \alpha \tag{1}
\end{equation*}


From conservation of mechanical energy,

\[
\begin{array}{lc} 
& \frac{1}{2} m v_{0}^{2}-\frac{\gamma m_{s} m}{r_{0}}=\frac{1}{2} m v_{1}^{2}-\frac{\gamma m_{s} m}{r_{1}} \\
\text { or, } & \frac{v_{0}^{2}}{2}-\frac{\gamma m_{s}}{r_{0}}=\frac{v_{0}^{2} r_{0}^{2} \sin ^{2} \alpha}{2 r_{1}^{2}}-\frac{\gamma m_{s}}{r_{1}}(\text { Using } 1) \\
\text { or, } & \left(v_{0}^{2}-\frac{2 \gamma m_{s}}{r_{0}}\right) r_{1}^{2}+2 \gamma m_{s} r_{1}-v_{0}^{2} r_{0}^{2} \sin \alpha=0 \\
\text { So, } & r_{1}=\frac{-2 \gamma m_{s} \pm \sqrt{4 \gamma^{2} m_{s}^{2}+4\left(v_{0}^{2} r_{0}^{2} \sin ^{2} \alpha\right)\left(v_{0}^{2}-\frac{2 \gamma m_{s}}{r_{0}}\right)}}{2\left(v_{0}^{2}-\frac{2 \gamma m_{s}}{r_{0}}\right)} \\
= & \frac{1 \pm \sqrt{1-\frac{v_{0}^{2} r_{0}^{2} \sin ^{2} \alpha}{\gamma m_{s}}\left(\frac{2}{r_{0}}-\frac{v_{0}^{2}}{r m_{s}}\right)}}{\left(\frac{2}{r_{0}}-\frac{v_{0}}{\gamma m_{s}}\right)}=\frac{r_{0}\left[1 \pm \sqrt{1-(2-\eta) \eta \sin ^{2} \alpha}\right]}{(2-\eta)}
\end{array}
\]

where \(\eta=v_{0}^{2} r_{0} / \gamma m_{s}\), ( \(m_{s}\) is the mass of the Sun).\\
1.210 At the minimum separation with the Sun, the cosmic body's velocity is perpendicular to its position vector relative to the Sun. If \(\boldsymbol{r}_{\text {min }}\) be the sought minimum distance, from conservation of angular momentum about the Sun (C).


\begin{equation*}
m v_{0} l=m v r_{\min } \text { or, } v=\frac{v_{0} l}{r_{\min }} \tag{1}
\end{equation*}


From conservation of mechanical energy of the system (sun + cosmic body),

\[
\frac{1}{2} m v_{0}^{2}=-\frac{\gamma m_{s} m}{r_{\min }}+\frac{1}{2} m v^{2}
\]

So,

\[
\frac{v_{0}^{2}}{2}=-\frac{\gamma m_{s}}{r_{\min }}+\frac{v_{0}^{2}}{2 r_{\min }^{2}} \quad(\text { using } 1)
\]

or,

\[
v_{0}^{2} r_{\min }^{2}+2 \gamma m_{s} r_{\min }^{\prime}-v_{0}^{2} l^{2}=0
\]

So,

\[
r_{\min }=\frac{-2 \gamma m_{s} \pm \sqrt{4 \gamma^{2} m^{2}+4 v_{0}^{2} v_{0}^{2} l^{2}}}{2 v_{0}^{2}}=\frac{-\gamma m_{s} \pm \sqrt{\gamma^{2} m_{s}^{2}+v_{0}^{4} l^{2}}}{v_{0}^{2}}
\]

Hence, taking positive root

\[
r_{\min }=\left(\gamma m_{s} / v_{0}^{2}\right)\left[\sqrt{1+\left(l v_{0}^{2} / \gamma m_{s}\right)^{2}}-1\right]
\]

1.211 Suppose that the sphere has a radius equal to \(a\). We may imagine that the sphere is made up of concentric thin spherical shells (layers) with radii ranging from 0 to \(a\), and each spherical layer is made up of elementry bands (rings). Let us first calculate potential due to an elementry band of a spherical layer at the point of location of the point mass \(m\) (say point \(P\) ) (Fig.). As all the points of the band are located at the distance \(l\) from the point \(P\), so,


\begin{equation*}
\partial \varphi=-\frac{\gamma \partial M}{l} \text { (where mass of the band) } \tag{1}
\end{equation*}



\begin{align*}
\partial M & =\left(\frac{d M}{4 \pi a^{2}}\right)(2 \pi a \sin \theta)(a d \theta) \\
& =\left(\frac{d M}{2}\right) \sin \theta d \theta \tag{2}
\end{align*}


And \(l^{2}=a^{2}+r^{2}-2 a r \cos \theta\)

Differentiating Eq. (3), we get


\begin{equation*}
l d l=\operatorname{ar} \sin \theta d \theta \tag{4}
\end{equation*}


Hence using above equations\\
\includegraphics[max width=\textwidth, alt={}, center]{9f3c5a8e-ac5d-42f4-8eb6-ea2af623f2d8-113_365_618_1617_731}


\begin{equation*}
\partial \varphi=-\left(\frac{\gamma d M}{2 a r}\right) d l \tag{5}
\end{equation*}


Now integrating this Eq. over the whole spherical layer


\begin{align*}
d \varphi=\int \partial \varphi & =-\frac{\gamma d M}{2 a r} \int_{r-a}^{\bar{r}+a} \\
d \varphi & =-\frac{\gamma d M}{r} \tag{6}
\end{align*}


So

This expression is similar to that of Eq. (6)\\
Hence the sought potential energy of gravitational interaction of the particle \(m\) and the sphere,

\[
U=m \varphi=-\frac{\gamma M m}{r}
\]

(b) Using the Eq.,

\[
G_{r}=-\frac{\partial \varphi}{\partial r}
\]

\[
G_{r}=-\frac{\gamma M}{r^{2}} \quad \text { (using Eq. } 7 \text { ) }
\]

So


\begin{equation*}
\vec{G}=-\frac{\gamma M}{r^{3}} \vec{r} \text { and } \vec{F}=m \vec{G}=-\frac{\gamma m M}{r^{3}} \vec{r} \tag{8}
\end{equation*}


1.212 (The problem has already a clear hint in the answer sheet of the problem book). Here we adopt a different method.\\
Let \(m\) be the mass of the spherical layer, wich is imagined to be made up of rings. At a point inside the spherical layer at distance \(r\) from the centre, the gravitational potential due to a ring element of radius \(a\) equals,\\
\(d \varphi=-\frac{\gamma m}{2 a r} d l\) (see Eq. (5) of solution of 1.211)\\
So, \(\varphi=\int d \varphi=-\frac{\gamma m}{2 a r} \int_{a-r}^{a+r} d l=-\frac{\gamma m}{a}\)\\
\includegraphics[max width=\textwidth, alt={}, center]{9f3c5a8e-ac5d-42f4-8eb6-ea2af623f2d8-114_358_357_1392_924}

Hence

\[
G_{r}=-\frac{\partial \varphi}{\partial r}=0 .
\]

Hence gravitational field strength as well as field force becomes zero, inside a thin sphereical layer.\\
1.213 One can imagine that the uniform hemisphere is made up of thin hemispherical layers of radii ranging from 0 to \(R\). Let us consider such a layer (Fig.). Potential at point O , due to this layer is,\\
\(d \varphi=-\frac{\gamma d m}{r}=-\frac{3 \gamma M}{R^{3}} r d r\), where \(d m=\frac{M}{(2 / 3) \pi R^{3}}\left(\frac{4 \pi r^{2}}{2}\right) d r\)\\
(This is because all points of each hemispherical shell are equidistant from \(O\).)\\
Hence, \(\quad \varphi=\int d \varphi=-\frac{3 \gamma M}{R^{3}} \int_{0}^{R} r d r=-\frac{3 \gamma M}{2 R}\)\\
Hence, the work done by the gravitational field force on the particle of mass \(m\), to remove it to infinity is given by the formula\\
\(A=m \varphi\), since \(\varphi=0\) at infinity.\\
Hence the sought work,

\[
A_{0 \rightarrow \infty}=-\frac{3 \gamma m M}{2 R}
\]

\includegraphics[max width=\textwidth, alt={}, center]{9f3c5a8e-ac5d-42f4-8eb6-ea2af623f2d8-115_430_576_265_784}\\
(The work done by the external agent is \(-A\).)\\
1.214 In the solution of problem 1.211, we have obtained \(\varphi\) and \(G\) due to a uniform shpere, at a distance \(r\) from it's centre outside it. We have from Eqs. (7) and (8) of 1.211,


\begin{equation*}
\varphi=-\frac{\gamma M}{r} \text { and } \vec{G}=-\frac{\gamma M}{r^{3}} \vec{r} \tag{A}
\end{equation*}


Accordance with the Eq. (1) of the solution of 1.212 , potential due to a spherical shell of radius \(a\), at any point, inside it becomes


\begin{equation*}
\varphi=\frac{\gamma M}{a}=\text { Const. and } G_{r}=-\frac{\partial \varphi}{\partial r}=0 \tag{B}
\end{equation*}


For a point (say \(P\) ) which lies inside the uniform solid sphere, the potential \(\varphi\) at that point may be represented as a sum.

\[
\varphi_{\text {inside }}=\varphi_{1}+\varphi_{2}
\]

where \(\varphi_{1}\) is the potential of a solid sphere having radius \(r\) and \(\varphi_{2}\) is the potential of the layer of radii \(r\) and \(R\). In accordance with equation (A)

\[
\varphi_{1}=-\frac{\gamma}{r}\left(\frac{M}{(4 / 3) \pi R^{3}} \frac{4}{3} \pi r^{3}\right)=-\frac{\gamma M}{R^{3}} r^{2}
\]

The potential \(\varphi_{2}\) produced by the layer (thick shell) is the same at all points inside it. The potential \(\varphi_{2}\) is easiest to calculate, for the point positioned at the layer's centre. Using Eq. (B)

\[
\varphi_{2}=-\gamma \int_{r}^{R} \frac{d M}{r}=-\frac{3}{2} \frac{\gamma M}{R^{3}}\left(R_{1}^{2}-r^{2}\right)
\]

where \(d M=\frac{M}{(4 / 3) \pi R^{3}} 4 \pi r^{2} d r=\left(\frac{3 M}{R^{3}}\right) r^{2} d r\) is the mass of a thin layer between the radii \(r\) and \(r+d r\).

Thus


\begin{equation*}
\varphi_{\text {inside }}=\varphi_{1}+\varphi_{2}=\left(\frac{\gamma M}{2 R}\right)\left(3-\frac{r^{2}}{R^{2}}\right) \tag{C}
\end{equation*}


From the Eq.\\
or


\begin{gather*}
G_{r}=\frac{-\partial \varphi}{\partial r} \\
G_{r}=\frac{\gamma M r}{R^{3}}  \tag{D}\\
\vec{G}=-\frac{\gamma M}{R^{3}} \vec{r}=-\gamma \frac{4}{3} \pi \rho \vec{r}
\end{gather*}


(where \(\rho=\frac{M}{\frac{4}{3} \pi R^{3}}\), is the density of the sphere)

The plots \(\varphi(r)\) and \(G(r)\) for a uniform sphere of radius \(R\) are shown in figure of answersheet. Alternate : Like Gauss's theorem of electrostatics, one can derive Gauss's theorem for gravitation in the form \(\oint \vec{G} \cdot d \vec{S}=-4 \pi \gamma m_{\text {inclosed }}\). For calculation of \(\vec{G}\) at a point inside the sphere at a distance \(r\) from its centre, let us consider a Gaussian surface of radius \(r\), Then,

\[
G_{r} 4 \pi r^{2}=-4 \pi \gamma\left(\frac{M}{R^{3}}\right) r^{3} \text { or, } G_{r}=-\frac{\gamma M}{R^{3}} r
\]

Hence,

\[
\vec{G}=-\frac{\gamma M}{R^{3}} \vec{r}=-\gamma \frac{4}{3} \pi \rho \vec{r}\left(\text { as } \rho=\frac{M}{(4 / 3) \pi R^{3}}\right)
\]

So,

\[
\varphi=\int_{r}^{\infty} G_{r} d r=\int_{r}^{R}-\frac{\gamma M}{R^{3}} r d r+\int_{R}^{\infty}-\frac{\gamma M}{r^{2}} d r
\]

Integrating and summing up, we get,

\[
\varphi=-\frac{\gamma M}{2 R}\left(3-\frac{r^{2}}{R^{2}}\right)
\]

And from Gauss's theorem for outside it :

\[
G_{r} 4 \pi r^{2}=-4 \pi \gamma M \quad \text { or } \quad G_{r}=-\frac{\gamma M}{r^{2}}
\]

Thus

\[
\varphi(r)=\int_{r}^{\infty} G_{r} d r=-\frac{\gamma M}{r}
\]

1.215 Treating the cavity as negative mass of density \(-\rho\) in a uniform sphere density \(+\rho\) and using the superposition principle, the sought field strength is :

\[
\vec{G}=\vec{G}_{1}+\vec{G}_{2}
\]

or \(\vec{G}=-\frac{4}{3} \pi \gamma \rho \overrightarrow{r_{+}}+-\frac{4}{3} \gamma \pi(-\rho) \overrightarrow{r_{-}}\)\\
(where \(\overrightarrow{r_{+}}\)and \(\overrightarrow{r_{-}}\)are the position vectors of an orbitrary point \(P\) inside the cavity with respect to centre of sphere and cavity respectively.)\\
Thus \(\quad \vec{G}=-\frac{4}{3} \pi \gamma \rho\left(\overrightarrow{r_{+}}-\overrightarrow{r_{-}}\right)=-\frac{4}{3} \pi \gamma \rho \vec{l}\)\\
\includegraphics[max width=\textwidth, alt={}, center]{9f3c5a8e-ac5d-42f4-8eb6-ea2af623f2d8-116_357_359_1697_855}\\
1.216 We partition the solid sphere into thin spherical layers and consider a layer of thickness \(d r\) lying at a distance \(r\) from the centre of the ball. Each spherical layer presses on the layers within it. The considered layer is attracted to the part of the sphere lying within it (the outer part does not act on the layer). Hence for the considered layer

\[
d p 4 \pi r^{2}=d F
\]

or, \(\quad d P 4 \pi r^{2}=\frac{\gamma\left(\frac{4}{3} \pi r^{3} \rho\right)\left(4 \pi r^{2} d r \rho\right)}{r^{2}}\) (where \(\rho\) is the mean density of sphere)\\
or, \(\quad d p=\frac{4}{3} \pi \gamma \rho^{2} r d r\)\\
Thus \(\quad p=\int_{r}^{R} d p=\frac{2 \pi}{3} \gamma \rho^{2}\left(R^{2}-r^{2}\right)\)\\
\includegraphics[max width=\textwidth, alt={}, center]{9f3c5a8e-ac5d-42f4-8eb6-ea2af623f2d8-117_420_408_357_868}\\
(The pressure.must vanish at \(r=R\).)\\
or, \(\quad p=\frac{3}{8}\left(1-\left(\dot{r}^{2} / R^{2}\right)\right) \gamma M^{2} / \pi R^{4}\), Putting \(\rho=M /(4 / 3) \pi R^{3}\)\\
Putting \(r=0\), we have the pressure at sphere's centre, and treating it as the Earth where mean density is equal to \(\rho=5.5 \times 10^{3} \mathrm{~kg} / \mathrm{m}^{3}\) and \(R=64 \times 10^{2} \mathrm{~km}\)\\
we have,

\[
p=1.73 \times 10^{11} \mathrm{~Pa} \text { or } 1.72 \times 10^{6} \mathrm{atms} .
\]

1.217 (a) Since the potential at each point of a spherical surface (shell) is constant and is equal to \(\varphi=-\frac{\gamma m}{R}\), [as we have in Eq. (1) of solution of problem 1.212]\\
We obtain in accordance with the equation

\[
\begin{aligned}
U & =\frac{1}{2} \int d m \varphi=\frac{1}{2} \varphi \int d m \\
& =\frac{1}{2}\left(-\frac{\gamma m}{R}\right) m=-\frac{\gamma m^{2}}{2 R}
\end{aligned}
\]

(The factor \(\frac{1}{2}\) is needed otherwise contribution of different mass elements is counted twice.)\\
(b) In this case the potential inside the sphere depends only on \(r\) (see Eq. (C) of the solution of problem 1.214)

\[
\varphi=-\frac{3 \gamma m}{2 R}\left(1-\frac{r^{2}}{3 R^{2}}\right)
\]

Here \(d m\) is the mass of an elementry spherical layer confined between the radii \(r\) and \(r+d r\) :

\[
d m=\left(4 \pi r^{2} d r \rho\right)=\left(\frac{3 m}{R^{3}}\right) r^{2} d r
\]

\[
\begin{gathered}
U=\frac{1}{2} \int d m \varphi \\
=\frac{1}{2} \int_{0}^{R}\left(\frac{3 m}{R^{3}}\right) r^{2} d r\left\{-\frac{3 \gamma m}{2 R}\left(1-\frac{r^{2}}{3 R^{2}}\right)\right\}
\end{gathered}
\]

After integrating, we get

\[
U=-\frac{3}{5} \frac{\gamma m^{2}}{R}
\]

1.218 Let \(\omega=\sqrt{\frac{\gamma M_{E}}{r^{3}}}=\) circular frequency of the satellite in the outer orbit,\\
\(\omega_{0}=\sqrt{\frac{\gamma M_{E}}{(r-\Delta r)^{3}}}=\) circular frequency of the satellite in the inner orbit.

So, relative angular velocity \(=\omega_{0} \pm \omega\) where - sign is to be taken when the satellites are moving in the same sense and + sign if they are moving in opposite sense. Hence, time between closest approaches

\[
=\frac{2 \pi}{\omega_{0} \pm \omega}=\frac{2 \pi}{\sqrt{\gamma M_{E}} / r^{3 / 2}} \frac{1}{\frac{3 \Delta r}{2 r}+\delta}=\left\{\begin{array}{c}
4.5 \text { days }(\delta=0) \\
0.80 \text { hour }(\delta=2)
\end{array}\right.
\]

where \(\delta\) is 0 in the first case and 2 in the second case.\\
\(1.219 \omega_{1}=\frac{\gamma M}{R^{2}}=\frac{6.67 \times 10^{-11} \times 5.96 \times 10^{24}}{\left(6.37 \times 10^{6}\right)^{2}}=9.8 \mathrm{~m} / \mathrm{s}^{2}\)\\
\(\omega_{2}=\omega^{2} R=\left(\frac{2 \pi}{T}\right)^{2} R=\left(\frac{2 \times 22}{24 \times 3600 \times 7}\right)^{2} 6.37 \times 10^{6}=0.034 \mathrm{~m} / \mathrm{s}^{2}\)\\
and \(\quad \omega_{3}=\frac{\gamma M_{S}}{R_{\text {mean }}^{2}}=\frac{6.67 \times 10^{-11} \times 1.97 \times 10^{30}}{\left(149.50 \times 10^{6} \times 10^{3}\right)^{2}}=5.9 \times 10^{-3} \mathrm{~m} / \mathrm{s}^{2}\)

Then

\[
\omega_{1}: \omega_{2}: \omega_{3}=1: 0.0034: 0.0006
\]

1.220 Let \(h\) be the sought height in the first case. so

\[
\begin{gathered}
\frac{99}{100} g=\frac{\gamma M}{(R+h)^{2}} \\
=\frac{\gamma M}{R^{2}\left(1+\frac{h}{R}\right)^{2}}=\frac{g}{\left(1+\frac{h}{R}\right)^{2}}
\end{gathered}
\]

or

\[
\frac{99}{100}=\left(1+\frac{h}{R}\right)^{-2}
\]

From the statement of the problem, it is obvious that in this case \(h \ll R\)\\
Thus \(\quad \frac{99}{100}=\left(1-\frac{2 h}{R}\right)\) or \(h=\frac{R}{200}=\left(\frac{6400}{200}\right) \mathrm{km}=32 \mathrm{~km}\)\\
In the other case if \(h^{\prime}\) be the sought height, than

\[
\frac{g}{2}=g\left(1+\frac{h^{\prime}}{R}\right)^{-2} \text { or } \frac{1}{2}=\left(1+\frac{h^{\prime}}{R}\right)^{-2}
\]

From the language of the problem, in this case \(\boldsymbol{h}^{\prime}\) is not very small in comparision with \(\boldsymbol{R}\). Therefore in this case we cannot use the approximation adopted in the previous case.\\
Here, \(\left(1+\frac{h^{\prime}}{R}\right)^{2}=2 \quad\) So, \(\frac{h^{\prime}}{R}= \pm \sqrt{2}-1\)\\
As - ve sign is not acceptable

\[
h^{\prime}=(\sqrt{2}-1) R=(\sqrt{2}-1) 6400 \mathrm{~km}=2650 \mathrm{~km}
\]

1.221 Let the mass of the body be \(m\) and let it go upto a height \(h\).

From conservation of mechanical energy of the system

\[
-\frac{\gamma M m}{R}+\frac{1}{2} m v_{0}^{2}=\frac{-\gamma M m}{(R+h)}+0
\]

Using \(\frac{\gamma M}{R^{2}}=g\), in above equation and on solving we get,

\[
h=\frac{R v_{0}^{2}}{2 g R-v_{0}^{2}}
\]

1.222 Gravitational pull provides the required centripetal acceleration to the satelite. Thus if \(h\) be the sought distance, we have\\
so, \(\quad \frac{m v^{2}}{(R+h)}=\frac{\gamma m M}{(R+h)^{2}}\) or, \((R+h) v^{2}=\gamma M\)\\
or, \(\quad R v^{2}+h v^{2}=g R^{2}\), as \(g=\frac{\gamma M}{R^{2}}\)\\
Hence

\[
h=\frac{g R^{2}-R v^{2}}{v^{2}}=R\left[\frac{g R}{v^{2}}-1\right]
\]

1.223 A satellite that hovers above the earth's equator and corotates with it moving from the west to east with the diurnal angular velocity of the earth appears stationary to an observer on the earth. It is called geostationary. For this calculation we may neglect the annual motion of the earth as well as all other influences. Then, by Newton's law,

\[
\frac{\gamma M m}{r^{2}}=m\left(\frac{2 \pi}{T}\right)^{2} r
\]

where \(M=\) mass of the earth, \(T=86400\) seconds = period of daily rotation of the earth and \(r=\) distance of the satellite from the centre of the earth. Then

\[
r=\sqrt[3]{\gamma M\left(\frac{T}{2 \pi}\right)^{2}}
\]

Substitution of \(M=5.96 \times 10^{24} \mathrm{~kg}\) gives

\[
r=4.220 \times 10^{4} \mathrm{~km}
\]

The instantaneous velocity with respect to an inertial frame fixed to the centre of the earth at that moment will be

\[
\left(\frac{2 \pi}{T}\right) r=3.07 \mathrm{~km} / \mathrm{s}
\]

and the acceleration will be the centripetal acceleration.

\[
\left(\frac{2 \pi}{T}\right)^{2} r=0.223 \mathrm{~m} / \mathrm{s}^{2}
\]

1.224 We know from the previous problem that a satellite moving west to east at a distance \(R=2.00 \times 10^{4} \mathrm{~km}\) from the centre of the earth will be revolving round the earth with an angular velocity faster than the earth's diurnal angualr velocity. Let

\[
\begin{gathered}
\omega=\text { angular velocity of the satellite } \\
\omega_{0}=\frac{2 \pi}{T}=\text { anuglar velocity of the earth. Then } \\
\omega-\omega_{0}=\frac{2 \pi}{\tau}
\end{gathered}
\]

as the relative angular velocity with respect to earth. Now by Newton's law

So,

\[
\begin{gathered}
\frac{\gamma M}{R^{2}}=\omega^{2} R \\
M=\frac{R^{3}}{\gamma}\left(\frac{2 \pi}{\tau}+\frac{2 \pi}{T}\right)^{2} \\
=\frac{4 \pi^{2} R^{3}}{\gamma T^{2}}\left(1+\frac{T}{\tau}\right)^{2}
\end{gathered}
\]

Substitution gives

\[
M=6.27 \times 10^{24} \mathrm{~kg}
\]

1.225 The velocity of the satellite in the inertial space fixed frame is \(\sqrt{\frac{\gamma M}{R}}\) east to west. With respect to the Earth fixed frame, from the \(\vec{v}_{1}^{\prime}=\vec{v}-(\vec{w} \times \vec{r})\) the velocity is

\[
v^{\prime}=\frac{2 \pi R}{T}+\sqrt{\frac{\gamma M}{R}}=7.03 \mathrm{~km} / \mathrm{s}
\]

Here \(M\) is the mass of the earth and \(T\) is its period of rotation about its own axis.

It would be \(-\frac{2 \pi R}{T}+\sqrt{\frac{Y M}{R}}\), if the satellite were moving from west to east.\\
To find the acceleration we note the formula

\[
m \overrightarrow{\omega^{\prime}}=\vec{F}+2 m\left(\overrightarrow{v^{\prime}} \times \vec{\omega}\right)+m \omega^{2} \vec{R}
\]

Here \(\vec{F}=-\frac{\gamma M m}{R^{3}} \vec{R}\) and \(\vec{v}^{\prime} \perp \vec{\omega}\) and \(\vec{v} \times \vec{\omega}\) is directed towards the centre of the Earth.

Thus

\[
w^{\prime}=\frac{\gamma M}{R^{2}}+2\left(\frac{2 \pi R}{T}+\sqrt{\frac{\gamma M}{R}}\right) \frac{2 \pi}{T}-\left(\frac{2 \pi}{T}\right)^{2} R
\]

toward the earth's rotation axis

\[
=\frac{\gamma M}{R^{2}}+\frac{2 \pi}{T}\left[\frac{2 \pi R}{T}+2 \sqrt{\frac{\gamma M}{R}}\right]=4.94 \mathrm{~m} / \mathrm{s}^{2} \text { on substitution. }
\]

1.226 From the well known relationship between the velocities of a particle w.r.t a space fixed frame \((K)\) rotating frame \(\left(K^{\prime}\right) \vec{v}=\vec{v}{ }^{\prime}+(\vec{w} \times \vec{r})\)

\[
v_{1}^{\prime}=v-\left(\frac{2 \pi}{T}\right) R
\]

Thus kinetic energy of the satellite in the earth's frame

\[
T_{1}^{\prime}=\frac{1}{2} m v_{1}^{\prime 2}=\frac{1}{2} m\left(v-\frac{2 \pi R}{T}\right)^{2}
\]

Obviously when the satellite moves in opposite sense comared to the rotation of the Earth its velocity relative to the same frame would be

And kinetic energy

\[
v_{2}^{\prime}=v+\left(\frac{2 \pi}{T}\right) R
\]


\begin{equation*}
T_{2}^{\prime}=\frac{1}{2} m v_{2}^{\prime 2}=\frac{1}{2} m\left(v+\frac{2 \pi R}{T}\right)^{2} \tag{2}
\end{equation*}


From (1) and (2)


\begin{equation*}
T^{\prime}=\frac{\left(v+\frac{2 \pi R}{T}\right)^{2}}{\left(v-\frac{2 \pi R}{T}\right)^{2}} \tag{3}
\end{equation*}


Now from Newton's second law


\begin{equation*}
\frac{\gamma M m}{R^{2}}=\frac{m v^{2}}{R} \quad \text { or } \quad v=\sqrt{\frac{\gamma M}{R}}=\sqrt{g R} \tag{4}
\end{equation*}


Using (4) and (3) •

\[
\frac{T_{2}^{\prime}}{T_{1}^{\prime}}=\frac{\left(\sqrt{g R}+\frac{2 \pi R}{T}\right)^{2}}{\left(\sqrt{g R}-\frac{2 \pi R}{T}\right)^{2}}=1.27 \quad \text { nearly (Using Appendices) }
\]

1.227 For a satellite in a circular orbit about any massive body, the following relation holds between kinetic, potential \& total energy :


\begin{equation*}
T=-E, U=2 E \tag{1}
\end{equation*}


Thus since total mechanical energy must decrease due to resistance of the cosmic dust, the kintetic energy will increase and the satellite will 'fall', We see then, by work energy theorm

\[
d T=-d E=-d A_{f r}
\]

So, \(\quad m v d v=\alpha v^{2} v d t \quad\) or, \(\quad \frac{\alpha d t}{m}=\frac{d v}{v^{2}}\)\\
Now from Netow's law at an arbitray radius \(r\) from the moon's centre.

\[
\frac{v^{2}}{r}=\frac{\gamma M}{r^{2}} \text { or } v=\sqrt{\frac{\gamma M}{r}}
\]

( \(M\) is the mass of the moon.) Then

\[
v_{i}=\sqrt{\frac{\gamma M}{\eta R}}, v_{f}=\sqrt{\frac{\gamma M}{R}}
\]

where \(R\) = moon's radius. So

\[
\int_{v_{1}}^{v_{f}} \frac{d v}{v^{2}}=\frac{\alpha}{m} \int_{0}^{\tau} d t=\frac{\alpha \tau}{m}
\]

or,

\[
\tau=\frac{m}{\alpha}\left(\frac{1}{v_{i}}-\frac{1}{v_{f}}\right)=\frac{m}{\alpha \sqrt{\frac{M}{R}}}(\sqrt{\eta}-1)=\frac{m}{\alpha \sqrt{g R}}(\sqrt{\eta}-1)
\]

where \(g\) is moon's gravity. The averaging implied by Eq. (1) (for noncircular orbits) makes the result approximate.\\
1.228 From Newton's second law


\begin{equation*}
\frac{\gamma M m}{R^{2}}=\frac{m v_{0}^{2}}{R} \text { or } v_{0}=\sqrt{\frac{\gamma M}{R}}=1.67 \mathrm{~km} / \mathrm{s} \tag{1}
\end{equation*}


From conservation of mechanical energy


\begin{equation*}
\frac{1}{2} m v_{e}^{2}-\frac{\gamma M m}{R}=0 \text { or, } v_{e}=\sqrt{\frac{2 \gamma \bar{M}}{R}}=2.37 \mathrm{~km} / \mathrm{s}^{2} \tag{2}
\end{equation*}


In Eq. (1) and (2), \(M\) and \(R\) are the mass of the moon and its radius. In Eq. (1) if \(M\) and \(R\) represent the mass of the earth and its radius, then, using appendices, we can easily get

\[
v_{0}=7.9 \mathrm{~km} / \mathrm{s} \text { and } v_{c}=11.2 \mathrm{~km} / \mathrm{s}
\]

1.229 In a parabolic orbit, \(E=0\)

So

\[
\frac{1}{2} m v_{i}^{2}-\frac{\gamma M m}{R}=0 \text { or, } v_{i}=\sqrt{2} \sqrt{\frac{\gamma M}{R}}
\]

where \(M=\) mass of the Moon, \(R=\) its radius. (This is just the escape velocity.)\\
On the other hand in orbit

\[
m v_{f}^{2} R=\frac{\gamma M m}{R^{2}} \text { or } v_{f}=\sqrt{\frac{\gamma M}{R}}
\]

Thus

\[
\Delta v=(1-\sqrt{2}) \sqrt{\frac{\gamma M}{R}}=-0.70 \mathrm{~km} / \mathrm{s} .
\]

1.230 From 1.228 for the Earth surface

\[
\dot{v}_{0}=\sqrt{\frac{\gamma M}{R}} \text { and } v_{e}=\sqrt{\frac{2 \gamma M}{R}}
\]

Thus the sought additional velocity

\[
\Delta v=v_{e}-v_{0}=\sqrt{\frac{\gamma M}{R}}(\sqrt{2}-1)=g R(\sqrt{2}-1)
\]

This 'kick' in velocity must be given along the direction of motion of the satellite in its orbit.\\
1.231 Let \(r\) be the sought distance, then\\
or

\[
\begin{gathered}
\frac{\eta \eta M}{(n R-r)^{2}}=\frac{\gamma M}{r^{2}} \text { or } \eta r^{2}=(n R-r)^{2} \\
\sqrt{\eta} r=(n R-r) \text { or } r=\frac{n R}{\sqrt{\eta}+1}=3.8 \times 10^{4} \mathrm{~km} .
\end{gathered}
\]

1.232 Between the earth and the moon, the potential energy of the spaceship will have a maximum at the point where the attractions of the earth and the moon balance each other. This maximum P.E. is approximately zero. We can also neglect the contribution of either body to the p.E. of the spaceship sufficiently near the other body. Then the minimum energy that must be imparted to the spaceship to cross the maximum of the P.E. is clearly (using \(E\) to denote the earth)

\[
\frac{\gamma M_{E} m}{R_{E}}
\]

With this energy the spaceship will cross over the hump in the P.E. and coast down the hill of p.E. towards the moon and crashland on it. What the problem seeks is the minimum energy reguired for softlanding. That reguies the use of rockets to loving about the braking of the spaceship and since the kinetic energy of the gases ejected from the rocket will always be positive, the total energy required for softlanding is greater than that required for crashlanding. To calculate this energy we assume that the rockets are used fairly close to the moon when the spaceship has nealy attained its terminal velocity on the moon \(\sqrt{\frac{2 \gamma M_{O}}{R_{0}}}\) where \(M_{0}\) is the mass of the moon and \(R_{0}\) is its radius. In general \(d E=v d p\) and since the speed of the ejected gases is not less than the speed of the rocket, and momentum transfered to the ejected gases must equal the momentum of the spaceship the energy \(E\) of the gass ejected is not less than the kinetic energy of spaceship

\[
\frac{\gamma M_{0} m}{R_{0}}
\]

Addding the two we get the minimum work done on the ejected gases to bring about the softlanding.

\[
A_{\min } \approx \gamma m\left(\frac{M_{E}}{R_{E}}+\frac{M_{0}}{R_{0}}\right)
\]

On substitution we get \(1.3 \times 10^{8} \mathrm{~kJ}\).

Assume first that the attraction of the earth can be neglected. Then the minimum velocity, that must be imparted to the body to escape from the Sun's pull, is, as in \(1 \cdot 230\), equal to

\[
(\sqrt{2}-1) v_{1}
\]

where \(v_{1}^{2}=\gamma M_{s} / r, r=\) radius of the earth's orbit, \(M_{s}=\) mass of the Sun.\\
In the actual case near the earth, the pull of the Sun is small and does not change much over distances, which are several times the radius of the Earth. The velocity \(v_{3}\) in question is that which overcomes the earth's pull with sufficient velocity to escape the Sun's pull. Thus

\[
\frac{1}{2} m v_{3}^{2}-\frac{\gamma M_{E}}{R}=\frac{1}{2} m(\sqrt{2}-1)^{2} v_{1}^{2}
\]

where \(R=\) radius of the earth, \(M_{E}=\) mass of the earth.\\
Writing \(v_{1}^{2}=\gamma M_{E} / R\), we get

\[
v_{3}=\sqrt{2 v_{2}^{2}+(\sqrt{2}-1)^{2} v_{1}^{2}}=16.6 \mathrm{~km} / \mathrm{s}
\]

\subsection*{1.5 DYNAMICS OF A SOLID BODY}
1.234 Since, motion of the rod is purely translational, net torque about the C.M. of the rod should be equal to zero.

Thus


\begin{equation*}
F_{1} \frac{l}{2}=F_{2}\left(\frac{l}{2}-a\right) \text { or, } \frac{F_{1}}{F_{2}}=1-\frac{a}{l / 2} \tag{1}
\end{equation*}


For the translational motion of rod.


\begin{equation*}
F_{2}-F_{1}=m w_{c} \text { or } 1-\frac{F_{1}}{F_{2}}=\frac{m w_{c}}{F_{2}} \tag{2}
\end{equation*}


From (1) and (2)

\[
\frac{a}{l / 2}=\frac{m w_{c}}{F_{2}} \text { or, } l=\frac{2 a F_{2}}{m w_{c}}=1 \mathrm{~m}
\]

1.235 Sought moment \(\vec{N}=\overrightarrow{r \times} \vec{F}=(a \vec{i}+b \vec{j}) \times(A \vec{i}+B \vec{j})\)

\[
=a B \vec{k}+A b(-\vec{k})=(a B-A b) \vec{k}
\]

and arm of the force

\[
l=\frac{N}{F}=\frac{a B-A b}{\sqrt{A^{2}+B^{2}}}
\]

1.236 Relative to point \(O\), the net moment of force :


\begin{align*}
\vec{N} & =\overrightarrow{r_{1}} \times \overrightarrow{F_{1}}+\overrightarrow{r_{2}} \times \overrightarrow{F_{2}}=(a \vec{i} \times A \vec{j})+(\overrightarrow{B j} \times \overrightarrow{B i}) \\
& =a b \vec{k}+A B(-\vec{k})=(a b-A B) \vec{k} \tag{1}
\end{align*}


Resultant of the external force


\begin{equation*}
\vec{F}=\vec{F}_{1}+\vec{F}_{2}=A \vec{j}+B \vec{i} \tag{2}
\end{equation*}


As \(\vec{N} \cdot \vec{F}=0(\) as \(\vec{N} \perp \vec{F})\) so the sought arm \(l\) of the force \(\vec{F}\)

\[
l=N / F=\frac{a b-A B}{\sqrt{A^{2}+B^{2}}}
\]

1.237 For coplanar forces, about any point in the same plane, \(\sum \overrightarrow{r_{i}} \times \overrightarrow{F_{i}}=\vec{r} \times \vec{F}_{n e t}\) (where \(\vec{F}_{n e t}=\sum \vec{F}_{i}=\) resultant force) or, \(\vec{N}_{n e t}=\vec{r} \times \vec{F}_{n e t}\)\\
Thus length of the arm, \(l=\frac{N_{\text {net }}}{F_{\text {net }}}\)\\
Here obviously \(\left|\vec{F}_{\text {net }}\right|=2 F\) and it is directed toward right along \(A C\). Take the origin at \(C\). Then about \(C\),\\
\(\vec{N}=\left(\sqrt{2} a F+\frac{a}{\sqrt{2}} F-\sqrt{2} a F\right)\) directed normally into the plane of figure.\\
(Here \(a=\) side of the square.)\\
Thus \(\vec{N}=F \frac{a}{\sqrt{2}}\) directed into the plane of the figure.\\
Hence

\[
l=\frac{F(a / \sqrt{2})}{2 F}=\frac{a}{2 \sqrt{2}}=\frac{a}{2} \sin 45^{\circ}
\]

Thus the point of application of force is at the mid point of the side \(B C\).\\
1.238 (a) Consider a strip of length \(d x\) at a perpendicular distance \(x\) from the axis about which we have to find the moment of inertia of the rod. The elemental mass of the rod equals

\[
d m=\frac{m}{l} d x
\]

Moment of inertia of this element about the axis

\[
d I=d m x^{2}=\frac{m}{l} d x \cdot x^{2}
\]

Thus, moment of inertia of the rod, as a whole about the given axis

\[
I=\int_{0}^{l} \frac{m}{l} x^{2} d x=\frac{m l^{2}}{3}
\]

(b) Let us imagine the plane of plate as \(x y\) plane taking the origin at the intersection point of the sides of the plate (Fig.).

Obviously\\
\includegraphics[max width=\textwidth, alt={}, center]{9f3c5a8e-ac5d-42f4-8eb6-ea2af623f2d8-126_367_397_576_830}

Similarly

\[
\begin{gathered}
I_{x}=\int d m y^{2} \\
=\int_{0}^{a}\left(\frac{m}{a b} b d y\right) y^{2} \\
=\frac{m a^{2}}{3} \\
I_{y}=\frac{m b^{2}}{3}
\end{gathered}
\]

Hence from perpendicular axis theorem

\[
I_{z}=I_{x}+I_{y}=\frac{m}{3}\left(a^{2}+b^{2}\right)
\]

which is the sought moment of inertia.\\
1.239 (a) Consider an elementry disc of thickness \(d x\). Moment of inertia of this element about the \(z\)-axis, passing through its C.M.

\[
d I_{z}=\frac{(d m) R^{2}}{2}=\rho S d x \frac{R^{2}}{2}
\]

where \(\rho=\) density of the material of the plate and \(S=\) area of cross section of the plate.\\
Thus the sought moment of inertia

\[
\begin{aligned}
I_{z} & =\frac{\rho S R^{2}}{2} \int_{0}^{b} d x=\frac{R^{2}}{2} \rho S b \\
& =\frac{\pi}{2} \rho b R^{4}\left(\text { as } S=\pi R^{2}\right)
\end{aligned}
\]

\includegraphics[max width=\textwidth, alt={}, center]{9f3c5a8e-ac5d-42f4-8eb6-ea2af623f2d8-126_417_535_1520_824}\\
putting all the vallues we get, \(I_{z}=2 \cdot \mathrm{gm} \cdot \mathrm{m}^{2}\) (b) Consider an element disc of radius \(r\) and thickness \(d x\) at a distance \(x\) from the point \(O\). Then \(r=x \tan \alpha\) and volume of the disc

\[
=\pi x^{2} \tan ^{2} \alpha d x
\]

Hence, its mass \(d m=\pi x^{2} \tan \alpha d x \cdot \rho\) (where \(\rho=\) density of the cone \(=m / \frac{1}{3} \pi R^{2} h\) )\\
Moment of inertia of this element, about the\\
\includegraphics[max width=\textwidth, alt={}, center]{9f3c5a8e-ac5d-42f4-8eb6-ea2af623f2d8-127_421_518_133_816}\\
axis \(O A\),

\[
\begin{gathered}
d I=d m \frac{r^{2}}{2} \\
=\left(\pi x^{2} \tan ^{2} \alpha d x\right) \frac{x^{2} \tan ^{2} \alpha}{2} \\
=\frac{\pi \rho}{2} x^{4} \tan ^{4} \alpha d x
\end{gathered}
\]

Thus the sought moment of inertia \(I=\frac{\pi \rho}{2} \tan ^{4} \alpha \int_{0}^{h} x^{4} d x\)

\[
=\frac{\pi \rho R^{4} \cdot h^{5}}{10 h^{4}}\left(\text { as } \tan \alpha=\frac{R}{h}\right)
\]

Hence

\[
I=\frac{3 m R^{2}}{10}\left(\text { putting } \rho=\frac{3 m}{\pi R^{2} h}\right)
\]

\includegraphics[max width=\textwidth, alt={}, center]{9f3c5a8e-ac5d-42f4-8eb6-ea2af623f2d8-127_543_336_692_1019}\\
L. 240 (a) Let us consider a lamina of an arbitrary shape and indicate by 1,2 and 3 , three axes coinciding with \(x, y\) and \(z\)-axes and the plane of lamina as \(x-y\) plane.\\
Now, moment of inertia of a point mass about \(x\) - axis, \(d I_{x}=d m y^{2}\)\\
Thus moment of inertia of the lamina about this axis, \(I_{x}=\int d m y^{2}\)\\
Similarly, \(I_{y}=\int d m x^{2}\) and \(I_{z}=\int d m r^{2}\)

\[
=\int d m\left(x^{2}+y^{2}\right) \text { as } r=\sqrt{x^{2}+y^{2}}
\]

\begin{center}
\includegraphics[max width=\textwidth, alt={}]{9f3c5a8e-ac5d-42f4-8eb6-ea2af623f2d8-127_483_570_1392_821}
\end{center}

Thus,

\[
I_{z}=I_{x}+I_{y} \text { or, } I_{3}=I_{1}+I_{2}
\]

(b) Let us take the plane of the disc as \(x-y\) plane and origin to the centre of the disc (Fig.) From the symmetry \(I_{x}=I_{y}\). Let us consider a ring element of radius \(r\) and thickness \(d r\), then the moment of inertia of the ring element about the \(y\)-axis.

\[
d I_{z}=d m r^{2}=\frac{m}{\pi R^{2}}(2 \pi r d r) r^{2}
\]

Thus the moment of inertia of the disc about \(z\) - axis\\
\(I_{z}=\frac{2 m}{R^{2}} \int_{0}^{R} r^{3} d r=\frac{m R^{2}}{2}\)\\
But we have

Thus

\[
\begin{gathered}
I_{z}=I_{x}+I_{y}=2 I_{x} \\
I_{x}=\frac{I_{z}}{2}=\frac{m R^{2}}{4}
\end{gathered}
\]

\includegraphics[max width=\textwidth, alt={}, center]{9f3c5a8e-ac5d-42f4-8eb6-ea2af623f2d8-128_488_521_111_837}\\
1.241 For simplicity let us use mathematical trick. We consider the portion of the given disc as the superposition of two complete discs (without holes), one of positive density and radius \(R\) and other of negative density but of same magnitude and radius \(R / 2\).\\
As (area) \(\alpha\) (mass), the respective masses of the considered discs are \((4 m / 3)\) and \((-m / 3)\) respectively, and these masses can be imagined to be situated at their respective centers (C.M). Let us take point \(O\) as origin and point \(x\)-axis towards right. Obviously the C.M. of the shaded position of given shape lies on the \(x\)-axis. Hence the C.M. (C) of the shaded portion is given by\\
\(x_{c}=\frac{(-m / 3)(-R / 2)+(4 m / 3) 0}{(-m / 3)+4 m / 3}=\frac{R}{6}\)\\
Thus C.M. of the shape is at a distance \(R / 6\) from point \(O\) toward \(x\)-axis\\
Using parallel axis theorem and bearing in mind that the moment of inertia of a complete homogeneous disc of radius \(m_{0}\) and radius \(r_{0}\) equals \(\frac{1}{2} m_{0} r_{0}^{2}\). The moment of inetia of the\\
\includegraphics[max width=\textwidth, alt={}, center]{9f3c5a8e-ac5d-42f4-8eb6-ea2af623f2d8-128_312_448_1033_824}\\
small disc of mass ( \(-m / 3\) ) and radius \(R / 2\) about the axis passing through point \(C\) and perpendicular to the plane of the disc

\[
\begin{gathered}
I_{2 C}=\frac{1}{2}\left(-\frac{m}{3}\right)\left(\frac{R}{2}\right)^{2}+\left(-\frac{m}{3}\right)\left(\frac{R}{2}+\frac{R}{6}\right)^{2} \\
=-\frac{m R^{2}}{24}-\frac{4}{27} m R^{2}
\end{gathered}
\]

Similarly

\[
\begin{gathered}
I_{1 C}=\frac{1}{2}\left(\frac{4 m}{3}\right) R^{2}+\left(\frac{4 m}{3}\right)\left(\frac{R}{6}\right)^{2} \\
=\frac{2}{3} m R^{2}+\frac{m R^{2}}{27}
\end{gathered}
\]

Thus the sought moment of inertia,

\[
I_{C}=I_{1 C}+I_{2 C}=\frac{15}{24} m R^{2}-\frac{3}{27} m R^{2}=\frac{37}{72} m R^{2}
\]

1.242 Moment of inertia of the shaded portion, about the axis passing through it's certre,

\[
\begin{gathered}
I=\frac{2}{5}\left(\frac{4}{3} \pi R^{3} \rho\right) R^{2}-\frac{2}{5}\left(\frac{4}{3} \pi r^{3} \rho\right) r^{3} \\
=\frac{2}{5} \frac{4}{3} \pi \rho\left(R^{5}-r^{5}\right)
\end{gathered}
\]

Now, if \(R=r+d r\), the shaded portion becomes a shell, which is the required shape to calculate the moment of inertia.

Now,

\[
\begin{aligned}
I & =\frac{2}{5}-\frac{4}{3} \pi \rho\left\{(r+d r)^{5}-r^{5}\right\} \\
& =\frac{2}{5} \cdot \frac{4}{3} \pi \rho\left(r^{5}+5 r^{4} d r+\ldots \ldots-r^{5}\right)
\end{aligned}
\]

\begin{center}
\includegraphics[max width=\textwidth, alt={}]{9f3c5a8e-ac5d-42f4-8eb6-ea2af623f2d8-129_333_325_323_874}
\end{center}

Neglecting higher terms.

\[
=\frac{2}{3}\left(4 \pi r^{2} d r \rho\right) r^{2}=\frac{2}{3} m r^{2}
\]

1.243 (a) Net force which is effective on the system (cylinder \(M+\) body \(m\) ) is the weight of the body \(m\) in a uniform gravitational field, which is a constant. Thus the initial acceleration of the body \(m\) is also constant.

From the conservation of mechanical energy of the said system in the uniform field of gravity at time \(t=\Delta t: \Delta T+\Delta U=0\)\\
or

\[
\frac{1}{2} m v^{2}+\frac{1}{2} \frac{M R^{2}}{2} \omega^{2}-m g \Delta h=0
\]

or, \(\quad \frac{1}{4}(2 m+M) v^{2}-m g \Delta h=0\) [ as \(v=\omega R\) at all times]

But


\begin{equation*}
v^{2}=2 w \Delta h \tag{1}
\end{equation*}


Hence using it in Eq. (1), we get

\[
\frac{1}{4}(2 m+M) 2 w \Delta h-m g \Delta h=0 \text { or } w=\frac{2 m g}{(2 m+M)}
\]

From the kinematical relationship, \(\beta=\frac{w}{R}=\frac{2 m g}{(2 m+M) R}\)\\
Thus the sought angular velocity of the cylinder

\[
\omega(t)=\beta t=\frac{2 m g}{(2 m+M) R} t=\frac{g t}{(1+M / 2 m) R}
\]

(b) Sought kinetic energy.

\[
T(t)=\frac{1}{2} m v^{2}+\frac{1}{2} \frac{M l^{2}}{2} \omega^{2}=\frac{1}{4}(2 m+M) R^{2} \omega^{2}
\]

1.244 For equilibrium of the disc and axle

\[
2 T=m g \text { or } T=m g / 2
\]

As the disc unwinds, it has an angular acceleration \(\beta\) given by

\[
I \beta=2 T r \text { or } \beta=\frac{2 T r}{I}=\frac{m g r}{I}
\]

The corresponding linear acceleration is\\
\(r \beta=w=\frac{m g r^{2}}{I}\)\\
Since the disc remains stationary under the combined action of this acceleration and the acceleration ( \(-w\) ) of the bar which is transmitted to the axle, we must have \(w=\frac{m g r^{2}}{I}\)\\
\includegraphics[max width=\textwidth, alt={}, center]{9f3c5a8e-ac5d-42f4-8eb6-ea2af623f2d8-130_321_494_335_846}\\
1.245 Let the rod be deviated through an angle \(\varphi\) ''from its initial position at an arbitrary instant of time, measured relative to the initial position in the positive direction. From the equation of the increment of the mechanical energy of the system.\\
or,

\[
\begin{aligned}
\Delta T & =A_{e x t} \\
\frac{1}{2} I \omega^{2} & =\int N_{z} d \varphi
\end{aligned}
\]

or,

\[
\frac{1}{2} \frac{M l^{2}}{3} \omega^{2}=\int_{n}^{\varphi} F l \cos \varphi d \varphi=F l \sin \varphi
\]

Thus,

\[
\omega=\sqrt{\frac{6 F \sin \varphi}{M l}}
\]

1.246 First of all, let us sketch free body diagram of each body. Since the cylinder is rotating and massive, the tension will be different in both the sections of threads. From Newton's law in projection form for the bodies \(m_{1}\) and \(m_{2}\) and noting that \(w_{1}=w_{2}=w=\beta R\), (as no thread slipping), we have ( \(m_{1}>m_{2}\) )


\begin{equation*}
m_{1} g-T_{1}=m_{1} w=m_{1} \beta R \tag{1}
\end{equation*}


and \(T_{2}-m_{2} g=m_{2} w\)

Now from the equation of rotational dynamics of a solid about stationary axis of rotation. i.e. \(N_{z}=I \beta_{z}\), for the cylinder.\\
or, \(\left(T_{1}-T_{2}\right) R=I \beta=m R^{2} \beta / 2\)

Similtaneous solution of the above equations yields :\\
\(\beta=\frac{\left(m_{1}-m_{2}\right) g}{R\left(m_{1}+m_{2}+\frac{m}{2}\right)}\) and \(\frac{T_{1}}{T_{2}}=\frac{m_{1}\left(m+4 m_{2}\right)}{m_{2}\left(m+4 m_{1}\right)}\)\\
\includegraphics[max width=\textwidth, alt={}, center]{9f3c5a8e-ac5d-42f4-8eb6-ea2af623f2d8-130_565_355_1437_922}\\
1.247 As the system ( \(m+m_{1}+m_{2}\) ) is under constant forces, the acceleration of body \(m_{1}\) an \(m_{2}\) is constant. In addition to it the velocities and accelerations of bodies \(m_{1}\) and \(m_{2}\) at equal in magnitude (say \(v\) and \(w\) ) because the length of the thread is constant.\\
From the equation of increament of mechanical energy i.e. \(\Delta T+\Delta U=A_{f r}\), at time \(t\) whe block \(m_{1}\) is distance \(h\) below from initial position corresponding to \(t=0\),


\begin{equation*}
\frac{1}{2}\left(m_{1}+m_{2}\right) v^{2}+\frac{1}{2}\left(\frac{m R^{2}}{2}\right) \frac{v^{2}}{R^{2}}-m_{2} g h=-k m_{1} g h \tag{(1}
\end{equation*}


(as angular velocity \(\omega=\nu / R\) for no slipping of thread.)\\
But

\[
v^{2}=2 w h
\]

So using it in (1), we get


\begin{equation*}
w=\frac{2\left(m_{2}-k m_{1}\right) g}{m+2\left(m_{1}+m_{2}\right)} . \tag{(2.}
\end{equation*}


Thus the work done by the friction force on \(m_{1}\)

\[
\begin{gathered}
A_{f r}=-k m_{1} g h=-k m_{1} g\left(\frac{1}{2} w t^{2}\right) \\
=-\frac{k m_{1}\left(m_{1}-k m_{1}\right) g^{2} t^{2}}{m+2\left(m_{1}+m_{2}\right)} \quad \text { (using 2). }
\end{gathered}
\]

1.248 In the problem, the rigid body is in translation equlibrium but there is an angular retardation. We first sketch the free body diagram of the cylinder. Obviously the friction forces, acting on the cylinder, are kinetic. From the condition of translational equlibrium for the cylinder,

Hence,

\[
\begin{aligned}
& m g=N_{1}+k N_{2} ; \quad N_{2}=k N_{1} \\
& N_{1}=\frac{m g}{1+k^{2}} ; \quad N_{2}=k \frac{m g}{1+k^{2}}
\end{aligned}
\]

For pure rotation of the cylinder about its rotation axis, \(N_{z}=I \beta_{z}\)\\
or, \(\quad-k N_{1} R-k N_{2} R=\frac{m R^{2}}{2} \beta_{z}\)\\
or, \(\quad-\frac{k m g R(1+k)}{1+k^{2}}=\frac{m R^{2}}{2} \beta_{z}\)\\
or, \(\quad \beta_{z}=-\frac{2 k(1+k) g}{\left(1+k^{2}\right) R}\)\\
\includegraphics[max width=\textwidth, alt={}, center]{9f3c5a8e-ac5d-42f4-8eb6-ea2af623f2d8-131_388_500_1400_797}

Now, from the kinematical equation,

\[
\begin{gathered}
\omega^{2}=\omega_{0}^{2}+2 \beta_{z} \Delta \varphi \text { we have, } \\
\Delta \varphi=\frac{\omega_{0}^{2}\left(1+k^{2}\right) R}{4 k(1+k) g}, \quad \text { because } \omega=0
\end{gathered}
\]

Hence, the sought number of turns,

\[
n=\frac{\Delta \varphi}{2 \pi}=\frac{\omega_{0}^{2}\left(1+k^{2}\right) R}{8 \pi k(1+k) g}
\]

1.249 It is the moment of friction force which brings the disc to rest. The force of friction is applied to each section of the disc, and since these sections lie at different distances from the axis, the moments of the forces of friction differ from section to section.\\
To find \(N_{z}\), where \(z\) is the axis of rotation of the disc let us partition the disc into thin rings (Fig.). The force of friction acting on the considered element\\
\(d f r=k(2 \pi r d r \sigma) g\), (where \(\sigma\) is the density of the disc)\\
The moment of this force of friction is

\[
d N_{z}=-r d f r=-2 \pi k \sigma g r^{2} d r
\]

Integrating with respect to \(r\) from zero to \(R\), we get

\[
N_{z}=-2 \pi k \sigma g \int_{0}^{R} r^{2} d r=-\frac{2}{3} \pi k \sigma g R^{3}
\]

For the rotation of the disc about the stationary axis \(z\), from the equation \(N_{z}=I \beta_{z}\)

\[
-\frac{2}{3} \pi k \sigma g R^{3}=\frac{\left(\pi R^{2} \sigma\right) R^{2}}{2} \beta_{z} \text { or } \beta_{z}=-\frac{4 k g}{3 R}
\]

Thus from the angular kinematical equation \(\omega_{z}=\omega_{0 z}+\beta_{z} t\)\\
\includegraphics[max width=\textwidth, alt={}, center]{9f3c5a8e-ac5d-42f4-8eb6-ea2af623f2d8-132_372_510_723_815}

\[
0=\omega_{0}+\left(-\frac{4 k g}{3 R}\right) t \text { or } t=\frac{3 R \omega_{0}}{4 k g}
\]

1.250 According to the question,

Integrating,

\[
I \frac{d \omega}{d t}=-k \sqrt{\omega} \quad \text { or, } I=\frac{d \omega}{\sqrt{\omega}}=-k d t
\]

or, \(\omega=\frac{k^{2} t^{2}}{4 I^{2}}-\frac{\sqrt{\omega_{0}} k t}{I}+\omega_{0}\), (Noting that at \(t=0, \omega=\omega_{0}\).)\\
Let the flywheel stops at \(t=t_{0}\) then from Eq. (1), \(t_{0}=\frac{2 I \sqrt{\omega_{0}}}{k}\)\\
Hence sought average angular velocity

\[
<\omega>=\frac{\int_{0}^{\frac{2 \sqrt{\omega_{0}}}{k}}\left(\frac{k^{2} t^{2}}{4 I^{2}}-\frac{\sqrt{\omega_{0}} k t}{I}+\omega_{0}\right) d t}{\frac{2 I \sqrt{\omega_{0}}}{k}}=\frac{\omega_{0}}{3}
\]

1.251 Let us use the equation \(\frac{d M_{z}}{d t}=N_{z}\) relative to the axis through \(O\)

For this purpose, let us find the angular momentum of the system \(\boldsymbol{M}_{\boldsymbol{z}}\) about the given rotation axis and the corresponding torque \(N_{z}\). The angular momentum is

\[
M_{z}=I \omega+m v R=\left(\frac{m_{0}}{2}+m\right) R^{2} \omega
\]

[where \(I=\frac{m_{0}}{2} R^{2}\) and \(v=\omega R\) (no cord slipping)]\\
So, \(\quad \frac{d M_{z}}{d t}=\left(\frac{M R^{2}}{2}+m R^{2}\right) \beta_{z}\)

The downward pull of gravity on the overhanging part is the only external force, which exerts a torque about the \(z\)-axis, passing through \(O\) and is given by,

\[
\begin{gathered}
N_{z}=\left(\frac{m}{l}\right) x g R \\
\frac{d M_{z}}{d t}=N_{z}
\end{gathered}
\]

Hence from the equation

\[
\left(\frac{M R^{2}}{2}+m R^{2}\right) \beta_{z}=\frac{m}{l} x g R
\]

Thus,

\[
\beta_{z}=\frac{2 m g x}{l R(M+2 m)}>0
\]

Note : We may solve this problem using conservation of mechanical energy of the system (cylinder + thread) in the uniform field of gravity.\\
1.252 (a) Let us indicate the forces acting on the sphere and their points of application. Choose positive direction of \(x\) and \(\varphi\) (rotation angle) along the incline in downward direction and in the sense of \(\vec{\omega}\) (for undirectional rotation) respectively. Now from equations of dynamics of rigid body i.e. \(F_{x}=m w_{c x}\) and \(N_{c z}=I_{c} \beta_{z}\) we get :


\begin{equation*}
m g \sin \alpha-f_{r}=m w \tag{1}
\end{equation*}


and


\begin{equation*}
f r R=\frac{2}{5} m R^{2} \beta \tag{2}
\end{equation*}


But


\begin{equation*}
f r \leq k m g \cos \alpha \tag{3}
\end{equation*}


In addition, the absence of slipping provides the kinematical realtionship between the accelerations :


\begin{equation*}
w=\beta R \tag{4}
\end{equation*}


The simultaneous solution of all the four equations yields :

\[
k \cos \alpha \geq \frac{2}{7} \sin \alpha, \text { or } k \geq \frac{2}{7} \tan \alpha
\]

\begin{center}
\includegraphics[max width=\textwidth, alt={}]{9f3c5a8e-ac5d-42f4-8eb6-ea2af623f2d8-133_357_465_1592_804}
\end{center}

\footnotetext{(b) Solving Eqs. (1) and (2) [of part (a)], we get :
}
\[
w_{c}=\frac{5}{7} g \sin \alpha
\]

As the sphere starts at \(t=0\) along positive \(x\) axis, for pure rolling


\begin{equation*}
v_{c}(t)=w_{c} t=\frac{5}{7} \mathrm{~g} \sin \alpha t \tag{5}
\end{equation*}


Hence the sought kinetic energy

\[
\begin{aligned}
T & =\frac{1}{2} m v_{c}^{2}+\frac{1}{2} \frac{2}{5} m R^{2} \omega^{2}=\frac{7}{10} m v_{c}^{2}\left(\text { as } \omega=v_{c} / R\right) \\
& =\frac{7}{10} m\left(\frac{5}{7} \mathrm{~g} \sin \alpha t\right)^{2}=\frac{5}{14} m g^{2} \sin ^{2} \alpha t^{2}
\end{aligned}
\]

1.253 (a) Let us indicate the forces and their points of application for the cylinder. Choosing the positive direction for \(x\) and \(\varphi\) as shown in the figure, we write the equation of motion of the cylinder axis and the equation of moments in the C.M. frame relative to that axis i.e. from equation \(F_{x}=m w_{c}\) and \(N_{z}=I_{c} \beta_{z}\)

\[
m g-2 T=m w_{c} ; 2 T R=\frac{m R^{2}}{2} \beta
\]

As there is no slipping of thread on the cylinder

\[
w_{c}=\beta R
\]

From these three equations\\
\(T=\frac{m g}{6}=13 \mathrm{~N}, \beta=\frac{2}{5} \frac{g}{R}=5 \times 10^{2} \mathrm{rad} / \mathrm{s}^{2}\)\\
(b) we have \(\beta=\frac{2}{3} \frac{g}{R}\)

So, \(w_{c}=\frac{2}{3} \mathrm{~g}>0\) or, in vector form \(\vec{w}_{c}=\frac{2}{3} \overrightarrow{\mathrm{~g}}\)\\
\includegraphics[max width=\textwidth, alt={}, center]{9f3c5a8e-ac5d-42f4-8eb6-ea2af623f2d8-134_427_520_937_868}\\
\(P=\vec{F} \cdot \vec{v}=\vec{F} \cdot\left(\vec{w}_{c} t\right)\)\\
\(=m \vec{g} \cdot\left(\frac{2}{3} \vec{g} t\right)=\frac{2}{3} m g^{2} t\)\\
1.254 Let us depict the forces and their points of application corresponding to the cylinder attached with the elevator. Newton's second law for solid in vector form in the frame of elevator, gives :


\begin{equation*}
2 \vec{T}+m \vec{g}+m\left(-\vec{w}_{0}\right)=m \vec{w} \tag{1}
\end{equation*}


The equation of moment in the C.M. frame relative to the cylinder axis i.e. from \(N_{z}=I_{c} \beta_{z}-\)\\
\(2 T R=\frac{m R^{2}}{2} \beta=\frac{m R^{2}}{2} \frac{w^{\prime}}{R}\)\\[0pt]
[as thread does not slip on the cylinder, \(w^{\prime}=\beta R\) ]\\
\includegraphics[max width=\textwidth, alt={}, center]{9f3c5a8e-ac5d-42f4-8eb6-ea2af623f2d8-134_337_491_1607_899}\\
or,

\[
T=\frac{m w^{\prime}}{4}
\]

As (1) \(\vec{T} \uparrow \downarrow \vec{w}\)\\
so in vector form


\begin{equation*}
\vec{T}=-\frac{m \vec{w}}{4} \tag{2}
\end{equation*}


Solving Eqs. (1) and (2), \(\vec{w}^{\prime}=\frac{2}{3}\left(\vec{g}-\vec{w}_{0}\right)\) and sought force

\[
\vec{F}=2 \vec{T}=\frac{1}{3} m\left(\vec{g}-\overrightarrow{w_{0}}\right) .
\]

1.255 Let us depict the forces and their points of application for the spool. Choosing the positive direction for \(x\) and \(\varphi\) as shown in the fig., we apply \(F_{x}=m w_{c x}\) and \(N_{c z}=I_{c} \beta_{z}\) and get

\[
m g \sin \alpha-T=m w ; T r=I \beta
\]

"Notice that if a point of a solid in plane motion is connected with a thread the projection of velocity vector of the solid's point of contact along the length of the thread equals the velocity of the other end of the thread (if it is not slacked)"\\
Thus in our problem, \(\nu_{p}=\nu_{0}\) but \(\nu_{0}=0\), hence point \(P\) is the instantaneous centre of rotation of zero velocity for the spool. Therefore \(v_{c}=\omega r\) and subsequently \(w_{c}=\beta r\).\\
Solving the equations simultaneously, we get\\
\(w=\frac{g \sin \alpha}{1+\frac{I}{m r^{2}}}=1.6 \mathrm{~m} / \mathrm{s}^{2}\)\\
\includegraphics[max width=\textwidth, alt={}, center]{9f3c5a8e-ac5d-42f4-8eb6-ea2af623f2d8-135_627_603_696_813}\\
1.256 Let us sketch the force diagram for solid cylinder and apply Newton's second law in projection form along \(x\) and \(y\) axes (Fig.) :


\begin{equation*}
f r_{1}+f r_{2}=m w_{c} \tag{1}
\end{equation*}


and

\[
N_{1}+N_{2}-m g-F=0
\]

or


\begin{equation*}
N_{1}+N_{2}=m g+F \tag{2}
\end{equation*}


Now choosing positive direction of \(\varphi\) as shown in the figure and using \(N_{c z}=I_{c} \beta_{z}\)\\
we get


\begin{equation*}
F R-\left(f r_{1}+f r_{2}\right) R=\frac{m R^{2}}{2} \beta=\frac{m R^{2}}{2} \frac{w_{c}}{R} \tag{3}
\end{equation*}


[as for pure rolling \(w_{c}=\beta R\) ]. In addition to, \(f r_{1}+f r_{2} \leq k\left(N_{1}+N_{2}\right)\)\\
\includegraphics[max width=\textwidth, alt={}, center]{9f3c5a8e-ac5d-42f4-8eb6-ea2af623f2d8-135_550_585_1446_799}

Solving the Eqs., we get\\
and

\[
\begin{gathered}
F \leq \frac{3 k m g}{(2-3 k)}, \quad \text { or } F_{\max }=\frac{3 k m g}{2-3 k} \\
w_{c(\max )}=\frac{k\left(N_{1}+N_{2}\right)}{m} \\
=\frac{k}{m}\left[m g+F_{\max }\right]=\frac{k}{m}\left[m g+\frac{3 k m g}{2-3 k}\right]=\frac{2 k g}{2-3 k}
\end{gathered}
\]

1.257 (a) Let us choose the positive direction of the rotation angle \(\varphi\), such that \(w_{c x}\) and \(\beta_{z}\) have identical signs (Fig.). Equation of motion, \(F_{x}=m w_{c x}\) and \(N_{c z}=I_{c} \beta_{z}\) gives :

\[
F \cos \alpha-f r=m w_{c x}: f r R-F r=I_{c} \beta_{z}=\gamma m R^{2} \beta_{z}
\]

In the absence of the slipping of the spool \(w_{c x}=\beta_{z} R\)\\
From the three equations \(w_{c x}=w_{c}=\frac{F[\cos \alpha-(r / R)]}{m(1+\gamma)}\), where \(\cos \alpha>\frac{r}{R}\)\\
(b) As static friction ( \(f r\) ) does not work on the spool, from the equation of the increment of mechanical energy \(A_{\text {ext }}=\Delta T\).

\[
\begin{gathered}
A_{e x t}=\frac{1}{2} m v_{c}^{2}+\frac{1}{2} \gamma m R^{2} \frac{v_{c}^{2}}{R^{2}}=\frac{1}{2} m(1+\gamma) v_{c}^{2} \\
=\frac{1}{2} m(1+\gamma) 2 w_{c} x=\frac{1}{2} m(1+\gamma) 2 w_{c}\left(\frac{1}{2} w_{c} t^{2}\right) \\
=\frac{F^{2}\left(\cos \alpha-\frac{r}{R}\right)^{2} t^{2}}{2 m(1+\gamma)}
\end{gathered}
\]

\begin{center}
\includegraphics[max width=\textwidth, alt={}]{9f3c5a8e-ac5d-42f4-8eb6-ea2af623f2d8-136_561_523_695_849}
\end{center}

Note that at \(\cos \alpha=r / R\), there is no rolling and for \(\cos \alpha<r / R, w_{c x}<0\), i.e. the spool will move towards negative \(x\)-axis and rotate in anticlockwise sense.\\
1.258 For the cylinder from the equation \(N_{z}=I \beta_{z}\) about its stationary axis of rotation.


\begin{equation*}
2 T r=\frac{m r^{2}}{2} \beta \text { or } \beta=\frac{4 T}{m r} \tag{1}
\end{equation*}


For the rotation of the lower cylinder from the equation \(N_{c z}=I_{c} \beta_{z}\)

\[
2 T r=\frac{m r^{2}}{2} \beta^{\prime} \quad \text { or }, \quad \beta^{\prime}=\frac{4 T}{m r}=\beta
\]

Now for the translational motion of lower cylinder from the Eq. \(F_{x}=m w_{c x}\) :


\begin{equation*}
m g-2 T=m w_{c} \tag{2}
\end{equation*}


As there is no slipping of threads on the cylinders :


\begin{equation*}
w_{c}=\beta^{\prime} r+\beta r=2 \beta r \tag{3}
\end{equation*}


\begin{center}
\includegraphics[max width=\textwidth, alt={}]{9f3c5a8e-ac5d-42f4-8eb6-ea2af623f2d8-136_603_581_1417_797}
\end{center}

Simultaneous solution of (1), (2) and (3) yields

\[
T=\frac{m g}{10} .
\]

1.259 Let us depict the forces acting on the pulley and weight \(A\), and indicate positive direction for \(x\) and \(\varphi\) as shown in the figure. For the cylinder from the equation \(F_{x}=m\) \& and \(N_{c z}=I_{c} \beta_{z}\) we get


\begin{equation*}
M g+T_{A}-2 T=M w_{c} \tag{1}
\end{equation*}


and \(2 T R+T_{A}(2 R)=I \beta=\frac{I w_{c}}{R}\)

For the weight \(A\) from the equation


\begin{gather*}
F_{x}=m w_{x} \\
m g-T_{A}=m w_{A} \tag{3}
\end{gather*}


As there is no slipping of the threads on the pulleys.\\
\includegraphics[max width=\textwidth, alt={}, center]{9f3c5a8e-ac5d-42f4-8eb6-ea2af623f2d8-137_703_559_215_844}


\begin{equation*}
w_{A}=w_{c}+2 \beta R=w_{c}+2 w_{c}=3 w_{c} \tag{4}
\end{equation*}


Simultaneous solutions of above four equations gives :

\[
w_{A}=\frac{3(M+3 m) g}{\left(M+9 m+\frac{I}{R^{2}}\right)}
\]

1.260 (a) For the translational motion of the system \(\left(m_{1}+m_{2}\right)\), from the equation : \(F_{x}=m w_{c x}\)


\begin{equation*}
F=\left(m_{1}+m_{2}\right) w_{c} \text { or, } w_{c}=F /\left(m_{1}+m_{2}\right) \tag{1}
\end{equation*}


Now for the rotational motion of cylinder from the equation : \(N_{c x}=I_{c} \beta_{z}\)


\begin{equation*}
F r=\frac{m_{1} r^{2}}{2} \beta \quad \text { or } \quad \beta r=\frac{2 F}{m_{1}} \tag{2}
\end{equation*}


But

\[
w_{K}=w_{c}+\beta r, \text { So }
\]


\begin{equation*}
w_{K}=\frac{F}{m_{1}+m_{2}}+\frac{2 F}{m_{1}}=\frac{F\left(3 m_{1}+2 m_{2}\right)}{m_{1}\left(m_{1}+m_{2}\right)} \tag{3}
\end{equation*}


\includegraphics[max width=\textwidth, alt={}, center]{9f3c5a8e-ac5d-42f4-8eb6-ea2af623f2d8-137_289_454_1419_868}\\
(b) From the equation of increment of mechanical energy : \(\Delta T=A_{\text {ext }}\)\\
Here

\[
\Delta T=T(t), \text { so, } T(t)=A_{e x t}
\]

As force \(F\) is constant and is directed along \(x\)-axis the sought work done.

\[
A_{e x t}=F x
\]

(where \(x\) is the displacement of the point of application of the force \(F\) during time interval \(t\) )

\[
=F\left(\frac{1}{2} w_{K} t^{2}\right)=\frac{F^{2} t^{2}\left(3 m_{1}+2 m_{2}\right)}{2 m_{1}\left(m_{1}+m_{2}\right)}=T(t)
\]

(using Eq. (3)\\
Alternate : \(T(t)=T_{\text {translation }}(t)+T_{\text {rotation }}(t)\)

\[
=\frac{1}{2}\left(m_{1}+m_{2}\right)\left(\frac{F t}{\left(m_{1}+m_{2}\right)}\right)^{2}+\frac{1}{2} \frac{m_{1} r^{2}}{2}\left(\frac{2 F t}{m_{1} r}\right)^{2}=\frac{F^{2} t^{2}\left(3 m_{1}+2 m_{2}\right)}{2 m_{1}\left(m_{1+} m_{2}\right)}
\]

1.261 Choosing the positive direction for \(x\) and \(\varphi\) as shown in Fig, let us we write the equation of motion for the sphere \(F_{x}=m w_{c x}\) and \(N_{c z}=I_{c} \beta_{z}\)\\
\(f r=m_{2} w_{2} ; f r r=\frac{2}{5} m_{2} r^{2} \beta\)\\
( \(w_{2}\) is the acceleration of the C.M. of sphere.)\\
For the plank from the Eq. \(F_{x}=m w_{x}\)\\
\(F-f_{r}=m_{1} w_{1}\)\\
In addition, the condition for the absence of slipping of the sphere yields the kinematical relation between the accelerations :

\[
w_{1}=w_{2}+\beta r
\]

\begin{center}
\includegraphics[max width=\textwidth, alt={}]{9f3c5a8e-ac5d-42f4-8eb6-ea2af623f2d8-138_341_420_687_908}
\end{center}

Simultaneous solution of the four equations yields :

\[
w_{1}=\frac{F}{\left(m_{1}+\frac{2}{7} m_{2}\right)} \text { and } w_{2}=\frac{2}{7} w_{1}
\]

1.262 (a) Let us depict the forces acting on the cylinder and their point of applications for the cylinder and indicate positive direction of \(x\) and \(\varphi\) as shown in the figure. From the equations for the plane motion of a solid \(F_{x}=m w_{c x}\) and \(N_{c z}=I_{c} \beta_{z}\) :


\begin{gather*}
k m g=m w_{c x} \text { or } w_{c x}=k g  \tag{1}\\
-k m g R=\frac{m R^{2}}{2} \beta^{z} \quad \text { or } \quad \beta_{z}=-2 \frac{k g}{R} \tag{2}
\end{gather*}


Let the cylinder starts pure rolling at \(t=t_{0}\) after releasing on the horizontal floor at \(t=0\).\\
From the angular kinematical equation


\begin{align*}
& \omega_{z}=\omega_{o z}+\beta_{z} t, \\
& \text { or } \quad \omega=\omega_{0}-2 \frac{k g}{R} t \tag{3}
\end{align*}


From the equation of the linear kinematics,\\
\includegraphics[max width=\textwidth, alt={}, center]{9f3c5a8e-ac5d-42f4-8eb6-ea2af623f2d8-138_353_449_1605_874}

\[
v_{c x}=v_{o c x}+w_{c x} t
\]

or


\begin{equation*}
v_{c}=0+k g t_{0} \tag{4}
\end{equation*}


But at the moment \(t=t_{0}\), when pure rolling starts \(v_{c}=\omega R\)\\
so,

\[
k g t_{0}=\left(\omega_{0}-2 \frac{k g}{R} t_{0}\right) R
\]

Thus

\[
t_{0}=\frac{\omega_{0} R}{3 k g}
\]

(b) As the cylinder pick, up speed till it starts rolling, the point of contact has a purely translatory movement equal to \(\frac{1}{2} w_{c} t_{0}^{2}\) in the forward directions but there is also a backward movement of the point of contact of magnitude \(\left(\omega_{0} \tau_{0}-\frac{1}{2} \beta t_{0}^{2}\right) R\). Because of slipping the net displacement is backwards. The total work done is then,

\[
\begin{aligned}
& A_{\mathrm{fr}}=k m g\left[\frac{1}{2} w_{c} t_{0}^{2}-\left(\omega_{0} t_{0}+\frac{1}{2} \beta t_{0}^{2}\right) R\right] \\
= & k m g\left[\frac{1}{2} k g t_{0}^{2}-\frac{1}{2}\left(-\frac{2 k g}{R}\right) t_{0}^{2} R-\omega_{0} t_{0} R\right] \\
= & k m g \frac{\omega_{0} R}{3 k g}\left[\frac{\omega_{0} R}{6}+\frac{\omega_{0} R}{3}-\omega_{0} R\right]=-\frac{m \omega_{0}^{2} R^{2}}{6}
\end{aligned}
\]

The same result can also be obtained by the work-energy theorem, \(A_{f r}=\Delta T\).\\
1.263 Let us write the equation of motion for the centre of the sphere at the moment of breaking-off:

\[
m v^{2} /(R+r)=m g \cos \theta,
\]

where \(\nu\) is the velocity of the centre of the sphere at that moment, and \(\theta\) is the corresponding angle (Fig.). The velocity \(v\) can be found from the energy conservation law :

\[
m g h=\frac{1}{2} m v^{2}+\frac{1}{2} I \omega^{2},
\]

where \(I\) is the moment of inertia of the sphere relative to the axis passing through the sphere's centre. i.e. \(I=\frac{2}{5} m r^{2}\). In addition,

\[
v=\omega r ; h=(R+r)(1-\cos \theta) .
\]

From these four equations we obtain

\[
\omega=\sqrt{10 g(R+r) 17 r^{2}} .
\]

\includegraphics[max width=\textwidth, alt={}, center]{9f3c5a8e-ac5d-42f4-8eb6-ea2af623f2d8-139_441_463_1250_835}\\
1.264 Since the cylinder moves without sliding, the centre of the cylinder rotates about the point \(O\), while passing through the common edge of the planes. In other words, the point \(O\) becomes the foot of the instantaneous axis of rotation of the cylinder.\\
It at any instant during this motion the velocity of the C.M. is \(\boldsymbol{v}_{1}\) when the angle (shown in the figure) is \(\beta\), we have

\[
\frac{m v_{1}^{2}}{R}=m g \cos \beta-N,
\]

where \(N\) is the normal reaction of the edge\\
or, \(\quad v_{1}^{2}=g R \cos \beta-\frac{N R}{m}\)

From the energy conservation law,

\[
\frac{1}{2} I_{0} \frac{v_{1}^{2}}{R^{2}}-\frac{1}{2} I_{0} \frac{v_{0}^{2}}{R^{2}}=m g R(1-\cos \beta)
\]

But \(\quad I_{0}=\frac{m R^{2}}{2}+m R^{2}=\frac{3}{2} m R^{2}\),\\
(from the parallel axis theorem)\\
\includegraphics[max width=\textwidth, alt={}, center]{9f3c5a8e-ac5d-42f4-8eb6-ea2af623f2d8-140_390_519_135_812}

Thus,


\begin{equation*}
v_{1}^{2}=v_{0}^{2}+\frac{4}{3} g R(1-\cos \beta) \tag{2}
\end{equation*}


From (1) and (2)

\[
v_{0}^{2}=\frac{g R}{3}(7 \cos \beta-4)-\frac{N R}{m}
\]

The angle \(\beta\) in this equation is clearly smaller than or equal to \(\alpha\) so putting \(\beta=\alpha\) we get

\[
v_{0}^{2}=\frac{g R}{3}(7 \cos \alpha-4)-\frac{N_{0} R}{M}
\]

where \(N_{0}\) is the corresponding reaction. Note that \(N \geq N_{0}\). No jumping occurs during this turning if \(N_{0}>0\). Hence, \(\nu_{0}\) must be less than

\[
v_{\max }=\sqrt{\frac{g R}{3}(7 \cos \alpha-4)}
\]

1.265 Clearly the tendency of bouncing of the hoop will be maximum when the small body \(A\), will be at the highest point of the hoop during its rolling motion. Let the velocity of C.M. of the hoop equal \(v\) at this position. The static friction does no work on the hoop, so from conservation of mechanical energy; \(E_{1}=E_{2}\)\\
or, \(\quad 0+\frac{1}{2} m v_{0}^{2}+\frac{1}{2} m R^{2}\left(\frac{v_{0}}{R}\right)^{2}-m g R=\frac{1}{2} m(2 v)^{2}+\frac{1}{2} m v^{2}+\frac{1}{2} m R^{2}\left(\frac{v}{R}\right)^{2}+m g R\)\\
or,


\begin{equation*}
3 v^{2}=v_{0}^{2}-2 g R \tag{1}
\end{equation*}


From the equation \(F_{n}=m w_{n}\) for body \(A\) at final position 2 :


\begin{equation*}
m g+N^{\prime}=m \omega^{2} R=m\left(\frac{v}{R}\right)^{2} R \tag{2}
\end{equation*}


\includegraphics[max width=\textwidth, alt={}, center]{9f3c5a8e-ac5d-42f4-8eb6-ea2af623f2d8-140_380_380_1619_284}\\
\includegraphics[max width=\textwidth, alt={}, center]{9f3c5a8e-ac5d-42f4-8eb6-ea2af623f2d8-140_450_389_1567_711}

As the hoop has no acceleration in vertical direction, so for the hoop,


\begin{equation*}
N+N^{\prime}=m g \tag{3}
\end{equation*}


From Eqs. (2) and (3),


\begin{equation*}
N=2 m g-\frac{m v^{2}}{R} \tag{4}
\end{equation*}


As the hoop does not bounce, \(N \geq 0\)

So from Eqs. (1), (4) and (5),

Hence


\begin{equation*}
\frac{8 g R-v_{0}^{2}}{3 R} \geq 0 \text { or } 8 g R \geq v_{0}^{2} \tag{5}
\end{equation*}


\[
v_{0} \leq \sqrt{8 g R}
\]

1.266 Since the lower part of the belt is in contact with the rigid floor, velocity of this part becomes zero. The crawler moves with velocity \(v\), hence the velocity of upper part of the belt becomes \(2 \nu\) by the rolling condition and kinetic energy of upper part \(=\frac{1}{2}\left(\frac{m}{2}\right)(2 v)^{2}=m v^{2}\), which is also the sought kinetic energy, assuming that the length of the belt is much larger than the radius of the wheels.\\
1.267 The sphere has two types of motion, one is the rotation about its own axis and the other is motion in a circle of radius \(R\). Hence the sought kinetic energy


\begin{equation*}
T=\frac{1}{2} I_{1} \omega_{1}^{2}+\frac{1}{2} I_{2} \omega_{2}^{2} \tag{1}
\end{equation*}


where \(I_{1}\) is the moment of inertia about its own axis, and \(I_{2}\) is the moment of inertia about the vertical axis, passing through \(O\),\\
But, \(I_{1}=\frac{2}{5} m r^{2}\) and \(I_{2}=\frac{2}{5} m r^{2}+m R^{2}\) (using parallel axis theorem,)

In addition to


\begin{equation*}
\omega_{1}=\frac{v}{r} \text { and } \omega_{2}=\frac{v}{R} \tag{3}
\end{equation*}


Using (2) and (3) in (1), we get \(T^{\prime}=\frac{7}{10} m v^{2}\left(1+\frac{2 r^{2}}{7 R^{2}}\right)\)\\
1.268 For a point mass of mass \(d m\), looked at from \(C\) rotating frame, the equation is

\[
d m \vec{w}=\vec{f}+d m \omega^{2} \vec{r}+2 d m\left(\vec{v}^{\prime} \times \vec{\omega}\right)
\]

where \(\vec{r}=\) radius vector in the rotating frame with respect to rotation axis and \(\overrightarrow{v^{\prime}}=\) velocity in the same frame. The total centrifugal force is clearly

\[
\vec{F}_{c f}=\sum d m \omega^{2} \vec{r}=m \omega^{2} \vec{R}_{c}
\]

\(\vec{R}_{c}\) is the radius vector of the C.M. of the body with respect to rotation axis, also

\[
\vec{F}_{c o r}=2 m \vec{v}_{c}^{\prime} \times \vec{\omega}
\]

where we have used the definitions

\[
m \overrightarrow{R_{c}}=\sum d m \vec{r} \text { and } m \overrightarrow{v_{c}}=\sum d m \overrightarrow{v^{\prime}}
\]

1.269 Consider a small element of length \(d x\) at a distance \(x\) from the point \(C\), which is rotating in a circle of radius \(r=x \sin \theta\)\\
Now, mass of the element \(=\left(\frac{m}{l}\right) d x\)\\
So, centrifugal force acting on this element \(=\left(\frac{m}{l}\right) d x \omega^{2} x \sin \theta\) and moment of this force about \(C\),

\[
\begin{aligned}
|d N| & =\left(\frac{m}{l}\right) d x \omega^{2} x \sin \theta \cdot x \cos \theta \\
& =\frac{m \omega^{2}}{2 l} \sin 2 \theta x^{2} d x
\end{aligned}
\]

and hence, total moment\\
l/ 2\\
\includegraphics[max width=\textwidth, alt={}, center]{9f3c5a8e-ac5d-42f4-8eb6-ea2af623f2d8-142_557_375_198_896}\\
\(N=2 \int_{0} \frac{m \omega^{2}}{2 l} \sin 2 \theta x^{2} d x=\frac{1}{24} m \omega^{2} l^{2} \sin 2 \theta\),\\
1.270 Let us consider the system in a frame rotating with the rod. In this frame, the rod is at rest and experiences not only the gravitational force \(m \vec{g}\) and the reaction force \(\vec{R}\), but also the centrifugal force \(\vec{F}_{c f}\).\\
In the considered frame, from the condition of equilibrium i.e. \(N_{0 z}=0\)\\
or,


\begin{equation*}
N_{c f}=m g \frac{l}{2} \sin \theta \tag{1}
\end{equation*}


where \(N_{c f}\) is the moment of centrifugal force about \(O\). To calculate \(N_{c f}\), let us consider an element of length \(d x\), situated at a distance \(x\) from the point \(O\). This element is subjected to a horizontal pseudo force \(\left(\frac{m}{l}\right) d x \omega^{2} x \sin \theta\). The moment of this pseudo force about the axis of rotation through the point \(O\) is\\
\includegraphics[max width=\textwidth, alt={}, center]{9f3c5a8e-ac5d-42f4-8eb6-ea2af623f2d8-142_413_513_1128_875}

\[
\begin{aligned}
d N_{c f} & =\left(\frac{m}{l}\right) d x \omega^{2} x \sin \theta x \cos \theta \\
& =\frac{m \omega^{2}}{l} \sin \theta \cos \theta x^{2} d x
\end{aligned}
\]

So


\begin{equation*}
N_{c f}=\int_{0}^{l} \frac{m \omega^{2}}{l} \sin \theta \cos \theta x^{2} d x=\frac{m \omega^{2} l^{2}}{3} \sin \theta \cos \theta \tag{2}
\end{equation*}


It follows from Eqs. (1) and (2) that,


\begin{equation*}
\cos \theta=\left(\frac{3 g}{2 \omega^{2} l}\right) \text { or } \theta=\cos ^{-1}\left(\frac{3 g}{2 \omega^{2} l}\right) \tag{3}
\end{equation*}


1.271 When the cube is given an initial velocity on the table in some direction (as shown) it acquires an angular momentum about an axis on the table perpendicular to the initial velocity and (say) just below the C.G.. This angular momentum will disappear when the cube stops and this can only by due to a torque. Frictional forces cannot do this by themselves because they act in the plain containing the axis. But if the force of normal reaction act eccentrically (as shown), their torque can bring about the vanishing of the angular momentum. We can calculate the distance \(\Delta x\) between the point of application of the normal reaction and the C.G. of the cube as follows. Take the moment about C.G. of all the forces. This must vanish because the cube does not turn or turnble on the table. Then if the force of friction is \(f r\)\\
\includegraphics[max width=\textwidth, alt={}, center]{9f3c5a8e-ac5d-42f4-8eb6-ea2af623f2d8-143_613_582_218_820}

\[
\text { fr } \frac{a}{2}=N \Delta x
\]

But \(N=m g\) and \(f r=k m g\), so

\[
\Delta x=k a / 2
\]

1.272 In the process of motion of the given system the kinetic energy and the angular momentum relative to rotation axis do not vary. Hence, it follows that

\[
\frac{1}{2} \frac{M l^{2}}{3} \omega_{0}^{2}=\frac{1}{2} m\left(\omega^{2} l^{2}+v^{\prime 2}\right)+\frac{1}{2} \frac{M l^{2}}{3} \omega^{2}
\]

( \(\omega\) is the final angular velocity of the rod)\\
and

\[
\frac{M l^{2}}{3} \omega_{0}=\frac{M l^{2}}{3} \omega+m l^{2} \omega
\]

From these equations we obtain

\[
\begin{aligned}
\omega & =\omega_{0} /\left(1+\frac{3 M}{M}\right) \text { and } \\
v^{\prime} & =\omega_{0} l / \sqrt{1+3 m / M}
\end{aligned}
\]

1.273 Due to hitting of the ball, the angular impulse received by the rod about the C.M. is equal to \(p \frac{1}{2}\). If \(\omega\) is the angular velocity acquired by the rod, we have


\begin{equation*}
\frac{m l^{2}}{12} \omega=\frac{p l}{2} \text { or } \omega=\frac{6 p}{m l} \tag{1}
\end{equation*}


In the frame of C.M., the rod is rotating about an axis passing through its mid point with the angular velocity \(\omega\). Hence the force exerted by one half on the other \(=\) mass of one half \(\times\) acceleration of C.M. of that part, in the frame of C.M.

\[
=\frac{m}{2}\left(\omega^{2} \frac{l}{4}\right)=m \frac{\omega^{2} l}{8}=\frac{9 p^{2}}{2 m l}=9 \mathrm{~N}
\]

1.274 (a) In the process of motion of the given system the kinetic energy and the angular momentum relative to rotation axis do not vary. Hence it follows that\\
and

\[
\begin{gathered}
\frac{1}{2} m v^{2}=\frac{1}{2} m v^{\prime 2}+\frac{1}{2}\left(\frac{M l^{2}}{3}\right) \omega^{2} \\
m v \frac{l}{2}=m v^{\prime} \frac{l}{2}+\frac{M l^{2}}{3} \omega
\end{gathered}
\]

From these equations we obtain

\[
v^{\prime}=\left(\frac{3 m-4 M}{3 m+4 M}\right) v . \quad \text { and } \quad \omega=\frac{4 v}{l(1+4 m / 3 M)}
\]

As \(\vec{v}^{\prime} \uparrow \uparrow \vec{v}\), so in vector form \(\vec{v}^{\prime}=\left(\frac{3 m-4 M}{3 m+4 M}\right) \vec{v}\)\\
(b) Obviously the sought force provides the centripetal acceleration to the C.M. of the rod and is

\[
\begin{gathered}
F_{n}=m w_{c n} \\
=M \omega^{2} \frac{l}{2}=\frac{8 M v^{2}}{l(1+4 M / 3 m)^{2}}
\end{gathered}
\]

1.275 (a) About the axis of rotation of the rod, the angular momentum of the system is conserved. Thus if the velocity of the flying bullet is \(v\).


\begin{gather*}
m v l=\left(m l^{2}+\frac{M l^{2}}{3}\right) \omega \\
\omega=\frac{m v}{\left(m+\frac{M}{3}\right) l} \approx \frac{3 m v}{M l} \text { as } m \ll M \tag{1}
\end{gather*}


Now from the conservation of mechanical energy of the system (rod with bullet) in the uniform field of gravity


\begin{equation*}
\frac{1}{2}\left(m l^{2}+\frac{M l^{2}}{3}\right) \omega^{2}=(M+m) g \frac{l}{2}(1-\cos \alpha) \tag{2}
\end{equation*}


[because C.M. of rod raises by the height \(\frac{l}{2}(1-\cos \alpha)\) ]\\
Solving (1) and (2), we get

\[
v=\left(\frac{M}{m}\right) \sqrt{\frac{2}{3} g l} \sin \frac{\alpha}{2} \text { and } \omega=\sqrt{\frac{6 g}{l}} \sin \frac{\alpha}{2}
\]

(b) Sought \(\Delta p=\left[m(\omega l)+M\left(\omega \frac{l}{2}\right)\right]-m v\)\\
where \(\omega l\) is the veloccity of the bullet and \(\omega \frac{l}{2}\) equals the velocity of C.M. of the rod after the impact. Putting the value of \(v\) and \(\omega\) we get

\[
\Delta p \sim \frac{1}{2} m v=M \sqrt{\frac{g l}{6}} \sin \frac{\alpha}{2}
\]

This is caused by the reaction at the hinge on the upper end.\\
(c) Let the rod starts swinging with angular velocity \(\omega^{\prime}\), in this case. Then, like part (a)

\[
m v x=\left(\frac{M l^{2}}{3}+m x^{2}\right) \omega^{\prime} \text { or } \omega^{\prime} \approx \frac{3 m v x}{M l^{2}}
\]

Final momentum is

\[
p_{f}=m x \omega^{\prime}+\int_{0}^{l} y \omega^{\prime} \frac{M}{l} d y \approx \frac{M}{2} \omega^{\prime} l \sim \frac{3}{2} m v \frac{x}{l}
\]

So,

\[
\Delta p=p_{f}-p_{i} \sim m v\left(\frac{3 x}{2 l}-1\right)
\]

This vanishes for

\[
x \approx \frac{2}{3} l
\]

1.276 (a) As force \(F\) on the body is radial so its angular momentum about the axis becomes zero and the angular momentum of the system about the given axis is conserved. Thus

\[
\frac{M R^{2}}{2} \omega_{0}+m \omega_{0} R^{2}=\frac{M R^{2}}{2} \omega \text { or } \omega=\omega_{0}\left(1+\frac{2 m}{M}\right)
\]

(b) From the equation of the increment of the mechanical energy of the system :

\[
\begin{gathered}
\Delta T=A_{e x t} \\
\frac{1}{2} \frac{M R^{2}}{2} \omega^{2}-\frac{1}{2}\left(\frac{M R^{2}}{2}+m R^{2}\right) \omega_{0}^{2}=A_{e x t}
\end{gathered}
\]

Putting the value of \(\omega\) from part (a) and solving we get

\[
A_{e x t}=\frac{m \omega_{0}^{2} R^{2}}{2}\left(1+\frac{2 m}{M}\right)
\]

1.277 (a) Let \(z\) be the rotation axis of disc and \(\varphi\) be its rotation angle in accordance with right-hand screw rule (Fig.). ( \(\varphi\) and \(\varphi^{\prime}\) are to be measured in the same sense algebraically.) As \(M_{z}\) of the system (disc + man) is conserved and \(M_{z(\text { initial })}=0\), we have at any instant,\\
or,

\[
0=-\frac{m_{2} R^{2}}{2} \frac{d \varphi}{d t}+m_{1}\left[\left(\frac{d \varphi^{\prime}}{d t}\right) R+\left(\frac{d \varphi}{d t}\right) R\right] R
\]

On integrating

\[
d \varphi=\left[-\frac{m_{1}}{m_{1}+\left(m_{2} / 2\right)}\right] d \varphi^{\prime}
\]

or,

\[
\int_{0}^{\varphi} d \varphi=-\int_{0}^{\varphi^{\prime}}\left(\frac{m_{1}}{m_{1}+\left(m_{2} / 2\right)}\right) d \varphi^{\prime}
\]


\begin{equation*}
\varphi=-\left(\frac{m_{1}}{m_{1}+\frac{m_{2}}{2}}\right) \varphi^{\prime} \tag{1}
\end{equation*}


This gives the total angle of rotation of the disc.\\
(b) From Eq. (1)

\[
\frac{d \varphi}{d t}=-\left(\frac{m_{1}}{m_{1}+\frac{m_{2}}{2}}\right) \frac{d \varphi^{\prime}}{d t}=-\left(\frac{m_{1}}{m_{1}+\frac{m_{2}}{2}}\right) \frac{v^{\prime}(t)}{R}
\]

Differentiating with respect to time

\[
\frac{d^{2} \varphi}{d t^{2}}=-\left(\frac{m_{1}}{m_{1}+\frac{m_{2}}{2}}\right) \frac{1}{R} \frac{d v^{\prime}(t)}{d t}
\]

Thus the sought force moment from the Eq. \(N_{z}=I \beta_{z}\)

\[
\begin{aligned}
& N_{z}=\frac{m_{2} R^{2}}{2} \frac{d^{2} \varphi}{d t^{2}}=-\frac{m_{2} R^{2}}{2}\left(\frac{m_{1}}{m_{1}+\frac{m_{2}}{2}}\right) \frac{1}{R} \frac{d v^{\prime}(t)}{d t} \\
& \text { Hence } \\
& \qquad N_{z}=-\frac{m_{1} m_{2} R}{2 m_{1}+m_{2}} \frac{d v^{\prime}(t)}{d t}
\end{aligned}
\]

\includegraphics[max width=\textwidth, alt={}, center]{9f3c5a8e-ac5d-42f4-8eb6-ea2af623f2d8-146_400_424_289_925}\\
1.278 (a) Frome the law of conservation of angular momentum of the system relative to vertical axis \(z\), it follows that:

Hence


\begin{gather*}
I_{1} \omega_{1 z}+I_{2} \omega_{2 z}=\left(I_{1}+I_{2}\right) \omega_{z} \\
\omega_{z}=\left(I_{1} \omega_{1 z}+I_{2} \omega_{2 z}\right) /\left(I_{1}+I_{2}\right) \tag{1}
\end{gather*}


Not that for \(\omega_{z}>0\), the corresponding vector \(\vec{\omega}\) coincides with the poitive direction to the \(z\) axis, and vice versa. As both discs rotates about the same vertical axis \(z\), thus in vector form.

\[
\vec{\omega}=I_{1} \overrightarrow{\omega_{1}}+I_{2} \overrightarrow{\omega_{2}} /\left(I_{1}+I_{2}\right)
\]

However, the problem makes sense only if \(\vec{\omega}_{1} \uparrow \uparrow \vec{\omega}_{2}\) or \(\vec{\omega}_{1} \uparrow \downarrow \vec{\omega}_{2}\)\\
(b) From the equation of increment of mechanical energy of a system: \(A_{f r}=\Delta T\).

\[
=\frac{1}{2}\left(I_{1}+I_{2}\right) \omega_{z}^{2}-\frac{1}{2} I_{1} \omega_{1 z}^{2}+\frac{1}{2} I_{2} \omega_{2 z}^{2}
\]

Using Eq. (1)

\[
A_{f r}=-\frac{I_{1} I_{2}}{2\left(I_{1}+I_{2}\right)}\left(\omega_{1 z}-\omega_{2 z}\right)^{2}
\]

1.279 For the closed system (disc + rod), the angular momentum is conserved about any axis. Thus from the conservation of angular momentum of the system about the rotation axis of rod passing through its C.M. gives :


\begin{equation*}
m v \frac{l}{2}=m v^{\prime} \frac{l}{2}+\frac{\eta m l^{2}}{12} \omega \tag{1}
\end{equation*}


( \(\nu^{\prime}\) is the final velocity of the disc and \(\omega\) angular velocity of the rod) For the closed system linear momentum is also conserved. Hence


\begin{equation*}
m v=m v^{\prime}+\eta m v_{c} \tag{2}
\end{equation*}


(where \(\boldsymbol{v}_{\boldsymbol{c}}\) is the velocity of C.M. of the rod)\\
From Eqs (1) and (2) we get

\[
v_{c}=\frac{l \omega}{3} \text { and } \quad v-v^{\prime}=\eta v_{c}
\]

Applying conservation of kinetic energy, as the collision is elastic


\begin{equation*}
\frac{1}{2} m v^{2}=\frac{1}{2} m v^{\prime 2}+\frac{1}{2} \eta m v_{c}^{2}+\frac{1}{2} \frac{\eta m l^{2}}{12} \omega^{2} \tag{3}
\end{equation*}


or

\[
v^{2}-v^{\prime 2}=4 \eta v_{c}^{2} \text { and hence } v+v^{\prime}=4 v_{c}
\]

Then

\[
v^{\prime}=\frac{4-\eta}{4+\eta} v \text { and } \omega=\frac{12 v}{(4+\eta) l}
\]

Vectorially, noting that we have taken \(\vec{v}^{\prime}\) parallel to \(\vec{v}\)

\[
\vec{u}^{\prime}=\left(\frac{4-\eta}{4+\eta}\right) \vec{v}
\]

So, \(\vec{u}^{\prime}=0\) for \(\eta=4\) and \(\vec{u}^{\prime} \downarrow \uparrow \vec{v}\) for \(\eta>4\)\\
\includegraphics[max width=\textwidth, alt={}, center]{9f3c5a8e-ac5d-42f4-8eb6-ea2af623f2d8-147_418_224_622_984}\\
1.280 See the diagram in the book (Fig. 1.72)\\
(a) When the shaft \(B B^{\prime}\) is turned through \(90^{\circ}\) the platform must start turning with angular velocity \(\Omega\) so that the angular momentum remains constant. Here

\[
\left(I+I_{0}\right) \Omega=I_{0} \omega_{0} \text { or, } \Omega=\frac{I_{0} \omega_{0}}{I+I_{0}}
\]

The work performed by the motor is therefore

\[
\frac{1}{2}\left(I+I_{0}\right) \Omega^{2}=\frac{1}{2} \frac{I_{0}^{2} \omega_{0}^{2}}{I+I_{0}}
\]

If the shaft is turned through \(180^{\circ}\), angular velocity of the sphere changes sign. Thus from conservation of angular momentum,

\[
I \Omega-I_{0} \omega_{0}=I_{0} \omega_{0}
\]

(Here \(-I_{0} \omega_{0}\) is the complete angular momentum of the sphere i. e. we assume that the angular velocity of the sphere is just \(-\omega_{0}\) ). Then

\[
\Omega=2 I_{0} \frac{\omega_{0}}{I}
\]

and the work done must be,

\[
\frac{1}{2} I \Omega^{2}+\frac{1}{2} I_{0} \omega_{0}^{2}-\frac{1}{2} I_{0} \omega_{0}^{2}=\frac{2 I_{0}^{2} \omega_{0}^{2}}{I}
\]

(b) In the case (a), first part, the angular momentum vector of the sphere is precessing with angular velocity \(\Omega\). Thus a torque,

\[
I_{0} \omega_{0} \Omega=\frac{I_{0}^{2} \omega_{0}^{2}}{I+I_{0}} \text { is needed. }
\]

1.281 The total centrifugal force can be calculated by,

\[
\int_{0}^{l_{0}} \frac{m}{l_{0}} \omega^{2} x d x=\frac{1}{2} m l_{0} \omega^{2}
\]

Then for equilibrium,

\[
\left(T_{2}-T_{1}\right) \frac{l}{2}=m g \frac{l_{0}}{2}
\]

and, \(T_{2}+T_{1}=\frac{1}{2} m l_{0} \omega^{2}\)\\
Thus \(T_{1}\) vanishes, when

\[
\omega^{2}=\frac{2 g}{l}, \omega=\sqrt{\frac{2 g}{l}}=6 \mathrm{rad} / \mathrm{s}
\]

Then \(T_{2}=m g \frac{l_{0}}{l}=25 \mathrm{~N}\)\\
\includegraphics[max width=\textwidth, alt={}, center]{9f3c5a8e-ac5d-42f4-8eb6-ea2af623f2d8-148_514_659_484_698}\\
1.282 See the diagram in the book (Fig. 1.71).\\
(a) The angular velocity \(\vec{\omega}\) about \(O O^{\prime}\) can be resolved into a component parallel to the rod and a component \(\omega \sin \theta\) perpendicular to the rod through \(C\). The component parallel to the rod does not contribute so the angular momentum

\[
M=I \omega \sin \theta=\frac{1}{12} m l^{2} \omega \sin \theta
\]

Also, \(\quad M_{z}=M \sin \theta=\frac{1}{12} m l^{2} \omega \sin ^{2} \theta\)\\
This can be obtained directly also,\\
(b) The modulus of \(\vec{M}\) does not change but the modulus of the change of \(\vec{M}\) is \(|\Delta \vec{M}|\).\\
\(|\Delta \vec{M}|=2 M \sin (90-\theta)=\frac{1}{12} m l^{2} \omega \sin 2 \theta\)\\
(c) Here \(M_{\perp}=M \cos \theta=I \omega \sin \theta \cos \theta\)\\
\includegraphics[max width=\textwidth, alt={}, center]{9f3c5a8e-ac5d-42f4-8eb6-ea2af623f2d8-148_534_325_1335_953}

Now \(\left|\frac{d \vec{M}}{d t}\right|=I \omega \sin \theta \cos \theta \frac{\omega d t}{d t}=\frac{1}{24} m l^{2} \omega^{2} \sin ^{2} \theta\)\\
as \(\vec{M}\) precesses with angular velocity \(\omega\).\\
1.283 Here \(M=I \omega\) is along the symmetry axis. It has two components, the part \(I \omega \cos \theta\) is constant and the part \(M_{\perp}=I \omega \sin \theta\) presesses, then

\[
\left|\frac{d \vec{M}}{d t}\right|=I \omega \sin \theta \omega^{\prime}=m g l \sin \theta
\]

or,

\[
\omega^{\prime}=\text { precession frequency }=\frac{m g l}{I \omega}=0.7 \mathrm{rad} / \mathrm{s}
\]

(b) This force is the centripetal force due to precession. It acts inward and has the magnitude

\[
|\vec{F}|=\left|\sum m_{i} \omega^{\prime 2} \overrightarrow{\rho_{i}}\right|=m \omega^{\prime 2} l \sin \theta=12 \mathrm{mN} .
\]

\(\vec{\rho}_{i}\) is the distance of the \(i\) th element from the axis. This is the force that the table will exert on the top. See the diagram in the answer sheet\\
\includegraphics[max width=\textwidth, alt={}, center]{9f3c5a8e-ac5d-42f4-8eb6-ea2af623f2d8-149_501_451_819_234}\\
\includegraphics[max width=\textwidth, alt={}, center]{9f3c5a8e-ac5d-42f4-8eb6-ea2af623f2d8-149_510_347_797_871}\\
1.284 See the diagram in the book (Fig. 1.73).

The moment of inertia of the disc about its symmentry axis is \(\frac{1}{2} m R^{2}\). If the angular velocity of the disc is \(\omega\) then the angular momentum is \(\frac{1}{2} m R^{2} \omega\). The precession frequency being \(2 \pi n\),\\
we have

\[
\left|\frac{\overrightarrow{d M}}{d t}\right|=\frac{1}{2} m R^{2} \omega \times 2 \pi n
\]

This must equal \(m(g+w) l\), the effective gravitational torques ( \(g\) being replaced by \(g+w\) in the elevator). Thus,

\[
\omega=\frac{(g+w) l}{\pi R^{2} n}=300 \mathrm{rad} / \mathrm{s}
\]

1.285 The effective \(g\) is \(\sqrt{g^{2}+w^{2}}\) inclined at angle \(\tan ^{-1} \frac{w}{g}\) with the vertical. Then with reference to the new " vertical" we proceed as in problem 1.283. Thus

\[
\omega^{\prime}=\frac{m l \sqrt{g^{2}+w^{2}}}{I \omega}=0.8 \mathrm{rad} / \mathrm{s} .
\]

The vector \(\vec{\omega}^{\prime}\) forms an angle \(\theta=\tan ^{-1} \frac{w}{g}=6^{\circ}\) with the normal vertical.\\
1.286 The moment of inertia of the sphere is \(\frac{2}{5} m R^{2}\) and hence the value of angular momentum is \(\frac{2}{5} m R^{2} \omega\). Since it precesses at speed \(\omega^{\prime}\) the torque required is

\[
\frac{2}{5} m R^{2} \omega \omega^{\prime}=F^{\prime} l
\]

So,

\[
F^{\prime}=\frac{2}{5} m R^{2} \omega \omega^{\prime} / l=300 \mathrm{~N}
\]

(The force \(F^{\prime}\) must be vertical.)\\
1.287 The moment of inertia is \(\frac{1}{2} m r^{2}\) and angular momentum is \(\frac{1}{2} m r^{2} \omega\). The axle oscillates about a horizontal axis making an instantaneous angle.

\[
\varphi=\varphi_{m} \sin \frac{2 \pi t}{T}
\]

This means that there is a variable precession with a rate of precession \(\frac{d \varphi}{d t}\). The maximum value of this is \(\frac{2 \pi \varphi_{m}}{T}\). When the angle between the axle and the axis is at its maximum value, a torque \(I \omega \Omega\)

\[
=\frac{1}{2} m r^{2} \omega \frac{2 \pi \varphi_{m}}{T}=\frac{\pi m r^{2} \omega \varphi_{m}}{T} \text { acts on it. }
\]

The corresponding gyroscopic force will be \(\frac{\pi m r^{2} \omega \varphi_{m}}{l T}=90 \mathrm{~N}\)\\
1.288 The revolutions per minute of the flywheel being \(n\), the angular momentum of the flywheel is \(l \times 2 \pi n\). The rate of precession is \(\frac{v}{R}\)\\
Thus \(N=2 \pi I N V / R=5.97 \mathrm{kN} . \mathrm{m}\).\\
1.289 As in the previous problem a couple \(2 \pi I n v / R\) must come in play. This can be done if a force, \(\frac{2 \pi I n v}{R l}\) acts on the rails in opposite directions in addition to the centrifugal and other forces. The force on the outer rail is increased and that on the inner rail decreased.\\
The additional force in this case has the magnitude 1.4 kN . m.

\subsection*{1.6 ELASTIC DEFORMATIONS OF A SOLID BODY}
1.290 Variation of length with temperature is given by


\begin{equation*}
l_{t}=l_{0}(1+\alpha \Delta t) \text { or } \frac{\Delta l}{l_{0}}=\alpha \Delta t=\varepsilon \tag{1}
\end{equation*}


But

\[
\varepsilon=\frac{\sigma}{E},
\]

Thus \(\sigma=\alpha \Delta t E\), which is the sought stress of pressure.\\
Putting the value of \(\alpha\) and \(E\) from Appendix and taking \(\Delta t=100^{\circ} \mathrm{C}\), we get

\[
\sigma=2.2 \times 10^{3} \mathrm{~atm} .
\]

1.291 (a) Consider a transverse section of the tube and concentrate on an element which subtends an angle \(\Delta \varphi\) at the centre. The forces acting on a portion of length \(\Delta l\) on the element are\\
(1) tensile forces side ways of magnitude \(\sigma \Delta r \Delta l\).

The resultant of these is

\[
2 \sigma \Delta r \Delta l \sin \frac{\Delta \varphi}{2} \approx \sigma \Delta r \Delta l \Delta \varphi
\]

radially towards the cente.\\
(2) The force due to fluid pressure \(=p r \Delta \varphi \Delta l\)

Since these balance, we get \(p_{\text {max }} \approx \sigma_{m} \frac{\Delta r}{r}\)\\
where \(\sigma_{m}\) is the maximum tensile force.\\
\includegraphics[max width=\textwidth, alt={}, center]{9f3c5a8e-ac5d-42f4-8eb6-ea2af623f2d8-151_346_312_730_1004}

Putting the values we get \(p_{\text {max }}=19.7\) atmos.\\
(b) Consider an element of area \(d S=\pi(r \Delta \theta / 2)^{2}\) about \(z\)-axis chosen arbitrarily. There are tangential tensile forces all around the ring of the cap. Their resultant is

\[
\sigma\left[2 \pi\left(r \frac{\Delta \theta}{2}\right) \Delta r\right] \sin \frac{\Delta \theta}{2}
\]

Hence in the limit

\[
p_{m} \pi\left(\frac{r \Delta \theta}{2}\right)^{2}=\sigma_{m} \pi\left(\frac{r \Delta \theta}{2}\right) \Delta r \Delta \theta
\]

or \(\quad p_{m}=\frac{2 \sigma_{m} \Delta r}{r}=39.5\) atmos.\\
\includegraphics[max width=\textwidth, alt={}, center]{9f3c5a8e-ac5d-42f4-8eb6-ea2af623f2d8-151_315_317_1307_965}\\
1.292 Let us consider an element of rod at a distance \(x\) from its rotation axis (Fig.). From Newton's second law in projection form directed towards the rotation axis

\[
-d T=(d m) \omega^{2} x=\frac{m}{l} \omega^{2} x d x
\]

On integrating

\[
-T=\frac{m \omega^{2}}{l} \frac{x^{2}}{2}+\mathrm{C}(\text { constant })
\]

But at

\[
x= \pm \frac{l}{2} \text { or free end, } T=0
\]

Thus

\[
0=\frac{m \omega^{2}}{2} \frac{l^{2}}{4}+C \text { or } C=-\frac{m \omega^{2} l}{8}
\]

Hence

\[
T=\frac{m \omega^{2}}{2}\left(\frac{l}{4}-\frac{x^{2}}{l}\right)
\]

Thus

\[
T_{\max }=\frac{m \omega^{2} l}{8} \text { (at mid point) }
\]

Condition required for the problem is\\
\includegraphics[max width=\textwidth, alt={}, center]{9f3c5a8e-ac5d-42f4-8eb6-ea2af623f2d8-152_420_446_302_893}

So, \(\frac{m \omega^{2} l}{8}=S \sigma_{m}\) or \(\omega=\frac{2}{l} \sqrt{\frac{2 \sigma_{m}}{\rho}}\)

\[
T_{\max }=S \sigma_{m}
\]

Hence the sought number of \(r p s\)

\[
n=\frac{\omega}{2 \pi}=\frac{1}{\pi l} \sqrt{\frac{2 \sigma_{m}}{\rho}} \quad\left[\text { using the table } n=0.8 \times 10^{2} \mathrm{rps}\right]
\]

1.293 Let us consider an element of the ring (Fig.). From Newton's law \(F_{n}=m w_{n}\) for this element, we get,

\[
T d \theta=\left(\frac{m}{2 \pi} d \theta\right) \omega^{2} r \quad \text { [see solution of 1.93or 1.92] }
\]

So,

\[
T=\frac{m}{2 \pi} \omega^{2} r
\]

Condition for the problem is :

\[
\frac{T}{\pi r^{2}} \leq \sigma_{m} \text { or, } \frac{m \omega^{2} r}{2 \pi^{2} r^{2}} \leq \sigma_{m}
\]

or, \(\quad \omega_{\max }^{2}=\frac{2 \pi^{2} \sigma_{m} r}{\pi r^{2}(2 \pi r \rho)}=\frac{\sigma_{m}}{\rho r^{2}}\)\\
Thus sought number of rps\\
\includegraphics[max width=\textwidth, alt={}, center]{9f3c5a8e-ac5d-42f4-8eb6-ea2af623f2d8-152_331_327_1191_961}

\[
n=\frac{\omega_{\max }}{2 \pi}=\frac{1}{2 \pi r} \sqrt{\frac{\sigma_{m}}{\rho}}
\]

Using the table of appendices \(n=23 \mathrm{rps}\)\\
1.294 Let the point \(O\) desend by the distance \(x\) (Fig.). From the condition of equilibrium of point \(O\).


\begin{equation*}
2 T \sin \theta=m g \text { or } T=\frac{m g}{2 \sin \theta}=\frac{m g}{2 x} \sqrt{(l / 2)^{2}+x^{2}} \tag{1}
\end{equation*}


Now,


\begin{equation*}
\frac{T}{\pi(d / 2)^{2}}=\sigma=\varepsilon E \text { or } T=\varepsilon E \pi \frac{d^{2}}{4} \tag{2}
\end{equation*}


( \(\sigma\) here is stress and \(\varepsilon\) is strain.)

In addition to it,


\begin{equation*}
\varepsilon=\frac{\sqrt{(l / 2)^{2}+x^{2}}-\frac{l}{2}}{l / 2}=\sqrt{1+\left(\frac{2 x}{l}\right)^{2}}-1 \tag{3}
\end{equation*}


From Eqs. (1), (2) and (3)\\
\(x-\frac{x}{\sqrt{1+\left(\frac{2 x}{l}\right)^{2}}}=\frac{m g l}{\pi E d^{2}}\) as \(x \ll l\)\\
So, \(\frac{4 x^{3}}{2 l^{2}} \approx \frac{m g l}{\pi E d^{2}}\)\\
\includegraphics[max width=\textwidth, alt={}, center]{9f3c5a8e-ac5d-42f4-8eb6-ea2af623f2d8-153_470_548_228_811}\\
or, \(x=l\left(\frac{m g}{2 \pi E d^{2}}\right)^{1 / 3}=2.5 \mathrm{~cm}\)\\
1.295 Let us consider an element of the rod at a distance \(x\) from the free end (Fig.). For the considered element ' \(T-T\) ' are internal restoring forces which produce elongation and \(d T\) provides the acceleration to the element. For the element from Newton's law :

\[
d T=(d m) w=\left(\frac{m}{l} d x\right) \frac{F_{o}}{m}=\frac{F_{o}}{l} d x
\]

As free end has zero tension, on integrating the above expression,

\[
\int_{o}^{T} d T=\frac{F_{o}}{l} \int_{o}^{x} d x \quad \text { or } \quad T=\frac{F_{o}}{l} x
\]

Elongation in the considered element of lenght \(d x\) :

\[
\partial \xi=\frac{\sigma}{E}(x) \quad d x=\frac{T}{S E} d x=\frac{F_{x o} x d x}{S E L}
\]

Thus total elengation \(\xi=\frac{F_{o}}{S E l} \int_{x}^{l} x d x=\frac{F_{o} l}{2 S E}\)\\
Hence the sought strain\\
\includegraphics[max width=\textwidth, alt={}, center]{9f3c5a8e-ac5d-42f4-8eb6-ea2af623f2d8-153_143_577_1483_832}

\[
\sigma=\frac{\xi}{l}=\frac{F_{o}}{2 S E}
\]

1.296 Let us consider an element of the rod at a distance \(r\) from it's rotation axis. As the element rotates in a horizontal circle of radius \(r\), we have from Newton's second law in projection form directed toward the axis of rotation :\\
or,

\[
\begin{gathered}
T-(T+d T)=(d m) \omega^{2} r \\
-d T=\left(\frac{m}{l} d r\right) \omega^{2} r=\frac{m}{l} \omega^{2} r d r
\end{gathered}
\]

At the free end tension becomes zero. Integrating the above experession we get, thus

\[
\begin{gathered}
-\int_{T}^{0} d T=\frac{m}{l} \omega^{2} \int_{r}^{l} r d r \\
T=\frac{m \omega^{2}}{l}\left(\frac{l^{2}-r^{2}}{2}\right)=\frac{m \omega^{2} l}{2}\left(1-\frac{r^{2}}{l^{2}}\right)
\end{gathered}
\]

Thus\\
Elongation in elemental length \(d r\) is given by :

\[
\partial \xi=\frac{\sigma(r)}{E} d r=\frac{T}{S E} d r
\]

(where \(S\) is the cross sectional area of the rod and \(T\) is the tension in the rod at the considered element)\\
or, \(\quad \partial \xi=\frac{m \omega^{2} l}{2 S E}\left(1-\frac{r^{2}}{l^{2}}\right) d r\)\\
Thus the sought elongation

\[
\xi=\int d \xi=\frac{m \omega^{2} l}{2 S E} \int_{0}^{l}\left(1-\frac{r^{2}}{l^{2}}\right) d r
\]

\includegraphics[max width=\textwidth, alt={}, center]{9f3c5a8e-ac5d-42f4-8eb6-ea2af623f2d8-154_297_395_651_885}\\
or, \(\quad \xi=\frac{m \omega^{2} l}{2 S E} \frac{2 l}{3}=\frac{(S l \rho)}{3 S E} \omega^{2} l^{3}\)

\[
=\frac{1}{3} \frac{\rho \omega^{2} l^{3}}{E} \text { (where } \rho \text { is the density of the copper.) }
\]

1.297 Volume of a solid cylinder

\[
V=\pi r^{2} l
\]

So, \(\quad \frac{\Delta V}{V}=\frac{\pi 2 r \Delta r l}{\pi r^{2} l}+\frac{\pi r^{2} \Delta l}{\pi r^{2} l}=\frac{2 \Delta r}{r}+\frac{\Delta l}{l}\)

But longitudinal strain \(\Delta l / l\) and accompanying lateral strain \(\Delta r / r\) are related as


\begin{equation*}
\frac{\Delta r}{r}=-\mu \frac{\Delta l}{l} \tag{2}
\end{equation*}


Using (2) in (1), we get :

But


\begin{gather*}
\frac{\Delta V}{V}=\frac{\Delta l}{l}(1-2 \mu)  \tag{3}\\
\frac{\Delta l}{l}=\frac{-F / \pi r^{2}}{E}
\end{gather*}


(Because the increment in the length of cylinder \(\Delta l\) is negative)\\
So,

\[
\frac{\Delta V}{V}=\frac{-F}{\pi r^{2} E}\left(1-2^{\prime} \mu\right)
\]

Thus,

\[
\Delta V=\frac{-F l}{E}(1-2 \mu)
\]

Negative sign means that the volume of the cylinder has decreased.\\
1.298 (a) As free end has zero tension, thus the tension in the rod at a vestical distance \(y\) from its lower end


\begin{equation*}
T=\frac{m}{l} g y \tag{1}
\end{equation*}


Let \(\partial l\) be the elongation of the element of length \(d y\), then

\[
\partial l=\frac{\sigma(y)}{E} d y
\]

\[
=\frac{T}{S E} d y=\frac{m g y d y}{S l E}=\rho g y d y / E \text { (where } \rho \text { is the density of the copper) }
\]

Thus the sought elongation


\begin{equation*}
\Delta l=\int \partial l=\rho g \int_{0}^{l} \frac{y d y}{E}=\frac{1}{2} \rho g l^{2} / E \tag{2}
\end{equation*}


(b) If the longitudinal (tensile) strain is \(\varepsilon=\frac{\Delta l}{l}\), the accompanying lateral (compressive) strain is given by


\begin{equation*}
\varepsilon^{\prime}=\frac{\Delta r}{r}=-\mu \varepsilon \tag{3}
\end{equation*}


Then since \(V=\pi r^{2} l\) we have

\[
\begin{gathered}
\frac{\Delta V}{V}=\frac{2 \Delta r}{r}+\frac{\Delta l}{l} \\
=(1-2 \mu) \frac{\Delta l}{l} \quad[\text { Using (3) }]
\end{gathered}
\]

where \(\frac{\Delta l}{l}\) is given in part (a), \(\mu\) is the Poisson ratio for copper.\\
1.299 Consider a cube of unit length before pressure is applied. The pressure acts on each face. The pressures on the opposite faces constitute a tensile stress producing longitudianl compression and lateral extension. The compressions is \(\frac{p}{E}\) and the lateral extension is \(\mu \frac{p}{E}\) The net result is a compression

\[
\frac{p}{E}(1-2 \mu) \text { in each side. }
\]

Hence \(\frac{\Delta V}{V}=-\frac{3 p}{E}(1-2 \mu)\) because from symmetry \(\frac{\Delta V}{V}=3 \frac{\Delta l}{l}\)\\
(b) Let us consider a cube under an equal compressive stress \(\sigma\), acting on all its faces.

Then,


\begin{equation*}
\text { volume strain }=-\frac{\Delta V}{V}=\frac{\sigma}{k}, \tag{1}
\end{equation*}


where \(k\) is the bulk modulus of elasticity.\\
So

\[
\frac{\sigma}{k}=\frac{3 \sigma}{E}(1-2 \mu)
\]

or,

\[
E=3 k(1-2 \mu)=\frac{3}{\beta}(1-2 \mu)\left(\text { as } k=\frac{1}{\beta}\right)
\]

\(\mu \leq \frac{1}{2}\) if \(E\) and \(\beta\) are both to remain positive.

A beam clamped at one end and supporting an applied load at the free end is called a cantilever. The theory of cantilevers is discussed in advanced text book on mechanics. The key result is that elastic forces in the beam generate a couple, whose moment, called the moment of resistances, balances the external bending moment due to weight of the beam, load etc. The moment of resistance, also called internal bending moment (I.B.M) is given by

\[
\text { I.B.M. }=\mathrm{EI} / \mathrm{R}
\]

Here \(R\) is the radius of curvature of the beam at the representative point \((x, y) . I\) is called the geometrical moment of inertia

\[
I=\int z^{2} d S
\]

of the cross section relative to the axis passing through the netural layer which remains unstretched. (Fig.1.). The section of the beam beyond \(P\) exerts the bending moment \(N(x)\) and we have,

\[
\frac{E I}{R}=N(x)
\]

If there is no load other than that due to the weight of the beam, then

\[
N(x)=\frac{1}{2} \rho g(l-x)^{2} b h
\]

where \(\rho=\) density of steel.\\
Hence, at

\[
\begin{gathered}
x=0 \\
\left(\frac{I}{R}\right)_{0}=\frac{\rho g l^{2} b h}{2 E I}
\end{gathered}
\]

Here \(b=\) width of the beam perpendicular to paper.\\
Also, \(\quad I=\int_{-h / 2}^{h / 2} z^{2} b d z=\frac{b h^{3}}{12}\).\\
Hence, \(\left(\frac{1}{R}\right)_{0}=\frac{6 \rho g l^{2}}{E h^{2}}=(0.121 \mathrm{~km})^{-1}\)\\
\includegraphics[max width=\textwidth, alt={}, center]{9f3c5a8e-ac5d-42f4-8eb6-ea2af623f2d8-156_319_524_1330_829}\\
\includegraphics[max width=\textwidth, alt={}, center]{9f3c5a8e-ac5d-42f4-8eb6-ea2af623f2d8-156_444_469_1652_877}\\
1.301 We use the equation given above and use the result that when \(y\) is small

\[
\frac{1}{R} \approx \frac{d^{2} y}{d x^{2}} . \text { Thus, } \frac{d^{2} y}{d x^{2}}=\frac{N(x)}{E I}
\]

(a) Here \(N(x)=N_{0}\) is a constant. Then integration gives,

\[
\frac{d y}{d x}=\frac{N_{0} x}{E I}+C_{1}
\]

But \(\quad\left(\frac{d y}{d x}\right)=0\) for \(x=0\), so \(C_{1}=0\). Integrating again,

\[
y=\frac{N_{0} x^{2}}{2 E I}
\]

where we have used \(y=0\) for \(x=0\) to set the constant of integration at zero. This is the equation of a parabola. The sag of the free end is

\[
\lambda=y(x=l)=\frac{N_{0} l^{2}}{2 E I}
\]

(b) In this case \(N(x)=F(l-x)\) because the load \(F\) at the extremity is balanced by a similar force at \(F\) directed upward and they constitute a couple. Then

\[
\frac{d^{2} y}{d x^{2}}=\frac{F(l-x)}{E I} .
\]

Integrating,

\[
\frac{d y}{d x}=\frac{F\left(l x-x^{2} / 2\right)}{E I}+C_{1}
\]

As before \(C_{1}=0\). Integrating again, using \(y=0\) for \(x=0\)

\[
y=\frac{F\left(\frac{l x^{2}}{2}-\frac{x^{3}}{6}\right)}{E I} \text { here } \lambda=\frac{F l^{3}}{3 E I}
\]

Here for a square cross section

\[
I=\int_{-a / 2}^{a / 2} z^{2} a d z=a^{4} / 12
\]

1.302 One can think of it as analogous to the previous case but with a beam of length \(l / 2\) loaded upward by a force \(F / 2\).

Thus

\[
\lambda=\frac{F l^{3}}{48 E I},
\]

\begin{center}
\includegraphics[max width=\textwidth, alt={}]{9f3c5a8e-ac5d-42f4-8eb6-ea2af623f2d8-157_226_569_1634_794}
\end{center}

On using the last result of the previous problem.\\
1.303 (a) In this case \(N(x)=\frac{1}{2} \rho g b h(l-x)^{2}\) where \(b=\) width of the girder.

Also

\[
I=b h^{3} / 12 \text {. Then, }
\]

\[
\frac{E b h^{2}}{12} \frac{d^{2} y}{d x^{2}}=\frac{\rho g b h}{2}\left(l^{2}-2 l x+x^{2}\right) .
\]

Integrating,

\[
\frac{d y}{d x}=\frac{6 \rho g}{E h^{2}}\left(l^{2} x-l x^{2}+\frac{x^{3}}{3}\right)
\]

using

\[
\frac{d y}{d x}=0 \text { for } x=0 . \text { Again integrating }
\]

\[
y=\frac{6 \rho g}{E h^{2}}\left(\frac{l^{2} x^{2}}{2}-\frac{l x^{3}}{3}+\frac{x^{4}}{12}\right)
\]

Thus

\[
\begin{aligned}
\lambda & =\frac{6 \rho g l^{4}}{E h^{2}}\left(\frac{1}{2}-\frac{1}{3}+\frac{1}{12}\right) \\
& =\frac{6 \rho g l^{4}}{E h^{2}} \frac{3}{12}=\frac{3 \rho g l^{4}}{2 E h^{2}}
\end{aligned}
\]

(b) As before, \(E I \frac{d^{2} y}{d x^{2}}=N(x)\) where \(N(x)\) is the bending moment due to section \(P B\).

This bending moment is clearly

\[
\begin{aligned}
N & =\int_{x}^{2 l} w d \xi(\xi-x)-w l(2 l-x) \\
& =w\left(2 l^{2}-2 x l+\frac{x^{2}}{2}\right)-w l(2 l-x)=w\left(\frac{x^{2}}{2}-x l\right)
\end{aligned}
\]

(Here \(\boldsymbol{w}=\boldsymbol{\rho} \boldsymbol{g} \boldsymbol{b} \boldsymbol{h}\) is weight of the beam per unit length)\\
Now integrating, \(E I \frac{d y}{d x}=w\left(\frac{x^{3}}{6}-\frac{x^{2} l}{2}\right)+c_{0}\)\\
or since \(\frac{d y}{d x}=0\) for \(x=l, c_{0}=w l^{3} / 3\)\\
\includegraphics[max width=\textwidth, alt={}, center]{9f3c5a8e-ac5d-42f4-8eb6-ea2af623f2d8-158_211_595_1280_790}

Integrating again, \(E I y=w\left(\frac{x^{4}}{24}-\frac{x^{3} l}{6}\right)+\frac{w l^{3} x}{3}+c_{1}\)\\
As \(y=0\) for \(x=0, c_{1}=0\). From this we find

\[
\lambda=y(x=l)=\frac{5 w l^{4}}{24} / E I=\frac{5 \rho g l^{4}}{2 E h^{2}}
\]

1.304 The deflection of the plate can be noticed by going to a co-rotating frame. In this frame each element of the plate experiences a pseudo force proportional to its mass. These forces have a moment which constitutes the bending moment of the problem. To calculate this moment we note that the acceleration of an element at a distance \(\xi\) from the axis is \(a=\xi \beta\) and the moment of the forces exerted by the section between \(x\) and \(l\) is

\[
N=\rho \ln \beta \int_{x}^{l} \xi^{2} d \xi=\frac{1}{3} \rho \ln \beta\left(l^{3}-x^{3}\right) .
\]

From the fundamental equation

\[
E I \frac{d^{2} y}{d x^{2}}=\frac{1}{3} \rho \ln \beta\left(l^{3}-x^{3}\right) .
\]

The moment of inertia \(I=\int_{-h / 2}^{+h / 2} z^{2} l d z=\frac{l h^{3}}{12 .}\)\\
Note that the neutral surface (i.e. the surface which contains lines which are neither stretched nor compressed) is a vertical plane here and \(z\) is perpendicular to it.

\[
\begin{gathered}
\frac{d^{2} y}{d x^{2}}=\frac{4 \rho \beta}{E h^{2}}\left(l^{3}-x^{3}\right) . \text { Integrating } \\
\frac{d y}{d x}=\frac{4 \rho \beta}{E h^{2}}\left(l^{3} x-\frac{x^{4}}{4}\right)+c_{1}
\end{gathered}
\]

Since

\[
\begin{array}{r}
\frac{d y}{d x}=0, \text { for } x=0, c_{1}=0 . \text { Integrating } \\
y=\frac{4 \rho \beta}{E h^{2}}\left(\frac{l^{3} x^{2}}{2}-\frac{x^{5}}{20}\right)+c_{2} \\
c_{2}=0 \text { because } y=0 \text { for } x=0 \\
\lambda=y(x=l)=\frac{9 \rho \beta l^{5}}{5 E h^{2}}
\end{array}
\]

Thus\\
1.305 (a) Consider a hollow cylinder of length \(l\), outer radius \(r+\Delta r\) inner radius \(r\), fixed at one end and twisted at the other by means of a couple of moment \(N\). The angular displacement \(\varphi\), at a distance \(l\) from the fixed end, is proportional to both \(l\) and \(N\). Consider an element of length \(d x\) at the twisted end. It is moved by an angle \(\varphi\) as shown. A vertical section is also shown and the twisting of the parallelopipe of length \(l\) and area \(\Delta r d x\) under the action of the twisting couple can be discussed by elementary means. If \(f\) is the tangential force generated then shearing stress is \(f / \Delta r d x\) and this must equal

\[
G \theta=G \frac{r \varphi}{l}, \text { since } \theta=\frac{r \varphi}{l} .
\]

Hence,

\[
f=G \Delta r d x \frac{r \varphi}{l} .
\]

The force \(f\) has moment \(f r\) about the axis and so the total moment is

\[
N=G \Delta r \frac{\varphi}{l} r^{2} \int d x=\frac{2 \pi r^{3} \Delta r \varphi}{l} G
\]

(b) For a solid cylinder we must integrate over \(r\). Thus\\
\includegraphics[max width=\textwidth, alt={}, center]{9f3c5a8e-ac5d-42f4-8eb6-ea2af623f2d8-160_682_1108_140_80}\\
1.306 Clearly \(N=\int_{d_{1} / 2} \frac{2 \pi r^{3} d r \varphi G}{l}=\frac{\pi}{32 l} G \varphi\left(d_{2}^{4}-d_{1}^{4}\right)\)\\
using

\[
\begin{gathered}
G=81 \mathrm{G} \mathrm{~Pa}=8.1 \times 10^{10} \frac{\mathrm{~N}}{\mathrm{~m}^{2}} \\
d_{2}=5 \times 10^{-2} \mathrm{~m}, d_{1}=3 \times 10^{-2} \mathrm{~m} \\
\varphi=2.0^{\circ}=\frac{\pi}{90} \text { radians, } l=3 \mathrm{~m} \\
N=\frac{\pi \times 8.1 \times \pi}{32 \times 3 \times 90}(625-81) \times 10^{2} \mathrm{~N} \cdot \mathrm{~m} \\
=0.5033 \times 10^{3} \mathrm{~N} \cdot \mathrm{~m} \approx 0.5 \mathrm{k} \mathrm{~N} \cdot \mathrm{~m}
\end{gathered}
\]

1.307 The maximum power that can be transmitted by means of a shaft rotating about its axis is clearly \(N \omega\) where \(N\) is the moment of the couple producing the maximum permissible torsion, \(\varphi\). Thus

\[
P=\frac{\pi r^{4} G \varphi}{2 l} \cdot \omega=16.9 \mathrm{kw}
\]

1.308 Consider an elementary ring of width \(d r\) at a distant \(r\) from the axis. The part outside exerts a couple \(N+\frac{d N}{d r} d r\) on this ring while the part inside exerts a couple \(N\) in the opposite direction. We have for equilibrium

\[
\frac{d N}{d r} d r=-d I \beta
\]

where \(d I\) is the moment of inertia of the elementary ring, \(\beta\) is the angular acceleration and minus sign is needed because the couple \(N(r)\) decreases, with distance vanshing at the outer radius, \(N\left(r_{2}\right)=0\). Now

\[
d I=\frac{m}{\pi\left(r_{2}^{2}-r_{1}^{2}\right)} 2 \pi r d r r^{2}
\]

Thus

\[
d N=\frac{2 m \beta}{\left(r_{1}^{2}-r_{1}^{2}\right)} r^{3} d r
\]

or, \(\quad N=\frac{1}{2} \frac{m \beta}{\left(r_{2}^{2}-r_{1}^{2}\right)}\left(r_{2}^{4}-r^{4}\right)\), on integration\\
\includegraphics[max width=\textwidth, alt={}, center]{9f3c5a8e-ac5d-42f4-8eb6-ea2af623f2d8-161_375_384_597_230}\\
\includegraphics[max width=\textwidth, alt={}, center]{9f3c5a8e-ac5d-42f4-8eb6-ea2af623f2d8-161_359_433_573_851}\\
1.309 We assume that the deformation is wholly due to external load, neglecting the effect of the weight of the rod (see next problem). Then a well known formula says, elastic energy per unit volume

\[
=\frac{1}{2} \text { stress } \text { × } \text { strain }=\frac{1}{2} \sigma \varepsilon
\]

This gives \(\frac{1}{2} \frac{m}{\rho} E \varepsilon^{2} \approx 0.04 \mathrm{~kJ}\) for the total deformation energy.\\
1.310 When a rod is deformed by its own weight the stress increases as one moves up, the stretching force being the weight of the portion below the element considered. The stress on the element \(d x\) is

\[
\rho \pi r^{2}(l-x) g / \pi r^{2}=\rho g(l-x)
\]

The extension of the element is

\[
\Delta d x=d \Delta x=\rho g(l-x) d x / E
\]

Integrating \(\Delta l=\frac{1}{2} \rho g l^{2} / E\) is the extension of the whole rod. The elastic energy of the element is

\[
\frac{1}{2} \rho g(l-x) \frac{\rho g(l-x)}{E} \pi r^{2} d x
\]

Integrating

\[
\Delta U=\frac{1}{6} \pi r^{2} \rho^{2} g^{2} l^{3} / E=\frac{2}{3} \pi r^{2} l E\left(\frac{\Delta l}{l}\right)^{2}
\]

\includegraphics[max width=\textwidth, alt={}, center]{9f3c5a8e-ac5d-42f4-8eb6-ea2af623f2d8-161_587_205_1416_1067}\\
1.311 The work done to make a loop out of a steel band appears as the elastic energy of the loop and may be calculated from the same.\\
If the length of the band is \(l\) the radius of the loop \(R=\frac{l}{2 \pi}\). Now consider an element \(A B C D\) of the loop. The elastic energy of this element can be calculated by the same sort of arguments as used to derive the formula for internal bending moment. Consider a fibre at a distance \(z\) from the neutral surface \(P Q\). This fibre experiences a force \(p\) and undergoes an extension \(d s\) where \(d s=Z d \varphi\), while \(P Q=s=R d \varphi\). Thus strain \(\frac{d s}{s}=\frac{Z}{R}\). If \(\alpha\) is the cross sectional area of the fibre, the elastic energy associated with it is

\[
\frac{1}{2} E\left(\frac{Z}{R}\right)^{2} R d \varphi \alpha
\]

Summing over all the fibres we get

\[
\frac{E I \varphi}{2 R} \sum \alpha Z^{2}=\frac{E l d \varphi}{2 R}
\]

For the whole loop this gives,\\
using \(\int d \varphi=2 \pi\),

\[
\frac{E l \pi}{R}=\frac{2 E I \pi^{2}}{l}
\]

Now

\[
I=\int_{-\delta / 2}^{\delta / 2} Z^{2} h d Z=\frac{h \delta^{3}}{12}
\]

\begin{center}
\includegraphics[max width=\textwidth, alt={}]{9f3c5a8e-ac5d-42f4-8eb6-ea2af623f2d8-162_478_548_695_816}
\end{center}

So the energy is

\[
\frac{1}{6} \frac{\pi^{2} E h \delta^{3}}{l}=0.08 \mathrm{~kJ}
\]

1.312 When the rod is twisted through an angle \(\theta\), a couple\\
\(N(\theta)=\frac{\pi r^{4} G}{2 l} \theta\) appears to resist this. Work done in twisting the rod by an angle \(\varphi\) is then

\[
\int_{0}^{\varphi} N(\theta) d \theta=\frac{\pi r^{4} G}{4 l} \varphi^{2}=7 \mathrm{~J} \text { on putting the values. }
\]

1.313 The energy between radii \(r\) and \(r+d r\) is, by differentiation, \(\frac{\pi r^{3} d r}{l} G \varphi^{2}\)

Its density is

\[
\frac{\pi r^{3} d r}{2 \pi r d r l} \frac{G \varphi^{2}}{l}=\frac{1}{2} \frac{G \varphi^{2} r^{2}}{l^{2}}
\]

1.314 The energy density is as usual \(1 / 2\) stress × strain. Stress is the pressure \(\rho g h\). Strain is \(\beta \times \rho g h\) by defination of \(\beta\). Thus

\[
u=\frac{1}{2} \beta(\rho g h)^{2}=23 \cdot 5 \mathrm{~kJ} / \mathrm{m}^{3} \text { on putting the values. }
\]

\subsection*{1.7 HYDRODYNAMICS}
1.315 Between 1 and 2 fluid particles are in nearly circular motion and therefore have centripetal acceleration. The force for this acceleration, like for any other situation in an ideal fluid, can only come from the pressure variation along the line joining 1 and 2 . This requires that pressure at 1 should be greater than the pressure at 2 i.e.

\[
p_{1}>p_{2}
\]

so that the fluid particles can have required acceleration. If there is no turbulence, the motion can be taken as irrotational. Then by considering

\[
\oint \vec{v} \cdot d l=0
\]

along the circuit shown we infer that

\[
v_{2}>v_{1}
\]

(The portion of the circuit near 1 and 2 are streamlines while the other two arms are at right angle to streamlines)\\
\includegraphics[max width=\textwidth, alt={}, center]{9f3c5a8e-ac5d-42f4-8eb6-ea2af623f2d8-163_339_317_492_1020}

In an incompressible liquid we also have div \(\vec{v}=0\)\\
By electrostatic analogy we then find that the density of streamlines is proportional to the velocity at that point.\\
1.316 From the conservation of mass


\begin{equation*}
v_{1} S_{1}=v_{2} S_{2} \tag{1}
\end{equation*}


But \(S_{1}<S_{2}\) as shown in the figure of the problem, therefore

\[
v_{1}>v_{2}
\]

As every streamline is horizontal between 1 \& 2 , Bernoull's theorem becomes

\[
\begin{aligned}
p+\frac{1}{2} \rho v^{2} & =\text { constant, which gives } \\
p_{1} & <p_{2} \text { as } v_{1}>v_{2}
\end{aligned}
\]

As the difference in height of the water column is \(\Delta h\), therefore


\begin{equation*}
p_{2}-p_{1}=\rho g \Delta h \tag{2}
\end{equation*}


From Bernoull's theorem between points 1 and 2 of a streamline\\
or,

\[
p_{1}+\frac{1}{2} \rho v_{1}^{2}=p_{2}+\frac{1}{2} \rho v_{2}^{2}
\]

or


\begin{equation*}
p_{2}-p_{1}=\frac{1}{2} \rho\left(v_{1}^{2}-v_{2}^{2}\right) \tag{3}
\end{equation*}


using (1) in (3), we get

\[
v_{1}=S_{2} \sqrt{\frac{2 g \Delta h}{S_{2}^{2}-S_{1}^{2}}}
\]

Hence the sought volume of water flowing per see\\
\includegraphics[max width=\textwidth, alt={}, center]{9f3c5a8e-ac5d-42f4-8eb6-ea2af623f2d8-163_210_547_1711_843}

\[
Q=v_{1} S_{1}=S_{1} S_{2} \sqrt{\frac{2 g \Delta h}{S_{2}^{2}-S_{1}^{2}}}
\]

1.317 Applying Bernoulli's theorem for the point \(A\) and \(B\),

\[
\begin{aligned}
& p_{A}=p_{B}+\frac{1}{2} \rho v^{2} \\
\text { or, } & \frac{1}{2} \rho v^{2} \\
\text { so, } & v_{A}=0 \\
\text { So, } & v \\
& =\sqrt{\frac{2 \Delta h \rho_{0} g}{\rho}}
\end{aligned}
\]

Thus, rate of flow of gas, \(Q=S v=S \sqrt{\frac{2 \Delta h \rho_{0} g}{\rho}}\)\\
\includegraphics[max width=\textwidth, alt={}, center]{9f3c5a8e-ac5d-42f4-8eb6-ea2af623f2d8-164_351_479_157_922}

The gas flows over the tube past it at \(B\). But at \(A\) the gas becomes stationary as the gas will move into the tube which already contains gas.\\
In applying Bernoulli's theorem we should remember that \(\frac{p}{\rho}+\frac{1}{2} v^{2}+g z\) is constant along a streamline. In the present case, we are really applying Bernoulli's theorem somewhat indirectly. The streamline at \(A\) is not the streamline at \(B\). Nevertheless the result is correct. To be convinced of this, we need only apply Bernoull's theorem to the streamline that goes through \(A\) by comparing the situation at \(A\) with that above \(B\) on the same level. In steady conditions, this agrees with the result derived because there cannot be a transverse pressure differential.\\
1.318 Since, the density of water is greater than that of kerosene oil, it will collect at the bottom. Now, pressure due to water level equals \(h_{1} \rho_{1} g\) and pressure due to kerosene oil level equals \(h_{2} \rho_{2} g\). So, net pressure becomes \(h_{1} \rho_{1} g+h_{2} \rho_{2} g\).

From Bernoulli's theorem, this pressure energy will be converted into kinetic energy while flowing through the whole \(A\).\\
i.e. \(\quad h_{1} \rho_{1} g+h_{2} \rho_{2} g=\frac{1}{2} \rho_{1} v^{2}\)\\
Hence \(\nu=\sqrt{2\left(h_{1}+h_{2} \frac{\rho_{2}}{\rho_{1}}\right) g}=3 \mathrm{~m} / \mathrm{s}\)\\
\includegraphics[max width=\textwidth, alt={}, center]{9f3c5a8e-ac5d-42f4-8eb6-ea2af623f2d8-164_374_507_1119_889}\\
1.319 Let, \(H\) be the total height of water column and the hole is made at a height \(h\) from the bottom.

Then from Bernoulli's theorem

\[
\frac{1}{2} \rho v^{2}=(H-h) \rho g
\]

or, \(v=\sqrt{(H-h) 2 g}\), which is directed horizontally.\\
For the horizontal range, \(l=v t\)\\
\(=\sqrt{2 g(H-h)} \cdot \sqrt{\frac{2 h}{g}}=2 \sqrt{\left(H h-h^{2}\right)}\)\\
\includegraphics[max width=\textwidth, alt={}, center]{9f3c5a8e-ac5d-42f4-8eb6-ea2af623f2d8-164_483_476_1556_926}

Now, for maximum \(l, \frac{d\left(H h-h^{2}\right)}{d h}=0\)\\
which yields

\[
h=\frac{H}{2}=25 \mathrm{~cm}
\]

1.320 Let the velocity of the water jet, near the orifice be \(v^{\prime}\), then applying Bernoullis theorem,\\
or,


\begin{align*}
\frac{1}{2} \rho v^{2} & =h_{0} \rho g+\frac{1}{2} \rho v^{2} \\
v^{\prime} & =\sqrt{v^{2}-2 g h_{0}} \tag{1}
\end{align*}


Here the pressure term on both sides is the same and equal to atmospheric pressure. (In the problem book Fig. should be more clear.)\\
Now, if it rises upto a height \(h\), then at this height, whole of its kinetic energy will be converted into potential energy. So,

\[
\begin{aligned}
& \frac{1}{2} \rho v^{\prime 2}=\rho g h \quad \text { or } \quad h=\frac{v^{\prime 2}}{2 g} \\
= & \frac{v^{2}}{2 g}-h_{0}=20 \mathrm{~cm},[\text { using Eq. }(1)]
\end{aligned}
\]

1.321 Water flows through the small clearance into the orifice. Let \(d\) be the clearance. Then from the equation of continuity\\
\(\left(2 \pi R_{1} d\right) v_{1}=(2 \pi r d) v=\left(2 \pi R_{2} d\right) v_{2}\)\\
or


\begin{equation*}
v_{1} R_{1}=v r=v_{2} R_{2} \tag{1}
\end{equation*}


where \(v_{1}, v_{2}\) and \(v\) are respectively the inward radial velocities of the fluid at 1,2 and 3.\\
Now by Bernoulli's theorem just before 2 and just after it in the clearance


\begin{equation*}
p_{0}+h \rho g=p_{2}+\frac{1}{2} \rho v_{2}^{2} \tag{2}
\end{equation*}


Applying the same theorem at 3 and 1 we find that this also equals\\
\includegraphics[max width=\textwidth, alt={}, center]{9f3c5a8e-ac5d-42f4-8eb6-ea2af623f2d8-165_452_621_1045_802}


\begin{equation*}
p+\frac{1}{2} \rho v^{2}=p_{0}+\frac{1}{2} \rho v_{1}^{2} \tag{3}
\end{equation*}


\section*{(since the pressure in the orifice is \(p_{0}\) )}
From Eqs. (2) and (3) we also hence\\
and


\begin{gather*}
v_{1}=\sqrt{2 g h}  \tag{4}\\
p=p_{0}+\frac{1}{2} \rho v_{1}^{2}\left(1-\left(\frac{v}{v_{1}}\right)^{2}\right) \\
=p_{0}+h \rho g\left(1-\left(\frac{R_{1}}{r}\right)^{2}\right) \quad[\text { Using (1) and (4) }]
\end{gather*}


1322 Let the force acting on the piston be F and the length of the cylinder be \(l\).\\
Then, work done \(=F l\)

Applying Bernoulli's theorem for points \(A\) and \(B, p=\frac{1}{2} \rho v^{2}\) where \(\rho\) is the density and \(v\) is the velocity at point \(B\). Now, force on the piston,


\begin{equation*}
F=p A=\frac{1}{2} \rho v^{2} A \tag{2}
\end{equation*}


where \(A\) is the cross section area of piston.\\
\includegraphics[max width=\textwidth, alt={}, center]{9f3c5a8e-ac5d-42f4-8eb6-ea2af623f2d8-166_320_543_236_852}

Also, discharge through the orifice during time interval \(t=S v t\) and this is equal to the volume of the cylinder, i.e.,


\begin{equation*}
V=S v t \text { or } v=\frac{V}{S t} \tag{3}
\end{equation*}


From Eq. (1), (2) and (3) work done

\[
=\frac{1}{2} \rho v^{2} A l=\frac{1}{2} \rho A \frac{V^{2}}{(S t)^{2}} l=\frac{1}{2} \rho V^{3} / S^{3} t^{2} \quad \text { (as } A l=V
\]

1.323 Let at any moment of time, water level in the vessel be \(H\) then speed of flow of water through the orifice, at that moment will be


\begin{equation*}
v=\sqrt{2 g H} \tag{1}
\end{equation*}


In the time interval \(d t\), the volume of water ejected through orifice,


\begin{equation*}
d V=s v d t \tag{2}
\end{equation*}


On the other hand, the volume of water in the vessel at time \(t\) equals

\[
V=S H
\]

Differentiating (3) with respect to time,


\begin{equation*}
\frac{d V}{d t}=S \frac{d H}{d t} \text { or } \quad d V=S d H \tag{4}
\end{equation*}


Eqs. (2) and (4)

Integrating,


\begin{gather*}
S d H=s v d t \text { or } d t=\frac{S}{s} \frac{d H}{\sqrt{2 g H}}, \text { from (2) }  \tag{2}\\
\int_{0}^{t} d t=\frac{S}{s \sqrt{2 h}} \int_{h}^{0} \frac{d h}{\sqrt{H}} \\
t=\frac{S}{s} \sqrt{\frac{2 h}{g}}
\end{gather*}


Thus,\\
1.324 In a rotating frame (with constant angular velocity) the Eulerian equation is

\[
-\vec{\nabla} p+\rho \vec{g}+2 \rho(\vec{v} \times \vec{\omega})+\rho \omega^{2} \vec{r}=\rho \frac{d \vec{v}^{\prime}}{d t}
\]

In the frame of rotating tube the liquid in the "column" is practically static because the orifice is sufficiently small. Thus the Eulerian Eq. in projection form along \(\vec{r}\) (which is\\
the position vector of an arbitrary liquid element of lenth \(d r\) relative to the rotation axis) reduces to\\
or, \(\quad d p=\rho \omega^{2} r d r\)\\
so, \(\quad \int_{P_{0}}^{P} d p=\rho \omega^{2} \int_{(l-h)}^{r} r d r\)\\
\includegraphics[max width=\textwidth, alt={}, center]{9f3c5a8e-ac5d-42f4-8eb6-ea2af623f2d8-167_367_666_105_645}

Thus


\begin{equation*}
p(r)=p_{0}+\frac{\rho \omega^{2}}{2}\left[r^{2}-(l-h)^{2}\right] \tag{1}
\end{equation*}


Hence the pressure at the end \(B\) just before the orifice i.e.


\begin{equation*}
p(l)=p_{0}+\frac{\rho \omega^{2}}{2}\left(2 l h-h^{2}\right) \tag{2}
\end{equation*}


Then applying Bernoull's theorem at the orifice for the points just inside and outside of the end \(B\)\\
\(p_{0}+\frac{1}{2} \rho \omega^{2}\left(2 l h-h^{2}\right)=p_{0}+\frac{1}{2} \rho v^{2}\) (where \(v\) is the sought velocity)

So,

\[
v=\omega h \sqrt{\frac{2 l}{h}-1}
\]

1.325 The Euler's equation is \(\rho \frac{d \vec{v}}{d t}=\vec{f}-\vec{\nabla} p=-\vec{\nabla}(p+\rho g z)\), where \(z\) is vertically upwards.

Now


\begin{equation*}
\frac{d \vec{v}}{d t}=\frac{\partial \vec{v}}{\partial t}+(\vec{v} \cdot \vec{\nabla}) \vec{v} \tag{1}
\end{equation*}


But


\begin{equation*}
(\vec{v} \cdot \vec{\nabla}) \vec{v}=\vec{\nabla}\left(\frac{1}{2} v^{2}\right)-\vec{v} \times \text { Curl } \vec{v} \tag{2}
\end{equation*}


we consider the steady (i.e. \(\partial \vec{v} \partial \partial t=0\) ) flow of an incompressible fluid then \(\rho=\) constant. and as the motion is irrotational Curl \(\vec{v}=0\)\\
So from (1) and (2)

\[
\begin{aligned}
& \rho \vec{\nabla}\left(\frac{1}{2} v^{2}\right)=-\vec{\nabla}(p+\rho g z) \\
& \vec{\nabla}\left(p+\frac{1}{2} \rho v^{2}+\rho g z\right)=0 \\
& p+\frac{1}{2} \rho v^{2}+\rho g z=\text { constant. }
\end{aligned}
\]

Hence\\
1.326 Let the velocity of water, flowing through \(A\) be \(v_{A}\) and that through \(B\) be \(v_{B}\), then discharging rate through \(A=Q_{A}=S v_{A}\) and similarly through \(B=S v_{B}\).\\
Now, force of reaction at \(A\),

\[
F_{A}=\rho Q_{A} v_{A}=\rho S v_{B}^{2}
\]

Hence, the net force,


\begin{equation*}
F=\rho S\left(v_{A}^{2}-v_{B}^{2}\right) \text { as } \vec{F}_{A} \uparrow \downarrow \vec{F}_{B} \tag{1}
\end{equation*}


Applying Bernoulli's theorem to the liquid flowing out of \(A\) we get

\[
\rho_{0}+\rho g h=\rho_{0}+\frac{1}{2} \rho v_{A}^{2}
\]

and similarly at \(B\)

Hence

\[
\rho_{0}+\rho g(h+\Delta h)=\rho_{0}+\frac{1}{2} \rho v_{B}^{2}
\]

Thus

\[
\begin{array}{r}
\left(v_{B}^{2}-v_{A}^{2}\right) \frac{\rho}{2}=\Delta h \rho g \\
F=2 \rho g S \Delta h=0.50 \mathrm{~N}
\end{array}
\]

\includegraphics[max width=\textwidth, alt={}, center]{9f3c5a8e-ac5d-42f4-8eb6-ea2af623f2d8-168_556_291_120_900}\\
1.327 Consider an element of height \(d y\) at a distance \(y\) from the top. The velocity of the fluid coming out of the element is

\[
v=\sqrt{2 g y}
\]

The force of reaction \(d F\) due to this is \(d F=\rho d A v^{2}\), as in the previous problem, \(=\rho(b d y) 2 g y\)

Integrating

\[
\begin{gathered}
F=\rho g b \int_{h-1}^{h} 2 y d y \\
=\rho g b\left[h^{2}-(h-l)^{2}\right]=\rho g b l(2 h-l)
\end{gathered}
\]

(The slit runs from a depth \(h-l\) to a depth \(h\) from the top.)\\
1.328 Let the velocity of water flowing through the tube at a certain instant of time be \(u\), then \(u=\frac{Q}{\pi r^{2}}\), where \(Q\) is the rate of flow of water and \(\pi r^{2}\) is the cross section area of the tube.\\
From impulse momentum theorem, for the stream of water striking the tube corner, in \(x\)-direction in the time interval \(d t\),

\[
F_{x} d t=-\rho Q u d t \quad \text { or } \quad F_{x}=-\rho Q u
\]

and similarly, \(F_{y}=\rho Q u\)\\
Therefore, the force exerted on the water stream by the tube,

\[
\vec{F}=-\rho Q u \vec{i}+\rho Q u \vec{j}
\]

According to third law, the reaction force on the tube's wall by the stream equals ( \(-\vec{F}\) )

\[
=\rho Q u \vec{i}-\rho Q u \vec{j}
\]

Hence, the sought moment of force about 0 becomes\\
\includegraphics[max width=\textwidth, alt={}, center]{9f3c5a8e-ac5d-42f4-8eb6-ea2af623f2d8-168_515_553_1356_851}

\[
\begin{aligned}
& \vec{N}=l(-\vec{i}) \times(\rho Q u \vec{i}-\rho Q u \vec{j})=\rho Q u l \vec{k}=\frac{\rho Q^{2}}{\pi r^{2}} l \vec{k} \\
& \text { and } \quad \\
& \qquad|\vec{N}|=\frac{\rho Q^{2} l}{\pi r^{2}}=0.70 \mathrm{~N} \cdot \mathrm{~m}
\end{aligned}
\]

1.329 Suppose the radius at \(A\) is \(R\) and it decreases uniformaly to \(r\) at \(B\) where \(S=\pi R^{2}\) and \(s=\pi r^{2}\). Assume also that the semi vectical angle at 0 is \(\alpha\). Then

So

\[
\begin{gathered}
\frac{R}{L_{2}}=\frac{r}{L_{1}}=\frac{y}{x} \\
y=r+\frac{R-r}{L_{2}-L_{1}}\left(x-L_{1}\right)
\end{gathered}
\]

where \(y\) is the radius at the point \(P\) distant \(x\) from the vertex \(O\). Suppose the velocity with which the liquid flows out is \(V\) at \(A, v\) at \(B\) and \(u\) at \(P\). Then by the equation of continuity

\[
\pi R^{2} V=\pi r^{2} v=\pi y^{2} u
\]

The velocity \(v\) of efflux is given by

\[
v=\sqrt{2 g h}
\]

and Bernoulli's theorem gives

\[
p_{p}+\frac{1}{2} \rho u^{2}=p_{0}+\frac{1}{2} \rho v^{2}
\]

where \(p_{p}\) is the pressure at \(P\) and \(p_{0}\) is the atmospheric pressure which is the pressure just outside of \(B\). The force on the nozzle tending to pull it out is then\\
\includegraphics[max width=\textwidth, alt={}, center]{9f3c5a8e-ac5d-42f4-8eb6-ea2af623f2d8-169_406_619_599_782}

\[
F=\int\left(p_{\rho}-p_{0}\right) \sin \theta 2 \pi y d s
\]

We have subtracted \(p_{0}\) which is the force due to atmosphenic pressure the factor \(\sin \theta\) gives horizontal component of the force and \(d s\) is the length of the element of nozzle surface, \(d s=d x \sec \theta\) and

\[
\tan \theta=\frac{R-r}{L_{2}-L_{1}}
\]

Thus

\[
\begin{aligned}
F & =\int_{L_{1}}^{L_{2}} \frac{1}{2}\left(v^{2}-u^{2}\right) \rho 2 \pi y \frac{R-r}{L_{2}-L_{1}} d x \\
& =\pi \rho \int_{r}^{R} v^{2}\left(1-\frac{r^{4}}{y^{4}}\right) y d y \\
& =\pi \rho v^{2} \frac{1}{2}\left(R^{2}-r^{2}+\frac{r^{4}}{R^{2}}-r^{2}\right)=\rho g h\left(\frac{\pi\left(R^{2}-r^{2}\right)^{2}}{R^{2}}\right) \\
& =\rho g h(S-s)^{2} / S=6.02 \mathrm{~N} \text { on putting the values. }
\end{aligned}
\]

Note : If we try to calculate \(F\) from the momentum change of the liquid flowing out w will be wrong even as regards the sign of the force.\\
There is of course the effect of pressure at \(S\) and \(s\) but quantitative derivation of \(F\) fron Newton's law is difficult.\\
1.330 The Euler's equation is \(\rho \frac{d \vec{v}}{d t}=\vec{f}-\vec{\nabla} p\) in the space fixed frame where \(\vec{f}=-\rho g \vec{k}\) downward. We assume incompressible fluid so \(\rho\) is constant.\\
Then \(\vec{f}=-\vec{\nabla}(\rho g z)\) where \(z\) is the height vertically upwards from some fixed origin. We go to rotating frame where the equation becomes

\[
\rho \frac{d \vec{v}^{\prime}}{d t}=-\vec{\nabla}(p+\rho g z)+\rho \omega^{2} \vec{r}+2 \rho\left(\vec{v}^{\prime} \times \vec{\omega}\right)
\]

the additional terms on the right are the well known coriolis and centrifugal forces. In the frame rotating with the liquid \(\vec{v}^{\prime}=0\) so\\
or

\[
\begin{array}{r}
\vec{v}\left(p+\rho g z-\frac{1}{2} \rho \omega^{2} r^{2}\right)=0 \\
p+\rho g z-\frac{1}{2} \rho \omega^{2} r^{2}=\text { constant }
\end{array}
\]

On the free surface \(p=\) constant, thus

\[
z=\frac{\omega^{2}}{2 g} r^{2}+\text { constant }
\]

If we choose the origin at point \(r=0\) (i.e. the axis) of the free surface then "cosntant" \(=0\) and

\[
z=\frac{\omega^{2}}{2 g} r^{2} \text { (The paraboloid of revolution) }
\]

At the bottom \(z=\) constant\\
So

\[
p=\frac{1}{2} \rho \omega^{2} r^{2}+\text { constant }
\]

If \(p=p_{0}\) on the axis at the bottom, then

\[
p=p_{0}+\frac{1}{2} \rho \omega^{2} r^{2} .
\]

1.331 When the disc rotates the fuild in contact with, corotates but the fluid in contact with the walls of the cavity does not rotate. A velocity gradient is then set up leading to viscous forces. At a distance \(r\) from the axis the linear velocity is \(\omega r\) so there is a velocity gradient \(\frac{\omega r}{h}\) both in the upper and lower clearance. The corresponding force on the element whose radial width is \(d r\) is

\[
\eta 2 \pi r d r \frac{\omega r}{h} \text { (from the formular } F=\eta A \frac{d v}{d x} \text { ) }
\]

The torque due to this force is

\[
\eta 2 \pi r d r \frac{\omega r}{h} r
\]

and the net torque considering both the upper and lower clearance is

\[
\begin{aligned}
& 2 \int_{0}^{R} \eta 2 \pi r^{3} d r \frac{\omega}{h} \\
& =\pi R^{4} \omega \eta / h
\end{aligned}
\]

So power developed is

\[
P=\pi R^{4} \omega^{2} \eta / h=9.05 \mathrm{~W} \text { (on putting the values). }
\]

(As instructed end effects i.e. rotation of fluid in the clearance \(r>R\) has been neglected.)\\
1.332 Let us consider a coaxial cylinder of radius \(r\) and thickness \(d r\), then force of friction or viscous force on this elemental layer, \(F=2 \pi r l \eta \frac{d \nu}{d r}\).\\
This force must be constant from layer to layer so that steady motion may be possible.

\[
\text { or, } \quad \frac{F d r}{r}=2 \pi l \eta d v .
\]

Integrating,

\[
\begin{aligned}
& F \int_{R_{2}}^{r} \frac{d r}{r}=2 \pi l \eta \int_{0}^{v} d v \\
& \text { or, } \quad F \ln \left(\frac{r}{R_{2}}\right)=2 \pi l \eta v
\end{aligned}
\]

Putting

\[
\begin{gathered}
r=R_{1}, \text { we get } \\
F \ln \frac{R_{1}}{R_{2}}=2 \pi l \eta v_{0}
\end{gathered}
\]

From (2) by (3) we get,

\[
v=v_{0} \frac{\ln r / R_{2}}{\ln R_{1} / R_{2}}
\]

Note : The force \(F\) is supplied by the agency which tries to carry the inner cylinder with velocity \(\nu_{0}\).\\
1.333 (a) Let us consider an elemental cylinder of radius \(r\) and thickness \(d r\) then from Newton's formula

\[
F=2 \pi r l \eta r \frac{d \omega}{d r}=2 \pi l \eta r^{2} \frac{d \omega}{d r}
\]

and moment of this force acting on the element,


\begin{equation*}
N=2 \pi r^{2} \ln \frac{d \omega}{d r} r=2 \pi r^{3} \ln \frac{d \omega}{d r} \tag{2}
\end{equation*}


or, \(\quad 2 \pi \ln d \omega=N \frac{d r}{r^{3}}\)

As in the previous problem \(N\) is constant when conditions are steady\\
\includegraphics[max width=\textwidth, alt={}, center]{9f3c5a8e-ac5d-42f4-8eb6-ea2af623f2d8-171_401_430_1187_936}

\[
\begin{array}{ll}
\text { Integrating, } & 2 \pi l \eta \int_{0}^{\omega} d \omega=N \int_{R_{1}}^{r} \frac{d r}{r^{3}} \\
\text { or, } & 2 \pi l \eta \omega=\frac{N}{2}\left[\frac{1}{R_{1}^{2}}-\frac{1}{r^{2}}\right] \\
\text { Putting } & r=R_{2}, \omega=\omega_{2}, \text { we get } \\
& 2 \pi l \eta \omega_{2}=\frac{N}{2}\left[\frac{1}{R_{1}^{2}}-\frac{1}{R_{2}^{2}}\right]
\end{array}
\]

From (3) and (4),

\[
\omega=\omega_{2} \frac{R_{1}^{2} R_{2}^{2}}{R_{2}^{2}-R_{1}^{2}}\left[\frac{1}{R_{1}^{2}}-\frac{1}{r^{2}}\right]
\]

(b) From Eq. (4),

\[
N_{1}=\frac{N}{l}=4 \pi \eta \omega_{2} \frac{R_{1}^{2} R_{2}^{2}}{R_{2}^{2}-R_{1}^{2}}
\]

1.334 (a) Let \(d V\) be the volume flowing per second through the cylindrical shell of thickness \(d r\) then,

\[
d V=-(2 \pi r d r) v_{0}\left(1-\frac{r^{2}}{R^{2}}\right)=2 \pi v_{0}\left(r-\frac{r^{3}}{R^{2}}\right) d r
\]

and the total volume,

\[
V=2 \pi v_{0} \int_{0}^{R}\left(r-\frac{r^{3}}{R^{2}}\right) d r=2 \pi v_{0} \frac{R^{2}}{4}=\frac{\pi}{2} R^{2} v_{0}
\]

\includegraphics[max width=\textwidth, alt={}, center]{9f3c5a8e-ac5d-42f4-8eb6-ea2af623f2d8-172_361_358_771_123}\\
(b) Let, \(d E\) be the kinetic energy, within the above cylindrical shell. Then

\[
\begin{gathered}
d T=\frac{1}{2}(d m) v^{2}=\frac{1}{2}(2 \pi r l d r \rho) v^{2} \\
=\frac{1}{2}(2 \pi l \rho) r d r v_{0}^{2}\left(1-\frac{r^{2}}{R^{2}}\right)=\pi l \rho v_{0}\left[r-\frac{2 r^{3}}{R^{2}}+\frac{r^{5}}{R^{4}}\right] d r
\end{gathered}
\]

Hence, total energy of the fluid,

\[
T=\pi l \rho v_{0}^{2} \int_{0}\left(r-\frac{2 r^{3}}{R^{2}}+\frac{r^{5}}{R^{4}}\right) d r=\frac{\pi R^{2} \rho l v_{0}^{2}}{6}
\]

(c) Here frictional force is the shearing force on the tube, exerted by the fluid, which equals \(-\eta S \frac{d \nu}{d t}\).

Given,

\[
v=v_{0}\left(1-\frac{r^{2}}{R^{2}}\right)
\]

So,

\[
\frac{d v}{d r}=-2 v_{0} \frac{r}{R^{2}}
\]

And at

\[
r=R, \frac{d v}{d r}=-\frac{2 v_{0}}{R}
\]

Then, viscous force is given by, \(F=-\eta(2 \pi R l)\left(\frac{d v}{d r}\right)_{r=R}\)

\[
=-2 \pi R \eta l\left(-\frac{2 v_{0}}{R}\right)=4 \pi \eta v_{0} l
\]

(d) Taking a cylindrical shell of thickness \(d r\) and radius \(r\) viscous force,

\[
F=-\eta(2 \pi r l) \frac{d v}{d r},
\]

Let \(\Delta p\) be the pressure difference, then net force on the element \(=\Delta p \pi r^{2}+2 \pi \eta l r \frac{d \nu}{d r}\) But, since the flow is steady, \(F_{\text {net }}=0\)\\
or, \(\quad \Delta p=\frac{-2 \pi \ln r \frac{d v}{d r}}{\pi r^{2}}=\frac{-2 \pi l \eta r\left(-2 v_{0} \frac{r}{R^{2}}\right)}{\pi r^{2}}=4 \eta v_{0} l / R^{2}\)\\
1.335 The loss of pressure head in travelling a distance \(l\) is seen from the middle section to be \(h_{2}-h_{1}=10 \mathrm{~cm}\). Since \(h_{2}-h_{1}=h_{1}\) in our problem and \(h_{3}-h_{2}=15 \mathrm{~cm}=5+h_{2}-h_{1}\), we see that a pressure head of 5 cm remains incompensated and must be converted into kinetic energy, the liquid flowing out. Thus

\[
\frac{\rho v^{2}}{2}=\rho g \Delta h \text { where } \Delta h=h_{3}-h_{2}
\]

Thus

\[
v=\sqrt{2 g \Delta h} \approx 1 \mathrm{~m} / \mathrm{s}
\]

1.336 We know that, Reynold's number \(\left(R_{e}\right)\) is defined as, \(R_{e}=\rho v l / \eta\), where \(v\) is the velocity \(l\) is the characteristic length and \(\eta\) the coefficient of viscosity. In the case of circular cross section the chracteristic length is the diameter of cross-section \(d\), and \(v\) is taken as average velocity of flow of liquid.\\
Now, \(R_{e_{1}}\) (Reynold's number at \(x_{1}\) from the pipe end) \(=\frac{\rho d_{1} v_{1}}{m \eta}\) where \(v_{1}\) is the velocity at distance \(\boldsymbol{x}_{1}\)\\
and similarly,

\[
R_{e_{2}}=\frac{\rho d_{2} v_{2}}{\eta} \text { so } \frac{R_{e_{1}}}{R_{e_{2}}}=\frac{d_{1} v_{1}}{d_{2} v_{2}}
\]

From equation of continuity, \(A_{1} v_{1}=A_{2} v_{2}\)\\
or,

\[
\pi r_{1}^{2} v_{1}=\pi r_{2}^{2} v_{2} \text { or } d_{1} v_{1} r_{1}=d_{2} v_{2} r_{2}
\]

\[
\frac{d_{1} v_{1}}{d_{2} v_{2}}=\frac{r_{2}}{r_{1}}=\frac{r_{0} e^{-\alpha x_{2}}}{r_{0} e^{-\alpha x_{1}}}=e^{-\alpha \Delta x} \quad\left(\text { as } x_{2}-x_{1}=\Delta x\right)
\]

Thus

\[
\frac{R_{e_{2}}}{R_{e_{1}}}=e^{\alpha \Delta x}=5
\]

1.337 We know that Reynold's number for turbulent flow is greater than that on laminar flow.

Now,

\[
\left(R_{e}\right)_{l}=\frac{\rho v d}{\eta}=\frac{2 \rho_{1} v_{1} r_{1}}{\eta_{1}} \text { and }\left(R_{e}\right)_{t}=\frac{2 \rho_{2} v_{2} r_{2}}{\eta}
\]

But, \(\left(R_{e}\right)_{t} \geq\left(R_{e}\right)_{l}\)\\
so \(\quad v_{2_{\max }}=\frac{\rho_{1} v_{1} r_{1} \eta_{2}}{\rho_{2} r_{2} \eta_{1}}=5 \mu \mathrm{~m} / \mathrm{s}\) on putting the values.\\
1.338 We have \(R=\frac{v \rho_{0} d}{\eta}\) and \(v\) is given by

\[
6 \pi \eta r v=\frac{4 \pi}{3} r^{2}\left(\rho-\rho_{0}\right) g
\]

( \(\rho=\) density of lead, \(\rho_{0}=\) density of glycerine.)

\[
v=\frac{2}{9 \eta}\left(\rho-\rho_{0}\right) g r^{2}=\frac{1}{18 \eta}\left(\rho-\rho_{0}\right) g d^{2}
\]

Thus

\[
\frac{1}{2}=\frac{1}{18 \eta^{2}}\left(\rho-\rho_{o}\right) g \rho_{0} d^{3}
\]

and

\[
d=\left[9 \eta^{2} / \rho_{0}\left(\rho-\rho_{0}\right) g\right]^{1 / 3}=5 \cdot 2 \mathrm{~mm} \text { on putting the values. }
\]

1.339

\[
m \frac{d v}{d t}=m g-6 \pi \eta r v
\]

or

\[
\frac{d v}{d t}+\frac{6 \pi \eta r}{m} v=g
\]

or

\[
\frac{d v}{d t}+k v=g, k=\frac{6 \pi \eta r}{m}
\]

or

\[
e^{k t} \frac{d v}{d t}+k e^{k t} v=g e^{k t} \text { or } \frac{d}{d t} e^{k t} v=g e^{k t}
\]

or \(\quad v e^{k t}=\frac{g}{k} e^{k t}+C\) or \(v=\frac{g}{k}+C e^{-k t}\) (where \(C\) is const.)\\
Since

\[
v=0 \text { for } t=0,0=\frac{g}{k}+C
\]

So

\[
C=-\frac{g}{k}
\]

Thus

\[
v=\frac{g}{k}\left(1-e^{-k t}\right)
\]

The steady state velocity is \(\frac{g}{k}\).\\
\(v\) differs from \(\frac{g}{k}\) by \(n\) where \(e^{-k t}=n\)\\
or

\[
t=\frac{1}{k} \ln n
\]

Thus

\[
\frac{1}{k}=-\frac{\frac{4 \pi}{3} r^{3} P}{6 \pi \eta r}=-\frac{4 r^{2} \rho}{18 \eta}=-\frac{d^{2} \rho}{18 \eta}
\]

We have neglected buoyancy in olive oil.

\subsection*{1.8 RELATIVISTIC MECHANICS}
1.340 From the formula for length contraction

\[
\left(l_{0}-l_{0} \sqrt{1-\frac{v^{2}}{c^{2}}}\right)=\eta l_{0}
\]

So,

\[
1-\frac{v^{2}}{c^{2}}=(1-\eta)^{2} \quad \text { or } \quad v=c \sqrt{\eta(2-\eta)}
\]

1.341 (a) In the frame in which the triangle is at rest the space coordinates of the vertices are (000), \(\left(a \frac{\sqrt{3}}{2},+\frac{a}{2}, 0\right)\left(a \frac{\sqrt{3}}{2},-\frac{a}{2}, 0\right)\), all measured at the same time \(t\). In the moving frame the corresponding coordinates at time \(t^{\prime}\) are\\
\(A:\left(v t^{\prime}, 0,0\right), B:\left(\frac{a}{2} \sqrt{3} \sqrt{1-\beta^{2}}+v t^{\prime}, \frac{a}{2}, 0\right)\) and \(C:\left(\frac{a}{2} \sqrt{3} \sqrt{1-\beta^{2}}+v t^{\prime},-\frac{a}{2}, 0\right)\)\\
The perimeter \(P\) is then

\[
P=a+2 a\left(\frac{3}{4}\left(1-\beta^{2}\right)+\frac{1}{4}\right)^{1 / 2}=a\left(1+\sqrt{4-3 \beta^{2}}\right)
\]

(b) The coordinates in the first frame are shown at time \(t\). The coordinates in the moving frame are,\\
\includegraphics[max width=\textwidth, alt={}, center]{9f3c5a8e-ac5d-42f4-8eb6-ea2af623f2d8-175_428_504_1039_174}\\
\includegraphics[max width=\textwidth, alt={}, center]{9f3c5a8e-ac5d-42f4-8eb6-ea2af623f2d8-175_408_572_1038_769}\\
\(A:\left(v t^{\prime}, 0,0\right), B:\left(\frac{a}{2} \sqrt{1-\beta^{2}}+v t^{\prime}, a \frac{\sqrt{3}}{2}, 0\right), C:\left(a \sqrt{1-\beta^{2}}+v t^{\prime}, 0,0\right)\)\\
The perimeter \(P\) is then\\
\(P=a \sqrt{1-\beta^{2}}+\frac{a}{2}\left[1-\beta^{2}+3\right]^{1 / 2} \times 2=a\left(\sqrt{1-\beta^{2}}+\sqrt{4-\beta^{2}}\right)\) here \(\beta=\frac{V}{c}\)\\
1.342 In the rest frame, the coordinates of the ends of the rod in terms of proper length \(l_{0}\)

\[
A:(0,0,0) B:\left(l_{0} \cos \theta_{0}, l_{0} \sin \theta_{0}, 0\right)
\]

at time \(t\). In the laboratory frame the coordinates at time \(t^{\prime}\) are

\[
A:\left(v t^{\prime}, 0,0\right), B:\left(l_{0} \cos \theta_{0} \sqrt{1-\beta^{2}}+v t^{\prime}, l_{0} \sin \theta_{0}, 0\right)
\]

Therefore we can write,

\[
\begin{aligned}
l \cos \theta_{0} & =l_{0} \cos \theta_{0} \sqrt{1-\beta^{2}} \text { and } l \sin \theta= \\
\text { Hence } l_{0}^{2} & =\left(l^{2}\right)\left(\frac{\cos ^{2} \theta+\left(1-\beta^{2}\right) \sin ^{2} \theta}{1-\beta^{2}}\right) \\
\text { or, } \quad & =\sqrt{\frac{1-\beta^{2} \sin ^{2} \theta}{1-\beta^{2}}}
\end{aligned}
\]

\includegraphics[max width=\textwidth, alt={}, center]{9f3c5a8e-ac5d-42f4-8eb6-ea2af623f2d8-176_363_534_82_860}\\
1.343 In the frame \(K\) in which the cone is at rest the coordinates of \(A\) are \((0,0,0)\) and of \(B\) are ( \(h, h \tan \theta, 0\) ). In the frame \(K^{\prime}\), which is moving with velocity \(v\) along the axis of the cone, the coordinates of \(A\) and \(B\) at time \(t^{\prime}\) are\\
\(A:\left(-v t^{\prime}, 0,0\right), B:\left(h \sqrt{1-\beta^{2}}-v t^{\prime}, h \tan \theta, 0\right)\)\\
Thus the taper angle in the frame \(K^{\prime}\) is

\[
\tan \theta^{\prime}=\frac{\tan \theta}{\sqrt{1-\beta^{2}}}\left(=\frac{y_{B}^{\prime}-y_{A}^{\prime}}{x_{B}^{\prime}-x_{A}^{\prime}}\right)
\]

and the lateral surface area is,

\[
\begin{gathered}
S=\pi h^{\prime 2} \sec \theta^{\prime} \tan \theta^{\prime} \\
=\pi h^{2}\left(1-\beta^{2}\right) \frac{\tan \theta}{\sqrt{1-\beta^{2}}} \sqrt{1+\frac{\tan ^{2} \theta}{1-\beta^{2}}}=S_{0} \sqrt{1-\beta^{2} \cos ^{2} \theta}
\end{gathered}
\]

Here \(S_{0}=\pi h^{2} \sec \theta \tan \theta\) is the lateral surface area in the rest frame and

\[
h^{\prime}=h \sqrt{1-\beta^{2}}, \beta=v / c .
\]

1.344 Because of time dilation, a moving clock reads less time. We write,

Thus,

\[
t-\Delta t=t \sqrt{1-\beta^{2}}, \beta=\frac{v}{c}
\]

or,

\[
\begin{gathered}
1-\frac{2 \Delta t}{t}+\left(\frac{\Delta t}{t}\right)^{2}=1-\beta^{2} \\
v=c \sqrt{\frac{\Delta t}{t}\left(2-\frac{\Delta t}{t}\right)}
\end{gathered}
\]

1.345 In the frame \(K\) the length \(l\) of the rod is related to the time of flight \(\Delta t\) by

\[
l=v \Delta t
\]

In the reference frame fixed to the rod (frame \(K^{\prime}\) )the proper length \(l_{0}\) of the rod is given by

\[
l_{0}=v \Delta t^{\prime}
\]

But

\[
l_{0}=\frac{l}{\sqrt{1-\beta^{2}}}=\frac{v \Delta t}{\sqrt{1-\beta^{2}}}, \beta=\frac{v}{c}
\]

Thus,

\[
v \Delta t^{\prime}=\frac{v \Delta t}{\sqrt{1-\beta^{2}}}
\]

So

\[
1-\beta^{2}=\left(\frac{\Delta t}{\Delta t^{\prime}}\right)^{2} \text { or } v=c \sqrt{1-\left(\frac{\Delta t}{\Delta t^{\prime}}\right)^{2}}
\]

and

\[
l_{0}=c \sqrt{\left(\Delta t^{\prime}\right)^{2}-(\Delta t)^{2}}=c \Delta t^{\prime} \sqrt{1-\left(\frac{\Delta t}{\Delta t^{\prime}}\right)^{2}}
\]

1.346 The distance travelled in the laboratory frame of reference is \(v \Delta t\) where \(v\) is the velocity of the particle. But by time dilation

\[
\Delta t=\frac{\Delta t_{0}}{\sqrt{1-v^{2} / c^{2}}} \text { So } v=c \sqrt{1-\left(\Delta t_{0} / \Delta t\right)^{2}}
\]

Thus the distance traversed is

\[
c \Delta t \sqrt{1-\left(\Delta t_{0} / \Delta t\right)^{2}}
\]

1.347 (a) If \(\tau_{0}\) is the proper life time of the muon the life time in the moving frame is

\[
\frac{\tau_{0}}{\sqrt{1-v^{2} / c^{2}}} \text { and hence } l=\frac{v \tau_{0}}{\sqrt{1-v^{2} / c^{2}}}
\]

Thus

\[
\tau_{0}=\frac{l}{v} \sqrt{1-v^{2} / c^{2}}
\]

(The words "from the muon's stand point" are not part of any standard terminology)\\
1.348 In the frame \(K\) in which the particles are at rest, their positions are \(A\) and \(B\) whose coordinates may be taken as,

\[
A:(0,0,0), B=\left(l_{0}, 0,0\right)
\]

In the frame \(K^{\prime}\) with respect to which \(K\) is moving with a velocity \(v\) the coordinates of \(A\) and \(B\) at time \(t^{\prime}\) in the moving frame are\\
\(A=\left(v t^{\prime}, 0,0\right) B=\left(l_{0} \sqrt{1-\beta^{2}}+v t^{\prime}, 0,0\right), \beta=\frac{v}{c}\)\\
Suppose \(B\) hits a stationary target in \(K^{\prime}\) after time \(\boldsymbol{t}_{\boldsymbol{B}}^{\boldsymbol{\prime}}\) while \(\boldsymbol{A}\) hits it after time \(\boldsymbol{t}_{\boldsymbol{B}}+\Delta \boldsymbol{t}\). Then,

\[
l_{0} \sqrt{1-\beta^{2}}+v t_{B}^{\prime}=v\left(t_{B}^{\prime}+\Delta t\right)
\]

\begin{center}
\includegraphics[max width=\textwidth, alt={}]{9f3c5a8e-ac5d-42f4-8eb6-ea2af623f2d8-177_341_443_1413_881}
\end{center}

So, \(\quad l_{0} \frac{v \Delta t}{\sqrt{1-v^{2} / c^{2}}}\)\\
1.349 In the reference frame fixed to the ruler the rod is moving with a velocity \(v\) and suffers Lorentz contraction. If \(l_{0}\) is the proper length of the rod, its measured length will be

\[
\Delta x_{1}=l_{0} \sqrt{1-\beta^{2}}, \beta=\frac{v}{c}
\]

In the reference frame fixed to the rod the ruler suffers Lorentz contraction and we must have

\[
\Delta x_{2} \sqrt{1-\beta^{2}}=l_{0} \text { thus } l_{0}=\sqrt{\Delta x_{1} \Delta x_{2}}
\]

and

\[
1-\beta^{2}=\frac{\Delta x_{1}}{\Delta x_{2}} \text { or } v=c \sqrt{1-\frac{\Delta x_{1}}{\Delta x_{2}}}
\]

1.350 The coordinates of the ends of the rods in the frame fixed to the left rod are shown. The points \(B\) and \(D\) coincides when

\[
l_{0}=c_{1}-v t_{0} \text { or } t_{0}=\frac{c_{1}-l_{0}}{v}
\]

The points \(\boldsymbol{A}\) and \(\boldsymbol{E}\) coincide when\\
\(0=c_{1}+l_{0} \sqrt{1-\beta^{2}}-v t_{1}, t_{1}=\frac{c_{1}+l_{0} \sqrt{1-\beta^{2}}}{v}\)\\
Thus \(\Delta t=t_{1}-t_{0}=\frac{l_{0}}{v}\left(1+\sqrt{1-\beta^{2}}\right)\)\\
or \(\left(\frac{v \Delta t}{l_{0}}-1\right)^{2}=1-\beta^{2}=1-\frac{v^{2}}{c^{2}}\)\\
\includegraphics[max width=\textwidth, alt={}, center]{9f3c5a8e-ac5d-42f4-8eb6-ea2af623f2d8-178_375_534_591_798}

From this

\[
v=\frac{2 c^{2} \Delta t / l_{0}}{1+c^{2} \Delta t^{2} / l_{0}^{2}}=\frac{2 l_{0} / \Delta t}{1+\left(l_{0} / c \Delta t\right)^{2}}
\]

1.351 In \(K_{0}\), the rest frame of the particles, the events corresponding to the decay of the particles are,

\[
A:(0,0,0,0) \text { and }\left(0, l_{0}, 0,0\right)=B
\]

In the reference frame \(K\), the corresponding coordintes are by Lorentz transformation

Now

\[
\begin{gathered}
A:(0,0,0,0), B:\left(\frac{v l_{0}}{c^{2} \sqrt{1-\beta^{2}}}, \frac{l_{0}}{\sqrt{1-\beta^{2}}}, 0,0\right) \\
l_{0} \sqrt{1-\beta^{2}}=l
\end{gathered}
\]

by Lorentz Fitzgerald contraction formula. Thus the time lag of the decay time of \(B\) is

\[
\Delta t=\frac{v l_{0}}{c^{2} \sqrt{1-\beta^{2}}}-\frac{v l}{c^{2}\left(1-\beta^{2}\right)}=\frac{v l}{c^{2}-v^{2}}
\]

\(B\) decays later ( \(B\) is the forward particle in the direction of motion)\\
1.352 (a) In the reference frame \(K\) with respect to which the rod is moving with velocity \(v\), the coordinates of \(A\) and \(B\) are

\[
\begin{aligned}
& A: t, x_{A}+v\left(t-t_{A}\right), 0,0 \\
& B: t, x_{B}+v\left(t-t_{B}\right), 0,0
\end{aligned}
\]

Thus \(l=x_{A}-x_{B}-v\left(t_{A}-t_{B}\right)=l_{0} \sqrt{1-\beta^{2}}\)\\
So \(\quad l_{0}=\frac{x_{A}-x_{B}-v\left(t_{A}-t_{B}\right)}{\sqrt{1-v^{2} / c^{2}}}\)\\
(b) \(\pm l_{0}-v\left(t_{A}-t_{B}\right)=l=l_{0} \sqrt{1-v^{2} / c^{2}}\)\\
(since \(x_{A}-x_{B}\) can be either \(+l_{0}\) or \(-l_{0}\) )\\
Thus \(v\left(t_{A}-t_{B}\right)=\left( \pm 1-\sqrt{1-v^{2} / c_{2}}\right) l_{o}\)\\
i.e. \(t_{A}-t_{B}=\frac{l_{0}}{v}\left(1-\sqrt{1-\frac{v^{2}}{c^{2}}}\right)\)\\
\includegraphics[max width=\textwidth, alt={}, center]{9f3c5a8e-ac5d-42f4-8eb6-ea2af623f2d8-179_412_507_247_812}\\
or \(t_{B}-t_{A}=\frac{l_{0}}{v^{*}}\left(1+\sqrt{1-v^{2} / c^{2}}\right)\)\\
1.353 At the instant the picture is taken the coordintes of \(A, B, A^{\prime}, B^{\prime}\) in the rest frame of \(A B\) are

\[
\begin{gathered}
A:(0,0,0,0) \\
B:\left(0, l_{0}, 0,0\right) \\
B^{\prime}:(0,0,0,0) \\
A^{\prime}:\left(0_{1}-l_{0} \sqrt{1-v^{2} / c^{2}}, 0,0\right)
\end{gathered}
\]

\begin{center}
\includegraphics[max width=\textwidth, alt={}]{9f3c5a8e-ac5d-42f4-8eb6-ea2af623f2d8-179_212_577_827_695}
\end{center}

In this frame the coordinates of \(B^{\prime}\) at other times are \(B^{\prime}:(t, v t, 0,0)\). So \(B^{\prime}\) is opposite to \(B\) at time \(t(B)=\frac{l_{0}}{v}\). In the frame in which \(B^{\prime}, A^{\prime}\) is at rest the time corresponding this is by Lorentz tranformation.

\[
t^{0}\left(B^{\prime}\right)=\frac{1}{\sqrt{1-\frac{v^{2}}{c^{2}}}}\left(\frac{l_{0}}{v}-\frac{v l_{0}}{c^{2}}\right)=\frac{l_{0}}{v} \sqrt{1-v^{2} / c^{2}}
\]

Similarly in the rest frame of \(A, B\), te coordinates of \(A\) at other times are

\[
A^{\prime}:\left(t,-l_{0} \sqrt{1-\frac{v^{2}}{c^{2}}}+v t, 0,0\right)
\]

\(A^{\prime}\) is opposite to \(A\) at time \(t(A)=\frac{l_{0}}{v} \sqrt{1-\frac{v^{2}}{c^{2}}}\)\\
The corresponding time in the frame in which \(A^{\prime}, B^{\prime}\) are at rest is

\[
t\left(A^{\prime}\right)=\gamma t(A)=\frac{l_{0}}{v}
\]

1.354 By Lorentz transformation \(t^{\prime}=\frac{1}{\sqrt{1-\frac{v^{2}}{c^{2}}}}\left(t-\frac{v x}{c^{2}}\right)\)

So at time

\[
t=0, t^{\prime}=\frac{v x}{c^{2}} \frac{1}{\sqrt{1-v^{2} / c^{2}}}
\]

If \(x>0 t^{\prime}<0\), if \(x<0, t^{\prime}>0\) and we get the diagram given below "in terms of the \(K\)-clock".\\
\includegraphics[max width=\textwidth, alt={}, center]{9f3c5a8e-ac5d-42f4-8eb6-ea2af623f2d8-180_261_761_290_353}

The situation in terms of the \(K^{\prime}\) clock is reversed.\\
1.355 Suppose \(x(t)\) is the locus of points in the frame \(K\) at which the readings of the clocks of both reference system are permanently identical, then by Lorentz transformation

\[
t^{\prime}=\frac{1}{\sqrt{1-V^{2} / c^{2}}}\left(t-\frac{V x(t)}{c^{2}}\right)=t
\]

So differentiating \(x(t)=\frac{c^{2}}{V}\left(1-\sqrt{1-\frac{V^{2}}{c^{2}}}\right)=\frac{c}{\beta}\left(1-\sqrt{1-\beta^{2}}\right), \beta=\frac{V}{c}\)\\
Let

\[
\beta=\tan h \theta, 0 \leq \theta<\infty \text {, Then }
\]

\[
\begin{aligned}
& x(t)=\frac{c}{\tan h \theta}\left(1-\sqrt{1-\tan h^{2} \theta}\right)=c \frac{\cos h \theta}{\sin h \theta}\left(1-\frac{1}{\cos h \theta}\right) \\
&=c \frac{\cos h \theta-1}{\sin h \theta}=c \sqrt{\frac{\cos h \theta-1}{\cos h \theta+1}}=c \tan h \frac{\theta}{2} \leq v
\end{aligned}
\]

\(\cdot(\tan h \theta\) is a monotonically increasing function of \(\theta)\)\\
1.356 We can take the coordinates of the two events to be

\[
A:(0,0,0,0) B:(\Delta t, a, 0,0)
\]

For \(B\) to be the effect and \(A\) to be cause we must have \(\Delta t>\frac{|a|}{c}\).\\
In the moving frame the coordinates of \(A\) and \(B\) become\\
\(A:(0,0,0,0), B:\left[\gamma\left(\Delta t-\frac{a V}{c^{2}}\right), \gamma(a-V \Delta t), 0,0\right]\) where \(\gamma=\frac{1}{\sqrt{1-\left(\frac{V^{2}}{c^{2}}\right)}}\)\\
Since\\
\(\left(\Delta t^{\prime}\right)^{2}-\frac{a^{\prime 2}}{c^{2}}=\gamma^{2}\left[\left(\Delta t-\frac{a V}{c^{2}}\right)^{2}-\frac{1}{c^{2}}(a-V \Delta t)^{2}\right)=(\Delta t)^{2}-\frac{a^{2}}{c^{2}}>0\)\\
we must have \(\Delta t^{\prime}>\frac{\left|a^{\prime}\right|}{c}\)\\
1.357 (a) The four-dimensional interval between \(A\) and \(B\) (assuming \(\Delta y=\Delta z=0\) ) is :

\[
5^{2}-3^{2}=16 \text { units }
\]

Therefore the time interval between these two events in the reference frame in which the events occurred at the same place is

\[
c\left(t_{B}^{\prime}-t_{A}^{\prime}\right)=\sqrt{16}=4 \mathrm{~m}
\]

or \(\quad t_{B}^{\prime}-t_{A}^{\prime}=\frac{4}{c}=\frac{4}{3} \times 10^{-8} \mathrm{~s}\)\\
(b) The four dimensional interval between \(A\) and \(C\) is (assuming \(\Delta y=\Delta z=0\) )

\[
3^{2}-5^{2}=-16
\]

So the distance between the two events in the frame in which they are simultaneous is 4 units \(=4 \mathrm{~m}\).\\
\includegraphics[max width=\textwidth, alt={}, center]{9f3c5a8e-ac5d-42f4-8eb6-ea2af623f2d8-181_444_596_245_858}\\
1.358 By the velocity addition formula\\
and

\[
\begin{gathered}
v_{x}^{\prime}=\frac{v_{x}-V}{1-\frac{V v_{x}}{c^{2}}}, v_{y}^{\prime}=\frac{v_{y} \sqrt{1-V^{2} / c^{2}}}{1-\frac{v_{x} V}{c^{2}}} \\
v^{\prime}=\sqrt{v_{x}^{\prime 2}+v_{y}^{\prime 2}}=\frac{\sqrt{\left(v_{x}-V\right)^{2}+v_{y}^{2}\left(1-V^{2} / c^{2}\right)}}{1-\frac{v_{x} V}{c^{2}}}
\end{gathered}
\]

1.359 (a) By definition the velocity of apporach is

\[
v_{\text {approach }}=\frac{d x_{1}}{d t}-\frac{d x_{2}}{d t}=v_{1}-\left(-v_{2}\right)=v_{1}+v_{2}
\]

in the reference frame \(K\).\\
(b) The relative velocity is obtained by the transformation law

\[
v_{r}=\frac{v_{1}-\left(-v_{2}\right)}{1-\frac{v_{1}\left(-v_{2}\right)}{c^{2}}}=\frac{v_{1}+v_{2}}{1+\frac{v_{1} v_{2}}{c^{2}}}
\]

1.360 The velocity of one of the rods in the reference frame fixed to the other rod is

\[
V=\frac{v+v}{1+\frac{v^{2}}{c^{2}}}=\frac{2 v}{1+\beta^{2}}
\]

The length of the moving rod in this frame is

\[
l=l_{0} \sqrt{1-\frac{4 v^{2} / c^{2}}{\left(1+\beta^{2}\right)^{2}}}=l_{0} \frac{1-\beta^{2}}{1+\beta^{2}}
\]

1.361 The approach velocity is defined by

\[
\vec{V}_{\text {approach }}=\frac{d \overrightarrow{r_{1}}}{d t}-\frac{d \overrightarrow{r_{2}}}{d t}=V_{1}-V_{2}
\]

in the laboratory frame. So \(V_{\text {approach }}=\sqrt{v_{1}^{2}+v_{2}^{2}}\)

On the other hand, the relative velocity can be obtained by using the velocity addition formula and has the components

\[
\left[-v_{1}, v_{2} \sqrt{1-\left(\frac{v_{1}^{2}}{c^{2}}\right)}\right] \text { so } V_{r}=\sqrt{v_{1}^{2}+v_{2}^{2}-\frac{v_{1} v_{2}^{2}}{c^{2}}}
\]

1.362 The components of the velocity of the unstable particle in the frame \(K\) are

\[
\left(V, v^{\prime} \sqrt{1-\frac{V^{2}}{c^{2}}}, 0\right)
\]

so the velocity relative to \(K\) is

\[
\sqrt{V^{2}+v^{\prime 2}-\frac{v^{\prime 2} V^{2}}{c^{2}}}
\]

The life time in this frame dilates to

\[
\Delta t_{0} / \sqrt{1-\frac{V^{2}}{c^{2}}-\frac{v^{\prime 2}}{c^{2}}+\frac{v^{\prime 2} V^{2}}{c^{4}}}
\]

and the distance traversed is

\[
\Delta t_{0} \frac{\sqrt{V^{2}+v^{\prime 2}-\left(v^{\prime 2} V^{2}\right) / c^{2}}}{\sqrt{1-V^{2} / c^{2}} \sqrt{1-v^{\prime 2} / c^{2}}}
\]

\includegraphics[max width=\textwidth, alt={}, center]{9f3c5a8e-ac5d-42f4-8eb6-ea2af623f2d8-182_443_443_444_739}\\
1.363 In the frame \(K^{\prime}\) the components of the velocity of the particle are

\[
\begin{gathered}
v_{x}^{\prime}=\frac{v \cos \theta-V}{1-\frac{v V \cos \theta}{c^{2}}} \\
v_{y}^{\prime}=\frac{v \sin \theta \sqrt{1-V^{2} / c^{2}}}{1-\frac{v V}{c^{2}} \cos \theta}
\end{gathered}
\]

Hence, \(\tan \theta^{\prime}=\frac{v_{y}^{\prime}}{v_{x}^{\prime}}=\frac{v \sin \theta}{v \cos \theta-V} \sqrt{\left(1-V^{2}\right) / c^{2}}\)\\
\includegraphics[max width=\textwidth, alt={}, center]{9f3c5a8e-ac5d-42f4-8eb6-ea2af623f2d8-182_399_393_1073_826}\\
1.364 In \(K^{\prime}\) the coordinates of \(A\) and \(B\) are

\[
A:\left(t^{\prime}, 0,-v^{\prime} t^{\prime}, 0\right) ; B:\left(t^{\prime}, l,-v^{\prime} t^{\prime}, 0\right)
\]

After performing Lorentz transformation to the frame \(K\) we get

\[
\begin{gathered}
A: t=\gamma t^{\prime} \quad B: t=\gamma\left(t^{\prime}+\frac{V l}{c^{2}}\right) \\
x=\gamma V t^{\prime}, x=\gamma\left(l+V t^{\prime}\right) \\
y=v^{\prime} t^{\prime} \quad y=-v^{\prime} t^{\prime} \\
z=0 \quad z=0
\end{gathered}
\]

By translating \(t^{\prime} \rightarrow t^{\prime}-\frac{V l}{c^{2}}\), we can write the coordinates of \(B\) as \(B: t=\gamma t^{\prime}\)

\[
\begin{aligned}
& x=\gamma l\left(1-\frac{V^{2}}{c^{2}}\right)+V t^{\prime} \gamma=l \sqrt{1-\frac{v^{2}}{c^{2}}}+V t^{\prime} \gamma \\
& y=-v^{\prime}\left(t^{\prime}-\frac{V l}{c^{2}}\right), z=0 \\
& \text { Thus } \\
& \Delta x=l \sqrt{1-\left(\frac{V^{2}}{c^{2}}\right)}, \Delta y=\frac{v^{\prime} V l}{c^{2}} \\
& \text { Hence } \\
& \tan \theta^{\prime}-\frac{v^{\prime} V}{c^{2} \sqrt{1-\frac{v^{\prime} V}{c^{2}}}}
\end{aligned}
\]

In \(K\) the velocities at time \(t\) and \(t+d t\) are respectively \(v\) and \(v+w d t\) along \(x\) - axis which is parallel to the vector \(\vec{V}\). In the frame \(K^{\prime}\) moving with velocity \(\vec{V}\) with respect to \(K\), the velocities are respectively,

\[
\frac{v-V}{1-\frac{v V}{c^{2}}} \text { and } \frac{v+w d t-V}{1-(v+w d t) \frac{V}{c^{2}}}
\]

The latter velocity is written as

\[
\frac{v-V}{1-v \frac{V}{c^{2}}}+\frac{w d t}{1-v \frac{V}{c^{2}}}+\frac{v-V}{\left(1-\frac{v V}{c^{2}}\right)} \frac{w V}{c^{2}} d t=\frac{v-V}{1-v \frac{V}{c^{2}}}+\frac{w d t\left(1-\frac{V^{2}}{c^{2}}\right)}{\left(1-\frac{v V}{c^{2}}\right)^{2}}
\]

Also by Lorentz transformation

\[
d t^{\prime}=\frac{d t-V d x / c^{2}}{\sqrt{1-V^{2} / c^{2}}}=d t \frac{1-v V / c^{2}}{\sqrt{1-V^{2} / c^{2}}}
\]

Thus the acceleration in the \(K^{\prime}\) frame is

\[
w^{\prime}=\frac{d v^{\prime}}{d t^{\prime}}=\frac{w}{\left(1-\frac{v V}{c^{2}}\right)^{3}}\left(1-\frac{V^{2}}{c^{2}}\right)^{3 / 2}
\]

(b) In the \(K\) frame the velocities of the particle at the time \(t\) and \(t+d t\) are repectively

\[
(0, v, 0) \text { and }(0, v+w d t, 0)
\]

where \(\vec{V}\) is along \(x\)-axis. In the \(K^{\prime}\) frame the velocities are\\
and

\[
\begin{gathered}
\left(-V, v \sqrt{1-V^{2} / c^{2}}, 0\right) \\
\left(-V,(v+w d t) \sqrt{1-V^{2} / c^{2}}, 0\right) \text { respectively }
\end{gathered}
\]

Thus the acceleration

\[
w^{\prime}=\frac{w d t \sqrt{\left(1-V^{2} / c^{2}\right)}}{d t^{\prime}}=w\left(1-\frac{V^{2}}{c^{2}}\right) \text { along the } y \text {-axis. }
\]

We have used \(d t^{\prime}=\frac{d t}{\sqrt{1-V^{2} / c^{2}}}\)\\
1.366 In the instantaneous rest frame \(v=V\) and

So,

\[
\begin{aligned}
w^{\prime}= & \frac{w}{\left(1-\frac{v^{2}}{c^{2}}\right)^{3 / 2}} \text { (from 1.365a) } \\
& =\frac{d v}{\left(1-\frac{V^{2}}{c^{2}}\right)^{3 / 2}}=w^{\prime} d t
\end{aligned}
\]

\(w^{\prime}\) is constant by assumption. Thus integration gives

\[
v=\frac{w^{\prime} t}{\sqrt{1+\left(\frac{w^{\prime} t}{c}\right)^{2}}}
\]

Integrating once again \(x=\frac{c^{2}}{w^{\prime}}\left(\sqrt{1+\left(\frac{w^{\prime} t}{c}\right)^{2}}-1\right)\)\\
1.367 The boost time \(\tau_{0}\) in the reference frame fixed to the rocket is related to the time \(\tau\) elapsed on the earth by

\[
\begin{aligned}
& \tau_{0}=\int_{0}^{\tau} \sqrt{1-\frac{v^{2}}{c^{2}}} d t=\int_{0}^{\tau}\left[1-\frac{\left(\frac{w^{\prime} t}{c}\right)^{2}}{1+\left(\frac{w^{\prime} t}{c}\right)^{2}}\right]^{1 / 2} d t \\
& =\int_{0}^{\tau} \frac{d t}{\sqrt{1+\left(\frac{w^{\prime} t}{c}\right)^{2}}}=\frac{c}{w^{\prime}} \int_{0}^{\left(w^{\prime}\right) / c} \frac{d \xi}{\sqrt{1+\xi^{2}}}=\frac{c}{w^{\prime}} \ln \left[\frac{w^{\prime} \tau}{c}+\sqrt{1+\left(\frac{w^{\prime} \tau}{c}\right)^{2}}\right] \\
& m=\frac{m_{0}}{\sqrt{1-\beta^{2}}}
\end{aligned}
\]

For

\[
\beta \approx 1, \frac{m}{m_{0}} \approx \frac{1}{\sqrt{2(1-\beta)}}=\frac{1}{\sqrt{2 \eta}}
\]

1.369 We define the density \(\rho\) in the frame \(K\) in such a way that \(\rho d x d y d z\) is the rest mass \(d m_{0}\) of the element. That is \(\rho d x d y d z=\rho_{0} d x_{0} d y_{0} d z_{0}\), where \(\rho_{0}\) is the proper density \(d x_{0}, d y_{0}, d z_{0}\) are the dimensions of the element in the rest frame \(K_{0}\). Now

\[
d y=d y_{0}, d z=d z_{0}, d x=d x_{0} \sqrt{1-\frac{v^{2}}{c^{2}}}
\]

if the frame \(K\) is moving with velocity, \(v\) relative to the frame \(K_{0}\). Thus

\[
\rho=\frac{\rho_{0}}{\sqrt{1-\frac{v^{2}}{c^{2}}}}
\]

Defining \(\eta\) by \(\rho=\rho_{0}(1+\eta)\)\\
We get \(1+\eta=\frac{1}{\sqrt{1-\frac{v^{2}}{c^{2}}}}\) or, \(\frac{v^{2}}{c^{2}}=1-\frac{1}{(1+\eta)^{2}}=\frac{\eta(2+\eta)}{(1+\eta)^{2}}\)\\
or

\[
v=c \sqrt{\frac{\eta(2+\eta)}{(1+\eta)^{2}}}=\frac{c \sqrt{\eta(2+\eta)}}{1+\eta}
\]

1.370 We have

\[
\frac{m_{0} v}{\sqrt{1-\frac{v^{2}}{c^{2}}}}=p \text { or, } \frac{m_{0}}{\sqrt{1-\frac{v^{2}}{c^{2}}}}=\sqrt{m_{0}^{2}+\frac{p^{2}}{c^{2}}}
\]

or

\[
1-\frac{v^{2}}{c^{2}}=\frac{m_{0}^{2} c^{2}}{m_{0}^{2} c^{2}+p^{2}}=1-\frac{p^{2}}{p^{2}+m_{0}^{2} c^{2}}
\]

or

\[
v=\frac{c_{\mathrm{p}}}{\sqrt{p^{2}+m_{0}^{2} c^{2}}}=\frac{c}{\sqrt{1+\left(\frac{m_{0} c}{p}\right)^{2}}}
\]

So \(\frac{c-v}{c}=\left[1-\left(1+\left(\frac{m_{0} c}{p}\right)^{2}\right)^{-1 / 2}\right] \times 100 \% \cdots \frac{1}{2}\left(\frac{m_{0} c}{p}\right)^{2} \times 100 \%\)\\
1.371 By definition of \(\eta\),

\[
\frac{m_{0} v}{\sqrt{1-\frac{v^{2}}{c^{2}}}}=\eta m_{0} v \text { or } 1-\frac{v^{2}}{c^{2}}=\frac{1}{\eta^{2}}
\]

or

\[
v=c \sqrt{1-\frac{1}{\eta^{2}}}=\frac{c}{\eta} \sqrt{\eta^{2}-1}
\]

1.372 The work done is equal to change in kinetic energy which is different in the two cases Classically i.e. in nonrelativistic mechanics, the change in kinetic energy is

\[
\frac{1}{2} m_{0} c^{2}\left((0.8)^{2}-(0.6)^{2}\right)=\frac{1}{2} m_{0} c^{2} 0.28=0.14 m_{0} c^{2}
\]

Relativistically it is, .

\[
\begin{aligned}
\frac{m_{0} c^{2}}{\sqrt{1-(0.8)^{2}}}-\frac{m_{0} c^{2}}{\sqrt{1-(0.6)^{2}}} & =\frac{m_{0} c^{2}}{0.6}-\frac{m_{0} c^{2}}{0.8}=m_{0} c^{2}(1.666-1.250) \\
& =0.416 m_{0} c^{2}=0.42 m_{0} c^{2}
\end{aligned}
\]

\(1.373 \frac{m_{0} c^{2}}{\sqrt{1-\frac{v^{2}}{c^{2}}}}=2 m_{0} c^{2}\)\\
or

\[
\sqrt{1-\frac{v^{2}}{c^{2}}}=\frac{1}{2} \quad \text { or } \quad 1-\frac{v^{2}}{c^{2}}=\frac{1}{4}
\]

or

\[
\frac{v}{c}=\frac{\sqrt{3}}{2} \cdot \text { i.e. } v=c \frac{\sqrt{3}}{2}
\]

1.374 Relativistically

So

\[
\beta_{r e l}^{2} \approx \frac{2 T}{m_{0} c^{2}}-\frac{3}{4}\left(\beta_{r e l}^{2}\right)^{2} \approx \frac{2 T}{m_{0} c^{2}}-\frac{3}{4}\left(\frac{2 T}{m_{0} c^{2}}\right)^{2}
\]

Thus

\[
-\beta_{r e l}=\left[\frac{2 T}{m_{0} c^{2}}-3 \frac{T^{2}}{m_{0}^{2} c^{4}}\right]^{1 / 2}=\sqrt{\frac{2 T}{m_{0} c^{2}}}\left(1-\frac{3}{4} \frac{T}{m_{0} c^{2}}\right)
\]

But Classically, \(\beta_{c l}=\sqrt{\frac{2 T}{m_{0} c^{2}}}\) so \(\frac{\beta_{\text {rel }}-\beta_{c l}}{\beta_{c l}}=\frac{3}{4} \frac{T}{m_{0} c^{2}}=\varepsilon\)\\
Hence if

\[
\frac{T}{m_{0} c^{2}}<\frac{4}{3} \varepsilon
\]

the velocity \(\beta\) is given by the classical formula with an error less than \(\varepsilon\).\\
1.375 From the formula

\[
E=\frac{m_{0} c^{2}}{\sqrt{1-\frac{v^{2}}{c^{2}}}}, p=\frac{m_{0} v}{\sqrt{1-v^{2} / c^{2}}}
\]

we find

\[
E^{2}=c^{2} p^{2}+m_{0}^{2} c^{4} \text { or }\left(m_{0} c^{2}+T\right)^{2}=c^{2} p^{2}+m_{0}^{2} c^{4}
\]

or

\[
T\left(2 m_{0} c^{2}+T\right)=c^{2} p^{2} \text { i.e. } p=\frac{1}{c} \sqrt{T\left(2 m_{0} c^{2}+T\right)}
\]

1.376 Let the total force exerted by the beam on the target surface be \(F\) and the power liberated there be \(P\). Then, using the result of the previous problem we see

\[
F=N p=\frac{N}{c} \sqrt{T\left(T+2 m_{0} c^{2}\right)}=\frac{I}{e c} \sqrt{T\left(T+2 m_{0} c^{2}\right)}
\]

since \(I=N e, N\) being the number of particles striking the target per second. Also,

\[
P=N\left(\frac{m_{0} c^{2}}{\sqrt{1-v^{2} / c^{2}}}-m_{0} c^{2}\right)=\frac{I}{e} T
\]

These will be, respectively, equal to the pressure and power developed per unit area of the target if \(I\) is current density.\\
1.377 In the frame fixed to the sphere :- The momentum transferred to the eastically scatterred particle is

\[
\frac{2 m v}{\sqrt{1-\frac{v^{2}}{c^{2}}}}
\]

The density of the moving element is, from \(1.369, n \frac{1}{\sqrt{1-\frac{v^{2}}{c^{2}}}}\)\\
and the momentum transferred per unit time per unit area is\\
\(p=\) the pressure \(=\frac{2 m v}{\sqrt{1-\frac{v^{2}}{c^{2}}}} n \frac{1}{\sqrt{1-\frac{v^{2}}{c^{2}}}} \cdot v=\frac{2 m n v^{2}}{1-\frac{v^{2}}{c^{2}}}\)\\
In the frame fixed to the gas :- When the sphere hits a stationary particle, the latter recoils with a velocity

\[
=\frac{v+v}{1+\frac{v^{2}}{c^{2}}}=\frac{2 v}{1+\frac{v^{2}}{c^{2}}}
\]

The momentum transferred is \(\frac{\frac{m \cdot 2 v}{1+v^{2} / c^{2}}}{\sqrt{1-\frac{4 v^{2} / c^{2}}{\left(1-v^{2} / c^{2}\right)^{2}}}}=\frac{2 m v}{1-\frac{v^{2}}{c^{2}}}\)\\
and the pressure is \(\frac{2 m v}{1-\frac{v^{2}}{c^{2}}} \cdot n \cdot v=\frac{2 m n v^{2}}{1-\frac{v^{2}}{c^{2}}}\)\\
1.378 The equation of motion is

\[
\frac{d}{d t}\left(\frac{m_{0} v}{\sqrt{1-\frac{v^{2}}{c^{2}}}}\right)=F
\]

Integrating \(=\frac{v / c}{\sqrt{1-\frac{v^{2}}{c^{2}}}}=\frac{\beta}{\sqrt{1-\beta^{2}}}=\frac{F_{t}}{m_{0} c}\), using \(v=0\) for \(t=0\)\\
\(\frac{\beta^{2}}{1-\beta^{2}}=\left(\frac{F t}{m_{0} c}\right)^{2}\) or, \(\beta^{2}=\frac{(F t)^{2}}{(F t)^{2}+\left(m_{0} c\right)^{2}}\) or, \(v=\frac{F c t}{\sqrt{\left(m_{0} c\right)^{2}+(F t)^{2}}}\)\\
or \(\quad x=\int \frac{F c t d t}{\sqrt{F^{2} t^{2}+m_{0}^{2} c^{2}}}=\frac{c}{F} \int \frac{\xi d \xi}{\sqrt{\xi^{2}+\left(m_{0} c\right)^{2}}}=\frac{c}{F} \sqrt{F^{2} t^{2}+m_{0}^{2} c^{2}}+\) constant\\
or using \(x=0\) at \(t=0\), we get, \(x=\sqrt{c^{2} t^{2}+\left(\frac{m_{0} c^{2}}{F}\right)^{2}}-\frac{m_{0} c^{2}}{F}\)\\
\(1.379 x=\sqrt{a^{2}+c^{2} t^{2}}\), so \(\dot{x}=v=\frac{c^{2} t}{a^{2}+c^{2} t^{2}}\)\\
or, \(\frac{v}{\sqrt{1-\frac{v^{2}}{c^{2}}}}=\frac{c^{2} t}{a}\). Thus \(\frac{d}{d t}\left(\frac{m_{0} v}{\sqrt{1-\frac{v^{2}}{c^{2}}}}\right)=\frac{m_{0} c^{2}}{a}=F\)\\
\(1.380 \vec{F}=\frac{d}{d t}\left(\frac{m_{0} \vec{v}}{\sqrt{1-\frac{v^{2}}{c^{2}}}}\right)=m_{0} \frac{\dot{\vec{v}}}{\sqrt{1-\frac{v^{2}}{c^{2}}}}+m_{0} \frac{\vec{v}}{c^{2}} \vec{v} \cdot \dot{\vec{v}} \frac{1}{\left(1-\frac{v^{2}}{c^{2}}\right)^{3 / 2}}\)\\
Thus

\[
\begin{aligned}
& \vec{F}_{\perp}=m_{0} \frac{\vec{w}}{\sqrt{1-\beta^{2}}}, \vec{w}=\dot{\vec{v}}, \vec{w}_{\perp} \vec{v} \\
& \vec{F}_{\|}=m_{0} \frac{\vec{w}}{\left(1-\beta^{2}\right)^{3 / 2}}, \vec{w}=\dot{\vec{v}}, \overrightarrow{w_{\|}} \vec{v}
\end{aligned}
\]

1.381 By definition,\\
\(E=m_{0} \frac{c^{2}}{\sqrt{1-\frac{v_{x}^{2}}{c^{2}}}}=\frac{m_{0} c^{3} d t}{d s}, p_{x}=m_{0} \frac{v_{x}}{\sqrt{1-\frac{v^{2}}{c^{2}}}}=\frac{c m_{0} d x}{d s}\)\\
where \(d s^{2}=c^{2} d t^{2}-d x^{2}\) is the invariant interval \((d y=d z=0)\)\\
Thus,

\[
\begin{aligned}
& p_{x}^{\prime}=c m_{0} \frac{d x^{\prime}}{d s}=c m_{0} \gamma \frac{(d x-V d t)}{d s}=\frac{p_{x}-V E / c^{2}}{\sqrt{1-V^{2} / c^{2}}} \\
& E^{\prime}=m_{0} c^{3} \frac{d t^{\prime}}{d s}-c^{3} m_{0} \gamma \frac{\left(d t-\frac{V d_{x}}{c^{2}}\right)}{d s}=\frac{E-V p_{x}}{\sqrt{1-\frac{V^{2}}{c^{2}}}}
\end{aligned}
\]

1.382 For a photon moving in the \(x\) direction

\[
\varepsilon=c p_{x}, p_{y}=p_{z}=0
\]

In the moving frame, \(\varepsilon^{\prime}=\frac{1}{\sqrt{1-\beta^{2}}}\left(\varepsilon-V \frac{\varepsilon}{c}\right)=\varepsilon \sqrt{\frac{1-V / c}{1+V / c}}\)\\
Note that

\[
\varepsilon^{\prime}=\frac{\varepsilon}{2} \text { if, } \frac{1}{4}=\frac{1-\beta}{1+\beta} \text { or } \beta=\frac{3}{5}, V=\frac{3 c}{5}
\]

1.383 As before

\[
E=m_{0} c^{3} \frac{d t}{d s}, p_{x}=m_{0} c \frac{d x}{d s}
\]

Similarly

\[
p_{y}=m_{0} c \frac{d y}{d s}, p_{z}=m_{0} c \frac{d z}{d s}
\]

Then \(\quad E^{2}-c^{2} p^{2}=E^{2}-c^{2}\left(p_{x}^{2}+p_{y}^{2}+p_{x}^{2}\right)\)

\[
=m_{0}^{2} c^{4} \frac{\left(c^{2} d t^{2}-d x^{2}-d y^{2}-d z^{2}\right)}{d s^{2}}=m_{0}^{2} c^{4} \text { is invariant }
\]

1.384 (b) \& (a) In the CM frame, the total momentum is zero, Thus

\[
\frac{V}{c}=\frac{c p_{1 x}}{E_{1}+E_{2}}=\frac{\sqrt{T\left(T+2 m_{0} c^{2}\right)}}{T+2 m_{0} c^{2}}=\sqrt{\frac{T}{T+2 m_{0} c^{2}}}
\]

where we have used the result of probiem (1.375)\\
Then

\[
\frac{1}{\sqrt{1-V^{2} / c^{2}}}=\frac{1}{\sqrt{1-\frac{T}{T+2 m_{0} c^{2}}}}=\sqrt{\frac{T+2 m_{0} c^{2}}{2 m_{0} c^{2}}}
\]

Total energy in the CM frame is

\[
\frac{2 m_{0} c^{2}}{\sqrt{1-V^{2} / c^{2}}}=2 m_{0} c^{2} \sqrt{\frac{T+2 m_{0} c^{2}}{2 m_{0} c^{2}}}=\sqrt{2 m_{0} c^{2}\left(T+2 m_{0} c^{2}\right)}=\tilde{T}+2 m_{0} c^{2}
\]

So

\[
\tilde{T}=2 m_{0} c^{2}\left(\sqrt{1+\frac{T}{2 m_{0} c^{2}}}-1\right)
\]

Also \(2 \sqrt{c^{2} \tilde{p}^{2}+m_{0}^{2} c^{4}}=\sqrt{2 m_{0} c^{2}\left(T+2 m_{0} c^{2}\right)}, 4 c^{2} \tilde{p}^{2}=2 m_{0} c^{2} T\), or \(\tilde{p}=\sqrt{\frac{1}{2} m_{0} T}\)\\
1.385

\[
\begin{aligned}
& M_{0} c^{2}=\sqrt{E^{2}-c^{2} p^{2}} \\
& \sqrt{\left(2 m_{0} c^{2}+T\right)^{2}-T\left(2 m_{0} c^{2}+T\right)}=\sqrt{2 m_{0} c^{2}\left(2 m_{0} c^{2}+T\right)}=c \sqrt{2 m_{0}\left(2 m_{0} c^{2}+T\right)}
\end{aligned}
\]

Also \(\quad c p=\sqrt{T\left(T+2 m_{0} c^{2}\right)}, v=\frac{c^{2} p}{E}=c \sqrt{\frac{T}{T+2 m_{0} c^{2}}}\)\\
1.386 Let \(T=\) kinetic energy of a proton striking another stationary particle of the same rest mass. Then, combined kinetic energy in the CM frame

\[
\begin{gathered}
=2 m_{0} c^{2}\left(\sqrt{1+\frac{T^{\prime}}{2 m_{0} c^{2}}}-1\right)=2 T,\left(\frac{T}{m_{0} c^{2}}+1\right)^{2}=1+\frac{T^{\prime}}{2 m_{0} c^{2}} \\
\frac{T^{\prime}}{2 m_{0} c^{2}}=\frac{T\left(2 m_{0} c^{2}+T\right)}{m_{0}^{2} c^{4}}, T^{\prime}=\frac{2 T\left(T+2 m_{o} c^{2}\right)}{m_{0} c^{2}}
\end{gathered}
\]

1.387 We have

\[
E_{1}+E_{2}+E_{3}=m_{0} c^{2}, \overrightarrow{p_{1}}+\overrightarrow{p_{2}}+\overrightarrow{p_{3}}=0
\]

Hence \(\quad\left(m_{0} c^{2}-E_{1}\right)^{2}-c^{2} \vec{p}_{1}^{2}=\left(E_{2}+E_{3}\right)^{2}-\left(\overrightarrow{p_{2}}+\overrightarrow{p_{3}}\right)^{2} c^{2}\)\\
The L.H.S. \(=\left(m_{0}^{2} c^{4}-E_{1}\right)^{2}-c^{2} \vec{p}_{1}^{2}=\left(m_{0}^{2}+m_{1}^{2}\right) c^{4}-2 m_{0} c^{2} E_{1}\)\\
The R.H.S. is an invariant. We can evaluate it in any frame. Choose the CM frame of the particles 2 and 3.\\
In this frame R.H.S. \(=\left(E_{2}^{\prime}+E_{3}^{\prime}\right)^{2}=\left(m_{2}+m_{3}\right)^{2} c^{4}\)\\
Thus \(\left(m_{0}^{2}+m_{1}^{2}\right) c^{4}-2 m_{0} c^{2} E_{1}=\left(m_{2}+m_{3}\right)^{2} c^{4}\)\\
or \(2 m_{0} c^{2} E_{1} \leq\left\{m_{0}^{2}+m_{1}^{2}-\left(m_{2}+m_{3}\right)^{2}\right\} c^{4}\), or \(E_{1} \leq \frac{m_{0}^{2}+m_{1}^{2}-\left(m_{2}+m_{3}\right)^{2}}{2 m_{0}} c^{2}\)\\
1.388 The velocity of ejected gases is \(u\) realtive to the rocket. In an earth centred frame it is

\[
\frac{v-u}{1-\frac{v u}{c^{2}}}
\]

in the direction of the rocket. The momentum conservation equation then reads

\[
(m+d m)(v+d v)+\frac{v-u}{1-\frac{u v}{c^{2}}}(-d m)=m v
\]

or

\[
m d v-\left(\frac{v-u}{1-\frac{u v}{c^{2}}}-v\right) d m=0
\]

Here - \(d m\) is the mass of the ejected gases. so

\[
\begin{gathered}
m d v-\frac{-u+\frac{u v^{2}}{c^{2}}}{1-\frac{u v}{c^{2}}} d m=0, \text { or } m d v+u\left(1-\frac{v^{2}}{c^{2}}\right) d m=0 \\
\text { (neglecting } 1-\frac{u v}{c^{2}} \text { since } u \text { is non- relativistic.) }
\end{gathered}
\]

Integrating \(\left(\beta=\frac{v}{c}\right), \int \frac{d \beta}{1-\beta^{2}}+\frac{u}{c} \int \frac{d m}{m}=0, \ln \frac{1+\beta}{1-\beta}+\frac{u}{c} \ln m=\) constant\\
The constant \(=\frac{u}{c} \ln m_{0}\) since \(\beta=0\) initially.

Thus

\[
\frac{1-\beta}{1+\beta}=\left(\frac{m}{m_{0}}\right)^{w / c} \quad \text { or } \quad \beta=\frac{1-\left(\frac{m}{m_{0}}\right)^{w / c}}{1+\left(\frac{m}{m_{0}}\right)^{w / c}}
\]

\section*{PART TWO}
\section*{THERMODYNAMICS AND MOLECULAR PHYSICS}
\subsection*{2.1 EQUATION OF THE GAS STATE • PROCESSES}
2.1 Let \(m_{1}\) and \(m_{2}\) be the masses of the gas in the vessel before and after the gas is releasedHence mass of the gas released,

\[
\Delta m=m_{1}-m_{2}
\]

Now from ideal gas equation

\[
p_{1} V=m_{1} \frac{R}{M} T_{0} \text { and } p_{2} V=m_{2} \frac{R}{M} T_{0}
\]

as \(V\) and \(T\) are same before and after the release of the gas.\\
so, \(\quad\left(p_{1}-p_{2}\right) V=\left(m_{1}-m_{2}\right) \frac{R}{M} T_{0}=\Delta m \frac{R}{M} T_{0}\)\\
or, \(\quad \Delta m=\frac{\left(p_{1}-p_{2}\right) V M}{R T_{0}}=\frac{\Delta p V M}{R T_{0}}\)

We also know \(p=\rho \frac{R}{M} T\) so, \(\frac{M}{R T_{0}}=\frac{\rho}{p_{0}}\)\\
(where \(p_{0}=\) standard atmospheric pressure and \(T_{0}=273 \mathrm{~K}\) )\\
From Eqs. (1) and (2) we get

\[
\Delta m=\rho V \cdot \frac{\Delta p}{p_{0}}=1.3 \times 30 \times \frac{0.78}{1}=30 \mathrm{~g}
\]

2.2 Let \(m_{1}\) be the mass of the gas enclosed.

Then,

\[
p_{1} V=v_{1} R T_{1}
\]

When heated, some gas, passes into the evacuated vessel till pressure difference becomes \(\Delta p\). Let \(p_{1}^{\prime}\) and \(p_{2}^{\prime}\) be the pressure on the two sides of the valve. Then \(p_{1}^{\prime} V=v_{1}^{\prime} R T_{2}\) and

\[
p_{2}^{\prime} V=v_{2}^{\prime} R T_{2}=\left(v_{1}-v_{1}^{\prime}\right) R T_{2}
\]

\[
p_{2}^{\prime} V=\left(\frac{p_{1} V}{R T_{1}}-\frac{p_{1}^{\prime} V}{R T_{2}}\right) \quad \text { or } \quad p_{2}^{\prime}=\left(\frac{p_{1}}{T_{1}}-\frac{p_{1}^{\prime}}{T_{2}}\right) T_{2}
\]

But,

\[
p_{1}^{\prime}-p_{2}^{\prime}=\Delta p
\]

So,

\[
\begin{aligned}
p_{2}^{\prime} & =\left(\frac{p_{1}}{T_{1}}-\frac{p_{2}^{\prime}+\Delta p}{T_{2}}\right) T_{2} \\
& =\frac{p_{1} T_{2}}{T_{1}}-p_{2}^{\prime}-\Delta p
\end{aligned}
\]

or,

\[
p_{2}^{\prime}=\frac{1}{2}\left(\frac{p_{1} T_{2}}{T_{1}}-\Delta p\right)=0.08 \mathrm{~atm}
\]

2.3 Let the mixture contain \(v_{1}\) and \(v_{2}\) moles of \(H_{2}\) and \(H_{e}\) respectively. If molecular weights of \(H_{2}\) and \(H_{\mathrm{e}}\) are \(M_{1}\) and \(M_{2}\), then respective masses in the mixture are equal to

\[
m_{1}=v_{1} M_{1} \text { and } m_{2}=v_{2} M_{2}
\]

Therefore, for the total mass of the mixture we get,


\begin{equation*}
m=m_{1}+m_{2} \text { or } m=v_{1} M_{1}+v_{2} M_{2} \tag{1}
\end{equation*}


Also, if \(v\) is the total number of moles of the mixture in the vessels, then we know,


\begin{equation*}
v=v_{1}+v_{2} \tag{2}
\end{equation*}


Solving (1) and (2) for \(v_{1}\) and \(v_{2}\), we get,

\[
v_{1}=\frac{\left(v M_{2}-m\right)}{M_{2}-M_{1}}, \quad v_{2}=\frac{m-v M_{1}}{M_{2}-M_{1}}
\]

Therefore, we get \(m_{1}=M_{1} \cdot \frac{\left(v M_{2}-m\right)}{M_{2}-M_{1}}\) and \(m_{2}=M_{2} \frac{\left(m-v M_{1}\right)}{M_{2}-M_{1}}\)\\
or,

\[
\frac{m_{1}}{m_{2}}=\frac{M_{1}}{M_{2}} \frac{\left(v M_{2}-m\right)}{\left(m-v M_{1}\right)}
\]

One can also express the above result in terms of the effective molecular weight \(M\) of the mixture, defined as,

\[
M=\frac{m}{v}=m \frac{R T}{p V}
\]

Thus,

\[
\frac{m_{1}}{m_{2}}=\frac{M_{1}}{M_{2}} \cdot \frac{M_{2}-M}{M-M_{1}}=\frac{1-M / M_{2}}{M / M_{1}-1}
\]

Using the data and table, we get :

\[
M=3.0 \mathrm{~g} \text { and }, \frac{m_{1}}{m_{2}}=0.50
\]

2.4 We know, for the mixture, \(\mathrm{N}_{2}\) and \(\mathrm{CO}_{2}\) (being regarded as ideal gases, their mixture too behaves like an ideal gas)

\[
p V=v R T, \quad \text { so } \quad p_{0} V=v R T
\]

where, \(v\) is the total number of moles of the gases (mixture) present and \(V\) is the volume of the vessel. If \(v_{1}\) and \(v_{2}\) are number of moles of \(N_{2}\) and \(C O_{2}\) respectively present in the mixture, then

\[
v=v_{1}+v_{2}
\]

Now number of moles of \(\mathrm{N}_{2}\) and \(\mathrm{CO}_{2}\) is, by definition, given by

\[
v_{1}=\frac{m_{1}}{M_{1}} \text { and, } v_{2}=\frac{m_{2}}{M_{2}}
\]

where, \(m_{1}\) is the mass of \(N_{2}\) (Moleculer weight \(=M_{1}\) ) in the mixture and \(m_{2}\) is the mass of \(\mathrm{CO}_{2}\) (Molecular weight \(=M_{2}\) ) in the mixture.\\
Therefore density of the mixture is given by

\[
\begin{gathered}
\rho=\frac{m_{1}+m_{2}}{V}=\frac{m_{1}+m_{2}}{\left(v R T / P_{0}\right)} \\
=\frac{p_{0}}{R T} \cdot \frac{m_{1}+m_{2}}{v_{1}+v_{2}}=\frac{p_{0}\left(m_{1}+m_{2}\right) M_{1} M_{2}}{R T\left(m_{1} M_{2}+m_{2} M_{1}\right)} \\
=1.5 \mathrm{~kg} / \mathrm{m}^{3} \text { on substitution }
\end{gathered}
\]

2.5 (a) The mixture contains \(v_{1}, v_{2}\) and \(v_{3}\) moles of \(\mathrm{O}_{2}, \mathrm{~N}_{2}\) and \(\mathrm{CO}_{2}\) respectively. Then the total number of moles of the mixture

\[
v=v_{1}+v_{2}+v_{3}
\]

We know, ideal gas equation for the mixture

\[
p V=v R T \quad \text { or } \quad p=\frac{v R T}{V}
\]

or, \(\quad p=\frac{\left(v_{1}+v_{2}+v_{3}\right) R T}{V}=1.968 \mathrm{~atm}\) on substitution\\
(b) Mass of oxygen ( \(O_{2}\) ) present in the mixture : \(m_{1}=v_{1} M_{1}\)

Mass of nitrogen \(\left(N_{2}\right)\) present in the mixture : \(m_{2}=v_{2} M_{2}\)\\
Mass of carbon dioxide ( \(C O_{2}\) ) present in the mixture : \(m_{3}=v_{3} M_{3}\)\\
So, mass of the mixture

\[
m=m_{1}+m_{2}+m_{3}=v_{1} M_{1}+v_{2} M_{2}+v_{3} M_{3}
\]

Moleculer mass of the mixture : \(M=\frac{\text { mass of the mixture }}{\text { total number of moles }}\)

\[
=\frac{v_{1} M_{1}+v_{2} M_{2}+v_{3} M_{3}}{v_{1}+v_{2}+v_{3}}=36.7 \mathrm{~g} / \mathrm{mol} . \text { on substitution }
\]

2.6 Let \(p_{1}\) and \(p_{2}\) be the pressure in the upper and lower part of the cylinder respectively at temperature \(T_{0}\). At the equilibrium position for the piston :\\
\(p_{1} S+m g=p_{2} S\) or, \(p_{1}+\frac{m g}{S}=p_{2}(m\) is the mass of the piston.)\\
But \(p_{1}=\frac{R T_{0}}{\eta V_{0}}\) (where \(V_{0}\) is the initial volume of the lower part)

So,


\begin{equation*}
\frac{R T_{o}}{\eta V_{0}}+\frac{m g}{S}=\frac{R T_{0}}{V_{0}} \text { or, } \frac{m g}{S}=\frac{R T_{0}}{V_{0}}\left(1-\frac{1}{\eta}\right) \tag{1}
\end{equation*}


Let \(T\) be the sought temperature and at this temperature the volume of the lower part becomes \(V^{\prime}\), then according to the problem the volume of the upper part becomes \(\eta^{\prime} V^{\prime}\)

Hence,


\begin{equation*}
\frac{m g}{S}=\frac{R T}{V^{\prime}}\left(1-\frac{1}{\eta^{\prime}}\right) \tag{2}
\end{equation*}


From (1) and (2).

\[
\frac{R T_{0}}{V_{0}}\left(1-\frac{1}{\eta}\right)=\frac{R T}{V^{\prime}}\left(1-\frac{1}{\eta^{\prime}}\right) \quad \text { or, } \quad T=\frac{T_{0}\left(1-\frac{1}{\eta}\right) V^{\prime}}{V_{0}\left(1-\frac{1}{\eta^{\prime}}\right)}
\]

As, the total volume must be constant,

\[
V_{0}(1+\eta)=V^{\prime}\left(1+\eta^{\prime}\right) \quad \text { or, } \quad V^{\prime}=\frac{V_{0}(1+\eta)}{\left(1+\eta^{\prime}\right)}
\]

Putting the value of \(V^{\prime}\) in Eq. (3), we get

\[
\begin{aligned}
& T=\frac{T_{0}\left(1-\frac{1}{\eta}\right) V_{0} \frac{(1+\eta)}{\left(1+\eta^{\prime}\right)}}{V_{0}\left(1-\frac{1}{\eta^{\prime}}\right)} \\
& =\frac{T_{0}\left(\eta^{2}-1\right) \eta^{\prime}}{\left(\eta^{\prime 2}-1\right) \eta}=0.42 \mathrm{kK}
\end{aligned}
\]

2.7 Let \(\rho_{1}\) be the density after the first stroke. The the mass remains constant

\[
V \rho=(V+\Delta V) \rho_{1}, \quad \text { or, } \quad \rho_{1}=\frac{V \rho}{(V+\Delta V)}
\]

Similarly, if \(\rho_{2}\) is the density after second stroke

\[
V \rho_{1}=(V+\Delta V) \rho_{2} \quad \text { or, } \quad \rho_{2}=\left(\frac{V}{V+\Delta V}\right) \rho_{1}=\left(\frac{V}{V+\Delta V}\right)^{2} \rho_{0}
\]

In this way after \(n\)th stroke.

\[
\rho_{n}=\left(\frac{V}{V+\Delta V}\right)^{n} \rho_{0}
\]

Since pressure \(\alpha\) density,

\[
p_{n}=\left(\frac{V}{V+\Delta V}\right)^{n} p_{0} \text { (because temperature is constant.) }
\]

It is required by \(\frac{p_{n}}{p_{0}}\) to be \(\frac{1}{\eta}\)\\
so, \(\quad \frac{1}{\eta}=\left(\frac{V}{V+\Delta V}\right)^{n}\) or, \(\eta=\left(\frac{V+\Delta V}{V}\right)^{n}\)

Hence

\[
n=\frac{\ln \eta}{\ln \left(1+\frac{\Delta V}{V}\right)}
\]

2.8 From the ideal gas equation \(p=\frac{m}{M} \frac{R T}{V}\)


\begin{equation*}
\frac{d p}{d t}=\frac{R T}{M V} \frac{d m}{d t} \tag{1}
\end{equation*}


In each stroke, volume \(v\) of the gas is ejected, where \(v\) is given by

\[
v=\frac{V}{m_{N}}\left[m_{N-1}-m_{N}\right]
\]

In case of continuous ejection, if ( \(m_{N-1}\) ) corresponds to mass of gas in the vessel at time \(t\), then \(m_{N}\) is the mass at time \(t+\Delta t\), where \(\Delta t\), is the time in which volume \(v\) of the gas has come out. The rate of evacuation is therefore \(\frac{v}{\Delta t}\) i.e.

\[
C=\frac{v}{\Delta t}=-\frac{V}{m(t+\Delta t)} \cdot \frac{m(t+\Delta t)-m(t)}{\Delta t}
\]

In the limit \(\Delta t-0\), we get


\begin{equation*}
C=\frac{V}{m} \frac{d m}{d t} \tag{2}
\end{equation*}


From (1) and (2)

\[
\frac{d p}{d t}=-\frac{C}{V} \frac{m R T}{M V}=-\frac{C}{V} p \quad \text { or } \quad \frac{d p}{p}=-\frac{C}{V} d t
\]

Integrating \(\int_{p}^{p_{0}} \frac{d p}{p}=-\frac{C}{V} \int_{t}^{0} d t\) or \(\ln \frac{p}{p_{0}}=-\frac{C}{V} t\)\\
Thus

\[
p=p_{0} e^{-C t / V}
\]

2.9 Let \(\rho\) be the instantaneous density, then instantaneous mass \(=V_{\rho}\). In a short interval \(d t\) the volume is increased by \(C d t\).\\
So, \(\quad V \rho=(V+C d t)(\rho+d \rho)\)\\
(because mass remains constant in a short interval \(d t\) )\\
so,

\[
\frac{d \rho}{\rho}=-\frac{C}{V} d t
\]

Since pressure \(\alpha\) density \(\frac{d p}{p}=-\frac{C}{V} d t\)\\
or

\[
\int_{p_{1}}^{p_{2}}-\frac{d p}{p}=\frac{C}{V} t,
\]

or

\[
t=\frac{V}{C} \ln \frac{p_{1}}{p_{2}}=\frac{V}{C} \ln \frac{1}{\eta} 1.0 \mathrm{~min}
\]

2.10 The physical system consists of one mole of gas confined in the smooth vertical tube. Let \(m_{1}\) and \(m_{2}\) be the masses of upper and lower pistons and \(S_{1}\) and \(S_{2}\) are their respective areas.\\
For the lower piston


\begin{equation*}
p S_{2}+m_{2} g=p_{0} S_{2}+T \tag{1}
\end{equation*}


or, \(\quad T=\left(p-p_{0}\right) S_{2}+m_{2} g\)

Similarly for the upper piston


\begin{equation*}
p_{0} S_{1}+T+m_{1} g=p S_{1}, \tag{2}
\end{equation*}


or, \(\quad T=\left(p-p_{0}\right) S_{1}-m_{1} g\)

From (1) and (2)

\[
\left(p-p_{0}\right)\left(S_{1}-S_{2}\right)=\left(m_{1}+m_{2}\right) g
\]

or, \(\quad\left(p-p_{0}\right) \Delta S=m g\)\\
\includegraphics[max width=\textwidth, alt={}, center]{9f3c5a8e-ac5d-42f4-8eb6-ea2af623f2d8-196_452_225_798_900}\\
so, \(\quad p=\frac{m g}{\Delta S}+p_{0}=\) constant\\
From the gas law, \(p V=v R T\)

\[
p \Delta V=v R \Delta T \text { (because } p \text { is constant) }
\]

So,

\[
\left(p_{0}+\frac{m g}{\Delta S}\right) \Delta S l=R \Delta T
\]

Hence,

\[
\Delta T=\frac{1}{R}\left(p_{0} \Delta S+m g\right) l=0.9 \mathrm{~K}
\]

2.11 (a) \(p=p_{0}-\alpha V^{2}=p_{0}-\alpha\left(\frac{R T}{p}\right)^{2}\)\\
(as, \(V=R T / p\) for one mole of gas)\\
Thus,


\begin{equation*}
T=\frac{1}{R \sqrt{\alpha}} p \sqrt{p_{0}-p}=\frac{1}{R \sqrt{\alpha}} \sqrt{p_{0} p^{2}-p^{3}} \tag{1}
\end{equation*}


For \(\mathrm{T}_{\max }, \quad \frac{d}{d p}\left(p_{0} p^{2}-p^{3}\right)\) must be zero\\
which yields,


\begin{equation*}
p=\frac{2}{3} p_{0} \tag{2}
\end{equation*}


Hence, \(\quad T_{\max }=\frac{1}{R \sqrt{\alpha}} \cdot \frac{2}{3} p_{0} \sqrt{p_{0}-\frac{2}{3} p_{0}}=\frac{2}{3}\left(\frac{p_{0}}{R}\right) \sqrt{\frac{p_{0}}{3 \alpha}}\)\\
(b) \(p=p_{0} e^{-\beta V}=p_{0} e^{-\beta R T / p}\)\\
so


\begin{equation*}
\frac{\beta R T}{p}=\ln \frac{p_{o}}{p}, \text { and } T=\frac{p}{\beta R} \ln \frac{p_{0}}{p} \tag{1}
\end{equation*}


For \(T_{\max }\) the condition is \(\frac{d T}{d p}=0\), which yields

\[
p=\frac{p_{0}}{e}
\]

Hence using this value of \(p\) in Eq. (1), we get

\[
T_{\max }=\frac{p_{o}}{e \beta R}
\]

\(2.12 T=T_{0}+\alpha V^{2}=T_{0}+\alpha \frac{R^{2} T^{2}}{p^{2}}\)

\[
\text { (as, } V=R T / p \text { for one mole of gas) }
\]

So,


\begin{equation*}
p=\sqrt{\alpha} R T\left(T-T_{0}\right)^{1 / 2} \tag{1}
\end{equation*}


For

\[
p_{\min }, \frac{d p}{d T}=0, \text { which gives }
\]


\begin{equation*}
T=2 T_{0} \tag{2}
\end{equation*}


From (1) and (2), we get,

\[
p_{\min }=\sqrt{\alpha} R 2 T_{0}\left(2 T_{0}-T_{0}\right)^{-1 / 2}=2 R \sqrt{\alpha T_{0}}
\]

2.13 Consider a thin layer at a height \(h\) and thickness \(d h\). Let \(p\) and \(d p+p\) be the pressure on the two sides of the layer. The mass of the layer is Sdhp. Equating vertical downward force to the upward force acting on the layer.


\begin{equation*}
S d h \rho g+(p+d p) S=p S \tag{1}
\end{equation*}


So, \(\quad \frac{d p}{d h}=-\rho g\)

But, \(p=\frac{\rho}{M} R T\), we have \(d p=\frac{\rho R}{M} d T\),

\[
\text { or, } \quad-\frac{\rho R}{M} d T=\prime \rho g d h
\]

\begin{center}
\includegraphics[max width=\textwidth, alt={}]{9f3c5a8e-ac5d-42f4-8eb6-ea2af623f2d8-197_328_323_1569_964}
\end{center}

So, \(\quad \frac{d T}{d h}=-\frac{g M}{R}=-34 \mathrm{~K} / \mathrm{km}\)\\
That means, temperature of air drops by \(34^{\circ} \mathrm{C}\) at a height of 1 km above bottom.\\
2.14 We have, \(\frac{d p}{d h}=-\rho g\) (See 2.13)

But, from \(\quad p=C \rho^{n}\) (where \(C\) is, \(a\) const) \(\frac{d p}{d \rho}=C n \rho^{n-1}\)

We have from gas low \(\quad p=\rho \frac{R}{M} T\), so using (2)

Thus,


\begin{gather*}
C \rho^{n}=\rho \frac{R}{M} \cdot T, \quad \text { or } \quad T=\frac{M}{R} C \rho^{n-1} \\
\frac{d T}{d \rho}=\frac{M}{R} \cdot C(n-1) \rho^{n-2}  \tag{3}\\
\frac{d T}{d h}=\frac{d T}{d \rho} \cdot \frac{d \rho}{d p} \cdot \frac{d p}{d h}
\end{gather*}


But,

\[
\frac{d T}{d h}=\frac{M}{R} C(n-1) \rho^{n-2} \frac{1}{C n \rho^{n-1}}(-\rho g)=\frac{-M g(n-1)}{n R}
\]

2.15 We have, \(d p=-\rho g d h\) and from gas law \(\rho=\frac{M}{R T} p\)

Thus


\begin{equation*}
\frac{d p}{p}=-\frac{M g}{R T} d h \tag{1}
\end{equation*}


Integrating, we get\\
or, \(\quad \int_{p_{0}}^{p} \frac{d p}{p}=-\frac{M g}{R T} \int_{0}^{h} d h \quad\) or, \(\quad \ln \frac{p}{p_{0}}=-\frac{M g}{R T} h\),\\
(where \(p_{0}\) is the pressure at the surface of the Earth.)

\[
p=p_{0} e^{-M g h / R T}
\]

[Under standard condition, \(p_{0}=1 \mathrm{~atm}, T=273 \mathrm{~K}\)\\
Pressure at a height of \(5 \mathrm{~atm}=1 \times e^{-28 \times 9.81 \times 5000 / 8314 \times 273}=0.5 \mathrm{~atm}\).\\
Pressure in a mine at a depth of \(5 \mathrm{~km}=1 \times e^{-28 \times 9.81 \times(-5000) / 8314 \times 273}=2 \mathrm{~atm}\).]\\
2.16 We have \(d p=-\rho g d h\) but from gas law \(p=\frac{\rho}{M} R T\),

Thus

\[
d p=\frac{d \rho}{M} R T \text { at const. temperature }
\]

So,

\[
\frac{d \rho}{\rho}=\frac{g M}{R T} d h
\]

Integrating within limits \(\int_{\rho_{0}}^{\rho} \frac{d \rho}{\rho}=\int_{0}^{h} \frac{g M}{R T} d h\)\\
or,

\[
\ln \frac{\rho}{\rho_{0}}=-\frac{g M}{R T} h
\]

So, \(\rho=\rho_{0} e^{-M g h / R T}\) and \(h=-\frac{R T}{M g} \ln \frac{\rho}{\rho_{0}}\)\\
(a) Given \(T=273^{\circ} \mathrm{K}, \frac{\rho_{0}}{\rho}=e\)

Thus

\[
h=-\frac{R T}{M g} \ln e^{-1}=8 \mathrm{~km}
\]

(b) \(T=273^{\circ} \mathrm{K}\) and

\[
\frac{\rho_{0}-\rho}{\rho_{0}}=0.01 \text { or } \frac{\rho}{\rho_{0}}=0.99
\]

Thus

\[
h=-\frac{R T}{M g} \ln \frac{\rho}{\rho_{0}}=0.09 \mathrm{~km} \text { on substitution }
\]

2.17 From the Barometric formula, we have

\[
p=p_{0} e^{-M g h / R T}
\]

and from gas law

\[
\rho=\frac{p M}{R T}
\]

So, at constant temperature from these two Eqs.


\begin{equation*}
\rho=\frac{M p_{0}}{R T} e^{-M g h / R T}=\rho_{0} e^{-M g h / R T} \tag{1}
\end{equation*}


Eq. (1) shows that density varies with height in the same manner as pressure. Let us consider the mass element of the gas contained in the coltumn.

\[
d m=\rho(S d h)=\frac{M p_{0}}{R T} e^{-M g h / R T} S d h
\]

Hence the sought mass,

\[
m=\frac{M p_{0} S}{R T} \int_{0}^{h} e^{-M g h / R T} d h=\frac{p_{0} S}{g}\left(1-e^{-M g h / R T}\right)
\]

2.18 As the gravitational field is constant the centre of gravity and the centre of mass are same. The location of C.M.

\[
h=\frac{\int_{0}^{\infty} h d m}{\int_{0}^{\infty} d m}=\frac{\int_{0}^{\infty} h \rho d h}{\int_{0}^{\infty} \rho d h}
\]

But from Barometric formula and gas law \(\rho=\rho_{0} e^{-M g h / R T}\)

So,

\[
h=\frac{\int_{0}^{\alpha} h\left(e^{-M g h / R T}\right) d h}{\int_{0}^{\alpha}\left(e^{-M g h / R T}\right) d h}=\frac{R T}{M g}
\]

2.19 (a) We know that the variation of pressure with height of a fluid is given by :

\[
d p=-\rho g d h
\]

But from gas law

\[
p=\frac{\rho}{M} R T \quad \text { or, } \quad \rho=\frac{p M}{R T}
\]

From these two Eqs.


\begin{equation*}
d p=-\frac{p M g}{R T} d h \tag{1}
\end{equation*}


or, \(\quad \frac{d p}{p}=\frac{-M g d h}{R T_{0}(1-a h)}\)\\
Integrating, \(\quad \int_{p_{0}}^{p} \frac{d p}{p}=\frac{-M g}{R T_{0}} \int_{0}^{h} \frac{d h}{(1-a h)^{2}}\), we get

\[
\ln \frac{p}{p_{0}}=\ln (1-a h)^{M g / a R T_{0}}
\]

Hence, \(\quad \quad p=p_{0}(1-a h)^{M g / a R T_{0}}\). Obvionsly \(h<\frac{1}{a}\)\\
(b) Proceed up to Eq. (1) of part (a), and then put \(T=T_{0}(1+a h)\) and proceed further in the same fashion to get

\[
p=\frac{p_{0}}{(1+a h)^{M g / a R T_{0}}}
\]

2.20 Let us consider the mass element of the gas (thin layer) in the cylinder at a distance \(r\) from its open end as shown in the figure.\\
Using Newton's second law for the element

\[
\begin{gathered}
F_{n}=m w_{n}: \\
(p+d p) S-p S=(\rho S d r) \omega^{2} r
\end{gathered}
\]

\includegraphics[max width=\textwidth, alt={}, center]{9f3c5a8e-ac5d-42f4-8eb6-ea2af623f2d8-200_294_393_1565_839}\\
or, \(\quad d p=\rho \omega^{2} r d r=\frac{p M}{R T} \omega^{2} r d r\)

So, \(\quad \frac{d p}{p}=\frac{M \omega^{2}}{R T} r d r\) or, \(\int_{p_{0}}^{p} \frac{d p}{p}=\frac{M \omega^{2}}{R T} \int_{0}^{r} r d r\),\\
Thus, \(\quad \ln \frac{p}{p_{0}}=\frac{M \omega^{2}}{2 R T} r^{2}\) or, \(p=p_{0} e^{M \omega^{2} r^{2} / 2 R T}\)\\
2.21 For an ideal gas law

\[
p=\frac{\rho}{M} R T
\]

So,

\[
p=0.082 \times 300 \times \frac{500}{44} \mathrm{atms}=279.5 \text { atmosphere }
\]

For Vander Waal gas Eq.\\
\(\left(p+\frac{v^{2} a}{V^{2}}\right)(V-v b)=v R T, \quad\) where \(V=v V_{M}\)\\
or,

\[
\begin{gathered}
p=\frac{v R T}{V-v b}-\frac{a v^{2}}{V^{2}}=\frac{m R T / M}{V-\frac{m b}{M}}-\frac{a m^{2}}{V^{2} M^{2}} \\
=\frac{\rho R T}{M-\rho b}-\frac{a \rho^{2}}{M^{2}}=79.2 \mathrm{~atm}
\end{gathered}
\]

2.22 (a) \(p=\left[\frac{R T}{V_{M}-b}-\frac{a}{V_{M}^{2}}\right](1+\eta)=\frac{R T}{V_{M}}\)\\
(The pressure is less for a Vander Waal gas than for an ideal gas)\\
or, \(\quad \frac{a(1+\eta)}{V_{M}^{2}}=R T\left[\frac{-1}{V_{M}}+\frac{1+\eta}{V_{M}-b}\right]=R T \frac{\eta V_{M}+b}{V_{M}\left(V_{M}-b\right)}\)\\
or, \(\quad T=\frac{a(1+\eta)\left(V_{M}-b\right)}{R V_{M}\left(\eta V_{M}+b\right)}\), (here \(V_{M}\) is the molar volume.)

\[
=\frac{1.35 \times 1.1 \times(1-0.039)}{0.082 \times(0.139)} \approx 125 \mathrm{~K}
\]

(b) The corresponding pressure is

\[
\begin{aligned}
p & =\frac{R T}{V_{M}-b}-\frac{a}{V_{M}^{2}}=\frac{a(1+\eta)}{V_{M}\left(\eta V_{M}+b\right)}-\frac{a}{V_{M}^{2}} \\
& =\frac{a}{V_{M}^{2}} \frac{\left(V_{M}+\eta V_{M}-\eta V_{M}-b\right)}{\left(\eta V_{M}+b\right)}=\frac{a}{V_{M}^{2}} \frac{\left(V_{M}-b\right)}{\left(V_{M}+b\right)} \\
& =\frac{1.35}{1} \times \frac{0.961}{0.139} \approx 9.3 \mathrm{~atm}
\end{aligned}
\]

\(2.23 p_{1}=R T_{1} \frac{1}{V-b}-\frac{a}{V^{2}}, p_{2}=R T_{2} \frac{1}{V-b}-\frac{a}{V^{2}}\)

So,

\[
p_{2}-p_{1}=\frac{R\left(T_{2}-T_{1}\right)}{V-b}
\]

or.

\[
V-b=\frac{R\left(T_{2}-T_{1}\right)}{p_{2}-p_{1}} \text { or, } b=V-\frac{R\left(T_{2}-T_{1}\right)}{p_{2}-p_{1}}
\]

Also,

\[
p_{1}=T_{1} \frac{p_{2}-p_{1}}{T_{2}-T_{1}}-\frac{a}{V^{2}}
\]

\[
\frac{a}{V^{2}}=\frac{T_{1}\left(p_{2}-p_{1}\right)}{T_{2}-T_{1}}-p_{1}=\frac{T_{1} p_{2}-p_{1} T_{2}}{T_{2}-T_{1}}
\]

or,

\[
a=V^{2} \frac{T_{1} p_{2}-p_{1} T_{2}}{T_{2}-T_{1}}
\]

Using \(T_{1}=300 \mathrm{~K}, p_{1}=90 \mathrm{atms}, T_{2}=350 \mathrm{~K}, p_{2}=110 \mathrm{~atm}, V=0.250\) litre

\[
a=1.87 \mathrm{~atm} . \mathrm{litre}^{2} / \mathrm{mole}^{2}, b=0.045 \text { litre } / \mathrm{mole}
\]

\(2.24 p=\frac{R T}{V-b}-\frac{a}{V^{2}}-V\left(\frac{\partial p}{\partial V}\right)_{T}=\frac{R T V}{(V-b)^{2}}-\frac{2 a}{V^{2}}\)\\
or,

\[
\kappa=\frac{-1}{V}\left(\frac{\partial V}{\partial p}\right)_{T}
\]

\[
=\left[\frac{R T V^{3}-2 a(V-b)^{2}}{V^{2}(V-b)^{2}}\right]^{-1}=\frac{V^{2}(V-b)}{\left[R T V^{3}-2 a(V-b)^{2}\right]}
\]

2.25 For an ideal gas \(\kappa_{0}=\frac{V}{R T}\)

Now \(\kappa=\frac{(V-b)^{2}}{R T V}\left\{1-\frac{2 a(V-b)^{2}}{R T V^{3}}\right\}^{-1}=\kappa_{0}\left(1-\frac{b}{V}\right)^{2}\left\{1-\frac{2 a}{R T V}\left(1-\frac{b}{V}\right)^{2}\right\}^{-1}\)

\[
=\kappa_{0}\left\{1-\frac{2 b}{V}+\frac{2 a}{R T V}\right\}, \text { to leading order in } a, b
\]

Now

\[
\kappa>\kappa_{0} \text { if } \frac{2 a}{R T V}>\frac{2 b}{V} \text { or } T<\frac{a}{b R}
\]

If \(a, b\) do not vary much with temperature, then the effect at high temperature is clearly determined by \(b\) and its effect is repulsive so compressibility is less.

\subsection*{2.2 THE FIRST LAW OF THERMODYNAMICS. HEAT CAPACITY}
2.26 Internal energy of air, treating as an ideal gas


\begin{equation*}
U=\frac{m}{M} C_{V} T=\frac{m}{M} \frac{R}{\gamma-1} T=\frac{p V}{\gamma-1} \tag{1}
\end{equation*}


Using \(\quad C_{V}=\frac{R}{\gamma-1}\), since \(C_{p}-C_{V}=R\) and \(\frac{C_{p}}{C_{V}}=\gamma\)\\
Thus at constant pressure \(U=\) constant, because the volume of the room is a constant.\\
Puting the value of \(p=p_{\text {atm }}\) and \(V\) in Eq. (1), we get \(U=10 \mathrm{MJ}\).\\
2.27 From energy conservation

\[
U_{t}+\frac{1}{2}(v M) v^{2}=U_{f}
\]

or,


\begin{equation*}
\Delta U=\frac{1}{2} v M v^{2} \tag{1}
\end{equation*}


But from \(U=v \frac{R T}{\gamma-1}, \quad \Delta U=\frac{v R}{\gamma-1} \Delta T\) (trom the previous problem)

Hence from Eqs. (1) and (2).

\[
\Delta T=\frac{M v^{2}(\gamma-1)}{2 R}
\]

2.28 On opening the valve, the air will flow from the vessel at heigher pressure to the vessel at lower pressure till both vessels have the same air pressure. If this air pressure is \(p\), the total volume of the air in the two vessels will be \(\left(V_{1}+V_{2}\right)\). Also if \(v_{1}\) and \(v_{2}\) be the number of moles of air initially in the two vessels, we have


\begin{equation*}
p_{1} V_{1}=v_{1} R T_{1} \text { and } p_{2} V_{2}=v_{2} R T_{2} \tag{1}
\end{equation*}


After the air is mixed up, the total number of moles are ( \(v_{1}+v_{2}\) ) and the mixture is at temperature \(T\).\\
Hence


\begin{equation*}
p\left(V_{1}+V_{2}\right)=\left(v_{1}+v_{2}\right) R T \tag{2}
\end{equation*}


Let us look at the two portions of air as one single system. Since this system is contained in a thermally insulated vessel, no heat exchange is involved in the process. That is, total heat transfer for the combined system \(Q=0\)\\
Moreover, this combined system does not perform mechanical work either. The walls of the containers are rigid and there are no pistons etc to be pushed, looking at the total system, we know \(A=0\).\\
Hence, internal energy of the combined system does not change in the process. Initially energy of the combined system is equal to the sum of internal energies of the two portions of air :


\begin{equation*}
U_{t}=U_{1}+U_{2}=\frac{v_{1} R T_{1}}{\gamma-1}+\frac{v_{2} R T_{2}}{\gamma-1} \tag{3}
\end{equation*}


Final internal energy of ( \(\boldsymbol{n}_{1}+\boldsymbol{n}_{2}\) ) moles of air at temperature \(\boldsymbol{T}\) is given by


\begin{equation*}
U_{f}=\frac{\left(v_{1}+v_{2}\right) R T}{\gamma-1} \tag{4}
\end{equation*}


Therefore, \(U_{i}=U_{f}\) implies :

\[
T=\frac{v_{1} T_{1}+v_{2} T_{2}}{v_{1}+v_{2}}=\frac{p_{1} V_{1}+p_{2} V_{2}}{\left(p_{1} V_{1} / T_{1}\right)+\left(p_{2} V_{2} / T_{2}\right)}=T_{1} T_{2} \frac{p_{1} V_{1}+p_{2} V_{2}}{p_{1} V_{1} T_{2}+p_{2} V_{2} T_{2}}
\]

From (2), therefore, final pressure is given by :

\[
p=\frac{v_{1}+v_{2}}{V_{1}+V_{2}} R T=\frac{R}{V_{1}+V_{2}}\left(v_{1} T_{1}+v_{2} T_{2}\right)=\frac{p_{1} V_{1}+p_{2} V_{2}}{V_{1}+V_{2}}
\]

This process in an example of free adiabatic expansion of ideal gas.\\
2.29 By the first law of thermodynamics,

\[
Q=\Delta U+A
\]

Here \(A=0\), as the volume remains constant,\\
So,

\[
Q=\Delta U=\frac{v R}{\gamma-1} \Delta T
\]

From gas law,

\[
p_{0} V=v R T_{0}
\]

So, \(\quad \Delta U=\frac{p_{0} V \Delta T}{T_{0}(\gamma-1)}=-0.25 \mathrm{~kJ}\)\\
Hence amount of heat lost \(=-\Delta U=0.25 \mathrm{~kJ}\)\\
2.30 By the first law of thermodynamics \(Q=\Delta U+A\)

But

\[
\begin{gathered}
\Delta U=\frac{p \Delta V}{\gamma-1}=\frac{A}{\gamma-1}(\text { as } p \text { is constant }) \\
Q=\frac{A}{\gamma-1}+A=\frac{\gamma \cdot A}{\gamma-1}=\frac{1 \cdot 4}{1 \cdot 4-1} \times 2=7 \mathrm{~J}
\end{gathered}
\]

2.31 Under isobaric process \(A=p \Delta V=R \Delta T\) (as \(v=1\) ) \(=0.6 \mathrm{~kJ}\)

From the first law of thermodynamics

\[
\Delta U=Q-A=Q-R \Delta T=1 \mathrm{~kJ}
\]

Again increment in internal energy \(\Delta U=\frac{R \Delta T}{\gamma-1}\), for \(v=1\)\\
Thus

\[
Q-R \Delta T=\frac{R \Delta T}{\gamma-1} \quad \text { or } \quad \gamma=\frac{Q}{Q-R \Delta T}=1.6
\]

2.32 Let \(v=2\) moles of the gas. In the first phase, under isochoric process, \(A_{1}=0\), therefore from gas law if pressure is reduced \(n\) times so that temperature i.e. new temperature becomes \(T_{0} / n\).\\
Now from first law of thermodynamics

\[
Q_{1}=\Delta U_{1}=\frac{v R \Delta T}{\gamma-1}
\]

\[
=\frac{v R}{\gamma-1}\left(\frac{T_{o}}{n}-T_{0}\right)=\frac{v R T_{0}(1-n)}{n(\gamma-1)}
\]

During the second phase (under isobaric process),

\[
A_{2}=p \Delta V=v R \Delta T
\]

Thus from first law of thermodynamics :

\[
\begin{aligned}
& Q_{2}=\Delta U_{2}+A_{2}=\frac{v R \Delta T}{\gamma-1}+v R \Delta T \\
& =\frac{v R\left(T_{0}-\frac{T_{0}}{n}\right) \gamma}{\gamma-1}=\frac{v R T_{0}(n-1) \gamma}{n(\gamma-1)}
\end{aligned}
\]

Hence the total amount of heat absorbed

\[
\begin{aligned}
Q & =Q_{1}+Q_{2}=\frac{v R T_{0}(1-n)}{n(\gamma-1)}+\frac{v R T_{0}(n-1) \gamma}{n(\gamma-1)} \\
& =\frac{v R T_{0}(n-1) \gamma}{n(\gamma-1)}(-1+\gamma)=v R T_{0}\left(1-\frac{1}{n}\right)
\end{aligned}
\]

2.33 Total no. of moles of the mixture \(v=v_{1}+v_{2}\)

At a certain temperature, \(U=U_{1}+U_{2}\) or \(v C_{V}=v_{1} C_{V_{1}}+v_{2} C_{V_{2}}\)

Thus

Similarly

\[
C_{V}=\frac{v_{1} C_{V_{1}}+v_{2} C_{V_{2}}}{v}=\frac{\left(v_{1} \frac{R}{\gamma_{1}-1}+v_{2} \frac{R}{\gamma_{2}-1}\right)}{v}
\]

Thus

\[
\begin{gathered}
C_{p}=\frac{v_{1} C_{p_{1}}+v_{2} C_{p_{2}}}{v} \\
=\frac{v_{1} \gamma_{1} C_{V_{1}}+v_{2} \gamma_{2} C_{V_{2}}}{v}=\frac{\left(v_{1} \frac{\gamma_{1} R}{\gamma_{1}-1}+v_{2} \frac{\gamma_{2} R}{\gamma_{2}-1}\right)}{v} \\
\gamma=\frac{C_{p}}{C_{V}}=\frac{v_{1} \frac{\gamma_{1}}{\gamma_{1}-1} R+v_{2} \frac{\gamma_{2}}{\gamma_{2}-1} R}{v_{1} \frac{R}{\gamma_{1}-1}+v_{2} \frac{R}{\gamma_{2}-1}} \\
=\frac{v_{1} \gamma_{1}\left(\gamma_{2}-1\right)+v_{2} \gamma_{2}\left(\gamma_{1}-1\right)}{v_{1}\left(\gamma_{2}-1\right)+v_{2}\left(\gamma_{1}-1\right)}
\end{gathered}
\]

2.34 From the previous problem

\[
C_{V}=\frac{v_{1} \frac{R}{\gamma_{1}-1}+v_{2} \frac{R}{\gamma_{2}-1}}{v_{1}+v_{2}}=15.2 \mathrm{~J} / \text { mole. } \mathrm{K}
\]

and

\[
C_{p}=\frac{v_{1} \frac{\gamma_{1} R}{\gamma_{1}-1}+v_{2} \frac{\gamma_{2} R}{\gamma_{2}-1}}{v_{1}+v_{2}}=23.85 \mathrm{~J} / \mathrm{mole} . \mathrm{K}
\]

Now molar mass of the mixture \((M)=\frac{\text { Total mass }}{\text { Total number of moles }}=\frac{20+7}{\frac{1}{2}+\frac{1}{4}}=36\)

Hence

\[
c_{V}=\frac{C_{V}}{M}=0.42 \mathrm{~J} / \mathrm{g} \cdot \mathrm{~K} \text { and } c_{p}=\frac{C_{p}}{M}=0.66 \mathrm{~J} / \mathrm{g} \cdot \mathrm{~K}
\]

2.35 Let \(S\) be the area of the piston and \(F\) be the force exerted by the external agent. Then, \(F+p S=p_{0} S\) (Fig.) at an arbitrary instant of time. Here \(p\) is the pressure at the instant the volume is \(V\). (Initially the pressure inside is \(p_{0}\) )

\[
\begin{aligned}
& A(\text { Work done by the agent })=\int_{V_{0}}^{\eta V_{0}} F d x \\
& =\int_{V_{0}}^{\eta V_{0}}\left(p_{0}-p\right) S \cdot d x=\int_{V_{0}}^{\eta V_{0}}\left(p_{0}-p\right) d V \\
& =p_{0}(\eta-1) V_{0}-\int_{V_{0}} p d V=p_{0}(\eta-1) V_{0}-\int_{V_{0}}^{\eta V_{0}} v R T \cdot \frac{d V}{V} \\
& =(\eta-1) p_{0} V_{0}-n R T \ln \eta=(\eta-1) v R T-v R T \ln \eta \\
& =v R T(\eta-1-\ln \eta)=R T(\eta-1-\ln \eta)(\text { For } v=1 \text { mole })
\end{aligned}
\]

\includegraphics[max width=\textwidth, alt={}, center]{9f3c5a8e-ac5d-42f4-8eb6-ea2af623f2d8-206_403_298_771_979}\\
2.36 Let the agent move the piston to the right by \(x\). In equilibirium position,

\[
p_{1} S+F_{\text {agent }}=p_{2} S, \quad \text { or, } \quad F_{\text {agent }}=\left(p_{2}-p_{1}\right) S
\]

Work done by the agent in an infinitesmal change \(d x\) is

\[
F_{\text {agent }} \cdot d x=\left(p_{2}-p_{1}\right) S d x=\left(p_{2}-p_{1}\right) d V
\]

By applying \(p V=\) constant, for the two parts,

\[
p_{1}\left(V_{0}+S x\right)=p_{0} V_{0} \text { and } p_{2}\left(V_{0}-S x\right)=p_{0} V_{0}
\]

So,

\[
p_{2}-p_{1}=\frac{p_{0} V_{0} 2 S x}{V_{0}^{2}-S^{2} x^{2}}=\frac{2 p_{0} V_{0} V}{V_{0}^{2}-V^{2}}(\text { where } S x=V)
\]

When the volume of the left end is \(\eta\) times the volume of the right end

\[
\left(V_{0}+V\right)=\eta\left(V_{0}-V\right), \text { or, } \quad V=\frac{\eta-1}{\eta+1} V_{0}
\]

\[
\begin{gathered}
A=\int_{0}^{V}\left(p_{2}-p_{1}\right) d V=\int_{0}^{V} \frac{2 p_{0} V_{0} V}{V_{0}^{2}-V^{2}} d V=-p_{0} V_{0}\left[\ln \left(V_{0}^{2}-V^{2}\right)\right]_{0}^{V} \\
=-p_{0} V_{0}\left[\ln \left(V_{0}^{2}-V^{2}\right)-\ln V_{0}^{2}\right] \\
=-p_{0} V_{0}\left[\ln \left\{V_{0}^{2}-\left(\frac{\eta-1}{\eta+1^{2}}\right) V_{0}^{2}\right\}-\ln V_{0}^{2}\right] \\
=-p_{0} V_{0}\left(\ln \frac{4 \eta}{(\eta+1)^{2}}\right)=p_{0} V_{0} \ln \frac{(\eta+1)^{2}}{4 \eta}
\end{gathered}
\]

2.37 In the isothermal process, heat transfer to the gas is given by

\[
Q_{1}=v R T_{0} \ln \frac{V_{2}}{V_{1}}=v R T_{0} \ln \eta \quad\left(\text { For } \eta=\frac{V_{2}}{V_{1}}=\frac{p_{1}}{p_{2}}\right)
\]

In the isochoric process, \(A=0\)\\
Thus heat transfer to the gas is given by

\[
Q_{2}=\Delta U=v C_{V} \Delta T=\frac{v R}{\gamma-1} \Delta T \quad\left(\text { for } \quad C_{V}=\frac{R}{\gamma-1}\right)
\]

But \(\quad \frac{p_{2}}{p_{1}}=\frac{T_{0}}{T}\), or, \(T=T_{0} \frac{p_{1}}{p_{2}}=\eta T_{0}\left(\right.\) for \(\left.\eta=\frac{p_{1}}{p_{2}}\right)\)\\
or, \(\quad \Delta T=\eta T_{0}-T_{0}=(\eta-1) T_{0}\) so, \(Q_{2}=\frac{v R}{\gamma-1} \cdot(\eta-1) T_{0}\)\\
Thus, net heat transfer to the gas

\[
Q=v R T_{0} \ln \eta+\frac{v R}{\gamma-1} \cdot(\eta-1) T_{0}
\]

or,

\[
\frac{Q}{v R T_{0}}=\ln \eta+\frac{\eta-1}{\gamma-1}, \quad \text { or }, \quad \frac{Q}{v R T_{0}}-\ln \eta=\frac{\eta-1}{\gamma-1}
\]

or, \(\quad \gamma=1+\frac{\eta-1}{\frac{Q}{v R T_{0}}-\ln \eta}=1+\frac{6-1}{\left(\frac{80 \times 10^{3}}{3 \times 8.314 \times 273}\right)-\ln 6}=1.4\)\\
2.38 (a) From ideal gas law \(p=\left(\frac{v R}{V}\right) T=k T\) (where \(k=\frac{v R}{V}\) )

For isochoric process, obviously \(k=\) constant, thus \(p=k T\), represents a straight line passing through the origin and its slope becomes \(k\).\\
For isobaric process \(p=\) constant, thus on \(p-T\) curve, it is a horizontal straight line parallel to \(T\)-axis, if \(T\) is along horizontal (or \(x\)-axis)\\
For isothermal process, \(T=\) constant, thus on \(p-T\) curve, it represents a vertical straight line if \(T\) is taken along horizontal (or \(x\)-axis)\\
For adiabatic process \(T^{\gamma} p^{1-\gamma}=\) constant\\
After diffrentiating, we get \((1-\gamma) p^{-\gamma} d p \cdot T^{\gamma}+\gamma p^{1-\gamma} \cdot T^{\gamma-1} \cdot d T=0\)

\[
\frac{d p}{d T}=\left(\frac{\gamma}{1-\gamma}\right)\left(\frac{p^{1-\gamma}}{p^{-\gamma}}\right)\left(\frac{T^{\gamma-1}}{T^{\gamma}}\right)=\left(\frac{\gamma}{\gamma-1}\right) \frac{p}{T}
\]

The approximate plots of isochoric, isobaric, isothermal, and adiabatic processess are drawn in the answersheet.\\
(b) As \(p\) is not considered as variable, we have from ideal gas law

\[
V=\frac{v R}{p} T=k^{\prime} T\left(\text { where } k^{\prime}=\frac{v R}{p}\right)
\]

On \(V-T\) co-ordinate system let us, take \(T\) along \(x\)-axis.\\
For isochoric process \(V=\) constant, thus \(k^{\prime}=\) constant and \(V=k^{\prime} T\) obviously represents a straight line pasing through the origin of the co-ordinate system and \(k^{\prime}\) is its slope.\\
For isothermal process \(T=\) constant. Thus on the stated co- ordinate system it represents a straight line parallel to the \(V\)-axis.\\
For adiabatic process \(T V^{\gamma-1}=\) constant\\
After differentiating, we get \((\gamma-1) V^{\gamma-2} d V \cdot T+V^{\mu-1} d T=0\)

\[
\frac{d V}{d T}=-\left(\frac{1}{\gamma-1}\right) \cdot \frac{V}{T}
\]

The approximate plots of isochoric, isobaric, isothermal and adiabatic processess are drawn in the answer sheet.\\
2.39 According to \(T-p\) relation in adiabatic process, \(T^{\gamma}=k p^{\gamma-1}\) (where \(k=\) constant)\\
and

\[
\left(\frac{T_{2}}{T_{1}}\right)^{\gamma}=\left(\frac{p_{2}}{p_{1}}\right)^{\gamma-1} \text { So, } \frac{T^{\gamma}}{T_{0}^{\gamma}}=\eta^{\gamma-1}\left(\text { for } \eta=\frac{p_{2}}{p_{1}}\right)
\]

Hence

\[
T=T_{0} \cdot \eta \frac{\gamma-1}{\gamma}=290 \times 10^{(1.4-1) 1.4}=0.56 \mathrm{kK}
\]

(b) Using the solution of part (a), sought work done

\[
A=\frac{v R \Delta T}{\gamma-1}=\frac{v R T_{0}}{\gamma-1}\left(\eta^{(\gamma-1)_{\gamma}}-1\right)=5.61 \mathrm{~kJ} \text { (on substitution) }
\]

2.40 Let \(\left(p_{0}, V_{0}, T_{0}\right)\) be the initial state of the gas.

We know \(A_{\text {adia }}=\frac{-v R \Delta T}{\gamma-1}\) (work done by the gas)\\
But from the equation \(T V^{\gamma-1}=\) constant, we get \(\Delta T=T_{0}\left(\eta^{\gamma-1}-1\right)\)

Thus

\[
A_{\mathrm{adia}}=\frac{-v R T_{0}\left(\eta^{\gamma-1}-1\right)}{\gamma-1}
\]

On the other hand, we know \(A_{\text {iso }}=v R T_{0} \ln \left(\frac{1}{\eta}\right)=-v R T_{0} \ln \eta\) (work done by the gas)\\
Thus

\[
\frac{A_{\text {adia }}}{A_{\text {iso }}}=\frac{\eta^{\gamma-1}-1}{(\gamma-1) \ln \eta}=\frac{5^{0.4}-1}{0.4 \times \ln 5}=1.4
\]

2.41 Since here the piston is conducting and it is moved slowly the temperature on the two sides increases and maintained at the same value.\\
Elementary work done by the agent \(=\) Work done in compression - Work done in expansion i.e. \(d A=p_{2} d V-p_{1} d V=\left(p_{2}-p_{1}\right) d V\)\\
where \(p_{1}\) and \(p_{2}\) are pressures at any instant of the gas on expansion and compression side respectively.\\
From the gas law \(p_{1}\left(V_{0}+S x\right)=v R T\) and \(p_{2}\left(V_{0}-S x\right)=v R T\), for each section ( \(x\) is the displacement of the piston towards section 2 )

So,

\[
p_{2}-p_{1}=v R T \frac{2 S x}{V_{0}^{2}-S^{2} x^{2}}=v R T \cdot \frac{2 V}{V_{0}^{2}-V^{2}}(\text { as } S x=V)
\]

So

\[
d A=v R T \frac{2 V}{V_{0}^{2}-V^{2}} d V
\]

Also, from the first law of thermodynamics\\
\(d A=-d U=-2 \nu \frac{R}{\gamma-1} d T \quad(\) as \(d Q=0)\)\\
So, work done on the gas \(=-d A=2 \nu \cdot \frac{R}{\gamma-1} d T\)\\
Thus \(2 v \frac{R}{\gamma-1} d T=v R T \frac{2 V \cdot d V}{V_{0}^{2}-V^{2}}\),\\
or, \(\frac{d T}{T}=\gamma-1 \frac{V d V}{V_{0}^{2}-V^{2}}\)\\
When the left end is \(\eta\) times the volume of the right end.

\[
\left(V_{0}+V\right)=\eta\left(V_{0}-V\right) \text { or } V=\frac{\eta-1}{\eta+1} V_{0}
\]

\begin{center}
\includegraphics[max width=\textwidth, alt={}]{9f3c5a8e-ac5d-42f4-8eb6-ea2af623f2d8-209_252_483_1096_908}
\end{center}

On integrating\\
or

\[
\begin{gathered}
\ln \frac{T}{T_{0}}=(\gamma-1)\left[-\frac{1}{2} \ln \left(V_{0}^{2}-V^{2}\right)\right]_{0}^{V} \\
\left.=-\frac{\gamma-1}{2}\left[\ln \left(V_{0}^{2}-V^{2}\right)-\ln V_{0}^{2}-V^{2}\right)-\ln V_{0}^{2}\right] \\
=\frac{\gamma-1}{2}\left[\ln V_{0}^{2}-\ln V_{0}^{2}\left\{1-\left(\frac{\eta-1}{\eta+1}\right)\right]\right]=\frac{\gamma-1}{2} \ln \frac{(\eta+1)^{2}}{4 \eta} \\
T=T_{0}\left(\frac{(\eta+1)^{2}}{4 \eta}\right)^{\frac{\gamma-1}{2}}
\end{gathered}
\]

Hence\\
2.42 From energy conservation as in the derivation of Bernoullis theorem it reads


\begin{equation*}
\frac{p}{\rho}+\frac{1}{2} v^{2}+g z+u+Q_{d}=\text { constant } \tag{1}
\end{equation*}


In the Eq. (1) \(u\) is the internal energy per unit mass and in this case is the thermal energy per unit mass of the gas. As the gas vessel is thermally insulated \(Q_{d}=0\), also in our case. Just inside the vessel \(u=\frac{C_{V} T}{M}=\frac{R T}{M(\gamma-1)}\) also \(\frac{p}{\rho}=\frac{R T}{M}\). Inside the vessel \(v=0\) also. Just outside \(p=0\), and \(u=0\). Ingeneral \(g z\) is not very significant for gases.\\
Thus applying Eq. (1) just inside and outside the hole, we get

\[
\begin{aligned}
\frac{1}{2} v^{2} & =\frac{p}{\rho}+u \\
& =\frac{R T}{M}+\frac{R T}{M(\gamma-1)}=\frac{\gamma R T}{M(\gamma-1)}
\end{aligned}
\]

Hence \(\quad v^{2}=\frac{2 \gamma R T}{M(\gamma-1)} \quad\) or, \(\quad v=\sqrt{\frac{2 \gamma R T}{M(\gamma-1)}}=3.22 \mathrm{~km} / \mathrm{s}\).\\
Note : The velocity here is the velocity of hydrodynamic flow of the gas into vaccum. This requires that the diameter of the hole is not too small ( \(D>\) mean free path \(l\) ). In the opposite case ( \(D \ll l\) ) the flow is called effusion. Then the above result does not apply and kinetic theory methods are needed.\\
2.43 The differential work done by the gas

So,

\[
\begin{gathered}
d A=p d V=\frac{v R T^{2}}{a}\left(-\frac{a}{T^{2}}\right) d T=-v R d T \\
\left(\begin{array}{c}
\text { as } \left.p V=v R T \text { and } V=\frac{a}{T}\right) \\
T+\Delta T
\end{array}\right.
\end{gathered}
\]

\[
A=-\int_{T}^{T+\Delta T} v R d T=-v R \Delta T
\]

From the first law of thermodynamics

\[
\begin{gathered}
Q=\Delta U+A=\frac{v R}{\gamma-1} \Delta T-v R \Delta T \\
=v R \Delta T \cdot \frac{2-\gamma}{\gamma-1}=R \Delta T \cdot \frac{2-\gamma}{\gamma-1}(\text { for } v=1 \text { mole })
\end{gathered}
\]

2.44 According to the problem : \(A \alpha U\) or \(d A=a U\) (where \(a\) is proportionality constant)\\
or,


\begin{equation*}
p d V=\frac{a v R d T}{\gamma-1} \tag{1}
\end{equation*}


From ideal gas law, \(p V=v R T\), on differentiating


\begin{equation*}
p d V+V d p=v R d T \tag{2}
\end{equation*}


Thus from (1) and (2)\\
\(p d V=\frac{a}{\gamma-1}(p d V+V d p)\)\\
or, \(p d V\left(\frac{a}{\gamma-1}-1\right)+\frac{a}{\gamma-1} V d p=0\)\\
or, \(p d V(k-1)+k V d p=0\) (where \(k=\frac{a}{\gamma-1}=\) another constant)\\
or, \(p d V \frac{k-1}{k}+V d p=0\)\\
or, \(p d V n+V d p=0\) (where \(\frac{k-1}{k}=n=\) ratio)\\
Dividing both the sides by \(p V\)

\[
n \frac{d V}{V}+\frac{d p}{p}=0
\]

On integrating \(n \ln V+\ln p=\ln C\) (where \(C\) is constant)\\
or,

\[
\ln \left(p V^{n}\right)=\ln C \quad \text { or, } \quad p V^{n}=C \text { (const.) }
\]

2.45 In the polytropic process work done by the gas

\[
A=\frac{v R\left[T_{i}-T_{f}\right]}{n-1}
\]

(where \(T_{i}\) and \(T_{f}\) are initial and final temperature of the gas like in adiabatic process)\\
and

\[
\Delta U=\frac{v R}{\gamma-1}\left(T_{f}-T_{i}\right)
\]

By the first law of thermodynamics \(Q=\Delta U+A\)

\[
\begin{gathered}
=\frac{v R}{\gamma-1}\left(T_{f}-T_{i}\right)+\frac{v R}{n-1}\left(T_{i}-T_{f}\right) \\
=\left(T_{f}-T_{i}\right) v R\left[\frac{1}{\gamma-1}-\frac{1}{n-1}\right]=\frac{v R[n-\gamma]}{(n-1)(\gamma-1)} \Delta T
\end{gathered}
\]

According to definition of molar heat capacity when number of moles \(v=1\) and \(\Delta T=1\) then \(Q=\) Molar heat capacity.

Here,

\[
C_{n}=\frac{R(n-\gamma)}{(n-1)(\gamma-1)}<0 \text { for } 1<n<\gamma
\]

2.46 Let the process be polytropic according to the law \(p V^{n}=\) constant

Thus,

\[
p_{f} V_{f}^{n}=p_{i} V_{i}^{n} \quad \text { or, } \quad\left(\frac{p_{i}}{p_{f}}\right)=\beta
\]

So,

\[
\alpha^{n}=\beta \quad \text { or } \quad \ln \beta=n \ln \alpha \quad \text { or } \quad n=\frac{\ln \beta}{\ln \alpha}
\]

In the polytropic process molar heat capacity is given by

\[
\begin{aligned}
C_{n} & =\frac{R(n-\gamma)}{(n-1)(\gamma-1)}=\frac{R}{\gamma-1}-\frac{R}{n-1} \\
& =\frac{R}{\gamma-1}-\frac{R \ln \alpha}{\ln \beta-\ln \alpha}, \quad \text { where } n=\frac{\ln \beta}{\ln \alpha}
\end{aligned}
\]

So,

\[
C_{n}=\frac{8.314}{1.66-1}-\frac{8.314 \ln 4}{\ln 8-\ln 4}=-42 \mathrm{~J} / \mathrm{mol} . \mathrm{K}
\]

2.47 (a) Increment of internal energy for \(\Delta T\), becomes

\[
\Delta U=\frac{v R \Delta T}{\gamma-1}=\frac{R \Delta T}{\gamma-1}=-324 \mathrm{~J}(\text { as } v=1 \mathrm{~mole})
\]

From first law of thermodynamics

\[
Q=\Delta U+A=\frac{R \Delta T}{\gamma-1}-\frac{R \Delta T}{n-1}=0.11 \mathrm{~kJ}
\]

(b) Sought work done, \(A_{n}=\int p d V=\int_{V_{1}} \frac{k}{V^{n}} d V\)

\[
\begin{gathered}
\text { (where } \left.p V^{n}=k=p_{i} V_{i}^{n}=p_{f} V_{f}^{n}\right) \\
=\frac{k}{1-n}\left(V_{f}^{1-n}-V_{i}^{1-n}\right)=\frac{\left(p_{f} V_{f}^{n} V_{f}^{1-n}-p_{i} V_{i}^{n} V_{i}^{1-n}\right)}{1-n} \\
=\frac{p_{f} V_{f}-p_{i} V_{i}}{1-n}=\frac{v R\left(T_{f}-T_{i}\right)}{1-n} \\
=\frac{v R \Delta T}{n-1}=-\frac{R \Delta T}{n-1}=0.43 \mathrm{~kJ}(\text { as } v=1 \text { mole })
\end{gathered}
\]

2.48 Law of the process is \(p=\alpha V\) or \(p V^{-1}=\alpha\)\\
so the process is polytropic of index \(n=-1\)\\
As \(p=\alpha V\) so, \(p_{i}=\alpha V_{0}\) and \(p_{f}=\alpha \eta V_{0}\)\\
(a) Increment of the internal energy is given by

\[
\Delta U=\frac{v R}{\gamma-1}\left[T_{f}-T_{i}\right]=\frac{p_{f} V_{f}-p_{i} V_{i}}{\gamma-1}
\]

(b) Work done by the gas is given by

\[
\begin{aligned}
A & =\frac{p_{i} V_{i}-p_{f} V_{f}}{n-1}=\frac{\alpha V_{0}^{2}-\alpha \eta V_{0} \cdot \eta V_{0}}{-1-1} \\
& =\frac{\alpha V_{0}^{2}\left(1-\eta^{2}\right)}{-2}=\frac{1}{2} \alpha V_{0}^{2}\left(\eta^{2}-1\right)
\end{aligned}
\]

(c) Molar heat capacity is given by

\[
C_{n}=\frac{R(n-\gamma)}{(n-1)(\gamma-1)}=\frac{R(-1-\gamma)}{(-1-1)(\gamma-1)}=\frac{R}{2} \frac{\gamma+1}{\gamma-1}
\]

2.49 (a) \(\Delta U=\frac{v R}{\gamma-1} \Delta T\) and \(Q=v C_{n} \Delta T\)\\
where \(C_{n}\) is the molar heat capacity in the process. It is given that \(Q=-\Delta U\)\\
So,

\[
C_{n} \Delta T=\frac{R}{\gamma-1} \Delta T, \quad \text { or } \quad C_{n}=-\frac{R}{\gamma-1}
\]

(b) By the first law of thermodynamics, \(d Q=d U+d A\),\\
or,

\[
2 d Q=d A(\text { as } d Q=-d U)
\]

\[
2 v C_{n} d T=p d V, \text { or, } \frac{2 R v}{\gamma-1} d T+p d V=0
\]

So,

\[
\frac{2 R V}{\gamma-1} d T+\frac{v R T}{V} d V=0, \quad \text { or, } \quad \frac{2 d T}{(\gamma-1) T}+\frac{d V}{V}=0
\]

or, \(\quad \frac{d T}{T}+\frac{\gamma-1}{2} \frac{d V}{V}=0\), or, \(\quad T V^{(\gamma-1) / 2}=\) constant.\\
(c) We know \(C_{n}=\frac{(n-\gamma) R}{(n-1)(\gamma-1)}\)

But from part (a), we have \(C_{n}=-\frac{R}{\gamma-1}\)\\
Thus

\[
\begin{gathered}
-\frac{R}{\gamma-1}=\frac{(n-\gamma) R}{(n-1)(\gamma-1)} \text { which yields } \\
n=\frac{1+\gamma}{2}
\end{gathered}
\]

From part (b); we know \(T V^{(\gamma-1) / 2}=\) constant\\
So, \(\quad \frac{T_{o}}{T}=\left(\frac{V}{V_{0}}\right)^{(\gamma-1) / 2}=\eta^{(\gamma-1) / 2}\) (where \(T\) is the final temperature)\\
Work done by the gas for one mole is given by

\[
A=R \frac{\left(T_{0}-T\right)}{n-1}=\frac{2 R T_{0}\left[1-\eta^{(1-\gamma) / 2}\right]}{\gamma-1}
\]

2.50 Given \(p=a T^{\alpha}\) (for one mole of gas)

So,

\[
p T^{-\alpha}=a \quad \text { or, } \quad p\left(\frac{p V}{R}\right)^{-\alpha}=a
\]

or,

\[
p^{1-\alpha} V^{-\alpha}=a R^{-\alpha} \text { or, } p V^{\alpha /(\alpha-1)}=\text { constant }
\]

Here polytropic exponent \(n=\frac{\alpha}{\alpha-1}\)\\
(a) In the polytropic process for one mole of gas :

\[
A=\frac{R \Delta T}{1-n}=\frac{R \Delta T}{\left(1-\frac{\alpha}{\alpha-1}\right)}=R \Delta T(1-\alpha)
\]

(b) Molar heat capacity is given by

\[
C=\frac{R}{\gamma-1}-\frac{R}{n-1}=\frac{R}{\gamma-1}-\frac{R}{\frac{\alpha}{\alpha-1}-1}=\frac{R}{\gamma-1}+R(1-\alpha)
\]

2.51 Given \(U=a V^{\alpha}\)\\
\(\begin{array}{lc}\text { or } & v C_{V} T=a V^{\alpha}, \text { or, } v C_{V} \frac{p V}{v R}=a V^{\alpha} \\ \text { or, } & a V^{\alpha} \cdot \frac{R}{C_{V}} \cdot \frac{1}{p V}=1, \text { or }, V^{\alpha-1} \cdot p^{-1}=\frac{C_{V}}{R a} \\ \text { or } & p V^{1-\alpha}=\frac{R a}{C_{V}}=\text { constant }=a(\gamma-1)\left[\text { as } C_{V}=\frac{R}{\gamma-1}\right]\end{array}\)\\
So polytropric index \(n=1-\alpha\).\\
(a) Work done by the gas is given by

Hence

\[
\begin{gathered}
A=\frac{-v R \Delta T}{n-1} \text { and } \Delta U=\frac{v R \Delta T}{\gamma-1} \\
A=\frac{-\Delta U(\gamma-1)}{n-1}=\frac{\Delta U(\gamma-1)}{\alpha}(\text { as } n=1-\alpha)
\end{gathered}
\]

By the first law of thermodynamics, \(Q=\Delta U+A\)

\[
=\Delta U+\frac{\Delta U(\gamma-1)}{\alpha}=\Delta U\left[1+\frac{\gamma-1}{\alpha}\right]
\]

(b) Molar heat capacity is given by

\[
\begin{aligned}
C= & \frac{R}{\gamma-1}-\frac{R}{n-1}=\frac{R}{\gamma-1}-\frac{R}{1-\alpha-1} \\
& =\frac{R}{\gamma-1}+\frac{R}{\alpha}(\text { as } n=1-\alpha)
\end{aligned}
\]

2.52 (a) By the first law of thermodynamics

\[
d Q=d U+d A=v C_{V} d T+p d V
\]

Molar specific heat according to definition

\[
\begin{gathered}
C=\frac{d Q}{v d T}=\frac{C_{V} d T+p d V}{v d T} \\
=\frac{v C_{V} d T+\frac{v R T}{V} d V}{v d T}=C_{V}+\frac{R T}{V} \frac{d V}{d T}
\end{gathered}
\]

We have

\[
T=T_{0} e^{\alpha V}
\]

After differentiating, we get \(d T=\alpha T_{0} e^{\alpha V} \cdot d V\)\\
So,

\[
\frac{d V}{D T}=\frac{1}{\alpha T_{0} e^{\alpha V}}
\]

Hence

\[
C=C_{V}+\frac{R T}{V} \cdot \frac{1}{\alpha T_{0} e^{\alpha V}}=C_{V}+\frac{R T_{0} e^{\alpha V}}{\alpha V T_{0} e^{\alpha \nu}}=C_{V}+\frac{R}{\alpha V}
\]

(b) Process is \(p=p_{0} e^{\alpha V}\)

\[
p=\frac{R T}{V}=p_{0} e^{\alpha V}
\]

\[
\text { or, } \quad T=\frac{p_{0}}{R} e^{\alpha V} \cdot V
\]

So, \(C=C_{V}+\frac{R T}{V} \frac{d V}{d T}=C_{V}+p_{0} e^{\alpha V} \cdot \frac{R}{p_{0} e^{\alpha V}(1+\alpha V)}=C_{V}+\frac{R}{1+\alpha V}\)\\
2.53 Using 2.52\\
(a) \(C=C_{V}+\frac{R T}{V} \frac{d V}{d T}=C_{V}+\frac{p d V}{d T}\) (for one mole of gas)

We have \(p=p_{0}+\frac{\alpha}{V}, \quad\) or, \(\quad \frac{R T}{V}=p_{0}+\frac{\alpha}{V}, \quad\) or, \(\quad R T=p_{0} V+\alpha\)\\
Therefore \(\quad R d T=p_{0} d V\), So, \(\frac{d V}{d T}=\frac{R}{p_{0}}\)\\
Hence

\[
\begin{gathered}
C=C_{V}+\left(p_{0}+\frac{\alpha}{V}\right) \cdot \frac{R}{p_{0}}=\frac{R}{\gamma-1}+\left(1+\frac{\alpha}{p_{0} V}\right) R \\
=\left(R+\frac{R}{\gamma-1}\right)+\frac{\alpha R}{p_{0} V}=\frac{\gamma R}{\gamma-1}+\frac{\alpha R}{p_{0} V}
\end{gathered}
\]

(b) Work done is given by

\[
\begin{gathered}
A=\int_{V_{1}}^{V_{2}}\left(p_{0}+\frac{\alpha}{V}\right) d V=p_{0}\left(V_{2}-V_{1}\right)+\alpha \ln \frac{V_{2}}{V_{1}} \\
\Delta U=C_{V}\left(T_{2}-T_{1}\right)=C_{V}\left(\frac{p_{2} V_{2}}{R}-\frac{p_{1} V_{1}}{R}\right) \text { (for one mole) } \\
=\frac{R}{(\gamma-1) R}\left(p_{2} V_{2}-p_{1} V_{1}\right) \\
=\frac{1}{\gamma-1}\left[\left(p_{0}+\alpha V_{2}\right) V_{2}-\left(p_{0}+\frac{\alpha}{V_{1}}\right) V_{1}\right]=\frac{p_{0}\left(V_{2}-V_{1}\right)}{\gamma-1}
\end{gathered}
\]

By the first law of thermodynamics \(Q=\Delta U+A\)

\[
\begin{gathered}
=p_{0}\left(V_{2}-V_{1}\right)+\alpha \ln \frac{V_{2}}{V_{1}}+\frac{p_{0}\left(V_{2}-V_{1}\right)}{(\gamma-1)} \\
=\frac{\gamma p_{0}\left(V_{2}-V_{1}\right)}{\gamma-1}+\alpha \ln \frac{V_{2}}{V_{1}}
\end{gathered}
\]

2.54 (a) Heat capacity is given by

\[
C=C_{V}+\frac{R T}{V} \frac{d V}{d T} \text { (see solution of 2.52) }
\]

We have

\[
T=T_{0}+\alpha V \quad \text { or, } \quad V=\frac{T}{\alpha}-\frac{T_{0}}{\alpha}
\]

After differentiating, we get, \(\frac{d V}{d T}=\frac{1}{\alpha}\)

Hence\\
(b) Given \(T=T_{0}+\alpha V\)

As \(T=\frac{p V}{R}\) for one mole of gas

Now

\[
\begin{gathered}
C=C_{V}+\frac{R T}{V} \cdot \frac{1}{\alpha}=\frac{R}{\gamma-1}+\frac{R\left(T_{0}+\alpha V\right)}{V} \cdot \frac{1}{\alpha} \\
=\frac{R}{\gamma-1}+R\left(\frac{T_{0}}{\alpha V}+1\right)=\frac{\gamma R}{\gamma-1}+\frac{R T_{0}}{\alpha V}=C_{V}+\frac{R T}{\alpha V}=C_{p}+\frac{R T_{0}}{\alpha V}
\end{gathered}
\]

\[
p=\frac{R}{V}\left(T_{0}+\alpha V\right)=\frac{R T}{V}=\alpha R
\]

\[
\begin{gathered}
A=\int_{V_{1}}^{V_{2}} p d V=\int_{V_{1}}^{V_{2}}\left(\frac{R T_{0}}{v}+\alpha R\right) d V(\text { for one mole }) \\
=R T_{0} \ln \frac{V_{2}}{V_{1}}+\alpha\left(V_{2}-V_{1}\right) \\
\Delta U=C_{V}\left(T_{2}-T_{1}\right) \\
=C_{V}\left[T_{0}+\alpha V_{2}-T_{0} \alpha V_{1}\right]=\alpha C_{V}\left(V_{2}-V_{1}\right)
\end{gathered}
\]

By the first law of thermodynamics \(Q=\Delta U+A\)

\[
\begin{gathered}
=\frac{\alpha R}{\gamma-1}\left(V_{2}-V_{1}\right)+R T_{0} \ln \frac{V_{2}}{V_{1}}+\alpha R\left(V_{2}-V_{1}\right) \\
=\alpha R\left(V_{2}-V_{1}\right)\left[1+\frac{1}{\gamma-1}\right]+R T_{0} \ln \frac{V_{2}}{V_{1}} \\
=\alpha C_{p}\left(V_{2}-V_{1}+R T_{0} \ln \frac{V_{2}}{V_{1}}\right. \\
=\alpha C_{p}\left(V_{2}-V_{1}\right)+R T_{0} \ln \frac{V_{2}}{V_{1}}
\end{gathered}
\]

2.55 Heat capacity is given by \(C=C_{V}+\frac{R T}{V} \frac{d V}{d T}\)\\
(a) Given \(C=C_{V}+\alpha T\)

So, \(\quad C_{V}+\alpha T=C_{V}+\frac{R T}{V} \frac{d V}{d T}\) or, \(\quad \frac{\alpha}{R} d T=\frac{d V}{V}\)\\
Integrating both sides, we get \(\frac{\alpha}{R} T=\ln V+\ln C_{0}=\operatorname{In} V C_{0}, C_{0}\) is a constant.\\
Or, \(\quad V \cdot C_{0}=e^{\alpha T / R}\) or \(V \cdot e^{\alpha T / R}=\frac{1}{C_{0}}=\) constant\\
(b) \(C=C_{V}+\beta V\)\\
and

\[
C=C_{V}+\frac{R T}{V} \frac{d V}{d T} \text { so, } C_{V} \frac{R T}{V} \frac{d V}{d T}=C_{V}+\beta V
\]

or, \(\quad \frac{R T}{V} \frac{d V}{d T}=\beta V\) or, \(\frac{d V}{V^{2}}=\frac{\beta}{R} \frac{d T}{T}\) or, \(V^{-2}=\frac{d T}{T}\)\\
Integrating both sides, we get \(\frac{R}{\beta} \frac{V^{-1}}{\beta-1}=\ln T+\ln C_{0}=\ln T \cdot C_{0}\)\\
So, \(\ln T \cdot C_{0}=-\frac{R}{\beta V} T \cdot C_{0}=e^{-R / \beta V}\) or, \(T e^{-R / \beta V}=\frac{1}{C_{0}}=\) constant\\
(c) \(C=C_{V}+a p\) and \(C=C_{V}+\frac{R T}{V} \frac{d V}{d T}\)

So, \(\quad C_{V}+a p=C_{V}+\frac{R T}{V} \frac{d V}{d T}\) so, \(a p=\frac{R T}{V} \frac{d V}{d T}\)\\
or, \(\quad a \frac{R T}{V}=\frac{R T}{V} \frac{d V}{d T}\) (as \(p=\frac{R T}{V}\) for one mole of gas)\\
or, \(\quad \frac{d V}{d T}=a\) or, \(d V=a d T\) or, \(d T=\frac{d V}{a}\)\\
So, \(T=\frac{V}{a}+\) constant or \(V-a T=\) constant\\
2.56 (a) By the first law of thermodynamics \(A=Q-\Delta U\)\\
or,

\[
=C d T-C_{V} d T=\left(C-C_{V}\right) d T \text { (for one mole) }
\]

Given

\[
C=\frac{\alpha}{T}
\]

So,

\[
\begin{aligned}
A & =\int_{T_{0}}^{\eta T_{0}}\left(\frac{\alpha}{T}-C_{V}\right) d T=\alpha \ln \frac{\eta T_{0}}{T_{0}}-C_{V}\left(\eta T_{0}-T_{0}\right) \\
& =\alpha \ln \eta-C_{V} T_{0}(\eta-1)=\alpha \ln \eta+\frac{R T}{\gamma-1}(\eta-1)
\end{aligned}
\]

(b) \(C=+\frac{d Q}{d T}=\frac{R T}{V} \frac{d V}{d T}+C_{V}\)

Given

\[
C=\frac{\alpha}{T}, \quad \text { so } \quad C_{V}+\frac{R T}{V} \frac{d V}{d T}=\frac{\alpha}{T}
\]

or,

\[
\frac{R}{\gamma-1} \frac{1}{R T}+\frac{d V}{V}=\frac{\alpha}{R T^{2}} d T
\]

or,

\[
\frac{d V}{V}=\frac{\alpha}{R T^{2}} d T-\frac{1}{\gamma-1} \cdot \frac{d T}{T}
\]

or,

\[
(\gamma-1) \frac{d V}{V}=\frac{\alpha(\gamma-1)}{R T^{2}} d T-\frac{d T}{T}
\]

Integrating both sides, we get\\
or,

\[
(\gamma-1) \ln V=-\frac{\alpha(\gamma-1)}{R T}-\ln T+\ln K
\]

or,

\[
\begin{aligned}
\ln V^{\gamma-1} \frac{T}{K} & =\frac{-\alpha(\gamma-1)}{R T} \\
\ln V^{\gamma-1} \cdot \frac{p V}{R K} & =\frac{-\alpha(\gamma-1)}{p V}
\end{aligned}
\]

or,

\[
\frac{p V^{\gamma}}{R K}=e^{-\alpha(\gamma-1)_{p} V}
\]

or,

\[
p V^{\gamma e^{\alpha(\gamma-1) / p V}}=R K=\text { constant }
\]

2.57 The work done is

\[
\begin{aligned}
A & =\int_{V_{1}}^{V_{2}} p d V=\int_{V_{1}}^{V_{2}}\left(\frac{R T}{v-b}-\frac{a}{V^{2}}\right) d V \\
& =R T \ln \frac{V_{2}-b}{V_{1}-b}+a\left(\frac{1}{V_{2}}-\frac{1}{V_{1}}\right)
\end{aligned}
\]

2.58 (a) The increment in the internal energy is

\[
\Delta U=\int_{V_{1}}^{V_{1}}\left(\frac{\partial U}{\partial V}\right)_{T} d V
\]

But from second law

\[
\left(\frac{\partial U}{\partial V}\right)_{T}=T\left(\frac{\partial S}{\partial V}\right)_{T}-p=T\left(\frac{\partial p}{\partial T}\right)_{V}-p
\]

On the other hand

\[
p=\frac{R T}{V-b}-\frac{a}{V^{2}}
\]

or,

\[
T\left(\frac{\partial p}{\partial T}\right)_{V}=\frac{R T}{V-b} \text { and }\left(\frac{\partial U}{\partial V}\right)_{t}=\frac{a}{V^{2}}
\]

So,

\[
\Delta U=a\left(\frac{1}{V_{1}}-\frac{1}{V_{2}}\right)
\]

(b) From the first law

\[
Q=A+\Delta U=R T \ln \frac{V_{2}-b}{V_{1}-b}
\]

2.59 (a) From the first law for an adiabatic

\[
d Q=d U+p d V=0
\]

From the previous problem

\[
d U=\left(\frac{\partial U}{\partial T}\right)_{V} d T+\left(\frac{\partial U}{\partial V}\right)_{T} d V=C_{V} d T+\frac{a}{V^{2}} d V
\]

So,

\[
0=C_{V} d T+\frac{R T d V}{V-b}
\]

This equation can be integrated if we assume that \(C_{V}\) and \(b\) are constant then \(\frac{R}{C_{V}} \frac{d V}{V-b}+\frac{d T}{T}=0\), or, \(\ln T+\frac{R}{C_{V}} \ln (V-b)=\) constant\\
or,

\[
T(V-b)^{R / C_{v}}=\text { constant }
\]

(b) We use

Now,

So along constant \(p\),

\[
C_{p}=C_{V}+\frac{R T}{V-b}\left(\frac{\partial V}{\partial T}\right)_{p}
\]

Thus

\[
C_{p}-C_{V}=\frac{R T}{V-b}\left(\frac{\partial V}{\partial T}\right)_{p}, \text { But } p=\frac{R T}{V-b}-\frac{a}{V^{2}}
\]

or,\\
and\\
2.60 From the first law\\
\(Q=U_{f}-U_{i}+A=0\), as the vessels are themally insulated.\\
As this is free expansion, \(A=0\), so, \(U_{f}=U_{i}\)\\
But

\[
U=v C_{V} T-\frac{a v^{2}}{V}
\]

So,

\[
\begin{gathered}
C_{V}\left(T_{f}-T_{i}\right)=\left(\frac{a}{V_{1}+V_{2}}-\frac{a}{V_{1}}\right) v=\frac{-a V_{2} v}{V_{1}\left(V_{1}+V_{2}\right)} \\
\Delta T=\frac{-a(\gamma-1) V_{2} v}{R V_{1}\left(V_{1}+V_{2}\right)}
\end{gathered}
\]

or,

\[
d U=C_{V} d T+\frac{a}{V^{2}} d V
\]

\[
d Q=C_{V} d T+\frac{R T}{V-b} d V
\]

So along constant \(p, \quad C_{p}=C_{V}+\frac{R T}{V-b}\left(\frac{\partial V}{\partial T}\right)_{p}\)\\
Thus\\
and

\[
\begin{gathered}
T\left(\frac{\partial V}{\partial T}\right)_{p}=\frac{R T / V-b}{\frac{R T}{(V-b)^{2}}-\frac{2 a}{V^{3}}}=\frac{V-b}{1-\frac{2 a(V-b)^{2}}{R T V^{3}}} \\
C_{p}-C_{V}=\frac{R}{1-\frac{2 a(V-b)^{2}}{R T V^{3}}}
\end{gathered}
\]

\subsection*{2.3 KINETIC THEORY OF GASES. BOLTZMANN'S LAW AND MAXWELL'S DISTRIBUTION}
2.62 From the formula \(p=n k T\)

\[
\begin{aligned}
n & =\frac{p}{k T}=\frac{4 \times 10^{-15} \times 1.01 \times 10^{5}}{1.38 \times 10^{-23} \times 300} \text { per } \mathrm{m}^{3} \\
& =1 \times 10^{11} \text { per } \mathrm{m}^{3}=10^{5} \text { per c.c }
\end{aligned}
\]

Mean distance between molecules

\[
\left(10^{-5} \text { c.c. }\right)^{1 / 3}=10^{1 / 3} \times 10^{-2} \mathrm{~cm}=0.2 \mathrm{~mm} .
\]

2.63 After dissociation each \(N_{2}\) molecule becomes two \(N\)-atoms and so contributes, \(2 \times 3\) degrees of freedom. Thus the number of moles becomes

\[
\frac{m}{M}(1+\eta) \text { and } p=\frac{m R T}{M V}(1+\eta)
\]

Here \(M\) is the molecular weight in grams of \(N_{2}\).\\
2.64 Let \(n_{1}=\) number density of He atoms, \(n_{2}=\) number density of \(N_{2}\) molecules

Then

\[
\rho=n_{1} m_{1}+n_{2} m_{2}
\]

where \(m_{1}=\) mass of He atom, \(m_{2}=\) mass of \(N_{2}\) molecule also \(p=\left(n_{1}+n_{2}\right) k T\)\\
From these two equations we get

\[
n_{1}=\left(\frac{p}{k T}-\frac{\rho}{m_{2}}\right) /\left(1-\frac{m_{1}}{m_{2}}\right)
\]

\(2.65 p=\frac{n v \times 2 m v \cos \theta \times d A \cos \theta}{d A}\)

\[
=2 m n v^{2} \cos ^{2} \theta
\]

2.66 From the formula

\[
v=\sqrt{\frac{\gamma p}{\rho}}, \gamma=\frac{\rho V^{2}}{p}
\]

\begin{center}
\includegraphics[max width=\textwidth, alt={}]{9f3c5a8e-ac5d-42f4-8eb6-ea2af623f2d8-220_381_348_1028_744}
\end{center}

If \(\boldsymbol{i}=\) number of degrees of freedom of the gas then

\[
\begin{aligned}
C_{p} & =C_{V}+R T \text { and } C_{V}=\frac{i}{2} R T \\
\gamma=\frac{C_{p}}{C_{v}} & =1+\frac{2}{i} \text { or } i=\frac{2}{\gamma-1}=\frac{2}{\frac{\rho v^{2}}{p}-1}
\end{aligned}
\]

\(2.67 v_{\text {sound }}=\sqrt{\frac{Y P}{\rho}}=\sqrt{\frac{Y R T}{M}}\), and \(v_{r m s}=\sqrt{\frac{3 k T}{m}}=\sqrt{\frac{3 R T}{M}}\)\\
so,

\[
\frac{v_{\mathrm{sound}}}{v_{\mathrm{rms}}}=\sqrt{\frac{\gamma}{3}}=\sqrt{\frac{i+2}{3 i}}
\]

(a) For monoatomic gases \(i=3\)

\[
\frac{v_{\text {sound }}}{v_{\mathrm{rms}}}=\sqrt{\frac{5}{9}}=0.75
\]

(b) For rigid diatomic molecules \(i=5\)

\[
\frac{v_{\text {sound }}}{v_{\mathrm{rms}}}=\sqrt{\frac{7}{15}}=0.68
\]

2.68 For a general noncollinear, nonplanar molecule\\
mean energy \(=\frac{3}{2} k T(\) translational \()+\frac{3}{2} k T(\) rotational \()+(3 N-6) k T(\) vibrational \()\)\\
\(=(3 N-3) k T\) per molecule\\
For linear molecules, mean energy \(=\frac{3}{2} k T\) (translational)\\
\(+k T\) (rotational) \(+(3 N-5) k T\) (vibrational)\\
\(=\left(3 N-\frac{5}{2}\right) k T\) per molecule\\
Translational energy is a fraction \(\frac{1}{2(N-1)}\) and \(\frac{1}{2 N-\frac{5}{3}}\) in the two cases.\\
2.69 (a) A diatomic molecule has 2 translational, 2 rotational and one vibrational degrees of freedom. The corresponding energy per mole is

\[
\begin{gathered}
\frac{3}{2} R T,(\text { for translational })+2 \times \frac{1}{2} R T,(\text { for rotational }) \\
+1 \times R T,(\text { for vibrational })=\frac{7}{2} R T
\end{gathered}
\]

Thus,

\[
C_{V}=\frac{7}{2} R, \text { and } \gamma=\frac{C_{p}}{C_{V}}=\frac{9}{7}
\]

(b) For linear \(N\) - atomic molecules energy per mole

\[
=\left(3 N-\frac{5}{2}\right) R T \text { as before }
\]

So,

\[
C_{V}=\left(3 N-\frac{5}{2}\right) R \text { and } \gamma=\frac{6 N-3}{6 N-5}
\]

(c) For noncollinear \(N\) - atomic molecules\\
\(C_{V}=3(N-1) R\) as betore (2.68) \(\gamma=\frac{3 N-2}{3 N-3}=\frac{N-2 / 3}{N-1}\)\\
2.70 In the isobaric process, work done is

\[
A=p d v=R d T \text { per mole. }
\]

On the other hand heat transferred \(Q=C_{p} d T\)\\
Now \(C_{p}=(3 N-2) R\) for non-collinear molecules and \(C_{p}=\left(3 N-\frac{3}{2}\right) R\) for linear molecules

Thus

\[
\frac{A}{Q}=\left\{\begin{array}{l}
\frac{1}{3 N-2} \text { non collinear } \\
\frac{1}{3 N-\frac{3}{2}} \text { linear }
\end{array}\right.
\]

For monoatomic gases, \(c_{p}=\frac{5}{2}\) and \(\frac{A}{Q}=\frac{2}{5}\)\\
2.71 Given specific heats \(c_{p}, c_{v}\) (per unit mass)

\[
M\left(c_{p}-c_{\nu}\right)=R \text { or, } M=\frac{R}{c_{p}-c_{\nu}}
\]

Also

\[
\gamma=\frac{c_{p}}{c_{v}}=\frac{2}{i}+1, \text { os, } i=\frac{2}{\frac{c_{p}}{c_{v}} 1}=\frac{2 c_{v}}{c_{p}-c_{v}}
\]

2.72 (a) \(C_{\nu}=29 \frac{J}{{ }^{\circ} \mathrm{K} \text { mole }}=\frac{29}{8.3} \mathrm{R}\)

\[
\begin{aligned}
C_{v}=\frac{20 \cdot 7}{8 \cdot 3} R, \quad \gamma & =\frac{29}{20 \cdot 7}=1 \cdot 4=\frac{7}{5} \\
i & =5
\end{aligned}
\]

(b) In the process \(p T=\) const.

\[
\frac{T^{2}}{V}=\text { const, So } 2 \frac{d T}{T}-\frac{d V}{V}=0
\]

Thus

\[
C d T=C_{V} d T+p d V=C_{V} d T+\frac{R T}{V} d V=C_{V} d T+\frac{2 R T}{T} d T
\]

or

\[
C=C_{V}+2 R=\left(\frac{29}{8 \cdot 3}\right) R \quad \text { So } \quad C_{V}=\frac{12 \cdot 4}{8 \cdot 3} R=\frac{3}{2} R
\]

Hence \(i=3\) (monoatomic)\\
2.73 Obviously

\[
\frac{1}{R} C_{V}=\frac{3}{2} \gamma_{1}+\frac{5}{2} \gamma_{2}
\]

(Since a monoatomic gas has \(C_{V}=\frac{3}{2} R\) and a diatomic gas has \(C_{V}=\frac{5}{2} R\). [The diatomic molecule is rigid so no vibration])

Hence

\[
\begin{aligned}
\frac{1}{R} C_{p} & =\frac{3}{2} \gamma_{1}+\frac{5}{2} \gamma_{2}+\gamma_{1}+\gamma_{2} \\
\gamma & =\frac{C_{p}}{C_{V}}=\frac{5 \gamma_{1}+7 \gamma_{2}}{3 \gamma_{1}+5 \gamma_{2}}
\end{aligned}
\]

2.74 The internal energy of the molecules are

\[
U=\frac{1}{2} m N<(\vec{u}-\vec{v})^{2}>=\frac{1}{2} m N<u^{2}-v^{2}>
\]

where \(\vec{v}=\) velocity of the vessel, \(N=\) number of molecules, each of mass \(m\). When the vessel is stopped, internal energy becomes \(\frac{1}{2} m N\left\langle u^{2}\right\rangle\)\\
So there is an increase in internal energy of \(\Delta U=\frac{1}{2} m N v^{2}\). This will give rise to a rise in temperature of

\[
\begin{aligned}
\Delta T & =\frac{\frac{1}{2} m N v^{2}}{\frac{i}{2} R} \\
& =\frac{m N v^{2}}{i R}
\end{aligned}
\]

there being no flow of heat. This change of temperature will lead to an excess pressure\\
and finally

\[
\begin{aligned}
& \Delta p=\frac{R \Delta T}{V}=\frac{m N v^{2}}{i V} \\
& \frac{\Delta p}{p}=\frac{M v^{2}}{i R T}=2.2 \%
\end{aligned}
\]

where \(M=\) molecular weight of \(N_{2}, i=\) number of degrees of freedom of \(N_{2}\)\\
2.75 (a) From the equipartition theorem\\
\(\bar{\varepsilon}=\frac{3}{2} k T=6 \times 10^{-21} \mathrm{~J} ;\) and \(v_{r m s}=\sqrt{\frac{3 k T}{m}}=\sqrt{\frac{3 R T}{M}}=0.47 \mathrm{~km} / \mathrm{s}\)\\
(b) In equilibrium the mean kinetic energy of the droplet will be equal to that of a molecule.

\[
\frac{1}{2} \frac{\pi}{6} d^{3} \rho v_{r m s}^{2}=\frac{3}{2} k T \text { or } v_{r m s}=3 \sqrt{\frac{2 k T}{\pi d^{3} \rho}}=0.15 \mathrm{~m} / \mathrm{s}
\]

2.76 Here \(i=5, C_{V}=\frac{5}{2} R, \gamma=\frac{7}{5}\) given

\[
v_{r m s}^{\prime}=\sqrt{\frac{3 R T}{M}}=\frac{1}{\eta} v_{r m s}=\frac{1}{\eta} \sqrt{\frac{3 R T}{M}} \quad \text { or } \quad T=\frac{1}{\eta^{2}} T
\]

Now in an adiabatic process

\[
\begin{gathered}
T V^{\gamma-1}=T V^{2 / i}=\text { constant or } V T^{i / 2}=\text { constant } \\
V^{\prime}\left(\frac{1}{\eta^{2}} T\right)^{i / 2}=V T^{i / 2} \text { or } V^{\prime} \eta^{-i}=V \text { or } V^{\prime}=\eta^{i} V
\end{gathered}
\]

The gas must be expanded \(\eta^{i}\) times, i.e 7.6 times.\\
2.77 Here

\[
C_{V}=\frac{5}{2} \frac{m}{M} R(i=5 \text { here })
\]

\(m=\) mass of the gas, \(M=\) molecular weight. If \(v_{r m s}\) increases \(\eta\) times, the temperature will have increased \(\eta^{2}\) times. This will require (neglecting expansion of the vessels) a heat flow of amount

\[
\frac{5}{2} \frac{m}{M} R\left(\eta^{2}-1\right) T=10 \mathrm{~kJ}
\]

2.78 The root mean square angular velocity is given by

\[
\begin{aligned}
\frac{1}{2} I \omega^{2} & =2 \times \frac{1}{2} k T(2 \text { degrees of rotations }) \\
\omega & =\sqrt{\frac{2 k T}{I}}=6.3 \times 10^{12} \mathrm{rad} / \mathrm{s}
\end{aligned}
\]

or\\
2.79 Under compression, the temperature will rise

\[
T V^{\gamma-1}=\text { constant, } T V^{2 / i}=\text { constant }
\]

or,

\[
T\left(\eta^{-1} V_{0}\right)^{2 / i}=T_{0} V_{0}^{2 / i} \quad \text { or, } \quad T=\eta^{+2 / i} T_{0}
\]

So mean kinetic energy of rotation per molecule in the compressed state

\[
=k T=k T_{0} \eta^{2 / i}=0.72 \times 10^{-20} \mathrm{~J}
\]

2.80 No. of collisions \(=\frac{1}{4} n\langle v\rangle=v\)

Now,

\[
\frac{v^{\prime}}{v}=\frac{n^{\prime}}{n} \frac{\left\langle v^{\prime}\right\rangle}{\langle v\rangle}=\frac{1}{\eta} \sqrt{\frac{T}{T}}
\]

(When the gas is expanded \(\eta\) times, \(n\) decreases by a factor \(\eta\) ). Also \(T^{\prime}(\eta V)^{2 / i}=T V^{2 / i}\) or \(T=\eta^{2 / i} T\) so, \(\frac{v^{\prime}}{v}=\frac{1}{\eta} \eta^{-1 / i}=\eta^{\frac{-i-1}{i}}\)\\
i.e. collisions decrease by a factor \(\eta^{\frac{i+1}{i}}, i=5\) here.\\
2.81 In a polytropic process \(p V^{n}=\) constant., where \(n\) is called the polytropic index. For this process

\[
\begin{aligned}
p V^{n}=\text { constant or } T V^{n-1} & =\text { constant } \\
\frac{d T}{T}+(n-1) \frac{d V}{V} & =0
\end{aligned}
\]

Then

\[
\begin{gathered}
d Q=C d T=d U+p d V=C_{V} d T+p d V \\
=\frac{i}{2} R d T+\frac{R T}{V} d V=\frac{i}{2} R d T-\frac{1}{n-1} R d T=\left(\frac{i}{2}-\frac{1}{n-1}\right) R d T
\end{gathered}
\]

Now

\[
C=R \text { so } \frac{i}{2}-\frac{1}{n-1}=1
\]

or,

\[
\frac{1}{n-1}=\frac{i}{2}-1=\frac{i-2}{2} \text { or } n=\frac{i}{i-2}
\]

Now

\[
\frac{v^{\prime}}{v}=\frac{n^{\prime}}{n} \frac{\left\langle v^{\prime}\right\rangle}{\langle v\rangle}=\frac{1}{\eta} \sqrt{\frac{T}{T}}=\frac{1}{\eta}\left(\frac{V}{V^{\prime}}\right)^{\frac{n-1}{2}}
\]

\[
=\frac{1}{\eta}\left(\frac{1}{\eta}\right)^{\frac{1}{i-2}}=\left(\frac{1}{\eta}\right)^{\frac{i-1}{i-2}}=\eta^{\frac{-i-1}{i-2}} \text { times }=\frac{1}{2 \cdot 52} \text { times }
\]

2.82 If \(\alpha\) is the polytropic index then

\[
p V^{\alpha}=\text { constant, } T V^{\alpha-1}=\text { constant } .
\]

Now

\[
\frac{v^{\prime}}{v}=\frac{n^{\prime}}{n} \frac{\left\langle v^{\prime}\right\rangle}{\langle v\rangle}=\frac{V}{V^{\prime}} \sqrt{\frac{T^{\prime}}{T}}=\frac{V T^{-1 / 2}}{V^{\prime} T^{-1 / 2}}=1
\]

Hence

\[
\frac{1}{\alpha-1}=-\frac{1}{2} \quad \text { or } \quad \alpha=-1
\]

Then

\[
C=\frac{i R}{2}+\frac{R}{2}=3 R
\]

\(2.83 v_{p}=\sqrt{\frac{2 k T}{m}}=\sqrt{\frac{2 R T}{M}}=\sqrt{\frac{2 p}{\rho}}=0.45 \mathrm{~km} / \mathrm{s}\),

\[
v_{a v}=\sqrt{\frac{8}{\pi} \frac{p}{\rho}}=\cdot 51 \mathrm{~km} / \mathrm{s} \quad \text { and } \quad v_{r m s}=\sqrt{\frac{3 p}{\rho}}=0.55 \mathrm{~km} / \mathrm{s}
\]

2.84 (a) The formula is

\[
d f(u)=\frac{4}{\sqrt{\pi}} u^{2} e^{-u^{2}} d u, \text { where } u=\frac{v}{v_{p}}
\]

Now Prob \(\left(\frac{\left|v-v_{p}\right|}{v_{p}}<\delta \eta\right)=\int_{1-\delta \eta}^{1+\delta \eta} d f(u)\)

\[
=\frac{4}{\sqrt{\pi}} e^{-1} \times 2 \delta \eta=\frac{8}{\sqrt{\pi} e} \delta \eta=0.0166
\]

(b) \(\operatorname{Prob}\left(\left|\frac{v-v_{r m s}}{v_{r m s}}\right|<\delta \eta\right)=\operatorname{Prob}\left(\left|\frac{v}{v_{p}}-\frac{v_{r m s}}{v_{p}}\right|<\delta \eta \frac{v_{r m s}}{v_{p}}\right)\)

\[
\begin{aligned}
= & \operatorname{Prob}\left(\left|u-\sqrt{\frac{3}{2}}\right|<\sqrt{\frac{3}{2}} \delta \eta\right) \\
& \sqrt{\frac{3}{2}}+\sqrt{\frac{3}{2}} \delta \eta \\
= & \int \frac{4}{\sqrt{\pi}} u^{2} e^{-u^{2}} d u \\
& \sqrt{\frac{3}{2}}-\sqrt{\frac{3}{2}} \delta \eta \\
= & \frac{4}{\sqrt{\pi}} \times \frac{3}{2} e^{-3 / 2} \times 2 \sqrt{\frac{3}{2}} \delta \eta=\frac{12 \sqrt{3}}{\sqrt{2 \pi}} e^{-3 / 2} \delta \eta=0.0185
\end{aligned}
\]

2.85 (a) \(v_{r m s}-v_{p}=(\sqrt{3}-\sqrt{2}) \sqrt{\frac{k T}{m}}=\Delta v\),

\[
T=\frac{m}{k}\left(\frac{\Delta v}{(\sqrt{3}-\sqrt{2})}\right)^{2} / \mathrm{K}=384 \mathrm{~K}
\]

(b) Clearly \(v\) is the most probable speed at this temperature. So

\[
\sqrt{\frac{2 k T}{m}}=v \quad \text { or } \quad T=\frac{m v^{2}}{2 k}=342 \mathrm{~K}
\]

2.86 (a) We have,\\
\(\frac{v_{1}^{2}}{v_{p}^{2}} e^{-v_{1}^{2} / v_{p}^{2}}=\frac{v_{2}^{2}}{v_{p}^{2}} e^{-v_{2}^{2} / v_{p}^{2}} \quad\) or \(\quad\left(\frac{v_{1}}{v_{2}}\right)^{2}=e^{v_{1}^{2}-v_{2}^{2} / v_{p}^{2}} \quad\) or \(\quad v_{p}^{2}=\frac{2 k T}{m}=\frac{v_{1}^{2}-v_{2}^{2}}{\left(\ln v_{1}^{2} / v_{2}^{2}\right)}\)

So

\[
T=\frac{m\left(v_{1}^{2}-v_{2}^{2}\right)}{2 k \ln \frac{v_{1}^{2}}{v_{2}^{2}}}=330 \mathrm{~K}
\]

(b) \(F(v)=\frac{4}{\sqrt{\pi}} \frac{v^{2}}{v_{p}^{2}} e^{-v^{2} / v_{p}^{2}} \times \frac{1}{v_{p}}\left(\frac{1}{v_{p}}\right.\) comes from \(\left.F(v) d v=d f(u), d u=\frac{d v}{v_{p}}\right)\)

Thus \(\frac{v^{2}}{v_{p_{1}}^{3}} e^{-v^{2} / v_{p_{1}}^{2}}=v^{2} / v_{p_{2}} e^{-v_{2} / v_{2_{2}}} v_{p_{1}}^{2}=\frac{2 k T_{0}}{m}, v_{p_{2}}^{2}=\frac{2 k T_{0}}{m} \eta\) now

\[
e^{-\frac{m v^{2}}{2 k T_{0}}\left(1-\frac{1}{\eta}\right)}=\frac{1}{\eta^{3 / 2}} \text { or } \frac{m v^{2}}{2 k T_{0}}\left(1-\frac{1}{\eta}\right)=\frac{3}{2} \ln \eta
\]

Thus

\[
v=\sqrt{\frac{3 k T_{0}}{m}} \sqrt{\frac{\ln \eta}{1-1 / \eta}}
\]

\(2.87 \quad v_{p N}=\sqrt{\frac{2 k T}{m_{N}}}=\sqrt{\frac{2 R T}{M_{N}}}, \quad v_{p_{0}}=\sqrt{\frac{2 R T}{M_{0}}}\)

\[
v_{p_{N}}-v_{p_{0}}=\Delta v=\sqrt{\frac{2 R T}{M_{N}}}\left(1-\sqrt{\frac{M_{N}}{M_{0}}}\right)
\]

\(T=\frac{M_{N}(\Delta v)^{2}}{2 R\left(1-\sqrt{\frac{M_{N}}{M_{0}}}\right)^{2}}=\frac{m_{N}(\Delta v)^{2}}{2 k\left(1-\sqrt{\frac{m_{N}}{M_{0}}}\right)^{2}}=363 \mathrm{~K}\)\\
2.88\\
\(\frac{v^{2}}{v_{P_{H}}^{3}} e^{-v^{2} / v_{P_{H}}^{2}}=\frac{v^{2}}{v_{P_{H e}}^{3}} e^{-v^{2} / v_{P_{H e}}^{2}}\) or \(e^{v^{2}\left(\frac{m_{H e}}{2 k T}-\frac{m_{H}}{2 k T}\right)}=\left(\frac{m_{H e}}{m_{H}}\right)^{3 / 2}\)\\
\(v^{2}=3 k T \frac{\ln m_{H_{e}} / m_{H}}{m_{H_{e}}-m_{H}}\), Putting the values we get \(v=1.60 \mathrm{~km} / \mathrm{s}\)\\
\(2.89 d N(v)=\frac{N 4}{\sqrt{\pi}} \frac{v^{2} d v}{v_{p}^{3}} e^{-v^{2} / v_{p}^{2}}\)\\
For a given range \(v\) to \(v+d v\) (i.e. given \(v\) and \(d v\) ) this is maximum when

\[
\frac{\delta}{\delta v_{p}} \frac{d N(v)}{N v^{2} d v}=0=\left(-3 v_{p}^{-4}+\frac{2 v^{2}}{v_{p}^{6}}\right) e^{-v^{2} / v_{p}^{2}}
\]

or,

\[
v^{2}=\frac{3}{2} v_{p}^{2}=\frac{3 k T}{m} . \text { Thus } T=\frac{1}{3} \frac{m v^{2}}{k}
\]

\(2.90 d^{3} v=2 \pi v_{\perp} d v_{\perp} d v_{x}\)\\
Thus \(\quad d n(v)=N\left(\frac{m}{2 \pi k T}\right)^{3 / 2} e^{-\frac{m}{2 k T}}\left(v_{x}^{2}+v_{\perp}^{2}\right) d v_{x} 2 \pi v_{\perp} d v_{\perp}\)\\
\(2.91\left\langle v_{x}\right\rangle=0\) by symmetry

\[
\begin{gathered}
\langle | v_{x}| \rangle=\int_{-\infty}^{\infty}\left|v_{x}\right| e^{-\frac{m v_{x}^{2}}{2 k T}} d v_{x} / \int_{0}^{\infty} e^{\frac{-m v_{x}^{2}}{2 k T}} d v_{x}=\int_{0}^{\infty} v_{x} e^{\frac{-m v_{x}^{2}}{2 k T}} d v_{x} / \int_{0}^{\infty} e^{\frac{-m v^{2}}{2 k T}} d v_{x} \\
=\sqrt{\frac{2 k T}{m}} \int_{0}^{\infty} u e^{-u^{2}} d u / \int_{0}^{\infty} e^{-u^{2}} d u \\
=\sqrt{\frac{2 k T}{m}} \int_{0}^{\infty} \frac{1}{2} e^{-x} d x / \int_{0}^{\infty} e^{-x} \frac{d x}{2 \sqrt{x}} \\
=\sqrt{\frac{2 k T}{m}} \Gamma(1) / \Gamma\left(\frac{1}{2}\right)=\sqrt{\frac{2 k T}{m \pi}}
\end{gathered}
\]

\(2.92\left\langle v_{x}^{2}\right\rangle=\int_{0}^{\infty} v_{x}^{2} e^{-\frac{m v_{x}^{2}}{2 k t}} d v_{x} / \int_{0}^{\infty} e^{-\frac{m v_{x}^{2}}{2 k t}} d v_{x}\)

\[
\begin{gathered}
=\frac{2 k T}{m} \int_{0}^{\infty} x e^{-x} \frac{d x}{2 \sqrt{x}} / \int_{0}^{\infty} e^{-x} \frac{d x}{2 \sqrt{x}} \\
=\frac{2 k T}{m} \Gamma\left(\frac{3}{2}\right) / \Gamma\left(\frac{1}{2}\right)=\frac{k T}{m}
\end{gathered}
\]

2.93 Here \(v d A=\) No. of molecules hitting an area \(d A\) of the wall per second

\[
=\int_{0}^{\infty} d N\left(v_{x}\right) v_{x} d A
\]

or,

\[
\begin{aligned}
& v= \int_{0}^{\infty} n\left(\frac{m}{2 \pi k T}\right)^{1 / 2} e^{-\frac{m v_{x}^{2}}{2 k T}} v_{x} d v_{x} \\
&= \int_{0}^{\infty} \frac{n}{\sqrt{\pi}}\left(\frac{2 k T}{m}\right)^{1 / 2} e^{-u^{2}} u d u \\
&=\frac{1}{2} n \sqrt{\frac{2 k T}{m \pi}}=n \sqrt{\frac{k T}{2 m \pi}}=\frac{1}{4} n<v>, \\
&\left(\text { where }<v>=\sqrt{\frac{8 k T}{m \pi}}\right)
\end{aligned}
\]

\includegraphics[max width=\textwidth, alt={}, center]{9f3c5a8e-ac5d-42f4-8eb6-ea2af623f2d8-228_185_356_531_991}\\
2.94 Let, \(d n\left(v_{x}\right)=n\left(\frac{m}{2 \pi k T}\right)^{1 / 2} e^{-m v_{x}^{2} / 2 k T} d v_{x}\)\\
be the number of molecules per unit volume with \(x\) component of velocity in the range \(v_{x}\) to \(v_{x}+d v_{x}\)

Then

\[
\begin{aligned}
p & =\int_{0}^{\infty} 2 m v_{x} \cdot v_{x} d n\left(v_{x}\right) \\
& =\int_{0}^{\infty} 2 m v_{x}^{2} n\left(\frac{m}{2 \pi k T}\right)^{1 / 2} e^{-m v_{x}^{2} / 2 k T} d v_{x} \\
& =2 m n \frac{1}{\sqrt{\pi}} \frac{2 k T}{m} \int_{0}^{\infty} u^{2} e^{-u^{2}} d u \\
& =\frac{4}{\sqrt{\pi}} n k T \cdot \int_{0}^{\infty} x e^{-x} \frac{d x}{2 \sqrt{x}}=n k T
\end{aligned}
\]

\(2.95<\frac{1}{v}>=\int_{0}^{\infty}\left(\frac{m}{2 \pi k T}\right)^{3 / 2} e^{-\frac{m v^{2}}{2 k T}} 4 \pi v^{2} d v \frac{1}{v}\)

\[
\begin{gathered}
=\left(\frac{m}{2 \pi k T}\right)^{3 / 2} 4 \pi \frac{1}{2} \frac{2 k T}{m} \int_{0}^{\infty} e^{-x} d x \\
=2\left(\frac{m}{2 \pi k T}\right)^{1 / 2}=\left(\frac{2 m}{\pi k T}\right)^{1 / 2}=\left(\frac{16}{\pi^{2}} \frac{m \pi}{8 k T}\right)^{1 / 2}=\frac{4}{\pi<v>}
\end{gathered}
\]

\(d N(v)=N\left(\frac{m}{2 \pi k T}\right)^{3 / 2} e^{-m v^{2} / 2 k T} 4 \pi v^{2} d v=d N(\varepsilon)=\frac{d N(\varepsilon)}{d \varepsilon} d \varepsilon\)\\
or, \(\quad \frac{d N(\varepsilon)}{d \varepsilon}=N\left(\frac{m}{2 \pi k T}\right)^{3 / 2} e^{-m v^{2} / 2 R T} 4 \pi v^{2} \frac{d v}{d \varepsilon}\)\\
Now,

\[
\varepsilon=\frac{1}{2} m v^{2} \text { so } \frac{d v}{d \varepsilon}=\frac{1}{m v}
\]

or,

\[
\begin{aligned}
\frac{d N(\varepsilon)}{d \varepsilon} & =N\left(\frac{m}{2 \pi k T}\right)^{3 / 2} e^{-\varepsilon / k T} 4 \pi \sqrt{\frac{2 \varepsilon}{m}} \frac{1}{m} \\
& =N \frac{2}{\sqrt{\pi}}(k T)^{-3 / 2} e^{-\varepsilon / k T} \varepsilon^{1 / 2}
\end{aligned}
\]

i.e.

\[
d N(\varepsilon)=N \frac{2}{\sqrt{\pi}}(k T)^{-3 / 2} e^{-\varepsilon / k T} \varepsilon^{1 / 2} d \varepsilon
\]

The most probable kinetic energy is given from\\
\(\frac{d}{d \varepsilon} \frac{d N(\varepsilon)}{d \varepsilon}=0\) or, \(\frac{1}{2} \varepsilon^{-1 / 2} e^{-\varepsilon / k T}-\frac{\varepsilon^{1 / 2}}{k T} e^{\varepsilon / k T}=0\) or \(\varepsilon=\frac{1}{2} k T=\varepsilon_{p r}\)\\
The corresponding velocity is \(v=\sqrt{\frac{k T}{m}} \neq v_{p r}\)\\
2.97 The mean kinetic energy is\\
\(<\varepsilon>=\int_{0}^{\infty} \varepsilon^{3 / 2} e^{-\varepsilon / k T} d \varepsilon / \int_{0}^{\infty} \varepsilon^{1 / 2} e^{-\varepsilon / k T} d \varepsilon=k T \frac{\Gamma(5 / 2)}{\Gamma(3 / 2)}=\frac{3}{2} k T\)\\
Thus

\[
\begin{aligned}
& \frac{\delta N}{N}=\int_{\frac{3}{2} k T(1-\delta \eta)}^{\frac{3}{2}(1+\delta \eta) k T} \frac{2}{\sqrt{\pi}}(k T)^{-3 / 2} e^{-\varepsilon / k T} \varepsilon^{1 / 2} d \varepsilon \\
& =\frac{2}{\sqrt{\pi}} e^{-3 / 2}\left(\frac{3}{2}\right)^{3 / 2} 2 \delta \eta=3 \sqrt{\frac{6}{\pi}} e^{3 / 2} \delta \eta
\end{aligned}
\]

If \(\delta \eta=1 \%\) this gives \(0.9 \%\)\\
\(2.98 \frac{\Delta N}{N}=\frac{2}{\sqrt{\pi}}(k T)^{-3 / 2} \int_{\varepsilon_{0}}^{\infty} \sqrt{\varepsilon} e^{-\varepsilon / k T} d \varepsilon\)

\[
\approx \frac{2}{\sqrt{\pi}}(k T)^{-3 / 2} \sqrt{\varepsilon_{0}} \int_{\varepsilon_{0}}^{\infty} e^{-\varepsilon / k T} d \varepsilon\left(\varepsilon_{0} \gg k T\right)
\]

\[
=\frac{2}{\sqrt{\pi}}(k T)^{-3 / 2} \sqrt{\varepsilon_{0}} k T e^{-\varepsilon_{0} / k T}=2 \sqrt{\frac{\varepsilon_{0}}{\pi k T}} e^{-\varepsilon_{0} / k T}
\]

(In evaluating the integral, we have taken out \(\sqrt{\varepsilon}\) as \(\sqrt{\varepsilon_{0}}\) since the integral is dominated by the lower limit.)\\
2.99 (a) \(F(v)=A v^{3} e^{-m v^{2} / 2 k T}\)

For the most probable value of the velocity \(\frac{d F(v)}{d v}=0\) or \(3 A v^{2} e^{-m v^{2} / 2 k T}-A v^{3} \frac{2 m v}{2 k T} e^{-m v^{2} / 2 k T}=0\)

So,

\[
v_{p r}=\sqrt{\frac{3 k T}{m}}
\]

This should be compared with the value \(v_{p r}=\sqrt{\frac{2 k T}{m}}\) for the Maxwellian distribution.\\
(b) In terms of energy, \(\varepsilon=\frac{1}{2} m v^{2}\)

\[
\begin{gathered}
F(\varepsilon)=A v^{3} e^{-m v^{2} / 2 k T} \frac{d v}{d \varepsilon} \\
=A\left(\frac{2 \varepsilon}{m}\right)^{3 / 2} e^{-\varepsilon / k T} \frac{1}{\sqrt{2 m \varepsilon}}=A \frac{2 \varepsilon}{m^{2}} e^{-\varepsilon / k T}
\end{gathered}
\]

From this the probable energy comes out as follows : \(F^{\prime}(\varepsilon)=0\) implies

\[
\frac{2 A}{m^{2}}\left(e^{-\varepsilon / k T}-\frac{\varepsilon}{k T} e^{-\varepsilon / k T}\right)=0, \text { or, } \quad \varepsilon_{p r}=k T
\]

2.100 The number of molecules reaching a unit area of wall at angle between \(\theta\) and \(\theta+d \theta\) to its normal per unit time is

\[
\begin{aligned}
& \quad d v=\int_{v=0}^{v=\infty} d n(v) \frac{d \Omega}{4 \pi} v \cos \theta \\
& =\int_{0}^{\infty} n\left(\frac{m}{2 \pi k T}\right)^{3 / 2} e^{-m v^{2} / 2 k T} v^{3} d v \sin \theta \cos \theta d \theta \times 2 \pi \\
& =n\left(\frac{2 k T}{m \pi}\right)^{1 / 2} \int_{0}^{\infty} e^{-x} x d x \sin \theta \cos \theta d \theta=n\left(\frac{2 k T}{m \pi}\right)^{1 / 2} \sin \theta \cos \theta d \theta
\end{aligned}
\]

\includegraphics[max width=\textwidth, alt={}, center]{9f3c5a8e-ac5d-42f4-8eb6-ea2af623f2d8-230_184_343_1314_993}\\
2.101 Similarly the number of molecules reaching the wall (per unit area of the wall with velocities in the interval \(v\) to \(v+d v\) per unit time is

\[
d v=\int_{\theta=0}^{\theta=\pi / 2} d n(v) \frac{d \Omega}{4 \pi} v \cos 0
\]

\[
\begin{gathered}
=\int_{\theta=0}^{\theta=\pi / 2} n\left(\frac{m}{2 \pi k T}\right)^{3 / 2} e^{-m v^{2} / 2 k T} v^{3} d v \sin \theta \cos \theta d \theta \times 2 \pi \\
=n \pi\left(\frac{m}{2 \pi k T}\right)^{3 / 2} e^{-m v^{2} / 2 k T} v^{3} d v
\end{gathered}
\]

2.102 If the force exerted is \(F\) then the law of variation of concentration with height reads \(n(Z)=n_{0} e^{-F Z / k T}\) So, \(\eta=e^{F \Delta h / k T} \quad\) or \(\quad F=\frac{k T \ln \eta}{\Delta h}=9 \times 10^{-20} \mathrm{~N}\)\\
2.103 Here \(F=\frac{\pi}{6} d^{3} \Delta \rho g=\frac{R T \ln \eta}{N_{a} h}\) or \(N_{a}=\frac{6 R T \ln \eta}{\pi d^{3} g \Delta \rho h}\)

In the problem, \(\frac{\eta}{\eta_{0}}=1.39\) here\\
\(T=290 \mathrm{~K}, \eta=2, h=4 \times 10^{-5} \mathrm{~m}, d=4 \times 10^{-7} \mathrm{~m}, g=9 \cdot 8 \mathrm{~m} / \mathrm{s}^{2}, \Delta \rho=0 \cdot 2 \times 10^{3} \mathrm{~kg} / \mathrm{m}^{3}\) and \(R=8.31 \mathrm{~J} / \mathrm{k}\)\\
Hence, \(\quad N_{a}=\frac{6 \times 8.31 \times 290 \times \ln 2}{\pi \times 64 \times 9.8200 \times 4} \times 10^{26}=6.36 \times 10^{23} \mathrm{~mole}^{-1}\)\\
\(2.104 \eta=\frac{\text { concetration of } H_{2}}{\text { concentration of } N_{2}}=\eta_{0} \frac{e^{-M_{H_{2}} g h / R T}}{e^{-M_{N_{2}}-g h / R T}}=\eta_{0} e^{\left(M_{N_{2}}-M_{H_{2}}\right)} g h / R T\)\\
So more \(N_{2}\) at the bottom, \(\left(\frac{\eta_{1}}{\eta_{0}}=1.39\right.\) here \()\)\\
\(2.105 n_{1}(h)=n_{1} e^{-m_{1} g h / k t}, n_{2}(h)=n_{2} e^{-m_{2} g h / k T}\)\\
They are equal at a height \(h\) where \(\frac{n_{1}}{n_{2}}=e^{g h\left(m_{1}-m_{2}\right) V k T}\)

\[
\text { or } \quad h=\frac{k T}{g} \frac{\ln n_{1}-\ln n_{2}}{m_{1}-m_{2}}
\]

2.106 At a temperature \(T\) the concentration \(n(z)\) varies with height according to

\[
n(z)=n_{0} e^{-m g z / k T}
\]

This means that the cylinder contains \(\int_{0}^{\infty} n(z) d z\)

\[
=\int_{0}^{\infty} n_{0} e^{-m g z / k T} d z=\frac{n_{0} k T}{m g}
\]

particles per unit area of the base. Clearly this cannot change. Thus \(n_{0} k T=p_{0}=\) pressure at the bottom of the cylinder must not change with change of temperature.

\[
<U>=\frac{\int_{0}^{\infty} m g z e^{-m g z / k t} d z}{\int_{0}^{\infty} e^{-m g z / k T} d z}=k T \frac{\int_{0}^{\infty} x e^{-x} d x}{\int_{0}^{\infty} e^{-x} d x}=k T \frac{\Gamma(2)}{\Gamma(1)}=k T
\]

When there are many kinds of molecules, this formula holds for each kind and the average energy

\[
<U>=\frac{\sum f_{i} k T}{\sum f_{i}}=k T
\]

where \(f_{i} \alpha\) fractional concentration of each kind at the ground level.\\
2.108 The constant acceleration is equivalent to a pseudo force wherein a concentration gradient is set up. Then\\
or

\[
\begin{gathered}
e^{-M_{A} w l / R T}=1-\eta \\
w=-\frac{R T \ln (1-\eta)}{M_{A} l}=\frac{\eta R T}{M_{A} l}=70 \mathrm{~g}
\end{gathered}
\]

2.109 In a centrifuge rotating with angular velocity \(\omega\) about an axis, there is a centrifugal acceleration \(\omega^{2} r\) where \(r\) is the radial distance from the axis. In a fluid if there are suspended colloidal particles they experience an additional force. If \(m\) is the mass of each particle then its volume is \(\frac{m}{\rho}\) and the excess force on this particle is\\
\(\frac{m}{\rho}\left(\rho-\rho_{0}\right) \omega^{2} r\) outward corresponding to a potential energy \(-\frac{m}{2 \rho}\left(\rho-\rho_{0}\right) \omega^{2} r^{2}\)\\
This gives rise to a concentration variation

\[
n(r)=n_{0} \exp \left(+\frac{m}{2 \rho k T}\left(\rho-\rho_{0}\right) \omega^{2} r^{2}\right)
\]

Thus

\[
\frac{n\left(r_{2}\right)}{n\left(r_{1}\right)}=\eta=\exp \left(+\frac{M}{2 \rho R T}\left(\rho-\rho_{0}\right) \omega^{2}\left(r_{2}^{2}-r_{1}^{2}\right)\right)
\]

where

\[
\frac{m}{k}=\frac{M}{R}, M=N_{A} m \text { is the molecular weight }
\]

Thus

\[
M=\frac{2 \rho R T \ln \eta}{\left(\rho-\rho_{0}\right) \omega^{2}\left(r_{2}^{2}-r_{1}^{2}\right)}
\]

2.110 The potential energy associated with each molecule is : \(-\frac{1}{2} m \omega^{2} r^{2}\) and there is a concentration variation

\[
n(r)=n_{0} \exp \left(\frac{m \omega^{2} r^{2}}{2 k T}\right)=n_{0} \exp \left(\frac{M \omega^{2} r^{2}}{2 R T}\right)
\]

Thus

\[
\eta=\exp \left(\frac{M \omega^{2} l^{2}}{2 R T}\right) \text { or } \omega=\sqrt{\frac{2 R T}{M l^{2}} \ln \eta}
\]

Using \(M=12+32=44 \mathrm{gm}, l=100 \mathrm{~cm}, R=8.31 \times 10^{7} \frac{\mathrm{erg}}{{ }^{\circ} \mathrm{K}}, T=300\), we get \(\omega=280\) radians per second.\\
2.111 Here \(n(r)=n_{0} \exp \left(-\frac{a r^{2}}{k T}\right)\)\\
(a) The number of molecules located at the distance between \(r\) and \(r+d r\) is

\[
4 \pi r^{2} d r n(r)=4 \pi n_{0} \exp \left(-\frac{a r^{2}}{k T}\right) r^{2} d r
\]

(b) \(r_{p r}\) is given by \(\frac{d}{d r} r^{2} n(r)=0\) or, \(2 r-\frac{2 a r^{3}}{k T}=0\) or \(r_{p r}=\sqrt{\frac{k T}{a}}\)\\
(c) The fraction of molecules lying between \(r\) and \(r+d r\) is

Thus

\[
\begin{gathered}
\frac{d N}{N}=\frac{4 \pi r^{2} d r n_{0} \exp \left(-a r^{2} / k T\right)}{\int_{0}^{\infty} 4 \pi r^{2} d r n_{0} \exp \left(-a r^{2} / k T\right)} \\
\int_{0}^{\infty} 4 \pi r^{2} d r \exp \left(-\frac{a r^{2}}{k T}\right)=\left(\frac{k T}{a}\right)^{3 / 2} 4 \pi \int_{0}^{\infty} x \frac{d x}{2 \sqrt{x}} \exp (-x) \\
=\left(\frac{k T}{a}\right)^{3 / 2} 2 \pi \Gamma\left(\frac{3}{2}\right)=\left(\frac{\pi k T}{a}\right)^{3 / 2}
\end{gathered}
\]

(d)

\[
\frac{d N}{N}=\left(\frac{a}{\pi k T}\right)^{3 / 2} 4 \pi r^{2} d r \exp \left(\frac{-a r^{2}}{k T}\right)
\]

So

\[
d N=N\left(\frac{a}{\pi k T}\right)^{3 / 2} 4 \pi r^{2} d r \exp \left(\frac{-a r^{2}}{k T}\right)
\]

\subsection*{2.112}
\[
n(r)=N\left(\frac{a}{\pi k T}\right)^{1 / 2} \exp \left(\frac{-a r^{2}}{k T}\right)
\]

When \(T\) decreases \(\eta\) times \(n(0)=n_{0}\) will increase \(\eta^{3 / 2}\) times.\\
Write \(U=a r^{2}\) or \(r=\sqrt{\frac{U}{a}}\), so \(d r=\sqrt{\frac{1}{a}} \frac{d U}{2 \sqrt{U}}=\frac{d U}{2 \sqrt{a U}}\)\\
so \(\quad d N=n_{0} 4 \pi \frac{U}{a} \frac{d U}{2 \sqrt{a U}} \exp \left(\frac{U}{k T}\right)\)\\
\(=2 \pi n_{0} a^{-3 / 2} U^{1 / 2} \exp \left(\frac{-U}{k T}\right) d U\)\\
The most probable value of \(U\) is given by

\[
\frac{d}{d U}\left(\frac{d N}{d U}\right)=0=\left(\frac{1}{2 \sqrt{u}}-\frac{U^{1 / 2}}{k T}\right) \exp \left(\frac{-U}{k T}\right) \quad \text { or, } \quad U_{p r}=\frac{1}{2} k T
\]

From 2.111 (b), the potential energy at the most probable distance is \(k T\).

\subsection*{2.4 THE SECOND LAW OF THERMODYNAMICS. ENTROPY}
2.113 The efficiency is given by

\[
\eta=\frac{T_{1}-T_{2}}{T_{1}}, T_{1}>T_{2}
\]

Now in the two cases the efficiencies are

\[
\begin{aligned}
\eta_{h} & =\frac{T_{1}+\Delta T-T_{2}}{T_{1}+\Delta T}, \quad T_{1} \text { increased } \\
\eta_{l} & =\frac{T_{1}-T_{2}+\Delta T}{T_{1}}, T_{2} \text { decreased }
\end{aligned}
\]

Thus

\[
\eta_{h}<n_{l}
\]

2.114 For \(H_{2}, \gamma=\frac{7}{5}\)

\[
\begin{aligned}
& p_{1} V_{1}=p_{2} V_{2}, p_{3} V_{3}=p_{4} V_{4} \\
& p_{2} V_{2}^{\gamma}=p_{3} V_{3}^{\gamma}, p_{1} V_{1}^{\gamma}=p_{4} V_{4}^{\gamma}
\end{aligned}
\]

Define \(n\) by \(V_{3}=n V_{2}\)\\
Then \(p_{3}=p_{2} n^{-\gamma}\) so\\
\(p_{4} V_{4}=p_{3} V_{3}=p_{2} V_{2} n^{1-\gamma}=p_{1} V_{1} n^{1-\gamma}\)\\
\includegraphics[max width=\textwidth, alt={}, center]{9f3c5a8e-ac5d-42f4-8eb6-ea2af623f2d8-234_422_545_680_831}\\
\(p_{4} V_{4}^{\gamma}=p_{1} V_{1}^{\gamma}\) so \(V_{4}^{1-\gamma}=V_{1}^{1-\gamma} n^{1-\gamma}\) or \(V_{4}=n V_{1}\)

Also

\[
Q_{1}=p_{2} V_{2} \ln \frac{V_{2}}{V_{1}}, Q_{2}^{\prime}=p_{3} V_{3} \ln \frac{V_{3}}{V_{4}} n^{1-\gamma}=p_{2} v_{2} \ln \frac{V_{3}}{v_{4}}
\]

Finally

\[
\eta=1-\frac{Q_{2}}{Q_{1}}=1-n^{1-\gamma}=0.242
\]

(b) Define \(n\) by \(p_{3}=\frac{p_{2}}{n}\)

\[
p_{2} V_{2}^{\gamma}=\frac{p_{2}}{n} V_{3}^{\gamma} \quad \text { or } \quad V_{3}=n^{1 / \gamma} V_{2}
\]

So we get the formulae here by \(n \rightarrow n^{1 / \gamma}\) in the previous case.

\[
\eta=1-n^{(1 / \gamma)-1}=1-n^{-\frac{2}{7}} \approx 0.18
\]

2.115 Used as a refrigerator, the refrigerating efficiency of a heat engine is given by\\
\(\varepsilon=\frac{Q_{2}{ }^{\prime}}{A}=\frac{Q_{2}{ }^{\prime}}{Q_{1}-Q_{2}}=\frac{Q_{2}{ }^{\prime} / Q_{1}}{1-\frac{Q_{2}{ }^{\prime}}{Q_{1}}}=\frac{1-\eta}{\eta}=9\) here, where \(\eta\) is the efficiency of the heat engine.\\
\includegraphics[max width=\textwidth, alt={}, center]{9f3c5a8e-ac5d-42f4-8eb6-ea2af623f2d8-234_437_435_1668_968}\\
2.116 Given \(V_{2}=n V_{1}, V_{4}=n V_{3}\)\\
\(Q_{1}=\) Heat taken at the upper temperature

\[
=R T_{1} \ln n+R T_{2} \ln n=R\left(T_{1}+T_{2}\right) \ln n
\]

Now \(\quad T_{1} V_{2}^{\gamma-1}=T_{2} V_{3}^{\gamma-1}\) or \(V_{3}=\left(\frac{T_{1}}{T_{2}}\right)^{\frac{1}{\gamma-1}} V_{2}\)\\
Similarly \(\quad V_{5}=\left(\frac{T_{2}}{T_{3}}\right)^{\frac{1}{\gamma-1}} V_{4}, V_{6}=\left(\frac{T_{1}}{T_{3}}\right)^{\frac{1}{\gamma-1}} V_{1}\)\\
\includegraphics[max width=\textwidth, alt={}, center]{9f3c5a8e-ac5d-42f4-8eb6-ea2af623f2d8-235_408_431_121_963}

Thus \(Q_{2}=\) heat ejected at the lower temperature \(=-R T_{3} \ln \frac{V_{6}}{V_{5}}\)

\[
\begin{aligned}
& =-R T_{3} \ln \left(\frac{T_{1}}{T_{2}}\right)^{\frac{1}{\gamma-1}} \frac{V_{1}}{V_{4}}=-R T_{3} \ln \left(\frac{T_{1}}{T_{2}}\right)^{\frac{1}{\gamma-1}} \frac{V_{2}}{n^{2} V_{3}} \\
& \qquad=-R T_{3} \ln \left(\frac{T_{1}}{T_{2}}\right)^{\frac{1}{\gamma-1}} \frac{1}{n^{2}}\left(\frac{T_{1}}{T_{2}}\right)^{-\frac{1}{\gamma-1}}=2 R T_{3} \ln n \\
& \text { Thus } \quad \eta=1-\frac{2 T_{3}}{T_{1}+T_{2}}
\end{aligned}
\]

\(2.117 Q_{2}^{\prime}=C_{V}\left(T_{2}-T_{3}\right)=\frac{C_{V}}{R} V_{2}\left(p_{2}-p_{3}\right)\)

\[
Q_{1}=\frac{C_{V}}{R} V_{1}\left(p_{1}-p_{4}\right)
\]

Thus \(\quad \eta=1-\frac{V_{2}\left(p_{2}-p_{3}\right)}{V_{1}\left(p_{1}-p_{4}\right)}\)\\
On the other hand,\\
\(6 p_{1} V_{1}^{\gamma}=p_{2} V_{2}^{\gamma}, p_{3} V_{2}^{\gamma}=p_{4} V_{1}^{\gamma}\) also \(V_{2}=n V_{1}\)\\
Thus \(p_{1}=p_{2} n^{\gamma}, p_{4}=p_{3} n^{\gamma}\)\\
\includegraphics[max width=\textwidth, alt={}, center]{9f3c5a8e-ac5d-42f4-8eb6-ea2af623f2d8-235_403_436_1072_911}\\
and \(\eta=1-n^{1-\gamma}\), with \(\gamma=\frac{7}{5}\) for \(N_{2}\) this is \(\eta=0.602\)\\
\(2.118 Q_{1}=\frac{C_{p}}{R} p_{1}\left(V_{2}-V_{1}\right), Q_{2}^{\prime}=\frac{C_{p}}{R} p_{2}\left(V_{3}-V_{4}\right)\)\\
So \(\quad \eta=1-\frac{p_{2}\left(V_{3}-V_{4}\right)}{p_{1}\left(V_{2}-V_{1}\right)}\)\\
Now \(p_{1}=n p_{2}, p_{1} V_{2}^{\gamma}\) or \(V_{3}=n^{\frac{1}{\gamma}} V_{2}\)\\
\(p_{2} V_{4}^{\gamma}=p_{1} V_{1}^{\gamma} \quad\) or \(\quad V_{4}=n^{\frac{1}{\gamma}} V_{1}\)\\
so \(\eta=1-\frac{1}{n} \cdot n^{\frac{1}{\gamma}}=1-n^{\frac{1}{\gamma}-1}\)\\
\includegraphics[max width=\textwidth, alt={}, center]{9f3c5a8e-ac5d-42f4-8eb6-ea2af623f2d8-235_378_559_1656_824}\\
2.119 Since the absolute temperature of the gas rises \(n\) times both in the isochoric heating and in the isobaric expansion

\[
\begin{gathered}
p_{1}=n p_{2} \text { and } V_{2}=n V_{1} . \text { Heat taken is } \\
Q_{1}=Q_{11}+Q_{12}
\end{gathered}
\]

where \(Q_{11}=C_{p}(n-1) T_{1}\) and \(Q_{12}=C_{V} T_{1}\left(1-\frac{1}{n}\right)\)

Heat rejected is

\[
\begin{aligned}
& Q_{2}^{\prime}=Q_{21}^{\prime}+Q_{22}^{\prime} \text { where } \\
& Q_{21}^{\prime}=C_{V} T_{1}(n-1), Q_{22}^{\prime}=C_{p} T_{1}\left(1-\frac{1}{n}\right) \\
& \text { Thus } \eta=1-\frac{Q_{2}^{\prime}}{Q_{1}}=1-\frac{C_{V}(n-1)+C_{p}\left(1-\frac{1}{n}\right)}{C_{p}(n-1)+C_{V}\left(1-\frac{1}{n}\right)} \\
& =1-\frac{n-1+\gamma\left(1-\frac{1}{n}\right)}{\gamma(n-1)+\left(1-\frac{1}{n}\right)}=1-\frac{1+\frac{\gamma}{n}}{\gamma+\frac{1}{n}}=1-\frac{n+\gamma}{1+n \gamma}
\end{aligned}
\]

\includegraphics[max width=\textwidth, alt={}, center]{9f3c5a8e-ac5d-42f4-8eb6-ea2af623f2d8-236_478_544_363_855}\\
2.120 (a) Here \(p_{2}=n p_{1}, p_{1} V_{1}=p_{0} V_{0}\),

\[
\begin{gathered}
n p_{1} V_{1}^{\gamma}=p_{0} V_{0}^{\gamma} \\
Q_{2}^{\prime}=R T_{0} \ln \frac{V_{0}}{V_{1}}, Q_{1}=C_{V} T_{0}(n-1)
\end{gathered}
\]

But \(n V_{1}^{\gamma-1}=V_{0}^{\gamma-1}\) or, \(V_{1}=V_{0} n^{\frac{-1}{\gamma-1}}\)

\[
Q_{2}^{\prime}=R T_{0} \ln n^{\frac{1}{\gamma-1}}=\frac{R T_{0}}{\gamma-1} \ln n
\]

Thus \(\eta=1-\frac{\ln n}{n-1}\), on using \(C_{V}=\frac{R}{\gamma-1}\)\\
(b) Here \(V_{2}=n V_{1}, p_{1} V_{1}=p_{0} V_{0}\)\\
and

\[
p_{1}\left(n V_{1}\right)^{\gamma}=p_{0} V_{0}^{\gamma}
\]

i.e. \(n^{\gamma} V_{1}^{\gamma-1}=V_{0}^{\gamma-1}\) or \(V_{1}=n^{-\frac{\gamma}{\gamma-1}} V_{0}\)\\
Also \(Q_{1}=C_{p} T_{0}(n-1), Q_{2}^{\prime}=R T_{0} \ln \frac{V_{0}}{V_{1}}\)\\
or \(Q_{2}^{\prime}=R T_{0} \ln n \frac{\gamma}{\gamma-1}=\frac{R \gamma}{\gamma-1} T_{0} \ln n=C_{p} T_{0} \ln n\)\\
Thus

\[
\eta=1-\frac{\ln n}{n-1}
\]

\includegraphics[max width=\textwidth, alt={}, center]{9f3c5a8e-ac5d-42f4-8eb6-ea2af623f2d8-236_324_545_1727_854}\\
2.121 Here the isothermal process proceeds at the maximum temperature instead of at the minimum temperature of the cycle as in 2.120.\\
\includegraphics[max width=\textwidth, alt={}, center]{9f3c5a8e-ac5d-42f4-8eb6-ea2af623f2d8-237_338_522_211_123}\\
\includegraphics[max width=\textwidth, alt={}, center]{9f3c5a8e-ac5d-42f4-8eb6-ea2af623f2d8-237_389_643_196_747}\\
(a) Here \(p_{1} V_{1}=p_{0} V_{0}, p_{2}=\frac{p_{1}}{n}\)

\[
\begin{aligned}
& p_{2} V_{1}^{\gamma}=p_{0} V_{0}^{\gamma} \text { or } p_{1} V_{1}^{\gamma}=n p_{0} V_{0}^{\gamma} \\
& V_{1}^{\gamma-1}=n V_{0}^{\gamma-1} \text { or } V_{1}=V_{0} n \frac{1}{\gamma-1}
\end{aligned}
\]

i.e.\\
\(Q_{2}^{\prime}=C_{V} T_{0}\left(1-\frac{1}{n}\right), Q_{1}=R T_{0} \ln \frac{V_{1}}{V_{0}}=\frac{R T_{0}}{\gamma-1} \ln n=C_{V} T_{0} \ln n\).\\
Thus

\[
\eta=1-\frac{Q_{2}^{\prime}}{Q_{1}}=1-\frac{n-1}{n \ln n}
\]

(b) Here \(V_{2}=\frac{V_{1}}{n}, p_{0} V_{0}=p_{1} V_{1}\)\\
\(p_{0} V_{0}^{\gamma}=p_{1} V_{2}^{\gamma}=p_{1} n^{-\gamma} V_{1}^{\gamma}=V_{0}^{\gamma-1} n^{-\gamma} V_{1}^{\gamma-1}\) or \(V_{1}=n^{(\gamma / \gamma-1)} V_{0}\)\\
\(Q_{2}^{\prime}=C_{p} T_{0}\left(1-\frac{1}{n}\right), Q_{1}=R T_{0} \ln \frac{V_{1}}{V_{0}}=\frac{R y}{\gamma-1} T_{0} \ln n=C_{p} T_{0} \ln n\)\\
Thus

\[
\eta=1-\frac{n-1}{n \ln n}
\]

2.122 The section from \(\left(\dot{p}_{1}, V_{1}, T_{0}\right)\) to \(\left(p_{2}, V_{2}, T_{0} / n\right)\) is a polytropic process of index \(\alpha\). We shall assume that the corresponding specific heat \(C\) is +ve .\\
Here, \(d Q=C d T=C_{V} d T+p d V\)\\
Now \(p V^{\alpha}=\) constant or \(T V^{\alpha-1}=\) constant.\\
so \(p d V=\frac{R T}{V} d V=-\frac{R}{\alpha-1} d T\)\\
Then \(C=C_{V}-\frac{R}{\alpha-1}=R\left(\frac{1}{\gamma-1}-\frac{1}{\alpha-1}\right)\)\\
We have \(p_{1} V_{1}=R T_{0}=p_{2} V_{2}=\frac{R T_{0}}{n}=\frac{p_{1} V_{1}}{n}\)\\
\includegraphics[max width=\textwidth, alt={}, center]{9f3c5a8e-ac5d-42f4-8eb6-ea2af623f2d8-237_401_507_1566_905}

\[
p_{0} V_{0}=p_{1} V_{1}=n p_{2} V_{2}, p_{0} V_{0}^{\gamma}=p_{2} V_{2}^{\gamma},
\]

\(p_{1} V_{1}^{\alpha}=p_{2} V_{2}^{\alpha} \quad\) or \(\quad V_{0}^{\gamma-1}=\frac{1}{n} V_{2}^{\gamma-1} \quad\) or \(\quad V_{2}=V_{0} n \frac{1}{\gamma-1}\)\\
\(V_{1}^{\alpha-1}=\frac{1}{n} V_{2}^{\alpha-1} \quad\) or \(\quad V_{1}=n^{-\frac{1}{\alpha-1}} V_{2}=n_{\gamma-1}^{\frac{1}{\alpha-1}}-\frac{1}{0}\)\\
Now \(Q_{2}^{\prime}=C T_{0}\left(1-\frac{1}{n}\right), Q_{1}=R T_{0} \ln \frac{V_{1}}{V_{0}}=R T_{0}\left(\frac{1}{\gamma-1}-\frac{1}{\alpha-1}\right) \ln n=C T_{0} \ln n\)\\
Thus

\[
\eta=1-\frac{n-1}{n \ln n}
\]

2.123\\
\includegraphics[max width=\textwidth, alt={}, center]{9f3c5a8e-ac5d-42f4-8eb6-ea2af623f2d8-238_411_439_507_176}\\
\includegraphics[max width=\textwidth, alt={}, center]{9f3c5a8e-ac5d-42f4-8eb6-ea2af623f2d8-238_382_464_539_749}\\
(a) Here \(Q_{2}^{\prime}=C_{p}\left(T_{1}-\frac{T_{1}}{n}\right)=C_{p} T_{1}\left(1-\frac{1}{n}\right), Q_{1}=C_{V}\left(T_{0}-\frac{T_{1}}{n}\right)\)

Along the adiabatic line

\[
T_{0} V_{0}^{\gamma-1}=T_{1}\left(n V_{0}\right)^{\gamma-1} \quad \text { or, } \quad T_{0}=T_{1} n^{\gamma-1} .
\]

so

\[
Q_{1}=C_{V} \frac{T_{1}}{n}\left(n^{\gamma}-1\right) . \text { Thus } \eta=1-\frac{\gamma(n-1)}{n^{\gamma-1}}
\]

(b) Here \(Q_{2}^{\prime}=C_{V}\left(n T_{1}-T_{0}\right), Q_{1}=C_{p} \cdot T_{1}(n-1)\)

Along the adiabatic line \(T V^{\gamma-1}=\) constant

Thus

\[
T_{0} V_{0}^{\gamma-1}=T_{1}\left(\frac{V_{0}}{n}\right)^{\gamma-1} \quad \text { or } \quad T_{1}=n^{\gamma-1} T_{0}
\]

2.124

\[
\eta=1-\frac{n^{\gamma}-1}{\gamma n^{\gamma-1}(n-1)}
\]

\begin{figure}[h]
\begin{center}
  \includegraphics[alt={},max width=\textwidth]{9f3c5a8e-ac5d-42f4-8eb6-ea2af623f2d8-238_416_1039_1626_195}
\captionsetup{labelformat=empty}
\caption{(a)}
\end{center}
\end{figure}

(b)\\
(a) \(Q_{2}^{\prime}=C_{p} T_{0}\left(1-\frac{1}{n}\right), Q_{1}^{\prime \prime}=R T_{0} \ln n, Q_{1}^{\prime}=C_{V} T_{0}\left(1-\frac{1}{n}\right), Q_{1}=Q_{1}^{\prime}+Q_{1}^{\prime \prime}\)

So

\[
\begin{gathered}
\eta=1-\frac{Q_{2}^{\prime}}{Q_{1}}=1-\frac{C_{p}\left(1-\frac{1}{n}\right)}{C_{V}\left(1-\frac{1}{n}\right)+R \ln n} \\
=1-\frac{\gamma}{1+\frac{R}{C_{V}} \frac{n \ln n}{n-1}}=1-\frac{\gamma(n-1)}{n-1+(\gamma-1) n \ln n}
\end{gathered}
\]

(b) \(Q_{1}=C_{p} T_{0}(n-1), Q^{\prime \prime}{ }_{2}=C_{V} T_{0}(n-1) \quad Q^{\prime \prime \prime}{ }_{2}=R T_{0} \ln n, Q_{2}^{\prime}=Q^{\prime \prime}{ }_{2}+Q^{\prime \prime \prime}{ }_{2}\)

So

\[
\eta=1-\frac{Q_{2}^{\prime}}{Q_{1}}=1-\frac{n-1+(\gamma-1) \ln n}{\gamma(n-1)}
\]

2.125 We have\\
\(Q_{1}^{\prime}=\tau R T_{0} \ln \nu, Q^{\prime \prime}{ }_{2}=C_{V} T_{0}(\tau-1) Q_{1}=Q_{1}^{\prime}+Q^{\prime \prime}{ }_{1}\) and \(Q^{\prime \prime \prime}{ }_{2}=R T_{0} \ln v, Q^{\prime \prime}{ }_{1}=C_{V} T_{0}(\tau-1)\)\\
as well as \(Q_{1}=Q_{1}{ }^{\prime}+Q_{1}{ }^{\prime \prime}\) and\\
\(Q_{2}{ }^{\prime}=Q_{2}{ }^{\prime \prime}+Q_{2}{ }^{\prime \prime \prime}\)\\
So \(\eta=1-\frac{Q_{2}^{\prime}}{Q_{1}}+1=\frac{C_{V}(\tau-1)+R \ln v}{C_{V}(\tau-1)+\tau R \ln v}\)

\[
=1-\frac{\frac{\tau-1}{\gamma-1}+\ln v}{\frac{\tau-1}{\gamma-1}+\tau \ln v}=\frac{(\tau-1) \ln v}{\tau \ln v+\frac{\tau-1}{\gamma-1}}
\]

\includegraphics[max width=\textwidth, alt={}, center]{9f3c5a8e-ac5d-42f4-8eb6-ea2af623f2d8-239_352_452_906_831}\\
2.126 Here \(Q_{1}^{\prime \prime}=C_{p} T_{0}(\tau-1), Q^{\prime} 1^{\prime \prime}=\tau R T_{0} \ln n\) and\\
\(Q_{2}{ }^{\prime \prime}=C^{p} T_{0}(\tau-1), Q_{2}{ }^{\prime \prime \prime}=R T_{0} \ln n\)\\
in addition to we have\\
\(Q_{1}=Q_{1}^{\prime}+Q_{1}^{\prime \prime}\) and\\
\(Q_{2}{ }^{\prime}=Q_{2}{ }^{\prime \prime}+Q_{2}{ }^{\prime \prime \prime}\)\\
So \(\quad \eta=1-\frac{Q_{2}^{\prime}}{Q_{1}}=1-\frac{C_{p}(\tau-1)+R \ln n}{C_{p}(\tau-1)+\tau R \ln n}\)\\
\(=1-\frac{\tau-1+\left(1-\frac{1}{\gamma}\right) \ln n}{\tau-1+\left(1-\frac{1}{\gamma}\right) \tau \ln n}\)\\
\includegraphics[max width=\textwidth, alt={}, center]{9f3c5a8e-ac5d-42f4-8eb6-ea2af623f2d8-239_422_563_1518_801}\\
\(=1-\frac{\tau-1+\left(1-\frac{1}{\gamma}\right) \ln n}{\tau-1+\left(1-\frac{1}{\gamma}\right) \tau \ln n}=\frac{(\tau-1) \ln n}{\tau \ln n+\frac{\gamma(\tau-1)}{\gamma-1}}\)\\
2.127 Because of the linearity of the section\\
\(B C\) whose equation is

\[
\frac{p}{p_{0}}=\frac{v V}{V_{0}}(m p=\alpha V)
\]

We have \(\frac{\tau}{v}=v\) or \(v=\sqrt{\tau}\)\\
Here \(\quad Q_{2}^{\prime \prime}=C_{V} T_{0}(\sqrt{\tau}-1)\),\\
\(Q^{\prime \prime \prime}{ }_{2}=C_{p} T_{0}\left(1-\frac{1}{\sqrt{\tau}}\right)=C_{p} \frac{T_{0}}{\sqrt{\tau}}(\sqrt{\tau}-1)\)\\
\includegraphics[max width=\textwidth, alt={}, center]{9f3c5a8e-ac5d-42f4-8eb6-ea2af623f2d8-240_418_462_135_877}

Thus \(Q^{\prime}{ }_{2}=Q^{\prime \prime}{ }_{2}+Q^{\prime \prime \prime}{ }_{2}=\frac{R T_{0}}{\gamma-1}(\sqrt{\tau}-1)\left(1+\frac{\gamma}{\sqrt{\tau}}\right)\)\\
Along \(B C\), the specific heat \(C\) is given by\\
\(C d T=C_{V} d T+p d V=C_{V} d T+d\left(\frac{1}{2} \alpha V^{2}\right)=\left(C_{V}+\frac{1}{2} R\right) d T\)\\
Thus

\[
Q_{1}=\frac{1}{2} R T_{0} \frac{\gamma+1}{\gamma-1} \frac{\tau-1}{\sqrt{\tau}}
\]

Finally

\[
\eta=1-\frac{Q_{2}^{\prime}}{Q_{1}}=1-2 \frac{\sqrt{\tau}+\gamma}{\sqrt{\tau}+1} \frac{1}{\gamma+1}=\frac{(\gamma-1)(\sqrt{\tau}-1)}{(\gamma+1)(\sqrt{\tau}+1)}
\]

2.128 We write Claussius inequality in the form

\[
\int \frac{\partial_{1} Q}{T}-\int \frac{\partial_{2} Q}{T} \leq 0
\]

where \(\pi Q\) is the heat transeferred to the system but \(\pi_{2} Q\) is heat rejected by the system, both are \(+v e\) and this explains the minus sign before \(t_{2} Q\),\\
In this inequality \(T_{\max }>T>T_{\min }\) and we can write

Thus

\[
\begin{gathered}
\int \frac{d_{1} Q}{T_{\max }}-\int \frac{d_{2} Q}{T_{\min }}<0 \\
\frac{Q_{1}}{T_{\max }}<\frac{Q_{2}^{\prime}}{T_{\min }} \text { or } \frac{T_{\min }}{T_{\max }}<\frac{Q_{2}^{\prime}}{Q_{1}} \\
\eta=1-\frac{Q_{2}^{\prime}}{Q_{1}}<1-\frac{T_{\min }}{T_{\max }}=\eta_{\text {carnot }}
\end{gathered}
\]

or\\
2.129 We consider an infinitesimal carnot cycle with isothermal process at temperatures \(T+d T\) and \(T\).\\
Let \(\delta A\) be the work done in the cycle and \(\delta Q\), be the heat received at the higher temperature. Then by Carnot's theorem

\[
\frac{\delta A}{\delta Q_{1}}=\frac{d T}{T}
\]

On the other hand \(\delta A=d p d V=\left(\frac{\partial p}{\partial T}\right)_{v} d T d V\)\\
while \(\delta Q_{1}=d U_{12}+p d V=\left[\left(\frac{\partial U}{\partial V}\right)_{T}+p\right] d V\)\\
Hence \(\left(\frac{\partial U}{\partial V}\right)_{T}+p=T\left(\frac{\partial p}{\partial T}\right)_{V}\)\\
\includegraphics[max width=\textwidth, alt={}, center]{9f3c5a8e-ac5d-42f4-8eb6-ea2af623f2d8-241_491_557_59_767}\\
2.130 (a) In an isochoric process the entropy change will be

\[
\Delta S=\int_{T_{1}}^{T_{f}} \frac{C_{V} d T}{T}=C_{V} \ln \frac{T_{f}}{T_{i}}=C_{V} \ln n=\frac{R \ln n}{\gamma-1}
\]

For carbon dioxide \(\gamma=1.30\)\\
so,

\[
\Delta S=19 \cdot 2 \mathrm{Joule} /{ }^{\circ} \mathrm{K}-\text { mole }
\]

(b) For an isobaric process,

\[
\begin{aligned}
\Delta S= & C_{p} \ln \frac{T_{f}}{T_{t}}=C_{p} \ln n=\frac{\gamma R \ln n}{\gamma-1} \\
& =25 \text { Joule } /{ }^{\circ} \mathrm{K}-\text { mole }
\end{aligned}
\]

2.131 In an isothermal expansion\\
\(\Delta S=v R \ln \frac{V_{f}}{V_{i}}\)\\
so, \(\frac{V_{f}}{V_{t}}=e^{\Delta S / V R}=2.0\) times

\begin{figure}[h]
\begin{center}
  \includegraphics[alt={},max width=\textwidth]{9f3c5a8e-ac5d-42f4-8eb6-ea2af623f2d8-241_373_329_940_1038}
\captionsetup{labelformat=empty}
\caption{2.132}
\end{center}
\end{figure}

2.132 The entropy change depends on the final \& initial states only, so we can calculate it directly along the isotherm, it is \(\Delta S=2 R \ln n=20 \mathrm{~J} /{ }^{\circ} \mathrm{K}\)\\
(assuming that the final volume is \(n\) times the inital volume)\\
2.133 If the initial temperature is \(T_{0}\) and volume is \(V_{0}\) then in adiabatic expansion.

\[
T V^{\gamma-1}=T_{0} V_{0}^{\gamma-1}
\]

so,

\[
T=T_{0} n^{1-\gamma}=T_{1} \quad \text { where } n=\frac{V_{1}}{V_{0}}
\]

\(V_{1}\) being the volume at the end of the adiabatic process. There is no entropy change in this process. Next the gas is compressed isobarically and the net entropy change is

\[
\Delta S=\left(\frac{m}{M} C_{p}\right) \ln \frac{T_{f}}{T_{1}}
\]

But

\[
\frac{V_{1}}{T_{1}}=\frac{V_{0}}{T_{f}}, \quad \text { or } \quad T_{f}=T_{1} \frac{V_{0}}{V_{1}}=T_{0} n^{-\gamma}
\]

So \(\quad \Delta S=\left(\frac{m}{M} C_{p}\right) \ln \frac{1}{n}=-\frac{m}{M} C_{p} \ln n=-\frac{m}{M} \frac{R \gamma}{\gamma-1} \ln n=-9.7 \mathrm{~J} / \mathrm{K}\)\\
2.134 The entropy change depends on the initial and final state only so can be calculated for any process whatsoever.\\
We choose to evaluate the entropy change along the pair of lines shown above. Then

\[
\begin{aligned}
\Delta S=\int_{T_{0}}^{\frac{T_{0}}{\beta}} \frac{v C_{V} d T}{T}+\int_{\frac{T_{0}}{\beta}}^{\frac{\alpha T_{0}}{\beta}} v C_{p} \frac{d T}{T} \quad P_{0}, V_{0}, \frac{T_{0}}{\beta} \frac{P_{0}}{\beta}, \alpha V_{0}, \frac{\mathcal{L}_{0} T_{0}}{\beta} \\
=\left(-C_{V} \ln \beta+C_{p} \ln \alpha\right) v=\frac{v R}{\gamma-1}(\gamma \ln \alpha-\ln \beta) \approx-11 \frac{\text { Joule }}{{ }^{\circ} \mathrm{K}}
\end{aligned}
\]

2.135 To calculate the required entropy difference we only have to calculate the entropy difference for a process in which the state of the gas in vessel 1 is changed to that in vessel 2 .

\[
\begin{aligned}
& \begin{aligned}
\Delta S & =v\left(\int_{T_{1}}^{\frac{T_{1}}{\alpha \beta}} C_{V} \frac{d T}{T}+\int_{\frac{T_{1}}{\alpha \beta}}^{\frac{T_{1}}{\beta}} C_{p} \frac{d T}{T}\right) \\
& =v\left(C_{p} \ln \alpha-C_{V} \ln \alpha \beta\right)
\end{aligned} \\
& =v\left(R \ln \alpha-\frac{R}{\gamma-1} \ln \beta\right)=v R\left(\ln \alpha-\frac{\ln \beta}{\gamma-1}\right) \\
& \text { With } \gamma=\frac{5}{3}, \alpha=2 \text { and } \beta=1.5, v=1.2, \quad P_{1}, V_{1}, T_{1} \\
& \text { this gives } \Delta S=0.85 \text { Joule } /{ }^{\circ} \mathrm{K}
\end{aligned}\left\{\begin{array}{l}
\frac{P_{1}}{\alpha \beta}, V_{1}, \frac{T_{1}}{\alpha \beta} \quad P_{2}=\frac{P_{1}}{\mathcal{L} \beta}, V_{1} \alpha_{2} \frac{T_{1}}{\beta}
\end{array}\right.
\]

\begin{center}
\includegraphics[max width=\textwidth, alt={}]{9f3c5a8e-ac5d-42f4-8eb6-ea2af623f2d8-242_483_623_288_792}
\end{center}

To calculate the required entropy difference we only have to calculate the entropy difference\\
for a process in which the state of the gas in vessel 1 is changed to that in vessel 2 .\\
2.136 For the polytropic process with index \(n\)

\[
p V^{n}=\text { constant }
\]

Along this process (See 2.122)

\[
\begin{gathered}
C=R\left(\frac{1}{\gamma-1}-\frac{1}{n-1}\right)=\frac{n-\gamma}{(\gamma-1)(n-1)} \cdot R \\
\tau T_{0} \\
\Delta S=\int_{T_{0}} C \frac{d T}{T}=\frac{n-\gamma}{(\gamma-1)(n-1)} R \ln \tau
\end{gathered}
\]

So\\
2.137 The process in question may be written as

\[
\frac{p}{p_{0}}=\alpha \frac{V}{V_{0}}
\]

where \(\alpha\) is a constant and \(p_{0}, V_{0}\) are some reference values. For this process (see 2.127) the specific heat is

\[
C=C_{V}+\frac{1}{2} R=R\left(\frac{1}{\gamma-1}+\frac{1}{2}\right)=\frac{1}{2} R \frac{\gamma+1}{\gamma-1}
\]

Along the line volume increases \(\alpha\) times then so does the pressure. The temperature must then increase \(\alpha^{2}\) times. Thus

\[
\begin{gathered}
\Delta S=\int_{T_{0}}^{\alpha^{2} T_{0}} v C \frac{d T}{T}=\frac{v R}{2} \frac{\gamma+1}{\gamma-1} \ln \alpha^{2}=v R \frac{\gamma+1}{\gamma-1} \ln \alpha \\
v=2, \gamma=\frac{5}{3}, \alpha=2, \Delta S=46 \cdot 1 \mathrm{Joule} /^{\circ} \mathrm{K}
\end{gathered}
\]

if\\
2.138 Let \(\left(p_{1}, V_{1}\right)\) be a reference point on the line

\[
p=p_{0}-\alpha V
\]

and let \((p, V)\) be any other point.\\
The entropy difference

\[
\begin{gathered}
\Delta S=S(p, V)-S\left(p_{1}, V_{1}\right) \\
=C_{V} \ln \frac{p}{p_{1}}+C_{p} \ln \frac{V}{V_{1}}=C_{V} \ln \frac{p_{0}-\alpha V}{p_{1}}+C_{p} \ln \frac{V}{V_{1}}
\end{gathered}
\]

For an exetremum of \(\Delta S\)\\
\(\frac{\partial \Delta S}{\partial V}=\frac{-\alpha C_{V}}{p_{0}-\alpha V}+\frac{C_{p}}{V}=0\)\\
or \(C_{p}\left(p_{0}-\alpha V\right)-\alpha V C_{V}=0\)\\
\includegraphics[max width=\textwidth, alt={}, center]{9f3c5a8e-ac5d-42f4-8eb6-ea2af623f2d8-243_403_411_985_911}\\
or \(\gamma\left(p_{0}-\alpha V\right)-\alpha V=0\) or \(V=V_{m}=\frac{\gamma p_{0}}{\alpha(\gamma+1)}\)\\
This gives a maximum of \(\Delta S\) because \(\frac{\partial^{2} \Delta S}{\partial V^{2}}<0\)\\
(Note :- a maximum of \(\Delta S\) is a maximum of \(S(p, V)\) )\\
2.139 Along the process line : \(S=a T+C_{V} \ln T\)\\
or the specific heat is : \(C=T \frac{d S}{d T}=a T+C_{V}\)\\
On the other hand : \(d Q=C d T=C_{V} d T+p d V\) for an ideal gas.\\
Thus,

\[
p d V=\frac{R T}{V} d V=a T d T
\]

or

\[
\frac{R}{a} \frac{d V}{V}=d T \text { or, } \frac{R}{a} \ln V+\text { constant }=T
\]

Using \(\quad T=T_{0}\) when \(V=V_{0}\), we get, \(T=T_{0}+\frac{R}{a} \ln \frac{V}{V_{0}}\)\\
2.140 For a Vander Waal gas

\[
\left(p+\frac{a}{V^{2}}\right)(V-b)=R T
\]

The entropy change along an isotherm can be calculated from

\[
\Delta S=\int_{V_{1}}^{V_{2}}\left(\frac{\partial S}{\partial V}\right)_{T} d V
\]

It follows from (2.129) that

\[
\left(\frac{\partial S}{\partial V}\right)_{T}=\left(\frac{\partial p}{\partial T}\right)_{V}=\frac{R}{V-b}
\]

assuming \(a, b\) to be known constants.\\
Thus

\[
\Delta S=R \ln \frac{V_{2}-b}{V_{1}-b}
\]

2.141 We use, \(\Delta S=\int_{V_{1}, T_{1}}^{V_{2}, T_{2}} d S(V, T)=\int_{T_{1}}^{T_{2}}\left(\frac{\partial S}{\partial T}\right)_{V_{1}} d T+\int_{V_{1}}^{V_{2}}\left(\frac{\partial S}{\partial V}\right)_{T=T_{2}} d V\)

\[
=\int_{T_{1}}^{T_{2}} \frac{C_{V} d T}{T}+\int_{V_{1}}^{V_{2}} \frac{R}{V-b} d V=C_{V} \ln \frac{T_{2}}{T_{1}}+R \ln \frac{V_{2}-b}{V_{1}-b}
\]

assuming \(C_{V}, a, b\) to be known constants.\\
2.142 We can take \(S \rightarrow 0\) as \(T \rightarrow 0\) Then

\[
S=\int_{0}^{T} C \frac{d T}{T}=\int_{0}^{T} a T^{2} d T=\frac{1}{3} a T^{3}
\]

\(2.143 \Delta S \equiv \int_{T_{1}}^{T_{2}} \frac{C d T}{T}=\int_{T_{1}}^{T_{2}} \frac{m(a+b T)}{T} d T=m b\left(T_{2}-T_{1}\right)+m a \ln \frac{T_{2}}{T_{1}}\)\\
2.144 Here \(T=a S^{n} \quad\) or \(S=\left(\frac{T}{a}\right)^{\frac{1}{n}}\)

Then

\[
\begin{aligned}
& C=T \frac{1}{n} \frac{T^{\frac{1}{n}-1}}{a^{1 / n}}=\frac{S}{n} \\
& C<0 \text { if } n<0 .
\end{aligned}
\]

Clearly\\
\includegraphics[max width=\textwidth, alt={}, center]{9f3c5a8e-ac5d-42f4-8eb6-ea2af623f2d8-245_472_565_356_823}\\
2.145 We know,

\[
S-S_{0}=\int_{T_{0}}^{T} \frac{C d T}{T}=C \ln \frac{T}{T_{0}}
\]

assuming \(C\) to be a known constant.\\
Then \(T=T_{0} \exp \left(\frac{S-S_{0}}{C}\right)\)\\
2.146 (a) \(C=T \frac{d S}{d T}=-\frac{\alpha}{T}\)\\
(b) \(Q=\int_{T_{1}} C d T=\alpha \ln \frac{T_{1}}{T_{2}}\)\\
(c) \(W=\Delta Q-\Delta U=\alpha \ln \frac{T_{1}}{T_{2}}+C_{V}\left(T_{1}-T_{2}\right)\)

Since for an ideal gas \(C_{V}\) is constant and \(\Delta U=C_{V}\left(T_{2}-T_{1}\right)\)\\
( \(U\) does not depend on \(V\) )\\
2.147 (a) We have from the definition

\[
\begin{aligned}
Q & =\int T d S=\text { area under the curve } \\
Q_{1} & =T_{0}\left(S_{1}-S_{0}\right) \\
Q_{2}^{\prime} & =\frac{1}{2}\left(T_{0}+T_{1}\right)\left(S_{1}-S_{0}\right)
\end{aligned}
\]

Thus, using \(T_{1}=\frac{T_{0}}{n}\),

\[
\eta=1-\frac{T_{0}+T_{1}}{2 T_{0}}=1-\frac{1+\frac{1}{n}}{2}=\frac{n-1}{2 n}
\]

(b) Here \(Q_{1}=\frac{1}{2}\left(S_{1}-S_{0}\right)\left(T_{1}+T_{0}\right)\)

\[
\begin{aligned}
& Q_{2}^{\prime}=T_{1}\left(S_{1}-S_{0}\right) \\
& \eta=1-\frac{2 T_{1}}{T_{1}+T_{0}}=\frac{T_{0}-T_{1}}{T_{0}-T_{1}}=\frac{n-1}{n+1}
\end{aligned}
\]

\includegraphics[max width=\textwidth, alt={}, center]{9f3c5a8e-ac5d-42f4-8eb6-ea2af623f2d8-245_514_550_1102_830}\\
\includegraphics[max width=\textwidth, alt={}, center]{9f3c5a8e-ac5d-42f4-8eb6-ea2af623f2d8-245_457_561_1600_832}\\
2.148 In this case, called free expansion no work is done and no heat is exchanged. So internal energy must remain unchanged \(U_{f}=U_{i}\). For an ideal gas this implies constant temperature \(T_{f}=T_{i}\). The process is irreversible but the entropy change can be calculated by considering a reversible isothermal process. Then, as before

\[
\Delta S=\int \frac{d Q}{T}=\int_{V_{1}}^{V_{2}} \frac{p d V}{T}=v R \ln n=20.1 \mathrm{~J} / \mathrm{K}
\]

2.149 The process consists of two parts. The first part is free expansion in which \(U_{f}=U_{i}\). The second part is adiabatic compression in which work done results in change of internal energy. Obviously,

\[
0=U_{F}-U_{f}+\int_{V_{f}}^{V_{0}} p d V, V_{f}=2 V_{0}
\]

Now in the first part \(p_{f}=\frac{1}{2} p_{0}, V_{f}=2 V_{\theta}\), because there is no change of temperature.\\
In the second part, \(p V^{\gamma}=\frac{1}{2} p_{0}\left(2 V_{0}\right)^{\gamma}=2^{\gamma-1} p_{0} V_{0}^{\gamma}\)

\[
\begin{gathered}
\int_{2 V_{0}}^{V_{0}} p d V=\int_{2 V_{0}}^{V_{0}} \frac{2^{\gamma-1} p_{0} V_{0}^{\gamma}}{V^{\gamma}} d V=\left[\frac{2^{\gamma-1} p_{0} V_{0}^{\gamma}}{-\gamma+1} V^{1-\gamma}\right]_{2 V_{0}}^{V_{0}} \\
=2^{\gamma-1} p_{0} V_{0}^{\gamma} V_{0}^{-\gamma+1} \frac{2^{-\gamma+1}-1}{\gamma-1}=-\frac{\left(2^{\gamma-1}-1\right.}{\gamma-1} R T \\
\Delta U=U_{F}-U_{i}=\frac{R T_{0}}{\gamma-1}\left(2^{\gamma-1}-1\right)
\end{gathered}
\]

Thus\\
The entropy change \(\Delta S=\Delta S_{I}+\Delta S_{I I}\)\\
\(\Delta S_{I}=R \ln 2\) and \(\Delta S_{I I}=0\) as the process is reversible adiabatic. Thus \(\Delta S=R \ln 2\).\\
2.150 In all adiabatic processes

\[
Q=U_{f}-U_{i}+A=0
\]

by virtue of the first law of thermodynamics. Thus,

\[
U_{f}=U_{i}-A
\]

For a slow process, \(A^{\prime}=\int_{V_{0}}^{V} p d V\) where for a quasistatic adiabatic process \(p V^{\gamma}=\) constant.\\
On the other hand for a fast process the external work done is \(A^{\prime \prime}<A^{\prime}\). In fact \(A^{\prime \prime}=0\) for free expansion. Thus \(U_{f}^{\prime}\) (slow) \(<U^{\prime \prime}{ }_{f}\) (fast)\\
Since \(U\) depends on temperature only, \(T_{f}<T^{\prime}{ }_{f}\)\\
Consequently, \(\boldsymbol{p}^{\prime \prime}{ }_{\boldsymbol{f}}>\boldsymbol{p}_{\boldsymbol{f}}\)\\
(From the ideal gas equation \(p V=R T\) )\\
2.151 Let \(V_{1}=V_{0}, V_{2}=n V_{0}\)

Since the temperature is the same, the required entropy change can be calculated by considering isothermal expansion of the gas in either parts into the whole vessel.

Thus

\[
\begin{aligned}
\Delta S & =\Delta S_{I}+\Delta S_{I I}=v_{1} R \ln \frac{V_{1}+V_{2}}{V_{1}}+v_{2} R \ln \frac{V_{1}+V_{2}}{V_{2}} \\
& =v_{1} R \ln (1+n)+v_{2} R \ln \frac{1+n}{n}=5 \cdot 1 \mathrm{~J} / \mathrm{K}
\end{aligned}
\]

2.152 Let \(c_{1}=\) specific heat of copper specific heat of water \(=c_{2}\)

\[
T_{0} \quad 97+273
\]

Then \(\Delta S=\int_{7+273} \frac{c_{2} m_{2} d T}{T}-\int_{T_{0}} \frac{m_{1} c_{1} d T}{T}=m_{2} c_{2} \ln \frac{T_{0}}{280}-m_{1} c_{1} \ln \frac{370}{T_{0}}\)\\
\(T_{0}\) is found from

\[
\begin{aligned}
& c_{2} m_{2}\left(T_{0}-280\right)=m_{1} c_{1}\left(370-T_{0}\right) \text { or } T_{0}=\frac{280 m_{2} c_{2}+370 m_{1} c_{1}}{c_{2} m_{2}+m_{1} c_{1}} \\
& \text { using } \\
& \qquad \begin{aligned}
& c_{1}=0.39 \mathrm{~J} / \mathrm{g}{ }^{\circ} \mathrm{K}, c_{2}=4.18 \mathrm{~J} / \mathrm{g}{ }^{\circ} \mathrm{K}, \\
& T_{0} \approx 300^{\circ} \mathrm{K} \text { and } \Delta S=28.4-24.5 \approx 3.9 \mathrm{~J} /{ }^{\circ} \mathrm{K}
\end{aligned}
\end{aligned}
\]

2.153 For an ideal gas the internal energy depends on temperature only. We can consider the process in question to be one of simultaneous free expansion. Then the total energy \(U=U_{1}+U_{2}\). Since\\
\(U_{1}=C_{V} T_{1}, U_{2}=C_{V} T_{2}, U=2 C_{V} \frac{T_{1}+T_{2}}{2}\) and \(\left(T_{1}+T_{2}\right) / 2\) is the final temperature. The entropy change is obtained by considering isochoric processes because in effect, the gas remains confined to its vessel.

\[
\Delta S=\int_{T_{1}}^{\left(T_{1}+T_{2}\right) / 2} \frac{C_{V} d T}{T}-\int_{\left(T_{1}+T_{2}\right) / 2}^{T_{2}} C_{V} \frac{d T}{T}=C_{V} \ln \frac{\left(T_{1}+T_{2}\right)^{2}}{4 T_{1} T_{2}}
\]

Since

\[
\left(T_{1}+T_{2}\right)^{2}=\left(T_{1}-T_{2}\right)^{2}+4 T_{1} T_{2}, \Delta S>0
\]

2.154 (a) Each atom has a probability \(\frac{1}{2}\) to be in either campartment. Thus

\[
p=2^{-N}
\]

(b) Typical atomic velocity at room temperature is \(\sim 10^{5} \mathrm{~cm} / \mathrm{s}\) so it takes an atom \(10^{-5} \mathrm{sec}\) to cross the vessel. This is the relevant time scale for our problem. Let \(T=10^{-5} \mathrm{sec}\), then in time \(t\) there will be \(t / T\) crossing or arrangements of the atoms. This will be large enough to produce the given arrangement if

\[
\frac{t}{\tau} 2^{-N} \sim 1 \text { or } N \sim \frac{\ln t / \tau}{\ln 2} \sim 75
\]

2.155 The statistical weight is

\[
N_{C_{N / 2}}=\frac{N!}{N / 2!\frac{N}{2}!}=\frac{10 \times 9 \times 8 \times 7 \times 6}{8 \times 4 \times 3 \times 2}=252
\]

The probability distribution is

\[
N_{C_{N / 2}} 2^{-N}=252 \times 2^{-10}=24.6 \%
\]

2.156 The probabilites that the half \(A\) contains \(n\) molecules is

\[
N_{C_{n}} \times 2^{-N}=\frac{N!}{n!(N-n)!} 2^{-N}
\]

2.157 The probability of one molecule being confined to the marked volume is

\[
p=\frac{V}{V_{0}}
\]

We can choose this molecule in many \(\left(N_{C_{1}}\right)\) ways. The probability that \(n\) molecules get confined to the marked volume is cearly

\[
N_{C_{n}} p^{n}(1-p)^{N-n}=\frac{N!}{n!(N-n)!} p^{n}(1-p)^{N-n}
\]

2.158 In a sphere of diameter \(\boldsymbol{d}\) there are

\[
N=\frac{\pi d^{3}}{6} n_{0} \quad \text { molecules }
\]

where \(n_{0}=\) Loschmidt's number \(=\) No. of molecules per unit volume \((1 \mathrm{cc})\) under NTP. The relative fluctuation in this number is

\[
\frac{\partial N}{N}=\frac{\sqrt{N}}{N}=\frac{1}{\sqrt{N}}=\eta
\]

or \(\frac{1}{\eta^{2}}=\frac{\pi}{6} d^{3} n_{0}\) or \(d^{3}=\frac{6}{\pi n_{0} \eta^{2}}\) or \(d=\left(\frac{6}{\pi \eta^{2} n_{0}}\right)^{1 / 3}=0.41 \mu \mathrm{~m}\)\\
The average number of molecules in this sphere is \(\frac{1}{\eta^{2}}=10^{6}\)\\
2.159 For a monoatomic gas \(C_{V}=\frac{3}{2} R\) per mole

The entropy change in the process is

\[
\Delta S=S-S_{0}=\int_{T_{0}}^{T_{0}+\Delta T} C_{V} \frac{d T}{T}=\frac{3}{2} R \ln \left(1+\frac{\Delta T}{T_{0}}\right)
\]

Now from the Boltzmann equation

\[
\begin{gathered}
S=k \ln \Omega \\
\frac{\Omega}{\Omega_{0}}=e^{\left(S-S_{0}\right) / k}=\left(1+\frac{\Delta T}{T_{0}}\right)^{\frac{3 N_{A}}{2}}=\left(1+\frac{1}{300}\right)^{\frac{3 \times 6}{2} \times 10^{23}}=10^{13} \times 10^{21}
\end{gathered}
\]

Thus the statistical weight increases by this factor.

\subsection*{2.5 LIQUIDS. CAPILLARY EFFECTS}
2.160

\[
\text { (a) } \begin{aligned}
\Delta p & =\alpha\left(\frac{1}{d / 2}+\frac{1}{d / 2}\right)=\frac{4 \alpha}{d} \\
& =\frac{4 \times 490 \times 10^{-3}}{1.5 \times 10^{-6}} \frac{\mathrm{~N}}{\mathrm{~m}^{2}}=1.307 \times 10^{6} \frac{\mathrm{~N}}{\mathrm{~m}^{2}}=13 \text { atmosphere }
\end{aligned}
\]

(b) The soap bubble has two surfaces\\
so

\[
\begin{gathered}
\Delta p=2 \alpha\left(\frac{1}{d / 2}+\frac{1}{d / 2}\right)=\frac{8 \alpha}{d} \\
=\frac{8 \times 45}{3 \times 10^{-3}} \times 10^{-3}=1.2 \times 10^{-3} \text { atomsphere } .
\end{gathered}
\]

2.161 The pressure just inside the hole will be less than the outside pressure by \(4 \alpha / d\). This can support a height \(h\) of Hg where

\[
\begin{gathered}
\rho g h=\frac{4 \alpha}{d} \text { or } h=\frac{4 \alpha}{\rho g d} \\
=\frac{4 \times 490 \times 10^{-3}}{13.6 \times 10^{3} \times 9.8 \times 70 \times 10^{-6}}=\frac{200}{13.6 \times 70} \approx 21 \mathrm{~m} \text { of } \mathrm{Hg}
\end{gathered}
\]

2.162 By Boyle's law\\
or

\[
\left(p_{0}+\frac{8 \alpha}{d}\right) \frac{4 \pi}{3}\left(\frac{d}{2}\right)^{3}=\left(\frac{p_{0}}{n}+\frac{8 \alpha}{\eta d}\right) \frac{4 \pi}{3}\left(\frac{\eta d}{2}\right)^{3}
\]

Thus

\[
\begin{aligned}
& p_{0}\left(1-\frac{\eta^{3}}{n}\right)=\frac{8 \alpha}{d}\left(\eta^{2}-1\right) \\
& \alpha=\frac{1}{8} p_{0} d\left(1-\frac{\eta^{3}}{n}\right)\left(\eta^{2}-1\right)
\end{aligned}
\]

2.163 The pressure has terms due to hydrostatic pressure and capillarity and they add

\[
\begin{gathered}
p=p_{0}+\rho g h+\frac{4 \alpha}{d} \\
=\left(1+\frac{5 \times 9.8 \times 10^{3}}{10^{5}}+\frac{4 \times .73 \times 10^{-3}}{4 \times 10^{-6}} \times 10^{-5}\right) \text { atoms }=2.22 \text { atom } .
\end{gathered}
\]

2.164 By Boyle's law

\[
\left(p_{0}+h g \rho+\frac{4 \alpha}{d}\right) \frac{\pi}{6} d^{3}=\left(p_{0}+\frac{4 \alpha}{n d}\right) \frac{\pi}{6} n^{3} d^{3}
\]

or

\[
\left[h g \rho-p_{0}\left(n^{3}-1\right)\right]=\frac{4 \alpha}{d}\left(n^{2}-1\right)
\]

or

\[
h=\left[p_{0}\left(n^{3}-1\right)+\frac{4 \alpha}{d}\left(n^{2}-1\right)\right] / g \rho=4.98 \text { meter of water }
\]

2.165 Clearly

\[
\begin{gathered}
\Delta h \rho g=4 \alpha|\cos \theta|\left(\frac{1}{d_{1}}-\frac{1}{d_{2}}\right) \\
\Delta h=\frac{4 \alpha|\cos \theta|\left(d_{2}-d_{1}\right)}{d_{1} d_{2} \rho g}=11 \mathrm{~mm}
\end{gathered}
\]

\includegraphics[max width=\textwidth, alt={}, center]{9f3c5a8e-ac5d-42f4-8eb6-ea2af623f2d8-250_310_472_79_887}\\
2.166 In a capillary with diameter \(d=0.5 \mathrm{~mm}\) water will rise to a height

\[
\begin{gathered}
\frac{2 \alpha}{\rho g r}=\frac{4 \alpha}{\rho g d} \\
=\frac{4 \times 73 \times 10^{-3}}{10^{3} \times 9.8 \times 0.5 \times 10^{-3}}=59.6 \mathrm{~mm}
\end{gathered}
\]

Since this is greater than the height ( \(=25 \mathrm{~mm}\) ) of the tube, a meniscus of radius \(R\) will be formed at the top of the tube, where\\
\(R=\frac{2 \alpha}{\rho g h}=\frac{2 \times 73 \times 10^{-3}}{10^{3} \times 9.8 \times 25 \times 10^{-3}} \approx 0.6 \mathrm{~mm}\)\\
\includegraphics[max width=\textwidth, alt={}, center]{9f3c5a8e-ac5d-42f4-8eb6-ea2af623f2d8-250_412_440_487_913}\\
2.167 Initially the pressure of air in the eppillary is \(p_{0}\) and it's length is \(l\). When submerged under water, the pressure of air in the portion above water must be \(p_{0}+4 \frac{\alpha}{d}\), since the level of water inside the capillary is the same as the level outside. Thus by Boyle's law\\
or

\[
\begin{gathered}
\left(p_{0}+\frac{4 \alpha}{d}\right)(l-x)=p_{0} l \\
\frac{4 \alpha}{d}(l-x)=p_{0} x \quad \text { or } \quad x=\frac{l}{1+\frac{p_{0} d}{4 \alpha}}
\end{gathered}
\]

2.168 We have by Boyle's law\\
or,

\[
\left(p_{0}-\rho g h+\frac{4 \alpha \cos \theta}{d}\right)(l-h)=p_{0} l
\]

Hence,

\[
\begin{aligned}
& \frac{4 \alpha \cos \theta}{d}=\rho g h+\frac{p_{0} h}{l-h} \\
& \alpha=\left(\rho g h+\frac{p_{0} h}{l-h}\right) \frac{d}{4 \cos \theta}
\end{aligned}
\]

2.169 Suppose the liquid rises to a height \(h\). Then the total energy of the liquid in the capillary is

\[
E(h)=\frac{\pi}{4}\left(d_{2}^{2}-d_{1}^{2}\right) h \times \rho g \times \frac{h}{2}-\pi\left(d_{2}-d_{1}\right) \alpha h
\]

Minimising \(\boldsymbol{E}\) we get

\[
h=\frac{4 \alpha}{\rho g\left(d_{2}-d_{1}\right)}=6 \mathrm{~cm}
\]

2.170 Let \(h\) be the height of the water level at a distance \(x\) from the edge. Then the total energy of water in the wedge above the level outside is.

\[
\begin{aligned}
& E=\int x \delta \varphi \cdot d x \cdot h \cdot \rho g \frac{h}{2}-2 \int d x \cdot h \cdot \alpha \cos 0 \\
& =\int d x \frac{1}{2} x \rho g \delta \varphi\left(h^{2}-2 \frac{2 \alpha \cos \theta}{x \rho g \delta \varphi} h\right) \\
& =\int d x \frac{1}{2} x \rho g \delta \varphi\left[\left(h-\frac{2 \alpha \cos \theta}{x \rho g \delta \varphi}\right)^{2}-\frac{4 \alpha^{2} \cos ^{2} \theta}{x^{2} \rho^{2} g^{2} \delta \varphi^{2}}\right]
\end{aligned}
\]

\begin{center}
\includegraphics[max width=\textwidth, alt={}]{9f3c5a8e-ac5d-42f4-8eb6-ea2af623f2d8-251_572_469_60_916}
\end{center}

This is minimum when \(h=\frac{2 \alpha \cos \theta}{x \rho g \delta \varphi}\)

\subsection*{2.171 From the equation of continuity}
\[
\frac{\pi}{4} d^{2} \cdot v=\frac{\pi}{4}\left(\frac{d}{n}\right)^{2} \cdot V \quad \text { or } \quad V=n^{2} v
\]

We then apply Bernoulli's theorem

\[
\frac{p}{\rho}+\frac{1}{2} v^{2}+\Phi=\text { constant }
\]

The pressure \(p\) differs from the atmospheric pressure by capillary effects. At the upper section

\[
p=p_{0}+\frac{2 \alpha}{d}
\]

neglecting the curvature in the vertical plane. Thus,\\
or

\[
\begin{gathered}
\frac{p_{0}+\frac{2 \alpha}{d}}{\rho}+\frac{1}{2} v^{2}+g l=\frac{p_{0}+\frac{2 n \alpha}{d}}{\rho}+\frac{1}{2} n^{4} v^{2} \\
v=\sqrt{\frac{2 g l-\frac{4 \alpha}{\rho d}(n-1)}{n^{4}-1}}
\end{gathered}
\]

Finally, the liquid coming out per second is,

\[
V=\frac{1}{4} \pi d^{2} \sqrt{\frac{2 g l-\frac{4 \alpha}{\rho d}(n-1)}{n^{4}-1}}
\]

2.172 The radius of curvature of the drop is \(R_{1}\) at the upper end of the drop and \(R_{2}\) at the lower end. Then the pressure inside the drop is \(p_{0}+\frac{2 \alpha}{R_{1}}\) at the top end and \(p_{0}+\frac{2 \alpha}{R_{2}}\) at the bottom end. Hence

\[
p_{0}+\frac{2 \alpha}{R_{1}}=p_{0}+\frac{2 \alpha}{R_{2}}+\rho g h \quad \text { or } \quad \frac{2 \alpha\left(R_{2}-R_{1}\right)}{R_{1} R_{2}}=\rho g h
\]

To a first approximation \(R_{1} \approx R_{2} \approx \frac{h}{2}\) so \(R_{2}-R_{1} \approx \frac{1}{8} \rho g h^{3} / \alpha . \approx 0.20 \mathrm{~mm}\)\\
if

\[
h=2.3 \mathrm{~mm}, \alpha=73 \mathrm{mN} / \mathrm{m}
\]

2.173 We must first calculate the pressure difference inside the film from that outside. This is

\[
p=\alpha\left(\frac{1}{r_{1}}+\frac{1}{r_{2}}\right)
\]

Here \(2 r_{1}|\cos \theta|=h\) and \(r_{2} \sim-R\) the radius of the tablet and can be neglected. Thus the total force exerted by mercury drop on the upper glass plate is

\[
\frac{2 \pi R^{2} \alpha|\cos \theta|}{h} \text { typically }
\]

\begin{center}
\includegraphics[max width=\textwidth, alt={}]{9f3c5a8e-ac5d-42f4-8eb6-ea2af623f2d8-252_358_455_118_957}
\end{center}

We should put \(h / n\) for \(h\) because the tablet is compresed \(n\) times. Then since \(H g\) is nearly, incompressible, \(\pi R^{2} h=\) constants so \(R \rightarrow R \sqrt{n}\). Thus,

\[
\text { total force }=\frac{2 \pi R^{2} \alpha|\cos \theta|}{h} n^{2}
\]

Part of the force is needed to keep the \(H g\) in the shape of a table rather than in the shape of infinitely thin sheet. This part can be calculated being putting \(n=1\) above. Thus\\
or

\[
\begin{gathered}
m g+\frac{2 \pi R^{2} \alpha|\cos \theta|}{h}=\frac{2 \pi R^{2} \alpha|\cos \theta|}{h} n^{2} \\
m=\frac{2 \pi R^{2} \alpha|\cos \theta|}{h g}\left(n^{2}-1\right)=0.7 \mathrm{~kg}
\end{gathered}
\]

2.174 The pressure inside the film is less than that outside by an amount \(\alpha\left(\frac{1}{r_{1}}+\frac{1}{r_{2}}\right)\) where \(r_{1}\) and \(r_{2}\) are the principal radii of curvature of the meniscus. One of these is small being given by \(h=2 r_{1} \cos \theta\) while the other is large and will be ignored. Then \(F \approx \frac{2 A \cos \theta}{h} \alpha\) where \(A=\) area of the water film between the plates.\\
Now \(A=\frac{m}{\rho h}\) so \(F=\frac{2 m \alpha}{\rho h^{2}}\) when \(\theta\) (the angle of contact) \(=0\)\\
2.175 This is analogous to the previous problem except that: \(A=\pi R^{2}\)

So

\[
F=\frac{2 \pi R^{2} \alpha}{h}=0.6 \mathrm{kN}
\]

2.176 The energy of the liquid between the plates is\\
\(E=l d h \rho g \frac{h}{2}-2 \alpha l h=\frac{1}{2} \rho g l d h^{2}-2 \alpha l h\)\\
\(=\frac{1}{2} \rho g l d\left(h-\frac{2 \alpha}{\rho g d}\right)^{2}-\frac{2 \alpha^{2} l}{\rho g d}\)\\
This energy is minimum when, \(h=\frac{2 \alpha}{\rho g d}\) and\\
\includegraphics[max width=\textwidth, alt={}, center]{9f3c5a8e-ac5d-42f4-8eb6-ea2af623f2d8-252_374_555_1330_844}\\
the minimum potential energy is then \(E_{\min }=-\frac{2 \alpha^{2} l}{\rho g d}\)\\
The force of attraction between the plates can be obtained from this as\\
\(F=\frac{-\partial E_{\text {min }}}{\partial d}=-\frac{2 \alpha^{2} l}{\rho g d^{2}}\) (minus sign means the force is attractive.)\\
Thus

\[
F=-\frac{\alpha l h}{d}=13 \mathrm{~N}
\]

2.177 Suppose the radius of the bubble is \(x\) at some instant. Then the pressure inside is \(p_{0}+\frac{4 \alpha}{x}\). The flow through the capillary is by Poiseuille's equation,

\[
Q=\frac{\pi r^{4}}{8 \eta l} \frac{4 \alpha}{x}=-4 \pi^{2} \frac{d x}{d t}
\]

Integrating \(\frac{\pi r^{4} \alpha}{2 \eta l} t=\pi\left(R^{4}-x^{4}\right)\) where we have used the fact that \(t=0\) where \(x=R\).\\
This gives \(t=\frac{2 \eta l R^{4}}{\alpha r^{4}}\) as the life time of the bubble corresponding to \(x=0\)\\
2.178 If the liquid rises to a height \(h^{\prime}\), the energy of the liquid column becomes

\[
E=\rho g \pi r^{2} h \cdot \frac{h}{2}-2 \pi r h \alpha=\frac{1}{2} \rho g \pi\left(r h-2 \frac{\alpha}{\rho g}\right)^{2}-\frac{2 \pi \alpha^{2}}{\rho g}
\]

This is minimum when \(r h=\frac{2 \alpha}{\rho g}\) and that is relevant height to which water must rise.\\
At this point, \(\quad E_{\text {min }}=-\frac{2 \pi \alpha^{2}}{\rho g}\)\\
Since \(E=0\) in the absence of surface tension a heat \(Q=\frac{2 \pi \alpha^{2}}{\rho g}\) must have been liberated.\\
2.179 (a) The free energy per unit area being \(\alpha\),

\[
F=\pi \alpha d^{2}=3 \mu \mathrm{~J}
\]

(b) \(F=2 \pi \alpha d^{2}\) because the soap bubble has two surfaces. Substitution gives \(F=10 \mu \mathrm{~J}\)\\
2.180 When two mercury drops each of diameter \(d\) merge, the resulting drop has diameter \(d_{1}\) where

\[
\frac{\pi}{6} d_{1}^{3}=\frac{\pi}{6} d^{3} \times 2 \quad \text { or, } \quad d_{1}=2^{1 / 3} d
\]

The increase in free energy is

\[
\Delta F=\pi 2^{2 / 3} d^{2} \alpha-2 \pi d^{2} \alpha=2 \pi d^{2} \alpha\left(2^{-1 / 3}-1\right)=-1.43 \mu \mathrm{~J}
\]

2.181 Work must be done to stretch the soap film and compress the air inside. The former is simply \(2 \alpha \times 4 \pi R^{2}=8 \pi R^{2} \alpha\), there being two sides of the film. To get the latter we note that the compression is isothermal and work done is

\[
\begin{gathered}
-\int_{V_{1}=V_{0}}^{V_{f}=V} p d V \text { where } V_{0} p_{0}=\left(p_{0}+\frac{4 \alpha}{R}\right) \cdot V, V=\frac{4 \pi}{3} R^{3} \\
V_{0}=\frac{p V}{p_{0}}, p=p_{0}+\frac{4 \alpha}{R}
\end{gathered}
\]

or\\
and minus sign is needed becaue we are calculating work done on the system. Thus since \(p V\) remains constants, the work done is

\[
p V \ln \frac{V_{0}}{V}=p V \ln \frac{p}{p_{0}}
\]

So

\[
A^{\prime}=8 \pi R^{2} \alpha+p V \ln \frac{p}{p_{0}}
\]

2.182 When heat is given to a soap bubble the temperature of the air inside rises and the bubble expands but unless the bubble bursts, the amount of air inside does not change. Further we shall neglect the variation of the surface tension with temperature. Then from the gas equations

Differentiating

\[
\left(p_{0}+\frac{4 \alpha}{r}\right) \frac{4 \pi}{3} r^{3}=v R T, v=\text { Constant }
\]

or

\[
\begin{array}{r}
\left(p_{0}+\frac{8 \alpha}{3 r}\right) 4 \pi r^{2} d r=v R d T \\
d V=4 \pi r^{2} d r=\frac{v R d T}{p_{0}+\frac{8 \alpha}{3 r}}
\end{array}
\]

Now from the first law\\
or

\[
\begin{gathered}
d Q=v C d T=v C_{V} d T+\frac{v R d T}{p_{0}+\frac{8 \alpha}{3 r}} \cdot\left(p_{0}+\frac{4 \alpha}{r}\right) \\
C=C_{V}+R \frac{p_{0}+\frac{4 \alpha}{r}}{p_{0} \frac{8 \alpha}{3 r}}
\end{gathered}
\]

using

\[
C_{p}=C_{V}+R, C=C_{p}+\frac{\frac{1}{2} R}{1+\frac{3 p_{0} r}{8 \alpha}}
\]

2.183 Consider an infinitesimal Carnot cycle with isotherms at \(T-d T\) and \(T\). Let \(A\) be the work done during the cycle. Then\\
\(A=[\alpha(T-d T)-\alpha(T)] \delta \sigma=-\frac{d \alpha}{d T} d T \delta \sigma\)\\
Where \(\delta \sigma\) is the change in the area of film (we are considering only one surface).\\
Then \(\eta=\frac{A}{Q_{1}}=\frac{d T}{T}\) by Carnot therom.\\
or \(\frac{-\frac{d \alpha}{d T} d T \delta \sigma}{q \delta \sigma}=\frac{d T}{T}\) or \(q=-T \frac{d \alpha}{d T}\)\\
\includegraphics[max width=\textwidth, alt={}, center]{9f3c5a8e-ac5d-42f4-8eb6-ea2af623f2d8-254_478_569_1153_815}\\
2.184 As before we can calculate the heat required. It, is taking into account two sides of the soap film

Thus

\[
\begin{aligned}
& \delta q=-T \frac{d \alpha}{d T} \delta \sigma \times 2 \\
& \Delta S=\frac{\delta q}{T}=-2 \frac{d \alpha}{d T} \delta \sigma
\end{aligned}
\]

Now

\[
\Delta F=2 \alpha \delta \sigma \text { so, } \Delta U=\Delta F+T \Delta S=2\left(\alpha-T \frac{d \alpha}{d T}\right) \delta \sigma
\]

\subsection*{2.6 PHASE TRANSFORMATIONS}
2.185 The condensation takes place at constant pressure and temperature and the work done is

\[
p \Delta V
\]

where \(\Delta V\) is the volume of the condensed vapour in the vapour phase. It is

\[
p \Delta V=\frac{\Delta m}{M} R T=120.6 \mathrm{~J}
\]

where \(M=18 \mathrm{gm}\) is the molecular weight of water.\\
2.186 The specific volume of water (the liquid) will be written as \(V_{l}^{\prime}\). Since \(V_{\nu}^{\prime} \gg V_{l}^{\prime}\), most of the weight is due to water. Thus if \(m_{l}\) is mass of the liquid and \(m_{\nu}\) that of the vapour then

\[
\begin{gathered}
m=m_{l}+m_{v} \\
V=m_{l} V_{l}^{\prime}+m_{v} V_{v}^{\prime} \text { or } V-m V_{l}^{\prime}=m_{v}\left(V_{v}^{\prime}-V_{l}^{\prime}\right)
\end{gathered}
\]

So \(m_{v}=\frac{V-m V_{l}^{\prime}}{V_{v}^{\prime}-V_{1}^{\prime}}=20 \mathrm{gm}\) in the present case. Its volume is \(m_{v} V_{v}^{\prime}=1.01\)\\
2.187 The volume of the condensed vapour was originally \(V_{0}-V\) at temperature \(T=373 \mathrm{~K}\). Its mass will be given by\\
\(p\left(V_{0}-V\right)=\frac{m}{M} R T\) or \(m=\frac{M p\left(V_{0}-V\right)}{R T}=2 \mathrm{gm}\) where \(p=\) atmospheric pressure\\
2.188 We let \(V_{l}^{\prime}=\) specific volume of liquid. \(V_{\nu}^{\prime}=N V_{l}^{\prime}=\) specific volume of vapour.

Let \(V=\) Original volume of the vapour. Then

\[
M \frac{p V}{R T}=m_{l}+m_{v}=\frac{V}{N V_{l}^{\prime}} \text { or } \frac{V}{n}=\left(m_{l}+N m_{v}\right) V_{l}^{\prime}
\]

So

\[
(N-1) m_{l} V_{l}^{\prime}=V\left(1-\frac{1}{n}\right)=\frac{V}{n}(n-1) \text { or } \eta=\frac{m_{l} V_{l}^{\prime}}{V / n}=\frac{n-1}{N-1}
\]

In the case when the final volume of the substance corresponds to the midpoint of a horizontal portion of the isothermal line in the \(p, v\) diagram, the final volume must be \((1+N) \frac{V_{l}^{\prime}}{2}\) per unit mass of the substance. Of this the volume of the liquid is \(V_{1}^{\prime} / 2\) per unit total mass of the substance.

Thus

\[
\eta=\frac{1}{1+N}
\]

2.189 From the first law of thermodynamics

\[
\Delta U+A=Q=m q
\]

where \(q\) is the specific latent heat of vaporization\\
Now

\[
A=p\left(V_{\nu}^{\prime}-V_{l}^{\prime}\right) m=m \frac{R T}{M}
\]

Thus

\[
\Delta U=m\left(q-\frac{R T}{M}\right)
\]

For water this gives \(\approx 2.08 \times 10^{6}\) Joules.\\
2.190 Some of the heat used in heating water to the boiling temperature

\[
\begin{aligned}
T=100^{\circ} \mathrm{C} & =373 \mathrm{~K} . \text { The remaining heat } \\
& =Q-m c \Delta T
\end{aligned}
\]

( \(c=\) specific heat of water, \(\Delta T=100 \mathrm{~K}\) ) is used to create vapour. If the piston rises to a height \(h\) then the volume of vapour will be \(\approx \operatorname{sh}\) (neglecting water). Its mass will be \(\frac{p_{0} s h}{R T} \times M\) and heat of vapourization will be \(\frac{p_{0} s h M q}{R T}\). To this must be added the work done in creating the saturated vapour \(\approx p_{0} s h\). Thus\\
\(Q-m c \Delta T \sim p_{0} S h\left(1+\frac{q M}{R T}\right)\) or \(h=\frac{Q-m c \Delta T}{p_{0} s\left(1+\frac{q M}{R T}\right)}=20 \mathrm{~cm}\)\\
2.191 A quantity \(\frac{m c\left(T-T_{0}\right)}{q}\) of saturated vapour must condense to heat the water to boiling point \(T=373^{\circ} \mathrm{K}\)\\
(Here \(c=\) specific heat of water, \(T_{0}=295 \mathrm{~K}=\) initial water temperature).\\
The work done in lowering the piston will then be

\[
\frac{m c\left(T-T_{0}\right)}{q} \times \frac{R T}{M}=25 \mathrm{~J},
\]

since work done per unit mass of the condensed vapour is \(p V=\frac{R T}{M}\)\\
2.192 Given \(\Delta P=\frac{\rho_{v}}{\rho_{l}} \frac{2 \alpha}{r}=\frac{\rho_{v}}{\rho_{l}} \times \frac{4 \alpha}{d}=\eta p_{\text {vap }}=\eta \frac{\frac{m}{M} R T}{V_{\text {vap }}}=\frac{\eta R T}{M} \rho_{v}\)\\
or

\[
d=\frac{4 \alpha M}{\rho_{1} R T \eta}
\]

For water \(\alpha=73\) dynes \(/ \mathrm{cm}, M=18 \mathrm{gm}, \rho_{l}=\mathrm{gm} / \mathrm{cc}, T=300 \mathrm{~K}\), and with \(\eta \approx 0.01\), we get

\[
d \approx 0.2 \mu \mathrm{~m}
\]

2.193 In equilibrium the number of "liquid" molecules evaporoting must equals the number of "vapour" molecules condensing. By kinetic theory, this number is

\[
\eta \times \frac{1}{4} n<v>=\eta \times \frac{1}{4} n \times \sqrt{\frac{8}{\pi} \frac{k T}{m}}
\]

Its mass is

\[
\begin{gathered}
\mu=m \times \eta \times n \times \sqrt{\frac{k T}{2 \pi m}}=\eta n k T \sqrt{\frac{m}{2 \pi k t}} \\
=\eta p_{0} \sqrt{\frac{M}{2 \pi R T}}=0.35 \mathrm{~g} / \mathrm{cm}^{2} \cdot \mathrm{~s} .
\end{gathered}
\]

where \(p_{0}\) is atmospheric pressure and \(T=373 \mathrm{~K}\) and \(M=\) molecular weight of water.\\
2.194 Here we must assume that \(\mu\) is also the rate at which the tungsten filament loses mass when in an atmosphere of its own vapour at this temperature and that \(\eta\) (of the previous problem) \(\approx 1\). Then

\[
p=\mu \sqrt{\frac{2 \pi R T}{M}}=0.9 \mathrm{n} \mathrm{~Pa}
\]

from the previous problem where \(p=\) pressure of the saturated vapour.\\
2.195 From the Vander Waals equation

\[
p=\frac{R T}{V-b}-\frac{a}{V^{2}}
\]

where \(V=\) Volume of one gm mole of the substances.\\
For water \(V=18 \mathrm{c} . \mathrm{c}\). per mole \(=1.8 \times 10^{-2}\) litre per mole

\[
a=5 \cdot 47 \mathrm{atmos} \cdot \frac{\mathrm{litre}^{2}}{\mathrm{~mole}^{2}}
\]

If molecular attraction vanished the equation will be

\[
p^{\prime}=\frac{R T}{V-b}
\]

for the same specific volume. Thus\\
\(\Delta p=\frac{a}{V^{2}}=\frac{5.47}{1.8 \times 1.8} \times 10^{4} \mathrm{atmos} \approx 1.7 \times 10^{4} \mathrm{atmos}\)\\
2.196 The internal pressure being \(\frac{a}{V^{2}}\), the work done in condensation is

\[
\int_{V_{l}}^{V_{g}} \frac{a}{V^{2}} d V=\frac{a}{V_{l}}-\frac{a}{V_{g}} \approx \frac{a}{V_{l}}
\]

This by assumption is \(M q, M\) being the molecular weight and \(V_{l}, V_{g}\) being the molar volumes of the liquid and gas.

Thus

\[
p_{l}=\frac{a}{V_{l}^{2}}=\frac{M q}{V_{l}}=\rho q
\]

where \(\rho\) is the density of the liquid. For water \(p_{i} \equiv 3.3 \times 10^{13} \mathrm{~atm}\)\\
2.197 The Vandar Waal's equation can be written as (for one mole)

\[
p(V)=\frac{R T}{V-b}-\frac{a}{V^{2}}
\]

At the critical point \(\left(\frac{\partial p}{\partial V}\right)_{T}\) and \(\left(\frac{\partial^{2} p}{\partial V^{2}}\right)_{T}\) vanish. Thus

\[
\begin{aligned}
& 0=\frac{R T}{(V-b)^{2}}+\frac{2 a}{V^{3}} \text { or } \frac{R T}{(V-b)^{2}}=\frac{2 a}{V^{3}} \\
& 0=\frac{2 R T}{(V-b)^{3}}-\frac{6 a}{V^{4}} \text { or } \frac{R T}{(V-b)^{3}}=\frac{3 a}{V^{4}}
\end{aligned}
\]

Solving these simultaneously we get on division

\[
V-b=\frac{2}{3} V, \quad \mathrm{~V}=3 b \cong V_{M C r}
\]

This is the critical molar volume. Putting this back

\[
\frac{R T_{C r}}{4 b^{2}}=\frac{2 a}{27 b^{3}} \quad \text { or } \quad T_{C r}=\frac{8 a}{27 b R}
\]

Finally

\[
p_{C r}=\frac{R T_{C r}}{V_{M C r}-b}-\frac{a}{V_{M C r}^{2}}=\frac{4 a}{27 b^{2}}-\frac{a}{9 b^{2}}=\frac{a}{27 b^{2}}
\]

From these we see that \(\frac{P_{C r} V_{M C r}}{R T_{C r}}=\frac{a / 9 b}{8 a / 27 b}=\frac{3}{8}\)\\
\(2.198 \frac{P_{C r}}{R T_{C r}}=\frac{a / 27 b^{2}}{8 a / 27 b}=\frac{1}{8 b}\)\\
Thus

\[
b=R \frac{T_{C r}}{8 p_{C r}}=\frac{0.082 \times 304}{73 \times 8}=0.043 \text { litre } / \mathrm{mol}
\]

and \(\quad \frac{\left(R T_{C r}\right)^{2}}{p_{C r}}=\frac{64 a}{27}\) or \(a=\frac{27}{64}\left(R T_{C r}\right)^{2} / p_{C r}=3.59 \frac{\mathrm{~atm} \cdot \mathrm{litre}^{2}}{(\mathrm{~mol})^{2}}\)\\
2.199 Specific volume is molar volume divided by molecular weight. Thus

\[
V_{C r}^{\prime}=\frac{V_{M C r}}{M}=\frac{3 R T_{C r}}{8 M P_{C r}}=\frac{3 \times .082 \times 562}{8 \times 78 \times 47} \frac{\text { litre }}{g}=4.71 \frac{\mathrm{cc}}{g}
\]

\(2.200\left(p+\frac{a}{V_{M}^{2}}\right)\left(V_{m}-b\right)=R T\)\\
or

\[
\frac{p+\frac{a}{V_{m}^{2}}}{P_{C r}} \times \frac{V_{m}-b}{V_{M C r}}=\frac{8}{3} \frac{T}{T_{C r}}
\]

or

\[
\left(\pi+\frac{a}{p_{C r} V_{m}^{2}}\right) \times\left(v-\frac{b}{V_{M C r}}\right)=\frac{8}{3} \tau,
\]

where

\[
\pi=\frac{p}{p_{C r}}, v=\frac{V_{m}}{V_{M C r}}, \tau=\frac{T}{T_{C r}}
\]

or

\[
\left(\pi+\frac{27 b^{2}}{V_{M}^{2}}\right)\left(v-\frac{1}{3}\right)=\frac{8}{3} \tau, \text { or }\left(\pi+\frac{3}{v^{2}}\right)\left(v-\frac{1}{3}\right)=\frac{8}{3} \tau
\]

When

\[
\pi=12 \text { and } v=\frac{1}{2}, \tau=\frac{3}{8} \times 24 \times \frac{1}{6}=\frac{3}{2}
\]

2.201 (a) The ciritical Volume \(V_{M C r}\) is the maximum volume in the liquid phase and the minimum volume in the gaseous. Thus

\[
V_{\max }=\frac{1000}{18} \times 3 \times 030 \text { litre } \approx 5 \text { litre }
\]

(b) The critical pressure is the maximum possible pressure in the vapour phase in equilibrium with liquid phase. Thus

\[
p_{\max }=\frac{a}{27 b^{2}}=\frac{5.47}{27 \times .03 \times .03}=225 \text { atmosphere }
\]

\(2.202 T_{C r}=\frac{8}{27} \frac{a}{b R}=\frac{8}{27} \times \frac{3.62}{.043 \times .082} \approx 304 \mathrm{~K}\)

\[
\rho_{c_{r}}=\frac{M}{3 b}=\frac{44}{3 \times 43} \mathrm{gm} / \text { c.c. }=0.34 \mathrm{gm} / \text { c.c. }
\]

2.203 The vessel is such that either vapour or liquid of mass \(m\) occupies it at critical point. Then its volume will be

\[
v_{C r}=\frac{m}{M} V_{M C r}=\frac{3}{8} \frac{R T_{C r}}{p_{C r}} \frac{m}{M}
\]

The corresponding volume in liquid phase at room temperature is

\[
V=\frac{m}{\rho}
\]

where \(\rho=\) density of liquid ether at room tmeperature. Thus

\[
\eta=\frac{V}{V_{C r}}=\frac{8 M p_{C r}}{3 R T_{C r} \rho} \approx 0.254
\]

using the given data (and \(\rho=720 \mathrm{gm}\) per litre)\\
2.204 We apply the relation ( \(T=\) constant)

\[
T \oint d S=\oint d U+\oint p d V
\]

to the cycle 1234531.\\
Here \(\oint d S=\oint d U=0\)\\
So \(\quad \oint p d V=0\)\\
This implies that the areas I and II are equal. This reasoning is inapplicable to the cycle 1231, for example. This cycle is irreversible because it involves the irreversible transition from a single phase to a two-phase state at the point 3 .\\
\includegraphics[max width=\textwidth, alt={}, center]{9f3c5a8e-ac5d-42f4-8eb6-ea2af623f2d8-259_478_576_1197_785}\\
2.205 When a portion of supercool water turns into ice some heat is liberated, which should heat it upto ice point. Neglecting the variation of specific heat of water, the fraction of water turning inot ice is clearly

\[
f=\frac{c|t|}{q}=0.25
\]

where \(c \boldsymbol{=}\) specific heat of water and \(q \boldsymbol{\sim}\) latent heat of fusion of ice, Clearly \(f=1\) at \(t=-80^{\circ} \mathrm{C}\)\\
2.206 From the Claussius-Clapeyron ( \(C-C\) )equations

\[
\frac{d T}{d p}=\frac{T\left(V_{2}^{\prime}-V_{1}^{\prime}\right)}{q_{12}}
\]

\(\boldsymbol{q}_{12}\) is the specific latent heat absorbed in \(1 \rightarrow 2(1=\) solid, \(2=\) liquid \()\)

\[
\begin{aligned}
& \Delta T=\frac{T\left(V_{w}^{\prime}-V_{\text {ice }}^{\prime}\right)}{q_{12}} \Delta p=-\frac{273 \times .091}{333} \times 1 \frac{\mathrm{~atm} \times \mathrm{cm}^{3} \times \mathrm{K}}{\text { joule }} \\
& 1 \frac{\mathrm{~atm} \times \mathrm{cm}^{3}}{\text { Joule }} \approx \frac{10^{5} \frac{\mathrm{~N}}{\mathrm{~m}^{2}} \times 10^{-6} \mathrm{~m}^{3}}{\text { Joule }}=10^{-1}, \Delta T=-.0075 \mathrm{~K}
\end{aligned}
\]

2.207 Here 1 = liquid, 2 = Steam

\[
\Delta T=\frac{T\left(V_{s}^{\prime}-V_{l i q}^{\prime}\right)}{q_{12}} \Delta P
\]

or

\[
V_{s}^{\prime} \approx \frac{q_{12}}{T} \frac{\Delta T}{\Delta p}=\frac{2250}{373} \times \frac{0.9}{3.2} \times 10^{-3} \dot{\mathrm{~m}}^{3} / \mathrm{g}=1.7 \mathrm{~m}^{3} / \mathrm{kg}
\]

2.208 From \(C-C\) equations

\[
\frac{d p}{d T}=\frac{q_{12}}{T\left(V_{2}^{\prime}-V_{1}^{\prime}\right)} \approx \frac{q_{12}}{T V_{2}^{\prime}}
\]

Assuming the saturated vapour to be ideal gas

\[
\frac{1}{V_{2}^{\prime}}=\frac{m p}{R T}, \text { Thus } \Delta p=\frac{M q}{R T^{2}} p \Delta T
\]

and

\[
p \approx p_{0}\left(1+\frac{M q}{R T^{2}} \Delta T\right) \cong 1 \cdot 04 \text { atmosphere }
\]

2.209 From \(C-C\) equation, neglecting the voolume of the liquid\\
or

\[
\begin{gathered}
\frac{d p}{d T} \cong \frac{q_{12}}{T V_{2}^{\prime}} \cong \frac{M q}{R T^{2}} p,\left(q=q_{12}\right) \\
\frac{d p}{p}=\frac{M q}{R T} \frac{d T}{T}
\end{gathered}
\]

Now

\[
p V=\frac{m}{M} R T \text { or } m=\frac{M p V}{R T} \text { for a perfect gas }
\]

So

\[
\begin{gathered}
\quad \frac{d m}{m}=\frac{d p}{p}-\frac{d T}{T}(V \text { is Const }=\text { specific volume }) \\
=\left(\frac{M q}{R T}-1\right) \frac{d T}{T}=\left(\frac{18 \times 2250}{8.3 \times 373}-1\right) \times \frac{1.5}{373} \approx 4.85 \%
\end{gathered}
\]

2.210 From C-C equation

\[
\frac{d p}{d T} \approx \frac{q}{T V_{2}^{\prime}}=\frac{M q}{R T^{2}} p
\]

Integrating \(\ln p=\) constant \(-\frac{M q}{R T}\)\\
So

\[
p=p_{0} \exp \left[\frac{M q}{R}\left(\frac{1}{T_{0}}-\frac{1}{T}\right)\right]
\]

This is reasonable for \(\left|T-T_{0}\right| \ll T_{0}\), and far below critical temperature.\\
2.211 As before (2.206) the lowering of melting point is given by

\[
\Delta T=-\frac{T \Delta V^{\prime}}{q} p
\]

The superheated ice will then melt in part. The fraction that will melt is

\[
\eta=\frac{C T \Delta V^{\prime}}{q^{2}} p \sim .03
\]

2.212 (a) The equations of the transition lines are

\[
\begin{aligned}
& \log p=9.05-\frac{1800}{T}: \text { Solid gas } \\
& \quad=6.78-\frac{1310}{T}: \text { Liquid gas }
\end{aligned}
\]

At the triple point they intersect. Thus

\[
2 \cdot 27=\frac{490}{T_{t r}} \text { or } T_{t r}=\frac{490}{2 \cdot 27}=216 \mathrm{~K}
\]

corresponding \(p_{t r}\) is 5.14 atmosphere.\\
In the formula \(\log p=a-\frac{b}{T}\), we compare \(b\) with the corresponding term in the equation in 2.210. Then\\
or,

\[
\begin{aligned}
& \ln p=a \times 2.303-\frac{2.303 b}{T} \text { So, } 2.303=\frac{M q}{R} \\
& q_{\text {sublimation }}=\frac{2.303 \times 1800 \times 8.31}{44}=783 \mathrm{~J} / \mathrm{gm} \\
& q_{\text {liquid-gas }}=\frac{2.303 \times 1310 \times 8.31}{44}=570 \mathrm{~J} / \mathrm{gm}
\end{aligned}
\]

Finally \(q_{\text {solid-liquid }}=213 \mathrm{~J} / \mathrm{gm}\) on subtraction

\[
2.213 \begin{aligned}
\Delta S=\int_{T_{1}} m c \frac{d T}{T}+\frac{m q}{T_{2}} & \cong m\left(c \ln \frac{T_{2}}{T_{1}}+\frac{q}{T_{2}}\right) \\
& =10^{3}\left(4.18 \ln \frac{373}{283}+\frac{2250}{373}\right) \approx 7.2 \mathrm{~kJ} / \mathrm{K}
\end{aligned}
\]

\(2.214 \Delta S=\frac{q_{m}}{T_{1}}+c \ln \frac{T_{2}}{T}+\frac{q_{v}}{T_{2}}\)

\[
=\frac{333}{273}+4 \cdot 18 \ln \frac{373}{283} \approx 8 \cdot 56 \mathrm{~J} /{ }^{\circ} \mathrm{K}
\]

\(2.215 c=\) specific heat of copper \(=0.39 \cdot \frac{\mathrm{~J}}{\mathrm{~g} \cdot \mathrm{~K}} \cdot\) Suppose all ice does not melt, then heat rejected \(=90 \times 0.39(90-0)=3159 \mathrm{~J}\)\\
heat gained by ice \(=50 \times 2.09 \times 3+x \times 333\)\\
Thus

\[
x=8.5 \mathrm{gm}
\]

The hypothesis is correct and final temperature will be \(T=273 \mathrm{~K}\).\\
Hence change in entropy of copper piece

\[
=m c \ln \frac{273}{363}=-10 \mathrm{~J} / \mathrm{K}
\]

2.216 (a) Here \(t_{2}=60^{\circ} \mathrm{C}\). Suppose the final temperature is \(t{ }^{\circ} \mathrm{C}\). Then

\[
\text { heat lost by water }=m_{2} c\left(t_{2}-t\right)
\]

heat gained by ice \(=m_{1} q_{m}+m_{1} c\left(t-t_{1}\right)\), if all ice melts\\
In this case \(m_{1} q_{m}=m_{2} \times 4 \cdot 18(60-t)\), for \(m_{1}=m_{2}\)\\
So the final temperature will be \(0^{\circ} \mathrm{C}\) and only some ice will melt.\\
Then

\[
100 \times 4 \cdot 18(60)=m_{1}^{\prime} \times 333
\]

\(m_{1}^{\prime}=75.3 \mathrm{gm}=\) amount of ice that will melt

Finally

\[
\begin{aligned}
\Delta S= & 75.3 \times \frac{333}{273}+100 \times 4.18 \ln \frac{273}{333} \\
& \Delta S=\frac{m_{1}^{\prime} q_{m}}{T_{1}}+m_{2} c \ln \frac{T_{1}}{T_{2}} \\
= & m_{2} c \frac{\left(T_{2}-T_{1}\right)}{T_{1}}-m_{2} \ln \frac{T_{2}}{T_{1}} \\
= & m_{2} C\left[\frac{T_{2}}{T_{1}}-1-\ln \frac{T_{2}}{T_{1}}\right]=8.8 \mathrm{~J} / \mathrm{K}
\end{aligned}
\]

(b) If \(m_{2} c t_{2}>m_{1} q_{m}\) then all ice will melt as one can check and the final temperature can be obtained like this\\
or

\[
\begin{aligned}
& m_{2} c\left(T_{2}-T\right)=m_{1} q_{m}+m_{1} c\left(T-T_{1}\right) \\
& \left(m_{2} T_{2}+m_{1} T_{1}\right) c-m_{1} q_{m}=\left(m_{1}+m_{2}\right) c T
\end{aligned}
\]

\[
T=\frac{m_{2} T_{2}+m_{1} T_{1}-\frac{m_{1} q_{m}}{c}}{m_{1}+m_{2}} \cong 280 \mathrm{~K}
\]

and

\[
\Delta S=\frac{m_{1} q}{T_{1}}+c\left(m_{1} \ln \frac{T}{T_{1}}-m_{2} \ln \frac{T_{2}}{T}\right)=19 \mathrm{~J} / \mathrm{K}
\]

\(2.217 \quad \Delta S=-\frac{m q_{1}}{T_{2}}-m c \ln \frac{T_{2}}{T_{1}}+\frac{M q_{\text {ice }}}{T_{1}}\)\\
where

\[
\begin{aligned}
& M q_{\text {ice }}=m\left(q_{2}+c\left(T_{2}-T_{1}\right)\right) \\
= & m q_{2}\left(\frac{1}{T_{1}}-\frac{1}{T_{2}}\right)+m c\left(\frac{T_{2}}{T_{1}}-1-\frac{T_{2}}{T_{1}}\right) \\
= & 0.2245+0.2564 \approx 0.48 \mathrm{~J} / \mathrm{K}
\end{aligned}
\]

2.218 When heat \(d Q\) is given to the vapour its temperature will change by \(d T\), pressure by \(d p\) and volume by \(d V\), it being assumed that the vapour remains saturated.\\
Then by C-C equation

\[
\frac{d p}{d T}=\frac{q}{T V^{\prime}}\left(V_{\text {vapour }}^{\prime} \gg V_{\mathrm{Lq}}^{\prime}\right), \text { or } d p=\frac{q}{T V^{\prime}} d T
\]

on the other hand, \(p V^{\prime}=\frac{R T}{M}\)\\
So

\[
p d V^{\prime}+V^{\prime} d p=\frac{R d T}{M}
\]

Hence

\[
p d V^{\prime}=\left(\frac{R}{M}-\frac{q}{T}\right) d T
\]

finally

\[
d Q=C d T=d U+p d V^{\prime}
\]

\[
=C_{V} d T+\left(\frac{T}{M}-\frac{q}{T}\right) d T=C_{p} d T-\frac{q}{T} d T
\]

( \(C_{p}, C_{V}\) refer to unit mass here). Thus

\[
C=C_{p}-\frac{q}{T}
\]

For water \(C_{p}=\frac{R \gamma}{\gamma-1} \cdot \frac{1}{M}\) with \(\gamma=1.32\) and \(M=18\)\\
So

\[
C_{p}=1.90 \mathrm{~J} / \mathrm{gm} \mathrm{~K}
\]

and

\[
C=-4 \cdot 13 \mathrm{~J} / \mathrm{gm}^{\circ} \mathrm{K}=-74 \mathrm{~J} / \text { mole } \mathrm{K}
\]

2.219 The required entropy change can be calculated along a process in which the water is heated from \(T_{1}\) to \(T_{2}\) and then allowed to evaporate. The entropy change for this is

\[
\Delta S=C_{p} \ln \frac{T_{2}}{T_{1}}+\frac{q M}{T_{2}}
\]

where \(q=\) specific latent heat of vaporization.

\subsection*{2.7 TRANSPORT PHENOMENA}
2.220 (a) The fraction of gas molecules which traverses distances exceeding the mean free path without collision is just the probability to traverse the distance \(s=\lambda\) without collision.

Thus

\[
P=e^{-1}=\frac{1}{e}=0.37
\]

(b) This probability is

\[
P=e^{-1}-e^{-2}=0.23
\]

2.221 From the formula

\[
\frac{1}{\eta}=e^{-\Delta l / \lambda} \text { or } \lambda=\frac{\Delta l}{\ln \eta}
\]

2.222 (a) Let \(P(t)=\) probability of no collision in the interval \((0, t)\). Then

\[
\begin{gathered}
P(t+d t)=P(t)(1-\alpha d t) \\
\frac{d P}{d t}=-\alpha P(t) \quad \text { or } P(t)=e^{-\alpha t}
\end{gathered}
\]

or\\
where we have used \(P(0)=1\)\\
(b) The mean interval between collision is also the mean interval of no collision. Then

\[
<t>=\frac{\int_{0}^{\infty} t e^{-\alpha t} d t}{\int_{0}^{\infty} e^{-\alpha t} d t}=\frac{1}{\alpha} \frac{\Gamma(2)}{\Gamma(1)}=\frac{1}{\alpha}
\]

2.223\\
(a) \(\lambda=\frac{1}{\sqrt{2} \pi d^{2} n}=\frac{k T}{\sqrt{2} \pi d^{2} p}\)

\[
\begin{gathered}
=\frac{1.38 \times 10^{-23} \times 273}{\sqrt{2} \pi\left(0.37 \times 10^{-9}\right)^{2} \times 10^{5}}=6.2 \times 10^{-8} \mathrm{~m} \\
\tau=\frac{\lambda}{\langle v\rangle}=\frac{6.2 \times 10^{-8}}{454} s=136 \mathrm{n} \mathrm{~s}
\end{gathered}
\]

\(\lambda=6.2 \times 10^{6} \mathrm{~m}\)\\
(b) \(\eta=1.36 \times 10^{4} \mathrm{~s}=3.8\) hours\\
2.224 The mean distance between molecules is of the order

\[
\left(\frac{22.4 \times 10^{-3}}{6.0 \times 10^{23}}\right)^{1 / 3}=\left(\frac{224}{6}\right)^{1 / 3} \times 10^{-9} \text { meters } \approx 3.34 \times 10^{-9} \text { meters }
\]

This is about 18.5 times smaller than the mean free path calculated in 2.223 (a) above.\\
2.225 We know that the Vander Waal's constant \(b\) is four times the molecular volume. Thus

Hence

\[
\begin{gathered}
b=4 N_{A} \frac{\pi}{6} d^{3} \text { or } d=\left(\frac{3 b}{2 \pi N_{A}}\right)^{1 / 3} \\
\lambda=\left(\frac{k T_{0}}{\sqrt{2} \pi p_{0}}\right)\left(\frac{2 \pi N_{A}}{3 b}\right)^{2 / 3}
\end{gathered}
\]

2.226 The volocity of sound in \(N_{2}\) is

\[
\begin{array}{cc}
\text { so, } & \sqrt{\frac{Y P}{\rho}}=\sqrt{\frac{Y R T}{M}} \\
\text { or, } & \frac{1}{v}=\sqrt{\frac{\gamma R T_{0}}{M}}=\frac{R T_{0}}{\sqrt{2} \pi d^{2} p_{0} N_{A}} \\
v=\pi d^{2} p_{0} N_{A} \sqrt{\frac{2 \gamma}{M R T_{0}}}
\end{array}
\]

2.227 (a) \(\lambda>l\) if \(p<\frac{k T}{\sqrt{2} \pi d^{2} l}\)

Now \(\frac{k T}{\sqrt{2} \pi d^{2} l}\) for \(O_{2}\) of \(O\) is 0.7 Pa .\\
(b) The corresponding \(n\) is obtained by dividing by \(k T\) and is \(1.84 \times 10^{20}\) per \(m^{3}=1.84^{14}\) per c.c. and the corresponding mean distance is \(\frac{l}{n^{1 / 3}}\).

\[
=\frac{10^{-2}}{(0.184)^{1 / 3} \times 10^{5}}=1.8 \times 10^{-} m \approx 0.18 \mu \mathrm{~m} .
\]

2.228\\
(a) \(\nu=\frac{1}{\tau}=\frac{1}{\lambda /\langle\nu\rangle}=\frac{\langle\nu\rangle}{\lambda}\)

\[
=\sqrt{2} \pi d^{2} n\langle\nu\rangle=.74 \times 10^{10} \mathrm{~s}^{-1}(\text { see } 2.223)
\]

(b) Total number of collisions is

\[
\frac{1}{2} n v \approx 1.0 \times 10^{29} \mathrm{~s} \mathrm{~cm}^{-3}
\]

Note, the factor \(\frac{1}{2}\). When two molecules collide we must not count it twice.\\
2.229 (a) \(\lambda=\frac{1}{\sqrt{2} \pi d^{2} n}\)\\
\(d\) is a constant and \(n\) is a constant for an isochoric process so \(\lambda\) is constant for an isochoric process.

\[
v=\frac{\langle v\rangle}{\lambda}=\frac{\sqrt{\frac{8 R T}{M \pi}}}{\lambda} \alpha \sqrt{T}
\]

(b) \(\lambda=\frac{1}{\sqrt{2} \pi d^{2}} \frac{k T}{p} \alpha T\) for an isobaric process.

\[
v=\frac{\langle v\rangle}{\lambda} \alpha \frac{\sqrt{T}}{T}=\frac{1}{\sqrt{T}} \text { for an isobaric process. }
\]

2.230 (a) In an isochoric process \(\lambda\) is constant and

\[
v \propto \sqrt{T} \propto \sqrt{p V} \propto \sqrt{p} \alpha \sqrt{n}
\]

(b) \(\lambda=\frac{k T}{\sqrt{2} \pi d^{2} p}\) must decrease \(n\) times in an isothermal process and \(v\) must increase \(n\) times because \(<v>\) is constant in an isothermal process.\\
2.231 (a) \(\lambda \alpha \frac{1}{n} \Rightarrow \frac{1}{N / V}=\frac{V}{N}\)

Thus

\[
\lambda \alpha V \text { and } v \alpha \frac{T^{1 / 2}}{V}
\]

But in an adiabatic process \(\left(\gamma=\frac{7}{5}\right.\) here \()\)

\[
T V^{\gamma-1}=\text { constant so } T V^{2 / 5}=\text { constant }
\]

or

\[
T^{1 / 2} \alpha V^{-1 / 5} \text { Thus } v \alpha V^{-6 / 5}
\]

(b) \(\lambda \alpha \frac{T}{p}\)

But \(p\left(\frac{T}{p}\right)^{\gamma}=\) constant or \(\frac{T}{p} \alpha p^{-1 / \gamma}\) or \(T \alpha p^{1-1 / \gamma}\)\\
Thus

\[
\begin{gathered}
\lambda \alpha p^{-1 / \gamma}=p^{-5 / 7} \\
v=\frac{\langle v\rangle}{\lambda} \alpha \frac{p}{\sqrt{T}} \alpha p^{1 / 2+\frac{1}{2 \gamma}}=p^{\frac{\gamma+1}{2 \gamma}}=p^{6 / 7}
\end{gathered}
\]

(c) \(\lambda \propto V\)

But

\[
T V^{2 / 5}=\text { constant or } V \alpha T^{-5 / 2}
\]

Thus

\[
\begin{gathered}
\lambda \alpha T^{-5 / 2} \\
v \alpha \frac{T^{1 / 2}}{V} \alpha T^{3}
\end{gathered}
\]

2.232 In the polytropic process of index \(n\)

\[
p V^{n}=\text { constant }, T V^{n-1}=\text { constant and } p^{1-n} T^{n}=\text { constant }
\]

(a) \(\lambda \alpha V\)

\[
v \alpha \frac{T^{1 / 2}}{V}=V^{\frac{1-n}{2}} V^{-1}=V^{\frac{-n+1}{2}}
\]

(b) \(\lambda \alpha \frac{T}{p}, T^{n} \alpha p^{n-1}\) or \(T \alpha p^{1-\frac{1}{n}}\)\\
so

\[
\begin{gathered}
\lambda \alpha p^{-1 / n} \\
v=\frac{\langle v\rangle}{\lambda} \alpha \frac{p}{\sqrt{T}} \alpha p^{1-\frac{1}{2}+\frac{1}{2 n}}=p^{\frac{n+1}{2 n}}
\end{gathered}
\]

(c) \(\lambda \alpha \frac{T}{p}, p \alpha T^{\frac{n}{n-1}}\)

\[
\begin{aligned}
& \lambda \alpha T^{1-\frac{n}{n 1}}=T^{-\frac{1}{n-1}}=T^{\frac{1}{1-n}} \\
& v \alpha \frac{p}{\sqrt{T}} \alpha T^{\frac{n}{n-1}-\frac{1}{2}}=T^{\frac{n+1}{2(n-1)}}
\end{aligned}
\]

2.233 (a) The number of collisions between the molecules in a unit volume is

\[
\frac{1}{2} n v=\frac{1}{\sqrt{2}} \pi d^{2} n^{2}<v>\alpha \frac{\sqrt{T}}{V^{2}}
\]

This remains constant in the poly process \(p V^{-3}=\) constant\\
Using (2.122) the molar specific heat for the polytropic process\\
is

Thus

\[
\begin{gathered}
p V^{\alpha}=\text { constant, } \\
C=R\left(\frac{1}{\gamma-1}-\frac{1}{\alpha-1}\right) \\
C=R\left(\frac{1}{\gamma-1}+\frac{1}{4}\right)=R\left(\frac{5}{2}+\frac{1}{4}\right)=\frac{11}{4} R
\end{gathered}
\]

It can also be written as \(\frac{1}{4} R(1+2 i)\) where \(i=5\)\\
(b) In this case \(\frac{\sqrt{T}}{V}=\) constant and so \(p V^{-1}=\) constant\\
so

\[
C=R\left(\frac{1}{\gamma-1}+\frac{1}{2}\right)=R\left(\frac{5}{2}+\frac{1}{2}\right)=3 R
\]

It can also be written as \(\frac{R}{2}(i+1)\)\\
2.234 We can assume that all molecules, incident on the hole, leak out. Then,

\[
-d N=-d(n V)=\frac{1}{4} n<v>S d t
\]

or

\[
d n=-n \frac{d t}{4 v / S\langle v\rangle}=-n \frac{d t}{\tau}
\]

Integrating \(\quad n=n_{0} e^{-t / \tau}\). Hence \(<v>=\sqrt{\frac{8 R T}{\pi M}}\)\\
2.235 If the temperature of the compartment 2 is \(\eta\) times more than that of compartment 1 , it must contain \(\frac{1}{\eta}\) times less number of molecules since pressure must be the same when the big hole is open. If \(M=\) mass of the gas in 1 than the mass of the gas in 2 must be \(\frac{M}{\eta}\). So immediately after the big hole is closed.

\[
n_{1}^{0}=\frac{M}{m V}, n_{2}^{0}=\frac{M}{m V \eta}
\]

where \(m=\) mass of each molecule and \(n_{1}^{0}, n_{2}^{0}\) are concentrations in 1 and 2 . After the big hole is closed the pressures will differ and concentration will become \(n_{1}\) and \(n_{2}\) where

\[
n_{1}+n_{2}=\frac{M}{m V \eta}(1+\eta)
\]

On the other hand

\[
n_{1}<v_{1}>=n_{2}<v_{2}>\text { i.e. } n_{1}=\sqrt{\eta} n_{2}
\]

Thus

\[
n_{2}(1+\sqrt{\eta})=\frac{m}{m V \eta}(1+\eta)=n_{2}^{0}(1+\eta)
\]

So

\[
n_{2}=n_{2}^{0} \frac{1+\eta}{1+\sqrt{\eta}}
\]

2.236 We know

\[
\eta=\frac{1}{3}<v>\lambda \rho=\frac{1}{3}<v>\frac{1}{\sqrt{2} \pi d^{2}} m \alpha \sqrt{T}
\]

Thus \(\eta\) changing \(\alpha\) times implies \(T\) changing \(\alpha^{2}\) times.\\
On the other hand

\[
D=\frac{1}{3}<v>\lambda=\frac{1}{3} \sqrt{\frac{8 k T}{\pi m}} \frac{k T}{\sqrt{2} \pi d^{2} p}
\]

Thus \(D\) changing \(\beta\) times means \(\frac{T^{3 / 2}}{p}\) changing \(\beta\) times\\
So \(p\) must change \(\frac{\alpha^{3}}{\beta}\) times\\
\(2.237 D \propto \frac{\sqrt{T}}{n} \propto V \sqrt{T}, \eta \propto \sqrt{T}\)\\
(a) \(D\) will increase \(n\) times\\
\(\eta\) will remain constant if \(T\) is constant\\
(b) \(D \alpha \frac{T^{3 / 2}}{p} \alpha \frac{(p V)^{3 / 2}}{p}=p^{1 / 2} V^{3 / 2}\)

\[
\eta \alpha \sqrt{p V}
\]

Thus \(D\) will increase \(n^{3 / 2}\) times, \(\eta\) will increase \(n^{1 / 2}\) times, if \(p\) is constant \(2.238 D \alpha V \sqrt{T}, \eta \alpha \sqrt{T}\)

In an adiabatic process

\[
T V^{\gamma-1}=\text { constant, or } T \propto V^{1-}
\]

Now \(V\) is decreased \(\frac{1}{n}\) times. Thus

\[
\begin{aligned}
& D \alpha V^{\frac{3-\gamma}{2}}=\left(\frac{1}{n}\right)^{\frac{3-\gamma}{2}}=\left(\frac{1}{n}\right)^{4 / 5} \\
& \eta \alpha \text { of } V^{\frac{1-\gamma}{2}}=\left(\frac{1}{n}\right)^{-1 / 5}=n^{1 / 5}
\end{aligned}
\]

So \(D\) decreases \(n^{4 / 5}\) times and \(\eta\) increase \(n^{1 / 5}\) times.\\
2.239 (a) \(D \alpha V \sqrt{T} \alpha \sqrt{p V^{3}}\)

Thus \(D\) remains constant in the process \(p V^{3}=\) constant\\
So polytropic index \(n=3\)\\
(b) \(\eta \propto \sqrt{T} \alpha \sqrt{p V}\)

So \(\eta\) remains constant in the isothermal process

\[
p V=\text { constant, } n=1, \text { here }
\]

(c) Heat conductivity \(\kappa=\eta C_{V}\)\\
and \(C_{V}\) is a constant for the ideal gas\\
Thus

\[
n=1 \text { here also, }
\]

\(2.240 \eta=\frac{1}{3} \sqrt{\frac{8 k T}{\pi m}} \frac{m}{\sqrt{2} \pi d^{2}}=\frac{2}{3} \sqrt{\frac{m k T}{\pi^{3}}} \frac{1}{d^{2}}\)\\
or \(d=\left(\frac{2}{3 \eta}\right)^{1 / 2}\left(\frac{m k T}{\pi^{3}}\right)^{1 / 4}=\left(\frac{2}{3 \times 18.9} \times 10^{6}\right)^{1 / 2}\left(\frac{4 \times 8.31 \times 273 \times 10^{-3}}{\pi^{3} \times 36 \times 10^{46}}\right)^{1 / 4}\)

\[
=10^{-10}\left(\frac{2}{3 \times 18.9}\right)^{1 / 2}\left(\frac{4 \times 83.1 \times 273}{\pi^{3} \times .36}\right)^{1 / 4} \approx 0.178 \mathrm{~nm}
\]

\(2.241 \kappa=\frac{1}{3}<v>\lambda \rho c_{v}\)

\[
=\frac{1}{3} \sqrt{\frac{8 k T}{m \pi}} \frac{1}{\sqrt{2} \pi d^{2} n} m n \frac{C_{V}}{M}
\]

\(\left(C_{V}\right.\) is the specific heat capacity which is \(\left.\frac{C_{V}}{M}\right)\). Now \(C_{V}\) is the same for all monoatomic gases such as He and \(A\). Thus\\
or

\[
\begin{gathered}
\kappa \alpha \frac{1}{\sqrt{M} d^{2}} \\
\frac{\kappa_{H_{e}}}{\kappa_{A}}=8.7=\frac{\sqrt{M_{A}} d_{A}^{2}}{\sqrt{M_{H_{e}}} d_{H_{e}}^{2}}=\sqrt{10} \frac{d_{A}^{2}}{d_{H_{e}}^{2}} \\
\frac{d_{A}}{d_{H_{e}}}=\sqrt{\frac{8.7}{\sqrt{10}}}=1.658 \approx 1.7
\end{gathered}
\]

2.242 In this case

\[
N_{1} \frac{r_{2}^{2}-r_{1}^{2}}{r_{1}^{2} r_{2}^{2}}=4 \pi \eta \omega
\]

or

\[
N_{1} \frac{2 R \Delta R}{R^{4}} \approx 4 \pi \eta \omega \text { or } N_{1}=\frac{2 \pi \eta \omega R^{3}}{\Delta R}
\]

To decrease \(N_{1}, n\) times \(\eta\) must be decreased \(n\) times. Now \(\eta\) does not depend on pressure until the pressure is so low that the mean free path equals, say, \(\frac{1}{2} \Delta R\). Then the mean free path is fixed and \(\eta\) decreases with pressure. The mean free path cquals \(\frac{1}{2} \Delta R\) when

\[
\frac{1}{\sqrt{2} \pi d^{2} n_{0}}=\Delta R\left(n_{0}=\text { concentration }\right)
\]

Corresponding pressure is \(p_{0}=\frac{\sqrt{2} k T}{\pi d^{2} \Delta R}\)\\
The sought pressure is \(n\) times less

\[
p=\frac{\sqrt{2} k T}{\pi d^{2} n \Delta R}=70.7 \times \frac{10^{-23}}{10^{-18} \times 10^{-3}} \approx 0.71 \mathrm{~Pa}
\]

The answer is qualitative and depends on the choice \(\frac{1}{2} \Delta R\) for the mean free path.\\
2.243 We neglect the moment of inertia of the gas in a shell. Then the moment of friction forces on a unit length of the cylinder must be a constant as a function of \(r\).\\
So, \(\quad 2 \pi r^{3} \eta \frac{d \omega}{d r}=N_{1}\) or \(\omega(r)=\frac{N_{1}}{4 \pi \eta}\left(\frac{1}{r_{1}^{2}}-\frac{1}{r^{2}}\right)\)\\
and

\[
\omega=\frac{N_{1}}{4 \pi \eta}\left(\frac{1}{r_{1}^{2}}-\frac{1}{r_{2}^{2}}\right) \quad \text { or } \quad \eta=\frac{N_{1}}{4 \pi \omega}\left(\frac{1}{r_{1}^{2}}-\frac{1}{r_{2}^{2}}\right)
\]

2.244 We consider two adjoining layers. The angular velocity gradient is \(\frac{\omega}{h}\). So the moment of the frictional force is

\[
N=\int_{0}^{a} r \cdot 2 \pi r d r \cdot \eta r \frac{\omega}{h}=\frac{\pi \eta a^{4} \omega}{2 h}
\]

2.245 In the ultrararefied gas we must determine \(\eta\) by taking \(\lambda=\frac{1}{2} h\). Then

\[
\begin{array}{rlrl}
\eta=\frac{1}{3} \sqrt{\frac{8 k T}{m \pi}} \times \frac{1}{2} h \times \frac{m p}{k T}=\frac{1}{3} \sqrt{\frac{2 M}{\pi R T}} h p & \\
\text { so, } & N & =\frac{1}{3} \omega a^{4} p \sqrt{\frac{\pi M}{2 R T}}
\end{array}
\]

2.246 Take an infinitesimal section of length \(d x\) and apply Poiseuilles equation to this. Then

From the formula\\
or

\[
\begin{aligned}
\frac{d V}{d t} & =\frac{-\pi a^{4}}{8 \eta} \frac{\partial p}{\partial x} \\
p V & =R T \cdot \frac{m}{M} \\
p d V & =\frac{R T}{M} d m \\
\frac{d m}{d t}=\mu & =-\frac{\pi a^{4} M}{8 \eta R T} \frac{p d p}{d x}
\end{aligned}
\]

This equation implies that if the flow is isothermal then \(p \frac{d p}{d x}\) must be a constant and so equals \(\frac{\left|p_{2}^{2}-p_{1}^{2}\right|}{2 l}\) in magnitude.

Thus,

\[
\mu=\frac{\pi a^{4} M}{16 \eta R T}\left|\frac{p_{2}^{2}-p_{1}^{2}}{l}\right|
\]

2.247 Let \(T=\) temperature of the interface.

Then heat flowing from left \(=\) heat flowing into right in equilibrium.\\
Thus, \(\kappa_{1} \frac{T_{1}-T}{l_{1}}=\kappa_{2} \frac{T-T_{2}}{l_{2}}\) or \(T=\frac{\left(\frac{\kappa_{1} T_{1}}{l_{1}}+\frac{\kappa_{2} T_{2}}{l_{2}}\right)}{\left(\frac{\kappa_{1}}{l_{1}}+\frac{\kappa_{2}}{l_{2}}\right)}\)\\
2.248 We have

\[
\kappa_{1} \frac{T_{1}-T}{l_{1}}=\kappa_{2} \frac{T-T_{2}}{l_{2}}=\kappa \frac{T_{1}-T_{2}}{l_{1}+l_{2}}
\]

or using the previous result

\[
\frac{\kappa_{1}}{l_{1}}\left(T_{1}-\frac{\frac{\kappa_{1} T_{1}}{l_{1}}+\frac{\kappa_{2} T_{2}}{l_{2}}}{\frac{\kappa_{1}}{l_{1}}+\frac{\kappa_{1}}{l_{2}}}\right)=\kappa \frac{T_{1}-T_{2}}{l_{1}+l_{2}}
\]

or \(\quad \frac{\kappa_{1}}{l_{1}} \frac{\frac{\kappa_{2}}{l_{2}}\left(T_{1}-T_{2}\right)}{\frac{\kappa_{1}}{l_{1}}+\frac{\kappa_{2}}{l_{2}}}=\kappa \frac{T_{1}-T_{2}}{l_{1}+l_{2}}\) or \(\kappa=\frac{l_{1}+l_{2}}{\frac{l_{1}}{\kappa_{1}}+\frac{l_{2}}{-\kappa_{2}}}\)\\
2.249 By definition the heat flux (per unit area) is

\[
\dot{Q}=-K \frac{d T}{d x}=-\alpha \frac{d}{d x} \ln T=\text { constant }=+\alpha \frac{\ln T_{1} / T_{2}}{l}
\]

Integrating

\[
\ln T=\frac{x}{l} \ln \frac{T_{2}}{T_{1}}+\ln T_{1}
\]

where \(T_{1}=\) temperature at the end \(x=0\)

So

\[
T=T_{1}\left(\frac{T_{2}}{T_{1}}\right)^{x / l} \text { and } \dot{Q}=\frac{\alpha \ln T_{1} / T_{2}}{l}
\]

2.250 Suppose the chunks have temperatures \(T_{1}, T_{2}\) at time \(t\) and \(T_{1}-d T_{1}, T_{2}+d T_{2}\) at time \(d t+t\).

Then

\[
C_{1} d T_{1}=C_{2} d T_{2}=\frac{\kappa S}{l}\left(T_{1}-T_{2}\right) d t
\]

Thus

\[
d \Delta T=-\frac{\kappa S}{l}\left(\frac{1}{C_{1}}+\frac{1}{C_{2}}\right) \Delta T d t \text { where } \Delta T=T_{1}-T_{2}
\]

Hence

\[
\Delta T=(\Delta T)_{0} e^{-t / \tau} \text { where } \frac{1}{\tau}=\frac{\kappa s}{l}\left(\frac{1}{C_{1}}+\frac{1}{C_{2}}\right)
\]

\(2.251 \dot{Q}=\kappa \frac{\partial T}{\partial x}=-A \sqrt{T} \frac{\partial T}{\partial x}\)

\[
\begin{gathered}
=-\frac{2}{3} A \frac{\partial T^{3 / 2}}{\partial x},(A=\text { constant }) \\
=\frac{2}{3} A \frac{\left(T_{1}^{3 / 2}-T_{2}^{3 / 2}\right)}{l}
\end{gathered}
\]

Thus

\[
T^{3 / 2}=\text { constant }-\frac{x}{l}\left(T_{1}^{3 / 2}-T_{2}^{3 / 2}\right)
\]

or using

\[
T=T_{1} \quad \text { at } x=0
\]

\[
\begin{gathered}
T^{3 / 2}=T_{1}^{3 / 2}+\frac{x}{l}\left(T_{2}^{3 / 2}-T_{1}^{3 / 2}\right) \text { or }\left(\frac{T}{T_{1}}\right)^{3 / 2}=1+\frac{x}{l}\left(\left(\frac{T_{2}}{T}\right)^{3 / 2}-1\right) \\
T=T_{1}\left[1+\frac{x}{l}\left\{\left(\frac{T_{2}}{T_{1}}\right)^{3 / 2}-1\right\}\right]^{2 / 3}
\end{gathered}
\]

\(2.252 \kappa=\frac{1}{3} \sqrt{\frac{8 R T}{\pi M}} \frac{1}{\sqrt{2} \pi d^{2} n} m n \frac{R \frac{i}{2}}{M}=\frac{R^{3 / 2} i T^{3 / 2}}{3 \pi^{3 / 2} d^{2} \sqrt{M} N_{A}}\)\\
Then from the previous problem

\[
q=\frac{2 i R^{3 / 2}\left(T_{2}^{3 / 2}-T_{1}^{3 / 2}\right)}{9 \pi^{3 / 2} d^{2} \sqrt{M} N_{A} l}, i=3 \text { here. }
\]

2.253 At this pressure and average temparature \(=27^{\circ} \mathrm{C}=300 \mathrm{~K}=T=\frac{\left(T_{1}+T_{2}\right)}{2}\)\\
\(\lambda=\frac{1}{\sqrt{2} \pi d^{2} p}=2330 \times 10^{-5} \mathrm{~m}=23.3 \mathrm{~mm} \gg 5.0 \mathrm{~mm}=l\)\\
The gas is ultrathin and we write \(\lambda=\frac{1}{2} l\) here

Then

\[
q=\kappa \frac{d T}{d x}=\kappa \frac{T_{2}-T_{1}}{l}
\]

where

\[
\kappa=\frac{1}{3}<v>\times \frac{1}{2} l \times \frac{M P}{R T} \times \frac{R}{\gamma-1} \times \frac{1}{M}=\frac{p<v>}{6 T(\gamma-1)} l
\]

and

\[
q=\frac{p<v>}{6 T(\gamma-1)}\left(T_{2}-T_{1}\right)
\]

where \(<v>=\sqrt{\frac{8 R T}{M \pi}}\). We have used \(T_{2}-T_{1} \ll \frac{T_{2}+T_{1}}{2}\) here.\\
2.254 In equilibrium \(2 \pi r \kappa \frac{d T}{d r}=-A=\) constant. So \(T=B-\frac{A}{2 \pi \kappa} \ln r\)

But \(T=T_{1}\) when \(r=R_{1}\) and \(T=T_{2}\). when \(r=R_{2}\).\\
From this we find \(T=T_{1}+\frac{T_{2}-T_{1}}{\ln \frac{R_{2}}{R_{1}}} \ln \frac{r}{r_{1}}\)\\
2.255 In equilibrium \(4 \pi r^{2} \kappa \frac{d T}{d r}=-A=\) constant

\[
T=B+\frac{A}{4 \pi \kappa} \frac{1}{r}
\]

Using \(T=T_{1}\) when \(r=R_{1}\) and \(T=T_{2}\) when \(r=R_{2}\),

\[
T=T_{1}+\frac{T_{2}-T_{1}}{\frac{1}{R_{2}}-\frac{1}{R_{1}}}\left(\frac{1}{r}-\frac{1}{R_{1}}\right)
\]

2.256 The heat flux vector is \(-\kappa\) grad \(T\) and its divergence equals \(w\). Thus

\[
\nabla^{2} T=-\frac{w}{\kappa}
\]

or

\[
\frac{1}{r} \frac{\partial}{\partial r}\left(r \frac{\partial T}{\partial r}\right)=-\frac{w}{\kappa} \text { in cylindrical coordinates. }
\]

or

\[
T=B+A \ln r-\frac{\omega}{2 x} r^{2}
\]

Since \(T\) is finite at \(r=0, A=0\). Also \(T=T_{0}\) at \(r=R\)\\
so

\[
B=T_{0}+\frac{w}{4 \kappa} R^{2}
\]

Thus

\[
T=T_{0}+\frac{w}{4 \kappa}\left(R^{2}-r^{2}\right)
\]

\(r\) here is the distance from the axis of wire (axial radius).\\
2.257 Here again

\[
\nabla^{2} T=-\frac{w}{\kappa}
\]

So in spherical polar coordinates,

\[
\frac{1}{r^{2}} \frac{\partial}{\partial r}\left(r^{2} \frac{\partial T}{\partial r}\right)=-\frac{w}{\kappa} \text { or } r^{2} \frac{\partial T}{\partial r}=-\frac{w}{3 \kappa} r^{3}+A
\]

or

\[
T=B-\frac{A}{r}-\frac{w}{6 \kappa} r^{2}
\]

Again

\[
A=0 \text { and } B=T_{0}+\frac{w}{6 \kappa} R^{2}
\]

so finally

\[
T=T_{0}+\frac{w}{6 \kappa}\left(R^{2}-r^{2}\right)
\]

\section*{PART THREE}
\section*{ELECTRODYNAMICS}
\subsection*{3.1 CONSTANT ELECTRIC FIELD IN VACUUM}
\(3.1 F_{d}\) (for electorns) \(=\frac{q^{2}}{4 \pi \varepsilon_{0} r^{2}}\) and \(F_{g r}=\frac{\gamma m^{2}}{r^{2}}\)

Thus

\[
\begin{aligned}
& \quad \frac{F_{e l}}{F_{g r}}(\text { for electrons })=\frac{q^{2}}{4 \pi \varepsilon_{0} \gamma m^{2}} \\
& =\frac{\left(1.602 \times 10^{-19} \mathrm{C}\right)^{2}}{\left(\frac{1}{9 \times 10^{9}}\right) \times 6.67 \times 10^{-11} \mathrm{~m}^{3} /\left(\mathrm{kg} \cdot \mathrm{~s}^{2}\right) \times\left(9.11 \times 10^{-31} \mathrm{~kg}\right)^{2}}=4 \times 10^{42} \\
& \frac{F_{e l}(\text { for proton })=\frac{q^{2}}{4 \pi \varepsilon_{0} \gamma m^{2}}}{\left(\frac{1}{9 \times 10^{9}}\right) \times 6.67 \times 10^{-11} \mathrm{~m}^{3} /\left(\mathrm{kg} \cdot \mathrm{~s}^{2}\right) \times\left(1.672 \times 10^{-27} \mathrm{~kg}\right)^{2}}=1 \times 10^{36} \\
& \quad \text { For } F_{e l}=F_{g r} \\
& =\sqrt{\frac{q^{2}}{4 \pi \varepsilon_{0} r^{2}}}=\frac{\gamma m^{2}}{r^{2}} \text { or } \frac{q}{m}=\sqrt{4 \pi \varepsilon_{0} \gamma} \\
& =\frac{\left(1.602 \times 10^{-19} \mathrm{C}\right)^{2}}{9 \times 10^{9}\left(\mathrm{~kg}-\mathrm{s}^{2}\right)} \\
& =0.86 \times 10^{-10} \mathrm{C} / \mathrm{kg}
\end{aligned}
\]

3.2 Total number of atoms in the sphere of mass \(1 \mathrm{gm}=\frac{1}{63.54} \times 6.023 \times 10^{23}\)

So the total nuclear charge \(\lambda=\frac{6.023 \times 10^{23}}{63.54} \times 1.6 \times 10^{-19} \times 29\)\\
Now the charge on the sphere \(=\) Total nuclear charge - Total electronic charge

\[
=\frac{6.023 \times 10^{23}}{63.54} \times 1.6 \times 10^{-19} \times \frac{29 \times 1}{100}=4.298 \times 10^{2} \mathrm{C}
\]

Hence force of interaction between these two spheres,

\[
F=\frac{1}{4 \pi \varepsilon_{0}} \cdot \frac{\left[4.398 \times 10^{2}\right]^{2}}{1^{2}} \mathrm{~N}=9 \times 10^{9} \times 10^{4} \times 19.348 \mathrm{~N}=1.74 \times 10^{15} \mathrm{~N}
\]

3.3 Let the balls be deviated by an angle \(\theta\), from the vertical when separtion between them equals \(x\).\\
Applying Newton's second law of motion for any one of the sphere, we get,


\begin{equation*}
T \cos \theta=m g \tag{1}
\end{equation*}


and


\begin{equation*}
T \sin \theta=F_{e} \tag{2}
\end{equation*}


From the Eqs. (1) and (2)


\begin{equation*}
\tan \theta=\frac{F_{e}}{m g} \tag{3}
\end{equation*}


But from the figure


\begin{equation*}
\tan \theta=\frac{x}{2 \sqrt{l^{2}-\left(\frac{x}{2}\right)^{2}}} \cong \frac{x}{2 l} \text { as } x \ll l \tag{4}
\end{equation*}


From Eqs. (3) and (4)\\
\includegraphics[max width=\textwidth, alt={}, center]{9f3c5a8e-ac5d-42f4-8eb6-ea2af623f2d8-275_438_546_615_862}

\[
F_{e}=\frac{m g x}{2 l} \text { or } \frac{q^{2}}{4 \pi \varepsilon_{0} x^{2}}=\frac{m g x}{2 l}
\]

Thus


\begin{equation*}
q^{2}=\frac{2 \pi \varepsilon_{0} m g x^{3}}{l} \tag{5}
\end{equation*}


Differentiating Eqn. (5) with respect to time

\[
2 q \frac{d q}{d t}=\frac{2 \pi \varepsilon_{0} m g}{l} 3 x^{2} \frac{d x}{d t}
\]

According to the problem \(\frac{d x}{d t}=v=a / \sqrt{x}\) (approach velocity is \(\frac{d x}{d t}\) )\\
so, \(\quad\left(\frac{2 \pi \varepsilon_{0} m g}{l} x^{3}\right)^{1 / 2} \frac{d q}{d t}=\frac{3 \pi \varepsilon_{0} m g}{l} x^{2} \frac{a}{\sqrt{x}}\)\\
Hence, \(\frac{d q}{d t}=\frac{3}{2} a \sqrt{\frac{2 \pi \varepsilon_{0} m g}{l}}\)\\
3.4 Let us choose coordinate axes as shown in the figure and fix three charges, \(q_{1}, q_{2}\) and \(q_{3}\) having position vectors \(\overrightarrow{r_{1}}, \overrightarrow{r_{2}}\) and \(\overrightarrow{r_{3}}\) respectively.\\
Now, for the equilibrium of \(q_{3}\)

\[
\frac{+q_{2} q_{3}\left(\overrightarrow{r_{2}}-\overrightarrow{r_{3}}\right)}{\left|\overrightarrow{r_{2}}-\overrightarrow{r_{3}}\right|^{3}}+\frac{\dot{q_{1}} q_{3}\left(\overrightarrow{r_{1}}-\overrightarrow{r_{3}}\right)}{\left|\overrightarrow{r_{1}}-\overrightarrow{r_{3}}\right|^{3}}=0
\]

\includegraphics[max width=\textwidth, alt={}, center]{9f3c5a8e-ac5d-42f4-8eb6-ea2af623f2d8-275_461_515_1614_813}\\
or,

\[
\frac{q_{2}}{\left|\overrightarrow{r_{2}}-\overrightarrow{r_{3}}\right|^{2}}=\frac{q_{1}}{\left|\overrightarrow{r_{1}}-\overrightarrow{r_{3}}\right|^{2}}
\]

because

\[
\frac{\overrightarrow{r_{2}}-\overrightarrow{r_{3}}}{\left|\overrightarrow{r_{2}}-\overrightarrow{r_{3}}\right|}=-\frac{\overrightarrow{r_{1}}-\overrightarrow{r_{3}}}{\left|\overrightarrow{r_{1}}-\overrightarrow{r_{3}}\right|}
\]

or,

\[
\sqrt{q_{2}}\left(\overrightarrow{r_{1}}-\overrightarrow{r_{3}}\right)=\sqrt{q_{1}}\left(\overrightarrow{r_{3}}-\overrightarrow{r_{2}}\right)
\]

or,

\[
\overrightarrow{r_{3}}=\frac{\sqrt{q_{2}} \overrightarrow{r_{1}}+\sqrt{q_{1}} \overrightarrow{r_{2}}}{\sqrt{q_{1}}+\sqrt{q_{2}}}
\]

Also for the equilibrium of \(\boldsymbol{q}_{1}\),

\[
\frac{q_{3}\left(\overrightarrow{r_{3}}-\overrightarrow{r_{1}}\right)}{\left|\overrightarrow{r_{3}}-\overrightarrow{r_{1}}\right|^{3}}+\frac{q_{2}\left(\overrightarrow{r_{2}}-\overrightarrow{r_{1}}\right)}{\left|\overrightarrow{r_{2}}-\overrightarrow{r_{1}}\right|^{3}}=0
\]

or,

\[
q_{3}=\frac{-q_{2}}{\left|\overrightarrow{r_{2}}-\overrightarrow{r_{1}}\right|^{2}}\left|\overrightarrow{r_{1}}-\overrightarrow{r_{3}}\right|^{2}
\]

Substituting the value of \(\overrightarrow{r_{3}}\), we get,

\[
q_{3}=\frac{-q_{1} q_{2}}{\left(\sqrt{q_{1}}+\sqrt{q_{2}}\right)^{2}}
\]

3.5 When the charge \(q_{0}\) is placed at the centre of the ring, the wire get stretched and the extra tension, produced in the wire, will balance the electric force due to the charge \(q_{0}\). Let the tension produced in the wire, after placing the charge \(\boldsymbol{q}_{0}\), be \(T\). From Newton's second law in projection form \(F_{n}=m w_{n}\).\\
\(T d \theta-\frac{1}{4 \pi \varepsilon_{0}} \frac{q_{0}}{r^{2}}\left(\frac{q}{2 \pi r} r d \theta\right)=(d m) 0\),\\
\includegraphics[max width=\textwidth, alt={}, center]{9f3c5a8e-ac5d-42f4-8eb6-ea2af623f2d8-276_419_391_1017_844}\\
or,

\[
T=\frac{q q_{0}}{8 \pi^{2} \varepsilon_{0} r^{2}}
\]

3.6 Sought field strength

\[
\vec{E}=\frac{1}{4 \pi \varepsilon_{0}} \frac{q}{\left|\overrightarrow{r-} \overrightarrow{r_{0}}\right|^{2}}
\]

\(=4.5 \mathrm{kV} / \mathrm{m}\) on putting the values.\\
3.7 Let us fix the coordinate system by taking the point of intersection of the diagonals as the origin and let \(\vec{k}\) be directed normally, emerging from the plane of figure. Hence the sought field strength :\\
\(\vec{E}=\frac{q}{4 \pi \varepsilon_{0}} \frac{\overrightarrow{l i+x \vec{k}}}{\left(l^{2}+x^{2}\right)^{3 / 2}}+\frac{-q}{4 \pi \varepsilon_{0}} \frac{l(-\vec{i})+x \vec{k}}{\left(l^{2}+x^{2}\right)^{3 / 2}}\)\\
\(+\frac{-q}{4 \pi \varepsilon_{0}} \frac{\overrightarrow{l j+x \vec{k}}}{\left(l^{2}+x^{2}\right)^{3 / 2}}+\frac{q}{4 \pi \varepsilon_{0}} \frac{l(-\vec{j})+x \vec{k}}{\left(l^{2}+x^{2}\right)^{3 / 2}}\)\\
\(=\frac{q}{4 \pi \varepsilon_{0}\left(l^{2}+x^{2}\right)^{3 / 2}}[2 \overrightarrow{l i}-2 l \vec{j}]\)\\
Thus \(E=\frac{q l}{\sqrt{2} \pi \varepsilon_{0}\left(l^{2}+x^{2}\right)^{3 / 2}}\)\\
\includegraphics[max width=\textwidth, alt={}, center]{9f3c5a8e-ac5d-42f4-8eb6-ea2af623f2d8-277_401_477_93_879}\\
3.8 From the symmetry of the problem the sought field.

\[
E=\int d E_{x}
\]

where the projection of field strength along \(x\)-axis due to an elemental charge is\\
\(d E_{x}=\frac{d q \cos \theta}{4 \pi \varepsilon_{0} R^{2}}=\frac{q R \cos \theta d \theta}{4 \pi^{2} \varepsilon_{0} R^{3}}\)\\
Hence

\[
E=\frac{q}{4 \pi^{2} \varepsilon_{0} R^{2}} \int_{\pi / 2}^{\pi / 2} \cos \theta d \theta \frac{q}{2 \pi^{2} \varepsilon_{0} R^{2}}
\]

\includegraphics[max width=\textwidth, alt={}, center]{9f3c5a8e-ac5d-42f4-8eb6-ea2af623f2d8-277_433_399_584_857}\\
3.9 From the symmetry of the condition, it is clear that, the field along the normal will be zero\\
i.e.

Now

\[
d E_{l}=\frac{d q}{4 \pi \varepsilon_{0}\left(R^{2}+l^{2}\right)} \cos \theta
\]

But

\[
d q=\frac{q}{2 \pi R} d x \text { and } \cos \theta=\frac{l}{\left(R^{2}+l^{2}\right)^{1 / 2}}
\]

Hence

\[
E=\int d E_{l}=\int_{0}^{2 \pi R} \frac{q l}{2 \pi R} \cdot \frac{d x}{4 \pi \varepsilon_{0}\left(R^{2}+l^{2}\right)^{3 / 2}}
\]

or \(E=\frac{1}{4 \pi \varepsilon_{0}} \frac{q l}{\left(l^{2}+R^{2}\right)^{3 / 2}}\)\\
and for \(l \gg R\), the ring behaves like a point charge, reducing the field to the value,

\[
E \approx \frac{1}{4 \pi \varepsilon_{0}} \frac{q}{l^{2}}
\]

\begin{center}
\includegraphics[max width=\textwidth, alt={}]{9f3c5a8e-ac5d-42f4-8eb6-ea2af623f2d8-277_458_598_1484_782}
\end{center}

For \(E_{\text {max }}\), we should have \(\frac{d E}{d l}=0\)\\
So, \(\left(l^{2}+R^{2}\right)^{3 / 2}-\frac{3}{2} l\left(l^{2}+R^{2}\right)^{1 / 2} 2 l=0\) or \(l^{2}+R^{2}-3 l^{2}=0\)\\
Thus \(l=\frac{R}{\sqrt{2}}\) and \(E_{\text {max }}=\frac{q}{6 \sqrt{3} \pi \varepsilon_{0} R^{2}}\)\\
3.10 The electric potential at a distance \(x\) from the given ring is given by,

\[
\varphi(x)=\frac{q}{4 \pi \varepsilon_{0} x}-\frac{q}{4 \pi \varepsilon_{0}\left(R^{2}+x^{2}\right)^{1 / 2}}
\]

Hence, the field strength along \(x\)-axis (which is the net field strength in our case),

\[
\begin{aligned}
E_{x} & =-\frac{d \varphi}{d x}=\frac{q}{4 \pi \varepsilon_{0}} \frac{1}{x^{2}}-\frac{q x}{4 \pi \varepsilon_{0}\left(R^{2}+x^{2}\right)^{3 / 2}} \\
& =\frac{\frac{q}{4 \pi \varepsilon_{0}} x^{3}\left[\left(1+\frac{R^{2}}{x^{2}}\right)^{3 / 2}-1\right]}{x^{2}\left(R^{2}+x^{2}\right)^{3 / 2}} \\
& =\frac{\frac{q}{4 \pi \varepsilon_{0}} x^{3}\left[1+\frac{3}{2} \frac{R^{2}}{x^{2}}+\frac{3}{8} \frac{R^{4}}{x^{4}}+\ldots\right]}{x^{2}\left(R^{2}+x^{2}\right)^{3 / 2}}
\end{aligned}
\]

Neglecting the higher power of \(R / x\), as \(x \gg R\).

\[
E=\frac{3 q R^{2}}{8 \pi \varepsilon_{0} x^{4}}
\]

Note : Instead of \(\varphi(x)\), we may write \(\boldsymbol{E}(\boldsymbol{x})\) directly using 3.9\\
3.11 From the solution of 3.9, the electric field strength due to ring at a point on its axis (say \(x\)-axis) at distance \(x\) from the centre of the ring is given by :

\[
E(x)=\frac{q x}{4 \pi \varepsilon_{0}\left(R^{2}+x^{2}\right)^{3 / 2}}
\]

And from symmetry \(\vec{E}\) at every point on the axis is directed along the \(x\)-axis (Fig.).\\
Let us consider an element ( \(d x\) ) on thread which carries the charge ( \(\lambda d x\) ). The electric force experienced by the element in the field of ring.\\
\(d F=(\lambda d x) E(x)=\frac{\lambda q x d x}{4 \pi \varepsilon_{0}\left(R^{2}+x^{2}\right)^{3 / 2}}\)\\
Thus the sought interaction\\
\(F=\int_{0}^{\infty} \frac{\lambda q x d x}{4 \pi \varepsilon_{0}\left(R^{2}+x^{2}\right)^{3 / 2}}\)\\
On integrating we get, \(F=\frac{\lambda q}{4 \pi \varepsilon_{0} R}\)\\
\includegraphics[max width=\textwidth, alt={}, center]{9f3c5a8e-ac5d-42f4-8eb6-ea2af623f2d8-278_323_529_1653_759}\\
3.12 (a) The given charge distribution is shown in Fig. The symmetry of this distribution implies that vector \(\vec{E}\) at the point \(O\) is directed to the right, and its magnitude is equal to the sum of the projection onto the direction of \(\vec{E}\) of vectors \(d \vec{E}\) from elementary charges \(d q\). The projection of vector \(d \vec{E}\) onto vector \(\vec{E}\) is

\[
d E \cos \varphi=\frac{1}{4 \pi \varepsilon_{0}} \frac{d q}{R^{2}} \cos \varphi,
\]

where \(d q=\lambda R d \varphi=\lambda_{0} R \cos \varphi d \varphi\).\\
Integrating (1) over \(\varphi\) between 0 and \(2 \pi\) we find the magnitude of the vector \(E\) :

\[
E=\frac{\lambda_{0}}{4 \pi \varepsilon_{0} R} \int_{0}^{2 \pi} \cos ^{2} \varphi d \varphi=\frac{\lambda_{0}}{4 \varepsilon_{0} R}
\]

It should be noted that this integral is evaluated in the most simple way if we take into account that \(\left\langle\cos ^{2} \varphi\right\rangle=1 / 2\). Then\\
\includegraphics[max width=\textwidth, alt={}, center]{9f3c5a8e-ac5d-42f4-8eb6-ea2af623f2d8-279_441_577_380_815}

\[
\int_{0} \cos ^{2} \varphi d \varphi=<\cos ^{2} \varphi>2 \pi=\pi
\]

(b) Take an element \(S\) at an azimuthal angle \(\varphi\) from the \(x\)-axis, the element subtending an angle \(d \varphi\) at the centre.\\
The elementary field at \(P\) due to the element is\\
\(\frac{\lambda_{0} \cos \varphi d \varphi R}{4 \pi \varepsilon_{0}\left(x^{2}+R^{2}\right)}\) along \(S P\) with components

\[
\frac{\lambda_{0} \cos \varphi d \varphi R}{4 \pi \varepsilon_{0}\left(x^{2}+R^{2}\right)} \times\{\cos \theta \text { along } O P, \sin \theta \text { along } O S\}
\]

where

\[
\cos \theta=\frac{x}{\left(x^{2}+R^{2}\right)^{1 / 2}}
\]

The component along \(O P\) vanishes on integration as \(\int_{0}^{2 \pi} \cos \varphi d \varphi=0\)\\
The component alon \(O S\) can be broken into the parts along \(O X\) and \(O Y\) with

\[
\frac{\lambda_{0} R^{2} \cos \varphi d \varphi}{4 \pi \varepsilon_{0}\left(x^{2}+R^{2}\right)^{3 / 2}} \times\{\cos \varphi \text { along } O X, \sin \varphi \text { along } O Y\}
\]

On integration, the part along \(O Y\) vanishes. Finally

For \(x \gg R\)

\[
E=E_{x}=\frac{\lambda_{0} R^{2}}{4 \varepsilon_{0}\left(x^{2}+R^{2}\right)^{3 / 2}}
\]

\[
E_{x}=\frac{p}{4 \pi \varepsilon_{0} x^{3}} \text { where } p=\lambda_{0} \pi R^{2}
\]

\includegraphics[max width=\textwidth, alt={}, center]{9f3c5a8e-ac5d-42f4-8eb6-ea2af623f2d8-279_334_638_1715_749}\\
3.13 (a) It is clear from symmetry considerations that vector \(\vec{E}\) must be directed as shown in the figure. This shows the way of solving this problem : we must find the component \(d E_{r}\) of the field created by the element \(d l\) of the rod, having the charge \(d q\) and then integrate the result over all the elements of the rod. In this case

\[
d E_{r}=d E \cos \alpha=\frac{1}{4 \pi \varepsilon_{0}} \frac{\lambda d l}{r_{0}^{2}} \cos \alpha,
\]

where \(\lambda=\frac{q}{2 a}\) is the linear charge density. Let us reduce this equation of the form convenient for integration. Figure shows that \(d l \cos \alpha=r_{0} d \alpha\) and \(r_{0}=\frac{r}{\cos \alpha}\);\\
Consequently,\\
\(d E_{r}=\frac{1}{4 \pi \varepsilon_{0}} \frac{\lambda r_{0} d \alpha}{r_{0}^{2}}=\frac{\lambda}{4 \pi \varepsilon_{0} r} \cos \alpha d \alpha\)\\
This expression can be easily integrated :\\
\(E=\frac{\lambda}{4 \pi \varepsilon_{0} r} 2 \int_{0}^{\alpha_{0}} \cos \alpha d \alpha=\frac{\lambda}{4 \pi \varepsilon_{0} r} 2 \sin \alpha_{0}\)\\
where \(\alpha_{0}\) is the maximum value of the angle \(\alpha\),\\
\(\sin \alpha_{0}=a / \sqrt{a^{2}+r^{2}}\)\\
\includegraphics[max width=\textwidth, alt={}, center]{9f3c5a8e-ac5d-42f4-8eb6-ea2af623f2d8-280_415_530_608_840}

Thus, \(E=\frac{q / 2 a}{4 \pi \varepsilon_{0} r} 2 \frac{a}{\sqrt{a^{2}+r^{2}}}=\frac{q}{4 \pi \varepsilon_{0} r \sqrt{a^{2}+r^{2}}}\)\\
Note that in this case also \(E \propto \frac{q}{4 \pi \varepsilon_{0} r^{2}}\) for \(r \gg a\) as of the field of a point charge.\\
(b) Let, us consider the element of length \(d l\) at a distance \(l\) from the centre of the rod, as shown in the figure.\\
Then field at \(P\), due to this element.

\[
d E=\frac{\lambda d l}{4 \pi \varepsilon_{0}(r-l)^{2}}
\]

\includegraphics[max width=\textwidth, alt={}, center]{9f3c5a8e-ac5d-42f4-8eb6-ea2af623f2d8-280_192_649_1323_719}\\
if the element lies on the side, shown in the\\
diagram, and \(d E=\frac{\lambda d l}{4 \pi \varepsilon_{0}(r+l)^{2}}\), if it lies on\\
other side.

Hence

\[
E=\int d E=\int_{0}^{a} \frac{\lambda d l}{4 \pi \varepsilon_{0}(r-l)^{2}}+\int_{0}^{a} \frac{\lambda d l}{4 \pi \varepsilon_{0}(r+l)^{2}}
\]

On integrating and putting \(\lambda=\frac{q}{2 a}\), we get, \(E=\frac{q}{4 \pi \varepsilon_{0}} \frac{1}{\left(r^{2}-a^{2}\right)}\)\\
For

\[
r \gg a, \quad E \sim \frac{q}{4 \pi \varepsilon_{0} r^{2}}
\]

3.14 The problem is reduced to finding \(E_{x}\) and \(E_{y}\) viz. the projections of \(\vec{E}\) in Fig, where it is assumed that \(\lambda>0\).\\
Let us start with \(E_{x}\). The contribution to \(E_{x}\) from the charge element of the segment \(d x\) is


\begin{equation*}
d E_{x}=\frac{1}{4 \pi \varepsilon_{0}} \frac{\lambda d x}{r^{2}} \sin \alpha \tag{1}
\end{equation*}


Let us reduce this expression to the form convenient for integration. In our case, \(d x=r d \alpha / \cos \alpha, r=y / \cos \alpha\). Then

\[
d E_{x}=\frac{\lambda}{4 \pi \varepsilon_{0} y} \sin \alpha d \alpha
\]

Integrating this expression over \(\alpha\) between 0 and \(\pi / 2\), we find

\[
E_{x}=\lambda / 4 \pi \varepsilon_{0} y .
\]

In order to find the projection \(E_{y}\) it is sufficient to recall that \(d E_{y}\) differs from \(d E_{x}\) in that \(\sin \alpha\) in (1) is simply replaced by \(\cos \alpha\).\\
This gives\\
\includegraphics[max width=\textwidth, alt={}, center]{9f3c5a8e-ac5d-42f4-8eb6-ea2af623f2d8-281_372_370_495_890}\\
\(d E_{y}=(\lambda \cos \alpha d \alpha) / 4 \pi \varepsilon_{0} y\) and \(E_{y}=\lambda / 4 \pi \varepsilon_{0} y\).\\
We have obtained an interesting result :\\
\(E_{x}=E_{y}\) independently of \(y\),\\
i.e. \(\vec{E}\) is oriented at the angle of \(45^{\circ}\) to the rod. The modulus of \(\vec{E}\) is

\[
E=\sqrt{E_{x}^{2}+E_{y}^{2}}=\lambda \sqrt{2} / 4 \pi \varepsilon_{0} y .
\]

3.15 (a) Using the solution of 3.14, the net electric field strength at the point \(O\) due to straight parts of the thread equals zero. For the curved part (arc) let us derive a general expression i.e. let us calculate the field strength at the centre of arc of radius \(R\) and linear charge density \(\lambda\) and which subtends angle \(\theta_{0}\) at the centre.\\
From the symmetry the sought field strength will be directed along the bisector of the angle \(\theta_{0}\) and is given by

\[
+\theta_{0} / 2
\]

\(E=\int_{-\theta_{0} / 2} \frac{\lambda(R d \theta)}{4 \pi \varepsilon_{0} R^{2}} \cos \theta=\frac{\lambda}{2 \pi \varepsilon_{0} R} \sin \frac{\theta_{0}}{2}\)\\
In our problem \(\theta_{0}=\pi / 2\), thus the field strength due to the turned part at the point \(E_{0}=\frac{\sqrt{2} \lambda}{4 \pi \varepsilon_{0} R}\) which is also the sought result.\\
\includegraphics[max width=\textwidth, alt={}, center]{9f3c5a8e-ac5d-42f4-8eb6-ea2af623f2d8-281_476_279_1353_965}\\
(b) Using the solution of 3.14 (a), net field strength at \(O\) due to stright parts equals \(\sqrt{2}\left(\frac{\sqrt{2} \lambda}{4 \pi \varepsilon_{0} R}\right)=\frac{\lambda}{2 \pi \varepsilon_{0} R}\) and is directed vertically down. Now using the solution of 3.15\\
(a), field strength due to the given curved part (semi-circle) at the point \(O\) becomes \(\frac{\lambda}{2 \pi \varepsilon_{0} R}\) and is directed vertically upward. Hence the sought net field strengh becomes zero.\\
3.16 Given charge distribution on the surface \(\underset{\rightarrow}{g}=\vec{a} \cdot \vec{r}\) is shown in the figure. Symmetry of this distribution implies that the sought \(\vec{E}\) at the centre \(O\) of the sphere is opposite to \(\vec{a} d q=\sigma(2 \pi r \sin \theta) r d \theta=(\vec{a} \cdot \vec{r}) 2 \pi r^{2} \sin \theta d \theta=2 \pi a r^{3} \sin \theta \cos \theta d \theta\)\\
Again from symmetry, field strength due to any ring element \(d E\) is also opposite to \(\vec{a}\) i.e. \(d \vec{E} \uparrow \downarrow \vec{a}\). Hence\\
\(d \vec{E}=\frac{d q r \cos \theta}{4 \pi \varepsilon_{0}\left(r^{2} \sin ^{2} \theta+r^{2} \cos ^{2} \theta\right)^{3 / 2}} \frac{-\vec{a}}{a}\) (Using the result of 3.9)\\
\(=\frac{\left(2 \pi a r^{3} \sin \theta \cos \theta d \theta\right) r \cos \theta}{4 \pi \varepsilon_{0} r^{3}} \frac{(-\vec{a})}{a}\)\\
\(=\frac{-\overrightarrow{a r}}{2 \varepsilon_{0}} \sin \theta \cos ^{2} d \theta\)\\
Thus \(\vec{E}=\int d \vec{E}=\frac{(-\vec{a}) r}{2 \varepsilon_{0}} \int_{0}^{\pi} \sin \theta \cos ^{2} \theta d \theta\)\\
\includegraphics[max width=\textwidth, alt={}, center]{9f3c5a8e-ac5d-42f4-8eb6-ea2af623f2d8-282_313_573_636_807}

Integrating, we get \(\vec{E}=-\frac{\overrightarrow{a r}}{2 \varepsilon_{0}} \frac{2}{3}=-\frac{\overrightarrow{a r}}{3 \varepsilon_{0}}\)\\
3.17 We start from two charged spherical balls each of radius \(R\) with equal and opposite charge densities \(+\rho\) and \(-\rho\). The centre of the balls are at \(+\frac{\vec{a}}{2}\) and \(-\frac{\vec{a}}{2}\) respectively so the equation of their surfaces are \(\left|\vec{r}-\frac{\vec{a}}{2}\right|=R\) or \(r-\frac{a}{2} \cos \theta \propto R\) and \(r+\frac{a}{2} \cos \theta \propto R\), considering \(a\) to be small. The distance between the two surfaces in the radial direction at angle \(\theta\) is \(|a \cos \theta|\) and does not depend on the azimuthal angle. It is seen from the diagram that the surface of the sphere has in effect a surface density \(\sigma=\sigma_{0} \cos \theta\) when

\[
\sigma_{0}=\rho a .
\]

Inside any uniformly charged spherical ball, the field is radial and has the magnitude given by Gauss's theorm\\
or

\[
\begin{aligned}
4 \pi r^{2} E & =\frac{4 \pi}{3} r^{3} \rho / \varepsilon_{0} \\
E & =\frac{\rho r}{3 \varepsilon_{0}}
\end{aligned}
\]

In vector notation, using the fact the \(V\) must be measured from the centre of the ball, we get, for the present case

\[
\vec{E}=\frac{\rho}{3 \varepsilon_{0}}\left(\vec{r}-\frac{\vec{a}}{2}\right)-\frac{\rho}{3 \varepsilon_{0}}\left(\vec{r}+\frac{\vec{a}}{2}\right)
\]

\begin{center}
\includegraphics[max width=\textwidth, alt={}]{9f3c5a8e-ac5d-42f4-8eb6-ea2af623f2d8-282_476_360_1615_888}
\end{center}

\[
=-\rho \vec{a} / 3 \varepsilon_{0}=\frac{\sigma_{0}}{3 \varepsilon_{0}} \vec{k}
\]

When \(\vec{k}\) is the unit vector along the polar axis from which \(\theta\) is measured.\\
3.18 Let us consider an elemental spherical shell of thickness \(d r\). Thus surface charge density of the shell \(\sigma=\rho d r=(\vec{a} \cdot \vec{r}) d r\).\\
Thus using the solution of 3.16 , field strength due to this sperical shell

\[
d \vec{E}=-\frac{\overrightarrow{a r}}{3 \varepsilon_{0}} d r
\]

Hence the sought field strength

\[
\vec{E}=-\frac{\vec{a}}{3 \varepsilon_{0}} \int_{0}^{R} r d r=-\frac{\vec{a} R^{2}}{6 \varepsilon_{0}} .
\]

3.19 From the solution of 3.14 field strength at a perpendicular distance \(r<R\) from its left end

\[
\vec{E}(r)=\frac{\lambda}{4 \pi \varepsilon_{0} r}(-\vec{i})+\frac{\lambda}{4 \pi \varepsilon_{0} r}\left(\hat{e}_{r}\right)
\]

Here \(\hat{e}_{r}\) is a unit vector along radial direction.\\
Let us consider an elemental surface, \(d S=d y d z=d z(r d \theta)\) a figure. Thus flux of \(\vec{E}(r)\) over the element \(d \vec{S}\) is given by

\[
\begin{aligned}
d \Phi & =\vec{E} \cdot d \vec{S}=\left[\frac{\lambda}{4 \pi \varepsilon_{0} r}(\overrightarrow{-i})+\frac{\lambda}{4 \pi \varepsilon_{0} r}\left(\hat{e}_{r}\right)\right] \cdot d r(r d \theta) \vec{i} \\
& =-\frac{\lambda}{4 \pi \varepsilon_{0}} d r d \theta\left(\text { as } \vec{e}_{r} \perp \vec{i}\right)
\end{aligned}
\]

The sought flux, \(\Phi=-\frac{\lambda}{4 \pi \varepsilon_{0}} \int_{0}^{R} d r \int_{0}^{2 \pi} d \theta=-\frac{\lambda R}{2 \varepsilon_{0}}\).\\
If we have taken \(d \vec{S} \uparrow \uparrow(-\vec{i})\), then \(\Phi\) were \(\frac{\lambda R}{2 \varepsilon_{0}}\)\\
\includegraphics[max width=\textwidth, alt={}, center]{9f3c5a8e-ac5d-42f4-8eb6-ea2af623f2d8-283_509_424_1088_970}

Hence

\[
|\Phi|=\frac{\lambda R}{2 \varepsilon_{0}}
\]

3.20 Let us consider an elemental surface area as shown in the figure. Then flux of the vector \(\vec{E}\) through the elemental area,

\[
\begin{array}{r}
d \Phi=\vec{E} \cdot d \vec{S}=E d S=2 E_{0} \cos \varphi d S(\text { as } \vec{E} \uparrow \uparrow d \vec{S}) \\
=\frac{2 q}{4 \pi \varepsilon_{0}\left(l^{2}+r^{2}\right)} \frac{l}{\left(l^{2}+r^{2}\right)^{1 / 2}}(r d \theta) d r=\frac{2 q l r d r d \theta}{4 \pi \varepsilon_{0}\left(r^{2}+l^{2}\right)^{3 / 2}}
\end{array}
\]

where \(E_{0}=\frac{q}{4 \pi \varepsilon_{0}\left(l^{2}+r^{2}\right)}\) is magnitude of field strength due to any point charge at the point of location of considered elemental area.

\[
\begin{aligned}
& \text { Thus } \Phi=\frac{2 q l}{4 \pi \varepsilon_{0}} \int_{0}^{R} \frac{r d r}{\left(r^{2}+l^{2}\right)^{3 / 2}} \int_{0}^{2 \pi} d \theta \\
& =\frac{2 q l \times 2 \pi}{4 \pi \varepsilon_{0}} \int_{0}^{R} \frac{r d r}{\left(r^{2}+l^{2}\right)^{3 / 2}}=\frac{q}{\varepsilon_{0}}\left[1-\frac{l}{\sqrt{l^{2}+R^{2}}}\right]
\end{aligned}
\]

\begin{center}
\includegraphics[max width=\textwidth, alt={}]{9f3c5a8e-ac5d-42f4-8eb6-ea2af623f2d8-284_529_568_70_832}
\end{center}

It can also be solved by considering a ring element or by using solid angle.\\
3.21 Let us consider a ring element of radius \(x\) and thickness \(d x\), as shown in the figure. Now, flux over the considered element,

\[
d \Phi=\vec{E} \cdot d \vec{S}=E_{r} d S \cos \theta
\]

But \(E_{r}=\frac{\rho r}{3 \varepsilon_{0}}\) from Gauss's theorem, and \(d S=2 \pi x d x, \cos \theta=\frac{r_{0}}{r}\)

Thus \(\quad d \Phi=\frac{\rho r}{3 \varepsilon_{0}} 2 \pi x d x \frac{r_{0}}{r}=\frac{\rho r_{0}}{3 \varepsilon_{0}} 2 \pi x d x\)\\
Hence sought flux

\[
\Phi=\frac{2 \pi \rho r_{0}}{3 \varepsilon_{0}} \int_{0}^{\sqrt{R^{2}-r_{0}^{2}}} x d x=\frac{2 \pi \rho r_{0}}{3 \varepsilon_{0}} \frac{\left(R^{2}-r_{0}^{2}\right)}{2}=\frac{\pi \rho r_{0}}{3 \varepsilon_{0}}\left(R^{2}-r_{0}^{2}\right)
\]

\includegraphics[max width=\textwidth, alt={}, center]{9f3c5a8e-ac5d-42f4-8eb6-ea2af623f2d8-284_468_582_766_809}\\
3.22 The field at \(P\) due to the threads at \(A\) and \(B\) are both of magnitude \(\frac{\lambda}{2 \pi \varepsilon_{0}\left(x^{2}+l^{2} / 4\right)^{1 / 2}}\) and directed along \(A P\) and \(B P\). The resultant is along \(O P\) with

\[
\begin{aligned}
E & =\frac{2 \lambda \cos \theta}{2 \pi \varepsilon_{0}\left(\pi^{2}+\pi^{1 / 2}\right)^{1 / 2}}=\frac{\lambda x}{\pi \varepsilon_{0}\left(x^{2}+l^{2 / 4}\right)} \\
& =\frac{\lambda}{\pi \varepsilon_{0}\left[x+\frac{l^{2}}{4 x}-2 \cdot \frac{l}{2 \sqrt{x}} \cdot \sqrt{x}+l\right]} \\
& =\frac{\lambda}{\pi \varepsilon_{0}\left[\left(\sqrt{x}-\frac{l}{2 \sqrt{x}}\right)^{2}+l\right]}
\end{aligned}
\]

\begin{center}
\includegraphics[max width=\textwidth, alt={}]{9f3c5a8e-ac5d-42f4-8eb6-ea2af623f2d8-284_438_591_1497_800}
\end{center}

This is maximum when \(x=l / 2\) and then \(E=E_{\max }=\frac{\lambda}{\pi \varepsilon_{0} l}\)\\
3.23 Take a section of the cylinder perpendicular to its axis through the point where the electric field is to be calculated. (All points on the axis are equivalent.) Consider an element \(S\) with azimuthal angle \(\varphi\). The length of the element is \(R, d \varphi, R\) being the radius of cross section of the cylinder. The element itself is a section of an infinite strip. The electric field at \(O\) due to this strip is

\[
\frac{\sigma_{0} \cos \varphi(R d \varphi)}{2 \pi \varepsilon_{0} R} \text { along } S O
\]

This can be resolved into\\
\(\frac{\sigma_{0} \cos \varphi d \varphi}{2 \pi \varepsilon_{0}}\left\{\begin{array}{c}\cos \varphi \text { along } O X \text { towards } O \\ \sin \varphi \text { along } Y O\end{array}\right.\)\\
\includegraphics[max width=\textwidth, alt={}, center]{9f3c5a8e-ac5d-42f4-8eb6-ea2af623f2d8-285_419_680_258_691}

On integration the component along \(Y O\) vanishes. What remains is

\[
\int_{0}^{2 \pi} \frac{\sigma_{0} \cos ^{2} \varphi d \varphi}{2 \pi \varepsilon_{0}}=\frac{\sigma_{0}}{2 \varepsilon_{0}} \text { along } X O \text { i.e. along the direction } \varphi=\pi
\]

3.24 Since the field is axisymmetric (as the field of a uniformly charged filament), we conclude that the flux through the sphere of radius \(R\) is equal to the flux through the lateral surface of a cylinder having the same radius and the height \(2 R\), as arranged in the figure.\\
Now, \(\quad \Phi=\oint \vec{E} \cdot d \vec{S}=E_{r} S\)

But

\[
E_{r}=\frac{a}{R}
\]

\begin{center}
\includegraphics[max width=\textwidth, alt={}]{9f3c5a8e-ac5d-42f4-8eb6-ea2af623f2d8-285_449_286_868_944}
\end{center}

Thus \(\Phi=\frac{a}{R} S=\frac{a}{R} 2 \pi R \cdot 2 R=4 \pi a R\)\\
3.25 (a) Let us consider a sphere of radius \(r<R\) then charge, inclosed by the considered sphere,


\begin{equation*}
q_{\text {inclosed }}=\int_{0}^{r} 4 \pi r^{2} d r \rho=\int_{0}^{r} 4 \pi r^{2} \rho_{0}\left(1-\frac{r}{R}\right) d r \tag{1}
\end{equation*}


Now, applying Gauss' theorem,\\
\(E_{r} 4 \pi r^{2}=\frac{q_{\text {inclosed }}}{\varepsilon_{0}}\), (where \(E_{r}\) is the projection of electric field along the radial line.)\\
or,

\[
\begin{gathered}
=\frac{\rho_{0}}{\varepsilon_{0}} \int_{0}^{r} 4 \pi r^{2}\left(1-\frac{r}{R}\right) d r \\
E_{r}=\frac{\rho_{0}}{3 \varepsilon_{0}}\left[r^{2}-\frac{3 r^{2}}{4 R}\right]
\end{gathered}
\]

And for a point. outside the sphere \(r>R\).\\
\(q_{\text {inclosed }}=\int_{0}^{R} 4 \pi r^{2} d r \rho_{0}\left(1-\frac{r}{R}\right)\) (as there is no charge outside the ball)\\
Again from Gauss' theorem,

\[
\begin{gathered}
E_{r} 4 \pi r^{2}=\int_{0}^{R} \frac{4 \pi r^{2} d r \rho_{0}\left(1-\frac{r}{R}\right)}{\varepsilon_{0}} \\
E_{r}=\frac{\rho_{0}}{r^{2} \varepsilon_{0}}\left[\frac{R^{3}}{3}-\frac{R^{4}}{4 R}\right]=\frac{\rho_{0} R^{3}}{12 r^{2} \varepsilon_{0}}
\end{gathered}
\]

(b) As ⋅ magnitude of electric field decreases with increasing \(r\) for \(r>R\), field will be maximum for \(r<R\). Now, for \(E_{r}\) to be maximum,

Hence

\[
\begin{gathered}
\frac{d}{d r}\left(r-\frac{3 r^{2}}{4 R}\right)=0 \quad \text { or } \quad 1-\frac{3 r}{2 R}=0 \quad \text { or } r=r_{m}=\frac{2 R}{3} \\
E_{\max }=\frac{\rho_{0} R}{9 \varepsilon_{0}}
\end{gathered}
\]

3.26 Let the charge carried by the sphere be \(q\), then using Gauss' theorem for a spherical surface having radius \(r>R\), we can write.\\
\(E 4 \pi r^{2}=\frac{q_{\text {inclosed }}}{\varepsilon_{0}}=\frac{q}{\varepsilon_{0}}+\frac{1}{\varepsilon_{0}} \int_{R}^{r} \frac{\alpha}{r} 4 \pi r^{2} d r\)\\
On integrating we get,\\
\(E 4 \pi r^{2}=\frac{\left(q-2 \pi \alpha R^{2}\right)}{\varepsilon_{0}}+\frac{4 \pi \alpha r^{2}}{2 \varepsilon_{0}}\)\\
The intensity \(E\) does not depend on \(r\) when the experession in the parentheses is equal to zero. Hence\\
\includegraphics[max width=\textwidth, alt={}, center]{9f3c5a8e-ac5d-42f4-8eb6-ea2af623f2d8-286_373_345_1138_907}

\[
q=2 \pi \alpha R^{2} \text { and } E=\frac{\alpha}{2 \varepsilon_{0}}
\]

3.27 Let us consider a spherical layer of radius \(r\) and thickness \(d r\), having its centre coinciding with the centre of the system. Then using Gauss' theorem for this surface,

\[
\begin{aligned}
& E_{r} 4 \pi r^{2}=\frac{q_{\text {inclosed }}}{\varepsilon_{0}}=\int_{0}^{r} \frac{\rho d V}{\varepsilon_{0}} \\
& \quad=\frac{1}{\varepsilon_{0}} \int_{0}^{r} \rho_{0} e^{-\alpha r^{3}} 4 \pi r^{2} d r
\end{aligned}
\]

After integration

\[
E_{r} 4 \pi r^{2}=\frac{\rho 4 \pi}{3 \varepsilon_{0} \alpha}\left[1-e^{-\alpha r^{3}}\right]
\]

or, \(E_{r}=\frac{\rho_{0}}{3 \varepsilon \alpha r^{2}}\left[1-e^{-\alpha r^{3}}\right]\)\\
Now when \(\alpha r^{3} \ll 1, E_{r} \approx \frac{\rho_{0} r}{3 \varepsilon_{0}}\)\\
And when \(\alpha r^{3} \gg 1, E_{r} \approx \frac{\rho_{0}}{3 \varepsilon_{0} \alpha r^{2}}\)\\
\includegraphics[max width=\textwidth, alt={}, center]{9f3c5a8e-ac5d-42f4-8eb6-ea2af623f2d8-287_427_375_103_848}\\
3.28 Using Gauss theorem we can easily show that the electric field strength within a uniformly charged sphere is \(\vec{E}=\left(\frac{\rho}{3 \varepsilon_{0}}\right) \vec{r}\)\\
The cavity, in our problem, may be considered as the superposition of two balls, one with the charge density \(\rho\) and the other with - \(\rho\).\\
Let \(P\) be a point inside the cavity such that its position vector with respect to the centre of cavity be \(\overrightarrow{r_{-}}\)and with respect to the centre of the ball \(\overrightarrow{r_{+}}\). Then from the principle of superposition, field inside the cavity, at an arbitrary point \(P\),\\
\includegraphics[max width=\textwidth, alt={}, center]{9f3c5a8e-ac5d-42f4-8eb6-ea2af623f2d8-287_412_401_756_864}

\[
\begin{gathered}
\vec{E}=\vec{E}_{+}+\vec{E}_{-} \\
=\frac{\rho}{3 \varepsilon_{0}}\left(\vec{r}_{+}-\vec{r}_{-}\right)=\frac{\rho}{3 \varepsilon_{0}} \vec{a}
\end{gathered}
\]

Note : Obtained expression for \(\vec{E}\) shows that it is valid regardless of the ratio between the radii of the sphere and the distance between their centres.\\
3.29 Let us consider a cylinderical Gaussian surface of radius \(r\) and height \(h\) inside an infinitely long charged cylinder with charge density \(\rho\). Now from Gauss theorem :\\
\(E_{r} 2 \pi r h=\frac{q_{\text {inclosed }}}{\varepsilon_{0}}\)\\
(where \(E_{r}\) is the field inside the cylinder at a distance \(r\) from its axis.)\\
or, \(E_{r} 2 \pi r h=\frac{\rho \pi r^{2} h}{\varepsilon_{0}}\) or \(E_{r}=\frac{\rho r}{2 \varepsilon_{0}}\)\\
Now, using the method of 3.28 field at a point \(P\), inside the cavity, is\\
\(\vec{E}=\vec{E}_{+}+\vec{E}_{-}=\frac{\rho}{2 \varepsilon_{0}}\left(\vec{r}_{+}-\vec{r}_{-}\right)=\frac{\rho}{2 \varepsilon_{0}} \vec{a}\)\\
\includegraphics[max width=\textwidth, alt={}, center]{9f3c5a8e-ac5d-42f4-8eb6-ea2af623f2d8-287_464_288_1580_987}\\
3.30 The arrangement of the rings are as shown in the figure. Now, potential at the point 1 , \(\varphi_{1}=\) potential at 1 due to the ring \(1+\) potential at 1 due to the ring 2 .

\[
=\frac{q}{4 \pi \varepsilon_{0} R}+\frac{-q}{4 \pi \varepsilon_{0}\left(R^{2}+a^{2}\right)^{1 / 2}}
\]

Similarly, the potential at point 2 ,

\[
\varphi_{2}=\frac{-q}{4 \pi \varepsilon_{0} R}+\frac{q}{4 \pi \varepsilon_{0}\left(R^{2}+a^{2}\right)^{1 / 2}}
\]

Hence, the sought potential difference,

\[
\begin{aligned}
\varphi_{1}-\varphi_{2}=\Delta \varphi & =2\left(\frac{q}{4 \pi \varepsilon_{0} R}+\frac{-q}{4 \pi \varepsilon_{0}\left(R^{2}+a^{2}\right)^{1 / 2}}\right) \\
& =\frac{q}{2 \pi \varepsilon_{0} R}\left(1-\frac{1}{\sqrt{1+(a / R)^{2}}}\right)
\end{aligned}
\]

\includegraphics[max width=\textwidth, alt={}, center]{9f3c5a8e-ac5d-42f4-8eb6-ea2af623f2d8-288_277_413_295_853}\\
3.31 We know from Gauss theorem that the electric field due to an infinietly long straight wire, at a perpendicular distance \(r\) from it equals, \(E_{r}=\frac{\lambda}{2 \pi \varepsilon_{0} r}\). So, the work done is

\[
\int_{1}^{2} E_{r} d r=\int_{x}^{\eta x} \frac{\lambda}{2 \pi \varepsilon_{0} r} d r
\]

(where \(x\) is perpendicular distance from the thread by which point 1 is removed from it.)\\
Hence

\[
\Delta \varphi_{12}=\frac{\lambda}{2 \pi \varepsilon_{0}} \ln \eta
\]

3.32 Let us consider a ring element as shown in the figure. Then the charge, carried by the element, \(d q=(2 \pi R \sin \theta) R d \theta \sigma\),\\
Hence, the potential due to the considered element at the centre of the hemisphere,

\[
d \varphi=\frac{1}{4 \pi \varepsilon_{0}} \frac{d q}{R}=\frac{2 \pi \sigma R \sin \theta d \theta}{4 \pi \varepsilon_{0}}=\frac{\sigma R}{2 \varepsilon_{0}} \sin \theta d \theta
\]

So potential due to the whole hemisphere\\
\(\varphi=\frac{R \sigma}{2 \varepsilon_{0}} \int_{0}^{\pi / 2} \sin \theta d \theta=\frac{\sigma R}{2 \varepsilon_{0}}\)\\
Now from the symmetry of the problem, net electric field of the hemisphere is directed towards the negative \(y\)-axis. We have

\[
d E_{y}=\frac{1}{4 \pi \varepsilon_{0}} \frac{d q \cos \theta}{R^{2}}=\frac{\sigma}{2 \varepsilon_{0}} \sin \theta \cos \theta d \theta
\]

\begin{center}
\includegraphics[max width=\textwidth, alt={}]{9f3c5a8e-ac5d-42f4-8eb6-ea2af623f2d8-288_493_632_1365_772}
\end{center}

Thus \(E=E_{y}{ }^{\prime}=\frac{\sigma}{2 \varepsilon_{0}} \int_{0}^{\pi / 2} \sin \theta \cos \theta d \theta=\frac{\sigma}{4 \varepsilon_{0}} \int_{0}^{\pi / 2} \sin 2 \theta d \theta=\frac{\sigma}{4 \varepsilon_{0}}\), along \(Y O\)\\
3.33 Let us consider an elementary ring of thickness \(d y\) and radius \(y\) as shown in the figure. Then potential at a point \(P\), at distance \(l\) from the centre of the disc, is

\[
d \varphi=\frac{\sigma 2 \pi y d y}{4 \pi \varepsilon_{0}\left(y^{2}+l^{2}\right)^{1 / 2}}
\]

Hence potential due to the whole disc,\\
\includegraphics[max width=\textwidth, alt={}, center]{9f3c5a8e-ac5d-42f4-8eb6-ea2af623f2d8-289_384_596_71_772}

\[
\varphi=\int_{0}^{R} \frac{\sigma 2 \pi y d y}{4 \pi \varepsilon_{0}\left(y^{2}+l^{2}\right)^{1 / 2}}=\frac{\sigma l}{2 \varepsilon_{0}}\left(\sqrt{1+(R / l)^{2}}-1\right)
\]

From symmetry

\[
\begin{aligned}
E=E_{l} & =-\frac{d \varphi}{d l} \\
=-\frac{\sigma}{2 \varepsilon_{0}}\left[\frac{2 l}{2 \sqrt{R^{2}+l^{2}}}-1\right] & =\frac{\sigma}{2 \varepsilon_{0}}\left[1-\frac{1}{\sqrt{1+(R / l)^{2}}}\right]
\end{aligned}
\]

when \(l \rightarrow 0, \varphi \approx \frac{\sigma R}{2 \varepsilon_{0}}, E=\frac{\sigma}{2 \varepsilon_{0}}\) and when \(l \gg R\),

\[
\varphi \approx \frac{\sigma R^{2}}{4 \varepsilon_{0} l}, E=\frac{\sigma R^{2}}{4 \varepsilon_{0} l^{2}}
\]

3.34 By definition, the potential in the case of a surface charge distribution is defined by integral \(\varphi=\frac{1}{4 \pi \varepsilon_{0}} \int \frac{\sigma d S}{r}\). In order to simplify integration, we shall choose the area element \(d S\) in the form of a part of the ring of radius \(r\) and width \(d r\) in (Fig.). Then \(d S=2 \theta r d r\), \(r=2 R \cos \theta\) and \(d r=-2 R \sin \theta d \theta\). After substituting these expressions into integral \(\varphi=\frac{1}{4 \pi \varepsilon_{0}} \int \frac{\sigma d S}{r}\), we obtain the expression for \(\varphi\) at the point \(O\) :\\
\(\varphi=-\frac{\sigma R}{\pi \varepsilon_{0}} \int_{\pi / 2}^{0} \theta \sin \theta d \theta\).\\
We integrate by parts, denoting \(\theta=u\) and \(\sin \theta d \theta=d v:\)\\
\(\int \theta \sin \theta d \theta=-\theta \cos \theta\)\\
\(+\int \cos \theta d \theta=-\theta \cos \theta+\sin \theta\)\\
which gives -1 after substituting the limits of integration. As a result, we obtain

\[
\varphi=\sigma R / \pi \varepsilon_{0} .
\]

\includegraphics[max width=\textwidth, alt={}, center]{9f3c5a8e-ac5d-42f4-8eb6-ea2af623f2d8-289_473_465_1515_796}\\
3.35 In accordance with the problem \(\varphi=\vec{a} \cdot \vec{r}\)

Thus from the equation : \(\vec{E}=-\vec{\nabla}_{\varphi}\)\\
\(\vec{E}=-\left[\frac{\partial}{\partial x}\left(a_{x} x\right) \vec{i}+\frac{\partial}{\partial_{y}}\left(a_{y} y\right) \vec{j}+\frac{\partial}{\partial_{z}}\left(a_{z} z\right) \vec{k}\right]=-\left[a_{x} \vec{i}+a_{y} \vec{j}+{ }_{z} \vec{k}\right]=-\vec{a}\)\\
3.36 (a) Given, \(\varphi=a\left(x^{2}-y^{2}\right)\)

So,

\[
\vec{E}=-\vec{\nabla} \varphi=-2 a(x \vec{i}-y \vec{j})
\]

The sought shape of field lines is as shown in the figure (a) of answersheet assuming \(a>0\) :\\
(b) Since \(\varphi=a x y\)

So,

\[
\vec{E}=-\vec{\nabla} \varphi=-a y \vec{i}-a x \vec{j}
\]

Plot as shown in the figure (b) of answersheet.\\
3.37 Given, \(\varphi=a\left(x^{2}+y^{2}\right)+b z^{2}\)

So,

\[
\begin{gathered}
\vec{E}=-\vec{\nabla} \varphi=-[2 a x \vec{i}+2 a y \vec{j}+2 b z \vec{k}] \\
|\vec{E}|=2 \sqrt{a^{2}\left(x^{2}+y^{2}\right)+b^{2} z^{2}}
\end{gathered}
\]

Hence\\
Shape of the equipotential surface :\\
Put

\[
\vec{\rho}=x \vec{i}+y \vec{j} \text { or } \rho^{2}=x^{2}+y^{2}
\]

Then the equipotential surface has the equation

\[
a \rho^{2}+b z^{2}=\text { constant }=\varphi
\]

If \(a>0, b>0\) then \(\varphi>0\) and the equation of the equipotential surface is

\[
\frac{\rho}{\varphi / a}+\frac{z^{2}}{\varphi / b}=1
\]

which is an ellipse in \(\rho, z\) coordinates. In three dimensions the surface is an ellipsoid of revolution with semi- axis \(\sqrt{\varphi / a}, \sqrt{\varphi / a}, \sqrt{\varphi / b}\).\\
If \(a>0, b<0\) then \(\varphi\) can be \(\geq 0\). If \(\varphi>0\) then the equation is

\[
\frac{\rho^{2}}{\varphi / a}-\frac{\chi^{2}}{\varphi /|b|}=1
\]

This is a single cavity hyperboloid of revolution about \(z\) axis. If \(\varphi=0\) then

\[
\begin{aligned}
& a \rho^{2}-|b| z^{2}=0 \\
& z= \pm \sqrt{\frac{a}{|b|}} \rho
\end{aligned}
\]

or\\
is the equation of a right circular cone.\\
If \(\varphi<0\) then the equation can be written as\\
or

\[
\begin{gathered}
|b| z^{2}-a \rho^{2}=|\varphi| \\
\frac{z^{2}}{|\varphi| /|b|}-\frac{\rho^{2}}{|\varphi| / a}=1
\end{gathered}
\]

This is a two cavity hyperboloid of revolution about \(z\)-axis.\\
3.38 From Gauss' theorem intensity at a point, inside the sphere at a distance \(r\) from the centre is given by, \(E_{r}=\frac{\rho r}{3 \varepsilon_{0}}\) and outside it, is given by \(E_{r}=\frac{1}{4 \pi \varepsilon_{0}} \frac{q}{r^{2}}\).\\
(a) Potential at the centre of the sphere,\\
as

\[
\begin{gathered}
\varphi_{0}=\int_{0}^{\infty} \vec{E} \cdot d \vec{r}=\int_{0}^{R} \frac{\rho r}{3 \varepsilon_{0}} d r+\int_{R}^{\infty} \frac{q}{4 \pi \varepsilon_{0} r^{2}} d r=\frac{\rho}{3 \varepsilon} \frac{R^{2}}{2}+\frac{q}{4 \pi \varepsilon_{0} R} \\
=\frac{q}{8 \pi \varepsilon_{0} R}+\frac{q}{4 \pi \varepsilon_{0} R}=\frac{3 q}{8 \pi \varepsilon_{0} R}\left(\text { as } \rho=\frac{3 q}{4 \pi R^{3}}\right)
\end{gathered}
\]

(b) Now, potential at any point, inside the sphere, at a distance \(r\) from it s centre.

\[
\varphi(r)=\int_{r}^{R} \frac{\rho}{3 \varepsilon_{0}} r d r \int_{r}^{\infty} \frac{q}{4 \pi \varepsilon_{0}} \frac{d r}{r^{2}}
\]

On integration : \(\varphi(r)=\frac{3 q}{8 \pi \varepsilon_{0} R}\left[1-\frac{r^{2}}{3 R^{2}}\right]=\varphi_{0}\left[1-\frac{r^{2}}{3 R^{2}}\right]\)\\
3.39 Let two charges \(+q\) and \(-q\) be separated by a distance \(l\). Then electric potential at a point at distance \(r \gg l\) from this dipole,

But


\begin{gather*}
\varphi(r)=\frac{+q}{4 \pi \varepsilon_{0} r_{+}}+\frac{-q}{4 \pi \varepsilon_{0} r_{-}}=\frac{q}{4 \pi \varepsilon_{0}}\left(\frac{r_{-}-r_{+}}{r_{+} r_{-}}\right)  \tag{1}\\
r_{-}-r_{+}=l \cos \theta \text { and } r_{+} r_{-} \cong r^{2}
\end{gather*}


From Eqs. (1) and (2),

\[
\varphi(r)=\frac{q l \cos \theta}{4 \pi \varepsilon_{0} r^{2}}=\frac{p \cos \theta}{4 \pi \varepsilon_{0} r^{2}} \varphi=\frac{\vec{p} \cdot \vec{r}}{4 \pi \varepsilon_{0} r^{3}}+q \vec{r}_{+}
\]

where \(p\) is magnitude of electric moment vector.\\
Now, \(E_{r}=-\frac{\partial \varphi}{\partial r}=\frac{2 p \cos \theta}{4 \pi \varepsilon_{0} r^{3}}\)\\
and \(E_{\theta}=-\frac{\partial \varphi}{r \partial \theta}=\frac{p \sin \theta}{4 \pi \varepsilon_{0} r^{3}}\)\\
\includegraphics[max width=\textwidth, alt={}, center]{9f3c5a8e-ac5d-42f4-8eb6-ea2af623f2d8-291_297_170_1290_811}

So \(E=\sqrt{E_{r}^{2}+E_{0}^{2}}=\frac{p}{4 \pi \varepsilon_{0} r^{3}} \sqrt{4 \cos ^{2} \theta+\sin ^{2} \theta}\)\\
3.40 From the results, obtained in the previous problem,

\[
E_{r}=\frac{2 p \cos \theta}{4 \pi \varepsilon_{0} r^{3}} \text { and } E_{\theta}=\frac{p \sin \theta}{4 \pi \varepsilon_{0} r^{3}}
\]

From the given figure, it is clear that,

\[
E_{z}=E_{r} \cos \theta-E_{\theta} \sin \theta=\frac{p}{4 \pi \varepsilon_{0} r^{3}}\left(3 \cos ^{2} \theta-1\right)
\]

and

\[
E_{\perp}=E_{r} \sin \theta+E_{0} \cos \theta=\frac{3 p \sin \theta \cos \theta}{4 \pi \varepsilon_{0} r^{3}}
\]

When

\[
\overrightarrow{E \perp} \vec{p},|\vec{E}|=E_{\perp} \text { and } E_{z}=0
\]

So

\[
3 \cos ^{2} \theta=1 \text { and } \cos \theta=\frac{1}{\sqrt{3}}
\]

Thus \(\vec{E}_{\perp} \vec{p}\) at the points located on the lateral surface of the cone, having its axis, coinciding with the direction of \(z\)-axis and semi vertex angle \(\theta=\cos ^{-1} 1 / \sqrt{3}\).\\
3.41 Let us assume that the dipole is at the centre of the one equipotential surface which is spherical (Fig.). On an equipotential surface the net electric field strength along the tangent of it becomes zero. Thus\\
\(-E_{0} \sin \theta+E_{\theta}=0\) or \(-E_{0} \sin \theta+\frac{p \sin \theta}{4 \pi \varepsilon_{0} r^{3}}=0\)

Hence

\[
r=\left(\frac{p}{4 \pi \varepsilon_{0} E_{0}}\right)^{1 / 3}
\]

Alternate : Potential at the point, near the dipole is given by,

\[
\begin{aligned}
& \varphi=\frac{\vec{p} \cdot \vec{r}}{4 \pi \varepsilon_{0} r^{3}}-\vec{E}_{0} \cdot \vec{r}+\text { constant }, \\
& =\left(\frac{p}{4 \pi \varepsilon_{0} r^{3}}-E_{0}\right) \cos \theta+\text { Const }
\end{aligned}
\]

For \(\varphi\) to be constant,\\
\includegraphics[max width=\textwidth, alt={}, center]{9f3c5a8e-ac5d-42f4-8eb6-ea2af623f2d8-292_294_415_873_928}

\[
\frac{p}{4 \pi \varepsilon_{0} r^{3}}-E_{0}=0 \text { or, } \frac{p}{4 \pi \varepsilon_{0} r^{3}}=E_{0}
\]

Thus

\[
r=\left(\frac{p}{4 \pi \varepsilon_{0} E_{0}}\right)^{1 / 3}
\]

3.42 Let \(P\) be a point, at distace \(r \gg l\) and at an angle to \(\theta\) the vector \(\vec{l}\) (Fig.).

Thus \(\vec{E}\) at \(P=\frac{\lambda}{2 \pi \varepsilon_{0}} \frac{\vec{r}+\frac{\vec{l}}{2}}{\left|\vec{r}+\frac{\vec{l}}{2}\right|^{2}}-\frac{\lambda}{2 \pi \varepsilon_{0}} \frac{\vec{r}-\frac{\vec{l}}{2}}{\left|\vec{r}-\frac{\vec{l}}{2}\right|^{2}}\)

\[
\begin{aligned}
& =\frac{\lambda}{2 \pi \varepsilon_{0}}\left[\frac{\vec{r}+\vec{l} / 2}{r^{2}+\frac{l^{2}}{4}+r l \cos \theta}-\frac{\vec{r}-\vec{l} / 2}{r^{2}+\frac{\dot{l}^{2}}{4}-r l \cos \theta}\right] \\
& \quad=\frac{\lambda}{2 \pi \varepsilon_{0}}\left(\frac{\vec{l}}{r^{2}}-\frac{2 l \vec{r}}{r^{3}} \cos \theta\right)
\end{aligned}
\]

Hence \(E=|\vec{E}|=\frac{\lambda l}{2 \pi \varepsilon_{0} r^{2}}, r \gg l\)\\
\includegraphics[max width=\textwidth, alt={}, center]{9f3c5a8e-ac5d-42f4-8eb6-ea2af623f2d8-292_490_531_1552_846}

Also,

\[
\begin{gathered}
\left.\left.\varphi=\frac{\lambda}{2 \pi \varepsilon_{0}} \ln |\vec{r}+\vec{l} / 2|-\frac{\lambda}{2 \pi \varepsilon_{0}} \ln \right\rvert\, \vec{r}-\vec{l} / 2\right] \\
=\frac{\lambda}{4 \pi \varepsilon_{0}} \ln \frac{r^{2}+r l \cos \theta+l^{2} / 4}{r^{2}-r l \cos \theta+l^{2} / 4}=\frac{\lambda l \cos \theta}{2 \pi \varepsilon_{0} r}, r \gg l
\end{gathered}
\]

3.43 The potential can be calculated by superposition. Choose the plane of the upper ring as \(x=l / 2\) and that of the lower ring as \(x=-l / 2\).

\[
\begin{aligned}
& \text { Then } \quad \varphi=\frac{q}{4 \pi \varepsilon_{0}\left[R^{2}+(x-l / 2)^{2}\right]^{1 / 2}}-\frac{q}{4 \pi \varepsilon_{0}\left[R^{2}+(x+l / 2)^{2}\right]^{1 / 2}} \\
& \quad \cong \frac{q}{4 \pi \varepsilon_{0}\left[R^{2}+x^{2}-l x\right]^{1 / 2}}-\frac{q}{4 \pi \varepsilon_{0}\left[R^{2}+x^{2}+l x\right]^{1 / 2}} \\
& \propto \frac{q}{4 \pi \varepsilon_{0}\left(R^{2}+x^{2}\right)^{1 / 2}}\left(1+\frac{l x}{2\left(R^{2}+x^{2}\right)}\right)-\frac{q}{4 \pi \varepsilon_{0}\left(R^{2}+x^{2}\right)^{1 / 2}}\left(1-\frac{l x}{2\left(R^{2}+x^{2}\right)}\right) \\
& \quad \propto \frac{q l x}{4 \pi \varepsilon_{0}\left(R^{2}+x^{2}\right)^{3 / 2}} \\
& \text { For } \quad|x| \gg R, \varphi \approx \frac{q l}{4 \pi x^{2}}
\end{aligned}
\]

The electric field is \(E=-\frac{\partial \varphi}{\partial x}\)

\[
=-\frac{q l}{4 \pi \varepsilon_{0}\left(R^{2}+x^{2}\right)^{3 / 2}}+\frac{3}{2} \frac{q l}{\left(R^{2}+x^{2}\right)^{5 / 2} 4 \pi \varepsilon_{0}} \times 2 x=\frac{q l\left(2 x^{2}-R^{2}\right)}{4 \pi \varepsilon_{0}\left(R^{2}+x^{2}\right)^{5 / 2}}
\]

For

\[
|x| \gg R, E \approx \frac{q l}{2 \pi \varepsilon_{0} x^{3}} . \quad \text { The plot is as given in the book. }
\]

3.44 The field of a pair of oppositely charged sheets with holes can by superposition be reduced to that of a pair of uniform opposite charged sheets and discs with opposite charges. Now the charged sheets do not contribute any field outside them. Thus using the result of the previous problem

\[
\begin{aligned}
& \varphi=\int_{0}^{R} \frac{(-\sigma) l 2 \pi r d r x}{4 \pi \varepsilon_{0}\left(r^{2}+x^{2}\right)^{3 / 2}} \\
& =-\frac{\sigma x l}{4 \varepsilon_{0}} \int_{x^{2}}^{R^{2}+x^{2}} \frac{d y}{y^{3 / 2}}=\frac{\sigma x l}{2 \varepsilon_{0} \sqrt{R^{2}+x^{2}}} \\
& E_{x}=-\frac{\partial \varphi}{\partial x}=-\frac{\sigma l}{2 \varepsilon_{0}}\left[\frac{1}{\sqrt{R^{2}+x^{2}}}-\frac{x^{2}}{\left(R^{2}+x^{2}\right)^{3 / 2}}\right]=-\frac{\sigma l R^{2}}{2 \varepsilon_{0}\left(R^{2}+x^{2}\right)^{3 / 2}} .
\end{aligned}
\]

\begin{center}
\includegraphics[max width=\textwidth, alt={}]{9f3c5a8e-ac5d-42f4-8eb6-ea2af623f2d8-293_397_418_1412_899}
\end{center}

The plot is as shown in the answersheet.\\
3.45 For \(x>0\) we can use the result as given above and write

\[
\varphi \approx \pm \frac{\sigma l}{2 \varepsilon_{0}}\left(1-\frac{|x|}{\left(R^{2}+x^{2}\right)^{1 / 2}}\right)
\]

for the solution that vanishes at \(\alpha\). There is a discontinuity in potential for \(|x|=0\). The solution for negative \(x\) is obtained by \(\sigma \rightarrow-\sigma\). Thus

\[
\varphi=-\frac{\sigma l x}{2 \varepsilon_{0}\left(R+x^{2}\right)^{1 / 2}}+\text { constant }
\]

Hence ignoring the jump

\[
E=-\frac{\partial \varphi}{\partial x}=\frac{\sigma l R^{2}}{2 \varepsilon_{0}\left(R^{2}+x^{2}\right)^{3 / 2}}
\]

for large \(|x| \quad \varphi \approx \pm \frac{p}{4 \pi \varepsilon_{0} x^{2}}\) and \(E \approx \frac{p}{2 \pi \varepsilon_{0}|x|^{3}}\left(\right.\) where \(\left.p=\pi R^{2} \sigma l\right)\)\\
3.46 Here \(E_{r}=\frac{\lambda}{2 \pi \varepsilon_{0} r}, E_{\theta}=E_{\varphi}=0\) and \(\vec{F}=p \frac{\partial \vec{E}}{\partial l}\)\\
(a) \(\vec{p}\) along the thread.\\
\(\vec{E}\) does not change as the point of observation is moved along the thread.

\[
\vec{F}=0
\]

(b) \(\vec{p}\) along \(\vec{r}\),

\[
\vec{F}=F_{r} \vec{e}_{r}=\frac{\lambda p}{2 \pi \varepsilon_{0} r^{2}} \vec{e}_{r}=-\frac{\lambda \vec{p}}{2 \pi \varepsilon_{0} r^{2}}\left(\text { On using } \frac{\partial}{\partial r} \vec{e}_{r}=0\right)
\]

(c) \(\vec{p}\) along \(\overrightarrow{e_{\theta}}\)

\[
\begin{gathered}
\vec{F}=p \frac{\partial}{r \partial \theta} \frac{\lambda}{2 \pi \varepsilon_{0} r} e_{r} \\
=\frac{p \lambda}{2 \pi \varepsilon_{0} r^{2}} \frac{\partial \overrightarrow{e_{r}}}{\partial \theta}=\frac{p \lambda}{2 \pi \varepsilon_{0} r^{2}} \overrightarrow{e_{\theta}}=\frac{\vec{p} \lambda}{2 \pi \varepsilon_{0} r^{2}} .
\end{gathered}
\]

3.47 Force on a dipole of moment \(\boldsymbol{p}\) is given by,

\[
F=\left|\varphi \frac{\partial \vec{E}}{\partial l}\right|
\]

In our problem, field, due to a dipole at a distance \(l\), where a dipole is placed,

\[
|\vec{E}|=\frac{p}{2 \pi \varepsilon_{0} l^{3}}
\]

Hence, the force of interaction,

\[
F=\frac{3 p^{2}}{2 \pi \varepsilon_{0} l^{4}}=2.1 \times 10^{-16} \mathrm{~N}
\]

\(3.48-d \varphi=\vec{E} \cdot d \vec{r}=a(y d x+x d y)=a d(x y)\)\\
On integrating,

\[
\varphi=-a x y+C
\]

\(3.49-d \varphi=\vec{E} \cdot d \vec{r}=\left[2 a x y \vec{i}+2\left(x^{2}-y^{2}\right) \vec{j}\right] \cdot[d x \vec{i}+d y \vec{j}]\)\\
or,

\[
d \varphi=2 a x y d x+a\left(x^{2}-y^{2}\right) d y=a d\left(x^{2} y\right)-a y^{2} d y
\]

On integrating, we get,

\[
\varphi=a y\left(\frac{y^{2}}{3}-x^{2}\right)+C
\]

3.50 Given, again

\[
\begin{aligned}
-d \varphi & =\vec{E} \cdot d \vec{r}=(a y \vec{i}+(a x+b z) \vec{j}+b y \vec{k}) \cdot(d x \vec{i}+d y \vec{j}+d x \vec{k}) \\
& =a(y d x+a x d y)+b(z d y+y d z)=a d(x y)+b d(y z)
\end{aligned}
\]

On integrating,

\[
\varphi=-(a x y+b y z)+C
\]

3.51 Field intensity along \(x\)-axis.


\begin{equation*}
E_{x}=-\frac{\partial \varphi}{\partial x}=3 a x^{2} \tag{1}
\end{equation*}


Then using Gauss's theorem in differential from

\[
\frac{\partial E_{x}}{\partial x}=\frac{\rho(x)}{\varepsilon_{0}} \text { so, } \rho(x)=6 a \varepsilon_{0} x .
\]

3.52 In the space between the plates we have the Poisson equation\\
or,

\[
\begin{gathered}
\frac{\partial^{2} \varphi}{\partial x^{2}}=-\frac{\rho_{0}}{\varepsilon_{0}} \\
\varphi=-\frac{\rho_{0}}{2 \varepsilon_{0}} x^{2}+A x+B
\end{gathered}
\]

where \(\rho_{0}\) is the constant space charge density between the plates.\\
We can choose

\[
\varphi(0)=0 \text { so } B=0
\]

Then

\[
\varphi(d)=\Delta \varphi=A d-\frac{\rho_{0} d^{2}}{2 \varepsilon_{0}} \quad \text { or, } \quad A=\frac{\Delta \varphi}{d}+\frac{\rho_{0} d}{2 \varepsilon_{0}}
\]

Now

\[
E=-\frac{\partial \varphi}{\partial x}=\frac{\rho_{0}}{\varepsilon_{0}} x-A=0 \text { for } x=0
\]

if

\[
A=\frac{\Delta \varphi}{d}+\frac{\rho_{0} d}{2 \varepsilon_{0}}=0
\]

then

\[
\rho_{0}=-\frac{2 \varepsilon_{0} \Delta \varphi}{d^{2}}
\]

Also

\[
E(d)=\frac{\rho_{0} d}{\varepsilon_{0}}
\]

3.53 Field intensity is along radial line and is


\begin{equation*}
E_{r}=-\frac{\partial \varphi}{\partial r}=-2 a r \tag{1}
\end{equation*}


From the Gauss' theorem,

\[
4 \pi r^{2} E_{r}=\int \frac{d q}{\varepsilon_{0}}
\]

where \(d q\) is the charge contained between the sphere of radii \(r\) and \(r+d r\).\\
Hence


\begin{equation*}
4 \pi r^{2} E_{r}=4 \pi r^{2} \times(-2 a r)=\frac{4 \pi}{\varepsilon_{0}} \int_{0}^{r} r^{\prime 2} \rho\left(r^{\prime}\right) d r^{\prime} \tag{2}
\end{equation*}


Differentiating (2) \(\rho=-6 \varepsilon_{0} a\)

\subsection*{3.2 CONDUCTORS AND DIELECTRICS IN AN ELECTRIC FIELD}
3.54 When the ball is charged, for the equilibrium of ball, electric force on it must counter balance the excess spring force, exerted, on the ball due to the extension in the spring.\\
Thus \(F_{e l}=F_{s p r}\)\\
or, \(\frac{q^{2}}{4 \pi \varepsilon_{0}(2 l)^{2}}=\kappa x\), (The force on the charge \(q\) might be considered as arised from attraction by the electrical image)\\
or,\(q=4 l \sqrt{\pi \varepsilon_{0} \kappa x}\),\\
sought charge on the sphere.\\
\includegraphics[max width=\textwidth, alt={}, center]{9f3c5a8e-ac5d-42f4-8eb6-ea2af623f2d8-296_501_331_184_960}\\
3.55 By definition, the work of this force done upon an elementry displacement \(d x\) (Fig.) is given by

\[
d A=F_{x} d x=-\frac{q^{2}}{4 \pi \varepsilon_{0}(2 x)^{2}} d x
\]

where the expression for the force is obtained with the help of the image method. Integrating this equation over \(x\) between \(l\) and \(\infty\), we find

\[
A=-\frac{q^{2}}{16 \pi \varepsilon_{0}} \int_{l}^{\infty} \frac{d x}{x^{2}}=-\frac{q^{2}}{16 \pi \varepsilon_{0} l}
\]

\includegraphics[max width=\textwidth, alt={}, center]{9f3c5a8e-ac5d-42f4-8eb6-ea2af623f2d8-296_383_557_783_850}\\
3.56 (a) Using the concept of electrical image, it is clear that the magnitude of the force acting on each charge,

\[
\begin{aligned}
|\vec{F}| & =\sqrt{2} \frac{q^{2}}{4 \pi \varepsilon_{0} l^{2}}-\frac{q^{2}}{4 \pi \varepsilon_{0}(\sqrt{2} l)^{2}} \\
& =\frac{q^{2}}{8 \pi \varepsilon_{0} l^{2}}(2 \sqrt{2}-1)
\end{aligned}
\]

(b) Also, from the figure, magnitude of electrical field strength at \(P\)

\[
E=2\left(1-\frac{1}{5 \sqrt{5}}\right) \frac{q}{\pi \varepsilon_{0} l^{2}}
\]

3.57 Using the concept of electrical image, it is easily seen that the force on the charge \(q\) is,

\[
\begin{aligned}
& F=\frac{\sqrt{2} q^{2}}{4 \pi \varepsilon_{0}(2 l)^{2}}+\frac{(-q)^{2}}{4 \pi \varepsilon_{0}(2 \sqrt{2} l)^{2}} \\
& =\frac{(2 \sqrt{2}-1) q^{2}}{32 \pi \varepsilon_{0} l^{2}} \text { (It is attractive) }
\end{aligned}
\]

\includegraphics[max width=\textwidth, alt={}, center]{9f3c5a8e-ac5d-42f4-8eb6-ea2af623f2d8-296_415_437_1242_920}\\
\includegraphics[max width=\textwidth, alt={}, center]{9f3c5a8e-ac5d-42f4-8eb6-ea2af623f2d8-296_355_383_1677_920}\\
3.58 Using the concept of electrical image, force on the dipole \(\vec{p}\),\\
\(\vec{F}=p \frac{\partial \vec{E}}{\partial l}\), where \(\vec{E}\) is field at the location of\\
\(\vec{p}\) due to \((-\vec{p})\)\\
or, \(|\vec{F}|=\left|\frac{\partial \vec{E}}{\partial l}\right| p=\frac{3 p^{2}}{32 \pi \varepsilon_{0} l^{4}}\)\\
\includegraphics[max width=\textwidth, alt={}, center]{9f3c5a8e-ac5d-42f4-8eb6-ea2af623f2d8-297_168_413_235_949}\\
as, \(|\vec{E}|=\frac{p}{4 \pi \varepsilon_{0}(2 l)^{3}}\)\\
3.59 To find the surface charge density, we must know the electric field at the point \(P\) (Fig.) which is at a distance \(r\) from the point \(O\).\\
Using the image mirror method, the field at \(P\),

\[
E=2 E \cos \alpha=2 \frac{q}{4 \pi \varepsilon_{0} x^{2}} \frac{l}{x}=\frac{q l}{2 \pi \varepsilon_{0}\left(l^{2}+r^{2}\right)^{3 / 2}}
\]

Now from Gauss' theorem the surface charge density on conductor is connected with the electric field near its surface (in vaccum) through the relation \(\sigma=\varepsilon_{0} E_{n}\), where \(E_{n}\) is the projection of \(\vec{E}\) onto the outward normal \(\vec{n}\) (with respect to the conductor).\\
As our field strength \(\vec{E} \uparrow \downarrow \vec{n}\), so\\
\includegraphics[max width=\textwidth, alt={}, center]{9f3c5a8e-ac5d-42f4-8eb6-ea2af623f2d8-297_455_509_585_862}

\[
\sigma=-\varepsilon_{0} E=-\frac{q l}{2 \pi\left(l^{2}+r^{2}\right)^{3 / 2}}
\]

3.60 (a) The force \(F_{1}\) on unit length of the thread is given by

\[
F_{1}=\lambda E_{1}
\]

where \(E_{1}\) is the field at the thread due to image charge :

\[
E_{1}=\frac{-\lambda}{2 \pi \varepsilon_{0}(2 l)}
\]

Thus

\[
F_{1}=\frac{-\lambda^{2}}{4 \pi \varepsilon_{0} l}
\]

minus singn means that the force is one of attraction.\\
(b) There is an image thread with charge density- \(\lambda\) behind the conducting plane. We\\
\includegraphics[max width=\textwidth, alt={}]{9f3c5a8e-ac5d-42f4-8eb6-ea2af623f2d8-297_407_473_1278_917} calculate the electric field on the conductor. It is

\[
E(x)=E_{n}(x)=\frac{\lambda l}{\pi \varepsilon_{0}\left(x^{2}+l^{2}\right)}
\]

on considering the thread and its image.\\
Thus

\[
\sigma(x)=\varepsilon_{0} E_{n}=\frac{\lambda l}{\pi\left(x^{2}+l^{2}\right)}
\]

3.61 (a) At \(O\),

\[
E_{n}(O)=2 \int_{l}^{\infty} \frac{\lambda d x}{4 \pi \varepsilon_{0} x^{2}}=\frac{\lambda}{2 \pi \varepsilon_{0} l}
\]

So

\[
\sigma(O)=\varepsilon_{0} E_{n}=\frac{\lambda}{2 \pi l}
\]

(b) \(E_{n}(r)=2 \int_{l}^{\infty} \frac{\lambda d x}{4 \pi \varepsilon_{0}\left(x^{2}+r^{2}\right)} \frac{x}{\left(x^{2}+r^{2}\right)^{1 / 2}}=\frac{\lambda}{2 \pi \varepsilon_{0}} \int_{l}^{\infty} \frac{x d x}{\left(x^{2}+r^{2}\right)^{3 / 2}}\)

\[
\begin{array}{r}
=\frac{\lambda}{4 \pi \varepsilon_{0}} \int_{l^{2}+r^{2}}^{\infty} \frac{d y}{y^{3 / 2}}, \text { on putting } \\
=\frac{\lambda}{2 \pi \varepsilon_{0} \sqrt{l^{2}+r^{2}}}
\end{array}
\]

Hence \(\sigma(r)=\varepsilon_{0} E_{n}=\frac{\lambda}{2 \pi \sqrt{l^{2}+r^{2}}}\)\\
\includegraphics[max width=\textwidth, alt={}, center]{9f3c5a8e-ac5d-42f4-8eb6-ea2af623f2d8-298_406_484_630_867}\\
3.62 It can be easily seen that in accordance with the image method, a charge \(-q\) must be located on a similar ring but on the other side of the conducting plane. (Fig.) at the same perpendicular distance. From the solution of 3.9 net electric field at \(O\),\\
\(\vec{E}=2 \frac{q l}{4 \pi \varepsilon_{0}\left(R^{2}+l^{2}\right)^{3 / 2}}(-\vec{n})\) where \(\vec{n}\) is\\
outward normal with respect to the conducting plane.

Now

\[
E_{n}=\frac{\sigma}{\varepsilon_{0}}
\]

Hence \(\quad \sigma=\frac{-q l}{2 \pi\left(R^{2}+l^{2}\right)^{3 / 2}}\)\\
where minus sign indicates that the induced carge is opposite in sign to that of charge\\
\includegraphics[max width=\textwidth, alt={}, center]{9f3c5a8e-ac5d-42f4-8eb6-ea2af623f2d8-298_484_404_1205_920}\\
\(q>0\).\\
3.63 Potential \(\varphi\) is the same for all the points of the sphere. Thus we calculate its value at the centre \(O\) of the sphere. Thus we can calculate its value at the centre \(O\) of the sphere, because only for this point, it can be calculated in the most simple way.


\begin{equation*}
\varphi=\frac{1}{4 \pi \varepsilon_{0}} \frac{q}{l}+\varphi^{\prime} \tag{1}
\end{equation*}


where the first term is the potential of the charge \(q\), while the second is the potential due to the charges induced on the surface of the sphere. But since all induced charges are at the same distance equal to the radius of the circle from the point \(C\) and the total induced charge is equal to zero, \(\varphi^{\prime}=0\), as well. Thus equation (1) is reduced to the form,

\[
\varphi=\frac{1}{4 \pi \varepsilon_{0}} \frac{q}{l}
\]

3.64 As the sphere has conducting layers, charge \(-q\) is induced on the inner surface of the sphere \(q\) and consequently charge \(+q\) is induced on the outer layer as the sphere as a whole is uncharged.\\
Hence, the potential at \(O\) is given by,

\[
\varphi_{0}=\frac{q}{4 \pi \varepsilon_{0} r}+\frac{(-q)}{4 \pi \varepsilon_{0} R_{1}}+\frac{q}{4 \pi \varepsilon_{0} R_{2}}
\]

It should be noticed that the potential can be found in such a simple way only at \(O\), since all the induced charges are at the same distance from this point, and their distribution, (which\\
\includegraphics[max width=\textwidth, alt={}, center]{9f3c5a8e-ac5d-42f4-8eb6-ea2af623f2d8-299_358_548_90_787}\\
\includegraphics[max width=\textwidth, alt={}, center]{9f3c5a8e-ac5d-42f4-8eb6-ea2af623f2d8-299_373_371_765_816}\\
is unknown to us), does not play any role.\\
3.65 Potential at the inside sphere,

\[
\varphi_{a}=\frac{q_{1}}{4 \pi \varepsilon_{0} a}+\frac{q_{2}}{4 \pi \varepsilon_{0} b}
\]

Obviously


\begin{equation*}
\varphi_{a}=0 \text { for } q_{2}=-\frac{b}{a} q_{1} \tag{1}
\end{equation*}


When

\[
r \geq b,
\]

\[
\varphi_{r}=\frac{q_{1}}{4 \pi \varepsilon_{0} r}+\frac{q_{2}}{4 \pi \varepsilon_{0} r}=\frac{q_{1}}{4 \pi \varepsilon_{0}}\left(1-\frac{b}{a}\right) / r, \text { using Eq. (1). }
\]

And when \(r \leq b\)

\[
\varphi_{r}=\frac{q_{1}}{4 \pi \varepsilon_{0} r}+\frac{q_{2}}{4 \pi \varepsilon_{0} b}=\frac{q_{1}}{4 \pi \varepsilon_{0}}\left(\frac{1}{r}-\frac{1}{a}\right)
\]

3.66 (a) As the metallic plates 1 and 4 are isolated and conncted by means of a conductor, \(\varphi_{1}=\varphi_{4}\). Plates 2 and 3 have the same amount of positive and negative charges and due to induction, plates 1 and 4 are respectively negatively and positively charged and in addition to it all the four plates are located a small but at equal distance \(d\) relative to each\\
other, the magnitude of electric field strength between 1-2 and 3-4 are both equal in magnitude and direction (say \(\vec{E}\) ). Let \(\overrightarrow{E^{\prime}}\) be the field strength between the plates 2 and 3 , which is directed form 2 to 3 . Hence \(\overrightarrow{E^{\prime}} \uparrow \downarrow \vec{E}\) (Fig.).\\
According to the problem


\begin{equation*}
E^{\prime} d=\Delta \varphi=\varphi_{2}-\varphi_{3} \tag{1}
\end{equation*}


In addition to

\[
\begin{aligned}
& \varphi_{1}-\varphi_{4}=0=\left(\varphi_{1}-\varphi_{2}\right)+\left(\varphi_{2}-\varphi_{3}\right)+\left(\varphi_{3}-\varphi_{4}\right) \\
& \text { or, } \quad 0=-E d+\Delta \varphi-E d
\end{aligned}
\]

\includegraphics[max width=\textwidth, alt={}, center]{9f3c5a8e-ac5d-42f4-8eb6-ea2af623f2d8-300_288_604_317_778}\\
or, \(\Delta \varphi=2 E d\) or \(E=\frac{\Delta \varphi}{2 d}\)

Hence


\begin{equation*}
E=\frac{E^{\prime}}{2}=\frac{\Delta \varphi}{2 d} \tag{2}
\end{equation*}


(b) Since \(E \alpha \sigma\), we can state that according to equation (2) for part (a) the charge on the plate 2 is divided into two parts; such that \(1 / 3 \mathrm{rd}\) of it lies on the upper side and \(2 / 3 \mathrm{rd}\) on its lower face.\\
Thus charge density of upper face of plate 2 or of plate 1 or plate 4 and lower face of \(3 \sigma=\varepsilon_{0} E=\frac{\varepsilon_{0} \Delta \varphi}{2 d}\) and charge density of lower face of 2 or upper face of 3

\[
\sigma^{\prime}=\varepsilon_{0} E^{\prime}=\varepsilon_{0} \frac{\Delta \varphi}{d}
\]

Hence the net charge density of plate 2 or 3 becomes \(\sigma+\sigma^{\prime}=\frac{3 \varepsilon_{0} \Delta \varphi}{2 d}\), which is obvious from the argument.\\
3.67 The problem of point charge between two conducting planes is more easily tackled (if we want only the total charge induced on the planes) if we replace the point charge by a uniformly charged plane sheet.\\
Let \(\sigma\) be the charge density on this sheet and \(E_{1}, E_{2}\) outward electric field on the two sides of this sheet.

Then

\[
E_{1}+E_{2}=\frac{\sigma}{\varepsilon_{0}}
\]

The conducting planes will be assumed to be grounded. Then \(E_{1} x=E_{2}(l-x)\).

Hence

\[
E_{1}=\frac{\sigma}{l \varepsilon_{0}}(l-x), E_{2}=\frac{\sigma}{l \varepsilon_{0}} x
\]

This means that the induced charge density on the plane conductors are

\[
\sigma_{1}=-\frac{\sigma}{l}(l-x), \sigma_{2}=-\frac{\sigma}{l} x
\]

\begin{center}
\includegraphics[max width=\textwidth, alt={}]{9f3c5a8e-ac5d-42f4-8eb6-ea2af623f2d8-300_428_218_1485_991}
\end{center}

Hence \(q_{1}=-\frac{q}{l}(l-x), q_{2}=-\frac{q}{l} x\)\\
3.68 Near the conductor \(E=E_{n}=\frac{\sigma}{\varepsilon_{0}}\)

This field can be written as the sum of two parts \(E_{1}\) and \(E_{2} . E_{1}\) is the electric field due to an infinitesimal area \(d S\).\\
Very near it \(E_{1}= \pm \frac{\sigma}{2 \varepsilon_{0}}\)\\
The remaining part contributes \(E_{2}=\frac{\sigma}{2 \varepsilon_{0}}\) on both sides. In calculating the force on the element \(d S\) we drop \(E_{1}\) (because it is a self-force.) Thus

\[
\frac{d F}{d S}=\sigma \cdot \frac{\sigma}{2 \varepsilon_{0}}=\frac{\sigma^{2}}{2 \varepsilon_{0}}
\]

\includegraphics[max width=\textwidth, alt={}, center]{9f3c5a8e-ac5d-42f4-8eb6-ea2af623f2d8-301_393_255_273_981}\\
3.69 The total force on the hemisphere is

\[
\begin{aligned}
F & =\int_{0}^{\pi / 2} \frac{\sigma^{2}}{2 \varepsilon_{0}} \cdot \cos \theta \cdot 2 \pi R \sin \theta R d \theta \\
& =\frac{2 \pi R^{2} \sigma^{2}}{2 \varepsilon_{0}} \int_{0}^{\pi / 2} \cos \theta \sin \theta d \theta \\
= & \frac{2 \pi R^{2}}{2 \varepsilon_{0}} \times \frac{1}{2} \times\left(\frac{q}{4 \pi R^{2}}\right)^{2}=\frac{q^{2}}{32 \pi \varepsilon_{0} R} .
\end{aligned}
\]

\includegraphics[max width=\textwidth, alt={}, center]{9f3c5a8e-ac5d-42f4-8eb6-ea2af623f2d8-301_419_421_796_872}\\
3.70 We know that the force acting on the area element \(d S\) of a conductor is,


\begin{equation*}
d \vec{F}=\frac{1}{2} \sigma \vec{E} d S \tag{1}
\end{equation*}


It follows from symmetry considerations that the resultant force \(F\) is directed along the \(z\)-axis, and hence it can be represented as the sum (integral) of the projection of elementary forces (1) onto the \(z\)-axis :


\begin{equation*}
d F_{z}=d F \cos \theta \tag{2}
\end{equation*}


For simplicity let us consider an element area \(d S=2 \pi R \sin \theta R d \theta\) (Fig.). Now considering that \(E=\sigma / \varepsilon_{0}\). Equation (2) takes the from\\
\includegraphics[max width=\textwidth, alt={}, center]{9f3c5a8e-ac5d-42f4-8eb6-ea2af623f2d8-301_422_334_1400_876}

\[
\begin{aligned}
d F_{z} & =\frac{\pi \sigma^{2} R^{2}}{\varepsilon_{0}} \sin \theta \cos \theta d \theta \\
& =-\left(\frac{\pi \sigma_{0}^{2} R^{2}}{\varepsilon_{0}}\right) \cos ^{3} \theta d \cos \theta
\end{aligned}
\]

Integrating this expression over the half sphere (i.e. with respect to \(\cos \theta\) between 1 and 0 ), we obtain

\[
F=F_{z}=\frac{\pi \sigma_{0}^{2} R^{2}}{4 \varepsilon_{0}}
\]

3.71 The total polarization is \(P=(\varepsilon-1) \varepsilon_{0} E\). This must equals \(\frac{n_{0} p}{N}\) where \(n_{0}\) is the concerntation of water molecules. Thus

\[
N=\frac{n_{0} p}{(\varepsilon-1) \varepsilon_{0} E}=2.93 \times 10^{3} \text { on putting the values }
\]

3.72 From the general formula

\[
\begin{gathered}
\vec{E}=\frac{1}{4 \pi \varepsilon_{0}} \frac{3 \vec{p} \cdot \vec{r} \vec{r}-\vec{p} r^{2}}{r^{5}} \\
\vec{E}=\frac{1}{4 \pi \varepsilon_{0}} \frac{2 \vec{p}}{l^{3}}, \text { where } r=l \text { and } \vec{r} \uparrow \uparrow \vec{p}
\end{gathered}
\]

This will cause the induction of a dipole moment.

\[
\vec{p}_{i n d}=\beta \frac{1}{4 \pi \varepsilon_{0}} \frac{2 \vec{p}}{l^{3}} \times \varepsilon_{0}
\]

Thus the force,

\[
\vec{F}=\frac{\beta}{4 \pi} \frac{2 p}{l^{3}} \frac{\partial}{\partial l} \frac{1}{4 \pi \varepsilon_{0}} \frac{2 p}{l^{3}}=\frac{3 \beta p^{2}}{4 \pi^{2} \varepsilon_{0} l^{7}}
\]

3.73 The electric field \(E\) at distance \(x\) from the centre of the ring is,

\[
E(x)=\frac{q x}{4 \pi \varepsilon_{0}\left(R^{2}+x^{2}\right)^{3 / 2}}
\]

The induced dipole moment is \(p=\beta \varepsilon_{0} E=\frac{q \beta x}{4 \pi\left(R^{2}+x^{2}\right)^{3 / 2}}\)\\
The force on this molecule is\\
\(F=p \frac{\partial}{\partial x} E=\frac{q \beta x}{4 \pi\left(R^{2}+x^{2}\right)^{3 / 2}} \frac{q}{4 \pi \varepsilon_{0}} \frac{\partial}{\partial x} \frac{x}{\left(R^{2}+x^{2^{3 / 2}}\right.}=\frac{q^{2} \beta}{16 \pi^{2} \varepsilon_{0}} \frac{x\left(R^{2}-2 x^{2}\right)}{\left(R^{2}+x^{2}\right)^{4}}\)\\
This vanishes for \(x=\frac{ \pm R}{\sqrt{2}}\) (apart from \(x=0, x=\infty\) )\\
It is maximum when

\[
\frac{\partial}{\partial x} \frac{x\left(R^{2}-x^{2} \times 2\right)}{\left(R^{2}+x^{2}\right)^{4}}=0
\]

or,

\[
\left(R^{2}-2 x^{2}\right)\left(R^{2}+x^{2}\right)-4 x^{2}\left(R^{2}+x^{2}\right)-8 x^{2}\left(R^{2}-2 x^{2}\right)=0
\]

or,

\[
R^{4}-13 x^{2} R^{2}+10 x^{4}=0 \text { or, } x^{2}=\frac{R^{2}}{20}(13 \pm \sqrt{129})
\]

or, \(x=\frac{R}{\sqrt{20}} \sqrt{13 \pm \sqrt{129}}\) (on either side), Plot of \(F_{x}(x)\) is as shown in the answersheet.\\
3.74 Inside the ball

\[
\vec{D}(\vec{r})=\frac{q}{4 \pi} \frac{\vec{r}}{r^{3}}=\varepsilon \varepsilon_{0} \vec{E} .
\]

Also

\[
\varepsilon_{0} \vec{E}+\vec{P}=\vec{D} \text { or } \vec{P}=\frac{\varepsilon-1}{\varepsilon} \vec{D}=\frac{\varepsilon-1}{\varepsilon} \frac{q}{4 \pi} \frac{\vec{r}}{r^{3}}
\]

Also,

\[
q^{\prime}=-\oint \vec{P} \cdot d \vec{S}=-\frac{\varepsilon-1}{\varepsilon} \frac{q}{4 \pi} \int d \Omega=-\frac{\varepsilon-1}{\varepsilon} q
\]

3.75

\[
\begin{aligned}
D_{\text {diel }}=\varepsilon \varepsilon_{0} E_{\text {diel }}=D_{\text {conductor }} & =\sigma \text { or, } E_{\text {diel }}=\frac{\sigma}{\varepsilon \varepsilon_{0}} \\
P_{n} & =(\varepsilon-1) \varepsilon_{0} E_{\text {diel }}=\frac{\varepsilon-1}{\varepsilon} \sigma \\
\sigma^{\prime} & =-P_{n}=-\frac{\varepsilon-1}{\varepsilon} \sigma
\end{aligned}
\]

This is the surface density of bound charges.\\
3.76 From the solution of the previous problem \(q_{\text {in }}^{\prime}=\) charge on the interior surface of the conductor

\[
=-(\varepsilon-1) / \varepsilon \int \sigma d S=-\frac{\varepsilon-1}{\varepsilon} q
\]

Since the dielectric as a whole is neutral there must be a total charge equal to \(q_{\text {outer }}^{\prime}=+\frac{\varepsilon-1}{\varepsilon} q\) on the outer surface of the dielectric.\\
3.77 (a) Positive extraneous charge is distributed uniformly over the internal surface layer. Let \(\sigma_{0}\) be the surface density of the charge.\\
Clearly,

\[
E=0, \text { for } r<a
\]

For \(a<r\)\\
\(\varepsilon_{0} E \times 4 \pi r^{2}=4 \pi a^{2} \sigma_{0}\) by Gauss theorem.\\
or, \(E=\frac{\sigma_{0}}{\varepsilon_{0} \varepsilon}\left(\frac{a}{r}\right)^{2}, a<r<b\)\\
For \(r>b\), similarly

\[
E=\frac{\sigma_{0}}{\varepsilon_{0}}\left(\frac{a}{r}\right)^{2}, r>b
\]

\begin{center}
\includegraphics[max width=\textwidth, alt={}]{9f3c5a8e-ac5d-42f4-8eb6-ea2af623f2d8-303_403_426_1331_925}
\end{center}

Now, \(\quad E=-\frac{\partial \varphi}{\partial r}\).\\
So by integration from infinity where \(\varphi(\infty)=0\),

\[
\varphi=\frac{\sigma_{0} a^{2}}{\varepsilon_{0} r} r>b
\]

\(a<r<b \quad \varphi=\frac{\sigma_{0} a^{2}}{\varepsilon \varepsilon r}+B, B\) is a constant\\
or by continuity, \(\varphi=\frac{\sigma_{0} a^{2}}{\varepsilon_{0} \varepsilon}\left(\frac{1}{r}-\frac{1}{b}\right)+\frac{\sigma_{0} a^{2}}{\varepsilon_{0} b}, a<r<b\)\\
For \(r<a . \varphi=A=\) Constant\\
By continuity, \(\varphi=\frac{\sigma_{0} a^{2}}{\varepsilon_{0} \varepsilon}\left(\frac{1}{a}-\frac{1}{b}\right)+\frac{\sigma_{0} a^{2}}{\varepsilon_{0} b}\)\\
(b) Positive extraneous charge is distributed uniformly over the internal volume of the dielectric\\
Let \(\rho_{0}=\) Volume density of the charge in the dielectric, for \(a<r<b\).

\[
\begin{array}{cc} 
& E=0, r<a \\
\varepsilon_{0} \varepsilon 4 \pi r^{2} E=\frac{4 \pi}{3}\left(r^{3}-a^{3}\right) \rho_{0},(a<r<b) \\
\text { or, } & E=\frac{\rho_{0}}{3 \varepsilon_{0} \varepsilon}\left(r-\frac{a^{3}}{r^{2}}\right) \\
E=\frac{4 \pi}{3}\left(b^{3}-a^{3}\right) \rho_{0} / \varepsilon_{0} 4 \pi r^{2}, r>b \\
\text { or, } & E=\frac{\left(b^{3}-a^{3}\right) \rho_{0}}{3 \varepsilon_{0} r^{2}} \text { for } r>b
\end{array}
\]

By integration,\\
or,

\[
\begin{gathered}
\varphi=\frac{\left(b^{3}-a^{3}\right) \rho_{0}}{3 \varepsilon_{0} r} \text { for } r>b \\
\varphi=B-\frac{\rho_{0}}{3 \varepsilon_{0} \varepsilon}\left(\frac{r^{2}}{2}+\frac{a^{3}}{r}\right), a<r<b
\end{gathered}
\]

By continuity\\
or,

\[
\begin{array}{r}
\frac{b^{3}-a^{3}}{3 \varepsilon_{0} b} \rho_{0}=B-\frac{\rho_{0}}{3 \varepsilon_{0} \varepsilon}\left(\frac{b^{2}}{2}+\frac{a^{3}}{b}\right) \\
B=\frac{\rho_{0}}{3 \varepsilon_{0} \varepsilon}\left\{\frac{\varepsilon\left(b^{3}-a^{3}\right)}{b}+\left(\frac{b^{2}}{2}+\frac{a^{3}}{b}\right)\right\}
\end{array}
\]

Finally

\[
\varphi=B-\frac{\rho_{0}}{3 \varepsilon_{0} \varepsilon}\left(\frac{a^{2}}{2}+a^{2}\right)=B-\frac{\rho_{0} a^{2}}{2 \varepsilon_{0} \varepsilon}, r<a
\]

On the basis of obtained expressions \(E(r)\) and \((\varphi)(r)\) can be plotted as shown in the answer-sheet.\\
3.78 Let the field in the dielectric be \(\vec{E}\) making an angle \(\alpha\) with \(\vec{n}\). Then we have the boundary conditions,\\
\(E_{0} \cos \alpha_{0}=\varepsilon E \cos \alpha\) and \(E_{0} \sin c_{0}=E \sin \alpha\)\\
So \(E=E_{0} \sqrt{\sin ^{2} \alpha_{0}+\frac{1}{\varepsilon^{2}} \cos ^{2} \alpha_{0}}\) and \(\tan \alpha=\varepsilon \tan \alpha_{0}\)\\
In the dielectric the normal component of the induction vector is\\
\(D_{n}=\varepsilon_{0} \varepsilon E_{n}=\varepsilon_{0} \varepsilon E \cos \alpha=\varepsilon_{0} E_{0} \cos \alpha_{0}\)\\
\(\sigma^{\prime}=P_{n}=D_{n}-\varepsilon_{0} E_{n}=\left(1-\frac{1}{\varepsilon}\right) \varepsilon_{0} E_{0} \cos \alpha_{0}\)\\
or, \(\quad \sigma^{\prime}=\frac{\varepsilon-1}{\varepsilon} \varepsilon_{0} E_{0} \cos \alpha_{0}\)\\
3.79 From the previous problem, \(\sigma^{\prime}=\varepsilon_{0} \frac{\varepsilon-1}{\varepsilon} E_{0} \cos \theta\)\\
\includegraphics[max width=\textwidth, alt={}, center]{9f3c5a8e-ac5d-42f4-8eb6-ea2af623f2d8-305_401_466_361_865}\\
\includegraphics[max width=\textwidth, alt={}, center]{9f3c5a8e-ac5d-42f4-8eb6-ea2af623f2d8-305_342_861_899_292}\\
(a) Then \(\oint \vec{E} \cdot \overrightarrow{d S}=\frac{1}{\varepsilon_{0}} Q=\pi R^{2} E_{0} \cos \theta \frac{\varepsilon-1}{\varepsilon}\)\\
(b) \(\oint \vec{D} \cdot \overrightarrow{d l}=\left(D_{1 t}-D_{2 t}\right) l=\left(\varepsilon_{0} E_{0} \sin \theta-\varepsilon \varepsilon_{0} E_{0} \sin \theta\right)=-(\varepsilon-1) \varepsilon_{0} E_{0} l \sin \theta\)\\
3.80 (a) \(\operatorname{div} \vec{D}=\frac{\partial D_{x}}{\partial x}=\rho\) and \(D=\rho l\)

\[
E_{x}=\frac{\rho l}{\varepsilon \varepsilon_{0}}, l<d \text { and } E_{x}=\frac{\rho d}{\varepsilon_{0}} \text { constant for } l>d
\]

\(\varphi(x)=-\frac{\rho l^{2}}{2 \varepsilon \varepsilon_{0}}, l<d\) and \(\varphi(x)=A-\frac{\rho l d}{\varepsilon_{0}}, l>d\) then \(\varphi(x)=\frac{\rho d}{\varepsilon_{0}}\left(d-\frac{d}{2 \varepsilon}-l\right)\), by continuity.

On the basis of obtained expressions \(E_{x}(x)\) and \(\varphi(x)\) can be plotted as shown in the figure of answersheet.\\
(b) \(\rho^{\prime}=-\operatorname{div} \vec{P}=-\operatorname{div}(\varepsilon-1) \varepsilon_{0} \vec{E}=-\rho \frac{(\varepsilon-1)}{\varepsilon}\)\\
\(\sigma^{\prime}=P_{1 n}-P_{2 n}\), where \(n\) is the normal from 1 to 2 .

\[
\begin{aligned}
& \left.=P_{1 n}, \overrightarrow{(P}_{2}=0 \text { as } 2 \text { is vacuum. }\right) \\
& =(\rho d-\rho d / \varepsilon)=. \rho d \frac{\varepsilon-1}{\varepsilon}
\end{aligned}
\]

\(3.81 \operatorname{div} \vec{D}=\frac{l}{r^{2}} \frac{\partial}{\partial r} r^{2} D_{r}=\rho\)\\
\(r^{2} D_{r}=\rho \frac{r^{3}}{3}+A \quad D_{r}=\frac{1}{3} \rho r+\frac{A}{r^{2}}, r<R\)\\
\(A=0\) as \(D_{r} \neq \infty\) at \(r=0\), Thus, \(E_{r}=\frac{\rho r}{3 \varepsilon \varepsilon_{0}}\)\\
For

\[
r>R, D_{r}=\frac{B}{r^{2}}
\]

By continuity of \(D_{r}\) at \(r=R ; B=\frac{\rho R^{3}}{3}\)\\
so,

\[
\begin{aligned}
E_{r} & =\frac{\rho R^{3}}{3 \varepsilon_{0} r^{2}}, r>R \\
\varphi & =\frac{\rho R^{3}}{3 \varepsilon_{0} r}, r>R \text { and } \varphi=-\frac{\rho r^{2}}{6 \varepsilon \varepsilon_{0}}+C, r<R \\
C & =+\frac{\rho R_{6}^{2}}{3 \varepsilon_{0}}+\frac{\rho R^{2}}{6 \varepsilon \varepsilon_{0}}, \text { by continuity of } \varphi .
\end{aligned}
\]

See answer sheet for graphs of \(E(r)\) and \(\varphi(r)\)\\
(b) \(\rho^{\prime}=\operatorname{div} \vec{P}=-\frac{1}{r^{2}} \frac{\partial}{\partial r}\left\{\frac{r^{3}}{3} \rho\left(1-\frac{1}{\varepsilon}\right)\right\}=-\frac{\rho(\varepsilon-1)}{\varepsilon}\)

\[
\sigma^{\prime}=P_{1 r}-P_{2 r}=P_{1 r}=\frac{1}{3} \rho R\left(1-\frac{1}{\varepsilon}\right)
\]

3.82 Because there is a discontinuity in polarization at the boundary of the dielectric disc, a bound surface charge appears, which is the source of the electric field inside and outside the disc.\\
We have for the electric field at the origin.

\[
\vec{E}=-\int \frac{\sigma^{\prime} d S}{4 \pi \varepsilon_{0} r^{3}} \vec{r},
\]

where \(\vec{r}=\) radius vector to the origin from the element \(d S\).\\
\(\sigma^{\prime}=P_{n}=P \cos \theta\) on the curved surface\\
( \(P_{n}=0\) on the flat surface.)\\
Here \(\theta=\) angle between \(\vec{r}\) and \(\vec{P}\)\\
By symmetry, \(\vec{E}\) will be parallel to \(\vec{P}\). Thus

\[
E \approx-\int_{0}^{2 \pi} \frac{P \cos \theta R d \theta \cdot \cos \theta}{4 \pi \varepsilon_{0} R^{2}} \cdot d
\]

\includegraphics[max width=\textwidth, alt={}, center]{9f3c5a8e-ac5d-42f4-8eb6-ea2af623f2d8-307_229_434_284_836}\\
where, \(\quad r=R\) if \(d \ll R\).\\
So, \(\quad E \approx-\frac{P d}{4 \varepsilon_{0} R}\) and \(\vec{E}=-\frac{\vec{P} d}{4 \varepsilon_{0} R}\)\\
3.83. Since there are no free extraneous charges anywhere

\[
\operatorname{div} \vec{D}=\frac{\partial D_{x}}{\partial x}=0 \text { or, } D_{x}=\text { Constant }
\]

But

\[
D_{x}=0 \text { at } \infty, \text { so, } D_{x}=0, \text { every where. }
\]

Thus,

\[
\vec{E}=-\frac{\vec{P}_{0}}{\varepsilon_{0}}\left(1-\frac{x^{2}}{d^{2}}\right) \quad \text { or, } \quad E_{x}=-\frac{P_{0}}{\varepsilon_{0}}\left(1-\frac{x^{2}}{d^{2}}\right)
\]

So,

\[
\varphi=\frac{P_{0} x}{\varepsilon_{0}}-\frac{P_{0} x^{3}}{3 \varepsilon_{0} d^{2}}+\text { constant }
\]

Hence,

\[
\varphi(+d)-\varphi(-d)=\frac{2 P_{0} d}{\varepsilon_{0}}-\frac{2 P_{0} d^{3}}{3 d^{2} \varepsilon_{0}}=\frac{4 P_{0} d}{3 \varepsilon_{0}}
\]

3.84 (a) We have \(D_{1}=D_{2}\), or, \(\varepsilon E_{2}=E_{1}\)

Also,

\[
E_{1} \frac{d}{2}+E_{2} \frac{d}{2}=E_{0} d \text { or, } E_{1}+E_{2}=2 E_{0}
\]

Hence, \(\quad E_{2}=\frac{2 E_{0}}{\varepsilon+1}\) and \(E_{1}=\frac{2 \varepsilon E_{0}}{\varepsilon+1}\) and \(D_{1}=D_{2}=\frac{2 \varepsilon \varepsilon_{0} E_{0}}{\varepsilon+1}\)\\
(b) \(D_{1}=D_{2}\), or, \(\varepsilon E_{2}=E_{1}=\frac{\sigma}{\varepsilon_{0}}=E_{0}\)

Thus,

\[
E_{1}=E_{0}, E_{2}=\frac{E_{0}}{\varepsilon} \text { and } D_{1}=D_{2}=\varepsilon_{0} E_{0}
\]

3.85 (a) Constant voltage acros the plates;

\[
E_{1}=E_{2}=E_{0}, D_{1}=\varepsilon_{0} E_{0}, D_{2}=\varepsilon_{0} \varepsilon E_{0}
\]

(b) Constant charge across the plates;

\[
\begin{gathered}
E_{1}=E_{2}, D_{1}=\varepsilon_{0} E_{1}, D_{2}=\varepsilon \varepsilon_{0} E_{2}=\varepsilon D_{1} \\
E_{1}(1+\varepsilon)=2 E_{0} \quad \text { or } \quad E_{1}=E_{2}=\frac{2 E_{0}}{\varepsilon+1}
\end{gathered}
\]

3.86 At the interface of the dielectric and vacuum,

\[
E_{1 t}=E_{2 t}
\]

The electric field must be radial and

\[
E_{1}=E_{2}=\frac{A}{\varepsilon_{0} \varepsilon r^{2}}, a<r<b
\]

Now, \(q=\frac{A}{R^{2}}\left(2 \pi R^{2}\right)+\frac{A}{\varepsilon R^{2}}\left(2 \pi R^{2}\right)\)

\[
=A\left(1+\frac{1}{\varepsilon}\right) 2 \pi
\]

\includegraphics[max width=\textwidth, alt={}, center]{9f3c5a8e-ac5d-42f4-8eb6-ea2af623f2d8-308_378_374_522_894}\\
or, \(E_{1}=E_{2}=\frac{q}{2 \pi \varepsilon_{0} r^{2}(1+\varepsilon)}\)\\
3.87 In air the forces are as shown. In \(K\)-oil, \(F \rightarrow F^{\prime}=F / \varepsilon\) and \(m g \rightarrow m g\left(1-\frac{\rho_{0}}{\rho}\right)\).\\
Since the inclinations do not change

\[
\frac{1}{\varepsilon}=1-\frac{\rho_{0}}{\rho}
\]

or, \(\frac{\rho_{0}}{\rho}=1-\frac{1}{\varepsilon}=\frac{\varepsilon-1}{\varepsilon}\)\\
\includegraphics[max width=\textwidth, alt={}, center]{9f3c5a8e-ac5d-42f4-8eb6-ea2af623f2d8-308_445_695_1013_697}\\
or, \(\quad \rho=\rho_{0} \frac{\varepsilon}{\varepsilon-1}\)\\
where \(\rho_{0}\) is the density of \(K\)-oil and \(\rho\) that of the material of which the balls are made.\\
3.88 Within the ball the electric field can be resolved into normal and tangential components.

\[
E_{n}=E \cos \theta, E_{t}=E \sin \theta
\]

Then, \(D_{n}=\varepsilon \varepsilon_{0} E \cos \theta\)\\
and \(P_{n}=(\varepsilon-1) \varepsilon_{0} E \cos \theta\)\\
or, \(\quad \sigma^{\prime}=(\varepsilon-1) \varepsilon_{0} E \cos \theta\)\\
so, \(\quad \sigma_{\max }=(\varepsilon-1) \varepsilon_{0} E\),\\
and total charge of one sign,\\
\includegraphics[max width=\textwidth, alt={}, center]{9f3c5a8e-ac5d-42f4-8eb6-ea2af623f2d8-308_361_545_1636_742}

\[
q^{\prime}=\int_{0}^{1}(\varepsilon-1) \varepsilon_{0} E \cos \theta 2 \pi R^{2} d(\cos \theta)=\pi R^{2} \varepsilon_{0}(\varepsilon-1) E
\]

(Since we are interested in the total charge of one sign we must intergrate \(\cos \theta\) from 0 to 1 only).\\
3.89 The charge is at \(A\) in the medium 1 and has an image point at \(A^{\prime}\) in the medium 2 . The electric field in the medium 1 is due to the actual charge \(q\) at \(A\) and the image charge \(q^{\prime}\) at \(A^{\prime}\). The electric field in 2 is due to a corrected charge \(q^{\prime}\) at \(A\). Thus on the boundary between 1 and 2,

\[
\begin{aligned}
& E_{1 n}=\frac{q^{\prime}}{4 \pi \varepsilon_{0} r^{2}} \cos \theta-\frac{q}{4 \pi \varepsilon_{0} r^{2}} \cos \theta \\
& E_{2 n}=\frac{-q^{\prime \prime}}{4 \pi \varepsilon_{0} r^{2}} \cos \theta \\
& E_{1 t}=\frac{q^{\prime}}{4 \pi \varepsilon_{0} r^{2}} \sin \theta+\frac{q}{4 \pi \varepsilon_{0} r^{2}} \sin \theta \\
& E_{2 t}=\frac{q^{\prime \prime}}{4 \pi \varepsilon_{0} r^{2}} \sin \theta
\end{aligned}
\]

The boundary conditions are\\
\includegraphics[max width=\textwidth, alt={}, center]{9f3c5a8e-ac5d-42f4-8eb6-ea2af623f2d8-309_488_373_525_966}

\[
\begin{gathered}
D_{1 n}=D_{2 n} \text { and } E_{1 t}=E_{2 t} \\
\varepsilon q^{\prime \prime}=q-q^{\prime} \\
q^{\prime \prime}=q+q^{\prime}
\end{gathered}
\]

So,

\[
q^{\prime \prime}=\frac{2 q}{\varepsilon+1}, q^{\prime}=-\frac{\varepsilon-1}{\varepsilon+1} q
\]

(a) The surface density of the bound charge on the surface of the dielectric

\[
\begin{aligned}
\sigma^{\prime} & =P_{2 n}=D_{2 n}-\varepsilon_{0} E_{2 n}=(\varepsilon-1) \varepsilon_{0} E_{2 n} \\
& =-\frac{\varepsilon-1}{\varepsilon+1} \frac{q}{2 \pi r^{2}} \cos \theta=-\frac{\varepsilon-1}{\varepsilon+1} \frac{q l}{2 \pi r^{3}}
\end{aligned}
\]

(b) Total bound charge is, \(-\frac{\varepsilon-1}{\varepsilon+1} q \int_{0}^{\infty} \frac{l}{2 \pi\left(l^{2}+x^{2}\right)^{3 / 2}} 2 \pi x d x=-\frac{\varepsilon-1}{\varepsilon+1} q\)\\
3.90 The force on the point charge \(q\) is due to the bound charges. This can be calculated from the field at this charge after extracting out the self field. This image field is

\[
E_{\text {image }}=\frac{\varepsilon-1}{\varepsilon+1} \frac{q}{4 \pi \varepsilon_{0}(2 l)^{2}}
\]

Thus,

\[
F=\frac{\varepsilon-1}{\varepsilon+1} \frac{q^{2}}{16 \pi \varepsilon_{0} l^{2}}
\]

\(3.91 E_{p}=\frac{q \overrightarrow{r_{1}}}{4 \pi \varepsilon_{0} r_{1}^{3}}+\frac{q^{\prime} \overrightarrow{r_{2}}}{4 \pi r_{2}^{3} \varepsilon_{0}}, P\) in 1\\
\(E_{p}=\frac{q^{\prime \prime} \overrightarrow{r_{1}}}{4 \pi \varepsilon_{0} r_{1}^{3}}, P\) in 2\\
where \(q^{\prime \prime}=\frac{2 q}{\varepsilon+1}, q^{\prime}=q^{\prime \prime}-q\)\\
In the limit \(\vec{l} \rightarrow 0\)\\
\includegraphics[max width=\textwidth, alt={}, center]{9f3c5a8e-ac5d-42f4-8eb6-ea2af623f2d8-310_456_401_75_821}\\
\(\vec{E}_{p}=\frac{\left(q+q^{\prime}\right) \vec{r}}{4 \pi \varepsilon_{0} r^{3}}=\frac{q \vec{r}}{2 \pi \varepsilon_{0}(1+\varepsilon) r^{3}}\), in either part.

Thus,

\[
\begin{gathered}
E_{p}=\frac{q}{2 \pi \varepsilon_{0}(1+\varepsilon) r^{2}} \\
\varphi=\frac{q}{2 \pi \varepsilon_{0}(1+\varepsilon) r} \\
D=\frac{q}{2 \pi \varepsilon_{0}(1+\varepsilon) r^{2}} \times\left\{\begin{array}{c}
1 \text { in vacuum } \\
\varepsilon \text { in dielectric }
\end{array}\right.
\end{gathered}
\]

\(3.92 \overrightarrow{E_{p}}=\frac{q \overrightarrow{r_{2}}}{4 \pi \varepsilon_{0} \varepsilon r_{2}^{3}}+\frac{q^{\prime} \overrightarrow{r_{1}}}{4 \pi \varepsilon_{0} r_{1}^{3}} ; P\) in 2

\[
\vec{E}_{p}=\frac{q^{\prime \prime} \vec{r}_{2}}{4 \pi \varepsilon_{0} r_{2}^{3}} ; P \text { in } 1
\]

Using the boundary conditions,

\[
E_{1 n}=\varepsilon E_{2 n}, E_{1 t}=E_{2 t}
\]

This implies\\
\(q-\varepsilon q^{\prime}=q^{\prime \prime}\) and \(q+\varepsilon q^{\prime}=\varepsilon q^{\prime \prime}\)\\
So, \(\quad q^{\prime \prime}=\frac{2 q}{\varepsilon+1}, q^{\prime}=\frac{\varepsilon-1}{\varepsilon+1} q\)\\
Then, as earlier,

\[
\sigma^{\prime}=\frac{q l}{2 \pi r^{3}} \cdot\left(\frac{\varepsilon-1}{\varepsilon+1}\right) \cdot \frac{1}{\varepsilon}
\]

\includegraphics[max width=\textwidth, alt={}, center]{9f3c5a8e-ac5d-42f4-8eb6-ea2af623f2d8-310_494_545_1389_732}\\
3.93 To calculate the electric field, first we note that an image charge will be needed to ensure that the electric field on the metal boundary is normal to the surface.

The image charge must have magnitude \(-\frac{q}{\varepsilon}\)\\
so that the tangential component of the electric field may vanish. Now,\\
\(E_{n}=\frac{1}{4 \pi \varepsilon_{0}}\left(\frac{q}{\varepsilon r^{2}}\right) 2 \cos \theta=\frac{q l}{2 \pi \varepsilon_{0} \varepsilon r^{3}}\)\\
Then \(P_{n}=D_{n}-\varepsilon_{0} E_{n}=\frac{(\varepsilon-1) q l}{2 \pi \varepsilon r^{3}}=\sigma^{\prime}\)\\
\includegraphics[max width=\textwidth, alt={}, center]{9f3c5a8e-ac5d-42f4-8eb6-ea2af623f2d8-311_403_452_63_858}

This is the density of bound charge on the surface.\\
3.94 Since the condenser plates are connected,

\[
E_{1} h+E_{2}(d-h)=0
\]

and \(\quad P+\varepsilon_{0} E_{1}=\varepsilon_{0} E_{2}\)\\
or, \(\quad E_{1}+\frac{P}{\varepsilon_{0}}=E_{2}\)\\
Thus, \(E_{2} d-\frac{P h}{\varepsilon_{0}}=0\), or, \(E_{2}=\frac{P h}{\varepsilon_{0} d}\)\\
\includegraphics[max width=\textwidth, alt={}, center]{9f3c5a8e-ac5d-42f4-8eb6-ea2af623f2d8-311_288_719_668_680}

\[
E_{1}=-\frac{P}{\varepsilon_{0}}\left(1-\frac{h}{d}\right)
\]

3.95 Given \(\vec{P}=\alpha \vec{r}\), where \(\vec{r}=\) distance from the axis. The space density of charges is given\\
by, \(\rho^{\prime}=-\operatorname{div} \vec{P}=-2 \alpha\)\\
On using.

\[
\operatorname{div} \vec{r}=\frac{1}{r} \frac{\partial}{\partial r}(\vec{r} \cdot \vec{r})=2
\]

3.96 In a uniformly charged sphere,\\
\(E_{r}=\frac{\rho_{0} r}{3 \varepsilon_{0}} \quad\) or, \(\quad \vec{E}=\frac{\rho_{0}}{3 \varepsilon_{0}} \vec{r}\)\\
The total electric field is

\[
\begin{gathered}
\vec{E}=\frac{1}{3 \varepsilon_{0}} \rho_{0} \vec{r}-\frac{1}{3 \varepsilon_{0}}(\vec{r}-\delta \vec{r}) \rho_{0} \\
=\frac{1}{3 \varepsilon_{0}} \rho_{0} \delta \vec{r}=-\frac{\vec{P}}{3 \varepsilon_{0}}
\end{gathered}
\]

where \(\rho \delta \vec{r}=-\vec{P}\) (dipole moment is defined with its direction being from the -ve charge to \(+v e\) charge.)\\
\includegraphics[max width=\textwidth, alt={}, center]{9f3c5a8e-ac5d-42f4-8eb6-ea2af623f2d8-311_400_427_1440_883}

\[
\varphi=\frac{1}{4 \pi \varepsilon_{0}}\left(\frac{Q}{r}-\frac{Q}{|\vec{r}-\delta \vec{r}|}\right)=\frac{\overrightarrow{p_{0}} \cdot \vec{r}}{4 \pi \varepsilon_{0} r^{3}}, r>R
\]

where \(\vec{p}_{0}=-\frac{4 \pi}{3} R^{3} \rho_{0} \partial \vec{r}\) is the total dipole moment.\\
3.97 The electric field \(\vec{E}_{0}\) in a spherical cavity in a uniform dielectric of permittivity \(\varepsilon\) is related to the far away field \(\vec{E}\), in the following manner. Imagine the cavity to be filled up with the dielectric. Then there will be a uniform field \(\vec{E}\) everywhere and a polarization \(\vec{P}\), given by,

\[
\vec{P}=(\varepsilon-1) \varepsilon_{0} \vec{E}
\]

Now take out the sphere making the cavity, the electric field inside the sphere will be \(-\frac{P}{3 \varepsilon_{0}}\)\\
By superposition. \(\vec{E}_{0}-\frac{\vec{P}}{3 \varepsilon_{0}}=\vec{E}\)\\
\includegraphics[max width=\textwidth, alt={}, center]{9f3c5a8e-ac5d-42f4-8eb6-ea2af623f2d8-312_302_467_282_889}\\
or, \(\overrightarrow{E_{0}}=\vec{E}+\frac{1}{3}(\varepsilon-1) \vec{E}=\frac{1}{3}(\varepsilon+2) \vec{E}\)\\
3.98 By superposition the field \(\vec{E}\) inside the ball is given by

\[
\vec{E}=\vec{E}_{0}-\frac{\vec{P}}{3 \varepsilon_{0}}
\]

On the other hand, if the sphere is not too small, the macroscopic equation \(\vec{P}=(\varepsilon-1) \varepsilon_{0} \vec{E}\) must hold. Thus,

Also

\[
\begin{gathered}
\vec{E}\left(1+\frac{1}{3}(\varepsilon-1)\right)=\vec{E}_{0} \text { or, } \vec{E}=\frac{3 \vec{E}_{0}}{\varepsilon+2} \\
\vec{P}=3 \varepsilon_{0} \frac{\varepsilon-1}{\varepsilon+2} \vec{E}_{0}
\end{gathered}
\]

3.99 This is to be handled by the same trick as in 3.96. We have effectively a two dimensional situation. For a uniform cylinder full of charge with charge density \(\rho_{0}\) (charge per unit volume), the electric field \(E\) at an inside point is along the (cylindrical) radius vector \(\vec{r}\) and equal to,

\[
\begin{aligned}
\vec{E} & =\frac{1}{2 \varepsilon_{0}} \rho \vec{r} \\
\left(\operatorname{div} \vec{E}=\frac{l}{r} \frac{\partial}{\partial r}\left(r E_{r}\right)\right. & \left.=\frac{\rho}{\varepsilon_{0}}, \quad \text { hence, } \quad E_{r}=\frac{\rho}{2 \varepsilon_{0}} r\right)
\end{aligned}
\]

Therefore the polarized cylinder can be thought of as two equal and opposite charge distributions displaced with respect to each other

\[
\vec{E}=\frac{1}{2 \varepsilon_{0}} \rho \vec{r}-\frac{1}{2 \varepsilon_{0}} \rho(\vec{r}-\delta \vec{r})=\frac{1}{2 \varepsilon_{0}} \rho \delta \vec{r}=-\frac{\vec{P}}{2 \varepsilon_{0}}
\]

Since \(\vec{P}=-\rho \delta \vec{r}\). (direction of electric dipole moment vector being from the negative charge to positive charge .)\\
3.100 As in 3.98 , we write \(\vec{E}=\overrightarrow{E_{0}}-\frac{\vec{P}}{2 \varepsilon_{0}}\) using here the result of the foregoing problem.\\
Also

\[
\vec{P}=(\varepsilon-1) \varepsilon_{0} \vec{E}
\]

So,

\[
\vec{E}\left(\frac{\varepsilon+1}{2}\right)=\vec{E}_{0}, \quad \text { or, } \quad \vec{E}=\frac{2 \vec{E}_{0}}{\varepsilon+1} \text { and } \vec{P}=2 \varepsilon_{0} \frac{\varepsilon-1}{\varepsilon+1} \vec{E}_{0}
\]

\subsection*{3.3 ELECTRIC CAPACITANCE ENERGY OF AN ELECTRIC FIELD}
3.101 Let us mentally impart a charge \(q\) on the conductor, then

\[
\begin{aligned}
\varphi_{+}-\varphi_{-} & =\int_{R_{1}}^{R_{2}} \frac{q}{4 \pi \varepsilon_{0} \varepsilon r^{2}} d r+\int_{R_{2}}^{\infty} \frac{q}{4 \pi \varepsilon_{0} r^{2}} d r \\
& =\frac{q}{4 \pi \varepsilon_{0} \varepsilon}\left[\frac{1}{R_{1}}-\frac{1}{R_{2}}\right]+\frac{q}{4 \pi \varepsilon_{0}} \frac{1}{R_{2}} \\
& =\frac{q}{4 \pi \varepsilon_{0} \varepsilon}\left[\frac{(\varepsilon-1)}{R_{2}}+\frac{1}{R_{1}}\right]
\end{aligned}
\]

\begin{center}
\includegraphics[max width=\textwidth, alt={}]{9f3c5a8e-ac5d-42f4-8eb6-ea2af623f2d8-313_326_372_293_888}
\end{center}

Hence the sought capacitance,

\[
C=\frac{q}{\varphi_{+}-\varphi_{-}}=\frac{q 4 \pi \varepsilon_{0} \varepsilon}{q\left[\frac{(\varepsilon-1)}{R_{2}}+\frac{1}{R_{1}}\right]}=\frac{4 \pi \varepsilon_{0} \varepsilon R_{1}}{(\varepsilon-1) \frac{R_{1}}{R_{2}}+1}
\]

3.102 From the symmetry of the problem, the voltage across each capacitor, \(\Delta \varphi=\xi / 2\) and charge on each capacitor \(q=C \xi / 2\) in the absence of dielectric.\\
Now when the dielectric is filled up in one of the capacitors, the equivalent capacitance of the system,

\[
C_{0}^{\prime}=\frac{C \varepsilon}{1+\varepsilon}
\]

and the potential difference across the capacitor, which is filled with dielectric,

\[
\Delta \varphi^{\prime}=\frac{q^{\prime}}{\varepsilon C}=\frac{C \varepsilon}{(1+\varepsilon)} \frac{\xi}{C \varepsilon}=\frac{\xi}{(1+\varepsilon)}
\]

But

\[
\varphi \alpha E
\]

So, as \(\varphi\) decreases \(\frac{1}{2}(1+\varepsilon)\) times, the field strength also decreases by the same factor and flow of charge, \(\Delta q=q^{\prime}-q\)\\
\(=\frac{C \varepsilon}{(1+\varepsilon)} \xi-\frac{C}{2} \xi=\frac{1}{2} C \xi \frac{(\varepsilon-1)}{(\varepsilon+1)}\)\\
3.103 (a) As it is series combination of two capacitors,

\[
\frac{1}{C}=\frac{d_{1}}{\varepsilon_{0} \varepsilon_{1} S}+\frac{d_{2}}{\varepsilon_{0} \varepsilon_{2} S} \quad \text { or, } \quad C=\frac{\varepsilon_{0} S}{\left(d_{1} / \varepsilon_{1}\right)+\left(d_{2} / \varepsilon_{2}\right)}
\]

(b) Let, \(\sigma\) be the initial surface charge density, then density of bound charge on the boundary plane.

\[
\sigma^{\prime}=\sigma\left(1-\frac{1}{\varepsilon_{1}}\right)-\sigma\left(1-\frac{1}{\varepsilon_{2}}\right)=\sigma\left(\frac{1}{\varepsilon_{2}}-\frac{1}{\varepsilon_{1}}\right)
\]

\begin{center}
\includegraphics[max width=\textwidth, alt={}]{9f3c5a8e-ac5d-42f4-8eb6-ea2af623f2d8-313_359_343_1583_917}
\end{center}

But,

\[
\sigma=\frac{q}{S}=\frac{C V}{S}=\frac{\varepsilon_{0} S \varepsilon_{1} \varepsilon_{2}}{\varepsilon_{2} d_{1}+\varepsilon_{1} d_{2}} \frac{V}{S}
\]

So,

\[
\sigma^{\prime}=\frac{\varepsilon_{0} V\left(\varepsilon_{1}-\varepsilon_{2}\right)}{\varepsilon_{2} d_{1}+\varepsilon_{1} d_{2}}
\]

3.104 (a) We point the \(x\)-axis lowards right and place the origin on the left hand side plate. The left plate is assumed to be positively charged.\\
Since \(\varepsilon\) varies linearly, we can write,

\[
\varepsilon(x)=a+b x
\]

where \(a\) and \(b\) can be determined from the boundary condition. We have \(\varepsilon=\varepsilon_{1}\) at \(x=0\) and \(\varepsilon=\varepsilon_{2}\) at \(x=d\),

Thus,

\[
\varepsilon(x)=\varepsilon_{1}+\left(\frac{\varepsilon_{2}-\varepsilon_{1}}{d}\right) x
\]

Now potential difference between the plates\\
\(\varphi_{+}-\varphi_{-}=\int_{0}^{d} \vec{E} \cdot d \vec{r}=\int_{0}^{d} \frac{\sigma}{\varepsilon_{0} \varepsilon(x)} d x\)

\[
=\int_{0}^{d} \frac{\sigma}{\varepsilon_{0}\left(\varepsilon_{1}+\frac{\varepsilon_{2}-\varepsilon_{1}}{d} x\right)} d x=\frac{\sigma d}{\left(\varepsilon_{2}-\varepsilon_{1}\right) \varepsilon_{0}} \ln \frac{\varepsilon_{2}}{\varepsilon_{1}}
\]

Hence, the sought capacitance, \(C=\frac{\sigma S}{\varphi_{+}-\varphi_{-}}=\frac{\left(\varepsilon_{2}-\varepsilon_{1}\right) \varepsilon_{0} S}{\left(\ln \varepsilon_{2} / \varepsilon_{1}\right) d}\)\\
\includegraphics[max width=\textwidth, alt={}, center]{9f3c5a8e-ac5d-42f4-8eb6-ea2af623f2d8-314_447_312_708_1070}\\
(b) \(D=\frac{q}{S}\) and \(P=\frac{q}{S}-\frac{q}{S \varepsilon(x)}\)\\
and the space density of bound charges is

\[
\rho^{\prime}=-\operatorname{div} P=-\frac{q\left(\varepsilon_{2}-\varepsilon_{1}\right)}{S d \varepsilon^{2}(x)}
\]

3.105 Let, us mentally impart a charge \(q\) to the conductor. Now potential difference between the plates,

\[
\begin{gathered}
\varphi_{+}-\varphi_{-}=\int_{R_{1}}^{R_{2}} \vec{E} \cdot d \vec{r} \\
=\int_{R_{1}}^{R_{2}} \frac{q}{4 \pi \varepsilon_{0} a / r} \frac{1}{r^{2}} d r=\frac{q}{4 \pi \varepsilon_{0} a} \ln R_{2} / R_{1}
\end{gathered}
\]

Hence, the sought capacitance,

\[
C=\frac{q}{\varphi_{+}-\varphi_{-}}=\frac{q 4 \pi \varepsilon_{0} a}{q \ln R_{2} / R_{1}}=\frac{4 \pi \varepsilon_{0} a}{\ln R_{2} / R_{1}}
\]

3.106 Let \(\lambda\) be the linear charge density then,


\begin{equation*}
E_{1 m}=\frac{\lambda}{2 \pi \varepsilon_{0} R_{1} \varepsilon_{1}} \tag{1}
\end{equation*}


and,


\begin{equation*}
E_{2 m}=\frac{\lambda}{2 \pi \varepsilon_{0} R_{2} \varepsilon_{2}} \tag{2}
\end{equation*}


The breakdown in either case will occur at the smaller value of \(r\) for a simultaneous breakdown of both dielectrics.\\
From (1) and (2)\\
\(E_{1 m} R_{1} \varepsilon_{1}=E_{2 m} R_{2} \varepsilon_{2}\), which is the sought relationship.\\
3.107 Let, \(\lambda\) be the linear charge density then, the sought potential difference,

\[
\begin{aligned}
& \varphi_{+}-\varphi_{-}=\int_{R_{1}}^{R_{2}} \frac{\lambda}{2 \pi \varepsilon_{0} \varepsilon_{1} r} d r+\int_{R_{2}}^{R_{3}} \frac{\lambda}{2 \pi \varepsilon_{0} \varepsilon_{2} r} d r \\
& =\frac{\lambda}{2 \pi \varepsilon_{0}}\left[\frac{1}{\varepsilon_{1}} \ln R_{2} / R_{1}+\frac{1}{\varepsilon_{2}} \ln R_{3} / R_{2}\right]
\end{aligned}
\]

Now, as, \(E_{1} R_{1} \varepsilon<E_{2} R_{2} \varepsilon_{2}\), so

\[
\frac{\lambda}{2 \pi \varepsilon_{0}}=E_{1} R_{1} \varepsilon_{1}
\]

is the maximum acceptable value, and for values greater than \(E_{1} R_{1} \varepsilon_{1}\), dielectric breakdows will take place,\\
Hence, the maximum potential difference between the plates,\\
\(\varphi_{+}-\varphi_{-}=E_{1} R_{1} \varepsilon_{1}\left[\frac{1}{\varepsilon_{1}} \ln R_{2} / R_{1}+\frac{1}{\varepsilon_{2}} \ln R_{3} / R_{2}\right]=E_{1} R_{1}\left[\ln R_{2} / R_{1}+\frac{\varepsilon_{1}}{\varepsilon_{2}} \ln R_{3} / R_{2}\right]\)\\
3.106 Let us suppose that linear charge density of the wires be \(\lambda\) then, the potential difference, \(\varphi_{+}-\varphi_{-}=\varphi-(-\varphi)=2 \varphi\). The intensity of the electric field created by one of the wires at a distance \(x\) from its axis can be easily found with the help of the Gauss's theorem, \(E=\frac{\lambda}{2 \pi \varepsilon_{\theta} x}\)\\
Then, \(\varphi=\int_{a}^{b-a} E d x=\frac{\lambda}{2 \pi \varepsilon_{0}} \ln \frac{b-a}{a}\)\\
\({ }^{\mathrm{t}}\) Hence, capacitance, per unit length,\\
\includegraphics[max width=\textwidth, alt={}, center]{9f3c5a8e-ac5d-42f4-8eb6-ea2af623f2d8-315_594_278_1342_971}\\
\(\frac{\lambda}{\varphi_{+}-\varphi_{-}} \cong \frac{2 \pi \varepsilon_{0}}{\ln b / a}\)\\
3.109 The field in the region between the conducting plane and the wire can be obtained by using an oppositely charged wire as an image on the other side.\\
Then the potential difference between the wire and the plane,

\[
\begin{aligned}
\Delta \varphi & =\int_{b} \vec{E} \cdot d \vec{r} \\
& =\int_{a}\left[\frac{\lambda}{2 \pi \varepsilon_{0} r}+\frac{\lambda}{2 \pi \varepsilon_{0}(2 b-r)}\right] d r \\
& =\frac{\lambda}{2 \pi \varepsilon_{0}} \ln \frac{b}{a}-\frac{\lambda}{2 \pi \varepsilon_{0}} \ln \frac{b}{2 b-a} \\
& =\frac{\lambda}{2 \pi \varepsilon_{0}} \ln \frac{2 b-a}{a} \\
& \cong \frac{\lambda}{2 \pi \varepsilon_{0}} \ln \frac{2 b}{a}, \text { as } b \gg a
\end{aligned}
\]

\begin{center}
\includegraphics[max width=\textwidth, alt={}]{9f3c5a8e-ac5d-42f4-8eb6-ea2af623f2d8-316_484_610_219_782}
\end{center}

Hence, the sought mutual capacitance of the system per unit length of the wire

\[
=\frac{\lambda}{\Delta \varphi}=\frac{2 \pi \varepsilon_{0}}{\ln 2 b / a}
\]

3.110 When \(b \gg a\), the charge distribution on each spherical conductor is practically unaffected by the presence of the other conductor. Then, the potential \(\varphi_{+}\left(\varphi_{-}\right)\)on the positive (respectively negative) charged conductor is \(+\frac{q}{4 \pi \varepsilon_{0} \varepsilon a}\left(-\frac{q}{4 \pi \varepsilon_{0} \varepsilon a}\right)\)\\
Thus \(\varphi_{+}-\varphi_{-}=\frac{q}{2 \pi \varepsilon \varepsilon_{0} a}\)\\
and \(C=\frac{q}{\varphi_{+}-\varphi_{-}}=2 \pi \varepsilon_{0} \varepsilon a\).\\
\includegraphics[max width=\textwidth, alt={}, center]{9f3c5a8e-ac5d-42f4-8eb6-ea2af623f2d8-316_291_607_1077_761}

Note : if we require terms which depend on \(\frac{a}{b}\), we have to take account of distribution of charge on the conductors.\\
3.111 As in 3.109 we apply the method of image. Then the potentical difference between the +vely charged sphere and the conducting plane is one half the nominal potential difference between the sphere and its image and is\\
\(\Delta \psi=\frac{1}{2}\left(\psi_{+}-\psi_{-}\right) \cong \frac{q}{4 \pi \varepsilon_{0} a}\)\\
Thus

\[
C=\frac{q}{\Delta \varphi}=4 \pi \varepsilon_{0} a . \text { for } l \gg a .
\]

\includegraphics[max width=\textwidth, alt={}, center]{9f3c5a8e-ac5d-42f4-8eb6-ea2af623f2d8-316_401_614_1594_797}\\
\includegraphics[max width=\textwidth, alt={}, center]{9f3c5a8e-ac5d-42f4-8eb6-ea2af623f2d8-317_288_1155_218_157}\\
(a) Since \(\varphi_{1}=\varphi_{B}\) and \(\varphi_{2}=\varphi_{A}\)

The arrangement of capacitors shown in the problem is equivalent to the arrangement shown in the Fig.\\
\includegraphics[max width=\textwidth, alt={}, center]{9f3c5a8e-ac5d-42f4-8eb6-ea2af623f2d8-317_429_378_724_195}\\
\includegraphics[max width=\textwidth, alt={}, center]{9f3c5a8e-ac5d-42f4-8eb6-ea2af623f2d8-317_73_95_908_684}\\
\includegraphics[alt={}]{smile-564219e94ed88610824606fd5449681aff525002}\\
and hence the capacitance between \(A\) and \(B\) is,

\[
C=C_{1}+C_{2}+C_{3}
\]

(B) From the symmetry of the problem, there is no P.d. between D and E.. So, the combination reduces to a simple arrangement shown in the Fig and hence the net capacitance,

\[
C_{0}=\frac{C}{2}+\frac{C}{2}=C
\]

3.113 (a) In the given arrangement, we have three capacitors of equal capacitance \(C=\frac{\varepsilon_{0} S}{d}\) and the first and third plates are at the same potential.

Hence, we can resolve the network into a simple form using series and parallel grouping of\\
\includegraphics[max width=\textwidth, alt={}, center]{9f3c5a8e-ac5d-42f4-8eb6-ea2af623f2d8-317_318_488_1485_875}\\
capacitors, as shown in the figure.\\
Thus the equivalent capacitance

\[
C_{0}=\frac{(C+C) C}{(C+C)+C}=\frac{2}{3} C
\]

(b) Let us mentally impart the charges \(+q\) and \(-q\) to the plates 1 and 2 and then distribute them to other plates using charge conservation and electric induction. (Fig.).\\
As the potential difference between the plates 1 and 2 is zero,

\[
\begin{aligned}
& -\frac{q_{1}}{C}+\frac{q_{2}}{C}-\frac{q_{1}}{C}=0, \quad\left(\text { where } C=\frac{\varepsilon_{0} S}{d}\right) \\
& \text { or, } \quad q_{2}=2 \cdot q_{1}
\end{aligned}
\]

The potential difference between \(A\) and \(B\),

\[
\varphi=\varphi_{A}-\varphi_{B}=q_{2} / C
\]

Hence the sought capacitane,

\[
C_{0}=\frac{q}{\varphi}=\frac{q_{1}+q_{2}}{q_{2} / C}=\frac{3 q_{1}}{2 q_{1} / C}=\frac{3}{2} C=\frac{3 \varepsilon_{0} S}{2 d}
\]

3.114 Amount of charge, that the capacitor of capacitance \(C_{1}\) can withstand, \(q_{1}=C_{1} V_{1}\) and similarly the charge, that the capacitor of capacitance \(C_{2}\) can withstand, \(q_{2}=C_{2} V_{2}\). But in series combination, charge on both the capacitors will be same, so, \(q_{\text {max }}\), that the combination can withstand = \(C_{1} V_{1}\),\\
as \(C_{1} V_{1}<C_{2} V_{2}\), from the numerical data, given.\\
Now, net capacitance of the system,

\[
C_{0}=\frac{C_{1} C_{2}}{C_{1}+C_{2}}
\]

and hence, \(V_{\max }=\frac{q_{\max }}{C_{0}}=\frac{C_{1} V_{1}}{C_{1} C_{2} / C_{1}+C_{2}}=V_{1}\left(1+\frac{C_{1}}{C_{2}}\right)=9 \mathrm{kV}\)\\
3.115 Let us distribute the charges, as shown in the figure.

Now, we know that in a closed circuit, \(-\Delta \varphi=0\)\\
-So, in the loop, \(D C F E D\),


\begin{equation*}
\frac{q_{1}}{C_{1}}-\frac{q_{2}}{C_{1}}-\frac{q_{2}}{C_{2}}=\xi \quad \text { or, } q_{1}=C_{1}\left[\xi+q_{2}\left(\frac{1}{C_{1}}+\frac{1}{C_{2}}\right)\right] \tag{1}
\end{equation*}


Again in the loop DGHED,


\begin{equation*}
\frac{q_{1}}{C_{1}}+\frac{q_{1}+q_{2}}{C_{2}}=\xi \tag{2}
\end{equation*}


\begin{center}
\includegraphics[max width=\textwidth, alt={}]{9f3c5a8e-ac5d-42f4-8eb6-ea2af623f2d8-318_328_780_1638_314}
\end{center}

Using Eqs. (1) and (2), we get

\[
q_{2}\left[\frac{1}{C_{1}}+\frac{3}{C_{2}}+\frac{C_{1}}{C_{2}}\right]=-\frac{\xi C_{1}}{C_{2}}
\]

Now,

\[
\varphi_{A}-\varphi_{B}=\frac{-q_{2}}{C_{2}}=\frac{\xi}{C_{2}^{2} / C_{1}} \frac{1}{\left[\frac{1}{C_{1}}+\frac{3}{C_{2}}+\frac{C_{1}}{C_{2}^{2}}\right]}
\]

or,

\[
\varphi_{A}-\varphi_{B}=\frac{\xi}{\left[\frac{C_{2}{ }^{2}}{C_{1}{ }^{2}}+\frac{3 C_{2}}{C_{1}}+1\right]}=\frac{\xi}{\eta^{2}+3 \eta+1}=10 \mathrm{~V}
\]

3.116 The infinite circuit, may be reduced to the circuit, shown in the Fig. where, \(C_{0}\) is the net capacitance of the combination.

So,

\[
\frac{1}{C+C_{0}}+\frac{1}{C}=\frac{1}{C_{0}}
\]

Solving the quadratic,

\[
C C_{0}+\dot{C}_{0}^{2}-C^{2}=0,
\]

we get,\\
\includegraphics[max width=\textwidth, alt={}, center]{9f3c5a8e-ac5d-42f4-8eb6-ea2af623f2d8-319_233_450_741_864}\\
\(C_{0}=\frac{(\sqrt{5}-1)}{2} C\), taking only \(+v e\) value as \(C_{0}\) can not be negative.\\
3.117 Let, us make the charge distribution, as shown in the figure.

Now,

\[
\varphi_{A}-\varphi_{B}=\frac{q}{C_{1}}-\xi+\frac{q}{C_{2}}
\]

or, \(\quad q=\frac{\left(\varphi_{A}-\varphi_{B}\right)+\xi}{C_{1}+C_{2}} C_{1} C_{2}\)\\
\includegraphics[max width=\textwidth, alt={}, center]{9f3c5a8e-ac5d-42f4-8eb6-ea2af623f2d8-319_123_553_1260_668}

Hence, voltage across the capacitor \(C_{1}\)

\[
=\frac{q}{C_{1}}=\frac{\left(\varphi_{A}-\varphi_{B}\right)+\xi}{C_{1}+C_{2}} C_{2}=10 \mathrm{~V}
\]

and voltage across the capacitor, \(C_{2}\)

\[
=\frac{q}{C_{2}}=\frac{\left(\varphi_{A}-\varphi_{B}\right)+\xi}{C_{1}+C_{2}} C_{1}=5 \mathrm{~V}
\]

3.118 Let \(\xi_{2}>\xi_{1}\), then using \(-\Delta \varphi=0\) in the closed circuit, (Fig.)

\[
\begin{gathered}
\frac{-q}{C_{1}}+\xi_{2}-\frac{q}{C_{2}}-\xi_{1}=0 \\
q=\frac{\left(\xi_{2}-\xi_{1}\right) C_{1} C_{2}}{\left(C_{1}+C_{2}\right)}
\end{gathered}
\]

\begin{center}
\includegraphics[max width=\textwidth, alt={}]{9f3c5a8e-ac5d-42f4-8eb6-ea2af623f2d8-319_428_495_1683_854}
\end{center}

Hence the P.D. accross the left and right plates of capacitors,

\[
\varphi_{1}=\frac{q}{C_{1}}=\frac{\left(\xi_{2}-\xi_{1}\right) C_{2}}{C_{1}+C_{2}}
\]

and similarly\\
\(\varphi_{2}=\frac{-q}{C_{2}}=\frac{\left(\xi_{1}-\xi_{2}\right) C_{1}}{C_{1}+C_{2}}\)\\
3.119 Taking benefit of the foregoing problem, the amount of charge on each capacitor

\[
|q|=\frac{\left|\xi_{2}-\xi_{1}\right| C_{1} C_{2}}{C_{1}+C_{2}}
\]

3.120 Make the charge distribution, as shown in the figure. In the circuit, 12561.

\[
\begin{aligned}
& -\Delta \varphi=0 \text { yields } \\
& \frac{q_{1}}{C_{4}}+\frac{q_{1}}{C_{3}}-\xi=0 \quad \text { or, } \quad q_{1}=\frac{\xi C_{3} C_{4}}{C_{3}+C_{4}}
\end{aligned}
\]

and in the circuit 13461,\\
\(\frac{q_{2}}{C_{2}}+\frac{q_{2}}{C_{1}}-\xi=0 \quad\) or, \(\quad q_{2}=\frac{\xi C_{1} C_{2}}{C_{1}+C_{2}}\)\\
Now \(\quad \varphi_{A}-\varphi_{B}=\frac{q_{2}}{C_{1}}-\frac{q_{1}}{C_{3}}\)\\
\includegraphics[max width=\textwidth, alt={}, center]{9f3c5a8e-ac5d-42f4-8eb6-ea2af623f2d8-320_435_598_736_747}

\[
=\xi\left[\frac{C_{2}}{C_{1}+C_{2}}-\frac{C_{4}}{C_{3}+C_{4}}\right]=\xi\left[\frac{C_{2} C_{3}-C_{1} C_{4}}{\left(C_{1}+C_{2}\right)\left(C_{3}+C_{4}\right)}\right]
\]

It becomes zero, when

\[
\left(C_{2} C_{3}-C_{1} C_{4}\right)=0 . \quad \text { or } \quad \frac{C_{1}}{C_{2}}=\frac{C_{3}}{C_{4}}
\]

3.121 Let, the charge \(q\) flows through the connecting wires, then at the state of equilibrium, charge distribution will be as shown in the Fig. In the closed circuit 12341, using \(-\Delta \varphi=0\)\\
\(-\frac{\left(C_{1} V-q\right)}{C_{1}}+\frac{q}{C_{2}}+\frac{q}{C_{3}}=0\)\\
or, \(q=\frac{V}{\left(1 / C_{1}+1 / C_{2}+1 / C_{3}\right)}=0.06 \mathrm{mC}\)\\
3.122 Initially, charge on the capacitor \(C_{1}\) or \(C_{2}\),\\
\includegraphics[max width=\textwidth, alt={}, center]{9f3c5a8e-ac5d-42f4-8eb6-ea2af623f2d8-320_301_658_1607_744}\\
\(q=\frac{\xi C_{1} C_{2}}{C_{1}+C_{2}}\), as they are in series combination (Fig.-a)\\
when the switch is closed, in the circuit CDEFC' from \(-\Delta \varphi=0\), (Fig. b )


\begin{equation*}
\xi-\frac{q_{2}}{C_{2}}=0 \text { or } q_{2}=C_{2} \xi \tag{1}
\end{equation*}


And in the closed loop BCFAB from \(-\Delta \varphi=0\)


\begin{equation*}
\frac{-q_{1}}{C_{1}}+\frac{q_{2}}{C_{2}}-\xi=0 \tag{2}
\end{equation*}


\begin{figure}[h]
\begin{center}
  \includegraphics[alt={},max width=\textwidth]{9f3c5a8e-ac5d-42f4-8eb6-ea2af623f2d8-321_370_622_364_64}
\captionsetup{labelformat=empty}
\caption{(a)}
\end{center}
\end{figure}

\begin{figure}[h]
\begin{center}
  \includegraphics[alt={},max width=\textwidth]{9f3c5a8e-ac5d-42f4-8eb6-ea2af623f2d8-321_364_606_370_727}
\captionsetup{labelformat=empty}
\caption{(b)}
\end{center}
\end{figure}

From (1) and (2) \(q_{1}=0\)\\
Now, charge flown through section \(1=\left(q_{1}+q_{2}\right)-0=C_{2} \xi\)\\
and charge flown through section \(2=-q_{1}-q=-\frac{\xi C_{1} C_{2}}{C_{1}+C_{2}}\)\\
3.123 When the switch is open, (Fig-a)

\[
q_{0}=\frac{2 \xi C_{1} C_{2}}{C_{1}+C_{2}}
\]

\includegraphics[max width=\textwidth, alt={}, center]{9f3c5a8e-ac5d-42f4-8eb6-ea2af623f2d8-321_390_573_1190_48}\\
\includegraphics[max width=\textwidth, alt={}, center]{9f3c5a8e-ac5d-42f4-8eb6-ea2af623f2d8-321_395_574_1193_782}\\
and when the switch is closed,

\[
q_{1}=\xi C_{1} \text { and } q_{2}=\xi C_{2}
\]

Hence, the flow of charge, due to the shortening of switch,\\
through section \(1=q_{1}-q_{0}=\xi C_{1}\left[\frac{C_{1}-C_{2}}{C_{1}+C_{2}}\right]=-24 \mu \mathrm{C}\)\\
through the section \(2=-q_{2}-\left(q_{0}\right)=\xi C_{2}\left[\frac{C_{1}-C_{2}}{C_{1}+C_{2}}\right]=-36 \mu \mathrm{C}\)\\
and through the section \(3=q_{2}-\left(q_{2}-q_{1}\right)-0=\xi\left(C_{2}-C_{1}\right)=-60 \mu \mathrm{C}\)\\
3.124 First of all, make the charge distribution, as shown in the figure.

In the loop 12341, using \(-\Delta \varphi=0\)


\begin{equation*}
\frac{q_{1}}{C_{1}}-\xi_{1}+\frac{q_{1}-q_{2}}{C_{3}}=0 \tag{1}
\end{equation*}


Similarly, in the loop 61456, using \(-\Delta \varphi=0\)


\begin{equation*}
\frac{q_{2}}{C_{2}}+\frac{q_{2}-q_{1}}{C_{3}}-\xi_{2}=0 \tag{2}
\end{equation*}


From Eqs. (1) and (2) we have

\[
q_{2}-q_{1}=\frac{\xi_{2} C_{2}-\xi_{1} C_{1}}{\frac{C_{2}}{C_{3}}+\frac{C_{1}}{C_{3}}+1}
\]

\begin{center}
\includegraphics[max width=\textwidth, alt={}]{9f3c5a8e-ac5d-42f4-8eb6-ea2af623f2d8-322_431_542_187_796}
\end{center}

Hence,

\[
\varphi_{A}-\varphi_{B}=\frac{q_{2}-q_{1}}{C_{3}}=\frac{\xi_{2} C_{2}-\xi_{1} C_{1}}{C_{1}+C_{2}+C_{3}}
\]

3.125 In the loop ABDEA, using \(-\Delta \varphi=0\)


\begin{equation*}
-\xi_{3}+\frac{q_{1}}{C_{3}}+\frac{q_{1}+q_{2}}{C_{1}}+\xi_{1}=0 \tag{1}
\end{equation*}


Similarly in the loop ODEF, O


\begin{equation*}
\frac{q_{1}+q_{2}}{C_{1}}+\xi_{1}-\xi_{2}+\frac{q_{2}}{C_{2}}=0 \tag{2}
\end{equation*}


Solving Eqs. (1) and (2), we get,\\
\(q_{1}+q_{2}=\frac{\xi_{2} C_{2}-\xi_{1} C_{2}-\xi_{1} C_{3}+\xi_{3} C_{3}}{\frac{C_{3}}{C_{1}}+\frac{C_{2}}{C_{1}}+1}\)\\
\includegraphics[max width=\textwidth, alt={}, center]{9f3c5a8e-ac5d-42f4-8eb6-ea2af623f2d8-322_513_597_783_781}

Now, \(\varphi_{1}-\varphi_{0}=\varphi_{1}=-\frac{\left(q_{1}+q_{2}\right)}{C_{1}}\), as ( \(\varphi_{0}=0\) )

\[
=\frac{\xi_{1}\left(C_{2}+C_{3}\right)-\xi_{2} C_{2}-\xi_{3} C_{3}}{C_{1}+C_{2}+C_{3}}
\]

And using the symmetry, \(\varphi_{2}=\frac{\xi_{2}\left(C_{1}+C_{3}\right)-\xi_{1} C_{1}-\xi_{3} C_{3}}{C_{1}+C_{2}+C_{3}}\)\\
and

\[
\varphi_{3}=\frac{\xi_{3}\left(C_{1}+C_{2}\right)-\xi_{1} C_{1}-\xi_{2} C_{2}}{C_{1}+C_{2}+C_{3}}
\]

The answers have wrong sign in the book.\\
3.126 Taking the advantage of symmetry of the problem charge distribution may be made, as shown in the figure.\\
In the loop, 12561, \(-\Delta \varphi=0\)\\
or \(\frac{q_{2}}{C_{2}}+\frac{q_{2}-q_{1}}{C_{3}}-\frac{q_{1}}{C_{1}}=0\)\\
or \(\frac{q_{1}}{q_{2}}=\frac{C_{1}\left(C_{3}+C_{2}\right)}{C_{2}\left(C_{1}+C_{3}\right)}\)

Now, capacitance of the network,


\begin{align*}
& C_{0}=\frac{q_{1}+q_{2}}{\varphi_{A}-\varphi_{B}}=\frac{q_{1}+q_{2}}{q_{2} / C_{2}+q_{1} / C_{1}} \\
& =\frac{\left(1+q_{1} / q_{2}\right)}{\left(\frac{1}{C_{2}}+\frac{q_{1}}{q_{2} C_{1}}\right)} \tag{2}
\end{align*}


\begin{center}
\includegraphics[max width=\textwidth, alt={}]{9f3c5a8e-ac5d-42f4-8eb6-ea2af623f2d8-323_433_741_158_635}
\end{center}

From Eqs. (1) and (2)

\[
C_{0}=\frac{2 C_{1} C_{2}+C_{3}\left(C_{1}+C_{2}\right)}{C_{1}+C_{2}+2 C_{3}}
\]

3.127 (a) Interaction energy of any two point charges \(q_{1}\) and \(q_{2}\) is given by \(\frac{q_{1} q_{2}}{4 \pi \varepsilon_{0} r}\) where \(r\) is the separation between the charges.\\
\includegraphics[max width=\textwidth, alt={}, center]{9f3c5a8e-ac5d-42f4-8eb6-ea2af623f2d8-323_391_397_1063_168}\\
\includegraphics[max width=\textwidth, alt={}, center]{9f3c5a8e-ac5d-42f4-8eb6-ea2af623f2d8-323_393_314_1063_635}\\
\includegraphics[max width=\textwidth, alt={}, center]{9f3c5a8e-ac5d-42f4-8eb6-ea2af623f2d8-323_381_308_1065_1030}

Hence, interaction energy of the system,

\[
\begin{aligned}
& U_{a}=4 \frac{q^{2}}{4 \pi \varepsilon_{0} a}+2 \frac{q^{2}}{4 \pi \varepsilon_{0}(\sqrt{2} a)} \\
& U_{b}=4 \frac{-q^{2}}{4 \pi \varepsilon_{0} a}+2 \frac{q^{2}}{4 \pi \varepsilon_{0}(\sqrt{2} a)}
\end{aligned}
\]

and

\[
U_{c}=2 \frac{q^{2}}{4 \pi \varepsilon_{0} a}-\frac{2 q^{2}}{4 \pi \varepsilon_{0} a}-\frac{2 q^{2}}{4 \pi \varepsilon_{0}(\sqrt{2} a)}=-\frac{\sqrt{2} q^{2}}{4 \pi \varepsilon_{0} a}
\]

3.128 As the chain is of infinite length any two charge of same sign will occur symmetrically to any other charge of opposite sign.\\
So, interaction energy of each charge with all the others,


\begin{equation*}
U=-2 \frac{q^{2}}{4 \pi \varepsilon_{0} a}\left[1-\frac{1}{2}+\frac{1}{3}-\frac{1}{4}+\ldots \ldots \ldots \text { up to } \infty\right] \tag{1}
\end{equation*}


But

\[
\ln (1+x)=x-\frac{1}{2} x^{2}+\frac{1}{3} x^{3} \ldots \ldots \ldots \text { up to } \infty
\]

and putting \(x=1\) we get \(\ln 2=1-\frac{1}{2}+\frac{1}{3}+\ldots \ldots \ldots\) up to \(\infty\)

From Eqs. (1) and (2),

\[
U=\frac{-2 q^{2} \ln 2}{4 \pi \varepsilon_{0} a}
\]

3.129 Using electrical image method, interaction energy of the charge \(q\) with those induced on the plane.

\[
U=\frac{-q^{2}}{4 \pi \varepsilon_{0}(2 l)}=-\frac{q^{2}}{8 \pi \varepsilon_{0} l}
\]

3.130 Consider the interaction energy of one of the balls (say 1) and thin spherical shell of the other. This interaction energy can be written as \(\int d \varphi q\)

\[
\begin{aligned}
& =\int \frac{q_{1}}{4 \pi \varepsilon_{0} R} \rho_{2}(r) 2 \pi r^{2} \sin \theta d \theta d r=\int_{0}^{\pi} \frac{\rho_{2}(r) q_{1} r^{2} \sin \theta d \theta d r}{2 \varepsilon_{0}\left(l^{2}+r^{2}+2 l r \cos \theta\right)^{1 / 2}} \\
& \quad=\frac{q_{1} r}{2 \varepsilon_{0} l} d r \int_{l-r}^{l+r} d x \rho_{2}(r) \\
& \quad=\frac{q_{1} r}{2 \varepsilon_{0} l} d r \cdot 2 r \rho_{2}(r) \cdot 2 \\
& \quad=\frac{q_{1}}{4 \pi \varepsilon_{0} l} 4 \pi r^{2} d r \rho_{2}(r) \\
& \text { Hence finally integrating }
\end{aligned}
\]

\(U_{\text {int }}=\frac{q_{1} q_{2}}{4 \pi \varepsilon_{0} l}\) where, \(q_{2}=\int_{0}^{\infty} 4 \pi r^{2} \rho_{2}(r) d r\)\\
3.131 Charge contained in the capacitor of capacitance \(C_{1}\) is \(q=C_{1} \varphi\) and the energy, stored in it :

\[
U_{i}=\frac{q^{2}}{2 C_{1}}=\frac{1}{2} C_{1} \varphi^{2}
\]

Now, when the capacitors are connected in parallel, equivalent capacitance of the system, \(C=C_{1}+C_{2}\) and hence, energy stored in the system :\\
\(U_{f}=\frac{C_{1}^{2} \varphi^{2}}{2\left(C_{1}+C_{2}\right)}\), as charge remains conserved during the process.\\
So, increment in the energy,

\[
\Delta U=\frac{C_{1}^{2} \varphi^{2}}{2}\left(\frac{1}{C_{1}+C_{2}}-\frac{1}{C_{1}}\right)=\frac{-C_{2} C_{1} \varphi^{2}}{2\left(C_{1}+C_{2}\right)}=-0.03 \mathrm{~mJ}
\]

3.132 The charge on the condensers in position 1 are as shown. Here

\[
\frac{q}{C}=\frac{q_{0}}{C_{0}}=\frac{q+q_{0}}{C+C_{0}}
\]

and

\[
\left(q+q_{0}\right)\left(\frac{1}{C+C_{0}}+\frac{1}{C}\right)=\xi \text { or, } q+q_{0}=\frac{C\left(C+C_{0}\right) \xi}{C_{0}+2 C}
\]

Hence,

\[
q=\frac{C^{2} \xi}{C_{0}+2 C} \text { and } q_{0}=\frac{C C_{0} \xi}{C_{0}+2 C}
\]

\includegraphics[max width=\textwidth, alt={}, center]{9f3c5a8e-ac5d-42f4-8eb6-ea2af623f2d8-325_438_693_823_25}\\
\includegraphics[max width=\textwidth, alt={}, center]{9f3c5a8e-ac5d-42f4-8eb6-ea2af623f2d8-325_437_596_824_749}

After the switch is thrown to position 2, the charges change as shown in (Fig-b). A charge \(q_{0}\) has flown in the right loop through the two condensers and a charge \(q_{0}\) through the cell, Because of the symmetry of the problem there is no change in the energy stored in the condensers. Thus\\
\(H\) (Heat produced) \(=\) Energy delivered by the cell

\[
=\Delta q \xi=q_{0} \xi=\frac{C C_{0} \xi^{2}}{C_{0}+2 C}
\]

3.133 Initially, the charge on the right plate of the capacitor, \(q=C\left(\xi_{1}-\xi_{2}\right)\) and finally, when switched to the position, 2 . charge on the same plate of capacitor ;

So,

\[
\begin{gathered}
q^{\prime}=C \xi_{1} \\
\Delta q=q^{\prime}=C \xi_{2}
\end{gathered}
\]

Now, from energy conservation,\\
\(\Delta U+\) Heat liberated \(=A_{\text {cell }}\), where \(\Delta U\) is the electrical energy.

\[
\frac{1}{2} C \xi_{1}^{2}-\frac{1}{2} C\left(\xi_{1}-\xi_{2}\right)^{2}+\text { Heat liberated }=\Delta q \xi_{1}
\]

as only the cell with e.m.f. \(\xi_{1}\) is responsible for redistribution of the charge. So,

\[
C \xi_{1} \xi_{2}-\frac{1}{2} C \xi_{2}^{2}+\text { Heat liberated }=C \xi_{2} \xi_{1}
\]

Hence heat liberated \(=\frac{1}{2} C \xi_{2}^{2}\)\\
3.134 Self energy of each shell is given by \(\frac{q \varphi}{2}\), where \(\varphi\) is the potential of the shell, created only by the charge \(q\), on it.\\
Hence, self energy of the shells 1 and 2 are :

\[
W_{1}=\frac{q_{1}^{2}}{8 \pi \varepsilon_{0} R_{1}} \text { and } W_{2}=\frac{q_{2}^{2}}{8 \pi \varepsilon_{0} R_{2}}
\]

The interaction energy between the charged shells equals charge \(q\) of one shell, multiplied by the potential \(\varphi\), created by other shell, at the point of location of charge \(q\).

So,

\[
W_{12}=q_{1} \frac{q_{2}}{4 \pi \varepsilon_{0} R_{2}}=\frac{q_{1} q_{2}}{4 \pi \varepsilon_{0} R_{2}}
\]

Hence, total enegy of the system,

\[
\begin{gathered}
U=W_{1}+W_{2}+W_{12} \\
=\frac{1}{4 \pi \varepsilon_{0}}\left[\frac{q_{1}^{2}}{2 R_{1}}+\frac{q_{2}^{2}}{2 R_{2}}+\frac{q_{1} q_{2}}{R_{2}}\right]
\end{gathered}
\]

3.135 Electric fields inside and outside the sphere with the help of Gauss theorem :

\[
E_{1}=\frac{q r}{4 \pi \varepsilon_{0} R^{2}}(r \leq R), E_{2}=\frac{q}{4 \pi \varepsilon_{0}} \frac{1}{r^{2}}(r>R)
\]

Sought self energy of the ball

\[
U=W_{1}+W_{2}
\]

\[
\begin{gathered}
=\int_{0}^{R} \frac{\varepsilon_{0} E_{1}^{2}}{2} 4 \pi r^{2} d r+\int_{R}^{\infty} \frac{\varepsilon_{0} E_{2}^{2}}{2} 4 \pi r^{2} d r=\frac{q^{2}}{8 \pi \varepsilon_{0} R}\left(\frac{1}{5}+1\right) \\
U=\frac{3 q^{2}}{4 \pi \varepsilon_{0} 5 R} \text { and } \frac{W_{1}}{W_{2}}=\frac{1}{5}
\end{gathered}
\]

Hence,\\
3.136 (a) By the expression \(\int \frac{1}{2} \varepsilon_{0} \varepsilon E^{2} d V=\int \frac{1}{2} \varepsilon \varepsilon_{0} E^{2} 4 \pi r^{2} d r\), for a spherical layer.

To find the electrostatic energy inside the dielectric layer, we have to integrate the upper expression in the limit \([a, b]\)

\[
U=\frac{1}{2} \varepsilon_{0} \varepsilon \int_{a}^{b}\left(\frac{q}{4 \pi \varepsilon_{0} \varepsilon r^{2}}\right)^{2} 4 \pi r^{2} d r=\frac{q^{2}}{8 \pi \varepsilon_{0} \varepsilon}\left[\frac{1}{a}-\frac{1}{b}\right]=27 \mathrm{~mJ}
\]

3.137 As the field is conservative total work done by the field force,

\[
\begin{gathered}
A_{f d}=U_{t}-U_{f}=\frac{1}{2} q\left(\varphi_{1}-\varphi_{2}\right) \\
=\frac{1}{2} \frac{q^{2}}{4 \pi \varepsilon_{0}}\left[\frac{1}{R_{1}}-\frac{1}{R_{2}}\right]=\frac{q^{2}}{8 \pi \varepsilon_{0}}\left[\frac{1}{R_{1}}-\frac{1}{R_{2}}\right]
\end{gathered}
\]

3.138 Initially, energy of the system, \(U_{i}=W_{1}+W_{12}\) where, \(W_{1}\) is the self energy and \(W_{12}\) is the mutual energy.

So,

\[
U_{i}=\frac{1}{2} \frac{q^{2}}{4 \pi \varepsilon_{0} R_{1}}+\frac{q q_{0}}{4 \pi \varepsilon_{0} R_{1}}
\]

and on expansion, energy of the system,

\[
\begin{aligned}
U_{f} & =W_{1}^{\prime}+W_{12}^{\prime} \\
& =\frac{1}{2} \frac{q^{2}}{4 \pi \varepsilon_{0} R_{2}}+\frac{q q_{0}}{4 \pi \varepsilon_{0} R_{2}}
\end{aligned}
\]

\begin{center}
\includegraphics[max width=\textwidth, alt={}]{9f3c5a8e-ac5d-42f4-8eb6-ea2af623f2d8-327_350_341_481_889}
\end{center}

Now, work done by the field force, \(A\) equals the decrement in the electrical energy, i.e.

\[
A=\left(U_{\imath}-U_{f}\right)=\frac{q\left(q_{0}+q / 2\right)}{4 \pi \varepsilon_{0}}\left(\frac{1}{R_{1}}-\frac{1}{R_{2}}\right)
\]

Alternate : The work of electric forces is equal to the decrease in electric energy of the system,

\[
A=U_{i}-U_{f}
\]

In order to find the difference \(U_{i}-U_{\rho}\) we note that upon expansion of the shell, the electric field and hence the energy localized in it, changed only in the hatched spherical layer consequently (Fig.).

\[
U_{i}-U_{f}=\int_{R_{1}}^{R_{2}} \frac{\varepsilon_{0}}{2}\left(E_{1}^{2}-E_{2}^{2}\right) \cdot 4 \pi r^{2} d r
\]

where \(E_{1}\) and \(E_{2}\) are the field intensities (in the hatched region at a distance \(r\) from the centre of the system) before and after the expansion of the shell. By using Gauss' theorem, we find

\[
E_{1}=\frac{1}{4 \pi \varepsilon_{0}} \frac{q+q_{0}}{r^{2}} \text { and } E_{2}=\frac{1}{4 \pi \varepsilon_{0}} \frac{q_{0}}{r^{2}}
\]

As a result of integration, we obtain

\[
A=\frac{q\left(q_{0}+q / 2\right)}{4 \pi \varepsilon_{0}}\left(\frac{1}{R_{1}}-\frac{1}{R_{2}}\right)
\]

3.139 Energy of the charged sphere of radius \(r\), from the equation

\[
U=\frac{1}{2} q \varphi=\frac{1}{2} q \frac{q}{4 \pi \varepsilon_{0} r}=\frac{q^{2}}{8 \pi \varepsilon_{0} r}
\]

If the radius of the shell changes by \(d r\) then work done is

\[
4 \pi r^{2} F_{u} d r=-d U=q^{2} / 8 \pi \varepsilon_{0} r^{2}
\]

Thus sought force per unit area,

\[
F_{u}=\frac{q^{2}}{4 \pi r^{2}\left(8 \pi \varepsilon_{0} r^{2}\right)}=\frac{\left(4 \pi r^{2} \sigma\right)^{2}}{4 \pi r^{2} \times 8 \pi \varepsilon_{0} r^{2}}=\frac{\sigma^{2}}{2 \varepsilon_{0}}
\]

3.140 Initially, there will be induced charges of magnitude \(-q\) and \(+q\) on the inner and outer surface of the spherical layer respectively. Hence, the total electrical energy of the system is the sum of self energies of spherical shells, having radii \(a\) and \(b\), and their mutual energies including the point charge \(q\).\\
or

\[
\begin{gathered}
U_{i}=\frac{1}{2} \frac{q^{2}}{4 \pi \varepsilon_{0} b}+\frac{1}{2} \frac{(-q)^{2}}{4 \pi \varepsilon_{0} a}+\frac{-q q}{4 \pi \varepsilon_{0} a}+\frac{q q}{4 \pi \varepsilon_{0} b}+\frac{-q q}{4 \pi \varepsilon_{0} b} \\
U_{i}=\frac{q^{2}}{8 \pi \varepsilon_{0}}\left[\frac{1}{b}-\frac{1}{a}\right]
\end{gathered}
\]

Finally, charge \(q\) is at infinity hence, \(U_{f}=0\)\\
Now, work done by the agent = increment in the energy

\[
=U_{f}-U_{i}=\frac{q^{2}}{8 \pi \varepsilon_{0}}\left[\frac{1}{a}-\frac{1}{b}\right]
\]

3.141 (a) Sought work is equivalent to the work performed against the electric field created by one plate, holding at rest and to bring the other plate away. Therefore the required work,

\[
A_{\mathrm{agent}}=q E\left(x_{2}-x_{1}\right),
\]

where \(E=\frac{\sigma}{2 \varepsilon_{0}}\) is the intensity of the field created by one plate at the location of other.\\
So,

\[
A_{\text {agent }}=q \frac{\sigma}{2 \varepsilon_{0}}\left(x_{2}-x_{1}\right)=\frac{q^{2}}{2 \varepsilon_{0} S}\left(x_{2}-x_{1}\right)
\]

Alternate : \(A_{\text {ext }}=\Delta U\) (as field is potential)

\[
=\frac{q^{2}}{2 \varepsilon_{0} S} x_{2}-\frac{q^{2}}{2 \varepsilon_{0} S} x=\frac{q^{2}}{2 \varepsilon_{0} S}\left(x_{2}-x_{1}\right)
\]

(b) When voltage is kept const., the force acing on each plate of capacitor will depend on the distance between the plates.\\
So, elementary work done by agent, in its displacement over a distance \(d x\), relative to the other,

\[
d A=-F_{x} d x
\]

But,

\[
F_{x}=-\underset{x_{2}}{\left(\frac{\sigma(x)}{2 \varepsilon_{0}}\right) S \sigma(x)} \text { and } \sigma(x)=\frac{\varepsilon_{0} V}{x}
\]

Hence,

\[
A=\int d A=\int_{x_{1}} \frac{1}{2} \varepsilon_{0} \frac{S V^{2}}{x^{2}} d x=\frac{\varepsilon_{0} S V^{2}}{2}\left[\frac{1}{x_{1}}-\frac{1}{x_{2}}\right]
\]

Alternate : From energy Conservation,

\[
\begin{aligned}
& U_{f}-U_{i}=A_{\text {cell }}+A_{\text {agent }} \\
& \text { or } \quad \frac{1}{2} \frac{\varepsilon_{0} S}{x_{2}} V^{2}-\frac{1}{2} \frac{\varepsilon_{0} S}{x_{1}} V^{2}=\left[\frac{\varepsilon_{0} S}{x_{2}}-\frac{\varepsilon_{0} S}{x_{1}}\right] V^{2}+A_{\text {agent }} \\
& \left(\text { as } A_{\text {cell }}=\left(q_{f}-q_{i}\right) V=\left(C_{f}-C_{i}\right) V^{2}\right) \\
& \text { So } \\
& \qquad A_{\text {agent }}=\frac{\varepsilon_{0} S V^{2}}{2}\left[\frac{1}{x_{1}}-\frac{1}{x_{2}}\right]
\end{aligned}
\]

3.142 (a) When metal plate of thickness \(\eta d\) is inserted inside the capacitor, capacitance of the system becomes \(C_{0}=\frac{\varepsilon_{0} S}{d(1-\eta)}\)\\
Now, initially, charge on the capacitor, \(q_{0}=C_{0} V=\frac{\varepsilon_{0} S V^{.}}{d(1-\eta)}\)\\
Finally, capacitance of the capacitor, \(C=\frac{\varepsilon_{0} S}{d}\)\\
As the source is disconnected, charge on the plates will remain same during the process.\\
Now, from energy conservation,\\
or,

\[
\begin{gathered}
U_{f}-U_{i}=A_{\text {agent }}(\text { as cell does no work }) \\
\frac{1}{2} \frac{q_{0}^{2}}{C}-\frac{1}{2} \frac{q_{0}^{2}}{C_{0}}=A_{\text {agent }}
\end{gathered}
\]

Hence \(A_{\text {agent }}=\frac{1}{2}\left[\frac{\varepsilon_{0} S V}{d(1-\eta)}\right]^{2}\left[\frac{1}{C}-\frac{(1-\eta)}{C}\right]=\frac{1}{2} \frac{C V^{2} \eta}{(1-\eta)^{2}}=1.5 \mathrm{~mJ}\)\\
(b) Initially, capacitance of the system is given by,\\
\(C_{0}=\frac{C \varepsilon}{\eta(1-\varepsilon)+\varepsilon}\) (this is the capacitance of two capacitors in series)\\
So, charge on the plate, \(q_{0}=C_{0} V\)\\
Capacitance of the capacitor, after the glass plate has been removed equals \(C\) From energy conservation,

\[
\begin{aligned}
A_{\text {agent }} & =U_{f}-U_{i} \\
& =\frac{1}{2} q_{0}^{2}\left[\frac{1}{C}-\frac{1}{C_{0}}\right]=\frac{1}{2} \frac{C V^{2} \varepsilon \eta(\varepsilon-1)}{[\varepsilon-\eta(\varepsilon-1)]^{2}}=0.8 \mathrm{~mJ}
\end{aligned}
\]

3.143 When the capactior which is immersed in water is connected to a constant voltage source, it gets charged. Suppose \(\sigma_{0}\) is the free charge density on the condenser plates. Because water is a dielectric, bound charges also appear in it. Let \(\sigma^{\prime}\) be the surface density of bound charges. (Because of homogeneity of the medium and uniformity of the field when we ignore edge effects no volume density of bound charges exists.) The electric field due to free charges only \(\frac{\sigma_{0}}{\varepsilon_{0}}\); that due to bound charges is \(\frac{\sigma^{\prime}}{\varepsilon_{0}}\) and the total electric field is \(\frac{\sigma_{0}}{\varepsilon \varepsilon_{0}}\). Recalling that the sign of bound charges is opposite of the free charges, we have

\[
\frac{\sigma_{0}}{\varepsilon \varepsilon_{0}}=\frac{\sigma_{0}}{\varepsilon_{0}}-\frac{\sigma^{\prime}}{\varepsilon_{0}} \quad \text { or, } \quad \sigma^{\prime}=\left(\frac{\varepsilon-1}{\varepsilon}\right) \sigma_{0}
\]

Because of the field that exists due to the free charges (not the total field; the field due to the bound charges must be excluded for this purpose as they only give rise to self energy effects), there is a force attracting the bound charges to the near by plates. This force is

\[
\frac{1}{2} \sigma^{\prime} \frac{\sigma_{o}}{\varepsilon_{0}}=\frac{(\varepsilon-1)}{2 \varepsilon \varepsilon_{0}} \sigma_{0}^{2} \text { per unit area. }
\]

The factor \(\frac{1}{2}\) needs an explanation. Normally the force on a test charge is \(q E\) in an electric field \(E\). But if the charge itself is produced by the electric filed then the force must be constructed bit by bit and is

\[
F=\int_{0}^{E} q\left(E^{\prime}\right) d E^{\prime}
\]

if

\[
\begin{aligned}
q\left(E^{\prime}\right) & \propto E^{\prime} \text { then we get } \\
F & =\frac{1}{2} q(E) E
\end{aligned}
\]

This factor of \(\frac{1}{2}\) is well known. For example the energy of a dipole of moment \(\vec{p}\) in an electric field \(\overrightarrow{E_{0}}\) is \(-\vec{p} \cdot \overrightarrow{E_{0}}\) while the energy per unit volume of a linear dielectric in an electric field is \(-\frac{1}{2} \vec{P} \cdot \overrightarrow{E_{0}}\) where \(\vec{P}\) is the polarization vector (i.e. dipole moment per unit volume). Now the force per unit area manifests ifself as excess pressure of the liquid.

Noting that

\[
\frac{V}{d}=\frac{\sigma_{0}}{\varepsilon \varepsilon_{0}}
\]

We get

\[
\Delta p=\frac{\varepsilon_{0} \varepsilon(\varepsilon-1) V^{2}}{2 d^{2}}
\]

Substitution, using \(\varepsilon=81\) for water, gives \(\Delta p=7.17 \mathrm{k} \mathrm{Pa}=0.07 \mathrm{~atm}\).\\
3.144 One way of doing this problem will be exactly as in the previous case so let us try an alternative method based on energy. Suppose the liquid rises by a distance \(h\). Then let us calculate the extra energy of the liquid as a sum of polarization energy and the ordinary gravitiational energy. The latter is

\[
\frac{1}{2} h \cdot \rho g \cdot S h=\frac{1}{2} \rho g S h^{2}
\]

If \(\sigma\) is the free charge surface density on the plate, the bound charge density is, from the previous problem,

\[
\sigma^{\prime}=\frac{\varepsilon-1}{\varepsilon} \sigma
\]

This is also the volume density of induced dipole moment i.e. Polarization. Then the energy is, as before

\[
-\frac{1}{2} \cdot \sigma^{\prime} E_{0}=\frac{-1}{2} \cdot \sigma^{\prime} \frac{\sigma}{\varepsilon_{0}}=\frac{-(\varepsilon-1) \sigma^{2}}{2 \varepsilon_{0} \varepsilon}
\]

and the total polarization energy is

\[
-S(a+h) \frac{(\varepsilon-1) \sigma^{2}}{2 \varepsilon_{0} \varepsilon}
\]

Then, total energy is\\
\includegraphics[max width=\textwidth, alt={}, center]{9f3c5a8e-ac5d-42f4-8eb6-ea2af623f2d8-331_179_651_120_704}

\[
U(h)=-S(a+h) \frac{(\varepsilon-1) \sigma^{2}}{2 \varepsilon_{0} \varepsilon}+\frac{1}{2} \rho g S h^{2}
\]

The actual height to which the liquid rises is determined from the formula

\[
\frac{d U}{d h}=U^{\prime}(h)=0
\]

This gives

\[
h=\frac{(\varepsilon-1) \sigma^{2}}{2 \varepsilon_{0} \varepsilon \rho g}
\]

3.145 We know that energy of a capacitor, \(U=\frac{q^{2}}{2 C}\).

Hence, from \(F_{x}=\left.\frac{\partial U}{\partial x}\right|_{q \text {-Const. }}\) we have, \(F_{x}=\frac{q^{2}}{2} \frac{\partial c}{\partial x} / C^{2}\)

Now, since \(d \ll R\), the capacitance of the given capacitor can be calculated by the formula of a parallel plate capacitor. Therefore, if the dielectric is introduced upto a depth \(x\) and the length of the capacitor is \(l\), we have,


\begin{equation*}
C=\frac{2 \pi \varepsilon_{0} \varepsilon R x}{d}+\frac{2 \pi R \varepsilon_{0}(l-x)}{d} \tag{2}
\end{equation*}


From (1) and (2), we get,

\[
F_{x}=\varepsilon_{0}(\varepsilon-1) \frac{\pi R V^{2}}{d}
\]

3.146 When the capacitor is kept at a constant potential difference \(V\), the work performed by the moment of electrostatic forces between the plates when the inner moveable plate is rotated by an angle \(d \varphi\) equals the increase in the potential energy of the system. This comes about because when charges are made, charges flow from the battery to keep the potential constant and the amount of the work done by these charges is twice in magnitude and opposite in sign to the change in the energy of the capacitor Thus

\[
N_{Z}=\frac{\partial U}{\partial \varphi}=\frac{1}{2} V^{2} \frac{\partial C}{\partial \varphi}
\]

Now the capacitor can be thought of as made up two parts (with and without the dielectric) in parallel.\\
Thus \(C=\frac{\varepsilon_{0} R^{2} \varphi}{2 d}+\frac{\varepsilon_{0} \varepsilon(\pi-\varphi) R^{2}}{2 d}\)\\
as the area of a sector of angle \(\varphi\) is \(\frac{1}{2} R^{2} \varphi\). Differentiation then gives

\[
N_{z}=-\frac{(\varepsilon-1) \varepsilon_{0} R^{2} V^{2}}{4 d}
\]

The negative sign of \(N_{z}\) indicates that the moment of the force is acting clockwise (i.e. trying to suck in the dielectric).

\subsection*{3.4 ELECTRIC CURRENT}
3.147 The convection current is


\begin{equation*}
I=\frac{d q}{d t} \tag{1}
\end{equation*}


here, \(d q=\lambda d x\), where \(\lambda\) is the linear charge density.\\
But, from the Gauss' theorem, electric field at the surface of the cylinder,

\[
E=\frac{\lambda}{2 \pi \varepsilon_{0} a}
\]

Hence, substituting the value of \(\lambda\) and subsequently of \(d q\) in Eqs. (1), we get

\[
\begin{gathered}
I=\frac{2 E \pi \varepsilon_{0} a d x}{d t} \\
=2 \pi \varepsilon_{0} E a v, \text { as } \frac{d x}{d t}=v
\end{gathered}
\]

3.148 Since \(d \ll r\), the capacitance of the given capacitor can be calculated using the formula for a parallel plate capacitor. Therefore if the water (permittivity \(\varepsilon\) ) is introduced up to a height \(x\) and the capacitor is of length \(l\), we have,

\[
C=\frac{\varepsilon \varepsilon_{0} \cdot 2 \pi r x}{d}+\frac{\varepsilon_{0}(l-x) 2 \pi r}{d}=\frac{\varepsilon_{0} 2 \pi r}{d}(\varepsilon x+l-x)
\]

Hence charge on the plate at that instant, \(q=C V\)\\
Again we know that the electric current intensity,

\[
\begin{gathered}
I=\frac{d q}{d t}=\frac{d(C V)}{d t} \\
=\frac{V \varepsilon_{0} 2 \pi r}{d} \frac{d(\varepsilon x+l-x)}{d t}=\frac{V 2 \pi r \varepsilon_{0}}{d}(\varepsilon-1) \frac{d x}{d t}
\end{gathered}
\]

But,

\[
\frac{d x}{d t}=v,
\]

So,

\[
I=\frac{2 \pi r \varepsilon_{0}(\varepsilon-1) V}{d} v=0.11 \mu \mathrm{~A}
\]

3.149 We have, \(R_{t}=R_{0}(1+\alpha t),(1)\)\\
where \(R_{t}\) and \(R_{0}\) are resistances at \(t^{\circ} \mathrm{C}\) and \(0^{\circ} \mathrm{C}\) respectively and \(\alpha\) is the mean temperature coefficient of resistance.\\
So,

\[
R_{1}=R_{0}\left(1+\alpha_{1} t\right) \text { and } R_{2}=\eta R_{0}\left(1+\alpha_{2} t\right)
\]

(a) In case of series combination, \(R=\Sigma R_{i}\)\\
so


\begin{align*}
R & =R_{1}+R_{2}=R_{0}\left[(1+\eta)+\left(\alpha_{1}+\eta \alpha_{2}\right) t\right]  \tag{1}\\
& =R_{0}(1+\eta)\left[1+2 \frac{\alpha_{1}+\eta \alpha_{2}}{1+\eta} t\right] \tag{2}
\end{align*}


Comparing Eqs. (1) and (2), we conclude that temperature co-efficient of resistance of the circuit,

\[
\alpha=\frac{\alpha_{1}+\eta \alpha_{2}}{1+\eta}
\]

(b) In parallel combination

\[
R=\frac{R_{0}\left(1+\alpha_{1} t\right) R_{0} \eta\left(1+\alpha_{2} t\right)}{R_{0}\left(1+\alpha_{1} t\right)+\eta R_{0}\left(1+\alpha_{2} t\right)}=R^{\prime}\left(1+\alpha^{\prime} t\right), \text { where } R^{\prime}=\frac{\eta R_{0}}{1+\eta}
\]

Now, neglecting the terms, proportional to the product of temperature coefficients, as being very small, we get,

\[
\alpha^{\prime} \approx \frac{\eta \alpha_{1}+\alpha_{2}}{1+\eta}
\]

3.150 (a) The currents are as shown. From Ohm's law applied between 1 and 7 via 1487 (say)\\
\includegraphics[max width=\textwidth, alt={}, center]{9f3c5a8e-ac5d-42f4-8eb6-ea2af623f2d8-333_650_650_614_98}\\
\includegraphics[max width=\textwidth, alt={}, center]{9f3c5a8e-ac5d-42f4-8eb6-ea2af623f2d8-333_592_655_618_778}

Thus,

\[
I R_{\mathrm{eq}}=\frac{I}{3} R+\frac{I}{6} R+\frac{I}{3} R=\frac{5}{6} R I
\]

(b)\\
(a)

\[
R_{\mathrm{eq}}=\frac{5 R}{6}
\]

(b) Between 1 and 2 from the loop 14321, \(I_{1} R=2 I_{2} R+I_{3} R\) or, \(I_{1}=I_{3}+2 I_{2}\)\\
From the loop 48734,

\[
\left(I_{2}-I_{3}\right) R+2\left(I_{2}-I_{3}\right) R+\left(I_{2}-I_{3}\right) R=I_{3} R
\]

or,

\[
4\left(I_{2}-I_{3}\right)=I_{3} \quad \text { or } \quad I_{3}=\frac{4}{5} I_{2}
\]

so

\[
I_{1}=\frac{14}{5} I_{2}
\]

Then, \(\left(I_{1}+2 I_{2}\right) R_{\text {eq }}=\frac{24}{5} I_{2} R_{\text {eq }}=I_{1} R=\frac{14}{5} I_{2} R\)\\
or

\[
R_{\mathrm{eq}}=\frac{7}{12} R
\]

(c) Between 1 and 3

From the loop 15621

\[
I_{2} R=I_{1} R+\frac{I_{1}}{2} R \quad \text { or }, \quad I_{2}=3 \frac{I_{1}}{2}
\]

Then,

\[
\begin{aligned}
& \left(I_{1}+2 I_{2}\right) R_{e q}=4 I_{1} R_{e q} \\
& =I_{2} R+I_{2} R=3 I_{1} R
\end{aligned}
\]

Hence, \(\quad R_{\text {eq }}=\frac{3}{4} R\).\\
\includegraphics[max width=\textwidth, alt={}, center]{9f3c5a8e-ac5d-42f4-8eb6-ea2af623f2d8-334_483_559_101_692}\\
3.151 Total resistance of the circuit will be independent of the number of cells,\\
\includegraphics[max width=\textwidth, alt={}, center]{9f3c5a8e-ac5d-42f4-8eb6-ea2af623f2d8-334_319_902_711_240}\\
if

\[
\begin{gathered}
R_{x}=\frac{\left(R_{x}+2 R\right) R}{R_{x}+2 R+R} \\
R_{x}^{2}+2 R R_{x}-2 R^{2}=0
\end{gathered}
\]

or,\\
On solving and rejecting the negative root of the quadratic equation, we have,

\[
R_{x}=R(\sqrt{3}-1)
\]

3.152 Let \(R_{0}\) be the resistance of the network,\\
\includegraphics[max width=\textwidth, alt={}, center]{9f3c5a8e-ac5d-42f4-8eb6-ea2af623f2d8-334_287_892_1419_244}\\
then, \(R_{0}=\frac{R_{0} R_{2}}{R_{0}+R_{2}}\) or \(R_{0}^{2}-R_{0} R_{1}-R_{1} R_{2}=0\)\\
On solving we get,

\[
R_{0}=\frac{R_{1}}{2}\left(1+\sqrt{1+4 \frac{R_{2}}{R_{1}}}\right)=6 \Omega
\]

3.153 Suppose that the voltage \(V\) is applied between the points \(A\) and \(B\) then

\[
V=I R=I_{0} R_{0}
\]

where \(R\) is resistance of whole the grid, \(I\), the current through the grid and \(I_{0}\), the current through the segment \(A B\). Now from symmetry, \(I / 4\) is the part of the current, flowing through all the four wire segments, meeting at the point \(A\) and similarly the amount of current flowing through the wires, meeting at \(B\) is also \(I / 4\). Thus a current \(I / 2\) flows through the conductor \(A B\), i.e.

\[
I_{0}=\frac{I}{2}
\]

Hence,

\[
R=\frac{R_{0}}{2}
\]

3.154 Let us mentally isolate a thin cylindrical layer of inner and outer radii \(r\) and \(r+d r\) respectively. As lines of current at all the points of this layer are perpendicular to it, such a layer can be treated as a cylindrical conductor of thickness \(d r\) and cross-sectional area \(2 \pi r l\). So, we have,

\[
d R=\rho \frac{d r}{S(r)}=\rho \frac{d r}{2 \pi r l}
\]

and integrating between the limits, we get,\\
\includegraphics[max width=\textwidth, alt={}, center]{9f3c5a8e-ac5d-42f4-8eb6-ea2af623f2d8-335_417_635_722_772}

\[
R=\frac{\rho}{2 \pi l} \ln \frac{b}{a}
\]

3.155 Let us mentally isolate a thin spherical layer of inner and outer radii \(r\) and \(r+d r\). Lines of current at all the points of the this layer are perpendicular to it and therefore such a layer can be treated as a spherical conductor of thickness \(d r\) and cross sectional area \(4 \pi r^{2}\). So


\begin{equation*}
d R=\rho \frac{d r}{4 \pi r^{2}} \tag{1}
\end{equation*}


And integrating (1) between the limits \([a, b]\), we get,

\[
R=\frac{\rho}{4 \pi}\left[\frac{1}{a}-\frac{1}{b}\right]
\]

Now, for \(b \rightarrow \infty\), we have

\[
R=\frac{\rho}{4 \pi a}
\]

3.156 In our system, resistance of the medium \(R=\frac{\rho}{4 \pi}\left[\frac{1}{a}-\frac{1}{b}\right]\),\\
where \(\rho\) is the resistivity of the medium\\
The current

\[
i=\frac{\varphi}{R}=\frac{\varphi}{\frac{\rho}{4 \pi}\left[\frac{1}{a}-\frac{1}{b}\right]}
\]

Also, \(\quad i=\frac{-d q}{d t}=-\frac{d(C \varphi)}{d t}=-C \frac{d \varphi}{d t}\), as capacitance is constant.

So, equating (1) and (2) we get,\\
or,

\[
\begin{gathered}
\frac{\varphi}{\frac{\rho}{4 \pi}\left[\frac{1}{a}-\frac{1}{b}\right]}=-C \frac{d \varphi}{d t} . \\
-\int \frac{d \varphi}{\varphi}=\frac{\Delta t}{\frac{C \rho}{4 \pi}\left[\frac{1}{a}-\frac{1}{b}\right]} \\
\text { in } \eta=\frac{\Delta t 4 \pi a b}{C \rho(b-a)}
\end{gathered}
\]

or,\\
Hence, resistivity of the medium,

\[
\rho=\frac{4 \pi \Delta t a b}{C(b-a) \ln \eta}
\]

3.157 Let us mentally impart the charge \(+q\) and \(-q\) to the balls respectively. The electric field strength at the surface of a ball will be determined only by its own charge and the charge can be considered to be uniformly distributed over the surface, because the other ball is at infinite distance. Magnitude of the field strength is given by,

\[
E=\frac{q}{4 \pi \varepsilon_{0} a^{2}}
\]

So, current density \(j=\frac{1}{\rho} \frac{q}{4 \pi \varepsilon_{0} a^{2}}\) and electric current

\[
I=\int \vec{j} \cdot d \vec{S}=j S=\frac{q}{\rho 4 \pi \varepsilon_{0} a^{2}} \cdot 4 \pi a^{2}=\frac{q}{\rho \varepsilon_{0}}
\]

But, potential difference between the balls,

\[
\varphi_{+}-\varphi_{-}=2 \frac{q}{4 \pi \varepsilon_{0} a}
\]

Hence, the sought resistance,

\[
R=\frac{\varphi_{+}-\varphi_{-}}{I}=\frac{2 q / 4 \pi \varepsilon_{0} a}{q / \rho \varepsilon_{0}}=\frac{\rho}{2 \pi a}
\]

3.158 (a) The potential in the unshaded region beyond the conductor as the potential of the given charge and its image and has the form

\[
\varphi=A\left(\frac{1}{r_{1}}-\frac{1}{r_{2}}\right)
\]

where \(r_{1}, r_{2}\) are the distances of the point from the charge and its image. The potential has been taken to be zero on the conducting plane and on the ball

\[
\varphi \approx A\left(\frac{1}{a}-\frac{1}{2 l}\right)=V
\]

\begin{center}
\includegraphics[max width=\textwidth, alt={}]{9f3c5a8e-ac5d-42f4-8eb6-ea2af623f2d8-336_396_594_1624_780}
\end{center}

So \(A \approx V a\). In this calculation the conditions \(a \ll l\) is used to ignore the variation of \(\varphi\) over the ball.

The electric field at \(P\) can be calculated similarly. The charge on the ball is\\
and

\[
\begin{gathered}
Q=4 \pi \varepsilon_{0} V a \\
E_{P}=\frac{V a}{r^{2}} \quad 2 \cos \theta=\frac{2 a I V}{r^{3}}
\end{gathered}
\]

Then \(j=\frac{1}{\rho} E=\frac{2 a l V}{\rho r^{3}}\) normal to the plane.\\
(b) The total current flowing into the conducting plane is

\[
I=\int_{0}^{\infty} 2 \pi x d x j=\int_{0}^{\infty} 2 \pi x d x \frac{2 a l V}{\rho\left(\pi^{2}+l^{2}\right)^{3 / 2}}
\]

(On putting \(y=x^{2}+l^{2}\) )

\[
\begin{gathered}
I=\frac{2 \pi a l V}{\rho} \int_{l^{2}}^{\infty} \frac{d y}{y^{3 / 2}}=\frac{4 \pi a V}{\rho} \\
R=\frac{V}{I}=\frac{\rho}{2 \pi a}
\end{gathered}
\]

Hence\\
3.159 (a) The wires themselves will be assumed to be perfect conductors so the resistance is entirely due to the medium. If the wire is of length \(L\), the resistance \(R\) of the medium is \(\alpha \frac{1}{L}\) because different sections of the wire are connected in parallel (by the medium) rather than in series. Thus if \(R_{1}\) is the resistance per unit length of the wire then \(R=R_{1} / L\). Unit of \(R_{1}\) is ohm-meter.\\
The potential at a point \(P\) is by symmetry and superposition\\
(for \(l \gg a\) )

\[
\begin{aligned}
\varphi \cong \frac{A}{2} \ln \frac{r_{1}}{a} & -\frac{A}{2} \ln \frac{r_{2}}{a} \\
& =\frac{A}{2} \ln \frac{r_{1}}{r_{2}}
\end{aligned}
\]

\begin{center}
\includegraphics[max width=\textwidth, alt={}]{9f3c5a8e-ac5d-42f4-8eb6-ea2af623f2d8-337_494_478_1343_842}
\end{center}

Then \(\varphi_{1}=\frac{V}{2}=\frac{A}{2} \ln \frac{a}{l}\) (for the potential of 1)\\
or,

\[
A=-V / \ln \frac{l}{a}
\]

and

\[
\varphi=-\frac{V}{2 \ln l / a} \ln r_{1} / r_{2}
\]

We then calculate the field at a point \(P\) which is equidistant from \(1 \& 2\) and at a distance \(r\) from both :

Then

\[
\begin{aligned}
E & =\frac{V}{2 \ln l / a}\left(\frac{1}{r}\right) \times 2 \sin \theta \\
& =\frac{V l}{2 \ln l / a} \frac{1}{r^{2}}
\end{aligned}
\]

and

\[
J=\sigma E=\frac{1}{\rho} \frac{V}{2 \ln l / a} \frac{1}{r^{2}}
\]

(b) Near either wire

\[
E=\frac{V}{2 \ln l / a} \frac{1}{a}
\]

and

\[
J=\sigma E=\frac{1}{\rho} \frac{V}{2 \ln l / a}
\]

Then

\[
I=\frac{V}{R}=L \frac{V}{R_{1}}=J 2 \pi a L
\]

Which gives

\[
R_{1}=\frac{\rho}{\pi} \ln l / a
\]

3.160 Let us mentally impart the charges \(+q\) and \(-q\) to the plates of the capacitor.

Then capacitance of the network,


\begin{equation*}
C=\frac{q}{\varphi}=\frac{\varepsilon \varepsilon_{0} \int E_{n} d S}{\varphi} \tag{1}
\end{equation*}


Now, electric current,


\begin{equation*}
i=\int \vec{j} \cdot d \vec{S}=\int \sigma E_{n} d S \text { as } \vec{j} \uparrow \uparrow \vec{E} \tag{2}
\end{equation*}


Hence, using (1) in (2), we get,

\[
i=\frac{C \varphi}{\varepsilon \varepsilon_{0}} \sigma=\frac{C \varphi}{\rho \varepsilon \varepsilon_{0}}=1.5 \mu \mathrm{~A}
\]

3.161 Let us mentally impart charges \(+q\) and \(-q\) to the conductors. As the medium is poorly conducting, the surfaces of the conductors are equipotential and the field configuration is same as in the absence of the medium.\\
Let us surround, for example, the positively charged conductor, by a closed surface \(S\), just containing the conductor,\\
then,

\[
R=\frac{\varphi}{i}=\frac{\varphi}{\int \vec{j} \cdot d \vec{S}}=\frac{\varphi}{\int \sigma E_{n} d S} ; \text { as } \vec{j} \uparrow \uparrow \vec{E}
\]

and

\[
C=\frac{q}{\varphi}=\frac{\varepsilon \varepsilon_{0} \int E_{n} d S}{\varphi}
\]

So,

\[
R C=\frac{\varepsilon \varepsilon_{0}}{\sigma}=\rho \varepsilon \varepsilon_{0}
\]

3.162 The dielectric ends in a conductor. It is given that on one side (the dielectric side) the electric displacement \(D\) is as shown. Within the conductor, at any point \(A\), there can be no normal component of electric field. For if there were such a field, a current will flow towards depositing charge there which in turn will set up countering electric field causing the normal component to vanish. Then by Gauss theorem, we easily derive \(\sigma=D_{n}=D \cos \alpha\) where \(\sigma\) is the surface charge density at \(A\).\\
The tangential component is determined from the circulation theorem

\[
\oint \vec{E} \cdot d \vec{r}=0
\]

It must be continuous across the surface of the conductor. Thus, inside the conductor there is a tangential electric field of magnitude,

\[
\frac{D \sin \alpha}{\varepsilon_{0} \varepsilon} \text { at } A
\]

This implies a current, by Ohm's law, of\\
\includegraphics[max width=\textwidth, alt={}, center]{9f3c5a8e-ac5d-42f4-8eb6-ea2af623f2d8-339_375_619_437_782}

\[
j=\frac{D \sin \alpha}{\varepsilon_{0} \varepsilon \rho}
\]

3.163 The resistance of a layer of the medium, of thickness \(d x\) and at a distance \(x\) from the first plate of the capacitor is given by,


\begin{equation*}
d R=\frac{1}{\sigma(x)} \frac{d x}{S} \tag{1}
\end{equation*}


Now, since \(\sigma\) varies linearly with the distance from the plate. It may be represented as, \(\sigma=\sigma_{1}+\left(\frac{\sigma_{2}-\sigma_{1}}{d}\right) x\), at a distance \(x\) from any one of the plate.\\
From Eq. (1)\\
or,

\[
\begin{gathered}
d R=\frac{1}{\sigma_{1}+\left(\frac{\sigma_{2}-\sigma_{1}}{d}\right) x} \frac{d x}{S} \\
R=\frac{1}{S} \int_{0}^{d} \frac{d x}{\sigma_{1}+\left(\frac{\sigma_{2}-\sigma_{1}}{d}\right) x}=\frac{d}{S\left(\sigma_{2}-\sigma_{1}\right)} \ln \frac{\sigma_{2}}{\sigma_{1}}
\end{gathered}
\]

Hence,

\[
i=\frac{V}{R}=\frac{S V\left(\sigma_{2}-\sigma_{1}\right)}{d \ln \frac{\sigma_{2}}{\sigma_{1}}}=5 \mathrm{nA}
\]

3.164 By charge conservation, current \(\mathbf{j}\), leaving the medium (1) must enter the medium (2). Thus

\[
j_{1} \cos \alpha_{1}=j_{2} \cos \alpha_{2}
\]

Another relation follows from

\[
E_{1 t}=E_{2 t},
\]

which is a consequence of \(\oint \vec{E} \cdot d \vec{r}=0\)\\
Thus \(\frac{1}{\sigma_{1}} j_{1} \sin \alpha_{1}=\frac{1}{\sigma_{2}} j_{2} \sin \alpha_{2}\)\\
or, \(\quad \frac{\tan \alpha_{1}}{\sigma_{1}}=\frac{\tan \alpha_{2}}{\sigma_{2}}\)\\
\includegraphics[max width=\textwidth, alt={}, center]{9f3c5a8e-ac5d-42f4-8eb6-ea2af623f2d8-340_400_453_116_722}\\
or, \(\quad \frac{\tan \alpha_{1}}{\tan \alpha_{2}}=\frac{\sigma_{1}}{\sigma_{2}}\)\\
3.165 The electric field in conductor 1 is

\[
E_{1}=\frac{\rho_{1} I}{\pi R^{2}}
\]

and that in 2 is \(E_{2}=\frac{\rho_{2} I}{\pi R^{2}}\)\\
Applying Gauss' theorem to a small cylindrical pill-box at the boundary.

\[
-\frac{\rho_{1} I}{\pi R^{2}} d S+\frac{\rho_{2} I}{\pi R^{2}} d S=\frac{\sigma d S}{\varepsilon_{0}}
\]

Thus,

\[
\sigma=\varepsilon_{0}\left(\rho_{2}-\rho_{1}\right) \frac{1}{\pi R^{2}}
\]

and charge at the boundary \(=\varepsilon_{0}\left(\rho_{2}-\rho_{1}\right) I\)\\
3.166 We have, \(E_{1} d_{1}+E_{2} d_{2}=V\)\\
and by current conservation

\[
\frac{1}{\rho_{1}} E_{1}=\frac{1}{\rho_{2}} E_{2}
\]

Thus, \(\quad E_{1}=\frac{\rho_{1} V}{\rho_{1} d_{1}+\rho_{2} d_{2}}\),\\
\includegraphics[max width=\textwidth, alt={}, center]{9f3c5a8e-ac5d-42f4-8eb6-ea2af623f2d8-340_238_552_1154_707}\\
\(d_{1} \quad d_{2}\)

\[
E_{2}=\frac{\rho_{2} V}{\rho_{1} d_{1}+\rho_{2} d_{2}}
\]

At the boundary between the two dielectrics,

\[
\begin{gathered}
\sigma=D_{2}-D_{1}=\varepsilon_{0} \varepsilon_{2} E_{2}-\varepsilon_{0} \varepsilon_{1} E_{1} \\
\frac{\varepsilon_{0} V}{\rho_{1} d_{1}+\rho_{2} d_{2}}\left(\varepsilon_{2} \rho_{2}-\varepsilon_{1} \rho_{1}\right)
\end{gathered}
\]

3.167 By current conservation \(\frac{E(x)}{\rho(x)}=\frac{E(x)+d E(x)}{\rho(x)+d \rho(x)}=\frac{d E(x)}{d \rho(x)}\)\\
This has the solution,\\
\(E(x)=C \rho(x)=\frac{I \rho(x)}{A}\)\\
\includegraphics[max width=\textwidth, alt={}, center]{9f3c5a8e-ac5d-42f4-8eb6-ea2af623f2d8-340_218_94_1807_754}\\
\includegraphics[max width=\textwidth, alt={}, center]{9f3c5a8e-ac5d-42f4-8eb6-ea2af623f2d8-340_276_210_1744_890}\\
\includegraphics[max width=\textwidth, alt={}, center]{9f3c5a8e-ac5d-42f4-8eb6-ea2af623f2d8-340_214_129_1804_1156}

Hence charge induced in the slice per unit area\\
\(d \sigma=\varepsilon_{0} \frac{I}{A}[\{\varepsilon(x)+d \varepsilon(x)\}\{\rho(x)+d \rho(x)\}-\varepsilon(x) \rho(x)]=\varepsilon_{0} \frac{I}{A} d[\varepsilon(x) \rho(x)]\)\\
Thus,

\[
d Q=\varepsilon_{0} I d[\varepsilon(x) \rho(x)]
\]

Hence total charge induced, is by integration,

\[
Q=\varepsilon_{0} I\left(\varepsilon_{2} \rho_{2}-\varepsilon_{1} \rho_{1}\right)
\]

3.168 As in the previous problem

\[
E(x)=C \rho(x)=C\left(\rho_{0}+\rho_{1} x\right)
\]

where

\[
\rho_{0}+\rho_{1} d=\eta \rho_{0} \quad \text { or, } \quad \rho_{1}=\frac{(\eta-1) \rho_{0}}{d}
\]

By integration \(V=\int_{0} C \rho(x) d x=C \rho_{0} d\left(1+\frac{\eta-1}{2}\right)=\frac{1}{2} C \rho_{0} d(\eta+1)\)\\
Thus

\[
C=\frac{2 V}{\rho_{0} d(\eta+1)}
\]

Thus volume density of charge present in the medium

\[
\begin{gathered}
=\frac{d Q}{S d x}=\varepsilon_{0} d E(x) / d x \\
=\frac{2 \varepsilon_{0} V}{\rho_{0} d(\eta+1)} \times \frac{(\eta-1) \rho_{0}}{d}=\frac{2 \varepsilon_{0} V(\eta-1)}{(\eta+1) d^{2}}
\end{gathered}
\]

3.169 (a) Consider a cylinder of unit length and divide it into shells of radius \(r\) and thickness \(d r\) Different sections are in parallel. For a typical section,

Integrating,

\[
d\left(\frac{1}{R_{1}}\right)=\frac{2 \pi r d r}{\left(\alpha / r^{2}\right)}=\frac{2 \pi r^{3} d r}{\alpha}
\]

or,

\[
\frac{1}{R_{1}}=\frac{\pi R^{4}}{2 \alpha}=\frac{S^{2}}{2 \pi \alpha}
\]

or, \(\quad R_{1}=\frac{2 \pi \alpha}{S^{2}}\), where \(S=\pi R^{2}\)\\
(b) Suppose the electric filed inside is \(E_{z}=E_{o}\) ( \(Z\) axis is along the axiz of the conductor). This electric field cannot depend on \(r\) in steady conditions when other components of \(E\) are absent, otherwise one violates the circulation theorem

\[
\oint \vec{E} \cdot d \vec{r}=0
\]

The current through a section between radii ( \(r+d r, r\) ) is

\[
2 \pi r d r \frac{1}{\propto / r^{2}} E=2 \pi r^{3} d r \frac{E}{\alpha}
\]

Thus

\[
I=\int_{0}^{R} 2 \pi r^{3} d r \frac{E}{\propto}=\frac{\pi R^{4} E}{2 \propto}
\]

Hence

\[
E=\frac{2 \propto \pi I}{S^{2}} \text { when } S=\pi R^{2}
\]

3.170 The formula is,

\[
q=C V_{0}\left(1-e^{-t / R C}\right)
\]

or,

\[
V=\frac{q}{C}=V_{0}\left(1-e^{-t / R C}\right) \text { or, } \frac{V}{V_{0}}=1-e^{-t / R C}
\]

or,

\[
e^{-t / R C}=1-\frac{V}{V_{0}}=\frac{V_{0}-V}{V_{0}}
\]

Hence, \(t=R C \ln \frac{V_{0}}{V_{0}-V}=R C \ln 10\), if \(V=0.9 V_{0}\).\\
Thus \(t=0.6 \mu \mathrm{~S}\).\\
3.171 The charge decays according to the foumula

\[
q=q_{0} e^{-t / R C}
\]

Here,

\[
R C=\text { mean life }=\text { Half-life } / \ln 2
\]

So, half life \(=T=R C \ln 2\)\\
But,

\[
C=\frac{\varepsilon \varepsilon_{0} A}{d}, R=\frac{\rho d}{A}
\]

Hence,

\[
\rho=\frac{T}{\varepsilon \varepsilon_{0} \ln 2}=1.4 \times 10^{13} \Omega \cdot \mathrm{~m}
\]

3.172 Suppose \(q\) is the charge at time \(t\). Initially \(q=C \xi\), at \(t=0\).

Then at time \(t\),

\[
\frac{\eta q}{C}-i R-\xi=0
\]

But, \(i=-\frac{d q}{d t}\) (- sign because charge decreases)\\
\includegraphics[max width=\textwidth, alt={}, center]{9f3c5a8e-ac5d-42f4-8eb6-ea2af623f2d8-342_293_513_954_873}

So

\[
\begin{aligned}
\frac{\eta q}{C}+R \frac{d q}{d t} & =\xi \\
\frac{d q}{d t}+\frac{\eta}{R C} q & =\frac{\xi}{R}
\end{aligned}
\]

or,

\[
\frac{d}{d t} q e^{t \eta / R C}=\frac{\xi}{R} e^{t \eta / R C}
\]

or,

\[
q=\frac{C \xi}{\eta}+A e^{-t \eta / R C}
\]

\[
A=C \xi\left(1-\frac{1}{\eta}\right), \text { from } q=C \xi \text { at } t=0
\]

Hence,

\[
q=C \xi\left(\frac{1}{\eta}+\left(1-\frac{1}{\eta}\right) e^{-\eta t / R C}\right)
\]

Finally,

\[
i=-\frac{d q}{d t}=\frac{\xi(\eta-1)}{R} e^{-t \eta / R C}
\]

3.173 Let \(r=\) internal resistance of the battery. We shall take the resistance of the ammeter to \(\mathrm{be}=0\) and that of voltmeter to be G .\\
Initially, \(V=\xi-I r, I=\frac{\xi}{r+G}\)

So, \(\quad V=\xi \frac{G}{r+G}\)

After the voltmeter is shunted


\begin{align*}
& \qquad \frac{V}{\eta}=\xi-\frac{\xi r}{r+\frac{R G}{R+G}} \text { (Voltmeter) }  \tag{2}\\
& \text { and } \frac{\xi}{r+\frac{R G}{R+G}}=\eta \frac{\xi}{r+G} \text { (Ammeter) }  \tag{3}\\
& \text { From (2) and (3) we have }
\end{align*}


From (1) and (4)

\[
\frac{G}{\eta}=r+G-\eta r \text { or } G=\eta r
\]

Then (1) gives the required reading

\[
\frac{V}{\eta}=\frac{\xi}{\eta+1}
\]

3.174 Assume the current flow, as shown. Then potentials are as shown. Thus,\\
or, \(\quad I=\frac{\xi_{1}-\xi_{2}}{R_{1}+R_{2}}\)\\
And \(\varphi_{2}=\varphi_{1}-I R_{1}+\xi_{1}\)\\
So, \(\quad \varphi_{1}-\varphi_{2}=-\xi_{1}+\frac{\xi_{1}-\xi_{2}}{R_{1}+R_{2}} R_{1}\)\\
\(=-\left(\xi_{1} R_{2}+\xi_{2} R_{1}\right) /\left(R_{1}+R_{2}\right)=-4 \mathrm{~V}\)\\
\includegraphics[max width=\textwidth, alt={}, center]{9f3c5a8e-ac5d-42f4-8eb6-ea2af623f2d8-343_423_719_939_679}\\
3.175 Let, us consider the current \(i\), flowing through the circuit, as shown in the figure.

Applying loop rule for the circuit, \(-\Delta \varphi=0\)

\[
-2 \xi+i R_{1}+i R_{2}+i R=0
\]

or, \(\quad i\left(R_{1}+R_{2}^{\prime \prime}+R\right)=2 \xi\)\\
\includegraphics[max width=\textwidth, alt={}, center]{9f3c5a8e-ac5d-42f4-8eb6-ea2af623f2d8-343_380_545_1475_793}

Now, if

\[
\begin{gathered}
\varphi_{1}-\varphi_{2}=0 \\
-\xi+i R_{1}=0
\end{gathered}
\]

or, \(\frac{2 \xi R_{1}}{R+R_{1}+R_{2}}=\xi\) and \(2 R_{1}=R_{2}+R+R_{1}\)\\
or, \(R=R_{1}-R_{2}\), which is not possible as \(R_{2}>R_{1}\)\\
Thus,

\[
\varphi_{2}-\varphi_{3}=-\xi+i R_{2}=0
\]

or,

\[
\frac{2 \xi R_{2}}{R+R_{1}+R_{2}}=\xi
\]

So, \(\quad R=R_{2}-R_{1}\), which is the required resistance.\\
3.176 (a) Current, \(i=\frac{N \xi}{N R}=\frac{N \alpha R}{N R}=\alpha\), as \(\xi=\alpha R\) (given)\\
(b) \(\varphi_{A}-\varphi_{B}=n \xi-n i R=n \alpha R-n \alpha R=0\)\\
3.177 As the capacitor is fully charged, no current flows through it. So, current

\[
i=\frac{\xi_{2}-\xi_{1}}{R_{1}+R_{2}}\left(\text { as } \xi_{2}>\xi_{1}\right)
\]

And hence, \(\varphi_{A}-\varphi_{B}=\xi_{1}-\xi_{2}+i R_{2}\)

\[
\begin{aligned}
& =\xi_{1}-\xi_{2}+\frac{\xi_{2}-\xi_{1}}{R_{1}+R_{2}} R_{2} \\
& =\frac{\left(\xi_{1}-\xi_{2}\right) R_{1}}{R_{1}+R_{2}}=-0.5 \mathrm{~V}
\end{aligned}
\]

\includegraphics[max width=\textwidth, alt={}, center]{9f3c5a8e-ac5d-42f4-8eb6-ea2af623f2d8-344_365_442_485_801}\\
3.178 Let us make the current distribution, as shown in the figure.

Current \(i=\frac{\xi}{R+\frac{R_{1} R_{2}}{R_{1}+R_{2}}}\) (using loop rule)\\
So, current through the resistor \(R_{1}\),

\[
\begin{aligned}
i_{1} & =\frac{\xi}{R+\frac{R_{1} R_{2}}{R_{1}+R_{2}}} \frac{R_{2}}{R_{1}+R_{2}} \\
& =\frac{\xi R_{2}}{R R_{1}+R R_{2}+R_{1} R_{2}}=1.2 \mathrm{~A}
\end{aligned}
\]

\includegraphics[max width=\textwidth, alt={}, center]{9f3c5a8e-ac5d-42f4-8eb6-ea2af623f2d8-344_253_479_1068_700}\\
and similary, current through the resistor \(R_{2}\),\\
\(i_{2}=\frac{\xi}{R+\frac{R_{1} R_{2}}{R_{1}+R_{2}}} \frac{R_{1}}{R_{1}+R_{2}}=\frac{\xi R_{1}}{R R_{1}+R_{1} R_{2}+R R_{2}}=0.8 \mathrm{~A}\)\\
3.179 Total resistance \(=\frac{l-x}{l} R_{0}+\frac{R \cdot \frac{x R_{0}}{l}}{R+\frac{x R_{0}}{l}}\)

\[
\begin{aligned}
= & \frac{l-x}{l} R_{0}+\frac{x R R_{0}}{l R+x R_{0}} \\
= & R_{0}\left[\frac{l-x}{l}+\frac{x R}{l R+x R_{0}}\right]
\end{aligned}
\]

\begin{center}
\includegraphics[max width=\textwidth, alt={}]{9f3c5a8e-ac5d-42f4-8eb6-ea2af623f2d8-344_277_447_1737_755}
\end{center}

Then \(\quad V=V_{0} \frac{x R}{l R+x R_{0}} /\left(1-\frac{x}{l}+\frac{x R}{x R_{0}+l R}\right)=V_{0} R x /\left\{l R+R_{0} x\left(1-\frac{x}{l}\right)\right\}\)\\
For \(\quad R \gg R_{0}, V \cong V_{0} \frac{x}{l}\)\\
3.180 Let us connect a load of resistance \(K\) between the points \(A\) and \(B\) (Fig.)

From the loop rule, \(\Delta \varphi=0\), we obtain\\
and


\begin{equation*}
i R=\xi_{1}-i_{1} R_{1} \tag{1}
\end{equation*}


or

\[
i R=\xi_{2}-\left(i-i_{1}\right) R_{2}
\]

or \(\quad i\left(R+R_{2)}=\xi+i_{1} R_{2}(2)\right.\)\\
Solving Eqs. (1) and (2), we get


\begin{equation*}
i=\frac{\xi_{1} R_{1}+\xi_{2} R_{2}}{R_{1}+R_{2}} / R+\frac{R_{1} R_{2}}{R_{1}+R_{2}}=\frac{\xi_{0}}{R+R_{0}} \tag{3}
\end{equation*}


\includegraphics[max width=\textwidth, alt={}, center]{9f3c5a8e-ac5d-42f4-8eb6-ea2af623f2d8-345_429_312_361_1035}\\
where \(\xi_{0}=\frac{\xi_{1} R_{1}+\xi_{2} R_{2}}{R_{1}+R_{2}}\) and \(R_{0}=\frac{R_{1} R_{2}}{R_{1}+R_{2}}\)\\
Thus one can replace the given arrangement of the cells by a single cell having the emf \(\xi_{0}\) and internal resistance \(R_{0}\).\\
3.181 Make the current distribution, as shown in the diagram.\\
Now, in the loop 12341, applying \(-\Delta \varphi=0\)


\begin{equation*}
i R+i_{1} R_{1}+\xi_{1}=0 \tag{1}
\end{equation*}


and in the loop 23562,


\begin{equation*}
i R-\xi_{2}+\left(i-i_{1}\right) R_{2}=0 \tag{2}
\end{equation*}


Solving (1) and (2), we obtain current through the resistance \(R\),\\
\includegraphics[max width=\textwidth, alt={}, center]{9f3c5a8e-ac5d-42f4-8eb6-ea2af623f2d8-345_338_569_1002_802}

\[
i=\frac{\left(\xi_{2} R_{1}-\xi_{1} R_{2}\right)}{R R_{1}+R R_{2}+R_{1} R_{2}}=0.02 \mathrm{~A}
\]

and it is directed from left to the right\\
3.182 At first indicate the currents in the branches using charge conservation (which also includes the point rule).\\
In the loops 1 BA 61 and B34AB from the loop rule, \(-\Delta \varphi=0\), we get, respectively


\begin{equation*}
-\xi_{2}+\left(i-i_{1}\right) R_{2}+\xi_{1}-i_{1} R_{1}=0 \tag{1}
\end{equation*}



\begin{equation*}
i R_{3}+\xi_{3}-\left(i-i_{1}\right) R_{2}+\xi_{2}=0 \tag{2}
\end{equation*}


On solving Eqs (1) and (2), we obtain

\[
i_{1}=\frac{\left(\xi_{1}-\xi_{2}\right) R_{3}+R_{2}\left(\xi_{1}+\xi_{3}\right)}{R_{1} R_{2}+R_{2} R_{3}+R_{3} R_{1}} \approx 0.06 \mathrm{~A}
\]

\begin{center}
\includegraphics[max width=\textwidth, alt={}]{9f3c5a8e-ac5d-42f4-8eb6-ea2af623f2d8-345_361_474_1592_915}
\end{center}

Thus \(\varphi_{A}-\varphi_{B}=\xi_{2}-i_{2} R_{2} \approx 0.9 \mathrm{~V}\)\\
3.183 Indicate the currents in all the branches using charge conservation as shown in the figure. Applying loop rule, \(-\Delta \varphi=0\) in the loops 1A781, 1B681 and B456B, respectively, we get


\begin{align*}
& \xi_{0}=\left(i_{0}-i_{1}\right) R_{1}  \tag{1}\\
& i_{3} R_{3}+i_{1} R_{2}-\xi_{0}=0  \tag{2}\\
& \left(i_{1}-i_{3}\right) R-\xi-i_{3} R_{3}=0 \tag{3}
\end{align*}


Solving Eqs. (1), (2) and (3), we get the sought current\\
\(\left(i_{1}-i_{3}\right)=\frac{\xi\left(R_{2}+R_{3}\right)+\xi_{0} R_{3}}{R\left(R_{2}+R_{3}\right)+R_{2} R_{3}}\)\\
\includegraphics[max width=\textwidth, alt={}, center]{9f3c5a8e-ac5d-42f4-8eb6-ea2af623f2d8-346_388_758_199_643}\\
3.184 Indicate the currents in all the branches using charge conservation as shown in the figure.

Applying the loop rule \((-\Delta \varphi=0)\) in the loops\\
12341 and 15781, we get


\begin{equation*}
-\xi_{1}+i_{3} R_{1}-\left(i_{1}-i_{3}\right) R_{2}=0 \tag{1}
\end{equation*}


and \(\left(i_{1}-i_{3}\right) R_{2}-\xi_{2}+i_{1} R_{3}=0\)

Solving Eqs. (1) and (2), we get

\[
i_{3}=\frac{\xi_{1}\left(R_{2}+R_{3}\right)+\xi_{2} R_{2}}{R_{1} R_{2}+R_{2} R_{3}+R_{3} R_{1}}
\]

Hence, the sought p.d.\\
\includegraphics[max width=\textwidth, alt={}, center]{9f3c5a8e-ac5d-42f4-8eb6-ea2af623f2d8-346_354_563_768_834}

\[
\begin{gathered}
\varphi_{A}-\varphi_{B}=\xi_{2}-i_{3} R_{1} \\
=\frac{\xi_{2} R_{3}\left(R_{1}+R_{2}\right)-\xi_{1} R_{1}\left(R_{2}+R_{3}\right)}{R_{1} R_{2}+R_{2} R_{3}+R_{3} R_{1}}=-1 \mathrm{~V}
\end{gathered}
\]

3.185 Let us distribute the currents in the paths as shown in the figure.\\
Now, \(\quad \varphi_{1}-\varphi_{2}=i R_{1}+i_{1} R_{2}\)\\
and \(\varphi_{1}-\varphi_{3}=i R_{1}+\left(i-i_{1}\right) R_{3}\)

Simplifying Eqs. (1) and (2) we get

\[
i=\frac{R_{3}\left(\varphi_{1}-\varphi_{2}\right)+R_{2}\left(\varphi_{1}-\varphi_{3}\right)}{R_{1} R_{2}+R_{2} R_{3}+R_{3} R_{1}}=0.2 \mathrm{~A}
\]

\includegraphics[max width=\textwidth, alt={}, center]{9f3c5a8e-ac5d-42f4-8eb6-ea2af623f2d8-346_471_571_1234_824}\\
3.186 Current is as shown. From Kirchhoff's Second law

\[
\begin{gathered}
i_{1} R_{1}=i_{2} R_{3}, \\
i_{1} R_{1}+\left(i_{1}-i_{3}\right) R_{2}=V, \\
i_{2} R_{3}+\left(i_{3}+i_{2}\right) R_{4}=V
\end{gathered}
\]

\begin{center}
\includegraphics[max width=\textwidth, alt={}]{9f3c5a8e-ac5d-42f4-8eb6-ea2af623f2d8-346_265_607_1719_790}
\end{center}

Eliminating \(\boldsymbol{i}_{2}\)

Hence\\
or,

\[
\begin{gathered}
i_{1}\left(R_{1}+R_{2}\right)-i_{3} R_{2}=V \\
i_{1} \frac{R_{1}}{R_{3}}\left(R_{3}+R_{4}\right)+i_{3} R_{4}=V \\
i_{3}\left[R_{4}\left(R_{1}+R_{2}\right)+\frac{R_{1} R_{2}}{R_{3}}\left(R_{3}+R_{4}\right)\right] \\
=V\left[\left(R_{1}+R_{2}\right)-\frac{R_{1}}{R_{3}}\left(R_{3}+R_{4}\right)\right] \\
i_{3}=\frac{R_{3}\left(R_{1}+R_{2}\right)-R_{1}\left(R_{3}+R_{4}\right)}{R_{3} R_{4}\left(R_{1}+R_{2}\right)+R_{1} R_{2}\left(R_{3}+R_{4}\right)}
\end{gathered}
\]

On substitution we get \(i_{3}=1.0 \mathrm{~A}\) from \(C\) to \(D\)\\
3.187 From the symmetry of the problem, current flow is indicated, as shown in the figure.\\
Now,


\begin{equation*}
\varphi_{A}-\varphi_{B}=i_{1} r+\left(i-i_{1}\right) R \tag{1}
\end{equation*}


In the loop 12561, from \(-\Delta \varphi=0\)


\begin{gather*}
\quad\left(i-i_{1}\right) R+\left(i-2 i_{1}\right) r-i_{1} r=0 \\
\text { or, } \quad i_{1}=\frac{(R+r)}{3 r+R} i \tag{2}
\end{gather*}


Equivalent resistance between the terminals\\
\includegraphics[max width=\textwidth, alt={}, center]{9f3c5a8e-ac5d-42f4-8eb6-ea2af623f2d8-347_330_571_821_798}\\
\(A\) and \(B\) using (1) and (2),\\
\(R_{0}=\frac{\varphi_{A}-\varphi_{B}}{i}=\frac{\left[\frac{R+r}{3 r+R}(r-R)+R\right] i}{i}=\frac{r(3 R+r)}{3 r+R}\)\\
3.188 Let, at any moment of time, charge on the plates be \(+q\) and \(-q\) respectively, then voltage\\
across the capacitor, \(\varphi=q / C\)

Now, from charge conservation,


\begin{equation*}
i=i_{1}+i_{2}, \text { where } i_{2}=\frac{d q}{d t} \tag{2}
\end{equation*}


In the loop 65146, using \(-\Delta \varphi=0\).


\begin{equation*}
\frac{q}{C}+\left(i_{1}+\frac{d q}{d t}\right) R-\xi=0 \tag{3}
\end{equation*}


[using (1) and (2)]\\
In the loop 25632, using \(-\Delta \varphi=0\)\\
\includegraphics[max width=\textwidth, alt={}, center]{9f3c5a8e-ac5d-42f4-8eb6-ea2af623f2d8-347_364_527_1525_823}


\begin{equation*}
-\frac{q}{C}+i_{1} R=0 \quad \text { or, } \quad i_{1} R=\frac{q}{C} \tag{4}
\end{equation*}


From (1) and (2),


\begin{equation*}
\frac{d q}{d t} R=\xi_{1}-\frac{2 q}{C} \quad \text { or, } \quad \frac{d q}{\xi-\frac{2 q}{C}}=\frac{d t}{R} \tag{5}
\end{equation*}


On integrating the expression (5) between suitable limits,

\[
\int_{n}^{q} \frac{d q}{\xi-\frac{2 q}{C}}=\frac{1}{R} \int_{0}^{t} d t \quad \text { or, } \quad-\frac{C}{2} \ln \frac{\xi-\frac{2 q}{C}}{\xi}=\frac{t}{R}
\]

Thus

\[
\frac{q}{C}=V=\frac{1}{2} \xi\left(1-e^{-2 t / R C}\right)
\]

3.189 (a) As current \(i\) is linear function of time, and at \(t=0\) and \(\Delta t\), it equals \(i_{0}\) and zero respectively, it may be represented as,

\[
i=i_{0}\left(1-\frac{t}{\Delta t}\right)
\]

Thus

\[
q=\int_{0}^{\Delta t} i d t=\int_{0}^{\Delta t} i_{0}\left(1-\frac{t}{\Delta t}\right) d t=\frac{i_{0} \Delta t}{2}
\]

So,

\[
i_{0}=\frac{2 q}{\Delta t}
\]

Hence,

\[
i=\frac{2 q}{\Delta t}\left(1-\frac{t}{\Delta t}\right)
\]

The heat generated.

\[
H=\int_{0}^{\Delta t} i^{2} R d t=\int_{0}^{\Delta t}\left[\frac{2 q}{\Delta t}\left(1-\frac{t}{\Delta t}\right)\right]^{2} R d t=\frac{4 q^{2} R}{3 \Delta t}
\]

(b) Obviously the current through the coil is given by

\[
i=i_{0}\left(\frac{1}{2}\right)^{t / \Delta t}
\]

Then charge

\[
q=\int_{0}^{\infty} i d t=\int_{0}^{\infty} i_{0} 2^{-t / \Delta t} d t=\frac{i_{0} \Delta t}{\ln 2}
\]

So,

\[
i_{0}=\frac{q \ln 2}{\Delta t}
\]

And hence, heat generated in the circuit in the time interval \(t[0, \infty]\),

\[
H=\int_{0}^{\infty} i^{2} R d t=\int_{0}^{\infty}\left[\frac{q \ln 2}{\Delta t} 2^{-t / \Delta}\right]^{2} R d t=-\frac{q^{2} \ln 2}{2 \Delta t} R
\]

6.190 The equivalent circuit may be drawn as in the figure.

Resistance of the network \(=R_{0}+(R / 3)\)\\
Let, us assume that e.m.f. of the cell is \(\xi\), then current

\[
i=\frac{\xi}{R_{0}+(R / 3)}
\]

Now, thermal power, generated in the circuit

\[
P=i^{2} R / 3=\frac{\xi^{2}}{\left(R_{0}+(R / 3)\right)^{2}}(R / 3)
\]

For \(P\) to be maximum, \(\frac{d P}{d R}=0\), which yields\\
\includegraphics[max width=\textwidth, alt={}, center]{9f3c5a8e-ac5d-42f4-8eb6-ea2af623f2d8-349_509_395_139_942}

\[
R=3 R_{0}
\]

3.191 We assume current conservation but not Kirchhoff's second law. Then thermal power dissipated is

\[
\begin{aligned}
P\left(i_{1}\right) & =i_{1}^{2} R_{1}+\left(i-i_{1}\right)^{2} R_{2} \\
& =i_{1}^{2}\left(R_{1}+R_{2}\right)-2 i i_{1} R_{2}+i^{2} R_{2} \\
& =\left[R_{1}+R_{2}\right]\left[i_{1}-\frac{R_{2}}{R_{1}+R_{2}} i\right]^{2}+i^{2} \frac{R_{1} R_{2}}{R_{1}+R_{2}}
\end{aligned}
\]

\begin{center}
\includegraphics[max width=\textwidth, alt={}]{9f3c5a8e-ac5d-42f4-8eb6-ea2af623f2d8-349_272_544_755_803}
\end{center}

The resistances being positive we see that the power dissipated is minimum when

\[
i_{1}=i \frac{R_{2}}{R_{1}+R_{2}}
\]

This corresponds to usual distribution of currents over resistance joined is parallel.\\
3.192 Let, internal resistance of the cell be \(r\), then


\begin{equation*}
V=\xi-i r \tag{1}
\end{equation*}


where \(i\) is the current in the circuit. We know that thermal power generated in the battery.


\begin{equation*}
Q=i^{2} r \tag{2}
\end{equation*}


Putting \(r\) from (1) in (2), we obtain,

\[
Q=(\xi-V) i=0.6 \mathrm{~W}
\]

In a battery work is done by electric forces (whose origin lies in the chemical processes going on inside the cell). The work so done is stored and used in the electric circuit outside.\\
\includegraphics[max width=\textwidth, alt={}, center]{9f3c5a8e-ac5d-42f4-8eb6-ea2af623f2d8-349_307_428_1456_790}\\
Its magnitude just equals the power used in the electric circuit. We can say that net power developed by the electric forces is

\[
P=-I V=-2 \cdot 0 \mathrm{~W}
\]

Minus sign means that this is generated not consumed.\\
3.193 As far as motor is concerned the power delivered is dissipated and \href{http://can.be}{can.be} represented by a load, \(R_{0}\). Thus

\[
I=\frac{V}{R+R_{0}}
\]

and \(P=I^{2} R_{0}=\frac{V^{2} R_{0}}{\left(R_{0}+R\right)^{2}}\)\\
This is maximum when \(R_{0}=R\) and the current\\
\includegraphics[max width=\textwidth, alt={}, center]{9f3c5a8e-ac5d-42f4-8eb6-ea2af623f2d8-350_253_356_194_1023}\\
\(I\) is then

\[
I=\frac{V}{2 R}
\]

The maximum power delivered is

\[
\frac{V^{2}}{4 R}=P_{\max }
\]

The power input is \(\frac{V^{2}}{R+R_{0}}\) and its value when \(P\) is maximum is \(\frac{V^{2}}{2 R}\)\\
The efficiency then is \(\frac{1}{2}=50 \%\)\\
3.194 If the wire diameter decreases by \(\delta\) then by the information given\\
\(P=\) Power input \(=\frac{V^{2}}{R}=\) heat lost through the surface, \(H\).\\
Now, \(H \propto(1-\delta)\) like the surface area and

\[
R \propto(1-\delta)^{-2}
\]

So, \(\quad \frac{V^{2}}{R_{0}}(1-\delta)^{2}=A(1-\delta)\) or, \(V^{2}(1-\delta)=\) constant\\
But \(\quad V \propto 1+\eta\) so \((1+\eta)^{2}(1-\delta)=\) Const \(=1\)\\
Thus \(\quad \delta=2 \eta=2 \%\)\\
3.195 The equation of heat balance is

\[
\frac{V^{2}}{R}-k\left(T-T_{0}\right)=C \frac{d T}{d t}
\]

Put

\[
T-T_{0}=\xi
\]

So,

\[
C \xi+k \xi=\frac{V^{2}}{R} \quad \text { or, } \quad \xi+\frac{k}{C} \xi=\frac{V^{2}}{C R}
\]

or,

\[
\frac{d}{d t}\left(\xi e^{k t / c}\right)=\frac{V^{2}}{C R} e^{k t / c}
\]

or,

\[
\xi e^{k t / c}=\frac{V^{2}}{k R} e^{k t / c}+A
\]

where \(A\) is a constant. Clearly

\[
\begin{aligned}
\xi=0 \text { at } t & =0, \text { so } A=-\frac{V^{2}}{k R} \text { and hence, } \\
T & =T_{0}+\frac{V^{2}}{k R}\left(1-e^{-k t / c}\right)
\end{aligned}
\]

3.196 Let, \(\varphi_{A}-\varphi_{B}=\varphi\)

Now, thermal power generated in the resistance \(\boldsymbol{R}_{\boldsymbol{x}}\),

\[
P=i^{\prime 2} R_{x}=\left[\frac{\varphi}{R_{1}+\frac{R_{2} R_{x}}{R_{2}+R_{x}}} \frac{R_{2}}{R_{2}+R_{x}}\right]^{2} R_{x}
\]

For \(P\) to be independent of \(R_{x}\),

\[
\begin{aligned}
& \frac{d P}{d R_{x}}=0, \text { which yeilds } \\
& R_{x}=\frac{R_{1} R_{2}}{R_{1}+R_{2}}=12 \Omega
\end{aligned}
\]

\includegraphics[max width=\textwidth, alt={}, center]{9f3c5a8e-ac5d-42f4-8eb6-ea2af623f2d8-351_261_540_292_775}\\
3.197 Indicate the currents in the circuit as shown in the figure.

Appying loop rule in the closed loop 12561, \(-\Delta \varphi=0\) we get


\begin{equation*}
i_{1} R-\xi_{1}+i R_{1}=0 \tag{1}
\end{equation*}


and in the loop 23452,


\begin{equation*}
\left(i-i_{1}\right) R_{2}+\xi_{2}-i_{1} R=0 \tag{2}
\end{equation*}


Solving (1) and (2), we get,

\[
i_{1}=\frac{\xi_{1} R_{2}+\xi_{2} R_{1}}{R R_{1}+R_{1} R_{2}+R R_{2}}
\]

So, thermal power, generated in the resistance \(R\),\\
\(P=i_{1}^{2} R=\left[\frac{\xi_{1} R_{2}+\xi_{2} R_{1}}{R R_{1}+R_{1} R_{2}+R R_{2}}\right]^{2} R\)\\
For \(P\) to be maximum, \(\frac{d P}{d R}=0\), which fields\\
\includegraphics[max width=\textwidth, alt={}, center]{9f3c5a8e-ac5d-42f4-8eb6-ea2af623f2d8-351_510_450_828_828}

Hence,

\[
\begin{gathered}
R=\frac{R_{1} R_{2}}{R_{1}+R_{2}} \\
P_{\max }=\frac{\left(\xi_{1} R_{2}+\xi_{2} R_{1}\right)^{2}}{4 R_{1} R_{2}\left(R_{1}+R_{2}\right)}
\end{gathered}
\]

3.198 Let, there are \(x\) number of cells, connected in series in each of the \(n\) parallel groups then, \(n x=N\) or, \(x=\frac{N}{n}\)

Now, for any one of the loop, consisting of \(x\) cells and the resistor \(R\), from loop rule

\[
i R+\frac{i}{n} x r-x \xi=0
\]

So, \(i=\frac{x \xi}{R+\frac{x r}{n}}=\frac{\frac{N}{n} \xi}{R+\frac{N r}{n^{2}}}\), using\\
\includegraphics[max width=\textwidth, alt={}, center]{9f3c5a8e-ac5d-42f4-8eb6-ea2af623f2d8-351_373_414_1602_824}

Heat generated in the resistor \(R\),


\begin{equation*}
Q=i^{2} R=\left(\frac{N n \xi}{n^{2} R+N R}\right)^{2} R \tag{2}
\end{equation*}


and for \(Q\) to be maximum, \(\frac{d Q}{d n}=0\), which yields

\[
n=\sqrt{\frac{N r}{R}}=3
\]

3.199 When switch 1 is closed, maximum charge accumulated on the capacitor,


\begin{equation*}
q_{\max }=C \xi, \tag{1}
\end{equation*}


and when switch 2 is closed, at any arbitrary instant of time,

\[
\left(R_{1}+R_{2}\right)\left(-\frac{d q}{d t}\right)=q / C
\]

because capacitor is discharging.\\
or, \(\int_{q_{\text {max }}}^{q} \frac{1}{q} d q=-\frac{1}{\left(R_{1}+R_{2}\right) C} \int_{0}^{t} d t\)\\
\includegraphics[max width=\textwidth, alt={}, center]{9f3c5a8e-ac5d-42f4-8eb6-ea2af623f2d8-352_385_339_552_1027}

Integrating, we get


\begin{equation*}
\ln \frac{q}{q_{\max }}=\frac{-t}{\left(R_{1}+R_{2}\right) C} \quad \text { or, } \quad q=q_{\max } e^{\frac{-t}{\left(R_{1}+R_{2}\right) C}} \tag{2}
\end{equation*}


Differentiating with respect to time,

\[
i(t)=\frac{d q}{d t}=q_{\max } e^{\frac{-t}{\left(R_{1}+R_{2}\right) C}}\left(-\frac{1}{\left(R_{1}+R_{2}\right) C}\right)
\]

or,

\[
i(t)=\frac{C \xi}{\left(R_{1}+R_{2}\right) C} e^{\frac{-t}{\left(R_{1}+R_{2}\right) C}}
\]

Negative sign is ignored, as we are not interested in the direction of the current.\\
thus,


\begin{equation*}
i(t)=\frac{\xi}{\left(R_{1}+R_{2}\right)} e^{\frac{-1}{\left(R_{1}+R_{2}\right) C}} \tag{3}
\end{equation*}


When the switch ( \(S w\) ) is at the position 1 , the charge (maximum) accumalated on the capacitor is,

\[
q=C \xi
\]

When the \(S w\) is thrown to position 2 , the capacitor starts discharging and as a result the electric energy stored in the capacitor totally turns into heat energy tho' the resistors \(R_{1}\) and \(R_{2}\) (during a very long interval of time). Thus from the energy conservation, the total heat liberated tho' the resistors.

\[
H=U_{i}=\frac{q^{2}}{2 C}=\frac{1}{2} C \xi^{2}
\]

During the process of discharging of the capacitor, the current tho' the resistors \(R_{1}\) and \(R_{2}\) is the same at all the moments of time, thus

\[
H_{1} \propto R_{1} \text { and } H_{2} \propto R_{2}
\]

So,

\[
H_{1}=\frac{H R_{1}}{\left(R_{1}+R_{2}\right)} \quad\left(\text { as } H=H_{1}+H_{2}\right)
\]

Hence

\[
H_{1}=\frac{1}{2} \frac{C R_{1}}{R_{1}+R_{2}} \xi^{2}
\]

3.200 When the plate is absent the capacity of the condenser is

\[
C=\frac{\varepsilon_{0} S}{d}
\]

When it is present, the capacity is

\[
C^{\prime}=\frac{\varepsilon_{0} S}{d(1-\eta)}=\frac{C}{1-\eta}
\]

(a) The energy increment is clearly.

\[
\Delta U=\frac{1}{2} C V^{2}-\frac{1}{2} C^{\prime} V^{2}=\frac{C \eta}{2(1-\eta)} V^{2}
\]

(b) The charge on the plate is

\[
q_{i}=\frac{C V}{1-\eta} \text { initially and } q_{f}=C V \text { finally. }
\]

A charge \(\frac{C V \eta}{1-\eta}\) has flown through the battery charging it and withdrawing \(\frac{C V^{2} \eta}{1-\eta}\) units of energy from the system into the battery. The energy of the capacitor has decreased by just half of this. The remaining half i.e. \(\frac{1}{2} \frac{C V^{2} \eta}{1-\eta}\) must be the work done by the external agent in withdrawing the plate. This ensures conservation of energy.\\
3.201 Initially, capacitance of the system \(=C \varepsilon\).

So, initial energy of the system : \(U_{i}=\frac{1}{2}(C \varepsilon) V^{2}\)\\
and finally, energy of the capacitor : \(U_{f}=\frac{1}{2} C V^{2}\)\\
Hence capacitance energy increment,

\[
\Delta U=\frac{1}{2} C V^{2}-\frac{1}{2}(C \varepsilon) V^{2}=-\frac{1}{2} C V^{2}(\varepsilon-1)=-0.5 \mathrm{~mJ}
\]

From energy conservation

\[
\Delta U=A_{\text {cell }}+A_{\text {agent }}
\]

(as there is no heat liberation)\\
But \(A_{\text {cell }}=\left(C_{f}-C_{i}\right) V^{2}=(C-C \varepsilon) V^{2}\)

Hence \(A_{\text {agent }}=\Delta U-A_{\text {cell }}\)

\[
=\frac{1}{2} C(1-\varepsilon) V^{2}=0.5 \mathrm{~m} \mathrm{~J}
\]

3.202 If \(C_{0}\) is the initial capacitance of the condenser before water rises in it then

\[
U_{i}=\frac{1}{2} C_{0} V^{2}, \text { where } C_{0}=\frac{\varepsilon_{0} 2 l \pi R}{d}
\]

( \(R\) is the mean radius and \(l\) is the length of the capacitor plates.)\\
Suppose the liquid rises to a height \(h\) in it. Then the capacitance of the condenser is

\[
C=\frac{\varepsilon \varepsilon_{0} h 2 \pi R}{d}+\frac{\varepsilon(l-h) 2 \pi R}{d}=\frac{\varepsilon_{0} 2 \pi R}{d}(l+(\varepsilon-1) h)
\]

and energy of the capacitor and the liquid (including both gravitational and electrosatic contributions) is

\[
\frac{1}{2} \frac{\varepsilon_{0} 2 \pi R}{d}(l+(\varepsilon-1) h) V^{2}+\rho g(2 \pi R h d) \frac{h}{2}
\]

If the capacitor were not connected to a battery this energy would have to be minimized. But the capacitor is connected to the battery and, in effect, the potential energy of the whole system has to be minimized. Suppose we increase \(h\) by \(\delta h\). Then the energy of the capacitor and the liquid increases by

\[
\delta h\left(\frac{\varepsilon_{0} 2 \pi R}{2 d}(\varepsilon-1) V^{2}+\rho g(2 \pi R d) h\right)
\]

and that of the cell diminishes by the quantity \(A_{\text {cell }}\) which is the product of charge flown and \(V\)

\[
\delta h \frac{\varepsilon_{0}(2 \pi R)}{d}(\varepsilon-1) V^{2}
\]

In equilibrium, the two must balance; so

\[
\rho g d h=\frac{\varepsilon_{0}(\varepsilon-1) V^{2}}{2 d}
\]

Hence

\[
h=\frac{\varepsilon_{0}(\varepsilon-1) V^{2}}{2 \rho g d^{2}}
\]

3.203 (a) Let us mentally islolate a thin spherical layer with inner and outer radii \(r\) and \(r+d r\) respectively. Lines of current at all the points of this layer are perpendicular to it and therefore such a layer can be treated as a spherical conductor of thickness \(d r\) and cross sectional area \(4 \pi r^{2}\). Now, we know that resistance,


\begin{equation*}
d R=\rho \frac{d r}{S(r)}=\rho \frac{d r}{4 \pi r^{2}} \tag{1}
\end{equation*}


Integrating expression (1) between the limits,


\begin{equation*}
\int_{0}^{R} d R=\int_{a}^{b} \rho \frac{d r}{4 \pi r^{2}} \quad \text { or, } \quad R=\frac{\rho}{4 \pi}\left[\frac{1}{a}-\frac{1}{b}\right] \tag{2}
\end{equation*}


Capacitance of the network, \(C=\frac{4 \pi \varepsilon_{0} \varepsilon}{\left[\frac{1}{a}-\frac{1}{b}\right]}\)\\
and

\[
q=C \varphi\left[\begin{array}{c}
\text { where } q \text { is the charge }  \tag{4}\\
\text { at any arbitrary moment }
\end{array}\right]
\]

also,


\begin{equation*}
\varphi=\left(\frac{-d q}{d t}\right) R, \text { as capacitor is discharging. } \tag{5}
\end{equation*}


From Eqs. (2), (3), (4) and (5) we get,

\[
q=\frac{4 \pi \varepsilon_{0} \varepsilon}{\left[\frac{1}{a}-\frac{1}{b}\right]} \frac{\left[-\frac{d q}{d t}\right] \rho\left[\frac{1}{a}-\frac{1}{b}\right]}{4 \pi} \text { or, } \frac{d q}{q}=\frac{d t}{\rho \varepsilon \varepsilon_{0}}
\]

Integrating

\[
\int_{q_{0}}^{q}-\frac{d q}{q}=\frac{1}{\rho \varepsilon_{0} \varepsilon} \int_{0}^{t} d t=\frac{d t}{\rho \varepsilon \varepsilon_{0}}
\]

Hence

\[
q=q_{0} e^{\frac{-t}{\rho \varepsilon_{0} \varepsilon}}
\]

(b) From energy conservation heat generated, during the spreading of the charge,

\[
\begin{gathered}
H=U_{i}-U_{f}\left[\text { because } A_{\text {coll }}=0\right] \\
=\frac{1}{2} \frac{q_{0}^{2}}{4 \pi \varepsilon_{0} \varepsilon}\left[\frac{1}{a}-\frac{1}{b}\right]-0=\frac{q_{0}^{2}}{8 \pi \varepsilon_{0} \varepsilon} \frac{b-a}{a b}
\end{gathered}
\]

3.204 (a) Let, at any moment of time, charge on the plates be \(\left(q_{0}-q\right)\) then current through the resistor, \(i=-\frac{d\left(q_{0}-q\right)}{d t}\), because the capacitor is discharging.\\
or, \(\quad i=\frac{d q}{d t}\)\\
Now, applying loop rule in the circuit,

\[
i R-\frac{q_{0}-q}{C}=0
\]

or, \(\quad \frac{d q}{d t} R=\frac{q_{0}-q}{C}\)\\
or, \(\quad \frac{d q}{q_{0}-q}=\frac{1}{R C} d t\)\\
\includegraphics[max width=\textwidth, alt={}, center]{9f3c5a8e-ac5d-42f4-8eb6-ea2af623f2d8-355_428_542_1605_794}

At \(t=0, q=0\) and at \(t=\tau, q=q\)

So,

\[
\ln \frac{q_{0}-q}{q_{0}}=\frac{-\tau}{R C}
\]

Thus

\[
q=q_{0}\left(1-e^{-\tau / R C}\right)=0.18 \mathrm{mC}
\]

(b) Amount of heat generated \(=\) decrement in capacitance energy

\[
\begin{gathered}
=\frac{1}{2} \frac{q_{0}^{2}}{C}-\frac{1}{2} \frac{\left[q_{0}-q_{0}\left(1-e^{-\tau / R C}\right)\right]^{2}}{C} \\
=\frac{1}{2} \frac{q_{0}^{2}}{C}\left[1-e^{-\frac{2 \tau}{R C}}\right]=82 \mathrm{~mJ}
\end{gathered}
\]

3.205 Let, at any moment of time, charge flown be \(q\) then current \(i=\frac{d q}{d t}\)

Applying loop rule in the circuit, \(-\Delta \varphi=0\), we get :

\[
\frac{d q}{d t} I R-\frac{\left(C V_{0}-q\right)}{C}+\frac{q}{C}=0
\]

or, \(\quad \frac{d q}{C V_{0}-2 q}=\frac{1}{R C} d t\)\\
So, \(\ln \frac{C V_{0}-2 q}{C V_{0}}=-2 \frac{t}{R C}\) for \(0 \leq t \leq t\)\\
or, \(\quad q=\frac{C V_{0}}{2}\left(1-e^{\frac{-2 t}{R C}}\right)\)\\
\includegraphics[max width=\textwidth, alt={}, center]{9f3c5a8e-ac5d-42f4-8eb6-ea2af623f2d8-356_367_655_902_749}

Hence,

\[
i=\frac{d q}{d t}=\frac{C V_{0}}{2} \frac{2}{R C} e^{-2 t / R C}=\frac{V_{0}}{R} e^{-2 t / R C}
\]

Now, heat liberaled,

\[
Q=\int_{0}^{\infty} i^{2} R d t=\frac{V_{0}^{2}}{R^{2}} R \int_{0}^{\infty} e^{\frac{-4 t}{R C}} d t=\frac{1}{4} C V_{0}^{2}
\]

3.206 In a rotating frame, to first order in \(\omega\), the main effect is a coriolis force \(2 m \overrightarrow{v^{\prime}} \times \vec{\omega}\). This unbalanced force will cause electrons to react by setting up a magnetic field \(\vec{B}\) so that the magnetic force \(e \vec{v} \times \vec{B}\) balances the coriolis force.

Thus

\[
-\frac{e}{2 m} \vec{B}=\vec{\omega} \text { or, } \vec{B}=-\frac{2 m}{e} \vec{\omega}
\]

The flux associated with this is

\[
\Phi=N \pi r^{2} B=N \pi r^{2} \frac{2 m}{e} \omega
\]

where \(N=\frac{l}{2 \pi r}\) is the number of turns of the ring. If \(\omega\) changes (and there is time for the electron to rearrange) then \(B\) also changes and so does \(\Phi\). An emf will be induced and a current will flow. This is

\[
I=N \pi r^{2} \frac{2 m}{e} \omega / R
\]

The total charge flowing through the ballastic galvanometer, as the ring is stopped, is

So,

\[
\begin{aligned}
& q=N \pi r^{2} / \frac{2 m}{e} \omega / R \\
& \frac{e}{m}=\frac{2 N \pi r^{2} \omega}{q R}=\frac{l \omega r}{q R}
\end{aligned}
\]

3.207 Let, \(n_{0}\) be the total number of electorns then, total momentum of electorns,


\begin{equation*}
p=n_{0} m_{e} v_{d} \tag{1}
\end{equation*}


Now,


\begin{equation*}
I=\rho S_{x} v_{d}=\frac{n_{0} e}{V} S_{x} v_{d}=\frac{n e}{l} v_{d} \tag{2}
\end{equation*}


Here \(S_{x}=\) Cross sectional area, \(\rho=\) electron charge density, \(V=\) volume of sample From (1) and (2)

\[
p=\frac{m_{e}}{e} I l=0.40 \mu \mathrm{Ns}
\]

3.208 By definition\\
\(n e v_{d}=j\) (where \(v_{d}\) is the drift velocity, \(n\) is number density of electrons.)\\
Then

\[
\tau=\frac{l}{v_{d}}=\frac{n e l}{j}
\]

So distance actually travelled

\[
S=\langle v\rangle \tau=\frac{n e l\langle v\rangle}{j}
\]

(<v>= mean velocity of thermal motion of an electron)\\
3.209 Let, \(n\) be the volume density of electrons, then from \(I=\rho S_{x} v_{d}\),

\[
I=n e S_{x}|<\vec{v}>|=n e S_{x} \frac{l}{t}
\]

So,

\[
t=\frac{n e S_{x} l}{I}=3 \mu \mathrm{~s} .
\]

(b) Sum of electric forces

\[
\begin{aligned}
=|(n v) e \vec{E}|= & \mid n S \text { le } \rho \vec{j} \mid, \text { where } \rho \text { is resistivity of the material. } \\
& =n S \text { le } \rho \frac{I}{S}=n \text { el } \rho I=1.0 \mu \mathrm{~N}
\end{aligned}
\]

3.210 From Gauss theorem field strength at a surface of a cylindrical shape equals, \(\frac{\lambda}{2 \pi \varepsilon_{0} r}\), where \(\lambda\) is the linear charge density.

Now,


\begin{equation*}
e V=\frac{1}{2} m_{e} v^{2} \quad \text { or, } \quad v=\sqrt{\frac{2 e V}{m_{e}}} \tag{1}
\end{equation*}


Also,

\[
d q=\lambda d x \quad \text { so, } \quad \frac{d q}{d t}=\lambda \frac{d x}{d t}
\]

or,

\[
I=\lambda v \quad \text { or, } \lambda=\frac{I}{v}=\frac{I}{\sqrt{\frac{2 e V}{m_{e}}}} \text {, using (1) }
\]

Hence

\[
E=\frac{I}{2 \pi \varepsilon_{0} r} \sqrt{\frac{m_{e}}{2 e V}}=32 \mathrm{~V} / \mathrm{m}
\]

(b) For the point, inside the solid charged cylinder, applying Gauss' theorem,

\[
2 \pi r h E=\pi r^{2} h \frac{q}{\varepsilon_{0} \pi R^{2} l}
\]

or,

\[
E=\frac{q / l}{2 \pi \varepsilon_{0} R^{2}} r=\frac{\lambda r}{2 \pi \varepsilon_{0} R^{2}}
\]

So,from

\[
E=-\frac{d \varphi}{d r},
\]

\[
\int_{\varphi_{1}}^{\varphi_{2}}-d \varphi=\int_{0}^{R} \frac{\lambda}{2 \pi \varepsilon_{0} R^{2}} r d r
\]

or,

\[
\varphi_{1}-\varphi_{2}=\frac{\lambda}{2 \pi \varepsilon_{0} R^{2}}\left[\frac{r^{2}}{2}\right]_{0}^{R}=\frac{\lambda}{4 \pi \varepsilon_{0}}
\]

Hence,

\[
\varphi_{1}-\varphi_{2}=\frac{V I}{4 \pi \varepsilon_{0}} \sqrt{\frac{m_{e}}{2 e V}}=0.80 \mathrm{~V}
\]

3.211 Between the plates \(\varphi=a x^{4 / 3}\)\\
or,

\[
\frac{\partial \varphi}{\partial x}=a \times \frac{4}{3} x^{1 / 3}
\]

\[
\frac{d^{2} \varphi}{\partial x^{2}}=\frac{4}{9} a x^{-2 / 3}=-\rho / \varepsilon_{0}
\]

or,

\[
\rho=-\frac{4 \varepsilon_{0} a}{9} x^{-2 / 3}
\]

Let the charge on the electron be \(-e\),\\
then

\[
\frac{1}{2} m v^{2}-e \varphi=\text { Const. }=0
\]

as the electron is initially emitted with neligible energy.

So,

\[
\begin{gathered}
v^{2}=\frac{2 e \varphi}{m}, \quad v=\sqrt{\frac{2 e \varphi}{m}} \\
j=-\rho v=\frac{4 \varepsilon_{0} a}{9} \sqrt{\frac{2 \varphi}{m}} x^{-2 / 3} .
\end{gathered}
\]

( \(j\) is measued from the anode to cathode, so the - ve sign.)\\
\(3.212 E=\frac{V}{d}\)\\
So by the definition of the mobility\\
and

\[
\begin{aligned}
v^{+} & =u_{0}^{+} \frac{V}{d}, v^{-}=u_{0}^{-}-\frac{V}{d} \\
j & =\left(n_{+} u_{0}^{+}+n_{-} u_{0}^{-}\right) \frac{e V}{d}
\end{aligned}
\]

(The negative ions move towards the anode and the positive ion towards the cathode and the total current is the sum of the currents due to them.)\\
On the other hand, in equilibrium \(n_{+}=n_{-}\)\\
So,

\[
\begin{gathered}
n_{+}=n_{-}=\frac{I}{S} /\left(u_{0}^{+}+u_{0}^{-}\right) \frac{e V}{d} \\
=\frac{I d}{e V S\left(u_{0}^{+}+u_{0}^{-}\right)}=2.3 \times 10^{8} \mathrm{~cm}^{-3}
\end{gathered}
\]

3.213 Velocity \(=\) mobility × field\\
or, \(v=u \frac{V_{0}}{l} \sin \omega t\), which is positive for \(0 \leq \omega t \leq \pi\)\\
So, maximum displacement in one direction is

\[
x_{\max }=\int_{0}^{\pi} u \frac{V_{0}}{l} \sin \omega t d t=\frac{2 u V_{0}}{l \omega}
\]

At \(\omega=\omega_{0}, x_{\text {max }}=l\), so, \(\frac{2 u V_{0}}{l \omega}=l\)\\
Thus

\[
u=\frac{\omega l^{2}}{2 V_{0}}
\]

3.214 When the current is saturated, all the ions, produced, reach the plate.

Then,

\[
\dot{n}_{i}=\frac{I_{\mathrm{sat}}}{e V}=6 \times 10^{9} \mathrm{~cm}^{-3} \mathrm{~s}^{-1}
\]

(Both positive ions and negative ions are counted here)\\
The equation of balance is, \(\frac{d n}{d t}=\dot{n}_{i}-r n^{2}\)

The first term on the right is the production rate and the second term is the recombination rate which by the usual statistical arguments is proportional to \(n^{2}(=\) no of positive ions × no. of -ve ion). In equilibrium,

\[
\frac{d n}{d t}=0
\]

so,

\[
n_{e q}=\sqrt{\frac{\dot{n}_{i}}{r}}=6 \times 10^{7} \mathrm{~cm}^{-3}
\]

3.215 Initially \(n=n_{0}=\sqrt{\dot{n}_{i} / r}\)

Since we can assume that the long exposure to the ionizer has caused equilibrium to be set up. Afer the ionizer is switched off,

\[
\frac{d n}{d t}=-r n^{2}
\]

or

\[
r d t=-\frac{d n}{n^{2}}, \quad \text { or }, \quad r t=\frac{1}{n}+\text { constant }
\]

But

\[
n=n_{0} \text { at } t=0, \text { so, } r t=\frac{1}{n}-\frac{1}{n_{0}}
\]

The concentration will decrease by a factor \(\eta\) when\\
or,

\[
\begin{aligned}
r t_{0} & =\frac{1}{n_{0} / \eta}-\frac{1}{n_{0}}=\frac{\eta-1}{n_{0}} \\
t_{0} & =\frac{\eta-1}{\sqrt{r \dot{n}_{i}}}=13 \mathrm{~ms}
\end{aligned}
\]

3.216 Ions produced will cause charge to decay. Clearly,\\
\(\eta C V=\) decrease of charge \(=\dot{n}_{i} e A d t=\frac{\varepsilon_{0} A}{d} V \eta\)\\
or,

\[
t=\frac{\varepsilon_{0} V \eta}{\dot{n}_{i} e d^{2}}=4.6 \text { days }
\]

Note, that \(n_{i}\), here, is the number of ion pairs produced.\\
3.217 If \(v=\) number of electrons moving to the anode at distance \(x\), then

\[
\frac{d v}{d x}=\alpha v \text { or } v=v_{0} e^{\alpha x}
\]

Assuming saturation, \(I=e v_{0} e^{\alpha d}\)\\
3.218 Since the electrons are produced uniformly through the volume, the total current attaining saturation is clearly,

\[
I=\int_{0}^{d} e\left(\dot{n}_{i} A d x\right) e^{\alpha x}=e \dot{n}_{i} A\left(\frac{e^{\alpha d}-1}{\alpha}\right)
\]

Thus,

\[
j=\frac{I}{A}=e \dot{n}_{i}\left(\frac{e^{\alpha d}-1}{\alpha}\right)
\]

\subsection*{3.5 CONSTANT MAGNETIC FIELD. MAGNETICS}
3.219 (a) From the Biot - Savart law,

\[
\begin{gathered}
d \vec{B}=\frac{\mu_{0}}{4 \pi} i \frac{d \vec{l} \times \vec{r}}{r^{3}}, \text { so } \\
d B=\frac{\mu_{0}}{4 \pi} i \frac{(R d \theta) R}{R^{3}}(\text { as } d \vec{l} \perp \vec{r})
\end{gathered}
\]

From the symmetry

\[
B=\int d B=\int_{0}^{2 \pi} \frac{\mu_{0}}{4 \pi} \frac{i}{R} d \theta=\frac{\mu_{0}}{2} \frac{i}{R}=6.3 \mu \mathrm{~T}
\]

(b) From Biot-Savart's law :\\
\(\vec{B}=\frac{\mu_{0}}{4 \pi} i \int \frac{d \vec{l} \times \vec{r}}{r^{3}}\) (here \(\left.\vec{r}=\vec{R}+\vec{x}\right)\)\\
\includegraphics[max width=\textwidth, alt={}, center]{9f3c5a8e-ac5d-42f4-8eb6-ea2af623f2d8-361_591_433_274_772}

So, \(\quad \vec{B}=\frac{\mu_{0}}{4 \pi} i[\oint d \vec{l} \times \vec{R}+\oint d \vec{l} \times \vec{x}]\)\\
Since \(\vec{x}\) is a constant vector and \(|\vec{R}|\) is also constant\\
So, \(\oint d \vec{l} \times \vec{x}=(\oint d \vec{l}) \times \vec{x}=0\) (because \(\oint d \vec{l}=0\) )\\
and

\[
\begin{aligned}
& \oint d \vec{l} \times \vec{R}=\oint R d l \vec{n} \\
& =\vec{n} R \oint d l=2 \pi R^{2} \vec{n}
\end{aligned}
\]

Here \(\vec{n}\) is a unit vector perpendicular to the plane containing the current loop (Fig.) and in the direction of \(\overrightarrow{\boldsymbol{x}}\)

Thus we get

\[
\vec{B}=\frac{\mu_{0}}{4 \pi} \frac{2 \pi R^{2} i}{\left(x^{2}+R^{2}-\right)^{3 / 2}} \vec{n}
\]

3.220 As \(\angle A O B=\frac{2 \pi}{n}, O C\) or perpendicular distance of any segment from centre equals \(R \cos \frac{\pi}{n}\). Now magnetic induction at \(O\), due to the right current carrying element \(A B\)

\[
=\frac{\mu_{0}}{4 \pi} \frac{i}{R \cos \frac{\pi}{n}} 2 \sin \frac{\pi}{n}
\]

(From Biot-Savart's law, the magnetic field at \(O\) due to any section such as \(A B\) is perpendicular to the plane of the figure and has the magnitude.)

\[
B=\int \frac{\mu_{0}}{4 \pi} i \frac{d x}{r^{2}} \cos \theta=\int_{-\frac{\pi}{n}}^{\frac{\pi}{n}} \frac{\mu_{0} i}{4 \pi} \frac{R \cos \frac{\pi}{n} \sec ^{2} \theta d \theta}{R^{2} \cos \frac{2 \pi}{n} \sec ^{2} \theta} \cos \theta=\frac{\mu_{0} i}{4 \pi} \frac{1}{R \cos \frac{\pi}{n}} 2 \sin \frac{\pi}{n}
\]

As there are \(n\) number of sides and magnetic induction vectors, due to each side at \(O\), are equal in magnitude and direction. So,

\[
\begin{aligned}
B_{0}= & \frac{\mu_{0}}{4 \pi} \frac{n i}{R \cos \frac{\pi}{n}} 2 \sin \frac{\pi}{n} \cdot n \\
= & \frac{\mu_{0}}{2 \pi} \frac{n i}{R} \tan \frac{\pi}{n} \text { and for } n \rightarrow \infty \\
B_{0}= & \frac{\mu_{0}}{2} \frac{i}{R} \underset{n \rightarrow \infty}{L t}\left(\frac{\tan \frac{\pi}{n}}{\pi / n}\right)=\frac{\mu_{0}}{2} \frac{i}{R}
\end{aligned}
\]

\includegraphics[max width=\textwidth, alt={}, center]{9f3c5a8e-ac5d-42f4-8eb6-ea2af623f2d8-362_337_377_482_604}\\
\includegraphics[max width=\textwidth, alt={}, center]{9f3c5a8e-ac5d-42f4-8eb6-ea2af623f2d8-362_474_418_377_999}\\
3.221 We know that magnetic induction due to a straight current carrying wire at any point, at a perpendicular distance from it is given by :

\[
B=\frac{\mu_{0}}{4 \pi} \frac{i}{r}\left(\sin \theta_{1}+\sin \theta_{2}\right),
\]

where \(r\) is the perpendicular distance of the wire from the point, considered, and \(\theta_{1}\) is the angle between the line, joining the upper point of straight wire to the considered point and the perpendicular drawn to the wire and \(\theta_{2}\) that from the lower point of the straight wire.

Here,

\[
B_{1}=B_{3}=\frac{\mu_{0}}{4 \pi} \frac{i}{(d / 2) \sin \frac{\varphi}{2}}\left\{\cos \frac{\varphi}{2}+\cos \frac{\varphi}{2}\right\}
\]

and

\[
B_{2}=B_{4}=\frac{\mu_{0}}{4 \pi} \frac{i}{(d / 2) \cos \frac{\varphi}{2}}\left(\sin \frac{\varphi}{2}+\sin \frac{\varphi}{2}\right)
\]

Hence, the magnitude of total magnetic induction at \(O\),

\[
\begin{gathered}
B_{0}=B_{1}+B_{2}+B_{3}+B_{4} \\
=\frac{\mu_{0}}{4 \pi} \frac{4 i}{d / 2}\left[\frac{\cos \frac{\varphi}{2}}{\sin \frac{\varphi}{2}}+\frac{\sin \frac{\varphi}{2}}{\cos \frac{\varphi}{2}}\right] \\
=\frac{4 \mu_{0} i}{\pi d \sin \varphi}=0.10 \mathrm{mT}
\end{gathered}
\]

\includegraphics[max width=\textwidth, alt={}, center]{9f3c5a8e-ac5d-42f4-8eb6-ea2af623f2d8-362_413_473_1611_892}\\
3.222 Magnetic induction due to the arc segment at \(O\),

\[
B_{\mathrm{arc}}=\frac{\mu_{0}}{4 \pi} \frac{i}{R}(2 \pi-2 \varphi)
\]

and magnetic induction due to the line segment at \(O\),

\[
B_{\text {line }}=\frac{\mu_{0}}{4 \pi} \frac{i}{R \cos \varphi}[2 \sin \varphi]
\]

So, total magnetic induction at \(O\),\\
\includegraphics[max width=\textwidth, alt={}, center]{9f3c5a8e-ac5d-42f4-8eb6-ea2af623f2d8-363_316_379_152_818}\\
\(B_{0}=B_{\text {arc }}+B_{\text {line }}=\frac{\mu_{0}}{2 \pi} \frac{i}{R}[\pi-\varphi+\tan \varphi]=28 \mu \mathrm{~T}\)\\
3.223 (a) From the Biot-Savart law,

\[
d B=\frac{\mu_{0}}{4 \pi} i \frac{(d \vec{l} \times \vec{r})}{r^{3}}
\]

So, magnetic field induction due to the segment 1 at \(O\),

\[
B_{1}=\frac{\mu_{0}}{4 \pi} \frac{i}{a}(2 \pi-\varphi)
\]

also

\[
B_{2}=B_{4}=0 \text {, as } d \vec{l} \uparrow \uparrow \vec{r}
\]

and

\[
B_{3}=\frac{\mu_{0}}{4 \pi} \frac{i}{b} \varphi
\]

Hence, \(B_{0}=B_{1}+B_{2}+B_{3}+B_{4}\)\\
\includegraphics[max width=\textwidth, alt={}, center]{9f3c5a8e-ac5d-42f4-8eb6-ea2af623f2d8-363_373_396_789_847}

\[
=\frac{\mu_{0}}{4 \pi} i\left[\frac{2 \pi-\varphi}{a}+\frac{\varphi}{b}\right],
\]

(b) Here, \(B_{1}=\frac{\mu_{0}}{4 \pi} \frac{i}{a} \frac{3 \pi}{a}, \vec{B}_{2}=0\),

\[
\begin{gathered}
B_{3}=\frac{\mu_{0}}{4 \pi} \frac{i}{b} \sin 45^{\circ} \\
B_{4}=\frac{\mu_{0}}{4 \pi} \frac{i}{b} \sin 45^{\circ} \\
B_{5}=0
\end{gathered}
\]

and\\
So, \(\quad B_{0}=B_{1}+B_{2}+B_{3}+B_{4}+B_{5}\)\\
\includegraphics[max width=\textwidth, alt={}, center]{9f3c5a8e-ac5d-42f4-8eb6-ea2af623f2d8-363_443_387_1279_779}

\[
\begin{gathered}
=\frac{\mu_{0}}{4 \pi} \frac{i}{a} \frac{3 \pi}{2}+0+\frac{\mu_{0}}{4 \pi} \frac{i}{b} \sin 45^{\circ}+\frac{\mu_{0}}{4 \pi} \frac{i}{b} \sin 45^{\circ}+0 \\
=\frac{\mu_{0}}{4 \pi} i\left[\frac{3 \pi}{2 a}+\frac{\sqrt{2}}{b}\right]
\end{gathered}
\]

3.224 The thin walled tube with a longitudinal slit can be considered equivalent to a full tube and a strip carrying the same current density in the opposite direction. Inside the tube, the former does not contribute so the total magnetic field is simply that due to the strip. It is

\[
B=\frac{\mu_{0}}{2 \pi} \frac{(I / 2 \pi R) h}{r}=\frac{\mu_{0} I h}{4 \pi^{2} R r}
\]

where \(r\) is the distance of the field point from the strip.\\
3.225 First of all let us find out the direction of vector \(\vec{B}\) at point \(O\). For this purpose, we divide the entire conductor into elementary fragments with current di. It is obvious that the sum of any two symmetric fragments gives a resultant along \(\vec{B}\) shown in the figure and consequently, vector \(\vec{B}\) will also be directed as shown

So,

\[
\begin{aligned}
& |\vec{B}|=\int d B \sin \varphi \\
& =\int \frac{\mu_{0}}{2 \pi R} d i \sin \varphi
\end{aligned}
\]

\(=\int_{0}^{\pi} \frac{\mu_{0}}{2 \pi^{2} R} i \sin \varphi d \varphi,\left(\right.\) as \(\left.d i=\frac{i}{\pi} d \varphi\right)\)\\
Hence

\[
B=\mu_{0} i / \pi^{2} R
\]

3.226 (a) From symmetry

\[
\begin{gathered}
B_{0}=B_{1}+B_{2}+B_{3} \\
=0+\frac{\mu_{0}}{4 \pi} \frac{i}{R} \pi+0=\frac{\mu_{0}}{4} \frac{i}{R}
\end{gathered}
\]

(b) From symmetry

\[
\begin{gathered}
B_{0}=B_{1}+B_{2}+B_{3} \\
=\frac{\mu_{0}}{4 \pi} \frac{i}{R}+\frac{\mu_{0}}{2 \pi} \frac{i}{R} \frac{3 \pi}{2}+0=\frac{\mu_{0}}{4 \pi} \frac{i}{R}\left[1+\frac{3 \pi}{2}\right]
\end{gathered} .
\]

(c) From symmetry

\[
\begin{gathered}
B_{0}=B_{1}+B_{2}+B_{3} \\
=\frac{\mu_{0}}{4 \pi} \frac{i}{R}+\frac{\mu_{0}}{4 \pi} \frac{i}{R} \pi+\frac{\mu_{0}}{4 \pi} \frac{i}{R}=\frac{\mu_{0}}{4 \pi} \frac{i}{R}(2+\pi)
\end{gathered}
\]

\(3.227 \overrightarrow{B_{0}}=\overrightarrow{B_{1}}+\overrightarrow{B_{2}}\)\\
or, \(\left|\overrightarrow{B_{0}}\right|=\frac{\mu_{0} i}{4 \pi l} \sqrt{2}=2.0 \mu \mathrm{~T}\), (using 3.221)\\
\includegraphics[max width=\textwidth, alt={}, center]{9f3c5a8e-ac5d-42f4-8eb6-ea2af623f2d8-364_317_497_678_859}\\
\includegraphics[max width=\textwidth, alt={}, center]{9f3c5a8e-ac5d-42f4-8eb6-ea2af623f2d8-364_144_486_1001_861}\\
\includegraphics[max width=\textwidth, alt={}, center]{9f3c5a8e-ac5d-42f4-8eb6-ea2af623f2d8-364_201_488_1157_865}\\
\includegraphics[max width=\textwidth, alt={}, center]{9f3c5a8e-ac5d-42f4-8eb6-ea2af623f2d8-364_188_486_1384_865}\\
\includegraphics[max width=\textwidth, alt={}, center]{9f3c5a8e-ac5d-42f4-8eb6-ea2af623f2d8-364_265_272_1636_986}\\
3.228 (a)

\[
\begin{aligned}
\vec{B}_{0} & =\vec{B}_{1}+\vec{B}_{2}+\vec{B}_{3} \\
& =\frac{\mu_{0}}{4 \pi} \frac{i}{R}(-\vec{k})+\frac{\mu_{0}}{4 \pi} \frac{i}{R} \pi(-\vec{i})+\frac{\mu_{0}}{4 \pi} \frac{i}{R}(-\vec{k}) \\
& =-\frac{\mu_{0}}{4 \pi} \frac{i}{R}[2 \vec{k}+\pi \vec{i}]
\end{aligned}
\]

So,

\[
\left|\vec{B}_{0}\right|=\frac{\mu_{0}}{4 \pi} \frac{i}{R} \sqrt{\pi^{2}+4}=0.30 \mu \mathrm{~T}
\]

(b) \(\overrightarrow{B_{0}}=\overrightarrow{B_{1}}+\overrightarrow{B_{2}}+\overrightarrow{B_{3}}\)

\[
\begin{aligned}
& =\frac{\mu_{0}}{4 \pi} \frac{i}{R}(-\vec{k})+\frac{\mu_{0}}{4 \pi} \frac{i}{R} \pi(-\vec{i})+\frac{\mu_{0}}{4 \pi} \frac{i}{R}(-\vec{i}) \\
& =-\frac{\mu_{0}}{4 \pi} \frac{i}{R}[\vec{k}+(\pi+1) \vec{i}]
\end{aligned}
\]

So,\\
\(\left|\vec{B}_{0}\right|=\frac{\mu_{0}}{4 \pi} \frac{i}{R} \sqrt{1+(\pi+1)^{2}}=0.34 \mu \mathrm{~T}\)\\
(c) Here using the law of parallel resistances

\[
i_{1}+i_{2}=i \text { and } \frac{i_{1}}{i_{2}}=\frac{1}{3},
\]

So,

\[
\frac{i_{1}+i_{2}}{i_{2}}=\frac{4}{3}
\]

Hence

\[
i_{2}=\frac{3}{4} i, \text { and } i_{1}=\frac{1}{4} i
\]

Thus \(\vec{B}_{0}=\frac{\mu_{0}}{4 \pi} \frac{i}{R}(-\vec{k})+\frac{\mu_{0}}{4 \pi} \frac{i}{R}(-\vec{j})+\left[\frac{\mu_{0}}{4 \pi}\left(\frac{3 \pi}{2}\right) \frac{i_{1}}{R}(-\vec{i})+\frac{\mu_{0}}{4 \pi} \frac{(\pi / 2) i_{2}}{R} \vec{i}\right]\)

Thus,

\[
\begin{aligned}
& =\frac{-\mu_{0}}{4 \pi} \frac{i}{R}(\vec{j}+\vec{k})+0 \\
\left|\vec{B}_{0}\right| & =\frac{\mu_{0}}{4 \pi} \frac{\sqrt{2} i}{R}=0.11 \mu \mathrm{~T}
\end{aligned}
\]

3.229 (a) We apply circulation theorem as shown. The current is vertically upwards in the plane and the magnetic field is horizontal and parallel to the plane.\\
\(\oint \vec{B} \cdot d \vec{l}=2 B l=\mu_{0} i l\) or, \(B=\frac{\mu_{0} i}{2}\)\\
(b) Each plane contributes \(\mu_{0} \frac{i}{2}\) between the planes and outside the plane that cancel.\\
Thus

\[
B=\left\{\begin{array}{c}
\mu_{0} i \text { between the plane } \\
0 \text { outside. }
\end{array}\right.
\]

\includegraphics[max width=\textwidth, alt={}, center]{9f3c5a8e-ac5d-42f4-8eb6-ea2af623f2d8-365_423_446_1651_839}\\
3.230 We assume that the current flows perpendicular to the plane of the paper, by circulation theorem,

\[
2 B d l=\mu_{0}(2 x d l) j
\]

or,

\[
B=\mu_{0} x j,|x| \leq d
\]

Outside, \(2 B d l=\mu_{0}(2 d d l) j\)\\
or,

\[
B=\mu_{0} d j|x| \geq d .
\]

\includegraphics[max width=\textwidth, alt={}, center]{9f3c5a8e-ac5d-42f4-8eb6-ea2af623f2d8-366_351_197_97_986}\\
3.231 It is easy to convince oneself that both in the regions. 1 and 2 , there can only be a circuital magnetic field (i.e. the component \(B_{\varphi}\) ). Any radial field in region 1 or any \(B_{Z}\) away from the current plane will imply a violation of Gauss' law of magnetostatics, \(B_{\varphi}\) must obviously be symmetrical about the straight wire. Then in 1,\\
or,

\[
B_{\varphi} 2 \pi r=\mu_{0} I
\]

\[
B_{\varphi}=\frac{\mu_{0}}{2 \pi} \frac{I}{r}
\]

In 2, \(\quad B_{\varphi} \cdot 2 \pi r=0\), or \(B_{\varphi}=0\)\\
\includegraphics[max width=\textwidth, alt={}, center]{9f3c5a8e-ac5d-42f4-8eb6-ea2af623f2d8-366_276_420_648_884}\\
3.232 On the axis, \(B=\frac{\mu_{0} I R^{2}}{2\left(R^{2}+x^{2}\right)^{3 / 2}}=B_{x}\) along the axis.

Thus,

\[
\begin{gathered}
\int \vec{B} \cdot d \vec{r}=\int_{-\infty}^{\infty} B_{x} d x=\frac{\mu_{0} I R^{2}}{2} \int_{-\infty}^{\infty} \frac{d x}{\left(R^{2}+x^{2}\right)^{3 / 2}} \\
=\frac{\mu_{0} I R^{2}}{2} \int_{-\pi / 2}^{\pi / 2} \frac{R \sec ^{2} \theta d \theta}{R^{3} \sec ^{3} \theta}, \text { on putting } x=R \tan \theta \\
=\mu_{0} I \frac{1}{2} \int_{-\pi / 2}^{\pi / 2} \cos \theta d \theta=\mu_{0} I
\end{gathered}
\]

The physical interpretation of this result is that \(\int_{-\infty}^{\infty} B_{x} d x\) can be thought of as the circulation of \(B\) over a closed loop by imaging that the two ends of the axis are connected, by a line at infinity (e.g. a semicircle of infinite radius).\\
3.233 By circulation theorem inside the conductor

\[
B_{\varphi} 2 \pi r=\mu_{0} j_{z} \pi r^{2} \quad \text { or, } \quad B_{\varphi}=\mu_{0} j_{z} r / 2
\]

i.e.,

\[
\vec{B}=\frac{1}{2} \mu_{0} \vec{j} \times \vec{r}
\]

Similarly outside the conductor,

\[
B_{\varphi} 2 \pi r=\mu_{0} j_{z} \pi R^{2} \quad \text { or, } \quad B_{\varphi}=\frac{1}{2} \mu_{0} j_{z} \frac{R^{2}}{r}
\]

So,

\[
\vec{B}=\frac{1}{2} \mu_{0}(\vec{j} \times \vec{r}) \frac{R^{2}}{r^{2}}
\]

2.234 We can think of the given current which will be assumed uniform, as arising due to a negative current, flowing in the cavity, superimposed on the true current, everywhere including the cavity. Then from the previous problem, by superposition.

\[
\vec{B}=\frac{1}{2} \mu_{0} \vec{j} \times(A \vec{P}-B \vec{P})=\frac{1}{2} \mu_{0} \vec{j} \times \vec{l}
\]

If \(\vec{l}\) vanishes so that the cavity is concentric\\
\includegraphics[max width=\textwidth, alt={}, center]{9f3c5a8e-ac5d-42f4-8eb6-ea2af623f2d8-367_337_332_371_832}\\
with the conductor, there is no magnetic field in the cavity.\\
3.235 By Circulation theorem,

\[
B_{\varphi} \cdot 2 \pi r=\mu_{0} \int_{0}^{r} j\left(r^{\prime}\right) \times 2 \pi r^{\prime} d r^{\prime}
\]

or using \(B_{\varphi}=b r^{\alpha}\) inside the stream,

\[
b r^{\alpha+1}=\mu_{0} \int_{0}^{r} j\left(r^{\prime}\right) r^{\prime} d r^{\prime}
\]

So by differentiation,

Hence,

\[
\begin{aligned}
& (\alpha+1) b r^{\alpha}=\mu_{0} j(r) r \\
& j(r)=\frac{b(\alpha+1)}{\mu_{0}} r^{\alpha-1}
\end{aligned}
\]

3.236 On the surface of the solenoid there is a surface current density

\[
\overrightarrow{j_{s}}=n I \hat{e}_{\varphi}
\]

Then,

\[
\vec{B}=-\frac{\mu_{0}}{4 \pi} n I \int R d \varphi d z \frac{\hat{e}_{\varphi} \times \vec{r}_{0}}{r_{0}^{3}}
\]

where \(\overrightarrow{r_{0}}\) is the vector from the current element to the field point, which is the centre of the solenoid,\\
Now,

\[
\begin{aligned}
& -\hat{e}_{\varphi} \times \vec{r}_{0}=R \hat{e}_{z} \\
& r_{0}=\left(z^{2}+R^{2}\right)^{1 / 2}
\end{aligned}
\]

Thus,

\[
B=B_{z}=\frac{\mu_{0} n I}{4 \pi} \times 2 \pi R^{2} \int_{-l / 2}^{-l / 2} \frac{d z}{\left(R^{2}+z^{2}\right)^{3 / 2}}
\]

\[
\begin{aligned}
= & \left.\frac{1}{2} \mu_{0} n I \int_{-\tan ^{-1} \frac{l}{2 R}}^{+\tan ^{-1} \frac{l}{2 R}} \cos \alpha d \alpha \quad \text { (on putting } z=R \tan \alpha\right) \\
= & \mu_{0} n I \sin \alpha=\mu_{0} n I \frac{l / 2}{\sqrt{(l / 2)^{2}+R^{2}}}=\mu_{0} n I / \sqrt{1+\left(\frac{2 R}{l}\right)^{2}}
\end{aligned}
\]

3.237 We proceed exactly as in the previous problem. Then (a) the magnetic induction on the axis at a distance \(x\) from one end is clearly,

\[
\begin{aligned}
B & =\frac{\mu_{n} n I}{4 \pi} \times 2 \pi R^{2} \int_{0}^{\infty} \frac{d z}{\left[R^{2}+(z-x)^{2}\right]^{3 / 2}}=\frac{1}{2} \mu_{0} n I R^{2} \int_{x}^{\infty} \frac{d z}{\left(z^{2}+R^{2}\right)^{3 / 2}} \\
& =\frac{1}{2} \mu_{0} n I \int_{\tan ^{-1} \frac{x}{R}}^{\pi / 2} \cos \theta d \theta=\frac{1}{2} \mu_{0} n I\left(1-\frac{x}{\sqrt{x^{2}+R^{2}}}\right)
\end{aligned}
\]

\(x>0\) menas that the field point is outside the solenoid. \(B\) then falls with \(x . x<0\) means that the field point gets more and more inside the solenoid. \(B\) then increases with \((x)\) and eventually becomes constant, equal to \(\mu_{0} n I\). The \(B-x\) graph is as given in the answer script.\\
(b) We have, \(\frac{B_{0}-\delta B}{B_{0}}=\frac{1}{2}\left[1-\frac{x_{0}}{\sqrt{R^{2}+x_{0}^{2}}}\right]=1-\eta\)\\
or,

\[
-\frac{x_{0}}{\sqrt{R^{2}+x_{0}^{2}}}=1-2 \eta
\]

Since \(\eta\) is small \((\approx 1 \%), x_{0}\) must be negative. Thus \(x_{0}=-\left|x_{0}\right|\)\\
and

\[
\begin{aligned}
\frac{\left|x_{0}\right|}{\sqrt{R^{2}+\left|x_{0}\right|^{2}}} & =1-2 \eta \\
\left|x_{0}\right|^{2} & =\left(1-4 \eta+4 \eta^{2}\right)\left(R^{2}+\left|x_{0}\right|^{2}\right) \\
0 & =(1-2 \eta)^{2} R^{2}-4 \eta(1-\eta)\left|x_{0}\right|^{2} \\
\left|x_{0}\right| & =\frac{(1-2 \eta) R}{2 \sqrt{\eta(1-\eta)}}
\end{aligned}
\]

3.238 If the strip is tightly wound, it must have a pitch of \(h\). This means that the current will flow obliquely, partly along \(\hat{e}_{\varphi}\) and partly along \(\hat{e}_{z^{\circ}}\) Obviously, the surface current density is,

\[
\vec{J}_{s}=\frac{I}{h}\left[\sqrt{1-(h / 2 \pi R)^{2}} \hat{e}_{\varphi}+\frac{h}{2 \pi R} \hat{e}_{z}\right]
\]

By comparision with the case of a solenoid and a hollow straight conductor, we see that field inside the coil

\[
=\mu_{0} \frac{I}{h} \sqrt{1-(h / 2 \pi R)^{2}}
\]

(Cf. \(B=\mu_{0} n I\) ).\\
Outside, only the other term contributes, so\\
or,

\[
\begin{gathered}
B_{\varphi} \times 2 \pi r=\mu_{0} \frac{I}{h} \times \frac{h}{2 \pi R} \times 2 \pi R \\
B_{\varphi}=\frac{\mu_{0}}{4 \pi} \cdot \frac{2 I}{r} .
\end{gathered}
\]

Note - Surface current density is defined as current flowing normally across a unit length over a surface.\\
3.239 Suppose \(a\) is the radius of cross section of the core. The winding has a pitch \(2 \pi R / N\), so the surface current density is

\[
\overrightarrow{J_{s}}=\frac{I}{2 \pi R / N} \overrightarrow{e_{1}}+\frac{I}{2 \pi a} \overrightarrow{e_{2}}
\]

where \(\overrightarrow{e_{1}}\) is a unit vector along the cross section of the core and \(\overrightarrow{e_{2}}\) is a unit vector along its length.\\
The magnetic field inside the cross section of the core is due to first term above, and is given by

\[
B_{\varphi} \cdot 2 \pi R=\mu_{0} N I
\]

( \(N I\) is total current due to the above surface current (first term.))\\
Thus,

\[
B_{\varphi}=\mu_{0} N I / 2 \pi R
\]

The magnetic field at the centre of the core can be obtained from the basic formula.

\[
d \vec{B}=\frac{\mu_{0}}{4 \pi} \frac{\overrightarrow{J_{s}} \times \overrightarrow{r_{0}}}{r_{0}^{3}} d S \text { and is due to the second term. }
\]

So,

\[
\vec{B}=B_{z} \vec{e}_{z}=\vec{e}_{z} \frac{\mu_{0}}{4 \pi} \frac{I}{2 \pi a} \int \frac{1}{R^{3}} R d \varphi \times 2 \pi a
\]

or, \(\quad B_{z}=\frac{\mu_{0} I}{2 R}\)\\
The ratio of the two magnetic field, is \(=\frac{N}{\pi}\)\\
3.240 We need the flux through the shaded area.

Now by Ampere's theorem',\\
or,

\[
\begin{gathered}
B_{\varphi} 2 \pi r=\mu_{0} \frac{I}{\pi R^{2}} \cdot \pi r^{2} \\
B_{\varphi}=\frac{\mu_{0}}{2 \pi} I \frac{r}{R^{2}}
\end{gathered}
\]

The flux through the shaded region is,\\
\includegraphics[max width=\textwidth, alt={}, center]{9f3c5a8e-ac5d-42f4-8eb6-ea2af623f2d8-369_489_340_1516_934}

\[
\begin{aligned}
\varphi_{1} & =\int_{0}^{R} 1 \cdot d r \cdot B_{\varphi}(r) \\
& =\int_{0}^{R} d r \frac{\mu_{0}}{2 \pi} I \frac{r}{R^{2}}=\frac{\mu_{0}}{4 \pi} I .
\end{aligned}
\]

3.241 Using 3.237, the magnetic field is given by,

\[
B=\frac{1}{2} \mu_{0} n I\left(1-\frac{x}{\sqrt{x^{2}+R^{2}}}\right)
\]

At the end, \(B=\frac{1}{2} \mu_{0} n I=\frac{1}{2} B_{0}\), where \(B_{0}=\mu_{0} n I\), is the field deep inside the solenoid. Thus,

\[
\Phi=\frac{1}{2} \mu_{0} n I S=\Phi_{0} / 2, \text { where } \Phi=\mu_{0} n I S
\]

is the flux of the vector \(B\) through the cross section deep inside the solenid.\\
3.242

\[
\begin{aligned}
B_{\varphi} 2 \pi r & =\mu_{0} N I \\
B_{\varphi_{1}} & =\frac{\mu_{0} N I}{2 \pi r}
\end{aligned}
\]

or,

\[
\Phi=\int_{a}^{b} B_{\varphi} h d r, a \leq r \leq b=\frac{\mu_{0}}{4 \pi} 2 N l h \ln \eta \text {, where } \eta=b / a
\]

3.243 Magnetic moment of a current loop is given by \(p_{m}=n i S\) (where \(n\) is the number of turns and \(S\), the cross sectional area.) In our problem, \(n=1, S=\pi R^{2}\) and \(B=\frac{\mu_{0}}{2} \frac{i}{R}\)

So,

\[
p_{m}=\frac{2 B R}{\mu_{0}} \pi R^{2}=\frac{2 \pi B R^{3}}{\mu_{0}}
\]

3.244 Take an element of length \(r d \theta\) containing \(\frac{N}{\pi r} \cdot r d \theta\) turns. Its magnetic moment is

\[
\frac{N}{\pi} d \theta \cdot \frac{\pi}{4} d^{2} I
\]

normal to the plane of cross section. We resolve it along \(O A\) and \(O B\). The moment along \(O A\) integrates to

\[
\int_{0}^{\pi} \frac{N}{4} d^{2} I d \theta \cos \theta=0
\]

while that along \(O B\) gives

\[
p_{m}=\int_{0}^{\pi} \frac{N d^{2} I}{4} \sin \theta d \theta=\frac{1}{2} N d^{2} I
\]

\includegraphics[max width=\textwidth, alt={}, center]{9f3c5a8e-ac5d-42f4-8eb6-ea2af623f2d8-370_352_487_1654_849}\\
3.245 (a) From Biot-Savart's law, the magnetic induction due to a circular current carrying wire loop at its centre is given by,

\[
B_{r}=\frac{\mu_{0}}{2 r} i
\]

The plane spiral is made up of concentric circular loops, having different radii, varying from \(a\) to \(b\). Therefore, the total magnetic induction at the centre,


\begin{equation*}
B_{0}=\int \frac{\mu_{0}}{2 r} d N \tag{1}
\end{equation*}


where \(\frac{\mu_{0}}{2 r} i\) is the contribution of one turn of radius \(r\) and \(d N\) is the number of turns in the interval ( \(r, r+d r\) )\\
i.e.

\[
d N=\frac{N}{b-a} d r
\]

Substituting in equation (1) and integrating the result over \(r\) between \(a\) and \(b\), we obtain,

\[
B_{0}=\int_{a}^{b} \frac{\mu_{0} i}{2 r} \frac{N}{(b-a)} d r=\frac{\mu_{0} i N}{2(b-a)} \ln \frac{b}{a}
\]

(b) The magnetic moment of a turn of radius \(r\) is \(p_{m}=i \pi r^{2}\) and of all turns,

\[
p=\int p_{m} d N=\int_{a}^{b} i \pi r^{2} \frac{N}{b-a} d r=\frac{\pi i N\left(b^{3}-a^{2}\right)}{3(b-a)}
\]

3.246 (a) Let us take a ring element of radius \(r\) and thickness \(d r\), then charge on the ring element., \(d q=\sigma 2 \pi r d r\)\\
and current, due to this element, \(d i=\frac{(\sigma 2 \pi r d r) \omega}{2 \pi}=\sigma \omega r d r\)\\
So, magnetic induction at the centre, due to this element : \(d B=\frac{\mu_{0}}{2} \frac{d i}{r}\) and hence, from symmetry : \(B=\int d B=\int_{0}^{R} \frac{\mu_{0} \sigma \omega r d r}{r}=\frac{\mu_{0}}{2} \sigma \omega R\)\\
(b) Magnetic moment of the element, considered,

\[
d p_{m}=(d i) \pi r^{2}=\sigma \omega d r \pi r^{2}=\sigma \pi \omega r^{3} d r
\]

Hence, the sought magnetic moment,

\[
p_{m}=\int d p_{m}=\int_{0}^{R} \sigma \pi \omega r^{3} d r=\sigma \omega \pi \frac{R^{4}}{4}
\]

3.247 As only the outer surface of the sphere is charged, consider the element as a ring, as shown in the figure.\\
The equivalent current due to the ring element,


\begin{equation*}
d i=\frac{\omega}{2 \pi}(2 \pi r \sin \theta r d \theta) \sigma \tag{1}
\end{equation*}


and magnetic induction due to this loop element at the centre of the sphere, \(O\),\\
\(d B=\frac{\mu_{0}}{4 \pi} d i \frac{2 \pi r \sin \theta r \sin \theta}{r^{3}}=\frac{\mu_{0}}{4 \pi} d i \frac{\sin ^{2} \theta}{r}\)\\[0pt]
[Using 3.219 (b) ]\\
Hence, the total magnetic induction due to the sphere at the centre, \(O\),\\
\includegraphics[max width=\textwidth, alt={}, center]{9f3c5a8e-ac5d-42f4-8eb6-ea2af623f2d8-372_448_528_211_845}

\[
\begin{gathered}
B=\int d B=\int_{0}^{\pi / 2} \frac{\mu_{0}}{4 \pi} \frac{\omega}{2 \pi} \frac{2 \pi r^{2} \sin \theta d \theta \sin ^{2} \theta \sigma}{r}[\text { using (1) }] \\
B=\int_{0}^{\pi / 2} \frac{\mu_{0} \sigma \omega r}{4 \pi} \sin ^{3} \theta d \theta=\frac{2}{3} \mu_{0} \sigma \omega r=29 \mathrm{pT}
\end{gathered}
\]

3.248 The magnetic moment must clearly be along the axis of rotation. Consider a volume element \(d V\). It contains a charge \(\frac{q}{4 \pi / 3 R^{3}} d V\). The rotation of the sphere causes this charge to revolve around the axis and constitute a current.

\[
\frac{3 q}{4 \pi R^{3}} d V \times \frac{\omega}{2 \pi}
\]

Its magnetic moment will be

\[
\frac{3 q}{4 \pi R^{3}} d V \times \frac{\omega}{2 \pi} \times \pi r^{2} \sin ^{2} \theta
\]

So the total magnetic moment is

\[
p_{m}=\int_{0}^{R} \int_{0}^{\pi} \frac{3 q}{2 R^{3}} r^{2} \sin \theta d \theta \times \frac{\omega r^{2} \sin ^{2} \theta}{2} d r=\frac{3 q}{2 R^{3}} \times \frac{\omega}{2} \times \frac{R^{5}}{5} \times \frac{4}{3}=\frac{1}{5} q R^{2} \omega
\]

The mechanical moment is

\[
M=\frac{2}{5} m R^{2} \omega, \text { So, } \quad \frac{p_{m}}{M}=\frac{q}{2 m}
\]

3.249 Because of polarization a space charge is present within the cylinder. It's density is

\[
\rho_{p}=-\operatorname{div} \vec{P}=-2 \alpha
\]

Since the cylinder as a whole is neutral a surface charge density \(\sigma_{p}\) must be present on the surface of the cylinder also. This has the magnitude (algebraically)

\[
\sigma_{p} \times 2 \pi R=2 \alpha \pi R^{2} \quad \text { or, } \quad \sigma_{p}=\alpha R
\]

When the cylinder rotates, currents are set up which give rise to magnetic fields. The contribution of \(\rho_{p}\) and \(\sigma_{p}\) can be calculated separately and then added.\\
For the surface charge the current is (for a particular element)

\[
\alpha R \times 2 \pi R d x \times \frac{\omega}{2 \pi}=\alpha R^{2} \omega d x
\]

Its contribution to the magnetic field at the centre is

\[
\frac{\mu_{0} R^{2}\left(\alpha R^{2} \omega d x\right)}{2\left(x^{2}+R^{2}\right)^{3 / 2}}
\]

and the total magnetic field is

\[
B_{s}=\int_{\infty}^{-\infty} \frac{\mu_{0} R^{2}\left(\alpha R^{2} \omega d x\right)}{2\left(x^{2}+R^{2}\right)^{3 / 2}}=\frac{\mu_{0} \alpha R^{4} \omega}{2} \int_{-\infty}^{\infty} \frac{d x}{\left(x^{2}+R^{2}\right)^{3 / 2}}=\frac{\mu_{0} \alpha R^{4} \omega}{2} \times \frac{2}{R^{2}}=\mu_{0} \alpha R^{2} \omega
\]

As for the volume charge density consider a circle of radius \(r\), radial thickness \(d r\) and length \(d x\).\\
The current is \(-2 \alpha \times 2 \pi r d r d x \times \frac{\omega}{2 \pi}=-2 \alpha r d r \omega d x\)\\
The total magnetic field due to the volume charge distribution is

\[
\begin{aligned}
B_{v} & =-\int_{0}^{R} d r \int_{-\infty}^{\infty} d x 2 \pi r \omega \frac{\mu_{0} r^{2}}{2\left(x^{2}+r^{2}\right)^{3 / 2}}=-\int_{0}^{R} \alpha \mu_{0} \omega r^{3} d r \int_{-\infty}^{\infty} d x\left(x^{2}+r^{2}\right)^{3 / 2} \\
& =-\int_{0} \alpha \mu_{0} \omega r d r \times 2=-\mu_{0} \alpha \omega R^{2} \quad \text { so, } \quad B=B_{s}+B_{v}=0
\end{aligned}
\]

3.250 Force of magnetic interaction, \(\vec{F}_{\text {mag }}=e(\vec{v} \times \vec{B})\)

Where,

\[
\vec{B}=\frac{\mu_{0}}{4 \pi} \frac{e(\vec{v} \times \vec{r})}{r^{3}}
\]

So,

\[
\vec{F}_{m a g}=\frac{\mu_{0}}{4 \pi} \frac{e^{2}}{r^{3}}[\vec{v} \times(\vec{v} \times \vec{r})]
\]

\[
=\frac{\mu_{0}}{4 \pi} \frac{e^{2}}{r^{3}}[(\vec{v} \cdot \vec{r}) \times \vec{v}-(\vec{v} \cdot \vec{v}) \times \vec{r}]=\frac{\mu_{0}}{4 \pi} \frac{e^{2}}{r^{3}}\left(-v^{2} \vec{r}\right)
\]

And

\[
\vec{F}_{e l e}=e \vec{E}=e \frac{1}{4 \pi \varepsilon_{0}} \frac{e \vec{r}}{|\vec{r}|^{3}}
\]

Hence,

\[
\frac{\left|\vec{F}_{\text {mag }}\right|}{\left|\vec{F}_{\text {electic }}\right|}=-v^{2} \mu_{0} \varepsilon_{0}=\left(\frac{v}{c}\right)^{2}=1.00 \times 10^{-6}
\]

3.251 (a) The magnetic field at \(O\) is only due to the curved path, as for the line element, \(d \vec{l} \uparrow \uparrow \overrightarrow{r .}\)

Hence, \(\quad \vec{B}=\frac{\mu_{0} i}{4 \pi R} \pi(-\vec{k})=\frac{\mu_{0} i}{4 R}(-\vec{k})\)\\
Thus \(\vec{F}_{u}=i B(-\vec{j})=\frac{\mu_{0} i^{2}}{4 R}(-\vec{j})\)\\
So, \(F_{u}=\frac{\pi_{0} i^{2}}{4 R}=0.20 \mathrm{~N} / \mathrm{m}\)\\
(b) In this part, magnetic induction \(\vec{B}\) at \(O\)\\
\includegraphics[max width=\textwidth, alt={}]{9f3c5a8e-ac5d-42f4-8eb6-ea2af623f2d8-374_455_534_192_826} will be effective only due to the two semi infinite segments of wire. Hence

\[
\begin{aligned}
\vec{B} & =2 \cdot \frac{\mu_{0} i}{4 \pi\left(\frac{l}{2}\right)} \sin \frac{\pi}{2}(-\vec{k}) \\
& =\frac{\mu_{0} i}{\pi l}(-\vec{k})
\end{aligned}
\]

Thus force per unit length,\\
\includegraphics[max width=\textwidth, alt={}, center]{9f3c5a8e-ac5d-42f4-8eb6-ea2af623f2d8-374_325_547_689_811}

\[
\vec{F}_{u}=\frac{\mu_{0} l^{2}}{\pi l}(-\vec{i})
\]

3.252 Each element of length \(d l\) experiences a force \(B I d l\). This causes a tension \(T\) in the wire. For equilibrium,

\[
T d \alpha=B I d l,
\]

where \(d \alpha\) is the angle subtended by the element at the centre.

Then,

\[
T=B I \frac{d l}{d \alpha}=B I R
\]

The wire experiences a stress

\[
\frac{B I R}{\pi d^{2} / 4}
\]

\includegraphics[max width=\textwidth, alt={}, center]{9f3c5a8e-ac5d-42f4-8eb6-ea2af623f2d8-374_366_296_1280_676}\\
\includegraphics[max width=\textwidth, alt={}, center]{9f3c5a8e-ac5d-42f4-8eb6-ea2af623f2d8-374_306_406_1340_1001}

This must equals the breaking stress \(\sigma_{m}\) for rupture. Thus,

\[
B_{\max }=\frac{\pi d^{2} \sigma_{m}}{4 I R}
\]

3.253 The Ampere forces on the sides \(O P\) and \(O^{\prime} P^{\prime}\) are directed along the same line, in opposite directions and have equal values, hence the net force as well as the net torque of these forces about the axis \(O O^{\prime}\) is zero. The Ampere-force on the segment \(P P^{\prime}\) and the corresponding moment of this force about the axis \(O O^{\prime}\) is effective and is deflecting in nature.

In equilibrium (in the dotted position) the deflecting torque must be equal to the restoring torque, developed due to the weight of the shape.\\
Let, the length of each side be \(l\) and \(\rho\) be the density of the material then,

\[
\begin{aligned}
i l B(l \cos \theta)=(S l \rho) g \frac{l}{2} \sin \theta & +(S l \rho) g \frac{l}{2} \sin \theta \\
& +(S l \rho) g l \sin \theta
\end{aligned}
\]

or, \(\quad i l^{2} B \cos \theta=2 S \rho g l^{2} \sin \theta\)\\
Hence,

\[
B=\frac{2 S \rho g}{i} \tan \theta
\]

3.254 We know that the torque acting on a magnetic dipole.

\[
\vec{N}=\overrightarrow{p_{m}} \times \vec{B}
\]

\begin{figure}[h]
\begin{center}
  \includegraphics[alt={},max width=\textwidth]{9f3c5a8e-ac5d-42f4-8eb6-ea2af623f2d8-375_624_473_87_913}
\captionsetup{labelformat=empty}
\caption{S2gg}
\end{center}
\end{figure}

But, \(\vec{p}_{m}=i S \hat{n}\), where \(\hat{n}\) is the normal on the plane of the loop and is directed in the direction of advancement of a right handed screw, if we rotate the screw in the sense of current in the loop.

On passing a current through the coil, this torque acting on the magnetic dipol, is counterbalanced by the moment of additional weight, about \(O\).

Hence, the direction of current in the loop must be in the direction, shown in the figure.

\[
\overrightarrow{p_{m}} \times \vec{B}=-\vec{l} \times \Delta m \vec{g}
\]

\includegraphics[max width=\textwidth, alt={}, center]{9f3c5a8e-ac5d-42f4-8eb6-ea2af623f2d8-375_322_607_957_776}\\
or, \(\quad N i S B=\Delta m g l\)\\
So, \(\quad B=\frac{\Delta m g l}{N i S}=0.4 \mathrm{~T}\) on putting the values.\\
3.255 (a) As is clear from the condition, Ampere's forces on the sides (2) and (4) are equal in magnitude but opposite in direction. Hence the net effective force on the frame is the resultant of the forces, experienced by the sides (1) and (3).\\
Now, the Ampere force on (1),

\[
F_{1}=\frac{\mu_{0}}{2 \pi} \frac{i i_{0}}{\left(\eta-\frac{1}{2}\right)}
\]

and that on (3),

\[
F_{3}=\frac{\mu_{0}}{2 \pi} \frac{i_{0} i}{\left(\eta+\frac{1}{2}\right)}
\]

So, the resultant force on the frame \(=F_{1}-F_{3}\), (as they are opposite in nature.)\\
\includegraphics[max width=\textwidth, alt={}, center]{9f3c5a8e-ac5d-42f4-8eb6-ea2af623f2d8-375_462_445_1581_849}

\[
=\frac{2 \mu_{0} \ddot{i}_{0}}{\pi\left(4 \eta^{2}-1\right)}=0.40 \mu \mathrm{~N} .
\]

(b) Work done in turning the frame through some angle, \(A=\int i d \Phi=i\left(\Phi_{f}-\Phi_{i}\right)\), where \(\Phi_{f}\) is the flux through the frame in final position, and \(\Phi_{i}\), that in the the initial position. Here,

\[
\left|\Phi_{f}\right|=\left|\Phi_{i}\right|=\Phi \text { and } \Phi_{i}=-\Phi_{f}
\]

so,

\[
\Delta \Phi=2 \Phi \text { and } A=i 2 \Phi
\]

Hence,

\[
A=2 i \int \vec{B} \cdot d \vec{S}
\]

\[
=2 i \int_{a\left(\eta-\frac{1}{2}\right)}^{a\left(\eta+\frac{1}{2}\right)} \frac{\mu_{0}}{2 \pi} \frac{i_{0} a}{r} d r=\frac{\mu_{0} i_{0} a}{\pi} \ln \left(\frac{2 \eta+1}{2 \eta-1}\right)
\]

3.256 There are excess surface charges on each wire (irrespective of whether the current is flowing through them or not). Hence in addition to the magnetic force \(\vec{F}_{m}\), we must take into account the electric force \(\vec{F}_{e}\). Suppose that an excess charge \(\lambda\) corresponds to a unit length of the wire, then electric force exerted per unit length of the wire by other wire can be found with the help of Gauss's theorem.


\begin{equation*}
F_{e}=\lambda E=\lambda \frac{1}{4 \pi \varepsilon_{0}} \frac{2 \lambda}{l}=\frac{2 \lambda^{2}}{4 \pi \varepsilon_{0} l}, \tag{1}
\end{equation*}


where \(l\) is the distance between the axes of the wires. The magnetic force acting per unit length of the wire can be found with the help of the theorem on circulation of vector \(\vec{B}\)


\begin{equation*}
F_{m}=\frac{\mu_{0}}{4 \pi} \frac{2 i^{2}}{l} \tag{2}
\end{equation*}


where \(i\) is the current in the wire.\\
\includegraphics[max width=\textwidth, alt={}, center]{9f3c5a8e-ac5d-42f4-8eb6-ea2af623f2d8-376_310_538_1230_815}\\
Now, from the relation, \(\lambda=C \varphi\), where \(C\) is the capacitance of the wires per unit lengths and is given in problem 3.108 and \(\varphi=i R\)


\begin{equation*}
\lambda=\frac{\pi \varepsilon_{0}}{\ln \eta} i R \quad \text { or, } \quad \frac{i}{\lambda}=\frac{\ln \eta}{\pi \varepsilon_{0} R} \tag{3}
\end{equation*}


Dividing (2) by (1) and then substuting the value of \(\frac{i}{\lambda}\) from (3), we get,

\[
\frac{F_{m}}{F_{e}}=\frac{\mu_{0}}{\varepsilon_{0}} \frac{(\ln \eta)^{2}}{\pi^{2} R^{2}}
\]

The resultant force of interaction vanishes when this ratio equals unity. This is possible when \(R=R_{0}\), where

\[
R_{0}=\sqrt{\frac{\mu_{0}}{\varepsilon_{0}}} \frac{\ln \eta}{\pi}=0.36 \mathrm{k} \Omega
\]

3.257 Use 3.225

The magnetic field due to the conductor with semicircular cross section is

\[
B=\frac{\mu_{0} I}{\pi^{2} R}
\]

Then

\[
\frac{\partial F}{\partial l}=B I=\frac{\mu_{0} I^{2}}{\pi^{2} R}
\]

3.258 We know that Ampere's force per unit length on a wire element in a magnetic field is given by.\\
\(d \vec{F}_{n}=i(\hat{n} \times \hat{B})\) where \(\hat{n}\) is the unit vector along the direction of current.

Now, let us take an element of the conductor \(i_{2}\), as shown in the figure. This wire element is in the magnetic field, produced by the current \(i_{1}\), which is directed normally into the sheet of the paper and its magnitude is given by,


\begin{equation*}
|\vec{B}|=\frac{\mu_{0} I_{1}}{2 \pi r} \tag{2}
\end{equation*}


From Eqs. (1) and (2)\\
\includegraphics[max width=\textwidth, alt={}, center]{9f3c5a8e-ac5d-42f4-8eb6-ea2af623f2d8-377_435_407_643_965}\\
\(d \vec{F}_{n}=\frac{I_{2}}{b} d r(\hat{n} \times \vec{B})\), (because the current through the element equals \(\frac{I_{2}}{b} d r\) )\\
So,

\[
d \vec{F}_{n}=\frac{\mu_{0}}{2 \pi} \frac{I_{1} I_{2}}{b} \frac{d r}{r}, \text { towards left }\left(\text { as } \hat{n} \perp^{r} \vec{B}\right)
\]

Hence the magnetic force on the conductor :\\
\(\vec{F}_{n}=\frac{\mu_{0}}{2 \pi} \frac{I_{1} I_{2}}{b} \int_{a}^{a+b} \frac{d r}{r}\) (towards left) \(=\frac{\mu_{0}}{2 \pi} \frac{I_{1} I_{2}}{b} \ln \frac{a+b}{a}\) (towards left).\\
Then according to the Newton's third law the magnitude of sought magnetic interaction force

\[
=\frac{\mu_{0} I_{1} I_{2}}{2 \pi b} \ln \frac{a+b}{a}
\]

3.259 By the circulation theorem \(B=\mu_{0} i\),\\
where \(i=\) current per unit length flowing along the plane perpendicular to the paper. Currents flow in the opposite sense in the two planes and produce the given field B by superposition.\\
The field due to one of the plates is just \(\frac{1}{2} B\). The force on the plate is,\\
\(\frac{1}{2} B \times i \times\) Length × Breadth \(=\frac{B^{2}}{2 \mu_{0}}\) per unit area.\\
(Recall the formula \(F=B I l\) on a straight wire)\\
3.260 (a) The external field must be \(\frac{B_{1}+B_{2}}{2}\), which when superposed with the internal field \(\frac{B_{1}-B_{2}}{2}\) (of opposite sign on the two sides of the plate) must give actual field. Now

\[
\begin{array}{ll} 
& \frac{B_{1}-B_{2}}{2}=\frac{1}{2} \mu_{0} i \\
\text { or, } & i=\frac{B_{1}-B_{2}}{\mu_{0}} \\
\text { Thus, } & F=\frac{B_{1}^{2}-B_{2}^{2}}{2 \mu_{0}}
\end{array}
\]

\begin{figure}[h]
\begin{center}
  \includegraphics[alt={},max width=\textwidth]{9f3c5a8e-ac5d-42f4-8eb6-ea2af623f2d8-378_224_124_302_1044}
\captionsetup{labelformat=empty}
\caption{\(B_{1}\)}
\end{center}
\end{figure}

\includegraphics[max width=\textwidth, alt={}, center]{9f3c5a8e-ac5d-42f4-8eb6-ea2af623f2d8-378_283_143_301_1188}\\
(b) Here, the external field must be \(\frac{B_{1}-B_{2}}{2}\) upward with an internal field, \(\frac{B_{1}+B_{2}}{2}\), upward on the left and dewnward on the right. Thus,

\[
i=\frac{B_{1}+B_{2}}{\mu_{0}} \text { and } F=\frac{B_{1}^{2}-B_{2}^{2}}{2 \mu_{0}} .
\]

(c) Our boundary condition following from

Gauss' law is, \(B_{1} \cos \theta_{1}=B_{2} \cos \theta_{2}\).\\
Also, \(\left(B_{1} \sin \theta_{1}+B_{2} \sin \theta_{2}\right)=\mu_{0} i \quad\) where \(i=\) current per unit length.

\begin{figure}[h]
\begin{center}
  \includegraphics[alt={},max width=\textwidth]{9f3c5a8e-ac5d-42f4-8eb6-ea2af623f2d8-378_299_162_748_1086}
\captionsetup{labelformat=empty}
\caption{\(B_{1}\)}
\end{center}
\end{figure}

\begin{center}
\includegraphics[max width=\textwidth, alt={}]{9f3c5a8e-ac5d-42f4-8eb6-ea2af623f2d8-378_359_121_745_1277}
\end{center}

The external field parallel to the plate must be \(\frac{B_{1} \sin \theta_{1}-B_{2} \sin \theta_{2}}{2}\)\\
(The perpendicular component \(B_{1} \cos \theta_{1}\), does not matter since the corresponding force is tangential)\\
Thus, \(F=\frac{B_{1}{ }^{2} \sin ^{2} \theta_{1}-B_{2}{ }^{2} \sin ^{2} \theta_{2}}{, 2 \mu_{0}}\) per unit area

\[
=\frac{B_{1}^{2}-B_{2}^{2}}{2 \mu_{0}} \text { per unit area. }
\]

The direction of the current in the plane conductor is perpendicular to the paper and\\
\includegraphics[max width=\textwidth, alt={}, center]{9f3c5a8e-ac5d-42f4-8eb6-ea2af623f2d8-378_308_447_1341_928}\\
beyond the drawing.\\
3.261 The Current density is \(\frac{I}{a L}\), where \(L\) is the length of the section. The difference in pressure produced must be,

\[
\Delta p=\frac{1}{a L} \times B \times(a b L) / a b=\frac{I B}{a}
\]

\includegraphics[max width=\textwidth, alt={}, center]{9f3c5a8e-ac5d-42f4-8eb6-ea2af623f2d8-378_347_368_1675_994}\\
3.262 Let \(t=\) thickness of the wall of the cylinder. Then, \(J=I / 2 \pi R t\) along \(z\) axis. The magnetic field due to this at a distance \(r \left(R-\frac{t}{2}<r<R+\frac{t}{2}\right)\), is given by,\\
\(B_{\varphi}(2 \pi r)=\mu_{0} \frac{I}{2 \pi R t} \pi\left\{r^{2}-\left(R-\frac{t}{2}\right)^{2}\right\}\) or, \(\quad B_{\varphi}=\frac{\mu_{0} I}{4 \pi R r t}\left\{r^{2}-(R-t / 2)^{2}\right\}\)

Now, \(\vec{F}=\int \vec{J} \times \vec{B} d V\)\\
\includegraphics[max width=\textwidth, alt={}, center]{9f3c5a8e-ac5d-42f4-8eb6-ea2af623f2d8-379_381_490_224_728}

\[
\begin{aligned}
& \text { and } p=\frac{F_{r}}{2 \pi R L}=\frac{1}{2 \pi R L} \int_{R-\frac{t}{2}}^{R+\frac{t}{2}} \frac{\mu_{0} I^{2}}{8 \pi^{2} R^{2} t^{2} r}\left\{r^{2}-\left(R-\frac{t}{2}\right)^{2}\right\} \times 2 \pi r L d r \\
& =\frac{\mu_{0} I^{2}}{8 \pi^{2} R^{3} t^{2}} \int_{R-\frac{t}{2}}^{R+\frac{t}{2}}\left\{r^{2}-\left(R-\frac{t}{2}\right)^{2}\right\} d r=\frac{\mu_{0} I^{2}}{8 \pi^{2} R^{3} t^{2}}\left[\frac{\left(R+\frac{t}{2}\right)^{3}-\left(R-\frac{t}{2}\right)^{3}}{3}-\left(R-\frac{t}{2}\right)^{2} t\right] \\
& =\frac{\mu_{0} I^{2}}{8 \pi^{2} R^{3} t}\left[R t+0\left(t^{2}\right)\right] \approx \frac{\mu_{0} I^{2}}{8 \pi^{2} R^{2}}
\end{aligned}
\]

3.263 When self-forces are involved, a typical factor of \(\frac{1}{2}\) comes into play. For example, the force on a current carrying straight wire in a magnetic induction \(B\) is \(B I l\). If the magnetic induction \(B\) is due to the current itself then the force can be written as,

\[
F=\int_{0}^{I} B\left(I^{\prime}\right) d I^{\prime} l
\]

If \(B\left(I^{\prime}\right) \alpha I^{\prime}\), then this becomes, \(F=\frac{1}{2} B(I) I l\).\\
In the present case, \(B(I)=\mu_{0} n I\) and this acts on \(n I\) ampere turns per unit length, so, pressure \(p=\frac{F}{\text { Area }}=\frac{1}{2} \mu_{0} n \frac{I \times n I \times 1 \times l}{1 \times l}=\frac{1}{2} \mu_{0} n^{2} I^{2}\)\\
3.264 The magnetic induction \(B\) in the solenoid is given by \(B=\mu_{0} n I\). The force on an element \(d l\) of the current carrying conductor is,

\[
d F=\frac{1}{2} \mu_{0} n I I d l=\frac{1}{2} \mu_{0} n I^{2} d l
\]

This is radially outwards. The factor \(\frac{1}{2}\) is explained above.

To relate \(d F\) to the tensile strength \(F_{\text {lim }}\) we proceed as follows. Consider the equilibrium of the element dl. The longitudinal forces \(F\) have a radial component equal to,

\[
d F=2 F \sin \frac{d \theta}{2}=F d \theta
\]

Thus using \(d l=R d \theta, F=\frac{1}{2} \mu_{0} n I^{2} R\)\\
This equals \(F_{\text {lim }}\) when, \(I=I_{\text {lim }}=\sqrt{\frac{2 F_{\text {lim }}}{\mu_{0} n R}}\)\\
\includegraphics[max width=\textwidth, alt={}, center]{9f3c5a8e-ac5d-42f4-8eb6-ea2af623f2d8-380_358_536_237_790}

Note that \(F_{\text {lim }}\), here, is actually a force and not a stress.\\
3.265 Resistance of the liquid between the plates \(=\frac{\rho d}{S}\)

Voltage between the plates \(=E d=v B d\),\\
Current through the plates \(=\frac{v B d}{R+\frac{\rho d}{S}}\)\\
Power, generated, in the external resistance \(R\),\\
\(P=\frac{v^{2} B^{2} d^{2} R}{\left(R+\frac{\rho d}{S}\right)^{2}}=\frac{v^{2} B^{2} d^{2}}{\left(\sqrt{R}+\frac{\rho d}{S \sqrt{R}}\right)^{2}}=\frac{v^{2} B^{2} d^{2}}{\left[\left\{R^{1 / 4}-\left(\frac{\rho d}{S \sqrt{R}}\right)^{1 / 2}\right\}+2 \sqrt{\frac{\rho d}{S}}\right]^{2}}\)\\
This is maximum when \(R=\frac{\rho D}{S}\) and \(P_{\text {max }}=\frac{v^{2} B^{2} S d}{4 \rho}\)\\
3.266 The electrons in the conductor are drifting with a speed of,

\[
v_{d}=\frac{J}{n e}=\frac{I}{\pi R^{2} n e},
\]

where \(e=\) magnitude of the charge on the electron, \(n=\) concentration of the conduction electorns.\\
The magnetic field inside the conductor due to this current is given by,

\[
B_{\varphi}(2 \pi r)=\pi r^{2} \frac{I}{\pi R^{2}} \mu_{0} \quad \text { or, } \quad B_{\varphi}=\frac{\mu_{0}}{2 \pi} \frac{I r}{R^{2}}
\]

A radial electric field \(v B_{\varphi}\) must come into being in equilibrium. Its P.D. is,

\[
\Delta \varphi=\int_{0}^{R} \frac{I}{\pi R^{2} n e} \frac{\mu_{0}}{2 \pi} \frac{I r}{R^{2}} d r=\frac{I}{\pi R^{2} n e}\left(\frac{\mu_{0}}{4 \pi} I\right)=\frac{\mu_{0} I^{2}}{4 \pi R^{2} n e}
\]

3.267 Here, \(v_{d}=\frac{E}{B}\) and \(j=n e v_{d}\)

\[
\text { so, } \begin{aligned}
n & =\frac{j B}{e E}=\frac{200 \times 10^{4} \frac{\mathrm{~A}}{\mathrm{~m}^{2}} \times 1 \mathrm{~T}}{1.6 \times 10^{-19} \mathrm{C} \times 5 \times 10^{-4}} \mathrm{~V} / \mathrm{m} \\
& =2.5 \times 10^{28} \text { per m }^{3}=2.5 \times 10^{22} \text { per c.c. }
\end{aligned}
\]

Atomic weight of Na being 23 and its density \(\approx 1\), molar volume is 23 c.c. Thus number of atoms per unit volume is \(\frac{6 \times 10^{23}}{23}=2.6 \times 10^{22}\) per c.c.\\
Thus there is almost one conduction electron per atom.\\
3.268 By definition, mobility \(=\frac{\text { dirft velocity }}{\text { Electric field component causing this drift }}\) or \(\mu=\frac{v}{E_{L}}\)

On other hand,

\[
E_{T}=\nu B=\frac{E_{L}}{\eta}, \text { as given } \text { so, } \mu=\frac{1}{\eta B}=3.2 \times 10^{-3} \mathrm{~m}^{2} /(\mathrm{V} \cdot \mathrm{~s})
\]

3.269 Due to the straight conductor, \(B_{\varphi}=\frac{\mu_{0} I}{2 \pi r}\)

We use the formula, \(\quad \vec{F}=\left(\overrightarrow{p_{m}} \cdot \vec{\nabla}\right) \vec{B}\)\\
(a) The vector \(\vec{p}_{m}\) is parallel to the straight conductor.

\[
\vec{F}=p_{m} \frac{\partial}{\partial Z} \vec{B}=0
\]

because neither the direction nor the magnitude of \(\vec{B}\) depends on \(z\)\\
(b) The vector \(\vec{p}_{m}\) is oriented along the radius vector \(\vec{r}\)

\[
\vec{F}=p_{m} \frac{\partial}{\partial r} \vec{B}
\]

The direction of \(\vec{B}\) at \(r+d r\) is parallel to the direction at \(r\). Thus only the \(\varphi\) component of \(\vec{F}\) will survive.

\[
F_{\varphi}=p_{m} \frac{\partial}{\partial r} \frac{\mu_{0} I}{2 \pi r}=-\frac{\mu_{0} I p_{m}}{2 \pi r^{2}}
\]

(c) The vector \(\vec{p}_{m}\) coincides in direction with the magnetic field, produced by the conductor carrying current \(I\)

\[
\begin{aligned}
& \vec{F}=p_{m} \frac{\partial}{r \partial \varphi} \frac{\mu_{0} I}{2 \pi} \vec{e}_{\varphi}=\frac{\mu_{0} I p_{m}}{2 \pi r^{2}} \frac{\partial \vec{e}_{\varphi}}{\partial \varphi} \\
& \vec{F}=-\frac{\mu_{0} I p_{m}}{2 \pi r^{2}} \overrightarrow{e_{r}} \text { As, } \frac{\partial \vec{e}_{\varphi}}{\partial_{\varphi}}=-\overrightarrow{e_{r}}
\end{aligned}
\]

So,\\
3.270\\
\(F_{x}=p_{m} \frac{\partial}{\partial x} B_{x}\)\\
But, \(B_{x}=\frac{\mu_{0} I}{4 \pi} \int \frac{R d l}{\left(x^{2}+R^{2}\right)^{3 / 2}}=\frac{\mu_{0} I R^{2}}{2\left(x^{2}+R^{2}\right)^{3 / 2}}\)\\
So, \(\quad F=\frac{\mu_{0}}{4 \pi} \frac{I \cdot 2 \pi R^{2}}{\left(x^{2}+R^{2}\right)^{5 / 2}} \frac{3}{2} \cdot 2 x \cdot p_{m}\)

\[
=\frac{\mu_{0}}{4 \pi} \frac{6 \pi R^{2} I p_{m} x}{\left(x^{2}+R^{2}\right)^{5 / 2}}
\]

\includegraphics[max width=\textwidth, alt={}, center]{9f3c5a8e-ac5d-42f4-8eb6-ea2af623f2d8-382_356_345_196_970}\\
3.271

\[
\begin{aligned}
F & =P_{2 m} \frac{\partial}{\partial l}\left[\frac{\mu_{0}}{4 \pi} \frac{3 \vec{p}_{1 m} \cdot \vec{r} \vec{r}-\vec{p}_{1 m} r^{2}}{r^{5}}\right] \\
& =P_{2 m} \frac{\partial}{\partial l}\left[\frac{\mu_{0}}{2 \pi} \frac{p_{1 m}}{l^{3}}\right]=\frac{-3}{2} \frac{\mu_{0} p_{1 m} p_{2 m}}{\pi l^{4}}=9 \mathrm{nN}
\end{aligned}
\]

3.272 From 3.270, for \(x \gg R\),

\[
B_{x} \approx \frac{\mu_{0} I^{\prime} R^{2}}{2 x^{3}}
\]

or, \(I^{\prime} \approx \frac{2 B_{x} x^{3}}{\mu_{0} R^{2}}=\frac{2 \times 3 \times 10^{-5} T \times\left(10^{-1} \mathrm{~m}\right)^{3}}{1.26 \times 10^{-6} \times\left(10^{-2} \mathrm{~m}\right)^{4}} \approx 0.5 \mathrm{kA}\)\\
3.273\\
\includegraphics[max width=\textwidth, alt={}, center]{9f3c5a8e-ac5d-42f4-8eb6-ea2af623f2d8-382_388_480_1090_833}\\
so,

\[
\begin{aligned}
B_{n}^{\prime} & =B \cos \alpha, \\
H_{t}^{\prime} & =\frac{1}{\mu_{0}} B \sin \alpha, \\
B_{t}^{\prime} & =\mu B \sin \alpha
\end{aligned}
\]

\[
B^{\prime}=B \sqrt{\mu^{2} \sin ^{2} \alpha+\cos ^{2} \alpha}
\]

3.274 (a) \(\oint \vec{H} \cdot d \vec{S}=\oint\left(\frac{\vec{B}}{\mu_{0}}-\vec{J}\right) \cdot d \vec{S}=-\oint \vec{J} \cdot d \vec{S}\), since \(\oint \vec{B} \cdot d \vec{S}=0\)

Now \(\vec{J}\) is nonvanishing only in the bottom half of the sphere.\\
\includegraphics[max width=\textwidth, alt={}, center]{9f3c5a8e-ac5d-42f4-8eb6-ea2af623f2d8-382_383_854_1663_243}

Here, \(B_{n}^{\prime}=B \cos \theta, H_{t}^{\prime}=\frac{1}{\mu_{0}} B \sin \theta, B_{t}^{\prime}=\mu B \sin \theta, H_{n}^{\prime}=\frac{B}{\mu \mu_{0}} \cos \theta\)

\[
J_{n}=\frac{B \cos \theta}{\mu_{0}}\left(1-\frac{1}{\mu}\right) \text { and } J_{t}=\frac{\mu-1}{\mu_{0}} B \sin \theta
\]

Only \(J_{n}\) contributes the surface integral and

\[
-\oint \vec{J} \cdot d \vec{S}=-\oint_{\text {lower }} \vec{J} \cdot d \vec{S}=\oint_{\text {lower }} J_{n} d S=\frac{\pi R^{2} B \cos \theta}{\mu_{0}}\left(1-\frac{1}{\mu}\right)
\]

(b) \(\oint_{\Gamma} \vec{B} \cdot d \vec{r}=\left(B_{t}-B_{t}^{\prime}\right) l=(1-\mu) B l \sin \theta\)\\
3.275 Inside the cylindrical wire there is an external current of density \(\frac{I}{\pi R^{2}}\). This gives a magnetic field \(H_{\varphi}\) with

\[
H_{\varphi} 2 \pi r=I \frac{r^{2}}{R^{2}} \quad \text { or, } \quad H_{\varphi}=\frac{I r}{2 \pi R^{2}}
\]

From this \(B_{\varphi}=\frac{\mu \mu_{0} I r}{2 \pi R^{2}}\) and \(J_{\varphi}=\frac{\mu-1}{2 \pi} \frac{I r}{R^{2}}=\frac{\chi I r}{2 \pi R^{2}}=\) Magnetization.\\
Hence total volume molecular current is,

\[
\oint_{r=R} \vec{J}_{\varphi} \cdot d \vec{r}=\int \frac{\chi I}{2 \pi R} d l=\chi I
\]

The surface current is obtained by using the equivalence of the surface current density to \(\vec{J} \times \vec{n}\), this gives rise to a surface current density in the \(z\)-direction of \(-\frac{\chi I}{2 \pi R}\)\\
The total molecular surface current is,

\[
I_{s}^{\prime}=-\frac{\chi I}{2 \pi R}(2 \pi R)=-\chi I .
\]

The two currents have opposite signs.\\
3.276 We can obtain the form of the curves, required here, by qualitative arguments.

From

\[
\oint \vec{H} \cdot d \vec{l}=I
\]

we get

\[
H(x \gg 0)=H(x \ll 0)=n I
\]

Then

\[
\begin{aligned}
& B(x \gg 0)=\mu \mu_{0} n I \\
& B(x<0)=\mu_{0} n I
\end{aligned}
\]

Also,

\[
\begin{aligned}
B(x<0) & =\mu_{0} H(x<0) \\
J(x<0) & =0
\end{aligned}
\]

\(B\) is continuous at \(x=0, H\) is not. These give the required curves as shown in the answersheet.\\
3.277 The lines of the \(B\) as well as \(H\) field are circles around the wire. Thus

\[
H_{1} \pi r+H_{2} \pi r=I \quad \text { or, } \quad H_{1}+H_{2}=\frac{I}{\pi r}
\]

Also \(\mu_{0} \mu_{1} H_{1}=\mu_{2} H_{2} \mu_{0}=B_{1}=B_{2}=B\)\\
Thus \(\quad H_{1}=\frac{\mu_{2}}{\mu_{1}+\mu_{2}} \frac{I}{\pi r}\),

\[
H_{2}=\frac{\mu_{1}}{\mu_{1}+\mu_{2}} \frac{I}{\pi r}
\]

and \(\quad B=\mu_{0} \frac{\mu_{1} \mu_{2}}{\mu_{1}+\mu_{2}} \frac{I}{\pi r}\).\\
\includegraphics[max width=\textwidth, alt={}, center]{9f3c5a8e-ac5d-42f4-8eb6-ea2af623f2d8-384_467_561_126_784}\\
3.278 The medium \(I\) is vacuum and contains a circular current carrying coil with current \(I\). The medium II is a magnetic with permeability \(\mu\). The boundary is the plane \(z=0\) and the coil is in the plane \(z=l\). To find the magnetic induction, we note that the effect of the magnetic medium can be written as due to an image coil in II as far as the medium I is concerned. On the other hand, the induction in II can be written as due to the coil in I, carrying a different current. It is sufficient to consider the far away fields and ensure that the boundary conditions are satisfied there. Now for actual coil in medium \(I\),\\
\includegraphics[max width=\textwidth, alt={}, center]{9f3c5a8e-ac5d-42f4-8eb6-ea2af623f2d8-384_537_356_911_240}\\
\includegraphics[max width=\textwidth, alt={}, center]{9f3c5a8e-ac5d-42f4-8eb6-ea2af623f2d8-384_554_459_879_725}

\[
B_{r}=-\frac{2 p_{m} \cos \theta^{\prime}}{r^{3}} \cdot\left(\frac{\mu_{0}}{4 \pi}\right), \quad B_{\theta}=\frac{p_{m} \sin \theta^{\prime}}{r^{3}}\left(\frac{\mu_{0}}{4 \pi}\right)
\]

so, \(\quad B_{2}=\frac{\mu_{0} p_{m}}{4 \pi}\left(2 \cos ^{2} \theta^{\prime}-\sin ^{2} \theta^{\prime}\right)\) and \(B_{x}=\frac{\mu_{0} p_{m}}{4 \pi}\left(-3 \sin \theta^{\prime} \cos \theta^{\prime}\right)\)\\
where

\[
p_{m}=I\left(\pi a^{2}\right), a=\text { radius of the coil. }
\]

Similarly due to the image coil,\\
\(B_{z}=\frac{\mu_{0} p^{\prime}{ }_{m}}{4 \pi}\left(2 \cos ^{2} \theta^{\prime}-\sin ^{2} \theta^{\prime}\right), B_{x}=\frac{\mu_{0} p^{\prime}{ }_{m}}{4 \pi}\left(3 \sin \theta^{\prime} \cos \theta^{\prime}\right), p^{\prime}{ }_{m}=I^{\prime}\left(\pi a^{2}\right)\)\\
As far as the medium II is concerned, we write similarly\\
\(B_{z}=\frac{\mu_{0} p^{\prime \prime}{ }_{m}}{4 \pi}\left(2 \cos ^{2} \dot{\theta}^{\prime}-\sin ^{2} \theta^{\prime}\right), B_{x}=\frac{\mu_{0} p^{\prime \prime}}{4 \pi}\left(-3 \sin \theta^{\prime} \cos \theta^{\prime}\right), p^{\prime \prime}{ }_{m}=I^{\prime \prime}\left(\pi a^{2}\right)\)

The boundary conditions are, \(p_{m}+p_{m}^{\prime}=p_{m}^{\prime \prime}\left(\right.\) from \(\left.B_{1 n}=B_{2 n}\right)\)\\
\(-p_{m}+p_{m}^{\prime}=-\frac{1}{\mu} p_{m}^{\prime \prime}\left(\right.\) from \(\left.H_{1 t}=H_{2 t}\right)\)\\
Thus,

\[
I^{\prime \prime}=\frac{2 \mu}{\mu+1} I, I^{\prime}=\frac{\mu-1}{\mu+1} I
\]

In the limit, when the coil is on the boundary, the magnetic field enverywhere can be obtained by taking the current to be \(\frac{2 \mu}{\mu+1} I\). Thus, \(\vec{B}=\frac{2 \mu}{\mu+1} \vec{B}_{0}\)\\
3.279 We use the fact that within an isolated uniformly magnetized ball,\\
\(\overrightarrow{H^{\prime}}=-\vec{J} / 3, \overrightarrow{B^{\prime}}=\frac{2 \mu_{0} \vec{J}}{3}\), where \(\vec{J}\) is the magnetization vector. Then in a uniform magnetic field with induction \(\overrightarrow{B_{0}}\), we have by superposition,

\[
\vec{B}_{i n}=\vec{B}_{0}+\frac{2 \mu_{0} \vec{J}}{3}, \vec{H}_{i n}=\frac{\vec{B}_{0}}{\mu_{0}}-\overrightarrow{J / 3}
\]

or,

\[
\begin{aligned}
\vec{B}_{i n}+2 \mu_{0} \vec{H}_{i n} & =3 \overrightarrow{B_{0}} \\
\overrightarrow{B_{i n}} & =\mu \mu_{0} \xrightarrow{H_{i n}}
\end{aligned}
\]

also,\\
Thus,

\[
\vec{H}_{i n}=\frac{3 \vec{B}_{0}}{\mu_{0}(\mu+2)} \text { and } \vec{B}_{i n}=\frac{3 \mu \vec{B}_{0}}{\mu+2}
\]

3.280 The coercive force \(H_{c}\) is just the magnetic field within the cylinder. This is by circulation theorem, \(H_{c}=\frac{N I}{l}=6 \mathrm{kA} / \mathrm{m}\)\\
(from \(\oint \vec{H} \cdot d \vec{r}=I\), total current, considering a rectangular contour.)\\
3.281 We use, \(\oint \vec{H} \cdot d \vec{l}=0\)

Neglecting the fringing of the lines of force, we write this as

\[
H(\pi d-b)+\frac{B}{\mu_{0}} b=0
\]

or, \(\quad H \approx \frac{-B b}{\mu_{0} \pi d}=101 \mathrm{~A} / \mathrm{m}\)\\
The sense of \(H\) is opposite to \(B\)\\
\includegraphics[max width=\textwidth, alt={}, center]{9f3c5a8e-ac5d-42f4-8eb6-ea2af623f2d8-385_437_462_1136_811}\\
3.282 Here, \(\oint \vec{H} \cdot d \vec{l}=N I\) or, \(H(2 \pi R)+\frac{B b}{\mu_{0}}=N I\), so, \(H=\frac{N I \mu_{0}-B b}{2 \pi R \mu_{0}}\)

Hence,

\[
\mu=\frac{B}{\mu_{0} H}=\frac{2 \pi R B}{\mu_{0} N I-B b}=3700
\]

3.283 One has to draw the graph of \(\mu=\frac{B}{\mu_{0} H}\) versus \(H\) from the given graph. The \(\mu-H\) graph starts out horizontally, and then rises steeply at about \(H=0.04 \mathrm{k} \mathrm{A} / \mathrm{m}\) before falling agian. It is easy to check that \(\mu_{\max } \approx 10,000\).\\
3.284 From the theorem on circulation of vector \(\vec{H}\).

\[
H \pi d+\frac{B b}{\mu_{0}}=N I \quad \text { or, } B=\frac{\mu_{0} N I}{b}-\frac{\mu_{0} \pi d}{b} H=(1.51-0.987) H,
\]

where B is in Tesla and H in \(\mathrm{kA} / \mathrm{m}\). Besides, \(B\) and \(H\) are interrelated as in the Fig. 3.76 of the text. Thus we have to solve for \(\mathrm{B}, \mathrm{H}\) graphically by simultaneously drawing the two curves (the hysterisis curve and the straight line, given above) and find the point of intersection. It is at

Then,

\[
\begin{gathered}
H \approx 0.26 \mathrm{kA} / \mathrm{m}, B=1.25 \mathrm{~T} \\
\mu=\frac{B}{\mu_{0} H} \approx 4000
\end{gathered}
\]

3.285 From the formula,

Thus

\[
\begin{gathered}
\vec{F}=\left(\overrightarrow{p_{m}} \cdot \vec{\nabla}\right) \vec{B} \rightarrow \vec{F}-P \int(\vec{J} \cdot \vec{\nabla}) \vec{B} d V \\
\vec{F}=\frac{\chi}{\mu \mu_{0}} \int(\vec{B} \cdot \vec{\nabla}) \vec{B} d V
\end{gathered}
\]

or since \(\vec{B}\) is predominantly along the \(x\)-axis,

\[
F_{x}=\frac{\chi}{\mu \mu_{0}} \int B_{x} \frac{\partial B_{x}}{\partial x} S d x=\frac{\chi S}{2 \mu \mu_{0}} \int_{x=0}^{x=L} d B_{x}^{2}=-\frac{\chi S B^{2}}{2 \mu \mu_{0}} \approx \frac{\chi S B^{2}}{2 \mu_{0}}
\]

3.286 The force in question is,

\[
\vec{F}=\left(\overrightarrow{p_{m}} \cdot \vec{\nabla}\right) \vec{B}=\frac{\chi B V}{\mu \mu_{0}} \frac{d B}{d x}
\]

since \(B\) is essentiatly in the x-direction.\\
So, \(\quad F_{x} \approx \frac{\chi V}{2 \mu_{0}} \frac{d B^{2}}{d x}=\frac{\chi B_{0}^{2} V}{2 \mu_{0}} \frac{d}{d x}\left(e^{-2 a x^{2}}\right)=-4 a x e^{-2 a x^{2}} \frac{\chi B_{0}^{2}}{2 \mu_{0}} V\)\\
This is maximum when its derivative vanishes i.e. \(\quad 16 a^{2} x^{2}-4 a=0\), or, \(x_{m}=\frac{1}{\sqrt{4 a}}\)\\
The maximum force is,

\[
F_{\max }=4 a \frac{1}{\sqrt{4 a}} e^{-1 / 2} \frac{\chi B_{0}^{2} V}{2 \mu_{0}}=\frac{\chi B_{0}^{2} V}{\mu_{0}} \sqrt{\frac{a}{e}}
\]

So,

\[
\chi \approx\left(\mu_{0} F_{\max } \sqrt{\frac{e}{a}}\right) / V B_{0}^{2}=3.6 \times 10^{-4}
\]

\(3.287 F_{x}=\left(\overrightarrow{p_{m}} \cdot \vec{\nabla}\right) B_{x}=\frac{\chi B V}{\mu \mu_{0}} \frac{d B}{d x} \approx \frac{\chi V}{2 \mu_{0}} \frac{d B^{2}}{d x}\)\\
This force is attractive and an equal force must be applied for balance. The work done by applied forces is,

\[
A=\int_{x=0}^{x=L}-F_{x} d x=\frac{\chi V}{2 \mu_{0}}\left(-B^{2}\right)_{x=0}^{x=L} \approx \frac{\chi V B^{2}}{2 \mu_{0}}
\]

\subsection*{3.6 ELECTROMAGNETIC INDUCTION. MAXWELL'S EQUATIONS}
3.288 Obviously, from Lenz's law, the induced current and hence the induced e.m.f. in the loop is anticlockwise.\\
From Faraday's law of electromagnetic indcution,

Here, \(\quad d \Phi=\vec{B} \cdot d \vec{S}=-2 B x d y\),\\
and from \(y=a x^{2}, x=\sqrt{\frac{y}{a}}\)\\
Hence,\(\quad \xi_{\text {in }}=2 B \sqrt{\frac{y}{a}} \frac{d y}{d t}\)

\[
=B y \sqrt{\frac{8 w}{a}}, \text { using } \frac{d y}{d t}=\sqrt{2 w y}
\]

\includegraphics[max width=\textwidth, alt={}, center]{9f3c5a8e-ac5d-42f4-8eb6-ea2af623f2d8-387_551_520_260_854}\\
3.289 Let us assume, \(\vec{B}\) is directed into the plane of the loop. Then the motional e.m.f.

\[
\xi_{i n}=\left|\int-(\vec{v} \times \vec{B}) \cdot d \vec{l}\right|=v B l
\]

and directed in the same of \((\vec{v} \times \vec{B})\) (Fig.)\\
So, \(\quad i=\frac{\xi_{i n}}{R+\frac{R_{1} R_{2}}{R_{1}+R_{2}}}=\frac{B v l}{R+R_{u}}\)\\
As \(R_{1}\) and \(R_{2}\) are in parallel connections.\\
\includegraphics[max width=\textwidth, alt={}, center]{9f3c5a8e-ac5d-42f4-8eb6-ea2af623f2d8-387_363_433_901_899}\\
3.290 (a) As the metal disc rotates, any free electron also rotates with it with same angular velocity \(\omega\), and that's why an electron must have an acceleration \(\omega^{2} r\) directed towards the disc's centre, where \(r\) is separation of the electron from the centre of the disc. We know from Newton's second law that if a particle has some acceleration then there must be a net effecetive force on it in the direction of acceleration. We also know that a charged particle can be influenced by two fields electric and magnetic. In our problem magnetic field is absent hence we reach at the conslusion that there is an electric field near any electron and is directed opposite to the acceleration of the electron.\\
If \(E\) be the electric field strength at a distance \(r\) from the centre of the disc, we have from Newton's second law.

\[
\begin{gathered}
F_{n}=m w_{n} \\
e E=m r \omega^{2}, \text { or, } E=\frac{m \omega^{2} r}{e},
\end{gathered}
\]

and the potential difference,

\[
\varphi_{c e n}-\varphi_{r i m}=\int_{0}^{a} \vec{E} \cdot d \vec{r}=\int_{0}^{a} \frac{m \omega^{2} r}{e} d r, \text { as } \vec{E} \uparrow \downarrow d \vec{r}
\]

Thus

\[
\varphi_{c e n}-\varphi_{r i m}=\Delta \varphi=\frac{m \omega^{2}}{e} \frac{a^{2}}{2}=3.0 \mathrm{nV}
\]

(b) When field \(\vec{B}\) is present, by definition, of motional e.m.f. :

\[
\varphi_{1}-\varphi_{2}=\int_{1}^{2}-(\vec{v} \times \vec{B}) \cdot d \vec{r}
\]

Hence the sought potential difference,

\[
\varphi_{c e n}-\varphi_{r i m}=\int_{0}^{a}-v B d r=\int_{0}^{a}-\omega r B d r,(\text { as } v=\omega r)
\]

Thus

\[
\varphi_{\text {rim }}-\varphi_{\text {cen }}=\varphi=\frac{1}{2} \omega B a^{2}=20 \mathrm{mV}
\]

(In general \(\omega<\frac{e B}{m}\) so we can neglect the effect discussed in (1) here).\\
3.291 By definition,

\[
\vec{E}=-(\vec{v} \times \vec{B})
\]

So,

\[
\int_{A}^{C} \vec{E} \cdot d \vec{r}=\int_{A}^{C}-(\vec{v} \times \vec{B}) \cdot d \vec{r}=\int_{0}^{d}-v B d r
\]

But, \(v=\omega r\), where \(r\) is the perpendicular distance of the point from \(A\).\\
Hence, \(\quad \int_{A}^{c} \vec{E} \cdot d \vec{r}=\int_{0}^{d}-\omega B r d r=-\frac{1}{2} \omega B d^{2}=-10 \mathrm{mV}\)\\
This result can be generalized to a wire \(A C\) of arbitary planar shape. We have

\[
\begin{aligned}
\int_{A}^{C} \vec{E} \cdot d \vec{r} & =-\int_{A}^{C}(\vec{v} \times \vec{B}) \cdot d r=-\int_{A}^{C}((\omega \times \vec{r}) \times \vec{B}) \cdot d \vec{r} \\
& =-\int_{A}^{C}(\vec{B} \cdot \overrightarrow{r \omega}-\vec{B} \cdot \vec{\omega} \vec{r}) \cdot d \vec{r} \\
& =-\frac{1}{2} B \omega d^{2},
\end{aligned}
\]

\(d\) being \(A C\) and \(\vec{r}\) being measured from \(A\).\\
3.292 Flux at any moment of time,

\[
\left|\Phi_{t}\right|=|\vec{B} \cdot d \vec{S}|=B\left(\frac{1}{2} R^{2} \varphi\right)
\]

where \(\varphi\) is the sector angle, enclosed by the field.\\
Now, magnitude of induced e.m.f. is given by,

\[
\xi_{i n}=\left|\frac{d \Phi_{t}}{d t}\right|=\left|\frac{B R^{2}}{2} \frac{d \varphi}{d t}\right|=\frac{B R^{2}}{2} \omega,
\]

where \(\omega\) is the angular velocity of the disc. But as it starts rotating from rest at \(t \equiv 0\) with an angular acceleration \(\beta\) its angular velocity \(\omega(t)=\beta t\). So,

\[
\xi_{i n}=\frac{B R^{2}}{2} \beta t .
\]

According to Lenz law the first half cycle current in the loop is in anticlockwise sense, and in subsequent half cycle it is in clockwise sense.\\
Thus in general, \(\xi_{i n}=(-1)^{n} \frac{B R^{2}}{2} \beta t\), where \(n\) in number of half revolutions.\\
The plot \(\xi_{i n}(t)\), where \(t_{n}=\sqrt{2 \pi n / \beta}\) is shown in the answer sheet.\\
3.293 Field, due to the current carrying wire in the region, right to it, is directed into the plane of the paper and its magnitude is given by,\\
\(B=\frac{\mu_{0}}{2 \pi} \frac{i}{r}\) where \(r\) is the perpendicular distance from the wire.\\
As \(B\) is same along the length of the rod thus motional e.m.f.

\[
\xi_{i n}=\left|-\int_{1}^{2}(\vec{v} \times \vec{B}) \cdot d \vec{l}\right|=v B l
\]

and it is directed in the sense of ( \(\vec{v} \times \vec{B}\) )\\
So, current (induced) in the loop,

\[
i_{i n}=\frac{\xi_{i n}}{R}=\frac{1}{2} \frac{\mu_{0} I v i}{\pi R r}
\]

3.294 Field, due to the current carrying wire, at a perpendicular distance \(x\) from it is given by,

\[
B(x)=\frac{\mu_{0}}{2 \pi} \frac{i}{x}
\]

Motional e.m.f is given by \(\left|\int-(\vec{v} \times \vec{B}) \cdot d \vec{l}\right|\)\\
There will be no induced e.m.f. in the segments (2) and (4) as, \(\vec{v} \uparrow \uparrow d \vec{l}\) and magnitude of e.m.f. induced in 1 and 3 , will be

\[
\xi_{1}=v\left(\frac{\mu_{0}}{2 \pi} \frac{i}{x}\right) a \text { and } \xi_{2}=v\left(\frac{\mu_{0}}{2 \pi} \frac{i}{(a+x)}\right) a,
\]

respectively, and their sense will be in the direction of \((\vec{v} \times \vec{B})\).\\
So, e.m.f., induced in the network \(=\xi_{1}-\xi_{2}\left[\right.\) as \(\left.\xi_{1}>\xi_{2}\right]\)

\[
=\frac{a v \mu_{0} i}{2 \pi}\left[\frac{1}{x}-\frac{1}{a+x}\right]=\frac{v a^{2} \mu_{0} i}{2 \pi x(a+x)}
\]

3.295 As the rod rotates, an emf.

\[
\frac{d}{d t} \frac{1}{2} a^{2} \theta \cdot B=\frac{1}{2} a^{2} B \omega
\]

is induced in it. The net current in the conductor is then \(\frac{\xi(t)-\frac{1}{2} a^{2} B \omega}{R}\)\\
A magnetic force will then act on the conductor of magnitude \(B I\) per unit length. Its direction will be normal to \(B\) and the rod and its torque will be

\[
\int_{o}^{a}\left(\frac{\xi(t)-\frac{1}{2} a^{2} B \omega}{R}\right) d x B x
\]

Obviously both magnetic and mechanical torque acting on the C.M. of the rod must be equal but opposite in sense. Then\\
for equilibrium at constant \(\omega\)

\[
\frac{\xi(t)-\frac{1}{2} a^{2} B \omega}{R} \cdot \frac{B a^{2}}{2}=\frac{1}{2} m g a \sin \omega t
\]

or, \(\quad \xi(t)=\frac{1}{2} a^{2} B \omega+\frac{m g R}{a B} \sin \omega t=\frac{1}{2 a B}\left(a^{3} B^{2} \omega+2 m g R \sin \omega t\right)\)\\
(The answer given in the book is incorrect dimensionally.)\\
3.296 From Lenz's law, the current through the connector is directed form \(A\) to \(B\). Here \(\xi_{\text {in }}=v B l\) between \(A\) and \(B\)\\
where \(v\) is the velocity of the rod at any moment.\\
For the rod, from \(F_{x}=m w_{x}\)\\
or, \(\quad m g \sin \alpha-i l B=m w\)\\
For steady state, acceleration of the rod must be equal to zero.\\
Hence, \(\quad m g \sin \alpha=i l B\)\\
\includegraphics[max width=\textwidth, alt={}, center]{9f3c5a8e-ac5d-42f4-8eb6-ea2af623f2d8-390_541_456_1062_914}

But, \(\quad i=\frac{\xi_{\text {in }}}{R}=\frac{v B l}{R}\)\\
From (1) and (2) \(v=\frac{m g \sin \alpha R}{B^{2} l^{2}}\)\\
3.297 From Lenz's law, the current through the copper bar is directed from 1 to 2 or in other words, the induced crrrent in the circuit is in clockwise sense.

Potential difference across the capacitor plates,

\[
\frac{q}{C}=\xi_{i n} \quad \text { or, } \quad q=C \xi_{i n}
\]

\begin{center}
\includegraphics[max width=\textwidth, alt={}]{9f3c5a8e-ac5d-42f4-8eb6-ea2af623f2d8-390_479_527_1617_861}
\end{center}

Hence, the induced current in the loop,

\[
i=\frac{d q}{d t}=C \frac{d \xi_{i n}}{d t}
\]

But the variation of magnetic flux through the loop is caused by the movement of the bar. So, the induced e.m.f. \(\xi_{\text {in }}=B l v\)\\
and,

\[
\frac{d \xi_{i n}}{d t}=B l \frac{d v}{d t}=B l w
\]

Hence,

\[
i=C \frac{d \xi}{d t}=C B l w
\]

Now, the forces acting on the bars are the weight and the Ampere's force, where \(F_{\text {amp }}=i l B(C B l w) B=C l^{2} B^{2} w\).\\
From Newton's second law, for the rod, \(F_{x}=m w_{x}\)\\
or,

\[
m g \sin \alpha-C l^{2} B^{2} w=m w
\]

Hence

\[
w=\frac{m g \sin \alpha}{C l^{2} B^{2}+m}=\frac{g \sin \alpha}{1+\frac{l^{2} B^{2} C}{m}}
\]

3.298 Flux of \(\vec{B}\), at an arbitrary moment of time \(t\) :

\[
\Phi_{t}=\vec{B} \cdot \vec{S}=B \frac{\pi a^{2}}{2} \cos \omega t,
\]

From Faraday's law, induced e.m.f., \(\boldsymbol{\xi}_{\text {in }}=-\frac{d \Phi}{d t}\)

\[
=-\frac{d\left(B \pi \frac{a^{2}}{2} \cos \omega t\right)}{d t}=\frac{B \pi a^{2} \omega}{2} \sin \omega t .
\]

and induced current, \(\quad i_{i n}=\frac{\xi_{i n}}{R}=\frac{B \pi a^{2}}{2 R} \omega \sin \omega t\).\\
Now thermal power, generated in the circuit, at the moment \(t=t\) :

\[
P(t)=\xi_{i n} \times i_{i n}=\left(\frac{B \pi a^{2} \omega}{2}\right)^{2} \frac{1}{R} \sin ^{2} \omega t
\]

and mean thermal power generated,

\[
<P>=\frac{\left[\frac{B \pi a^{2} \omega}{2}\right]^{2} \frac{1}{R} \int_{0}^{T} \sin ^{2} \omega t d t}{\int_{0}^{T} d t}=\frac{1}{2 R\left(\frac{\pi \omega a^{2} B}{2}\right)^{2}}
\]

Note : The claculation of \(\xi_{\text {in }}\) which can also be checked by using motional emf is correct even though the conductor is not a closed semicircle, for the flux linked to the rectangular part containing the resistance \(R\) is not changing. The answer given in the book is off by a factor \(1 / 4\).\\
3.299 The flux through the coil changes sign. Initially it is \(B S\) per turn.

Finally it is - BS per turn. Now if flux is \(\Phi\) at an intermediate state then the current at that moment will be

\[
i=\frac{-N \frac{d \Phi}{d t}}{R}
\]

So charge that flows during a sudden turning of the coil is

Hence,

\[
\begin{aligned}
& q=\int i d t=-\frac{N}{R}[\Phi-(-\Phi)]=2 N B S / R \\
& B=\frac{1}{2} \frac{q R}{N S}=0.5 \mathrm{~T} \text { on putting the values. }
\end{aligned}
\]

3.300 According to Ohm's law and Faraday's law of induction, the current \(i_{0}\) appearing in the frame, during its rotation, is determined by the formula,

\[
i_{0}=-\frac{d \Phi}{d t}=-\frac{L d i_{0}}{d t}
\]

Hence, the required amount of electricity (charge) is,

\[
q=\int i_{0} d t=-\frac{1}{R} \int\left(d \Phi+L d i_{0}\right)=-\frac{1}{R}\left(\Delta \Phi+L \Delta i_{0}\right)
\]

Since the frame has been stopped after rotation, the current in it vanishes, and hence \(\Delta i_{0}=0\). It remains for us to find the increment of the flux \(\Delta \Phi\) through the frame \(\left(\Delta \Phi=\Phi_{2}-\Phi_{1}\right)\).

Let us choose the normal \(\vec{n}\) to the plane of the frame, for instance, so that in the final position, \(\vec{n}\) is directed behind the plane of the figure (along \(\vec{B}\) ).\\
\includegraphics[max width=\textwidth, alt={}, center]{9f3c5a8e-ac5d-42f4-8eb6-ea2af623f2d8-392_400_561_968_840}

Then it can be easily seen that in the final pos.tion, \(\Phi_{2}>0\), while in the initial position, \(\Phi_{1}<0\) (the normal is opposite to \(\vec{B}\) ), and \(\Delta \Phi\) turns out to be simply equal to the fulx through the surface bounded by the final and initial positions of the frame :

\[
\Delta \Phi=\Phi_{2}+\left|\Phi_{1}\right|=\int_{b-a}^{b+a} B a d r
\]

where \(B\) is a function of \(r\), whose form can be easily found with the help of the theorem of circulation. Finally omitting the minus sign, we obtain,

\[
q=\frac{\Delta \Phi}{R}=\frac{\mu_{0} a i}{2 \pi R} \ln \frac{b+a}{b-a}
\]

3.301 As \(\vec{B}\), due to the straight current carrying wire, varies along the rod (connector) and enters linerarly so, to make the calculations simple, \(\vec{B}\) is made constant by taking its average value in the range \([a, b]\).\\
\includegraphics[max width=\textwidth, alt={}, center]{9f3c5a8e-ac5d-42f4-8eb6-ea2af623f2d8-393_432_1229_70_117}\\
(a) The flux of \(\vec{B}\) changes through the loop due to the movement of the connector. According to Lenz's law, the current in the loop will be anticlockwise. The magnitude of motional e.m.f.,

\[
\begin{gathered}
\xi_{\text {in }}=v<B>(b-a) \\
=\frac{\mu_{0}}{2 \pi} \frac{i_{0}}{(b-a)} \ln \frac{b}{a}(b-a) \frac{d x}{d t}=\frac{\mu_{0}}{2 \pi} i_{0} \ln \frac{b}{a} v
\end{gathered}
\]

So, induced current

\[
i_{i n}=\frac{\xi_{i n}}{R}=\frac{\mu_{0}}{2 \pi} \frac{i_{0} v}{R} \ln \frac{b}{a}
\]

(b) The force required to maintain the constant velocity of the connector must be the magnitude equal to that of Ampere's acting on the connector, but in opposite direction.\\
So, \(\quad F_{e x t}=i_{i n} l<B>=\left(\frac{\mu_{0}}{2 \pi} \frac{i_{0}}{R} v \ln \frac{b}{a}\right)(b-a)\left(\frac{\mu_{0}}{2 \pi} \frac{i_{0}}{(b-a)} \ln \frac{b}{a}\right)\)\\
\(=\frac{v}{R}\left(\frac{\mu_{0}}{2 \pi} i_{0} \ln \frac{b}{a}\right)^{2}\), and will be directed as shown in the (Fig.)\\
3.302 (a) The flux through the loop changes due to the movement of the \(\operatorname{rod} A B\). According to Lenz's law current should be anticlockwise in sense as we have assumed \(\vec{B}\) is directed into the plane of the loop. The motion e.m.f \(\xi_{\text {in }}(t)=B l v\)\\
and induced current \(i_{i n}=\frac{v B l}{R}\)\\
From Newton's law in projection form \(F_{x}=m w_{x}\)

\[
-F_{a m p}=m \cdot \frac{v d v}{d x}
\]

But \(F_{\text {amp }}=i_{\text {in }} l B=\frac{v B^{2} l^{2}}{R}\)\\
\includegraphics[max width=\textwidth, alt={}, center]{9f3c5a8e-ac5d-42f4-8eb6-ea2af623f2d8-393_341_457_1520_874}

So, \(\quad-\frac{v B^{2} l^{2}}{R}=m v \frac{d v}{d x}\)\\
or,

\[
\int_{0}^{x} d x=-\frac{m R}{B^{2} l^{2}} \int_{v_{0}}^{0} d v \text { or, } x=\frac{m R v_{0}}{B^{2} l^{2}}
\]

(b) From equation of energy conservation; \(E_{f}-E_{i}+\) Heat liberated \(=A_{\text {cell }}+A_{\text {ext }}\)

\[
\left[0-\frac{1}{2} m v_{0}^{2}\right]+\text { Heat liberated }=0+0
\]

So, heat liberated \(=\frac{1}{2} m v_{0}^{2}\)\\
3.303 With the help of the calculation, done in the previous problem, Ampere's force on the connector,\\
\(\vec{F}_{a m p}=\frac{\nu B^{2} l^{2}}{R}\) directed towards left.\\
Now from Newton's second law,

\[
F-F_{a m p}=m \frac{d v}{d t}
\]

So, \(\quad F=\frac{v B^{2} l^{2}}{R}+m \frac{d v}{d t}\)\\
or, \(\quad \int_{0}^{t} d t=m \int_{0}^{v} \frac{d v}{F-\frac{v B^{2} l^{2}}{R}}\)\\
\includegraphics[max width=\textwidth, alt={}, center]{9f3c5a8e-ac5d-42f4-8eb6-ea2af623f2d8-394_369_484_714_907}\\
or,

\[
\frac{t}{m}=-\frac{R}{B^{2} l^{2}} \ln \left(\frac{F-\frac{v B^{2} l^{2}}{R}}{F}\right)
\]

Thus

\[
v=\left(1-e \frac{-t B^{2} l^{2}}{R m}\right) \frac{R F}{B^{2} l^{2}}
\]

3.304 According to Lenz, the sense of induced e.m.f. is such that it opposes the cause of change of flux. In our problem, magnetic field is directed away from the reader and is diminishing.

\begin{figure}[h]
\begin{center}
  \includegraphics[alt={},max width=\textwidth]{9f3c5a8e-ac5d-42f4-8eb6-ea2af623f2d8-394_198_196_1613_127}
\captionsetup{labelformat=empty}
\caption{(a)}
\end{center}
\end{figure}

\begin{figure}[h]
\begin{center}
  \includegraphics[alt={},max width=\textwidth]{9f3c5a8e-ac5d-42f4-8eb6-ea2af623f2d8-394_222_204_1583_394}
\captionsetup{labelformat=empty}
\caption{(b)}
\end{center}
\end{figure}

\begin{figure}[h]
\begin{center}
  \includegraphics[alt={},max width=\textwidth]{9f3c5a8e-ac5d-42f4-8eb6-ea2af623f2d8-394_240_233_1576_674}
\captionsetup{labelformat=empty}
\caption{(C)}
\end{center}
\end{figure}

\begin{figure}[h]
\begin{center}
  \includegraphics[alt={},max width=\textwidth]{9f3c5a8e-ac5d-42f4-8eb6-ea2af623f2d8-394_248_370_1567_983}
\captionsetup{labelformat=empty}
\caption{(d)}
\end{center}
\end{figure}

So, in figure (a), in the round conductor, it is clockwise and there is no current in the connector\\
In figure (b) in the outside conductor, clockwise.

In figure (c) in both the conductor, clockwise; and there is no current in the connector to obey the charge conservation.\\
In figure (d) in the left side of the figure, clockwise.\\
3.305 The loops are connected in such a way that if the current is clockwise in one, it is anticlockwise in the other. Hence the e.m.f. in loop \(b\) opposes the e.m.f. in loop \(a\).\\
e.m.f. in loop \(a=\frac{d}{d t}\left(a^{2} B\right)=a^{2} \frac{d}{d t}\left(B_{0} \sin \omega t\right)\)

Similarly, e.m.f. in loop \(b=b^{2} B_{0} \omega \cos \omega t\).\\
Hence, net e.m.f. in the circuit \(=\left(a^{2}-b^{2}\right) B_{0} \omega \cos \omega t\), as both the e.m.f's are in opposite sense, and resistance of the circuit \(=4(a+b) \rho\)\\
Therefore, the amplitude of the current

\[
=\frac{\left(a^{2}-b^{2}\right) B_{0} \omega}{4(a+b) \rho}=0.5 \mathrm{~A} .
\]

3.306 The flat shape is made up of concentric loops, having different radii, varying from 0 to \(a\).\\
Let us consider an elementary loop of radius \(r\), then e.m.f. induced due to this loop \(=\frac{-d(\vec{B} \cdot \vec{S})}{d t}=\pi r^{2} B_{0} \omega \cos \omega t\).\\
and the total induced e.m.f.,


\begin{equation*}
\xi_{i n d}=\int_{0}^{a}\left(\pi r^{2} B_{0} \omega \cos \omega t\right) d N \tag{1}
\end{equation*}


where \(\pi r^{2} \omega \cos \omega t\) is the contribution of one turn of radius \(r\) and \(d N\) is the number of turns in the interval \([r, r+d r]\).

So,


\begin{equation*}
d N=\left(\frac{N}{a}\right) d r \tag{2}
\end{equation*}


From (1) and (2), \(\xi=\int_{0}^{a}-\left(\pi r^{2} B_{0} \omega \cos \omega t\right) \frac{N}{a} d r=\frac{\pi B_{0} \omega \cos \omega t N a^{2}}{3}\)\\
Maximum value of e.m.f. amplitude \(\xi_{\max }=\frac{1}{3} \pi B_{0} \omega N a^{2}\)\\
3.307 The flux through the loop changes due to the variation in \(\vec{B}\) with time and also due to the movement of the connector.

So,

\[
\xi_{i n}=\left|\frac{d(\vec{B} \cdot \vec{S})}{d t}\right|=\left|\frac{d(B S)}{d t}\right| \text { as } \vec{S} \text { and } \vec{B} \text { are colliniear }
\]

But, \(B\), after \(t\) sec. of beginning of motion \(=B t\), and \(S\) becomes \(=l \frac{1}{2} w t^{2}\), as connector starts moving from rest with a constant acceleration \(w\).

So,

\[
\xi_{i n d}=\frac{3}{2} B l w t^{2}
\]

3.308 We use \(B=\mu_{0} n I\)

Then, from the law of electromagnetic induction

\[
\oint \vec{E} \cdot d \vec{l}=-\frac{d \Phi}{d t}
\]

So, for \(r<a\)

\[
E_{\varphi} 2 \pi r=-\pi r^{2} \mu_{0} n \dot{I} \text { or, } E_{\varphi}=-\frac{1}{2} \mu_{0} n r \dot{I} .(\text { where } I=d I / d t)
\]

For \(r>a\)

\[
E_{\varphi} 2 \pi r=-\pi a^{2} \mu_{0} n \dot{I} \text { or, } E_{\varphi}=-\mu_{0} n \dot{I} a^{2} / 2 r
\]

The meaning of minus sign can be deduced from Lenz's law.\\
3.309 The e.m.f. induced in the turn is \(\mu_{0} n \dot{I} \pi \frac{d^{2}}{4}\)

The resistance is \(\frac{\pi d}{S} \rho\).\\
So, the current is \(\frac{\mu_{0} n \dot{I} S d}{4 \rho}=2 \mathrm{~mA}\), where \(\rho\) is the resistivity of copper.\\
3.310 The changing magnetic field will induce an e.m.f. in the ring, which is obviously equal, in the two parts by symmetry (the e.m.f. induced by electromagnetic induction does not depend on resistance). The current, that will flow due to this, will be different in the two parts. This will cause an acceleration of charge, leading to the setting up of an electric field \(E\) which has opposite sign in the two parts. Thus,

\[
\frac{\xi}{2}-\pi a E=r I \text { and, } \frac{\xi}{2}+\pi a E=\eta r I
\]

where \(\xi\) is the total induced e.m.f. From this,\\
and

\[
\begin{gathered}
\xi=(\eta+1) r I, \\
E=\frac{1}{2 \pi a}(\eta-1) r I=\frac{1}{2 \pi a} \frac{\eta-1}{\eta+1} \xi
\end{gathered}
\]

But by Faraday's law, \(\xi=\pi a^{2} b\)\\
so,

\[
E=\frac{1}{2} a b \frac{\eta-1}{\eta+1}
\]

3.311 Go to the rotating frame with an instantaneous angular velocity \(\vec{\omega}(t)\). In this frame, a Coriolis force, \(2 m \overrightarrow{v^{\prime}} \times \vec{\omega}(t)\)\\
acts which must be balanced by the magnetic force, \(e \vec{v} \times \vec{B}(t)\)\\
Thus,

\[
\vec{\omega}(t)=-\frac{e}{2 m} \vec{B}(t)
\]

(It is assumed that \(\vec{\omega}\) is small and varies slowly, so \(\omega^{2}\) and \(\dot{\omega}\) can be neglected.)\\
3.312 The solenoid has an inductance,

\[
L=\mu_{0} n^{2} \pi b^{2} l,
\]

where \(n=\) number of turns of the solenoid per unit length. When the solenoid is connected to the source an e.m.f. is set up, which, because of the inductance and resistance, rises slowly, according to the equation,

\[
R I+L U=V
\]

This has the well known solution,

\[
I=\frac{V}{R}\left(1-e^{-t R / L}\right) .
\]

Corresponding to this current, an e.m.f. is induced in the ring. Its magnetic field \(B=\mu_{0} n I\) in the solenoid, produces a force per unit length, \(\frac{d F}{d l}=B i=\mu_{0}^{2} n^{2} \pi a^{2} I I / r\)

\[
=\frac{\mu_{0}^{2} \pi a^{2} V^{2}}{r}\left(\frac{n^{2}}{R L}\right) e^{-t R / L}\left(1-e^{-t R / L}\right),
\]

acting on each segment of the ring. This force is zero initially and zero for large \(t\). Its maximum value is for some finite \(t\). The maximum value of

So

\[
\begin{gathered}
e^{-t R / L}\left(1-e^{-t R / L}\right)=\frac{1}{4}-\left(\frac{1}{2}-e^{-t R / L}\right)^{2} \text { is } \frac{1}{4} \\
\frac{d F_{\max }}{d l}=\frac{\mu_{0}^{2} \pi a^{2} V^{2}}{r} \frac{n^{2}}{4 R L}=\frac{\mu_{0} a^{2} V^{2}}{4 r R l b^{2}}
\end{gathered}
\]

3.313 The amount of heat generated in the loop during a small time interval \(d t\),

\[
d Q=\xi^{2} / R d t, \text { but, } \xi=-\frac{d \Phi}{d t}=2 a t-a \tau
\]

So,

\[
d Q=\frac{(2 a t-a \tau)^{2}}{R} d t
\]

and hence, the amount of heat, generated in the loop during the time interval 0 to \(\tau\).\\
3.314 Take an elementary ring of radius \(r\) and width dr. The e.m.f. induced in this elementary ring is \(\pi r^{2} \beta\). Now the conductance of this ring is.

\[
d\left(\frac{1}{R}\right)=\frac{h d r}{\rho 2 \pi r} \text { so } d I=\frac{h r d r}{2 \rho} \beta
\]

Integrating we get the total current,

\[
I=\int_{a}^{b} \frac{h r d r}{2 \rho} \beta=\frac{h \beta\left(b^{2}-a^{2}\right)}{4 \rho}
\]

\includegraphics[max width=\textwidth, alt={}, center]{9f3c5a8e-ac5d-42f4-8eb6-ea2af623f2d8-397_503_467_1520_862}\\
3.315 Given \(L=\mu_{0} n^{2} V=\mu_{0} n^{2} l_{0} \pi R^{2}\), where \(R\) is the radius of the solenoid.

Thus,

\[
n=\sqrt{\frac{L}{\mu_{0} l_{0} \pi}} \frac{1}{R} .
\]

So, length of the wire required is,

\[
l=n l_{0} 2 \pi R=\sqrt{\frac{4 \pi L l_{0}}{\mu_{0}}}=0.10 \mathrm{~km}
\]

3.316 From the previous problem, we know that,\\
\(l^{\prime}=\) length of the wire needed \(=\sqrt{\frac{L l 4 \pi}{\mu_{0}}}\), where \(l=\) length of solenoid here.\\
Now, \(R=\frac{\rho_{0} l^{\prime}}{S}\), (where \(S=\) area of crossection of the wire. Also \(m=\rho S l^{\prime}\) )\\
Thus, \(\quad l^{\prime}=\frac{R S}{\rho_{0}}=\frac{R m}{\rho \rho_{0} l^{\prime}} \quad\) or \(\quad l^{\prime}=\sqrt{\frac{R m}{\rho \rho_{0}}}\)\\
where \(\rho_{0}=\) resistivity of copper and \(\rho=\) its density.\\
Equating,

\[
\begin{aligned}
\frac{R m}{\rho \rho_{0}} & =\frac{L l}{\mu_{0} / 4 \pi} \\
L & =\frac{\mu_{0}}{4 \pi} \frac{m R}{\rho \rho_{0} l}
\end{aligned}
\]

3.317 The current at time \(t\) is given by,

\[
I(t)=\frac{V}{R}\left(1-e^{-t R L}\right)
\]

The steady state value is, \(I_{0}=\frac{V}{R}\)\\
and

\[
\frac{I(t)}{I_{0}}=\eta=1-e^{-t R / L} \text { or }, e^{-t R / L}=1-\eta
\]

or,

\[
t_{o} \frac{R}{L}=\ln \frac{1}{1-\eta} \quad \text { or, } \quad t_{o}=\frac{L}{R} \ln \frac{1}{1-\eta}=1.49 \mathrm{~s}
\]

3.318 The time constant \(\tau\) is given by

\[
\tau=\frac{L}{R}=\frac{L}{\rho_{0} \frac{l_{0}}{S}},
\]

where, \(\rho_{0}=\) resistivity, \(l_{0}=\) length of the winding wire, \(S=\) cross section of the wire.\\
But

\[
m=l \rho_{0} S
\]

So eliminating \(S, \tau=\frac{L}{\frac{\rho_{0} l_{0}}{m / \rho l_{0}}}=\frac{m L}{\rho \rho_{0} l_{0}^{2}}\)

From problem \(3.315 l_{0}=\sqrt{\frac{4 \pi l L}{\mu_{0}}}\)\\
(note the interchange of \(l\) and \(l_{0}\) because of difference in notation here.)\\
Thus,

\[
\tau=\frac{m L}{\rho \rho_{0} \frac{4 \pi}{\mu_{0}} L l}=\mu_{0} 4 \pi \frac{m}{\rho \rho_{0} l}=0.7 \mathrm{~ms},
\]

3.319 Between the cables, where \(a<r<b\), the magnetic field \(\vec{H}\) satisfies

\[
H_{\varphi} 2 \pi r=I \quad \text { or, } \quad H_{\varphi}=\frac{I}{2 \pi r}
\]

So

\[
B_{\varphi}=\frac{\mu \mu_{0} I}{2 \pi r}
\]

The associated flux per unit length is, \(\Phi=\int_{r=a}^{r=b} \frac{\mu \mu_{0} I}{2 \pi r} \times 1 \times d r=\frac{\mu \mu_{0} I}{2 \pi} \ln \frac{b}{a}\)\\
Hence, the inductance per unit length \(L_{1}=\frac{\Phi}{I}=\frac{\mu \mu_{0}}{2 \pi} \ln \eta\), where \(\eta=\frac{b}{a}\)\\
We get \(L_{1}=0.26 \frac{\mu H}{m}\)\\
3.320 Within the solenoid, \(H_{\varphi} \cdot 2 \pi r=N I\) or \(H_{\varphi}=\frac{N I}{2 \pi r}, B_{\varphi}=\mu \mu_{0} \frac{N I}{2 \pi r}\)\\
and the flux, \(\Phi=N \Phi_{1},=N \frac{\mu \mu_{0}}{2 \pi} N I \int_{b}^{a+b} \frac{a d r}{r}\)\\
Finally,

\[
L=\frac{\mu \mu_{0}}{2 \pi} N^{2} a \ln \left(1+\frac{a}{b}\right)
\]

3.321 Neglecting end effects the magnetic field \(B\), between the plates, which is mainly parallel to the plates, is \(B=\mu_{0} \frac{I}{b}\)\\
(For a derivation see 3.229 b)\\
Thus, the associated flux per unit length of the plates is,

\[
\Phi=\mu_{0} \frac{I}{b} \times h \times 1=\left(\mu_{0} \frac{h}{b}\right) \times I .
\]

\begin{center}
\includegraphics[max width=\textwidth, alt={}]{9f3c5a8e-ac5d-42f4-8eb6-ea2af623f2d8-399_462_622_1460_801}
\end{center}

So, \(L_{1}=\) inductance per unit length \(=\mu_{0} \frac{h}{b}=25 \mathrm{nH} / \mathrm{m}\).\\
3.322 For a single current carrying wire, \(B_{\varphi}=\frac{\mu_{0} I}{2 \pi r}(r>a)\). For the double line cable, with current, flowing in opposite directions, in the two conductors, \(B_{\varphi} \approx \frac{\mu_{0} I}{\pi r}\), between the cables, by superposition. The associated flux is,

\[
\begin{gathered}
\Phi=\int_{a}^{d-a} \frac{\mu_{0} I}{\pi} \frac{d r \times 1}{r} \approx \frac{\mu_{0} I}{\pi} \ln \frac{d}{a}=\frac{\mu_{0}}{\pi} \ln \eta \times I, \text { per unit length } \\
L_{1}=\frac{\mu_{0}}{\pi} \ln \eta
\end{gathered}
\]

Hence,\\
is the inductance per unit length.\\
3.323 In a superconductor there is no resistance, Hence,

So integrating,

\[
L \frac{d I}{d t}=+\frac{d \Phi}{d t}
\]

because

\[
\begin{aligned}
I & =\frac{\Delta \Phi}{L}=\frac{\pi a^{2} B}{L} \\
\Delta \Phi & =\Phi_{f}-\Phi_{i}, \Phi_{f}=\pi a^{2} B, \Phi_{i}=0
\end{aligned}
\]

Also, the work done is, \(A=\int \xi I d t=\int I d t \frac{d \Phi}{d t}=\frac{1}{2} L I^{2}=\frac{1}{2} \frac{\pi^{2} a^{4} B^{2}}{L}\)\\
3.324 In a solenoid, the inductance \(L=\mu \mu_{0} n^{2} V=\mu \mu_{0} \frac{N^{2} S}{l}\),\\
where \(S=\) area of cross section of the solenoid, \(l=\) its length, \(V=S l, N=n l=\) total number of turns.\\
When the length of the solenoid is increased, for example, by pulling it, its inductance will decrease. If the current remains unchanged, the flux, linked to the solenoid, will also decrease. An induced e.m.f. will then come into play, which by Lenz's law will try to oppose the decrease of flux, for example, by increasing the current. In the superconducting state the flux will not change and so,

\[
\frac{I}{l}=\text { constant }
\]

Hence,

\[
\frac{I}{l}=\frac{I_{0}}{l_{0}}, \quad \text { or }, \quad I=I_{0} \frac{l}{l_{0}}=I_{0}(1+\eta)
\]

3.325 The flux linked to the ring can not change on transition to the superconduction state, for reasons, similar to that given above. Thus a current \(I\) must be induced in the ring, where,

\[
I=\frac{\Phi}{L}=\frac{\pi a^{2} B}{\mu_{0} a\left(\ln \frac{8 a}{b}-2\right)}=\frac{\pi a B}{\mu_{0}\left(\ln \frac{8 a}{b}-2\right)}
\]

3.326 We write the equation of the circuit as,

\[
R i+\frac{L}{\eta} \frac{d i}{d t}=\xi
\]

for \(t \geq 0\). The current at \(t=0\) just after inductance is changed, is \(i=\eta \frac{\xi}{R}\), so that the flux through the inductance is unchanged.\\
We look for a solution of the above equation in the form

\[
i=A+B e^{-t / C}
\]

Substituting \(C=\frac{L}{\eta R}, B=\eta-1, A=\frac{\xi}{R}\)\\
Thus,

\[
i=\frac{\xi}{R}\left(1+(\eta-1) e^{-\eta R t / L}\right)
\]

3.327 Clearly, \(L \frac{d i}{d t}=R(I-i)=\xi-R I\)

So, \(2 L \frac{d i}{d t}=\xi-R i\)\\
This equation has the solution (as in 3.312)\\
\includegraphics[max width=\textwidth, alt={}, center]{9f3c5a8e-ac5d-42f4-8eb6-ea2af623f2d8-401_293_500_617_864}\\
\(i=\frac{\xi}{R}\left(1-e^{-t R / 2 L}\right)\)\\
3.328 The equations are,

\[
L_{1} \frac{d i_{1}}{d t}=L_{2} \frac{d i_{2}}{d t}=\xi-R\left(i_{1}+i_{2}\right)
\]

Then, \(\frac{d}{d t}\left(L_{1} i_{1}-L_{2} i_{2}\right)=0\)\\
or, \(L_{1} i_{1}-L_{2} i_{2}=\) constant\\
But initially at \(t=0, i_{1}=i_{2}=0\)\\
so constant must be zero and at all times,\\
\includegraphics[max width=\textwidth, alt={}, center]{9f3c5a8e-ac5d-42f4-8eb6-ea2af623f2d8-401_395_453_1009_888}

\[
L_{1} i_{1}=L_{2} i_{2}
\]

In the final steady atate, current must obviously be \(i_{1}+i_{2}=\frac{\xi}{R}\). Thus in steady state,

\[
i_{1}=\frac{\xi L_{2}}{R\left(L_{1}+L_{2}\right)} \text { and } i_{2}=\frac{\xi L_{1}}{R\left(L_{1}+L_{2}\right)}
\]

3.329 Here, \(B=\frac{\mu_{0} I}{2 \pi r}\) at a distance \(r\) from the wire. The flux through the frame is obtained as,

\[
\begin{aligned}
& \Phi_{12}=\int_{l}^{a+l} \frac{\mu_{0} I}{2 \pi r} b d r=\frac{\mu_{0} b}{2 \pi} I \ln \left(1+\frac{a}{l}\right) \\
& \text { Thus, } L_{12}=\frac{\Phi_{12}}{I}=\frac{\mu_{0} b}{2 \pi} \ln \left(1+\frac{a}{l}\right)
\end{aligned}
\]

\includegraphics[max width=\textwidth, alt={}, center]{9f3c5a8e-ac5d-42f4-8eb6-ea2af623f2d8-401_342_366_1729_893}\\
3.330 Here also, \(B=\frac{\mu_{0} I}{2 \pi r}\) and \(\Phi=\mu_{0} \mu \frac{I}{2 \pi} \int_{a}^{b} \frac{h d r}{r} N\).

Thus,

\[
L_{12}=\frac{\mu \mu_{0} h N}{2 \pi} \ln \frac{b}{a}
\]

3.331 The direct calculation of the flux \(\Phi_{2}\) is a rather complicated problem, since the configuration of the field itself is complicated. However, the application of the reciprocity theorem simplifies the solution of the problem. Indeed, let the same current \(i\) flow through loop 2. Then the magnetic flux created by this current through loop 1 can be easily found.

Magnetic induction at the centre of the loop, \(: B=\frac{\mu_{0} i}{2 b}\)\\
So, flux throug loop 1 : \(\Phi_{12}=\pi a^{2} \frac{\mu_{0} i}{2 b}\) and from reciprocity theorem,\\
\(\varphi_{12}=\Phi_{21}, \Phi_{21}=\frac{\mu_{0} \pi a^{2} i}{2 b}\)\\
\includegraphics[max width=\textwidth, alt={}, center]{9f3c5a8e-ac5d-42f4-8eb6-ea2af623f2d8-402_213_447_713_867}

So, \(L_{12}=\frac{\Phi_{21}}{i}=\frac{1}{2} \mu_{0} \pi a^{2} / b\)\\
3.332 Let \(\vec{p}_{m}\) be the magnetic moment of the magnet \(M\). Then the magnetic field due to this magnet is,

\[
\frac{\mu_{0}}{4 \pi}\left[\frac{3\left(\vec{p}_{m} \cdot \vec{r}\right) \vec{r}}{r^{5}}-\frac{\vec{p}_{m}}{r^{3}}\right]
\]

The flux associated with this, when the magnet is along the axis at a distance \(x\) from the centre, is

\[
\Phi=\frac{\mu_{0}}{4 \pi} \int\left[\frac{3\left(\vec{p}_{m} \cdot \vec{r}\right) \vec{r}}{r^{5}}-\frac{\vec{p}_{m}}{r^{3}}\right] \cdot d \vec{S}=\Phi_{1}-\Phi_{2}
\]

where, \(\Phi_{2}=\frac{\mu_{0}}{4 \pi} p_{m} \int_{0}^{a} \frac{2 \pi \rho d \rho}{\left(x^{2}+\rho^{2}\right)^{3 / 2}}=\frac{\mu_{0} p_{m}}{2}\left(\frac{1}{x}-\frac{1}{\sqrt{x^{2}+a^{2}}}\right)\)\\
and \(\Phi_{1}=\frac{3 \mu_{0} p_{m} x^{2}}{4 \pi} \int_{0}^{a} \frac{2 \pi \rho d \rho}{\left(x^{2}+\rho^{2}\right)^{5 / 2}}\)\\
\(=\frac{\mu_{0} p_{m} x^{2}}{2}\left(\frac{1}{x^{3}}-\frac{1}{\left(x^{2}+a^{2}\right)^{3 / 2}}\right)\)\\
So, \(\Phi=\frac{-\mu_{0} p_{m} a^{2}}{2\left(x^{2}+a^{2}\right)^{3 / 2}}\)\\
\includegraphics[max width=\textwidth, alt={}, center]{9f3c5a8e-ac5d-42f4-8eb6-ea2af623f2d8-402_322_585_1676_742}

When the flux changes, an e.m.f. \(-N \frac{d \Phi}{d t}\) is induced and a current \(-\frac{N}{R} \frac{d \Phi}{d t}\) flows. The total charge \(q\), flowing, as the magnet is removed to infinity from \(x=0\) is,\\
or,

\[
\begin{gathered}
q=\frac{N}{R} \Phi(x=0)=\frac{N}{R} \cdot \frac{\mu_{0} p_{m}}{2 a} \\
p_{m}=\frac{2 a q R}{N \mu_{0}}
\end{gathered}
\]

3.333 If a current \(I\) flows in one of the coils, the magnetic field at the centre of the other coil is,

\[
B=\frac{\mu_{0} a^{2} I}{2\left(l^{2}+a^{2}\right)^{3 / 2}}=\frac{\mu_{0} a^{2} I}{2 l^{3}}, \text { as } l \gg a .
\]

The flux associated with the second coil is then approximately \(\mu_{0} \pi a^{4} I / 2 l^{3}\)\\
Hence,

\[
L_{12}=\frac{\mu_{0} \pi a^{4}}{2 l^{3}}
\]

3.334 When the current in one of the loop is \(I_{1}=\alpha t\), an e.m.f. \(L_{12} \frac{d I_{1}}{d t}=L_{12} \alpha\), is induced in the other loop. Then if the current in the other loop is \(I_{2}\) we must have,

\[
L_{2} \frac{d I_{2}}{d t}+R I_{2}=L_{12} \alpha
\]

This familiar equation has the solution, \(I_{2}=\frac{L_{12} \alpha}{R}\left(1-e^{\frac{-t R}{L_{2}}}\right)\) which is the required current\\
3.335 Initially, after a steady current is set up, the current is flowing as shown.

In steady condition \(i_{20}=\frac{\xi}{R}, i_{10}=\frac{\xi}{R_{0}}\).\\
When the switch is disconnected, the current through \(R_{0}\) changes from \(i_{10}\) to the right, to \(i_{20}\) to the left. (The current in the inductance cannot change suddenly.). We then have the equation,

\[
L \frac{d i_{2}}{d t}+\left(R+R_{0}\right) i_{2}=0
\]

This equation has the solution \(i_{2}=i_{20} e^{-t\left(R+R_{0}\right) / L}\)\\
\includegraphics[max width=\textwidth, alt={}, center]{9f3c5a8e-ac5d-42f4-8eb6-ea2af623f2d8-403_463_543_1267_843}

The heat dissipated in the coil is,

\[
\begin{aligned}
& Q=\int_{0}^{\infty} i_{2}^{2} R d t=i_{20}^{2} R \int_{0}^{\infty} e^{-2 t\left(R+R_{0}\right) \nu L} d t \\
= & R i_{20}^{2} \times \frac{L}{2\left(R+R_{0}\right)}=\frac{L \xi^{2}}{2 R\left(R+R_{0}\right)}=3 \mu \mathrm{~J}
\end{aligned}
\]

3.336 To find the magnetic field energy we recall that the flux varies linearly with current. Thus, when the flux is \(\Phi\) for current \(i\), we can write \(\phi=A i\). The total energy inclosed in the field, when the current is \(I\), is

\[
\begin{aligned}
W & =\int \xi i d t=\int N \frac{d \Phi}{d t} i d t \\
& =\int N d \Phi i=\int_{0}^{I} N A i d i=\frac{1}{2} N A I^{2}=\frac{1}{2} N \Phi I
\end{aligned}
\]

The characteristic factor \(\frac{1}{2}\) appears in this way.\\
3.337 We apply circulation theorem,

\[
H \cdot 2 \pi b=N I, \quad \text { or, } \quad H=N I / 2 \pi b .
\]

Thus the total energy,

\[
W=\frac{1}{2} B H \cdot 2 \pi b \cdot \pi a^{2}=\pi^{2} a^{2} b B H .
\]

Given \(N, I, b\) we know \(H\), and can find out \(B\) from the \(B-H\) curve. Then \(W\) can be calculated.\\
3.338 From \(\oint \vec{H} \cdot d \vec{r}=N I\),

\[
H \cdot \pi d+\frac{B}{\mu_{0}} \cdot b \approx N I,(d \gg b)
\]

Also,

\[
B=\mu \mu_{0} H . \text { Thus, } H=\frac{N I}{\pi d+\mu b} .
\]

Since \(B\) is continuous across the gap, \(B\) is given by, \(B=\mu \mu_{0} \frac{N I}{\pi d+\mu b}\), both in the magnetic and the gap.\\
(a) \(\frac{W_{\text {gap }}}{W_{\text {magnetic }}}=\frac{\frac{B^{2}}{2 \mu_{0}} \times S \times b}{\frac{B^{2}}{2 \mu \mu_{0}} \times S \times \pi d}=\frac{\mu b}{\pi d}\).\\
(b) The flux is \(N \int \vec{B} \cdot d \vec{S}=N \mu \mu_{0} \frac{N I}{\pi d+\mu b} \cdot S=\mu_{0} \frac{S N^{2} I}{b+\frac{\pi d}{\mu}}\),

So,

\[
L \approx \frac{\mu_{0} S N^{2}}{b+\frac{\pi d}{\mu}}
\]

Energy wise; total energy

\[
=\frac{B^{2}}{2 \mu_{0}}\left(\frac{\pi d}{\mu}+b\right) S=\frac{1}{2} \frac{\mu_{0} N^{2} S}{b+\frac{\pi d}{\mu}} \cdot I^{2}=\frac{1}{2} L I^{2}
\]

The \(L\), found in the one way, agrees with that, found in the other way. Note that, in calculating the flux, we do not consider the field in the gap, since it is not linked to the winding. But the total energy includes that of the gap.\\
3.339 When the cylinder with a linear charge density \(\lambda\) rotates with a circular frequency \(\omega\), a surface current density (charge \(/\) length × time) of \(i=\frac{\lambda \omega}{2 \pi}\) is set up.

The direction of the surface current is normal to the plane of paper at \(Q\) and the contribution of this current to the magnetic field at \(P\) is \(d \vec{B}=\frac{\mu_{0}}{4 \pi} \frac{i(\vec{e} \times \vec{r})}{r^{3}} d S\) where \(\vec{e} \quad\) is the direction of the current. In magnitude, \(|\vec{e} \times \vec{r}|=r\), since \(\vec{e}\) is normal to \(\vec{r}\) and the direction of \(d \vec{B}\) is as shown.\\
It's component, \(d \overrightarrow{B_{\|}}\)cancels out by cylindrical\\
\includegraphics[max width=\textwidth, alt={}, center]{9f3c5a8e-ac5d-42f4-8eb6-ea2af623f2d8-405_579_505_309_869}\\
symmetry. The component that survives is,

\[
\left|\vec{B}_{\perp}\right|=\frac{\mu_{0}}{4 \pi} \int \frac{i d S}{r^{2}} \cos \theta=\frac{\mu_{0} i}{4 \pi} \int d \Omega=\mu_{0} i
\]

where we have used \(\frac{d S \cos \theta}{r^{2}}=d \Omega\) and \(\int d \Omega=4 \pi\), the total solid angle around any point.\\
The magnetic field vanishes outside the cylinder by similar argument.\\
The total energy per unit length of the cylinder is,

\[
W_{1}=\frac{1}{2 \mu_{0}} \mu_{0}^{2}\left(\frac{\lambda \omega}{2 \pi}\right)^{2} \times \pi a^{2}=\frac{\mu_{0}}{8 \pi} a^{2} \lambda^{2} \omega^{2}
\]

\(3.340 w_{E}=\frac{1}{2} \varepsilon_{0} E^{2}\), for the electric field,\\
\(w_{B}=\frac{1}{2 \mu_{0}} B^{2}\) for the magnetic field.\\
Thus,

\[
\frac{1}{2 \mu_{0}} B^{2}=\frac{1}{2} \varepsilon_{0} E^{2}
\]

when

\[
E=\frac{B}{\sqrt{\varepsilon_{0} \mu_{0}}}=3 \times 10^{8} \mathrm{~V} / \mathrm{m}
\]

3.341 The electric field at \(P\) is,

\[
E_{p}=\frac{q l}{4 \pi \varepsilon_{0}\left(a^{2}+l^{2}\right)^{3 / 2}}
\]

To get the magnetic field, note that the rotating ring constitutes a current \(i=q \omega / 2 \pi\), and the corresponding magnetic field at \(P\) is,

\[
B_{p}=\frac{\mu_{0} a^{2} i}{2\left(a^{2}+l^{2}\right)^{3 / 2}}
\]

Thus, \(\frac{w_{E}}{w_{M}}=\frac{\varepsilon_{0} \mu_{0} E^{2}}{B^{2}}=\varepsilon_{0} \mu_{0}\left(\frac{q l \times 2}{4 \pi \varepsilon_{0} \mu_{0} a^{2} i}\right)^{2}\)

\[
=\frac{1}{\varepsilon_{0} \mu_{0}}\left(\frac{l}{a^{2} \omega}\right)^{2}
\]

or, \(\quad \frac{w_{M}}{w_{E}}=\varepsilon_{0} \mu_{0} \omega^{2} a^{4} / l^{2}\)\\
\includegraphics[max width=\textwidth, alt={}, center]{9f3c5a8e-ac5d-42f4-8eb6-ea2af623f2d8-406_393_469_220_822}\\
3.342 The total energy of the magnetic field is,

\[
\begin{aligned}
& \frac{1}{2} \int(\vec{B} \cdot \vec{H}) d V=\frac{1}{2} \int \vec{B} \cdot\left(\frac{\vec{B}}{\mu_{0}}-\vec{J}\right) d V \\
& =\frac{1}{2 \mu_{0}} \int \vec{B} \cdot \vec{B} d V-\frac{1}{2} \int \vec{J} \cdot \vec{B} d V
\end{aligned}
\]

The second term can be interpreted as the energy of magnetization, and has the density

\[
-\frac{1}{2} \vec{J} \cdot \vec{B}
\]

3.343 (a) In series, the current I flows through both coils, and the total e.m.f. induced. when the current changes is,\\
or,

\[
\begin{aligned}
-2 L \frac{d I}{d t} & =-L^{\prime} \frac{d I}{d t} \\
L^{\prime} & =2 L
\end{aligned}
\]

(b) In parallel, the current flowing through either coil is, \(\frac{I}{2}\) and the e.m.f. induced is

\[
-L\left(\frac{1}{2} \frac{d I}{d t}\right)
\]

Equating this. to \(-L^{\prime} \frac{d I}{d t}\), we find \(L^{\prime}=\frac{1}{2} L\).\\
3.344 We use \(L_{1}=\mu_{0} n_{1}^{2} V, L_{2}=\mu_{0} n_{2}^{2} V\)

So,

\[
L_{12}=\mu_{0} n_{1} n_{2} V=\sqrt{L_{1} L_{2}}
\]

3.345 The interaction energy is

\[
\begin{gathered}
\frac{1}{2 \mu_{0}} \int\left|\overrightarrow{B_{1}}+\overrightarrow{B_{2}}\right|^{2} d V-\frac{1}{2 \mu_{0}} \int\left|\overrightarrow{B_{1}}\right|^{2} d V-\frac{1}{2 \mu_{0}} \int\left|\overrightarrow{B_{2}}\right|^{2} d V \\
=\frac{1}{\mu_{0}} \int \overrightarrow{B_{1}} \cdot \overrightarrow{B_{2}} d V
\end{gathered}
\]

Here, if \(\overrightarrow{B_{1}}\) is the magnetic field produced by the first of the current carrying loops, and \(\overrightarrow{B_{2}}\), that of the second one, then the magnetic field due to both the loops will be \(\overrightarrow{B_{1}}+\overrightarrow{B_{2}}\).\\
3.346 We can think of the smaller coil as constituting a magnet of dipole moment,

\[
p_{m}=\pi a^{2} I_{1}
\]

Its direction is normal to the loop and makes an angle \(\theta\) with the direction of the magnetic field, due to the bigger loop. This magnetic field is,

\[
B_{2}=\frac{\mu_{0} I_{2}}{2 b}
\]

The interaction energy has the magnitude,

\[
|W|=\frac{\mu_{0} I_{1} I_{2}}{2 b} \pi a^{2} \cos \theta
\]

Its sign depends on the sense of the currents.\\
3.347 (a) There is a radial outward conduction current. Let \(Q\) be the instantaneous charge on the inner sphere, then,

\[
j \times 4 \pi r^{2}=-\frac{d Q}{d t} \text { or, } \vec{j}=-\frac{1}{4 \pi r^{2}} \frac{d Q}{d t} \hat{r}
\]

On the other hand \(\overrightarrow{j_{d}}=\frac{\partial \vec{D}}{\partial t}=\frac{d}{d t}\left(\frac{Q}{4 \pi r^{2}} \hat{r}\right)=-\vec{j}\)\\
(b) At the given moment, \(\vec{E}=\frac{q}{4 \pi \varepsilon_{0} \varepsilon r^{2}} \hat{r}\)\\
and by Ohm's law, \(\vec{j}=\frac{\vec{E}}{\rho}=\frac{q}{4 \pi \varepsilon_{0} \varepsilon \rho r^{2}} \hat{r}\)\\
Then,

\[
\overrightarrow{j_{d}}=-\frac{q}{4 \pi \varepsilon_{0} \varepsilon \rho r^{2}} \hat{r}
\]

and

\[
\oint \overrightarrow{j_{d}} \cdot d \vec{S}=-\frac{q}{4 \pi \varepsilon_{0} \varepsilon \rho} \int \frac{d S \cos \theta}{r^{2}}=-\frac{q}{\varepsilon_{0} \varepsilon \rho}
\]

The surface integral must be \(-v e\) because \(\overrightarrow{j_{d}}\), being opposite of \(\vec{j}\), is inward.\\
3.348 Here also we see that neglecting edge effects, \(\overrightarrow{j_{d}}=-\vec{j}\). Thus Maxwell's equations reduce to,div \(\vec{B}=0\), Curl \(\vec{H}=0, \vec{B}=\mu \vec{H}\)\\
A general solution of this equation is \(\vec{B}=\) constant \(=\overrightarrow{B_{0}} \cdot \overrightarrow{B_{0}}\) can be thought of as an extraneous magnetic field. If it is zero, \(\vec{B}=0\).\\
3.349 Given \(I=I_{m} \sin \omega t\). We see that

\[
j=\frac{I_{m}}{S} \sin \omega t=-j_{d}=-\frac{\partial D}{\partial t}
\]

or, \(D=\frac{I_{m}}{\omega S} \cos \omega t\), so, \(E_{m}=\frac{I_{m}}{\varepsilon_{0} \omega S}\) is the amplitude of the electric field and is \(7 \mathrm{~V} / \mathrm{cm}\)\\
3.350 The electric field between the plates can be written as,

\[
E=\operatorname{Re} \frac{V_{m}}{d} e^{i \omega t}, \text { instead of } \frac{V_{m}}{d} \cos \omega t
\]

This gives rise to a conduction current,

\[
j_{c}=\sigma E=\operatorname{Re} \frac{\sigma}{d} V_{m} e^{i \omega t}
\]

and a displacement current,

\[
j_{d}=\frac{\partial D}{\partial t}=\operatorname{Re} \varepsilon_{0} \varepsilon i \omega \frac{V_{m}}{d} e^{i \omega t}
\]

The total current is,

\[
j_{T}=\frac{V_{m}}{d} \sqrt{\sigma^{2}+\left(\varepsilon_{0} \varepsilon \omega\right)^{2}} \cos (\omega t+\alpha)
\]

where, \(\quad \tan \alpha=\frac{\sigma}{\varepsilon_{0} \varepsilon \omega}\) on taking the real part of the resultant.\\
The corresponding magnetic field is obtained by using circulation theorem,

\[
H \cdot 2 \pi r=\pi r^{2} j_{T}
\]

or, \(H=H_{m} \cos (\omega t+\alpha)\), where, \(H_{m}=\frac{r V_{m}}{2 d} \sqrt{\sigma^{2}+\left(\varepsilon_{0} \varepsilon \omega\right)^{2}}\)\\
3.351 Inside the solenoid, there is a magnetic field,

\[
B=\mu_{0} n I_{m} \sin \omega t
\]

Since this varies in time there is an associated electric field. This is obtained by using,

\[
\oint_{C} \vec{E} \cdot d \vec{r}=-\frac{d}{d t} \int_{S} \vec{B} \cdot d \vec{S}
\]

For

\[
r<R, 2 \pi r E=-\dot{B} \cdot \pi r^{2}, \quad \text { or }, \quad E=-\frac{\dot{B} r}{2}
\]

For

\[
r>R, E=-\frac{\dot{B} R^{2}}{2 r}
\]

The associated displacement current density is,

\[
j_{d}=\varepsilon_{0} \frac{\partial E}{\partial t}=\left[\begin{array}{l}
-\varepsilon_{0} \ddot{B} r / 2 \\
-\varepsilon_{0} \ddot{B} R^{2} / 2 r .
\end{array}\right.
\]

The answer, given in the book, is dimensionally incorrect without the factor \(\varepsilon_{0}\).\\
3.352 In the non-relativistic limit.

\[
\vec{E}=\frac{q}{4 \pi \varepsilon_{0} r^{3}} \vec{r}
\]

(a) On a straight line coinciding with the charge path,

\[
\overrightarrow{j_{d}}=\varepsilon_{0} \frac{\partial \vec{E}}{\partial t}=\frac{q}{4 \pi}\left[\frac{-\vec{V}}{r^{3}}-\frac{3 \vec{r} r}{r^{4}}\right],\left(\text { using, } \frac{d \vec{r}}{d t}=-\overrightarrow{v_{.}}\right)
\]

But in this case, \(\dot{r}=-v\) and \(v \frac{\vec{r}}{r}=\overrightarrow{v,}\) so, \(j_{d}=\frac{2 q \vec{v}}{4 \pi r^{3}}\)\\
(b) In this case, \(\dot{r}=0\), as, \(\vec{r} \perp \overrightarrow{v .}\) Thus,

\[
j_{d}=-\frac{q \vec{v}}{4 \pi r^{3}}
\]

3.353 We have, \(E_{p}=\frac{q x}{4 \pi \varepsilon_{0}\left(a^{2}+x^{2}\right)^{3 / 2}}\)\\
then

\[
j_{d}=\frac{\partial D}{\partial t}=\varepsilon_{0} \frac{\partial E}{\partial t}=\frac{q v}{4 \pi\left(a^{2}+x^{2}\right)^{5 / 2}}\left(a^{2}-2 x^{2}\right)
\]

This is maximum, when \(x=x_{m}=0\), and minimum at some other value. The maximum displacement current density is

\[
\left(j_{d}\right)_{\max }=\frac{q v}{4 \pi a^{3}}
\]

To check this we calculate \(\frac{\partial j_{d}}{\partial x}\);

\[
\frac{\partial j_{d}}{\partial x}=\frac{q v}{4 \pi}\left[\left(-4 x\left(a^{2}+x^{2}\right)-5 x\left(a^{2}-2 x^{2}\right)\right]\right.
\]

This vanishes for \(x=0\) and for \(x=\sqrt{\frac{3}{2}} a\). The latter is easily shown to be a smaller local minimum (negative maximum).\\
3.354 We use Maxwell's equations in the form,

\[
\oint \vec{B} \cdot d \vec{r}=\varepsilon_{0} \mu_{0} \frac{\partial}{\partial t} \int \vec{E} \cdot d \vec{S}
\]

when the conduction current vanishes at the site.\\
We know that,

\[
\begin{aligned}
\int \vec{E} \cdot d \vec{S} & =\frac{q}{4 \pi \varepsilon_{0}} \int \frac{d \vec{S} \cdot \hat{r}}{r^{2}} \\
& =\frac{q}{4 \pi \varepsilon_{0}} \int d \Omega=\frac{q}{4 \pi \varepsilon_{0}} 2 \pi(1-\cos \theta)
\end{aligned}
\]

\includegraphics[max width=\textwidth, alt={}, center]{9f3c5a8e-ac5d-42f4-8eb6-ea2af623f2d8-409_288_561_1229_805}\\
where, \(2 \pi(1-\cos \theta)\) is the solid angle, formed by the disc like surface, at the charge.\\
Thus,

\[
\oint \vec{B} \cdot d \vec{r}=2 \pi a B=\frac{1}{2} \mu_{0} q \cdot \sin \theta \cdot \theta
\]

On the other hand, \(x=a \cot \theta\)\\
differentiating and using \(\frac{d x}{d t}=-v\),

Thus,

\[
\begin{aligned}
v & =a \operatorname{cosec}^{2} \theta \dot{\theta} \\
B & =\frac{\mu_{0} q v r \sin \theta}{4 \pi r^{3}}
\end{aligned}
\]

This can be written as, \(\vec{B}=\frac{\mu_{0} q(\vec{v} \times \vec{r})}{4 \pi r^{3}}\)\\
and \(\vec{H}=\frac{q}{4 \pi} \frac{\vec{v} \times \vec{r}}{r^{3}}\) (The sense has to be checked independently.)\\
3.355 (a) If \(\vec{B}=\vec{B}(t)\), then,

Curl \(\vec{E}=\frac{-\partial \vec{B}}{\partial t} \neq 0\).\\
So, \(\vec{E}\) cannot vanish.\\
(b) Here also, curl \(\vec{E} \neq 0\), so \(\vec{E}\) cannot be uniform.\\
(c) Suppose for instance, \(\vec{E}=\vec{a} f(t)\)\\
where \(\vec{a}\) is spatially and temporally fixed vector. Then \(-\frac{\partial \vec{B}}{d t}=\operatorname{curl} \vec{E}=0\). Generally speaking this contradicts the other equation curl \(\vec{H}=\frac{\partial \vec{D}}{\partial t} \neq 0\) for in this case the left hand side is time independent but RHS. depends on time. The only exception is when \(f(t)\) is linear function. Then a uniform field \(\vec{E}\) can be time dependent.\\
3.356 From the equation Curl \(\vec{H}-\frac{\partial \vec{D}}{\partial t}=\vec{j}\)

We get on taking divergence of both sides

\[
-\frac{\partial}{\partial t} \operatorname{div} \vec{D}=\operatorname{div} \vec{j}
\]

But \(\operatorname{div} \vec{D}=\rho\) and hence \(\operatorname{div} \vec{j}+\frac{\partial \rho}{\partial t}=0\)\\
3.357 From \(\vec{\nabla} \times \vec{E}=-\frac{\partial \vec{B}}{\partial t}\)\\
we get on taking divergence

\[
0=-\frac{\partial}{\partial t} \operatorname{div} \vec{B}
\]

This is compatible with \(\operatorname{div} \vec{B}=0\)\\
3.358 A rotating magnetic field can be represented by,\\
\(B_{x}=B_{0} \cos \omega t ; B_{y}=B_{0} \sin \omega t\) and \(B_{z}=B_{z o}\)\\
Then curl, \(\quad \vec{E}=-\frac{\partial \vec{B}}{\partial t}\).\\
So, \(\quad-(\operatorname{Curl} \vec{E})_{x}=-\omega B_{0} \sin \omega t=-\omega B_{y}\)

\[
-(\operatorname{Curl} \vec{E})_{y}=\omega B_{0} \cos \omega t=\omega B_{x} \text { and }-(\operatorname{Curl} \vec{E})_{z}=0
\]

Hence, Curl \(\vec{E}=-\vec{\omega} \times \vec{B}\),\\
where, \(\quad \vec{\omega}=\overrightarrow{e_{3}} \omega\).\\
3.359 Consider a particle with charge \(e\), moving with velocity \(\vec{v}\), in frame \(K\). It experiences a force \(\vec{F}=e \vec{v} \times \vec{B}\)

In the frame \(K^{\prime}\), moving with velocity \(\vec{v}\), relative to \(K\), the particle is at rest. This means that there must be an electric field \(\vec{E}\) in \(K^{\prime}\), so that the particle experinces a force,

Thus,

\[
\begin{aligned}
\overrightarrow{F^{\prime}}=e \overrightarrow{E^{\prime}} & =\vec{F}=e \vec{v} \times \vec{B} \\
\overrightarrow{E^{\prime}} & =\vec{v} \times \vec{B}
\end{aligned}
\]

3.360 Within the plate, there will appear a \((\vec{v} \times \vec{B})\) force, which will cause charges inside the plate to drift, until a countervailing electric field is set up. This electric field is related to \(B\), by \(E=e B\), since \(v \& B\) are mutually perpendicular, and \(E\) is perpendicular to both. The charge density \(\pm \sigma\), on the force of the plate, producing this electric field, is given by

\[
E=\frac{\sigma}{\varepsilon_{0}} \text { or } \sigma=\varepsilon_{0} \nu B=0.40 \mathrm{pC} / \mathrm{m}^{2}
\]

3.361 Choose \(\vec{\omega} \uparrow \uparrow \vec{B}\) along the \(z\)-axis, and choose \(\vec{r}\), as the cylindrical polar radius vector of a reference point (perpendicular distance from the axis). This point has the velocity,

\[
\vec{v}=\vec{\omega} \times \vec{r},
\]

and experiences a \((\vec{v} \times \vec{B})\) force, which must be counterbalanced by an electric field,

\[
\vec{E}=-(\vec{\omega} \times \vec{r}) \times \vec{B}=-(\vec{\omega} \cdot \vec{B}) \vec{r}
\]

There must appear a space charge density,

\[
\rho=\varepsilon_{0} \operatorname{div} \vec{E}=-2 \varepsilon_{0} \vec{\omega} \cdot \vec{B}=-8 \mathrm{pC} / \mathrm{m}^{3}
\]

Since the cylinder, as a whole is electrically neutral, the surface of the cylinder must acquire a positive charge of surface density,

\[
\sigma=+\frac{2 \varepsilon_{0}(\vec{\omega} \cdot \vec{B}) \pi a^{2}}{2 \pi a}=\varepsilon_{0} a \vec{\omega} \cdot \vec{B}=+\cdot 2 \mathrm{pC} / \mathrm{m}^{2}
\]

3.362 In the reference frame \(K^{\prime}\), moving with the particle,

\[
\begin{gathered}
\overrightarrow{E^{\prime}} \cong \vec{E}+\overrightarrow{v_{0}} \times \vec{B}=\frac{q \vec{r}}{4 \pi \varepsilon_{0} r^{3}} \\
\overrightarrow{B^{\prime}} \cong \vec{B}-\overrightarrow{v_{0}} \times \vec{E} / c^{2}=0
\end{gathered}
\]

Here, \(\overrightarrow{v_{0}}=\) velocity of \(K^{\prime}\), relative to the \(K\) frame, in which the particle has velocity \(\vec{v}\).\\
Clearly, \(\overrightarrow{v_{0}}=\vec{v}\). From the second equation,

\[
\vec{B} \equiv \frac{\vec{v} \times \vec{E}}{c^{2}}=\varepsilon_{0} \mu_{0} \times \frac{q}{4 \pi \varepsilon_{0}} \frac{\vec{v} \times \vec{r}}{r^{3}}=\frac{\mu_{0}}{4 \pi} \frac{q(\vec{v} \times \vec{r})}{r^{3}}
\]

3.363 Suppose, there is only electric field \(\vec{E}\), in \(K\). Then in \(K^{\prime}\), considering nonrelativistic velocity \(\vec{v}, \overrightarrow{E^{\prime}}=\vec{E}, \vec{B}=-\frac{\vec{v} \times \vec{E}}{c^{2}}\),\\
So,

\[
\overrightarrow{E^{\prime}} \cdot \overrightarrow{B^{\prime}}=0
\]

In the relativistic case,

\[
\left.\begin{array}{c}
\vec{E}_{\|}^{\prime}=\vec{E}_{\|} \\
E_{\perp}^{\prime}=\frac{\overrightarrow{E \perp}}{\sqrt{1-v^{2} / c^{2}}}
\end{array}\right\} \begin{gathered}
\vec{B}_{\|}^{\prime}=\vec{B}_{\|}=0 \\
\vec{B}_{\perp}^{\prime}=\frac{-\vec{v} \times \vec{E} / c^{2}}{\sqrt{1-v^{2} / c^{2}}}
\end{gathered}
\]

Now,

\[
\overrightarrow{E^{\prime}} \cdot \overrightarrow{B^{\prime}}=\overrightarrow{E_{\|}^{\prime}} \cdot \overrightarrow{B_{\|}^{\prime}}+\overrightarrow{E_{\perp}^{\prime}} \cdot \overrightarrow{B_{\perp}^{\prime}}=0, \text { since }
\]

\(E_{\perp}^{\prime} \cdot B_{\perp}^{\prime}=-\overrightarrow{E_{\perp}} \cdot(\vec{v} \times \vec{E}) /\left(1-v^{2} / c^{2}\right)=-\vec{E}_{\perp} \cdot\left(\vec{v} \times \vec{E}_{\perp}\right) /\left(1-\frac{v^{2}}{c^{2}}\right)=0\)\\
3.364 \(\operatorname{In} \cdot K, \vec{B}=b \frac{y \hat{i}-x \hat{j}}{x^{2}+y^{2}}, b=\) constant.

\[
\text { In } \quad K^{\prime}, \overrightarrow{E^{\prime}}=\vec{v} \times \vec{B}=b v \frac{y \hat{j}-x \hat{i}}{x^{2}+y^{2}}=b v \frac{\vec{r}}{r^{2}}
\]

The electric field is radial \((\vec{r}=x \hat{i}+y \hat{j})\).\\
3.365 In \(K, \vec{E}=a \frac{\vec{r}}{r^{2}}, \vec{r}=(x \hat{i}+y \hat{j})\)

In

\[
K^{\prime}, \overrightarrow{B^{\prime}}=-\frac{\vec{v} \times \vec{E}}{c^{2}}=\frac{a \vec{r} \times \vec{v}}{c^{2} r^{2}}
\]

The magnetic lines are circular.\\
3.366 In the non relativistic limit, we neglect \(v^{2} / c^{2}\) and write,

\[
\left.\begin{array}{l}
\overrightarrow{E_{1}^{\prime}}=\overrightarrow{E_{1}} \\
\overrightarrow{E_{\perp}} \cong \overrightarrow{E_{\perp}}+\vec{v} \times \vec{B}
\end{array}\right\} \underset{\overrightarrow{B_{1}^{\prime}}}{\overrightarrow{B_{1}^{\prime}}} \cong \overrightarrow{B_{1}}=\overrightarrow{B_{\perp}}-\vec{v} \times \vec{E} / c^{2}
\]

These two equations can be combined to give,

\[
\overrightarrow{E^{\prime}}=\vec{E}+\vec{v} \times \vec{B}, \overrightarrow{B^{\prime}}=\vec{B}-\vec{v} \times \vec{E} / c^{2}
\]

3.367 Choose \(\vec{E}\) in the direction of the \(z\)-axis, \(\vec{E}=(0,0, E)\). The frame \(K^{\prime}\) is moving with velocity \(\vec{v}=(v \sin \alpha, 0, v \cos \alpha)\), in the \(x-z\) plane. Then in the frame \(K^{\prime}\),

\[
\begin{gathered}
\vec{E}_{\|}^{\prime}=\vec{E}_{\|} B_{\|}^{\prime}=0 \\
\vec{E}_{\perp}^{\prime}=\frac{\vec{E}_{\perp}}{\sqrt{1-v^{2} / c^{2}}} \vec{B}_{\perp}^{\prime}=\frac{-\vec{v} \times \vec{E} / c^{2}}{\sqrt{1-v^{2} / c^{2}}}
\end{gathered}
\]

The vector along \(\vec{v}\) is \(\vec{e}=(\sin \alpha, 0, \cos \alpha)\) and the perpendicular vector in the \(x-z\) plane is,

\[
\vec{f}=(-\cos \alpha, 0, \sin \alpha),
\]

(a) Thus using \(\vec{E}=E \cos \alpha \vec{e}+E \sin \alpha \vec{f}\),

\[
E_{\|}^{\prime}=E \cos \alpha \text { and } E_{\|}^{\prime}=\frac{E \sin \alpha}{\sqrt{1-v^{2} / c^{2}}},
\]

So

\[
E^{\prime}=E \sqrt{\frac{1-\beta^{2} \cos ^{2} \alpha}{1-\beta^{2}}} \text { and } \tan \alpha^{\prime}=\frac{\tan \alpha}{\sqrt{1-v^{2} / c^{2}}}
\]

(b) \(B_{\|}^{\prime}=0, \vec{B}_{\perp}^{\prime}=\frac{\vec{v} \times \vec{E} / c^{2}}{\sqrt{1-v^{2} / c^{2}}}\)

\[
B^{\prime}=\frac{\beta E \sin \alpha}{c \sqrt{1-\beta^{2}}}
\]

3.368 Choose \(\vec{B}\) in the \(z\) direction, and the velocity \(\vec{v}=(v \sin \alpha, 0, v \cos \alpha)\) in the \(x-z\) plane, then in the \(K^{\prime}\) frame,

\[
\left.\begin{aligned}
& \vec{E}_{\|}^{\prime}=\vec{E}_{\|}=0 \\
& \vec{E}_{\perp}^{\prime}
\end{aligned}=\frac{\vec{v} \times \vec{B}}{\sqrt{1-v^{2} / c^{2}}} \right\rvert\, \begin{aligned}
& \vec{B}_{\|}^{\prime}=\vec{B}_{\|} \\
& \vec{B}_{\perp}^{\prime}=\frac{\vec{B}_{\perp}}{\sqrt{1-v^{2} / c^{2}}}
\end{aligned}
\]

We find similarly, \(E^{\prime}=\frac{c \beta B \sin \alpha}{\sqrt{1-\beta^{2}}}\)

\[
B^{\prime}=B \sqrt{\frac{1-\beta^{2} \cos ^{2} \alpha}{1-\beta^{2}}} \tan \alpha^{\prime}=\frac{\tan \alpha}{\sqrt{1-\beta^{2}}}
\]

3.369 (a) We see that, \(\overrightarrow{E^{\prime}} \cdot \overrightarrow{B^{\prime}}=\overrightarrow{E^{\prime}} \cdot \overrightarrow{B^{\prime}}+\overrightarrow{E_{\perp}^{\prime}} \cdot \overrightarrow{B_{\perp}^{\prime}}\)

\[
\begin{aligned}
& =\overrightarrow{E_{\|}} \cdot \overrightarrow{B_{\|}}+\frac{\left(\overrightarrow{E_{\perp}}+\vec{v} \times \vec{B}\right) \cdot\left(\overrightarrow{B_{\perp}}-\frac{\vec{v} \times \vec{E}}{c^{2}}\right)}{1-\frac{v^{2}}{c^{2}}} \\
& =\overrightarrow{E_{\|}} \cdot \overrightarrow{B_{\|}}+\frac{\overrightarrow{E_{\perp}} \cdot \overrightarrow{B_{\perp}}-(\vec{v} \times \vec{B}) \cdot(\vec{v} \times \vec{E}) / c^{2}}{1-v^{2} / c^{2}} \\
& =\overrightarrow{E_{\|}} \cdot \overrightarrow{B_{\|}}+\frac{\overrightarrow{E_{\perp}} \cdot \overrightarrow{B_{\perp}}-\left(\vec{v} \times \overrightarrow{B_{\perp}}\right) \cdot\left(\vec{v} \times \overrightarrow{E_{\perp}}\right) / c^{2}}{1-\frac{v^{2}}{c^{2}}} \\
& \rightarrow \rightarrow
\end{aligned}
\]

But,

\[
\vec{A} \times \vec{B} \cdot \vec{C} \times \vec{D}=A \cdot C B \cdot D-A \cdot D B \cdot C
\]

so,

\[
\overrightarrow{E^{\prime}} \cdot \overrightarrow{B^{\prime}}=\overrightarrow{E_{\|}} \cdot \overrightarrow{B_{\|}}+\overrightarrow{E_{\perp}} \cdot \overrightarrow{B_{\perp}} \frac{\left(1-\frac{v^{2}}{c^{2}}\right)}{1-\frac{v^{2}}{c^{2}}}=\vec{E} \cdot \vec{B}
\]

(b) \(E^{\prime 2}-c^{2} B^{\prime 2}=E_{\|}^{\prime 2}-c^{2} B_{\|}^{\prime 2}+E_{\perp}^{\prime 2}-c^{2} B_{\perp}^{\prime 2}\)

\[
\begin{aligned}
& =E_{\|}^{2}-c^{2} B_{\|}^{2}+\frac{1}{1-\frac{v^{2}}{c^{2}}}\left[\left(\vec{E}_{\perp}+\vec{v} \times \vec{B}\right)^{2}-c^{2}\left(\vec{B}_{\perp}-\frac{\vec{v} \times \vec{E}}{c^{2}}\right)^{2}\right] \\
& =E_{\|}^{2}-c^{2} B_{\|}^{2}+\frac{1}{1-\frac{v^{2}}{c^{2}}}\left[E_{\perp}^{2}-c^{2} B_{\perp}^{2}+\left(\vec{v} \times \vec{B}_{\perp}\right)^{2}-\frac{1}{c^{2}}\left(\vec{v} \times E_{1}\right)^{2}\right] \\
& =E_{\|}^{2}-c^{2} B_{\|}^{2}+\frac{1}{1-\frac{v^{2}}{c^{2}}}\left[E_{\perp}^{2}-c^{2} B_{\perp}^{2}\right]\left(1-\frac{v^{2}}{c^{2}}\right)=E^{2}-c^{2} B^{2},
\end{aligned}
\]

since, \(\left(\vec{v} \times \overrightarrow{A_{\perp}}\right)^{2}=v^{2} A_{\perp}^{2}\)\\
3.370 In this case, \(\vec{E} \cdot \vec{B}=0\), as the fields are mutually perpendicular. Also,

\[
E^{2}-c^{2} B^{2}=-20 \times 10^{8}\left(\frac{\mathrm{~V}}{\mathrm{~m}}\right)^{2} \text { is - ve. }
\]

Thus, we can find a frame, in which \(E^{\prime}=0\), and\\
\(B^{\prime}=\frac{1}{c} \sqrt{c^{2} B^{2}-E^{2}}=B \sqrt{1-\frac{E^{2}}{c^{2} B^{2}}}=0.20 \sqrt{1-\left(\frac{4 \times 10^{4}}{3 \times 10^{8} \times 2 \times 10^{-4}}\right)^{2}}=0.15 \mathrm{mT}\)\\
3.371 Suppose the charge \(q\) moves in the positive direction of the \(x\)-axis of the frame \(K\). Let us go over to the moving frame \(K^{\prime}\), at whose origin the charge is at rest. We take the \(x\) and \(x^{\prime}\) axes of the two frames to be coincident, and the \(y \& y^{\prime}\) axes, to be parallel.\\
In the \(K^{\prime}\) frame, \(\vec{E}=\frac{1}{4 \pi \varepsilon_{0}} \frac{q \vec{r}^{\prime}}{r^{\prime 3}}\),\\
and this has the following components,

\[
E_{x}^{\prime}=\frac{1}{4 \pi \varepsilon_{0}} \frac{q x^{\prime}}{r^{\prime 3}}, E_{y}^{\prime}=\frac{1}{4 \pi \varepsilon_{0}} \frac{q y^{\prime}}{r^{\prime 3}} .
\]

Now let us go back to the frame \(K\). At the moment, when the origins of the two frames coincide, we take \(t=0\). Then,

\[
x=r \cos \theta=x^{\prime} \sqrt{1-\frac{v^{2}}{c^{2}}}, y=r \sin \theta=y^{\prime}
\]

Also,

\[
E_{x}=E_{x}^{\prime}, E_{y}=E_{y}^{\prime} / \sqrt{1-v^{2} / c^{2}}
\]

From these equations, \(r^{\prime 2}=\frac{r^{2}\left(1-\beta^{2} \sin ^{2} \theta\right)}{1-\beta^{2}}\)

\[
\begin{gathered}
\vec{E}=\frac{q}{4 \pi \varepsilon_{0}} \frac{1}{r^{3}\left(1-\beta^{2} \sin ^{2} \theta\right)^{3 / 2}}\left[\left(1-\beta^{2}\right)^{3 / 2}\left(x^{\prime} \hat{i}+\frac{y^{\prime}}{\sqrt{1-\beta^{2}}} \hat{j}\right)\right] \\
=\frac{q \vec{r}\left(1-\beta^{2}\right)}{4 \pi \varepsilon_{0} r^{3}\left(1-\beta^{2} \sin ^{2} \theta\right)^{3 / 2}}
\end{gathered}
\]

\subsection*{3.7 MOTION OF CHARGED PARTICLES IN ELECTRIC AND MAGNETIC FIELDS}
3.372 Let the electron leave the negative plate of the capacitor at time \(t=0\)

As,

\[
E_{x}=-\frac{d \varphi}{d x}, E=\frac{\varphi}{l}=\frac{a t}{l},
\]

and, therefore, the acceleration of the electron,

\[
w=\frac{e E}{m}=\frac{e a t}{m l} \quad \text { or, } \quad \frac{d v}{d t}=\frac{e a t}{m l}
\]

or,


\begin{equation*}
\int_{0}^{v} d v=\frac{e a}{m l} \int_{0}^{t} t d t, \quad \text { or, } \quad v=\frac{1}{2} \frac{e a}{m l} t^{2} \tag{1}
\end{equation*}


But, from \(s=\int v d t\),

\[
l=\frac{1}{2} \frac{e a}{m l} \int_{0}^{t} t^{2} d t=\frac{e a t^{3}}{6 m l} \text { or, } t=\left(\frac{6 m l^{2}}{e a}\right)^{\frac{1}{3}}
\]

Putting the value of \(t\) in (1),

\[
v=\frac{1}{2} \frac{e a}{m l}\left(\frac{6 m l^{2}}{e a}\right)^{\frac{2}{3}}=\left(\frac{9}{2} \frac{a l e}{m}\right)^{\frac{1}{3}}=16 \mathrm{~km} / \mathrm{s}
\]

3.373 The electric field inside the capacitor varies with time as,

\[
E=a t .
\]

Hence, electric force on the proton,

\[
F=e a t
\]

and subsequently, acceleration of the proton,

\[
w=\frac{e a t}{m}
\]

Now, if \(t\) is the time elapsed during the motion of the proton between the plates, then \(t=\frac{l}{v_{\|}}\), as no acceleration is effective in this direction. (Here \(v_{\|}\)is velocity along the length of the plate.)\\
From kinematics, \(\frac{d v_{\perp}}{d t}=w\)\\
so,

\[
\int_{0}^{v_{1}} d v_{\perp}=\int_{0}^{t} w d t
\]

(as initially, the component of velocity in the direction, \(\perp\) to plates, was zero.)\\
or

Now,

\[
\begin{gathered}
v_{\perp}=\int_{0}^{t} \frac{e a}{m} \frac{t^{2}}{2 m}=\frac{e a}{2 m} \frac{l^{2}}{v_{\|}^{2}} \\
\tan \alpha=\frac{v_{\perp}}{v_{\|}}=\frac{e a l^{2}}{2 m v_{\|}^{3}} \\
=\frac{e a l^{2}}{2 m\left(\frac{2 e V}{m}\right)^{\frac{3}{2}}}, \text { as } v_{\|}=\left(\frac{2 e V}{m}\right)^{\frac{1}{2}}, \text { from energy conservation. } \\
=\frac{a l^{2}}{4} \sqrt{\frac{m}{2 e V^{3}}}
\end{gathered}
\]

3.374 The equation of motion is,

\[
\frac{d v}{d t}=v \frac{d v}{d x}=\frac{q}{m}\left(E_{0}-a x\right)
\]

Integrating

\[
\frac{1}{2} v^{2}-\frac{q}{m}\left(E_{0} x-\frac{1}{2} a x^{2}\right)=\text { constant. }
\]

But initially \(\nu=0\) when \(x=0\), so "constant" \(=0\)\\
Thus,

\[
v^{2}=\frac{2 q}{m}\left(E_{0} x-\frac{1}{2} a x^{2}\right)
\]

Thus, \(v=0\), again for \(x=x_{m}=\frac{2 E_{0}}{a}\)\\
The corresponding acceleration is,

\[
\left(\frac{d v}{d t}\right)_{x_{m}}=\frac{q}{m}\left(E_{0}-2 E_{0}\right)=-\frac{q E_{0}}{m}
\]

3.375 From the law of relativistic conservation of energy

\[
\frac{m_{0} c^{2}}{\sqrt{1-\left(v^{2} / c^{2}\right)}}-e E x=m_{0} c^{2}
\]

as the electron is at rest ( \(v=0\) for \(x=0\) ) initially.\\
Thus clearly

\[
T=e E x .
\]

On the other hand, \(\sqrt{1-\left(v^{2} / c^{2}\right)}=\frac{m_{0} c^{2}}{m_{0} c^{2}+e E x}\)\\
or,

\[
\frac{v}{c}=\frac{\sqrt{\left(m_{0} c^{2}+e E x\right)^{2}-m_{0}^{2} c^{4}}}{m_{0} c^{2}+e E x}
\]

or,

\[
c t=\int c d t=\int \frac{\left(m_{0} c^{2}+e E x\right) d x}{\sqrt{\left(m_{0} c^{2}+e E x\right)^{2}-m_{0}^{2} c^{4}}}
\]

\[
=\frac{1}{2 e E} \int \frac{d y}{\sqrt{y-m_{0}^{2} c^{4}}}=\frac{1}{e E} \sqrt{\left(m_{0} c^{2}+e E x\right)^{2}-m_{0}^{2} c^{4}}+\text { constant }
\]

The "constant" \(=0\), at \(t=0\), for \(x=0\),\\
So,

\[
c t=\frac{1}{e E} \sqrt{\left(m_{0} c^{2}+e E x\right)^{2}-m_{0}^{2} c^{4}}
\]

Finally, using \(T=e E x\),

\[
c e E t_{0}=\sqrt{T\left(T+2 m_{0} c^{2}\right)} \text { or, } t_{0}=\frac{\sqrt{T\left(T+2 m_{0} c^{2}\right)}}{e E c}
\]

3.376 As before, \(T=e E x\)

Now in linear motion,

\[
\begin{gathered}
\frac{d}{d t} \frac{m_{0} v}{\sqrt{1-v^{2} / c^{2}}}=\frac{m_{0} w}{\sqrt{1-v^{2} / c^{2}}}+\frac{m_{0} w}{\left(1-v^{2} / c^{2}\right)^{3 / 2}} \frac{v}{c^{2}} w \\
=\frac{m_{0}}{\left(1-v^{2} / c^{2}\right)^{3 / 2}} w=\frac{\left(T+m_{0} c^{2}\right)^{3}}{m_{0}^{2} c^{6}} w=e E,
\end{gathered}
\]

So, \(\quad w=\frac{e E m_{0}^{2} c^{6}}{\left(T+m_{0} c^{2}\right)^{3}}=\frac{e E}{m_{0}}\left(1+\frac{T}{m_{0} c^{2}}\right)^{-3}\)\\
3.377 The equations are,

\[
\frac{d}{d t}\left(\frac{m_{0} v_{x}}{\sqrt{1-\left(v^{2} / c^{2}\right)}}\right)=0 \quad \text { and } \quad \frac{d}{d t}\left(\frac{m_{0} v_{y}}{\sqrt{1-v^{2} / c^{2}}}\right)=e E
\]

Hence,

\[
\frac{v_{x}}{\sqrt{1-v^{2} / c^{2}}}=\text { constant }=\frac{v_{0}}{\sqrt{1-\left(v_{0}^{2} / c^{2}\right)}}
\]

Also, by energy conservation,

\[
\frac{m_{0} c^{2}}{\sqrt{1-\left(v^{2} / c^{2}\right)}}=\frac{m_{0} c^{2}}{\sqrt{1-\left(v_{0}^{2} / c^{2}\right)}}+e E y
\]

Dividing

\[
v_{x}=\frac{v_{0} \varepsilon_{0}}{\varepsilon_{0}+e E y}, \quad \varepsilon_{0}=\frac{m_{0} c^{2}}{\sqrt{1-\left(v_{0}^{2} / c^{2}\right)}}
\]

Also,

\[
\frac{m_{0}}{\sqrt{1-\left(v^{2} / c^{2}\right)}}=\frac{\varepsilon_{0}+e E y}{c^{2}}
\]

Thus,

\[
\left(\varepsilon_{0}+e E y\right) v_{y}=c^{2} e E t+\text { constant } .
\]

"constant" \(=0\) as \(v_{y}=0\) at \(t=0\).\\
Integrating again,

\[
\varepsilon_{0} y+\frac{1}{2} e E y^{2}=\frac{1}{2} c^{2} E t^{2}+\text { constant. }
\]

"constant" \(=0\), as \(y=0\), at \(t=0\).\\
Thus, \((c e E t)^{2}=(e y E)^{2}+2 \varepsilon_{0} e E y+\varepsilon_{0}^{2}-\varepsilon_{0}^{2}\)\\
or,

\[
c e E t=\sqrt{\left(\varepsilon_{0}+e E y\right)^{2}-\varepsilon_{0}^{2}}
\]

or,

\[
\varepsilon_{0}+e E y=\sqrt{\varepsilon_{0}^{2}+c^{2} e^{2} E^{2} t^{2}}
\]

Hence,

\[
v_{x}=\frac{v_{0} \varepsilon_{0}}{\sqrt{\varepsilon_{0}^{2}+c^{2} e^{2} E^{2} t^{2}}} \text { also, } v_{y}=\frac{c^{2} e E t}{\sqrt{\varepsilon_{0}^{2}+c^{2} e^{2} E^{2} t^{2}}}
\]

and

\[
\tan \theta=\frac{v_{y}}{v_{x}}=\frac{e E t}{m_{0} v_{0}} \sqrt{1-\left(v_{0}^{2} / c^{2}\right)} .
\]

3.378 From the figure,

\[
\sin \alpha=\frac{d}{R}=\frac{d q B}{m v},
\]

As radius of the \(\operatorname{arc} R=\frac{m v}{q B}\), where \(v\) is the velocity of the particle, when it enteres into the field. From initial condition of the problem,

\[
q V=\frac{1}{2} m v^{2} \quad \text { or, } \quad v=\sqrt{\frac{2 q V}{m}}
\]

Hence, \(\sin \alpha=\frac{d q B}{m \sqrt{\frac{2 q V}{m}}}=d B \sqrt{\frac{q}{2 m V}}\)\\
and \(\alpha=\sin ^{-1}\left(d B \sqrt{\frac{q}{2 m V}}\right)=30^{\circ}\), on putting the values.\\
3.379 (a) For motion along a circle, the magnetic force acted on the particle, will provide the centripetal force, necessary for its circular motion.\\
i.e.

\[
\frac{m v^{2}}{R}=e v B \quad \text { or, } \quad v=\frac{e B R}{m}
\]

and the period of revolution, \(T=\frac{2 \pi}{\omega}=\frac{2 \pi R}{v}=\frac{2 \pi m}{e B}\)\\
(b) Generally, \(\frac{d \vec{p}}{d t}=\vec{F}\)

But, \(\frac{d \vec{p}}{d t}=\frac{d}{d t} \frac{m_{0} \vec{v}}{\sqrt{1-\left(v^{2} / c^{2}\right)}}=\frac{m_{0} \dot{\vec{v}}}{\sqrt{1-\left(v^{2} / c^{2}\right)}}+\frac{m_{0}}{\left(1-\left(v^{2} / c^{2}\right)\right)^{3 / 2}} \frac{\vec{v}(\vec{v} \cdot \dot{\vec{v}})}{c^{2}}\)\\
For transverse motion, \(\vec{v} \cdot \dot{\vec{v}}=0\) so,

\[
\frac{d \vec{p}}{d t}=\frac{m_{0} \dot{\vec{v}}}{\sqrt{1-\left(v^{2} / c^{2}\right)}}=\frac{m_{0}}{\sqrt{1-\left(v^{2} / c^{2}\right)}} \frac{v^{2}}{r}, \text { here. }
\]

Thus, \(\quad \frac{m_{0} v^{2}}{r \sqrt{1-\left(v^{2} / c^{2}\right)}}=B\) ev or, \(\frac{v / c}{\sqrt{1-\left(v^{2} / c^{2}\right)}}=\frac{B e r}{m_{0} c}\)\\
or, \(\quad \frac{v}{c}=\frac{B e r}{\sqrt{B^{2} e^{2} r^{2}+m_{0}^{2} c^{2}}}\)\\
Finally,\(\quad T=\frac{2 \pi r}{v}=\frac{2 \pi m_{0}}{e B \sqrt{1-v^{2} / c^{2}}}=\frac{2 \pi}{c B e} \sqrt{B^{2} e^{2} r^{2}+m_{0}^{2} c^{2}}\)\\
3.380 (a) As before, \(p=B q r\).\\
(b) \(T=\sqrt{c^{2} p^{2}+m_{0}^{2} c^{4}}=\sqrt{c^{2} B^{2} q^{2} r^{2}+m_{0}^{2} c^{4}}\)\\
(c) \(w=\frac{v^{2}}{r}=\frac{c^{2}}{r\left[1+\left(m_{0} c / B q r\right)^{2}\right]}\)\\
using the result for \(v\) from the previous problem.\\
3.381 From (3.279),

\[
T=\frac{2 \pi \varepsilon}{c^{2} e B} \text { (relativistic) }, \quad T_{0}=\frac{2 \pi m_{0} c^{2}}{c^{2} e B} \text { (nonrelativistic), }
\]

Here,

\[
m_{0} c^{2} / \sqrt{1-v^{2} / c^{2}}=E
\]

Thus,

\[
\delta T=\frac{2 \pi T}{c^{2} e B},(T=K . E .)
\]

Now,

\[
\frac{\delta T}{T_{0}}=\eta=\frac{T}{m_{0} c^{2}}, \quad \text { so, } \quad T=\eta m_{0} c^{2}
\]

\(3.382 T=e V=\frac{1}{2} m v^{2}\)\\
(The given potential difference is not large enough to cause significant deviations from the nonrelativistic formula).

Thus,

\[
v=\sqrt{\frac{2 e V}{m}}
\]

So,

\[
v_{\|}=\sqrt{\frac{2 e V}{m}} \cos \alpha, v_{\perp}=\sqrt{\frac{2 e V}{m}} \sin \alpha
\]

Now,

\[
\frac{m v_{\perp}^{2}}{r}=B e v_{\perp} \quad \text { or, } \quad r=\frac{m v_{\perp}}{B e}
\]

and

\[
T=\frac{2 \pi r}{v_{\perp}}=\frac{2 \pi m}{B e}
\]

Pitch

\[
p=v_{\|} T=\frac{2 \pi m}{B e} \sqrt{\frac{2 e V}{m}}, \cos \alpha=2 \pi \sqrt{\frac{2 m V}{e B^{2}}} \cos \alpha
\]

3.383 The charged particles will traverse a helical trajectory and will be focussed on the axis after traversing a number of turms. Thus

\[
\frac{l}{v_{0}} \approx n \cdot \frac{2 \pi m}{q B_{1}}=(n+1) \frac{2 \pi m}{q B_{2}}
\]

So,

\[
\frac{n}{B_{1}}=\frac{n+1}{B_{2}}=\frac{1}{B_{2}-B_{1}}
\]

Hence,

\[
\frac{l}{v_{0}}=\frac{2 \pi m}{q\left(B_{2}-B_{1}\right)}
\]

or,

\[
\frac{l^{2}}{2 q V / m}=\frac{(2 \pi)^{2}}{\left(B_{2}-B_{1}\right)^{2}} \times \frac{1}{(q / m)^{2}}
\]

or,

\[
\frac{q}{m}=\frac{8 \pi^{2} V}{l^{2}\left(B_{2}-B_{1}\right)^{2}}
\]

3.384 Let us take the point \(A\) as the origin \(O\) and the axis of the solenoid as \(z\)-axis. At an arbitrary moment of time let us resolve the velocity of electron into its two rectangular components, \(\overrightarrow{v_{\|}}\)along the axis and \(\overrightarrow{v_{\perp}}\) to the axis of solenoid. We know the magnetic force does no work, so the kinetic energy as well as the speed of the electron \(\left|\vec{v}_{\perp}\right|\) will remain constant in the \(x-y\) plane. Thus \(\vec{v}_{\perp}\) can change only its direction as shown in the Fig.. \(\overrightarrow{v_{\|}}\)will remain constant as it is parallel to \(\vec{B}\).\\
Thus at \(t=t\)

\[
\begin{aligned}
& v_{x}=v_{\perp} \cos \omega t=v \sin \alpha \cos \omega t \\
& v_{y}=v_{\perp} \sin \omega t=v \sin \alpha \sin \omega t
\end{aligned}
\]

and

\[
v_{z}=v \cos \alpha, \text { where } \omega=\frac{e B}{m}
\]

As at \(t=0\), we have \(x=y=z=0\), so the motion law of the electron is.

\[
\left.\begin{array}{l}
z=v \cos \alpha t \\
x=\frac{v \sin \alpha}{\omega} \sin \omega t \\
y=\frac{v \sin \alpha}{\omega}(\cos \omega t-1)
\end{array}\right]
\]

(The equation of the helix)\\
On the screen,

\[
z=l, \text { so } t=\frac{l}{v \cos \alpha}
\]

Then,

\[
\begin{gathered}
r^{2}=x^{2}+y^{2}=\frac{2 v^{2} \sin ^{2} \alpha}{\omega^{2}}\left(1-\cos \frac{\omega l}{v \cos \alpha}\right) \\
r=\frac{2 v \sin \alpha}{\omega} \cdot\left|\sin \frac{\omega l}{2 v \cos \alpha}\right|=2 \frac{m v}{e B} \sin \alpha\left|\sin \frac{l e B}{2 m v \cos \alpha}\right|
\end{gathered}
\]

3.385 Choose the wire along the \(z\)-axis, and the initial direction of the electron, along the \(x\)-axis. Then the magnetic field in the \(x-z\) plane is along the \(y\)-axis and outside the wire it is,

\[
B=B_{y}=\frac{\mu_{0} I}{2 \pi x},\left(B_{x}=B_{z}=0, \text { if } y=0\right)
\]

The motion must be confined to the \(x-z\) plane. Then the equations of motion are,

\[
\begin{aligned}
& \frac{d}{d t} m v_{x}=-e v_{z} B_{y} \\
& \frac{d\left(m v_{z}\right)}{d t}=+e v_{x} B_{y}
\end{aligned}
\]

Multiplying the first equation by \(v_{x}\) and the second by \(v_{z}\) and then adding,

\[
v_{x} \frac{d v_{x}}{d t}+v_{z} \frac{d v_{z}}{d t}=0
\]

or,

\[
v_{x}^{2}+v_{z}^{2}=v_{0}^{2}, \text { say, or, } v_{z}=\sqrt{v_{0}^{2}-v_{x}^{2}}
\]

Then,

\[
v_{x} \frac{d v_{x}}{d x}=-\frac{e}{m} \sqrt{v_{0}^{2}-v_{x}^{2}} \frac{\mu_{0} I}{2 \pi x}
\]

or,

\[
\begin{aligned}
& -\frac{v_{x} d v_{x}}{\sqrt{v_{0}^{2}-v_{x}^{2}}}=\frac{\mu_{0} I e}{2 \pi m} \frac{d x}{x} \\
& \sqrt{v_{0}^{2}-v_{x}^{2}}=\frac{\mu_{0} I e}{2 \pi m} \ln \frac{x}{a}
\end{aligned}
\]

on using, \(\nu_{x}=\nu_{0}\), if \(x=a\) (i.e. initially).\\
Now,

\[
v_{x}=0, \text { when } x=x_{m},
\]

so,

\[
x_{m}=a e^{v_{0} / b}, \text { where } b=\frac{\mu_{0} I e}{2 \pi m}
\]

3.386 Inside the capacitor, the electric field follows a \(\frac{1}{r}\) law, and so the potential can be written as

\[
\varphi=\frac{V \ln r / a}{\ln b / a}, E=\frac{-V}{\ln b / a} \frac{1}{r}
\]

Here \(r\) is the distance from the axis of the capacitor.\\
Also,

\[
\frac{m v^{2}}{r}=\frac{q V}{\ln b / a r} \frac{1}{r} \text { or } m v^{2}=\frac{q V}{\ln b / a}
\]

On the other hand,

\[
m v=q B r \text { in the magnetic field. }
\]

Thus,

\[
v=\frac{V}{B r \ln b / a} \text { and } \frac{q}{m}=\frac{v}{B r}=\frac{V}{B^{2} r^{2} \ln (b / a)}
\]

3.387 The equations of motion are,

\[
m \frac{d v_{x}}{d t}=-q B v_{z}, m \frac{d v_{y}}{d t}=q E \text { and } m \frac{d v_{z}}{d t}=q v_{x} B
\]

These equations can be solved easily.\\
First,

\[
v_{y}=\frac{q E}{m} t, y=\frac{q E}{2 m} t^{\prime 2}
\]

Then,

\[
v_{x}^{2}+v_{z}^{2}=\text { constant }=v_{0}^{2} \text { as before }
\]

In fact, \(v_{x}=v_{0} \cos \omega t\) and \(v_{z}=v_{0} \sin \omega t\) as one can check.\\
Integrating again and using \(x=z=0\), at \(t=0\)

\[
x=\frac{v_{0}}{\omega} \sin \omega t, z=\frac{v_{0}}{\omega}(1-\cos \omega t)
\]

Thus,

\[
x=z=0 \text { for } t=t_{n}=n \frac{2 \pi}{\omega}
\]

At that instant, \(y_{n}=\frac{q E}{2 m} \times \frac{2 \pi}{q B / m} \times n^{2} \times \frac{2 \pi}{q B / m}=\frac{2 \pi^{2} m E n^{2}}{q B^{2}}\)

Also,

\[
\begin{aligned}
& \tan \alpha_{n}=\frac{v_{x}}{v_{y}},\left(v_{z}=0 \text { at this moment }\right) \\
= & \frac{m v_{0}}{q E t_{n}}=\frac{m v_{0}}{q E} \times \frac{q B}{m} \times \frac{1}{2 \pi n}=\frac{B v_{0}}{2 \pi E n} .
\end{aligned}
\]

3.388 The equation of the trajectory is,

\[
x=\frac{v_{0}}{\omega} \sin \omega t, z=\frac{v_{0}}{\omega}(1-\cos \omega t), y=\frac{q E}{2 m} t^{2} \text { as before see (3.384). }
\]

Now on the screen \(x=l\), so

\[
\sin \omega t=\frac{\omega l}{v_{0}} \quad \text { or, } \quad \omega t=\sin ^{-1} \frac{\omega l}{v_{0}}
\]

At that moment,\\
so,

\[
\frac{\omega l}{v_{0}}=\sin \sqrt{\frac{2 m \omega^{2} y}{q E}}=\sin \sqrt{\frac{2 q B^{2} y}{E m}}
\]

and

\[
\begin{gathered}
z=\frac{v_{0}}{\omega} 2 \sin ^{2} \frac{\omega t}{2}=l \tan \frac{\omega t}{2} \\
=l \tan \frac{1}{2}\left[\sin ^{-1} \frac{\omega l}{v_{0}}\right]=l \tan \sqrt{\frac{q B^{2} y}{2 m E}}
\end{gathered}
\]

For small

\[
z, \frac{q B^{2} y}{2 m E}=\left(\tan ^{-1} \frac{z}{l}\right)^{2} \approx \frac{z^{2}}{l^{2}}
\]

or,

\[
y=\frac{2 m E}{q B^{2} l^{2}} \cdot z^{2} \text { is a parabola. }
\]

3.389 In crossed field,

\[
e E=e v B, \text { so } v=\frac{E}{B}
\]

Then, \(F=\) force exerted on the plate \(=\frac{I}{e} \times m \frac{E}{B}=\frac{m I E}{e B}\)\\
3.390 When the electric field is switched off, the path followed by the particle will be helical. and pitch, \(\Delta l=v_{\|} T\), (where \(v_{\|}\)is the velocity of the particle, parallel to \(\vec{B}\), and \(T\), the time period of revolution.)

\[
\begin{array}{r}
=v \cos (90-\varphi) T=v \sin \varphi T \\
=v \sin \varphi \frac{2 \pi m}{q B}\left(\text { as } T=\frac{2 \pi}{q B}\right) \tag{1}
\end{array}
\]

Now, when both the fields were present, \(q E=q v B \sin (90-\varphi)\), as no net force was effective on the system.\\
or,


\begin{equation*}
v=\frac{E}{B \cos \varphi} \tag{2}
\end{equation*}


From (1) and (2), \(\quad \Delta l=\frac{E}{B} \frac{2 \pi m}{q B}\) tan \(\varphi=6 \mathrm{~cm}\).\\
3.391 When there is no deviation,


\begin{equation*}
-q \vec{E}=q(\vec{v} \times \vec{B}) \tag{1}
\end{equation*}


or, in scalar from, \(E=v B\) (as \(\vec{v} \perp \vec{B}\) ) or, \(v=\frac{E}{B}\)

Now, when the magnetic field is switched on, let the deviation in the field be \(x\). Then,

\[
x=\frac{l}{2}\left(\frac{q v B}{m}\right) t^{2}
\]

where \(t\) is the time required to pass through this region.\\
also,

\[
t=\frac{a}{v}
\]

Thus


\begin{equation*}
x=\frac{1}{2}\left(\frac{q v B}{m}\right)\left(\frac{a}{v}\right)^{2}=\frac{1}{2} \frac{q}{m} \frac{a^{2} B^{2}}{E} \tag{2}
\end{equation*}


For the region where the field is absent, velocity in upward direction


\begin{equation*}
=\left(\frac{q v B}{m}\right) t=\frac{q}{m} a B \tag{3}
\end{equation*}


Now,

\[
\Delta x-x=\frac{q a B}{m} t^{\prime}
\]


\begin{equation*}
=\frac{q}{m} \frac{a B^{2} b}{E} \text { when } t^{\prime}=\frac{b}{v}=\frac{b B}{E} \tag{4}
\end{equation*}


From (2) and (4),\\
or,

\[
\begin{gathered}
\Delta x-\frac{1}{2} \frac{q}{m} \frac{a^{2} B^{2}}{E}=\frac{q}{m} \frac{a B^{2} b}{E} \\
\frac{q}{m}=\frac{2 E \Delta x}{a B^{2}(a+2 b)}
\end{gathered}
\]

3.392 (a) The equation of motion is,

Now,

\[
\begin{aligned}
m \frac{d^{2} \vec{r}}{d t^{2}}=q(\vec{E}+\vec{v} \times \vec{B}) \\
\vec{v} \times \vec{B}=\left|\begin{array}{ccc}
\vec{i} & \vec{j} & \vec{k} \\
\dot{x} & \dot{y} & \dot{z} \\
0 & 0 & B
\end{array}\right|=\overrightarrow{i B} \dot{y}-\overrightarrow{j B} \dot{x}
\end{aligned}
\]

So, the equation becomes,

\[
\frac{d v_{x}}{d t}=\frac{q B v_{y}}{m}, \frac{d v_{y}}{d t}=\frac{q E}{m}-\frac{q B}{m} v_{x}, \text { and } \frac{d v_{z}}{d t}=0
\]

Here, \(v_{x}=\dot{x}, v_{y}=\dot{y}, v_{z}=\dot{z}\). The last equation is easy to integrate;

\[
v_{z}=\text { constant }=0
\]

since \(\boldsymbol{v}_{\boldsymbol{z}}\) is zero initially. Thus integrating again,

\[
z=\text { constant }=0,
\]

and motion is confined to the \(x-y\) plane. We now multiply the second equation by \(i\) and add to the first equation.

\[
\xi=v_{x}+i v_{y}
\]

we get the equation,

\[
\frac{d \xi}{d t}=i \omega \frac{E}{B}-i \omega \xi, \omega=\frac{q B}{m} .
\]

This equation after being multiplied by \(e^{i \omega t}\) can be rewritten as,

\[
\frac{d}{d t}\left(\xi e^{i \omega t}\right)=i \omega e^{i \omega t} \frac{E}{B}
\]

and integrated at once to give,

\[
\xi=\frac{E}{B}+C e^{-i \omega t-i \alpha},
\]

where \(C\) and \(\alpha\) are two real constants. Taking real and imaginary parts.

\[
v_{x}=\frac{E}{B}+C \cos (\omega t+\alpha) \text { and } v_{y}=-C \sin (\omega t+\alpha)
\]

Since \(v_{y}=0\), when \(t=0\), we can take \(\alpha=0\), then \(v_{x}=0\) at \(t=0\) gives, \(C=-\frac{E}{B}\) and we get,

\[
v_{x}=\frac{E}{B}(1-\cos \omega t) \text { and } v_{y}=\frac{E}{B} \sin \omega t
\]

Integrating again and using \(x=y=0\), at \(t=0\), we get

\[
x(t)=\frac{E}{B}\left(t-\frac{\sin \omega t}{\omega}\right), y(t)=\frac{E}{\omega B}(1-\cos \omega t)
\]

This is the equation of a cycloid.\\
(b) The velocity is zero, when \(\omega t=2 n \pi\). We see that

\[
v^{2}=v_{x}^{2}+v_{y}^{2}=\left(\frac{E}{B}\right)^{2}(2-2 \cos \omega t)
\]

or,

\[
v=\frac{d s}{d t}=\frac{2 E}{B}\left|\sin \frac{\omega t}{2}\right|
\]

The quantity inside the modulus is positive for \(0<\omega t<2 \pi\). Thus we can drop the modulus and write for the distance traversed between two successive zeroes of velocity.

\[
S=\frac{4 E}{\omega B}\left(1-\cos \frac{\omega t}{2}\right)
\]

Putting

\[
\begin{aligned}
\omega t & =2 \pi, \text { we get } \\
S & =\frac{8 E}{\omega B}=\frac{8 m E}{q B^{2}}
\end{aligned}
\]

(c) The drift velocity is in the \(x\)-direction and has the magnitude,

\[
\left\langle v_{x}\right\rangle=<\frac{E}{B}(1-\cos \omega t)>=\frac{E}{B}
\]

3.393 When a current \(I\) flows along the axis, a magnetic field \(B_{\varphi}=\frac{\mu_{0} I}{2 \pi \rho}\) is set up where \(\rho^{2}=x^{2}+y^{2}\). In terms of components,

\[
B_{x}=-\frac{\mu_{0} I y}{2 \pi \rho^{2}}, B_{y}=\frac{\mu_{0} I x}{2 \pi \rho^{2}} \text { and } B_{z}=0
\]

Suppose a p.d. \(V\) is set up between the inner cathode and the outer anode. This means a potential function of the form

\[
\varphi=V \frac{\ln \rho / b}{\ln a / b}, a>\rho>b
\]

as one can check by solving Laplace equation.\\
The electric field corresponding to this is,\\
\(E_{x}=-\frac{V x}{\rho^{2} \ln a / b}, E_{y}=-\frac{V y}{\rho^{2} \ln a / b}, E_{z}=0\).\\
The equations of motion are,

\[
\begin{aligned}
& \frac{d}{d t} m v_{x}=+\frac{|e| V z}{\rho^{2} \ln a / b}+\frac{|e| \mu_{0} I}{2 \pi \rho^{2}} x \dot{z} \\
& \frac{d}{d t} m v_{y}=+\frac{|e| V y}{\rho^{2} \ln a / b}+\frac{|e| \mu_{0} I}{2 \pi \rho^{2}} y \dot{z} \\
& \text { and } \quad \frac{d}{d t} m v_{z}=-|e| \frac{\mu_{0} I}{2 \pi \rho^{2}}(x \dot{x}+y \dot{y})=-|e| \frac{\mu_{0} I}{2 \pi} \frac{d}{d t} \ln \rho
\end{aligned}
\]

\includegraphics[max width=\textwidth, alt={}, center]{9f3c5a8e-ac5d-42f4-8eb6-ea2af623f2d8-425_525_381_1077_931}\\
\((-|e|)\) is the charge on the electron.\\
Integrating the last equation,

\[
m v_{z}=-|e| \frac{\mu_{0} I}{2 \pi} \ln \rho / a=m \dot{z}
\]

since \(\nu_{z}=0\) where \(\rho=a\). We now substitute this \(\dot{z}\) in the other two equations to get

\[
\begin{gathered}
\frac{d}{d t}\left(\frac{1}{2} m v_{x}^{2}+\frac{1}{2} m v_{y}^{2}\right) \\
=\left[\frac{|e| V}{\ln a / b}-\frac{|e|^{2}}{m}\left(\frac{\mu_{0} I}{2 \pi}\right)^{2} \ln \rho / b\right] \cdot \frac{x \dot{x}+y \dot{y}}{\rho^{2}} \\
=\left[\frac{|e| V}{\ln \frac{a}{b}}-\frac{|e|^{2}}{m}\left(\frac{\mu_{0} I}{2 \pi}\right)^{2} \ln \frac{\rho}{b}\right] \cdot \frac{1}{2 \rho^{2}} \frac{d}{d t} \rho^{2} \\
=\left[\frac{|e| V}{\ln \frac{a}{b}}-\frac{|e|^{2}}{m}\left(\frac{\mu_{0} I}{2 \pi}\right)^{2} \ln \frac{\rho}{b}\right] \frac{d}{d t} \ln \frac{\rho}{b}
\end{gathered}
\]

Integrating and using \(v^{2}=0\), at \(\rho=b\), we get,

\[
\frac{1}{2} m v^{2}=\left[\frac{|e| V}{\ln \frac{a}{b}} \ln \frac{\rho}{b}-\frac{1}{2 m}|e|^{2}\left(\frac{\mu_{0} I}{2 \pi}\right)^{2}\left(\ln \frac{\rho}{b}\right)\right]
\]

The RHS must be positive, for all \(a>\rho>b\). The condition for this is,

\[
V \geq \frac{1}{2} \frac{|e|}{m}\left(\frac{\mu_{0} I}{2 \pi}\right)^{2} \ln \frac{a}{b}
\]

3.394 This differs from the previous problem in ( \(a \leftrightarrow b\) ) and the magnetic field is along the \(z\)-direction. Thus \(B_{x}=B_{y}=0, B_{z}=B\)\\
Assuming as usual the charge of the electron to be \(-|e|\), we write the equation of motion\\
and

\[
\begin{gathered}
\frac{d}{d t} m v_{x}=\frac{|e| V_{x}}{\rho^{2} \ln \frac{b}{a}}-|e| B \dot{y}, \frac{d}{d t} m v_{y}=\frac{|e| V_{y}}{\rho^{2} \ln \frac{b}{a}}+|e| B \dot{x} \\
\frac{d}{d t} m v_{z}=0 \Rightarrow z=0
\end{gathered}
\]

The motion is confined to the plane \(z=0\). Eliminating \(B\) from the first two equations,\\
or,

\[
\begin{aligned}
\frac{d}{d t}\left(\frac{1}{2} m v^{2}\right) & =\frac{|e| V}{\ln b / a} \frac{x \dot{x}+y \dot{y}}{\rho^{2}} \\
\frac{1}{2} m v^{2} & =|e| V \frac{\ln \rho / a}{\ln b / a}
\end{aligned}
\]

so, as expected, since magnetic forces do not work,

\[
v=\sqrt{\frac{2|e| V}{m}} \text {, at } \rho=b \text {. }
\]

On the other hand, eliminating \(V\), we also get,

\[
\frac{d}{d t} m\left(x v_{y}-y v_{x}\right)=|e| B(x \dot{x}+y \dot{y})
\]

i.e.

\[
\left(x v_{y}-y v_{x}\right)=\frac{|e| B}{2 m} \rho^{2}+\text { constant }
\]

The constant is easily evaluated, since \(v\) is zero at \(\rho=a\). Thus,

At

\[
\left(x v_{y}-y v_{x}\right)=\frac{|e| B}{2 m}\left(\rho^{2}-a^{2}\right)>0
\]

Thus,

\[
\rho=b,\left(x v_{y}-y v_{x}\right) \leq v b
\]

Thus,

\[
v b \geq \frac{|e| B}{2 m}\left(b^{2}-a^{2}\right)
\]

or,

\[
B \leq \frac{2 m b}{b^{2}-a^{2}} \sqrt{\frac{2|e| V}{m}} \times \frac{1}{|e|}
\]

or,

\[
B \leq \frac{2 b}{b^{2}-a^{2}} \sqrt{\frac{2 m B}{|e|}}
\]

3.395 The equations are as in 3.392.

\[
\frac{d v_{x}}{d t}=\frac{q B}{m} v_{y}, \frac{d v_{y}}{d t}=\frac{q E_{m}}{m} \cos \omega t-\frac{q B}{m} v_{x} \text { and } \frac{d v_{z}}{d t}=
\]

with

\[
\begin{gathered}
\omega=\frac{q B}{m}, \xi=v_{x}+i v_{y}, \text { we get, } \\
\frac{d \xi}{d t}=i \frac{E_{m}}{B} \omega \cos \omega t-i \omega \xi
\end{gathered}
\]

or multiplying by \(e^{i \omega t}\),

\[
\frac{d}{d t}\left(\xi e^{i \omega t}\right)=i \frac{E_{m}}{2 B} \omega\left(e^{2 i \omega t}+1\right)
\]

or integrating,

\[
\xi e^{i \omega t}=\frac{E_{m}}{4 B} e^{2 i \omega t}+\frac{E_{m}}{2 B} i \omega t
\]

or,

\[
\xi=\frac{E_{m}}{4 B}\left(e^{i \omega t}+2 i \omega t e^{i \omega t}\right)+C e^{i \omega t}
\]

since

\[
\xi=0 \text { at } t=0, C=-\frac{E_{m}}{4 B} .
\]

Thus,

\[
\xi=i \frac{E_{m}}{2 B} \sin \omega t+i \frac{E_{m}}{2 B} \omega t e^{i \omega t}
\]

or,

\[
v_{x}=\frac{E_{m}}{2 B} \omega t \sin \omega t \text { and } v_{y}=\frac{E_{m}}{2 B} \sin \omega t+\frac{E_{m}}{2 B} \omega t \cos \omega t
\]

Integrating again,

\[
x=\frac{a}{2 \omega^{2}}(\sin \omega t-\omega t \cos \omega t), y=\frac{a}{2 \omega} t \sin \omega t
\]

where \(a=\frac{q E_{m}}{m}\), and we have used \(x=y=0\), at \(t=0\).\\
The trajectory is an unwinding spiral.\\
3.396 We know that for a charged particle (proton) in a magnetic field,

\[
\frac{m v^{2}}{r}=B e v \text { or } m v=B e r
\]

But,

\[
\omega=\frac{e B}{m},
\]

Thus

\[
E=\frac{1}{2} m v^{2}=\frac{1}{2} m \omega^{2} r^{2} .
\]

So,

\[
\Delta E=m \omega^{2} r \Delta r=4 \pi^{2} v^{2} m r \Delta r
\]

On the other hand \(\Delta E=2 \mathrm{eV}\), where \(V\) is the effective acceleration voltage, across the Dees, there being two crossings per revolution. So,

\[
V \geq 2 \pi^{2} v^{2} m r \Delta r / e
\]

3.397 (a) From \(\frac{m v^{2}}{r}=B e v\), or, \(m v=B e r\)\\
and

\[
T=\frac{(B e r)^{2}}{2 m}=\frac{1}{2} m v^{2}=12 \mathrm{MeV}
\]

(b) From \(\frac{2 \pi}{\omega}=\frac{2 \pi r}{v}\)\\
we get,

\[
f_{\min }=\frac{v}{2 \pi r}=\frac{1}{\pi r} \sqrt{\frac{T}{2 m}}=15 \mathrm{MHz}
\]

3.398 (a) The total time of acceleration is,

\[
t=\frac{1}{2 v} \cdot n,
\]

where \(n\) is the number of passages of the Dees.\\
But,

\[
T=n e V=\frac{B^{2} e^{2} r^{2}}{2 m}
\]

or,

\[
n=\frac{B^{2} e r^{2}}{2 m V}
\]

So,

\[
t=\frac{\pi}{e B / m} \times \frac{B^{2} e r^{2}}{2 m V}=\frac{\pi B r^{2}}{2 V}=\frac{\pi^{2} m v r^{2}}{e V}=30 \mu \mathrm{~s}
\]

(b) The distance covered is, \(s=\sum v_{n} \cdot \frac{1}{2 v}\)

But,

\[
v_{n}=\sqrt{\frac{2 e V}{m}} \sqrt{n}
\]

So, \(\quad s=\sqrt{\frac{e V}{2 m v^{2}}} \sum \sqrt{n}=\sqrt{\frac{e V}{2 m v^{2}}} \int \sqrt{n} d n=\sqrt{\frac{e V}{2 m v^{2}}} \frac{2}{3} n^{3 / 2}\)

But,

\[
n=\frac{B^{2} e^{2} r^{2}}{2 e V m}=\frac{2 \pi^{2} m v^{2} r^{2}}{e V}
\]

Thus,

\[
s \approx \frac{4 \pi^{3} v^{2} m r^{2}}{3 e V}=1.24 \mathrm{~km}
\]

3.399 In the nth orbit, \(\frac{2 \pi r_{n}}{v_{n}}=n T_{0}=\frac{n}{v}\). We ignore the rest mass of the electron and write \(\nu_{n} \approx c\). Also \(W \approx c p=c B e r_{n}\).

Thus,

\[
\frac{2 \pi W}{B e c^{2}}=\frac{n}{v}
\]

or,

\[
n=\frac{2 \pi W v}{B e c^{2}}=9
\]

3.400 The basic condition is the relativistic equation,

\[
\frac{m v^{2}}{r}=B q v, \quad \text { or }, \quad m v=\frac{m_{0} v}{\sqrt{1-v^{2} / c^{2}}}=B q r .
\]

Or calling,

\[
\omega=\frac{B q}{m},
\]

we get,

\[
\omega=\frac{\omega_{0}}{\sqrt{1+\frac{\omega_{0}^{2} r^{2}}{c^{2}}}}, \quad \omega_{0}=\frac{B q}{m_{0}} r
\]

is the radius of the instantaneous orbit.\\
The time of acceleration is,

\[
t=\sum_{n=1}^{N} \frac{1}{2 v_{n}}=\sum_{n=1}^{N} \frac{\pi}{\omega_{n}}=\sum_{n}^{N} \frac{\pi W_{n}}{q B c^{2}} .
\]

\(N\) is the number of crossing of either Dee.\\
But, \(W_{n}=m_{0} c^{2}+\frac{n \Delta W}{2}\), there being two crossings of the Dees per revolution.

So,

\[
\begin{gathered}
t=\sum \frac{\pi m_{0} c^{2}}{q B c^{2}}+\sum \frac{\pi \Delta W_{n}}{2 q B c^{2}} \\
=N \frac{\pi}{\omega_{0}}+\frac{N(N+1)}{4} \frac{\pi \Delta W}{q B c^{2}} \approx N^{2} \frac{\pi \Delta W}{4 q B c^{2}}(N \gg 1) \\
r=r_{N} \frac{v_{N}}{\omega_{N}} \approx \frac{c}{\pi} \frac{\partial t}{\partial N} \cong \frac{\Delta W}{2 q B c} N
\end{gathered}
\]

Also,

Hence finally,

\[
\begin{gathered}
\omega=\frac{\omega_{0}}{\sqrt{1+\frac{q^{2} B^{2}}{m_{0}^{2} c^{2}} \times \frac{\Delta W^{2}}{4 q^{2} B^{2} c^{2}}} N^{2}} \\
=\frac{\omega_{0}}{\sqrt{1+\frac{(\Delta W)^{2}}{4 m_{0}^{2} c^{4}} \times \frac{4 q B c^{2}}{\pi \Delta W} t}=\frac{\omega_{0}}{\sqrt{1+a t}}} ; \\
a=\frac{q B \Delta W}{\pi m_{0}^{2} c^{2}}
\end{gathered}
\]

3.401 When the magnetic field is being set up in the solenoid, and electric field will be induced in it, this will accelerate the charged particle. If \(B\) is the rate, at which the magnetic field is increasing, then.

\[
\pi r^{2} \dot{B}=2 \pi r E \text { or } E=\frac{1}{2} r \dot{B}
\]

Thus,

\[
m \frac{d v}{d t}=\frac{1}{2} r \dot{B} q, \text { or } v=\frac{q B r}{2 m}
\]

After the field is set up, the particle will execute a circular motion of radius \(\rho\), where

\[
m v=B q \rho, \quad \text { or } \rho=\frac{1}{2} r
\]

3.402 The increment in energy per revolution is \(\boldsymbol{e} \Phi\), so the number of revolutions is,

\[
N=\frac{W}{e \Phi}
\]

The distance traversed is, \(s=2 \pi r N\)\\
3.403 On the one hand,

\[
\frac{d p}{d t}=e E=\frac{e}{2 \pi r} \frac{d \Phi}{d t}=\frac{e}{2 \pi r} \frac{d}{d t} \int_{0}^{r} 2 \pi r^{\prime} B\left(r^{\prime}\right) d r^{\prime}
\]

On the other ,

\[
p=B(r) \text { er, } r=\text { constant. }
\]

so,

\[
\frac{d p}{d t}=e r \frac{d}{d t} B(r) \equiv e r \dot{B}(r)
\]

Hence, \(e r \dot{B}(r)=\frac{e}{2 \pi r} \pi r^{2} \frac{d}{d t}\langle B\rangle\)

So,

\[
\dot{B}(r)=\frac{1}{2} \frac{d}{d t}<B>
\]

This equations is most easily satisfied by taking \(B\left(r_{0}\right)=\frac{1}{2}\langle B\rangle\).\\
3.404 The condition, \(B\left(r_{0}\right)=\frac{1}{2}<B>=\frac{1}{2} \int_{0}^{r_{0}} B \cdot 2 \pi r d r / \pi r_{0}^{2}\)\\
or,

\[
B\left(r_{0}\right)=\frac{1}{r_{0}^{2}} \int_{0}^{r_{0}} B r d r
\]

This gives \(r_{0}\).\\
In the present case,

\[
B_{0}-a r_{0}^{2}=\frac{1}{r_{0}^{2}} \int_{0}^{r_{0}}\left(B-a r^{2}\right) r d r=\frac{1}{2}\left(B_{0}-\frac{1}{2} a r_{0}^{2}\right)
\]

or,

\[
\frac{3}{4} a r_{0}^{2}=\frac{1}{2} B_{0} \quad \text { or } \quad r_{0}=\sqrt{\frac{2 B_{0}}{3 a}}
\]

3.405 The induced electric field (or eddy current field) is given by,

\[
E(r)=\frac{1}{2 \pi r} \frac{d}{d t} \int_{0}^{r} 2 \pi r^{\prime}\left(r^{\prime}\right) B\left(r^{\prime}\right) d r^{\prime}
\]

Hence,

\[
\begin{gathered}
\frac{d E}{d r}=-\frac{1}{2 \pi r^{2}} \frac{d}{d t} \int_{0}^{r} 2 \pi r^{\prime} B\left(r^{\prime}\right) d r^{\prime}+\frac{d B(r)}{d t} \\
=-\frac{1}{2} \frac{d}{d t}<B>+\frac{d B(r)}{d t}
\end{gathered}
\]

This vanishes for \(r=r_{0}\) by the betatron condition, where \(r_{0}\) is the radius of the equilibrium orbit.\\
3.406 From the betatron condition,

\[
\frac{1}{2} \frac{d}{d t}<B>=\frac{d B}{d t}\left(r_{0}\right)=\frac{B}{\Delta t}
\]

Thus,

\[
\frac{d}{d t}<B>=\frac{2 B}{\Delta t}
\]

and

\[
\frac{d \Phi}{d t}=\pi r^{2} \frac{d<B>}{d t}=\frac{2 \pi r^{2} B}{\Delta t}
\]

So, energy increment per revolution is,

\[
e \frac{d \Phi}{d t}=\frac{2 \pi r^{2} e B}{\Delta t}
\]

3.407 (a) Even in the relativistic case, we know that : \(p=B e r\)

Thus,

\[
W=\sqrt{c^{2} p^{2}+m_{0}^{2} c^{4}}-m_{0} c^{2}=m_{0} c^{2}\left(\sqrt{1+\left(\operatorname{Ber} / m_{0} c\right)^{2}}-1\right)
\]

(b) The distance traversed is,

\[
2 \pi r \frac{W}{e \Phi}=2 \pi r \frac{W}{2 \pi r^{2} e B / \Delta t}=\frac{W \Delta t}{B e r}
\]

on using the result of the previous problem.


\end{document}


\section*{CONTENTS}
Preface ..... vii\\
PART FOUR\\
OSCILLATIONS AND WAVES\\
4.1 Mechanical Oscillations ..... 1\\
4.2 Electric Oscillations ..... 54\\
4.3 Elastic Waves. Acoustics ..... 82\\
4.4 Electromagnetic Waves. Radiation ..... 103\\
PART FIVE\\
OPTICS\\
5.1 Photometry and Geometrical Optics ..... 115\\
5.2 Interference of Light ..... 149\\
5.3 Diffraction of Light ..... 162\\
5.4 Polarization of Light ..... 196\\
5.5 Dispersion and Absorption of Light ..... 219\\
5.6 Optics of Moving Sources ..... 229\\
5.7 Thermal Radiation. Quantum Nature of Light ..... 238\\
PART SIX\\
ATOMIC AND NUCLEAR PHYSICS\\
6.1 Scattering of Particles. Rutherford-Bohr Atom ..... 259\\
6.2 Wave Properties of Particles. Schrödinger Equation ..... 285\\
6.3 Properties of Atoms. Spectra ..... 310\\
6.4 Molecules and Crystals ..... 337\\
6.5 Radioactivity ..... 360\\
6.6 Nuclear Reactions ..... 372\\
6.7 Elementary Particles ..... 387

\section*{OSCILLATIONS AND WAVES}
\subsection*{4.1 MECHANICAL OSCILLATIONS}
4.1 (a) Given, \(x=a \cos \left(\omega t-\frac{\pi}{4}\right)\)

So, \(v_{x}=\dot{x}=-a \omega \sin \left(\omega t-\frac{\pi}{4}\right)\) and \(w_{x}=\ddot{x}=-a \omega^{2} \cos \left(\omega t-\frac{\pi}{4}\right)\)

On-the basis of obtained expressions plots \(x(t), v_{x}(t)\) and \(w_{x}(t)\) can be drawn as shown in the answersheet, (of the problem book).\\
(b) From Eqn (1)


\begin{equation*}
v_{x}=-a \omega \sin \left(\omega t-\frac{\pi}{4}\right) \text { So, } v_{x}^{2}=a^{2} \omega^{2} \sin ^{2}\left(\omega t-\frac{\pi}{4}\right) \tag{2}
\end{equation*}


But from the law \(x=a \cos (\omega t-\pi / 4)\), so, \(x^{2}=a^{2} \cos ^{2}(\omega t-\pi / 4)\)\\
or,


\begin{equation*}
\cos ^{2}(\omega t-\pi / 4)=x^{2} / a^{2} \text { or } \sin ^{2}(\omega t-\pi / 4)=1-\frac{x^{2}}{a^{2}} \tag{3}
\end{equation*}


Using (3) in (2),


\begin{equation*}
v_{x}^{2}=a^{2} \omega^{2}\left(1-\frac{x^{2}}{a^{2}}\right) \quad \text { or } \quad v_{x}^{2}=\omega^{2}\left(a^{2}-x^{2}\right) \tag{4}
\end{equation*}


Again from Eqn (4), \(w_{x}=-a \omega^{2} \cos (\omega t-\pi / 4)=-\omega^{2} x\)\\
4.2 (a) From the motion law of the particle

\[
x=a \sin ^{2}(\omega t-\pi / 4)=\frac{a}{2}\left[1-\cos \left(2 \omega t-\frac{\pi}{2}\right)\right]
\]

or,

\[
x-\frac{a}{2}=-\frac{a}{2} \cos \left(2 \omega t-\frac{\pi}{2}\right)=-\frac{a}{2} \sin 2 \omega t=\frac{a}{2} \sin (2 \omega t+\pi)
\]

i.e.


\begin{equation*}
x-\frac{a}{2}=\frac{a}{2} \sin (2 \omega t+\pi) . \tag{1}
\end{equation*}


Now compairing this equation with the general equation of harmonic oscillations :\\
\(X=A \sin \left(\omega_{0} t+\alpha\right)\)\\
Amplitude, \(A=\frac{a}{2}\) and angular frequency, \(\omega_{0}=2 \omega\).\\
Thus the period of one full oscillation, \(T=\frac{2 \pi}{\omega_{0}}=\frac{\pi}{\omega}\)\\
(b) Differentiating Eqn (1) w.r.t. time


\begin{equation*}
\mathrm{v}_{\mathrm{x}}=a \omega \cos (2 \omega t+\pi) \text { or } \mathrm{v}_{\mathrm{x}}^{2}=a^{2} \omega^{2} \cos ^{2}(2 \omega t+\pi)=a^{2} \omega^{2}\left[1-\sin ^{2}(2 \omega t+\pi)\right] \tag{2}
\end{equation*}


From Eqn (1)

\[
\left(x-\frac{a}{2}\right)^{2}=\frac{a^{2}}{4} \sin ^{2}(2 \omega t+\pi)
\]

or,


\begin{equation*}
4 \frac{x^{2}}{a^{2}}+1-\frac{4 x}{a}=\sin ^{2}(2 \omega t+\pi) \text { or } 1-\sin ^{2}(2 \omega t+\pi) \doteq \frac{4 x}{a}\left(1-\frac{x}{a}\right) \tag{3}
\end{equation*}


From Eqns (2) and (3), \(\quad v_{x}=a^{2} \omega^{2} \frac{4 x}{a}\left(1-\frac{x}{a}\right)=4 \omega^{2} x(a-x)\)\\
Plot of \(\mathrm{v}_{\mathrm{x}}(x)\) is as shown in the answersheet.\\
4.3 Let the general equation of S.H.M. be


\begin{equation*}
x=a \cos (\omega t+\alpha) \tag{1}
\end{equation*}


So,


\begin{equation*}
v_{x}=-a \omega \sin (\omega t+\alpha) \tag{2}
\end{equation*}


Let us assume that at \(t=0, x=x_{0}\) and \(\mathrm{v}_{\mathrm{x}}=\mathrm{v}_{\mathrm{x}_{0}}\).\\
Thus from Eqns (1) and (2) for \(t=0, x_{0}=a \cos \alpha\), and \(v_{x_{0}}=-a \omega \sin \alpha\)\\
Therefore \(\quad \tan \alpha=-\frac{\mathrm{v}_{\mathrm{x}_{0}}}{\omega \mathrm{x}_{0}}\) and \(a=\sqrt{x_{0}^{2}+\left(\frac{\mathrm{v}_{\mathrm{x}_{0}}}{\omega}\right)^{2}}=35.35 \mathrm{~cm}\)\\
Under our assumption Eqns (1) and (2) give the sought \(x\) and \(v_{x}\) if

\[
t=t=2.40 \mathrm{~s}, a=\sqrt{x_{0}^{2}+\left(\mathrm{v}_{\mathrm{x}_{0}} / \omega\right)^{2}} \text { and } \alpha=\tan ^{-1}\left(-\frac{\mathrm{v}_{\mathrm{x}}}{\omega \mathrm{x}_{0}}\right)=-\frac{\pi}{4}
\]

Putting all the given numerical values, we get :

\[
x=-29 \mathrm{~cm} \quad \text { and } \quad v_{x}=-81 \mathrm{~cm} / \mathrm{s}
\]

4.4 From the Eqn, \(v_{x}^{2}=\omega^{2}\left(a^{2}-x^{2}\right)\) (see Eqn. 4 of 4.1)

\[
v_{1}^{2}=\omega^{2}\left(a^{2}-x_{1}^{2}\right) \text { and } v_{2}^{2}=\omega^{2}\left(a^{2}-x_{2}^{2}\right)
\]

Solving these Eqns simultaneously, we get

\[
\omega=\sqrt{\left(v_{1}^{2}-v_{2}^{2}\right) /\left(x_{2}^{2}-x_{1}^{2}\right)}, a=\sqrt{\left(v_{1} x_{2}^{2}-v_{2}^{2} x_{1}^{2}\right) /\left(v_{1}^{2}-v_{2}^{2}\right)}
\]

4.5 (a) When a particle starts from an extreme position, it is useful to write the motion law as


\begin{equation*}
x=a \cos \omega t \tag{1}
\end{equation*}


(However \(x\) is the displacement from the equlibrium position)\\
It \(t_{1}\) be the time to cover the distence \(a / 2\) then from (1)

\[
a-\frac{a}{2}=\frac{a}{2}=a \cos \omega t_{1} \quad \text { or } \quad \cos \omega t_{1}=\frac{1}{2}=\cos \frac{\pi}{3}\left(\text { as } t_{1}<T / 4\right)
\]

Thus

\[
t_{1}=\frac{\pi}{3 \omega}=\frac{\pi}{3(2 \pi / T)}=\frac{T}{6}
\]

As

\[
x=a \cos \omega t, \text { so, } \mathrm{v}_{\mathrm{x}}=-a \omega \sin \omega t
\]

Thus

\[
v=\left|v_{x}\right|=-v_{x}=a \omega \sin \omega t, \text { for } t \leq t_{1}=T / 6
\]

Hence sought mean velocity

\[
\langle v\rangle=\frac{\int v d t}{\int d t}=\int_{0}^{T / 6} a(2 \pi / T) \sin \omega t d t / T / 6=\frac{3 a}{T}=0.5 \mathrm{~m} / \mathrm{s}
\]

(b) In this case, it is easier to write the motion law in the form :


\begin{equation*}
x=a \sin \omega t \tag{2}
\end{equation*}


If \(t_{2}\) be the time to cover the distance \(a / 2\), then from Eqn (2)

\[
a / 2=a \sin \frac{2 \pi}{T} t_{2} \quad \text { or } \quad \sin \frac{2 \pi}{T} t_{2}=\frac{1}{2}=\sin \frac{\pi}{6}\left(\text { as } t_{2}<T / 4\right)
\]

Thus

\[
\frac{2 \pi}{T} t_{2}=\frac{\pi}{6} \quad \text { or }, t_{2}=\frac{T}{12}
\]

Differentiating Eqn (2) w.r.t time, we get

\[
v_{x}=a \omega \cos \omega t=a \frac{2 \pi}{T} \cos \frac{2 \pi}{T} t
\]

So,

\[
v=\left|v_{\lambda}\right|=a \frac{2 \pi}{T} \cos \frac{2 \pi}{T} t, \text { for } t \leq t_{2}=T / 12
\]

Hence the sought mean velocity

\[
\langle v\rangle=\frac{\int v d t}{\int d t}=\frac{1}{(T / 12)} \int_{0}^{T / 12} a \frac{2 \pi}{T} \cos \frac{2 \pi}{T} t d t=\frac{6 a}{T}=1 \mathrm{~m} / \mathrm{s}
\]

4.6\\
(a) As \(x=a \sin \omega t\)\\
Thus \(\left\langle\mathrm{v}_{\mathrm{x}}\right\rangle=\int \mathrm{v}_{\mathrm{x}} d t / \int d t=\frac{\int a \omega \cos (2 \pi / T) t d t}{\frac{3}{8} T}=\frac{2 \sqrt{2} a \omega}{3 \pi}\left(\right.\) using \(\left.T=\frac{2 \pi}{\omega}\right)\)\\
(b) In accordance with the problem

\[
\vec{v}=v_{x} \vec{i}, \text { so, }\left|<\vec{v}>\left|=\left|<v_{x}>\right|\right.\right.
\]

Hence, using part (a), \(|\langle\vec{v}\rangle|=\left|\frac{2 \sqrt{2} a \omega}{3 \pi}\right|=\frac{2 \sqrt{2} a \omega}{3 \pi}\)\\
(c) We have got, \(\mathrm{v}_{\mathrm{x}}=a \omega \cos \omega t\)

So,

\[
\left.\begin{array}{rl}
v & =\left|v_{x}\right|=a \omega \cos \omega t, \text { for } t \leq T / 4 \\
& =-a \omega \cos \omega t, \text { for } T / 4 \leq t \leq \frac{3}{8} T
\end{array}\right]
\]

Hence,

\[
<\mathrm{v}>=\frac{\int \mathrm{v} d t}{\int d t}=\frac{\int_{0}^{T / 4} a \omega \cos \omega t d t+\int_{T / 4}^{3 T / 8}-a \omega \cos \omega t d t}{3 T / 8}
\]

Using \(\omega=2 \pi / T\), and on evaluating the integral we get

\[
<v>=\frac{2(4-\sqrt{2}) a \omega}{3 \pi}
\]

4.7 From the motion law, \(x=a \cos \omega t\),, it is obvious that the time taken to cover the distance equal to the amplitude (a), starting from extreme position equals \(T / 4\).\\
Now one can write

\[
t=n \frac{T}{4}+t_{0}, \quad\left(\text { where } t_{0}<\frac{T}{4} \text { and } n=0,1,2, \ldots\right)
\]

As the particle moves according to the law, \(x=a \cos \omega t\),\\
so at \(n=1,3,5 \ldots\) or for odd \(n\) values it passes through the mean positon and for even numbers of \(n\) it comes to an extreme position (if \(t_{0}=0\) ).\\
Case (1) when \(n\) is an odd number :\\
In this case, from the equation\\
\(x=x a \sin \omega t\), if the t is counted from \(n T / 4\) and the distance covered in the time interval to becomes, \(s_{1}=a \sin \omega t_{0}=a \sin \omega\left(t-n \frac{T}{4}\right)=a \sin \left(\omega t-\frac{n \pi}{2}\right)\)\\
Thus the sought distance covered for odd \(n\) is

\[
s=n a+s_{1}=n a+a \sin \left(\omega t-\frac{n \pi}{2}\right)=a\left[n+\sin \left(\omega t-\frac{n \pi}{2}\right)\right]
\]

Case (2), when \(n\) is even, In this case from the equation \(x=a \cos \omega t\), the distance covered ( \(s_{2}\) ) in the interval \(t_{0}\), is given by\\
or,

\[
\begin{aligned}
a-s_{2}=a \cos \omega t_{0} & =a \cos \omega\left(t-n \frac{T}{4}\right)=a \cos \left(\omega t-n \frac{\pi}{2}\right) \\
s_{2} & =a\left[1-\cos \left(\omega t-\frac{n \pi}{2}\right)\right]
\end{aligned}
\]

Hence the sought distance for \(n\) is even

\[
s=n a+s_{2}=n a+a\left[1-\cos \left(\omega t-\frac{n \pi}{2}\right)\right]=a\left[n+1-\cos \left(\omega t-\frac{n \pi}{2}\right)\right]
\]

In general

\[
s=\left\{\begin{array}{l}
a\left[n+1-\cos \left(\omega t-\frac{n \pi}{2}\right)\right], n \text { is even } \\
a\left[n+\sin \left(\omega t-\frac{n \pi}{2}\right)\right], n \text { is odd }
\end{array}\right.
\]

4.8 Obviously the motion law is of the from, \(x=a \sin \omega t\). and \(v_{x}=\omega a \cos \omega t\).

Comparing \(\mathrm{v}_{\mathrm{x}}=\omega a \cos \omega t\) with \(\mathrm{v}_{\mathrm{x}}=35 \cos \pi t\), we get

\[
\omega=\pi, a=\frac{35}{\pi}, \text { thus } T=\frac{2 \pi}{\omega}=2 \text { and } T / 4=0.5 \mathrm{~s}
\]

Now we can write

\[
t=2.8 \mathrm{~s}=5 \times \frac{T}{4}+0.3\left(\text { where } \frac{T}{4}=0.5 \mathrm{~s}\right)
\]

As \(n=5\) is odd, like ( \(4 \cdot 7\) ), we have to basically find the distance covered by the particle starting from the extreme position in the time interval 0.3 s .\\
Thus from the Eqn.

\[
\begin{gathered}
x=a \cos \omega t=\frac{35}{\pi} \cos \pi(0.3) \\
\frac{35}{\pi}-s_{1}=\frac{35}{\pi} \cos \pi(0.3) \text { or } s_{1}=\frac{35}{\pi}\{1-\cos 0.3 \pi\}
\end{gathered}
\]

Hence the sought distance

\[
\begin{gathered}
s=5 \times \frac{35}{\pi}+\frac{35}{\pi}\{1-\cos 0.3 \pi\} \\
=\frac{35}{\pi}\{6-\cos 0.3 \pi\}=\frac{35}{22} \times 7\left(6-\cos 54^{\circ}\right) \approx 60 \mathrm{~cm}
\end{gathered}
\]

4.9 As the motion is periodic the particle repeatedly passes through any given region in the range \(-a \leq x \leq a\). The probability that it lies in the range \((x, x+d x)\) is defined as the fraction \(\frac{\Delta t}{t}\) (as \(t \rightarrow \infty\) ) where \(\Delta t\) is the time that the particle lies in the range \((x, x+d x)\) out of the total time \(t\). Because of periodicity this is

\[
d P=\frac{d P}{d x} d x=\frac{d t}{T}=\frac{2 d x}{v T}
\]

where the factor 2 is needed to take account of the fact that the particle is in the range ( \(x, x+d x\) ) during both up and down phases of its motion. Now in a harmonic oscillator.

\[
v=\dot{x}=\omega a \cos \omega t=\omega \sqrt{a^{2}-x^{2}}
\]

Thus since \(\omega T=2 \pi\) ( \(T\) is the time period)\\
We get

\[
d P=\frac{d P}{d x} d x=\frac{1}{\pi} \frac{d x}{\sqrt{a^{2}-x^{2}}}
\]

Note that

\[
\int_{-a}^{+a} \frac{d P}{d x} d x=1
\]

so

\[
\frac{d P}{d x}=\frac{1}{\pi} \frac{1}{\sqrt{a^{2}-x^{2}}} \text { is properly normalized. }
\]

4.10 (a) We take a graph paper and choose an axis ( \(X\)-axis) and an origin. Draw a vector of magnitude 3 inclined at an angle \(\frac{\pi}{3}\) with the \(X\)-axis. Draw another vector of magnitude 8 inclined at an angle \(-\frac{\pi}{3}\) (Since \(\sin (\omega t+\pi / 6)=\cos (\omega t-\pi / 3)\) ) with the \(X\)-axis. The magnitude of the resultant of both these vectors (drawn from the origin) obtained using parallelogram law is the resultant, amplitude.\\
\includegraphics[max width=\textwidth, alt={}, center]{da0156b6-4133-434e-90d4-ec01b2070afe-014_437_476_85_1017}

Clearly

\[
\begin{aligned}
R^{2} & =3^{2}+8^{2}+2 \cdot 3 \cdot 8 \cdot \cos \frac{2 \pi}{3}=9+64-48 \times \frac{1}{2} \\
& =73-24=49
\end{aligned}
\]

Thus

\[
R=7 \text { units }
\]

(b) One can follow the same graphical method here but the result can be obtained more quickly by breaking into sines and cosines and adding :

Resultant

\[
\begin{aligned}
x & =\left(3+\frac{5}{\sqrt{2}}\right) \cos \omega t+\left(6-\frac{5}{\sqrt{2}}\right) \sin \omega t \\
& =A \cos (\omega t+\alpha)
\end{aligned}
\]

Then

\[
\begin{aligned}
A^{2} & =\left(3+\frac{5}{\sqrt{2}}\right)^{2}+\left(6-\frac{5}{\sqrt{2}}\right)^{2}=9+25+\frac{30-60}{\sqrt{2}}+36 \\
& =70-15 \sqrt{2}=70-21 \cdot 2
\end{aligned}
\]

So,

\[
A=6.985 \approx 7 \text { units }
\]

Note- In using graphical method convert all oscillations to either sines or cosines but do not use both.\\
4.11 Given, \(x_{1}=a \cos \omega t \quad\) and \(\quad x_{2}=a \cos 2 \omega t\)\\
so, the net displacement,

\[
x=x_{1}+x_{2}=a\{\cos \omega t+\cos 2 \omega t\}=a\left\{\cos \omega t+2 \cos ^{2} \omega t-1\right\}
\]

and

\[
v_{x}=\dot{x}=a\{-\omega \sin \omega t-4 \omega \cos \omega t \sin \omega t\}
\]

For \(\dot{x}\) to be maximum,

\[
\ddot{x}=a \omega^{2} \cos \omega t-4 a \omega^{2} \cos ^{2} \omega t+4 a \omega^{2} \sin ^{2} \omega t=0
\]

or, \(\quad 8 \cos ^{2} \omega t+\cos \omega t-4=0\), which is a quadratic equation for \(\cos \omega t\).\\
Solving for acceptable value\\
thus

\[
\cos \omega t=0.644
\]

and

\[
v_{\max }=\left|v_{x_{\max }}\right|=+a \omega[0.765+4 \times 0.765 \times 0.644]=+2.73 a \omega
\]

4.12 We write :

\[
a \cos 2 \cdot 1 t \cos 50 \cdot 0 t=\frac{a}{2}\{\cos 52 \cdot 1 t+\cos 47 \cdot 9 t\}
\]

Thus the angular frequencies of constituent oscillations are

\[
52 \cdot 1 \mathrm{~s}^{-1} \text { and } 47 \cdot 9 \mathrm{~s}^{-1}
\]

To get the beat period note that the variable amplitude \(a \cos 2 \cdot 1 t\) becomes maximum (positive or negative), when

\[
2 \cdot 1 t=n \pi
\]

Thus the interval between two maxima is

\[
\frac{\pi}{2 \cdot 1}=1.5 \text { s nearly. }
\]

4.13 If the frequency of \(A\) with respect to \(K^{\prime}\) is \(v_{0}\) and \(K^{\prime}\) oscillates with frequency \(\bar{v}\) with respect to \(K\), the beat frequency of the point \(A\) in the \(K\)-frame will be \(v\) when

\[
\bar{v}=v_{0} \pm v
\]

In the present case \(\bar{v}=20\) or 24 . This means

\[
v_{0}=22 . \& v=2
\]

Thus beats of \(2 v=4\) will be heard when \(\bar{v}=26\) or 18 .\\
4.14 (a) From the Eqn : \(x=a \sin \omega t\)


\begin{equation*}
\sin ^{2} \omega t=x^{2} / a^{2} \text { or } \cos ^{2} \omega t=1-\frac{x^{2}}{a^{2}} \tag{1}
\end{equation*}


And from the equation : \(\quad y=b \cos \omega t\)


\begin{equation*}
\cos ^{2} \omega t=y^{2} / b^{2} \tag{2}
\end{equation*}


From Eqns (1) and (2), we get :

\[
1-\frac{x^{2}}{a^{2}}=\frac{y^{2}}{b^{2}} \quad \text { or } \quad \frac{x^{2}}{a^{2}}+\frac{y^{2}}{b^{2}}=1
\]

which is the standard equation of the ellipse shown in the figure.\\
we observe that,\\
at

\[
\begin{gathered}
t=0, x=0 \text { and } y=b \\
t=\frac{\pi}{2 \omega}, x=+a \text { and } y=0
\end{gathered}
\]

and at

Thus we observe that at \(t=0\), the point is at point 1 (Fig.) and at the following moments, the co-ordinate \(y\) diminishes and \(x\) becomes positive. Consequently the motion is clockwise.\\
(b) As \(x=a \sin \omega t\) and \(y=b \cos \omega t\)

So we may write \(\vec{r}=a \sin \omega t \overrightarrow{i+b \cos \omega t \vec{j}}\)\\
Thus \(\ddot{\mathbf{r}}=\vec{w}=-\omega^{2} \vec{r}\)\\
4.15 (a) From the Eqn. : \(x=a \sin \omega t\), we have

\[
\cos \omega t=\sqrt{1-\left(x^{2} / a^{2}\right)}
\]

and from the Eqn. : \(y=a \sin 2 \omega t\)

\[
y=2 a \sin \omega t \cos \omega t=2 x \sqrt{1-\left(x^{2} / a^{2}\right)} \text { or } y^{2}=4 x^{2}\left(1-\frac{x^{2}}{a^{2}}\right)
\]

(b) From the Eqn. : \(x=a \sin \omega t\);

\[
\sin ^{2} \omega t=x^{2} / a^{2}
\]

Fromy \(=a \cos 2 \omega t\)

\[
y=a\left(1-2 \sin ^{2} \omega t\right)=a\left(1-2 \frac{x^{2}}{a^{2}}\right)
\]

For the plots see the plots of answersheet of the problem book.\\
4.16 As \(U(x)=U_{0}(1-\cos a x)\)

So,


\begin{equation*}
F_{x}=-\frac{d U}{d x}=-U_{0} a \sin a x \tag{1}
\end{equation*}


or, \(\quad F_{x}=-U_{0} a a x\) (because for small angle of oscillations \(\sin a x \cong a x\) )\\
or,


\begin{equation*}
F_{x}=-U_{0} a^{2} x \tag{1}
\end{equation*}


But we know \(F_{x}=-m \omega_{0}^{2} x\), for small oscillation

Thus

\[
\omega_{0}^{2}=\frac{U_{0} a^{2}}{m} \quad \text { or } \quad \omega_{0}=a \sqrt{\frac{U_{0}}{m}}
\]

Hence the sought time period

\[
T=\frac{2 \pi}{\omega_{0}}=\frac{2 \pi}{a} \sqrt{\frac{m}{U_{0}}}=2 \pi \sqrt{\frac{m}{a^{2} U_{0}}}
\]

4.17 If \(U(x)=\frac{a}{x^{2}}-\frac{b}{x}\)\\
then the equilibrium position is \(x=x_{0}\) when \(U^{\prime}\left(x_{0}\right)=0\)\\
or

\[
-\frac{2 a}{x_{0}^{3}}+\frac{b}{x_{0}^{2}}=0 \Rightarrow x_{0}=\frac{2 a}{b}
\]

Now write :

\[
x=x_{0}+y
\]

Then

\[
U(x)=\frac{a}{x_{0}^{2}}-\frac{b}{x_{0}}+\left(x-x_{0}\right) U^{\prime}\left(x_{0}\right)+\frac{1}{2}\left(x-x_{0}\right)^{2} U^{\prime \prime}\left(x_{0}\right)
\]

But

\[
U^{\prime \prime}\left(x_{0}\right)=\frac{6 a}{x_{0}^{4}}-\frac{2 b}{x_{0}^{3}}=(2 a / b)^{-3}(3 b-2 b)=b^{4} / 8 a^{3}
\]

So finally :

\[
U(x)=U\left(x_{0}\right)+\frac{1}{2}\left(\frac{b^{4}}{8 a^{3}}\right) y^{2}+\ldots
\]

We neglect remaining terms for small oscillations and compare with the P.E. for a harmonic, oscillator :

\[
\begin{aligned}
\frac{1}{2} m \omega^{2} y^{2} & =\frac{1}{2}\left(\frac{b^{4}}{8 a^{3}}\right) y^{2}, \text { so } \omega=\frac{b^{2}}{\sqrt{8 a^{3} m}} \\
T & =2 \pi \frac{\sqrt{8 m a^{3}}}{b^{2}}
\end{aligned}
\]

Thus

Note : Equilibrium position is generally a minimum of the potential energy. Then \(U^{\prime}\left(x_{0}\right)=0, U^{\prime \prime}\left(x_{0}\right)>0\). The equilibrium position can in principle be a maximum but then \(U^{\prime \prime}\left(x_{0}\right)<0\) and the frequency of oscillations about this equilibrium position will be imaginary.

The answer given in the book is incorrect both numerically and dimensionally.\\
4.18 Let us locate and depict the forces acting on the ball at the position when it is at a distance \(x\) down from the undeformed position of the string.\\
At this position, the unbalanced downward force on the ball

\[
=m g-2 F \sin \theta
\]

By Newton's law, \(m \ddot{x}=m g-2 F \sin \theta\)

\[
\begin{aligned}
& =m g-2 F \theta(\text { when θis small }) \\
& =m g-2 F \frac{x}{l / 2}=m g-\frac{4 F}{l} x
\end{aligned}
\]

Thus

\[
\ddot{x}=g-\frac{4 F}{m l} x=-\frac{4 F}{m l}\left(x-\frac{m g l}{4 F}\right)
\]

putting \(x^{\prime}=x-\frac{m g \cdot l}{T}\), we get\\
\includegraphics[max width=\textwidth, alt={}, center]{da0156b6-4133-434e-90d4-ec01b2070afe-017_364_655_821_783}

\[
\ddot{x}^{\prime}=-\frac{4 T}{m l} x^{\prime}
\]

Thus \(T=\pi \sqrt{\frac{m l}{F}}=0.2 \mathrm{~s}\)\\
4.19 Let us depict the forces acting on the oscillating ball at an arbitraty angular position \(\theta\). (Fig.), relative to equilibrium position where \(F_{B}\) is the force of buoyancy. For the ball from the equation :\\
\(N_{Z}=I \beta_{Z}\), (where we have taken the positive sense of \(Z\) axis in the direction of angular velocity i.e. \(\dot{\theta}\) of the ball and passes through the point of suspension of the pendulum \(O\) ), we get :


\begin{equation*}
-m g l \sin \theta+F_{B} l \sin \theta=m l^{2} \ddot{\theta} \tag{1}
\end{equation*}


\begin{center}
\includegraphics[max width=\textwidth, alt={}]{da0156b6-4133-434e-90d4-ec01b2070afe-017_444_420_1444_946}
\end{center}

Using \(m=\frac{4}{3} \pi r^{3} \sigma, F_{B}=\frac{4}{3} \pi r^{3} \rho\) and \(\sin \theta \propto \theta\) for small \(\theta\), in Eqn (1), we get :

\[
\ddot{\theta}=-\frac{g}{l}\left(1-\frac{\rho}{\sigma}\right) \theta
\]

Thus the sought time period

Hence

\[
\begin{gathered}
T=2 \pi \frac{1}{\sqrt{\frac{g}{l}\left(1-\frac{\rho}{\sigma}\right)}}=2 \alpha \sqrt{\frac{l / g}{1-\frac{1}{\eta}}} \\
T=2 \alpha \sqrt{\frac{\eta l}{g(\eta-1)}}=1.1 \mathrm{~s}
\end{gathered}
\]

4.20 Obviously for small \(\beta\) the ball execute part of S.H.M. Due to the perfectly elastic collision the velocity of ball simply reversed. As the ball is in S.H.M. \((|\theta|<\alpha\) on the left)its motion law in differential from can be written as


\begin{equation*}
\ddot{\theta}=-\frac{g}{l} \theta=-\omega_{0}^{2} \theta \tag{1}
\end{equation*}


If we assume that the ball is released from the extreme position, \(\theta=\beta\) at \(t=0\), the solution of differential equation would be taken in the form


\begin{equation*}
\theta=\beta \cos \omega_{0} t=\beta \cos \sqrt{\frac{g}{l}} t \tag{2}
\end{equation*}


If \(t^{\prime}\) be the time taken by the ball to go from the extreme position \(\theta=\beta\) to the wall i.e. \(\theta=-\alpha\), then Eqn. (2) can be rewritten as

\[
\begin{gathered}
-\alpha=\beta \cos \sqrt{\frac{g}{l}} t^{\prime} \\
\text { or } \quad t^{\prime}=\sqrt{\frac{l}{g}} \cos ^{-1}\left(-\frac{\alpha}{\beta}\right)=\sqrt{\frac{l}{g}}\left(\pi-\cos ^{-1} \frac{\alpha}{\beta}\right)
\end{gathered}
\]

Thus the sought time \(T=2 t^{\prime}=2 \sqrt{\frac{l}{g}}\left(\pi-\cos ^{-1} \frac{\alpha}{\beta}\right)\)

\[
=2 \sqrt{\frac{l}{g}}\left(\frac{\pi}{2}+\sin ^{-1} \frac{\alpha}{\beta}\right), \quad\left[\text { because } \sin ^{-1} x+\cos ^{-1} x=\pi / 2\right]
\]

4.21 Let the downward acceleration of the elevator car has continued for time \(t^{\prime}\), then the sought time\\
\(t=\sqrt{\frac{2 h}{w}}+t^{\prime}\), where obviously \(\sqrt{\frac{2 h}{w}}\) is the time of upward acceleration of the elevator.\\
One should note that if the point of suspension of a mathematical pendulum moves with an acceleration \(\vec{w}\), then the time period of the pendulum becomes


\begin{equation*}
2 \pi \sqrt{\frac{l}{|\vec{g}-\vec{w}|}} \tag{see4.30}
\end{equation*}


In this problem the time period of the pendulum while it is moving upward with acceleration \(w\) becomes\\
\(2 \pi \sqrt{\frac{l}{g+w}}\) and its time period while the elevator moves downward with the same magnitude of acceleration becomes

\[
2 \pi \sqrt{\frac{l}{g-w}}
\]

As the time of upward acceleration equals \(\sqrt{\frac{2 h}{w}}\), the total number of oscillations during this time equals

\[
\frac{\sqrt{2 h / w}}{2 \pi \sqrt{l /(g+w)}}
\]

Thus the indicated time \(=\frac{\sqrt{2 h / w}}{2 \pi \sqrt{l /(g+w)}} \cdot 2 \pi \sqrt{l / g}=\sqrt{2 h / w} \sqrt{(g+w) / g}\)\\
Similarly the indicated time for the time interval \(t^{\prime}\)

\[
=\frac{t^{\prime}}{2 \pi \sqrt{l /(g-w)}} 2 \pi \sqrt{l / g}=t^{\prime} \sqrt{(g-w) / g}
\]

we demand that

\[
\begin{gathered}
\sqrt{2 h / w} \sqrt{(g+w) / g}+t^{\prime} \sqrt{(g-w) / g}=\sqrt{2 h / w}+t^{\prime} \\
\text { or, } \quad t^{\prime}=\sqrt{2 h / w} \frac{\sqrt{g+w}-\sqrt{g}}{\sqrt{g}-\sqrt{g-w}}
\end{gathered}
\]

Hence the sought time

\[
\begin{aligned}
t & =\sqrt{\frac{2 h}{w}}+t^{\prime}=\sqrt{\frac{2 h}{w}} \frac{\sqrt{g+w}-\sqrt{g-w}}{\sqrt{g}-\sqrt{g-w}} \\
& =\sqrt{\frac{2 h}{w}} \frac{\sqrt{1+\beta}-\sqrt{1-\beta}}{1-\sqrt{1-\beta}}, \text { where } \beta=w / g
\end{aligned}
\]

4.22 If the hydromoter were in equlibrium or floating, its weight will be balanced by the buoyancy force acting on it by the fluid. During its small oscillation, let us locate the hydrometer when it is at a vertically downward distance \(x\) from its equilibrium position. Obviously the net unbalanced force on the hydrometer is the excess buoyancy force directed upward and equals \(\pi r^{2} x \rho g\). Hence for the hydrometer.\\
or,

\[
\begin{aligned}
m \ddot{x} & =-\pi r^{2} \rho g x \\
\ddot{x} & =-\frac{\pi r^{2} \rho g}{m} x
\end{aligned}
\]

Hence the sought time period

\[
T=2 \alpha \sqrt{\frac{m}{\pi r^{2} \rho g}}=2.5 \mathrm{~s} .
\]

4.23 At first let us calculate the stiffness \(\kappa_{1}\) and \(\kappa_{2}\) of both the parts of the spring. If we subject the original spring of stiffness \(\kappa\) having the natural length \(l_{0}\) (say), under the deforming forces \(F-F\) (say) to elongate the spring by the amount \(x\), then


\begin{equation*}
F=\kappa x \tag{1}
\end{equation*}


Therefore the elongation per unit length of the spring is \(x / l_{0}\). Now let us subject one of the parts of the spring of natural length \(\eta l_{0}\) under the same deforming forces \(F-F\). Then the elongation of the spring will be

Thus


\begin{align*}
& \frac{x}{l_{0}} \eta l_{0}=\eta x \\
& F=\kappa_{1}(\eta x) \tag{2}
\end{align*}


Hence from Eqns (1) and (2)


\begin{equation*}
\kappa=\eta \kappa_{1} \text { or } \kappa_{1}=\kappa / \eta \tag{3}
\end{equation*}


Similarly

\[
\kappa_{2}=\frac{\kappa}{1-\eta}
\]

The position of the block \(m\) when both the parts of the spring are non-deformed, is its equilibrium position \(O\). Let us displace the block \(m\) towards right or in positive \(x\) axis by the small distance \(x\). Let us depict the forces acting on the block when it is at a distance \(x\) from its equilibrium position (Fig.). From the second law of motion in projection form i.e. \(F_{x}=m w_{x}\)

\[
-\kappa_{1} x-\kappa_{2} x=m \ddot{x}
\]

or, \(\quad-\left(\frac{\kappa}{\eta}+\frac{\kappa}{1-\eta}\right) x=m \ddot{x}\)\\
Thus \(\quad \ddot{x}=-\frac{\kappa}{m} \frac{1}{\eta(m)} x\)\\
\includegraphics[max width=\textwidth, alt={}, center]{da0156b6-4133-434e-90d4-ec01b2070afe-020_184_617_990_843}

Hence the sought time period

\[
T=2 \pi \sqrt{\eta(1 \div \eta) m / \kappa}=0.13 \mathrm{~s}
\]

4.24 Similar to the Soln of 4.23, the net unbalanced force on the block \(m\) when it is at a small horizontal distance \(x\) from the equilibrium position becomes \(\left(\kappa_{1}+\kappa_{2}\right) x\).\\
From \(F_{x}=m w_{x}\) for the block :

\[
-\left(\kappa_{1}+\kappa_{2}\right) x=m \ddot{x}
\]

Thus

\[
\ddot{x}=-\left(\frac{\kappa_{1}+\kappa_{2}}{m}\right) x
\]

Hence the sought time period \(T=2 \pi \sqrt{\frac{m}{\kappa_{1}+\kappa_{2}}}\)\\
Alternate : Let us set the block \(m\) in motion to perform small oscillation. Let us locate the block when it is at a distance \(x\) from its equilibrium position.\\
As the spring force is restoring conservative force and deformation of both the springs are same, so from the conservation of mechanical energy of oscillation of the spring-block system :

\[
\frac{1}{2} m\left(\frac{d x}{d t}\right)^{2}+\frac{1}{2} \kappa_{1} x^{2}+\frac{1}{2} \kappa_{2} x^{2}=\text { Constant }
\]

Differentiating with respect to time

\[
\begin{gathered}
\frac{1}{2} m 2 \dot{x} \ddot{x}+\frac{1}{2}\left(\kappa_{1}+\kappa_{2}\right) 2 x \dot{x}=0 \\
\ddot{x}=-\frac{\left(\kappa_{1}+\kappa_{2}\right)}{m} x
\end{gathered}
\]

or,

\[
T=2 \pi \sqrt{\frac{m}{\kappa_{1}+\kappa_{2}}}
\]

4.25 During the vertical oscillation let us locate the block at a vertical down distance \(x\) from its equilibrium position. At this moment if \(x_{1}\) and \(x_{2}\) are the additional or further elongation of the upper \& lower springs relative to the equilibrium position, then the net unbalanced force on the block will be \(\kappa_{2} x_{2}\) directed in upward direction. Hence


\begin{equation*}
-\kappa_{2} x_{2}=m \ddot{x} \tag{1}
\end{equation*}


We also have


\begin{equation*}
x=x_{1}+x_{2} \tag{2}
\end{equation*}


As the springs are massless and initially the net force on the spring is also zero so for the spring


\begin{equation*}
\kappa_{1} x_{1}=\kappa_{2} x_{2} \tag{3}
\end{equation*}


Solving the Eqns (1), (2) and (3) simultaneously, we get

\[
\begin{gathered}
-\frac{\kappa_{1} \kappa_{2}}{\kappa_{1}+\kappa_{2}} x=m \ddot{x} \\
\ddot{x}=-\frac{\left(\kappa_{1} \kappa_{2} / \kappa_{1}+\kappa_{2}\right)}{m} x
\end{gathered}
\]

Thus\\
Hence the sought time period \(T=2 \pi \sqrt{m \frac{\left(\kappa_{1} \kappa_{2}\right)}{\kappa_{1} \kappa_{2}}}\)\\
4.26 The force \(F\), acting on the weight deflected from the position of equilibrium is \(2 T_{0} \sin \theta\).

Since the angle \(\theta\) is small, the net restoring force, \(F=2 T_{0} \frac{x}{l}\)\\
or, \(F=k x\), where \(k=\frac{2 T_{0}}{l}\)\\
So, by using the formula,\\
\(\omega_{0}=\sqrt{\frac{k}{m}}, \quad \omega_{0}=\sqrt{\frac{2 T_{0}}{m l}}\)\\
\includegraphics[max width=\textwidth, alt={}, center]{da0156b6-4133-434e-90d4-ec01b2070afe-021_256_595_1508_792}\\
4.27 If the mercury rises in the left arm by \(x\) it must fall by a slanting length equal to \(x\) in the other arm. Total pressure difference in the two arms will then be

\[
\rho g x+\rho g x \cos \theta=\rho g x(1+\cos \theta)
\]

This will give rise to a restoring force

\[
-\rho g S x(1+\cos \theta)
\]

This must equal mass times acceleration which can be obtained from work energy principle.\\
\includegraphics[max width=\textwidth, alt={}, center]{da0156b6-4133-434e-90d4-ec01b2070afe-022_380_931_63_265}

The K.E. of the mercury in the tube is clearly : \(\frac{1}{2} m \dot{x}^{2}\)\\
So mass times acceleration must be : \(m \ddot{x}\)\\
Hence

\[
m \ddot{x}+\rho g S(1+\cos \theta) x=0
\]

This is S.H.M. with a time period

\[
T=2 \pi \sqrt{\frac{m}{\rho g S(1+\cos \theta)}} .
\]

4.28 In the equilibrium position the C.M. of the rod lies nid way between the two rotating wheels. Let us displace the rod horizontally by some small distance and then release it. Let us depict the forces acting on the rod when its C.M. is at distance \(x\) from its equilibrium position (Fig.). Since there is no net vertical force acting on the rod, Newton's second law gives :\\
\includegraphics[max width=\textwidth, alt={}, center]{da0156b6-4133-434e-90d4-ec01b2070afe-022_201_468_1041_315}\\
\includegraphics[max width=\textwidth, alt={}, center]{da0156b6-4133-434e-90d4-ec01b2070afe-022_294_446_1013_846}


\begin{equation*}
N_{1}+N_{2}=m g \tag{1}
\end{equation*}


For the translational motion of the rod from the Eqn. : \(F_{x}=m w_{c x}\)


\begin{equation*}
k N_{1}-k N_{2}=m \ddot{x} \tag{2}
\end{equation*}


As the rod experiences no net torque about an axis perpendicular to the plane of the Fig. through the C.M. of the rod.


\begin{equation*}
N_{1}\left(\frac{l+x}{2}\right)=N_{2}\left(\frac{l-x}{2}\right) \tag{3}
\end{equation*}


Solving Eqns. (1), (2) and (3) simultaneously we get

\[
\ddot{x}=-k \frac{2 g}{l} x
\]

Hence the sought time period

\[
T=2 \pi \sqrt{\frac{l}{2 k g}}=\pi \sqrt{\frac{2 l}{k g}}=1.5 \mathrm{~s}
\]

4.29 (a) The only force acting on the ball is the gravitational force \(\vec{F}\), of magnitude \(\gamma \frac{4}{3} \pi \rho m r\), where \(\gamma\) is the gravitational constant \(\rho\), the density of the Earth and \(r\) is the distance of the body from the centre of the Earth.\\
But, \(g=\gamma \frac{4 \pi}{3} \rho R\), so the expression for \(\vec{F}\) can be written as,\\
\(\vec{F}=-m g \frac{\vec{r}}{R}\), here \(R\) is the radius of the Earth and the equation of motion in projection form has the form, or, \(m \ddot{x}+\frac{m g}{R} x=0\)\\
(b) The equation, obtained above has the form of an equation of S.H.M. having the time period,

\[
T=2 \pi \sqrt{\frac{R}{g}},
\]

Hence the body will reach the other end of the shaft in the time,

\[
t=\frac{T}{2}=\pi \sqrt{\frac{R}{g}}=42 \mathrm{~min}
\]

(c) From the conditions of S.H.M., the speed of the body at the centre of the Earth will be maximum, having the magnitude,

\[
v=R \omega=R \sqrt{g / R}=\sqrt{g R}=7.9 \mathrm{~km} / \mathrm{s}
\]

4.30 In the frame of point of suspension the mathematical pendulum of mass \(m\) (say) will oscillate. In this frame, the body \(m\) will experience the inertial force \(m(-\vec{w})\) in addition to the real forces during its oscillations. Therefore in equilibrium position \(m\) is deviated by some angle say \(\alpha\). In equilibrium position

\[
T_{0} \cos \alpha=m g+m w \cos (\pi-\beta) \text { and } T_{0} \sin \alpha=m w \sin (\pi-\beta)
\]

So, from these two Eqns

\[
\left.\begin{array}{l}
\tan \alpha=\frac{g-w \cos \beta}{w \sin \beta} \\
\text { and } \cos \alpha=\sqrt{\frac{m^{2} w^{2} \sin ^{2} \beta+(m g-m w \cos \beta)^{2}}{m g-m w \cos \beta}} \tag{1}
\end{array}\right]
\]

\includegraphics[max width=\textwidth, alt={}, center]{da0156b6-4133-434e-90d4-ec01b2070afe-023_599_426_1512_89}\\
\includegraphics[max width=\textwidth, alt={}, center]{da0156b6-4133-434e-90d4-ec01b2070afe-023_579_482_1532_776}

Let us displace the bob \(m\) from its equilibrium position by some small angle and then release it. Now locate the ball at an angular position \((\alpha+\theta)\) from vertical as shown in the figure.\\
From the Eqn. :

\[
N_{0 z}=I \beta_{z}
\]

\(-m g l \sin (\alpha+\theta)-m w \cos (\pi-\beta) l \sin (\alpha+\theta)+m w \sin (\pi-\beta) l \cos (\alpha+\theta)=m l^{2} \ddot{\theta}\)\\
or, \(-g(\sin \alpha \cos \theta+\cos \alpha \sin \theta)-w \cos (\pi-\beta)(\sin \alpha \cos \theta+\cos \alpha \sin \theta)+w \sin \beta\)\\
( \(\cos \alpha \cos \theta-\sin \alpha \sin \theta\) )

\[
=l \ddot{\theta}
\]

But for small

\[
\theta, \sin \theta \propto \theta \cos \theta=1
\]

So,

\[
\begin{gathered}
-g(\sin \alpha+\cos \alpha \theta)-w \cos (\pi-\beta)(\sin \alpha+\cos \alpha \theta)+w \sin \beta(\cos \alpha-\sin \alpha \theta) \\
=l \ddot{\theta}
\end{gathered}
\]

or,


\begin{equation*}
(\tan \alpha+\theta)(w \cos \beta-g)+w \sin \beta(1-\tan \alpha \theta)=\frac{l}{\cos \alpha} \ddot{\theta} \tag{2}
\end{equation*}


Solving Eqns (1) and (2) simultaneously we get

\[
-\left(g^{2}-2 w g \cos \beta+w^{2}\right) \theta=l \sqrt{g^{2}+w^{2}-2 w g \cos \beta} \ddot{\theta}
\]

Thus

\[
\ddot{\theta}=-\frac{|\vec{g}-\vec{w}|_{\theta}}{l} \theta
\]

Hence the sought time period \(T=\frac{2 \pi}{\omega_{0}}=2 \pi \sqrt{\frac{l}{|\vec{g}-\vec{w}|}}\)\\
4.31 Obviously the sleeve performs small oscillations in the frame of rotating rod. In the rod's frame let us depict the forces acting on the sleeve along the length of the rod while the sleeve is at a small distance \(x\) towards right from its equilibrium position. The free body diagram of block does not contain Coriolis force, because it is perpendicualr to the length of the rod. From \(F_{x}=m w_{x}\) for the sleeve in the frame of rod


\begin{equation*}
-\kappa x+m \omega^{2} x=m \ddot{x} \tag{1}
\end{equation*}


or, \(\quad \ddot{x}=-\left(\frac{\kappa}{m}-\omega^{2}\right) x\)

Thus the sought time period

\[
T=\frac{2 \pi}{\sqrt{\frac{\kappa}{m}-\omega^{2}}}=0.7 \mathrm{~s}
\]

\begin{center}
\includegraphics[max width=\textwidth, alt={}]{da0156b6-4133-434e-90d4-ec01b2070afe-024_212_623_1278_808}
\end{center}

It is obvious from Eqn (1) that the sleeve will not perform small oscillations if \(\omega \geq \sqrt{\frac{\kappa}{m}} 10 \mathrm{rod} / \mathrm{s}\).\\
4.32 When the bar is about to start sliding along the plank, it experiences the maximum restoring force which is being provided by the limiting friction,\\
Thus

\[
k N=m \omega_{0}^{2} a \quad \text { or, } k m g=m \omega_{0}^{2} a
\]

or,

\[
k=\frac{\omega_{0}^{2} a}{g}=\frac{a}{g}\left(\frac{2 \pi}{T}\right)^{2}=4 \mathrm{~s} .
\]

4.33 The natural angular frequency of a mathematical pendulum equals \(\omega_{0}=\sqrt{g / l}\)\\
(a) We have the solution of S.H.M. equation in angular form :

\[
\theta=\theta_{m} \cos \left(\omega_{0} t+\alpha\right)
\]

If at the initial moment i.e. at \(t=0, \theta=\theta_{m}\) than \(\alpha=0\).\\
Thus the above equation takes the form

\[
\begin{aligned}
\theta & =\theta_{m} \cos \omega_{0} t \\
& =\theta_{m} \cos \sqrt{\frac{g}{l}} t=3^{\circ} \cos \sqrt{\frac{9 \cdot 8}{0 \cdot 8}} t \\
\theta & =3^{\circ} \cos 3 \cdot 5 t
\end{aligned}
\]

Thus\\
(b) The S.H.M. equation in angular form :

\[
\theta=\theta_{m} \sin \left(\omega_{0} t+\alpha\right)
\]

If at the initial moment \(t=0, \theta=0\), then \(\alpha=0\). Then the above equation takes the form

\[
\theta=\theta_{m} \sin \omega_{0} t
\]

Let \(v_{0}\) be the velocity of the lower end of pendulam at \(\theta=0\), then from conserved of mechanical energy of oscillaton\\
or,

\[
\begin{gathered}
E_{\text {mean }}=E_{\text {extreme }} \text { or } \quad T_{\text {mean }}=U_{\text {extrem }} \\
\frac{1}{2} m v_{0}^{2}=m g l\left(1-\cos \theta_{m}\right)
\end{gathered}
\]

Thus

\[
\theta_{m}=\cos ^{-1}\left(1-\frac{v_{0}^{2}}{2 g l}\right)=\cos ^{-1}\left[1-\frac{(0 \cdot 22)^{2}}{2 \times 9 \cdot 8 \times 0 \cdot 8}\right]=4 \cdot 5^{c}
\]

Thus the sought equation becomes

\[
\theta=\theta_{m} \sin \omega_{0} t=4.5^{\circ} \sin 3.5 t
\]

(c) Let \(\theta_{0}\) and \(v_{0}\) be the angular deviation and linear velocity at \(t=0\).

As the mechanical energy of oscillation of the mathematical pendulum is conservation

\[
\frac{1}{2} m v_{0}^{2}+m g l\left(1-\cos \theta_{0}\right)=m g l\left(1-\cos \theta_{m}\right)
\]

or,

\[
\frac{v_{0}^{2}}{2}=g l\left(\cos \theta_{0}-\cos \theta_{m}\right)
\]

Thus \(\theta_{m}=\cos ^{-1}\left\{\cos \theta_{0}-\frac{v_{0}^{2}}{2 g l}\right\}=\cos ^{-1}\left\{\cos 3^{\circ}-\frac{(0.22)^{2}}{2 \times 9.8 \times 0.8}\right\}=5.4^{\circ}\)

Then from \(\theta=5.4^{\circ} \sin (3.5 t+\alpha)\), we see that \(\sin \alpha=\frac{3}{5.4}\) and \(\cos \alpha<0\) because the velovity is directed towards the centre. Thus \(\alpha=\frac{\pi}{2}+1.0\) radians and we get the answer.\\
4.34 While the body \(A\) is at its upper extreme position, the spring is obviously elongated by the amount

\[
\left(a-\frac{m_{1} g}{\kappa}\right) .
\]

If we indicate \(y\)-axis in vertically downward direction, Newton's second law of motion in projection form i.e. \(F_{y}=m w_{y}\) for body \(A\) gives :


\begin{equation*}
m_{1} g+\kappa\left(a-\frac{m_{1} g}{\kappa}\right)=m_{1} \omega^{2} a \quad \text { or, } \kappa\left(a-\frac{m_{1} g}{\kappa}\right)=m_{1}\left(\omega^{2} a-g\right) \tag{1}
\end{equation*}


(Because at any extreme position the magnitude of acceleration of an oscillating body equals \(\omega^{2} a\) and is restoring in nature.)\\
If \(N\) be the normal force exerted by the floor on the body \(B\), while the body \(A\) is at its upper extreme position, from Newton's second law for body \(B\)

\[
N+\kappa\left(a-\frac{m_{1} g}{\kappa}\right)=m_{2} g
\]

or,

\[
N=m_{2} g-\kappa\left(a-\frac{m_{1} g}{\kappa}\right)=m_{2} g-m_{1}\left(\omega^{2} a-g\right)(\text { using Eqn. } 1)
\]

Hence \(N=\left(m_{1}+m_{2}\right) g-m_{1} \omega^{2} a\)\\
When the body \(A\) is at its lower extreme position, the spring is compresed by the distance \(\left(a+\frac{m_{1} g}{\kappa}\right)\).\\
From Newton's second law in projeciton forn i.e. \(F_{y}=m w_{y}\) for body \(A\) at this state:


\begin{equation*}
m_{1} g-\kappa\left(a+\frac{m_{1} g}{\kappa}\right)=m_{1}\left(-\omega^{2} a\right) \text { or, } \kappa\left(a+\frac{m_{1} g}{\kappa}\right)=m_{1}\left(g+\omega^{2} a\right) \tag{3}
\end{equation*}


In this case if \(N^{\prime}\) be the normal force exerted by the floor on the body \(B\), From Newton's second law\\
for body \(B\) we get: \(N^{\prime}=\kappa\left(a+\frac{m_{1} g}{\kappa}\right)+m_{2} g=m_{1}\left(g+\omega^{2} a\right)+m_{2} g\) (using Eqn. 3 )\\
Hence

\[
N^{\prime}=\left(m_{1}+m_{2}\right) g+m_{1} \omega^{2} a
\]

From Newton's third law the magnitude of sought forces are \(N^{\prime}\) and \(N\), respectively.\\
4.35 (a) For the block from Newton's second law in projection form \(F_{y}=m w_{y}\)


\begin{equation*}
N-m g=m \ddot{y} \tag{1}
\end{equation*}


But from

\[
y=a(1-\cos \omega t)
\]

We get


\begin{equation*}
\ddot{y}=\omega^{2} a \cos \omega t \tag{(2}
\end{equation*}


From Eqns (1) and (2)


\begin{equation*}
N=m g\left(1+\frac{\omega^{2} a}{g} \cos \omega t\right) \tag{3}
\end{equation*}


From Newtons's third law the force by which the body \(m\) exerts on the block is directed vertically downward and equls \(N=m g\left(1+\frac{\omega^{2} a}{g} \cos \omega t\right)\)\\
(b) When the body \(m\) starts, falling behind the plank or loosing contact, \(N=0\), (because the normal reaction is the contact force). Thus from Eqn. (3)

Hence

\[
\begin{aligned}
m g\left(1+\frac{\omega^{2} a}{g} \cos \omega t\right) & =0 \quad \text { for some } t . \\
a_{\mathrm{min}}=g / \omega^{2} & \approx 8 \mathrm{~cm} .
\end{aligned}
\]

(c) We observe that the motion takes place about the mean position \(y=a\). At the initial instant \(y=0\). As shown in (b) the normal reaction vanishes at a height ( \(g / \omega^{2}\) ) above the position of equilibrium and the body flies off as a free body. The speed of the body at a distance \(\left(g / \omega^{2}\right)\) from the equilibrium position is \(\omega \sqrt{a^{2}-\left(g / \omega^{2}\right)^{2}}\), so that the condition of the problem gives

\[
\frac{\left[\omega \sqrt{a^{2}-\left(g / \omega^{2}\right)^{2}}\right]^{2}}{2 g}+\frac{g}{\omega^{2}}+a=h
\]

Hence solving the resulting quadratic equation and taking the positive roof,

\[
a=-\frac{g}{\omega^{2}}+\sqrt{\frac{2 h g}{\omega^{2}}} \cong 20 \mathrm{~cm} .
\]

4.36 (a) Let \(y(t)=\) displacement of the body from the end of the unstreched position of the spring (not the equilibrium position). Then

\[
m \ddot{y}=-\kappa y+m g
\]

This equation has the solution of the form

\[
y=A+B \cos (\omega t+\alpha)
\]

if \(\quad-m \omega^{2} B \cos (\omega t+\alpha)=-\kappa[A+B \cos (\omega t+\alpha)]+m g\)\\
Then

\[
\omega^{2}=\frac{\kappa}{m} \quad \text { and } \quad A=\frac{m g}{\kappa}
\]

we have

\[
\begin{array}{r}
y=0 \text { and } \dot{y}=0 \text { at } \\
-\omega B \sin \alpha=0 \\
A+B \cos \alpha=0
\end{array}
\]

Since \(B>0\) and \(A>0\) we must have \(\alpha=\pi\)

\[
B=A=\frac{m g}{\kappa}
\]

and

\[
y=\frac{m g}{\kappa}(1-\cos \omega t)
\]

(b) Tension in the spring is\\
so

\[
\begin{gathered}
T=\kappa y=m g(1-\cos \omega t) \\
T_{\max }=2 m g, T_{\min }=0
\end{gathered}
\]

4.37 In accordance with the problem

\[
\vec{F}=-\alpha m \vec{r}
\]

So,

\[
m(\ddot{x} \vec{i}+\ddot{y} \vec{j})=-\alpha m(x \vec{i}+y \vec{j})
\]

Thus

\[
\ddot{x}=-\alpha x \text { and } \ddot{y}=-\alpha y
\]

Hence the solution of the differential equation


\begin{equation*}
\ddot{x}=-\alpha x \text { becomes } x=a \cos \left(\omega_{0} t+\delta\right), \text { where } \omega_{0}^{2}=\alpha \tag{1}
\end{equation*}


So,


\begin{equation*}
\dot{x}=-a \omega_{0} \sin \left(\omega_{0} t+\alpha\right) \tag{2}
\end{equation*}


From the initial conditions of the problem, \(\mathrm{v}_{\mathrm{x}}=0\) and \(x=r_{0}\) at \(t=0\)\\
So from Eqn. (2) \(\alpha=0\), and Eqn takes the form


\begin{equation*}
x=r_{0} \cos \omega_{0} t \quad \text { so, } \cos \omega_{0} t=x / r_{0} \tag{3}
\end{equation*}


One of the solution of the other differential Eqn \(\ddot{y}=-\alpha y\), becomes


\begin{equation*}
y=a^{\prime} \sin \left(\omega_{0} t+\delta^{\prime}\right), \text { where } \omega_{0}^{2}=\alpha \tag{4}
\end{equation*}


From the initial condition, \(y=0\) at \(t=0\), so \(\delta^{\prime}=0\) and Eqn (4) becomes :\\
\(y=a^{\prime} \sin \omega_{0} t(5)\)\\
Differentiating w.r.t. time we get


\begin{equation*}
\dot{y}=a^{\prime} \omega_{0} \cos \omega_{0} t \tag{6}
\end{equation*}


But from the initial condition of the problem, \(\dot{y}=v_{0}\) at \(t=0\),\\
So, from Eqn (6)

\[
v_{0}=a^{\prime} \omega_{0} \quad \text { or, } \quad a^{\prime}=v_{0} / \omega_{0}
\]

Using it in Eqn (5), we get


\begin{equation*}
y=\frac{v_{0}}{\omega_{0}} \sin \omega_{0} t \text { or } \sin \omega_{0} t=\frac{\omega_{0} y}{v_{0}} \tag{7}
\end{equation*}


Squaring and adding Eqns (3) and (7) we get :

\[
\sin ^{2} \omega_{0} t+\cos ^{2} \omega_{0} t=\frac{\omega_{0}^{2} y^{2}}{v_{0}^{2}}+\frac{x^{2}}{r_{0}^{2}}
\]

or,

\[
\left(\frac{x}{r_{0}}\right)^{2}+\alpha\left(\frac{y}{v_{0}}\right)^{2}=1\left(\text { as } \alpha,=\omega_{0}^{2}\right)
\]

4.38 (a) As the elevator car is a translating non-inertial frame, therefore the body \(m\) will experience an inertial force \(m w\) directed downward in addition to the real forces in the elevator's frame. From the Newton's second law in projection form \(F_{y}=m w_{y}\) for the body in the frame of elevator car:


\begin{equation*}
-\kappa\left(\frac{m g}{\kappa}+y\right)+m g+m w=m \ddot{y} \tag{A}
\end{equation*}


( Because the initial elongation in the spring is \(m g / \kappa\) )\\
so,

\[
m \ddot{y}=-\kappa y+m w=-\kappa\left(y-\frac{m w}{\kappa}\right)
\]

or,


\begin{equation*}
\frac{d^{2}}{d t^{2}}\left(y-\frac{m w}{\kappa}\right)=-\frac{\kappa}{m}\left(y-\frac{m w}{\kappa}\right) \tag{1}
\end{equation*}


Eqn. (1) shows that the motion of the body \(m\) is S.H.M. and its solution becomes


\begin{equation*}
y-\frac{m w}{\kappa}=a \sin \left(\sqrt{\frac{\kappa}{m}} t+\alpha\right) \tag{2}
\end{equation*}


Differentiating Eqn (2) w.r.t. time


\begin{equation*}
\dot{y}=a \sqrt{\frac{\kappa}{m}} \cos \left(\sqrt{\frac{\kappa}{m}} t+\alpha\right) \tag{3}
\end{equation*}


Using the initial condition \(y(0)=0\) in Eqn (2), we get :

\[
a \sin \alpha=-\frac{m w}{\kappa}
\]

and using the other initial condition \(\dot{y}(0)=0\) in Eqn (3)\\
We get

\[
a \sqrt{\frac{\kappa}{m}} \cos \alpha=0
\]

Thus

\[
\alpha=-\alpha / 2 \text { and } a=\frac{m w}{\kappa}
\]

Hence using these values in Eqn (2), we get

\[
y=\frac{m w}{\kappa}\left(1-\cos \sqrt{\frac{\kappa}{m}} t\right)
\]

(b) Proceed up to Eqn.(1). The solution of this differential Eqn be of the form :

\[
y-\frac{m w}{\kappa}=a \sin \left(\sqrt{\frac{\kappa}{m}} t+\delta\right)
\]

or,

\[
y-\frac{\alpha t}{\kappa / m}=a \sin \left(\sqrt{\frac{\kappa}{m}} t+\delta\right)
\]

or,


\begin{equation*}
y-\frac{\alpha t}{\omega_{0}^{2}}=a \sin \left(\omega_{0} t+\delta\right)\left(\text { wher } \omega_{0}=\sqrt{\frac{\kappa}{m}}\right) \tag{4}
\end{equation*}


From the initial condition that at \(t=0, y(0)=0\), so \(0=a \sin \delta\) or \(\delta=0\)\\
Thus Eqn.(4) takes the from : \(y-\frac{\alpha t}{\omega_{0}^{2}}=a \sin \omega_{0} t\)

Differentiating Eqn. (5) we get : \(\dot{y}-\frac{\alpha}{\omega_{0}^{2}}=a \omega_{0} \cos \omega t\)

But from the other initial condition \(\dot{y}(0)=0\) at \(t=0\).\\
So, from Eqn.(6)

\[
-\frac{\alpha}{\omega_{0}^{2}}=a \omega_{0} \quad \text { or } \quad a=-\alpha / \omega_{0}^{3}
\]

Putting the value of \(a\) in Eqn. (5), we get the sought \(y(t)\). i.e.

\[
y-\frac{\alpha t}{\omega_{0}^{2}}=-\frac{\alpha}{\omega_{0}^{3}} \sin \omega_{0} t \text { or } y=\frac{\alpha}{\omega_{0}^{3}}\left(\omega_{0} t-\sin \omega_{0} t\right)
\]

4.39 There is an important difference between a rubber cord or steel coire and a spring. A spring can be pulled or compressed and in both cases, obey's Hooke's law. But a rubber cord becomes loose when one tries to compress it and does not then obey Hooke's law. Thus if we suspend a body by a rubber cord it stretches by a distance \(m g / \kappa\) in reaching the equilibrium configuration. If we further strech it by a distance \(\Delta h\) it will execute harmonic oscillations when released if \(\Delta h \leq m g / \kappa\) because only in this case will the cord remain taut and obey Hooke's law.\\
Thus

\[
\Delta h_{\max }=m g / \kappa
\]

The energy of oscillation in this case is

\[
\frac{1}{2} \kappa\left(\Delta h_{\max }\right)^{2}=\frac{-1}{2} \frac{m^{2} g^{2}}{\kappa}
\]

4.40 As the pan is of negligible mass, there is no loss of kinetic energy even though the collision is inelastic. The mechanical energy of the body \(m\) in the field generated by the joint action of both the gravity force and the elastic force is conserved i.e. \(\Delta E=0\). During the motion of the body \(m\) from the initial to the final (position of maximum compression of the spring) position \(\Delta T=0\), and therefore \(\Delta U=\Delta U_{g r}+\Delta U_{s p}=0\)\\
or

\[
-m g(h+x)+\frac{1}{2} \kappa x^{2}=0
\]

On solving the quadratic equation :

\[
x=\frac{m g}{\kappa} \pm \sqrt{\frac{m^{2} g^{2}}{\kappa^{2}}+\frac{2 m g h}{\kappa}}
\]

As minus sign is not acceptable

\[
x=\frac{m g}{\kappa}+\sqrt{\frac{m^{2} g^{2}}{\kappa^{2}}+\frac{2 m g h}{\kappa}}
\]

If the body \(m\) were at rest on the spring, the corresponding position of \(m\) will be its equilibrium position and at this position the resultant force on the body \(m\) will be zero. Therefore the equilibrium compression \(\Delta x\) (say) due to the body \(m\) will be given by

\[
\kappa \Delta x=m g \text { or } \Delta x=m g / \kappa
\]

Therefore seperation between the equilibrium position and one of the extreme position i.e. the sought amplitude

\[
a=x-\Delta x=\sqrt{\frac{m^{2} g^{2}}{\kappa^{2}}+\frac{2 m g h}{\kappa}}
\]

The mechanical energy of oscillation which is conserved equals \(E=U_{\text {extreme }}\),because at the extreme position kinetic energy becomes zero.\\
Although the weight of body \(m\) is a conservative force , it is not restoring in this problem, hence \(U_{\text {extreme }}\) is only concerned with the spring force. Therefore

\[
E=U_{\text {extreme }}=\frac{1}{2} \kappa a^{2}=m g h+\frac{m^{2} g^{2}}{2 \kappa}
\]

4.41 Unlike the previous \((4.40)\) problem the kinetic energy of body \(m\) decreases due to the perfectly inelastic collision with the pan.Obviously the body \(m\) comes to strike the pan with velocity \(v_{0}=\sqrt{2 g h}\).If \(v\) be the common velocity of the "body \(m+\) pan" system due to the collision then from the conservation of linear momentum\\
or


\begin{gather*}
m \mathrm{v}_{0}=(M+m) \mathrm{v} \\
\mathrm{v}=\frac{m \mathrm{v}_{0}}{(M+m)}=\frac{m \sqrt{2 g h}}{(M+m)} \tag{1)}
\end{gather*}


At the moment the body \(m\) strikes the pan, the spring is compressed due to the weight of 'he pan by the amount \(M g / \kappa\). If \(l\) be the further compression of the spring due to the verocity acquired by the "pan - body \(m\) " system, then from the conservation of mechanical energy of the said system in the field generatad by the joint action of both the gravity and spring forces


\begin{align*}
& \qquad \frac{1}{2}(M+m) \mathrm{v}^{2}+(M+m) g l=\frac{1}{2} \kappa\left(\frac{M g}{\kappa}+l\right)^{2}-\frac{1}{2} \kappa\left(\frac{M g}{\kappa}\right)^{2} \\
& \text { or, } \frac{1}{2}(M+m) \frac{m^{2} 2 g h}{(M+m)}+(M+m) g l=\frac{1}{2} \kappa\left(\frac{M g}{\kappa}\right)^{2}+\frac{1}{2} \kappa l^{2}+M g l-\frac{1}{2} \kappa\left(\frac{M g}{\kappa}\right)^{2}  \tag{Using1}\\
& \text { or, } \quad \frac{1}{2} \kappa l^{2}-m g l-\frac{m^{; 2} g h}{(m+M)}=0 \\
& \text { Thus } l=\frac{m g \pm \sqrt{m^{2} g^{2}+\frac{2 \kappa g h m^{2}}{M+m}}}{\kappa}
\end{align*}


As minus sign is not acceptable

\[
l=\frac{m g}{\kappa}+\frac{1}{\kappa} \sqrt{m^{2} g^{2}+\frac{2 \kappa m n^{2} g h}{(M+m)}}
\]

If the oscillating "pan + body \(m\) " system were at rest it correspond to their equilbrium position i.e. the spring were compressed by \(\frac{(M+m) g}{\kappa}\) therefore the amplitude of oscillation

\[
a=l-\frac{m g}{\kappa}=\frac{m g}{\kappa} \sqrt{1+\frac{2 h \kappa}{m g}}
\]

The mechanical energy of oscillation which is only conserved with the restoring forces becomes \(E=U_{\text {extreme }}=\frac{1}{2} \kappa a^{2}\) (Because spring force is the only restoring force not the weight of the body)

Alternately

\[
E=T_{\text {mean }}=\frac{1}{2}(M+m) a^{2} \omega^{2}
\]

thus

\[
E=\frac{1}{2}(M+m) a^{2}\left(\frac{\kappa}{M+m}\right)=\frac{1}{2} \kappa a^{2}
\]

4.42 We have \(\vec{F}=a(\dot{y} \vec{i}-\dot{x} \vec{j})\)\\
or,

\[
m(\ddot{x} \overrightarrow{i+} \ddot{y} \vec{j})=a(\dot{y} \vec{i}-\dot{x} \vec{j})
\]

So,


\begin{equation*}
m \ddot{x}=a \dot{y} \text { and } m \ddot{y}=-a \dot{x} \tag{1}
\end{equation*}


From the initial condition,at \(t=0, \dot{x}=0\) and \(y=0\)\\
So, integrating Eqn, \(m \ddot{x}=a \dot{y}\)\\
we get


\begin{equation*}
:=a y \text { or } \dot{x}=\frac{a}{m} y \tag{2}
\end{equation*}


Using Eqn (2) in the Eqn \(m:=-a \dot{x}\), we get


\begin{equation*}
m \ddot{y}=-\frac{a^{2}}{m} y \quad \text { or } \ddot{y}=-\left(\frac{a}{m}\right)^{2} y \tag{3}
\end{equation*}


one of the solution of differential Eqn (3) is

\[
y=A \sin \left(\omega_{0} t+\alpha\right), \text { where } \omega_{0}=a / m
\]

As at \(t=0, y=0\), so the solution takes the form \(y=A \sin \omega_{0} t\)\\
On differentiating w.r.t. time \(\dot{y}=A \omega_{0} \cos \omega_{0} t\)\\
From the initial condition of the problem, at \(t=.0, \dot{y}=v_{0}\)\\
So,


\begin{gather*}
\mathrm{v}_{0}=A \omega_{0} \text { or } A=\mathrm{v}_{0} / \omega_{0} \\
y=\left(\mathrm{v}_{0} / \omega_{0}\right) \sin \omega_{0} t \tag{4}
\end{gather*}


Thus from (2) \(\dot{x}=v_{0} \sin \omega_{0} t\) so integrating


\begin{equation*}
x=B-\frac{\mathrm{v}_{0}}{\omega_{0}} \cos \omega_{0} t \tag{5}
\end{equation*}


On using

\[
x=0 \text { at } t=0, B=\frac{\mathrm{v}_{0}}{\omega_{0}}
\]

Hence finally


\begin{equation*}
x=\frac{v_{0}}{\omega_{0}}\left(1-\cos \omega_{0} t\right) \tag{6}
\end{equation*}


Hence from Eqns (4) and (6) we get

\[
\left[x-\left(v_{0} / \omega_{0}\right)\right]^{2}+y^{2}=\left(v_{0} / \omega_{0}\right)^{2}
\]

which is the equation of a circle of radius \(\left(v_{0} / \omega_{0}\right)\) with the centre at the point \(x_{0}=v_{0} / \omega_{0}, y_{0}=0\)\\
4.43 If water has frozen, the system consisting of the light rod and the frozen water in the hollow sphere constitute a compound (physical) pendulum to a very good approximation because we can take the whole system to be rigid. For such systems the time period is given by\\
\(T_{1}=2 \pi \sqrt{\frac{l}{g}} \sqrt{1+\frac{k^{2}}{l^{2}}} \quad\) where \(\quad k^{2}=\frac{2}{5} R^{2}\) is the radius of gyration of the sphere.\\
The situation is different when water is unfrozen. When dissipative forces (viscosity) are neglected, we are dealing with ideal fluids. Such fluids instantaneously respond to (unbalanced) internal stresses. Suppose the sphere with liquid water actually executes small rigid oscillations. Then the portion of the fluid above the centre of the sphere will have a greater acceleration than the portion below the centre because the linear acceleration of any element is in this case, equal to angular acceleration of the element multiplied by the distance of the element from the centre of suspension (Recall that we are considering small oscillations). Then, as is obvious in a frame moving with the centre of mass, there will appear an unbalanced couple (not negated by any pseudoforces) which will cause the fluid to move rotationally so as to destroy differences in acceleration. Thus for this case of ideal fluids the pendulum must move in such a way that the elements of the fluid all undergo the same acceleration. This implies that we have a simple (mathematical) pendulum with the time period :

\[
T_{0}=2 \pi \sqrt{\frac{l}{g}}
\]

Thus

\[
T_{1}=T_{0} \sqrt{1+\frac{2}{5}\left(\frac{R}{l}\right)^{2}}
\]

(One expects that a liquid with very small viscosity will have a time period close \(T_{0}\) while one with high viscosity will have a time period closer to \(T_{1}\).)\\
4.44 Let us locate the rod at the position when it makes an angle \(\theta\) from the vertical. In this problem both, the gravity and spring forces are restoring conservative forces, thus from the conservation of mechanical energy of oscillation of the oscillating system :

\[
\frac{1}{2} \frac{m l^{2}}{3}(\dot{\theta})^{2}+m g \frac{l}{2}(1-\cos \theta)+\frac{1}{2} \kappa(l \theta)^{2}=\text { constant }
\]

Differentiating w.r.t. time, we get :

\[
\frac{1}{2} \frac{m l^{2}}{3} 2 \dot{\theta} \ddot{\theta}+\frac{m g l}{2} \sin \theta \dot{\theta}+\frac{1}{2} \kappa l^{2} 2 \theta \dot{\theta}=0
\]

Thus for very small \(\theta\)

Hence,

\[
\begin{aligned}
& \ddot{\theta}=-\frac{3 g}{2 l}\left(1+\frac{\kappa l}{m g}\right) \theta \\
& \omega_{0}=\sqrt{\frac{3 g}{2 l}\left(1+\frac{\kappa l}{m g}\right)}
\end{aligned}
\]

4.45 (a) Let us locate the system when the threads are deviated through an angle \(\alpha^{\prime}<\alpha\), during the oscillations of the system (Fig.). From the conservation of mechanical energy of the system :


\begin{equation*}
\frac{1}{2} \frac{m L^{2}}{12} \dot{\theta}^{2}+m g l\left(1-\cos \alpha^{\prime}\right)=\text { constant } \tag{1}
\end{equation*}


Where \(L\) is the length of the rod, \(\theta\) is the angular deviation of the rod from its equilibrium position i.e. \(\theta=0\).\\
Differentiating Eqn. (1) w.r.t. time


\begin{equation*}
\frac{1}{2} \frac{m L^{2}}{12} 2 \dot{\theta} \ddot{\theta}+m g l \sin \alpha^{\prime} \dot{\alpha}^{\prime}=0 \tag{2}
\end{equation*}


\begin{center}
\includegraphics[max width=\textwidth, alt={}]{da0156b6-4133-434e-90d4-ec01b2070afe-034_370_383_197_1022}
\end{center}

So, \(\quad \frac{L^{2}}{12} \dot{\theta} \ddot{\theta}+g l \alpha^{\prime} \dot{\alpha}^{\prime}=O\left(\right.\) for small \(\left.\alpha^{\prime}, \sin \alpha^{\prime} \cong \alpha^{\prime}\right)\)

But from the Fig.

So,

\[
\begin{gathered}
\frac{L}{2} \theta=l \alpha^{\prime} \text { or } \alpha^{\prime}=\frac{L}{2 l} \theta \\
\dot{\alpha}^{\prime}=\frac{L}{2 l} \dot{\theta}
\end{gathered}
\]

Putting these values of \(\alpha^{\prime}\) and \(\frac{d \alpha^{\prime}}{d t}\) in Eqn. (2) we get

\[
\frac{d^{2} \theta}{d t^{2}}=-\frac{3 g}{l} \theta
\]

Thus the sought time period

\[
T=\frac{2 \pi}{\omega_{0}}=2 \pi \sqrt{\frac{l}{3 g}}
\]

(b) The sought oscillation energy

\[
\begin{gathered}
E=U_{\text {extreme }}=m g l(1-\cos \alpha)=m g l 2 \sin ^{2} \frac{\alpha}{2} \\
\cong m g l 2 \frac{\alpha^{2}}{4}=\frac{m g l \alpha^{2}}{2}(\text { because for small angle } \sin \theta \cong \theta)
\end{gathered}
\]

4.46 The \(\mathcal{K}\). of the disc is \(\frac{1}{2} I \dot{\varphi}^{2}=\frac{1}{2}\left(\frac{m R^{2}}{2}\right) \dot{\varphi}^{2}=\frac{1}{4} m R^{2} \dot{\varphi}^{2}\)

The torsional potential energy is \(\frac{1}{2} k \varphi^{2}\). Thus the total energy is :

\[
\frac{1}{4} m R^{2} \dot{\varphi}^{2}+\frac{1}{2} k \varphi^{2}=\frac{1}{4} m R^{2} \dot{\varphi}_{0}^{2}+\frac{1}{2} k \varphi_{0}^{2}
\]

By definition of the amplitude \(\varphi_{m}, \dot{\varphi}=0\) when \(\varphi=\varphi_{m}\). Thus total energy is

\[
\begin{gathered}
\frac{1}{2} k \varphi_{m}^{2}=\frac{1}{4} m R^{2} \dot{\varphi}_{0}^{2}+\frac{1}{2} k \varphi_{0}^{2} \\
\varphi_{m}=\varphi_{0} \sqrt{1+\frac{m R^{2}}{2 k} \frac{\varphi_{0}^{2}}{\varphi_{0}^{2}}}
\end{gathered}
\]

or\\
4.47 Moment of inertia of the rod equals \(\frac{m l^{2}}{3}\) about its one end and perpendicular to its length Thus rotational kinetic energy of the rod \(=\frac{1}{2}\left(\frac{m l^{2}}{3}\right) \dot{\theta}^{2}=\frac{m l^{2}}{6} \dot{\theta}^{2}\)\\
when the rod is displaced by an angle \(\theta\) its C.G. goes up by a distance \(\frac{l}{2}(1-\cos \theta) \equiv \frac{l \theta^{2}}{4}\) for small \(\theta\).\\
Thus the P.E. becomes : \(m g \frac{l \theta^{2}}{4}\)\\
As the mechanical energy of oscillation of the rod is conserved.

\[
\frac{1}{2}\left(\frac{m l^{2}}{3}\right) \dot{\theta}^{2}+\frac{1}{2}\left(\frac{m g l}{2}\right) \theta^{2}=\text { Constant }
\]

on differentiating w.r.t. time and for the simplifies we get : \(\ddot{\theta}=-\frac{3 g}{2 l} \theta\) for small \(\theta\).\\
we see that the angular frequency \(\omega\) is

\[
=\sqrt{3 g / 2 l}
\]

we write the general solution of the angular oscillation as :

\[
\theta=A \cos \omega t+B \sin \omega t
\]

But

\[
\theta=\theta_{0} \quad \text { at } t=0, \quad \text { so } A=\theta_{0}
\]

and

\[
\theta=\theta_{0} \quad \text { at } t=0, \text { so }
\]

\[
B=\dot{\theta}_{0} / \omega
\]

Thus

\[
\theta=\theta_{0} \cos \omega t+\frac{\theta_{0}}{\omega} \sin \omega t
\]

Thus the KE. of the rod

\[
\begin{gathered}
T=\frac{m l^{2}}{6} \dot{\theta}^{2}=\left[-\omega \theta_{0} \sin \omega t+\hat{0}_{0} \cos \omega t\right]^{2} \\
=\frac{m l^{2}}{6}\left[\dot{\theta}_{0}^{2} \cos ^{2} \omega t+\omega^{2} \theta_{0}^{2} \sin ^{2} \omega t-2 \omega 0_{0} \dot{0}_{0} \sin \omega t \cos \omega t\right]
\end{gathered}
\]

On averaging over one time period the last term vanishes and \(\left\langle\sin ^{2} \omega t\right\rangle=\left\langle\cos ^{2} \omega t\right\rangle=1 / 2\). Thus

\[
<T>=\frac{1}{12} m l^{2} \dot{\theta}_{0}^{2}+\frac{1}{8} m g l^{2} 0_{0}^{2}\left(\text { where } \omega^{2}=3 g / 2 l\right)
\]

4.48 Let \(l=\) distance between the C.G. ( \(C\) ) of the pendulum and its point of suspension \(C\) Originally the pendulum is in inverted position and its C.G. is above \(O\). When it falls to the normal (stable) position of equilibrium its C.G. has fallen by a distance \(2 l\). In the equilibriun position the total energy is equal to K.E. \(=\frac{1}{2} I \omega^{2}\) and we have from energy conservation :

\[
\frac{1}{2} I \omega^{2}=m g 2 l \text { or } I=\frac{4 m g l}{\omega^{2}}
\]

Angular frequency of oscillation for a physical pendulum is given by \(\omega_{0}^{2}=m g l / I\)

Thus

\[
T=2 \pi \sqrt{\frac{I}{m g l}}=2 \pi \sqrt{\frac{4 m g l / \omega^{2}}{m g l}}=\frac{4 \pi}{3}
\]

4.49 Let, moment of inertia of the pendulum, about the axis, concerned is \(I\), then writing \(N_{z}=I \beta_{z}\), for the pendulum,


\begin{equation*}
-m g x \sin a \theta=I \dot{\theta} \quad \text { or, } \quad \dot{\theta}=-\frac{m g x}{I} \theta \tag{(For small \(\theta\) )}
\end{equation*}


which is the required equation for S.H.M. So, the frequency of oscillation,


\begin{equation*}
\omega_{1}=\sqrt{\frac{M g x}{I}} \quad \text { or, } \quad x=\frac{I}{M g} \sqrt{\omega_{1}^{2}} \tag{1}
\end{equation*}


Now, when the mass \(m\) is attached to the pendulum, at a distance \(l\) below the oscillating axis,\\
or,

\[
\begin{aligned}
& -M g x \sin \theta^{\prime}-m g l \sin \theta^{\prime}=\left(I+m l^{2}\right) \frac{d^{2} \theta^{\prime}}{d t^{2}} \\
& -\frac{g(M x+m l)}{\left(I+m l^{2}\right)} \theta^{\prime}=\frac{d^{2} \theta}{d t^{2}},(\text { For small } \theta)
\end{aligned}
\]

which is again the equation of S.H.M., So, the new frequency,


\begin{equation*}
\omega_{2}=\sqrt{\frac{g(M x+m l)}{\left(I+m l^{2}\right)}} \tag{2}
\end{equation*}


Solving Eqns. (1) and (2),\\
or,

\[
\begin{aligned}
& \omega_{2}=\sqrt{\frac{g\left((I / g) \omega_{1}^{2}+m l\right)}{\left(I+m l^{2}\right)}} \\
& \omega_{2}^{2}=\frac{I \omega_{1}^{2}+m g l}{I+m l^{2}} \\
& I\left(\omega_{2}^{2}-\omega_{1}^{2}\right)=m g l-m \omega_{2}^{2} l^{2}
\end{aligned}
\]

or,

\[
I=m l^{2}\left(\omega_{2}^{2}-g / l\right) /\left(\omega_{1}^{2}-\omega_{2}^{2}\right)=0.8 \mathrm{~g} \cdot \mathrm{~m}^{2}
\]

4.50 When the two pendulums are joined rigidly and set to oscillate, each exert wrques on the other, these torques are equal and opposite. We write the law of motion for the two pendulums as

\[
\begin{aligned}
& I_{1} \ddot{\theta}=-\omega_{1}^{2} I_{1} \theta+G \\
& I_{2} \ddot{\theta}=-\omega_{2}^{2} I_{2} \theta-G
\end{aligned}
\]

where \(\boldsymbol{\pm} \boldsymbol{G}\) is the torque of mutual interactions. We have written the restoring forces on each pendulum in the absence of the other as \(-\omega_{1}^{2} I_{1} \theta\) and \(-\omega_{2}^{2} I_{2} \theta\) respectively. Then

\[
\begin{gathered}
\ddot{\theta}=-\frac{I_{1} \omega_{1}^{2}+I_{2} \omega_{2}^{2}}{I_{1}+I_{2}} \theta=-\omega^{2} \theta \\
\omega=\sqrt{\frac{I_{1} \omega_{1}^{2}+I_{2} \omega_{2}^{2}}{I_{1}+I_{2}}}
\end{gathered}
\]

4.51 Let us locate the rod when it is at small angular position \(\theta\) relative to its equilibrium position.

If \(a\) be the sought distance, then from the conservation of mechanical energy of oscillation

\[
m g a(1-\cos \theta)+\frac{1}{2} I_{O O},(\dot{\theta})^{2}=\text { constant }
\]

Differentiating w.r.t, time we get :

\[
m g a \sin \theta \dot{\theta}+\frac{1}{2} I_{O O^{\prime}} 2 \dot{\theta} \ddot{\theta}=0
\]

But \(\quad I_{O O^{\prime}}=\frac{m l^{2}}{12}+m a^{2}\) and for small \(\theta, \sin \theta=\theta\), we get

\[
\ddot{\theta}=-\left(\frac{g a}{\frac{l^{2}}{12}+a^{2}}\right) \theta
\]

Hence the time period of one full osscillation becomes

\[
T=2 \pi \sqrt{\frac{l^{2}}{\frac{12}{a g}+a^{2}}} \text { or } T^{2}=\frac{4 \pi^{2}}{g}\left(\frac{l^{2}}{12 a}+a\right)
\]

For

\[
T_{\min } \text {, obviously } \frac{d}{d a}\left(\frac{l^{2}}{12 a}+a\right)=0
\]

So,

\[
-\frac{l^{2}}{12 a^{2}}+1=0 \quad \text { or } \quad a=\frac{l}{2 \sqrt{3}}
\]

Hence

\[
T_{\min }=2 \pi \sqrt{\frac{l}{g \sqrt{3}}}
\]

4.52 Consider the moment of inertia of the triangular plate about \(A B\).

\[
\begin{aligned}
& I=\iint x^{2} d m=\iint x^{2} \rho d x d y \\
& =\int_{0}^{h} x^{2} \rho d x \frac{h-x}{h} \cdot \frac{2 h}{\sqrt{3}}=\int_{0}^{h} x^{2} \frac{2 \rho}{\sqrt{3}}(h-x) d x \\
& =\frac{2 \rho}{\sqrt{3}}\left(\frac{h^{4}}{3}-\frac{h^{4}}{4}\right)=\frac{\rho h^{4}}{6 \sqrt{3}}=\frac{m h^{2}}{6}
\end{aligned}
\]

On using the area of the triangle \(\triangle A B C=\frac{h^{2}}{\sqrt{3}}\) and \(m=\rho \Delta\).\\
\includegraphics[max width=\textwidth, alt={}, center]{da0156b6-4133-434e-90d4-ec01b2070afe-038_295_356_181_1056}

Thus KE.

\[
=\frac{1}{2} \frac{m h^{2}}{6} \dot{\theta}^{2}
\]

P.E.

\[
=m g \frac{h}{3}(1-\cos \theta)=\frac{1}{2} m g h \frac{\theta^{2}}{3}
\]

Here \(\theta\) is the angle that the instantaneous plane of the plate makes with the equilibrium position which is vertical. (The plate rotates as a rigid body)

Thus

\[
E=\frac{1}{2} \frac{m h^{2}}{6} \dot{\theta}^{2}+\frac{1}{2} \frac{m g h}{3} \theta^{2}
\]

Hence

\[
\omega^{2}=\frac{2 g}{h}=\frac{m g h}{3} / \frac{m h^{2}}{6}
\]

So

\[
T=2 \pi \sqrt{\frac{h}{2 g}}=\pi \sqrt{\frac{2 h}{g}} . \text { and } \quad l_{\text {reduced }}=h / 2 .
\]

Let us go to the rotating frame, in which the disc is stationary. In this frame the rod is subjected to coriolis and centrifugal forces, \(\mathbf{F}_{c o r}\) and \(\mathbf{F}_{c f}\), where

\[
\mathbf{F}_{c o r}=\int 2 d m\left(\mathbf{v}^{\prime} \times \vec{\omega}_{0}\right) \text { and } \mathbf{F}_{c f}=\int d m \omega_{0}^{2} \mathbf{r},
\]

where \(r\) is the position of an elemental mass of the rod (Fig.) with respect to point \(O\) (disc's centre) and

\[
\mathbf{v}^{\prime}=\frac{d \mathbf{r}^{\prime}}{d t}
\]

As

\[
\mathbf{r}=\mathbf{O P}=\mathbf{O A}+\mathbf{A P}
\]

So,

\[
\frac{d \mathbf{r}}{d t}=\frac{d(\mathbf{A} \mathbf{P})}{d t}=\mathbf{v}^{\prime}
\]

\begin{figure}[h]
\begin{center}
  \includegraphics[alt={},max width=\textwidth]{da0156b6-4133-434e-90d4-ec01b2070afe-038_417_448_1207_988}
\captionsetup{labelformat=empty}
\caption{(as OA is constant)}
\end{center}
\end{figure}

As the rod is vibrating transversely, so \(\mathbf{v}^{\prime}\) is directed perpendicular to the length of the rod. Hence \(2 \operatorname{dim}\left(\mathbf{v}^{\prime} \times \vec{\omega}\right)\) for each elemental mass of the rod is directed along \(\mathbf{P A}\). Therefore the net torque of coriolis about \(A\) becomes zero. The not torque of centrifulgal force about point A :

Now,

\[
\vec{\tau}_{c f(A)}=\int \mathbf{A} \mathbf{P} \times d m \omega_{0}^{2} \mathbf{r}=\int \mathbf{A} \mathbf{P} \times\left(\frac{m}{l}\right) d s \omega_{0}^{2}(\mathbf{O} \mathbf{A}+\mathbf{A} \mathbf{P})
\]

\[
\begin{aligned}
& =\int \mathbf{A P} \times\left(\frac{m}{l} d s\right) \omega_{0}^{2} \mathbf{O A}=\int \frac{m}{l} d s \omega_{0}^{2} s a \sin \theta(-\mathbf{k}) \\
& \quad=\frac{m}{l} \omega_{0}^{2} a \sin \theta(-\mathbf{k}) \int_{0}^{t} s d s=m \omega_{0}^{2} a \frac{l}{2} \sin \theta(-\mathbf{k})
\end{aligned}
\]

So,

\[
\tau_{c f(Z)}=\vec{\tau}_{c f(A)} \cdot \mathbf{k}=-m \omega_{0}^{2} a \frac{l}{2} \sin \theta
\]

According to the equation of rotational dynamics : \(\tau_{A(Z)}=l_{A} \alpha_{Z}\)\\
or,

\[
-m \omega_{0}^{2} a \frac{l}{2} \sin \theta=\frac{m l^{2}}{3} \ddot{\theta}
\]

or,

\[
\ddot{\theta}=-\frac{3}{2} \frac{\omega_{0}^{2} a}{l} \sin \theta
\]

Thus, for small \(\theta\),

\[
\ddot{\theta} \approx-\frac{3}{2} \frac{\omega^{2} a}{2 l} \theta
\]

This implies that the frequency \(\omega_{0}\) of oscillation is \(\omega_{0}=\sqrt{\frac{3 \omega^{2} a}{2 l}}\)\\
4.54 The physical system consists with a pulley and the block. Choosing an intertial frame, let us direct the \(x\)-axis as shown in the figure.\\
\includegraphics[max width=\textwidth, alt={}, center]{da0156b6-4133-434e-90d4-ec01b2070afe-039_679_736_992_329}

Initially the system is in equilibrium position. Now from the condition of translation equilibrium for the block


\begin{equation*}
T_{0}=m g \tag{1}
\end{equation*}


Similarly for the rotational equilibrium of the pulley

\[
\kappa \Delta / R=T_{0} R
\]

or.


\begin{equation*}
T_{0}=\kappa \Delta l \tag{2}
\end{equation*}


from Eqns. (1) and (2)


\begin{equation*}
\Delta l=\frac{m g}{\kappa} \tag{3}
\end{equation*}


Now let us disturb the equilibrium of the system no matter in which way to analyse its motion. At an arbitrary position shown in the figure, from Newton's second law of motion for the block


\begin{gather*}
F_{x}=m w_{x} \\
m g-T=m w=m \ddot{x} \tag{4}
\end{gather*}


Similarly for the pulley

\[
N_{z}=I \beta_{z}
\]

But


\begin{equation*}
w=\beta R \quad \text { or, } \quad \ddot{x}=R \ddot{\theta} \tag{5}
\end{equation*}


from (5) and (6)


\begin{equation*}
T R-\kappa(\Delta l+x) R=\frac{I}{R} \ddot{x} \tag{6}
\end{equation*}


Solving (4) and (7) using the initial condition of the problem\\
or,

\[
\begin{gathered}
-\kappa R x=\left(m R+\frac{I}{R}\right) \ddot{x} \\
\ddot{x}=-\left(\frac{\kappa}{m+\frac{l}{R^{2}}}\right) x
\end{gathered}
\]

Hence the sought time period, \(T=\frac{2 \pi}{\omega_{0}}=2 \pi \sqrt{\frac{m+l / R^{2}}{\kappa}}\)\\
Note : we may solve this problem by using the conservation of mechanical energy also\\
4.55 At the equilbrium position, \(N_{o z}=0\) (Net torque about 0 )

So,


\begin{equation*}
m_{A} g R-m g R \sin \alpha=0 \quad \text { or } \quad m_{A}=m \sin \alpha \tag{1}
\end{equation*}


From the equation of rotational dynamics of a solid body about the stationary axis (say \(z\)-axis) of rotation i.e. from \(N_{z}=I \beta_{z}\)\\
when the pulley is rotated by the small angular displacement \(\theta\) in clockwise sense relative to the equilibrium position (Fig.), we get :

\[
\begin{array}{rl}
m_{A} & g R-m g R \sin (\alpha+\theta) \\
& =\left[\frac{M R^{2}}{2}+m R^{2}+m_{A} R^{2}\right] \ddot{\theta}
\end{array}
\]

Using Eqn. (1)\\
\(m g \sin \alpha-m g(\sin \alpha \cos \theta+\cos \alpha \sin \theta)\)\\
\includegraphics[max width=\textwidth, alt={}, center]{da0156b6-4133-434e-90d4-ec01b2070afe-040_363_267_1530_1007}

\[
=\left\{\frac{M R+2 m(1+\sin \alpha) R}{2}\right\} \ddot{\theta}
\]

But for small \(\theta\), we may write \(\cos \theta \cong 1\) and \(\sin \theta \leqq \theta\)\\
Thus we have

\[
m g \sin \alpha-m g(\sin \alpha+\cos \alpha \theta)=\frac{\{M R+2 m(1+\sin \alpha) R\}}{2} \ddot{\theta}
\]

Hence,

\[
\ddot{\theta}=-\frac{2 m g \cos \alpha}{[M R+2 m(1+\sin \alpha) R]} \theta
\]

Hence the sought angular frequency \(\omega_{0}=\sqrt{\frac{2 m g \cos \alpha}{M R+2 m R(1+\sin \alpha)}}\)\\
4.56 Let us locate solid cylinder when it is displaced from its stable equilibrium position by the small angle \(\theta\) during its oscillations (Fig.). If \(v_{c}\) be the instantaneous speed of the C.M. ( \(C\) ) of the solid cylinder which is in pure rolling, then its angular velocity about its own centre \(C\) is


\begin{equation*}
\omega=v_{c} / r \tag{1}
\end{equation*}


\begin{center}
\includegraphics[max width=\textwidth, alt={}]{da0156b6-4133-434e-90d4-ec01b2070afe-041_256_391_472_953}
\end{center}

Since \(C\) moves in a circle of radius ( \(R-r\) ), the speed of \(C\) at the same moment can be written as


\begin{equation*}
v_{c}=\dot{\theta}(R-r) \tag{2}
\end{equation*}


Thus from Eqns (1) and (2)


\begin{equation*}
\omega=\dot{\theta} \frac{(R-r)}{r} \tag{3}
\end{equation*}


As the mechanical energy of oscillation of the solid cylinder is conserved, i.e. \(E=T+U=\) constant

So, \(\quad \frac{1}{2} m v_{\mathrm{c}}^{2}+\frac{1}{2} I_{c} \omega^{2}+m g(R-r)(1-\cos \theta)=\) constant\\
(Where \(m\) is the mass of solid cylinder and \(I_{c}\) is the moment of inertia of the solid cylinder about an axis passing through its C.M. ( \(C\) ) and perpendicular to the plane of Fig. of solid cylinder)\\
or, \(\quad \frac{1}{2} m \omega^{2} r^{2}+\frac{1}{2} \frac{m r^{2}}{2} \omega^{2}+m g(R-r)(1-\cos \theta)=\) constant \(\quad(\) using \(\quad\) Eqn \(\quad(1) \quad\) and \(\left.I_{c}=m r, D\right)\)

\[
\frac{3}{4} r^{2}(\dot{\theta})^{2} \frac{(R-r)^{2}}{r^{2}}+g(R-r)(1-\cos \theta)=\text { constant, (using Eqn. 3) }
\]

Differentiating w.r.t. time

\[
\frac{3}{4}(R-r) 2 \dot{\theta} \ddot{\theta}+g \sin \theta \dot{\theta}=0
\]

So,

\[
\ddot{\theta}=-\frac{2 g}{3(R-r)} \theta,(\text { because fer small } \theta, \sin \theta \cong \theta)
\]

Thus

\[
\omega_{0}=\sqrt{\frac{2 g}{3(R-r)}}
\]

Hence the sought time period

\[
T=\frac{2 \pi}{\omega_{0}}=2 \pi \sqrt{\frac{3(R-r)}{2 g}}
\]

4.57 Let \(\kappa_{1}\) and \(\kappa_{2}\) be the spring constant of left and right sides springs. As the rolling of the solid cylinder is pure its lowest point becomes the instanteneous centre of rotation. If \(\theta\) be the small angular displacement of its upper most point relative to its equilibrium position, the deformation of each spring becomes ( \(2 R \theta\) ). Since the mechanical energy of oscillation of the solid cylinder is conserved, \(E=T+U=\) constant\\
i.e.

\[
\frac{1}{2} I_{p}(\dot{\theta})^{2}+\frac{1}{2} \kappa_{1}(2 R \theta)^{2}+\frac{1}{2} \kappa_{2}(2 R \theta)^{2}=\text { constant }
\]

Differentianting w.r.t. time\\
or,

\[
\begin{gathered}
\frac{1}{2} I_{p} 2 \dot{\theta} \ddot{\theta}+\frac{1}{2}\left(\kappa_{1}+\kappa_{2}\right) 4 R^{2} 2 \theta \dot{\theta}=0 \\
\left(\frac{m R^{2}}{2}+m R^{2}\right) \ddot{\theta}+4 R^{2} \kappa \theta=0 \\
\left(\text { Because } I_{P}=I_{C}+m R^{2}=\frac{m R^{2}}{2}+m R^{2}\right)
\end{gathered}
\]

Hence

\[
\ddot{\theta}=-\frac{8}{3} \frac{\kappa}{m} \theta
\]

Thus \(\omega_{0}=\frac{8 \kappa}{3 m}\) and sought time period

\[
T=\frac{2 \pi}{\omega_{0}}=2 \pi \sqrt{\frac{3 m}{8 \kappa}}=\pi \sqrt{\frac{3 m}{2 \kappa}}
\]

4.58 In the C.M. frame (which is rigidly attached with the centre of mass of the two cubes) the cubes oscillates. We know that the kinetic energy of two body system equals \(\frac{1}{2} \mu v_{\text {rel }}^{2}\), where \(\mu\) is the reduced mass and \(\nu_{\text {rel }}\) is the modulus of velocity of any one body particle relative to other. From the conservation of mechanical energy of oscillation :

\[
\frac{1}{2} \kappa x^{2}+\frac{1}{2} \mu\left\{\frac{d}{d t}\left(l_{0}+x\right)\right\}^{2}=\text { constant }
\]

Here \(l_{0}\) is the natural length of the spring.\\
Differenting the above equation w.r.t time, we get :\\
\(\frac{1}{2} \kappa 2 x \dot{x}+\frac{1}{2} \mu 2 \dot{x} \ddot{x}=0\left[\right.\) becomes \(\left.\frac{d\left(l_{0}+x\right)}{d t}=\dot{x}\right]\)\\
Thus \(\ddot{x}=-\frac{\kappa}{\mu} x\) (where \(\mu=\frac{m_{1} m_{2}}{m_{1}+m_{2}}\) )\\
Hence the natural frequency of oscillation : \(\omega_{0}=\sqrt{\frac{\kappa}{\mu}}\) where \(\mu=\frac{m_{1} m_{2}}{m_{1}+m_{2}}\)\\
4.59 Suppose the balls 1 \& 2 are displaced by \(x_{1}, x_{2}\) from their initial position. Then the energy is : \(E=\frac{1}{2} m_{1} \dot{x}_{1}^{2}+m_{2} \dot{x}^{2}+\frac{1}{2} k\left(x_{1}-x_{2}\right)^{2}=\frac{1}{2} m_{1} v_{1}^{2}\)\\
Also total momentum is : \(m_{1} \dot{x}_{1}+m_{2} \dot{x}_{2}=m_{1} v_{1}\)\\
Define

\[
X=\frac{m_{1} x_{1}+m_{1} x_{2}}{m_{1}+m_{2}}, x=x_{1}-x_{2}
\]

Then

\[
\begin{aligned}
& x_{1}=X+\frac{m_{2}}{m_{1}+m_{2}} x, x_{2}=X-\frac{m_{1}}{m_{1}+m_{2}} x \\
& E=\frac{1}{2}\left(m_{1}+m_{2}\right) \dot{X}^{2}+\frac{1}{2} \frac{m_{1} m_{2}}{m_{1}+m_{2}} \dot{x}^{2}+\frac{1}{2} k x^{2}
\end{aligned}
\]

Hence

\[
\dot{X}=\frac{m_{1} v_{1}}{m_{1}+m_{2}}
\]

So

\[
\frac{1}{2} \frac{m_{1} m_{2}}{m_{1}+m_{2}} \dot{x}^{2}+\frac{1}{2} k x^{2}=\frac{1}{2} m_{1} v_{1}^{2}-\frac{1}{2} \frac{m_{1}^{2} v_{1}^{2}}{m_{1}+m_{2}}=\frac{1}{2} \frac{m_{1} m_{2}}{m_{1}+m_{2}} v_{1}^{2}
\]

(a) From the above equation

We see \(\omega=\sqrt{\frac{k}{\mu}}=\sqrt{\frac{3 \times 24}{2}}=6 \mathrm{~s}^{-1}\), when \(\mu=\frac{m_{1} m_{2}}{m_{1}+m_{2}}=\frac{2}{3} \mathrm{~kg}\).\\
(b) The energy of oscillation is

\[
\frac{1}{2} \frac{m_{1} m_{2}}{m_{2}+m_{2}} v_{1}^{2}=\frac{1}{2} \frac{2}{3} \times(0.12)^{2}=48 \times 10^{-4}=4.8 \mathrm{~mJ}
\]

We have \(x=\mathrm{a} \sin (\omega t+\alpha)\)\\
Initially

\[
x=0 \text { at } t=0 \text { so } \alpha=0
\]

Then

\[
x=a \sin \omega t . \text { Also } x=v_{1} \text { at } t=0 .
\]

So

\[
\omega a=v_{1} \text { and hence } a=\frac{v_{1}}{\omega}=\frac{12}{6}=2 \mathrm{~cm} .
\]

4.60 Suppose the disc 1 rotates by angle \(\theta_{1}\) and the disc 2 by angle \(\theta_{2}\) in the opposite sense. Then total torsion of the rod \(=\boldsymbol{\theta}_{1}+\boldsymbol{\theta}_{2}\)\\
and torsional P.E. \(=\frac{1}{2} \kappa\left(\theta_{1}+\theta_{2}\right)^{2}\)\\
(1)

The K.E. of the system (neglecting the moment of inertia of the rod) is

\[
\frac{1}{2} I_{1} \dot{\theta}_{1}^{2}+\frac{1}{2} I_{2} \dot{\theta}_{2}^{2}
\]

So total energy of the rod\\
\includegraphics[max width=\textwidth, alt={}, center]{da0156b6-4133-434e-90d4-ec01b2070afe-043_231_447_1555_945}

\[
E=\frac{1}{2} I_{1} \dot{\theta}_{1}^{2}+\frac{1}{2} I_{2} \dot{\theta}_{2}^{2}+\frac{1}{2} \kappa\left(\theta_{1}+\theta_{2}\right)^{2}
\]

We can put the total angular momentum of the rod equal to zero since the frequency associated with the rigid rotation of the whole system must be zero (and is known).

Thus

\[
I_{1} \theta_{1}=I_{2} \dot{\theta}_{2} \quad \text { or } \quad \frac{\dot{\theta}_{1}}{1 / I_{1}}=\frac{\dot{\theta}_{2}}{1 / I_{2}}=\frac{\dot{\theta}_{1}+\dot{\theta}_{2}}{1 / I_{1}+1 / I_{2}}
\]

So

\[
\dot{\theta}_{1}=\frac{I_{2}}{I_{1}+I_{2}}\left(\dot{\theta}_{1}+\dot{\theta}_{2}\right) \quad \text { and } \quad \dot{\theta}_{2}=\frac{I_{1}}{I_{1}+I_{2}}\left(\dot{\theta}_{1}+\dot{\theta}_{2}\right)
\]

and

\[
E=\frac{1}{2} \frac{I_{1} I_{2}}{I_{1}+I_{2}}\left(\dot{\theta}_{1}+\dot{\theta}_{2}\right)^{2}+\frac{1}{2} \kappa\left(\theta_{1}+\theta_{2}\right)^{2}
\]

The angular oscillation, frequency corresponding to this is

\[
\omega^{2}=\kappa / \frac{I_{1} I_{2}}{I_{1}+I_{2}}=\kappa / I^{\prime} \text { and } T=2 \pi \sqrt{\frac{I^{\prime}}{\kappa}}, \text { where } I^{\prime}=\frac{I_{1} I_{2}}{I_{1}+I_{2}}
\]

4.61 In the first mode the carbon atom remains fixed and the oxygen atoms move in equal \& opposite steps. Then total energy is\\
(1)\\
\includegraphics[max width=\textwidth, alt={}, center]{da0156b6-4133-434e-90d4-ec01b2070afe-044_181_677_693_451}

\[
\frac{1}{2} 2 m_{0} \dot{x}^{2}+\frac{1}{2} 2 \kappa x^{2}
\]

where \(x\) is the displacement of one of the 0 atom (say left one). Thus


\begin{equation*}
\omega_{1}^{2}=\kappa / m_{0} . \tag{2}
\end{equation*}


\begin{center}
\includegraphics[max width=\textwidth, alt={}]{da0156b6-4133-434e-90d4-ec01b2070afe-044_216_663_1077_444}
\end{center}

In this mode the oxygen atoms move in equal steps in the same direction but the carbon atom moves in such a way as to keep the centre of mass fixed.

Thus

\[
2 m_{0} x+m_{c} y=0 \text { or, } y=-\frac{2 m_{0}}{m_{c}} x
\]

KE. \(=\frac{1}{2} 2 m_{0} \dot{x}^{2}+\frac{1}{2} m_{c}\left(\frac{2 m_{0}}{m_{c}} \dot{x}\right)^{2}=\frac{1}{2} 2 m_{0} \dot{x}^{2}+\frac{1}{2} 2 m_{0} \frac{2 m_{0}}{m_{c}} \dot{x}^{2} \quad=\frac{1}{2} 2 m_{0}\left(1+\frac{2 m_{0}}{m_{c}}\right) \dot{x}^{2}\)\\
P.E. \(=\frac{1}{2} k\left(1+\frac{2 m_{0}}{m_{c}}\right)^{2} x^{2}+\frac{1}{2} \kappa\left(1+\frac{2 m_{0}}{m_{c}}\right)^{2} x^{2}=\frac{1}{2} 2 \kappa\left(1+\frac{2 m_{0}}{m_{c}}\right)^{2} x^{2}\)

Thus

\[
\omega_{2}^{2}=\frac{\kappa}{m_{0}}\left(1+\frac{2 m_{0}}{m_{c}}\right) \text { and } \omega_{2}=\omega_{1} \sqrt{1+\frac{2 m_{0}}{m_{c}}}
\]

Hence,

\[
\omega_{2}=\omega_{1} \sqrt{1+\frac{32}{12}}=\omega_{1} \sqrt{\frac{11}{3}} \approx 1.91 \omega_{1}
\]

4.62 Let, us displace the piston through small distance \(x\), towards right, then from \(F_{x}=m w_{x}\)\\
\includegraphics[max width=\textwidth, alt={}, center]{da0156b6-4133-434e-90d4-ec01b2070afe-045_296_769_131_293}\\
or,


\begin{equation*}
\left(p_{1}-p_{1}\right) S=-m \ddot{x} \tag{1}
\end{equation*}


But, the process is adiabatic, so from \(P V^{\Upsilon}=\) const.

\[
p_{2}=\frac{p_{0} V_{0}^{\gamma}}{\left(V_{0}-S x\right)^{\gamma}} \quad \text { and } p_{1}=\frac{p_{0} V_{0}^{\gamma}}{\left(V_{0}+S x\right)^{\gamma}}
\]

as the new volumes of the left and the right parts are now \(\left(V_{0}+S x\right)\) and \(\left(V_{0}-S x\right)\) respectively. So, the Eqn (1) becomes.\\
or,

\[
\begin{aligned}
\frac{p_{0} V_{0}^{\gamma} S}{m}\left\{\frac{1}{\left(V_{0}-S x\right)^{\gamma}}-\frac{1}{\left(V_{0}+S x\right)^{\gamma}}\right\} & =-\ddot{x} \\
\frac{p_{0} V_{0}^{\gamma} S}{m}\left\{\frac{\left(V_{0}+S x\right)^{\gamma}-\left(V_{0}-S x\right)^{\gamma}}{\left(V_{0}^{2}-S^{2} x^{2}\right)^{\gamma}}\right\} & =-\ddot{x} \\
\frac{p_{0} V_{0}^{\gamma} S}{m}\left\{\frac{\left(1+\frac{\gamma S x}{V_{0}}\right)-\left(1-\frac{\gamma S x}{V_{0}}\right)}{V_{0}^{\gamma}\left(1-\frac{\gamma S^{2} x^{2}}{V_{0}^{2}}\right)}\right\} & =-\ddot{x}
\end{aligned}
\]

or,

Neglecting the term \(\frac{\gamma S^{2} x^{2}}{V_{0}^{2}}\) in the denominator, as it is very small, we get,

\[
\ddot{x}=-\frac{2 p_{0} S^{2} \gamma x}{m V_{0}},
\]

which is the equation for S.H.M. and hence the oscillating frequency.

\[
\omega_{0}=s \sqrt{\frac{2 p_{0} \gamma}{m V_{0}}}
\]

4.63 In the absence of the charge, the oscillation period of the ball

\[
T=2 \pi \sqrt{l / g}
\]

when we impart the charge \(q\) to the ball, it will be influenced by the induced charges on the conducting plane. From the electric image method the electric force on the ball by the plane equals \(\frac{q^{2}}{4 \pi \varepsilon_{0}(2 h)^{2}}\) and is directed downward. Thus in this case the effective acceleration of the ball

\[
g^{\prime}=g+\frac{q^{2}}{16 \pi \varepsilon_{0} m h^{2}}
\]

and the corresponding time period

\[
T^{\prime}=2 \pi \sqrt{\frac{l}{g^{\prime}}}=2 \pi \sqrt{\frac{l}{g+\frac{q^{2}}{16 \pi \varepsilon_{0} m h^{2}}}}
\]

From the conditon of the problem

\[
T=\eta T^{\prime}
\]

So,

\[
T^{2}=\eta^{2} T^{\prime 2} \text { or } \frac{1}{g}=\eta^{2}\left(\frac{1}{g+\frac{q^{2}}{16 \pi \varepsilon_{0} m h^{2}}}\right)
\]

Thus on solving

\[
q=4 h \sqrt{\pi \varepsilon_{0} m g\left(\eta^{2}-1\right)}=2 \mu C
\]

4.64 In a magnetic field of induction \(B\) the couple on the magnet is \(-M B \sin \theta=-M B \theta\) equating this to \(I \theta\) we get

\[
I \ddot{\theta}+M B \theta=0
\]

or

\[
\omega^{2}=\frac{M B}{I} \quad \text { or } \quad T=2 \pi \sqrt{\frac{I}{M B}} .
\]

Given

\[
T_{2}=T_{1} / \eta
\]

o.

\[
\sqrt{\frac{1}{B_{2}}}=\sqrt{\frac{1}{B_{1}}} \cdot \frac{1}{\eta} \text { or } \frac{1}{B_{2}}=\frac{1}{B_{1}} \cdot \frac{1}{\eta^{2}}
\]

or

\[
B_{2}=\eta^{2} B_{1}
\]

The induction of the field increased \(\eta^{2}\) times.\\
4.65 We have in the circuit at a certain instant of time \((t)\), from Faraday's law of electromagnetic induction :

\[
L \frac{d i}{d t}=B l \frac{d x}{d t} \text { or } L d i=B l d x
\]

As at

\[
t=0, x=0, \text { so } L i=B l x \text { or } i=\frac{B l}{L} x(1)
\]

For the rod from the second law of motion \(F_{x}=m w_{x}\)

\[
-i l B=m \ddot{x}
\]

Using Eqn. (1), we get : \(\quad \ddot{x}=-\left(\frac{l^{2} B^{2}}{m L}\right) x=-\omega_{0}^{2} x\)\\
where


\begin{equation*}
\omega_{0}=l B / \sqrt{m L} \tag{2}
\end{equation*}


The solution of the above differential equation is of the form

\[
x=a \sin \left(\omega_{0} t+\alpha\right)
\]

From the initial condition, at \(t=0, x=0\), so \(\alpha=0\)\\
Hence,


\begin{equation*}
x=a \sin \omega_{0} t \tag{3}
\end{equation*}


Differentiating w.r.t. time, \(\dot{x}=a \omega_{0} \cos \omega_{0} t\)\\
But from the initial condition of the problem at \(t=0, \dot{x}=\mathrm{v}_{0}\)\\
Thus


\begin{equation*}
v_{0}=a \omega_{0} \quad \text { or } \quad a=v_{0} / \omega_{0} \tag{4}
\end{equation*}


Putting the value of \(a\) from Eqn. (4) into Eqn. (3), we obtained

\[
x=\frac{v_{0}}{\omega_{0}} \sin \omega_{0} t\left(\text { where } \omega_{0}=\frac{l B}{\sqrt{m L}}\right)
\]

4.66 As the connector moves, an emf is set up in the circuit and a current flows, since the emf is \(\xi=-B l \dot{x}\), we must have \(:-B l \dot{x}+L \frac{d I}{d t}=0\)\\
so, \(\quad I=B l x / L\)\\
provided \(x\) is measured from the initial position.\\
We then have

\[
m \ddot{x}=-\frac{B l x}{L} \cdot B \cdot l+m g
\]

for by Lenz's law the induced current will oppose downward sliding. Finally\\
\includegraphics[max width=\textwidth, alt={}, center]{da0156b6-4133-434e-90d4-ec01b2070afe-047_434_397_612_965}

\[
\ddot{x}+\frac{(B l)^{2}}{m L} x=g
\]

on putting

\[
\begin{aligned}
& \omega_{0}=\frac{B l}{\sqrt{m L}} \\
& \ddot{x}+\omega_{0}^{2} x=g
\end{aligned}
\]

A solution of this equation is \(x=\frac{g}{\omega_{0}^{2}}+A \cos \left(\omega_{0} t+\alpha\right)\)\\
But

\[
x=0 \text { and } \dot{x}=0 \text { at } t=0 \text {. This gives }
\]

\[
x=\frac{g}{\omega_{0}^{2}}\left(1-\cos \omega_{0} t\right)
\]

4.67 We are given \(x=a_{0} e^{-\beta t} \sin \omega t\)\\
(a) The velocity of the point at \(t=0\) is obtained from

\[
v_{0}=(\dot{x})_{t=0}=\omega a_{0}
\]

The term "oscillation amplitude at the moment \(t=0\) " is meaningless. Probably the implication is the amplitude for \(t \ll \frac{1}{\beta}\). Then \(x \sim a_{0} \sin \omega t\) and amplitude is \(a_{0}\).\\
(b) \(\dot{x}=\left(-\beta a_{0} \sin \omega t+\omega a_{0} \cos \omega t\right) e^{-\beta t}=0\)\\
when the displacement is an extremum. Then

\[
\tan \omega t=\frac{\omega}{\beta}
\]

or

\[
\omega t=\tan ^{-1} \frac{\omega}{\beta}+n \pi, n=0,1,2, \ldots
\]

4.68 Given \(\varphi=\varphi_{0} e^{-\beta t} \cos \omega t\)\\
we have \(\dot{\varphi}=-\beta \varphi-\omega \varphi_{0} e^{-\beta t} \sin \omega t\)

\[
\begin{aligned}
\ddot{\varphi} & =-\beta \dot{\varphi}+\beta \omega \varphi_{0} e^{-\beta t} \sin \omega t-\omega^{2} \varphi_{0} e^{-\beta t} \cos \omega t \\
& =\beta^{2} \varphi+2 \beta \omega \varphi_{0} e^{-\beta t} \sin \omega t-\omega^{2} \varphi
\end{aligned}
\]

so\\
(a) \((\dot{\varphi})_{0}=-\beta \varphi_{0},(\ddot{\varphi})_{0}=\left(\beta^{2}-\omega^{2}\right) \varphi_{0}\)\\
(b) \(\dot{\varphi}=-\varphi_{0} e^{-\beta t}(\beta \cos \omega t+\omega \sin \omega t)\) becomes maximum (or minimum) when

\[
\ddot{\varphi}=\varphi_{0}\left(\beta^{2}-\omega^{2}\right) e^{-\beta t} \cos \omega t+2 \beta \omega \varphi_{0} e^{-\beta t} \sin \omega t=0
\]

or

\[
\tan \omega t=\frac{\omega^{2}-\beta^{2}}{2 \beta \omega}
\]

and

\[
t_{n}=\frac{1}{\omega}\left[\tan ^{-1} \frac{\omega^{2}-\beta^{2}}{2 \beta \omega}+n \pi\right], n=0,1,2, \ldots
\]

4.69 We write \(x=a_{0} e^{-\beta t} \cos (\omega t+\alpha)\).\\
(a) \(x(0)=0 \Rightarrow \alpha= \pm \frac{\pi}{2} \Rightarrow x=\mp a_{0} e^{-\beta t} \sin \omega t\)

\[
\dot{x}(0)=(\dot{x})_{t=0}=\mp \omega a_{0}
\]

Since \(a_{0}\) is \(+v e\), we must choose the upper sign if \(\dot{x}(C)<0\) and the lower sign if \(\dot{x}(0)>0\). Thus

\[
a_{0}=\frac{|\dot{x}(0)|}{\omega} \text { and } \alpha= \begin{cases}+\frac{\pi}{2} & \text { if } \dot{x}(0)<0 \\ -\frac{\pi}{2} & \text { if } \dot{x}(0)>0\end{cases}
\]

(b) we write \(x=\operatorname{Re} A e^{-\beta t+i \omega t}, A=a_{0} e^{i \alpha}\)

Then

\[
\dot{x}=v_{x}=\operatorname{Re}(-\beta+i \omega) A e^{-\beta t+i \omega t}
\]

From \(v_{x}(0)=0\) we get \(\operatorname{Re}(-\beta+i \omega) A=0\)\\
This implies \(A= \pm i(\beta+i \omega) B\) where \(B\) is real and positive. Also

\[
x_{0}=\operatorname{Re} A=\mp \omega B
\]

Thus

\[
B=\frac{\left|x_{0}\right|}{\omega} \text { with }+\operatorname{sign} \text { in } A \text { if } x_{0}<0
\]

So

\[
A= \pm i \frac{\beta+i \omega}{\omega}\left|x_{0}\right|=\left(\mp 1+ \pm \frac{i \beta}{\omega}\right)\left|x_{0}\right|
\]

Finally

\[
\begin{aligned}
a_{0} & =\sqrt{1+\left(\frac{\beta}{\omega}\right)^{2}}\left|x_{0}\right| \\
\alpha & =\frac{-\beta}{\omega}, \alpha=\tan ^{-1}\left(\frac{-\beta}{\omega}\right)
\end{aligned}
\]

\(\alpha\) is in the \(4^{\text {th }}\) quadrant \(\left(-\frac{\pi}{2}<\alpha<0\right)\) if \(x_{0}>0\) and \(\alpha\) is in the \(2^{\text {nd }}\) quadrant \(\left(\frac{\pi}{2}<\alpha<\pi\right)\) if \(x_{0}<0\).\\
\(4.70 x=a_{0} e^{-\beta t} \cos (\omega t+\alpha)\)\\
Then

\[
(\dot{x})_{t-0}=-\beta a_{0} \cos \alpha-\omega a_{0} \sin \alpha=0
\]

or

\[
\tan \alpha=-\frac{\beta}{\omega} .
\]

Also

\[
\begin{gathered}
(x)_{t-0}=a_{0} \cos \alpha=\frac{a_{0}}{\eta} \\
\sec ^{2} \alpha=\eta^{2}, \tan \alpha=-\sqrt{\eta^{2}-1}
\end{gathered}
\]

Thus

\[
\beta=\omega \sqrt{\eta^{2}-1}
\]

(We have taken the amplitude at \(t=0\) to be \(a_{0}\) ).\\
4.71 We write \(x=a_{0} e^{-\beta t} \cos (\omega t+\alpha)\)

\[
\begin{aligned}
& =\operatorname{Re} A e^{-\beta t+i \omega t}, A=a_{0} e^{i \alpha} \\
\dot{x} & =\operatorname{Re} A(-\beta+i \omega) e^{-\beta t+i \omega t}
\end{aligned}
\]

Velocity amplitude as a function of time is defined in the following manner. Put \(t=t_{0}+\tau\), then

\[
\begin{gathered}
x=\operatorname{Re} A e^{-\beta\left(t_{0}+\tau\right)} e^{i \omega\left(t_{0}+\tau\right)} \\
\approx \operatorname{Re} A e^{-\beta t_{0}} e^{i \omega t_{0}+i \omega \tau} \approx \operatorname{Re} A e^{-\beta t_{0}} e^{i \omega t}
\end{gathered}
\]

for \(\tau \ll \frac{1}{\beta}\). This means that the displacement amplitude around the time \(t_{0}\) is \(a_{0} e^{-\beta t_{0}}\) and we can say that the displacement amplitude at time \(t\) is \(a_{0} e^{-\beta t}\). Similarly for the velocity amplitude.\\
Clearly\\
(a) Velocity amplitude at time \(t=a_{0} \sqrt{\beta^{2}+\omega^{2}} e^{-\beta t}\)

Since

\[
A(-\beta+i \omega)=a_{0} e^{i \alpha}(-\beta+i \omega)
\]

\[
=a_{0} \sqrt{\beta^{2}+\omega^{2}} e^{i \gamma}
\]

where \(\gamma\) is anotner constant.\\
(b) \(x(0)=0 \Rightarrow \operatorname{Re} A=0\) or \(A= \pm i a_{0}\)\\
where \(a_{0}\) is real and positive.\\
Also

\[
\begin{aligned}
v_{x}(0) & =\dot{x}_{0}=\operatorname{Re} \pm i a_{0}(-\beta+i \omega) \\
& =\mp \omega a_{0}
\end{aligned}
\]

Thus \(a_{0}=\frac{\left|\dot{x}_{0}\right|}{\omega}\) and we take \(-(+)\) sign if \(x_{0}\) is negative (positive). Finally the velocity amplitude is obtained as

\[
\frac{\left|\dot{x}_{0}\right|}{\omega} \sqrt{\beta^{2}+\omega^{2}} e^{-\beta t} .
\]

4.72 The first oscillation decays faster in time. But if one takes the natural time scale, the period \(T\) for each oscillation, the second oscillation attenuates faster during that period.\\
4.73 By definition of the logarithemic decrement \(\left(\lambda=\beta \frac{2 \pi}{\omega}\right)\) we get for the original decrement \(\lambda_{0}\)

Now

\[
\begin{aligned}
& \lambda_{0}=\beta \frac{2 \pi}{\sqrt{\omega_{0}^{2}-\beta^{2}}} \text { and finally } \lambda=\frac{2 \pi n \beta}{\sqrt{\omega_{0}^{2}-n^{2} \beta^{2}}} \\
& \frac{\beta}{\sqrt{\omega_{0}^{2}-\beta^{2}}}=\frac{\lambda_{0}}{2 \pi} \text { or } \frac{\beta}{\omega_{0}}=\frac{\lambda_{0} / 2 \pi}{\sqrt{1+\left(\frac{\lambda_{0}}{2 \pi}\right)^{2}}}
\end{aligned}
\]

so

\[
\frac{\lambda / 2 \pi}{\sqrt{1+\left(\frac{\lambda}{2 \pi}\right)^{2}}}=\frac{n \frac{\lambda_{0}}{2 \pi}}{\sqrt{1+\left(\frac{\lambda_{0}}{2 \pi}\right)^{2}}}
\]

Hence

\[
\frac{\lambda}{2 \pi}=\frac{n \lambda_{0} / 2 \pi}{\sqrt{1-\left(n^{2}-1\right)\left(\frac{\lambda_{0}}{2 \pi}\right)^{2}}}
\]

For critical damping

\[
\omega_{0}=n_{c} \beta
\]

\[
\frac{1}{n_{c}}=\frac{\beta}{\omega_{0}}=\frac{\lambda_{0} / 2 \pi}{\sqrt{1+\left(\frac{\lambda_{0}}{2 \pi}\right)^{2}}} \quad \text { or } \quad n_{c}=\sqrt{1+\left(\frac{2 \pi}{\lambda_{0}}\right)^{2}}
\]

4.74 The Eqn of the dead weight is

\[
m \ddot{x}+2 \beta m \dot{x}+m \omega_{0}^{2} x=m g
\]

so

\[
\Delta x=\frac{g}{\omega_{0}^{2}} \quad \text { or } \quad \omega_{0}^{2}=\frac{g}{\Delta x}
\]

Now

\[
\lambda=\frac{2 \pi \beta}{\omega}=\frac{2 \pi \beta}{\sqrt{\omega_{0}^{2}-\beta^{2}}} \quad \text { or } \quad \frac{\omega_{0}}{\sqrt{\omega_{0}^{2}-\beta^{2}}}=\sqrt{1+\left(\frac{\lambda}{2 \pi}\right)^{2}}
\]

Thus

\[
\begin{gathered}
T=\frac{2 \pi}{\sqrt{\omega_{0}^{2}-\beta^{2}}}=\frac{2 \pi}{\omega_{0}} \sqrt{1+\left(\frac{\lambda}{2 \pi}\right)^{2}} \\
=2 \pi \sqrt{\frac{\Delta x}{g}} \sqrt{1+\left(\frac{\lambda}{2 \pi}\right)^{2}}=\sqrt{\frac{\Delta x}{g}\left(4 \pi^{2}+\lambda^{2}\right)}=0.70 \mathrm{sec} .
\end{gathered}
\]

4.75 The displacement amplitude decrease \(\eta\) times every \(n\) oscillations. Thus\\
or

So

\[
\begin{gathered}
\frac{1}{\eta}=e^{-\beta \cdot \frac{2 \pi}{\omega} \cdot n} \\
\frac{2 \pi n \beta}{\omega}=\ln \eta \quad \text { or } \frac{\beta}{\omega}=\frac{\ln \eta}{2 \pi n} . \\
Q=\frac{\omega}{2 \beta}=\frac{\pi n}{\ln \eta} \approx 499 .
\end{gathered}
\]

4.76 From \(x=a_{0} e^{-\beta t} \cos (\omega t+\alpha \dot{\alpha})\), we get using

\[
(x)_{t=0}=l=a_{0} \cos \alpha
\]

\(0=(\dot{x})_{t=0}=-\beta a_{0} \cos \alpha-\omega a_{0} \sin \alpha\)\\
Then \(\tan \alpha=-\frac{\beta}{\omega} \quad\) or \(\quad \cos \alpha=\frac{\omega}{\sqrt{\omega^{2}+\beta^{2}}}\)\\
and \(x=\frac{l \sqrt{\omega^{2}+\beta^{2}}}{\omega} e^{-\beta t} \cos \left(\omega t-\tan ^{-1} \frac{\beta}{\omega}\right)\)

\[
x=0 \text { at } t=\frac{1}{\omega}\left(n \pi+\frac{\pi}{2}+\tan ^{-1} \frac{\beta}{\omega}\right)
\]

Total distance travelled in the first lap \(=l\)\\
\includegraphics[max width=\textwidth, alt={}, center]{da0156b6-4133-434e-90d4-ec01b2070afe-051_492_553_1073_791}

To get the maximum displacement in the second lap we note that

\[
\begin{gathered}
\dot{x}=\left[-\beta \cos \left(\omega t-\tan ^{-1} \frac{\beta}{\omega}\right)-\omega \sin \left(\omega t-\tan ^{-1} \frac{\beta}{\omega}\right)\right] \\
x \frac{l \sqrt{\omega^{2}+\beta^{2}}}{\omega} e^{-\beta t}=0
\end{gathered}
\]

when \(\omega t=\pi, 2 \pi, 3 \pi, \ldots\) etc.

Thus

\[
\ddot{x}_{\max }=-a_{0} e^{-\pi \beta / \omega} \cos \alpha=-l e^{-\pi \beta / \omega} \text { for } t=\pi / \omega
\]

so, distance traversed in the \(2^{\text {nd }}\) lap \(=2 l e^{-\pi \beta / \omega}\)\\
Continuing total distance traversed \(=l+2 l e^{-\pi \beta / \omega}+2 l e^{-2 \pi \beta / \omega}+\ldots .\).

\[
\begin{aligned}
& \quad=l+\frac{2 l e^{-\pi \beta / \omega}}{1-e^{-\beta \pi / \omega}}=l+\frac{2 l}{e^{\beta \pi / \omega}-1} \\
& =l \frac{e^{\beta \pi / \omega}+1}{e^{\beta \pi / \omega}-1}=l \frac{1+e^{\lambda / 2}}{e^{\lambda / 2}-1}
\end{aligned}
\]

where \(\lambda=\frac{2 \pi \beta}{\omega}\) is the logarithemic decrement. Substitution gives 2 metres.\\
4.77 For an undamped oscillator the mechanical energy \(E=\frac{1}{2} m \dot{x}^{2}+\frac{1}{2} m \omega_{0}^{2} x^{2}\) is conserved. For a damped oscillator.\\
and

\[
\begin{gathered}
x=a_{0} e^{-\beta t} \cos (\omega t+\alpha), \omega=\sqrt{\omega_{0}^{2}-\beta^{2}} \\
E(t)=\frac{1}{2} m \dot{x}^{2}+\frac{1}{2} m \omega_{0}^{2} x^{2}
\end{gathered}
\]

\[
\begin{aligned}
= & \frac{1}{2} m a_{0}^{2} e^{-2 \beta t}\left[\beta^{2} \cos ^{2}(\omega t\right. \\
& \left.+\alpha)+2 \beta \omega \cos (\omega t+\alpha) \times \sin (\omega t+\alpha)+\omega^{2} \sin ^{2}(\omega t+\alpha)\right] \\
& +\frac{1}{2} m a_{0}^{2} \omega_{0}^{2} e^{-2 \beta t} \cos ^{2}(\omega t+\alpha) \\
= & \frac{1}{2} m a_{0}^{2} \omega_{0}^{2} e^{-2 \beta t}+\frac{1}{2} m a_{0}^{2} \beta^{2} e^{-2 \beta t} \cos (2 \omega t+2 \alpha)+\frac{1}{2} m a_{0}^{2} \beta \omega e^{-2 \beta t} \sin (2 \omega t+2 \alpha)
\end{aligned}
\]

If \(\beta \ll \omega\), then the average of the last two terms over many oscillations about the time \(t\) will vanish and

\[
<E(t)>\cong \frac{1}{2} m a_{0}^{2} \omega_{0}^{2} e^{-2 \beta t}
\]

and this is the relevant mechanical energy.\\
In time \(\tau\) this decreases by a factor \(\frac{1}{\eta}\) so

\[
\begin{aligned}
e^{-2 \beta \tau} & =\frac{1}{\eta} \text { or } \tau=\frac{\ln \eta}{2 \beta} . \\
\beta & =\frac{\ln \eta}{2 \tau}
\end{aligned}
\]

and

\[
\lambda=\frac{2 \pi \beta}{\sqrt{\omega_{0}^{2}-\beta^{2}}}=\frac{2 \pi}{\sqrt{\left(\frac{\omega_{0}}{\beta}\right)^{2}-1}}=\frac{2 \pi}{\sqrt{\frac{4 g \tau^{2}}{l \ln ^{2} \eta}-1}} \quad \text { since } \omega_{0}^{2}=\frac{g}{l} .
\]

and

\[
Q=\frac{\pi}{\lambda}=\frac{1}{2} \sqrt{\frac{4 g \tau^{2}}{l \ln ^{2} \eta}-1} \cong 130 .
\]

4.78 The restoring couple is

\[
\Gamma=-m g R \sin \varphi \cong-m g R \varphi
\]

The moment of inertia is

\[
I=\frac{3 m R^{2}}{2}
\]

Thus for undamped oscillations

\[
\frac{3 m R^{2}}{2} \ddot{\varphi}+m g R \varphi=0
\]

so, \(\quad \omega_{0}^{2}=\frac{2 g}{3 R}\)\\
\includegraphics[max width=\textwidth, alt={}, center]{da0156b6-4133-434e-90d4-ec01b2070afe-053_416_386_134_1022}

Also

\[
\lambda=\frac{2 \pi \beta}{\omega}=\frac{2 \pi \beta}{\sqrt{\omega_{0}^{2}-\beta^{2}}}
\]

Hence

\[
\frac{\beta}{\sqrt{\omega_{0}^{2}-\beta^{2}}}=\frac{\lambda}{2 \pi} \quad \text { or } \quad \frac{\omega_{0}}{\sqrt{\omega_{0}-\beta^{2}}}=\sqrt{1+\left(\frac{\lambda}{2 \pi}\right)^{2}}
\]

Hence finally the period. \(T\) of small oscillation comes to

\[
\begin{aligned}
T=\frac{2 \pi}{\omega} & =\frac{2 \pi}{\omega_{0}} \times \frac{\omega_{0}}{\sqrt{\omega_{0}^{2}-\beta^{2}}} 2 \pi \sqrt{\frac{3 R}{2 g}\left(1+\left(\frac{\lambda}{2 \pi}\right)^{2}\right)} \\
& =\sqrt{\frac{3 R}{2 g}\left(4 \pi^{2}+\lambda^{2}\right)}=0.90 \mathrm{sec} .
\end{aligned}
\]

4.79 Let us calculate the moment \(G_{1}\) of all the resistive forces on the disc. When the disc rotates an element ( \(r d r d \theta\) ) with coordinates ( \(r, \theta\) ) has a velocity \(r \dot{\varphi}\), where \(\varphi\) is the instantaneous angle of rotation from the equilibrium position and \(r\) is measured from the centre. Then

\[
\begin{aligned}
G_{1} & =\int_{0}^{2 \pi} d \theta \int_{0}^{R} d r \cdot r \cdot\left(F_{1} \times r\right) \\
& =\int_{0}^{R} \eta r \dot{\varphi} r^{2} d \gamma \times 2 \pi=\frac{\eta \pi R^{4}}{2} \dot{\varphi}
\end{aligned}
\]

Also moment of inertia \(=\frac{m R^{2}}{2}\)

Thus

\[
\frac{m R^{2}}{2} \ddot{\varphi}+\frac{\pi \eta R^{4}}{2} \dot{\varphi}+\alpha \varphi=0
\]

or

\[
\ddot{\varphi}+2 \frac{\pi \eta R^{2}}{2 m} \dot{\varphi}+\frac{2 \alpha}{m R^{2}} \varphi=0
\]

Hence

\[
\omega_{0}^{2}=\frac{2 \alpha}{m R^{2}} \quad \text { and } \quad \beta=\frac{\pi \eta R^{2}}{2 m}
\]

and angular frequency

\[
\omega=\sqrt{\left(\frac{2 \alpha}{m R^{2}}\right)-\left(\frac{\pi \eta R^{2}}{2 m}\right)^{2}}
\]

Note :- normally by frequency we mean \(\frac{\omega}{2 \pi}\).\\
4.80 From the law of viscosity, force per unit area \(=\eta \frac{d \nu}{d x}\)\\
so when the disc executes torsional oscillations the resistive couple on it is

\[
=\int_{0}^{R} \eta \cdot 2 \pi r \cdot \frac{r \varphi}{h} \cdot r \cdot d r \times 2=\frac{\eta \pi R^{4}}{h} \dot{\varphi}
\]

(factor 2 for the two sides of the disc; see the figure in the book)\\
where \(\varphi\) is torsion. The equation of motion is

\[
I \ddot{\varphi}+\frac{\eta \pi R^{4}}{h} \dot{\varphi}+c \varphi=0
\]

Comparing with

\[
\begin{gathered}
\ddot{\varphi}+2 \beta \dot{\varphi}+\omega_{0}^{2} \varphi=0 \text { we get } \\
\beta=\eta \pi R^{4} / 2 h I
\end{gathered}
\]

Now the logarithmic decrement \(\lambda\) is given by \(\lambda=\beta T, T=\) time period Thus

\[
\eta=2 \lambda h I / \pi R^{4} T
\]

4.81 If \(\varphi=\) angle of deviation of the frame from its normal position, then an e.m.f.

\[
\varepsilon=B a^{2} \dot{\varphi}
\]

is induced in the frame in the displaced position and a current \(\frac{\varepsilon}{R}=\frac{B a^{2} \dot{\varphi}}{R}\) flows in it. A couple

\[
\frac{B a^{2} \dot{\varphi}}{R} \cdot B \cdot a \cdot a=\frac{B^{2} a^{4}}{R} \dot{\varphi}
\]

then acts on the frame in addition to any elastic restoring couple \(\boldsymbol{c} \boldsymbol{\varphi}\). We write the equation of the frame as

\[
I \ddot{\varphi}+\frac{B^{2} a^{4}}{R} \dot{\varphi}+c \varphi=0
\]

Thus \(\beta=\frac{B^{2} a^{4}}{2 I R}\) where \(\beta\) is defined in the book.\\
Amplitude of oscillation die out according to \(e^{-\beta t}\) so time required for the oscillations to decrease to \(\frac{1}{e}\) of its value is

\[
\frac{1}{\beta}=\frac{2 I R}{B^{2} a^{4}}
\]

4.82 We shall denote the stiffness constant by \(\kappa\). Suppose the spring is stretched by \(x_{0}\). The bar in then subject to two horizontal forces (1) restoring force \(-\kappa x\) and (2) friction kmg opposing motion. If

\[
x_{0}>\frac{k m g}{\kappa}=\Delta
\]

the bar will come back.\\
(If \(x_{0} \leq \Delta\), the bar will stay put.)\\
The equation of the bar when it is moving to the left is

\[
m \ddot{x}=-\kappa x+k m g
\]

\begin{center}
\includegraphics[max width=\textwidth, alt={}]{da0156b6-4133-434e-90d4-ec01b2070afe-055_238_548_305_833}
\end{center}

This equation has the solution

\[
x=\Delta+\left(x_{0}-\Delta\right) \cos \sqrt{\frac{k}{m}} t
\]

where we have used \(x=x_{0}, \dot{x}=0\) at \(t=0\). This solution is only valid till the bar comes to rest. This happens at

\[
t_{1}=\pi / \sqrt{\frac{k}{m}}
\]

and at that time \(x=x_{1}=2 \Delta-x_{0}\). if \(x_{0}>2 \Delta\) the tendency of the rod will now be to move to the right .(if \(\Delta<x_{0}<2 \Delta\) the rod will stay put now) Now the equation for rightward motion becomes \({ }^{\text {® }}\)

\[
m \ddot{x}=-\kappa x-k m g
\]

(the friction force has reversed).\\
We notice that the rod will move to the right only if

\[
\kappa\left(x_{0}-2 \Delta\right)>k m g \text { i.e. } x_{0}>3 \Delta
\]

In this case the solution is

Since

\[
\begin{gathered}
x=-\Delta+\left(x_{0}-3 \Delta\right) \cos \sqrt{\frac{k}{m}} t \\
x=2 \Delta-x_{0} \text { and } \dot{x}=0 \text { at } t=t_{1}=\pi / \sqrt{\frac{k}{m}} .
\end{gathered}
\]

The rod will next come to rest at

\[
t=t_{2}=2 \pi / \sqrt{\frac{k}{m}}
\]

and at that instant \(x=x_{3}=x_{0}-4 \Delta\). However the rod will stay put unless \(x_{0}>5 \Delta\). Thus\\
(a) time period of one full oscillation \(=2 \pi / \sqrt{\frac{k}{m}}\).\\
(b) There is no oscillation if \(0<x_{0}<\Delta\)

One half oscillation if \(\quad \Delta<x_{0}<3 \Delta\)

2 half oscillation if \(3 \Delta<x_{0}<5 \Delta\) etc.\\
We can say that the number of full oscillations is one half of the integer \(n\)\\
where

\[
n=\left[\frac{x_{0}-\Delta}{2 \Delta}\right]
\]

where \([x]=\) smallest non-negative integer greater than \(x\).\\
4.83 The equation of motion of the ball is

\[
m\left(\ddot{x}+\omega_{0}^{2} x\right)=F_{0} \cos \omega t
\]

This equation has the solution

\[
x=A \cos \left(\omega_{0} t+\alpha\right)+B \cos \omega t
\]

where \(A\) and \(\alpha\) are arbitrary and \(B\) is obtained by substitution in the above equation

\[
B=\frac{F_{0} / m}{\omega_{0}^{2}-\omega^{2}}
\]

The conditions \(x=0, \dot{x}=0\) at \(t=0\) give

\[
A \cos \alpha+\frac{F_{0} / m}{\omega_{0}^{2}-\omega^{2}}=0 \text { and }-\omega_{0} A \sin \alpha=0
\]

This gives \(\alpha=0\),

\[
A=-\frac{F_{0} / m}{\omega_{0}^{2}-\omega^{2}}=\frac{F_{0} / m}{\omega^{2}-\omega_{0}^{2}}
\]

Finally,

\[
x=\frac{F_{0} / m}{\omega^{2}-\omega_{0}^{2}}\left(\cos \omega_{0} t-\cos \omega t\right)
\]

4.84 We have to look for solutions of the equation

\[
\begin{gathered}
m \ddot{x}+k x=F, 0<t_{1}<\tau, \\
m \ddot{x}+k x=0, t>\tau
\end{gathered}
\]

subject to \(x(0)=\dot{x}(0)=0\) where \(F\) is constant.\\
The solution of this equation will be sought in the form

\[
\begin{aligned}
x & =\frac{F}{k}+A \cos \left(\omega_{0} t+\alpha\right), 0 \leq t \leq \tau \\
x & =B \cos \left(\omega_{0}(t-\tau)+\beta\right), t>\tau
\end{aligned}
\]

\(A\) and \(\alpha\) will be determined from the boundary condition at \(t=0\).

\[
\begin{aligned}
& 0=\frac{F}{k}+A \cos \alpha \\
& 0=-\omega_{0} A \sin \alpha
\end{aligned}
\]

Thus

\[
\alpha=0 \text { and } A=-\frac{F}{k} \quad \text { and } \quad x=\frac{F}{k}\left(1-\cos \omega_{0} t\right) 0 \leq t<\tau .
\]

\(B\) and \(\beta\) will be determined by the continuity of \(x\) and \(\dot{x}\) at \(t=\tau\). Thus

\[
\frac{F}{k}\left(1-\cos \omega_{0} \tau\right)=B \cos \beta \text { and } \phi_{0} \frac{F}{k} \sin \omega_{0} \tau=-\phi_{0} \quad B \sin \beta
\]

Thus

\[
B^{2}=\left(\frac{F}{k}\right)^{2}\left(2-2 \cos \omega_{0} \tau\right)
\]

or

\[
B=2 \frac{F}{k}\left|\sin \frac{\omega_{0} \tau}{2}\right|
\]

\(x(t)\)\\
\includegraphics[max width=\textwidth, alt={}, center]{da0156b6-4133-434e-90d4-ec01b2070afe-057_335_749_311_359}\\
4.85 For the spring \(m g=\kappa \Delta l\)\\
where \(\kappa\) is its stifness coefficient. Thus

\[
\omega_{0}^{2}=\frac{\kappa}{m}=\frac{g}{\Delta l},
\]

The equation of motion of the ball is

Here

\[
\begin{gathered}
\ddot{x}+2 \beta \dot{x}+\omega_{0}^{2} x=\frac{F_{0}}{m} \cos \omega t \\
\lambda=\frac{2 \pi \beta}{\sqrt{\omega_{0}^{2}-\beta^{2}}} \text { or } \frac{\beta}{\omega}=\frac{\lambda / 2 \pi}{\sqrt{1+(\lambda / 2 \pi)^{2}}}
\end{gathered}
\]

To find the solution of the above equation we look for the solution of the auxiliary equation

\[
\ddot{z}+2 \beta \dot{z}+\omega_{0}^{2} z=\frac{F_{0}}{m} e^{i \omega t}
\]

Clearly we can take \(\operatorname{Re} \mathcal{z}=x\). Now we look for a particular integral for \(\mathcal{z}\) of the form

\[
z=A e^{i \omega t}
\]

Thus, substitution gives \(\boldsymbol{A}\) and we get

\[
z=\frac{\left(F_{0} / m\right) e^{i \omega t}}{\omega_{0}^{2}-\omega^{2}+2 i \beta \omega}
\]

so taking the real part

\[
\begin{aligned}
x & =\frac{\left(F_{0} / m\right)\left[\left(\omega_{0}^{2}-\omega^{2}\right) \cos \omega t+2 \beta \omega \sin \omega t\right]}{\left(\omega_{0}^{2}-\omega^{2}\right)^{2}+4 \beta^{2} \omega^{2}} \\
= & \frac{F_{0}}{m} \frac{\cos (\omega t-\varphi)}{\sqrt{\left(\omega_{0}^{2}-\omega^{2}\right)^{2}+4 \beta^{2} \omega^{2}}}, \varphi=\tan ^{-1} \frac{2 \beta \omega}{\omega_{0}^{2}-\omega^{2}}
\end{aligned}
\]

The amplitude of this oscillation is maximum when the denominator is minimum. Tthis happens when\\
\(\omega^{4}-2 \omega_{0}^{2} \omega^{2}+4 \beta^{2} \omega^{2}+\omega_{0}^{4}=\left(\omega^{2}-\omega_{0}^{2}+2 \beta^{2}\right)+4 \beta^{2} \omega_{0}^{2}-4 \beta^{4}\) is minimum. i.e for \(\omega^{2}=\omega_{0}^{2}-2 \beta^{2}\)

Thus

\[
\begin{aligned}
\omega_{\text {res }}^{2} & =\omega_{0}^{2}\left(1-\frac{2 \beta^{2}}{\omega_{0}^{2}}\right) \\
& =\frac{g}{\Delta l}\left[1-\frac{2\left(\frac{\lambda}{2 \pi}\right)^{2}}{1+\left(\frac{\lambda}{2 \pi}\right)^{2}}\right]=\frac{g}{\Delta l} \frac{1-\left(\frac{\lambda}{2 \pi}\right)^{2}}{1+\left(\frac{\lambda}{2 \pi}\right)^{2}}
\end{aligned}
\]

and

\[
\begin{aligned}
a_{r e s} & =\frac{F_{0} / m}{\sqrt{4 \beta^{2} \omega_{0}^{2}-4 \beta^{4}}}=\frac{F_{0} / m}{2 \beta \sqrt{\omega_{0}^{2}-\beta^{2}}}=\frac{F_{0} / m}{2 \beta^{2}} \cdot \frac{\lambda}{2 \pi} \\
& =\frac{F_{0}}{2 m \omega_{0}^{2}} \cdot \frac{1+\left(\frac{\lambda}{2 \pi}\right)^{2}}{\lambda / 2 \pi}=\frac{F_{0} \Delta l \lambda}{4 \pi m g}\left(1+\frac{4 \pi^{2}}{\lambda^{2}}\right)
\end{aligned}
\]

4.86 Since \(a=\frac{F_{0} / m}{\sqrt{\left(\omega^{2}-\omega_{0}^{2}+2 \beta^{2}\right)^{2}+4 \beta^{2}\left(\omega_{0}^{2}-\beta^{2}\right)}}\)\\
we must have

\[
\omega_{1}^{2}-\omega_{0}^{2}+2 \beta^{2}=-\left(\omega_{2}^{2}-\omega_{0}^{2}+2 \beta^{2}\right)
\]

or

\[
\omega_{0}^{2}-2 \beta^{2}=\frac{\omega_{1}^{2}+\omega_{2}^{2}}{2}=\omega_{\mathrm{res}}^{2}
\]

4.87

\[
\begin{aligned}
x & =\frac{F_{0}}{m} \frac{\left(\omega_{0}^{2}-\omega^{2}\right) \cos \omega t+2 \beta \omega \sin \omega t}{\sqrt{\left(\omega^{2}-\omega_{0}^{2}\right)^{2}+4 \beta^{2} \omega^{2}}} \\
\dot{x} & =\frac{F_{0} \omega}{m} \frac{2 \beta \omega \cos \omega t+\left(\omega^{2}-\omega_{0}^{2}\right) \sin \omega t}{\left(\omega_{0}^{2}-\omega^{2}\right)^{2}+4 \beta^{2} \omega^{2}}
\end{aligned}
\]

Then

\[
\begin{aligned}
V_{0} & =\frac{F_{0} \omega}{m \sqrt{\left(\omega_{0}^{2}-\omega^{2}\right)^{2}+4 \beta^{2} \omega^{2}}} \\
& =\frac{F_{0}}{\sqrt{\left(\frac{\omega_{0}^{2}}{\omega}-\omega\right)^{2}+4 \beta^{2}}}
\end{aligned}
\]

This is maximum when

\[
\omega^{2}=\omega_{0}^{2}=\omega_{r e s}^{2}
\]

and then

\[
V_{0 \mathrm{res}}=\frac{F_{0}}{2 m \beta} .
\]

Now at half maximum\\
or

\[
\begin{gathered}
\left(\frac{\omega_{0}^{2}}{\omega}-\omega\right)^{2}=12 \beta^{2} \\
\omega^{2} \pm 2 \sqrt{3} \beta \omega-\omega_{0}^{2}=0 \\
\omega=\mp \beta \sqrt{3}+\sqrt{\omega_{0}^{2}+3 \beta^{2}}
\end{gathered}
\]

where we have rejected a solution with - ve sign before there dical. Writing

\[
\omega_{1}=\sqrt{\omega_{0}^{2}+3 \beta^{2}}+\beta \sqrt{3}, \omega_{2}=\sqrt{\omega_{0}^{2}+3 \beta^{2}}-\beta \sqrt{3}
\]

we get (a) \(\omega_{\text {res }}=\omega_{0}=\sqrt{\omega_{1} \omega_{2}}\) (Velocity resonance frequency)\\
(b) \(\beta=\frac{\left|\omega_{1}-\omega_{2}\right|}{2 \sqrt{3}}\) and damped oscillation frequency

\[
\sqrt{\omega_{0}^{2}-\beta^{2}}=\sqrt{\omega_{1} \omega_{2}-\frac{\left(\omega_{1}-\omega_{2}\right)^{2}}{12}}
\]

4.88 In general for displacement amplitude

Thus

\[
\begin{gathered}
a=\frac{F_{0}}{m} \frac{1}{\sqrt{\left(\omega_{0}^{2}-\omega^{2}\right)^{2}+4 \beta^{2} \omega^{2}}} \\
=\frac{F_{0}}{m} \frac{1}{\sqrt{\left(\omega^{2}-\omega_{0}^{2}+2 \beta^{2}\right)^{2}+4 \beta^{2}\left(\omega_{0}^{2}-\beta^{2}\right)}}
\end{gathered}
\]

But

\[
\eta=\frac{a_{\text {res }}}{a_{\text {low }}}=\frac{\omega_{0}^{2}}{\sqrt{4 \beta^{2}\left(\omega_{0}^{2}-\beta^{2}\right)}}=\frac{\omega_{0}^{2}}{2 \beta \sqrt{\omega_{0}^{2}-\beta^{2}}}
\]

Hence

\[
\begin{aligned}
\frac{\beta}{\omega_{0}} & =\frac{\lambda / 2 \pi}{\sqrt{1+(\lambda / 2 \pi)^{2}}}, \frac{\lambda}{2 \pi}=\frac{\beta}{\sqrt{\omega_{0}^{2}-\beta^{2}}} \\
\eta & =\frac{\omega_{0}^{2}}{2 \beta^{2}} \cdot \frac{\lambda}{2 \pi}=\frac{1}{2} \frac{1+\left(\frac{\lambda}{2 \pi}\right)^{2}}{\frac{\lambda}{2 \pi}}=2.90
\end{aligned}
\]

4.89 The work done in one cycle is

\[
\begin{aligned}
A & =\int^{T} F d x=\int_{0}^{T} F \mathbf{v} d t=\int_{0}^{T} F_{0} \cos \omega t(-\omega a \sin (\omega t-\varphi)) d t \\
& =\int_{0}^{T} F_{0} \omega a\left(-\cos \omega t \sin \omega t \cos \varphi+\cos ^{2} \omega t \sin \varphi\right) d t \\
& =\frac{1}{2} F_{0} \omega a \frac{T}{2} \sin \varphi=\pi a F_{0} \sin \varphi
\end{aligned}
\]

4.90 In the formula \(x=a \cos (\omega t-\varphi)\)\\
we have

\[
\begin{gathered}
a=\frac{F_{0}}{m} \frac{1}{\sqrt{\left(\omega_{0}^{2}-\omega^{2}\right)^{2}+4 \beta^{2} \omega^{2}}} \\
\tan \varphi=\frac{2 \beta \omega}{\omega_{0}^{2}-\omega^{2}} \\
\beta=\frac{\left(\omega_{0}^{2}-\omega^{2}\right) \tan \varphi}{2 \omega} . \\
\omega_{0}=\sqrt{K / m}=20 \mathrm{~s}^{-1} .
\end{gathered}
\]

Thus

Hence\\
a\\
and (a) the quality factor

\[
Q=\frac{\pi}{\beta T}=\frac{\sqrt{\omega_{0}^{2}-\beta^{2}}}{2 \beta}=\frac{1}{2} \sqrt{\frac{4 \omega^{2} \omega_{0}^{2}}{\left(\omega_{0}^{2}-\omega^{2}\right)^{2} \tan ^{2} \varphi}-1}=2.17
\]

(b) work done is \(A=\pi a F_{0} \sin \varphi\)

\[
\begin{aligned}
& =\pi m a^{2} \sqrt{\left(\omega_{0}^{2}-\omega^{2}\right)^{2}+4 \beta^{2} \omega^{2}} \sin \varphi=\pi m a^{2} \times 2 \beta \omega . \\
& =\pi m a^{2}\left(\omega_{0}^{2}-\omega^{2}\right) \tan \varphi=6 \mathrm{~mJ} .
\end{aligned}
\]

4.91 Here as usual \(\tan \varphi=\frac{2 \beta \omega}{\omega_{0}^{2}-\omega^{2}}\) where \(\varphi\) is the phase lag of the displacement

\[
x=a \cos (\omega t-\varphi), a=\frac{F_{0}}{m} \frac{1}{\sqrt{\left(\omega_{0}^{2}-\omega^{2}\right)^{2}+4 \beta^{2} \omega^{2}}}
\]

(a) Mean power developed by the force over one oscillation period

\[
\begin{aligned}
& =\frac{\pi F_{0} a \sin \varphi}{T}=\frac{1}{2} F_{0} a \omega \sin \varphi \\
& =\frac{F_{0}^{2}}{m} \frac{\beta \omega^{2}}{\left(\omega_{0}^{2}-\omega^{2}\right)^{2}+4 \beta^{2} \omega^{2}}=\frac{F_{0}^{2} \beta}{m} \frac{1}{\left(\frac{\omega_{0}^{2}}{\omega}-\omega\right)^{2}+4 \beta^{2}}
\end{aligned}
\]

(b) Mean power \(\langle P\rangle\) is maximum when \(\omega=\omega_{0}\) (for the denominator is then minimum Also

\[
<P>_{\max }=\frac{F_{0}^{2}}{4 m \beta}
\]

4.92 Given \(\beta=\omega_{0} / \eta\). Then from the previous problem

\[
<P>=\frac{F_{0}^{2} \omega_{0}}{\eta m} \cdot \frac{1}{\left(\frac{\omega_{0}^{2}}{\omega}-\omega\right)^{2}+4 \frac{\omega_{0}^{2}}{\eta^{2}}}
\]

At displacement resonance \(\omega=\sqrt{\omega_{0}^{2}-2 \beta^{2}}\)

\[
\begin{aligned}
<P>_{\text {res }} & =\frac{F_{0}^{2} \omega_{0}}{\eta m} \frac{1}{\frac{4 \beta^{4}}{\omega_{0}^{2}-2 \beta^{2}}+\frac{4 \omega_{0}^{2}}{\eta^{2}}}=\frac{F_{0}^{2} \omega_{0}}{\eta m} \frac{1}{\frac{4 \omega_{0}^{4} / \eta^{4}}{\omega_{0}^{2}\left(1-\frac{2}{\eta^{2}}\right)}+4 \frac{\omega_{0}^{2}}{\eta^{2}}} \\
& =\frac{F_{0}^{2}}{4 \eta m \omega_{0}} \frac{\eta^{2}}{\frac{1}{\eta^{2}-2}+1}=\frac{F_{0}^{2} \eta}{4 m \omega_{0}} \frac{\eta^{2}-2}{\eta^{2}-1}
\end{aligned}
\]

while

\[
<P>_{\max }=\frac{F_{0}^{2} \eta}{4 m \omega_{0}} .
\]

Thus

\[
\frac{\langle P\rangle_{\max }-\langle P\rangle_{\mathrm{res}}}{\langle P\rangle_{\max }}=\frac{100}{\eta^{2}-1} \%
\]

4.93 The equation of the disc is \(\ddot{\varphi}+2 \beta \dot{\varphi}+\omega_{0}^{2} \varphi=\frac{N_{m} \cos \omega t}{I}\)

Then as before

\[
\varphi=\varphi_{m} \cos (\omega t-\alpha)
\]

where

\[
\varphi_{m}=\frac{N_{m}}{I\left[\left(\omega_{0}^{2}-\omega^{2}\right)^{2}+4 \beta^{2} \omega^{2}\right]^{1 / 2}}, \tan \alpha=\frac{2 \beta \omega}{\omega_{0}^{2}-\omega^{2}}
\]

(a) Work performed by frictional forces

\[
\begin{aligned}
& =-\int N_{r} d \varphi \quad \text { where } N_{r}=-2 I \beta \dot{\varphi}=-\int_{0}^{T} 2 \beta I \dot{\varphi}^{2} d t=-2 \pi \beta \omega I \varphi_{m}^{2} \\
& =-\pi I \varphi_{m}^{2}\left[\left(\omega_{0}^{2}-\omega^{2}\right)^{2}+4 \beta^{2} \omega^{2}\right]^{1 / 2} \sin \alpha=-\pi N_{m} \varphi_{m} \sin \alpha
\end{aligned}
\]

(b) The quality factor

\[
\begin{gathered}
Q=\frac{\pi}{\lambda}=\frac{\pi}{\beta T}=\frac{\sqrt{\omega_{0}^{2}-\beta^{2}}}{2 \beta}=\frac{\omega \sqrt{\omega_{0}^{2}-\beta^{2}}}{\left(\omega_{0}^{2}-\omega^{2}\right) \tan \alpha}=\frac{1}{2 \tan \alpha}\left\{\frac{4 \omega^{2} \omega_{0}^{2}}{\left(\omega_{0}^{2}-\omega^{2}\right)^{2}}-\frac{4 \beta^{2} \omega^{2}}{\left(\omega_{0}^{2}-\omega^{2}\right)^{2}}\right\}^{1 / 2} \\
=\frac{1}{2 \tan \alpha}\left\{\frac{4 \omega^{2} \omega_{0}^{2} I^{2} \varphi_{m}^{2}}{N_{m}^{2} \cos ^{2} \alpha}-\tan ^{2} \alpha\right\}^{1 / 2} \operatorname{since} \omega_{0}^{2}=\omega^{2}+\frac{N_{m}}{I \varphi_{m}} \cos \alpha \\
=\frac{1}{2 \sin \alpha}\left\{\frac{4 \omega^{2} \omega_{0}^{2} I^{2} \varphi_{m}^{2}}{N_{m}^{2}}-\sin ^{2} \alpha\right\}^{1 / 2} \\
=\frac{1}{2 \sin \alpha}\left\{\frac{4 \omega^{2} I^{2} \varphi_{m}^{2}}{N_{m}^{2}}\left(\omega^{2}+\frac{N_{m} \cos \alpha}{I \varphi_{m}}\right)+1-\cos ^{2} \alpha\right\}^{1 / 2} \\
=\frac{1}{2 \sin \alpha}\left\{\frac{4 I^{2} \varphi_{m}^{2}}{N_{m}^{2}} \omega^{4}+\frac{4 I \varphi_{m}}{N_{m}} \omega^{2} \cos \alpha+\cos ^{2} \alpha-1\right\}^{1 / 2}=\frac{1}{2 \sin \alpha}\left\{\left(\frac{2 I \varphi_{m} \omega^{2}}{N_{m}}+\cos \alpha\right)^{2}-1\right\}^{1 / 2}
\end{gathered}
\]

\subsection*{4.2 ELECTRIC OSCILLATIONS}
4.94 If the electron (charge of each electron \(=-e\) ) are shifted by a sman distance \(x\), a net \(+v e\) charge density (per unit area) is induced on the surface. This will result in an electric field \(E=n e x / \varepsilon_{0}\) in the direction of \(x\) and a restoring force on an electron of \(-\frac{n e^{2} x}{\varepsilon_{0}}\),

Thus\\
or

\[
\begin{gathered}
m \ddot{x}=-\frac{n e^{2} x}{\varepsilon_{0}} \\
\ddot{x}+\frac{n e^{2}}{m \varepsilon_{0}} x=0 \\
\omega_{p}=\sqrt{\frac{n e^{2}}{m \varepsilon_{0}}}=1.645 \times 10^{16} \mathrm{~s}^{-1} .
\end{gathered}
\]

This gives\\
as the plasma frequency for the problem.\\
4.95 Since there are no sources of emf in the circuit, Ohm's 1 law reads

\[
\frac{q}{c}=-L \frac{d I}{d t}
\]

where \(q=\) change on the capacitor, \(I=\frac{d q}{d t}=\) current through the coil. Then

\[
\frac{d^{2} q}{d t^{2}}+\omega_{0}^{2} q=0, \omega_{0}^{2}=\frac{1}{L C}
\]

The solution fo this equation is

\[
q=q_{m} \cos \left(\omega_{0} t+\alpha\right)
\]

From the problem \(V_{m}=\frac{q_{m}}{C}\). Then

\[
I=-\omega_{0} C V_{m} \sin \left(\omega_{0} t+\alpha\right)
\]

and

\[
\begin{gathered}
V=V_{m} \cos \left(\omega_{0} t-\alpha\right) \\
V^{2}+\frac{I^{2}}{\omega_{0}^{2} C^{2}}=V_{m}^{2} \\
V^{2}+\frac{L I^{2}}{C}=V_{m}^{2}
\end{gathered}
\]

or

\[
\frac{1}{2} L I^{2}+\frac{q^{2}}{2 C}=\text { constant }
\]

By energy conservation\\
When the P.D. across the capacitor takes its maximum value \(V_{m}\), the current \(I\) must be zero.\\
Thus "constant" \(=\frac{1}{2} C V_{m}^{2}\)

Hence

\[
\frac{L I^{2}}{C}+V^{2}=V_{m}^{2} \text { once again. }
\]

4.96 After the switch was closed, the circuit satisfies\\
or

\[
\begin{gathered}
-L \frac{d I}{d t}=\frac{q}{C} \\
\frac{d^{2} q}{d t^{2}}+\omega_{0}^{2} q=0 \Rightarrow q=C V_{m} \cos \omega_{0} t
\end{gathered}
\]

where we have used the fact that when the switch is closed we must have

\[
V=\frac{q}{C}=V_{m}, I=\frac{d q}{d t}=0 \text { at } t=0 .
\]

Thus (a)

\[
\begin{aligned}
I & =\frac{a q}{d t}=-C V_{m} \omega_{0} \sin \omega_{0} t \\
& =-V_{m} \sqrt{\frac{C}{L}} \sin \omega_{0} t
\end{aligned}
\]

(b) The electrical energy of the capacitor is \(\frac{q^{2}}{2 C} \alpha \cos ^{2} \omega_{0} t\) and of the inductor is

\[
\frac{1}{2} L I^{2} \alpha \sin ^{2} \omega_{0} t
\]

The two are equal when

\[
\omega_{0} t=\frac{\pi}{4}
\]

At that instant the emf of the self-inductance is

\[
-L \frac{d i}{d t}=V_{m} \cos \omega_{0} t=V_{m} / \sqrt{2}
\]

4.97 In the oscillating circuit, let

\[
q=q_{m} \cos \omega t
\]

be the change on the condenser where\\
\(\omega^{2}=\frac{1}{L C}\) and \(C\) is the instantaneous capacity of the condenser ( \(S=\) area of plates)

\[
C=\frac{\varepsilon_{0} S}{y}
\]

\(y=\) distance between the plates. Since the oscillation frequency increases \(\eta\) fold, the quantity

\[
\omega^{2}=\frac{y}{\varepsilon_{0} S L}
\]

changes \(\eta^{2}\) fold and so does \(y\) i.e. changes from \(y_{0}\) initially to \(\eta^{2} y_{0}\) finally. Now the P.D. across the condenser is

\[
V=\frac{q_{m}}{C} \cos \omega t=\frac{y q_{m}}{\varepsilon_{0} S} \cos \omega t
\]

and hence the electric field between the plates is

\[
E=\frac{q_{m}}{\varepsilon_{0} S} \cos \omega t
\]

Thus, the charge on the plate being \(q_{m} \cos \omega t\), the force on the plate is

\[
F=\frac{q_{m}^{2}}{\varepsilon_{0} S} \cos ^{2} \omega t
\]

Since this force is always positive and the plate is pulled slowly we can use the average force

\[
\bar{F}=\frac{q_{m}^{2}}{2 \varepsilon_{0} S}
\]

and work done is

\[
A=\bar{F}\left(\eta^{2} y_{0}-y_{0}\right)=\left(\eta^{2}-1\right) \frac{q_{m}^{2} y_{0}}{2 \varepsilon_{0} S}
\]

But \(\frac{q_{m}^{2} y_{0}}{2 \varepsilon_{0} S}=\frac{q_{m}^{2}}{2 C_{0}}=W\) the initial stored energy. Thus .

\[
A=\left(\eta^{2}-1\right) W .
\]

4.98 The equations of the \(L-C\) circuit are

\[
L \frac{d}{d t}\left(I_{1}+I_{)}=\frac{C_{1} V-\int I_{1} d t}{C_{1}}=\frac{C_{2} V-\int I_{2} d t}{C_{2}}\right.
\]

Differentiating again

\[
L\left(I_{1}+I_{2}\right)=-\frac{1}{C_{1}} I_{1}=-\frac{1}{C_{2}} I_{2}
\]

Then

\[
\begin{gathered}
I_{1}=\frac{C_{1}}{C_{1}-C_{2}} I, I_{2}=\frac{C_{2}}{C_{1}+C_{2}} I \\
I=I_{1}+I_{2}
\end{gathered}
\]


\begin{equation*}
L\left(C_{1}+C_{2}\right) I+I=0 \tag{so}
\end{equation*}


or

\[
I=I_{0} \sin \left(\omega_{0} t+\alpha\right)
\]

where \(\quad \omega_{0}^{2}=\frac{1}{L\left(C_{1}+C_{2}\right)}\) (Part a )\\
(Hence \(T=\frac{2 \pi}{\omega_{0}}=0.7 \mathrm{~ms}\) )\\
\includegraphics[max width=\textwidth, alt={}, center]{da0156b6-4133-434e-90d4-ec01b2070afe-064_279_480_1126_916}

At \(\quad t=0, I=0\) so \(\alpha=0\)

\[
I=I_{0} \sin \omega_{0} t
\]

The peak value of the current is \(I_{0}\) and it is related to the voltage \(V\) by the first equation\\
or

\[
\begin{gathered}
L I=V-\int I d t /\left(C_{1}+C_{2}\right) \\
+L \omega_{0} I_{0} \cos \omega_{0} t=V-\frac{1}{C_{1}+C_{2}} \int_{0}^{t} I_{0} \sin \omega_{0} t d t
\end{gathered}
\]

(The P.D. across the inductance is \(V\) at \(t=0\) )

\[
=V+\frac{1}{C_{1}+C_{2}} \cdot \frac{I_{0}}{\omega_{0}}\left(\cos \omega_{0} t-1\right)
\]

Hence

\[
I_{0}=\left(C_{1}+C_{2}\right) \omega_{0} V=V \sqrt{\frac{C_{1}+C_{2}}{L}}=8.05 \mathrm{~A} .
\]

4.99 Initially \(q_{1}=C V_{0}\) and \(q_{2}=0\). After the switch is closed change flows and we get


\begin{gather*}
q_{1}+q_{2}=C V_{0} \\
\frac{q_{1}}{C}+L \frac{d I}{d t}-\frac{q_{2}}{C}=0 \tag{1}
\end{gather*}


Also \(I=\dot{q}_{1}=-\dot{q}_{2}\). Thus

\[
L \ddot{I}+\frac{2 I}{C}=0
\]

Hence \(\ddot{I}+\omega_{0}^{2} I=0 \quad \omega_{0}^{2}=\frac{2}{L C}\),\\
The soloution of this equation subject to

\[
I=0 \text { at } t=0
\]

\includegraphics[max width=\textwidth, alt={}, center]{da0156b6-4133-434e-90d4-ec01b2070afe-065_349_589_252_780}\\
is \(\quad I=I_{0} \sin \omega_{0} t\).\\
Integrating

\[
\begin{aligned}
& q_{1}=A-\frac{I_{0}}{\omega_{0}} \cos \omega_{0} t \\
& q_{2}=B+\frac{I_{0}}{\omega_{0}} \cos \omega_{0} t
\end{aligned}
\]

Finally substituting in (1)

\[
\frac{A-B}{C}-\frac{2 I_{0}}{\omega_{0} C} \cos \omega_{0} t+L I_{0} \omega_{0} \cos \omega_{0} t=0
\]

Thus

\[
\begin{gathered}
A=B=\frac{C V_{0}}{2} \text { and } \\
\frac{C V_{0}}{2}+\frac{I_{0}}{\omega_{0}}=0
\end{gathered}
\]

so

\[
\begin{aligned}
& q_{1}=\frac{C V_{0}}{2}\left(1+\cos \omega_{0} t\right) \\
& q_{2}=\frac{C V_{0}}{2}\left(1-\cos \omega_{0} t\right)
\end{aligned}
\]

4.100 The flux in the coil is

\[
\Phi(t)=\left\{\begin{array}{cc}
\Phi & t<0 \\
0 & t>0
\end{array}\right.
\]

The equation of the current is \(\quad-L \frac{d I}{d t}=\frac{\int_{0}^{t} I d t}{C}\)

This mean that


\begin{equation*}
L C \frac{d^{2} I}{d t^{2}}+I=0 \tag{1}
\end{equation*}


or with

\[
\omega_{0}^{2}=\frac{1}{L C} \quad I=I_{0} \sin \left(\omega_{0} t+\alpha\right)
\]

Putting in (1) \(-L I_{0} \omega_{0} \cos \left(\omega_{0} t+\alpha\right)=-\frac{I_{0}}{\omega_{0} C}\left[\cos \left(\omega_{0} t+\alpha\right)-\cos \alpha\right]\)\\
This implies \(\cos \alpha=0 \quad \therefore \quad I= \pm I_{0} \cos \omega_{0} t\). From Faraday's law

\[
\varepsilon=-\frac{d \Phi}{d t}=-L \frac{d I}{d t}
\]

or integrating from \(t=-\varepsilon\) to \(-\varepsilon\) where \(\varepsilon \rightarrow 0\)

\[
\begin{gathered}
\Phi=L I_{0} \text { with }+\operatorname{sign} \text { in } I \\
I=\frac{\Phi}{L} \cos \omega_{0} t .
\end{gathered}
\]

4.101 Given \(V=V_{m} e^{-\beta t} \cos \omega t\)\\
(a) The phrase 'peak values' is not clear. The answer is obtained on taking \(|\cos \omega t|=1\)\\
i.e

\[
t=\frac{\pi n}{\omega} .
\]

(b) For extrema \(\frac{d V}{d t}=0\)\\
so,\\
Given \(V=V_{m} e^{-\beta t} \cos \omega t\)\\
(a) The phrase 'peak values' is not\\
i.e clea

\[
\begin{array}{lr}
\text { or } & -\beta \cos \omega t-\omega \sin \omega t=0 \\
\tan \omega t=-\beta / \omega \\
\text { i.e. } & \omega t=n \pi+\tan ^{-1}\left(\frac{-\beta}{\omega}\right) .
\end{array}
\]

4.102 The equation of the circuit is

\[
L \frac{d^{2} Q}{d t^{2}}+R \frac{d Q}{d t}+\frac{Q}{C}=0
\]

where \(Q=\) charge on the capacitor,\\
This has the solution

\[
Q=Q_{m} e^{-\beta t} \sin (\omega t+\alpha)
\]

where

\[
\beta=\frac{R}{2 L}, \omega=\sqrt{\omega_{0}^{2}-\beta^{2}}, \omega_{0}^{2}=\frac{1}{L C} .
\]

Now

\[
I=\frac{d Q}{d t}=0 \quad \text { at } t=0
\]

so,

\[
Q_{m} e^{-\beta t}(-\beta \sin (\omega t+\alpha)+\omega \cos (\omega t+\alpha))=0 \text { at } t=0
\]

Thus

\[
\omega \cos \alpha=\beta \sin \alpha \quad \text { or } \alpha=\tan ^{-1} \frac{\omega}{\beta}
\]

Now

\[
\begin{gathered}
V_{m}=\frac{Q_{m}}{C} \text { and } V_{0}=\text { P.D. at } t=0=\frac{Q_{m}}{C} \sin \alpha \\
\therefore \quad \frac{V_{0}}{V_{m}}=\sin \alpha=\frac{\omega}{\sqrt{\omega^{2}+\beta^{2}}}=\frac{\omega}{\omega_{0}}=\sqrt{1-\beta^{2} / \omega_{0}^{2}}=\sqrt{1-\frac{R^{2} C}{4 L^{2}}}
\end{gathered}
\]

4.103 We write

\[
\begin{aligned}
-\frac{d Q}{d t}=I & =I_{m} e^{-\beta t} \sin \omega t \\
& =g m I_{m} e^{-\beta t+i \omega t} \quad \text { (gm means imaginary part) }
\end{aligned}
\]

Then

\[
\begin{aligned}
Q & =g m I_{m} \frac{e^{-\beta t+i \omega t}}{-\beta+i \omega} \\
Q & =g m I_{m} \frac{e^{-\beta t+i \omega t}}{\beta-i \omega} \\
& =g m I_{m} \frac{(\beta+i \omega) e^{-\beta t+i \omega t}}{\beta^{2}+\omega^{2}} \\
& =I_{m} e^{-\beta t} \frac{\beta \sin \omega t+\omega \cos \omega t}{\beta^{2}+\omega^{2}} \\
& =I_{m} e^{-\beta t} \frac{\sin (\omega t+\delta)}{\sqrt{\beta^{2}+\omega^{2}}}, \tan \delta=\frac{\omega}{\beta} .
\end{aligned}
\]

( An arbitrary constant of integration has been put equal to zero.)\\
Thus

\[
\begin{aligned}
V & =\frac{Q}{C}=I_{m} \sqrt{\frac{L}{C}} e^{-\beta t} \sin (\omega t+\delta) \\
V(0) & =I_{m} \sqrt{\frac{L}{C}} \sin \delta=I_{m} \sqrt{\frac{L}{C}} \frac{\omega}{\sqrt{\omega^{2}+\beta^{2}}} \\
& =I_{m} \sqrt{\frac{L}{C\left(1+\beta^{2} / \omega^{2}\right)}}
\end{aligned}
\]

4.104 \(I=I_{m} e^{-\beta t} \sin \omega t\)

\[
\beta=\frac{R}{2 L}, \omega_{0}=\sqrt{\frac{1}{L C}}, \omega=\sqrt{\omega_{0}^{2}-\beta^{2}}
\]

\(I=-\dot{q}, q=\) charge on the capacitor\\
Then

\[
q=I_{m} e^{-\beta t} \frac{\sin (\omega t+\delta)}{\sqrt{\omega^{2}+\beta^{2}}}, \tan \delta=\frac{\omega}{\beta} .
\]

Thus

\[
W_{M}=\frac{1}{2} L I_{m}^{2} e^{-2 \beta t} \sin ^{2} \omega t
\]

\[
W_{E}=\frac{I_{m}^{2}}{2 C} \frac{e^{-2 \beta t} \sin ^{2}(\omega t+\delta)}{\omega^{2}+\beta^{2}}=\frac{L I_{m}^{2}}{2} e^{-2 \beta t} \sin ^{2}(\omega t+\delta)
\]

Current is maximum when

\[
\frac{d}{d t} e^{-\beta t} \sin \omega t=0
\]

Thus

\[
-\beta \sin \omega t+\omega \cos \omega t=0
\]

or

\[
\tan \omega t=\frac{\omega}{\beta}=\tan \delta
\]

i.e.

\[
\omega t=n \pi+\delta
\]

and hence

\[
\begin{aligned}
\frac{W_{M}}{W_{E}} & =\frac{\sin ^{2}(\omega t)}{\sin ^{2}(\omega t+\delta)}=\frac{\sin ^{2} \delta}{\sin ^{2} 2 \delta}=\frac{1}{4 \cos ^{2} \delta} \\
& =\frac{1}{4 \beta^{2} / \omega_{0}^{2}}=\frac{\omega_{0}^{2}}{4 \beta^{2}}=\frac{1}{L C} \times \frac{L^{2}}{R^{2}}=\frac{L}{C R^{2}}=5 .
\end{aligned}
\]

( \(W_{M}\) is the magnetic energy of the inductance coil and \(W_{E}\) is the electric energy of U capacitor.)\\
4.105 Clearly

\[
L=L_{1}+L_{2}, R=R_{1}+R_{2}
\]

4.106

\[
Q=\frac{\pi}{\beta T} \quad \text { or } \quad \beta=\frac{\pi}{Q T}
\]

Now

\[
\begin{aligned}
\beta t=\ln \eta \quad \text { so } t & =\frac{\ln \eta}{\pi} Q T \\
& =\frac{Q \ln \eta}{\pi v}=0.5 \mathrm{~ms}
\end{aligned}
\]

4.107

Current decreases e fold in time

\[
\begin{aligned}
t=\frac{1}{\beta} & =\frac{2 L}{R} \sec =\frac{2 L}{R T} \text { oscillations } \\
& =\frac{2 L}{R} \frac{\omega}{2 \pi} \\
& =\frac{L}{\pi R} \sqrt{\frac{1}{L C}-\frac{R^{2}}{4 L^{2}}}=\frac{1}{2 \pi} \sqrt{\frac{4 L}{R^{2} C}-1}=15.9 \text { oscillations }
\end{aligned}
\]

\(4.108 Q=\frac{\pi}{\beta T}=\frac{\omega}{2 \beta}\)

\[
\therefore \omega=2 \beta Q, \beta=\frac{\omega}{2 Q}
\]

Now

\[
\omega_{0}=\omega \sqrt{1+\frac{1}{4 Q^{2}}} \quad \text { or } \quad \omega=\frac{\omega_{0}}{\sqrt{1+\frac{1}{4 Q^{2}}}}
\]

\[
\left|\frac{\omega_{0}-\omega}{\omega_{0}}\right| \times 100 \% \approx \frac{1}{8 Q^{2}} \times 100 \%=0.5 \%
\]

\begin{center}
\includegraphics[max width=\textwidth, alt={}]{da0156b6-4133-434e-90d4-ec01b2070afe-069_235_633_79_355}
\end{center}

At \(t=0\) current through the coil \(=\frac{\varepsilon}{R+r}\)\\
P.D. across the condenser \(=\frac{\varepsilon}{R+r}\)\\
(a) At \(t=0\), energy stored \(=W_{0}\)

\[
=\frac{1}{2} L\left(\frac{\varepsilon}{R+r}\right)^{2}+\frac{1}{2} C\left(\frac{\varepsilon R}{R+r}\right)^{2}=\frac{1}{2} \varepsilon^{2} \frac{\left(L+C \dot{R}^{2}\right)}{(R+r)^{2}}=2.0 \mathrm{~mJ} .
\]

(b) The current and the change stored decrease as \(e^{-t R / 2 L}\) so energy decreases as \(e^{-t R / L}\)

\[
\therefore \quad W=W_{0} e^{-t R / L}=0.10 \mathrm{~mJ} .
\]

\(4.110 Q=\frac{\pi}{\beta T}=\frac{\pi v}{\beta}=\frac{\omega}{2 \beta}=\frac{\sqrt{\omega_{0}^{2}-\beta^{2}}}{2 \beta}\)\\
or

\[
\frac{\omega_{0}}{\beta}=\sqrt{1+4 Q^{2}} \quad \text { or } \beta=\frac{\omega_{0}}{\sqrt{1+Q^{2}}}
\]

Now

\[
W=W_{0} e^{-2 \beta t}
\]

Thus energy decreases \(\eta\) times in \(\frac{\ln \eta}{\mathbf{2 \beta}}\) sec.

\[
=\ln \eta \frac{\sqrt{1+4 \dot{Q}^{2}}}{2 \omega_{0}} \approx \frac{Q \ln \eta}{2 \pi v_{n}} \sec .=1.033 \mathrm{~ms} .
\]

4.111 In a leaky condenser

\[
\frac{d q}{d t}=I-I^{\prime} \quad \text { where } I^{\prime}=\frac{V}{R}=\text { leak current }
\]

Now

\[
\begin{aligned}
V=\frac{q}{C} & =-L \frac{d I}{d t}=-L \frac{d}{d t}\left(\frac{d q}{d t}+\frac{V}{R}\right) \\
& =-L \frac{d^{2} q}{d t^{2}}-\frac{L}{R C} \frac{d q}{d t}
\end{aligned}
\]

or

\[
\ddot{q}+\frac{1}{R C} \frac{d Q}{d t}+\frac{1}{L C} q=0
\]

Then

\[
q=q_{m} e^{-\beta t} \sin (\omega t+\alpha)
\]

(a) \(\beta=\frac{1}{2 R C}, \omega_{0}^{2}=\frac{1}{L C}, \omega=\sqrt{\omega_{0}^{2}-\beta^{2}}\)

\[
=\sqrt{\frac{1}{L C}-\frac{1}{4 R^{2} C^{2}}}
\]

(b) \(Q=\frac{\omega}{2 \beta}=R C \sqrt{\frac{1}{L C}-\frac{1}{4 R^{2} C^{2}}}\)

\[
=\frac{1}{2} \sqrt{\frac{4 C R^{2}}{L}-1}
\]

4.112 Given \(V=V_{m} e^{-\beta t} \sin \omega t, \omega \cong \omega_{0} \beta T \ll 1\)

\[
\begin{aligned}
\text { Power loss } & =\frac{\text { Energy loss per cycle }}{T} \\
& \propto \frac{1}{2} C V_{m}^{2} \times 2 \beta
\end{aligned}
\]

(energy decreases as \(W_{0} e^{-2 \beta t}\) so loss per cycle is \(W_{0} \times 2 \beta T\) )\\
Thus

\[
\begin{aligned}
<P> & =\frac{1}{2} C V_{m}^{2} \times \frac{R}{L} \\
R & =\frac{2<P>}{V_{m}^{2}} \frac{L}{C}
\end{aligned}
\]

or

Hence

\[
Q=\frac{1}{R} \sqrt{\frac{L}{C}}=\sqrt{\frac{C}{L}} \frac{V_{m}^{2}}{2\langle P\rangle}=100 \text { on putting the vales. }
\]

4.113 Energy is lost across the resistance and the mean power lass is

\[
\langle P\rangle=R\left\langle I^{2}\right\rangle=\frac{1}{2} R I_{m}^{2}=0.2 \mathrm{~mW} .
\]

This power should be fed to the circuit to maintain undamped oscillations.\\
4.114 \(\langle P\rangle=\frac{R C V_{m}^{2}}{2 L}\) as in (4.112). We get \(\langle P\rangle=5 \mathrm{~mW}\).\\
4.115 Given \(q=q_{1}+q_{2}\)

\[
\begin{aligned}
I_{1} & =-\dot{q}_{1}, I_{2}=-\dot{q}_{2} \\
L I_{1} & =R I_{2}=\frac{q}{C}
\end{aligned}
\]

Thus \(C L \ddot{q}_{1}+\left(q_{1}+q_{2}\right)=0\)\\
\(R C \dot{q}_{2}+q_{1}+q_{2}=0\)\\
\includegraphics[max width=\textwidth, alt={}, center]{da0156b6-4133-434e-90d4-ec01b2070afe-070_315_409_1527_899}

Putting \(q_{1}=A e^{i \omega t} \quad q_{2}=B e^{+i \omega t}\)

\[
\left(1-\omega^{2} L C\right) A+B=0
\]

\[
A+(1+i \omega R C) B=0
\]

A solution exists only if\\
or

\[
\left(1-\omega^{2} L C\right)(1+i \omega R C)=1
\]

o

\[
i \omega R C-\omega^{2} L C-i \omega^{3} L R C^{2}=0
\]

or

\[
L R C^{2} \omega^{2}-i \omega L C-R C=0
\]

\[
\omega^{2}-i \omega \frac{1}{R C}-\frac{1}{L C}=0
\]

\[
\omega=\frac{i}{2 R C} \pm \sqrt{\frac{1}{L C}-\frac{1}{4 R^{2} C^{2}}} \approx i \beta \pm \omega_{0}
\]

Thus

\[
q_{1}=\left(A_{1} \cos \omega_{0} t+A_{2} \sin \omega_{0} t\right) e^{-\beta t} \text { etc. }
\]

\(\omega_{0}\) is the oscillation frequency. Oscillations are possible only if \(\omega_{0}^{2}>0\)\\
i.e.

\[
\frac{1}{4 R^{2}}<\frac{C}{L}
\]

4.116 We have

\[
\begin{aligned}
L_{1} \dot{I}_{1}+R_{1} I_{1} & =L_{2} \dot{I}_{2}+R_{2} I_{2} \\
& =-\frac{\int I d t}{C} \\
I & =I_{1}+I_{2}
\end{aligned}
\]

Then differentiating we have the equations\\
\includegraphics[max width=\textwidth, alt={}, center]{da0156b6-4133-434e-90d4-ec01b2070afe-071_331_559_829_755}

\[
\begin{aligned}
& L_{1} C \ddot{I}_{1}+R_{1} C \dot{I}_{1}+\left(I_{1}+I_{2}\right)=0 \\
& L_{2} C \ddot{I}_{2}+R_{2} C \dot{I}_{2}+\left(I_{1}+I_{2}\right)=0
\end{aligned}
\]

Look for a solution

Then

\[
\begin{gathered}
I_{1}=A_{1} e^{\alpha t}, I_{2}=A_{2} e^{\alpha t} \\
\left(1+\alpha^{2} L_{1} C+\alpha R_{1} C\right) A_{1}+A_{2}=0 \\
A_{1}+\left(1+\alpha^{2} L_{2} C+\alpha R_{2} C\right) A_{2}=0
\end{gathered}
\]

This set of simultaneous equations has a nontrivial solution only if\\
or

\[
\begin{gathered}
\left(1+\alpha^{2} L_{1} C+\alpha R_{1} C\right)\left(1+\alpha^{2} L_{2} C+\alpha R_{2} C\right)=1 \\
\alpha^{3}+\alpha^{2} \frac{L_{1} R_{2}+L_{2} R_{1}}{L_{1} L_{2}}+\alpha \frac{L_{1}+L_{2}+R_{1} R_{2} C}{L_{1} L_{2} C}+\frac{R_{1}+R_{2}}{L_{1} L_{2} C}=0
\end{gathered}
\]

This cubic equation has one real root which we ignore and two complex conjugate roots. We require the condition that this pair of complex conjugate roots is identical with the roots of the equation

\[
\alpha^{2} L C+\alpha R C+1=0
\]

The general solution of this problem is not easy. We look for special cases. If \(R_{1}=R_{2}=0\), then \(R=0\) and \(L=\frac{L_{1} L_{2}}{L_{1}+L_{2}}\). If \(L_{1}=L_{2}=0\), then\\
\(L=0\) and \(R=R_{1} R_{2} /\left(R_{1}+R_{2}\right)\). These are the quoted solution but they are misleading.\\
We shall give the solution for small \(R_{1}, R_{2}\). Then we put \(\alpha=-\beta+i \omega\) when \(\beta\) is small\\
We get \(\left(1-\omega^{2} L_{1} C-2 i \beta \omega L_{1} C-\beta \not p_{1} C+i \omega R_{1} C\right)\)

\[
\left(1-\omega^{2} L_{2} C-2 i \beta \omega L_{2} C-\beta \not p_{2} C+i \omega R_{2} C\right)=1
\]

(we neglect \(\left.\beta^{2} \& \beta R_{1}, \beta R_{2}\right)\). Then

\[
\left(1-\omega^{2} L_{1} C\right)\left(1-\omega^{2} L_{2} C\right)=1 \Rightarrow \omega^{2}=\frac{L_{1}+L_{2}}{L_{1} L_{2} C}
\]

This is identical with \(\omega^{2}=\frac{1}{L C}\) if \(L=\frac{L_{1} L_{2}}{L_{1}+L_{2}}\).\\
also \(\quad\left(2 \beta L_{1}-R_{1}\right)\left(1-\omega^{2} L_{2} C\right)+\left(2 \beta L_{2}-R_{2}\right)\left(1-\omega^{2} L_{2} C\right)=0\)\\
This gives \(\beta=\frac{R}{2 L}=\frac{R_{1} L_{2}^{2}+R_{2} L_{1}^{2}}{2 L_{1} L_{2}\left(L_{1}+L_{2}\right)} \Rightarrow R=\frac{R_{1} L_{2}^{2}+R_{2} L_{1}^{2}}{\left(L_{1}+L_{2}\right)^{2}}\).\\
\(4.117 o=\frac{q}{C}+L \frac{d I}{d t}+R I, I=+\frac{d q}{d t}\)\\
For the critical case \(R=2 \sqrt{\frac{L}{C}}\)\\
Thus \(L C \ddot{q}+2 \sqrt{L C} \dot{q}+q=0\)\\
\includegraphics[max width=\textwidth, alt={}, center]{da0156b6-4133-434e-90d4-ec01b2070afe-072_256_541_1068_807}

Look for a solution with \(q \alpha e^{\alpha t}\)

\[
\alpha=-\frac{1}{\sqrt{L C}} .
\]

An independent solution is \(t e^{\alpha t}\). Thus

\[
q=(A+B t) e^{-t / \sqrt{L C}},
\]

At

\[
t=0 q=C V_{0} \text { thus } A=C V_{0}
\]

Also at

\[
t=0 \dot{q}=I=0
\]

\[
0=B-A \frac{1}{\sqrt{l C}} \Rightarrow B=V_{0} \sqrt{\frac{C}{L}}
\]

Thus finally

\[
\begin{aligned}
I & =\frac{d q}{d t}=V_{0} \sqrt{\frac{C}{L}} e^{-t / \sqrt{L C}} \\
-\frac{1}{\sqrt{L C}} & \left(C V_{0}+V_{0} \sqrt{\frac{C}{L}} t\right) e^{-t / \sqrt{L C}} \\
& =-\frac{V_{0}}{L} t e^{-t / \sqrt{L C}}
\end{aligned}
\]

The current has been defined to increase the charge. Hence the minus sign.\\
The current is maximum when

\[
\frac{d I}{d t}=-\frac{V_{0}}{L} e^{-t / \sqrt{L C}}\left(1-\frac{t}{\sqrt{L C}}\right)=0
\]

This gives \(t=\sqrt{L C}\) and the magnitude of the maximum current is

\[
\left|I_{\max }\right|=\frac{V_{0}}{e} \sqrt{\frac{C}{L}} .
\]

4.118 The equation of the circuit is ( \(I\) is the current)

\[
L \frac{d I}{d t}+R I=V_{m} \cos \omega t
\]

From the theory of differential equations

\[
I=I_{P}+I_{C}
\]

where \(I_{P}\) is a particular integral and \(I_{C}\) is the complementary function (Solution of the differential equation with the RHS \(=0\) ). Now

\[
I_{C}=I_{C O} e^{-t R / L}
\]

and for \(I_{P}\) we write \(I_{P}=I_{m} \cos (\omega t-\varphi)\)\\
Substituting we get

\[
I_{m}=\frac{V_{m}}{\sqrt{R^{2}+\omega^{2} L^{2}}}, \varphi=\tan ^{-1} \frac{\omega L}{R}
\]

Thus

\[
I_{m}=\frac{V_{m}}{\sqrt{R^{2}+\omega^{2} L^{2}}} \cos (\omega t-\varphi)+I_{C O} e^{-t R / L}
\]

Now in an inductive circuit \(I=0\) at \(t=0\)\\
because a current cannot change suddenly.\\
Thus

\[
I_{C O}=-\frac{V_{m}}{\sqrt{R^{2}+\omega^{2} L^{2}}} \cos \varphi
\]

and so

\[
I=\frac{V_{m}}{\sqrt{R^{2}+\omega^{2} L^{2}}}\left[\cos (\omega t-\varphi)-\cos \varphi e^{-t R / L}\right]
\]

4.119 Here the equation is ( \(Q\) is charge on the capacitor)

\[
\frac{Q}{C}+R \frac{d Q}{d t}=V_{n} \cos \omega t
\]

A solution subject to \(Q=0\) at \(t=0\) is of the form (as in the previous problem)

\[
Q=Q_{m}\left[\cos (\omega t-\bar{\varphi})-\cos \bar{\varphi} e^{-t / R C}\right]
\]

Substituting back\\
so

\[
\begin{gathered}
\frac{Q_{m}}{C} \cos (\omega t-\bar{\varphi})-\omega R Q_{m} \sin (\omega t-\bar{\varphi}) \\
=V_{m} \cos \omega t \\
=V_{m}\{\cos \bar{\varphi} \cos (\omega t-\bar{\varphi})-\sin \bar{\varphi} \sin (\omega t-\bar{\varphi})\} \\
Q_{m}=C V_{m} \cos \bar{\varphi} \\
\omega R Q_{m}=V_{m} \sin \bar{\varphi}
\end{gathered}
\]

This leads to

\[
Q_{m}=\frac{C V_{m}}{\sqrt{1+(\omega R C)^{2}}}, \tan \bar{\varphi}=\omega R C
\]

Hence

\[
I=\frac{d Q}{d t}=\frac{V_{m}}{\sqrt{R^{2}+\left(\frac{1}{\omega C}\right)^{2}}}\left[-\sin (\omega t-\bar{\varphi})+\frac{\cos ^{2} \bar{\varphi}}{\sin \bar{\varphi}} e^{-t / R C}\right]
\]

The solution given in the book satisfies \(I=0\) at \(t=0\). Then \(Q=0\) at \(t=0\) but this will not satisfy the equation at \(t=0\). Thus \(I \neq 0\), (Equation will be satisfied with \(I=0\) only if \(Q \neq 0\) at \(t=0\) )

With our I,

\[
I(t=0)=\frac{V_{m}}{R}
\]

4.120 The current lags behind the voltage by the phase angle

\[
\varphi=\tan ^{-1} \frac{\omega L}{R}
\]

Now \(L=\mu_{0} n^{2} \pi a^{2} l, \quad l=\) length of the solenoid\\
\(R=\frac{\rho \cdot 2 \pi a n \cdot l}{\pi b^{2}}, 2 b=\) diameter of the wire\\
But

\[
2 b n=1 \quad \therefore \quad b=\frac{1}{2 n}
\]

Then

\[
\begin{aligned}
\varphi & =\tan ^{-1} \frac{\mu_{0} n^{2} l \pi a^{2} \cdot 2 \pi v}{\rho \cdot 2 \pi a n l} \times \pi \frac{1}{4 n^{2}} \\
& =\tan ^{-1} \frac{\mu_{0} \pi^{2} a v}{4 \rho n}
\end{aligned}
\]

\[
\text { Here } \begin{aligned}
V & =V_{m} \cos \omega t \\
I & =I_{m} \cos (\omega t+\varphi)
\end{aligned}
\]

where

\[
\begin{gathered}
I_{m}=\frac{V_{m}}{\sqrt{R^{2}+\left(\frac{1}{\omega C}\right)^{2}}}, \tan \varphi=\frac{1}{\omega R C} \\
R^{2}+\frac{1}{(\omega C)^{2}}=\left(\frac{V_{m}}{I_{m}}\right)^{2} \\
\frac{1}{\omega R C}=\sqrt{\left(\frac{V_{m}}{R I_{m}}\right)^{2}-1}
\end{gathered}
\]

Thus the current is ahead of the voltage by

\[
\varphi=\tan ^{-1} \frac{1}{\omega R C}=\tan ^{-1} \sqrt{\left(\frac{V_{m}}{R I_{m}}\right)^{2}-1}=60^{\circ}
\]

\[
\begin{gathered}
V=I R+\frac{\int_{0}^{t} I d t}{C}: \\
R \dot{I}+\frac{1}{C} I=\dot{V}=-\omega V_{0} \sin \omega t
\end{gathered}
\]

Ignoring transients, a solution has the form

\[
I=I_{0} \sin (\omega t-\alpha)
\]

\begin{center}
\includegraphics[max width=\textwidth, alt={}]{da0156b6-4133-434e-90d4-ec01b2070afe-075_225_574_1005_763}
\end{center}

\[
\begin{aligned}
& \omega R I_{0} \cos (\omega t-\alpha)+\frac{I_{0}}{C} \sin (\omega t-\alpha)=-\omega V_{0} \sin \omega t \\
& =-\omega V_{0}\{\sin (\omega t-\alpha) \cos \alpha+\cos (\omega t-\alpha) \sin \alpha\}
\end{aligned}
\]

so

\[
R I_{0}=-V_{0} \sin \alpha
\]

\[
\begin{aligned}
& \frac{I_{0}}{\omega C}=-V_{0} \cos \alpha \quad \alpha=\pi+\tan ^{-1}(\omega R C) \\
& I_{0}=\frac{V_{0}}{\sqrt{R^{2}+\left(\frac{1}{\omega C}\right)^{2}}} \\
& I=I_{0} \sin \left(\omega t-\tan ^{-1} \omega R C-\pi\right)=-I_{0} \sin \left(\omega t-\tan ^{-1} \omega R C\right) \\
& Q=\int_{0}^{t} I d t=Q_{0}+\frac{I_{o}}{\omega} \cos \left(\omega t-\tan ^{-1} \omega R C\right)
\end{aligned}
\]

Then

\[
V_{0}(1+\cos \omega t)=R \frac{d Q}{d t}+\frac{Q}{C}
\]

It satisfies\\
if

\[
\begin{aligned}
V_{0}(1+\cos \omega t)= & -R I_{0} \sin \left(\omega t-\tan ^{-1} \omega R C\right) \\
& +\frac{Q_{0}}{C}+\frac{I_{0}}{\omega C} \cos \left(\omega t-\tan ^{-1} \omega R C\right)
\end{aligned}
\]

Thus

\[
Q_{0}=C V_{0}
\]

and

\[
\left.\begin{array}{rl}
\frac{I_{0}}{\omega C} & =V_{0} / \sqrt{1+(\omega R C)^{2}} \\
R I_{0} & =\frac{V_{0} \omega R C}{\sqrt{1+(\omega R C)^{2}}}
\end{array}\right\} \text { checks }
\]

Hence

\[
V^{\prime}=\frac{Q}{C}=V_{0}+\frac{V_{0}}{\sqrt{1+(\omega R C)^{2}}} \cos (\omega t-\alpha)
\]

(b) \(\frac{V_{0}}{\eta}=\frac{V_{0}}{\sqrt{1+(\omega R C)^{2}}}\)\\
or\\
or

\[
\begin{gathered}
\eta^{2}-1=\omega^{2}(R C)^{2} \\
R C=\sqrt{\eta^{2}-1} / \omega=22 \mathrm{~ms} .
\end{gathered}
\]

4.123

\begin{figure}[h]
\begin{center}
  \includegraphics[alt={},max width=\textwidth]{da0156b6-4133-434e-90d4-ec01b2070afe-076_383_475_900_188}
\captionsetup{labelformat=empty}
\caption{(a)}
\end{center}
\end{figure}

\begin{figure}[h]
\begin{center}
  \includegraphics[alt={},max width=\textwidth]{da0156b6-4133-434e-90d4-ec01b2070afe-076_419_431_879_948}
\captionsetup{labelformat=empty}
\caption{(b)}
\end{center}
\end{figure}

(b)\\
\(\tan \varphi=\frac{\omega L-\frac{1}{\omega C}}{R}=-v e\)\\
as

\[
\omega^{2}<\frac{1}{L C}
\]

4.124 (a) \(I_{m}=\frac{V_{m}}{\sqrt{R^{2}+\left(\omega L-\frac{1}{\omega C}\right)^{2}}} \approx 4.48 \mathrm{~A}\)\\
(b) \(\tan \varphi=\frac{\omega L-\frac{1}{\omega C}}{R}, \varphi \approx-60^{\circ}\)

Current lags behind the voltage \(V\) by \(\varphi\)\\
(c) \(V_{C}=\frac{I_{m}}{\omega C} \cong 0.65 \mathrm{kV}\)

\[
V_{L / R}=I_{m} \sqrt{R^{2}+\omega^{2} L^{2}}=0.5 \mathrm{kV}
\]

4.125 (a) \(V_{C}=\frac{1}{\omega C} \frac{V_{m}}{\sqrt{R^{2}+\left(\omega L-\frac{1}{\omega C}\right)^{2}}}\)

\[
\begin{aligned}
& =\frac{V_{m}}{\sqrt{(\omega R C)^{2}+\left(\omega^{2} L C-1\right)^{2}}}=\frac{V_{m}}{\sqrt{\left(\frac{\omega^{2}}{\omega_{0}^{2}}-1\right)^{2}+4 \beta^{2} \omega^{2} / \omega_{0}^{4}}} \\
& =\frac{V_{m}}{\sqrt{\left(\frac{\omega^{2}}{\omega_{0}^{2}}-1+\frac{2 \beta^{2}}{\omega_{0}^{2}}\right)^{2}+\frac{4 \beta^{2}}{\omega_{0}^{2}}-\frac{4 \beta^{4}}{\omega_{0}^{4}}}}
\end{aligned}
\]

This is maximum when

\[
\omega^{2}=\omega_{0}^{2}-2 \beta^{2}=\frac{1}{L C}-\frac{R^{2}}{2 L^{2}}
\]

(b) \(V_{L}=I_{m} \omega L=V_{m} \frac{\omega L}{\sqrt{R^{2}+\left(\omega L-\frac{1}{\omega C}\right)^{2}}}\)

\[
\begin{aligned}
& =\frac{V_{m} L}{\sqrt{\frac{R^{2}}{\omega^{2}}+\left(L-\frac{1}{\omega^{2} C}\right)^{2}}}=\frac{V_{m} L}{\sqrt{L^{2}-\frac{1}{\omega^{2}}\left(\frac{2 L}{C}-R^{2}\right)+\frac{1}{\omega^{4} C^{2}}}} \\
& =\frac{V_{m} L}{\sqrt{\left(\frac{1}{\omega^{2} C}-\left(L-\frac{C R^{2}}{2}\right)\right)^{2}+L^{2}-\left(L-\frac{1}{2} C R^{2}\right)^{2}}}
\end{aligned}
\]

This is maximum when\\
or

\[
\begin{aligned}
\omega^{2} & =\frac{\frac{1}{\omega^{2} C}=L-\frac{1}{2} C R^{2}}{L C-\frac{1}{2} C^{2} R^{2}}=\frac{1}{\frac{1}{\omega_{0}^{2}}-\frac{2 \beta^{2}}{\omega_{0}^{4}}} \\
& =\frac{\omega_{0}^{4}}{\omega_{0}^{2}-2 \beta^{2}} \quad \text { or } \omega=\frac{\omega_{0}^{2}}{\sqrt{\omega_{0}^{2}-2 \beta^{2}}} .
\end{aligned}
\]

4.126.

\[
\begin{aligned}
V_{L} & =I_{m} \sqrt{R^{2}+\omega^{2} L^{2}} \\
& =\frac{V_{m} \sqrt{R^{2}+\omega^{2} L^{2}}}{\sqrt{R^{2}+\left(\omega L-\frac{1}{\omega C}\right)^{2}}}
\end{aligned}
\]

for a given \(\omega, L, R\), this is maximum when

\[
\frac{1}{\omega C}=\omega L \quad \text { or } \quad C=\frac{1}{\omega^{2} L}=28.2 \mu \mathrm{~F}
\]

For that \(C, \quad V_{L}=\frac{V \sqrt{R^{2}+\omega^{2} L^{2}}}{R}=V \sqrt{1+(\omega L / R)^{2}}=0.540 \mathrm{kV}\)\\
At this \(C\)

\[
V_{C}=\frac{1}{\omega C} \frac{V_{m}}{R}=\frac{V_{m} \omega L}{R}=.509 \mathrm{kV}
\]

4.127

\[
\left|\begin{array}{llll}
0 & 0 & 0 & 0 \\
0 & 0 & 0 & 0 \\
0 & \text { poor cond } \\
0 & 0 & 0 & 0 \\
0 & 0 & 0 & 0
\end{array}\right|=\vec{I}\left|\overrightarrow{I^{\prime}}\right|
\]

We use the complex voltage \(V=V_{m} e^{i \omega t}\). Then the voltage across the capacitor is

\[
\left(I-I^{\prime}\right) \frac{1}{i \omega C}
\]

and that across the resistance \(R I^{\prime}\) and both equal \(V\). Thus

\[
I^{\prime}=\frac{V_{m}}{R} e^{i \omega t}, I-I^{\prime}=i \omega C V_{m} e^{i \omega t}
\]

Hence

\[
I=\frac{V_{m}}{R}(1+i \omega R C) e^{i \omega t}
\]

The actual voltage is obtained by taking the real part. Then

Where

\[
I=\frac{V_{m}}{R} \sqrt{1+(\omega R C)^{2}} \cos (\omega t+\varphi)
\]

Note → A condenser with poorly conducting material (dielectric of high resistance) be the plates is equalent to an an ideal condenser with a high resistance joined in \(p\) between its plates.\\
4.128

\[
\begin{gathered}
L_{1} \frac{d I_{1}}{d t}+\frac{\int I_{1} d t}{C}=-L_{12} \frac{d I_{2}}{d t} \\
L_{2} \frac{d I_{2}}{d t}=-L_{12} \frac{d I_{1}}{d t}
\end{gathered}
\]

from the second equation

\[
L_{2} I_{2}=-L_{12} I_{1}
\]

\begin{center}
\includegraphics[max width=\textwidth, alt={}]{da0156b6-4133-434e-90d4-ec01b2070afe-079_312_549_65_848}
\end{center}

Then

\[
\left(L_{1}-\frac{L_{12}^{2}}{L_{2}}\right) \ddot{I}{ }_{1}+\frac{I_{1}}{C}=0
\]

Thus the current oscillates with frequency

\[
\omega=\frac{1}{\sqrt{C\left(L_{1}-\frac{L_{12}^{2}}{L_{2}}\right)}}
\]

4.129 Given \(V=V_{m} \cos \omega t\)

\[
I=I_{m} \cos (\omega t-\varphi)
\]

where\\
\includegraphics[max width=\textwidth, alt={}, center]{da0156b6-4133-434e-90d4-ec01b2070afe-079_229_627_700_780}

Then,

\[
\begin{aligned}
V_{C} & =\frac{\int I d t}{C}=\frac{I_{m} \sin (\omega t-\varphi)}{\omega C} \\
& =\frac{V_{m}}{\sqrt{\left(1-\omega^{2} L C\right)^{2}+(\omega R C)^{2}}} \sin (\omega t-\varphi)
\end{aligned}
\]

As resonance the voltage amplitude across the capacitor

\[
=\frac{V_{m}}{R C \frac{1}{\sqrt{L C}}}=\sqrt{\frac{L}{C R^{2}}} V_{m}=n V_{m}
\]

So

\[
\frac{L}{C R^{2}}=n^{2}
\]

Now

\[
Q=\sqrt{\frac{L}{C R^{2}}-\frac{1}{4}}=\sqrt{n^{2}-\frac{1}{4}}
\]

4.130 For maximum current amplitude

\[
I_{m}=\frac{V_{m}}{\sqrt{R^{2}+\left(\omega L-\frac{1}{\omega C}\right)^{2}}}
\]

\[
L=\frac{1}{\omega^{2} C} \text { and then } I_{m 0}=\frac{V_{m}}{R}
\]

Now

\[
\frac{I_{m 0}}{\eta}=\frac{V_{m}}{\sqrt{R^{2}+\frac{(n-1)^{2}}{\omega^{2} C^{2}}}}
\]

So

\[
\begin{aligned}
\eta & =\sqrt{1+\frac{(n-1)^{2}}{(\omega R C)^{2}}} \\
\omega R C & =\frac{n-1}{\sqrt{\eta^{2}-1}}
\end{aligned}
\]

Now \(\quad Q=\sqrt{\left(\frac{L}{C R^{2}}\right)^{2}-\frac{1}{4}}=\sqrt{\left(\frac{1}{\omega R C}\right)^{2}-\frac{1}{4}}=\sqrt{\frac{\eta^{2}-1}{(n-1)^{2}}-\frac{1}{4}}\)\\
4.131 At resonance

\[
\omega_{0} L=\left(\omega_{0} C\right)^{-1} \quad \text { or } \quad \omega_{0}=\frac{1}{\sqrt{L C}},
\]

and

\[
\left(I_{m}\right)_{\mathrm{res}}=\frac{V_{m}}{R}
\]

Now

\[
\frac{V_{m}}{n R}=\frac{V_{m}}{\sqrt{R^{2}+\left(\omega_{1} L-\frac{1}{\omega_{1} C}\right)^{2}}}=\frac{V_{m}}{\sqrt{R^{2}+\left(\omega_{2} L-\frac{1}{\omega_{2} C}\right)^{2}}}
\]

Then

\[
\omega_{1} L-\frac{1}{\omega_{1} C}=\sqrt{n^{2}-1} R
\]

\[
\left.\omega_{2} L-\frac{1}{\omega_{2} C}=+\sqrt{n^{2}-1} R \quad \text { (assuming } \omega_{2}>\omega_{1}\right)
\]

or

\[
\omega_{1}-\frac{\omega_{0}^{2}}{\omega_{1}}=-\omega_{2}+\frac{\omega_{0}^{2}}{\omega_{2}}=-\sqrt{n^{2}-1} \frac{R}{L}
\]

or

\[
\omega_{1}+\omega_{2}=\frac{\omega_{0}^{2}}{\omega_{1} \omega_{2}}\left(\omega_{1}+\omega_{2}\right) \Rightarrow \omega_{0}=\sqrt{\omega_{1} \omega_{2}}
\]

and

\[
\begin{aligned}
& \omega_{2}-\omega_{1}=\sqrt{n^{2}-1} \frac{R}{L} \\
& \beta=\frac{R}{2 L}=\frac{\omega_{2}-\omega_{1}}{2 \sqrt{n^{2}-1}}
\end{aligned}
\]

and

\[
Q=\sqrt{\frac{\omega_{0}^{2}}{4 \beta^{2}}-\frac{1}{4}}=\sqrt{\frac{\left(n^{2}-1\right) \omega_{1} \omega_{2}}{\left(\omega_{2}-\omega_{1}\right)^{2}}-\frac{1}{4}}
\]

\(4.132 Q=\frac{\omega}{2 \beta} \approx \frac{\omega_{0}}{2 \beta}\) for low damping.\\
Now \(\frac{I_{m}}{\sqrt{2}}=\frac{R I_{m}}{\sqrt{R^{2}+\left(\omega L-\frac{1}{\omega C}\right)^{2}}}, I_{m}=\) current amplitude at resouance\\
or

\[
\omega-\frac{\omega_{0}^{2}}{\omega}= \pm \frac{R}{L}= \pm 2 \beta
\]

Thus

\[
\omega \cong \omega_{0} \pm \beta
\]

So

\[
\Delta \omega=2 \beta \text { and } Q=\frac{\omega_{0}}{\Delta \omega} .
\]

4.133 At resonance \(\omega=\omega_{0}\)

\[
I_{m}\left(\omega_{0}\right)=\frac{V_{m}}{R}
\]

Then \(I_{m}\left(\eta \omega_{0}\right)=\frac{V_{m}}{\sqrt{R^{2}+\left(\eta \omega_{0} L-\frac{1}{\eta \omega_{0} C}\right)^{2}}}\)

\[
=\frac{V_{m}}{\sqrt{R^{2}+\left(\eta-\frac{l}{\eta}\right)^{2} \frac{L}{C}}}=\frac{V_{m}}{\sqrt{1+\left(Q^{2}+\frac{l}{4}\right)\left(\eta-\frac{1}{\eta}\right)^{2} \frac{L}{C}}}
\]

4.134 The a.c. current must be

\[
I=I_{0} \sqrt{2} \sin \omega t
\]

Then D.C. component of the rectified current is

\[
\begin{aligned}
\left\langle I^{\prime}\right\rangle & =\frac{1}{T} \int_{0}^{T / 2} I_{0} \sqrt{2} \sin \omega t d t \\
& =I_{0} \sqrt{2} \frac{1}{2 \pi} \int_{0}^{\pi} \sin \theta d \theta \\
& =\frac{I_{0} \sqrt{2}}{\pi}
\end{aligned}
\]

Since the charge deposited must be the same

\[
I_{0} t_{0}=\frac{I_{0} \sqrt{2}}{\pi} t \quad \text { or } t=\frac{\pi t_{0}}{\sqrt{2}}
\]

The answer is incorrect.\\
4.135 (a) \(I(t)=I_{1} \frac{t}{T} 0 \leq t<T\)

\[
I(t \pm T)=I(t)
\]

Now mean current

Then

\[
<I>=\frac{1}{T} \int_{0}^{T} I_{1} \frac{t}{T} d t=I_{1} \frac{T^{2} / 2}{T^{2}}=I_{1} / 2
\]

\[
I_{1}=2 I_{0} \text { since }\langle I\rangle=I_{0} .
\]

Now mean square current \(\left\langle I^{2}\right\rangle\).

\[
=4 I_{0}^{2} \frac{1}{T} \int_{0}^{T} \frac{t^{2}}{T^{2}} d t=\frac{4 I_{0}^{2}}{3}
\]

so effective current \(=\frac{2 I_{0}}{\sqrt{3}}\).\\
(b) In this case \(I=I_{1}|\sin \omega t|\)\\
and

\[
\begin{aligned}
I_{0} & =\frac{1}{T} \int_{0}^{T} I_{1}|\sin \omega t| d t \\
& =\frac{1}{2 \pi} I_{1} \int_{0}^{2 \pi}|\sin \theta| d \theta=\frac{I_{1}}{\pi} \int_{0}^{\pi} \sin \theta d \theta=\frac{2 I_{1}}{\pi}
\end{aligned}
\]

So

\[
I_{1}=\frac{\pi I_{0}}{2}
\]

Then, mean square current \(=\left\langle I^{2}\right\rangle=\frac{\pi^{2} I_{0}^{2}}{4 T} \int_{0}^{T} \sin ^{2} \omega t d t\)

\[
=\frac{\pi^{2} I_{0}^{2}}{4} \times \frac{1}{2 \pi} \int_{0}^{2 \pi} \sin ^{2} \theta d \theta=\frac{\pi^{2} I_{0}^{2}}{8}
\]

so effective current \(=\frac{\pi I_{0}}{\sqrt{8}}\).\\
\(4.136 P_{d . c .}=\frac{V_{0}^{2}}{R}\)

\[
P_{a . c .}=\frac{V_{0}^{2}}{\sqrt{R^{2}+\omega^{2} L^{2}}} \cdot \frac{R}{\sqrt{R^{2}+\omega^{2} L^{2}}}=\frac{v_{0}^{2} / R}{1+\left(\frac{\omega L}{R}\right)^{2}}=\frac{P_{d . c .}}{\eta}
\]

Thus

\[
\frac{\omega L}{R}=\sqrt{\eta-1}
\]

or

\[
\omega=\frac{R}{L} \sqrt{\eta-1}
\]

\[
v=\frac{R}{2 \pi L} \sqrt{\eta-1}=2 \mathrm{kH} \text { of on putting the values. }
\]

\(4.137 Z=\sqrt{R^{2}+X_{L}^{2}}\) or \(R_{0}=\sqrt{Z^{2}-X_{L}^{2}}\)\\
The

So

\[
\begin{gathered}
\tan \theta=\frac{X_{L}}{\sqrt{z^{2}-X_{L}^{2}}} \\
\cos \varphi=\frac{\sqrt{z^{2}-X_{L}^{2}}}{z}=\sqrt{1-\left(\frac{X_{L}}{z}\right)^{2}} \\
\varphi=\cos ^{-1} \sqrt{1-\left(\frac{X_{L}}{z}\right)^{2}}=37^{\circ}
\end{gathered}
\]

The current lags by \(\varphi\) behind the voltage.\\
also

\[
P=V I \cos \varphi=\frac{V^{2}}{Z^{2}} \sqrt{Z^{2}-X_{L}^{2}}=.160 \mathrm{~kW} .
\]

\(4.138 P=\frac{V^{2}(R+r)}{(R+r)^{2}+\omega^{2} L^{2}}\)\\
This is maximum when \(R+r=\omega L\) for

\[
P=\frac{V^{2}}{R+r+\frac{(\omega L)^{2}}{R+r}}=\frac{V^{2}}{\left[\sqrt{R+r}-\frac{\omega L}{\sqrt{R+r}}\right]^{2}+2 \omega L}
\]

Thus \(R=\omega L-r\) for maximum power and \(P_{\text {max }}=\frac{V^{2}}{2 \omega L}\).\\
Substituting the values, we get \(R=200 \Omega\) and \(P_{\max }=.114 \mathrm{~kW}\).\\
\(4.139 P=\frac{V^{2} R}{R^{2}+\left(X_{L}-X_{C}\right)^{2}}\)\\
Varying the capacitor does not change \(R\) so if \(P\) increases \(n\) times

\[
z=\sqrt{R^{2}+\left(X_{L}-X_{C}\right)^{2}} \text { must decreases } \sqrt{n} \text { times }
\]

Thus \(\cos \varphi=\frac{R}{Z}\) increases \(\sqrt{n}\) times\\
\(\therefore \quad \%\) increase in \(\cos \varphi=(\sqrt{n}-1) \times 100 \%=30.4 \%\).\\
4.140 \(P=\frac{V^{2} R}{R^{2}+\left(X_{L}-X_{C}\right)^{2}}\)

At resonance \(X_{L}=X_{C} \Rightarrow \omega_{0}=\frac{1}{\sqrt{L C}}\).\\
Power generated will decrease \(n\) times when\\
or

\[
\begin{aligned}
& \left(X_{L}-X_{C}\right)^{2}=\left(\omega L-\frac{1}{\omega C}\right)^{2}=(n-1) R^{2} \\
& \omega-\frac{\omega_{0}^{2}}{\omega}= \pm \sqrt{n-1} \frac{R}{L}= \pm \sqrt{n-1} 2 \beta
\end{aligned}
\]

Thus\\
or

\[
\begin{gathered}
\omega^{2} \mp 2 \sqrt{n-1} \beta \omega-\omega_{0}^{2}=0 \\
(\omega \mp \sqrt{n-1} \beta)^{2}=\omega_{0}^{2}+(n-1) \beta^{2} \\
\frac{\omega}{\omega_{0}}=\sqrt{1+(n-1) \beta^{2} / \omega_{0}^{2}} \pm \sqrt{n-1} \beta / \omega_{0}
\end{gathered}
\]

(taking only the positive sign in the first term to ensure positive value for \(\frac{\omega}{\omega_{0}}\).)

Now

Thus

\[
\begin{gathered}
Q=\frac{\omega}{2 \beta}=\frac{1}{2} \sqrt{\left(\frac{\omega_{0}}{\beta}\right)^{2}-1} \\
\frac{\omega_{0}}{\beta}=\sqrt{1+4 Q^{2}} \\
\frac{\omega}{\omega_{0}}=\sqrt{1+\frac{n-1}{\left(1+4 Q^{2}\right)}} \pm \sqrt{n-1} / \sqrt{1+4 Q^{2}}
\end{gathered}
\]

For large Q

\[
\left|\frac{\omega-\omega_{0}}{\omega_{0}}\right|=\frac{\sqrt{e n-1}}{2 Q}=\frac{\sqrt{e n-1}}{2 Q} \times 100 \%=0.5 \%
\]

4.141 We have\\
so

\[
V_{1}=\frac{V R}{\sqrt{\left(R+R_{1}\right)^{2}+X_{L}^{2}}}, \quad V_{2}=\frac{V \sqrt{R_{1}^{2}+X_{L}^{2}}}{\sqrt{\left(R+R_{1}\right)^{2}+X_{L}^{2}}}
\]

Hence

\[
\left(R+R_{1}\right)^{2}+X_{L}^{2}=\left(\frac{V R}{V_{1}}\right)_{-}^{2}, R_{1}^{2}+X_{L}^{2}=\left(\frac{V_{2} R}{V_{1}}\right)^{2}
\]

or

\[
\begin{aligned}
R^{2}+2 R R_{1} & =\frac{R^{2}}{V_{1}^{2}}\left(V^{2}-V_{2}^{2}\right) \\
R_{1} & =\frac{R}{2 V_{1}^{2}}\left(V^{2}-V_{2}^{2}-V_{1}^{2}\right)
\end{aligned}
\]

Heat generated in the coil \(=\frac{V^{2} R_{1}}{\left(R_{1}+R_{2}\right)^{2}+X_{L}^{2}}=\frac{V_{1}^{2}}{R^{2}} \times R_{1}=\frac{V_{1}^{2}}{R^{2}} \times \frac{R^{2}}{2 V_{1}^{2}}\left(V^{2}-V_{1}^{2}-V_{2}^{2}\right)\)

\[
=\frac{V^{2}-V_{1}^{2}-V_{2}^{2}}{2 R}=30 \mathrm{~W}
\]

4.142 Here \(I_{2}=\frac{V}{R}, V=\) effective voltage

\[
I_{1}=\frac{V}{\sqrt{R^{2}+X_{L}^{2}}}
\]

and \(I=\frac{V \sqrt{\left(R+R_{1}\right)^{2}+X_{L}^{2}}}{R \sqrt{R_{1}^{2}+X_{L}^{2}}}=\frac{V}{R_{\text {eff }}}\)\\
\(R_{\text {eff }}\) is the impedance of the coil \& the resistance in parallel.\\
\includegraphics[max width=\textwidth, alt={}, center]{da0156b6-4133-434e-90d4-ec01b2070afe-085_344_451_362_958}

Now

\[
\begin{gathered}
\frac{I^{2}-I_{2}^{2}}{I_{2}^{2}}=\frac{R^{2}+2 R R_{1}}{R_{1}^{2}+X_{L}^{2}}=\left(\frac{I_{1}}{I_{2}}\right)^{2}+\frac{2 R R_{1}}{R^{2}+X_{L}^{2}} \\
\frac{I^{2}-I_{2}^{2}-I_{1}^{2}}{I_{2}^{2}}=\frac{2 R R_{1}}{R^{2}+X_{L}^{2}}
\end{gathered}
\]

Now mean power consumed in the coil

\[
=I_{1}^{2} R_{1}=\frac{V^{2} R_{1}}{R_{2}+X_{L}^{2}}=I_{2}^{2} R \cdot \frac{I^{2}-I_{1}^{2}-I_{2}^{2}}{2 I_{2}^{2}}=\frac{1}{2} R\left(I^{2}-I_{1}^{2}-I_{2}^{2}\right)=2.5 \mathrm{~W} .
\]

4.143\\
\(\frac{1}{Z}=\frac{1}{R}+\frac{1}{\frac{1}{i \omega C}}=\frac{1}{R}+i \omega C=\frac{1+i \omega R C}{R}\)

\[
|z|=\frac{R}{\sqrt{1+(\omega R C)^{2}}}=40 \Omega
\]

4.144 (a) For the resistance, the voltage and the current are in phase. For the coil the voltage is ahead of the current by less than \(90^{\circ}\). The current is obtained by addition because the elements are in parallel.

\begin{figure}[h]
\begin{center}
  \includegraphics[alt={},max width=\textwidth]{da0156b6-4133-434e-90d4-ec01b2070afe-085_304_521_1565_161}
\captionsetup{labelformat=empty}
\caption{(a)}
\end{center}
\end{figure}

\begin{figure}[h]
\begin{center}
  \includegraphics[alt={},max width=\textwidth]{da0156b6-4133-434e-90d4-ec01b2070afe-085_289_449_1526_884}
\captionsetup{labelformat=empty}
\caption{(b)}
\end{center}
\end{figure}

(b) \(I_{C}\) is ahead of the voltage by \(90^{\circ}\).\\
(c) The coil has no resistance so \(I_{L}\) is \(90^{\circ}\) behind the voltage.\\
\includegraphics[max width=\textwidth, alt={}, center]{da0156b6-4133-434e-90d4-ec01b2070afe-086_413_460_175_488}\\
4.145 When the coil and the condenser are in parallel, the equation is

\[
\begin{aligned}
& L \frac{d I_{1}}{d t}+R I_{1}=\frac{\int I_{2} d t}{C}=V_{m} \cos \omega t \\
& I=I_{1}+I_{2}
\end{aligned}
\]

Using complex voltages

\[
I_{1}=\frac{V_{m} e^{i \omega t}}{R+i \omega L}, I_{2}=i \omega C V_{m} e^{i \omega t}
\]

and\\
\includegraphics[max width=\textwidth, alt={}, center]{da0156b6-4133-434e-90d4-ec01b2070afe-086_339_497_648_720}

\[
I=\left(\frac{1}{R+i \omega L}+i \omega C\right) V_{m} e^{i \omega t}=\left[\frac{R-i \omega L+i \omega C\left(R^{2}+\omega^{2} L^{2}\right)}{R^{2}+\omega^{2} L^{2}}\right] V_{m} e^{i \omega t}
\]

Thus, taking real parts \(\quad I=\frac{V_{m}}{|Z(\omega)|} \cos (\omega t-\varphi)\)\\
where

\[
\begin{gathered}
\frac{1}{|Z(\omega)|}=\frac{\left[R^{2}+\left\{\omega C\left(R^{2}+\omega^{2} L^{2}\right)-\omega L\right\}^{2}\right]}{\left(R^{2}+\omega^{2} L^{2}\right)^{1 / 2}} \\
\tan \varphi=\frac{\omega L-\omega C\left(R^{2}+\omega^{2} L^{2}\right)}{R}
\end{gathered}
\]

and\\
(a) To get the frequency of resonance we must define what we mean by resouance. One definition requires the extremum (maximum or minimum) of current amplitude. The other definition requires rapid change of phase with \(\varphi\) passing through zero at resonance. For the series circuit.

\[
I_{m}=\frac{V_{m}}{\left\{R^{2}+\left(\omega L-\frac{1}{\omega C}\right)^{2}\right\}^{1 / 2}} \quad \text { and } \quad \tan \varphi=\frac{\omega L-\frac{1}{\omega C}}{R}
\]

both definitions give \(\omega^{2}=\frac{1}{L C}\) at resonance. In the present case the two definitions do not agree (except when \(R=0\) ). The definition that has been adopted in the answer given in the book is the vanishing of phase. This requires

\[
C\left(R^{2}+\omega^{2} L^{2}\right)=L
\]

or

\[
\omega^{2}=\frac{1}{L C}-\frac{R^{2}}{L^{2}}=\omega_{\mathrm{res}}^{2}, \quad \omega_{\mathrm{res}}=31.6 \times 10^{3} \mathrm{rad} / \mathrm{s}
\]

Note that for small \(R, \varphi\) rapidly changes from \(\approx-\frac{\pi}{2}\) to \(+\frac{\pi}{2}\) as \(\omega\) passes through \(\omega_{\text {res }}\) from \(<\omega_{\text {res }}\) to \(>\omega_{\text {res }}\).\\
(b) At resonance

\[
I_{m}=\frac{V_{m} R}{L / C}=V_{m} \frac{C R}{L}
\]

so \(I=\) effective value of total current \(=V \frac{C R}{L}=3.1 \mathrm{~mA}\).\\
similarly

\[
\begin{gathered}
I_{L}=\frac{V}{\sqrt{L / C}}=V \sqrt{\frac{C}{L}}=0.98 \mathrm{~A} . \\
I_{C}=\omega C V=V \sqrt{\frac{C}{L}-\frac{R^{2} C^{2}}{L^{2}}} \equiv 0.98 \mathrm{~A} .
\end{gathered}
\]

Note :- The vanishing of phase (its passing through zero) is considered a more basic definition of resonance.\\
4.146 We use the method of complex voltage

\[
V=V_{0} e^{i \omega t}
\]

Then \(I_{C}=\frac{V_{0} e^{i \omega t}}{\frac{1}{i \omega C}}=i \omega C V_{0} e^{i \omega t}\)\\
\includegraphics[max width=\textwidth, alt={}, center]{da0156b6-4133-434e-90d4-ec01b2070afe-087_330_427_1107_966}

\[
\begin{aligned}
& I_{L, R}=\frac{V_{0} e^{i \omega t}}{R+i \omega L} \\
& \quad I=I_{C}+I_{L, R}=V_{0} \frac{R-i \omega L+i \omega C\left(R^{2}+\omega^{2} L^{2}\right)}{R^{2}+\omega^{2} L^{2}} e^{i \omega t}
\end{aligned}
\]

Then taking the real part

\[
I=\frac{V_{0} \sqrt{R^{2}+\left\{\omega C\left(R^{2}+\omega^{2} L^{2}\right)-\omega L\right\}^{2}}}{R^{2}+\omega^{2} L^{2}} \cos (\omega t-\varphi)
\]

where

\[
\tan \varphi=\frac{\omega L-\omega C\left(R^{2}+\omega^{2} L^{2}\right)}{R}
\]

4.147 From the previous problem

\[
\begin{aligned}
\boldsymbol{z} & =\frac{R^{2}+\omega^{2} L^{2}}{\sqrt{R^{2}+\left\{\omega C\left(R^{2}+\omega^{2} L^{2}\right)-\omega L\right\}^{2}}} \\
& =\frac{R^{2}+\omega^{2} L^{2}}{\sqrt{\left(R^{2}+\omega^{2} L^{2}\right)\left(1-2 \omega^{2} L C\right)+\omega^{2} C^{2}\left(R^{2}+\omega^{2} L^{2}\right)^{2}}} \\
& =\frac{\sqrt{R^{2}+\omega^{2} L^{2}}}{\sqrt{\left(1-2 \omega^{2} L C\right)+\omega^{2} C^{2}\left(R^{2}+\omega^{2} L^{2}\right)}}=\frac{\sqrt{R^{2}+\omega^{2} L^{2}}}{\sqrt{\left(1-\omega^{2} L C\right)^{2}+(\omega R C)^{2}}}
\end{aligned}
\]

4.148 (a) We have

\[
\varepsilon=-\frac{d \Phi}{d t}=\omega \Phi_{0} \sin \omega t=L \dot{I}+R I
\]

Put

\[
I=I_{m} \sin (\omega t-\varphi) . \text { Then }
\]

\[
\begin{aligned}
\omega \Phi_{0} \sin \omega t & =\omega \Phi_{0}\{\sin (\omega t-\varphi) \cos \varphi+\cos (\omega t-\varphi) \sin \varphi\} \\
& =L I_{m} \omega \cos (\omega t-\varphi)+R I_{m} \sin (\omega t-\varphi)
\end{aligned}
\]

so

\[
R I_{m}=\omega \Phi_{0} \cos \varphi \quad \text { and } \quad L I_{m}=\Phi_{0} \sin \varphi
\]

or

\[
I_{m}=\frac{\omega \Phi_{0}}{\sqrt{R^{2}+\omega^{2} L^{2}}} \quad \text { and } \quad \tan \varphi=\frac{\omega L}{R}
\]

(b) Mean mechanical power required to maintain rotation \(=\) energy loss per unit time

\[
=\frac{1}{T} \int_{0}^{T} R I^{2} d t=\frac{1}{2} R I_{m}^{2}=\frac{1}{2} \frac{\omega^{2} \Phi_{0}^{2} R}{R^{2}+\omega^{2} L^{2}}
\]

4.149 We consider the force \(\vec{F}_{12}\) that a circuit 1 exerts on another closed circuit 2 :-

\[
\overrightarrow{F_{12}}=\oint \quad I_{\tau} d \overrightarrow{l_{2}} \times \overrightarrow{B_{12}}
\]

Here \(\vec{B}_{12}=\) magnetic field at the site of the current element \(d \overrightarrow{l_{2}}\) due to the current \(I_{1}\) flowing in 1.

\[
=\frac{\mu_{0}}{4 \pi} \int \frac{I_{1} d \overrightarrow{l_{1}} \times \overrightarrow{r_{12}}}{r_{12}^{3}}
\]

where \(\overrightarrow{r_{12}}=\overrightarrow{r_{2}}-\overrightarrow{r_{1}}=\) vector from current element \(d \overrightarrow{l_{1}}\) to the current element \(d \overrightarrow{l_{2}}\)\\
Now

\[
\overrightarrow{F_{12}}=\frac{\mu_{0}}{4 \pi} \iint I_{1} I_{2} \frac{d \overrightarrow{l_{2}} \times\left(d \overrightarrow{l_{1}} \times \overrightarrow{r_{12}}\right)}{r_{12}^{3}}=\frac{\mu_{0}}{4 \pi} \iint I_{1} I_{2} \frac{d \overrightarrow{l_{1}}\left(d \overrightarrow{l_{2}} \cdot \overrightarrow{r_{12}}\right)-\left(d \overrightarrow{l_{1}} \cdot d \overrightarrow{l_{2}}\right) \overrightarrow{r_{12}}}{r_{12}^{3}}
\]

In the first term, we carry out the integration over \(d \overrightarrow{l_{\tau}}\) first. Then

\[
\iint \frac{d \overrightarrow{l_{1}}\left(d \overrightarrow{l_{2}} \cdot \overrightarrow{r_{12}}\right)}{r_{12}^{3}}=\int d \overrightarrow{l_{1}} \oint \frac{d \overrightarrow{l_{2}} \cdot \overrightarrow{r_{12}}}{r_{12}^{3}}=-\int d \overrightarrow{l_{1}} \oint d \overrightarrow{l_{2}} \cdot \nabla_{2} \frac{1}{r_{12}}=0
\]

because

\[
\oint d l_{2} \cdot \nabla_{2} \frac{1}{r_{12}}=\int d \vec{S}_{2} \text { curl }\left(\nabla \frac{1}{r_{12}}\right)=0
\]

Thus

\[
F_{12}=-\frac{\mu_{0}}{4 \pi} \iint I_{1} I_{2} d \overrightarrow{l_{1}} \cdot d \overrightarrow{l_{2}} \frac{\overrightarrow{r_{12}}}{r_{12}^{3}}
\]

The integral involved will depend on the vector \(\vec{a}\) that defines the separation of the (suitably chosen )centre of the coils. Let \(C_{1}\) and \(C_{2}\) be the centres of the two coil suitably defined.\\
Write

\[
\overrightarrow{r_{12}}=\overrightarrow{r_{2}}-\overrightarrow{r_{1}}=\overrightarrow{\rho_{2}}-\overrightarrow{\rho_{1}}+\vec{a}
\]

where \(\overrightarrow{\rho_{1}}\left(\overrightarrow{\rho_{2}}\right)\) is the distance of \(d \overrightarrow{l_{1}}\left(d \overrightarrow{l_{2}}\right)\) from \(C_{1}\left(C_{2}\right)\) and \(\vec{a}\) stands for the vector \(\overrightarrow{C_{1}} \vec{C}_{2}\).\\
Then\\
and

\[
\begin{gathered}
\frac{\overrightarrow{r_{12}}}{r_{12}^{3}}=-\vec{\nabla}_{\vec{a}} \frac{1}{r_{12}} \\
\vec{F}_{12}=\vec{\nabla}_{a}\left[I_{1} I_{2} \frac{\mu_{0}}{4 \pi} \iint \frac{d \overrightarrow{l_{1}} \cdot d \overrightarrow{l_{2}}}{r_{12}}\right]
\end{gathered}
\]

The bracket defines the mutual inductance \(L_{12}\). Thus noting the definition of \(x\)

\[
\left\langle F_{x}\right\rangle=\frac{\partial L_{12}}{\partial x}\left\langle I_{1} I_{2}\right\rangle
\]

where \(<>\) denotes time average. Now

\[
I_{1}=I_{0} \cos \omega t=\text { Real part of } I_{0} e^{i \omega t}
\]

The current in the coil 2 satisfies \(R I_{2}+L_{2} \frac{d I_{2}}{d t}=-L_{12} \frac{d I_{1}}{d t}\)\\
or

\[
I_{2}=\frac{-i \omega L_{12}}{R+i \omega L_{2}} I_{0} e^{i \omega t}(\text { in the complex case })
\]

taking the real part\\
\(I_{2}=-\frac{\omega L_{12} I_{0}}{R^{2}+\omega^{2} L_{2}^{2}}\left(\omega L_{2} \cos \omega t-R \sin \omega t\right)=-\frac{\omega L_{12}}{\sqrt{R^{2}+\omega^{2} L_{\tau}^{2}}} I_{0} \cos (\omega t+\varphi)\)\\
Where \(\tan \varphi=\frac{R}{\omega L_{2}}\). Taking time average, we get

\[
<F_{x}>=\frac{\partial L_{12}}{\partial x} I_{0} \frac{\omega L_{12} I_{0}}{\sqrt{R^{2}+\omega^{2} L_{2}^{2}}} \cdot \frac{1}{2} \cos \varphi=\frac{\omega^{2} L_{2} L_{12} I_{0}^{2}}{2\left(R^{2}+\omega^{2} L_{2}^{2}\right)} \frac{\partial L_{12}}{\partial x}
\]

The repulsive nature of the force is also consistent with Lenz's law, assuming, of comse, that \(L_{12}\) decreases with \(x\).

\subsection*{4.3 ELASTIC WAVES. ACOUSTICS}
4.150 Since the temperature varies linearly we can write the temperature as a function of \(x\), which is, the distance from the point \(A\) towards \(B\).\\
i.e.,

\[
T=T_{1}+\frac{T_{2}-T_{1}}{l} x,[0<x<l]
\]

hence,


\begin{equation*}
d T=\left(\frac{T_{2}-T_{1}}{l}\right) d x \tag{1}
\end{equation*}


In order to travel an elemental distance of \(d x\) which is at a distance of \(x\) from \(A\) it will take a time


\begin{equation*}
d t=\frac{d x}{\alpha \sqrt{T}} \tag{2}
\end{equation*}


From Eqns (1) and (2), expressing \(d x\) in terms of \(d T\), we get

\[
d t=\frac{l}{\alpha \sqrt{T}}\left(\frac{l d T}{T_{2}-T_{1}}\right)
\]

Which on integration gives\\
or,

\[
\begin{aligned}
\int_{0}^{t} d t & =\frac{l}{\alpha\left(T_{2}-T_{1}\right)} \int_{T_{1}}^{T_{2}} \frac{d T}{\sqrt{T}} \\
t & =\frac{2 l}{\left(T_{2}-T_{1}\right)}\left(\sqrt{T_{2}}-\sqrt{T_{1}}\right)
\end{aligned}
\]

Hence the sought time \(t=\frac{2 l}{\alpha\left(\sqrt{T_{1}}+\sqrt{T_{2}}\right)}\)\\
4.151 Equation of plane wave is given by\\
\(\xi(r, t)=a \cos (\omega t-\vec{k} \cdot \vec{r})\), where \(\vec{k}=\frac{\omega}{v} \hat{n}\) called the wave vector and \(\hat{n}\) is the unit vector normal to the wave surface in the direction of the propagation of wave.\\
\includegraphics[max width=\textwidth, alt={}, center]{da0156b6-4133-434e-90d4-ec01b2070afe-090_498_563_1418_509}\\
or,

\[
\begin{aligned}
\xi(x, y, z) & =a \cos \left(\omega t-k_{x} x-k_{y} y-k_{z} z\right) \\
& =a \cos (\omega t-k x \cos \alpha-k y \cos \beta-k z \cos \gamma)
\end{aligned}
\]

Thus

\[
\xi\left(x_{1}, y_{1}, z_{1}, t\right)=a \cos \left(\omega t-k x_{1} \cos \alpha-k y_{1} \cos \beta-k z_{1} \cos \gamma\right)
\]

and

\[
\xi\left(x_{2}, y_{2}, z_{2}, t\right)=a \cos \left(\omega t-k x_{2} \cos \alpha-k y_{2} \cos \beta-k z_{2} \cos \gamma\right)
\]

Hence the sought wave phase difference

\[
\begin{aligned}
\varphi_{2}-\varphi_{1} & =k\left[\left(x_{1}-x_{2}\right) \cos \alpha+\left(y_{1}-y_{2}\right) \cos \beta+\left(z_{1}-z_{2}\right) \cos \gamma\right] \\
\text { or } \Delta \varphi & =\left|\varphi_{2}-\varphi_{1}\right|=k\left|\left[\left(x_{1}-x_{2}\right) \cos \alpha+\left(y_{1}-y_{2}\right) \cos \beta+\left(z_{1}-z_{2}\right) \cos \gamma\right]\right| \\
& =\frac{\omega}{v}\left|\left[\left(x_{1}-x_{2}\right) \cos \alpha+\left(y_{1}-y_{2}\right) \cos \beta+\left(z_{1}-z_{2}\right) \cos \gamma\right]\right|
\end{aligned}
\]

4.152 The phase of the oscillation can be written as

\[
\Phi=\omega t-\vec{k} \cdot \vec{r}
\]

When the wave moves along the \(x\)-axis

\[
\left.\Phi=\omega t-k_{x} x \quad \text { (On putting } k_{y}=k_{z}=0\right) .
\]

Since the velocity associated with this wave is \(\mathbf{v}_{1}\)

We have

\[
k_{x}=\frac{\omega}{v_{1}}
\]

Similarly

\[
\begin{aligned}
k_{y} & =\frac{\omega}{v_{2}} \text { and } k_{z}=\frac{\omega}{v_{3}} \\
\vec{k} & =\frac{\omega}{v_{1}} \hat{e}_{x}+\frac{\omega}{v_{2}} \hat{e}_{y}+\frac{\omega}{v_{3}} \hat{e}_{3} .
\end{aligned}
\]

4.153 The wave equation propagating in the direction of +ve \(x\) axis in medium \(K\) is give as

\[
\xi=a \cos (\omega t-k x)
\]

So, \(\quad \xi=a \cos k(v t-x), \quad\) where \(k=\frac{\omega}{v}\) and \(v\) is the wave velocity\\
In the refrence frame \(K^{\prime}\), the wave velocity will be \((v-V)\) propagating in the direction of \(+\mathrm{ve} x\) axis and \(x\) will be \(x^{\prime}\). Thus the sought wave equation.\\
\(\xi=a \cos k\left[(v-V) t-x^{\prime}\right]\)\\
or,

\[
\xi=a \cos \left[\left(\omega-\frac{\omega}{v} V\right) t-k x^{\prime}\right]=a \cos \left[\omega t\left(1-\frac{V}{v}\right)-k x^{\prime}\right]
\]

4.154 This follows on actually putting\\
in the wave equation

\[
\begin{aligned}
& \xi=f(t+\alpha x) \\
& \frac{\partial^{2} \xi}{\partial x^{2}}=\frac{1}{v^{2}} \frac{\partial^{2} \xi}{\partial t^{2}}
\end{aligned}
\]

(We have written the one dimensional form of the wave equation.) Then

\[
\frac{1}{v^{2}} f^{\prime \prime}(t+\alpha x)=\alpha^{2} f^{\prime \prime}(t+\alpha x)
\]

so the wave equation is satisfied if

\[
\alpha= \pm \frac{1}{v^{\prime}}
\]

That is the physical meaning of the constant \(\alpha\).\\
4.155 The given wave equation

\[
\xi=60 \cos (1800 t-5 \cdot 3 x)
\]

is of the type

\[
\begin{gathered}
\xi=a \cos (\omega t-k x), \quad \text { where } a=60 \times 10^{-6} \mathrm{~m} \\
\omega=1800 \text { per sec and } k=5.3 \text { per metre }
\end{gathered}
\]

As

\[
k=\frac{2 \pi}{\lambda}, \text { so } \lambda=\frac{2 \pi}{k}
\]

and also

\[
k=\frac{\omega}{v}, \quad \text { so } \quad v=\frac{\omega}{k}=340 \mathrm{~m} / \mathrm{s}
\]

(a) Sought ratio \(=\frac{a}{\lambda}=\frac{a k}{2 \pi}=5 \cdot 1 \times 10^{-5}\)\\
(b) Since \(\xi=a \cos (\omega t-k x)\)

\[
\frac{\partial \xi}{\partial t}=-a \omega \sin (\omega t-k x)
\]

So velocity oscillation amplitude


\begin{equation*}
\left(\frac{\partial \xi}{\partial t}\right)_{m} \text { or } v_{m}=a \omega=0.11 \mathrm{~m} / \mathrm{s} \tag{1}
\end{equation*}


and the sought ratio of velocity oscillation amplitude to the wave propagation velocity

\[
=\frac{v_{m}}{v}=\frac{0.11}{340}=3.2 \times 10^{-4}
\]

(c) Relative deformation \(=\frac{\partial \xi}{\partial x}=a k \sin (\omega t-k x)\)

So, relative deformation amplitude


\begin{equation*}
=\left(\frac{\partial \xi}{\partial x}\right)_{m}=a k=\left(60 \times 10^{-6} \times 5.3\right) \mathrm{m}=3.2 \times 10^{-4} \mathrm{~m} \tag{2}
\end{equation*}


From Eqns (1) and (2)

\[
\left(\frac{\partial \xi}{\partial x}\right)_{m}=a k=\frac{a \omega}{v}=\frac{1}{v}\left(\frac{\partial \xi}{\partial t}\right)_{m}
\]

Thus \(\left(\frac{\partial \xi}{\partial x}\right)_{m}=\frac{1}{v}\left(\frac{\partial \xi}{\partial t}\right)_{m}\), where \(v=340 \mathrm{~m} / \mathrm{s}\) is the wave velocity.\\
4.156 (a) The given equation is,

\[
\xi=a \cos (\omega t-k x)
\]

So at

Now,

\[
\begin{gathered}
t=o, \\
\xi=a \cos k x
\end{gathered}
\]

and

Also,\\
and at

\[
\frac{d \xi}{d t}=-a \omega \sin (\omega t-k x)
\]

\[
\frac{d \xi}{d t}=a \omega \sin k x, \text { at } t=0
\]

\begin{center}
\includegraphics[max width=\textwidth, alt={}]{image-not-found}
\end{center}

\[
\frac{d \xi}{d x}=+a k \sin (\omega t-k x)
\]

\[
\begin{gathered}
t=0 \\
\frac{d \xi}{d x}=-a k \sin k x
\end{gathered}
\]

Hence all the graphs are similar having different amplitudes, as shown in the answersheet of the problem book.\\
(b) At the points, where \(\xi=0\), the velocity direction is positive, i.e., along + ve \(x\)-axis in the case of longitudinal and \(+v e y\)-axis in the case of transverse waves, where \(\frac{d \xi}{d t}\) is positive and vice versa.

For sought plots see the answer-sheet of the problem book.\\
4.157 In the given wave equation the particle's displacement amplitude \(=a e^{-\gamma x}\)

Let two points \(x_{1}\) and \(x_{2}\), between which the displacement amplitude differ by \(\eta=1 \%\)\\
So,

\[
a e^{-\gamma x_{1}}-a e^{-\gamma x_{2}}=\eta a e^{-\gamma x_{1}}
\]

or

\[
e^{-\gamma x_{1}}(1-\eta)=e^{-\gamma x_{2}}
\]

or

\[
\ln (1-\eta)-\gamma x_{1}=-\gamma x_{2}
\]

or,

\[
x_{2}-x_{1}=-\frac{\ln (1-\eta)}{\gamma}
\]

So path difference \(=-\frac{\ln (1-\eta)}{\gamma}\)\\
and phase difference \(=\frac{2 \pi}{\lambda} \times\) path difference

\[
=-\frac{2 \pi}{\lambda} \frac{\ln (1-\eta)}{\gamma} \approx \frac{2 \pi \eta}{\lambda \cdot \gamma}=0.3 \mathrm{rad}
\]

4.158 Let \(S\) be the source whose position vector relative to the reference point \(O\) is \(\vec{r}\). Since intensities are inversely proportional to the square of distances,

\[
\frac{\text { Intensity at } P\left(I_{1}\right)}{\text { Intensity at } Q\left(I_{2}\right)}=\frac{d_{2}^{2}}{d_{1}^{2}}
\]

where \(d_{1}=P S\) and \(d_{2}=Q S\).\\
But intensity is proportional to the square of amplitude.\\
So, \(\frac{a_{1}^{2}}{a_{2}^{2}}=\frac{d_{2}^{2}}{d_{1}^{2}}\) or \(a_{1} d_{1}=a_{2} d_{2}=k\) (say)\\
Thus \(\quad d_{1}=\frac{k}{a_{1}}\) and \(d_{2}=\frac{k}{a_{2}}\)\\
Let \(\hat{n}\) be the unit vector along \(P Q\) directed from \(P\) to \(Q\).\\
Then

\[
\begin{aligned}
& \overrightarrow{P S}=d_{1} \hat{n}=\frac{k}{a_{1}} \hat{n} \\
& \overrightarrow{S Q}=d_{2} \hat{n}=\frac{k}{a_{2}} \hat{n}
\end{aligned}
\]

and

From the triangle law of vector addition.\\
\includegraphics[max width=\textwidth, alt={}, center]{da0156b6-4133-434e-90d4-ec01b2070afe-094_391_497_419_952}\\
or


\begin{gather*}
\overrightarrow{O P}+\overrightarrow{P S}=\overrightarrow{O S} \quad \text { or } \quad \overrightarrow{r_{1}}+\frac{k}{a_{1}} \hat{n}=\vec{r} \\
a_{1} \overrightarrow{r_{1}}+k \hat{n}=a_{1} \vec{r} \tag{1}
\end{gather*}


Similarly


\begin{equation*}
\vec{r}+\frac{k}{a_{2}} \hat{n}=\overrightarrow{r_{2}} \quad \text { or } \quad a_{2} \overrightarrow{r_{2}}-k \hat{n}=a_{2} \vec{r} \tag{2}
\end{equation*}


Adding (1) and (2),

Hence

\[
\begin{aligned}
a_{1} \overrightarrow{r_{1}}+a_{2} \overrightarrow{r_{2}} & =\left(a_{1}-a_{2}\right) \vec{r} \\
\vec{r} & =\frac{a_{1} \overrightarrow{r_{1}}+a_{2} \overrightarrow{r_{2}}}{a_{1}+a_{2}}
\end{aligned}
\]

4.159 (a) We know that the equation of a spherical wave in a homogeneous absorbing medium of wave damping coefficient \(\gamma\) is :

\[
\xi=\frac{a_{0}^{\prime} e^{-\gamma r}}{r} \cos (\omega t-k r)
\]

Thus particle's displacement amplitude equals

\[
\frac{a_{0}^{\prime} e^{-\gamma r}}{r}
\]

According to the conditions of the problem,\\
and when

\[
\text { at } \begin{array}{r}
r=r_{0}, a_{0}=\frac{a_{0}^{\prime} e^{-\gamma r_{0}}}{r_{0}} \\
r=r, \quad \frac{a_{0}}{\eta}=\frac{a_{0}^{\prime-\gamma r}}{r} \tag{2}
\end{array}
\]

Thus from Eqns (1) and (2)

\[
\begin{array}{cc} 
& e^{\gamma\left(r-r_{0}\right)}=\eta \frac{r_{0}}{r} \\
\text { or, } & \gamma\left(r-r_{0}\right)=\ln \left(\eta r_{0}\right)-\ln r \\
\text { or, } & \gamma=\frac{\ln \eta+\ln r_{0}-\ln r}{r-r_{0}}=\frac{\ln 3+\ln 5-\ln 10}{5}=0.08 \mathrm{~m}^{-1}
\end{array}
\]

(b) As \(\xi=\frac{a_{0}^{\prime} e^{-\gamma r}}{r} \cos (\omega t-k r)\)

So,

\[
\begin{gathered}
\frac{\partial \xi}{\partial t}=-\frac{a_{0}^{\prime} e^{-\gamma r}}{r} \omega \sin (\omega t-k r) \\
\left(\frac{\partial \xi}{\partial t}\right)_{n}=\frac{a_{0}^{\prime} e^{-\gamma r}}{r} \omega
\end{gathered}
\]

But at point A, \(\frac{a_{0}^{\prime} e^{-\gamma r}}{r}=\frac{a_{0}}{\eta}\)\\
So, \(\left(\frac{\partial \xi}{\partial t}\right)_{m}=\frac{a_{0} \omega}{\eta}=\frac{a_{0} 2 \pi}{\eta}=\frac{50 \times 10^{-6}}{3} \times 2 \times \frac{22}{7} \times 1.45 \times 10^{3}=15 \mathrm{~m} / \mathrm{s}\)\\
4.160 (a) Equation of the resultant wave,

\[
\begin{aligned}
& \xi=\xi_{1}+\xi_{2}=2 a \cos k\left(\frac{y-x}{2}\right) \cos \left\{\omega t-\frac{k(x+y)}{2}\right\}, \\
= & a^{\prime} \cos \left\{\omega t-\frac{k(x+y)}{2}\right\}, \text { where } a^{\prime}=2 a \cos k^{\prime}\left(\frac{y-x}{2}\right)
\end{aligned}
\]

Now, the equation of wave pattern is,

\[
x+y=k,(\text { a Const } .)
\]

For sought plots see the answer-sheet of the problem book.\\
For antinodes, i.e. maximum intensity\\
or,

\[
\begin{gathered}
\cos \frac{k(y-x)}{2}= \pm 1=\cos n \pi \\
\pm(x-y)=\frac{2 n \pi}{k}=n \lambda \\
y=x \pm n \lambda, n=0,1,2, \ldots
\end{gathered}
\]

or,\\
Hence, the particles of the medium at the points, lying on the solid straight lines ( \(y=x \pm n \lambda\) ), oscillate with maximum amplitude.\\
For nodes, i.e. minimum intensity,\\
or,

\[
\begin{gathered}
\cos \frac{k(y-x)}{2}=0 \\
\pm \frac{k(y-x)}{2}=(2 n+1) \frac{\pi}{2}
\end{gathered}
\]

or,

\[
y=x \pm(2 n+1) \lambda / 2,
\]

and hence the particles at the points, lying on dotted lines do not oscillate.\\
(b) When the waves are longitudinal,

For sought plots see the answer-sheet of the problem book.\\
or,


\begin{align*}
k & (y-x)=\cos ^{-1} \frac{\xi_{1}}{a}-\cos ^{-1} \frac{\xi_{2}}{a} \\
\frac{\xi_{1}}{a} & =\cos \left\{k(y-x)+\cos ^{-1} \frac{\xi_{2}}{a}\right\} \\
& =\frac{\xi_{2}}{a} \cos k(y-x)-\sin k(y-x) \sin \left(\cos ^{-1} \frac{\xi_{2}}{a}\right) \\
& =\frac{\xi_{2}}{a} \cos k(y-x)-\sin k(y-x) \sqrt{1-\frac{\xi_{2}^{2}}{a^{2}}} \tag{1}
\end{align*}


from (1),\\
if

\[
\begin{gathered}
\sin k(y-x)=0 \sin (n \pi) \\
\xi_{1}=\xi_{2}(-1)^{n}
\end{gathered}
\]

thus, the particles of the medium at the points lying on the straight lines, \(y=x \pm \frac{n \lambda}{2}\) will oscillate along those lines (even \(n\) ), or at right angles to them (odd \(n\) ).

Also from (1),\\
if

\[
\begin{gathered}
\cos k(y-x)=0=\cos (2 n+1) \frac{\pi}{2} \\
\frac{\xi_{1}^{2}}{a^{2}}=1-\xi_{2}^{2} / a^{2}, \text { a circle }
\end{gathered}
\]

Thus the particles, at the points, where \(y=x \pm(n \pm 1 / 4) \lambda\), will oscillate along circles. In general, all other particles will move along ellipses.\\
4.161 The displacement of oscillations is given by \(\xi=a \cos (\omega t-k x)\)

Without loss of generality, we confine ourselves to \(x=0\). Then the displacement maxima occur at \(\omega t=n \pi\). Concentrate at \(\omega t=0\). Now the energy density is given by

\[
w=\rho a^{2} \cdot \omega^{2} \sin ^{2} \omega t \quad \text { at } x=0
\]

\(T / 6\) time later (where \(T=\frac{2 \pi}{\omega}\) is the time period) than \(t=0\).

Thus

\[
\begin{gathered}
w=\rho a^{2} \omega^{2} \sin ^{2} \frac{\pi}{3}=\frac{3}{4} \rho a^{2} \omega^{2}=w_{0} \\
<w>=\frac{1}{2} \rho a^{2} \omega^{2}=\frac{2 w_{0}}{3} .
\end{gathered}
\]

4.162 The power output of the source much be

\[
4 \pi l^{2} I_{0}=Q \text { Watt. }
\]

The required flux of accoustic power is then : \(Q=\frac{\Omega}{4 \pi}\)\\
Where \(\Omega\) is the solid angle subtended by the disc enclosed by the ring at \(S\). This solid angle is\\
\includegraphics[max width=\textwidth, alt={}, center]{da0156b6-4133-434e-90d4-ec01b2070afe-097_227_424_141_983}

So flux

\[
\Omega=2 \pi(1-\cos \alpha)
\]

\[
\Phi=I_{0} I_{0}\left(1-\frac{l}{\sqrt{r^{2}+R^{2}}}\right) 2 \pi l^{2}
\]

Substitution gives \(\Phi=2 \Pi \times 30\left(1-\frac{l}{\sqrt{1+\frac{1}{4}}}\right) \mu \mathrm{W}=1.99 \mu \mathrm{~W}\).\\
Eqn. (1) is a well known result stich is derived as follows; Let \(S O\) be the polar axis. Then the required solid angle is the area of that part of the surface a sphere of much radius whose colatitude is \(\leq \alpha\).

Thus

\[
\Omega=\int_{0}^{\alpha} 2 \pi \sin \theta d \theta=2 \pi(1-\cos \alpha) .
\]

4.163 From the result of 4.162 power flowing out through anyone of the opening

\[
\begin{gathered}
=\frac{P}{2}\left(1-\frac{h / 2}{\sqrt{R^{2}+(h / 2)^{2}}}\right) \\
=\frac{P}{2}\left(1-\frac{h}{\sqrt{4 R^{2}+h^{2}}}\right)
\end{gathered}
\]

As total power output equals \(P\), so the power reaching the lateral surface must be.

\[
=P-2 \cdot \frac{P}{2}\left(1-\frac{h}{\sqrt{4 R^{2}+h^{2}}}\right)=\frac{p h}{\sqrt{4 R^{2}+h^{2}}}=0.07 \mathrm{~W}
\]

4.164 We are given\\
so

\[
\frac{\partial \xi}{\partial x}=-a k \sin k x \cos \omega t \quad \text { and } \quad \frac{\partial \xi}{\partial t}=-a \omega \cos k x \sin \omega t
\]

Thus

\[
\begin{aligned}
(\xi)_{t-0} & =a \cos k x,(\xi)_{t-T / 2}=-a \cos k x \\
\left(\frac{\partial \xi}{\partial x}\right)_{t=0} & =a k \sin k x,\left(\frac{\partial \xi}{\partial x}\right)_{t=T / 2}=-a k \sin k x
\end{aligned}
\]

(a) The graphs of \((\xi)\) and \(\left(\frac{\partial \xi}{\partial x}\right)\) are as shown in Fig. (35) of the book (p.332).\\
(b) We can calculate the density as follows :

Take a parallelopiped of cross section unity and length \(d x\) with its edges at \(x\) and \(x+d x\).

After the oscillation the edge at \(x\) goes to \(x+\xi(x)\) and the edge at \(x+d x\) goes to \(x+d x+\xi(x+d x)\)\\
\(=x+d x+\xi(x)+\frac{\partial \xi}{\partial x} d x\). Thus the volume of the element (originally \(\left.d x\right)\) becomes

\[
\left(1+\frac{\partial \xi}{\partial x}\right) d x
\]

and hence the density becomes \(\rho=\frac{\rho_{0}}{1+\frac{\partial \xi}{\partial x}}\).\\
I On substituting we get for the density \(\rho(x)\) the curves shown in Fig.(35). referred to above.\\
(c) The velocity \(v(x)\) at time \(t=T / 4\) is

\[
\left(\frac{\partial \xi}{\partial t}\right)_{t-T / 4}=-a \omega \cos k x
\]

On plotting we get the figure (36).\\
4.165 Given \(\xi=a \cos k x \cos \omega t\)\\
(a) The potential energy density (per unit volume) is the energy of longitudinal strain. This is

\[
\begin{aligned}
& \begin{aligned}
w_{p}=\left(\frac{1}{2} \text { stress } \text { × } \text { strain }\right) & =\frac{1}{2} E\left(\frac{\partial \xi}{\partial x}\right)^{2}, \quad\left(\frac{\partial \xi}{\partial x} \text { is the longitudinal strain }\right) \\
w_{p} & =\frac{1}{2} E a^{2} k^{2} \sin ^{2} k x \cos ^{2} \omega t
\end{aligned} \\
& \text { But } \quad \begin{aligned}
\frac{\omega^{2}}{k^{2}} & =\frac{E}{\rho} \text { or } E k^{2}=\rho \omega^{2}
\end{aligned}
\end{aligned}
\]

Thus

\[
w_{p}=\frac{1}{2} \rho a^{2} \omega^{2} \sin ^{2} k x \cos ^{2} \omega t
\]

(b) The kinetic energy density is

\[
=\frac{1}{2} \rho\left(\frac{\partial \xi}{\partial t}\right)^{2}=\frac{1}{2} \rho a^{2} \omega^{2} \cos ^{2} k x \sin ^{2} \omega t .
\]

On plotting we get Fig. 37 given in the book (p. 332). For example at \(t=0\)

\[
w=w_{p}+w_{k}=\frac{1}{2} \rho a^{2} \omega^{2} \sin ^{2} k x
\]

and the displacement nodes are at \(x= \pm \frac{\pi}{2 k}\) so we do get the figure.\\
4.166 Let us denote the displacement of the elements of the string by

\[
\xi=a \sin k x \cos \omega t
\]

since the string is 120 cm long we must have \(k \cdot 120=n \pi\)\\
If \(x_{1}\) is the distance at which the displacement amplitude first equals 3.5 mm then

\[
a \sin k x_{1}=3 \cdot 5=a \sin \left(k x_{1}+15 k\right)
\]

Then

\[
k x_{1}+15 k=\pi-k x_{1} \text { or } k x_{1}=\frac{\pi-15 k}{2}
\]

One can convince ourself that the string has the form shown below\\
\includegraphics[max width=\textwidth, alt={}, center]{da0156b6-4133-434e-90d4-ec01b2070afe-099_259_837_716_313}

It shows that

\[
k \times 120=4 \pi, \text { so } k=\frac{\pi}{30} \mathrm{~cm}^{-1}
\]

Thus we are dealing with the third overtone\\
Also

\[
k x_{1}=\frac{\pi}{4} \quad \text { so } \quad a=3.5 \sqrt{2} \mathrm{~mm} \approx 4.949 \mathrm{~mm} .
\]

4.167 We have \(n=\frac{1}{2 l} \sqrt{\frac{T}{m}}=\frac{1}{2 l} \sqrt{\frac{T l}{M}}\) Where \(M=\) total mass of the wire. When the wire is stretched, total mass of the wire remains constant. For the first wire the new length \(=l+\eta_{1} l\) and for the record wire, the length is \(l+\eta_{2} l\). Also \(T_{1}=\alpha\left(\eta_{1} l\right)\) where \(\alpha\) is a constant and \(T_{2}=\alpha\left(\eta_{2} l\right)\). Substituting in the above formula.

\[
\begin{gathered}
v_{1}=\frac{1}{2\left(l+\eta_{1} l\right)} \sqrt{\frac{\left(\alpha \eta_{1} l\right)\left(l+\eta_{1} l\right)}{M}} \\
v_{2}=\frac{1}{2\left(l+\eta_{2} l\right)} \sqrt{\frac{\left(\alpha \eta_{2} l\right)\left(l+\eta_{2} l\right)}{M}} \\
\therefore \frac{v_{2}}{v_{1}}=\frac{1+\eta_{1}}{1+\eta_{2}} \sqrt{\frac{\eta_{2}}{\eta_{1}} \cdot \frac{1+\eta_{2}}{1+\eta_{1}}} \\
\frac{v_{2}}{v_{1}}=\sqrt{\frac{\eta_{2}\left(1+\eta_{1}\right)}{\eta_{2}\left(1+\eta_{2}\right)}}=\sqrt{\frac{0.04(1+0.02)}{0.02(1+0.04)}}=1.4
\end{gathered}
\]

4.168 Let initial length and tension be \(l\) and \(T\) respectively.

So,

\[
v_{1}=\frac{1}{2 l} \sqrt{\frac{T}{\rho_{1}}}
\]

In accordance with the problem, the new length

\[
l^{\prime}=l-\frac{l \times 35}{100}=0.65 l
\]

and new tension, \(T^{\prime}=T+\frac{T \times 70}{100}=1.7 T\)\\
Thus the new frequency

Hence

\[
\begin{gathered}
v_{2}=\frac{1}{2 l^{\prime}} \sqrt{\frac{T^{\prime}}{\rho_{1}}}=\frac{1}{2 \times 0.65 l} \sqrt{\frac{1.7 T}{\rho_{1}}} \\
\frac{v_{2}}{v_{1}}=\frac{\sqrt{1.7}}{0.65}=\frac{1.3}{0.65}=2
\end{gathered}
\]

4.169 Obviously in this case the velocityof sound propagation

\[
v=2 v\left(l_{2}-l_{1}\right)
\]

where \(l_{2}\) and \(l_{1}\) are consecutive lengths at which resonance occur\\
In our problem, \(\left(l_{2}-l_{1}\right)=l\)\\
So

\[
v=2 v l=2 \times 2000 \times 8.5 \mathrm{~cm} / \mathrm{s}=0.34 \mathrm{~km} / \mathrm{s} .
\]

4.170 (a) When the tube is closed at one end

Thus for

\[
\begin{aligned}
v & =\frac{v}{41}(2 n+1), \text { where } n=0,1,2, \ldots \\
& =\frac{340}{4 \times 0.85}(2 n+1)=100(2 n+1)
\end{aligned}
\]

\[
\begin{gathered}
n=0,1,2,3,4,5,6, \ldots, \text { we get } \\
n_{1}=1001 H_{z}, n_{2}=3001 H_{z}, n_{3}=5001 H_{z}, n_{4}=7001 H_{z} \\
n_{5}=9001 H_{z}, n_{6}=11001 H_{z}, n_{7}=13001 H_{z}
\end{gathered}
\]

Since \(v\) should be \(<v_{0}=12501 H_{z}\), we need not go beyond \(n_{6}\).\\
Thus 6 natural oscillations are possible.\\
(b) Organ pipe opened from both ends vibrates with all harmonics of the fundamental frequency. Now, the fundamental mode frequency is given as

\[
v=v / \lambda
\]

or,

\[
v=v / 2 l
\]

Here, also, end correction has been neglected. So, the frequencies of higher modes of vibrations are given by


\begin{equation*}
v=n(v / 2 l) \tag{1}
\end{equation*}


or,

\[
v_{1}=v / 2 l, v_{2}=2(v / 2 l), v_{3}=3(v / 2 l)
\]

It may be checked by putting the values of \(n\) in the equation (1) that below 1285 Hz , there are a total of six possible natural oscillation frequencies of air column in the open pipe.\\
4.171 Since the copper rod is clamped at mid point, it becomes a mode and the two free ends will be anitinodes. Thus the fundamental mode formed in the rod is as shown in the Fig. (a).

\begin{figure}[h]
\begin{center}
  \includegraphics[alt={},max width=\textwidth]{da0156b6-4133-434e-90d4-ec01b2070afe-101_278_546_414_86}
\captionsetup{labelformat=empty}
\caption{4.171 (a)}
\end{center}
\end{figure}

\begin{figure}[h]
\begin{center}
  \includegraphics[alt={},max width=\textwidth]{da0156b6-4133-434e-90d4-ec01b2070afe-101_274_723_422_705}
\captionsetup{labelformat=empty}
\caption{4.171 (b)}
\end{center}
\end{figure}

In this case,

\[
l=\frac{\lambda}{2}
\]

So,

\[
v_{0}=\frac{\mathrm{v}}{2 l}=\frac{1}{2 l} \sqrt{\frac{E}{\rho}} \sqrt{\frac{E}{e}}
\]

where \(E=\) Young's modules and \(\rho\) is the density of the copper\\
Similarly the second mode or the first overtone in the rod is as shown above in Fig. (b).\\
Here

Hence

\[
\begin{gathered}
l=\frac{3 \lambda}{2} \\
v_{1}=\frac{3 v}{2 l}=\frac{3}{2 l} \sqrt{\frac{E}{\rho}} \\
v=\frac{2 n+1}{2 l} \sqrt{\frac{E}{\rho}} \text { where } n=0,1,2 \ldots
\end{gathered}
\]

Putting the given values of \(E\) and \(\rho\) in the general equation

\[
v=3.8(2 n+1) \mathrm{k} \mathrm{~Hz}
\]

Hence \(v_{0}=3.8 \mathrm{k} . \mathrm{Hz}, v_{1}=(3.8 \times 3) \mathrm{k} \mathrm{Hz}, v_{2}=(3.8) \times 5=19 \mathrm{k} \mathrm{Hz}\),

\[
\begin{aligned}
& v_{3}=(3.8 \times 7)=26.6 \mathrm{k} \mathrm{~Hz}, v_{4}=(3.8 \times 9)=34.2 \mathrm{k} \mathrm{~Hz} \\
& v_{5}=(3.8 \times 11)=41.8 \mathrm{k} \mathrm{~Hz}, v_{6}=(3.8) \times 13 \mathrm{k} \mathrm{~Hz}=49.4 \mathrm{k} \mathrm{~Hz} \text { and } \\
& v_{7}=(3.8) \times 14 \mathrm{k} \mathrm{~Hz}>50 \mathrm{k} \mathrm{~Hz}
\end{aligned}
\]

Hence the sought number of frequencies between 20 to 50 k Hz equals 4 .\\
4.172 Let two waves \(\xi_{1}=a \cos (\omega t-k x)\) and \(\xi_{2}=a \cos (\omega t+k x)\), superpose and as a result,\\
we have a standing wave (the resultant wave) in the string of the form \(\xi=2 a \cos k x \cos \omega t\).\\
According to the problem \(2 a=a_{m}\).

Hence the standing wave excited in the string is


\begin{align*}
\xi & =a_{m} \cos k x \cos \omega t  \tag{1}\\
\frac{\partial \xi}{\partial t} & =-\omega a_{m} \cos k x \sin \omega t \tag{2}
\end{align*}


or,

So the kinetic energy confined in the string element of length \(d x\), is given by :

\[
\text { or, } \begin{aligned}
d T & =\frac{1}{2}\left(\frac{m}{l} d x\right)\left(\frac{\partial \xi}{\partial t}\right)^{2} \\
d T & =\frac{1}{2}\left(\frac{m}{l} d x\right) a_{m}^{2} \omega^{2} \cos ^{2} k x \sin ^{2} \omega t \\
\text { or, } & d T \\
d T & =\frac{m a_{m}^{2} \omega^{2}}{2 l} \sin ^{2} \omega t \cos ^{2} \frac{2 \pi}{\lambda} x d x
\end{aligned}
\]

Hence the kinetic energy confined in the string corresponding to the fundamental tone

\[
T=\int d T=\frac{m a_{m}^{2} \omega^{2}}{2 l} \sin ^{2} \omega t \int_{0}^{\lambda / 2} \cos ^{2} \frac{2 \pi}{\lambda} x d x
\]

Because, for the fundamental tone, length of the string \(l=\frac{\lambda}{2}\)\\
Integrating we get, \(\quad: \quad T=\frac{1}{4} m a_{m}^{2} \omega^{2} \sin ^{2} \omega t\)\\
Hence the sought maximum kinetic energy equals, \(T_{\max }=\frac{1}{4} m a_{m}^{2} \omega^{2}\), because for \(T_{\text {max }}, \sin ^{2} \omega t=1\)\\
(ii) Mean kinetic energy averaged over one oscillation period

\[
\begin{gathered}
<T>=\frac{\int T d t}{\int d t}=\frac{1}{4} m a_{m}^{2} \omega^{2} \frac{\int_{0}^{2 \pi / \omega} \sin ^{2} \omega t d t}{\int_{0}^{2 \pi / \omega} d t} \\
<T>=\frac{1}{8} m a_{m}^{2} \omega^{2} .
\end{gathered}
\]

or,\\
4.173 We have a standing wave given by the equation

So,


\begin{equation*}
\frac{\partial \xi}{\partial t}=-a \omega \sin k x \sin \omega t \tag{1}
\end{equation*}


and


\begin{equation*}
\frac{\partial \xi}{\partial t}=a k \cos k x \cos \omega t \tag{2}
\end{equation*}


The kinetic energy confined in an element of length \(d x\) of the rod

\[
d T=\frac{1}{2}(\rho S d x)\left(\frac{\partial \xi}{\partial t}\right)^{2}=\frac{1}{2} \rho S a^{2} \omega^{2} \sin ^{2} \omega t \sin ^{2} k x d x
\]

So total kinetic energy confined into rod\\
or,


\begin{gather*}
T=\int d T=\frac{1}{2} \rho S a^{2} \omega^{2} \sin ^{2} \omega t \int_{0}^{\lambda / 2} \sin ^{2} \frac{2 \pi}{\lambda} x d x \\
T=\frac{\pi S a^{2} \omega^{2} \rho \sin ^{2} \omega t}{4 k} \tag{3}
\end{gather*}


The potential energy in the above rod element\\
or,

\[
d U=\int \partial U=-\int_{0}^{\xi} F_{\xi} d \xi, \text { where } F_{\xi}=(\rho S d x) \frac{\partial^{2} \xi}{\partial t^{2}}
\]

so,

\[
F_{\xi}=-(\rho S d x) \omega^{2} \xi
\]

or,

\[
d U=\frac{\rho \omega^{2} S \xi^{2}}{2} d x=\frac{\rho \omega^{2} S a^{2} \cos ^{2} \omega t \sin ^{2} k x d x}{2}
\]

Thus the total potential energy stored in the rod \(U=\int d U\)\\
or,

\[
\begin{gathered}
U=\rho \omega^{i 2} S a^{2} \cos ^{2} \omega t \int_{0}^{\lambda / 2} \sin ^{2} \frac{2 \pi}{\lambda} x d x \\
U=\frac{\pi \rho S a^{2} \omega^{2} \cos ^{2} \omega t}{4 k}
\end{gathered}
\]

So,

\[
d U=\omega^{2} \rho S d x \int_{0}^{\xi} \xi d \xi
\]

\(\lambda / 2\)\\
r,

To find the potential energy stored in the rod element we may adopt an easier way. We know that the potential energy density confined in a rod under elastic force equals :

\[
\begin{aligned}
U_{D} & =\frac{1}{2}(\text { stress } \times \text { strain })=\frac{1}{2} \sigma \varepsilon=\frac{1}{2} Y \varepsilon^{2} \\
& =\frac{1}{2} \rho v^{2} \varepsilon^{2}=\frac{1}{2} \frac{\rho \omega^{2}}{k^{2}} \varepsilon^{2} \\
& =\frac{1}{2} \frac{\rho \omega^{2}}{k^{2}}\left(\frac{\partial \xi}{\partial x}\right)^{2}=\frac{1}{2} \rho a^{2} \omega^{2} \cos ^{2} \omega t \cos ^{2} k x
\end{aligned}
\]

Hence the total potential energy stored in the rod


\begin{align*}
U & =\int U_{D} d V=\int_{0}^{\lambda / 2} \frac{1}{2} \rho a^{2} \omega^{2} \cos ^{2} \omega t \cos ^{2} k x S d x \\
& =\frac{\pi \rho S a^{2} \omega^{2} \cos ^{2} \omega t}{4 k} \tag{4}
\end{align*}


Hence the sought mechanical energy confined in the rod between the two adjacent nodes

\[
E=T+U=\frac{\pi \rho \omega^{2} a^{2} S}{4 k}
\]

4.174 Receiver \(R_{1}\) registers the beating, due to the sound waves reaching directly to it from source and the other due to the reflection from the wall.\\
Frequency of sound reaching directly from \(S\) to \(R_{1}\)

\[
v_{S \rightarrow R_{1}}=v_{0} \frac{\mathrm{v}}{\mathrm{v}-\mathrm{u}} \text { when } S \text { moves towards } R_{1}
\]

and

\[
v_{S \rightarrow R_{1}}^{\prime}=v_{0} \frac{\mathrm{v}}{\mathrm{v}+\mathrm{u}} \text { when } S \text { moves towards the wall }
\]

Now frequency reaching to \(R_{1}\) after reflection from wall

\[
v_{W \rightarrow R_{1}}=v_{0} \frac{\mathrm{v}}{\mathrm{v}+\mathrm{u}}, \text { when } S \text { moves towards } R_{1}
\]

and \(v_{W \rightarrow R_{1}}^{\prime}=v_{0} \frac{v}{v-u}\), when \(S\) moves towards the wall\\
\includegraphics[max width=\textwidth, alt={}, center]{da0156b6-4133-434e-90d4-ec01b2070afe-104_274_487_880_993}

Thus the sought beat frequency

\[
\begin{aligned}
\Delta v & =\left(v_{S \rightarrow R_{1}}-v_{W \rightarrow R_{1}}\right) \quad \text { or } \quad\left(v_{W \rightarrow R_{1}}^{\prime}-v_{S \rightarrow R_{1}}\right) \\
& =v_{0} \frac{v}{v-u}-v_{0} \frac{v}{v+u}=\frac{2 v_{0} v u}{v^{2}-u^{2}} \cong \frac{2 u v_{0}}{v}=1 \mathrm{~Hz}
\end{aligned}
\]

4.175 Let the velocity of tuning fork is \(u\). Thus frequency reaching to the observer due to the tuning fork that approaches the observer

\[
v^{\prime}=v_{0} \frac{v}{v-u}[v=\text { velocity of sound }]
\]

Frequency reaching the observer due to the tunning fork that recedes from the observer

\[
v^{\prime \prime}=v_{0} \frac{v}{v+u}
\]

So, Beat frequency \(v-v^{\prime \prime}=v=v_{0} v\left(\frac{1}{v-u}-\frac{1}{v+u}\right)\)\\
or,

\[
v=\frac{2 v_{0} v u}{v^{2}-u^{2}}
\]

So,

\[
v u^{2}+\left(2 v v_{0}\right) u-v^{2} v=0
\]

Hence

\[
u=\frac{-2 v v_{0} \pm \sqrt{4 v_{0}^{2} v^{2}+4 v^{2} v^{2}}}{2 v}
\]

Hence the sought value of \(u\), on simplifying and noting that \(u>0\)

\[
u=\frac{v v_{0}}{v}\left(\sqrt{1+\left(\frac{v}{v_{0}}\right)^{2}}-1\right)
\]

4.176 Obviously the maximum frequency will be heard when the source is moving with maximum velocity towards the receiver and minimum frequency will be heard when the source recedes with maximum velocity. As the source swing harmonically its maximum velocity equals \(a \omega\). Hence

\[
v_{\max }=v_{0} \frac{v}{v-a \omega} \text { and } v_{\min }=v_{0} \frac{v}{v+a \omega}
\]

So the frequency band width \(\Delta v=v_{\text {max }}-v_{\text {min }}=v_{0} v\left(\frac{2 a \omega}{v^{2}-a^{2} \omega^{2}}\right)\)\\
or,

\[
\left(\Delta v a^{2}\right) \omega^{2}+\left(2 v_{0} v a\right) \omega-\Delta v v^{2}=0
\]

So,

\[
\omega=\frac{-2 v_{0} v a \pm \sqrt{4 v_{0}^{2} v^{2} a^{2}+\Delta v^{2} a^{2} v^{2}}}{2 \Delta v a^{2}}
\]

On simplifying (and taking + sign as \(\omega \rightarrow 0\) if \(\Delta \nu \rightarrow 0\) )

\[
\omega,=\frac{v v_{0}}{\Delta v a}\left(\sqrt{1+\left(\frac{\Delta v}{v_{0}}\right)^{2}}-1\right)
\]

4.177 It should be noted that the frequency emitted by the source at time \(t\) could not be received at the same moment by the receiver, becouse till that time the source will cover the distance \(\frac{1}{2} w t^{2}\) and the sound wave will take the further time \(\frac{1}{2} w t^{2} / \mathrm{v}\) to reach the receiver. Therefore the frequency noted by the receiver at time \(t\) should be emitted by the source at the time \(t_{1}<t\). Therefore


\begin{equation*}
t_{1}+\left(\frac{1}{2} w t_{1}^{2} / \mathrm{v}\right)=t \tag{1}
\end{equation*}


and the frequency noted by the receiver


\begin{equation*}
v=v_{0} \frac{v}{v+w t_{1}} \tag{2}
\end{equation*}


Solving Eqns (1) and (2), we get

\[
v=\frac{v_{0}}{\sqrt{1+\frac{2 w t}{v}}}=1.35 \mathrm{kHz} .
\]

4.178 (a) When the observer receives the sound, the source is closest to him. It means, that frequency is emitted by the source sometimes before (Fig.) Figure shows that the source approaches the stationary observer with velocity \(v_{s} \cos \theta\).\\
Hence the frequency noted by the observer


\begin{align*}
& \quad v=v_{0}\left(\frac{v}{v-v_{s} \cos \theta}\right) \\
& =v_{0}\left(\frac{v}{v-\eta v \cos \theta}\right)=\frac{v_{0}}{1-\eta \cos \theta}  \tag{1}\\
& \text { But } \quad \frac{x}{v_{s}}=\frac{\sqrt{l^{2}+x^{2}}}{v}, \text { So, } \frac{x}{\sqrt{l^{2}+x^{2}}}=\frac{v_{s}}{v}=\eta \\
& \text { or, } \quad \cos \theta=\eta
\end{align*}


Hence from Eqns. (1) and (2) the sought frequency

\[
v=\frac{v_{0}}{1-\eta^{2}}=5 \mathrm{kHz}
\]

(b) When the source is right in front of \(O\), the sound emitted by it will not be Doppler shifted because \(\theta=90^{\circ}\). This sound will be received at \(O\) at time \(t=\frac{l}{v}\) after the source has passed it. The source will by then have moved ahead by a distance \(v_{s} t=l \eta\). The distance between the source and the observer at this time will be \(l \sqrt{1+\eta^{2}}=0.32 \mathrm{~km}\).\\
4.179 Frequency of sound when it reaches the wall

\[
v^{\prime}=v \frac{v+u}{v}
\]

wall will reflect the sound with same frequency \(v^{\prime}\). Thus frequency noticed by a stationary observer after reflection from wall\\
\(v^{\prime \prime}=v^{\prime} \frac{v}{v-u}\), since wall behaves as a source of frequency \(v^{\prime}\).\\
Thus,

\[
v^{\prime \prime}=v \frac{v+u}{v} \cdot \frac{v}{v-u}=v \frac{v+u}{v-u}
\]

or,

\[
\lambda^{\prime \prime}=\lambda \frac{\mathrm{v}-u}{\mathrm{v}+u} \quad \text { or } \quad \frac{\lambda^{\prime \prime}}{\lambda}=\frac{\mathrm{v}-u}{\mathrm{v}+u}
\]

So,

\[
1-\frac{\lambda^{\prime \prime}}{\lambda}=1-\frac{\mathrm{v}-u}{\mathrm{v}+u}=\frac{2 u}{\mathrm{v}+u}
\]

Hence the sought percentage change in wavelength

\[
=\frac{\lambda-\lambda^{\prime \prime}}{\lambda} \times 100=\frac{2 u}{v+u} \times 100 \%=0.2 \% \text { decrease. }
\]

4.180 Frequency of sound reaching the wall.


\begin{equation*}
v=v_{0}\left(\frac{v-u}{v}\right) \tag{1}
\end{equation*}


Now for the observer the wall becomes a source of frequency \(v\) receding from it with velocity \(u\) Thus, the frequency reaching the observer


\begin{equation*}
v^{\prime}=v\left(\frac{v}{v+u}\right)=v_{0}\left(\frac{v-u}{v+u}\right) \tag{1}
\end{equation*}


Hence the beat frequency registered by the receiver (observer)

\[
\Delta v=v_{0}-v^{\prime}=\frac{2 u v_{0}}{v+u}=0.6 \mathrm{~Hz}
\]

4.181 Intensity of a spherical sound wave emitted from a point source in a homogeneous absorbing medium of wave damping coefficient \(\gamma\) is given by

\[
I=\frac{1}{2} \rho a^{2} e^{-2 \gamma r} \omega^{2} v
\]

So, Intensity of sound at a distance \(r_{1}\) from the source

\[
=\frac{I_{1}}{r_{1}^{2}}=\frac{1 / 2 \rho a^{2} e^{-2 \gamma r_{1}} \omega^{2} v}{r_{1}^{2}}
\]

and intensity of sound at a distance \(r_{2}\) from the source

\[
=I_{2} / r_{2}^{2}=\frac{1 / 2 \rho a^{2} e^{-2 \gamma r_{2}} \omega^{2} v}{r_{2}^{2}}
\]

But according to the problem

\[
\frac{1}{\eta} \frac{I_{1}}{r_{1}^{2}}=\frac{I_{2}}{r_{2}^{2}}
\]

So,

\[
\frac{\eta r_{1}^{2}}{r_{2}^{2}}=e^{2 \gamma\left(r_{2}-r_{1}\right)} \quad \text { or } \quad \ln \frac{\eta r_{2}^{2}}{r_{1}^{2}}=2 \gamma\left(r_{2}-r_{1}\right)
\]

or,

\[
\gamma=\frac{\ln \left(\eta r_{2}^{2} / r_{1}^{2}\right)}{2\left(r_{2}-r_{1}\right)}=6 \times 10^{-3} \mathrm{~m}^{-1}
\]

4.182 (a)Loudness ievel in bells \(=\log \frac{I}{I_{0}}\). ( \(I_{0}\) is the theshold of audibility.)

So, loudness level in decibells, \(L=10 \log \frac{I}{I_{0}}\)\\
Thus loudness level at \(x=x_{1}=L_{x_{1}}=10 \log \frac{I_{x_{1}}}{I_{0}}\)

Similarly

\[
L_{x_{2}}=10 \log \frac{I_{x_{2}}}{I_{0}}
\]

Thus

\[
L_{x_{2}}-L_{x_{1}}=10 \log \frac{I_{x_{2}}}{I_{x_{1}}}
\]

or, \(L_{x_{2}}=L_{x_{1}}+10 \log \frac{1 / 2 \rho a^{2} \omega^{2} v e^{-2 \gamma x_{2}}}{1 / 2 \rho a^{2} \omega^{2} v e^{-2 \gamma x_{1}}}=L_{x_{1}}+10 \log e^{-2 \gamma\left(x_{2}-x_{1}\right)}\)

\[
L_{x_{2}}=L_{x_{1}}-20 \gamma\left(x_{2}-x_{1}\right) \log e
\]

Hence

\[
\begin{aligned}
L^{\prime} & =L-20 \gamma x \log e\left[\text { since }\left(x_{2}-x_{1}\right)=x\right] \\
& =20 d B-20 \times 0 \cdot 23 \times 50 \times 0 \cdot 4343 d B \\
& =60 d B-10 d B=50 d B
\end{aligned}
\]

(b) The point at which the sound is not heard any more, the loudness level should be zero.

Thus

\[
0=L-20 \gamma x \log e \quad \text { or } \quad x=\frac{L}{20 \gamma \log e}=\frac{60}{20 \times 0.23 \times 0.4343}=300 \mathrm{~m}
\]

4.183 (a) As there is no damping, so

\[
L_{r_{0}}=10 \log \frac{I}{I_{0}}=10 \log \frac{1 / 2 \rho a^{2} \omega^{2} v / r_{0}^{2}}{1 / 2 \rho a^{2} \omega^{2} v}=-20 \log r_{0}
\]

Similarly \(L_{r}=-20 \log r\)\\
So,

\[
L_{r}-L_{r_{0}}=20 \log \left(r_{0} / r\right)
\]

or, \(L_{r}=L_{r_{0}}+20 \log \left(\frac{r_{0}}{r}\right)=30+20 \times \log \frac{20}{10}=36 d B\)\\
(b) Let \(r\) be the sought distance at which the sound is not heard.

So, \(L_{r}=L_{r_{0}}+20 \log \frac{r_{0}}{r}=0\) or, \(L_{r_{0}}=20 \log \frac{r}{r_{0}} \quad\) or \(\quad 30=20 \log \frac{r}{20}\)\\
So,

\[
\log _{10} \frac{r}{20}=3 / 2 \text { or } 10^{(3 / 2)}=r / 20
\]

Thus

\[
r=200 \sqrt{10}=0.63 \mathrm{Km} .
\]

Thus for \(r>0.63 \mathrm{~km}\) no sound will be heard.\\
4.184 We treat the fork as a point source. In the absence of damping the oscillation has the form

\[
\frac{\text { Const. }}{r} \cos (\omega t-k r)
\]

Because of the damping of the fork the amplitude of oscillation decreases exponentially with the retarded time (i.e. the time at which the wave started from the source.). Thus we write for the wave amplitude.

This means that

\[
\begin{gathered}
\xi=\frac{\text { Const. }}{r} e^{-\beta\left(t-\frac{r}{v}\right)} \\
\frac{e^{-\beta\left(t+\tau-\frac{r}{v}\right)}}{r_{A}}=\frac{e^{-\beta\left(t+\tau-\frac{r_{B}}{v}\right)}}{r_{B}}
\end{gathered}
\]

\begin{figure}[h]
\begin{center}
  \includegraphics[alt={},max width=\textwidth]{da0156b6-4133-434e-90d4-ec01b2070afe-108_258_147_1671_1023}
\captionsetup{labelformat=empty}
\caption{x}
\end{center}
\end{figure}

\begin{figure}[h]
\begin{center}
  \includegraphics[alt={},max width=\textwidth]{da0156b6-4133-434e-90d4-ec01b2070afe-108_259_164_1677_1307}
\captionsetup{labelformat=empty}
\caption{\(x+d x\)}
\end{center}
\end{figure}

Thus

\[
e^{-\beta\left(\tau+\frac{r_{B}-r_{A}}{v}\right)}=\frac{r_{A}}{r_{B}} \quad \text { or } \beta=\frac{\ln \frac{r_{B}}{r_{A}}}{\tau+\frac{r_{B}-r_{A}}{v}}=0.12 . \mathrm{s}^{-1}
\]

4.185 (a) Let us consider the motion of an element of the medium of thickness \(d x\) and unit area of cross-section. Let \(\xi=\) displacement of the particles of the medium at location \(x\). Then by the equation of motion

\[
\rho d x \dot{\xi}=-d p
\]

where \(d p\) is the pressure increment over the length \(d x\)\\
Recalling the wave equation

\[
\ddot{\xi}=v^{2} \frac{\partial^{2} \xi}{\partial x^{2}}
\]

we can write the foregoing equation as

\[
\rho \mathrm{v}^{2} \frac{\partial^{2} \xi}{\partial x^{2}} d x=-d p
\]

Integrating this equation, we get

\[
\Delta p=\text { surplus pressure }=-\rho \mathrm{v}^{2} \frac{\partial \xi}{\partial x}+\text { Const. }
\]

In the absence of a deformation (a wave), the surplus pressure is \(\Delta p=0\). So 'Const' \(=\) 0 and

\[
\Delta p=-\rho v^{2} \frac{\partial \xi}{\partial x}
\]

(b) We have found earlier that

\[
\begin{gathered}
w=w_{k}+w_{p}=\text { total energy density } \\
w_{k}=\frac{1}{2} \rho\left(\frac{\partial \xi}{\partial t}\right)^{2}, w_{p}=\frac{1}{2} E\left(\frac{\partial \xi}{\partial x}\right)^{2}=\frac{1}{2} \rho \mathrm{v}^{2}\left(\frac{\partial \xi}{\partial x}\right)^{2}
\end{gathered}
\]

It is easy me see that the space-time average of both densities is the same and the space time average of total energy density is then

\[
<w>=<\rho v^{2}\left(\frac{\partial \xi}{\partial x}\right)^{2}>
\]

The intensity of the wave is

\[
I=\mathrm{v}\langle w\rangle=\left\langle\frac{(\Delta p)^{2}}{\rho \mathrm{v}}\right\rangle
\]

Using

\[
<(\Delta p)^{2}>=\frac{1}{2}(\Delta p)_{m}^{2} \quad \text { we get } \quad I=\frac{(\Delta p)_{m}^{2}}{2 \rho \mathrm{v}}
\]

4.186 The intensity of the sound wave is

\[
I=\frac{(\Delta p)_{m}^{2}}{2 \rho v}=\frac{(\Delta p)_{m}^{2}}{2 \rho v \lambda}
\]

Using \(v=v \lambda, \rho\) is the density of air.\\
Thus the mean energy flow reaching the ball is

\[
\pi R^{2} I=\pi R^{2} \frac{(\Delta p)_{m}^{2}}{2 \rho v \lambda}
\]

\includegraphics[max width=\textwidth, alt={}, center]{da0156b6-4133-434e-90d4-ec01b2070afe-110_319_619_113_817}\\
\(\pi R^{2}\) being the effective area (area of cross section) of the ball.\\
Substitution gives 10.9 mW .\\
4.187 We have \(\frac{P}{4 \pi r^{2}}=\) intensity \(=\frac{(\Delta p)_{m}^{2}}{2 \rho v}\)\\
or

\[
(\Delta p)_{m}=\sqrt{\frac{\rho v P}{2 \pi r^{2}}}
\]

\[
\begin{gathered}
=\sqrt{\frac{1.293 \mathrm{~kg} / \mathrm{m}^{3} \times 340 \mathrm{~m} / \mathrm{s} \times 0.80 \mathrm{w}}{2 \pi \times 1.5 \times 1.5 \mathrm{~m}^{2}}}=\sqrt{\frac{1.293 \times 340 \times .8}{2 \pi \times 1.5 \times 1.5}}\left(\frac{\mathrm{~kg} \mathrm{~kg} \mathrm{~m}^{2} \mathrm{~s}^{-3} \mathrm{~m} \mathrm{~s}^{-1}}{\mathrm{~m}^{5}}\right)^{\frac{1}{2}} \\
=4.9877\left(\mathrm{~kg} \mathrm{~m}^{-1} \mathrm{~s}^{-2}\right)=5 \mathrm{~Pa} \\
\frac{(\Delta p)_{m}}{p}=5 \times 10^{-5}
\end{gathered}
\]

(b) We have

\[
\begin{aligned}
\Delta p & =-\rho v^{2} \frac{\partial \xi}{\partial x} \\
p)_{m} & =\rho v^{2} k \xi_{m}=\rho v 2 \pi v \xi_{m}
\end{aligned}
\]

\[
\begin{gathered}
\xi_{m}=a=\frac{(\Delta p)_{m}}{2 \pi \rho v v}=\frac{5}{2 \pi \times 1.293 \times 340 \times 600}=3 \mu \mathrm{~m} \\
\frac{\xi_{m}}{\lambda}=\frac{3 \times 10^{-6}}{340 / 600}=\frac{1800}{340} \times 10^{-6}=5 \times 10^{-6}
\end{gathered}
\]

4.188 Express \(L\) in bels. (i.e. \(L=5\) bels).

Then the intensity at the relevant point (at a distance \(r\) from the source) is \(: I_{0} \cdot 10^{L}\)\\
Had there been no damping the intensity would have been \(: e^{2 \gamma r} I_{0} \cdot 10^{L}\)\\
Now this must equal the quantity\\
\(\frac{P}{4 \pi r^{2}}\), where \(P=\) sonic power of the source.\\
Thus

\[
\frac{P}{4 \pi r^{2}}=e^{2 \gamma r} I_{0} \cdot 10^{L}
\]

or

\[
P=4 \pi r^{2} e^{2 \gamma r} I_{0} \cdot 10^{L}=1.39 \mathrm{~W} .
\]

\subsection*{4.4 ELECTROMAGNETIC WAVES. RADIATION}
4.189 The velocity of light in a medium of relative permittivity \(\varepsilon\) is \(\frac{c}{\sqrt{\varepsilon}}\). Thus the change in wavelength of light (from its value in vaccum to its value in the medium) is

\[
\Delta \lambda=\frac{c / \sqrt{\varepsilon}}{v}-\frac{c}{v}=\frac{c}{v}\left(\frac{1}{\sqrt{\varepsilon}}-1\right)=-50 \mathrm{~m} .
\]

4.190 From the data of the problem the relative permittivity of the medium varies as

\[
\varepsilon(x)=\varepsilon_{1} e^{-(x / l) \ln \frac{s_{1}}{\varepsilon_{2}}}
\]

Hence the local velocity of light

\[
v(x)=\frac{c}{\sqrt{\varepsilon(x)}}=\frac{c}{\sqrt{\varepsilon_{1}}} e^{\frac{x}{2 l} \ln \frac{\varepsilon_{1}}{\varepsilon_{2}}}
\]

Thus the required time \(t=\int_{0}^{l} \frac{d x}{V(x)}=\frac{\sqrt{\varepsilon_{1}}}{c} \int_{0}^{l} e^{-\frac{x}{2 l} \ln \frac{\varepsilon_{1}}{\varepsilon_{2}}}\)

\[
d x=\frac{\sqrt{\varepsilon_{1}}}{c} \frac{-e^{-\frac{1}{2} \ln \frac{\varepsilon_{1}}{\varepsilon_{2}}}+1}{\frac{1}{2 l} \ln \frac{\varepsilon_{1}}{\varepsilon_{2}}}=\frac{2 l}{c} \frac{\sqrt{\varepsilon_{1}}-\sqrt{\varepsilon_{2}}}{\ln \frac{\varepsilon_{1}}{\varepsilon_{2}}}
\]

4.191 Conduction current density \(=\sigma \vec{E}\)

Displacement current density \(=\frac{\partial \vec{D}}{\partial t}=\varepsilon \varepsilon_{0} \frac{\partial \vec{E}}{\partial t}=i \omega \varepsilon \varepsilon_{0} \vec{E}\)\\
Ratio of magnitudes \(=\frac{\sigma}{\omega \varepsilon \varepsilon_{0}}=\frac{j_{c}}{j_{\text {dis }}}=2\), on puting the values.\\
\(4.192 \vec{\nabla} \times \vec{E}=-\frac{\partial \vec{B}}{\partial t}=-\mu_{0} \frac{\partial \vec{H}}{\partial t}\)

\[
=\nabla \cos (\omega t-\vec{k} \cdot \vec{r}) \times \vec{E}_{m}=\vec{k} \times \vec{E}_{m} \sin (\omega t-\vec{k} \cdot \vec{r})
\]

At

\[
\begin{gathered}
\vec{r}=0 \\
\frac{\partial \vec{H}}{\partial t}=-\frac{\vec{k} \times \vec{E}_{m}}{\mu_{0}} \sin \omega t
\end{gathered}
\]

So integrating (ignoring a constant) and using \(c=\frac{1}{\sqrt{\varepsilon_{0} \mu_{0}}}\)

\[
\vec{H}=\frac{\vec{k} \times \vec{E}_{m}}{\mu_{0}} \cos c k t \times \frac{1}{c k}=\sqrt{\frac{\varepsilon_{0}}{\mu_{0}}} \frac{\vec{k} \times \vec{E}_{m}}{k} \cos c k t
\]

4.193 As in the previous problem

\[
\begin{aligned}
\vec{H} & =\frac{\vec{k} \times \overrightarrow{E_{m}}}{\mu_{0} \omega} \cos (\omega t-\vec{k} \cdot \vec{r})=\frac{E_{m}}{\mu_{0} c} \hat{e}_{z} \cos (k x-\omega t) \\
& =\sqrt{\frac{\varepsilon_{0}}{\mu_{0}}} E_{m} \hat{e}_{z} \cos (k x-\omega t)
\end{aligned}
\]

Thus\\
(a) at \(t=0 \vec{H}=\sqrt{\frac{\varepsilon_{0}}{\mu_{0}}} E_{m} \hat{e}_{\varepsilon} \cos k x\)\\
(b)at \(t=t_{0}, \vec{H}=\sqrt{\frac{\varepsilon_{0}}{\mu_{0}}} E_{m} \hat{e}_{z} \cos \left(k x-\omega t_{0}\right)\)\\
\(4.194 \xi_{i n d}=\oint \vec{E} \cdot d \vec{l}=E_{m} l(\cos \omega t-\cos (\omega t-k l)\)

\[
=-2 E_{m} l \sin \frac{\omega l}{2 c} \sin \left(\omega t-\frac{\omega l}{2 c}\right)
\]

Putting the values \(E_{m}=50 \mathrm{~m} V / m, l=\frac{1}{2}\) metre

\[
\begin{aligned}
& \frac{\omega l}{c}=\frac{2 \pi v l}{c}=\frac{\pi \times 10^{8}}{3 \times 10^{8}}=\frac{\pi}{3} \\
\xi_{i n d} & =50 \mathrm{~m} V\left(-\sin \frac{\pi}{6}\right) \sin \left(\omega t-\frac{\pi}{6}\right) \\
& =-25 \sin \left(\omega t+\frac{\pi}{6}-\frac{\pi}{2}\right)=25 \cos \left(\omega t-\frac{\pi}{3}\right) \mathrm{m} V
\end{aligned}
\]

\(4.195 \vec{E}=\hat{j} E(t, x)\)\\
and

\[
\vec{B}=\hat{k} B\left(t_{1} x\right)
\]

\(\operatorname{Curl} \vec{E}=\hat{k} \frac{\partial E}{\partial x}=-\frac{\partial \vec{B}}{\partial t}=-\hat{k} \frac{\partial B}{\partial t}\)\\
so

\[
-\frac{\partial E}{\partial x}=\frac{\partial B}{\partial t}
\]

Also

\[
\text { Curl } \vec{B}=\varepsilon_{0} \mu_{0} \frac{\partial \vec{E}}{\partial t}=\frac{1}{c^{2}} \frac{\partial \vec{E}}{\partial t}
\]

and

\[
\text { Curl } \vec{B}=-\hat{j} \frac{\partial B}{\partial x} \quad \text { so } \quad \frac{\partial B}{\partial x}=-\frac{1}{c^{2}} \frac{\partial E}{\partial t} .
\]

\(4.196 \vec{E}=\vec{E}_{m} \cos (\omega t-\vec{k} \cdot \vec{r})\) then as before

\[
\vec{H}=\sqrt{\frac{\varepsilon_{0}}{\mu_{0}}} \frac{\vec{k} \times \vec{E}_{m}}{k} \cos (\omega t-\vec{k} \cdot \vec{r})
\]

so

\[
\begin{aligned}
\vec{S}=\vec{E} \times \vec{H} & =\sqrt{\frac{\varepsilon_{0}}{\mu_{0}}} \vec{E}_{m} \times\left(\vec{k} \times \vec{E}_{m}\right) \frac{1}{k} \cos ^{2}(\omega t-\vec{k} \cdot \vec{r}) \\
& =\sqrt{\frac{\varepsilon_{0}}{\mu_{0}}} \vec{E}_{m}^{2} \frac{\vec{k}^{2}}{k} \cos ^{2}(\omega t-\vec{k} \cdot \vec{r}) \\
\langle\vec{S}\rangle & =\frac{1}{2} \sqrt{\frac{\varepsilon_{0}}{\mu_{0}}} \overrightarrow{E_{m}^{2}} \frac{\vec{k}}{k}
\end{aligned}
\]

4.197 \(E=E_{m} \cos (2 \pi v t-k x)\)\\
(a) \(j_{\text {dis }}=\frac{\partial D}{\partial t}=-2 \pi \varepsilon_{0} \vee E_{m} \sin (\omega t-k x)\)

Thus

\[
\begin{aligned}
\left(j_{d i s}\right)_{r m s} & =\left\langle j_{d i s}^{2}\right\rangle^{1 / 2} \\
& =\sqrt{2} \pi \varepsilon_{0} v E_{m}=0.20 \mathrm{~mA} / \mathrm{m}^{2}
\end{aligned}
\]

(b) \(\left\langle S_{x}\right\rangle=\frac{1}{2} \sqrt{\frac{\varepsilon_{0}}{\mu_{0}}} E_{m}^{2}\) as in (196). Thus \(\left\langle S_{x}\right\rangle=3.3 \mu \mathrm{~W} / \mathrm{m}^{2}\)\\
4.198 For the Poynting vector we can derive as in (196)\\
\(<S>=\frac{1}{2} \sqrt{\frac{\varepsilon \varepsilon_{0}}{\mu_{0}}} E_{m}^{2}\) along the direction of propagation.\\
Hence in time \(t\) (which is much longer than the time period \(T\) of the wave), the energy reaching the ball is

\[
\pi R^{2} \times \frac{1}{2} \sqrt{\frac{\varepsilon \varepsilon_{0}}{\mu_{0}}} E_{m}^{2} \times t=5 \mathrm{~kJ}
\]

4.199 Here \(\vec{E}=\vec{E}_{m} \cos k x \cos \omega t\)

From \(\operatorname{div} \vec{E}=0\) we get \(E_{m x}=0\) so \(\vec{E}_{m}\) is in the \(y-z\) plane.\\
Also\\
so

\[
\begin{aligned}
\frac{\partial \vec{B}}{\partial t} & =-\vec{\nabla} \times \vec{E}=-\nabla \cos k x \times \vec{E}_{m} \cos \omega t \\
& =\vec{k} \times \vec{E}_{m} \sin k x \cos \omega t
\end{aligned}
\]

\[
\begin{gathered}
=k \times E_{m} \sin k x \cos \omega t \\
\vec{B}=\frac{\vec{k} \times \vec{E}_{m}}{\omega} \sin k x \sin \omega t=\vec{B}_{m} \sin k x \sin \omega t
\end{gathered}
\]

Where

\[
\left|\vec{B}_{m}\right|=\frac{E_{m}}{c} \text { and } \vec{B}_{m} \perp \vec{E}_{m} \text { in the } y-z \text { plane. }
\]

At

\[
t=0, \vec{B}=0, E=E_{m} \cos k x
\]

At

\[
t=T / 4 \quad \vec{E}=0, B=B_{m} \sin k x
\]

106\\
\(4.200 \vec{E}=\vec{E}_{m} \cos k x \omega t\)\\
\(\vec{H}=\frac{\vec{k} \times \overrightarrow{E_{m}}}{\mu_{0} \omega} \sin k x \sin \omega t\) (exactly as in 199)

Thus

\[
\vec{S}=\vec{E} \times \vec{H}=\frac{\vec{E}_{m} \times\left(\vec{k} \times \vec{E}_{m}\right)}{\mu_{0} \omega} \frac{1}{4} \sin 2 k x \sin 2 \omega t
\]

\[
\begin{gathered}
S_{x}=\frac{1}{4} \varepsilon_{0} c E_{m}^{2} \sin 2 k x \sin 2 \omega t\left(\text { as } \frac{1}{\mu_{0} c}=\varepsilon_{0} c\right) \\
<S_{x}>=0
\end{gathered}
\]

4.201 Inside the condenser the peak electrical energy \(W_{e}=\frac{1}{2} C V_{m}^{2}\)

\[
=\frac{1}{2} V_{m}^{2} \frac{\varepsilon_{0} \pi R^{2}}{d}
\]

( \(d=\) separation between the plates, \(\pi R^{2}=\) area of each plate.).\\
\(V=V_{m} \sin \omega t, V_{m}\) is the maximum voltage\\
Changing electric field causes a displacement current

\[
\begin{aligned}
j_{d u s}=\frac{\partial D}{\partial t} & =\varepsilon_{0} E_{m} \omega \cos \omega t \\
& =\frac{\varepsilon_{0} \omega V_{m}}{d} \cos \omega t
\end{aligned}
\]

This gives rise to a magnetic field \(B(r)\) ( at a radial distance \(r\) from the centre of the plate)

\[
\begin{aligned}
B(r) \cdot 2 \pi r & =\mu_{0} \pi r^{2} j_{d i s}=\mu_{0} \pi r^{2} \frac{\varepsilon_{0} \omega V_{m}}{d} \cos \omega t \\
B & =\frac{1}{2} \varepsilon_{0} \mu_{0} \omega \frac{r}{d} V_{m} \cos \omega t
\end{aligned}
\]

Energy associated with this field is

\[
\begin{aligned}
=\int d^{3} r \frac{B^{2}}{2 \mu_{0}} & =\frac{1}{8} \varepsilon_{0}^{2} \mu_{0} \frac{\omega^{2}}{d^{2}} 2 \pi \int_{0}^{R} r^{2} r d r \times d \times V_{m}^{2} \cos ^{2} \omega t \\
& =\frac{1}{16} \pi \varepsilon_{0}^{2} \mu_{0} \frac{\omega^{2} R^{4}}{d} V_{m}^{2} \cos ^{2} \omega t
\end{aligned}
\]

Thus the maximum magnetic energy

Hence

\[
\begin{gathered}
W_{m}=\frac{\varepsilon_{0}^{2} \mu_{0}}{16}(\omega R)^{2} \frac{\pi R^{2}}{d} V_{m}^{2} \\
\frac{W_{m}}{W_{e}}=\frac{1}{8} \varepsilon_{0} \mu_{0}(\omega R)^{2}=\frac{1}{8}\left(\frac{\omega R}{c}\right)^{2}=5 \times 10^{-15}
\end{gathered}
\]

The approximation are valid only if \(\omega R \ll c\).\\
4.202 Here \(I=I_{m} \cos \omega t\), then the peak magnetic energy is

\[
W_{m}=\frac{1}{2} L I_{m}^{2}=\frac{1}{2} \mu_{0} n^{2} I_{m}^{2} \pi R^{2} d
\]

Changing magnetic field induces an electric field which by Faraday's law is given by

\[
\begin{gathered}
E \cdot 2 \pi r=-\frac{d}{d t} \int \vec{B} \cdot d \vec{S}=\pi r^{2} \mu_{0} n I_{m} \omega \sin \omega t \\
E=\frac{1}{2} r \mu_{0} n I_{m} \omega \sin \omega t
\end{gathered}
\]

The associated peak electric energy is

Hence

\[
\begin{gathered}
W_{e}=\int \frac{1}{2} \varepsilon_{0} E^{2} d^{3} r=\frac{1}{8} \varepsilon_{0} \mu_{0}^{2} n^{2} I_{m}^{2} \omega^{2} \sin ^{2} \omega t \times \frac{\pi R^{4}}{2} d \\
\frac{W_{e}}{W_{m}}=\frac{1}{8} \varepsilon_{0} \mu_{0}(\omega R)^{2}=\frac{1}{8}\left(\frac{\omega R}{c}\right)^{2}
\end{gathered}
\]

Again we expect the results to be valid if and only if

\[
\left(\frac{\omega R}{c}\right) \ll 1
\]

4.203 If the charge on the capacitor is \(Q\), the rate of increase of the capacitor's energy

\[
=\frac{d}{d t}\left(\frac{1}{2} \frac{Q^{2}}{C}\right)=\frac{Q \dot{Q}}{C}=\frac{d}{\varepsilon_{0} \pi R^{2}} Q \dot{Q}
\]

Now electric field between the plates (inside it) is, \(E=\frac{Q}{\pi R^{2} \varepsilon_{0}}\).\\
So displacement current \(=\frac{\partial D}{\partial t}=\frac{\dot{Q}}{\pi R^{2}}\)\\
This will lead to a magnetic field, (circuital) inside the plates. At a radial distance \(r\)

\[
2 \pi r H_{\theta}(r)=\pi r^{2} \frac{\dot{Q}}{\pi R^{2}} \quad \text { or } \quad H_{\theta}=\frac{\dot{Q} r}{2 \pi R^{2}}
\]

Hence \(H_{\theta}(R)=\frac{Q}{2 \pi R}\) at the edge.\\
Thus inward Poynting vector \(=S=\frac{Q}{2 \pi R} \times \frac{Q}{\pi R^{2} \varepsilon_{0}}\)\\
Total flow \(=2 \pi R d \times S=\frac{Q \dot{Q} d}{\pi R^{2} \varepsilon_{0}}\) Proved\\
4.204 Suppose the radius of the conductor is \(R_{0}\). Then the conduction current density is

\[
j_{c}=\frac{I}{\pi R_{0}^{2}}=\sigma E \quad \text { or } \quad E=\frac{I}{\pi R_{0}^{2} \sigma}=\frac{\rho I}{\pi R_{0}^{2}}
\]

where \(\rho=\frac{1}{\sigma}\) is the resistivity.

Inside the conductor there is a magnetic field given by

\[
H \cdot 2 \pi R_{0}=I \quad \text { or } \quad H=\frac{I}{2 \pi R_{0}} \text { at the edge }
\]

∴ Energy flowing in per second in a section of length \(l\) is

\[
E H \times 2 \pi R_{0} l=\frac{\rho I^{2} l}{\pi R_{0}^{2}}
\]

But the resistance \(R=\frac{\rho l}{\pi R_{0}^{2}}\)\\
Thus the energy flowing into the conductor \(=I^{2} R\).\\
4.205 Here \(n\) ev \(=I / \pi R^{2}\)\\
where \(R=\) radius of cross section of the conductor and \(n=\) charge density (per unit volume)\\
Also

\[
\frac{1}{2} m v^{2}=e U \quad \text { or } \quad v=\sqrt{\frac{2 e U}{m}}
\]

Thus, the moving protons have a charge per unit length

\[
=n e \pi R^{2}=I \sqrt{\frac{m}{2 e U}} .
\]

This gives rise to an electric field at a distance \(r\) given by

\[
E=\frac{I}{\varepsilon_{0}} \sqrt{\frac{m}{2 e U}} / 2 \pi r
\]

The magnetic field is \(H=\frac{I}{2 \pi r}(\) for \(r>R)\)\\
Thus\\
\(S=\frac{I^{2}}{\varepsilon_{0} 4 \pi^{2} r^{2}} \sqrt{\frac{m}{2 e U}}\) radially outward from the axis\\
This is the Poynting vector.\\
4.206 Within the solenoid \(B=\mu_{0} n I\) and the rate of change of magnetic energy

\[
=\dot{W}_{m}=\frac{d}{d t}\left(\frac{1}{2} \mu_{0} n^{2} I^{2} \pi R^{2} l\right)=\mu_{0} n^{2} \pi R^{2} l I \dot{I}
\]

where \(R=\) radius of cross section of the solenoid \(l=\) length.\\
Also \(H=B / \mu_{0}=n I\) along the axis within the solenoid.\\
By Faraday's law, the induced electric field is\\
or

\[
\begin{gathered}
E_{\theta} 2 \pi r=\pi r^{2} \dot{B}=\pi r^{2} \mu_{0} n \dot{I} \\
E_{\theta}=\frac{1}{2} \mu_{0} n \dot{I} r
\end{gathered}
\]

so at the edge \(E_{\theta}(R)=\frac{1}{2} \mu_{0} n \dot{I} R\) (circuital)\\
Then

\[
S_{r}=E_{\theta} H_{z} \text { (radially inward) }
\]

and

\[
\dot{W}_{m}=\frac{1}{2} \mu_{0} n^{2} I \dot{I} R \times 2 \pi R l=\mu_{0} n^{2} \pi R^{2} l I I \text { as before. }
\]

4.207 Given \(\varphi_{2}>\varphi_{1}\)

The electric field is as shown by the dashed lines ( ⇒ → - ---\()\).\\
The magnetic field is as shown\\
( © ) emerging out of the paper.\\
\(\vec{S}=\vec{E} \times \vec{H}\) is parallel to the wires\\
\includegraphics[max width=\textwidth, alt={}, center]{da0156b6-4133-434e-90d4-ec01b2070afe-117_230_635_387_779}\\
and towards right.\\
Hence source must be on the left.\\
4.208 The electric field \((\rightarrow \rightarrow)\) and the magnetic field \((H \rightarrow)\) are as shown.

The electric field by Gauss's theorem is like

Integrating

\[
E_{r}=\frac{A}{r}
\]

so

\[
A=\frac{V}{\ln \frac{r_{2}}{r_{1}}}\left(r_{2}>r_{1}\right)
\]

Then

\[
E=\frac{V}{r \ln \frac{r_{2}}{r_{1}}}
\]

\begin{center}
\includegraphics[max width=\textwidth, alt={}]{da0156b6-4133-434e-90d4-ec01b2070afe-117_402_368_915_1091}
\end{center}

Magnetic field is

\[
H_{\theta}=\frac{I}{2 \pi r}
\]

The Poynting vector \(S\) is along the \(Z\) axis and non zero between the two wires \(\left(r_{1}<r<r_{2}\right)\). The total power flux is

\[
=\int_{r_{1}}^{r_{2}} \frac{I V}{2 \pi r^{2} \ln \frac{r_{2}}{r_{1}}} \cdot 2 \pi r d r=I V
\]

4.209 As in the previous problem

\[
E_{r}=\frac{V_{0} \cos \omega t}{r \ln \frac{r_{2}}{r_{1}}} \quad \text { and } \quad H_{\theta}=\frac{I_{0} \cos (\omega t-\varphi)}{2 \pi r}
\]

Hence time averaged power flux (along the \(z\) axis) \(=\frac{1}{2} V_{0} I_{0} \cos \varphi\)

On using

\[
<\cos \omega t \cos (\omega t-\varphi)>=\frac{1}{2} \cos \varphi .
\]

4.210 Let \(\vec{n}\) be along the z axis. Then\\
and

\[
\begin{aligned}
& S_{1 \mathrm{ln}}=E_{1 x} H_{1 y}-E_{1 y} H_{1 x} \\
& S_{2 n}=E_{2 x} H_{2 y}-E_{2 y} H_{2 x}
\end{aligned}
\]

Using the boundary condition \(E_{1 t}=E_{2 t}, H_{1 t}=H_{2 t}\) at the boundary ( \(t=x\) or \(y\) ) we see that

\[
S_{1 n}=S_{2 n} .
\]

\(4.211 P \cdot \alpha|\dot{\vec{p}}|^{2}\) when\\
if

\[
\vec{p}=\sum e_{i} \overrightarrow{r_{i}}=\sum \frac{e_{i}}{m_{i}} m_{i} \overrightarrow{r_{i}}=\frac{e}{m} \sum m_{i} \overrightarrow{r_{i}}
\]

\[
\frac{e_{i}}{m_{i}}=\frac{e}{m}=\text { fixed }
\]

But

\[
\frac{d^{2}}{d t^{2}} \sum m_{i} \overrightarrow{r_{i}}=0 \text { for a closed system }
\]

Hence

\[
P=0 .
\]

\(1.212 P=\frac{1}{4 \pi \varepsilon_{0}} \frac{2(\dot{\vec{p}})^{2}}{3 c^{3}}\)

\[
|\dot{\vec{p}}|^{2}=\left(e \omega^{2} a\right)^{2} \cos ^{2} \omega t
\]

Thus

\[
<P>=\frac{1}{4 \pi \varepsilon_{0}} \frac{2}{3 c^{3}}\left(e \omega^{2} a\right)^{2} \times \frac{1}{2}=\frac{e^{2} \omega^{4} a^{2}}{12 \pi \varepsilon_{0} c^{3}}=5.1 \times 10^{-15} \mathrm{~W} .
\]

1.213 Here

\[
\dot{\vec{p}}=\frac{e}{m} \times \text { force }=\frac{e^{2} q}{m R^{2}} \frac{1}{4 \pi \varepsilon_{0}} .
\]

Thus

\[
P=\frac{1}{\left(4 \pi \varepsilon_{0}\right)^{3}}\left(\frac{e^{2} q}{m R^{2}}\right)^{2} \frac{2}{3 c^{3}}
\]

.214 Most of the radiation occurs when the moving particle is closest to the stationary particle. In that region, we can write

\[
R^{2}=b^{2}+v^{2} t^{2}
\]

and apply the previous problem's formula

Thus

\[
\Delta W \approx \frac{1}{\left(4 \pi \varepsilon_{0}\right)^{3}} \frac{2}{3 c^{3}} \int_{-\infty}^{\infty}\left(\frac{q e^{2}}{m}\right)^{2} \frac{d t}{\left(b^{2}+v^{2} t^{2}\right)^{2}}
\]

(the integral can be taken between \(\pm \infty\) with little error.)

Now

\[
\begin{gathered}
\int_{-\infty}^{\infty} \frac{d t}{\left(b^{2}+v^{2} t^{2}\right)^{2}}=\frac{1}{v} \int_{-\infty}^{\infty} \frac{d x}{\left(b^{2}+x^{2}\right)^{2}}=\frac{\pi}{2 v b^{3}} \\
\Delta W=\frac{1}{\left(4 \pi \varepsilon_{0}\right)^{3}} \frac{\pi q^{2} e^{4}}{3 c^{3} m^{2} v b^{3}}
\end{gathered}
\]

Hence,\\
\includegraphics[max width=\textwidth, alt={}, center]{da0156b6-4133-434e-90d4-ec01b2070afe-119_363_427_429_874}

Hence energy radiated \(=\Delta W\)\\
4.215 For the semicircular path on the right

\[
\frac{m \mathrm{v}^{2}}{R}=B e \mathrm{v} \quad \text { or } \quad \mathrm{v}=\frac{B e R}{m}
\]

Thus K.E. \(=T=\frac{1}{2} m v^{2}=\frac{B^{2} e^{2} R^{2}}{2 m}\).\\
Power radiated \(=\frac{1}{4 \pi \varepsilon_{0}} \frac{2}{3 c^{3}}\left(\frac{e v^{2}}{R}\right)^{2}\)

\[
=\frac{1}{4 \pi \varepsilon_{0}} \frac{2}{3 c^{3}}\left(\frac{B^{2} e^{3} R}{m^{2}}\right)^{2} \cdot \frac{\pi R}{B e R} m=\frac{B^{3} e^{5} R^{2}}{6 \varepsilon_{0} m^{3} c^{3}}
\]

So

\[
\frac{\Delta W}{T}=\frac{B e^{3}}{3 \varepsilon_{0} c^{3} m^{2}}=2.06 \times 10^{-18}
\]

(neglecting the change in \(v\) due to radiation, correct if \(\Delta W / T \ll 1\) ).\\
\(4.216 R=\frac{m \mathrm{~V}}{e B}\).

Then

\[
\begin{aligned}
P & =\frac{1}{4 \pi \varepsilon_{0}} \frac{2}{3 c^{3}}\left(\frac{e v^{2}}{R}\right)^{2}=\frac{1}{4 \pi \varepsilon_{0}} \frac{2}{3 c^{3}}\left(\frac{e^{2} B v}{m}\right)^{2} \\
& =\frac{1}{3 \pi \varepsilon_{0} c^{3}}\left(\frac{B^{2} e^{4}}{m^{3}}\right) T
\end{aligned}
\]

This is the radiated power so

\[
\frac{d T}{d t}=-\frac{B^{2} e^{4}}{3 \pi \varepsilon_{0} m^{3} \cdot c^{3}} T
\]

Integrating, \(T=T_{0} e^{-t / \tau}\)

\[
\tau=\frac{3 \pi \varepsilon_{0} m^{3} c^{3}}{B^{2} e^{4}}
\]

\(\tau\) is \((1836)^{3} \approx 10^{10}\) times less for an electron than for a proton so electrons radiate away their energy much faster in a magnetic field.\\
4.217 \(P\) is a fixed point at a distance \(l\) from the equilibrium position of the particle. Because \(l>a\), to first order in \(\frac{a}{l}\) the distance between \(P\) and the instantaneous position of the particle is still \(l\). For the first case \(y=0\) so \(t=T / 4\)\\
The corresponding retarded time is \(t^{\prime}=\frac{T}{4}-\frac{l}{c}\)\\
Now

\[
\ddot{y}\left(t^{\prime}\right)=-\omega^{2} a \cos \omega\left(\frac{I}{4}-\frac{l}{c}\right)=-\omega^{2} a \sin \frac{\omega l}{c}
\]

For the second case \(y=a\) at \(t=0\) so at the retarded time \(t^{\prime}=-\frac{\omega l}{c}\)\\
Thus

\[
\ddot{y}\left(t^{\prime}\right)=-\omega^{2} a \cos \frac{\omega l}{c}
\]

The radiation fluxes in the two cases are proportional to \(\left(\ddot{y}\left(t^{\prime}\right)\right)^{2}\) so

\[
\frac{S_{1}}{S_{2}}=\tan ^{2} \frac{\omega l}{c}=3.06 \text { on substitution. }
\]

Note : The radiation received at \(P\) at time \(t\) depends on the acceleration of the charge at the retarded time.\\
4.218 Along the circle \(x=R \sin \omega t, y=R \cos \omega t\) where \(\omega=\frac{\mathrm{V}}{\mathrm{R}}\). If t is the parameter in \(x(t), y(t)\) and \(t^{\prime}\) is the observer time then

\[
t^{\prime}=t+\frac{l-x(t)}{c}
\]

where we have neglected the effect of the \(y\)-cordinate which is of second order. The observed cordinate are\\
\includegraphics[max width=\textwidth, alt={}, center]{da0156b6-4133-434e-90d4-ec01b2070afe-120_400_522_930_925}

\[
x^{\prime}\left(t^{\prime}\right)=x(t), y^{\prime}\left(t^{\prime}\right)=y(t)
\]

Then

\[
\frac{d y^{\prime}}{d t^{\prime}}=\frac{d y}{d t^{\prime}}=\frac{d t}{d t^{\prime}} \frac{d y}{d t}=\frac{-\omega R \sin \omega t}{1-\frac{\omega R}{c} \cos \omega t}=\frac{-\omega x}{1-\frac{\omega y}{c}}=\frac{-\mathrm{v} x / R}{1-\frac{\mathrm{v} y}{c R}}
\]

and

\[
\begin{gathered}
\frac{d^{2} y^{\prime}}{d t^{\prime 2}}=\frac{d t}{d t^{\prime}} \frac{d}{d t}\left(\frac{-\mathrm{v} x / R}{1-\frac{\mathrm{v} y}{c R}}\right) \\
=\frac{1}{1-\frac{\mathrm{v} y}{c R}}\left\{\frac{-\frac{\mathrm{v}^{2}}{\mathrm{R}^{2}} y}{1-\frac{\mathrm{v} y}{c R}}+\frac{\frac{\mathrm{v} x}{R}\left(\frac{\mathrm{v}^{2}}{\mathrm{cR}^{2}} x\right)}{\left(1-\frac{\mathrm{v} y}{c R}\right)^{2}}\right\}=\frac{\frac{v^{2}}{R}\left(\frac{\mathrm{v}}{\mathrm{c}}-\frac{y}{R}\right)}{\left(1-\frac{\mathrm{v} y}{c R}\right)^{3}} .
\end{gathered}
\]

This is the observed acceleration.\\
(b) Energy flow density of \(E M\) radiation \(S\) is proportional to the square of the \(y\)-projection of the observed acceleration of the particle \(\left(\right.\) i.e. \(\left.\frac{d^{2} y^{\prime}}{d t^{\prime 2}}\right)\).

Thus

\[
\frac{S_{1}}{S_{2}}=\left[\frac{\left(\frac{v}{c}-1\right)}{\left(1-\frac{v}{c}\right)^{3}} / \frac{\left(\frac{v}{c}+1\right)}{\left(1+\frac{v}{c}\right)^{3}}\right]^{2}=\frac{\left(1+\frac{v}{c}\right)^{4}}{\left(1-\frac{v}{c}\right)^{4}}
\]

4.219 We know that \(S_{0}(r) \propto \frac{1}{r^{2}}\)

At other angles \(S(r, \theta) \propto \sin ^{2} \theta\)\\
Thus \(S(r, \theta)=S_{0}(r) \sin ^{2} \theta=S_{0} \sin ^{2} \theta\)\\
Average power radiated\\
\includegraphics[max width=\textwidth, alt={}, center]{da0156b6-4133-434e-90d4-ec01b2070afe-121_222_587_536_860}

\[
\begin{aligned}
=S_{0} \times 4 \pi r^{2} \times \frac{2}{3}= & \frac{8 \pi}{3} S_{0} r^{2} \\
& \left(\text { Average of } \sin ^{2} \theta \text { over whole sphere is } \frac{2}{3}\right)
\end{aligned}
\]

4.220 From the previous problem.\\
or

Thus

\[
\begin{aligned}
P_{0} & =\frac{8 \pi S_{0} r^{2}}{3} \\
S_{0} & =\frac{3 P_{0}}{8 \pi r^{2}} \\
\langle w\rangle & =\frac{S_{0}}{c}=\frac{3 P_{0}}{8 \pi c r^{2}}
\end{aligned}
\]

(Poynting flux vector is the energy contained is a box of unit cross section and length \(c\) ).\\
4.221 The rotating dispole has moments

Thus

\[
\begin{gathered}
p_{x}=p \cos \omega t, p_{y}=p \sin \omega t \\
P=\frac{1}{4 \pi \varepsilon_{0}} \frac{2}{3 c^{3}} \omega^{4} p^{2}=\frac{p^{2} \omega^{4}}{6 \pi \varepsilon_{0} c^{3}}
\end{gathered}
\]

4.222 If the electric field of the wave is

\[
\vec{E}=\vec{E}_{0} \cos \omega t
\]

then this induces a dipole moment whose second derivative is

\[
\ddot{\vec{p}}=\frac{e^{2} \overrightarrow{E_{0}}}{m} \cos \omega t
\]

Hence radiated mean power \(<P>=\frac{1}{4 \pi \varepsilon_{0}} \frac{2}{3 c^{3}}\left(\frac{e^{2} E_{0}}{m}\right)^{2} \times \frac{1}{2}\)

On the other hand the mean Poynting flux of the incident radiation is

Thus

\[
\begin{aligned}
&\left\langle S_{i n c}\right\rangle=\sqrt{\frac{\varepsilon_{0}}{\mu_{0}}} \times \frac{1}{2} E_{0}^{2} \\
& \frac{P}{\left\langle S_{i n c}\right\rangle}= \frac{1}{4 \pi \varepsilon_{0}} \cdot \frac{2}{3}\left(\varepsilon_{0} \mu_{0}\right)^{3 / 2}\left(\frac{e^{2}}{m}\right)^{2} \times \sqrt{\frac{\mu_{0}}{\varepsilon_{0}}} \\
&= \frac{\mu_{0}^{2}}{6 \pi}\left(\frac{e^{2}}{m}\right)^{2}
\end{aligned}
\]

4.223 For the elastically bound electron

\[
m \ddot{\vec{r}}+m \omega_{0}^{2} \vec{r}=e \overrightarrow{E_{0}} \cos \omega t
\]

This equation has the particular integral\\
(i.e. neglecting the part which does not have the frequency of the impressed force)

\[
\vec{r}=\frac{e \vec{E}_{0}}{m} \frac{\cos \omega t}{\omega_{0}^{2}-\omega^{2}} \text { so and } \quad \ddot{p}=-\frac{e^{2} \vec{E}_{0} \omega^{2}}{\left(\omega_{0}^{2}-\omega^{2}\right) m} \cos \omega t
\]

Hence \(\boldsymbol{P}=\) mean radiated power

\[
=\frac{1}{4 \pi \varepsilon_{0}} \frac{2}{3 c^{3}}\left(\frac{e^{2} \omega^{2}}{m\left(\omega_{0}^{2}-\omega^{2}\right)}\right)^{2} \frac{1}{2} E_{0}^{2}
\]

The mean incident poynting flux is

Thus

\[
\begin{gathered}
\left\langle S_{i n c}\right\rangle=\sqrt{\frac{\varepsilon_{0}}{\mu_{0}}} \frac{1}{2} E_{0}^{2} \\
\frac{P}{\left\langle S_{i n c}\right\rangle}=\frac{\mu_{0}^{2}}{6 \pi}\left(\frac{e^{2}}{m}\right)^{2} \frac{\omega^{4}}{\left(\omega_{0}^{2}-\omega^{2}\right)^{2}} .
\end{gathered}
\]

4.224 Let \(r=\) radius of the ball\\
\(R=\) distance between the ball \& the Sun \(\quad(r \ll R)\).

Then

\[
\begin{gathered}
M=\text { mass of the Sun } \\
\gamma=\text { gravitational constant } \\
\frac{\gamma M}{R^{2}} \frac{4 \pi}{3} r^{3} \rho=\frac{P}{4 \pi R^{2}} \pi r^{2} \cdot \frac{1}{c}
\end{gathered}
\]

(the factor \(\frac{1}{c}\) converts the energy received on the right into momentum received. Then the right hand side is the momentum received per unit time and must equal the negative of the impressed force for equilibrium).

Thus

\[
r=\frac{3 P}{16 \pi \gamma M c \rho}=0.606 \mu \mathrm{~m} .
\]

\section*{PART FIVE}
\section*{OPTICS}
\subsection*{5.1 PHOTOMETRY AND GEOMETRICAL OPTICS}
5.1 (a) The relative spectral response \(V(\lambda)\) shown in Fig. (5.11) of the book is so defined that \(A / V(\lambda)\) is the energy flux of light of wave length \(\lambda\) needed to produce a unit luminous flux at that wavelength. (A is the conversion factor defined in the book.)\\
At \(\quad \lambda=0.51 \mu \mathrm{~m}\), we read from the figure

\[
V(\lambda)=0.50 \mathrm{so}
\]

energy flux corresponding to a luminous flux of 1 lumen \(=\frac{1.6}{0.50}=3.2 \mathrm{~mW}\)\\
At

\[
\begin{gathered}
\lambda=0.64 \mu \mathrm{~m}, \text { we read } \\
V(\lambda)=0.17
\end{gathered}
\]

and energy flux corresponding to a luminous flux of 1 lumen \(=\frac{1 \cdot 6}{\cdot 17}=9 \cdot 4 \mathrm{~mW}\)\\
(b) Here \(d \Phi_{e}(\lambda)=\frac{\Phi_{e}}{\lambda_{2}-\lambda_{1}} d \lambda, \quad \lambda_{1} \leq \lambda \leq \lambda_{2}\)\\
since energy is distributed uniformly. Then

\[
\Phi=\int_{\lambda_{1}}^{\lambda_{2}} V(\lambda) d \Phi_{2}(\lambda) / A=\frac{\Phi_{e}}{A\left(\lambda_{2}-\lambda_{1}\right)} \int_{\lambda_{1}}^{\lambda_{2}} V(\lambda) d \lambda
\]

since \(V(\lambda)\) is assumed to vary linearly in the interval \(\lambda_{1} \leq \lambda \leq \lambda_{2}\), we have

\[
\frac{1}{\lambda_{1}-\lambda_{2}} \int_{\lambda_{1}}^{\lambda_{2}} V(\lambda) d \lambda=\frac{1}{2}\left(V\left(\lambda_{1}\right)+V\left(\lambda_{2}\right)\right)
\]

Thus

\[
\Phi=\frac{\Phi_{e}}{2 A}\left(V\left(\lambda_{1}\right)+V\left(\lambda_{2}\right)\right)
\]

Using

\[
\begin{aligned}
& V(0.58 \mu \mathrm{~m})=0.85 \\
& V(0.63 \mu \mathrm{~m})=0.25
\end{aligned}
\]

Thus

\[
\Phi=\frac{\Phi_{e}}{2 \times 1.6} \times 1.1=1.55 \text { lumen } .
\]

5.2 We have \(\Phi_{e}=\frac{\Phi A}{V(\lambda)}\)

But

\[
\Phi_{e}=\frac{\frac{1}{2} \sqrt{\frac{\varepsilon_{0}}{\mu_{0}}}}{\substack{\downarrow \\ \text { mean energy } \\ \text { flux vector }}} \times \frac{4 \pi r^{2}}{\text { area }} \quad \text { or } \quad E_{m}^{2}=\frac{\Phi A}{2 \pi r^{2} V(\lambda)} \sqrt{\frac{\mu_{0}}{\varepsilon_{0}}}
\]

For

Also

\[
\begin{gathered}
\lambda=0.59 \mu \mathrm{~m} V(\lambda)=0.74 \text { Thus } \\
E_{m}=1.14 \mathrm{~V} / \mathrm{m} \\
H_{m}=\sqrt{\frac{\varepsilon_{0}}{\mu_{0}}} E_{m}=3.02 \mathrm{~mA} / \mathrm{m}
\end{gathered}
\]

5.3 (a) Mean illuminance

\[
=\frac{\text { Total luminous flux incident }}{\text { Total area illuminated }} .
\]

Now, to calculate the total luminous flux incident on the sphere, we note that the illuminance at the point of normal incidence is \(E_{0}\). Thus the incident flux is \(E_{o} \cdot \pi R^{2}\). Thus

\[
\begin{aligned}
\text { Mean illuminance } & =\frac{\pi R^{2} \cdot E_{0}}{2 \pi R^{2}} \\
\text { or } \quad<E> & =\frac{1}{2} E_{0}
\end{aligned}
\]

(b) The sphere subtends a solid angle\\
\includegraphics[max width=\textwidth, alt={}, center]{da0156b6-4133-434e-90d4-ec01b2070afe-124_336_654_932_795}

\[
2 \pi(1-\cos \alpha)=2 \pi\left(1-\frac{\sqrt{l^{2}-R^{2}}}{l}\right)
\]

at the point source and therefore receives a total flux of

The area irradiated is : \(\quad 2 \pi R^{2} \int_{0} \sin \theta d \theta=2 \pi R^{2}(1-\sin \alpha)=2 \pi R^{2}\left(1-\frac{R}{l}\right)\)

Thus

\[
\langle E\rangle=\frac{I}{R^{2}} \frac{1-\sqrt{1-(R / l)^{2}}}{1-\frac{R}{l}}
\]

Substituting we get \(<E>=50\) lux .\\
5.4 Luminance \(L\) is the light energy emitted per unit area of the emitting surface in a given direction per unit solid angle divided by \(\cos \theta\). Luminosity \(M\) is simply energy emitted per unit area.

Thus

\[
M=\int L \cdot \cos \theta \cdot d \Omega
\]

where the integration must be in the forward hemisphere of the emitting surface (assuming light is being emitted in only one direction say outward direction of the surface.) But

\[
\begin{gathered}
L=L_{0} \cos \theta \\
\pi / 2
\end{gathered}
\]

Thus

\[
M=\int L_{0} \cos ^{2} \theta \cdot d \Omega=2 \pi \int_{0}^{\pi / 2} L_{0} \cos ^{2} \theta \sin \theta d \theta=\frac{2}{3} \pi L_{0}
\]

5.5 (a) For a Lambert source \(L=\) Const

The flux emitted into the cone is

\[
\Phi=L \Delta S \cos \alpha d \Omega
\]

\includegraphics[max width=\textwidth, alt={}, center]{da0156b6-4133-434e-90d4-ec01b2070afe-125_567_652_737_375}\\
(b) The luminosity is obtained from the previous formula for \(\theta=90^{\circ}\)

\[
M=\frac{\Phi\left(\theta=90^{\circ}\right)}{\Delta S}=\pi L
\]

5.6 The equivalent luminous intensity in the direction \(O P\) is

\[
L S \cos \theta
\]

and the illuminance at \(P\) is

\[
\begin{aligned}
& \frac{L S \cos \theta}{\left(R^{2}+h^{2}\right)} \cos \theta=\frac{L S h^{2}}{\left(R^{2}+h^{2}\right)^{2}} \\
= & \frac{L S}{\left(\frac{R^{2}}{h}+h\right)^{2}}=\frac{L S}{\left[\left(\frac{R}{\sqrt{h}}-\sqrt{h}\right)^{2}+2 R\right]^{2}}
\end{aligned}
\]

\begin{center}
\includegraphics[max width=\textwidth, alt={}]{da0156b6-4133-434e-90d4-ec01b2070afe-125_494_436_1531_697}
\end{center}

This is maximum when

\[
R=h
\]

and the maximum illuminance is

\[
\frac{L S}{4 R^{2}}=\frac{1.6 \times 10^{2}}{4}=40 \operatorname{lux}
\]

5.7 The illiminance at \(P\) is

\[
E_{p}=\frac{I(\theta)}{\left(x^{2}+h^{2}\right)} \cos \theta=\frac{I(\theta) \cos ^{3} \theta}{h^{2}}
\]

since this is constant at all \(x\), we must have\\
\(I(\theta) \cos ^{3} \theta=\) const \(=I_{0}\)\\
\(\operatorname{or} I(\theta)=I_{0} / \cos ^{3} \theta\)\\
The luminous flux reaching the table is\\
\(\Phi=\pi R^{2} \times \frac{I_{0}}{h^{2}}=314\) lumen\\
\includegraphics[max width=\textwidth, alt={}, center]{da0156b6-4133-434e-90d4-ec01b2070afe-126_402_378_375_1012}\\
5.8 The illuminated area acts as a Lambert source of luminosity \(M=\pi L\) where\\
\(M S=\rho E S=\) total reflected light\\
Thus, the luminance

\[
L=\frac{\rho E}{\pi}
\]

The equivalent luminous intensity in the direction making an angle \(\theta\) from the vertical is

\[
L S \cos \theta=\frac{\rho E S}{\pi} \cos \theta
\]

and the illuminance at the point \(P\) is

\[
\frac{\rho E S}{\pi} \cos \theta \sin \theta / R^{2} \operatorname{cosec}^{2} \theta=\frac{\rho E S}{\pi R^{2}} \cos \theta \sin ^{3} \theta
\]

This is maximum when\\
\includegraphics[max width=\textwidth, alt={}, center]{da0156b6-4133-434e-90d4-ec01b2070afe-126_400_512_1034_941}

\[
\frac{d}{d \theta}\left(\cos \theta \sin ^{3} \theta\right)=-\sin ^{4} \theta+3 \sin ^{2} \theta \cos ^{2} \theta=0
\]

or

\[
\tan ^{2} \theta=3 \Rightarrow \tan \theta=\sqrt{3}
\]

Then the maximum illuminance is

\[
\frac{3 \sqrt{3}}{16 \pi} \frac{\rho E S}{R^{2}}
\]

This illuminance is obtained at a distance \(R \cot \theta=R / \sqrt{3}\) from the ceiling. Substitution gives the value\\
5.9 From the definition of luminance, the energy emitted in the radial direction by an element \(d S\) of the surface of the dome is

\[
d \Phi=L d S d \Omega
\]

Here \(L=\) constant. The solid angle \(d \Omega\) is given by

\[
d \Omega=\frac{d A \cos \theta}{R^{2}}
\]

\includegraphics[max width=\textwidth, alt={}, center]{da0156b6-4133-434e-90d4-ec01b2070afe-127_312_448_87_943}\\
where \(d A\) is the area of an element on the plane illuminated by the radial light. Then

\[
d \Phi=\frac{L d S d A}{R^{2}} \cos \theta
\]

The illuminance at 0 is then

\[
E=\int \frac{d \Phi}{d A}=\int_{0}^{\pi / 2} \frac{L}{R^{2}} 2 \pi R^{2} \sin \theta d \theta \cos \theta=2 \pi L \int_{0}^{l} x d x=\pi L
\]

5.10 Consider an element of area \(d S\) at point \(P\).

It emits light of flux

\[
\begin{aligned}
d \Phi & =L d S d \Omega \cos \theta \\
& =L d S \frac{d A}{h^{2} \sec ^{2} \theta} \cdot \cos ^{2} \theta \\
& =\frac{L d S d A}{h^{2}} \cos ^{4} \theta
\end{aligned}
\]

\includegraphics[max width=\textwidth, alt={}, center]{da0156b6-4133-434e-90d4-ec01b2070afe-127_373_626_785_759}\\
in the direction of the surface element \(d A\) at \(O\).\\
The total illuminance at \(O\) is then

\[
E=\int \frac{L d S}{h^{2}} \cos ^{4} \theta
\]

But

\[
\begin{aligned}
d S & =2 \pi r d r=2 \pi h \tan \theta d(h \tan \theta) \\
& =2 \pi h^{2} \sec ^{2} \theta \tan \theta d \theta
\end{aligned}
\]

5.11 Consider an angular element of area

\[
2 \pi x d x=2 \pi h^{2} \tan \theta \sec ^{2} \theta d \theta
\]

Light emitted from this ring is

\[
d \Phi=L d \Omega\left(2 \pi h^{2} \tan \theta \sec ^{2} \theta d \theta\right) \cdot \cos \theta
\]

\begin{center}
\includegraphics[max width=\textwidth, alt={}]{da0156b6-4133-434e-90d4-ec01b2070afe-127_295_388_1645_918}
\end{center}

Now

\[
d \Omega=\frac{d A \cos \theta}{h^{2} \sec ^{2} \theta}
\]

where \(d A=\) an element of area of the table just below the untre of the illuminant. Then the illuminance at the element \(\boldsymbol{d} \boldsymbol{A}\) will be

\[
E_{0}=\int_{\theta=0}^{\theta-\alpha} 2 \pi L \sin \theta \cos \theta d \theta
\]

where

\[
\sin \alpha=\frac{R}{\sqrt{h^{2}+R^{2}}} \text {. Finally using luminosity } M=\pi L
\]

\[
E_{0}=M \sin ^{2} \alpha=M \frac{R^{2}}{h^{2}+R^{2}}
\]

or

\[
M=E_{0}\left(1+\frac{h^{2}}{R^{2}}\right)=700 \mathrm{~lm} / \mathrm{m}^{2} *\left(1 \mathrm{x}=1 \frac{1 \mathrm{~m}}{\mathrm{~m}^{2}} \text { dimensionally }\right)
\]

5.12 See the figure below. The light emitted by an element of the illuminant towards the point \(O\) under consideration is

\[
d \Phi=L d S d \Omega \cos (\alpha+\beta)
\]

The element \(d S\) has the area

\[
d S=2 \pi R^{2} \sin \alpha d \alpha
\]

The distance

\[
O A=\left[h^{2}+R^{2}-2 h R \cos \alpha\right]^{1 / 2}
\]

we also have

\[
\frac{O A}{\sin \alpha}=\frac{h}{\sin (\alpha+\beta)}=\frac{R}{\sin \beta}
\]

From the diagram

\[
\begin{aligned}
\cos (\alpha+\beta) & =\frac{h \cos \alpha-R}{O A} \\
\cos \beta & =\frac{h-R \cos \alpha}{O A}
\end{aligned}
\]

If we imagine a small area \(d \Sigma\) at \(O\) then

\[
\frac{d \Sigma \cos \beta}{O A^{2}}=d \Omega
\]

\begin{center}
\includegraphics[max width=\textwidth, alt={}]{da0156b6-4133-434e-90d4-ec01b2070afe-128_584_436_1111_959}
\end{center}

Hence, the illuminance at \(O\) is

\[
\int \frac{d \Phi}{d \Sigma}=\int L 2 \pi R^{2} \sin \alpha d \alpha \frac{(h \cos \alpha-R)(h-R \cos \alpha)}{(O A)^{4}}
\]

The limit of \(\alpha\) is \(\alpha=0\) to that value for which \(\alpha+\beta=90^{\circ}\), for then light is emitted tangentially. Thus

Thus

\[
E=\int_{0} L \cdot 2 \pi R^{2} \sin \alpha d \alpha \frac{(h-R \cos \alpha)(h \cos \alpha-R)}{\left(h^{2}+R^{2}-2 h R \cos \alpha\right)^{2}}
\]

we put

\[
y=h^{2}+R^{2}-2 h R \cos \alpha
\]

So,

\[
\begin{gathered}
d y=2 h R \sin \alpha d \alpha \\
E=\int_{(h-R)^{2}}^{h^{2}-R^{2}} L \cdot 2 \pi R^{2} \frac{d y}{2 h R} \frac{\left(h-\frac{h^{2}+R^{2}-y}{2 h}\right)\left(\frac{h^{2}+R^{2}-y}{2 R}-R\right)}{y^{2}} \\
=\frac{L \cdot 2 \pi R^{2}}{8 h^{2} R^{2}} \int_{(h-R)^{2}}^{h^{2}-R^{2}} \frac{\left(h^{2}-R^{2}+y\right)\left(h^{2}-R^{2}-y\right)}{y^{2}} d y \\
=\frac{\pi L}{4 h^{2}} \int_{(h-R)^{2}}^{h^{2}-R^{2}}\left[\frac{\left(h^{2}-R^{2}\right)^{2}}{y^{2}}-1\right] d y=\frac{\pi L}{4 h^{2}}\left[-\frac{\left(h^{2}-R^{2}\right)^{2}}{y}-y\right]_{(h-R)^{2}}^{h^{2}-R^{2}} \\
=\frac{\pi L}{4 h^{2}}\left[(h+R)^{2}-\left(h^{2}-R^{2}\right)-\left(h^{2}-R^{2}\right)+(h-R)^{2}\right] \\
=\frac{\pi L}{4 h^{2}}\left[2 h^{2}+2 R^{2}-2 h^{2}+2 R^{2}\right]=\frac{\pi L R^{2}}{h^{2}}
\end{gathered}
\]

Substitution gives :

\[
E=25 \cdot 1 \text { lux }
\]

5.13 We see from the diagram that because of the law of reflection, the component of the incident unit vector \(\vec{e}\) along \(\cdot \vec{n}\) changes sign on reflection while the component \(\|\) to the mirror remains unchanged. Writing \(\vec{e}=\overrightarrow{e_{\|}}+\overrightarrow{e_{1}}\) where \(\overrightarrow{e_{\perp}}=\vec{n}(\overrightarrow{e \cdot n})\)

\[
\overrightarrow{e_{\|}}=\overrightarrow{e^{\prime}}-\vec{n}(\overrightarrow{e \cdot n})
\]

we see that the reflected unit vector is\\
\includegraphics[max width=\textwidth, alt={}, center]{da0156b6-4133-434e-90d4-ec01b2070afe-129_335_521_1543_844}

\[
\overrightarrow{e^{\prime}}=\overrightarrow{e_{\|}}-\overrightarrow{e_{\perp}}=\overrightarrow{e^{\prime}}-2 \vec{n}\left(\overrightarrow{e_{\cdot} \cdot n}\right)
\]

5.14 We choose the unit vectors perpendicular to the mirror as the \(x, y, z\) axes in space. Then after reflection from the mirror with normal along \(x\) axis

\[
\overrightarrow{e^{\prime}}=\vec{e}-2 \hat{i}(\hat{i} \cdot \vec{e})=-\overrightarrow{e_{x}} \hat{i}+e_{y} \hat{j}+e_{z} \hat{k}
\]

where \(\hat{i}, \hat{j}, \hat{k}\) are the basic unit vectors. After a second reflection from the 2 nd mirror say along \(y\) axis.

\[
\vec{e}{ }^{\prime \prime}=\overrightarrow{e^{\prime}}-2 \hat{j}\left(\hat{j} \cdot \overrightarrow{e^{\prime}}\right)=-e_{x} \hat{i}-e_{y} \hat{h}+e_{z} \hat{k}
\]

Finally after the third reflection

\[
\vec{e}^{\prime \prime \prime}=-e_{x} \hat{i}-e_{y} \hat{j}-e_{z} \hat{k}=-\vec{e} .
\]

5.15 Let \(P Q\) be the surface of water and \(n\) be the R.I. of water. Let \(A O\) is the shaft of light with incident angle \(\theta_{1}\) and \(O B\) and \(O C\) are the reflected and refracted light rays at angles \(\theta_{1}\) and \(\theta_{2}\) respectively (Fig.). From the figure \(\theta_{2}=\frac{\pi}{2}-\theta_{1}\)

From the law of refraction at the interface \(P Q\)\\
\(n=\frac{\sin \theta_{1}}{\sin \theta_{2}}=\frac{\sin \theta_{1}}{\sin \left(\frac{\pi}{2}-\theta_{1}\right)}\)\\
or, \(n=\frac{\sin \theta_{1}}{\cos \theta_{1}}=\tan \theta_{1}\)\\
\includegraphics[max width=\textwidth, alt={}, center]{da0156b6-4133-434e-90d4-ec01b2070afe-130_385_457_648_997}

Hence \(\theta_{1}=\tan ^{-1} n\)\\
5.16 Let two optical mediums of R.I. \(n_{1}\) and \(n_{2}\) respectively be such that \(n_{1}>n_{2}\). In the case when angle of incidence is \(\theta_{1 c r}\) (Fig.), from the law of refraction


\begin{equation*}
n_{1} \sin \theta_{1 c r}=n_{2} \tag{1}
\end{equation*}


In the case, when the angle of incidence is \(\theta_{1}\), from the law of refraction at the interface of mediums 1 and 2 .

\[
n_{1} \sin \theta_{1}=n_{2} \sin \theta_{2}
\]

But in accordance with the problem \(\theta_{2}=\left(\pi / 2-\theta_{1}\right)\)\\
\includegraphics[max width=\textwidth, alt={}, center]{da0156b6-4133-434e-90d4-ec01b2070afe-130_389_455_1140_1008}\\
so,


\begin{equation*}
n_{1} \sin \theta_{1}=n_{2} \cos \theta_{1} \tag{2}
\end{equation*}


Dividing Eqn (1) by (2)\\
or,


\begin{gather*}
\frac{\sin \theta_{1 c r}}{\sin \theta_{1}}=\frac{1}{\cos \theta_{1}} \\
\eta=\frac{1}{\cos \theta_{1}}, \quad \text { so } \quad \cos \theta_{1}=\frac{1}{\eta} \text { and } \sin \theta_{1}=\frac{\sqrt{\eta^{2}-1}}{\eta} \\
\frac{n_{1}}{n_{2}}=\frac{\cos \theta_{1}}{\sin \theta_{1}} \tag{3}
\end{gather*}


But

So,

Thus

\[
\frac{n_{1}}{n_{2}}=\frac{1}{\eta} \frac{\eta}{\sqrt{\eta^{2}-1}}
\]

(Using 3)

But from the law of refraction

\[
\frac{n_{1}}{n_{2}}=\frac{1}{\sqrt{\eta^{2}-1}}
\]

\subsection*{5.17 From the Fig. the sought lateral shift \\
 5.17 From the Fig. the sought lateral shift}

\begin{align*}
& x=O M \sin (\theta-\beta) \\
& =d \sec \beta \sin (\theta-\beta) \\
& =d \sec \beta(\sin \theta \cos \beta-\cos \theta \sin \beta) \\
& =d(\sin \theta-\cos \theta \tan \beta)
\end{align*}


\(\sin \theta=n \sin \beta \quad\) or, \(\sin \beta=\frac{\sin \theta}{n}\)\\
\includegraphics[max width=\textwidth, alt={}, center]{da0156b6-4133-434e-90d4-ec01b2070afe-131_461_657_328_703}

So, \(\quad \cos \beta=\frac{\sqrt{n^{2}-\sin ^{2} \theta}}{n}\) and \(\tan \beta=\frac{\sin \theta}{\sqrt{n^{2}-\sin ^{2} \theta}}\)\\
Thus

\[
\begin{aligned}
x & =d(\sin \theta-\cos \theta \tan \beta)=d\left(\sin \theta-\cos \theta \frac{\sin \theta}{\sqrt{n^{2}-\sin ^{2} \theta}}\right) \\
& =d \sin \theta\left[1-\sqrt{\frac{1-\sin ^{2} \theta}{n^{2}-\sin ^{2} \theta}}\right]
\end{aligned}
\]

5.18 From the Fig.

\[
\sin d \alpha=\frac{M P}{O M}=\frac{M N \cos \alpha}{h \sec (\alpha+d \alpha)}
\]

As \(d \alpha\) is very small, so


\begin{equation*}
d \alpha \approx \frac{M N \cos \alpha}{h \sec \alpha}=\frac{M N \cos ^{2} \alpha}{h} \tag{1}
\end{equation*}


Similarly


\begin{equation*}
d \theta=\frac{M N \cos ^{2} \theta}{h^{\prime}} \tag{2}
\end{equation*}


\begin{center}
\includegraphics[max width=\textwidth, alt={}]{da0156b6-4133-434e-90d4-ec01b2070afe-131_424_489_1210_746}
\end{center}

From Eqns (1) and (2)


\begin{equation*}
\frac{d \alpha}{d \theta}=\frac{h^{\prime} \cos ^{2} \alpha}{h \cos ^{2} \theta} \quad \text { or }, \quad h^{\prime}=\frac{h \cos ^{2} \theta}{\cos ^{2} \alpha} \frac{d \alpha}{d \theta} \tag{3}
\end{equation*}


From the law of refraction


\begin{equation*}
n \sin \alpha=\sin \theta \tag{A}
\end{equation*}



\begin{equation*}
\sin \alpha=\frac{\sin \theta}{n}, \text { so, } \cos \alpha=\sqrt{\frac{n^{2}-\sin ^{2} \theta}{n^{2}}} \tag{B}
\end{equation*}


Differentiating Eqn.(A)


\begin{equation*}
n \cos \alpha d \alpha=\cos \theta d \theta \quad \text { or } \quad \frac{d \alpha}{d \theta}=\frac{\cos \theta}{n \cos \alpha} \tag{4}
\end{equation*}


Using (4) in (3), we get


\begin{equation*}
h^{\prime}=\frac{h \cos ^{3} \theta}{n \cos ^{3} \alpha} \tag{5}
\end{equation*}


Hence \(\quad h^{\prime}=\frac{h \cos ^{3} \theta}{n\left(\frac{n^{2}-\sin ^{2} \theta}{n^{2}}\right)^{3 / 2}}=\frac{n^{2} h^{\prime} \cos ^{3} \theta}{\left(n^{2}-\sin ^{2} \theta\right)^{3 / 2}}\) [Using Eqn.(B)]\\
5.19 The figure shows the passage of a monochromatic ray through the given prism, placed in air medium.\\
From the figure, we have


\begin{equation*}
\theta=\beta_{1}+\beta_{2} \tag{A}
\end{equation*}


and


\begin{align*}
& \alpha=\left(\alpha_{1}+\alpha_{2}\right)-\left(\beta_{1}+\beta_{2}\right) \\
& \alpha=\left(\alpha_{1}+\alpha_{2}\right)-\theta \tag{1}
\end{align*}


From the Snell's law

\[
\sin \alpha_{1}=n \sin \beta_{1}
\]

\includegraphics[max width=\textwidth, alt={}, center]{da0156b6-4133-434e-90d4-ec01b2070afe-132_365_482_761_818}\\
or


\begin{equation*}
\alpha_{1}=n \beta_{1} \text { (for small angles) } \tag{2}
\end{equation*}


and

\[
\sin \alpha_{2}=n \sin \beta_{2}
\]

or,


\begin{equation*}
\alpha_{2}=n \beta_{2} \text { (for small angles) } \tag{3}
\end{equation*}


From Eqns (1), (2) and (3), we get

\[
\alpha=n\left(\beta_{1}+\beta_{2}\right)-\theta
\]

So,

\[
\alpha=n(\theta)-\theta=(n-1) \theta \text { [Using Eqn.A] }
\]

5.20 (a) In the general case, for the passage of a monochromatic ray through a prism as shown in the figure of the soln. of 5.19,


\begin{equation*}
\alpha=\left(\alpha_{1}+\alpha_{2}\right)-\theta \tag{1}
\end{equation*}


And from the Snell's law,


\begin{equation*}
\sin \alpha_{1}=n \sin \beta_{1} \quad \text { or } \quad \alpha_{1}=\sin ^{-1}\left(n \sin \beta_{1}\right) \tag{2}
\end{equation*}


Similarly \(\alpha_{2}=\sin ^{-1}\left(n \sin \beta_{2}\right)=\sin ^{-1}\left[n \sin \left(\theta-\beta_{1}\right)\right]\) (As \(\theta=\beta_{1}+\beta_{2}\) )\\
Using (2) in (1)

\[
\alpha=\left[\sin ^{-1}\left(n \sin \beta_{1}\right)+\sin ^{-1}\left(n \sin \left(\theta-\beta_{1}\right)\right)\right]-\theta
\]

For \(\alpha\) to be minimum, \(\frac{d \alpha}{d \beta_{1}}=0\)\\
or,

\[
\frac{n \cos \beta_{1}}{\sqrt{1-n^{2} \sin ^{2} \beta_{1}}}-\frac{n \cos \left(\theta-\beta_{1}\right)}{\sqrt{1-n^{2} \sin ^{2}\left(\theta-\beta_{1}\right)}}=0
\]

or,

\[
\frac{\cos ^{2} \beta_{1}}{\left(1-n^{2} \sin ^{2} \beta_{1}\right)}=\frac{\cos ^{2}\left(\theta-\beta_{1}\right)}{1-n^{2} \sin ^{2}\left(\theta-\beta_{1}\right)}
\]

or,

\[
\cos ^{2} \beta_{1}\left(1-n^{2} \sin ^{2}\left(\theta-\beta_{1}\right)\right)=\cos ^{2}\left(\theta-\beta_{1}\right)\left(1-n^{2} \sin ^{2} \beta_{1}\right)
\]

or, \(\left(1-\sin ^{2} \beta_{1}\right)\left(1-n^{2} \sin ^{2}\left(\theta-\beta_{1}\right)\right)=\left(1-\sin ^{2}\left(\theta-\beta_{1}\right)\right)\left(1-n^{2} \sin ^{2} \beta_{1}\right)\)\\
or,

\[
\begin{gathered}
1-n^{2} \sin ^{2}\left(\theta-\beta_{1}\right)-\sin ^{2} \beta_{1}+\sin ^{2} \beta_{1} n^{2} \sin ^{2}\left(\theta-\beta_{1}\right) \\
=1-n^{2} \sin ^{2} \beta_{1}-\sin ^{2}\left(\theta-\beta_{1}\right)+\sin ^{2} \beta_{1} n^{2} \sin ^{2}\left(\theta-\beta_{1}\right)
\end{gathered}
\]

or,

\[
\sin ^{2}\left(\theta-\beta_{1}\right)-n^{2} \sin ^{2}\left(\theta-\beta_{1}\right)=\sin ^{2} \beta_{1}\left(1-n^{2}\right)
\]

or,

\[
\sin ^{2}\left(\theta-\beta_{1}\right)\left(1-n^{2}\right)=\sin ^{2} \beta_{1}\left(1-n^{2}\right)
\]

or,

\[
\theta-\beta_{1}=\beta_{1} \quad \text { or } \quad \beta_{1}=\theta / 2
\]

But

\[
\beta_{1}+\beta_{2}=\theta, \quad \text { so, } \quad \beta_{2}=\theta / 2=\beta_{1}
\]

which is the case of symmetric passage of ray.\\
In the case of symmetric passage of ray\\
\(\alpha_{1}=\alpha_{2}=\alpha^{\prime}\) (say)\\
and \(\beta_{1}=\beta_{2}=\beta=\theta / 2\)\\
Thus the total deviation\\
\(\alpha=\left(\alpha_{1}+\alpha_{2}\right)-\theta\)\\
\includegraphics[max width=\textwidth, alt={}, center]{da0156b6-4133-434e-90d4-ec01b2070afe-133_370_531_937_872}


\begin{equation*}
\alpha=2 \alpha^{\prime}-\theta \quad \text { or } \quad \alpha^{\prime}=\frac{\alpha+\theta}{2} \tag{1}
\end{equation*}


But from the Snell's law \(\sin \alpha=n \sin \beta\)\\
So,

\[
\sin \frac{\alpha+\theta}{2}=n \sin \frac{\theta}{2}
\]

5.21 In this case we have

\[
\sin \frac{\alpha+\theta}{2}=n \sin \frac{\theta}{2}(\text { see soln. of } 5.20)
\]

In our problem \(\alpha=\theta\)\\
So, \(\quad \sin \theta=n \sin (\theta / 2) \quad\) or \(\quad 2 \sin (\theta / 2) \cos (\theta / 2)=n \sin (\theta / 2)\)\\
Hence \(\cos (\theta / 2)=\frac{n}{2}\) or \(\theta=2 \cos ^{-1}(n / 2)=83^{\circ}\), where \(n=1.5\)\\
5.22 In the case of minimum deviation

\[
\sin \frac{\alpha+\theta}{2}=n \sin \frac{\theta}{2}
\]

So, \(\alpha=2 \sin ^{-1}\left\{n \sin \frac{\theta}{2}\right\}-\theta=37^{\circ}\), for \(n=1.5\)\\
Passage of ray for grazing incidence and grazing imergence is the condition for maximum deviation (Fig.). From Fig.

\[
\alpha=\pi-\theta=\pi-2 \theta_{c r}
\]

(where \(\theta_{c r}\) is the critical angle)\\
So, \(\quad \alpha=\pi-2 \sin (1 / n)=58^{\circ}\), for \(n=1.5=\) R.I. of glass.\\
5.23 The least deflection angle is given by the formula,\\
\includegraphics[max width=\textwidth, alt={}, center]{da0156b6-4133-434e-90d4-ec01b2070afe-134_534_486_120_968}\\
\(\delta=2 \alpha-\theta\), where \(\alpha\) is the angle of incidence at first surface and \(\theta\) is the prism angle.\\
Also from Snell's law, \(n_{1} \sin \alpha=n_{2} \sin (\theta / 2)\), as the angle of refraction at first surface is equal to half the angle of prism for least deflection\\
so,

\[
\sin \alpha=\frac{n_{2}}{n_{1}} \sin (\theta / 2)=\frac{1.5}{1.33} \sin 30^{\circ}=.5639
\]

or,

\[
\alpha=\sin ^{-1}(\cdot 5639)=34 \cdot 3259^{\circ}
\]

Substituting in the above (1), we get, \(\delta=8.65^{\circ}\)\\
5.24 From the Cauchy's formula, and also experimentally the R.I. of a medium depends upon the wavelength of the mochromatic ray i.e. \(n=f(\lambda)\). In the case of least deviation of \(a\) monochromatic ray the passage a prism, we have: .


\begin{equation*}
n \sin \frac{\theta}{2}=\sin \frac{\alpha+\theta}{2} \tag{1}
\end{equation*}


The above equation tales us that we have \(n=n(\alpha)\), so we may write


\begin{equation*}
\Delta n=\frac{d n}{d \alpha} \Delta \alpha \tag{2}
\end{equation*}


From Eqn. (1)\\
or,


\begin{align*}
d n \sin \frac{\theta}{2} & =\frac{1}{2} \cos \frac{\alpha+\theta}{2} d \alpha \\
\frac{d n}{d \alpha} & =\frac{\cos \frac{\alpha+\theta}{2}}{2 \sin \frac{\theta}{2}} \tag{3}
\end{align*}


From Eqns (2) and (3)\\
\includegraphics[max width=\textwidth, alt={}, center]{da0156b6-4133-434e-90d4-ec01b2070afe-134_361_441_1465_954}

\[
\Delta n=\frac{\cos \frac{\alpha+\theta}{2}}{2 \sin \frac{\theta}{2}} \Delta \alpha
\]

or, \(\Delta n=\frac{\sqrt{1-\sin ^{2}\left(\frac{\alpha+\theta}{2}\right)}}{2 \sin \frac{\theta}{2}} \Delta \alpha=\frac{\sqrt{1-n^{2} \sin ^{2} \frac{\theta}{2}}}{2 \sin \frac{\theta}{2}} \Delta \alpha\) (Using Eqn. 1.)

Thus

\[
\Delta \alpha=\frac{2 \sin \frac{\theta}{2}}{\sqrt{1-n^{2} \sin ^{2} \frac{\theta}{2}}} \Delta n=0 \cdot 44
\]

5.25 Fermat's principle : " The actual path of propagation of light (trajectory of a light ray) is the path which can be followed by light with in the lest time, in comparison with all other hypothetical paths between the same two points. "\\
"Above statement is the original wordings of Fermat (A famous French scientist of 17th century)"\\
Deduction of the law of refraction from Fermat's principle :\\
Let the plane \(S\) be the interface between medium 1 and medium 2 with the refractive indices \(n_{1}=c / v_{1}\) and \(n_{2}=c / v_{2}\) Fig. (a). Assume, as usual, that \(n_{1}<n_{2}\). Two points are given- one above the plane \(S\) (point \(A\) ), the other under plane \(S\) (point \(B\) ). The various distances are :\\
\(A A_{1}=h_{1}, B B_{1}=h_{2}, A_{1} B_{1}=l\). We must find the path from \(A\) to \(B\) which can be covered by light faster than it can cover any other hypothetical path. Clearly, this path must consist of two straight lines, viz, \(A O\) in medium 1 and \(O B\) in medium 2; the point \(O\) in the plane \(S\) has to be found.\\
First of all, it follows from Fermat's priniciple that the point \(O\) must lie on the intersection of \(S\) and a plane \(P\), which is perpendicular to \(S\) and passes through \(A\) and \(B\).

\begin{figure}[h]
\begin{center}
  \includegraphics[alt={},max width=\textwidth]{da0156b6-4133-434e-90d4-ec01b2070afe-135_343_362_1276_173}
\captionsetup{labelformat=empty}
\caption{(a)}
\end{center}
\end{figure}

\begin{figure}[h]
\begin{center}
  \includegraphics[alt={},max width=\textwidth]{da0156b6-4133-434e-90d4-ec01b2070afe-135_325_394_1288_548}
\captionsetup{labelformat=empty}
\caption{(b)}
\end{center}
\end{figure}

\begin{figure}[h]
\begin{center}
  \includegraphics[alt={},max width=\textwidth]{da0156b6-4133-434e-90d4-ec01b2070afe-135_309_297_1292_1031}
\captionsetup{labelformat=empty}
\caption{(c)}
\end{center}
\end{figure}

Indeed, let us assume that this point does not lie in the plane \(P\); let this be point \(O_{1}\) in Fig. (b). Drop the perpendicular \(O_{1} O_{2}\) from \(O_{1}\) onto \(P\). Since \(A O_{2}<A O_{1}\) and \(B O_{2}<B O_{1}\), it is clear that the time required to traverse \(A O_{2} B\) is less than that needed to cover the path \(A O_{1} B\). Thus, using Fermat's principle, we see that the first law of refraction is observed : the incident and the refracted rays lie in the same plane as the perpendicular to the interface at the point\\
where the ray is refracted. This plane is the plane \(P\) in Fig. (b); it is called the plane of incidence.

Now let us consider light rays in the plane of incidence Fig. (c). Designate \(A_{1} O\) as \(x\) and \(O B_{1}=l-x\). The time it takes a ray to travel from \(A\) to \(O\) and then from \(O\) to \(B\) is


\begin{equation*}
T=\frac{A O}{\mathrm{v}_{1}}+\frac{O B}{\mathrm{v}_{2}}=\frac{\sqrt{h_{1}^{2}+x^{2}}}{\mathrm{v}_{1}}+\frac{\sqrt{h_{2}^{2}+(l-x)^{2}}}{\mathrm{v}_{2}} \tag{1}
\end{equation*}


The time depends on the value of \(x\). According to Fermat's principle, the value of \(x\) must minimize the time \(T\). At this value of \(x\) the derivative \(d T / d x\) equals zero :


\begin{equation*}
\frac{d T}{d x}=\frac{x}{v_{1} \sqrt{h_{1}^{2}+x^{2}}}-\frac{l-x}{v_{2} \sqrt{h_{2}^{2}+(l-x)^{2}}}=0 . \tag{2}
\end{equation*}


Now,

\[
\frac{x}{\sqrt{h_{1}^{2}+x^{2}}}=\sin \alpha, \text { and } \frac{l-x}{\sqrt{h_{2}^{2}+(l-x)^{2}}}=\sin \beta,
\]

Consequently,

So,

\[
\begin{gathered}
\frac{\sin \alpha}{\mathrm{v}_{1}}-\frac{\sin \beta}{\mathrm{v}_{2}}=0, \quad \text { or } \quad \frac{\sin \alpha}{\sin \beta}=\frac{\mathrm{v}_{1}}{\mathrm{v}_{2}} \\
\frac{\sin \alpha}{\sin \beta}=\frac{c / n_{1}}{c / n_{2}}=\frac{n_{2}}{n_{1}}
\end{gathered}
\]

Note : Fermat himself could not use Eqn. 2. as mathematical analysis was developed later by Newton and Leibniz. To deduce the law of the refraction of light, Fermat used his own maximum and minimum method of calculus, which, in fact, corresponded to the subsequently developed method of finding the minimum (maximum) of a function by differentiating it and equating the derivative to zero.\\
5.26 (a) Look for a point \(O^{\prime}\) on the axis such that \(O^{\prime} P^{\prime}\) and \(O^{\prime} P\) make equal angles with \(O^{\prime} O\). This determines the position of the mirror. Draw a ray from \(P\) parallel to the axis. This must on reflection pass through \(P^{\prime}\). The intersection of the reflected ray with principal axis determines the focus.\\
\includegraphics[max width=\textwidth, alt={}, center]{da0156b6-4133-434e-90d4-ec01b2070afe-136_472_1251_1496_133}\\
(b) Suppose \(P\) is the object and \(P^{\prime}\) is the image. Then the mirror is convex because the image is virtual, crect \& diminished. Look for a point \(X\) (between \(P\) \& \(P^{\prime}\) ) on the axis such that \(P X\) and \(P^{\prime} X\) make equal angle with the axis.

\begin{figure}[h]
\begin{center}
  \includegraphics[alt={},max width=\textwidth]{da0156b6-4133-434e-90d4-ec01b2070afe-137_352_744_205_108}
\captionsetup{labelformat=empty}
\caption{(a)}
\end{center}
\end{figure}

\begin{figure}[h]
\begin{center}
  \includegraphics[alt={},max width=\textwidth]{da0156b6-4133-434e-90d4-ec01b2070afe-137_362_565_199_874}
\captionsetup{labelformat=empty}
\caption{(b)}
\end{center}
\end{figure}

5.27 (a) From the mirror formula,


\begin{equation*}
\frac{1}{s^{\prime}}+\frac{1}{s}=\frac{1}{f} \text { we get } f=\frac{s^{\prime} s}{s^{\prime}+s} \tag{1}
\end{equation*}


In accordance with the problem \(s-s^{\prime}=l \frac{s^{\prime}}{s}=\beta\),\\
From these two relations, we get : \(s=\frac{l}{1-\beta}, \quad s^{\prime}=-\frac{l \beta}{1+\beta}\)\\
Substituting it in the Eqn. (1),

\[
f=\frac{\beta\left(\frac{l}{1-\beta}\right)^{2}}{l\left(\frac{1-\beta}{1-\beta}\right)}=\frac{l \beta}{\left(1-\beta^{2}\right)}=-10 \mathrm{~cm}
\]

(b) Again we have, \(\quad \frac{1}{s^{\prime}}+\frac{1}{s}=\frac{1}{f}\) or, \(\frac{s}{s^{\prime}}+1=\frac{s}{f}\)\\
or,

\[
\frac{1}{\beta_{1}}=\frac{s}{f}-1=\frac{s-f}{f}
\]

or,


\begin{equation*}
\beta_{1}=\frac{f}{s-f} \tag{2}
\end{equation*}


Now, it is clear from the above equation, that for smaller \(\beta, s\) must be large, so the object is displaced away from the mirror in second position.\\
i.e.


\begin{equation*}
\beta_{2}=\frac{f}{s+l-f} \tag{3}
\end{equation*}


Eliminating \(s\) from the Eqn. (2) and (3), we get,

\[
f=\frac{l \beta_{1} \beta_{2}}{\left(\beta_{2}-\beta_{1}\right)}=-2 \cdot 5 \mathrm{~cm}
\]

5.28 For a concave mirror as usual \(\frac{1}{s^{\prime}}+\frac{1}{s}=\frac{1}{f}\) so \(s^{\prime}=\frac{s f}{s-f}\)\\
(In coordinate convention \(s=-s\) is negative \& \(f=-|f|\) is also negative.)\\
If \(A\) is the area of the mirror (assumed small) and the object is on the principal axis, then the light incident on the mirror per second is \(I_{0} \frac{A}{s^{2}}\).\\
This follows from the definition of luminous intensity as light emitted per second per unit solid angle in a given direction and the fact that \(\frac{A}{s^{2}}\) is the solid angle subtended by the mirror at the source. Of this a fraction \(\rho\) is reflected so if \(I\) is the luminous intensity of the image, then \(I \frac{A}{s^{\prime 2}}=\rho I_{0} \frac{A}{s^{2}}\)

Hence

\[
I=\rho I_{0}\left(\frac{|f|}{|f|-s}\right)^{2}
\]

(Because our convention makes \(f\)-ve for a concave mirror, we have to write \(|f|\).)\\
Substitution gives

\[
I=2.0 \times 10^{3} \mathrm{~cd} .
\]

5.29 For \(O_{1}\) to be the image, the optical paths of all rays \(O A O_{1}\) must be equal upto terms of leading order in \(h\). Thus

\[
n_{1} O A+n_{2} A O_{1}=\text { constant }
\]

But, \(O P=|s|, O_{1} P=\left|s^{\prime}\right|\) and so\\
\(O A=\sqrt{h^{2}+(|s|+\delta)^{2}} \cong|s|+\delta+\frac{h^{2}}{2|s|}\)\\
\(O_{1} A=\sqrt{h^{2}+\left(\left|s^{\prime}\right|-\delta\right)^{2}} \cong\left|s^{\prime}\right|-\delta+\frac{h^{2}}{2\left|s^{\prime}\right|}\)\\
(neglecting products \(h^{2} \delta\) ). Then\\
\includegraphics[max width=\textwidth, alt={}, center]{da0156b6-4133-434e-90d4-ec01b2070afe-138_342_651_1133_832}\\
\(n_{1}|s|+n_{2}\left|s^{\prime}\right|+n_{1} \delta-n_{2} \delta+\frac{h^{2}}{2}\left(\frac{n_{1}}{|s|}+\frac{n_{2}}{\left|s^{\prime}\right|}\right)=\) Const.\\
Now \(\quad(r-\delta)^{2}+h^{2}=r^{2}\)\\
or \(\quad h^{2}=2 \dot{r} \delta \quad\) or \(\quad \delta=\frac{h^{2}}{2 r}\)\\
Here \(\quad r=C P\).\\
Hence \(n_{1}|s|+n_{2}\left|s^{\prime}\right|+\frac{h^{2}}{2}\left\{\frac{n_{1}-n_{2}}{r}+\frac{n_{1}}{|s|}+\frac{n_{2}}{\left|s^{\prime}\right|}\right\}=\) Constant\\
Since this must hold for all \(h\), we have

\[
\frac{n_{2}}{\left|s^{\prime}\right|}+\frac{n_{1}}{|s|}=\frac{n_{2}-n_{1}}{r}
\]

From our sign convention, \(s^{\prime}>0, s<0\) so we get

\[
\frac{n_{2}}{s^{\prime}}-\frac{n_{1}}{s}=\frac{n_{2}-n_{1}}{r} .
\]

5.30 All rays focusing at a point must have traversed the same optical path. Thus\\
\includegraphics[max width=\textwidth, alt={}, center]{da0156b6-4133-434e-90d4-ec01b2070afe-139_389_549_297_233}\\
\includegraphics[max width=\textwidth, alt={}, center]{da0156b6-4133-434e-90d4-ec01b2070afe-139_409_260_301_830}

\[
x+n \sqrt{r^{2}+(f-x)^{2}}=n f \quad \text { or } \quad(n f-x)^{2}=n^{2} r^{2}+n^{2}(f-x)^{2}
\]

or,

\[
\begin{aligned}
n^{2} r^{2}=(n f-x)^{2}-[n(f-x)]^{2} & =(n f-x+n f-n x)(n f-x-n f+n x) \\
& =x(n-1)(2 n f-(n+1) x) \\
& =2 n(n-1) f x-(n+1)(n-1) x^{2}
\end{aligned}
\]

Thus,

\[
(n+1)(n-1) x^{2}-2 n(n-1) f x+n^{2} r^{2}=0
\]

so,

\[
\begin{aligned}
x & =\frac{n(n-1) f \pm \sqrt{n^{2}(n-1)^{2} f^{2}-n^{2} r^{2}(n+1)(n-1)}}{(n+1)(n-1)} \\
& =\frac{n f}{n+1}\left[1 \pm \sqrt{1-\frac{n+1}{n-1} \frac{r^{2}}{f^{2}}}\right]
\end{aligned}
\]

Ray must move forward so \(x<f\), for \(+\operatorname{sign} x>f\) for small \(r\), so -sign.\\
(Also \(x \rightarrow 0\) as \(r \rightarrow 0\) )\\
( \(x>f\) means ray turning back in the direction of incidence. (see Fig.)

Hence

\[
x=\frac{n f}{n+1}\left[1-\sqrt{1-\frac{n+1}{n-1} \frac{r^{2}}{f^{2}}}\right]
\]

For the maximum value of \(r\),


\begin{equation*}
\sqrt{1-\frac{n+1}{n-1} \frac{r^{2}}{f^{2}}}=0 \tag{A}
\end{equation*}


because the expression under the radical sign must be non-negative, which gives the maximum value of \(r\).\\
Hence from Eqn. (A), \(\quad r_{\text {max }}=f \sqrt{(n-1) /(n+1)}\)\\
5.31 As the given lense has significant thickness, the thin lense, formula cannote be used. For refraction at the front surface from the formula \(\frac{n^{\prime}}{s^{\prime}}-\frac{n}{s}=\frac{n^{\prime}-n}{R}\)

\[
\frac{1 \cdot 5}{s^{\prime}}-\frac{1}{-20}=\frac{1 \cdot 5-1}{5}
\]

On simplifying we get , \(\boldsymbol{s}^{\boldsymbol{\prime}}=30 \mathrm{~cm}\).\\
Thus the image \(I^{\prime}\) produced by the front surface behaves as a virtual source for the rear surface at distance 25 cm from it, because the thickness of the lense is 5 cm . Again from the refraction formula at cerve surface

\[
\begin{aligned}
& \frac{n^{\prime}}{s^{\prime}}-\frac{n}{s}=\frac{n^{\prime}-n}{R} \\
& \frac{1}{s^{\prime}}-\frac{1 \cdot 5}{25}=\frac{1-1 \cdot 5}{-5}
\end{aligned}
\]

On simplifying, \(s^{\prime}=+6 \cdot 25 \mathrm{~cm}\)\\
\includegraphics[max width=\textwidth, alt={}, center]{da0156b6-4133-434e-90d4-ec01b2070afe-140_241_728_526_692}

Thus we get a real image \(I\) at a distance 6.25 cm beyond the rear surface (Fig.).\\
5.32 (a) The formation of the image of a source \(S\), placed at a distance \(u\) from the pole of the convex surface of plano-convex lens of thickness \(d\) is shown in the figure.

On applying the formula for refraction through spherical surface, we get\\
\(\frac{n}{s^{\prime}}-\frac{1}{s}=(n-1) / R,\left(\right.\) here \(n_{2}=n\) and \(\left.n_{1}=1\right)\)\\
or, \(\frac{n}{d}-\frac{1}{s}=(n-1) / R \quad\) or, \(\frac{1}{s}=\frac{n}{d}-\frac{(n-1)}{R}\)\\
or,

\[
\frac{s^{\prime}}{s}=s^{\prime}\left\{\frac{n}{d}-\frac{(n-1)}{R}\right\}
\]

\begin{center}
\includegraphics[max width=\textwidth, alt={}]{da0156b6-4133-434e-90d4-ec01b2070afe-140_398_533_1032_933}
\end{center}

But in this case optical path of the light, corresponding to the distance \(v\) in the medium is \(v / n\), so the magnification produced will be,

\[
\beta=\frac{s^{\prime}}{n s}=\frac{s^{\prime}}{n}\left\{\frac{n}{d}-\frac{(n-1)}{R}\right\}=\frac{d}{n}\left\{\frac{n}{d}-\frac{(n-1)}{R}\right\}=1-\frac{d(n-1)}{n R}
\]

Substituting the values, we get magnification \(\beta=-0.20\).\\
(b) If the transverse area of the object is \(A\) (assumed small), the area of the image is \(\beta^{2} A\). We shall assume that \(\frac{\pi D^{2}}{4}>A\). Then light falling on the lens is : \(L A \frac{\pi D^{2} / 4}{s^{2}}\)\\
from the definition of luminance (See Eqn. (5.1c) of the book; here \(\cos \theta \approx 1\) if \(D^{2} \ll s^{2}\) and \(d \Omega=\frac{\pi D^{2} / 4}{s^{2}}\) ). Then the illuminance of the image is

\[
L A \frac{\pi D^{2} / 4}{s^{2}} / \beta^{2} A=L n^{2} \pi D^{2} / 4 d^{2}
\]

Substitution gives 421 x .\\
5.33 (a) Optical power of a thin lens of R.I. \(n\) in a medium with R.I. \(n_{0}\) is given by :


\begin{equation*}
\Phi=\left(n-n_{0}\right)\left(\frac{1}{R_{1}}-\frac{1}{R_{2}}\right) \tag{A}
\end{equation*}


From Eqn.(A), when the lens is placed in air :


\begin{equation*}
\Phi_{0}=(n-1)\left(\frac{1}{R_{1}}-\frac{1}{R_{2}}\right) \tag{1}
\end{equation*}


Similarly from Eqn.(A), when the lens is placed in liquid :


\begin{equation*}
\Phi=\left(n-n_{0}\right)\left(\frac{1}{R_{1}}-\frac{1}{R_{2}}\right) \tag{2}
\end{equation*}


Thus from Eqns (1) and (2)

\[
\Phi=\frac{n-n_{0}}{n-1} \Phi_{0}=2 D
\]

The second focal length, is given by \(f^{\prime}=\frac{n^{\prime}}{\Phi}\), where \(n^{\prime}\) is the R.I. of the medium in which it is placed.

\[
f^{\prime}=\frac{n_{0}}{\Phi}=85 \mathrm{~cm}
\]

(b) Optical power of a thin lens of R.I. \(n\) placed in a medium of R.I. \(n_{0}\) is given by :


\begin{equation*}
\Phi=\left(n-n_{0}\right)\left(\frac{1}{R_{1}}-\frac{1}{R_{2}}\right) \tag{A}
\end{equation*}


For a biconvex lens placed in air medium from Eqn. (A)


\begin{equation*}
\left.\Phi_{0}=(n-1)\left(\frac{1}{R}-\frac{1}{-R}\right)=\frac{2(n-1)}{R}\right) \tag{1}
\end{equation*}


where \(R\) is the radius of each curve surface of the lens Optical power of a spherical refractive surface is given by :


\begin{equation*}
\Phi=\frac{n^{\prime}-n}{R} \tag{B}
\end{equation*}


For the rear surface of the lens which divides air and glass medium

\[
\Phi_{0}=\frac{n-1}{R} \text { (Here } n \text { is the R.I. (2) of glass) }
\]

Similarly for the front surface which divides wate nd glass medium


\begin{equation*}
\Phi_{l}=\frac{n-n}{-R}=\frac{n-n_{0}}{R} \tag{3}
\end{equation*}


Hence the optical power of the given optical system


\begin{equation*}
\Phi=\Phi_{a}+\Phi_{l}=\frac{n-1}{R}+\frac{n-n_{0}}{R}=\frac{2 n-n_{0}-1}{R} \tag{4}
\end{equation*}


From Eqns (1) and (4)

\[
\frac{\Phi}{\Phi_{0}}=\frac{2 n-n_{0}-1}{2(n-1)} \text { So } \Phi=\frac{\left(2 n-n_{0}-1\right)}{2(n-1)} \Phi_{0}
\]

Focal length in air, \(f=\frac{1}{\Phi}=15 \mathrm{~cm}\)\\
and focal length in water \(=\frac{n_{0}}{\Phi}=20 \mathrm{~cm}\) for \(n_{0}=\frac{4}{3}\).\\
5.34 (a) Clearly the media on the sides are different. The front focus \(F\) is the position of the object (virtual or real) for which the image is formated at infinity. The rear focus \(F^{\prime}\) is the position of the image (virtual or real) of the object at infinity. (a) Figures 5.7 (a) \& (b). This geometrical construction ensines that the second of the equations (5.1g) is obeyed.\\
\includegraphics[max width=\textwidth, alt={}, center]{da0156b6-4133-434e-90d4-ec01b2070afe-142_332_779_1041_102}\\
(a) Convex lens\\
(b) Figure 5.5 (a) \& (b) with lens\\
\includegraphics[max width=\textwidth, alt={}, center]{da0156b6-4133-434e-90d4-ec01b2070afe-142_391_581_1470_136}\\
(a) Convex lens\\
( \(P\) is the object)\\
\includegraphics[max width=\textwidth, alt={}, center]{da0156b6-4133-434e-90d4-ec01b2070afe-142_484_599_933_883}\\
\includegraphics[max width=\textwidth, alt={}, center]{da0156b6-4133-434e-90d4-ec01b2070afe-142_436_562_1433_870}\\
(b) Concave lens\\
(c) Figure (5.8) (a) \& (b).

Clearly, the important case is that when the rays\\
(1) \& (2) are not symmetric about the principal axis, otherwise the figure can be completed by reflection in the principal axis. Knowing one path we know the path of all rays connecting the two points. For a different object. We proceed as shown below, we use the fact that a ray incident at a given height above the optic centre suffers a definite deviation.\\
\includegraphics[max width=\textwidth, alt={}, center]{da0156b6-4133-434e-90d4-ec01b2070afe-143_425_667_148_806}

The concave lens can be discussed similarly.\\
5.35 Since the image is formed on the screen, it is real, so for a conversing lens object is in the incident side.\\
Let \(s_{1}\) and \(s_{2}\) be the magnitudes of the object distance in the first and second case respectively. We have the lens formula


\begin{equation*}
\frac{1}{s^{\prime}}-\frac{1}{s}=\frac{1}{f} \tag{1}
\end{equation*}


In the first case from Eqn. (1)

\[
\frac{1}{(+l)}-\frac{1}{\left(-s_{1}\right)}=\frac{1}{f} \quad \text { or, } \quad s_{1}=\frac{f(l)}{l-f}=26 \cdot 31 \mathrm{~cm} .
\]

Similarly from Eqn.(1) in the second case

\[
\frac{1}{(l-\Delta l)}-\frac{1}{\left(-s_{2}\right)}=\frac{1}{f} \quad \text { or, } s_{2}=\frac{l f .}{(l-\Delta l)-f}=26.36 \mathrm{~cm} .
\]

Thus the sought distance \(\Delta x=s_{2}-s_{1}=0.5 \mathrm{~mm} \approx \Delta l f^{2} /\left(l-f^{2}\right)\)\\
5.36 The distance between the object and the image is \(l\). Let \(x=\) distance between the object and the lens. Then, since the image is real, we have in our convention, \(u=-x, \mathrm{v}=l-x\)

So\\
or

\[
\begin{gathered}
\frac{1}{x}+\frac{1}{l-x}=\frac{1}{f} \\
x(l-x)=l f \quad \text { or } x^{2}-x l+l f=0
\end{gathered}
\]

Solving we get the roots

\[
x=\frac{1}{2}\left[l \pm \sqrt{l^{2}-4 l f}\right]
\]

(We must have \(l>4 f\) for real roots.)\\
(a) If the distance between the two positions of the lens is \(\Delta l\), then clearly \(\Delta l=x_{2}-x_{1}=\) difference between roots \(=\sqrt{l^{2}-4 l f}\)\\
so

\[
f=\frac{l^{2}-\Delta l^{2}}{4 l}=20 \mathrm{~cm} .
\]

(b) The two roots are conjugate in the sense that if one gives the object distance the other gives the corresponding image distance (in both cases). Thus the magnifications are

\[
-\frac{l+\sqrt{l^{2}-4 l f}}{l-\sqrt{l^{2}-4 l f}}(\text { enlarged }) \text { and }-\frac{l-\sqrt{l^{2}-4 l f}}{l+\sqrt{l^{2}-4 l f}} \text { (diminished). }
\]

The ratio of these magnification being \(\eta\) we have\\
or

\[
\begin{gathered}
\frac{l-\sqrt{l^{2}-4 l f}}{l-\sqrt{l^{2}-4 l f}}=\sqrt{\eta} \quad \text { or } \frac{\sqrt{l^{2}-4 l f}}{l}=\frac{\sqrt{\eta}-1}{\sqrt{\eta}+1} \\
1-\frac{4 f}{l}=\left(\frac{\sqrt{\eta}-1}{\sqrt{\eta}+1}\right)^{2}=1-4 \frac{\sqrt{\eta}}{(1+\sqrt{\eta})^{2}} \\
f=l \frac{\sqrt{\eta}}{(1+\sqrt{\eta})^{2}}=20 \mathrm{~cm} .
\end{gathered}
\]

Hence\\
5.37 We know from the previous problem that the two magnifications are reciprocals of each other ( \(\beta^{\prime} \beta^{\prime \prime}=1\) ). If \(h\) is the size of the object then \(h^{\prime}=\beta^{\prime} h\) and

Hence

\[
\begin{gathered}
h^{\prime \prime}=\beta^{\prime \prime} h \\
h=\sqrt{h^{\prime} h^{\prime \prime}} .
\end{gathered}
\]

5.38 Refer to problem \(5.32(\mathrm{~b})\). If \(A\) is the area of the object, then provided the angular diameter of the object at the lens is much smaller than other relevant angles like \(\frac{D}{f}\) we calculate the light falling on the lens as \(L A \frac{\pi D^{2}}{4 s^{2}}\)\\
where \(u^{2}\) is the object distance squared. If \(\beta\) is the transverse magnification \(\left(\beta=\frac{s^{\prime}}{u}\right)\) then the area of the image is \(\beta^{2} A\). Hence the illuminance of the image (also taking account of the light lost in the lens)

\[
E=(1-\alpha) L A \frac{\pi D^{2}}{4 s^{2}} \frac{1}{\beta^{2} A}=\frac{(1-\alpha) \pi D^{2} L}{4 f^{2}}
\]

since \(s^{\prime}=f\) for a distant object. Substitution gives \(E=151 \mathrm{x}\).\\
5.39 (a) If \(s=\) object distance, \(s^{\prime}=\) average distance, \(L=\) luminance of the sounce, \(\Delta S=\) area of the source assumed to be a plane surface held normal to the principal axis, then we find for the flux \(\Delta \Phi\) incident on the lens

\[
\begin{array}{r}
\Delta \Phi=\int L \Delta S \cos \theta d \Omega \\
\approx L \Delta S \int_{0}^{\infty} \cos \theta 2 \pi \sin \theta d \theta=L \Delta S \pi \sin ^{2} \alpha \cong L \Delta S \frac{\pi D^{2}}{4 s^{2}}
\end{array}
\]

Here we are assuming \(D \ll s\), and ignoring the variation of \(L\) since \(\alpha\) is small\\
Then if \(L^{\prime}\) is the luminance of the image, and \(\Delta S^{\prime}=\left(\frac{S^{\prime}}{S}\right)^{2} \Delta S\) is the are of the image then similarly

\[
\begin{aligned}
& L^{\prime} \Delta S^{\prime} \frac{D^{2}}{4 s^{\prime 2}} \pi=L^{\prime} \Delta S \frac{D^{2}}{4 s^{2}} \pi=L \Delta S \frac{D^{2}}{4 s^{2}} \pi \\
& \text { or } \quad L^{\prime}=L \text { irrespective of } D
\end{aligned}
\]

\includegraphics[max width=\textwidth, alt={}, center]{da0156b6-4133-434e-90d4-ec01b2070afe-145_206_633_242_813}\\
(b) In this case the image on the white screen from a Lambert source. Then if its luminance is \(L_{0}\) its luminosity will be the \(\pi L_{0}\) and\\
or

\[
\begin{aligned}
\pi L_{0} \frac{s^{\prime 2}}{s^{2}} \Delta S & =L \Delta S \frac{D^{2}}{4 s^{2}} \pi \\
L_{0} & \propto D^{2}
\end{aligned}
\]

since \(\boldsymbol{s}^{\prime}\) depends on \(f, s\) but not on \(D\).\\
5.40 Focal length of the converging lens, when it is submerged in water of R.I. \(n_{0}\) (say) :


\begin{equation*}
\frac{1}{f_{1}}=\left(\frac{n_{1}}{n_{0}}-1\right)\left(\frac{1}{R}-\frac{1}{R}\right), \frac{2\left(n_{1}-n_{0}\right)}{n_{0} R} \tag{1}
\end{equation*}


Simillarly, the focal length of diverging lens in water.


\begin{equation*}
\frac{1}{f_{2}}=\left(\frac{n_{2}}{n_{0}}-1\right)\left(\frac{1}{-R}-\frac{1}{R}\right) \frac{-2\left(n_{2}-n_{0}\right)}{n_{0} R} \tag{2}
\end{equation*}


Now, when they are put together in the water, the focal length of the system,\\
or,

\[
\begin{gathered}
\frac{1}{f}=\frac{1}{f_{1}}+\frac{1}{f_{2}} \\
=\frac{2\left(n_{1}-n_{2}\right)}{n_{0} R}-\frac{2\left(n_{2}-n_{0}\right)}{n_{0} R}=\frac{2\left(n_{1}-n_{2}\right)}{n_{0} R} \\
f=\frac{-n_{0} R}{2\left(n_{1}-n_{2}\right)}=35 \mathrm{~cm}
\end{gathered}
\]

5.41 \(C\) is the centre of curvature of the silvered surface and \(O\) is the effective centre of the equivalent mirror in the sence that an object at \(O\) forms a coincident image. From the figure, using the formula for refraction at a spherical surface, we have\\
\(\frac{n}{-R}-\frac{1}{2 f}=\frac{n-1}{R} \quad\) or \(\quad f=\frac{-R}{2(2 n-1)}\)\\
(In our convention \(f\) is \(-v e\) ).\\
\includegraphics[max width=\textwidth, alt={}, center]{da0156b6-4133-434e-90d4-ec01b2070afe-145_355_517_1534_886}

Substitution gives \(f=-10 \mathrm{~cm}\).\\
5.42 (a) Path of a ray, as it passes through the lens system is as shown below.

Focal length of all the three lenses, \(f=\frac{1}{10} \mathrm{~m}=10 \mathrm{~cm}\), neglecting their signs.

Applying lens formula for the first lens, considering a ray coming from infinity, \(\frac{1}{s^{\prime}}-\frac{1}{\infty}=\frac{1}{f} \quad\) or, \(s^{\prime}=f=10 \mathrm{~cm}\),

\begin{figure}[h]
\begin{center}
  \includegraphics[alt={},max width=\textwidth]{da0156b6-4133-434e-90d4-ec01b2070afe-146_323_671_123_791}
\captionsetup{labelformat=empty}
\caption{(a)}
\end{center}
\end{figure}

and so the position of the image is 5 cm to the right of the second lens, when only the first one is present, but the ray again gets refracted while passing through the second, so,

\[
\frac{1}{s^{\prime}}-\frac{1}{5}=\frac{1}{f}=\frac{1}{-10}
\]

or, \(s^{\prime}=10 \mathrm{~cm}\), which is now 5 cm left to the third lens so for this lens,\\
or,

\[
\begin{aligned}
& \frac{1}{s^{\prime \prime}}-\frac{1}{5}=\frac{1}{10} \quad \text { or } \quad \frac{1}{s^{\prime \prime}}=\frac{3}{10} \\
& s^{\prime \prime}=10 / 3=3.33 \mathrm{~cm} . \text { from the last lens. }
\end{aligned}
\]

(b) This means that if the object is \(x \mathrm{~cm}\) to be left of the first lens on the axis \(O O^{\prime}\) then the image is \(x\) on to the right of the 3rd (last) lens. Call the lenses \(1,2,3\) from the left and let \(O\) be the object, \(O_{1}\) its image by the first lens, \(O_{2}\) the image of \(O_{1}\) by the 2nd lens and \(O_{3}\), the image of \(O_{2}\) by the third lens.\\
\(O_{1}\) and \(O_{2}\) must be symmetrically located with respect to the lens \(L_{2}\) and since this lens is concave, \(O_{1}\) must be at a distance \(2\left|f_{2}\right|\) to be the right of \(L_{2}\) and \(O_{2}\) must be \(2\left|f_{2}\right|\) to be the left of \(L_{2}\). One can check that this satisfies lens

\begin{figure}[h]
\begin{center}
  \includegraphics[alt={},max width=\textwidth]{da0156b6-4133-434e-90d4-ec01b2070afe-146_263_531_1139_915}
\captionsetup{labelformat=empty}
\caption{(b)}
\end{center}
\end{figure}

Hence

\[
\begin{gathered}
u=-\left(2\left|f_{2}\right|+5\right)=-25 \mathrm{~cm} \\
s^{\prime}=x, \quad f_{3}=10 \mathrm{~cm} \\
\frac{1}{x}+\frac{1}{25}=\frac{1}{10} \quad \text { so } x=16.67 \mathrm{~cm}
\end{gathered}
\]

5.43 (a) Angular magnification for Galilean telescope in normal adjustment is given as.


\begin{gather*}
\Gamma=f_{o} / f_{e} \\
\text { or, } \quad 10=f_{o} y f_{e} \quad \text { or } f_{o}=10 f_{e}
\end{gather*}


The length of the telescope in this case.

\[
l=f_{o}-f_{e}=45 \mathrm{~cm} . \text { given },
\]

So, using (1), we get,

\[
f_{e}=+5 \text { and } f_{o}=+50 \mathrm{~cm} .
\]

(b) Using lens formula for the objective, \(\frac{1}{s_{o}^{\prime}}-\frac{1}{s_{o}}=\frac{1}{f_{o}}\) or, \(s_{o}^{\prime}=\frac{s_{o} f_{o}}{s_{o}+f_{o}}=3.5 \mathrm{~cm}\)

From the figure, it is clear that, \(s_{o}^{\prime}=l^{r}+f_{e}\) where \(l^{\prime}\) is the new tube length.\\
\includegraphics[max width=\textwidth, alt={}, center]{da0156b6-4133-434e-90d4-ec01b2070afe-147_350_481_186_863}\\
or, \(l^{\prime}=v_{o}-f_{e}=50.5-5=45.5 \mathrm{~cm}\).\\
So, the displacement of ocular is,

\[
\Delta l=l^{\prime}-l=45 \cdot 5-45=0 \cdot 5 \mathrm{~cm}
\]

5.44 In the Keplerian telescope, in normal adjustment, the distance between the objective and eyepiece is \(f_{0}+f_{e}\). The image of the mounting produced by the eyepiece is formed at a distance \(v\) to the right where

But

\[
\begin{gathered}
\frac{1}{s^{\prime}}-\frac{1}{s}=\frac{1}{f_{e}} \\
s=-\left(f_{0}+f_{e}\right), \\
\frac{1}{s^{\prime}}=\frac{1}{f_{e}}-\frac{1}{f_{0}+f_{e}}=\frac{f_{0}}{f_{e}\left(f_{0}+f_{e}\right)}
\end{gathered}
\]

so\\
The linear magnification produced by the eyepiece of the mounting is, in magnitude,

\[
|\beta|=\left|\frac{s^{\prime}}{s}\right|=\frac{f_{e}}{f_{0}}
\]

This equals \(\frac{d}{D}\) according to the problem so

\[
\Gamma=\frac{f_{e}}{f_{e}}=\frac{D}{d} .
\]

5.45 It is clear from the figure that a parallel beam\\
of light, originally of intensity \(I_{0}\) has, on emerging from the telescope, an intensity.

\[
I=I_{0}\left(\frac{f_{0}}{f_{e}}\right)^{2}
\]

because it is concentrated over a section whose diameter is \(f_{e} / f_{0}\) of the diameter of the cross section of the incident beam.\\
\includegraphics[max width=\textwidth, alt={}, center]{da0156b6-4133-434e-90d4-ec01b2070afe-147_224_691_1726_760}

Thus

\[
\eta=\left(\frac{f_{0}}{f_{e}}\right)^{2}
\]

So

\[
\Gamma=\frac{f_{0}}{f_{e}}=\sqrt{\eta}
\]

Now

\[
\Gamma=\frac{\tan \Psi^{\prime}}{\tan \Psi}=\frac{\Psi^{\prime}}{\Psi}
\]

Hence \(\Psi=\Psi^{\prime} / \sqrt{\eta}=0.6^{\prime}\) on substitution.\\
5.46 When a glass lens is immersed in water its focal length increases approximately four times.

We check this as follows as :

\[
\begin{gathered}
\frac{1}{f_{a}^{\prime}}=(n-1)\left(\frac{1}{R_{1}}-\frac{1}{R_{2}}\right) \\
\frac{1}{f_{w}}=\left(\frac{n}{n_{0}}-1\right)\left(\frac{1}{R_{1}}-\frac{1}{R_{2}}\right)=\frac{\frac{n}{n_{0}}-1}{n-1} \cdot \frac{1}{f_{a}}=\frac{n-n_{0}}{n_{0}(n-1)} \frac{1}{f_{a}}
\end{gathered}
\]

Now back to the problem. Originally in air

\[
\Gamma=\frac{f_{0}}{f_{e}}=15 \text { so } l=f_{0}+f_{e}=f_{e}(\Gamma+1)
\]

In water,

\[
f_{e}^{\prime}=\frac{n_{0}(n-1)}{n-n_{0}} f_{e}
\]

and the focal length of the replaced objective is given by the condition\\
or

\[
f_{0}^{\prime}+f_{e}^{\prime}=l=(\Gamma+1) f_{e}
\]

Hence

\[
f_{0}^{\prime}=(\Gamma+1) f_{e}-f_{e}^{\prime}
\]

Hence

\[
\Gamma^{\prime}=\frac{f_{0}^{\prime}}{f_{e}^{\prime}}=(\Gamma+1) \frac{n-n_{0}}{n_{0}(n-1)}-1
\]

Substitution gives ( \(n=1.5, n_{0}=1.33\) ), \(\Gamma^{\prime}=3.09\)\\
5.47 If \(L\) is the luminance of the object, \(A\) is its area, \(s=\) distance of the object then light falling on the objective is

\[
\frac{L \pi D^{2}}{4 s^{2}} A
\]

The area of the image formed by the telescope (assuming that the image coincides with the object) is \(\Gamma^{2} \mathrm{~A}\) and the area of the final image on the retina is

\[
=\left(\frac{f}{s}\right)^{2} \Gamma^{2} A
\]

Where \(f=\) focal length of the eye lens. Thus the illuminance of the image on the retin (when the object is observed through the telescope) is

\[
\frac{L \pi D^{2} A}{4 u^{2}\binom{f}{s}^{2} \Gamma^{2} A}=\frac{L \pi D^{2}}{4 f^{2} \Gamma^{2}}
\]

When the object is viewed directly, the illuminance is, similarly, \(\frac{L \pi d_{0}^{2}}{4 f^{2}}\)

We want

\[
\frac{L \pi D^{2}}{4 f^{2} \Gamma^{2}} \geq \frac{L \pi d_{0}^{2}}{4 f^{2}}
\]

So, \(\Gamma \leq \frac{D}{d_{0}}=20\) on substitution of the values.\\
5.48 Obviously, \(f_{o}=+1 \mathrm{~cm}\) and \(f_{e}=+5 \mathrm{~cm}\)

Now, we know that, magnification of a microscope,\\
or,

\[
\begin{aligned}
& \Gamma=\left(\frac{s_{o}^{\prime}}{f_{o}}-1\right) \frac{D}{f_{e}}, \text { for distinct vision } \\
& 50=\left(\frac{s_{o}^{\prime}}{1}-1\right) \frac{25}{5} \quad \text { or, } v_{o}=11 \mathrm{~cm} .
\end{aligned}
\]

Since distance between objective and ocular has increased by 2 cm , hence it will cause the increase of tube length by 2 cm .\\
so,

\[
s_{o}^{\prime}=s_{o}^{\prime}+2=13
\]

and hence, :

\[
\Gamma^{\prime}=\left(\frac{s_{o}^{\prime}}{f_{o}}-1\right) \frac{D}{f_{e}}=60
\]

5.49 It is implied in the problem that final image of the object is at infinity (otherwise light coming out of the eyepiece will not have a definite diameter).\\
(a) We see that \(s_{0}^{\prime} 2 \beta=\left|s_{0}\right| 2 \alpha\), then\\
\(\beta=\frac{\left|s_{0}\right|}{s_{0}^{\prime}} \alpha\)\\
Then, from the figure\\
\includegraphics[max width=\textwidth, alt={}, center]{da0156b6-4133-434e-90d4-ec01b2070afe-149_207_661_1345_762}\\
\(d=2 f_{e} \beta=2 f_{e} \alpha / \frac{s_{0}^{\prime}}{\left|s_{0}\right|}\)\\
But when the final image is at infinity, the magnification \(\Gamma\) in a microscope is given by\\
\(\Gamma=\frac{s_{0}^{\prime}}{\left|s_{0}\right|} \cdot \frac{l}{f_{e}}(l=\) least distance of distinct vision \()\) So \(d=2 l \alpha / \Gamma\)\\
So \(d=d_{0}\) when \(\Gamma=\Gamma_{0}=\frac{2 l \alpha}{d_{0}}=15\) on putting the values.\\
(b) If \(\Gamma\) is the magnification produced by the microscope, then the area of the image produced on the retina (when we observe an object through a microscope) is : \(\Gamma^{2}\left(\frac{f}{s}\right)^{2} A\)

Where \(u=\) distance of the image produced by the microscope from the eye lens, \(f=\) focal length of the eye lens and \(A=\) area of the object. If \(\Phi=\) luminous flux reaching the objective from the object and \(d \leq d_{0}\) so that the entire flux is admitted into the eye), then the illuminance of the final image on the retina

\[
=\frac{\Phi}{\Gamma^{2}(f / s)^{2} A}
\]

But if \(d \geq d_{0}^{2}\), then only a fraction \(\left(d_{0} \mid d\right)^{2}\) of light is admitted into the eye and the illuminance becomes

\[
\frac{\Phi}{A\left(\frac{f}{s}\right)^{2} \Gamma^{2}}\left(\frac{d_{0}}{d}\right)^{2}=\frac{\Phi d_{0}^{2}}{A\left(\frac{f}{s}\right)^{2}(2 l \alpha)^{2}}
\]

independent of \(\Gamma\). The condition for this is then

\[
d \geq d_{0} \text { or } \Gamma \leq \Gamma_{0}=15
\]

5.50 The primary and secondary focal length of a thick lens are given as,\\
and

\[
\begin{aligned}
f & =-(n / \Phi)\left\{1-\left(d / n^{\prime}\right) \Phi_{2}\right\} \\
f^{\prime} & =+\left(n^{\prime \prime} / \Phi\right)\left\{1-\left(d / n^{\prime}\right) \Phi_{1}\right\}
\end{aligned}
\]

where \(\Phi\) is the lens power \(n, n^{\prime}\) and \(n^{\prime \prime}\) are the refractive indices of first medium, lens material and the second medium beyond the lens. \(\Phi_{1}\) and \(\Phi_{2}\) are the powers of first and second spherical surface of the lens.\\
Here,

\[
n=1, \text { for lens, } \quad n^{\prime}=n, \text { for air }
\]

and

\[
n^{\prime \prime}=n_{0}, \text { for water. }
\]

So,

\[
\text { and } \left.\begin{array}{rl}
f & =-1 / \Phi_{1}  \tag{1}\\
f^{\prime} & =+n_{0} / \Phi
\end{array}\right\}, \text { as } d \approx 0
\]

Now, power of a thin lens,

\[
\Phi=\Phi_{1}+\Phi_{2},
\]

where,

\[
\Phi_{1}=\frac{(n-1)}{R}
\]

and

\[
\Phi_{2}=\frac{\left(n_{0}-n\right)}{-R}
\]

So,


\begin{equation*}
\Phi=\left(2 n-n_{0}-1\right) / R \tag{2}
\end{equation*}


From equations (1) and (2), we get,

\[
f=\frac{-R}{(2 n-n-1)}=-11.2 \mathrm{~cm}
\]

and

\[
f=\frac{n_{0} R}{\left(2 n-n_{0}-1\right)}=+14.9 \mathrm{~cm} .
\]

Since the distance between the primary principal point and primary nodal point is given as,

So, in this case,

\[
\begin{gathered}
x=f\left\{\left(n^{\prime \prime}-n\right) / n^{\prime \prime}\right\} \\
x=\left(n_{0} / \Phi\right)\left(n_{0}-1\right) / n_{0}=\left(n_{0}-1\right) / \Phi \\
=\frac{n_{0}}{\Phi}-\frac{1}{\Phi}=f^{\prime}+f=3.7 \mathrm{~cm}
\end{gathered}
\]

5.51 See the answersheet of problem book.\\
5.52 (a) Draw \(P^{\prime} X\) parallel to the axis \(O O^{\prime}\) and let \(P F\) interest it at \(X\). That determines the principal point \(H\). As the medium on both sides of the system is the same, the principal point coincides with the nodal point. Draw a ray parallel to \(P H\) through \(P^{\prime}\). That determines \(H^{\prime}\). Draw a ray \(P X^{\prime}\) parallel to the axis and join \(P^{\prime} X^{\prime}\). That gives \(\boldsymbol{F}^{\boldsymbol{\prime}}\).\\
(b) We let \(H\) stand for the principal point (on the axis). Determine \(H^{\prime}\) by drawing a ray \(P^{\prime} H^{\prime}\) passing through \(P^{\prime}\) and parallel to \(P H\). One ray (conjugate to \(S H\) ) can be obtained from this. To get the other ray one needs to know \(F\) or \(F^{\prime}\). This is easy because \(P\) and \(P^{\prime}\) are known. Finally we get \(S^{\prime}\).\\
(c) From the incident ray we determine \(Q\). A line parallel to \(O O^{\prime}\) through \(Q\) determines \(Q\) and hence \(H^{\prime}\) : \(H\) and \(H^{\prime}\) are then also the nodal points. A ray parallel to the incident ray through \(H\) will emerge parallel to itself through \(H^{\prime}\). That determines \(F^{\prime}\). Similarly a ray parallel to the emergent ray through \(H\) determines \(F\).

\begin{figure}[h]
\begin{center}
  \includegraphics[alt={},max width=\textwidth]{da0156b6-4133-434e-90d4-ec01b2070afe-151_273_504_415_959}
\captionsetup{labelformat=empty}
\caption{(a)}
\end{center}
\end{figure}

\begin{figure}[h]
\begin{center}
  \includegraphics[alt={},max width=\textwidth]{da0156b6-4133-434e-90d4-ec01b2070afe-151_277_490_747_963}
\captionsetup{labelformat=empty}
\caption{(b)}
\end{center}
\end{figure}

\begin{figure}[h]
\begin{center}
  \includegraphics[alt={},max width=\textwidth]{da0156b6-4133-434e-90d4-ec01b2070afe-151_323_541_1068_942}
\captionsetup{labelformat=empty}
\caption{(C)}
\end{center}
\end{figure}

5.53 Here we do not assume that the media on the two sides of the system are the same.

\begin{figure}[h]
\begin{center}
  \includegraphics[alt={},max width=\textwidth]{da0156b6-4133-434e-90d4-ec01b2070afe-151_215_609_1494_93}
\captionsetup{labelformat=empty}
\caption{(a)}
\end{center}
\end{figure}

\begin{center}
\includegraphics[max width=\textwidth, alt={}]{da0156b6-4133-434e-90d4-ec01b2070afe-151_325_641_1730_122}
\end{center}

\begin{figure}[h]
\begin{center}
  \includegraphics[alt={},max width=\textwidth]{da0156b6-4133-434e-90d4-ec01b2070afe-151_226_616_1499_805}
\captionsetup{labelformat=empty}
\caption{(b)}
\end{center}
\end{figure}

\begin{figure}[h]
\begin{center}
  \includegraphics[alt={},max width=\textwidth]{da0156b6-4133-434e-90d4-ec01b2070afe-151_286_579_1741_824}
\captionsetup{labelformat=empty}
\caption{(d)}
\end{center}
\end{figure}

5.54 (a) Optical power of the system of combination of two lenses,

\[
\Phi=\Phi_{1}+\Phi_{2}-d \Phi_{1} \Phi_{2}
\]

on putting the values,\\
or,

\[
\begin{gathered}
\Phi=4 \mathrm{D} \\
f=\frac{1}{\Phi}=25 \mathrm{~cm}
\end{gathered}
\]

Now, the position of primary principal plane with respect to the vertex of converging lens,

\[
X=\frac{d \Phi_{2}}{\Phi}=10 \mathrm{~cm}
\]

Similarly, the distance of secondary principal plane with respect to the vertex of diverging lens.

\[
' X^{\prime}=-\frac{d \Phi_{1}}{\Phi}=-10 \mathrm{~cm}, \text { i.e. } 10 \mathrm{~cm} \text { left to it. }
\]

(b) The distance between the rear principal focal point \(F^{\prime}\) and the vertex of converging lens,\\
and

\[
l=d+\left(\frac{1}{\Phi}\right)\left(-d \Phi_{1}\right)=\frac{\Phi d}{\Phi}+\left(\frac{-d \Phi_{1}}{\Phi}\right)
\]

\[
\begin{aligned}
& f / l=\left(\frac{1}{\Phi}\right) / \frac{\Phi d}{\Phi}-\frac{d \Phi_{1}}{\Phi}, \text { as } f=\frac{1}{\Phi} \\
&=1 / d \Phi-d \Phi_{1} \\
&\left(\Phi_{1}+\Phi_{2}-d \Phi_{1} \Phi_{2}\right)-d \Phi_{1}=1 / d \Phi_{2}-d^{2} \Phi_{1} \Phi_{2}
\end{aligned}
\]

Now, if \(f / l\) is maximum for certain value of \(d\) then \(l / f\) will be minimum for the same value of \(d\). And for minimum \(l / f\),

\[
\begin{array}{lc} 
& d(l / f) / d d=\Phi_{2}-2 d \Phi_{1} \Phi_{2}=0 \\
\text { or, } & d=\Phi_{2} / 2 \Phi_{1} \Phi_{2} \\
\text { or, } & d=1 / 2 \Phi_{1}=5 \mathrm{~cm}
\end{array}
\]

So, the required maximum ratio of \(f / l=4 / 3\).\\
5.55 The optical power of first convex surface is,

\[
\Phi=\frac{P(n-1)}{R_{1}}=5 \mathrm{D}, \quad \text { as } R_{1}=10 \mathrm{~cm}
\]

and the optical power of second concave surface is,

\[
\Phi_{2}=\frac{(1-n)}{R_{2}}=-10 \mathrm{D}
\]

So, the optical power of the system,

\[
\Phi=\Phi_{1}+\Phi_{2}-\frac{d}{n} \Phi_{1} \Phi_{2}=-4 \mathrm{D}
\]

Now, the distance of the primary principal plane from the vertex of convex surface is given as,

\[
\begin{aligned}
x & =\left(\frac{1}{\Phi}\right)\left(\frac{d}{n}\right) \Phi_{2}, \text { here } n_{1}=1 \text { and } n_{2}=n \\
& =\frac{d \Phi_{2}}{\Phi n}=5 \mathrm{~cm}
\end{aligned}
\]

and the distance of secondary principal plane from the vertex of second concave surface,

\[
x^{\prime}=-\left(\frac{1}{\Phi}\right)\left(\frac{d}{n}\right) \Phi_{1}=-\frac{d \Phi_{1}}{\Phi n}=2.5 \mathrm{~cm}
\]

5.56 The optical power of the system of two thin lenses placed in air is given as,

\[
\Phi=\Phi_{1}+\Phi_{2}-d \Phi_{1} \Phi_{2}
\]

or, \(\quad \frac{1}{f}=\frac{1}{f_{1}}+\frac{1}{f_{2}}-\frac{d}{f_{1} f_{2}}\), where \(f\) is the equivalent focal length\\
So,


\begin{align*}
& \frac{1}{f}=\frac{f_{2}+f_{1}-d}{f_{1} f_{2}} \\
& f=\frac{f_{1} f_{2}}{f_{1}+f_{2}-d} \tag{1}
\end{align*}


This equivalent focal length of the system of two lenses is measured from the primary principal plane.\\
As clear from the figure, the distance of the primary principal plane from the optical centre of the first is\\
\(O_{1} H=x=+(n / \Phi)\left(d / n^{\prime}\right) \Phi_{1}\)\\
\(=\frac{d \Phi_{1}}{\Phi}\), as \(n=n^{\prime}=1\), for air .

\[
\begin{aligned}
& =\frac{d f}{f_{1}} \\
& =\left(\frac{d}{f_{1}}\right)\left(\frac{f_{1} f_{2}}{f_{1}+f_{2}-d}\right) \\
& =\frac{d f_{2}}{f_{1}+f_{2}-d}
\end{aligned}
\]

\begin{center}
\includegraphics[max width=\textwidth, alt={}]{da0156b6-4133-434e-90d4-ec01b2070afe-153_413_719_1239_676}
\end{center}

Now, if we place the equivalent lens at the primary principal plane of the lens system, it will provide the same transverse magnification as the system. So, the distance of equivalent lens from the vertex of the first lens is,

\[
x=\frac{d f_{2}}{f_{1}+f_{2}-d}
\]

5.57 The plane mirror forms the image of the lens, and water, filled in the space between the two, behind the mirror, as shown in the figure.

So, the whole optical system is equivalent to two similar lenses, seperated by a distance \(2 l\) and thus, the power of this system,\\
\(\Phi=\Phi_{1}+\Phi_{2}-\frac{d \Phi_{1} \Phi_{2}}{n_{0}}\), where \(\Phi_{1}=\Phi_{2}=\Phi_{1}^{\prime}\)\\
\(=\) optical power of individual lens and \(n_{0}=\) R.I. of water.\\
\includegraphics[max width=\textwidth, alt={}, center]{da0156b6-4133-434e-90d4-ec01b2070afe-154_381_503_105_963}

Now, \(\Phi^{\prime}=\) optical power of first convex surface + optical power of second concave surface.\\
\(=\frac{(n-1)}{R}+\frac{n_{0}-n}{R}, n\) is the refractive index of glass.


\begin{equation*}
\frac{\left(2 n-n_{0}-1\right)}{R} \tag{1}
\end{equation*}


and so, the optical power of whole system,

\[
\Phi=2 \Phi^{\prime}-\frac{2 d \Phi^{\prime 2}}{n_{0}}=3.0 \mathrm{D} \text {, substituting the values. }
\]

5.58 (a) A telescope in normal adjustment is a zero power conbination of lenses. Thus we require

\[
\Phi=O=\Phi_{1}+\Phi_{2}-\frac{d}{n} \Phi_{1} \Phi_{2}
\]

But \(\Phi_{1}=\) Power of the convex surface \(=\frac{n-1}{R_{0}+\Delta R}\)\\
\(\Phi_{2}=\) Power of the concave surface \(=-\frac{n-1}{R_{0}}\)\\
Thus,

\[
O=\frac{(n-1) \Delta R}{R_{0}\left(R_{0}+\Delta R\right)}+\frac{d}{n} \frac{(n-1)^{2}}{R_{0}\left(R_{0}+\Delta R\right)}
\]

So

\[
d=\frac{n \Delta R}{n-1}=4.5 \mathrm{~cm} . \text { on putting the values. }
\]

(b) Here, \(\Phi=-1=\frac{.5}{.1}-\frac{.5}{.075}+\frac{d}{1.5} \times \frac{.5 \times .5}{.1 \times .075}\)

\[
\begin{aligned}
& =5-\frac{20}{3}+\frac{d \times 2}{3} \times \frac{5 \times 20}{3}=-\frac{5}{3}+\frac{200^{\prime} d}{9} \\
& =\frac{200 d}{9}=\frac{2}{3} \text { or } d=(3 / 100) \mathrm{m}=3 \mathrm{~cm}
\end{aligned}
\]

5.59 (a) The power of the lens is (as in the previous problem)

\[
\Phi=\frac{n-1}{R}-\frac{n-1}{R}-\frac{d}{n}\left(\frac{n-1}{R}\right)\left(-\frac{n-1}{R}\right)=\frac{d(n-1)^{2}}{n R^{2}}>0 .
\]

The principal planes are located on the side of the convex surface at a distance \(d\) from each other, with the front principal plane being removed from the convex surface of the lens by a distance \(R /(n-1)\).\\
(b) Here \(\Phi=-\frac{n-1}{R_{1}}+\frac{n+1}{R_{2}}+\frac{R_{2}-R_{1}}{n} \frac{(n-1)^{2}}{R_{1} R_{2}}\)

\[
\begin{aligned}
& =\frac{(n-1)\left(R_{2}-R_{1}\right)}{R_{2} R_{1}}\left[-1+\frac{n-1}{n}\right] \\
& =-\frac{n-1}{n}\left(\frac{1}{R_{1}}-\frac{1}{R_{2}}\right)<0
\end{aligned}
\]

\begin{center}
\includegraphics[max width=\textwidth, alt={}]{da0156b6-4133-434e-90d4-ec01b2070afe-155_263_475_267_904}
\end{center}

Both principal planes pass through the common centre\\
of curvature of the surfaces of the lens.\\
5.60 Let the optical powers of the first and second surfaces of the ball of radius \(R_{1}\) be \(\Phi_{1}^{\prime}\) and \(\Phi_{1}^{\prime \prime}\), then

\[
\Phi_{1}^{\prime}=(n-1) / R_{1} \quad \text { and } \quad \Phi_{1}^{\prime \prime}=(1-n) /\left(-R_{1}\right)=\frac{(n-1)}{R_{1}}
\]

This ball may be treated as a thick spherical lens of thickness \(2 R_{1}\). So the optical power of this sphere is,


\begin{equation*}
\Phi=\Phi_{1}^{\prime}-\frac{2 R_{1} \Phi_{1}^{\prime} \Phi_{1}^{\prime \prime}}{n}=2(n-1) / n R_{1} \tag{1}
\end{equation*}


Similarly, the optical power of second ball,

\[
\Phi_{2}=2(n-1) / n R_{2}
\]

If the distance between the centres of these balls be \(d\). Then the optical power of whole system,

\[
\begin{gathered}
\Phi=\Phi_{1}+\Phi_{2}-d \Phi_{1} \Phi_{2} \\
=\frac{2(n-1)}{n R_{1}}+\frac{2(n-1)}{n R_{2}}-\frac{4 d(n-1)^{2}}{n^{2} R_{1} R_{2}} \\
=\frac{2(n-1)}{n R_{1} R_{2}}\left[\left(R_{1}+R_{2}\right)-\frac{2 d(n-1)}{n}\right]
\end{gathered}
\]

Now, since this system serves as telescope, the optical power of the system must be equal to zero.\\
or,

\[
\begin{aligned}
\left(R_{1}+R_{2}\right) & =\frac{2 d(n-1)}{n}, \quad \text { as } \quad \frac{2(n-1)}{n R_{1} R_{2}} \neq 0 \\
d & =\frac{n\left(R_{1}+R_{2}\right)}{2(n-1)}=9 \mathrm{~cm}
\end{aligned}
\]

Since the diameter \(D\) of the objective is \(2 R_{1}\) and that of the eye-piece is \(d=2 R_{2}\)\\
So, the magnification,

\[
\Gamma=D / d=\frac{2 R_{1}}{2 R_{2}}=R_{1} R_{2}=5
\]

5.61 Optical powers of the two surfaces of the lens are

\[
\Phi_{1}=(n-1) / R \text { and } \Phi_{2}=(1-n) /-R=\frac{n-1}{R}
\]

So, the power of the lens of thickness \(d\),

\[
\Phi^{\prime}=\Phi_{1}+\Phi_{2}-\frac{d \Phi_{1} \Phi_{2}}{n}=\frac{n-1}{R}+\frac{n-1}{R}-\frac{d(n-1)^{2} / R^{2}}{n^{2}}=\frac{n^{2}-1}{n R}
\]

and optical power of the combination of these two thick lenses,

\[
\Phi=\Phi^{\prime}+\Phi^{\prime}=2 \Phi^{\prime}=\frac{2\left(n^{2}-1\right)}{n R}
\]

so, power of this system in air is, \(\Phi_{0}=\frac{\Phi}{n}=\frac{2\left(n^{2}-1\right)}{n^{2} R}=37 \mathrm{D}\).\\
5.62 We consider a ray \(Q P R\) in a medium of gradually varying refractive index \(n\). At \(P\), the gradient of \(n\) is a vector with the given direction while is nearly the same at neighbouring points \(Q, R\). The arc length \(Q R\) is \(d s\). We apply Snell's formula \(n \sin \theta=\) constant where \(\theta\) is to be measured from the direction \(\nabla n\). The refractive indices at \(Q, R\) whose mid point is \(P\) are

\[
\eta \pm \frac{1}{2}|\nabla \eta| d \theta \cos \theta
\]

so

\[
\left(\eta-\frac{1}{2}|\nabla n| d \theta \cos \theta\right)\left(\sin \theta+\frac{1}{2} \cos \theta d \theta\right)
\]

\(=\left(\eta+\frac{1}{2}|\nabla n| d \theta \cos \theta\right)\left(\sin \theta-\frac{1}{2} \cos \theta d \theta\right)\) or \(n \cos \theta d \theta=|\nabla n| d s \cos \theta \sin \theta\)\\
(we have used here \(\sin \left(\theta \pm \frac{1}{2} d \theta\right)=\sin \theta \pm \frac{1}{2} \cos \theta d \theta\) )\\
Now using the definition of the radius of curvature \(\frac{1}{\rho}=\frac{d \theta}{d s}\)

\[
\frac{1}{\rho}=\frac{1}{\eta}|\nabla n| \sin \theta
\]

The quantity \(|\nabla n| \sin \theta\) can be called \(\frac{\delta n}{\delta N}\) i.e. the derivative of \(n\) along the normal \(N\) to the ray. Then

\[
\frac{1}{\rho}=\frac{\delta}{\delta N} \ln n .
\]

5.63 From the above problem\\
\(\frac{1}{\rho}=\frac{1}{n} \hat{p} \cdot \vec{\nabla} n \approx \hat{p} \cdot \nabla n \approx|\nabla n|=3 \times 10^{-8} \mathrm{~m}^{-1}\)\\
( since \(\hat{p} \| \vec{\nabla} n\) both being vertical ). So \(\rho=3.3 \times 10^{7} \mathrm{~m}\)\\
For the ray of light to propagate all the way round the earth we must have\\
\includegraphics[max width=\textwidth, alt={}, center]{da0156b6-4133-434e-90d4-ec01b2070afe-156_199_363_1622_1037}\\
\(\rho=R=6400 \mathrm{~km}=6.4 \times 10^{6} \mathrm{~m}\)\\
Thus \(|\nabla n|=1.6 \times 10^{-7} \mathrm{~m}^{-1}\)

\subsection*{5.2 INTERFERENCE OF LIGHT}
5.64 (a) In this case the net vibration is given by

\[
x=a_{1} \cos \omega t+a_{2} \cos (\omega t+\delta)
\]

where \(\delta\) is the phase difference between the two vibrations which varies rapidly and randomly in the interval \((0,2 \pi)\). (This is what is meant by incoherence.)\\
Then

\[
x=\left(a_{1}+a_{2} \cos \delta\right) \cos \omega t+a_{2} \sin \delta \sin \omega t
\]

The total energy will be taken to be proportional to the time average of the square of the displacement.

Thus

\[
E=\left\langle\left(a_{1}+a_{2} \cos \delta\right)^{2}+a_{2}^{2} \sin ^{2} \delta\right\rangle=a_{1}^{2}+a_{2}^{2}
\]

as \(\langle\cos \delta\rangle=0\) and we have put \(\left\langle\cos ^{2} \omega t\right\rangle=\left\langle\sin ^{2} \omega t\right\rangle=\frac{1}{2}\) and has been absorbed in the overall constant of proportionality.

In the same units the energies of the two oscillations are \(\boldsymbol{a}_{1}^{2}\) and \(\boldsymbol{a}_{2}^{2}\) respectively so the proposition is proved.\\
(b) Here \(\vec{r}=a_{1} \cos \omega t \hat{i}+a_{2} \cos (\omega t+\delta) \hat{j}\)\\
and the mean square displacement is \(\boldsymbol{\alpha} \boldsymbol{a}_{\mathbf{1}}^{\mathbf{2}}+\boldsymbol{a}_{\mathbf{2}}^{\mathbf{2}}\)\\
if \(\delta\) is fixed but arbitrary. Then as in (a) we see that \(E=E_{1}+E_{2}\).\\
5.65 It is easier to do it analytically.

\[
\begin{gathered}
\xi_{1}=a \cos \omega t, \xi_{2}=2 a \sin \omega t \\
\xi_{3}=\frac{3}{2} a\left(\cos \frac{\pi}{3} \cos \omega t-\sin \frac{\pi}{3} \sin \omega t\right)
\end{gathered}
\]

\[
\begin{aligned}
& \text { Resultant vibration is } \\
& \qquad \xi=\frac{7 a}{4} \cos \omega t+a\left(2-\frac{3 \sqrt{3}}{4}\right) \sin \omega t \\
& \text { This has an amplitude }=\frac{a}{4} \sqrt{49+(8-3 \sqrt{3})^{2}}=1.89 a
\end{aligned}
\]

5.66 We use the method of complex amplitudes. Then the amplitudes are

\[
A_{1}=a, A_{2}=a e^{i \varphi}, \ldots A_{N}=a e^{i(N-1) \varphi}
\]

and the resultant complex amplitude is

\[
\begin{aligned}
A & =A_{1}+A_{2}+\ldots+A_{N}=a\left(1+e^{i \varphi}+e^{2 i \varphi}+\ldots+e^{i(N-1) \varphi}\right) \\
& =a \frac{1-e^{i N \varphi}}{1-e^{i \varphi}}
\end{aligned}
\]

The corresponding ordinary amplitude is

\[
\begin{aligned}
|A| & =a\left|\frac{1-e^{i N \varphi}}{1-e^{i \varphi}}\right|=a\left[\frac{1-e^{i N \varphi}}{1-e^{i \varphi}} \times \frac{1-e^{-i N \varphi}}{1-e^{-i \varphi}}\right]^{1 / 2} \\
& =a\left[\frac{2-2 \cos N \varphi}{2-2 \cos \varphi}\right]^{1 / 2}=a \frac{\sin \frac{N \varphi}{2}}{\sin \frac{\varphi}{2}} .
\end{aligned}
\]

5.67 (a) With dipole moment \(\perp^{r}\) to plane there is no variation with \(\theta\) of individual radiation amplitude. Then the intensity variation is due to interference only.

In the direction given by angle \(\theta\) the phase difference is\\
\(\frac{2 \pi}{\lambda}(d \cos \theta)+\varphi=2 k \pi \quad\) for maxima\\
Thus \(\quad d \cos \theta=\left(k-\frac{\varphi}{2 \pi}\right) \lambda\)\\
\includegraphics[max width=\textwidth, alt={}, center]{da0156b6-4133-434e-90d4-ec01b2070afe-158_497_480_462_983}

\[
k=0, \pm 1, \pm 2, \ldots
\]

We have added \(\varphi\) to \(\frac{2 \pi}{\lambda} d \cos \theta\) because the extra path that the wave from 2 has to travel in going to \(P\) (as compared to 1 ) makes it lag more than it already is (due to \(\varphi\) ).\\
(b) Maximum for \(\theta=\pi\) gives \(-d=\left(k-\frac{\varphi}{2 \pi}\right) \lambda\)

Minimum for \(\theta=0\) gives \(d=\left(k^{\prime}-\frac{\varphi}{2 \pi}+\frac{1}{2}\right) \lambda\)\\
Adding we get

\[
\left(k+k^{\prime}-\frac{\varphi}{\pi}+\frac{1}{2}\right) \lambda=0
\]

This can be true only if

\[
k^{\prime}=-k, \varphi=\frac{\pi}{2}
\]

since \(0<\varphi<\pi\)\\
Then

\[
-\dot{d}=\left(k-\frac{1}{4}\right) \lambda
\]

Here

\[
k=0,-1,-2,-3, \ldots
\]

(Otherwise R.H.S. will become \(+v e\) ).\\
Putting \(k=-\bar{k}, \bar{k}=0,+1,+2,+3, \ldots\)

\[
d=\left(\bar{k}+\frac{1}{4}\right) \lambda .
\]

5.68 If \(\Delta \varphi\) is the phase difference between neighbouring radiators then for a maximum in the direction \(\theta\) we must have\\
\(\frac{2 \pi}{\lambda} d \cos \theta+\Delta \varphi=2 \pi k\)\\
\includegraphics[max width=\textwidth, alt={}, center]{da0156b6-4133-434e-90d4-ec01b2070afe-159_184_786_227_346}

For scanning \(\theta=\omega t+\beta\)\\
Thus

\[
\frac{d}{\lambda} \cos (\omega t+\beta)+\frac{\Delta \varphi}{2 \pi}=k
\]

or

\[
\Delta \varphi=2 \pi\left[k-\frac{d}{\lambda} \cos (\omega t+\beta)\right]
\]

To get the answer of the book, put \(\beta=\alpha-\pi / 2\).\\
5.69 From the general formula\\
we find that

\[
\begin{gathered}
\Delta x=\frac{l \lambda}{d} \\
\frac{\Delta x}{\eta}=\frac{l \lambda}{d+2 \Delta h}
\end{gathered}
\]

since \(d\) increases to \(d+2 \Delta h\) when the source is moved away from the mirror plane by \(\Delta h\).\\
Thus

\[
\eta d=d+2 \Delta h \quad \text { or } \quad d=2 \Delta h /(\eta-1)
\]

Finally

\[
\lambda=\frac{2 \Delta h \Delta x}{(\eta-1) l}=0.6 \mu \mathrm{~m} .
\]

5.70 We can think of the two coherent plane waves as emitted from two coherent point sources very far away. Then

But

\[
\Delta x=\frac{l \lambda}{d}=\frac{\lambda}{d / l}
\]

so

\[
\frac{d}{l}=\psi(\text { if } \psi<1)
\]

\[
\Delta x=\frac{\lambda}{\psi} .
\]

5.71 (a) Here \(S^{\prime} S^{\prime \prime}=d=2 r \alpha\)

Then \(\Delta x=\frac{(b+r) \lambda}{2 \alpha}\)\\
Putting \(b=1.3\) metre , \(r=1\) metre\\
\(\lambda=0.55 \mu \mathrm{~m}, \alpha=12^{\prime}=\frac{1}{5 \times 57}\) radian\\
we get \(\boldsymbol{\Delta} \boldsymbol{x}=\mathbf{1} \cdot \mathbf{1} \mathrm{mm}\)\\
Number of possible maxima \(=\frac{2 b \alpha}{\Delta x}+1 \approx 8.3+1 \approx 9\)\\
( \(2 b \alpha\) is the length of the spot on the screen which gets light after reflection from both mirror. We add 1 above to take account of the fact that in a distance \(\Delta x\) there are two maxima).\\
(b) When the silt moves by \(\delta l\) along the arc of radius \(r\) the incident ray on the mirror rotates by \(\frac{\delta l}{r}\); this is the also the rotation of the reflected ray. There is then a shift of the fringe of magnitude.

\[
b \frac{\delta l}{r}=13 \mathrm{~mm} .
\]

(c) If the width of the slit is \(\delta\) then we can imagine the slit to consist of two narrow slits with separation \(\delta\). The fringe pattern due to the wide slit is the superposition of the pattern due to thése two narrow slits. The full pattern will not be sharp at all if the pattern due to the two narrow slits are \(\frac{1}{2} \Delta x\) apart because then the maxima due to one will fill the minima due to the other. Thus we demand that

\[
\frac{b \delta_{\max }}{r}=\frac{1}{2} \Delta x=\frac{(b+r) \lambda}{4 r \alpha}
\]

or

\[
\delta_{\max }=\left(1+\frac{r}{b}\right) \frac{\lambda}{4 \alpha}=42 \mu \mathrm{~m} .
\]

5.72 To get this case we must let \(r \rightarrow \infty\) in the formula for \(\Delta x\) of the last example.

So

\[
\Delta x=\frac{(b+r) \lambda}{2 \alpha r} \rightarrow \frac{\lambda}{2 \alpha} .
\]

(A plane wave is like light emitted from a point source at \(\infty\) ).\\
Then

\[
\lambda=2 \alpha \Delta x=0.64 \mu \mathrm{~m} .
\]

5.73\\
\includegraphics[max width=\textwidth, alt={}, center]{da0156b6-4133-434e-90d4-ec01b2070afe-160_329_565_1369_519}\\
(a) We show the upper half or the lens. The emergent light is at an angle \(\frac{a}{2 f}\) from the axis.

Thus the divergence angle of the two incident light beams is

\[
\psi=\frac{a}{f}
\]

When they interfere the fringes produced have a width

\[
\Delta x=\frac{\lambda}{\psi}=\frac{f \lambda}{a}=0.15 \mathrm{~mm} .
\]

The patch on the screen illuminated by both light has a width \(b \psi\) and this contains

\[
\frac{b \Psi}{\Delta x}=\frac{b a^{2}}{f^{2} \lambda} \text { fringes }=13 \text { fringes }
\]

(if we ignore 1 in comparis on to \(\frac{b \psi}{\Delta x}\) (if 5.71 (a))\\
(b) We follow the logic of \((5.71 \mathrm{c})\). From one edge of the slit to the other edge the distance is of magnitude \(\delta\) (i.e. \(\frac{a}{2}\) to \(\frac{a}{2}+\delta\) ).\\
If we imagine the edge to shift by this distance, the angle \(\psi / 2\) will increase by \(\frac{\Delta \psi}{2}=\frac{\delta}{2 f}\) and the light will shift \(\pm b \frac{\delta}{2 f}\)\\
The fringe pattern will therefore shift by \(\frac{\delta \cdot b}{f}\)\\
Equating this to \(\frac{\Delta x}{2}=\frac{f \lambda}{2 a}\) we get \(\delta_{\max }=\frac{f^{2} \lambda}{2 a b}=37.5 \mu \mathrm{~m}\).\\
\(5.74 \Delta x=\frac{l \lambda}{d}\)

\[
\begin{aligned}
l & =a+b \\
d & =2(n-1) \theta a \\
\delta & =(n-1) \theta \\
d & =2 \delta \cdot a \\
n & =\text { R.I. of glass }
\end{aligned}
\]

\begin{center}
\includegraphics[max width=\textwidth, alt={}]{da0156b6-4133-434e-90d4-ec01b2070afe-161_276_781_995_482}
\end{center}

Thus

\[
\lambda=\frac{2(n-1) \theta a \Delta x}{a+b} .=0.64 \mu \mathrm{~m} .
\]

5.75 It will be assumed that the space between the biprism and the glass plate filled with benzene constitutes complementary prisms as shown.

Then the two prisms being oppositely placed, the net deviation produced by them is\\
\(\delta=(n-1) \theta-\left(n^{\prime}-1\right) \theta=\left(n-n^{\prime}\right) \theta\)\\
Hence as in the previous problem

\[
d=2 a \delta=2 a \theta\left(n-n^{\prime}\right)
\]

So

\[
\Delta x=\frac{(a+b) \lambda}{2 a \theta\left(n-n^{\prime}\right)}
\]

\begin{center}
\includegraphics[max width=\textwidth, alt={}]{da0156b6-4133-434e-90d4-ec01b2070afe-161_504_322_1461_970}
\end{center}

For plane incident wave we let \(a \rightarrow \infty\)\\
so

\[
\Delta x=\frac{\lambda}{2 \theta\left(n-n^{\prime}\right)}=0.2 \mathrm{~mm}
\]

5.76 Extra phase difference introduced by the glass plate is

\[
\frac{2 \pi}{\lambda}(n-1) h
\]

This will cause a shift equal to \((n-1) \frac{h}{\lambda}\) fringe widths\\
i.e. by

\[
(n-1) \frac{h}{\lambda} \times \frac{l \lambda}{d}=\frac{(n-1) h l}{d}=2 \mathrm{~mm} .
\]

The fringes move down if the lower slit is covered by the plate to compensate for the extra phase shift introduced by the plate.\\
5.77 No. of fringes shifted \(=\left(n^{\prime}-n\right) \frac{l}{\lambda}=N\)\\
so

\[
n^{\prime}=n+\frac{N \lambda}{l}=1.000377
\]

5.78 (a) Suppose the vector \(\vec{E}, \overrightarrow{E^{\prime}}, \overrightarrow{E^{\prime \prime}}\) correspond to the incident, reflected and the transmitted wave. Due to the continuity of the tangential component of the electric field across the interface, it follows that


\begin{equation*}
E_{\tau}+E_{\tau}^{\prime}=E_{\tau}^{\prime \prime} \tag{1}
\end{equation*}


where the subscript \(\tau\) means tangential.\\
The energy flux density is \(\vec{E} \times \vec{H}=\vec{S}\).\\
Since

\[
\begin{gathered}
H \sqrt{\mu \mu_{0}}=E \sqrt{\varepsilon \varepsilon_{0}} \\
H=E \sqrt{\frac{\varepsilon_{0}}{\mu_{0}}} \sqrt{\varepsilon / \mu}=n \sqrt{\frac{\varepsilon_{0}}{\mu_{0}}} E
\end{gathered}
\]

Now \(S_{\sim} n E^{2}\) and since the light is incident normally

\[
\begin{array}{lc} 
& n_{1} E_{\tau}^{2}=n_{1} E_{\tau}^{\prime 2}+n_{2} E_{\tau}^{\prime \prime 2} \\
\text { or } & n_{1}\left(E_{\tau}^{2}-E_{\tau}^{\prime 2}\right)=n_{2} E_{\tau}^{\prime \prime 2} \\
\text { so } & n_{1}\left(E_{\tau}-E_{\tau}^{\prime}\right)=n_{2} E_{\tau}^{\prime \prime}  \tag{3}\\
\text { so } & E_{\tau}^{\prime \prime}=\frac{.2 n_{1}}{n_{1}+n_{2}} E_{\tau}
\end{array}
\]

Since \(E_{\tau}{ }^{\prime \prime}\) and \(E_{\tau}\) have the same sign, there is no phase change involved in this case.\\
(b) From (1) \& (3)\\
or

\[
\begin{gathered}
\left(n_{2}+n_{1}\right) E_{\tau}^{\prime}+\left(n_{2}-n_{1}\right) E_{\tau}=0 \\
E_{\tau}^{\prime}=\frac{n_{1}-n_{2}}{n_{1}+n_{2}} E_{\tau} .
\end{gathered}
\]

If \(n_{2}>n_{1}\), then \(E_{\tau}{ }^{\prime} \& E_{\tau}\) have opposite signs. Thus the reflected wave has an abrupt change of phase by \(\pi\) if \(n_{2}>n_{1}\) i.e. on reflection from the interface between two media when light is incident from the rater to denser medium.\\
5.79\\
\includegraphics[max width=\textwidth, alt={}, center]{da0156b6-4133-434e-90d4-ec01b2070afe-163_443_681_240_415}

Path difference between (1) \& (2) is

\[
\begin{aligned}
& 2 n d \sec \theta_{2}-2 d \tan \theta_{2} \sin \theta_{1} \\
= & 2 d \frac{n-\frac{\sin ^{2} \theta_{1}}{n}}{\sqrt{1-\frac{\sin ^{2} \theta_{1}}{n^{2}}}}=2 d \sqrt{n^{2}-\sin ^{2} \theta_{1}}
\end{aligned}
\]

For bright fringes this must equal \(\left(k+\frac{1}{2}\right) \lambda\) where \(\frac{1}{2}\) comes from the phase cnange of \(\pi\) for (1).

Here

\[
k=0,1,2, \ldots
\]

Thus

\[
4 d \sqrt{n^{2}-\sin ^{2} \theta_{1}}=(2 k+1) \lambda
\]

or

\[
d=\frac{\lambda(1+2 k)}{4 \sqrt{n^{2}-\sin ^{2} \theta_{1}}}=0.14(1+2 k) \mu \mathrm{m} . .
\]

5.80 Given

\[
\begin{aligned}
& 2 d \sqrt{n^{2}-1 / 4}=\left(k+\frac{1}{2}\right) \times 0.64 \mu \mathrm{~m} \text { ( bright fringe) } \\
& 2 d \sqrt{n^{2}-1 / 4}=k \times 0.40 \mu \mathrm{~m} \quad \text { (dark fringe) }
\end{aligned}
\]

where \(k, k^{\prime}\) are integers.\\
Thus

\[
64\left(k+\frac{1}{2}\right)=40 k^{\prime} \text { or } 4(2 k+1)=5 k^{\prime}
\]

This means, for the smallest integer solutions

\[
k=2, k^{\prime}=4
\]

Hence

\[
d=\frac{4 \times 0.40}{2 \sqrt{n^{2}-1 / 4}}=0.65 \mu \mathrm{~m} .
\]

5.81 When the glass surface is coated with a material of R.I. \(n^{\prime}=\sqrt{n}\) ( \(n=\) R.I. of glass) of appropriate thickness, reflection is zero because of interference between various multiply reflected waves. We show this below.

Let a wave of unit amplitude be normally incident from the left. The reflected amplitude is \(-r\) where

\[
r=\frac{\sqrt{n}-1}{\sqrt{n}+1}
\]

Its phase is -ve so we write the reflected wave as \(-r\). The transmitted wave has amplitude \(t\)

\[
t=\frac{2}{1+\sqrt{n}}
\]

This wave is reflected at the second face and has amplitude\\
-tr

\[
\left(\text { because } \frac{n-\sqrt{n}}{n+\sqrt{n}}=\frac{\sqrt{n}-1}{\sqrt{n}+1} .\right)
\]

\begin{center}
\includegraphics[max width=\textwidth, alt={}]{da0156b6-4133-434e-90d4-ec01b2070afe-164_483_534_343_920}
\end{center}

The emergent wave has amplitude \(-t t^{\prime} r\).\\
We prove below that \(-t t^{\prime}=1-r^{2}\). There is also a reflected part of emplitude \(t r r^{\prime}=-t r^{2}\), where \(r^{\prime}\) is the reflection coefficient for a ray incident from the coating towards air. After reflection from the second face a wave of amplitude

\[
+t t^{\prime} r^{3}=+\left(1-r^{2}\right) r^{3}
\]

emerges. Let \(\delta\) be the phase of the wave after traversing the coating both ways.\\
Then the complete reflected wave is

\[
\begin{gathered}
-r-\left(1-r^{2}\right) r e^{i \delta}+\left(1-r^{2}\right) r^{3} e^{2 i \delta} \\
-\left(1-r^{2}\right) r^{5} e^{3 i \delta} \cdots \cdots \\
=-r-\left(1-r^{2}\right) r e^{i \delta} \frac{1}{1+r^{2} e^{i \delta}} \\
=-r\left[1+r^{2} e^{i \delta}+\left(1-r^{2}\right) e^{i \delta}\right] \frac{1}{1+r^{2} e^{i \delta}} \\
=-r \frac{1+e^{i \delta}}{1+r^{2} e^{i \delta}}
\end{gathered}
\]

This vanishes if \(\delta=(2 k+1) \pi\). But

\[
\begin{aligned}
\delta & =\frac{2 \pi}{\lambda} 2 \sqrt{n} d \text { so } \\
d & =\frac{\lambda}{4 \sqrt{n}}(2 k+1)
\end{aligned}
\]

We now deduce \(t t^{\prime}=1-r^{2}\) and \(r^{\prime}=+r\). This follows from the principle of reversibility of light path as shown in the figure below.

\[
\begin{aligned}
& \quad t t^{\prime}+r^{2}=1 \\
& -r t+r^{\prime} t=0 \\
& \therefore \quad t t^{\prime}=1-r^{2} \\
& r^{\prime}=+r .
\end{aligned}
\]

( \(-r\) is the reflection ratio for the wave entering a denser medium ).\\
\includegraphics[max width=\textwidth, alt={}, center]{da0156b6-4133-434e-90d4-ec01b2070afe-165_471_459_80_977}\\
5.82 We have the condition for maxima

\[
2 d \sqrt{n^{2}-\sin ^{2} \theta_{1}}=\left(k+\frac{1}{2}\right) \lambda
\]

This must hold for angle \(\theta \pm \frac{\delta \theta}{2}\) with successive values of \(k\). Thus

Thus

\[
\begin{gathered}
2 d \sqrt{n^{2}-\sin ^{2}\left(\theta+\frac{\delta \theta}{2}\right)}=\left(k-\frac{1}{2}\right) \lambda \\
2 d \sqrt{n^{2}-\sin ^{2}\left(\theta-\frac{\delta \theta}{2}\right)}=\left(k+\frac{1}{2}\right) \lambda \\
\lambda=2 d\left\{\sqrt{n^{2}-\sin ^{2} \theta+\delta \theta \sin \theta \cos \theta}\right. \\
\left.-\sqrt{n^{2}-\sin ^{2} \theta-\delta \theta \sin \theta \cos \theta}\right\} \\
=2 d \frac{\delta \theta \sin \theta \cos \theta}{\sqrt{n^{2}-\sin ^{2} \theta}} \\
d=\frac{\sqrt{n^{2}-\sin ^{2} \theta \lambda}}{\sin 2 \theta \delta \theta}=15 \cdot 2 \mu \mathrm{~m}
\end{gathered}
\]

Thus\\
5.83 For small angles \(\theta\) we write for dark fringes

\[
2 d \sqrt{n^{2}-\sin ^{2} \theta}=2 d\left(n-\frac{\sin ^{2} \theta}{2 n}\right)=(k+0) \lambda
\]

For the first dark fringe \(\theta \cong 0\) and

\[
2 d n=\left(k_{0}+0\right) \lambda
\]

For the \(i^{\text {th }}\) dark fringe\\
or

\[
\begin{gathered}
2 d\left(n-\frac{\sin ^{2} \theta_{i}}{2 n}\right)=\left(k_{0}-i+1\right) \lambda \\
\sin ^{2} \theta_{i}=\frac{n \lambda}{d}(i-1)=\frac{r_{i}^{2}}{4 l^{2}}
\end{gathered}
\]

\begin{center}
\includegraphics[max width=\textwidth, alt={}]{da0156b6-4133-434e-90d4-ec01b2070afe-166_348_856_57_298}
\end{center}

Finally

\[
\frac{n \lambda}{d}(i-k)=\frac{r_{i}^{2}-r_{k}^{2}}{4 l^{2}}
\]

so

\[
\lambda=\frac{d\left(r_{i}^{2}-r_{k}^{2}\right)}{4 l^{2} n(i-k)}
\]

5.84 We have the usual equation for maxima

\[
2 h_{k} \propto \sqrt{n^{2}-\sin ^{2} \theta_{1}}=\left(k+\frac{1}{2}\right) \lambda
\]

Here \(\boldsymbol{h}_{\boldsymbol{k}}=\) distance of the fringe from top\\
\(h_{k} \alpha \cong d_{k}=\) thickness of the film\\
Thus on the screen placed at right angles to the reflected light

\[
\begin{aligned}
\Delta x & =\left(h_{k}-h_{k-1}\right) \cos \theta_{1} \\
& =\frac{\lambda \cos \theta_{1}}{2 \alpha \sqrt{n^{2}-\sin ^{2} \theta_{1}}}
\end{aligned}
\]

5.85 (a) For normal incidence we have using the above formula\\
\includegraphics[max width=\textwidth, alt={}, center]{da0156b6-4133-434e-90d4-ec01b2070afe-166_528_495_729_947}\\
so\\
(b) In a distance \(l\) on the wedge there are \(N=\frac{l}{\Delta x}\) fringes.

If the fringes disappear there, it must be due to the fact that the maxima due to the component of wavelength \(\lambda\) coincide with the mainima due to the component of wavelength \(\lambda+\Delta \lambda\). Thus\\
so

\[
\begin{gathered}
\Delta x=\frac{\lambda}{2 n \alpha} \\
\alpha=\frac{\lambda}{2 n \Delta x}=3^{\prime} \text { on putting the values }
\end{gathered}
\]

\[
\begin{gathered}
N \dot{\lambda}=\left(N-\frac{1}{2}\right)(\lambda+\Delta \lambda) \text { or } \Delta \lambda=\frac{\lambda}{2 N} \\
\frac{\Delta \lambda}{\lambda}=\frac{1}{2 N}=\frac{\Delta x}{2 l}=\frac{0.21}{30}=0.007 .
\end{gathered}
\]

The answer given in the book is off by a factor 2 .\\
5.86 We have

\[
r^{2}=\frac{1}{2} k \lambda R
\]

So for \(k\) differing by \(1(\Delta k=1)\)\\
or

\[
\begin{gathered}
2 r \Delta r=\frac{1}{2} \Delta k \lambda R=\frac{1}{2} \lambda R \\
\Delta r=\frac{\lambda R}{4 r} .
\end{gathered}
\]

5.87 The path traveresed in air film of the wave constituting the \(k^{\text {th }}\) ring is

\[
\frac{r^{2}}{R}=\frac{1}{2} k \lambda
\]

when the lens is moved a distance \(\Delta h\) the ring radius changes to \(r^{\prime}\) and the path length becomes

Thus

\[
\begin{gathered}
\frac{r^{\prime 2}}{R}+2 \Delta h=\frac{1}{2} k \lambda \\
r^{\prime}=\sqrt{r^{2}-2 R \Delta h}=1.5 \mathrm{~mm} .
\end{gathered}
\]

5.88 In this case the path difference is \(\frac{r^{2}-r_{0}^{2}}{R}\) for \(r>r_{0}\) and zero for \(r \leq r_{0}\). This must equals \((k-1 / 2) \lambda\) ( where \(k=6\) for the six \({ }^{\text {th }}\) bright ring.)\\
Thus \(r=\sqrt{r_{0}{ }^{2}+\left(k-\frac{1}{2}\right) \lambda R}=3.8 \mathrm{~mm}\)\\
5.89 From the formula for Newton's rings we derive for dark rings\\
so

\[
\begin{gathered}
\frac{d_{1}^{2}}{4}=k_{1} R \lambda, \frac{d_{2}^{2}}{4}=k_{2} R \lambda \\
\frac{d_{2}^{2}-d_{1}^{2}}{4\left(k_{2}-k_{1}\right) R}=\lambda
\end{gathered}
\]

Substituting the values, \(\lambda=0.5 \mu \mathrm{~m}\).\\
5.90 Path difference between waves reflected by the two convex surfaces is

\[
r^{2}\left(\frac{1}{R_{1}}+\frac{1}{R_{2}}\right)
\]

Taking account of the phase change at the \(2^{\text {nd }}\) surface we write the condition of bright rings as

\[
r^{2}\left(\frac{1}{R_{1}}+\frac{1}{R_{2}}\right)=\frac{2 k+1}{2} \lambda
\]

\includegraphics[max width=\textwidth, alt={}, center]{da0156b6-4133-434e-90d4-ec01b2070afe-167_344_215_1528_1023}\\
\(k=4\) for the fifth bright ring.

Thus

\[
\frac{1}{R_{1}}+\frac{1}{R_{2}}=\frac{9}{2} \lambda \cdot \frac{4}{d^{2}}=\frac{18 \lambda}{d^{2}}
\]

Now

\[
\frac{1}{f_{1}}=(n-1) \frac{1}{R_{1}}, \frac{1}{f_{2}}=(n-1) \frac{1}{R_{2}}
\]

so

\[
\frac{1}{f}=\frac{1}{f_{1}}+\frac{1}{f_{2}}=(n-1) \frac{18 \lambda}{d^{2}}=\Phi=2 \cdot 40 \mathrm{D}
\]

Here \(n=\) R.I. of glass \(=1.5\).\\
5.91 Here \(\varphi=(n-1)\left(\frac{2}{R_{1}}-\frac{2}{R_{2}}\right)\)\\
so

\[
\frac{1}{R_{1}}-\frac{1}{R_{2}}=\frac{\Phi}{2(n-1)}
\]

As in the previous example, for the dark rings we have

\[
r_{k}^{2}\left(\frac{1}{R_{1}}-\frac{1}{R_{2}}\right)=\frac{\Phi}{2(n-1)} r_{k}^{2}=k \lambda
\]

\(k=0\) is dark spot; excluding it, we take \(k=10\) hre.

Then

\[
r=\sqrt{\frac{20 \lambda(n-1)}{\Phi}}=3.49 \mathrm{~mm}
\]

(b) Path difference in water film will be

\[
n_{0} \bar{r}^{2}\left(\frac{1}{R_{1}}-\frac{1}{R_{2}}\right)
\]

where \(\bar{r}=\) new radius of the ring. Thus\\
or

\[
\begin{aligned}
n_{0} \bar{r}^{2} & =r^{2} \\
\bar{r}=r / \sqrt{n_{0}} & =3.03 \mathrm{~mm} .
\end{aligned}
\]

Where \(n_{0}=\) R.I. of water \(=1 \cdot 33\).\\
5.92 The condition for minima are

\[
\frac{r^{2}}{R} n_{2}=\left(k+\frac{1}{2}\right) \lambda
\]

(There occur phase changes at both surfaces on reflection, hence minima when path difference is half integer multiple of \(\lambda\) ).

In this case \(k=4\) for the fifth dark ring\\
(Counting from \(k=0\) for the first dark ring).\\
Thus, we can write

\[
r=\sqrt{(2 K-1) \lambda R / 2 n_{2}}, K=5
\]

Substituting we get \(r=1.17 \mathrm{~mm}\).\\
5.93 Sharpness of the fringe pattern is the worst when the maxima and minima intermingle :-

\[
n_{1} \lambda_{1}=\left(n_{1}-\frac{1}{2}\right) \lambda_{2}
\]

or putting

\[
\lambda_{1}=\lambda, \lambda_{2}=\lambda+\Delta \lambda
\]

we get

\[
n_{1} \Delta \lambda \cong \frac{\lambda}{2}
\]

or

\[
n_{1} \equiv \frac{\lambda}{2 \Delta \lambda} \approx \frac{\lambda_{1}}{2\left(\lambda_{2}-\lambda_{1}\right)}=140
\]

5.94 Interference pattern vanishes when the maxima due to one wavelength mingle with the minima due to the other. Thus

\[
2 \Delta h=k \lambda_{2}=(k+1) \lambda_{1}
\]

where \(\Delta \boldsymbol{h}=\) displacement of the mirror between the sharpest patterns of rings\\
Thus

\[
k\left(\lambda_{2}-\lambda_{1}\right)=\lambda_{1}
\]

or

\[
k=\frac{\lambda_{1}}{\lambda_{2}-\lambda_{1}}
\]

So

\[
\Delta h=\frac{\lambda_{1} \lambda_{2}}{2\left(\lambda_{2}-\lambda_{1}\right)} \approx \frac{\lambda^{2}}{2 \Delta \lambda} \approx 29 \mathrm{~mm}
\]

5.95 The path difference between (1) \& (2) can be seen to be

\[
\begin{aligned}
\Delta & =2 d \sec \theta-2 d \tan \theta \sin \theta \\
& =2 d \cos \theta=k \lambda
\end{aligned}
\]

for maxima. Here \(\boldsymbol{k}=\) half-integer.\\
The order of interference decreases as \(\theta\) increases i.e. as the radius of the rings increases.\\
(b) Differentiating\\
\includegraphics[max width=\textwidth, alt={}, center]{da0156b6-4133-434e-90d4-ec01b2070afe-169_384_745_1153_625}\\
\(\delta \theta\) decreases as \(\theta\) increases.\\
on putting \(\delta k=-1\). Thus\\
5.96\\
(a) We have \(k_{\max }=\frac{2 d}{\lambda}\). for \(\theta=0 .=10^{5}\).\\
(b) We must have

\[
2 d \cos \theta=k \lambda=(k-1)(\lambda+\Delta \lambda)
\]

Thus \(\frac{1}{k} \approx \frac{\lambda}{2 d}\). and \(\Delta \lambda=\frac{\lambda}{k}=\frac{\lambda^{2}}{2 d}=5 \mathrm{pm}\). on putting the values.

\subsection*{5.3 DIFFRACTION OF LIGHT}
5.97 The radius of the periphery of the \(N^{\text {th }}\) Fresnel zone is

\[
r_{N}=\sqrt{N b \lambda}
\]

Then by conservation of energy

\[
I_{0} \pi(\sqrt{N b \lambda})^{2}=\int_{0}^{\infty} 2 \pi r d r I(r)
\]

Here \(r\) is the distance from the point \(P\).\\
Thus

\[
I_{0}=\frac{2}{N b \lambda} \int_{0}^{\infty} r d r I(r)
\]

5.98 By definition

\[
r_{k}^{2}=k \frac{a b \lambda}{a+b}
\]

for the periphery of the \(\mathbf{k}^{\text {th }}\) zone. Then

So

\[
\begin{gathered}
a r_{k}^{2}+b r_{k}^{2}=k a b \lambda \\
b=\frac{a r_{k}^{2}}{k a \lambda-r_{k}^{2}}=\frac{a r^{2}}{k a \lambda-r^{2}}=2 \text { metre } .
\end{gathered}
\]

on putting the values. (It is given that \(r=r_{k}\) ) for \(k=3\) ).\\
5.99 Suppose maximum intensity is obtained when the aperture contains \(k\) zones. Then a minimum will be obtained for \(k+1\) zones. Another maximum will be obtained for \(k+2\) zones. Hence

Thus

\[
\begin{gathered}
r_{1}^{2}=k \lambda \frac{a b .}{a+b} \\
r_{2}^{2}=(k+2) \lambda \frac{a b}{a+b} \\
\lambda=\frac{a+b}{2 a b}\left(r_{2}^{2}-r_{1}^{2}\right)=0.598 \mu \mathrm{~m}
\end{gathered}
\]

On putting the values.\\
5.100 (a) When the aperture is equal to the first Fresnel Zone :-

The amplitude is \(A_{1}\) and should be compared with the amplitude \(\frac{A}{2}\) when the aperture is very wide. If \(I_{0}\) is the intensity in the second case the intensity in the first case will be \(4 I_{0}\).\\
When the aperture is equal to the internal half of the first zone :- Suppose \(A_{\text {in }}\) and \(A_{\text {out }}\) are the amplitudes due to the two halves of the first Fresnel zone. Clearly \(A_{\text {in }}\) and \(A_{\text {out }}\) differ in phase by \(\frac{\pi}{2}\) because only half the Fresnel zone in involved. Also in magnitude \(\left|A_{\text {in }}\right|=\left|A_{\text {out }}\right|\). Then

\[
A_{1}^{2}=2\left|A_{i n}\right|^{2} \quad \text { so } \quad\left|A_{i n}\right|^{2}=\frac{A_{1}^{2}}{2}
\]

Hence following the argument of the first case. \(I_{i n}=2 I_{0}\)\\
(b) The aperture was made equal to the first Fresnel zone and then half of it was closed along a diameter. In this case the amplitude of vibration is \(\frac{A_{1}}{2}\). Thus

\[
I=I_{0} .
\]

5.101 (a) Suppose the disc does not obstruct light at all. Then

\[
A_{d i s c}+A_{r e m a i n d e r}=\frac{1}{2} A_{d i s c}
\]

(because the disc covers the first Fresnel zone only).\\
So \(A_{\text {remainder }}=-\frac{1}{2} A_{\text {disc }}\)\\
Hence the amplitude when half of the disc is removed along a diameter\\
\(=\frac{1}{2} A_{\text {disc }}+A_{\text {remainder }}=\frac{1}{2} A_{\text {disc }}-\frac{1}{2} A_{\text {disc }} \approx 0\)\\
Hence \(I=0\).\\
(b) In this case

We write

\[
\begin{aligned}
A & =\frac{1}{2} A_{\text {external }}+A_{\text {remainder }} \\
& =\frac{1}{2} A_{\text {external }}-\frac{1}{2} A_{\text {disc }}
\end{aligned}
\]

\[
A_{d i s c}=A_{i n}+i A_{o w e}
\]

where \(A_{\text {in }}\left(A_{\text {out }}\right)\) stands for \(A_{\text {internal }}\left(A_{\text {external }}\right)\). The factor \(i\) takes account of the \(\frac{\pi}{2}\) phase difference between two halves of the first Fresnel zone. Thus

\[
A=-\frac{1}{2} A_{i n} \quad \text { and } \quad I=\frac{1}{4} A_{i n}^{2}
\]

On the other hand

\[
I_{0}=\frac{1}{4}\left(A_{i n}^{2}+A_{\text {out }}^{2}\right)=\frac{1}{2} A_{i n}^{2}
\]

so

\[
I=\frac{1}{2} I_{0}
\]

5.102 When the screen is fully transparent, the amplitude of vibrations is \(\frac{1}{2} A_{1}\) (with intensity \(\left.I_{0}=\frac{1}{4} A_{1}^{2}\right)\).\\
(a) (1) In this case \(A=\frac{3}{4}\left(\frac{1}{2} A_{1}\right)\) so squaring \(I=\frac{9}{16} I_{0}\)\\
(2) In this case \(\frac{1}{2}\) of the plane is blacked out so

\[
A=\frac{1}{2}\left(\frac{1}{2} A_{1}\right) \text { and } I=\frac{1}{4} I_{0}
\]

(3) In this case \(A=\frac{1}{4}\left(A_{1} / 2\right)\) and \(I=\frac{1}{16} I_{0}\).\\
(4) In this case \(A=\frac{1}{2}\left(\frac{1}{2} A_{1}\right)\) again and \(I=\frac{1}{4} I_{0}\) so \(I_{4}=\frac{I}{2}\)

In general/we get \(I(\varphi)=I_{0}\left(1-\left(\frac{\varphi}{2 \pi}\right)\right)^{2}\) where \(\varphi\) is the total angle blocked out by the screen.\\
(b) (5) Here \(A=\frac{3}{4}\left(\frac{1}{2} A_{1}\right)+\frac{1}{4} A_{1}\)\\
\(A_{1}\) being the contribution of the first Fresnel zone.\\
Thus

\[
A=\frac{5}{8} A_{1} \text { and } I=\frac{25}{16} I_{0}
\]

(6) \(A=\frac{1}{2}\left(\frac{1}{2} A_{1}\right)+\frac{1}{2} A_{1}=\frac{3}{4} A_{1}\) and \(I=\frac{9}{4} I_{0}\)\\
(7) \(A=\frac{1}{4}\left(\frac{1}{2} A_{1}\right)+\frac{3}{4} A_{1}=\frac{7}{8} A_{1}\) and \(I=\frac{49}{16} I_{0}\)\\
(8) \(A=\frac{1}{2}\left(\frac{1}{2} A_{1}\right)+\frac{1}{2} A_{1}=\frac{3}{4} A_{1}\) and \(I=\frac{9}{4} I_{0}\left(I_{8}=I_{6}\right)\)

In 5 to 8 the first term in the expression for the amplitude is the contribution of the plane part and the second term gives the expression for the Fresnel zone part. In general in (5) to (8) \(I=I_{0}\left(1+\left(\frac{\varphi}{2 \pi}\right)\right)^{2}\) when \(\varphi\) is the angle covered by the screen.\\
5.103 We would require the contribution to the amplitude of a wave at a point from half a Fresnel zone. For this we proceed directly from the Fresnel Huyghens principle. The complex amplitude is written as

\[
E=\int K(\varphi) \frac{a_{0}}{r} e^{-i k r} d S
\]

Here \(K(\varphi)\) is a factor which depends on the angle \(\varphi\) between a normal \(\vec{n}\) to the area \(d S\) and the direction from\\
\includegraphics[max width=\textwidth, alt={}]{da0156b6-4133-434e-90d4-ec01b2070afe-172_335_470_1417_992} \(d S\) to the point \(P\) and \(r\) is the distance from the element \(d S\) to \(P\).\\
We see that for the first Fresnal zone

\[
\begin{gathered}
\left(\text { using } r \propto b+\frac{\rho^{2}}{2 b}\left(\text { for } \sqrt{\rho^{2}+b^{2}}\right)\right) \\
\sqrt{b \lambda} \\
E \approx \frac{a_{0}}{b} \int_{0} e^{-i k b-i k \rho^{2} / 2 b} 2 \pi \rho d \rho \quad(K(\varphi) \approx 1)
\end{gathered}
\]

For the first Fresnel zone \(r=b+\lambda / 2\) so \(r^{2} \cong b^{2}+b \lambda\) and \(\rho^{2}=b \lambda\).

Thus

\[
\begin{aligned}
E & \propto \frac{a_{0}}{b} e^{-i k b} 2 \pi \int_{0}^{\frac{b \lambda}{2}} e^{-i \frac{k x}{b}} d x \\
& =\frac{a_{0}}{b} 2 \pi e^{-i k b} \frac{e^{-i k \lambda / 2}-1}{-i k / b} \\
& =\frac{a_{0}}{k} 2 \pi i e^{-i k b}(-2)=-\frac{4 \pi}{k} i a_{0} e^{-i k b} \cong A_{1}
\end{aligned}
\]

For the next half zone

\[
\begin{aligned}
E & =\frac{a_{0}}{b} e^{-i k b} 2 \pi \int_{\frac{b \lambda}{2}}^{\frac{3 b \lambda}{4}} e^{-i k x / b} d x x \\
& =\frac{a_{0}}{k} 2 \pi i e^{-i k b}\left(e^{-i \frac{3 k \lambda}{4}}-e^{-i k \lambda / 2}\right) \\
& =\frac{a_{0}}{k} 2 \pi i e^{-i k b}(+1+i)=-\frac{A_{1}(1+i)}{2}
\end{aligned}
\]

If we calculate the contribution of the full \(2^{\text {nd }}\) Fresnel zone we will get - \(A_{1}\). If we take account of the factors \(K(\varphi)\) and \(\frac{1}{r}\) which decrease monotonically we expect the contribution to change to \(-A_{2}\). Thus we write for the contribution of the half zones in the \(2^{\text {nd }}\) Fresnel zone as

\[
-\frac{A_{2}(1+i)}{2} \text { and }-\frac{A_{2}(1-i)}{2}
\]

The part lying in the recess has an extra phase difference equal to \(-\delta=-\frac{2 \pi}{\lambda}(n-1) h\). Thus the full amplitude is (note that the correct form is \(e^{-i k n}\) )

\[
\begin{gathered}
\left(A_{1}-\frac{A_{2}}{2}(1+i)\right) e^{+i \delta}-\frac{A_{2}}{2}(1-i)+A_{3}-A_{4}+\ldots \\
\sim\left(\frac{A_{1}}{2}(1-i)\right) e^{+i \delta}-\frac{A_{2}}{2}(1-i)+\frac{A_{3}}{2}
\end{gathered}
\]

\(\approx\left(\frac{A_{1}}{2}(1-i)\right) e^{+i \delta}+i \frac{A_{1}}{2}\left(\right.\) as \(\left.A_{2} \approx A_{3} \approx A_{1}\right)\) and \(A_{3}-A_{4}+A_{5} \ldots=\frac{A_{3}}{2}\).\\
The corresponding intensity is

\[
\begin{aligned}
I & =\frac{A_{1}^{2}}{4}\left[(1-i) e^{+i \delta}+\frac{i}{e}\right]\left[(1+i) e^{-i \delta}-i\right] \\
& =I_{0}[3-2 \cos \delta+2 \sin \delta]=I_{0}\left[3+2 \sqrt{2} \sin \left(\delta-\frac{\pi}{4}\right)\right]
\end{aligned}
\]

(a) For maximum intensity \(\sin \left(\delta-\frac{\pi}{4}\right)=+1\)\\
or

\[
\begin{gathered}
\delta-\frac{\pi}{4}=2 k \pi+\frac{\pi}{2}, k=0,1,2, \ldots \\
\delta=2 k \pi+\frac{3 \pi}{4}=\frac{2 \pi}{\lambda}(n-1) h
\end{gathered}
\]

so

\[
h=\frac{\lambda}{n-1}\left(k+\frac{3}{8}\right)
\]

(b) For minimum intensity

\[
\sin \left(\delta-\frac{\pi}{4}\right)=1
\]

\[
\delta-\frac{\pi}{4}=2 k \pi+\frac{3 \pi}{2} \text { or } \delta=2 k \pi+\frac{7 \pi}{4}
\]

so

\[
h=\frac{\lambda}{n-1}\left(k+\frac{7 \pi}{8}\right)
\]

(c) For \(\left.I=I_{0}, \begin{array}{c}\cos \delta=0 \\ \sin \delta=-1\end{array}\right\}\) or \(\left\{\begin{array}{l}\sin \delta=0 \\ \cos \delta=+1\end{array}\right.\)

Thus

\[
\delta=2 k \pi \quad h=\frac{k \lambda}{n-1}
\]

or

\[
\delta=2 k \pi+\frac{3 \pi}{2}, h=\frac{\lambda}{n-1}\left(k+\frac{3 \pi}{4}\right)
\]

5.104 The contribution to the wave amplitude of the inner half-zone is

\[
\begin{aligned}
\begin{aligned}
\frac{2 \pi a_{0} e^{-i k b}}{b} \int_{0}^{\sqrt{b \lambda / 2}} e^{-i k \rho^{2} / 2 b} \rho d \rho & \longrightarrow \\
=\frac{2 \pi a_{0} e^{-i k b}}{b} \int_{0}^{b \lambda / 4} e^{-i k x / b d x} & \longrightarrow \\
=\frac{2 \pi a_{0} e^{-i k b}}{b}\left(e^{-i k \lambda / 4}-1\right) \times \frac{1}{b} & \longrightarrow
\end{aligned}
\end{aligned}
\]

\[
=\frac{2 \pi i a_{0} e^{-i k b}}{k}(-i-1)=+\frac{A_{1}}{2}(1+i)
\]

With phase factor this becomes \(\frac{A_{1}}{2}(1+i) e^{i \delta}\) where \(\delta=\frac{2 \pi}{\lambda}(n-1) h\). The contribution of the remaining aperture is \(\frac{A_{1}}{2}(1-i)\)\\
(so that the sum of the two parts when \(\delta=0\) is \(A_{1}\) )\\
Thus the complete amplitude is

\[
\frac{A_{1}}{2}(1+i) e^{i \delta}+\frac{A_{1}}{2}(1-i)
\]

and the intensity is

\[
\begin{aligned}
I & =I_{0}\left[(1+i) e^{i \delta}+(1-i)\right]\left[(1-i) e^{-i \delta}+(1+i)\right] \\
& =I_{0}\left[2+2+(1-i)^{2} e^{-i \delta}+(1+i)^{2} e^{i \delta}\right] \\
& =I_{0}\left[4-2 i e^{-i \delta}+2 i e^{i \delta}\right]=I_{0}(4-4 \sin \delta)
\end{aligned}
\]

Here \(I_{0}=\frac{A_{1}^{2}}{4}\) is the intensity of the incident light which is the same as the intensity due to an aperture of infinite extent (and no recess). Now \(I\) is maximum when \(\sin \delta=-1\)\\
or

\[
\delta=2 k \pi+\frac{3 \pi}{2}
\]

so

\[
h=\frac{\lambda}{n-1}\left(k+\frac{3}{4}\right) \quad \text { and (b) } \quad I_{\max }=8 I_{0}
\]

5.105 We follow the argument of 5.103 . we find that the contribution of the first Fresnel zone is

\[
A_{1}=-\frac{4 \pi i}{k} a_{0} e^{-i k b}
\]

For the next half zone it is \(-\frac{A_{2}}{2}(1+i)\)\\
(The contribution of the remaining part of the \(2^{\text {nd }}\) Fresnel zone will be \(-\frac{A_{2}}{2}(1-i)\) )\\
If the disc has a thickness \(h\), the extra phase difference suffered by the light wave in passing through the disc will be

\[
\delta=\frac{2 \pi}{\lambda}(n-1) h
\]

Thus the amplitude at \(P\) will be

\[
\begin{aligned}
E_{P} & =\left(A_{1}-\frac{A_{2}}{2}(1+i)\right) e^{-i \delta}-\frac{A_{2}}{2}(1-i)+A_{3}-A_{4}-A_{5}+\ldots \\
& \sim\left(\frac{A_{1}(1-i)}{2}\right) e^{-i \delta}+\frac{i A_{1}}{2}=\frac{A_{1}}{2}\left[(1-i) e^{-i \delta}+i\right]
\end{aligned}
\]

The corresponding intensity will be

\[
I=I_{0}(3-2 \cos \delta-2 \sin \delta)=I_{0}\left(3-2 \sqrt{2} \sin \left(\delta+\frac{\pi}{4}\right)\right)
\]

The intensity will be a maximum when\\
or\\
i.e.\\
so

\[
\begin{gathered}
\sin \left(\delta+\frac{\pi}{4}\right)=-1 \\
\delta+\frac{\pi}{4}=2 k \pi+\frac{3 \pi}{2} \\
\delta=\left(k+\frac{5}{8}\right) \cdot 2 \pi \\
h=\frac{\lambda}{n-1}\left(k+\frac{5}{8}\right), k=0,1,2, \ldots
\end{gathered}
\]

Note :- It is not clear why \(k=2\) for \(h_{\min }\). The normal choice will be \(k=0\). If we take \(k=0\) we get \(h_{\text {min }}=0.59 \mu \mathrm{~m}\).\\
5.106 Here the focal point acts as a virtual source of light. This means that we can take spherical waves converging towards \(F\). Let us divide these waves into Fresnel zones just after they emerge from the stop. We write \(r^{2}=f^{2}-(f-h)^{2}=(b-m \lambda / 2)^{2}-(b-h)^{2}\)

Here \(r\) is the radius of the \(m^{\text {th }}\) fresnel zone and \(h\) is the distance to the left of the foot of the perpendicular. Thus

\[
\begin{aligned}
& & r^{2} \cong 2 f h=-b m \lambda+2 b h & \\
\text { So } & & h & =b m \lambda / 2(b-f) \\
\text { and } & & r^{2} & =f b m \lambda /(b-f) .
\end{aligned}
\]

\begin{center}
\includegraphics[max width=\textwidth, alt={}]{da0156b6-4133-434e-90d4-ec01b2070afe-176_287_532_892_963}
\end{center}

The intensity maxima are observed when an odd number of Fresnel zones are exposed by the stop. Thus

\[
r_{k}=\sqrt{\frac{k b f \lambda}{b-f}} \text { where } \quad k=1,3,5, \ldots
\]

5.107 For the radius of the periphery of the \(k^{\text {th }}\) zone we have

\[
r_{k}=\sqrt{k \lambda \frac{a b}{a+b}}=\sqrt{k \lambda b} \quad \text { if } a=\infty
\]

If the aperture diameter is reduced \(\eta\) times it will produce a similar deffraction pattern (reduced \(\boldsymbol{\eta}\) times) if the radii of the Fresnel zones are also \(\boldsymbol{\eta}\) times less. Thus\\
\(r^{\prime}{ }_{k}=r_{k} / \eta\)

\[
\text { This requires } b^{\prime}=b / \eta^{2} \text {. }
\]

5.108 (a) If a point source is placecd before an opaque ball, the diffraction pattern consists of a bright spot inside a dark disc followed by fringes. The bright spot is on the line joining the point source and the centre of the ball. When the object is a finite source of transverse\\
diamension \(y\), every point of the source has its corresponding image on the line joining that point and the centre of the ball. Thus the transverse dimension of the image is given by

\[
y^{\prime}=\frac{b}{a} y=9 \mathrm{~mm} .
\]

(b) The minimum height of the irregularities covering the surface of the ball at random, at which the ball obstructs light is, according to the note at the end of the problem, comparable with the width of the Fresnel zone along which the edge of opaque screen passes.\\
So

\[
h_{\min } \approx \Delta r
\]

To find \(\Delta r\) we note that

\[
\begin{gathered}
r^{2}=\frac{k \lambda a b}{a+b} \\
\text { or } \quad 2 r \Delta r=D \Delta r=\frac{\lambda a b}{a+b} \Delta k
\end{gathered}
\]

Where \(D=\) diameter of the disc ( \(=\) diameter of the last Fresnel zone) and \(\Delta k=1\)\\
Thus \(h_{\text {min }}=\frac{\lambda a b}{D(a+b)}=0.099 \mathrm{~mm}\).\\
5.109 In a zone plate an undarkened circular disc is followed by a number of alternately undarkened and darkened rings. For the proper case, these correspond to \(1^{\text {st }}, 2^{\text {nd }}, 3^{\text {rd }} \ldots\). Fresnel zones.

Let \(r_{1}=\) radius of the central undarkened circle. Then for this to be first Fresnel zone in the present case, we must have

\[
S L+L I-S I=\lambda / 2
\]

Thus if \(r_{1}\) is the radius of the periphery of the first zone\\
or

\[
\begin{aligned}
& \sqrt{a^{2}+r_{1}^{2}}+\sqrt{b^{2}+r_{1}^{2}}-(a+b)=\frac{\lambda}{2} \\
& \frac{r_{1}^{2}}{2}\left(\frac{1}{a}+\frac{1}{b}\right)=\frac{\lambda}{2} \quad \text { or } \quad \frac{1}{a}+\frac{1}{b}=\frac{1}{r_{1}^{2} / \lambda}
\end{aligned}
\]

\begin{center}
\includegraphics[max width=\textwidth, alt={}]{da0156b6-4133-434e-90d4-ec01b2070afe-177_232_568_1300_894}
\end{center}

It is clear that the plate is acting as a lens of focal length

\[
f_{1}=\frac{r_{1}^{2}}{\lambda}=\frac{a b}{a+b}=.6 \text { metre }
\]

This is the principle focal length.\\
Other maxima are obtained when

\[
S L+L I-S I=3 \frac{\lambda}{2}, 5 \frac{\lambda}{2}, \ldots
\]

These focal lengths are also

\[
\frac{r_{1}^{2}}{3 \lambda}, \frac{r_{1}^{2}}{5 \lambda}, \ldots
\]

5.110 Just below the edge the amplitude of the wave is given by

\[
\begin{aligned}
& A=\frac{1}{2}\left(A_{1}-A_{2}+A_{3}-A_{4}+\ldots\right) e^{-i \delta} \\
& +\frac{1}{2}\left(A_{1}-A_{2}+A_{3}-A_{4}+\ldots\right)
\end{aligned}
\]

\begin{center}
\includegraphics[max width=\textwidth, alt={}]{da0156b6-4133-434e-90d4-ec01b2070afe-178_299_490_41_962}
\end{center}

Here the quantity in the brackets is the contribution of various Fresnel zones; the factor \(\frac{1}{2}\) is to take account of the division of the plate into two parts by the ledge; the phase factor \(\delta\) is given by

\[
\delta=\frac{2 \pi}{\lambda} h(n-1)
\]

and takes into account the extra length traversed by the waves on the left.\\
Using

\[
A_{1}-A_{2}+A_{3}-A_{4}+\ldots \approx \frac{A_{1}}{2}
\]

we get

\[
A=\frac{A_{1}}{4}\left(1+e^{i \delta}\right)
\]

and the corresponding intensity is

\[
I=I_{0} \frac{1+\cos \delta}{2} . \text { where } I_{0} \propto\left(\frac{A_{1}}{2}\right)^{2}
\]

(a) This is minimum when

\[
\cos \delta=-1
\]

So

\[
\delta=(2 k+1) \pi
\]

and

\[
h=(2 k+1) \frac{\lambda}{2(n-1)}, \quad k=0,1,2, \ldots
\]

using \(n=1.5, \lambda=0.60 \mu m\)

\[
h=0.60(2 k+1) \mu \mathrm{m} .
\]

(b) \(I=I_{0} / 2\) when \(\cos \delta=0\)\\
or

\[
\delta=k \pi+\frac{\pi}{2}=(2 k+1) \frac{\pi}{2}
\]

Thus in this case

\[
h=0.30(2 k+1) \mu \mathrm{m} .
\]

5.111 (a) From the Cornu's spiral, the intensity of the first maximum is given as

\[
I_{\max , 1}=1.37 I_{0}
\]

and the intensity of the first minimum is given by

\[
I_{\min }=0.78 I_{0}
\]

so the required ratio is

\[
\frac{I_{\max }}{I_{\min }}=1.76
\]

(b) The value of the distance \(x\) is related to the parameter \(v\) in Fresnel's integral by

\[
v=x \sqrt{\frac{2}{b \lambda}} .
\]

For the first two maxima the distances \(x_{1}, x_{2}\) are related to the parameters \(v_{1}, v_{2}\) by

Thus

\[
\begin{gathered}
x_{1}=\sqrt{\frac{b \lambda}{2}} v_{1}, x_{2}=\sqrt{\frac{b \lambda}{2}} v_{2} \\
\left(v_{2}-v_{1}\right) \sqrt{\frac{b \lambda}{2}}=x_{2}-x_{1}=\Delta x \\
\lambda=\frac{2}{b}\left(\frac{\Delta x}{v_{2}-v_{1}}\right)^{2}
\end{gathered}
\]

or\\
From the Cornu's spiral the positions of the maxima are

\[
\begin{gathered}
v_{1}=1.22, v_{2}=2.34, v_{3}=3.08 \text { etc } \\
\lambda=\frac{2}{b}\left(\frac{\Delta x}{1.12}\right)^{2} \equiv 0.63 \mu \mathrm{~m} .
\end{gathered}
\]

Thus\\
5.112 We shall use the equation written down in 5.103, the Fresnel-Huyghens formula.\\
\includegraphics[max width=\textwidth, alt={}, center]{da0156b6-4133-434e-90d4-ec01b2070afe-179_426_622_942_448}

Suppose we want to find the intensity at \(P\) which is such that the coordinates of the edges ( \(x\)-coordinates) with respect to \(P\) are \(x_{2}\) and \(-x_{1}\). Then, the amplitude at \(P\) is

\[
E=\int K(\varphi) \frac{a_{0}}{r} e^{-i k r} d S
\]

We write \(d S=d x d y, y\) is to integrated from \(-\infty+0+\infty\). We write


\begin{equation*}
r \approx b+\frac{x^{2}+y^{2}}{2 b} \tag{1}
\end{equation*}


( \(r\) is the distance of the element of surface on \(I\) from \(P\). It is \(\sqrt{b^{2}+x^{2}+y^{2}}\) and hence approximately (1)). We then get

\[
\begin{aligned}
E & =A_{0}(b)\left[\int_{x_{2}}^{\infty} e^{-i k x^{2} / 2 b} d x+\int_{-\infty}^{-x_{1}} e^{-i k x^{2} / 2 b} d x\right] \\
& =A_{0}^{\prime}(b)\left[\int_{v_{2}}^{\infty} e^{-i \frac{\pi u^{2}}{2}} d u+\int_{-\infty}^{-v_{1}} e^{-i \pi u^{2} / 2} d u\right] \\
v_{2} & =\sqrt{\frac{2}{b \lambda}} x_{2}, v_{1}=\sqrt{\frac{2}{b \lambda}} x_{1}
\end{aligned}
\]

where\\
The intensity is the square of the amplitude. In our case, at the centre\\
\(v_{1}=v_{2}=\sqrt{\frac{2}{b \lambda}} \cdot \frac{a}{2}=\sqrt{\frac{a^{2}}{2 b \lambda}}=0.64\)\\
( \(a=\) width of the strip \(=0.7 \mathrm{~mm}, b=100 \mathrm{~cm}, \lambda=0.60 \mu \mathrm{~m}\) )\\
At, say, the lower edge

\[
v_{1}=0, v_{2}=1.28
\]

Thus\\
\(\frac{I_{\text {centre }}}{I_{\text {edge }}}=\left|\frac{\int_{0.64}^{\infty} e^{-i \pi u^{2} / 2} d u+\int_{-\infty}^{-0.64} e^{-i \pi u^{2} / 2} d u}{\int_{1.28}^{\infty} e^{-i \pi u^{2} / 2} d u+\int_{-\infty}^{0} e^{-i \pi u^{2} / 2} d u}\right|^{2}=4 \frac{\left(\frac{1}{2}-C(0.64)\right)^{2}+\left(\frac{1}{2}-S(0.64)\right)^{2}}{(1-C(1.28))^{2}+(1-S(1.28))^{2}}\)\\
where

\[
\begin{aligned}
& C(v)=\int_{0}^{v} \cos \frac{\pi u^{2}}{2} d u \\
& S(v)=\int_{0}^{v} \sin \frac{\pi u^{2}}{2} d u
\end{aligned}
\]

Rough evaluation of the integrals using cornu's spiral gives

\[
\frac{I_{\text {centre }}}{I_{\text {edge }}} \approx 2.4
\]

(We have used \(\int_{0}^{\infty} \cos \frac{\pi u^{2}}{2} d u=\int_{0}^{\infty} \sin \frac{\pi u^{2}}{2} d u=\frac{1}{2}\)

\[
\begin{aligned}
& C(0.64)=0.62, S(0.64)=0.15 \\
& C(1.28)=0.65 S(1.28)=0.67
\end{aligned}
\]

5.113 If the aperture has width \(h\) then the parameters \((v,-v)\) associated with \(\left(h / 2,-\frac{h}{2}\right)\) are given by

\[
v=\frac{h}{2} \sqrt{\frac{2}{b \lambda}}=h / \sqrt{2 b \lambda}
\]

The intensity of light at \(O\) on the screen is obtained as the square of the amplitude \(A\) of the wave at \(O\) which is\\
\(A \sim\) const \(\int_{-v}^{v} e^{-i \pi u^{2} / 2} d u\)

\[
\frac{\sum^{3}+h / 2}{\sum^{3}-h / 2} \quad b \quad 0
\]

Thus

\[
I=2 I_{0}\left((C(v))^{2}+(S(v))^{2}\right)
\]

where \(C(v)\) and \(S(v)\) have been defined above and \(I_{0}\) is the intensity at \(O\) due to an infinitety wide \((v=\infty)\) aperture for then

\[
I=2 I_{0}\left(\left(\frac{1}{2}\right)^{2}+\left(\frac{1}{2}\right)^{2}\right)=2 I_{0} \times \frac{1}{2}=I_{0}
\]

By definition \(v\) corresponds to the first minimum of the intensity. This means

\[
v=v_{1} \approx .90
\]

when we increase \(h\) to \(h+\Delta h\), the corresponding \(\nu_{2}=\frac{h+\Delta h}{\sqrt{2 b \lambda}}\) relates to the second minimum of intensity. From the cornu's spiral \(\nu_{2} \approx 2.75\)\\
Thus

\[
\Delta h=\sqrt{2 b \lambda}\left(v_{2}-v_{1}\right)=0.85 \sqrt{2 b \lambda}
\]

or

\[
\lambda=\left(\frac{\Delta h}{0.85}\right)^{2} \frac{1}{2 b}=\left(\frac{0.70}{0.85}\right)^{2} \frac{1}{2 \times 0.6} \mu \mathrm{~m}=0.565 \mu \mathrm{~m}
\]

5.114 Let \(a=\) width of the recess and

\[
v=\frac{a}{2} \sqrt{\frac{2}{b \lambda}}=\frac{a}{\sqrt{2 b \lambda}}=\frac{0.6}{\sqrt{2 \times 0.77 \times 0.65}} \approx 0.60
\]

be the parameter along Cornu's spiral corresponding to the half-width of the recess. The amplitude of the diffracted wave is given by

\[
\sim \operatorname{const}\left[e^{i \delta} \int_{-v}^{v} e^{-i \pi u^{2} / 2} d u+\int_{v}^{\infty} e^{-i \pi u^{2} / 2} d u+\int_{-\infty}^{v} e^{-i \pi u^{2} / 2} d u\right]
\]

where

\[
\delta=\frac{2 \pi}{\lambda}(n-1) h
\]

is the extra phase due to the recess. (Actually an extra phase \(e^{-i \delta}\) appears outside the recess. When we take it out and absorb it in the constant we get the expression written).\\
Thus the amplitude is\\
\includegraphics[max width=\textwidth, alt={}, center]{da0156b6-4133-434e-90d4-ec01b2070afe-181_265_500_1693_951}

\[
\sim \operatorname{const}\left[(C(v)-i S(v)) e^{i \delta}+\left(\frac{1}{2}-C(v)\right)-i\left(\frac{1}{2}-S(v)\right)\right]
\]

From the Cornu's spiral, the coordinates corresponding to the parameter \(v=0.60\) are

\[
C(v)=0.57, S(v)=0.13
\]

so the intensity at \(O\) is proportional to

\[
\begin{aligned}
& \left|\left[(0.57-0.13 i) e^{i \delta}-0.07-i 0.37\right]\right|^{2} \\
& \quad=\left(0.57^{2}+0.13^{2}\right)+0.07^{2}+0.37^{2} \\
& +(0.57-0.13 i)(-0.07+0.37 i) e^{i \delta} \\
& +(0.57+0.13 i)(-0.07-i 0.37 i) e^{-i \delta}
\end{aligned}
\]

We write

\[
\begin{aligned}
0.57 \mp 0.13 i & =0.585 e^{\mp i \alpha} \alpha=12.8^{\circ} \\
-0.07 \pm 0.37 i & =0.377 e^{ \pm i \beta} \beta=100.7^{\circ}
\end{aligned}
\]

Thus the cross term is

\[
\begin{array}{r}
2 \times 0.585 \times 0.377 \cos \left(\delta+88^{\circ}\right) \\
\approx 2 \times 0.585 \times 0.377 \cos \left(\delta+\frac{\pi}{2}\right)
\end{array}
\]

For maximum intensity

\[
\begin{aligned}
\delta+\frac{\pi}{2}= & 2 k^{\prime} \pi, \quad k^{\prime}=1,2,3,4, \ldots \\
= & 2(k+1) \pi, k=0,1,2,3, \ldots \\
& \delta=2 k \pi+\frac{3 \pi}{2} \\
h= & \frac{\lambda}{n-1}\left(k+\frac{3}{4}\right)
\end{aligned}
\]

or\\
so\\
5.115\\
\includegraphics[max width=\textwidth, alt={}, center]{da0156b6-4133-434e-90d4-ec01b2070afe-182_414_616_1315_411}

Using the method of problem 5.103 we can immediately write down the amplitudes at 1 and 2. We get :

At 1

\[
\text { amplitude } A_{1} \sim \operatorname{const}\left[\int_{-\infty}^{0} e^{-i \pi u^{2} / 2} d u+e^{-i \delta} \int_{v}^{\infty} e^{-i \pi u^{2} /^{2}} d u\right]
\]

At 2 amplitude \(A_{2} \sim\) const \(\left[\int_{-\infty}^{-\nu} e^{-i \pi u^{2} / 2} d u+e^{-i \delta} \int_{0}^{\infty} e^{-i \pi u^{2} /^{2}} d u\right]\)\\
where

\[
v=a \sqrt{\frac{2}{b \lambda}}
\]

is the parameter of Cornu's spiral and constant factor is common to 1 and 2 .\\
With the usual notation\\
and the result

\[
\begin{gathered}
C=C(v)=\int_{0}^{v} \cos \frac{\pi u^{2}}{2} d u \\
S=S(v)=\int_{0}^{v} \sin \frac{\pi u^{2}}{2} d u \\
\int_{0}^{\infty} \cos \frac{\pi u^{2}}{2} d u=\int_{0}^{\infty} \sin \frac{\pi u^{2}}{2} d u=\frac{1}{2}
\end{gathered}
\]

We find the ratio of intensities as

\[
\frac{I_{2}}{I_{1}}=\left|\frac{\left(\frac{1}{2}-C\right)-i\left(\frac{1}{2}-S\right)+e^{-i \delta} \frac{(1-i)}{2}}{: \left\lvert\, \frac{1-i}{2}+e^{-i \delta}\left\{\left(\frac{1}{2}-C\right)-i\left(\frac{1}{2}-S\right)\right\}\right.}\right|^{2}
\]

(The constants in \(A_{1}\) and \(A_{2}\) must be the same by symmetry)\\
In our case, \(a=0.30 \mathrm{~mm}, \lambda=0.65 \mu \mathrm{~m}, b=1.1 \mathrm{~m}\)

\[
\begin{gathered}
v=0.30 \times \sqrt{\frac{2}{1.1 \times 0.65}}=0.50 \\
C(0.50)=0.48 \quad S(0.50)=0.06 \\
\frac{I_{2}}{I_{1}}=\left|\frac{0.02-0.44 i+e^{-i \delta} \frac{(1-i)}{2}}{\frac{1-i}{2} e^{i \delta}+0.02-0.44 i}\right|^{2}=\left|\frac{1+(0.02-0.44 i) \sqrt{2} e^{i \delta+\frac{i \pi}{4}}}{1+(0.02-0.44 i) \sqrt{2} e^{-i \delta+\frac{i \pi}{4}}}\right|^{2}
\end{gathered}
\]

But \(0.02-0.44 i=0.44 e^{i \alpha}, \alpha=1.525 \mathrm{rad}\left(\approx 87.4^{\circ}\right)\)\\
So \(\frac{I_{2}}{I_{1}}=\left|\frac{1+0.44 \times \sqrt{2} \times e^{i(\delta-0.740)}}{1+0.44 \times \sqrt{2} \times e^{-i(\delta+0.740)}}\right|^{2}=\frac{1+2(0.44)^{2}+2 \sqrt{2} \times 0.44 \cos (\delta-0.740)}{1+2(0.44)^{2}+2 \sqrt{2} \times 0.44 \cos (\delta+0.740)}\)\\
\(I_{2}\) is maximum when

\[
\delta-0 \cdot 740=0(\text { modulo } 2 \pi)
\]

Thus in that case \(\frac{I_{2}}{I_{1}}=\frac{1.387+1.245}{1.387+1.245 \cos (1.48)}=\frac{2.632}{1.5} \approx 1.75\)\\
5.116 We apply the formula of problem 5.103 and calculate

\[
\int_{\text {aperture }} \frac{a_{0}}{r} e^{-i k r} d S=\int_{\text {Semicircle }}+\int_{\text {Slit }}
\]

The contribution of the full \(1^{\text {st }}\) Fresnel zone has been evaluated in 5.103 . The contribution of the semi-circle is one half of it and is

\[
-\frac{2 \pi}{k} i a_{0} e^{-i k b}=-i a_{0} \lambda e^{-i k b}
\]

The contribution of the slit is

Now

\[
\begin{gathered}
\frac{a_{0}}{b} \int_{0}^{0.90 \sqrt{b \lambda}} e^{-i k b} e^{-i k \frac{x^{2}}{2 b}} d x \int_{-\infty}^{\infty} e^{-i k y^{2} / 2 b} d y \\
\int_{-\infty}^{\infty} e^{-i k y^{2} / 2 b} d y=\int_{-\infty}^{\infty} e^{-i \frac{\pi y^{2}}{b \lambda}} d y \\
\sqrt{\frac{b \lambda}{2}} \int_{-\infty}^{\infty} e^{-i \pi u^{2} / 2} d u=\sqrt{b \lambda} e^{-i \pi / 4}
\end{gathered}
\]

Thus the contribution of the slit is

\[
\begin{gathered}
\frac{a_{0}}{b} \sqrt{b \lambda} e^{-i k b-i \pi / 4} \int_{0}^{0.9 \times \sqrt{2}} e^{-i \pi u^{2} / 2} d u \sqrt{\frac{b \lambda}{2}} \\
\quad=a_{0} \lambda e^{-i k b-i \pi / 4} \frac{1}{\sqrt{2}} \int_{0}^{1.27} e^{-i \pi u^{2} / 2} d u
\end{gathered}
\]

Thus the intensity at the observation point \(P\) on the screen is

\[
\begin{aligned}
a_{0}^{2} \lambda^{2} \left\lvert\,-i+\frac{1-i}{2}(C(1.27)\right. & -i S(1.27))\left.\right|^{2}=a_{0}^{2} \lambda^{2}\left|-i+\frac{(1-i)(0.67-0.65 i)}{2}\right|^{2} \\
\text { (on using } C(1.27)= & 0.67 \text { and } S(1.27)=0.65) \\
& =a_{0}^{2} \lambda^{2}|-i+0.01-0.66 i|^{2} \\
& =a_{0}^{2} \lambda^{2}|0.01-1.66 i|^{2} \\
& =2.76 a_{0}^{2} \lambda^{2}
\end{aligned}
\]

Now \(a_{0}^{2} \lambda^{2}\) is the intensity due to half of \(1^{\text {st }}\) Fresnel zone and is therefore equal to \(I_{0}\). (It can also be obtained by doing the \(x\)-integral over- \(\infty\) to \(+\infty\) ).\\
Thus

\[
I=2.76 I_{0}
\]

5.117 From the statement of the problem we know that the width of the slit \(=\) diameter of the first Fresnel zone \(=2 \sqrt{b \lambda}\) where \(b\) is the distance of the observation point from the slit.

We calculate the amplitudes by evaluating the integral of problem 5.103 We get\\
\includegraphics[max width=\textwidth, alt={}, center]{da0156b6-4133-434e-90d4-ec01b2070afe-185_164_605_184_859}

\[
\begin{aligned}
A_{1} & =\frac{a_{0}}{b} \int_{-\sqrt{b \lambda}}^{\sqrt{b \lambda}} e^{-i k b} e^{-i k \frac{x^{2}}{2 b}} d x \int_{0}^{\infty} e^{-i k \frac{y^{2}}{2 b}} d y \\
& =\frac{a_{0}}{b} e^{-i k b} \frac{b \lambda}{2} \int_{-\sqrt{2}}^{\infty} e^{-i \pi u^{2} / 2} d u \times \int_{0}^{\infty} e^{-i \pi u^{2} / 2} d u \\
& =\frac{a_{0} \lambda}{2}(1-i) e^{-i k b}(C(\sqrt{2})-i S(\sqrt{2})) \\
A_{2} & =\frac{a_{0}}{b} \int_{-\sqrt{b \lambda}} e^{-i k b} e^{-i k \frac{x^{2}}{2 b}} d x \int_{-\infty}^{\infty} e^{-i k y^{2} / 2 b} d y \\
& =2 A_{1} \\
A_{3} & =-i a_{0} \lambda e^{-i k b}+\frac{a_{0} \lambda(1-i)}{2}(C(\sqrt{2})-i S(\sqrt{2})) e^{-i k b}
\end{aligned}
\]

where the contribution of the \(1^{\text {st }}\) half Fresnel zone (in \(A_{3}\), first term) has been obtained from the last problem.

Thus

\[
I_{1}=a_{0}^{2} \lambda^{2}\left|\frac{(1-i)(0.53-0.72 i)}{2}\right|^{2}
\]

(on using \(C(\sqrt{2})=0.53, S(\sqrt{2})=0.72\) )

\[
\begin{aligned}
=a_{0}^{2} \lambda^{2}|-0.095-0.625 i|^{2}=0.3996 a_{0}^{2} \lambda^{2} & \\
I_{2} & =4 I_{1} \\
I_{3} & =a_{0}^{2} \lambda^{2}|-0.095-0.625 i-i|^{2} \\
& =a_{0}^{2} \lambda^{2}|-0.095-1.625 i|^{2} \\
& =2.6496 a_{0}^{2} \lambda^{2}
\end{aligned}
\]

So

\[
I_{3}=6.6 I_{1}
\]

Thus

\[
I_{1}: I_{2}: I_{3} \approx 1: 4: 7
\]

5.118 The radius of the first half Fresnel zone is \(\sqrt{b \lambda / 2}\) and the amplitude at \(P\) is obtained using problem 5.103.

\[
\begin{aligned}
& \begin{aligned}
A=\frac{a_{0}}{b}\left[\int_{-\infty}^{-\eta \sqrt{b \lambda / 2}}+\int_{\eta \sqrt{b \lambda / 2}}^{\infty} e^{-i k b-\frac{k x^{2}}{2 b}} d x\right. \\
\int_{-\infty}^{\infty} e^{-i k y^{2} / 2 b} d y+\frac{a_{0}}{b} e^{-i k b} \int_{0} e^{-i k \rho^{2} / 2 b} 2 \pi \rho d \rho
\end{aligned} \\
& \text { We use } \\
& \\
& =\int_{\eta \sqrt{b \lambda / 2}}^{\infty} e^{-i k x^{2} / 2 b} d x=\int_{\eta \sqrt{b \lambda / 2}}^{\infty} e^{-i \frac{\pi x^{2}}{b \lambda}} d x \\
& =\int_{\eta}^{\infty} e^{-i \pi x^{2} / 2 b} d x \\
& =\sqrt{\frac{b \lambda}{2}}\left(\left(\frac{1}{2}-C(\eta)\right)-i\left(\frac{1}{2}-S(\eta)\right)\right. \\
& \int_{0}^{\infty}-\int_{0}^{\eta} e^{-i \pi u^{2} / 2} d u
\end{aligned}
\]

Thus

\[
\begin{aligned}
A= & a_{0} \frac{\lambda}{2} \times 2 \times(1-i) e^{-\imath k b}\left[\left(\frac{1}{2}-C(\eta)\right)\right. \\
& \left.-i\left(\frac{1}{2}-S(\eta)\right)\right]+a_{0} \lambda(1-i) e^{-\imath k b}
\end{aligned}
\]

where we have used

\[
\int_{0}^{\sqrt{b \lambda / 2}} e^{-i k \rho^{2} / 2 b} 2 \pi \rho d \rho=\frac{2 \pi i b}{k}(-1-i)=\frac{2 \pi b}{k}(1-i)=\lambda b(1-i)
\]

Thus the intensity is

\[
I=|A|^{2}=a_{0}^{2} \lambda^{2} \times 2\left[(3 / 2-C(\eta))^{2}+\left(\frac{1}{2}-S(\eta)\right)^{2}\right]
\]

From Cornu's Spiral,

\[
\begin{aligned}
C(\eta)=C(1.07) & =0.76 \\
S(\eta)=S(1.07) & =0.50 \\
I=a_{0}^{2} \lambda^{2} \times 2 \times(0.74)^{2} & =1.09 a_{0}^{2} \lambda^{2}
\end{aligned}
\]

As before

\[
I_{0}=a_{0}^{2} \lambda^{2} \text { so } I \approx I_{0} .
\]

5.119 If a plane wave is incident normally from the left on a slit of width \(b\) and the diffracted wave is observed at a large distance, the resulting pattern is called Fraunhofer diffraction. The condition for this is \(b^{2} \ll l \lambda\) where \(l\) is the distance between the slit and the screen. In practice light may be focussed on the screen with the help of a lens (or a telescope).

Consider an element of the slit which is an infinite strip of width \(d x\). We use the formula of problem 5.103 with\\
\includegraphics[max width=\textwidth, alt={}, center]{da0156b6-4133-434e-90d4-ec01b2070afe-187_437_483_344_973}\\
the following modifications.\\
The factor \(\frac{1}{r}\) characteristic of spherical waves will be omitted. The factor \(K(\varphi)\) will also be dropped if we confine overselves to not too large \(\varphi\). In the direction defined by the angle \(\varphi\) the extra path difference of the wave emitted from the element at \(x\) relative to the wave emitted from the centre is

\[
\Delta=-x \sin \varphi
\]

Thus the amplitude of the wave is given by

\[
\begin{aligned}
\alpha \int_{-b / 2}^{+b / 2} e^{i k \sin \varphi} d x & =\left(e^{i \frac{1}{2} k b \sin \varphi}-e^{-i \frac{1}{2} k b \sin \varphi}\right) / i k \sin \varphi \\
& =\frac{\sin \left(\frac{\pi b}{\lambda} \sin \varphi\right)}{\frac{\pi b}{\lambda} \sin \varphi}
\end{aligned}
\]

Thus

\[
I=I_{0} \frac{\sin ^{2} \alpha}{\alpha^{2}}
\]

where

\[
\alpha=\frac{\pi b}{\lambda} \sin \varphi \text { and }
\]

\(I_{0}\) is a constant\\
Minima are observed for \(\sin \alpha=0\) but \(\alpha \neq 0\)\\
Thus we find minima at angles given by

\[
b \sin \varphi=k \lambda, k= \pm 1, \pm 2, \pm 3, \ldots
\]

5.120 Since \(I(\alpha)\) is \(+v e\) and vanishes for \(b \sin \varphi=k \lambda\) i.e for \(\alpha=k \pi\), we expect maxima of \(I(\alpha)\) between \(\alpha=+\pi \& \alpha=+2 \pi\), etc. We can get these values by.

\[
\begin{aligned}
& \frac{d}{d \alpha}(I(\alpha))=I_{0} 2 \frac{\sin \alpha}{\alpha} \frac{d}{d \alpha} \frac{\sin \alpha}{\alpha}=0 \\
& \frac{\alpha \cos \alpha-\sin \alpha}{\alpha^{2}}=0 \quad \text { or } \quad \tan \alpha=\alpha
\end{aligned}
\]

Solutions of this transcendental equation can be obtained graphically.\\
The first three solutions are

\[
\alpha_{1}=1.43 \pi, \alpha_{2}=2.46 \pi, \alpha_{3}=3.47 \pi
\]

on the + ve side. (On the negative side the solution are \(-\alpha_{1},-\alpha_{2},-\alpha_{3}, \ldots\) )

Thus

\[
\begin{aligned}
b \sin \varphi_{1} & =1.43 \lambda \\
b \sin \varphi_{2} & =2.46 \lambda \\
b \sin \varphi_{3} & =3.47 \lambda
\end{aligned}
\]

Asymptotically the solutions are

\[
b \sin \varphi_{m} \cong\left(M+\frac{1}{2}\right) \lambda
\]

5.121 The relation \(b \sin \theta=k \lambda\)\\
for minima (when light is incident normally on the slit) has a simple interpretation : \(b \sin \theta\) is the path difference between extreme wave normals emitted at angle \(\theta\)\\
\includegraphics[max width=\textwidth, alt={}, center]{da0156b6-4133-434e-90d4-ec01b2070afe-188_375_377_1040_261}\\
\includegraphics[max width=\textwidth, alt={}, center]{da0156b6-4133-434e-90d4-ec01b2070afe-188_383_420_1024_815}

When light is incident at an angle \(\theta_{0}\) the path difference is

\[
b\left(\sin \theta-\sin \theta_{0}\right)
\]

and the condition of minima is

\[
b\left(\sin \theta-\sin \theta_{0}\right)=k \lambda
\]

For the first minima

\[
b\left(\sin \theta-\sin \theta_{0}\right)= \pm \lambda \quad \text { or } \sin \theta=\sin \theta_{0} \pm \frac{\lambda}{b}
\]

Putting in numbers \(\theta_{0}=30^{\circ}, \lambda=0.50 \mu \mathrm{~m} b=10 \mu \mathrm{~m}\)

\[
\begin{gathered}
\sin \theta=\frac{1}{2} \pm \frac{1}{20}=0.55 \text { or } 0.45 \\
\theta_{+1}=33^{\circ}-20^{\prime} \text { and } \theta_{-1}=26^{\circ} 44^{\prime}
\end{gathered}
\]

5.122 (a) This case is analogous to the previous one except that the incident wave moves in glass of RI n. Thus the expression for the path difference for light diffracted at angle \(\theta\) from the normal to the hypotenuse of the wedge is

\[
b(\sin \theta-n \sin \Theta)
\]

we write

\[
\theta=\Theta+\Delta \theta
\]

Then for the direction of principal Fraunhofer maximum

\[
b(\sin (\Theta+\Delta \theta)-n \sin \Theta)=0
\]

or

\[
\Delta \theta=\sin ^{-1}(n \sin \Theta)-\Theta
\]

Using

\[
\begin{gathered}
\Theta=15^{\circ}, n=1.5 \text { we get } \\
\Delta \theta=7.84^{\circ}
\end{gathered}
\]

\includegraphics[max width=\textwidth, alt={}, center]{da0156b6-4133-434e-90d4-ec01b2070afe-189_352_424_242_1032}\\
(b) The width of the central maximum is obtained from ( \(\lambda=0.60 \mu \mathrm{~m}, b=10 \mu \mathrm{~m}\) )

\[
b\left(\sin \theta_{1}-n \sin \Theta\right)= \pm \lambda
\]

Thus

\[
\begin{aligned}
\theta_{+1} & =\sin ^{-1}\left(n \sin \Theta+\frac{\lambda}{b}\right)=26.63^{\circ} \\
\theta_{-1} & =\sin ^{-1}\left(n \sin \Theta \frac{\lambda}{b}\right)=19.16^{\circ} \\
\therefore \quad \delta \theta & =\theta_{+1}-\theta_{-1}=7.47^{\circ}
\end{aligned}
\]

\begin{center}
\includegraphics[max width=\textwidth, alt={}]{da0156b6-4133-434e-90d4-ec01b2070afe-189_285_799_1064_311}
\end{center}

The path difference between waves reflected at \(A\) and \(B\) is

\[
d\left(\cos \alpha_{0}-\cos \alpha\right)
\]

and for maxima

\[
d\left(\cos \alpha_{0}-\cos \alpha\right)=k \lambda, k=0, \pm 1, \pm 2, \ldots
\]

In our case, \(k=2\) and \(\alpha_{0}, \alpha\) are small in radiaus. Then

\[
2 \lambda=d\left(\frac{\alpha^{2}-\alpha_{0}^{2}}{2}\right)
\]

Thus

\[
\lambda-\frac{\left(\alpha^{2}-\alpha_{0}^{2}\right) d}{4}=0.61 \mu \mathrm{~m}
\]

for

\[
\alpha=\frac{3 \pi}{180}, \alpha_{0}=\frac{\pi}{180}, d=10^{-3} \mathrm{~m}
\]

5.124 The general formula for diffraction from \(N\) slits is

\[
\begin{gathered}
I=I_{0} \frac{\sin ^{2} \alpha}{\alpha^{2}} \frac{\sin ^{2} N \beta}{\sin ^{2} \beta} \\
\alpha=\frac{\pi a \sin \theta}{\lambda} \\
\beta=\frac{\pi(a+b) \sin \theta}{\lambda}
\end{gathered}
\]

where\\
and \(N=3\) in the cases here.\\
(a) In this case \(a+b=2 a\)\\
so

\[
\beta=2 \alpha \text { and } I=I_{0} \frac{\sin ^{2} \alpha}{\alpha^{2}}\left(3-4 \sin ^{2} 2 \alpha\right)^{2}
\]

On plotting we get a curve that qualitatively looks like the one below\\
\includegraphics[max width=\textwidth, alt={}, center]{da0156b6-4133-434e-90d4-ec01b2070afe-190_606_881_686_202}\\
(b) In this case \(\boldsymbol{a}+\boldsymbol{b}=\mathbf{3} \boldsymbol{a}\)\\
so\\
and

\[
\begin{gathered}
\beta=3 \alpha \\
I=I_{0} \frac{\sin ^{2} \alpha}{\alpha^{2}}\left(2-4 \sin ^{2} 3 \alpha\right)^{2}
\end{gathered}
\]

This has 3 minima between the principal maxima\\
5.125 From the formula \(\quad d \sin \theta=m \lambda\)\\
we have \(d \sin 45^{\circ}=2 \lambda_{1}=2 \times 0.65 \mu \mathrm{~m}\)\\
or

\[
d=2 \sqrt{2} \times 0.65 \mu \mathrm{~m}
\]

Then for \(\lambda_{2}=0.50\) in the third order

\[
\begin{aligned}
& 2 \sqrt{2} \times 0.65 \sin \theta=3 \times 0.50 \\
& \sin \theta=\frac{1.5}{1.3 \times \sqrt{2}}=0.81602
\end{aligned}
\]

This gives \(\theta=54.68^{\circ} \approx 55^{\circ}\)\\
5.126 The diffraction formula is

\[
d \sin \theta_{0}=n_{0} \lambda
\]

where \(\theta_{0}=35^{\circ}\) is the angle of diffraction corresponding to order \(n_{0}\) (which is not yet known).\\
Thus

\[
d=\frac{n_{0} \lambda}{\sin \theta_{0}}=n_{0} \times 0.9327 \mu \mathrm{~m}
\]

on using \(\lambda=0.535 \mu \mathrm{~m}\)\\
For the \(n^{\text {th }}\) order we get

\[
\sin \theta=\frac{n}{n_{0}} \sin \theta_{0}=\frac{n}{n_{0}}(0.573576)
\]

If \(n_{0}=1\), then \(n>n_{0}\) is at least 2 and \(\sin \theta>1\) so \(n=1\) is the highest order of diffraction.\\
If \(n_{0}=2\) then \(n=3,4\), but \(\sin \theta>1\) for \(n=4\) thus the highest order of diffraction is 3 .\\
If \(\boldsymbol{n}_{\mathbf{0}}=\mathbf{3}\),\\
then

\[
n=4,5,6 .
\]

For \(n=6, \sin \theta=2 \times 0.57>1\), so not allowed while for

\[
n=5, \sin \theta=\frac{5}{3} \times 0.573576<1
\]

is allowed. Thus in this case the highest order of diffraction is five as given. Hence

\[
n_{0}=3
\]

and

\[
d=3 \times 0.9327=2.7981 \approx 2.8 \mu \mathrm{~m} .
\]

5.127 Given that

\[
\begin{gathered}
d \sin \theta_{1}=\lambda \\
d \sin \theta_{2}=d \sin \left(\theta_{1}+\Delta \theta\right)=2 \lambda
\end{gathered}
\]

Thus \(\sin \theta_{1} \cdot \cos \Delta \theta+\cos \theta_{1} \sin \Delta \theta=2 \sin \theta_{1}\)\\
or

\[
\sin \theta_{1}(2-\cos \Delta \theta)=\cos \theta_{1} \sin \Delta \theta
\]

or

\[
\tan \theta_{1}=\frac{\sin \Delta \theta}{2-\cos \Delta \theta}
\]

or

\[
\begin{aligned}
\sin \theta_{1} & =\frac{\sin \Delta \theta}{\sqrt{\sin ^{2} \Delta \theta+(2-\cos \Delta \theta)^{2}}} \\
& =\frac{\sin \Delta \theta}{\sqrt{5-4 \cos \Delta \theta}}
\end{aligned}
\]

Finally

\[
\lambda=\frac{d \sin \Delta \theta}{\sqrt{5-4 \cos \Delta \theta}} .
\]

Substitution gives \(\lambda \approx 0.534 \mu \mathrm{~m}\)\\
5.128 (a) Here the simple formula

\[
d \sin \theta=m_{1} \lambda \text { holds. }
\]

Thus

\[
1.5 \sin \theta=m \times 0.530 \quad \sin \theta=\frac{m \times 0.530}{1.5}
\]

Highest permissible \(m\) is \(m=2\) because \(\sin \theta>1\) if \(m=3\). Thus

\[
\sin \theta=\frac{1.06}{1.50} \text { for } m=2, \text { This gives } \theta=45^{\circ} \text { nearby. }
\]

(b) Here \(d\left(\sin \theta_{0}-\sin \theta\right)=n \lambda\)

\[
\text { Thus } \begin{aligned}
\sin \theta & =\sin \theta_{0}-\frac{n \lambda}{d} \\
& =\sin 60^{\circ}-n \times \frac{0.53}{1.5} \\
& =0.86602-n \times 0.353333 .
\end{aligned}
\]

For \(n=5, \sin \theta=-0.900645\)\\
for \(\quad n=6, \sin \theta<-1\).\\
Thus the highest order is \(n=5\) and we get

\[
\theta=\sin ^{-1}(-0.900645) \cong-64^{\circ}
\]

\includegraphics[max width=\textwidth, alt={}, center]{da0156b6-4133-434e-90d4-ec01b2070afe-192_572_521_214_748}\\
5.129 For the lens

For the grating

\[
\frac{1}{f}=(n-1)\left(\frac{1}{R}-\frac{1}{\infty}\right) \quad \text { or } \quad f=\frac{R}{n-1}
\]

\[
\begin{aligned}
d \sin \theta_{1} & =\lambda \text { or } \sin \theta_{1}=\frac{\lambda}{d} \\
\operatorname{cosec} \theta_{1} & =\frac{d}{\lambda}, \cot \theta_{1}=\sqrt{\left(\frac{d}{\lambda}\right)^{2}-1} \\
\tan \theta_{1} & =\frac{1}{\sqrt{\left(\frac{d}{\lambda}\right)^{2}-1}}
\end{aligned}
\]

Hence the distance between the two symmetrically placed first order maxima

\[
=2 f \tan \theta_{1}=\frac{2 R}{(n-1) \sqrt{\left(\frac{d}{\lambda}\right)^{2}-1}}
\]

On putting \(\quad R=20, n=1 \cdot 5, d=6 \cdot 0 \mu \mathrm{~m}\)

\[
\lambda=0.60 \mu \mathrm{~m} \text { we get } 8.04 \mathrm{~cm} .
\]

5.130 The diffraction formula is easily obtained on taking account of the fact that the optical path in the glass wedge acquires a factor \(n\) (refractive index). We get

\[
d\left(n \sin \Theta-\sin \left(\Theta \theta_{k}\right)\right)=k \lambda
\]

Since \(n>0, \Theta-\theta_{0}>\Theta\) and so \(\theta_{0}\) must be negative. We get, using \(\Theta=30^{\circ}\)

\[
\frac{3}{2} \times \frac{1}{2}=\sin \left(30^{\circ}-\theta_{0}\right)=\sin 48.6^{\circ}
\]

Thus

\[
\theta_{0}=-18 \cdot 6^{\circ}
\]

Also for \(k=1\)

Thus

\[
\begin{gathered}
\frac{3}{4}-\sin \left(30^{\circ}-\theta_{+1}\right)=\frac{\lambda}{d}=\frac{0 \cdot 5}{2 \cdot 0}=\frac{1}{4} \\
\theta_{+1}=0^{\circ}
\end{gathered}
\]

We calculate \(\theta_{k}\) for various \(k\) by the above formula. For \(k=6\).

\[
\sin \left(\theta_{k}-30^{\circ}\right)=\frac{3}{4} \Rightarrow \theta_{k}=78.6^{\circ}
\]

For \(k=7\)\\
\includegraphics[max width=\textwidth, alt={}, center]{da0156b6-4133-434e-90d4-ec01b2070afe-193_347_399_150_1058}

\[
\sin \left(\theta_{k}-30^{\circ}\right)=+1 \Rightarrow \theta_{k}=120^{\circ}
\]

This is in admissible. Thus the highest order that can be observed is

\[
\begin{gathered}
k=6 \\
\theta_{k}=78 \cdot 6^{\circ}
\end{gathered}
\]

corresponding to\\
(for \(k=7\) the diffracted ray will be grazing the wedge).\\
5.131 The intensity of the central Fraunhofer maximum will be zero if the waves from successive grooves (not in the same plane) differ in phase by an odd multiple of \(\pi\). Then since the phase difference is

\[
\delta=\frac{2 \pi}{\lambda}(n-1) h
\]

for the central ray we have\\
or

\[
\begin{aligned}
\frac{2 \pi}{\lambda}(n-1) h & =\left(k-\frac{1}{2}\right) 2 \pi, k=1,2,3, \ldots \\
h & =\frac{\lambda}{n-1}\left(k-\frac{1}{2}\right) .
\end{aligned}
\]

The path difference between the rays \(1 \& 2\) is\\
approximately (neglecting terms of order \(\boldsymbol{\theta}^{\mathbf{2}}\) )\\
\(\boldsymbol{a} \sin \theta+\boldsymbol{a}-\boldsymbol{n} \boldsymbol{a}\)\\
\(=a \sin \theta-(n-1) a\)\\
Thus for a maximum\\
\(a \sin \theta-\left(k^{\prime}+\frac{1}{2}\right) \lambda=m \lambda\)\\
or \(a \sin \theta=\left(m+k^{\prime}+\frac{1}{2}\right) \lambda\),\\
\includegraphics[max width=\textwidth, alt={}, center]{da0156b6-4133-434e-90d4-ec01b2070afe-193_439_628_1338_772}

\[
\begin{gathered}
k^{\prime}=0,1,2, \ldots \\
m=0, \pm 1, \pm 2, \ldots
\end{gathered}
\]

The first maximum after the central minimum is obtained when \(m+k^{\prime}=0\)\\
We'get'

\[
a \sin \theta_{1}=\frac{1}{2} \lambda
\]

5.132 When standing ultra sonic waves are sustained in the tank it behaves like a grating whose grating element is

\[
d=\frac{v}{v}=\text { wavelength of the ultrasonic }
\]

\(v=\) velocity of ultrasonic. Thus for maxima

\[
\frac{v}{v} \sin \theta_{m}=m \lambda
\]

On the other hand

\[
f \tan \theta_{\mathrm{m}}=m \Delta x
\]

Assuming \(\theta_{\mathrm{m}}\) to be small (because \(\lambda \ll \frac{\mathrm{v}}{\mathrm{v}}\) )\\
we get

\[
\Delta x=\frac{f \tan \theta_{\mathrm{m}}}{m}=\frac{f \tan \theta_{\mathrm{m}}}{\frac{\mathrm{v}}{v \lambda} \sin \theta_{\mathrm{m}}}=\frac{\lambda v f}{v}
\]

or

\[
v=\frac{\lambda v f}{\Delta x}
\]

Putting the values \(\lambda=0.55 \mu \mathrm{~m}, v=4.7 \mathrm{MHz} f=0.35 \mathrm{~m}\) and \(\Delta x=0.60 \times 10^{-3} \mathrm{~m}\) we easily get

\[
v=1.51 \mathrm{~km} / \mathrm{sec} .
\]

5.133 Each star produces its own diffraction pattern in the focal plane of the objective and these patterns are separated by angle \(\psi\). As the distance \(d\) decreases the angle \(\theta\) between the neighbouring maxima in either diffraction pattern increases \((\sin \theta=\lambda / d)\). When \(\theta\) becomes equal to \(2 \psi\) the first deterioration of visibility occurs because the maxima of one system of fringes coincide with the minima of the other system. Thus from the condition \(\theta=2 \psi\) and \(\sin \theta=\frac{\lambda}{d}\) we get\\
\(\psi=\frac{1}{2} \theta \approx \frac{\lambda}{2 d}\) (radians)\\
\includegraphics[max width=\textwidth, alt={}, center]{da0156b6-4133-434e-90d4-ec01b2070afe-194_529_524_1061_965}

Putting the values we get \(\psi \cong 0.06^{\prime \prime}\)\\
5.134 (a) For normal incidence, the maxima are given by

\[
d \sin \theta=n \lambda
\]

so

\[
\sin \theta=n \frac{\lambda}{d}=n \times \frac{0.530}{1.500}
\]

Clearly

\[
n \leq 2 \text { as } \sin \theta>1 \text { for } n=3 .
\]

Thus the highest order is \(n=2\). Then

\[
D=\frac{d \theta}{d \lambda}=\frac{k}{d \cos \theta}=\frac{k}{d} \frac{1}{\sqrt{1-\left(\frac{k \lambda}{d}\right)^{2}}}
\]

Putting \(k=2, \lambda=0.53 \mu \mathrm{~m}, d=1.5 \mu \mathrm{~m}=1500 \mathrm{~nm}\)\\
we get \(D=\frac{2}{1500} \frac{1}{\sqrt{1-\left(\frac{1.06}{1.5}\right)^{2}}} \times \frac{180}{\pi} \times 60=6.47 \mathrm{ang} . \mathrm{min} / \mathrm{nm}\).\\
(b) We write the diffraction formula as

\[
d\left(\sin \theta_{0}+\sin \theta\right)=k \lambda
\]

so

\[
\sin \theta_{0}+\sin \theta=k \frac{\lambda}{d}
\]

Here

\[
\theta_{0}=45^{\circ} \text { and } \sin \theta_{0}=0.707
\]

so

\[
\sin \theta_{0}+\sin \theta \leq 1.707 . \text { Since }
\]

\[
\begin{gathered}
\frac{\lambda}{d}=\frac{0.53}{1.5}=0.353333, \text { we see that } \\
k \leq 4
\end{gathered}
\]

Thus highest order corresponds to \(k=4\).\\
Now as before \(D=\frac{d \theta}{d \lambda}\) so

\[
\begin{gathered}
D=\frac{k}{d \cos \theta}=\frac{k / d}{\sqrt{1-\left(\frac{k \lambda}{d}-\sin \theta_{0}\right)^{2}}} \\
=12.948 \text { ang. min } / \mathrm{n} \mathrm{~m}
\end{gathered}
\]

5.135 We have

\[
d \sin \theta=k \lambda
\]

so

\[
\frac{d \theta}{d \lambda}=D=\frac{k}{d \cos \theta}=\frac{\tan \theta}{\lambda}
\]

5.136 For the second order principal maximum\\
or

\[
\begin{aligned}
& d \sin \theta_{2}=2 \lambda=k \lambda \\
& \frac{N \pi}{\lambda} d \sin \theta_{2}=2 N \pi
\end{aligned}
\]

minima adjacent to this maximum occur at\\
or

\[
\begin{aligned}
\frac{N \pi}{\lambda} d \sin \left(\theta_{2} \pm \Delta \theta\right) & =(2 N \pm 1) \pi \\
d \cos \theta_{2} \Delta \theta & =\frac{\lambda}{N}
\end{aligned}
\]

Finally angular width of the \(2^{\text {nd }}\) principal maximum is

\[
2 \Delta \theta=\frac{2 \lambda}{N d \cos \theta_{2}}=\frac{2 \lambda}{N d \sqrt{1-(k \lambda / d)^{2}}}=\frac{\tan \theta_{2}}{N}
\]

On putting the values we get \(11.019^{\prime \prime}\) of arc\\
5.137 Using

\[
\begin{aligned}
R=\frac{\lambda}{\delta \lambda} & =k N=\frac{N d \sin \theta}{\lambda} \\
& =\frac{l \sin \theta}{\lambda} \leq \frac{l}{\lambda}
\end{aligned}
\]

5.138 For the just resolved waves the frequency difference

\[
\begin{aligned}
\delta v & =\frac{c \delta \lambda}{\lambda}=\frac{c}{\lambda R}=\frac{c}{\lambda k N} \\
& =\frac{c}{N d \sin \theta}=\frac{1}{\delta t}
\end{aligned}
\]

since \(N d \sin \theta\) is the path difference between waves emitted by the extremities of the grating.

\[
\delta \lambda=.050 \mathrm{~nm}
\]

\[
\begin{gathered}
R=\frac{\lambda}{\delta \lambda} \approx \frac{600}{.05}=12000 \text { (nearly) } \\
=k N
\end{gathered}
\]

On the other hand

Thus

\[
\begin{aligned}
d \sin \theta & =k \lambda \\
\frac{l}{k N} \sin \theta & =\lambda
\end{aligned}
\]

where \(l=10^{-2}\) metre is the width of the grating\\
Hence

\[
\begin{aligned}
\sin \theta & =12000 \times \frac{\lambda}{l} \\
& =12000 \times 600 \times 10^{-7}=0.72 \\
\theta & =46^{\circ}
\end{aligned}
\]

or\\
5.140 (a) We see that

\[
N=6.5 \times 10 \times 200=13000
\]

Now to resolve lines with \(\delta \lambda=0.015 \mathrm{~nm}\) and \(\lambda \approx 670.8 \mathrm{~nm}\) we must have

\[
R=\frac{670 \cdot 8}{0 \cdot 015}=44720
\]

Since \(3 N<R<4 N\) one must go to the fourth order to resolve the said components.\\
(b) we have \(d=\frac{1}{200} \mathrm{~mm}=5 \mu \mathrm{~m}\)\\
so

\[
\sin \theta=\frac{k \lambda}{d}=\frac{k \times 0.670}{5}
\]

since \(|\sin \theta| \leq 1\) we must have \(\quad k \leq 7.46\)\\
so

\[
k_{\max }=7 \approx \frac{d}{\lambda}
\]

Thus

\[
R_{\max }=k_{\max } N=91000 \approx \frac{N d}{\lambda}=\frac{l}{\lambda}
\]

where \(l=6.5 \mathrm{~cm}\) is the grating width.\\
Finally

\[
\delta \lambda_{\min }=\frac{\lambda}{R_{\max }}=\frac{670}{91000}=.007 \mathrm{~nm}=7 \mathrm{pm} \approx \frac{\lambda^{2}}{l} .
\]

5.141 Here\\
so

\[
\begin{gathered}
R=\frac{\lambda}{\delta \lambda}=\frac{589.3}{0.6}=\mathrm{kN}=5 \mathrm{~N} \\
N=\frac{589.3}{3}=\frac{10^{-2}}{d} \\
d=\frac{3 \times 10^{-2}}{589.3} \mathrm{~m}=.0509 \mathrm{~mm}
\end{gathered}
\]

(b) To resolve a doublet with \(\lambda=460.0 \mathrm{~nm}\) and \(\delta \lambda=0.13 \mathrm{~nm}\) in the third order we must have

\[
N=\frac{R}{3}=\frac{460}{3 \times 0.13}=1179
\]

This means that the grating is

\[
N d=1179 \times 0.0509=60.03 \mathrm{~mm}
\]

wide \(=6 \mathrm{~cm}\) wide.\\
5.142 (a) From \(d \sin \theta=k \lambda\)\\
we get

\[
\delta \theta=\frac{k \delta \lambda}{d \cos \theta}
\]

On the other hand

\[
x=f \sin \theta
\]

so

\[
\delta x=f \cos \theta \delta \theta=\frac{k f}{d} \delta \lambda
\]

For

\[
f=0.80 \mathrm{~m}, \delta \lambda=0.03 \mathrm{~nm} \text { and }
\]

\[
d=\frac{1}{250} \mathrm{~mm}
\]

we get

\[
\delta x=\left\{\begin{aligned}
6 \mu \mathrm{~m} & \text { if } k=1 \\
12 \mu \mathrm{~m} & \text { if } k=2
\end{aligned}\right.
\]

(b) Here \(N=25 \times 250=6250\)\\
and

\[
\frac{\lambda}{\delta \lambda}=\frac{310 \cdot 169}{0 \cdot 03}=10339 \cdot \cdot>N
\]

and so to resolve we need \(k=2\) For \(k=1\) gives an R.P. of only 6250 .\\
\includegraphics[max width=\textwidth, alt={}, center]{da0156b6-4133-434e-90d4-ec01b2070afe-198_306_823_86_308}

Suppose the incident light consists of two wavelengths \(\lambda\) and \(\lambda+\delta \lambda\) which are just resolved by the prism. Then by Rayleigh's criterion, the maximum of the line of wavelength \(\lambda\) must coincide with the first minimum of the line of wavelength \(\lambda+\delta \lambda\). Let us write both conditions in terms of the optical path differences for the extreme rays :\\
For the light of wavelength \(\lambda\)

\[
b n-(D C+C E)=0
\]

For the light of wavelength \(\lambda+\delta \lambda\)

\[
b(n+\delta n)-(D C+C E)=\lambda+\delta \lambda
\]

because the path difference between extreme rays equals \(\lambda\) for the first minimum in a single slit diffraction (from the formula \(a \sin \theta=\lambda\) ).\\
Hence

\[
b \delta n \approx \lambda
\]

and

\[
R=\frac{\lambda}{\delta \lambda}=b\left|\frac{\delta n}{\delta \lambda}\right|=b\left|\frac{d n}{d \lambda}\right|
\]


\begin{align*}
& \begin{array}{l}
\lambda \\
\delta \lambda \\
\text { For }
\end{array} \quad \begin{array}{l}
R=b\left|\frac{d n}{d \lambda}\right|=2 B b / \lambda^{3} \\
\text { for }
\end{array} \begin{array}{l}
R_{1}=1.223 \times 10^{4} \\
\lambda_{2}=0.656 \mu \mathrm{~m} \\
R_{2}=0.3542 \times 10^{4}
\end{array}  \tag{a}\\
& \text { To resolve the } D \text {-lines we require } \\
& \qquad R=\frac{5893}{6}=982 \\
& \text { Thus } \\
& \qquad 982=\frac{0.02 \times b}{(0.5893)^{3}} \\
& \qquad b=\frac{982 \times(0.5893)^{3}}{0.02} \mu \mathrm{~m}=1.005 \times 10^{4} \mu \mathrm{~m}=1.005 \mathrm{~cm}
\end{align*}


\(5.145 b\left|\frac{d n}{d \lambda}\right|=k N=2 \times 10,000\)

\[
\begin{gathered}
b \times 0.10 \mu \mathrm{~m}^{-1}=2 \times 10^{4} \\
b=2 \times 10^{5} \mu \mathrm{~m}=0.2 \mathrm{~m}=20 \mathrm{~cm}
\end{gathered}
\]

5.146 Resolving power of the objective

\[
=\frac{D}{1.22 \lambda}=\frac{5 \times 10^{-2}}{1.22 \times 0.55 \times 10^{-6}}=7.45 \times 10^{4}
\]

Let \((\Delta y)_{\min }\) be the minimum distance between two points at a distance of 3.0 km which the telescope can resolve. Then\\
or

\[
\begin{gathered}
\frac{(\Delta y)_{\min }}{3 \times 10^{3}}=\frac{1.22 \lambda}{D}=\frac{1}{7.45 \times 10^{4}} \\
(\Delta y)_{\min }=\frac{3 \times 10^{3}}{7.45 \times 10^{4}}=0.04026 \mathrm{~m}=4.03 \mathrm{~cm} .
\end{gathered}
\]

5.147 The limit of resolution of a reflecting telescope is determined by diffraction from the mirror and obeys a formula similar to that from a refracting telescope. The limit of resolution is

\[
\frac{1}{R}=\frac{1 \cdot 22 \lambda}{D}=\frac{(\Delta y)_{\min }}{L}
\]

where \(L=\) distance between the earth and the moon \(=384000 \mathrm{~km}\)\\
Then putting the values \(\lambda=0.55 \mu \mathrm{~m}, D=5 \mathrm{~m}\)\\
we get

\[
(\Delta y)_{\min }=51.6 \text { metre }
\]

5.148 By definition, the magnification

\[
\Gamma=\frac{\text { angle subtended by the image at the eye }}{\text { angle subtended by the object at the eye }}=\frac{\psi^{\prime}}{\psi}
\]

At the limit of resolution

\[
\psi=\frac{1 \cdot 22 \lambda}{D}
\]

where \(D=\) diameter of the objective\\
On the other hand to be visible to the eye \(\psi^{\prime} \geq \frac{1 \cdot 22 \lambda}{d_{0}}\)\\
where \(d_{0}=\) diameter of the pupil\\
Thus to avail of the resolution offered by the telescope we must have

Hence

\[
\begin{aligned}
& \Gamma \geq \frac{1.22 \lambda}{d_{0}} / \frac{1.22 \lambda}{D}=\frac{D}{d_{0}} \\
& \Gamma_{\min }=\frac{D}{d_{0}}=\frac{50 \mathrm{~mm}}{4 \mathrm{~mm}}=12.5
\end{aligned}
\]

\begin{center}
\includegraphics[max width=\textwidth, alt={}]{da0156b6-4133-434e-90d4-ec01b2070afe-199_417_922_1591_275}
\end{center}

Let \(A\) and \(B\) be two points in the field of a microscope which is represented by the lens C D. Let \(A^{\prime}, B^{\prime}\) be their image points which are at equal distances from the axis of the lens CD . Then all paths from \(A\) to \(A^{\prime}\) are equal and the extreme difference of paths from \(A\) to \(B^{\prime}\) is equal to

\[
\begin{gathered}
A D B^{\prime}-A C B^{\prime} \\
=A D+D B^{\prime}-\left(A C+C B^{\prime}\right) \\
=A D+D B^{\prime}-B D-D B^{\prime} \\
+B C+C B^{\prime}-A C-C B^{\prime} \\
\left(\text { as } B D+D B^{\prime}=B C+C B^{\prime}\right) \\
=A D-B D+B C-A C \\
=2 A B \cos \left(90^{\circ}-\alpha\right)=2 A B \sin \alpha
\end{gathered}
\]

From the theory of diffraction by circular apertures this distance must be equal to \(1.22 \lambda\)\\
when \(B^{\prime}\) coincides with the minimum of the diffraction due to \(A\) and \(A^{\prime}\) with the minimum of the diffraction due to \(B\). Thus

\[
A B=\frac{1.22 \lambda}{2 \sin \alpha}=0.61 \frac{\lambda}{\sin \alpha}
\]

Here \(2 \alpha\) is the angle subtended by the objective of the microscope at the object.\\
Substituting the values

\[
A B=\frac{0.61 \times 0.55}{0.24} \mu \mathrm{~m}=1.40 \mu \mathrm{~m} .
\]

5.150 Suppose \(d_{\min }=\) minimum separation resolved by the microscope\\
\(\psi=\) angle subtended at the eye by this object when the object is at the least distance of distinct vision \(l_{0}(=25 \mathrm{~cm})\).\\
\(\psi^{\prime}=\) minimum angular separation resolved by the eye \(=\frac{1 \cdot 22 \lambda}{d_{0}}\)\\
From the previous problem\\
and

\[
\begin{gathered}
d_{\min }=\frac{0.61 \lambda}{\sin \alpha} \\
\psi=\frac{d_{\min }}{l_{0}}=\frac{0.61 \lambda}{l_{0} \sin \alpha}
\end{gathered}
\]

Now

\[
\Gamma=\text { magnifying power }=\frac{\text { angle subtended at the eye by the image }}{\text { angle subtended at the eye by the object }}
\]

when the object is at the least distance of distinct vision

\[
\geq \frac{\psi^{\prime}}{\psi}=2\left(\frac{l_{0}}{d_{0}}\right) \sin \alpha
\]

Thus

\[
\Gamma_{\min }=2\left(\frac{l_{0}}{d_{0}}\right) \sin \alpha=2 \times \frac{25}{0.4} \times 0.24=30
\]

5.151 Path difference\\
\(=B C-A D\)\\
\(=a\left(\cos 60^{\circ}-\cos \alpha\right)\)\\
For diffraction maxima\\
\(a\left(\cos 60^{\circ}-\cos \alpha\right)=k \lambda\), since \(\lambda=\frac{2}{5} a\), we get\\
\includegraphics[max width=\textwidth, alt={}, center]{da0156b6-4133-434e-90d4-ec01b2070afe-201_321_669_61_584}\\
and we get

\[
\begin{gathered}
\cos \alpha=\frac{1}{2}-\frac{2}{5} k \\
k=-1, \cos \alpha=\frac{1}{2}+\frac{2}{5}=0.9, \alpha=26^{\circ} \\
k=0, \cos \alpha=\frac{1}{2}=0.5, \alpha=60^{\circ} \\
k=1, \cos \alpha=\frac{1}{2}-\frac{2}{5}=0.1, \alpha=84^{\circ} \\
k=2, \cos \alpha=\frac{1}{2}-\frac{4}{5}=-0.3, \alpha=107.5^{\circ} \\
k=3, \cos \alpha=\frac{1}{2}-\frac{6}{5}=-0.7, \alpha=134.4^{\circ}
\end{gathered}
\]

Other values of \(k\) are not allowed as they lead to \(|\cos \alpha|>1\).\\
5.152 We give here a simple derivation of the condition for diffraction maxima, known as Laue equations. It is easy to see form the above figure that the path difference between waves scattered by nearby scattering centres \(P_{1}\) and \(P_{2}\) is

\[
\begin{array}{r}
P_{2} A-P_{1} B=\overrightarrow{r \cdot} \overrightarrow{s_{0}}-\vec{r} \cdot \vec{s} \\
=\vec{r} \cdot\left(\overrightarrow{s_{0}}-\vec{s}\right)=\vec{r} \cdot \vec{S} .
\end{array}
\]

Here \(\vec{r}\) is the radius vector \(\overrightarrow{P_{1} P_{2}}\). For maxima this path difference must be an integer multiple of \(\lambda\) for any two neighbouring atoms. In the present case of two dunensional lattice with \(X\) - rays incident normally \(\vec{r} \cdot \vec{s}=0\). Taking successively - nearest neighbours in the \(x\) - \& \(y\) - directions\\
We get the equations\\
\includegraphics[max width=\textwidth, alt={}, center]{da0156b6-4133-434e-90d4-ec01b2070afe-201_467_677_1178_772}

\[
\begin{aligned}
a \cos \alpha & =h \lambda \\
b \cos \beta & =k \lambda
\end{aligned}
\]

Here \(\cos \alpha\) and \(\cos \beta\) are the direction cosines of the ray with respect to the \(x \& y\) axes of the two dimensional crystal.

\[
\cos \alpha=\frac{\Delta x}{\sqrt{(\Delta x)^{2}+4 l^{2}}}=\sin \left(\tan ^{-1} \frac{\Delta x}{2 l}\right)=0.28735
\]

so using

\[
\begin{gathered}
h=k=2 \text { we get } \\
a=\frac{40 \times 2}{.28735} \mathrm{pm}=0.278 \mathrm{~nm}
\end{gathered}
\]

Similarly

\[
\begin{gathered}
\cos \beta=\frac{\Delta y}{\sqrt{(\Delta y)^{2}+4 l^{2}}}=\sin \left(\tan ^{-1} \frac{\Delta y}{2 l}\right)=0.19612 \\
b=\frac{80}{\cos \beta} \mathrm{pm}=0.408 \mathrm{~nm}
\end{gathered}
\]

5.153 Suppose \(\alpha, \beta\), and \(\gamma\) are the angles between the direction to the diffraction maximum and the directions of the array along the periods \(a, b\), and \(c\) respectively (call them \(x, y, \& z\) axes). Then the value of these angles can be found from the following familiar conditions

\[
\begin{gathered}
a(1-\cos \alpha)=k_{1} \lambda \\
b \cos \beta=k_{2} \lambda \text { and } c \cos \gamma=k_{3} \lambda
\end{gathered}
\]

where \(k_{1}, k_{2}, k_{3}\) are whole numbers \((+,-\), or 0\()\)\\
(These formulas are, in effect, Laue equations, see any text book on modern physics). Squaring and adding we get on using \(\cos ^{2} \alpha+\cos ^{2} \beta+\cos ^{2} \gamma=1\)

Thus

\[
\begin{gathered}
2-2 \cos \alpha=\left[\left(\frac{k_{1}}{a}\right)^{2}+\left(\frac{k_{2}}{b}\right)^{2}+\left(\frac{k_{3}}{c}\right)^{2}\right] \lambda^{2}=\frac{2 k_{1} \lambda}{a} \\
\lambda=\frac{2 k_{1} / a}{\left[\left(k_{1} / a\right)^{2}+\left(k_{2} / a\right)^{2}+\left(k_{3} / a\right)^{2}\right]} .
\end{gathered}
\]

Knowing \(a, b, c\) and the integer \(k_{1}, k_{2}, k_{3}\) we can find \(\alpha, \beta, \gamma\) as well as \(\lambda\).\\
5.154 The unit cell of NaCl is shown below. In an infinite crystal, there are four \(\mathrm{Na}^{+}\)and four \(\mathrm{Cl}^{-}\)ions per unit cell. (Each ion on the middle of the edge is shared by four unit cells; each ion on the face centre by two unit cells, the ion in the middle of the cell by one cell only and finally each ion on the corner by eight unit cells.) Thus

\[
4 \frac{M}{N_{A}}=\rho \cdot a^{3}
\]

where \(M=\) molecular weight of NaCl in gms \(\quad=58.5 \mathrm{gms}\)\\
\includegraphics[max width=\textwidth, alt={}, center]{da0156b6-4133-434e-90d4-ec01b2070afe-202_527_618_1159_816}\\
\(N_{A}=\) Avogadro number \(=6.023 \times 10^{23}\)\\
Thus

\[
\frac{1}{2} a=\sqrt{\frac{M}{2 N_{A} \rho}}=2.822 \AA
\]

The natural facet of the crystal is one of the faces of the unit cell. The interplanar distance

\[
d=\frac{1}{2} a=2.822 \AA
\]

Thus

\[
2 d \sin \alpha=2 \lambda
\]

So

\[
\lambda=d \sin \alpha=2.822 \AA \times \frac{\sqrt{3}}{2}=244 \mathrm{pm} .
\]

5.155 When the crystal is rotated, the incident monochromatic beam is diffracted from a given crystal plane of interplanar spacing \(d\) whenever in the course of rotation the value of \(\theta\) satisfies the Bragg equation.\\
We have the equations \(2 d \sin \theta_{1}=k_{1} \lambda\) and \(2 d \sin \theta_{2}=k_{2} \lambda\)\\
But

\[
\pi-2 \theta_{1}=\pi-2 \theta_{2}+\alpha \text { or } 2 \theta_{1}=2 \theta_{2}-\alpha
\]

so

\[
\theta_{2}=\theta_{1}+\frac{\alpha}{2} .
\]

Thus

\[
2 d\left\{\sin \theta_{1} \cos \frac{\alpha}{2}+\cos \theta_{1} \sin \frac{\alpha}{2}\right\}=k_{2} \lambda
\]

\begin{center}
\includegraphics[max width=\textwidth, alt={}]{da0156b6-4133-434e-90d4-ec01b2070afe-203_318_799_618_645}
\end{center}

Hence

\[
2 d \sin \frac{\alpha}{2} \cos \theta_{1}=\left(k_{2}-k_{1} \cos \frac{\alpha}{2}\right) \lambda
\]

also

\[
2 d \sin \frac{\alpha}{2} \sin \theta_{1}=k_{1} \lambda \sin \frac{\alpha}{2}
\]

Squaring and adding \(2 d \sin \frac{\alpha}{2}=\left(k_{1}^{2}+k_{2}^{2}-2 k_{1} k_{2} \cos \frac{\alpha}{2}\right)^{1 / 2} \lambda\)\\
Hence

\[
d=\frac{\lambda}{2 \sin \frac{\alpha}{2}}\left[k_{1}^{2}+k_{2}^{2}-2 k_{1} k_{2} \cos \frac{\alpha}{2}\right]^{1 / 2}
\]

Substituting

\[
\alpha=60^{\circ}, k_{1}=2, k_{2}=3, \lambda=174 \mathrm{p} \mathrm{~m}
\]

we get

\[
d=281 \mathrm{p} \mathrm{~m}=2 \cdot 81 \AA
\]

(and not 0.281 p m as given in the book.)\\
(Lattice parameters are typically in \(\AA^{\prime} s\) and not in fractions of a pm.)\\
5.156 In a polycrystalline specimen, microcrystals are oriented at various angles with respect to one another. The microcrystals which are oriented at certain special angles with respect to the incident beam produce diffraction maxima that appear as rings.\\
The radial of these rings are given by

\[
r=l \tan 2 \alpha
\]

where the Bragg's law gives

\[
2 d \sin \alpha=k \lambda
\]

In our case \(k=2, d=155 \mathrm{pm}, \lambda=17.8 \mathrm{pm}\)\\
so \(\alpha=\sin ^{-1} \frac{17.8}{155}=6.6^{\circ}\) and \(r=3.52 \mathrm{~cm}\).\\
\includegraphics[max width=\textwidth, alt={}, center]{da0156b6-4133-434e-90d4-ec01b2070afe-203_290_577_1661_842}

\subsection*{5.4 POLARIZATION OF LIGHT}
5.157 Natural light can be considered to be an incoherent mixture of two plane polarized light of intensity \(I_{0} / 2\) with mutually perpendicular planes of vibration. The screen consisting of the two polaroid half-planes acts as an opaque half-screen for one or the other of these light waves. The resulting diffraction pattern has the alterations in intensity (in the illuminated region) characteristic of a straight edge on both sides of the boundary.

At the boundary the intensity due to either component is\\
\includegraphics[max width=\textwidth, alt={}, center]{da0156b6-4133-434e-90d4-ec01b2070afe-204_402_490_163_996}

\[
\frac{\left(I_{0} / 2\right)}{4}
\]

and the total intensity is \(\frac{I_{0}}{4}\). (Recall that when light of intensity \(I_{0}\) is incident on a straight edge, the illuminance in front of the edge is \(I_{0} / 4\) ).\\
5.158 (a) Assume first that there is no polaroid and the amplitude due to the entire hole which extends over the first Fresnel zone is \(A_{1}\)\\
Then, we know, as usual, \(I_{0}=\frac{A_{1}^{2}}{4}\),\\
When the polaroid is introduced as shown above, each half transmits only the corresponding polarized light. If the full hole were covered by one polaroid the\\
\includegraphics[max width=\textwidth, alt={}, center]{da0156b6-4133-434e-90d4-ec01b2070afe-204_299_300_900_1080}\\
amplitude transmitted will be \(\left(A_{1} / \sqrt{2}\right)\).

Therefore the amplitude transmitted in the present case will be \(\frac{A_{1}}{2 \sqrt{2}}\) through either half. Since these transmitted waves are polarized in mutually perpendicular planes, the total intensity will be

\[
\left(\frac{A_{1}}{2 \sqrt{2}}\right)^{2}+\left(\frac{A_{1}}{2 \sqrt{2}}\right)^{2}=\frac{A_{1}^{2}}{4}=I_{0} .
\]

(b) We interpret the problem to mean that the two polaroid pieces are separated along the circumference of the circle limiting the first half of the Fresnel zone. (This however is inconsistent with the polaroids being identical in shape; however no other interpretation makes sense.)\\
From (5.103) and the previous problems we see that the amplitudes of the waves transmitted through the two parts is

\[
\frac{A_{1}}{2 \sqrt{2}}(1+i) \text { and } \frac{A_{1}}{2 \sqrt{2}}(1-i)
\]

and the intensity is

\[
\begin{gathered}
\left|\frac{A_{1}^{2}}{2 \sqrt{2}}(1+i)\right|^{2}+\left|\frac{A_{1}}{2 \sqrt{2}}(1-i)\right|^{2} \\
=\frac{A_{1}^{2}}{2}=2 I_{0}
\end{gathered}
\]

5.159 When the polarizer rotates with angular velocity \(\omega\) its instantaneous principal direction makes angle \(\omega t\) from a reference direction which we choose to be along the direction of vibration of the plane polarized incident light. The transmitted flux at this instant is

\[
\Phi_{0} \cos ^{2} \omega t
\]

and the total energy passing through the polarizer per revolution is

\[
\begin{gathered}
\int_{0}^{T} \Phi_{0} \cos ^{2} \omega t d t, T=2 \pi / \omega \\
=\Phi_{0} \frac{\pi}{\omega}=0.6 \mathrm{~mJ}
\end{gathered}
\]

5.160 Let \(I_{0}=\) intensity of the incident beam.

Then the intensity of the beam transmitted through the first Nicol prism is

\[
I_{1}=\frac{1}{2} I_{0}
\]

and through the \(2^{\text {nd }}\) prism is

Through the \(\boldsymbol{N}^{\text {th }}\) prism it will be

\[
\begin{aligned}
I_{2} & =\left(\frac{1}{2} I_{0}\right) \cos ^{2} \varphi \\
I_{N} & =I_{N-1} \cos ^{2} \varphi \\
& =\frac{1}{2} I_{0} \cos ^{2(N-1)} \varphi
\end{aligned}
\]

\begin{center}
\includegraphics[max width=\textwidth, alt={}]{da0156b6-4133-434e-90d4-ec01b2070afe-205_380_383_1147_998}
\end{center}

Hence fraction transmitted\\
and

\[
\begin{aligned}
=\frac{I_{N}}{I_{0}}=\eta=\frac{1}{2} \cos ^{(2 N-1)} \varphi & =0.12 \text { for } N=6 . \\
\varphi & =30^{\circ}
\end{aligned}
\]

5.161 When natural light is incident on the first polaroid, the fraction transmitted will be \(\frac{1}{2} \tau\) (only the component polarized parallel to the principal direction of the polaroid will go).

The emergent light will be plane polarized and on passing through the second polaroid will be polarized in a different direction (corresponding to the principal direction of the \(2^{\text {nd }}\) polaroid) and the intensity will have decreased further by \(\tau \cos ^{2} \varphi\).\\
In the third polaroid the direction of polarization will again have to change by \(\varphi\) thus only a fraction \(\tau \cos ^{2} \varphi\) will go through.

Finally

\[
I=I_{0} \times \frac{1}{2} \tau^{3} \cos ^{4} \varphi
\]

Thus the intensity will have decreased\\
for

\[
\begin{gathered}
\frac{I_{o}}{I}=\frac{2}{\tau^{3} \cos ^{4} \varphi}=60.2 \text { times } \\
\tau=0.81, \varphi=60^{\circ}
\end{gathered}
\]

5.162 Suppose the partially polarized light consists of natural light of intensity \(I_{1}\) and plane polarized light of intensity \(I_{2}\) with direction of vibration parallel to, say, \(x\)-axis.\\
Then when a polaroid is used to transmit it, the light transmitted will have a maximum intensity

\[
\frac{1}{2} I_{1}+I_{2}
\]

when the principal direction of the polaroid is parallel to \(x\)-axis, and will have a minimum intensity \(\frac{1}{2} I_{1}\) when the principal direction is \(\perp^{r}\) to \(x\)-axis.

Thus

\[
\begin{aligned}
P & =\frac{I_{\max }-I_{\min }}{I_{\max }+I_{\min }}=\frac{I_{2}}{I_{1}+I_{2}} \\
\frac{I_{2}}{I_{1}} & =\frac{P}{1-P}=\frac{0.25}{0.75}=\frac{1}{3} .
\end{aligned}
\]

so\\
5.163 If, as above,\\
\(I_{1}=\) intensity of natural component\\
\(I_{2}=\) intensity of plane polarized component\\
then

\[
I_{\max }=\frac{1}{2} I_{1}+I_{2}
\]

and

\[
I=\frac{I_{\max }}{\eta}=\frac{1}{2} I_{1}+I_{2} \cos ^{2} \varphi
\]

so

\[
I_{2}=I_{\max }\left(1-\frac{1}{\eta}\right) \operatorname{cosec}^{2} \varphi
\]

\[
I_{1}=2 I_{\max }\left[1-\left(1-\frac{1}{\eta}\right) \operatorname{cosec}^{2} \varphi\right]=\frac{2 I_{\max }}{\sin ^{2} \varphi}\left[\frac{1}{\eta}-\cos ^{2} \varphi\right]
\]

Then

\[
P=\frac{I_{2}}{I_{1}+I_{2}}=\frac{1-\frac{1}{\eta}}{2\left(\frac{1}{\eta}-\cos ^{2} \varphi\right)+1-\frac{1}{\eta}}=\frac{\eta-1}{1-\eta \cos 2 \varphi}
\]

On putting\\
we get

\[
\begin{gathered}
\eta=3.0, \varphi=60^{\circ} \\
P=\frac{2}{1+3 \times \frac{1}{2}}=\frac{4}{5}=0.8
\end{gathered}
\]

5.164 Let us represent the natural light as a sum of two mutually perpendicular components, both with intensity \(I_{0}\). Suppose that each polarizer transmits a fraction \(\alpha_{1}\) of the light with oscillation plane parallel to the prinicipal direction of the polarizer and a fraction \(\alpha_{2}\) with oscillation plane perpendicular to the principal direction of the polaizer. Then the intensity of light transmitted through the two polarizers is equal to

\[
I_{\|}=\alpha_{1}^{2} I_{0}+\alpha_{2}^{2} I_{0}
\]

when their principal direction are parallel and

\[
I_{\perp}=\alpha_{1} \alpha_{2} I_{0}+\alpha_{2} \alpha_{1} I_{0}=2 \alpha_{1} \alpha_{2} I_{0}
\]

when they are crossed. But\\
so

\[
\begin{aligned}
& \frac{I_{\perp}}{I_{\|}}=\frac{2 \alpha_{1} \alpha_{2}}{\alpha_{1}^{2}+\alpha_{2}^{2}}=\frac{1}{\eta} \\
& \frac{\alpha_{1}-\alpha_{2}}{\alpha_{1}+\alpha_{2}}=\sqrt{\frac{\eta-1}{\eta+1}}
\end{aligned}
\]

(a) Now the degree of polarization produced by either polarizer when used singly is

\[
P_{0}=\frac{I_{\max }-I_{\min }}{I_{\max }+I_{\min }}=\frac{\alpha_{1}-\alpha_{2}}{\alpha_{1}+\alpha_{2}}
\]

(assuming, of course, \(\alpha_{1}>\alpha_{2}\) )

Thus

\[
P_{0}=\sqrt{\frac{\eta-1}{\eta+1}}=\sqrt{\frac{9}{11}}=0.905
\]

(b) When both polarizer are used with their principal directions parallel, the transmitted light, when analysed, has maximum intensity, \(I_{\max }=\alpha_{1}^{2} I_{0}\) and minimum intensity, \(I_{\min }=\alpha_{2}^{2} I_{0}\)\\
so

\[
\begin{aligned}
P & =\frac{\alpha_{1}^{2}-\alpha_{2}^{2}}{\alpha_{1}^{2}+\alpha_{2}^{2}}=\frac{\alpha_{1}-\alpha_{2}}{\alpha_{1}+\alpha_{2}} \cdot \frac{\left(\alpha_{1}+\alpha_{2}\right)^{2}}{\alpha_{1}^{2}+\alpha_{2}^{2}} \\
& =\sqrt{\frac{\eta-1}{\eta+1}} \cdot\left(1+\frac{2 \alpha_{1} \alpha_{2}}{\alpha_{1}^{2}+\alpha_{2}^{2}}\right) \\
& =\sqrt{\frac{\eta-1}{\eta+1}}\left(1+\frac{1}{\eta}\right)=\frac{\sqrt{\eta^{2}-1}}{\eta}=\sqrt{1-\frac{1}{\eta^{2}}}=0.995 .
\end{aligned}
\]

5.165 If the principal direction \(N\) of the Nicol is along \(A\) or \(B\), the intensity of light transmitted is the same whether the light incident is one with oscillation plane \(N_{1}\) or one with \(N_{2}\). If \(N\) makes an angle \(\delta \varphi\) with \(A\) as shown then the fractional difference in intensity transmitted (when the light incident is \(N_{1}\) or \(N_{2}\) ) is

\[
\begin{aligned}
\left(\frac{\Delta I}{I}\right)_{A} & =\frac{\cos ^{2}\left(90^{\circ}-\frac{\varphi}{2}-\delta \varphi\right)-\cos ^{2}\left(90^{\circ}+\frac{\varphi}{2}-\delta \varphi\right)}{\cos ^{2}\left(90^{\circ}-\frac{\varphi}{2}\right)} \\
& =\frac{\sin ^{2}\left(\frac{\varphi}{2}+\delta \varphi\right)-\sin ^{2}\left(\frac{\varphi}{2}-\delta \varphi\right)}{\sin ^{2} \frac{\varphi}{2}} \\
& =\frac{2 \sin \frac{\varphi}{2} \cdot 2 \cos \frac{\varphi}{2} \delta \varphi}{\sin ^{2} \frac{\varphi}{2}}=4 \cot \frac{\varphi}{2} \delta \varphi
\end{aligned}
\]

\begin{center}
\includegraphics[max width=\textwidth, alt={}]{da0156b6-4133-434e-90d4-ec01b2070afe-208_499_620_446_849}
\end{center}

If \(N\) makes an angle \(\delta \varphi(\ll \varphi)\) with \(B\) then\\
\(\left(\frac{\Delta I}{I}\right)_{B}=\frac{\cos ^{2}(\varphi / 2-\delta \varphi)-\cos ^{2}(\varphi / 2+\delta \varphi)}{\cos ^{2} \varphi / 2}=\frac{2 \cos \frac{\varphi}{2} \cdot 2 \sin \varphi / 2 \delta \varphi}{\cos ^{2} \varphi / 2}=4 \tan \varphi / 2 \delta \varphi\)

Thus

\[
\eta=\left(\frac{\Delta I}{I}\right)_{A} /\left(\frac{\Delta I}{I}\right)_{B}=\cot ^{2} \varphi / 2
\]

or

\[
\varphi=2 \tan ^{-1} \frac{1}{\sqrt{\eta}}
\]

This gives \(\varphi=11 \cdot 4^{\circ}\) for \(\eta=100\).\\
5.166 Fresnel equations read

\[
I_{\perp}^{\prime}=I_{\perp} \frac{\sin ^{2}\left(\theta_{1}-\theta_{2}\right)}{\sin ^{2}\left(\theta_{1}+\theta_{2}\right)} \text { and } I_{\|}^{\prime}=I_{\|} \frac{\tan ^{2}\left(\theta_{1}-\theta_{2}\right)}{\tan ^{2}\left(\theta_{1}+\theta_{2}\right)}
\]

At the boundary between vacuum and a dielectric \(\theta_{1} \neq \theta_{2}\) since by Snell's law

\[
\sin \theta_{1}=n \sin \theta_{2}
\]

Thus \(I_{1}^{\prime} / I_{1}\) cannot be zero. However, if \(\theta_{1}+\theta_{2}=90^{\circ}, I_{\|}^{\prime}=0\) and the reflected light is polarized in this case. The condition for this is

\[
\sin \theta_{1}=n \sin \theta_{2}, \quad=n \sin \left(90^{\circ}-\theta_{1}\right)
\]

or

\[
\tan \theta_{1}=n \quad \theta_{1} \text { is called Brewsta's angle. }
\]

The angle between reflected light and refracted light is \(90^{\circ}\) in this case.\\
5.167 (a) From Fresnel's equations

\[
\begin{aligned}
I_{\perp}^{\prime} & =I_{\perp} \frac{\sin ^{2}\left(\theta_{1}-\theta_{2}\right)}{\sin ^{2}\left(\theta_{1}+\theta_{2}\right)} \\
I_{\|}^{\prime} & =0 \\
I_{\perp}^{\prime} & =I_{\perp} \sin ^{2}\left(\theta_{1}-\theta_{2}\right) \\
& =\frac{1}{2} I\left(\sin \theta_{1} \cos \theta_{2}-\cos \theta_{1} \sin \theta_{2}\right)^{2}
\end{aligned}
\]

Now

\[
\begin{gathered}
\tan \theta_{1}=n, \sin \theta_{1}=\frac{n}{\sqrt{n^{2}+1}} \\
\cos \theta_{1}=\frac{1}{\sqrt{n^{2}+1}}, \sin \theta_{2}=\cos \theta_{1} \\
\cos \theta_{2}=\sin \theta_{1} \\
I_{\perp}^{\prime}=\frac{1}{2} I\left(\frac{n^{2}-1}{n^{2}+1}\right)^{2}
\end{gathered}
\]

Thus reflection coefficient \(=\rho=\frac{I_{\perp}^{\prime}}{I}\)

\[
=\frac{1}{2}\left(\frac{n^{2}-1}{n^{2}+1}\right)^{2}=0.074
\]

on putting \(n=1.5\)\\
(b) For the refracted light

\[
\begin{aligned}
I_{\perp}^{\prime \prime} & =I_{\perp}-I_{\perp}^{\prime}=\frac{1}{2} I\left\{1-\left(\frac{n^{2}-1}{n^{2}+1}\right)^{2}\right\} \\
& =\frac{1}{2} I \frac{4 n^{2}}{\left(n^{2}+1\right)^{2}} \\
I_{\|}^{\prime} & =\frac{1}{2} I
\end{aligned}
\]

at the Brewster's angle.\\
Thus the degree of polarization of the refracted light is

\[
\begin{aligned}
P & =\frac{I_{\|}^{\prime \prime}-I_{\perp}^{\prime \prime}}{I_{\|}^{\prime \prime}+I_{\perp}^{\prime \prime}}=\frac{\left(n^{2}+1\right)^{2}-4 n^{2}}{\left(n^{2}+1\right)^{2}+4 n^{2}} \\
& =\frac{\left(n^{2}-1\right)^{2}}{2\left(n^{2}+1\right)^{2}-\left(n^{2}-1\right)^{2}}=\frac{\rho}{1-\rho}
\end{aligned}
\]

On putting \(\rho=0.074\) we get \(P=0.080\).\\
5.168 The energy transmitted is, by conservation of energy, the difference between incident energy and the reflected energy. However the intensity is affected by the change of the cross section of the beam by refraction. Let \(A_{i}, A_{r}, A_{t}\) be the cross sections of the incident, reflected and transmitted beams. Then

\[
\begin{aligned}
& A_{i}=A_{r} \\
& A_{t}=A_{i} \frac{\cos r}{\cos i}
\end{aligned}
\]

\begin{center}
\includegraphics[max width=\textwidth, alt={}]{da0156b6-4133-434e-90d4-ec01b2070afe-210_450_659_81_837}
\end{center}

But at Brewster's angle \(r=90-i\)\\
so\\
Thus

\[
\begin{gathered}
A_{t}=A_{i} \tan i=n A_{i} \\
I_{t}=\frac{(1-\rho) I_{0}}{n}
\end{gathered}
\]

5.169 The amplitude of the incident component whose oseillation vector is perpendicular to the plane of incidence is\\
and similarly

\[
\begin{aligned}
& A_{\perp}=A_{0} \sin \varphi \\
& A_{\|}=A_{0} \cos \varphi
\end{aligned}
\]

Then

\[
\begin{aligned}
I_{\perp}^{\prime} & =I_{0} \frac{\sin ^{2}\left(\theta_{1}-\theta_{2}\right)}{\sin ^{2}\left(\theta_{1}+\theta_{2}\right)} \sin ^{2} \varphi \\
& =I_{0}\left[\frac{\sin \theta_{1} \cos \theta_{2}-\cos \theta_{1} \sin \theta_{2}}{\sin \theta_{1} \cos \theta_{2}+\cos \theta_{1} \sin \theta_{2}}\right]^{2} \sin ^{2} \varphi \\
& =I_{0}\left[\frac{n^{2}-1}{n^{2}+1}\right]^{2} \sin ^{2} \varphi
\end{aligned}
\]

Hence

\[
\rho=\frac{I_{\perp}^{\prime}}{I_{0}}=\left[\frac{n^{2}-1}{n^{2}+1}\right]^{2} \sin ^{2} \varphi
\]

Putting \(n=1.33\) for water we get \(\rho=0.0386\)\\
5.170 Since natural light is incident at the Brewster's angle, the reflected light 1 is completely polarized and \(P_{1}=1\).

Similarly the ray 2 is incident on glass air surface at\\
Brewster's angle \(\left(\tan ^{-1} \frac{1}{n}\right)\) sc 3 is also completely polarized. Thus \(P_{3}=1\)\\
Now as in 5.167 (b)

\[
P_{2}=\frac{\rho}{1-\rho}=0.087 \text { if } \rho=0.080
\]

\begin{center}
\includegraphics[max width=\textwidth, alt={}]{da0156b6-4133-434e-90d4-ec01b2070afe-211_403_581_255_81}
\end{center}

Finally as shown in the figure

\[
P_{4}=\frac{\frac{1}{2}-\frac{1}{2}(1-2 \rho)^{2}}{\frac{1}{2}+\frac{1}{2}(1-2 \rho)^{2}}=\frac{2 \rho(1-\rho)}{1-2 \rho(1-\rho)}=0.173
\]

\includegraphics[max width=\textwidth, alt={}, center]{da0156b6-4133-434e-90d4-ec01b2070afe-211_680_607_105_746}\\
5.171 (a) In this case from Fresnel's equations

\[
I_{\perp}^{\prime}=I_{\perp} \frac{\sin ^{2}\left(\theta_{1}-\theta_{2}\right)}{\sin ^{2}\left(\theta_{1}+\theta_{2}\right)}
\]

we get

\[
I_{1}=\left(\frac{n^{2}-1}{n^{2}+1}\right)^{2} I_{0}=\rho I_{0} \text { say }
\]

then

\[
I_{2}=(1-\rho) I_{0}, I_{3}=\rho(1-\rho) I_{0}
\]

( \(\rho\) is invariant under the substitution \(n \rightarrow \frac{1}{n}\) )\\
finally

\[
I_{4}=(1-\rho)^{2} I_{0}=\frac{16 n^{4}}{\left(n^{2}+1\right)^{4}} I_{0}=0.726 I_{0} .
\]

(b) Suppose \(\rho^{\prime}=\) coefficient of reflection for the component of light whose electric vector oscillates at right angles to the incidence plane.

From Fresnel's equations

\[
\rho^{\prime}=\left(\frac{n^{2}-1}{n^{2}+1}\right)^{2}
\]

Then in the transmitted beam we have a partially polarized beam which is a superposition of two ( \(||\&|\) ) components with intensities

Thus

\[
P=\frac{1-\left(1-\rho^{\prime}\right)^{2}}{1+\left(1-\rho^{\prime}\right)^{2}}=\frac{\frac{1}{2} I_{0} \& \frac{1}{2} I_{0}\left(1-\rho^{\prime}\right)^{2}}{\left(n^{2}+1\right)^{4}+16 n^{4}}=\frac{1-0.726}{1+0.726} \approx 0.158
\]

5.172 (a) When natural light is incident on a glass plate at Brewster's angle, the transmitted light has

\[
H_{\|}=I_{0} / 2 \text { and } I_{1}^{\prime \prime}=\frac{16 n^{4}}{\left(n^{2}+1\right)^{4}} I_{0} / 2=\alpha^{4} I_{0} / 2
\]

where \(I_{0}\) is the incident intensity (see 5.171 a)\\
After passing through the \(2^{\text {nd }}\) plate we find

Thus after \(N\) plates

\[
\begin{gathered}
I_{\|}^{\prime \prime \prime \prime}=\frac{1}{2} I_{0} \quad \text { and } I_{1}^{\prime \prime \prime \prime}=\left(\alpha^{4}\right)^{2} \frac{1}{2} I_{0} \\
I_{\|}^{\text {trans }}=\frac{1}{2} I_{0} \\
I_{\perp}^{\text {trans }}=\alpha^{4 N} \frac{1}{2} I_{0} \\
P=\frac{1-\alpha^{4 N}}{1+\alpha^{4 N}} \text { where } \alpha=\frac{2 n}{1+n^{2}}
\end{gathered}
\]

Hence\\
(b) \(\alpha^{4}=0.726\) for \(n=\frac{3}{2}\).

Thus

\[
\begin{gathered}
P(N=1)=0.158, P(N=2)=0.310 \\
P(N=5)=0.663, P(N=10)=0.922
\end{gathered}
\]

5.173 (a) We decompose the natural light into two components with intensity \(I_{\|}=\frac{1}{2} I_{0}=I_{\perp}\) where || has its electric vector oscillating parallel to the plane of incidence and ⟂ has the same \(\perp^{r}\) to it.\\
By Fresnel's equations for normal incidence

\[
\frac{I_{\perp}^{\prime}}{I_{\perp}^{\prime}}=\lim _{\theta_{1} \rightarrow 0} \frac{\sin ^{2}\left(\theta_{1}-\theta_{2}\right)}{\sin ^{2}\left(\theta_{1}+\theta_{2}\right)}=\lim _{\theta_{1} \rightarrow 0}\left(\frac{\theta_{1}-\theta_{2}}{\theta_{1}+\theta_{2}}\right)^{2}=\left(\frac{n-1}{n+1}\right)^{2}=\rho
\]

similarly

Thus

\[
\begin{gathered}
\frac{I_{\|}^{\prime}}{I_{\|}}=\rho=\left(\frac{n-1}{n+1}\right)^{2} \\
\frac{I^{\prime}}{I}=\rho=\left(\frac{0.5}{2.5}\right)^{2}=\frac{1}{25}=0.04
\end{gathered}
\]

(b) The reflected light at the first surface has the intensity \(I_{1}=\rho I_{0}\)\\
Then the transmitted light has the intensity \(I_{2}=(1-\rho) I_{0}\)\\
At the second surface where light emerges from glass into air, the reflection coefficient is again \(\rho\) because \(\rho\) is invariant under the substitution \(n \rightarrow \frac{1}{n}\).\\
\includegraphics[max width=\textwidth, alt={}, center]{da0156b6-4133-434e-90d4-ec01b2070afe-212_270_357_1593_1076}

Thus \(I_{3}=\rho(1-\rho) I_{0}\) and \(I_{4}=(1-\rho)^{2} I_{0}\).

For \(N\) lenses the loss in luminous flux is then

\[
\frac{\Delta \Phi}{\Phi}=1-(1-\rho)^{2 N}=0.335 \text { for } N=5
\]

5.174 Suppose the incident light can be decomposed into waves with intensity \(I_{\|} \& I_{\perp}\) with oscillations of the electric vectors parallel and perpendicular to the plane of incidence.\\
For normal incidence we have from Fresnel equations

\[
I_{\perp}^{\prime}=I_{1}\left(\frac{\theta_{1}-\theta_{2}}{\theta_{1}+\theta_{2}}\right)^{2} \rightarrow I_{1}\left(\frac{n-1}{n+1}\right)^{2}
\]

where we have used \(\sin \theta \approx \theta\) for small \(\theta\).

Similarly

\[
I_{\|}^{\prime}=I_{\|}\left(\frac{n^{\prime}-1}{n^{\prime}+1}\right)^{2}
\]

\begin{center}
\includegraphics[max width=\textwidth, alt={}]{da0156b6-4133-434e-90d4-ec01b2070afe-213_391_435_329_988}
\end{center}

Then the refracted wave will be

\[
I^{\prime \prime} \|=I_{\|} \frac{4 n^{\prime}}{\left(n^{\prime}+1\right)^{2}} \text { and } I_{\perp}^{\prime \prime}=I_{\perp} \frac{4 n^{\prime}}{\left(n^{\prime}+1\right)^{2}}
\]

At the interface with glass

\[
I_{\perp}^{\prime \prime \prime}=I_{\perp}^{\prime \prime}\left(\frac{n^{\prime}-n}{n^{\prime}+n}\right)^{2} \text {, similarly for } I_{\|}^{\prime \prime \prime}
\]

we see that

\[
\frac{I_{\perp}^{\prime}}{I_{\perp}}=\frac{I_{\perp}^{\prime \prime \prime}}{I_{\perp}^{\prime \prime}} \quad \text { if } \quad n^{\prime}=\sqrt{n} \text {, similarly for } \| \text { component. }
\]

This shows that the light reflected as a fraction of the incident light is the same on the two surfaces if \(\boldsymbol{n}^{\prime}=\sqrt{\boldsymbol{n}}\).\\
Note:- The statement of the problem given in the book is incorrect. Actual amplitudes are not equal; only the reflectance is equal.\\
5.175 Here \(\theta_{1}=45^{\circ}\)

Hence

\[
\begin{gathered}
\sin \theta_{2}=\frac{1}{\sqrt{2}} \times \frac{1}{n}=\frac{2}{3 \sqrt{2}}=\frac{\sqrt{2}}{3}=0.4714 \\
\theta_{2}=\sin ^{-1} 0.4714=28.1^{\circ} \\
I_{\perp}^{\prime}=I_{\perp} \frac{\sin ^{2}\left(\theta_{1}-\theta_{2}\right)}{\sin ^{2}\left(\theta_{1}+\theta_{2}\right)} \\
=\frac{1}{2} I_{0}\left(\frac{\sin 16.9^{\circ}}{\sin 73.1^{\circ}}\right)^{2}=\frac{1}{2} I_{0} \times 0.0923
\end{gathered}
\]

\[
I_{\|}^{\prime}=\frac{1}{2} I_{0}\left(\frac{\tan 16.9}{\tan 73.1}\right)^{2}=\frac{1}{2} I_{0} \times 0.0085
\]

Thus\\
(a) Degree of polarization \(P\) of the reflected light

\[
=\frac{0.0838}{0.1008}=0.831
\]

(b) By conservation of energy

Thus

\[
\begin{aligned}
I_{\perp}^{\prime \prime} & =\frac{1}{2} I_{0} \times 0.9077 \\
I_{\|}^{\prime \prime} & =\frac{1}{2} I_{0} \times 0.9915 \\
P & =\frac{0.0838}{1.8982}=0.044
\end{aligned}
\]

5.176 The wave surface of a uniaxial crystal consists of two sheets of which one is a sphere while the other is an ellipsoid of revolution.\\
\includegraphics[max width=\textwidth, alt={}, center]{da0156b6-4133-434e-90d4-ec01b2070afe-214_271_277_838_503}\\
\includegraphics[max width=\textwidth, alt={}, center]{da0156b6-4133-434e-90d4-ec01b2070afe-214_275_277_840_904}

The optic axis is the line joining the points of contact.\\
To makes the appropriate Huyghen's construction we must draw the relevant section of the wave surface inside the crystal and determine the directions of the ordinary and extraordinary rays. The result is as shown in Fig. \(42(\mathrm{a}, \mathrm{b} \& \mathrm{c})\) of the answers\\
5.177 In a uniaxial crystal, an unpolarized beam of light (or even a polarized one) splits up into \(O\) (for ordinary) and \(E\) (for extraordinary) light waves. The direction of vibration in the \(O\) and \(E\) waves are most easily specified in terms of the \(O\) and \(E\) principal planes. The principal plane of the ordinary wave is defined as the plane containing the \(O\) ray and the optic axis. Similarly the principal plane of the \(E\) wave is the plane containing the \(E\) ray and the optic axis. In terms of these planes the following is true : The \(O\) vibrations are perpendicular to the principal plane of the \(O\) ray while the E vibrations are in the principal plane of\\
\includegraphics[max width=\textwidth, alt={}, center]{da0156b6-4133-434e-90d4-ec01b2070afe-214_373_490_1413_981}\\
the \(E\) ray.\\
When we apply this definition to the wollaston prism we find the following :

When unpolarized light enters from the left the \(O\) and \(E\) waves travel in the same direction but with different speeds. The \(O\) ray on the left has its vibrations normal to the plane of the paper and it becomes \(E\) ray on crossing the diagonal boundary of the two prism similarly the \(E\) ray on the left becomes \(O\) ray on the right. In this case Snell's law is applicable only approximately. The two rays are incident on the boundary at an angle \(\theta\) and in the right prism the ray which we have called \(O\) ray on the right emerges at

\[
\sin ^{-1} \frac{n_{e}}{n_{0}} \sin \theta=\sin ^{-1} \frac{1.658}{1.486} \times \frac{1}{2}=33.91^{\circ}
\]

where we have used

\[
n_{e}=1 \cdot 1658, n_{0}=1 \cdot 486 \text { and } \theta=30^{\circ} .
\]

Similarly the \(E\) ray on the right emerges within the prism at

\[
\sin ^{-1} \frac{n_{0}}{n_{e}} \sin \theta=26.62^{\circ}
\]

This means that the \(O\) ray is incident at the boundary between the prism and air at

\[
33 \cdot 91-30^{\circ}=3 \cdot 91^{\circ}
\]

and will emerge into air with a deviation of

\[
\begin{gathered}
\sin ^{-1} n_{0} \sin 3.91^{\circ} \\
=\sin ^{-1}\left(1.65 \delta \sin 3.91^{\circ}\right)=6.40^{\circ}
\end{gathered}
\]

The \(E\) ray will emerge with an opposite deviation of

Hence

\[
\begin{gathered}
\sin ^{-1}\left(n_{e} \sin \left(30^{\circ}-26.62^{\circ}\right)\right) \\
=\sin ^{-1}\left(1.486 \sin 3.38^{\circ}\right)=5.03^{\circ} \\
\delta \approx 6.49^{\circ}+5.03^{\circ}=11.52^{\circ}
\end{gathered}
\]

This result is accurate to first order in ( \(n_{e}-n_{0}\) ) because Snell's law holds when \(n_{e}=n_{0}\).\\
5.178 The wave is moving in the direction of \(z\) - axis\\
(a) Here \(E_{x}=E \cos (\omega t-k z), E_{y}=E \sin (\omega t-k z)\)

\[
\frac{E_{x}^{2}}{E^{2}}+\frac{E_{y}^{2}}{E^{2}}=1
\]

so the \(2 \times\) of the electric vector moves along a circle. For the right handed coordinate system this represents circular anticlockwise polarization when observed towards the incoming wave.\\
(b) \(E_{x}=E \cos (\omega t-k z), E_{y}=E \cos \left(\omega t-k z+\frac{\pi}{4}\right)\)\\
so

\[
\frac{E_{y}}{E}=\frac{1}{\sqrt{2}} \cos (\omega t-k z)-\frac{1}{\sqrt{2}} \sin (\omega t-k z)
\]

or

\[
\left(\frac{E_{y}}{E}-\frac{1}{\sqrt{2}} \frac{E_{x}}{E}\right)^{2}=\frac{1}{2}\left(1-\frac{E_{x}^{2}}{E^{2}}\right)
\]

or

\[
\frac{E_{y}^{2}}{E^{2}}+\frac{E_{x}^{2}}{E^{2}}-\sqrt{2} \frac{E_{y} E_{x}}{E^{2}}=\frac{1}{2}
\]

This is clearly an ellipse. By comparing with the previous case (compare the phase of \(E_{y}\) in the two cases) we see this represents elliptical clockwise polarization when viewed towards the incoming wave.\\
We write the equations as

Thus

\[
\begin{gathered}
E_{x}+E_{y}=2 E \cos \left(\omega t-k z+\frac{\pi}{8}\right) \cos \frac{\pi}{8} \\
E_{x}-E_{y}=+2 E \sin \left(\omega t-k z+\frac{\pi}{8}\right) \sin \frac{\pi}{8} \\
\left(\frac{E_{x}+E_{y}}{2 E \cos \frac{\pi}{8}}\right)^{2}+\left(\frac{E_{x}-E_{y}}{2 E \sin \frac{\pi}{8}}\right)^{2}=1
\end{gathered}
\]

Since \(\cos \frac{\pi}{8}>\sin \frac{\pi}{8}\), the major axis is in the direction of the straight line \(y=x\).\\
(c) \(E_{x}=E \cos (\omega t-k z)\)\\
\(E_{y}=E \cos (\omega t-k z+\pi)=-E \cos (\omega t-k z)\)\\
Thus the top of the electric vector traces the curve

\[
E_{y}=-E_{x}
\]

which is a straight line ( \(y=-x\) ). It corresponds to plane polarization.\\
5.179 For quartz\\
\(\left.\begin{array}{l}n_{e}=1.553 \\ n_{0}=1.544\end{array}\right\}\) for \(\lambda=589 \mathrm{~nm}\).\\
In a quartz plate cut parallel to its optic axis, plane polarized light incident normally from the left divides itself into \(O\) and \(E\) waves which move in the same direction with different speeds and as a result acquire a phase difference. This phase difference is

\[
\delta=\frac{2 \pi}{\lambda}\left(n_{e}-n_{0}\right) d
\]

where \(d=\) thickness of the plate. In general this makes the emergent light elliptically polarized.\\
(a) For emergent light to experience only rotation of polarization plane \(\delta=(2 k+1) \pi, k=0,1,2,3 \ldots\)

For this

\[
\begin{aligned}
d & =(2 k+1) \frac{\lambda}{2\left(n_{e}-n_{0}\right)} \\
& =(2 k+1) \frac{.589}{2 \times .009} \mu \mathrm{~m}=(2 k+1) \frac{.589}{18} \mathrm{~mm}
\end{aligned}
\]

The maximum value of \((2 k+1)\) for which this is less than 0.50 is obtained from

\[
\frac{0.50 \times 18}{0.589}=15.28
\]

Then we must take \(k=7\) and \(d=15 \times \frac{.589}{18}=0.4908 \mathrm{~mm}\)\\
(b) For circular polarization \(\delta=\frac{\pi}{2}\)\\
modulo

\[
2 \pi \text { i.e. } \delta=(4 k+1) \frac{\pi}{2}
\]

so

\[
d=(4 k+1) \frac{\lambda}{4\left(n_{e}-n_{0}\right)}=(4 k+1) \frac{0.589}{36}
\]

Now

\[
\frac{0.50 \times 36}{0.589}=30.56
\]

The nearest integer less than this which is of the form \(4 k+1\) is 29 for \(k=7\). For this

\[
d=0.4749 \mathrm{~mm}
\]

5.180 As in the previous problem the quartz plate introduces a phase difference \(\delta\) between the \(O \& E\) components. When \(\delta=\pi / 2\) (modulo \(\pi)\) the resultant wave is circularly polarized. In this case intensity is independent of the rotation of the rear prism. Now

\[
\begin{aligned}
\delta & =\frac{2 \pi}{\lambda}\left(n_{e}-n_{0}\right) d \\
& =\frac{2 \pi}{\lambda} 0.009 \times 0.5 \times 10^{-3} \mathrm{~m} \\
& =\frac{9 \pi}{\lambda}, \lambda \text { in } \mu \mathrm{m}
\end{aligned}
\]

For \(\lambda=0.50 \mu m . \delta=18 \pi\). The relevant values of \(\delta\) have to be chosen in the form

\[
\begin{gathered}
\left(k+\frac{1}{2}\right) \pi . \text { For } k=17,16,15 \text { we get } \\
\lambda=0.5143 \mu \mathrm{~m}, 0.5435 \mu \mathrm{~m} \text { and } 0.5806 \mu \mathrm{~m}
\end{gathered}
\]

These are the values of \(\lambda\) which lie between \(0.50 \mu \mathrm{~m}\) and \(0.60 \mu \mathrm{~m}\).\\
5.181 As in the previous two problems the quartz plate will introduce a phase difference \(\delta\). The light on passing through the plate will remain plane polarized only for \(\delta=2 k \pi\) or \((2 k+1) \pi\). In the latter case the plane of polarization of the light incident on the plate will be rotated by \(90^{\circ}\) by it so light passing through the analyser (which was originally crossed) will be a maximum. Thus dark bands will be observed only for those \(\lambda\) for which

\[
\delta=2 k \pi
\]

Now

\[
\begin{aligned}
\delta & =\frac{2 \pi}{\lambda}\left(n_{\mathrm{e}}-n_{0}\right) d=\frac{2 \pi}{\lambda} \times .009 \times 1.5 \times 10^{-3} \mathrm{~m} \\
& =\frac{27 \pi}{\lambda}(\lambda \text { in } \mu \mathrm{m})
\end{aligned}
\]

For

\[
\lambda=0.55 \text { we get } \delta=49.09 \pi
\]

Choosing \(\delta=48 \pi, 46 \pi, 44 \pi, 42 \pi\) we get \(\lambda=0.5625 \mu \mathrm{~m}, \lambda=0.5870 \mu \mathrm{~m}, \lambda=0.6136 \mu \mathrm{~m}\) and \(\lambda=0.6429 \mu \mathrm{~m}\). These are the only values between \(0.55 \mu \mathrm{~m}\) and \(0.66 \mu \mathrm{~m}\). Thus there are four bands.\\
5.182 Here

\[
\begin{aligned}
\delta & =\frac{2 \pi}{\lambda} \times 0.009 \times 0.25 \mathrm{~m} \\
& =\frac{4.5 \pi}{\lambda}, \lambda \text { in } \mu \mathrm{m}
\end{aligned}
\]

We check that for

\[
\begin{array}{ll}
\lambda=428.6 \mathrm{~nm} & \delta=10.5 \pi \\
\lambda=529.4 \mathrm{~nm} & \delta=8.5 \pi \\
\lambda=692.3 \mathrm{~nm} & \delta=6.5 \pi
\end{array}
\]

These are the only values of \(\lambda\) for which the plate acts as a quarter wave plate.\\
5.183 Between crossed Nicols, a quartz plate, whose optic axis makes \(45^{\circ}\) with the principal directions of the Nicols, must introduce a phase difference of \((2 k+1) \pi\) so as to transmit the incident light ( of that wavelength) with maximum intensity. For in this case the plane of polarization of the light emerging from the polarizer will be rotated by \(90^{\circ}\) and will go through the analyser undiminished. Thus we write for light of wavelengths 643 nm


\begin{align*}
\delta & =\frac{2 \pi \times 0.009}{0.643 \times 10^{-6}} \times d(\mathrm{~mm}) \times 10^{-3} \\
& =\frac{18 \pi d}{0.643}=(2 k+1) \pi \tag{1}
\end{align*}


To nearly block light of wavelength 564 nm we require


\begin{equation*}
\frac{18 \pi d}{0.564}=\left(2 k^{\prime}\right) \pi \tag{2}
\end{equation*}


We must have \(2 k^{\prime}>2 k+1\). For the smallest value of d we take \(2 k^{\prime}=2 k+2\).\\
Thus

\[
\begin{gathered}
0.643(2 k+1)=0.564 \times(2 k+2) \\
0.079 \times 2 k=0.564 \times 2-0.643
\end{gathered}
\]

This is not quite an integer but is close to one. This means that if we take \(2 k=6\) equations (1) can be satisfied exactly while equation (2) will hold approximately. Thus

\[
d=\frac{7 \times 0.643}{18}=0.250 \mathrm{~mm}
\]

5.184 If a ray traverses the wedge at a distance \(x\) below the joint, then the distance that the ray moves in the wedge is \(2 x \tan \frac{\Theta}{2}\) and this cause a phase difference

\[
\delta=\frac{2 \pi}{\lambda}\left(n_{e}-n_{0}\right) 2 x \tan \frac{\Theta}{2}
\]

between the \(E\) and \(O\) wave components of the ray. For a general \(x\) the resulting light is elliptically polarized and is not completely quenched by the analyser polaroid. The condition for complete quenching is\\
\includegraphics[max width=\textwidth, alt={}, center]{da0156b6-4133-434e-90d4-ec01b2070afe-218_468_297_1451_1059}

\[
\delta=2 k \pi \text { - dark fringe }
\]

That for maximum brightness is

\[
\delta=(2 k+1) \pi \text { - bright fringe. }
\]

The fringe width is given by

\[
\Delta x=\frac{\lambda}{2\left(n_{e}-n_{0}\right) \tan \frac{\Theta}{2}}
\]

Hence

\[
\left(n_{e}-n_{0}\right)=\frac{\lambda}{2 \Delta x \tan \Theta / 2}
\]

using

\[
\tan (\Theta / 2)=\tan 175^{\circ}=0.03055,
\]

\[
\lambda=0.55 \mu \mathrm{~m} \text { and } \Delta x=1 \mathrm{~mm} \text {, we get }
\]

\[
n_{e}-n_{0}=9.001 \times 10^{-3}
\]

5.185 Light emerging from the first polaroid is plane polarized with amplitude \(A\) where \(N_{1}\) is the principal direction of the polaroid and a vibration of amplitude can be resolved into two vibration : \(E\) wave with vibration along the optic axis of amplitude \(A \cos \varphi\) and the \(O\) wave with vibration perpendicular to the optic axis and having an amplitude \(A \sin \varphi\). These acquire a phase difference \(\delta\) on passing through the plate. The second polaroid transmits the components

\[
A \cos \varphi \cos \varphi^{\prime}
\]

and

\[
A \sin \varphi \sin \varphi^{\prime}
\]

\[
\mathrm{A}^{2}=\frac{I_{0}}{2}
\]

\begin{center}
\includegraphics[max width=\textwidth, alt={}]{da0156b6-4133-434e-90d4-ec01b2070afe-219_326_492_769_985}
\end{center}

What emerges from the second polaroid is a set of two plane polarized waves in the same direction and same plane of polarization but phase difference \(\boldsymbol{\delta}\). They interfere and produce a wave of amplitude squared\\
using

\[
R^{2}=A^{2}\left[\cos ^{2} \varphi \cos ^{2} \varphi^{\prime}+\sin ^{2} \varphi \sin ^{2} \varphi^{\prime}+2 \cos \varphi \cos \varphi^{\prime} \sin \varphi \sin \varphi^{\prime} \cos \delta\right],
\]

\[
\begin{aligned}
\cos ^{2}\left(\varphi-\varphi^{\prime}\right) & =\left(\cos \varphi \cos \varphi^{\prime}+\sin \varphi \sin \varphi^{\prime}\right)^{2} \\
& =\cos ^{2} \varphi \cos ^{2} \varphi^{\prime}+\sin ^{2} \varphi \sin ^{2} \varphi^{\prime}+2 \cos \varphi \cos \varphi^{\prime} \sin \varphi \sin \varphi^{\prime}
\end{aligned}
\]

we easily find

\[
R^{2}=A^{2}\left[\cos ^{2}\left(\varphi-\varphi^{\prime}\right)-\sin 2 \varphi \sin 2 \varphi^{\prime} \sin ^{2} \frac{\delta}{2}\right]
\]

Now

\[
A^{2}=I_{0} / 2 \text { and } R^{2}=I \text { so the result is }
\]

\[
I=\frac{1}{2} I_{0}\left[\cos ^{2}\left(\varphi-\varphi^{\prime}\right)-\sin 2 \varphi \sin 2 \varphi^{\prime} \sin ^{2} \frac{\delta}{2}\right]
\]

Special cases :- Crossed polaroids : Here \(\varphi-\varphi^{\prime}=90^{\circ}\) or \(\varphi^{\prime}=\varphi-90^{\circ}\) and \(2 \varphi^{\prime}=2 \varphi-180^{\circ}\)

Thus in this case

\[
I=I_{\perp}=\frac{1}{2} I_{0} \sin ^{2} 2 \varphi \sin ^{2} \frac{\delta}{2}
\]

Parallel polaroids : Here \(\varphi=\varphi^{\prime}\) and

\[
I=I_{\|}=\frac{1}{2} I_{0}\left(1-\sin ^{2} 2 \varphi \sin ^{2} \frac{\delta}{2}\right)
\]

With \(\delta=\frac{2 \pi}{\lambda} \Delta\), the conditions for the maximum and minimum are easily found to be that shown in the answer.\\
5.186 Let the circularly polarized light be resolved into plane polarized components of amplitude \(A_{0}\) with a phase difference \(\frac{\pi}{2}\) between then.

On passing through the crystal the phase difference becomes \(\delta+\frac{\pi}{2}\) and the components of the \(E\) and \(O\) wave in the direction \(N\) are respectively\\
\(A_{0} \cos \varphi\) and \(A_{0} \sin \varphi\)\\
\includegraphics[max width=\textwidth, alt={}, center]{da0156b6-4133-434e-90d4-ec01b2070afe-220_526_501_321_960}

They interfere to produce the amplitude squared

\[
\begin{aligned}
R^{2} & =A_{0}^{2} \cos ^{2} \varphi+A_{0}^{2} \sin ^{2} \varphi+2 A_{0}^{2} \cos \varphi \sin \varphi \cos \left(\delta+\frac{\pi}{2}\right) \\
& =A_{0}^{2}(1+\sin 2 \varphi \sin \delta)
\end{aligned}
\]

Hence

\[
I=I_{0}(1+\sin 2 \varphi \sin \delta)
\]

Here \(I_{0}\) is the intensity of the light transmitted by the polaroid when there is no crystal plate.\\
5.187 (a) The light with right circular polarization (viewed against the oncoming light, this means that the light vector is moving clock wise.) becomes plane polarized on passing through a quarter-wave plate. In this case the direction of oscillations of the electric vector of the electromagnetic wave forms an angle of \(+45^{\circ}\) with the axis of the crystal \(\mathrm{OO}^{\prime}\) (see Fig (a) below). In the case of left hand circular polarizations, this angle will be \(-45^{\circ}\) ( Fig (b)).\\
\includegraphics[max width=\textwidth, alt={}, center]{da0156b6-4133-434e-90d4-ec01b2070afe-220_374_390_1448_383}\\
\includegraphics[max width=\textwidth, alt={}, center]{da0156b6-4133-434e-90d4-ec01b2070afe-220_363_380_1444_827}\\
(b) If for any position of the plate the rotation of the polaroid (located behind the plate) does not bring about any variation in the intensity of the transmitted light, the incident light\\
is unpolarized (i.e. natural). If the intensity of the transmitted light can drop to zero on rotating the analyzer polaroid for some position of the quarter wave plate, the incident light is circularly polarized. If it varies but does not drop to zero, it must be a mixture of natural and circularly polarized light.\\
5.188 The light from \(P\) is plane polarized with its electric vector vibrating at \(45^{\circ}\) with the plane of the paper. At first the sample \(S\) is absent. Light from \(P\) can be resolved into components vibrating in and perpendicular to the plane of the paper. The former is the \(E\) ray in the left half of the Babinet compensator and the latter is the \(O\) ray. In the right half the nomenclature is the opposite. In the compensator the two components acquire a pahse difference which depends on the relative\\
\includegraphics[max width=\textwidth, alt={}, center]{da0156b6-4133-434e-90d4-ec01b2070afe-221_378_585_308_869}\\
position of the ray. If the ray is incident at a distance \(x\) above the central line through the compensator then the \(E\) ray acquires a phase\\
\(\frac{2 \pi}{\lambda}\left(n_{\mathrm{E}}(l-x)+n_{0}(l+x)\right) \tan \Theta\)\\
while the O ray acquires\\
\(\frac{2 \pi}{\lambda}\left(n_{\mathrm{o}}(l-x)+n_{\mathrm{E}}(l+x)\right) \tan \Theta\)\\
so the phase difference between the two reays is

\[
\frac{2 \pi}{\lambda}\left(n_{\mathrm{o}}-n_{\mathrm{E}}\right) 2 x \tan \Theta=\delta
\]

\includegraphics[max width=\textwidth, alt={}, center]{da0156b6-4133-434e-90d4-ec01b2070afe-221_351_392_836_964}\\
we get dark fringes when ever \(\delta=2 k \pi\)\\
because then the emergent light is the same as that coming from the polarizer and is quenched by the analyser. (If \(\delta=(2 k+1) \pi\), we get bright fringes because in this case, the plane of polarizaton of the emergent hight has rotated by \(90^{\circ}\) and is therefore fully transmitted by the analyser.\}\\
If follows that the fringe width \(\Delta x\) is given by

\[
\Delta x=\frac{\lambda}{2\left|n_{0}-n_{E}\right| \tan \Theta}
\]

(b) If the fringes are displaced upwards by \(\delta x\), then the path difference introduced by the sample between the \(O\) and the \(E\) rays must be such as to be exactly cancelled by the compensator. Thus

\[
\frac{2 \pi}{\lambda}\left[d\left(n_{\mathrm{O}}^{\prime}-n_{\mathrm{E}}^{\prime}\right)+\left(n_{\mathrm{E}}-n_{\mathrm{O}}\right) 2 \delta x \tan \Theta\right]=0
\]

or

\[
d\left(n_{\mathrm{O}}^{\prime}-n_{\mathrm{E}}^{\prime}\right)=-2\left(n_{\mathrm{E}}-n_{\mathrm{O}}\right) \delta x \Theta
\]

using \(\tan \Theta \approx \Theta\).\\
5.189 Light polarized along the \(x\)-direction (i.e. one whose electric vector has only an \(x\) component) and propagating along the \(z\)-direction can be decomposed into left and right circularly polarized light in accordance with the formula

\[
E_{x}=\frac{1}{2}\left(E_{x}+i E_{y}\right)+\frac{1}{2}\left(E_{x}-i E_{y}\right)
\]

On passing through a distance \(l\) of an active medium these acquire the phases \(\delta_{\mathrm{R}}=\frac{2 \pi}{\lambda} n_{\mathrm{R}} l\) and \(\delta_{1}=\frac{2 \pi}{\lambda} n_{l} l\) so we get for the complex amplitude

\[
\begin{aligned}
E^{\prime} & =\frac{1}{2}\left(E_{\mathrm{x}}+i E_{\mathrm{y}}\right) e^{\imath \delta_{R}}+\frac{1}{2}\left(E_{\mathrm{x}}-i E_{\mathrm{y}}\right) e^{\imath \delta_{L}} \\
& =e^{\imath \frac{\delta_{\mathrm{R}}+\delta_{l}}{2}}\left[\frac{1}{2}\left(E_{\mathrm{x}}+i E_{\mathrm{y}}\right) e^{\imath \delta / 2}+\frac{1}{2}\left(E_{\mathrm{x}}-i E_{\mathrm{y}}\right) e^{-\imath \delta / 2}\right] \\
& =e^{\imath \frac{\delta_{\mathrm{R}}+\delta_{l}}{2}}\left[E_{\mathrm{x}} \cos \frac{\delta}{2}-E_{\mathrm{y}} \sin \frac{\delta}{2}\right], \delta=\delta_{R}-\delta_{L}
\end{aligned}
\]

Apart from an over all phase \(\left(\delta_{\mathrm{R}}+\delta_{l}\right) / 2\) (which is irrelevant) this represents a wave whose plane of polarization has rotated by

\[
\frac{\delta}{2}=\frac{\pi}{\lambda}(\Delta n) l, \quad \Delta n=\left|n_{\mathrm{R}}-n_{l}\right|
\]

By definition this equals \(\alpha l\) so

\[
\begin{aligned}
\Delta n & =\frac{\alpha \lambda}{\pi} \\
& =\frac{589.5 \times 10^{-6} \mathrm{~m} \mathrm{~m} \times 21.72 \mathrm{deg} / \mathrm{m} \mathrm{~m}}{\pi} \times \frac{\pi}{180}(\mathrm{rad}) \\
& =\frac{.5895 \times 21.72}{180} \times 10^{-3} \\
& =0.71 \times 10^{-4}
\end{aligned}
\]

5.190 Plane polarized light on entering the wedge decomposes into right and left circularly polarized light which travel with different speeds in \(P\) and the emergent light gets its plane of polarization rotated by an angle which depends on the distance travelled.\\
Given that \(\Delta x=\) fringe width\\
\(\Delta x \tan \theta=\) difference in the path length traversed by two rays which form successive bright or dark fringes.

Thus

\[
\frac{2 \pi}{\lambda}\left|n_{\mathrm{R}}-n_{l}\right| \Delta x \tan \theta=2 \pi
\]

\begin{center}
\includegraphics[max width=\textwidth, alt={}]{da0156b6-4133-434e-90d4-ec01b2070afe-222_450_471_1387_964}
\end{center}

Thus

\[
\begin{aligned}
\alpha=\frac{\pi \Delta n}{\lambda} & =\pi / \Delta x \tan \theta \\
& =20.8 \mathrm{ang} \mathrm{deg} / \mathrm{m} \mathrm{~m}
\end{aligned}
\]

Let \(x=\) distance on the polaroid Pol as measured from a maximum. Then a ray that falls at this distance traverses an extra distance equal to

\[
\pm x \tan \theta
\]

and hence a rotation of \(\quad \pm \alpha x \tan \theta= \pm \frac{\pi x}{\Delta x}\)\\
By Malus' law the intensity at this point will be \(\sim \cos ^{2}\left(\frac{\pi x}{\Delta x}\right)\).\\
5.191 If \(I_{0}=\) intensity of natural light then\\
\(\frac{1}{2} I_{0}=\) intensity of light emerging from the polarizer nicol.\\
Suppose the quartz plate rotates this light by \(\varphi\), then the analyser will transmit

\[
\begin{gathered}
\frac{1}{2} I_{0} \cos ^{2}(90-\varphi) \\
=\frac{1}{2} I_{0} \sin ^{2} \varphi
\end{gathered}
\]

of this intensity. Hence

\[
\eta I_{0}=\frac{1}{2} I_{0} \sin ^{2} \varphi
\]

or

\[
\varphi=\sin ^{-1} \sqrt{2 \eta}
\]

But

\[
\varphi=\alpha d \text { so }
\]

\[
d_{\min }=\frac{1}{\alpha} \sin ^{-1} \sqrt{2 \eta}
\]

\begin{center}
\includegraphics[max width=\textwidth, alt={}]{da0156b6-4133-434e-90d4-ec01b2070afe-223_402_552_815_896}
\end{center}

For minimum \(d\) we must take the principal value of inverse sine. Thus using \(\alpha=17\) ang deg \(/ \mathrm{m} \mathrm{m}\).

\[
d_{\min }=2.99 \mathrm{~m} \mathrm{~m} .
\]

5.192 For light of wavelength 436 nm

\[
41 \cdot 5^{\circ} \times d=k \times 180^{\circ}=2 k \times 90^{\circ}
\]

(Light will be completely cut off when the quartz plate rotates the plane of polarization by a multiple of \(180^{\circ}\).) Here \(d=\) thickness of quartz plate in mm .\\
For natural incident light, half the light will be transmitted when the quartz rotates light by an odd multiple of \(90^{\circ}\). Thus

\[
31 \cdot 1^{\circ} \times d=\left(2 k^{\prime}+1\right) \times 90^{\circ}
\]

Now

\[
\frac{41 \cdot 5}{31 \cdot 1}=1 \cdot 3344 \approx \frac{4}{3}
\]

Thus

\[
\begin{aligned}
k & =2 \text { and } k^{\prime}=1 \text { and } \\
d & =\frac{4 \times 90}{41 \cdot 5}=8.67 \mathrm{~m} \mathrm{~m}
\end{aligned}
\]

5.193 Two effects are involved here : rotation of plane of polarizatin by sugar solution and the effect of that rotation on the scattering of light in the transverse direction. The latter is shown in the figure given below. It is easy to see from the figure that there will be no scattering of light in this transverse direction if the incident light has its electric vector parallel to the line of sight. In such a situation, we expect fringes to occur in the given experiment. From the given data we see that in a distance of 50 cm , the rotation of plane of polarization must be \(180^{\circ}\). Thus the specific rotation constant of sugar\\
\includegraphics[max width=\textwidth, alt={}, center]{da0156b6-4133-434e-90d4-ec01b2070afe-224_353_522_204_970}

\[
\begin{aligned}
& =\frac{\text { rotation constant }}{\text { concentration }} \\
& =\frac{180 / 50}{500 \frac{g}{l}} \mathrm{ang} / \mathrm{deg} / \mathrm{cm}=\frac{180}{5 \cdot 0 \mathrm{~d} \mathrm{~m} \times(\cdot 500 \mathrm{gm} / \mathrm{cc})} \\
& =72^{\circ} \mathrm{ang} \mathrm{deg} /(\mathrm{dm} \cdot \mathrm{gm} / \mathrm{cc})(1 \mathrm{~d} \mathrm{~m}=10 \mathrm{~cm})
\end{aligned}
\]

5.144 (a) in passing through the Kerr cell the two perpendicular components of the electric field will acquire a phase difference. When this phase difference equals \(90^{\circ}\) the emergent light will be circularly polarized because the two perpendicular components \(O \& E\) have the same magnitude since it is given that the direction of electric field \(E\) in the capacitor forms an angle of \(45^{\circ}\) with the principal directions of then icols. In this case the intensity of light that emerges from this system will be independent of the rotation of the analyser prism.\\
Now the phase difference introduced is given by

\[
\delta=\frac{2 \pi}{\lambda}\left(n_{\mathrm{e}}-n_{0}\right) l
\]

In the present case \(\delta=\frac{\pi}{2}\) (for minimum electric field)

\[
n_{\mathrm{e}}-n_{0}=\frac{\lambda}{4 l}
\]

Now

\[
n_{e}-n_{0}=B \lambda E^{2}
\]

so

\[
E_{\min }=\sqrt{\frac{1}{4 B l}}=10^{5} / \sqrt{88}=10.66 \mathrm{kV} / \mathrm{cm} .
\]

(b) If the applied electric field is

\[
E=E_{\mathrm{m}} \sin \omega t, \omega=2 \pi v
\]

than the Kerr cell introduces a time varying phase difference

\[
\begin{aligned}
\delta & =2 \pi B \mid E_{\mathrm{m}}^{2} \sin ^{2} \omega t \\
& =2 \pi \times 2 \cdot 2 \times 10^{-10} \times 10 \times\left(50 \times 10^{3}\right)^{2} \sin ^{2} \omega t \\
& =11 \pi \sin ^{2} \omega t
\end{aligned}
\]

In one half-cycle (i.e. in time \(\frac{\pi}{\omega}=T / 2=\frac{1}{2 v}\) )\\
this reaches the value \(2 k \pi\) when

\[
\begin{array}{r}
\sin ^{2} \omega t=0, \frac{2}{11}, \frac{4}{11}, \frac{6}{11}, \frac{8}{11}, \frac{10}{11} \\
\frac{2}{11}, \frac{4}{11}, \frac{6}{11}, \frac{8}{11}, \frac{10}{11}
\end{array}
\]

i.e. 11 times. On each of these occasions light will be interrupted. Thus light will be interrupted

\[
2 v \times 11=2.2 \times 10^{8} \text { times per second }
\]

(Light will be interrupted when the Kerr cell (placed between crossed Nicols) introduces a phase difference of \(2 k \pi\) and in no other case.)\\
5.195 From problem 189, we know that

\[
\Delta n=\frac{\alpha \lambda}{\pi}
\]

where \(\alpha\) is the rotation constant. Thus

\[
\Delta n=\frac{2 \alpha}{2 \pi / \lambda}=\frac{2 \alpha c}{\omega}
\]

On the other hand

\[
\alpha_{\mathrm{mag}}=V H
\]

Thus for the magnetic rotations \(\quad \Delta n=\frac{2 c V H}{\omega}\).\\
5.196 Part of the rotations is due to Faraday effect and part of it is ordinary optical rotation. The latter does not change sign when magnetic field is reversed. Thus

\[
\begin{aligned}
& \varphi_{1}=\alpha l+V l H \\
& \varphi_{2}=\alpha l-V l H
\end{aligned}
\]

Hence

\[
2 V l H=\left(\varphi_{1}-\varphi_{2}\right)
\]

or

\[
V=\left(\frac{\varphi_{1}-\varphi_{2}}{2}\right) / l H
\]

Putting the values

\[
V=\frac{510 \text { ang min }}{2 \times .3 \times 56.5} \times 10^{-3} \text { per } A=0.015 \text { ang min } / \mathrm{A}
\]

5.197 We write

\[
\varphi=\varphi_{\text {chemical }}+\varphi_{\text {magnetic }}
\]

We look against the transmitted beam and count the positive direction clockwise. The chemical part of the rotation is annulled by reversal of wave vector upon reflection.\\
Thus

\[
\varphi_{\text {chemical }}=\alpha l
\]

Since in effect there is a single transmission.

On the other hand

\[
\varphi_{\mathrm{mag}}=-N H V l
\]

To get the signs right recall that dextro rotatory compounds rotate the plane of vibration in a clockwise direction on looking against the oncoming beam. The sense of rotation of light vibration in Faraday effect is defined in terms of the direction of the field, positive rotation being that of a right handed screw advancing in the direction of the field. This is the opposite of the definition of \(\varphi_{\text {chemical }}\) for the present case. Finally

\[
\varphi=(\alpha-V N H) l
\]

(Note : If plane polarized light is reflected back \& forth through the same active medium in a magnetic field, the Faraday rotation increases with each traveresal.)\\
5.198 There must be a Faraday rotation by \(45^{\circ}\) in the opposite direction so that light could pass through the second polaroid. Thus\\
or

\[
\begin{aligned}
V l H_{\min } & =\pi / 4 \\
H_{\min } & =\frac{\pi / 4}{V l}=\frac{45 \times 60}{2.59 \times 0.26} \frac{A}{\mathrm{~m}} \\
& =4.01 \frac{\mathrm{k} \mathrm{~A}}{\mathrm{~m}}
\end{aligned}
\]

If the direction of magnetic field is changed then the sense of rotation will also change. Light will be completely quenched in the above case.\\
5.199 Let \(r=\) radius of the disc\\
then its moment of inertia about its axis \(=\frac{1}{2} m r^{2}\)\\
In time \(t\) the disc will acquire an angular momentum

\[
t \cdot \pi r^{2} \cdot \frac{I}{\omega}
\]

when circularly polarized light of intensity \(\boldsymbol{I}\) falls on it. By conservation of angular momentum this must equal

\[
\frac{1}{2} m r^{2} \cdot \omega_{0}
\]

where \(\omega_{0}=\) final angular velocity.\\
Equating

\[
t=\frac{m \omega \omega_{0}}{2 \pi I} .
\]

But

\[
\frac{\omega}{2 \pi}=v=\frac{c}{\lambda} \text { so } t=\frac{m c \omega_{0}}{I \lambda} .
\]

Substituting the values of the various quantities we get

\[
t=11.9 \text { hours }
\]

\subsection*{5.5 DISPERSION AND ABSORPTION OF LIGHT}
5.200 In a travelling plane electromagnetic wave the intensity is simply the time averaged magnitude of the Poynting vector :-

\[
I=\langle | \vec{E} \times \vec{H}| \rangle=\left\langle\sqrt{\frac{\varepsilon_{0}}{\mu_{0}}} E^{2}\right\rangle=\left\langle c \varepsilon_{0} E^{2}\right\rangle
\]

on using

\[
c=\frac{1}{\sqrt{\varepsilon_{0} \mu_{0}}}, \quad E \sqrt{\varepsilon_{0}}=H \sqrt{\mu_{0}} .
\]

(see chapter 4.4 of the book).\\
Now time averaged value of \(E^{2}\) is \(E_{0}^{2} / 2\) so

\[
I=\frac{1}{2} c \varepsilon_{0} E_{0}^{2} \quad \text { or } E_{0}=\sqrt{\frac{2 I}{c \varepsilon_{0}}},
\]

(a) Represent the electric field at any point by \(E=E_{0} \sin \omega t\). Then for the electron we have the equation.\\
so

\[
\begin{aligned}
m \ddot{x} & =e E_{0} \sin \omega t \\
x & =-\frac{e E_{0}}{m \omega^{2}} \sin \omega t
\end{aligned}
\]

The ampitude of the forced oscillation is

\[
\frac{e E_{0}}{m \omega^{2}}=\frac{e}{m \omega^{2}} \sqrt{\frac{2 I}{c \varepsilon_{0}}}=5.1 \times 10^{-16} \mathrm{~cm}
\]

The velocity amplitude is clearly

\[
\frac{e E_{0}}{m \omega}=5.1 \times 10^{-16} \times 3.4 \times 10^{15}=1.73 \mathrm{~cm} / \mathrm{sec}
\]

(b) For the electric force\\
\(F_{e}=\) amplitude of the electric force

\[
=e E_{0}
\]

For the magnetic force (which we have neglected above), it is

\[
\begin{gathered}
(e \vee B)=\left(e \vee \mu_{0} H\right) \\
=e v E \sqrt{\varepsilon_{0} \mu_{0}}=e v \frac{E}{c}
\end{gathered}
\]

writing \(\mathrm{v}=-\mathrm{v}_{0} \cos \omega t\)\\
where

\[
\mathrm{v}_{0}=\frac{e E_{0}}{m \omega}
\]

we see that the magnetic force is apart from a sign

\[
\frac{e v_{0} E_{0}}{2 c} \sin 2 \omega t
\]

Hence \(\frac{F_{m}}{F_{e}}=\) Ratio of amplitudes of the two forces

\[
=\frac{v_{o}}{2 c}=2.9 \times 10^{-11}
\]

This is negligible and justifies the neglect of magnetic field of the electromagnetic wave in calculating \(\nu_{0}\).\\
5.201 (a) It turns out that one can neglect the spatial dependence of the electric field as well as the magnetic field. Thus for a typical electron

\[
m \overrightarrow{\vec{r}}=e \overrightarrow{E_{0}} \sin \omega t
\]

so \(\vec{r}=-\frac{e \overrightarrow{E_{0}}}{m \omega^{2}} \sin \omega t\) (neglecting any nonsinusoidal part).\\
The ions will be practically unaffected. Then\\
and

\[
\begin{gathered}
\vec{P}=n_{o} e \vec{r}=-\frac{n_{0} e^{2}}{m \omega^{2}} \vec{E} \\
\vec{D}=\varepsilon_{0} \vec{E}+\vec{P}=\varepsilon_{0}\left(1-\frac{n_{0} e^{2}}{\varepsilon_{0} m \omega^{2}}\right) \vec{E}
\end{gathered}
\]

Hence the permittivity

\[
\varepsilon=1-\frac{n_{0} e^{2}}{\varepsilon_{0} m \omega^{2}}
\]

(b) The phase velocity is given by

So

Thus

\[
\begin{gathered}
\mathrm{v}=\omega / K=\frac{c}{\sqrt{\varepsilon}} \\
c k=\omega \sqrt{1-\frac{\omega_{P}^{2}}{\omega^{2}}}, \omega_{P}^{2}=\frac{n_{0} e^{2}}{\varepsilon_{0} m} \\
\omega^{2}=c^{2} k^{2}+\omega_{P}^{2} \\
\mathrm{v}=c \sqrt{1+\frac{\omega_{P}^{2}}{c^{2} k^{2}}}=c \sqrt{1+\left(\frac{n_{0} e^{2}}{4 \pi^{2} m c^{2} \varepsilon_{0}}\right) \lambda^{2}}
\end{gathered}
\]

5.202 From the previous problem

\[
\begin{aligned}
n^{2} & =1-\frac{n_{0} e^{2}}{\varepsilon_{0} m \omega^{2}} \\
& =1-\frac{n_{0} e^{2}}{4 \pi^{2} \varepsilon_{0} m v^{2}}
\end{aligned}
\]

Thus \(n_{0}=\left(4 \pi^{2} v^{2} m \varepsilon_{0} / e^{2}\right)\left(1-n^{2}\right)=2.36 \times 10^{7} \mathrm{~cm}^{-3}\)\\
5.203 For hard \(x\) - rays, the electrons in graphite will behave as if nearly free and the formula of previous problem can be applied. Thus\\
and

\[
\begin{aligned}
& n^{2}=1-\frac{n_{0} e^{2}}{\varepsilon_{0} m \omega^{2}} \\
& n \approx 1-\frac{n_{0} e^{2}}{2 \varepsilon_{0} m \omega^{2}}
\end{aligned}
\]

on taking square root and neglecting higher order terms.\\
So

\[
n-1=-\frac{n_{0} e^{2}}{2 \varepsilon_{0} m \omega^{2}}=-\frac{n_{0} e^{2} \lambda^{2}}{8 \pi^{2} \varepsilon_{0} m e^{2}}
\]

We calculate \(n_{0}\) as follows : There are \(6 \times 6.023 \times 10^{23}\) electrons in 12 gms of graphite of density \(1.6 \mathrm{gm} /\) c.c. Thus

\[
n_{0}=\frac{6 \times 6.023 \times 10^{23}}{(12 / 1.6)} \text { per c.c }
\]

Using the values of other constants and \(\lambda=50 \times 10^{-12}\) metre we get

\[
n-1=-5 \cdot 4 \times 10^{-7}
\]

5.204 (a) The equation of the electron can (under the stated conditions) be written as

\[
m \ddot{x}+\gamma \dot{x}+k x=e E_{0} \cos \omega t
\]

To solve this equation we shall find it convenient to use complex displacements. Consider the equation

\[
m \ddot{z}+\gamma \dot{z}+k z=e E_{0} e^{-i \omega t}
\]

Its solution is

\[
z=\frac{e E_{0} e^{-i \omega t}}{-m \omega^{2}-i \gamma \omega+k}
\]

(we ignore transients.)\\
Writing

\[
\beta=\frac{\gamma}{2 m}, \omega_{0}^{2}=\frac{k}{m}
\]

we find

\[
z=\frac{e E_{0}}{m} e^{-i \omega t} /\left(\omega_{o}^{2}-\omega^{2}-2 i \beta \omega\right)
\]

Now \(x=\) Real part of \(z\)

\[
=\frac{e E_{0}}{m} \cdot \frac{\cos (\omega t+\varphi)}{\sqrt{\left(\omega_{0}^{2}-\omega^{2}\right)^{2}+4 \beta^{2} \omega^{2}}}=a \cos (\omega t+\varphi)
\]

where

\[
\begin{gathered}
\tan \varphi=\frac{2 \beta \omega}{\omega^{2}-\omega_{0}^{2}} \\
\left.\sin \varphi=-\frac{2 \beta \omega}{\sqrt{\left(\omega^{2}-\omega_{0}^{2}\right)^{2}+4 \beta^{2} \omega^{2}}}\right)
\end{gathered}
\]

(b) We calculate the power absorbed as

\[
\begin{aligned}
P & =\langle F \dot{x}\rangle=\left\langle e E_{0} \cos \omega t(-\omega a \sin (\omega t+\varphi))\right\rangle \\
& =e E_{0} \cdot \frac{e E_{0}}{m} \frac{1}{2} \cdot \frac{2 \beta \omega}{\left(\omega_{0}^{2}-\omega^{2}\right)^{2}+4 \beta^{2} \omega^{2}} \cdot \omega=\left(\frac{e E_{0}}{m}\right)^{2} \frac{\beta m \omega^{2}}{\left(\omega_{0}^{2}-\omega^{2}\right)^{2}+4 \beta^{2} \omega^{2}}
\end{aligned}
\]

This is clearly maximum when \(\omega_{0}=\omega\) because \(P\) can be written as\\
and

\[
\begin{aligned}
P & =\left(\frac{e E_{0}}{m}\right)^{2} \frac{\beta m}{\left(\frac{\omega_{0}^{2}}{\omega}-\omega\right)^{2}+4 \beta^{2}} \\
P_{\max } & =\frac{m}{4 \beta}\left(\frac{e E_{0}}{m}\right)^{2} \text { for } \omega=\omega_{0} .
\end{aligned}
\]

\(P\) can also be calculated from \(P=\langle\gamma \dot{x} \cdot \dot{x}\rangle\)

\[
=\left(\gamma \omega^{2} a^{2} / 2\right)=\frac{\beta m \omega^{2}\left(e E_{0} / m\right)^{2}}{\left(\omega_{0}^{2}-\omega^{2}\right)^{2}+4 \beta^{2} \omega^{2}} .
\]

5.205 Let us write the solutions of the wave equation in the form

\[
A=A_{0} e^{i(\omega t-k x)}
\]

where \(k=\frac{2 \pi}{\lambda}\) and \(\lambda\) is the wavelength in the medium. If \(n^{\prime}=n+i \chi\), then

\[
k=\frac{2 \pi}{\lambda_{0}} n^{\prime}
\]

( \(\lambda_{0}\) is the wavelength in vacuum) and the equation becomes

\[
A=A_{0} e^{\chi^{\prime} x} \exp \left(i\left(\omega t_{1}-k^{\prime} x\right)\right)
\]

where \(\chi^{\prime}=\frac{2 \pi}{\lambda_{0}} \chi\) and \(k^{\prime}=\frac{2 \pi}{\lambda_{0}} n\). In real form,

\[
A=A_{0} e^{\chi^{\prime} x} \cos \left(\omega t-k^{\prime} x\right)
\]

This represents a plane wave whose amplitude diminishes as it propagates to the right (provided \(\chi^{\prime}<0\) ).\\
when \(n^{\prime}=i \chi\), then similarly

\[
A=A_{0} e^{\chi^{\prime} x} \cos \omega t
\]

(on putting \(n=0\) in the above equation).\\
This represents a standing wave whose amplitude diminishes as one goes to the right (if \(\chi^{\prime}<0\) ). The wavelength of the wave is infinite ( \(k^{\prime}=0\) ).\\
Waves of the former type are realized inside metals as well as inside dielectrics when there is total reflection. (penetration of wave).\\
5.206 In the plasma radio waves with wavelengths exceeding \(\lambda_{0}\) are not propagated. We interpret this to mean that the permittivity becomes negative for such waves. Thus

Hence

\[
0=1-\frac{n_{0} e^{2}}{\varepsilon_{0} m \omega^{2}} \text { if } \omega=\frac{2 \pi}{\lambda_{0}} \frac{c}{2}
\]

\[
\frac{n_{0} e^{2} \lambda_{0}^{2}}{4 \pi^{2} \varepsilon_{0} m c^{2}}=1
\]

or

\[
n_{0}=\frac{4 \pi^{2} \varepsilon_{0} m c^{2}}{e^{2} \lambda_{0}^{2}}=1.984 \times 10^{9} \text { per c.c }
\]

5.207 By definition

\[
u=\frac{d \omega}{d k}=\frac{d}{d k}(\mathrm{v} k) \quad \text { as } \omega=\mathrm{v} k=\mathrm{v}+k \frac{d \mathrm{v}}{d k}
\]

Now

\[
k=\frac{2 \pi}{\lambda} \quad \text { so } \quad d k=-\frac{2 \pi}{\lambda^{2}} d \lambda
\]

Thus

\[
u=v-\lambda \frac{d v}{d \lambda}
\]

Its interpretation is the following :\\
\(\left(\frac{d v}{d \lambda}\right)_{\lambda=\lambda^{\prime}}\) is the slope of the \(v-\lambda\) curve at \(\lambda=\lambda^{\prime} . V^{\prime}\)\\
\includegraphics[max width=\textwidth, alt={}, center]{da0156b6-4133-434e-90d4-ec01b2070afe-231_424_543_664_928}

Thus as is obvious from the diagram\\
\(\mathrm{v}^{\prime}=\mathrm{v}\left(\lambda^{\prime}\right)-\lambda^{\prime}\left(\frac{d \mathrm{v}}{d \lambda}\right)_{\lambda=\lambda^{\prime}}\) is the group velocity for \(\lambda=\lambda^{\prime}\).\\
5.208 (a) \(v=a / \sqrt{\lambda}, a=\) constant

Then

\[
\begin{aligned}
u & =v-\lambda \frac{d v}{d \lambda} \\
& =\frac{a}{\sqrt{\lambda}}-\lambda\left(-\frac{1}{2} a \lambda^{-3 / 2}\right)=\frac{3}{2} \cdot \frac{a}{\sqrt{\lambda}}=\frac{3}{2} v
\end{aligned}
\]

(b) \(\mathrm{v}=b k=\omega k, b=\) constant\\
so

\[
\omega=b k^{2} \quad \text { and } \quad u=\frac{d \omega}{d k}=2 b k=2 \mathrm{v} .
\]

(c) \(\mathrm{v}=\frac{c}{\omega^{2}}, c=\) constant \(=\frac{\omega}{k}\).\\
so

\[
\omega^{3}=c k \quad \text { or } \omega=c^{1 / 3} k^{1 / 3}
\]

Thus

\[
u=\frac{d \omega}{d k}=c^{1 / 3} \frac{1}{3} k^{-2 / 3}=\frac{1}{3} \frac{\omega}{k}=\frac{1}{3} \mathrm{v}
\]

5.209 We have

\[
u \mathrm{v}=\frac{\omega}{k} \frac{d \omega}{d k}=c^{2}
\]

Integrating we find\\
so

\[
\omega^{2}=A+c^{2} k^{2}, A \text { is a constant. }
\]

and

\[
\begin{aligned}
k & =\frac{\sqrt{\omega^{2}-A}}{c} \\
v & =\frac{\omega}{k}=\frac{c}{\sqrt{1-\frac{A}{\omega^{2}}}}
\end{aligned}
\]

writing this as \(c / \sqrt{\varepsilon(\omega)}\) we get \(\varepsilon(\omega)=1-\frac{A}{\omega^{2}}\)\\
(A can be +ve or negative)\\
5.210 The phase velocity of light in the vicinity of \(\lambda=534 n m=\lambda_{0}\) is obtained as

\[
v\left(\lambda_{0}\right)=\frac{c}{n\left(\lambda_{0}\right)}=\frac{3 \times 10^{8}}{1.640}=1.829 \times 10^{8} \mathrm{~m} / \mathrm{s}
\]

To get the group velocity we need to calculate

\[
\left(\frac{d n}{d \lambda}\right)_{\lambda=\lambda_{0}} . \text { We shall use linear }
\]

interpolation in the two intervals. Thus

\[
\begin{aligned}
& \left(\frac{d n}{d \lambda}\right)_{\lambda=521.5}=-\frac{.007}{25}=-28 \times 10^{-5} \text { per nm } \\
& \left(\frac{d n}{d \lambda}\right)_{\lambda=561.5}=-\frac{.01}{55}=-18.2 \times 10^{-5} \text { per nm }
\end{aligned}
\]

There \((d n / d \lambda)\) values have been assigned to the mid-points of the two intervals. Interpolating again we get

\[
\left(\frac{d n}{d \lambda}\right)_{\lambda=534}=\left[-28+\frac{9 \cdot 8}{40} \times 12 \cdot 5\right] \times 10^{-5} \text { per mm }=-24 \cdot 9 \times 10^{-5} \text { per n m. }
\]

Finally

At \(\lambda=534\)

\[
u=\frac{c}{n}-\lambda \frac{d}{d \lambda}\left(\frac{c}{n}\right)=\frac{c}{n}\left[1+\frac{\lambda}{n}\left(\frac{d n}{d \lambda}\right)\right]
\]

\[
u=\frac{3 \times 10^{8}}{1.640}\left[1-\frac{534}{1.640} \times 24.9 \times 10^{-5}\right] \mathrm{m} / \mathrm{s}=1.59 \times 10^{8} \mathrm{~m} / \mathrm{s}
\]

\subsection*{5.211 We write}
(since \(k=\frac{2 \pi}{\lambda}\) ). Suppose a wavetrain at time \(t=0\) has the form

\[
F(x, 0)=\int f(k) e^{i k x} d k
\]

Then at time \(t\) it will have the form

\[
\begin{aligned}
F(x, t) & =\int f(k) e^{i k x-i \omega t} d k \\
& =\int f(k) e^{i k x-i(2 \pi b+a k) t}=\int f(k) e^{i k(x-a t)} e^{-i 2 \pi b t} d k
\end{aligned}
\]

At \(t=\frac{1}{b}=\tau\)

\[
F(x, \tau)=F(x-a \tau, 0)
\]

so at time \(t=\tau\) the wave train has regained its shape though it has advanced by \(a \tau\).

On passing through the first (polarizer) Nicol the intensity of light becomes \(\frac{1}{2} I_{0}\) because one of the components has been cut off. On passing through the solution the plane of polarization of the light beam will rotate by\\
\(\varphi=V l H\)\\
and its intensity will also decrease by a factor \(e^{-x l}\). The plane of vibraton of the light wave will then make an angle \(90^{\circ}-\varphi\) with the principal direction of the analyzer Nicol. Thus by Malus' law the intensity of light coming out of the second Nicol will be

\[
\begin{gathered}
\frac{1}{2} I_{0} \cdot e^{-\chi l} \cdot \cos ^{2}\left(90^{\circ}-\varphi\right) \\
=\frac{1}{2} I_{0} e^{-\chi l} \sin ^{2} \varphi .
\end{gathered}
\]

5.213 (a) The multiple reflections are shown below. Transmission gives a factor ( \(1-\rho\) ) while reflections give factors of \(\rho\). Thus the transmitted intensity assuming incoheren light is

\[
\begin{aligned}
& (1-\rho)^{2} I_{0}+(1-\rho)^{2} \rho^{2} I_{0}+(1-\rho)^{2} \rho^{4} I_{0}+\ldots \\
& =(1-\rho)^{2} I_{0}\left(1+\rho^{2}+\rho^{4}+\rho^{6}+\ldots\right) \\
& =(1-\rho)^{2} I_{0} \times \frac{1}{1-\rho^{2}}=I_{0} \frac{1-\rho}{1+\rho}
\end{aligned}
\]

\includegraphics[max width=\textwidth, alt={}, center]{da0156b6-4133-434e-90d4-ec01b2070afe-233_668_497_1275_975}\\
(b) When there is absorption, we pick up a factor \(\sigma=e^{-\chi d}\) in each traversal of the plate. Thus we get

\[
\begin{aligned}
(1-\rho)^{2} & \sigma I_{0}+(1-\rho)^{2} \sigma^{3} \rho^{2} I_{0}+(1-\rho)^{2} \sigma^{5} \rho^{4} I_{0}+\ldots \\
& =(1-\rho)^{2} \sigma I_{0}\left(1+\sigma^{2} \rho^{2}+\sigma^{4} \rho^{4}+\ldots\right) \\
& =I_{0} \frac{\sigma(1-\rho)^{2}}{1-\sigma^{2} \rho^{2}}
\end{aligned}
\]

5.214 We have

\[
\begin{aligned}
& \tau_{1}=e^{-\chi d_{1}}(1-\rho)^{2} \\
& \tau_{2}=e^{-\chi d_{2}}(1-\rho)^{2}
\end{aligned}
\]

where \(\rho\) is the reflectivity; see previous problem, multiple reflection have been ignored.\\
Thus

\[
\begin{gathered}
\frac{\tau_{1}}{\tau_{2}}=e^{\chi\left(d_{2}-d_{1}\right)} \\
\chi=\frac{\ln \left(\frac{\tau_{1}}{\tau_{2}}\right)}{d_{2}-d_{1}}=0.35 \mathrm{~cm}^{-1}
\end{gathered}
\]

or\\
5.215 On each surface we pick up a factor \((1-\rho)\) from reflection and a factor \(e^{-\chi l}\) due to absorption in each plate.\\
Thus

Thus

\[
\begin{gathered}
\tau=(1-\rho)^{2 N} e^{-\chi N l} \\
\chi=\frac{1}{N l} \ln \frac{(1-\rho)^{2 N}}{\tau}=0.034 \mathrm{~cm}^{-1}
\end{gathered}
\]

5.216 Apart from the factor \((1-\rho)\) on each end face of the plate, we shall get a factor due to absorptions. This factor can be calculated by assuming the plate to consist of a large number of very thin slab within each of which the absorption coefficient can be assumed to be constant. Thus we shall get a product like

\[
\ldots \ldots e^{-\chi(x) d x} e^{-\chi(x+d x) d x} e^{-\chi(x+2 d x) d x} \ldots \ldots
\]

This product is nothing but

\[
e^{-\int_{0}^{1} x(x) d x}
\]

Now \(\chi(0)=\chi_{1}, \chi(l)=\chi_{2}\) and variation\\
with \(x\) is linear so \(\chi(x)=\chi_{1}+\frac{x}{l}\left(\chi_{2}-\chi_{1}\right)\)\\
Thus the factor becomes

\[
e^{-\int_{0}^{l}\left[x_{1}+\frac{x}{l}\left(x_{2}-x_{1}\right)\right] d x}=e^{-\frac{1}{2}\left(x_{1}+x_{2}\right) l}
\]

5.217 The spectral density of the incident beam (i.e. intensity of the components whose wave length lies in the interval \(\lambda \& \lambda+d \lambda)\) is

\[
\frac{I_{0}}{\lambda_{2}-\lambda_{1}} d \lambda, \lambda_{1} \leq \lambda \leq \lambda_{2}
\]

The absorption factor for this component is

\[
e^{-\left[x_{1}+\frac{\lambda+\lambda_{1}}{\lambda_{2}-\lambda_{1}}\left(x_{2}-x_{1}\right)\right] l}
\]

and the transmission factor due to reflection at the surface is \((1-\rho)^{2}\). Thus the intensity of the transmitted beam is

\[
\begin{gathered}
(1-\rho)^{2} \frac{I_{0}}{\lambda_{2}-\lambda_{1}} \int_{\lambda_{1}}^{\lambda_{2}} d \lambda e^{-l\left[x_{1}+\frac{\lambda-\lambda_{1}}{\lambda_{2}-\lambda_{1}}\left(x_{2}-x_{1}\right)\right]} \\
=(1-\rho)^{2} \frac{I_{0}}{\lambda_{2}-\lambda_{1}} e^{-x_{1} l}\left(\frac{1-e^{-\left(x_{2}-x_{1}\right) l}}{\left(x_{2}-x_{1}\right)^{\prime} m e}\right) \times\left(\lambda_{2}-\lambda_{1}\right)=(1-\rho)^{2} I_{0} \frac{e^{-x_{1} l}-e^{-x_{2} l}}{\left(x_{2}-x_{1}\right) l}
\end{gathered}
\]

5.218 At the wavelength \(\lambda_{0}\), the absorption coefficient vanishes and loss in transmission is entirely due to reflection. This factor is the same at all wavelengths and therefore cancels out in calculating the pass band and we need not worry about it. Now

\[
\begin{aligned}
T_{0} & =\left(\text { transmissivity at } \lambda=\lambda_{0}\right)=(1-\rho)^{2} \\
T & =\text { transmissivity at } \lambda=(1-\rho)^{2} e^{-\chi(\lambda) d}
\end{aligned}
\]

The edges of the passband are \(\lambda_{0} \pm \frac{\Delta \lambda}{2}\) and at the edge

Thus

\[
\frac{T}{T_{0}}=e^{-\alpha d\left(\frac{\Delta \lambda}{2 \lambda_{0}}\right)^{2}}=\eta
\]

or

\[
\begin{aligned}
& \frac{\Delta \lambda}{2 \lambda_{0}}=\sqrt{\left(\ln \frac{1}{\eta}\right) / \propto d} \\
& \Delta \lambda=2 \lambda_{0} \sqrt{\frac{1}{\alpha d}\left(\ln \frac{1}{\eta}\right)}
\end{aligned}
\]

5.219 We have to derive the law of decrease of intensity in ań absorbing medium taking in to account the natural geometrical fall-off (inverse sequare law) as well as absorption.\\
Consider a thin spherical shell of thickness \(d x\) and internal radius \(x\). Let \(I(x)\) and \(I(x+d x)\) be the intersities at the inner and outer surfaces of this shell.

Then

\[
4 \pi x^{2} I(x) e^{-x d x}=4 \pi(x+d x)^{2} I(x+d x)
\]

Except for the factor \(e^{-\chi d x}\) this is the usual equation. We rewrite this as

\[
x^{2} I(x)=I(x+d x)(x+d x)^{2}(1+\chi d x)
\]

\[
\begin{gathered}
=\left(I+\frac{d I}{d x} d x\right)\left(x^{2}+2 x d x\right)(1+\chi d x) \\
x^{2} \frac{d I}{d x}+\chi x^{2} I+2 x I=0 \\
\frac{d}{d x}\left(x^{2} I\right)+\chi\left(x^{2} I\right)=0 \\
x^{2} I=C e^{-\chi x}
\end{gathered}
\]

or

Hence\\
so\\
where \(C\) is a constant of integration.\\
In our case we apply this equation for \(a \leq x \leq b\)\\
For \(x \leq a\) the usual inverse square law gives

\[
I(a)=\frac{\Phi}{4 \pi a^{2}}
\]

Hence

\[
C=\frac{\Phi}{4 \pi} e^{\chi a}
\]

and

\[
I(b)=\frac{\Phi}{4 \pi b^{2}} e^{-\chi(b-a)}
\]

This does not take into account reflections. When we do that we get

\[
I(b)=\frac{\Phi}{4 \pi b^{2}}(1-\rho)^{2} e^{-\chi(b-a)}
\]

5.220 The transmission factor is \(e^{-\mu d}\) and so the intensity will decrease

\[
=e^{3.6 \times 11.3 \times 0.1}=58.4 \text { timestimes }
\]

(we have used \(\mu=(\mu / \rho) \times \rho\) and used the known value of density of lead).\\
5.221 We require \(\mu_{P b} d_{P b}=\mu_{A l} d_{A l}\)\\
or

\[
\begin{aligned}
\left(\frac{\mu_{P b}}{\rho_{P b}}\right) \rho_{P b} d_{P b} & =\left(\frac{\mu_{A l}}{\rho_{A l}}\right) \rho_{A l} d_{A l} \\
72.0 \times 11.3 \times d_{P b} & =3.48 \times 2.7 \times 2.6 \\
d_{P b} & =0.3 \mathrm{~m} \mathrm{~m}
\end{aligned}
\]

\(5.222 \quad \frac{1}{2}=e^{-\mu d}\)\\
or

\[
d=\frac{\ln 2}{\mu}=\frac{\ln 2}{\left(\frac{\mu}{\rho}\right) \rho}=0.80 \mathrm{~cm}
\]

5.223 We require \(N\) plates where

\[
\left(\frac{1}{2}\right)^{N}=\frac{1}{50} \quad \text { So } \quad N=\frac{\ln 50}{\ln 2}=5.6
\]

\subsection*{5.6 OPTICS OF MOVING SOURCES}
5.224 In the Fizean experiment, light disappears when the wheel rotates to bring a tooth in the position formerly occupied by a gap in the time taken by light to go from the wheel to the mirror and back. Thusdistance travelled \(=2 l\) Suppose the \(m^{\text {th }}\) tooth after the gap has come in place of the latter. Then time taken

\[
\begin{aligned}
& =\frac{2(m-1)+1}{2 z n_{1}} \text { sec. in the first case } \\
& =\frac{2 m+1}{2 z n_{2}} \text { sec in the second case }=\frac{1}{z\left(n_{2}-n_{1}\right)}
\end{aligned}
\]

Then

\[
c=2 l z\left(n_{2}-n_{1}\right)=3.024 \times 10^{8} \mathrm{~m} / \mathrm{sec}
\]

5.225 When \(v \ll c\) time dilation effect of relativity can be neglected (i.e. \(t^{\prime} \approx t\) ) and we can use time in the reference frame fixed to the observer. Suppose the source emits short pulses with intervals \(T_{0}\). Then in the reference frame fixed to the receiver the distance between two seccessive pulses is \(\lambda=c T_{0}-v_{\mathrm{r}} T_{0}\) when measured along the observation line.\\
Here \(v_{r}=v \cos \theta\) is the projection of the source velocity on the observation line. The frequency of the pulses received by the observer is\\
\(v=\frac{c}{\lambda}=\frac{v_{0}}{1-\frac{v_{\underline{\mathrm{r}}}}{\mathrm{c}}} \approx v_{0}\left(1+\frac{v_{\mathrm{r}}}{\mathrm{c}}\right)\)\\
(The formula is accurate to first order only)\\
Thus \(\frac{v-v_{0}}{v_{0}}=\frac{v_{\mathrm{r}}}{\mathrm{c}}=\frac{\mathrm{v} \cos \theta}{c}\)\\
\includegraphics[max width=\textwidth, alt={}, center]{da0156b6-4133-434e-90d4-ec01b2070afe-237_369_757_850_685}

The frequency increases when the source is moving towards the observer.\\
\(5.226 \frac{\Delta v}{v}=\frac{v}{c} \cos \theta\) from the previous problem\\
But \(v \lambda=c\) gives an differentiation\\
\includegraphics[max width=\textwidth, alt={}, center]{da0156b6-4133-434e-90d4-ec01b2070afe-237_118_472_1344_894}

\[
\frac{\Delta v}{v}=-\frac{\Delta \lambda}{\lambda}
\]

So

\[
\Delta \lambda=-\lambda \sqrt{\frac{v^{2}}{c^{2}}} \cos \theta=-\lambda \sqrt{\frac{2 T}{m c^{2}}} \cos \theta
\]

on using

\[
T=\frac{1}{2} m \mathrm{v}^{2}, m=\text { mass of } \mathrm{He}^{+} \text {ion }
\]

We use \(m c^{2}=4 \times 938 \mathrm{MeV}\). Putting other values

\[
\Delta \lambda=-26 \mathrm{~nm}
\]

5.227 One end of the solar disc is moving towards us while the other end is moving away from us. The angle \(\theta\) between the direction in which the edges of the disc are moving and the line of observation is small \((\cos \theta \approx 1)\). Thus

\[
\frac{\Delta \lambda}{\lambda}=\frac{2 \omega R}{c}
\]

where \(\omega=\frac{2 \pi}{T}\) is the angular velocity of the Sun. Thus

So

\[
\begin{aligned}
\omega & =\frac{c \Delta \lambda}{2 R \lambda} \\
T & =\frac{4 \pi R \lambda}{c \Delta \lambda}
\end{aligned}
\]

Putting the values ( \(R=6.95 \times 10^{8} \mathrm{~m}\) )\\
we get

\[
T=24.85 \text { days }
\]

5.228 Maximum splitting of the spectral lines\\
will occur when both of the stars are moving in the direction of line of observation as shown.\\
We then have the equations

\[
\begin{aligned}
\left(\frac{\Delta \lambda}{\lambda}\right)_{m} & =\frac{2 v}{i} \\
\frac{m v^{2}}{R} & =\frac{\gamma m^{2}}{4 R^{2}} \\
\tau & =\frac{\pi R}{v}
\end{aligned}
\]

\begin{center}
\includegraphics[max width=\textwidth, alt={}]{da0156b6-4133-434e-90d4-ec01b2070afe-238_322_287_861_1000}
\end{center}

From these we get

\[
\begin{gathered}
d=2 R=\left(\frac{\Delta \lambda}{\lambda}\right)_{m} c \tau / \pi=2.97 \times 10^{7} \mathrm{~km} \\
m=\left(\frac{\Delta \lambda}{\lambda}\right)_{m}^{3} c^{3} \tau / 2 \pi \gamma=2.9 \times 10^{29} \mathrm{~kg}
\end{gathered}
\]

5.229 We define the frame \(S\) (the lab frame) by the condition of the problem. In this frame the mirror is moving with velocity v (along say \(x\)-axis) towards left and light of frequency \(\omega_{0}\) is approaching it from the left. We introduce the frame \(S^{\prime}\) whose axes are parallel to those of \(S\) but which is moving with velocity v along \(x\) axis towards left (so that the mirror is at rest in \(S^{\prime}\) ). In \(S^{\prime}\) the frequency of the incident light is

\[
\omega_{1}=\omega_{0}\left(\frac{1+\mathrm{v} / c}{1-\mathrm{v} / c}\right)^{1 / 2}
\]

In \(S\).' the reflected light still has frequency \(\omega_{1}\) but it is now moving towards left. When we transform back to \(S\) this reflected light has the frequency

\[
\omega=\omega_{1}\left(\frac{1+\mathrm{v} / \mathrm{c}}{1-\mathrm{v} / \mathrm{c}}\right)^{1 / 2}=\omega_{0}\left(\frac{1+\mathrm{v} / \mathrm{c}}{1-\mathrm{v} / c}\right)
\]

In the nonrelativistic limit

\[
\omega \approx \omega_{0}\left(1+\frac{2 v}{c}\right)
\]

5.230 From the previous problem, the beat frequency is clearly

\[
\Delta v=v_{0} \frac{2 v}{c}=2 v /\left(c / v_{0}\right)=2 v / \lambda_{9}
\]

Hence

\[
v=\frac{1}{2} \lambda \Delta v=\frac{10^{+3}}{2} \times 50 \mathrm{~cm} / \mathrm{sec}=900 \mathrm{~km} / \mathrm{hour}
\]

5.231 From the invariance of phase under Lorentz transformations we get

\[
\omega t-k x=\omega^{\prime} t^{\prime}-k^{\prime} x^{\prime}
\]

Here \(\omega=c k\). The primed coordinates refer to the frame \(S^{\prime}\) which is moving to the right with velocity v :-

Then \(x^{\prime}=\gamma(x-\mathrm{v} t)\)

\[
t^{\prime}=\gamma\left(t-\frac{v x}{c^{2}}\right)
\]

where

\[
\gamma=\left(\sqrt{1-v^{2} / c^{2}}\right)^{-1}
\]

Substituting and equating the coefficients of \(t \& x\)

\[
\begin{aligned}
& \omega=\gamma \omega^{\prime}+\gamma k^{\prime} v=\omega^{\prime} \frac{1+v / c}{\sqrt{1-v^{2} / c^{2}}} \\
& k=\gamma \frac{\omega^{\prime} v}{c^{2}}+\gamma k^{\prime}=k^{\prime} \frac{1+v / c}{\sqrt{1-v^{2} / c^{2}}}
\end{aligned}
\]

\includegraphics[max width=\textwidth, alt={}, center]{da0156b6-4133-434e-90d4-ec01b2070afe-239_480_639_945_790}\\
5.232 From the previous problem using\\
we get

\[
k=\frac{2 \pi}{\lambda}
\]

\[
\lambda^{\prime}=\lambda \sqrt{\frac{1+v / c}{1-v / c}}
\]

Thus

\[
\frac{1+\frac{v}{c}}{1-v / c}=\frac{\lambda^{\prime 2}}{\lambda^{2}}
\]

or

\[
v / c=\frac{\lambda^{\prime 2}-\lambda^{2}}{\lambda^{\prime 2}+\lambda^{2}}=\frac{564^{2}-434^{2}}{564^{2}+434^{2}}=0.256
\]

5.233 As in the previous problem

\[
\begin{gathered}
\frac{v}{c}=\frac{\lambda^{2}-\lambda^{\prime 2}}{\lambda^{2}+\lambda^{\prime 2}} \\
v=c \frac{\left(\frac{\lambda}{\lambda^{\prime}}\right)^{2}-1}{\left(\frac{\lambda}{\lambda^{\prime}}\right)^{2}+1}=7.1 \times 10^{4} \mathrm{~km} / \mathrm{s}
\end{gathered}
\]

so\\
5.233 We go to the frame in which the observer is at rest. In this frame the velocity of the source of light is, by relativistic velocity addition foumula,

\[
v=\frac{\frac{3}{4} c-\frac{1}{2} c}{1-\left(\frac{3}{4} c \cdot \frac{1}{2} c / c^{2}\right)}=\frac{2 c}{5} .
\]

When this source emits light of proper frequency \(\omega_{0}\), the frequency recorded by observer will be

\[
\omega=\omega_{0} \sqrt{\frac{1-v / c}{1+v / c}+v / c}=\sqrt{\frac{3}{7}} \omega_{0}
\]

Note that \(\omega<\omega_{0}\) as the source is moving away from the observer (red shift).\\
5.235 In transverse Doppler effect.

So

\[
\begin{gathered}
\omega=\omega_{0} \sqrt{1-\beta^{2}} \cong \omega_{0}\left(1-\frac{1}{2} \beta^{2}\right) \\
\lambda=\frac{c}{\omega}=\frac{c}{\omega_{0}}\left(1+\frac{1}{2} \beta^{2}\right)=\lambda_{0}\left(1+\frac{1}{2} \beta^{2}\right) \\
\Delta \lambda=\frac{1}{2} \beta^{2} \lambda
\end{gathered}
\]

Hence\\
Using \(\beta^{2}=\frac{\mathrm{v}^{2}}{\mathrm{c}^{2}}=\frac{2 T}{m c^{2}}\), where \(T=\mathrm{K.E}\) of \(H\) atoms

\[
\Delta \lambda=\frac{T}{m c^{2}} \lambda=\frac{1}{938} \times 656.3 \mathrm{~nm}=0.70 \mathrm{~nm}
\]

, 236 (a) If light is received by the observer at \(P\) at the moment when the source is at \(O\), it must have been emitted by the source when it was at \(O^{\prime}\) and travelled along \(O^{\prime} P\). Then if \(O^{\prime} P=c t, O^{\prime} O=v t\) and \(\cos \theta=\frac{\mathrm{v}}{\mathrm{c}}=\beta\)\\
In the frame of the observer, the frequency of the light is \(\omega\) while its wave vector is\\
\includegraphics[max width=\textwidth, alt={}, center]{da0156b6-4133-434e-90d4-ec01b2070afe-240_360_484_1589_979}

\[
\frac{\omega}{c}(\cos \theta, \sin \theta, 0)
\]

we can calculate the value of \(\omega\) by relating it to proper frequency \(\omega\). The relation is

\[
\omega_{0}=\frac{\omega}{\sqrt{1-\beta^{2}}}(1-\beta \cos \theta)
\]

(To derive the formula is this form it is easiest to note that \(\frac{\omega}{\sqrt{1-v / c^{2}}}-\frac{\vec{k} \cdot \vec{v}}{\sqrt{1-v^{2} / c^{2}}}\) is an invariant which takes the value \(\omega_{0}\) in the rest frame of the source).

Thus

\[
\omega=\frac{\omega_{0} \sqrt{1-\beta^{2}}}{1-\beta^{2}}=\frac{\omega_{0}}{\sqrt{1-\beta^{2}}}=5 \times 10^{10} \mathrm{sec}^{-1}
\]

(b) For the light to be received at the instant observer sees the source at \(O\), light must be emitted when the observer is at \(O\) at \(90^{\circ}=\theta\)

\[
\cos \theta=0
\]

Then as before \(\omega_{0}=\frac{\omega}{\sqrt{1-\beta^{2}}}\) or \(\omega=\omega_{0} \sqrt{1-\beta^{2}}=1.8 \times 10^{10} \mathrm{sec}^{-1}\)\\
In this case the observer will receive light along \(O P\) and he will "see" that the source is at \(O\) even though the source will have moved ahead at the instant the light is received.\\
5.237 An electron moving in front of a metal mirror sees an image charge of equal and opposite type. The two together constitute a dipole. Let us look at the problem in the rest frame of the electron. In this frame the grating period is Lorenz contracted to

\[
d^{\prime}=d \sqrt{1-v^{2} / c^{2}}
\]

Because the metal has etchings the dipole moment of electron-image pair is periodically disturbed with a period \(\frac{d^{\prime}}{v}\)

The corresponding frequency is \(\frac{\mathbf{v}}{\mathbf{d}^{\prime}}\) which is also the proper frequency of radiation emitted.\\
Due to Doppler effect the frequency observed at an angle \(\theta\) is

\[
v=v^{\prime} \frac{\sqrt{1-(v / c)^{2}}}{1-\frac{v}{c} \cos \theta}=\frac{v / d}{1-\frac{v}{c} \cos \theta}
\]

The corresponding wave length is \(\lambda=\frac{c}{v}=d\left(\frac{c}{v}-\cos \theta\right)\)\\
Putting

\[
\begin{gathered}
c \approx \mathrm{v}, \theta=45^{\circ}, d=2 \mu \mathrm{~m} \text { we get } \\
\lambda=0.586 \mu \mathrm{~m}
\end{gathered}
\]

5.238 (a) Let \(v_{x}\) be the projection of the velocity vector of the radiating atom on the observation direction. The number of atoms with projections falling within the interval \(v_{x}\) and \(\mathrm{v}_{\mathrm{x}}+d \mathrm{v}_{\mathrm{x}}\) is

\[
n\left(v_{\mathrm{x}}\right) d v_{\mathrm{x}} \sim \exp \left(-m v_{\mathrm{x}}^{2} / 2 k T\right) d v_{\mathrm{x}}
\]

The frequency of light emitted by the atoms moving with velocity \(v_{x}\) is \(\omega=\omega_{0}\left(1+\frac{v_{x}}{c}\right)\). From the expressions the frequency distribution of atoms can be found \(: n(\omega) d \omega=n\left(\mathrm{v}_{\mathrm{x}}\right) d \mathrm{v}_{\mathrm{x}}\). Now using\\
we get

Now the spectral radiation density \(I_{\omega} \propto n_{\omega}\)

Hence

\[
\begin{gathered}
\mathrm{v}_{\mathrm{x}}=c \frac{\omega-\omega_{0}}{\omega_{0}} \\
n(\omega) d \omega_{\sim} \exp \left(-\frac{m c^{2}}{2 k T}\left(1-\frac{\omega}{\omega_{0}}\right)^{2}\right) \frac{c d \omega}{\omega_{0}}
\end{gathered}
\]

(The constant of proportionality is fixed by \(I_{0}\).)\\
(b) On putting \(\omega=\omega_{0} \pm \frac{1}{2} \Delta \omega\) and demanding

\[
I_{\omega}=I_{0} e^{-a\left(1-\frac{\omega}{\omega_{0}}\right)^{2}}, a=\frac{m c^{2}}{2 k T}
\]

The constant of proportionality is fixed by \(I_{0}\).)

\[
\begin{array}{lc} 
& I_{\omega}=I_{0} / 2 \\
\text { we get } & \frac{1}{2}=e^{-a\left(\frac{\Delta \omega}{2 \omega_{0}}\right)^{2}} \\
\text { so } & a\left(\frac{\Delta \omega}{2 \omega_{0}}\right)^{2}=\ln 2 \\
\text { Hence } & \frac{\Delta \omega}{2 \omega_{0}}=\sqrt{(2 \ln 2) k T / m c^{2}} \\
\text { and } & \frac{\Delta \omega}{\omega}=2 \sqrt{(2 \ln 2) k T / m c^{2}}
\end{array}
\]

5.239 In vacuum inertial frames are all equivalent; the velocity of light is \(c\) in any frame. This equivalence of inertial frames does not hold in material media and here the frame in which the medium is at rest is singled out. It is in this frame that the velocity of light is \(\frac{c}{n}\) where \(n\) is the refractive in desc of light for that medium.\\
The velocity of light in the frame in which the medium is moving is then by the law of addition of velocities

\[
\begin{gathered}
\frac{\frac{c}{n}+\mathrm{v}}{1+\frac{c}{n} \cdot \mathrm{v} / c^{2}}=\frac{\frac{c}{n}+\mathrm{v}}{1+\frac{\mathrm{v}}{\mathrm{cn}}}=\left(\frac{c}{n}+\mathrm{v}\right)\left(1-\frac{\mathrm{v}}{c n}+\ldots\right) \\
=\frac{c}{n}+\mathrm{v}-\frac{\mathrm{v}}{\mathrm{n}^{2}}+\ldots \approx \frac{c}{n}+\mathrm{v}\left(1-\frac{1}{n^{2}}\right)
\end{gathered}
\]

This is the velocity of light in the medium in a frame in which the medium is moving with velocity \(\mathrm{v} \ll c\).\\
5.240 Although speed of light is the same in all inertial frames of reference according to the principles of relativity, the direction of a light ray can appear different in different frames. This phenomenon is called aberration and to first order in \(\frac{v}{c}\), can be calculated by the elementary law of addition of velocities applied to light waves. The angle of aberration is

\[
\tan ^{-1} \frac{v}{c}
\]

and in the present case it equals \(\frac{1}{2} \delta \theta\) on either side. Thus equating \(\frac{v}{c}=\tan \frac{1}{2} \delta \theta \approx \frac{1}{2} \delta \theta(\delta \theta\) radians \()\)\\
or

\[
\begin{aligned}
v=\frac{c}{2} \delta \theta & =\frac{3 \times 10^{8}}{2} \times \frac{41}{3600} \times \frac{\pi}{180} \\
& =\frac{3 \times 4.1 \times \pi}{3.6 \times 3.6} \times 10^{4} \mathrm{~m} / \mathrm{s}=29.8 \mathrm{k} \mathrm{~m} / \mathrm{sec}
\end{aligned}
\]

\includegraphics[max width=\textwidth, alt={}, center]{da0156b6-4133-434e-90d4-ec01b2070afe-243_425_249_639_1163}\\
5.241 We consider the invariance of the phase of a wave moving in the \(x-y\) plane.

We write

\[
\omega^{\prime} t^{\prime}-k_{\mathrm{x}}^{\prime} \chi^{\prime}-k_{\mathrm{y}}^{\prime} y^{\prime}=\omega t-k_{\mathrm{x}} \chi-k_{\mathrm{y}} y
\]

From Lorcntz transformations, L.H.S.

\[
=\omega^{\prime} \gamma\left(t-\frac{\mathrm{v} x}{c^{2}}\right)-k_{\mathrm{x}}^{\prime}(x-\mathrm{v} t) \gamma-k_{\mathrm{y}}^{\prime} y
\]

so equating

\[
\begin{aligned}
& \omega=\gamma\left(\omega^{\prime}+\mathrm{v} k_{\mathrm{x}}^{\prime}\right) \\
& k_{\mathrm{x}}=\gamma\left(k_{\mathrm{x}}^{\prime}+\frac{\mathrm{v} \omega^{\prime}}{c^{2}}\right) \text { and } k_{\mathrm{y}}=k_{\mathrm{y}}^{\prime}
\end{aligned}
\]

so inverting

\[
\begin{aligned}
& \omega^{\prime}=\gamma\left(\omega-\mathrm{v} k_{\mathrm{x}}\right)^{\prime} \\
& k_{\mathrm{x}}^{\prime}=\gamma\left(k_{\mathrm{x}}-\frac{\mathrm{v} \omega}{c^{2}}\right) \\
& k_{\mathrm{y}}^{\prime}=k_{\mathrm{y}}
\end{aligned}
\]

writing

\[
\begin{aligned}
& k_{\mathrm{x}}^{\prime}=k^{\prime} \cos \theta^{\prime}, k_{\mathrm{x}}=k \cos \theta \\
& k_{\mathrm{y}}^{\prime}=k^{\prime} \sin \theta^{\prime}, k_{\mathrm{y}}=k \sin \theta
\end{aligned}
\]

we get on using \(c k^{\prime}=\omega^{\prime}, c k=\omega\)

\[
\cos \theta^{\prime}=\frac{\cos \theta-\beta}{1-\beta \cos \theta}
\]

where \(\beta=\frac{v}{c}\) and the primed frame is moving with velocity \(v\) in the \(x\)-direction w.r.t the unprimed frame. For small \(\beta \ll 1\), the situation is as shown.

We see that if

Then\\
\(\cos \boldsymbol{\theta}^{\prime} \approx-\boldsymbol{\beta}\)

\[
\theta=-\pi / 2 .
\]

\[
\theta^{\prime}=-\left(\frac{\pi}{2}+\sin ^{-1} \beta\right)
\]

\begin{center}
\includegraphics[max width=\textwidth, alt={}]{da0156b6-4133-434e-90d4-ec01b2070afe-244_429_520_97_919}
\end{center}

This is exactly what we get from elementary nonrelativistic law of addition of velocities,\\
5.242 The statement of the problem is not quite properly worked and is in fact misleading. The correct situation is described below. We consider, for simplicity, stars in the \(x-z\) plane. Then the previous formula is applicable,\\
and we have

\[
\cos \theta^{\prime}=\frac{\cos \theta-\beta}{1-\beta \cos \theta}=\frac{\cos \theta-0.99}{1-0.99 \cos \theta}
\]

The distribution of \(\theta^{\prime}\) is given in the diagram below\\
The light that appears to come from the forward quadrant in the frame \(K(\theta=-\pi\) to \(\theta=-\pi / 2)\) is compressed into an angle of magnitude \(+8 \cdot 1^{\circ}\) in the forward direction while the remaining stars are\\
\includegraphics[max width=\textwidth, alt={}, center]{da0156b6-4133-434e-90d4-ec01b2070afe-244_441_691_785_755}\\
spread out.

The three dimensional distribution can also be found out from the three dimensional generalization of the formula in the previous problems.

243 The field induced by a charged particle moving with velocity \(V\) excites the atoms of the medium turning them into sources of light waves. Let us consider two arbitrary points \(A\) and \(B\) along the path of the particle. The light waves emitted from these points when the particle passes them reach the point \(P\) simultaneously and reinforce each other provided they are in phase which is the case is general if the time taken by the light wave to propagate from the point \(A\) to the point \(C\) is equal to that taken by\\
\includegraphics[max width=\textwidth, alt={}, center]{da0156b6-4133-434e-90d4-ec01b2070afe-244_347_455_1459_970}\\
the particle to fly over the distance \(A B\). Hence we obtain

\[
\cos \theta=\frac{v}{V}
\]

where \(\mathrm{v}=\frac{c}{n}\) is the phase velocity of light. It is evident that the radiation is possible only if \(V>\mathrm{v}\) i.e. when the velocity of the particle exceeds the phase velocity of light in the medium.\\
5.244 We must have

\[
V \geq \frac{c}{n}=\frac{3}{1.6} \times 10^{8} \mathrm{~m} / \mathrm{s} \quad \text { or } \quad \frac{V}{c} \geq \frac{1}{1.6}
\]

For electrons this means a K.E. greater than

\[
\begin{gathered}
T_{e} \geq \frac{m_{e} c^{2}}{\sqrt{1-\left(\frac{1}{1.6}\right)^{2}}}-m_{e} c^{2}=m_{e} c^{2}\left[\frac{n}{\sqrt{n^{2}-1}}-1\right] \\
=0.511\left[\frac{1}{\sqrt{1\left(\frac{1}{1.6}\right)^{2}}}-1\right] \quad \text { using } m_{e} c^{2}=0.511 \mathrm{MeV}=0.144 \mathrm{MeV}
\end{gathered}
\]

For protons with \(m_{p} c^{2}=938 \mathrm{MeV}\)

Also

\[
\begin{gathered}
T_{p} \geq 938\left[\frac{1}{\sqrt{1-\left(\frac{1}{1.6}\right)^{2}}}-1\right]=264 \mathrm{MeV}=0.264 \mathrm{GeV} \\
T_{\min }=29.6 \mathrm{MeV}=m c^{2}\left[\frac{1}{\sqrt{1-\left(\frac{1}{1.6}\right)^{2}}}-1\right]
\end{gathered}
\]

Then \(m c^{2}=105 \cdot 3 \mathrm{MeV}\). This is very nearly the mass of means.\\
5.245

From \(\cos \theta=\frac{\mathrm{v}}{\mathrm{V}}\)\\
we get

\[
V=v \sec \theta
\]

so

\[
\frac{V}{c}=\frac{v}{c} \sec \theta=\frac{\sec \theta}{n}=\frac{\sec 30^{\circ}}{1 \cdot 5}=\frac{2 / \sqrt{3}}{3 / 2}=\frac{4}{3 \sqrt{3}}
\]

Thus for electorns

\[
T_{\mathrm{e}}=0.511\left[\frac{1}{\sqrt{1-\frac{16}{27}}}-1\right]=0.511\left[\sqrt{\frac{27}{11}}-1\right]=0.289 \mathrm{MeV}
\]

Generally

\[
T=m c^{2}\left[\frac{1}{\sqrt{1-\frac{1}{n^{2} \cos ^{2} \theta}}}-1\right]
\]

\subsection*{5.7 THERMAL RADIATION. QUANTUM NATURE OF LIGHT}
5.246 (a) The most probable radiation frequency \(\omega_{p r}\) is the frequency for which

\[
\frac{d}{d \omega} u_{\omega}=3 \omega^{2} F(\omega / T)+\frac{\omega^{3}}{T} F^{\prime}(\omega / T)=0
\]

The maximum frequency is the root other than \(\omega=0\) of this equation. It is

\[
\omega=-\frac{3 T F(\omega / T)}{F^{\prime}(\omega / T)}
\]

or \(\omega_{\mathrm{pr}}=x_{0} T\) where \(x_{0}\) is the solution of the transcendental equation

\[
3 F\left(x_{0}\right)+x_{0} F^{\prime}\left(x_{0}\right)=0
\]

(b) The maximum spectral density is the density corresponding to most probable frequency. It is

\[
\left(u_{\omega}\right)_{\max }=x_{0}^{3} F\left(x_{0}\right) T^{3} \alpha T^{3}
\]

where \(x_{0}\) is defined above.\\
(c) The radiosity is

\[
M_{e}=\frac{c}{4} \int_{0}^{\infty} \omega^{3} F\left(\frac{\omega}{T}\right) d \omega=T^{4}\left[\frac{c}{4} \int_{0}^{\infty} x^{3} F(x) d x\right] \alpha T^{4}
\]

5.247 For the first black body

Then

\[
\left(\lambda_{m}\right)_{1}=\frac{b}{T_{1}}
\]

Hence

\[
\left(\lambda_{m}\right)_{2}=\frac{b}{T_{1}}+\Delta \lambda=\frac{b}{T_{2}}
\]

\[
T_{2}=\frac{b}{\frac{b}{T_{1}}+\Delta \lambda}=\frac{b T_{1}}{b+T_{1} \Delta \lambda}=1.747 \mathrm{kK}
\]

5.248 From the radiosity we get the temperature of the black body. It is

\[
T=\left(\frac{M_{e}}{\sigma}\right)^{1 / 4}=\left(\frac{3.0 \times 10^{4}}{5.67 \times 10^{-8}}\right)^{1 / 4}=852.9 K
\]

Hence the wavelength corresponding to the maximum emissive capacity of the body is

\[
\begin{aligned}
& \frac{b}{T}=\frac{0.29}{852.9} \mathrm{~cm}=3.4 \times 10^{-4} \mathrm{~cm}=3.4 \mu \mathrm{~m} \\
& \text { (Note that } 3.0 \mathrm{~W} / \mathrm{cm}^{2}=3.0 \times 10^{4} \mathrm{~W} / \mathrm{m}^{2} . \text { ) }
\end{aligned}
\]

5.249 The black body temperature of the sun may be taken as

\[
T_{\odot}=\frac{0.29}{0.48 \times 10^{-4}}=6042 \mathrm{~K}
\]

Thus the radiosity is

\[
M_{e \odot}=5.67 \times 10^{-8}(6042)^{4}=0.7555 \times 10^{8} \mathrm{~W} / \mathrm{m}^{2}
\]

Energy lost by sun is

\[
4 \pi(6.95)^{2} \times 10^{16} \times 0.7555 \times 10^{8}=4.5855 \times 10^{26} \text { watt }
\]

This corresponds to a mass loss of

\[
\frac{4.5855 \times 10^{26}}{9 \times 10^{16}} \mathrm{~kg} / \mathrm{sec}=5.1 \times 10^{9} \mathrm{~kg} / \mathrm{sec}
\]

The sun loses \(1 \%\) of its mass in

\[
\frac{1.97 \times 10^{30} \times 10^{-2}}{5.1 \times 10^{9}} \sec \approx 1.22 \times 10^{11} \text { years }
\]

5.250 For an ideal gas \(p=n k T\) where \(n=\) number density of the particles and \(k=\frac{R}{N_{A}}\) is Boltzman constant. In a fully ionized hydrogen plasma, both \(H\) ions (protons) and electrons contribute to pressure but since the mass of electrons is quite small ( \(\approx m_{p} / 1836\) ), only protons contribute to mass density. Thus\\
and

\[
\begin{gathered}
n=\frac{2 \rho}{m_{H}} \\
p=\frac{2 \rho R}{N_{A} m_{H}} T
\end{gathered}
\]

where \(m_{H} \approx m_{p}\) is the proton or hydrogen mass.\\
Equating this to thermal radiation pressure

Then

\[
\begin{gathered}
\frac{2 \rho R}{N_{A} m_{H}} T=\frac{u}{3}=\frac{M_{e}}{3} \times \frac{4}{c}=\frac{4 \sigma T^{4}}{3 c} \\
T^{3}=\frac{3 c \rho R}{2 \sigma N_{A} m_{H}}=\frac{3 c \rho R}{\sigma M}
\end{gathered}
\]

where \(=2 N_{A} m_{H}=\) molecular weight of hydrogen \(=2 \times 10^{-3} \mathrm{~kg}\).

Thus

\[
T=\left(\frac{3 c \rho R}{\sigma M}\right)^{1 / 3} \approx 1.89 \times 10^{7} \mathrm{~K}
\]

5.251 In time \(d t\) after the instant \(t\) when the temperature of the ball is \(T\), it loses

\[
\pi d^{2} \sigma T^{4} d t
\]

Joules of energy. As a result its temperature falls by \(-\boldsymbol{d T}\) and

\[
\pi d^{2} \sigma T^{4} d t=-\frac{\pi}{6} d^{3} \rho C d T
\]

where \(\rho=\) density of copper, \(C=\) its sp.heat

Thus

\[
d t=-\frac{C \rho d}{6 \sigma} \frac{d T}{T^{4}}
\]

or

\[
t_{0}=\frac{C \rho d}{6 \sigma} \int_{T_{0}}^{T_{0} / \eta}-\frac{d T}{T^{4}}=\frac{C \rho d}{18 \sigma T_{0}^{3}}\left(\eta^{3}-1\right)=2.94 \text { hours. }
\]

5.252 Taking account of cosine low of emission we write for the energy radiated per second by the hole in cavity \# 1 as

\[
d I(\Omega)=A \cos \theta d \Omega
\]

where \(A\) is an constant, \(d \Omega\) is an element of solid angle around some direciton defined by the symbol \(\Omega\). Integrating over the whole forward hemisphere we get

\[
I=A \int_{0}^{\pi / 2} \cos \theta 2 \pi \sin \theta d \theta=\pi A
\]

\begin{center}
\includegraphics[max width=\textwidth, alt={}]{da0156b6-4133-434e-90d4-ec01b2070afe-248_292_783_702_385}
\end{center}

We find \(A\) by equating this to the quantity \(\sigma T_{1}^{4} \cdot \frac{\pi d^{2}}{4} \sigma\) is stefan-Boltzman constant and \(d\) is the diameter of th hole.

Then

\[
A=\frac{1}{4} \sigma d^{2} T_{1}^{4}
\]

Now energy reaching 2 from 1 is \((\cos \theta \sim 1)\)

\[
\frac{1}{4} \sigma d^{2} T_{1}^{4} \cdot \Delta \Omega
\]

where \(\Delta \Omega=\frac{\left(\pi d^{2} / 4\right)}{l^{2}}\) is the solid angle subtended by the hole of 2 at 1 . \{We are assuming \(d \ll l\) so \(\Delta \Omega=\) area of hole \(\left./(\text { distance })^{2}\right\}\).\\
This must equal

\[
\boldsymbol{\sigma} T_{2}^{4} \pi d^{2} / 4
\]

which is the energy emitted by 2 . Thus equating\\
or

\[
\begin{gathered}
\frac{1}{4} \sigma d^{2} T_{1}^{4} \frac{\pi d^{2}}{4 l^{2}}=\sigma T_{2}^{4} \frac{\pi d^{2}}{4} \\
T_{2}=T_{1} \sqrt{\frac{d}{2 l}} \\
T_{2}=0.380 k K=380 \mathrm{~K} .
\end{gathered}
\]

Substituting we get\\
5.253 (a) The total internal energy of the cavity is

Hence

\[
\begin{aligned}
& U=\frac{4 \sigma}{c} T^{4} V \\
& C_{V}=\left(\frac{\partial U}{\partial T}\right)_{V}=\frac{16 \sigma}{C} T^{3} V \\
&=\frac{16 \times 5.67 \times 10^{8}}{3 \times 10^{8}} \times 10^{9} \times 10^{-3} \mathrm{Joule} /{ }^{\circ} \mathrm{K} \\
&=\frac{1.6 \times 5.67}{3} n \mathrm{~J} / \mathrm{K}=3.024 n \mathrm{~J} / \mathrm{K}
\end{aligned}
\]

(b) From first law

\[
\begin{array}{lrl}
T d S & =d U+p d V \\
& =V d U+U d V+\frac{u}{3} d V\left(p=\frac{U}{3}\right) \\
& =V d U+\frac{4 U}{3} d V \\
& =\frac{16 \sigma}{C} V T^{3} d T+\frac{16 \sigma}{3 C} T^{4} d V \\
& d S=\frac{16 \sigma}{C} V T^{2} d T+\frac{16 \sigma}{3 C} T^{3} d V \\
& =d\left(\frac{16 \sigma}{3 C} V T^{3}\right) \\
\text { so } & & \\
\text { Hence } & S=\frac{16 \sigma}{3 C} V T^{3}=\frac{1}{3} C_{V}=1.008 n \mathrm{~J} / \mathrm{K} .
\end{array}
\]

5.254 We are given

\[
u(\omega, T)=A \omega^{3} \exp (-a \omega / T)
\]

(a) Then

\[
\frac{d u}{d \omega}=\left(\frac{3}{\omega}-\frac{a}{T}\right) u=0
\]

so

\[
\omega_{p r}=\frac{3 T}{a}=\frac{6000}{7.64} \times 10^{12} \mathrm{~s}^{-1}
\]

(b) We determine the spectral distribution in wavelength.

But

\[
-\tilde{u}(\lambda, T) d \lambda=u(\omega, T) d \omega
\]

o

\[
\cdots=\frac{2 \pi c}{\lambda} \quad \text { or } \quad \lambda=\frac{2 \pi c}{\omega}=\frac{C^{\prime}}{\omega}
\]

so

\[
d \lambda=-\frac{C^{\prime}}{\omega^{2}} d \omega, d \omega=-\frac{C^{\prime}}{\lambda^{2}} d \lambda
\]

(we have put a minus sign before \(d \lambda\) to subsume just this fact \(d \lambda\) is -ve where \(d \omega\) is +ve.)

\[
\tilde{u}(\lambda, T)=\frac{C^{\prime}}{\lambda^{2}} u\left(\frac{C^{\prime}}{\lambda}, T\right)=\frac{C^{\prime 4} A}{\lambda^{5}} \exp \left(-\frac{a C^{\prime}}{\lambda T}\right)
\]

This is maximum when\\
or

\[
\begin{gathered}
\frac{\partial \tilde{u}}{\partial \lambda}=6=\tilde{u}\left[\frac{-5}{\lambda}+\frac{a C^{\prime}}{\lambda^{2} T}\right] \\
\lambda_{p r}=\frac{a C^{\prime}}{5 T}=\frac{2 \pi c a}{5 T}=1.44 \mu \mathrm{~m}
\end{gathered}
\]

5.255 From Planek's formula

\[
u_{\omega}=\frac{\hbar \omega^{3}}{\pi^{2} c^{3}} \frac{1}{e^{\hbar \omega / k T}-1}
\]

(a) In a range \(\hbar \omega \ll k T\) (long wavelength or high temperature).

\[
\begin{gathered}
u_{\omega} \rightarrow \frac{\hbar \omega^{3}}{\pi^{2} c^{3}} \frac{1}{\hbar \omega} \\
=\frac{\omega^{2}}{k T} \\
\pi^{2} c^{3} k T, \text { using } e^{x}=1+x \quad \text { for small } x .
\end{gathered}
\]

(b) In the range \(\hbar \omega \gg T\) (high frequency or low temperature) :\\
and

\[
\begin{aligned}
& \frac{\hbar \omega}{k T} \gg 1 \text { so } e^{\frac{\hbar \omega}{k T}} \gg 1 \\
& u_{\omega}=\frac{\hbar \omega^{3}}{\pi^{2} c^{3}} e^{-\hbar \omega / k T}
\end{aligned}
\]

5.256 We write

\[
u_{\omega} d \omega=\tilde{u}_{v} d v \text { where } \omega=2 \pi v
\]

Then

\[
\tilde{u}_{v}=\frac{2 \pi \hbar(2 \pi v)^{3}}{\pi^{2} c^{3}} \frac{1}{e^{2 \pi h v / k T}-1}=\frac{16 \pi^{2} \hbar v^{3}}{c^{3}} \frac{1}{e^{2 \pi h v / k T}-1}
\]

Also

\[
\begin{gathered}
-\tilde{u}(\lambda, T) d \lambda=u_{\omega} d \omega \quad \text { where } \quad \lambda=\frac{2 \pi c}{\omega} \\
d \omega=-\frac{2 \pi c}{\lambda^{2}} d \lambda \\
\tilde{u}(\lambda, T)=\frac{2 \pi c}{\lambda^{2}} u\left(\frac{2 \pi c}{\lambda}, T\right) \\
\left(\frac{2 \pi c}{\lambda}\right)^{3} \frac{\hbar}{\pi^{2} c^{3}} \frac{1}{e^{2 \pi c h c / \lambda k T}-1}=\frac{16 \pi^{2} c h}{\lambda^{5}} \frac{1}{e^{2 \pi c h c / \lambda k T}-1}
\end{gathered}
\]

5.257 We write the required power in terms of the radiosity by considering ónly the energy radiated in the given range. Then from the previous problem

\[
\begin{aligned}
\Delta P & =\frac{c}{4} \tilde{u}\left(\lambda_{m}, T\right) \Delta \lambda \\
& =\frac{4 \pi^{2} c^{2} \hbar}{\lambda_{m}^{5}} \frac{\Delta \lambda}{e^{2 \pi c h / k \lambda_{m} T}-1}
\end{aligned}
\]

But

\[
\lambda_{m} T=b
\]

so

\[
\Delta P=\frac{4 \pi^{2} c^{2} \hbar}{\lambda_{m}^{5}} \frac{\Delta \lambda}{e^{2 \pi c h / k b}-1} \Delta \lambda
\]

Using the data\\
and

\[
\begin{gathered}
\frac{2 \pi c h}{k b}=\frac{2 \pi \times 3 \times 10^{8} \times 1.05 \times 0^{-34}}{1.38 \times 10^{-23} \times 2.9 \times 10^{-3}}=4.9643 \\
\frac{1}{e^{2 \pi c h / k b}-1}=7.03 \times 10^{-3} \\
\Delta P=0.312 \mathrm{~W} / \mathrm{cm}^{2}
\end{gathered}
\]

5.258 (a) From the curve of the function \(y(x)\) we see that \(y=0.5\) when \(x=1.41\)

Thus

\[
\lambda=1.41 \times \frac{0.29}{3700} \mathrm{~cm}=1.105 \mu \mathrm{~m}
\]

(b) At 5000 K

\[
\lambda=\frac{0.29}{.5} \times 10^{-6} \mathrm{~m}=0.58 \mu \mathrm{~m}
\]

So the visible range \((0.40\) to 0.70\() \mu \mathrm{m}\) corresponds to a range \((0.69\) to 1.31\()\) of \(x\).\\
From the curve

\[
\begin{aligned}
& y(0.69)=0.07 \\
& y(1.31)=0.44
\end{aligned}
\]

so the fraction is 0.37\\
(c) The value of \(x\) corresponding to 0.76 are

\[
\begin{aligned}
& x_{1}=0.76 / \frac{0.29}{0.3}=0.786 \text { at } 3000 \mathrm{~K} \\
& x_{2}=0.76 / \frac{0.29}{0.5}=1.31 \text { at } 5000 \mathrm{~K}
\end{aligned}
\]

The requisite fraction is then

\[
\begin{aligned}
\left(\frac{P_{2}}{P_{1}}\right) & =\underset{\substack{\text { ratio of } \\
\text { total power } \\
T_{1} \\
\text { tration } \\
\text { required wavelengths } \\
\text { in the radiated power } \\
\text { fraction of } \\
\frac{T_{2}}{1-y_{1}}}}{\frac{1-y_{2}}{1-0.12}}=4.91
\end{aligned}
\]

5.259 We use the formula \(\varepsilon=\hbar \omega\)

Then the number of photons in the spectral interval ( \(\omega, \omega+d \omega\) ) is

\[
n(\omega) d \omega=\frac{u(\omega, T) d \omega}{\hbar \omega}=\frac{\omega^{2}}{\pi^{2} c^{3}} \frac{1}{e^{\hbar \omega / k T}-1} d \omega
\]

using

\[
\begin{aligned}
n(\omega) & d \omega=-\tilde{n}(\lambda) d \lambda \\
d \lambda \tilde{n}(\lambda) & =n\left(\frac{2 \pi c}{\lambda}\right) \frac{2 \pi c}{\lambda^{2}} d \lambda \\
& =\frac{(2 \pi c)^{3}}{\pi^{2} c^{3} \lambda^{4}} \frac{1}{e^{2 \pi h c / k \lambda T}-1} d \lambda \\
& =\frac{8 \pi}{\lambda^{4}} \frac{d \lambda}{e^{2 \pi h c / \lambda k T}-1}
\end{aligned}
\]

we get\\
5.260 (a) The mean density of the flow of photons at a distance \(r\) is

\[
\begin{aligned}
\langle j\rangle & =\frac{P}{4 \pi r^{2}} / \frac{2 \pi \hbar c}{\lambda}=\frac{P \lambda}{8 \pi^{2} \hbar c r^{2}} \mathrm{~m}^{-2} \mathrm{~s}^{-2} \\
& =\frac{10 \times .589 \times 10^{-6}}{8 \pi^{2} \times 1.054 \times 10^{-34} \times 10^{8} \times 4} \mathrm{~m}^{-2} \mathrm{~s}^{-1} \\
& =\frac{10 \times .589}{8 \pi^{2} \times 1.054 \times 12} \times 10^{16} \mathrm{~cm}^{-2} \mathrm{~s}^{-1} \\
& =5.9 \times 10^{13} \mathrm{~cm}^{-2} \mathrm{~s}^{-1}
\end{aligned}
\]

(b) If \(n(r)\) is the mean concentration (number per unit volume) of photons at a distance \(r\) form the source, then, since all photons are moving outwards with a velocity \(c\), there is an outward flux of \(c n\) which is balanced by the flux from the source. In steady state, the two are equal and so\\
so

\[
\begin{aligned}
n(r) & =\frac{\langle j\rangle}{c}=\frac{P \lambda}{8 \pi^{2} \hbar c^{2} r^{2}}=n \\
r & =\frac{1}{2 \pi c} \sqrt{\frac{P \lambda}{2 \hbar n}} \\
& =\frac{1}{6 \pi \times 10^{8}} \sqrt{\frac{10 \times .589 \times 10^{-6}}{2 \times 1.054 \times 10^{-34} \times 10^{2} \times 10^{6}}} \\
& =\frac{10^{2}}{6 \pi} \sqrt{\frac{5.89}{2.108}}=8.87 \text { metre }
\end{aligned}
\]

5.261 The statement made in the question is not always correct. However it is correct in certain cases, for example, when light is incident on a perfect reflector or perfect absorber.

Consider the former case. If light is incident at an angle \(\theta\) and reflected at the angle \(\theta\), then momentum transfered by each photon is

\[
2 \frac{h v}{c} \cos \theta
\]

If there are \(n(v) d v\) photons in frequency interval ( \(v, v+d v\) ), then total momentum transfered is

\[
\begin{gathered}
\int_{0}^{\infty} 2 n(v) \frac{h v}{c} \cos \theta d v \\
=\frac{2 \Phi_{e}}{c} \cos \theta
\end{gathered}
\]

\includegraphics[max width=\textwidth, alt={}, center]{da0156b6-4133-434e-90d4-ec01b2070afe-253_325_598_887_832}\\
5.262 The mean pressure \(<p>\) is related to the force Fexerted by the beam by

\[
<p>\times \frac{\pi d^{2}}{4}=F
\]

The force \(F\) equals momentum transfered per second. This is (assuming that photons, not reflected, are absorbed)

\[
2 \rho \frac{E}{c \tau}+(1-\rho) \frac{E}{c \tau}=(1+\rho) \frac{E}{c \tau}
\]

The first term is the momentum transfered on reflection (see problem (261)); the second on absorption.

\[
<p>=\frac{4(1+\rho) E}{\pi d^{2} c \tau}
\]

Substituting the values we get

\[
\langle p\rangle=48.3 \text { atmosphére. }
\]

5.263 The momentum transfered to the plate is\\
\(=\frac{E}{c}(1-\rho)\{\sin \theta \hat{i}-\cos \theta \hat{j}\}\)

\[
\begin{array}{cc}
\underset{\substack{\text { (momentum transfered } \\
\text { on absorption) }}}{\downarrow} & +\frac{E}{c} \rho\{-2 \cos \theta \hat{j}\} \\
& \downarrow \\
& \begin{array}{c}
\downarrow \\
\text { momentum transtered } \\
\text { on reflection ) }
\end{array} \\
= & \frac{E}{c}(1-\rho) \sin \theta \hat{j}-\frac{E}{c}(1+P) \cos \theta \hat{j} \quad(\hat{i})
\end{array}
\]

\begin{center}
\includegraphics[max width=\textwidth, alt={}]{da0156b6-4133-434e-90d4-ec01b2070afe-254_426_466_99_953}
\end{center}

Its magnitude is

\[
\frac{E}{c} \sqrt{(1-\rho)^{2} \sin ^{2} \theta+(1+\rho)^{2} \cos ^{2} \theta}=\frac{E}{c} \sqrt{1+\rho^{2}+2 \rho \cos 2 \theta}
\]

Substitution gives \(35 n\) N.s as the answer.\\
5.264 Suppose the mirror has a surtace area \(A\).

The incident bean then has a cross section of \(A \cos \theta\) and the incident energy is \(I A \cos \theta\) : Then the momentum transfered per second ( \(=\) Force ) is from the last problem \(-\frac{I A \cos \theta}{c}(1+\rho) \cos \theta \hat{j}+\frac{I A \cos \theta}{c}(1-\rho) \sin \theta \hat{i}\)

The normal pressure is then \(p=\frac{I}{c}(1+\rho) \cos ^{2} \theta\)\\
( \(\hat{j}\) is the unit vector \(\perp^{r}\) to the plane mirror.)\\
Putting in the values

\[
p=\frac{0.20 \times 10^{4}}{3 \times 10^{8}} \times 1.8 \times \frac{1}{2}=0.6 n \mathrm{~N} \mathrm{~cm}^{-2}
\]

\includegraphics[max width=\textwidth, alt={}, center]{da0156b6-4133-434e-90d4-ec01b2070afe-254_609_357_756_981}\\
5.265 We consider a strip defined by the angular range \((\theta, \theta+d \theta)\). From the previous problem the normal pressure exerted on this strip is

\[
\frac{2 I}{c} \cos ^{2} \theta
\]

This pressure gives rise to a force whose resultant, by symmetry is in the direction of the incident light. Thus\\
\includegraphics[max width=\textwidth, alt={}, center]{da0156b6-4133-434e-90d4-ec01b2070afe-254_390_505_1420_842}\\
\(F=\frac{2 I}{c} \int_{0}^{\pi / 2} \cos ^{2} \theta \cdot \cos \theta \cdot 2 \pi R^{2} \sin \theta d \theta=\pi R^{2} \frac{I}{c}\)

Putting in the values

\[
F=\pi \times 25 \times 10^{-4} \frac{0.70 \times 10^{4}}{3 \times 10^{8}} \mathrm{~N}=0.183 \mu \mathrm{~N}
\]

5.266 Consider a ring of radius \(x\) on the plate. The normal pressure on this ring is, by problem (264),

\[
\begin{aligned}
& \frac{2}{c} \frac{P}{4 \pi\left(x^{2}+\eta^{2} R^{2}\right)} \cdot \cos ^{2} \theta \\
= & \frac{P}{2 \pi c} \frac{\eta^{2} R^{2}}{\left(x^{2}+\eta^{2} R^{2}\right)^{2}}
\end{aligned}
\]

The total force is then

\[
\begin{gathered}
\int_{0}^{R} \frac{P}{2 \pi c} \frac{\eta^{2} R^{2}}{\left(x^{2}+\eta^{2} R^{2}\right)^{2}} 2 \pi x d x \\
=\frac{P \eta^{2} R^{2}}{2 c} \int_{\eta^{2} R^{2}}^{R^{2}\left(1+\eta^{2}\right)} \frac{d y}{y^{2}} \\
=\frac{P \eta^{2} R^{2}}{2 c}\left[\frac{1}{\eta^{2} R^{2}}-\frac{1}{R^{2}\left(1+\eta^{2}\right)}\right]=\frac{P}{2 c\left(1+\eta^{2}\right)}
\end{gathered}
\]

\includegraphics[max width=\textwidth, alt={}, center]{da0156b6-4133-434e-90d4-ec01b2070afe-255_387_527_350_859}\\
5.267 (a) In the reference frame fixed to the mirror, the frequency of the photon is, by the Doppler shift formula

\[
\bar{\omega}=\omega \sqrt{\frac{1+\beta}{1-\beta}}\left(=\omega \frac{\sqrt{1-\beta^{2}}}{1-\beta}\right) .
\]

(see Eqn. (5.6b) of the book.)\\
In this frame momentum imparted to the mirrof is

\[
\frac{2 \hbar \bar{\omega}}{c}=\frac{2 \hbar \omega}{c} \sqrt{\frac{1+\beta}{1-\beta}},
\]

(b) In the \(K\) frame, the incident particle carries a momentum of \(\hbar \omega / c\) and returns with momentum

\[
\frac{\hbar \omega}{c} \frac{1+\beta}{1-\beta}
\]

(see problem 229). The momentum imparted to the mirror, then, has the magnitude

Here

\[
\begin{gathered}
\frac{\hbar \omega}{c}\left[\frac{1+\beta}{1-\beta}+1\right]=\frac{2 \hbar \omega}{c} \frac{1}{1-\beta} \\
\beta=\frac{V}{c} .
\end{gathered}
\]

5.268 When light falls on a small mirror and is reflected by it, the mirror recoils. The energy of recoil is obtained from the incident beam photon and the frequency of reflected photons is less than the frequency of the incident photons. This shift of frequency can however be neglected in calculating quantities related to recoil (to a first approximation.) Thus, the momentum acquired by the mirror as a result of the laser pulse is

\[
\left|\overrightarrow{p_{f}}-\overrightarrow{p_{1}}\right|=\frac{2 E}{c}
\]

Or assuming \(\overrightarrow{p_{1}}=0\), we get

\[
\left|\overrightarrow{p_{f}}\right|=\frac{2 E}{c}
\]

Hence the kinetic energy of the mirror is

\[
\frac{p_{f}^{2}}{2 m}=\frac{2 E^{2}}{m c^{2}}
\]

Suppose the mirror is deflected by an angle \(\theta\). Then by conservation of energy

\[
\begin{aligned}
\text { final P.E. }=m g l(1-\cos \theta)= & \text { Initial K.E. }=\frac{2 E^{2}}{m c^{2}} \\
\text { or } & m g l 2 \sin ^{2} \frac{\theta}{2}=\frac{2 E^{2}}{m c^{2}} \\
\text { or } & \sin \frac{\theta}{2}=\left(\frac{E}{m c}\right) \frac{1}{\sqrt{g l}}
\end{aligned}
\]

Using the data. \(\quad \sin \frac{\theta}{2}=\frac{13}{10^{-5} \times 3 \times 10^{8} \sqrt{9 \cdot 8 \times \cdot 1}}=4.377 \times 10^{-3}\)\\
This gives \(\theta=0.502\) degrees .\\
5.269 We shall only consider stars which are not too compact so that the gravitational field at their surface is weak :

\[
\frac{\gamma M}{c^{2} R} \ll 1
\]

We shall also clarify the problem by making clear the meaning of the (slightly changed) notation.\\
Suppose the photon is emitted by some atom whose total rclativistic energies (including the rest mass) are \(E_{1} \& E_{2}\) with \(E_{1}<E_{2}\). These energies are defined in the absence of gravitational field and we have

\[
\omega_{0}=\frac{E_{2}-E_{1}}{\hbar}
\]

as the frequency at infinity of the photon that is emitted in \(2 \rightarrow 1\) transition. On the surface of the star, the energies have the values

\[
\begin{aligned}
& E_{2}^{\prime}=E_{2}-\frac{E_{2}}{c^{2}} \cdot \frac{\gamma M}{R}=E_{2}\left(1-\frac{\gamma M}{c^{2} R}\right) \\
& E_{1}^{\prime}=E_{1}\left(1-\frac{\gamma M}{c^{2} R}\right)
\end{aligned}
\]

Thus, from \(\hbar \omega=E_{2}^{\prime}-E_{2}^{\prime}\) we get

\[
\omega=\omega_{0}\left(1-\frac{\gamma M}{c^{2} R}\right)
\]

Here \(\omega\) is the frequency of the photon emitted in the transition \(2 \rightarrow 1\) when the atom is on the surface of the star. In shows that the frequency of spectral lines emitted by atoms on the surface of some star is less than the frequency of lines emitted by atoms here on earth (where the gravitational effect is quite small).

Finally

\[
\frac{\Delta \omega}{\omega_{0}} \cong-\frac{\gamma M}{c^{2} R} .
\]

The answer given in the book is incorrect in general though it agrees with the above result for \(\frac{\gamma M}{c^{2} R} \ll 1\).\\
5.270 The general formula is

\[
\frac{2 \pi \hbar c}{\lambda}=e V
\]

Thus

\[
\lambda=\frac{2 \pi \hbar c}{e V}
\]

Now

\[
\Delta \lambda=\frac{2 \pi \hbar c}{e V}\left(1-\frac{1}{\eta}\right)
\]

Hence

\[
V=\frac{2 \pi \hbar c}{e \Delta \lambda}\left(\frac{\eta-1}{\eta}\right)=15.9 \mathrm{kV}
\]

5.271 We have as in the above problem

\[
\frac{2 \pi \hbar c}{\lambda}=e V
\]

On the other hand, from Bragg's law

\[
2 d \sin \alpha=k \lambda=\lambda
\]

since \(k=1\) when \(\alpha\) takes its smallest value.\\
Thus

\[
V=\frac{\pi \hbar c}{e d \sin \alpha}=30.974 k V \sim 31 \mathrm{kV}
\]

5.272 The wavelength of \(X\) - rays is the least when all the K.E. of the electrons approaching the anticathode is converted into the energy of \(X\) - rays.\\
But the K.E. of electron is

\[
\begin{aligned}
& \quad T_{\mathrm{m}}=m c^{2}\left[\frac{1}{\sqrt{1-\mathrm{v}^{2} / c^{2}}}-1\right] \\
& \left(m c^{2}=\text { rest mass energy of electrons }=0.511 \mathrm{MeV}\right) \\
& \text { Thus } \\
& \text { or } \begin{aligned}
& \quad \lambda=\frac{2 \pi \hbar c}{\lambda}=T_{m} \\
& \quad= \frac{2 \pi c}{m c(\gamma-1)}, \gamma=\frac{2 \pi \hbar}{m c}\left[\frac{1}{\sqrt{1-\mathrm{v}^{2} / c^{2}}}-1\right]^{-1} \\
& \frac{1-\mathrm{v}^{2} / c^{2}}{m c}=2.70 \mathrm{pm}
\end{aligned}
\end{aligned}
\]

5.273 The work function of zinc is

\[
A=3.74 \text { romane }=3.74 \times 1.602 \times 10^{-19} \text { Joule }
\]

The threshold wavelength for photoelectric effect is given by\\
or

\[
\begin{gathered}
\frac{2 \pi \hbar c}{\lambda_{0}}=A \\
\lambda_{0}=\frac{2 \pi \hbar c}{A}=331.6 \mathrm{~nm}
\end{gathered}
\]

The maximum velocity of photoelectrons liberated by light of wavelength \(\lambda\) is given by

So

\[
\begin{gathered}
\frac{1}{2} m v_{\max }^{2}=2 \pi \hbar c\left(\frac{1}{\lambda}-\frac{1}{\lambda_{0}}\right) \\
v_{\max }=\sqrt{\frac{4 \pi \hbar c}{m}\left(\frac{1}{\lambda}-\frac{1}{\lambda_{0}}\right)}=6.55 \times 10^{5} \mathrm{~m} / \mathrm{s}
\end{gathered}
\]

5.274 From the last equation of the previous problem, we find

Thus

\[
\eta=\frac{\left(v_{1}\right)_{\max }}{\left(v_{2}\right)_{\max }}=\sqrt{\frac{\frac{1}{\lambda_{1}}-\frac{1}{\lambda_{0}}}{\frac{1}{\lambda_{2}}-\frac{1}{\lambda_{0}}}}
\]

or

\[
\eta^{2}\left(\frac{1}{\lambda_{2}}-\frac{1}{\lambda_{0}}\right)=\frac{1}{\lambda_{1}}-\frac{1}{\lambda_{0}}
\]

\[
\frac{1}{\lambda_{0}}\left(\eta^{2}-1\right)=\frac{\eta^{2}}{\lambda_{2}}-\frac{1}{\lambda_{1}}
\]

and

So

\[
\begin{gathered}
\frac{1}{\lambda_{0}}=\left(\frac{\eta^{2}}{\lambda_{2}}-\frac{1}{\lambda_{1}}\right) /\left(\eta^{2}-1\right) \\
A=\frac{2 \pi \hbar c}{\lambda_{0}}=\frac{2 \pi \hbar c}{\lambda_{2}} \frac{\eta^{2}-\frac{\lambda_{2}}{\lambda_{1}}}{\eta^{2}-1}=1.88 \mathrm{eV}
\end{gathered}
\]

5.275 When light of sufficiently short wavelength falls on the ball, photoelectrons are ejected and the copper ball gains positive change. The charged ball tends to resist further emission of electrons by attracting them. When the copper ball has enough charge even the most energetic electrons are unable to leave it. We can calculate this final maximum potential of the copper ball. It is obviously equal in magnitude (in volt) to the maximum K.F of electrons (in electron volts) initially emitted. Hence

\[
\begin{aligned}
\varphi_{\max } & =\frac{2 \pi \hbar c}{\lambda e}-A_{c u} \\
& =8.86-4.47=4.39 \text { volts }
\end{aligned}
\]

( \(\boldsymbol{A}_{\boldsymbol{c} \boldsymbol{u}}\) is the work function of copper.)\\
5.276 We write

\[
\begin{aligned}
E & =a(1+\cos \omega t) \cos \omega_{0} t \\
& =a \cos \omega_{0} t+\frac{a}{2}\left[\cos \left(\omega_{0}-\omega\right) t+\cos \left(\omega_{0}+\omega\right) t\right]
\end{aligned}
\]

It is obvious that light has three frequencies and the maximum K.E. of photo electrons ejected is

\[
h\left(\omega+\omega_{0}\right)-A_{L i}
\]

where \(A_{L i}=2.39 \mathrm{eV}\). Substituting we get 0.37 eV .\\
5.277 Suppose \(N\) photons fall on the photocell per sec. Then the power incident is

\[
N \frac{2 \pi \hbar c}{\lambda}
\]

This will give rise to a photocurrent of \(N \frac{2 \pi \hbar c}{\lambda} \cdot \mathrm{~J}\)\\
which means that

\[
N \frac{2 \pi \hbar c}{e \lambda} \cdot \mathrm{~J}
\]

electrons have been emitted. Thus the number of photoelectrons produced by each photon is

\[
w=\frac{2 \pi \hbar c J}{e \lambda}=0.0198 \approx 0.02
\]

A simple application of Einstein's equation

\[
\frac{1}{2} m v_{\max }^{2}=h v-h v_{0}=\frac{2 \pi \hbar c}{\lambda}-A_{c s}
\]

gives incorrect result in this case because the photoelectrons emitted by the Cesium electrode are retarded by the small electric field that exists between the cesium electrode and the Copper electrode even in the absence of external emf. This small electric field is caused by the contact potential difference whose magnitude equals the difference of work functions

\[
\frac{1}{e}\left(A_{c u}-A_{c s}\right) \text { volts. }
\]

Its physical origin is explained below.\\
The maximum velocity of the photoelectrons reaching the copper electrode is then

\[
\frac{1}{2} m v_{\mathrm{m}}^{2}=\frac{1}{2} m v_{0}^{2}-\left(A_{c u}-A_{c s}\right)=\frac{2 \pi \hbar c}{\lambda}-A_{c u}
\]

Here \(\mathbf{v}_{0}\) is the maximum velocity of the photoelectrons immediately after emission. Putting the values we get, on using \(A_{c u}=4.47 \mathrm{eV}, \lambda=0.22 \mu \mathrm{~m}\),

\[
v_{m}=6.41 \times 10^{5} \mathrm{~m} / \mathrm{s}
\]

The origin of contact potential difference is the following. Inside the metals free electrons can be thought of as a Fermi gas which occupy energy levels upto a maximum called the Fermi energy \(E_{F}\). The work function \(A\) measures the depth of the Fermi level.\\
\includegraphics[max width=\textwidth, alt={}, center]{da0156b6-4133-434e-90d4-ec01b2070afe-260_505_453_868_193}\\
\includegraphics[max width=\textwidth, alt={}, center]{da0156b6-4133-434e-90d4-ec01b2070afe-260_394_617_900_798}

When two metals 1 \& 2 are in contact, electrons flow from one to the other till their Fermi levels are the same. This requires the appearance of contact potential difference of \(A_{1}-A_{2}\) between the two metals externally.\\
5.279 The maximum K.E. of the photoelectrons emitted by the Zn cathode is

\[
E_{\max }=\frac{2 \pi \hbar c}{\lambda}-A_{2 n}
\]

On calculating this comes out to be \(0.993 \mathrm{eV} \approx 1.0 \mathrm{eV}\)\\
Since an external decelerating voltage of 1.5 V is required to cancel this current, we infer that a contact potential difference of \(1.5-1.0=0.5 \mathrm{~V}\) exists in the circuit whose polarity is opposite of the decelerating voltage.\\
5.280 The unit of \(\hbar\) is Joule-sec. Since \(m c^{2}\) is the rest mass energy, \(\frac{\hbar}{m c^{2}}\) has the dimension of time and multiplying by \(c\) we get a quantity

\[
\lambda_{c}=\frac{\hbar}{m c} .
\]

whose dimension is length. This quantity is called reduced compton wavelength.\\
(The name compton wavelength is traditionally reserved for \(\frac{2 \pi \hbar}{m c}\) ).\\
5.281 We consider the collision in the rest frame of the initial electron. Then the reaction is

\[
\gamma+e(\text { rest }) \rightarrow e(\text { moving })
\]

Energy momentum conservation gives

\[
\begin{aligned}
\hbar \omega+m_{0} c^{2} & =m_{0} c^{2} / \sqrt{1-\beta^{2}} \\
\frac{\hbar \omega}{c} & =\frac{m_{0} c \beta}{\sqrt{1-\beta^{2}}}
\end{aligned}
\]

where \(\omega\) is the angular frequency of the photon.\\
Eliminating \(h \omega\) we get

\[
m_{0} c^{2}=m_{0} c^{2} \frac{1-\beta}{\sqrt{1-\beta^{2}}}=m_{0} c^{2} \sqrt{\frac{1-\beta}{1+\beta}}
\]

This gives \(\beta=0\) which implies \(\hbar \omega=0\).\\
But a zero energy photon means no photon.\\
5.282 (a) Compton scattering is the scattering of light by free electrons. (The free electrons are the electrons whose binding is much smaller than the typical energy transfer to the electrons). For this reason the increase in wavelength \(\Delta \lambda\) is independent of the nature of the scattered substance.\\
(b) This is because the effective number of free electrons increases in both cases. With increasing angle of scattering, the energy transfered to electrons increases. With diminishing atomic number of the substance the binding energy of the electrons decreases.\\
(c) The presence of a non-displaced component in the scattered radiation is due to scattering from strongy bound (inner) electrons as well as nuclei. For scattering by these the atom essentially recoils as a whole and there is very little energy transfer.\\
5.283 Let \(\lambda_{0}=\) wavelength of the incident radiation.

Then\\
wavelength of the radiation scattened at \(\boldsymbol{\theta}_{1}=60^{\circ}\)

\[
=\lambda_{1}=\lambda_{0}+2 \pi \lambda_{c}\left(1-\cos \theta_{1}\right) \quad \text { where } \lambda_{c}=\frac{\hbar}{m c} .
\]

and similarly

\[
\lambda_{2}=\lambda_{0}+2 \pi \lambda_{c}\left(1-\cos \theta_{2}\right)
\]

From the data \(\theta_{1}=60^{\circ}, \theta_{2}=120^{\circ}\) and

\[
\lambda_{2}=\eta \lambda_{1}
\]

Thus

\[
\begin{aligned}
(\eta-1) \lambda_{0} & =2 \pi \lambda_{c}\left[1-\cos \theta_{2}-\eta\left(1-\cos \theta_{1}\right)\right] \\
& =2 \pi \lambda_{c}\left[1-\eta+\eta \cos \theta_{1}-\cos \theta_{2}\right]
\end{aligned}
\]

Hence

\[
\begin{aligned}
\lambda_{0} & =2 \pi \lambda_{c}\left[\frac{\eta \cos \theta_{1}-\cos \theta_{2}}{\eta-1}-1\right] \\
& =4 \pi \lambda_{c}\left[\frac{\sin ^{2} \theta_{2} / 2-\eta \sin ^{2} \theta_{1} / 2}{\eta-1}\right]=1.21 \mathrm{pm} .
\end{aligned}
\]

The expression \(\lambda_{0}\) given in the book contains misprints.\\
5.284 The wave lengths of the photon has increased by a fraction \(\eta\) so its final wavelength is

\[
\begin{gathered}
\lambda_{f}=(2+\eta) \lambda_{i} \\
\frac{\hbar \omega}{1+\eta}
\end{gathered}
\]

and its energy is\\
The K.E. of the compton electron is the energy lost by the photon and is

\[
T=\hbar \omega\left(1-\frac{1}{1+\eta}\right)=\hbar \omega \frac{\eta}{1+\eta}
\]

5.285 (a) From the Compton formula

Thus

\[
\begin{gathered}
\lambda^{\prime}=2 \pi \lambda_{c}(1-\cos 90)+\lambda \\
\omega^{\prime}=\frac{2 \pi c}{\lambda^{\prime}}=\frac{2 \pi c}{\lambda+2 \pi \lambda_{c}} \quad \text { where } 2 \pi \lambda_{c}=\frac{h}{m c} .
\end{gathered}
\]

Substituting the values. we get \(\omega^{\prime}=2.24 \times 10^{20} \mathrm{rad} / \mathrm{sec}\)\\
(b) The kinetic energy of the scattered electron (in the frame in which the initial electron was stationary) is simply

\[
T=\hbar \omega-\hbar \omega^{\prime}
\]

\[
\begin{aligned}
& =\frac{2 \pi \hbar c}{\lambda}-\frac{2 \pi \hbar c}{\lambda+2 \pi \lambda_{c}} \\
& =\frac{4 \pi^{2} \hbar c \lambda_{c}}{\lambda\left(\lambda+2 \pi \lambda_{c}\right)}=\frac{2 \pi \hbar c / \lambda}{1+\lambda / 2 \pi \lambda_{c}}=59.5 \mathrm{kV}
\end{aligned}
\]

5.286 The wave length of the incident photon is

\[
\lambda_{0}=\frac{2 \pi c}{\omega}
\]

Then the wavelength of the final photon is

\[
\frac{2 \pi c}{\omega}+2 \pi \lambda_{c}(1-\cos \theta)
\]

and the energy of the final photon is

\[
\begin{aligned}
\hbar \omega^{\prime} & =\frac{2 \pi \hbar c}{\frac{2 \pi c}{\omega}+2 \pi \lambda_{c}(1-\cos \theta)}=\frac{h \omega}{1+\frac{h \omega}{m c^{2}}(1-\cos \theta)} \\
& =\frac{h \omega}{1+2\left(\frac{h \omega}{m c^{2}}\right) \sin ^{2}(\theta / 2)}=144 \cdot 2 \mathrm{kV}
\end{aligned}
\]

5.287 We use the equation \(\lambda=\frac{h}{p}=\frac{2 \pi \hbar}{p}\).

Then from Compton formula\\
so

\[
\begin{gathered}
\frac{2 \pi \hbar}{p^{\prime}}=\frac{2 \pi \hbar}{p}+2 \pi \frac{\hbar}{m c}(1-\cos \theta) \\
\frac{1}{p^{\prime}}=\frac{1}{p}+\frac{1}{m c} \cdot 2 \sin ^{2} \theta / 2
\end{gathered}
\]

Hence

\[
\sin ^{2} \frac{\theta}{2}=\frac{m c}{2}\left(\frac{1}{p^{\prime}}-\frac{1}{p}\right)
\]

\[
=\frac{m c\left(p-p^{\prime}\right)}{2 p p^{\prime}}
\]

or

\[
\sin \frac{\theta}{2}=\sqrt{\frac{m c\left(p-p^{\prime}\right)}{2 p p^{\prime}}},
\]

Substituting from the data

\[
\sin \frac{\theta}{2}=\sqrt{\frac{m c^{2}\left(c p-c p^{\prime}\right)}{2 c p \cdot c p^{\prime}}}=\sqrt{\frac{0.511(1.02-0.255)}{2 \times 1.02 \times 0.255}}
\]

This gives \(\theta=120 \cdot 2\) degrees .\\
5.288 From the Compton formula

\[
\lambda=\lambda_{0}+\frac{2 \pi \hbar}{m c}(1-\cos \theta)
\]

From conservation of energy\\
or

\[
\begin{gathered}
\frac{2 \pi \hbar c}{\lambda_{0}}=\frac{2 \pi \hbar c}{\lambda}+T=\frac{2 \pi \hbar c}{\lambda_{0}+\frac{2 \pi \hbar}{m c}(1-\cos \theta)}+T \\
\frac{4 \pi \hbar}{m c} \sin ^{2} \frac{\theta}{2}=\frac{T}{2 \pi \hbar c} \lambda_{0}\left(\lambda_{0}+\frac{4 \pi \hbar}{m c} \sin ^{2} \frac{\theta}{2}\right)
\end{gathered}
\]

or introducing \(\hbar \omega_{0}=2 \pi \hbar c / \lambda_{0}\)

Hence

\[
\begin{gathered}
\frac{2 \sin ^{2} \theta / 2}{m c^{2}}=\frac{T}{\hbar \omega_{0}}\left(\frac{1}{\hbar \omega_{0}}+\frac{2}{m c^{2}} \sin ^{2} \frac{\theta}{2}\right) \\
\left(\frac{1}{\hbar \omega_{0}}\right)^{2}+2 \frac{1}{\hbar \omega_{0}} \frac{\sin ^{2} \frac{\theta}{2}}{m c^{2}}-\frac{2 \sin ^{2} \frac{\theta}{2}}{m c^{2} T}=0 \\
\left(\frac{1}{\hbar \omega_{0}}+\frac{\sin ^{2} \frac{\theta}{2}}{m c^{2}}\right)^{2}=\frac{2 \sin ^{2} \frac{\theta}{2}}{m c^{2} T}+\left(\frac{\sin ^{2} \frac{\theta}{2}}{m c^{2}}\right)^{2} \\
\frac{1}{\hbar \omega_{0}}=\frac{\sin ^{2} \frac{\theta}{2}}{m c^{2}}\left[\sqrt{1+\frac{2 m c^{2}}{T \sin ^{2} \theta / 2}}-1\right] \\
\hbar \omega_{0}=\frac{m c^{2} / \sin ^{2} \theta / 2}{\sqrt{1+\frac{2 m c^{2}}{T \sin ^{2} \frac{\theta}{2}}-1}} \\
=\frac{T}{2}\left[\sqrt{1+\frac{2 m c^{2}}{T \sin ^{2} \theta / 2}}+1\right] \\
\hbar \omega_{0}=0 \cdot 677 \mathrm{MeV}
\end{gathered}
\]

Substituting we get\\
5.289 We see from the previous problem that the electron gains the maximum K.E. when the photon is scattered backwards \(\theta=180^{\circ}\). Then

Hence

\[
\begin{gathered}
\omega_{0}=\frac{m c^{2} / h}{\sqrt{1+\frac{2 m c^{2}}{T_{\text {max }}}-1}} \\
\lambda_{0}=\frac{2 \pi c}{\omega_{0}}=\frac{2 \pi \hbar}{m c}\left[\sqrt{1+\frac{2 m c^{2}}{T_{\text {max }}}}-1\right]
\end{gathered}
\]

Substituting the values we get \(\quad \lambda_{0}=3.695 \mathrm{pm}\).\\
5.290 Refer to the diagram. Energy momentum conservation gives

\[
\begin{aligned}
& \frac{\hbar \omega^{\prime}}{c}-\frac{\hbar \omega^{\prime}}{c} \cos \theta=p \cos \varphi \\
& \frac{h \omega^{\prime}}{c} \sin \theta=p \sin \varphi \\
& h \omega+m c^{2}=h \omega^{\prime}+E \\
& \text { here } \quad E^{2}=c^{2} p^{2}+m^{2} c^{4} . \text { we see }
\end{aligned}
\]

\begin{center}
\includegraphics[max width=\textwidth, alt={}]{da0156b6-4133-434e-90d4-ec01b2070afe-265_362_599_604_695}
\end{center}

\[
\begin{aligned}
\tan \varphi & =\frac{\omega^{\prime} \sin \theta}{\omega-\omega^{\prime} \cos \theta}=\frac{\frac{1}{\lambda^{\prime}} \sin \theta}{\frac{1}{\lambda}-\frac{1}{\lambda^{\prime}} \cos \theta} \\
& =\frac{\lambda \sin \theta}{\lambda^{\prime}-\lambda \cos \theta}=\frac{\sin \theta}{\frac{\Delta \lambda}{\lambda}+2 \sin ^{2} \frac{\theta}{2}}
\end{aligned}
\]

where

\[
\Delta \lambda=\lambda^{\prime}-\lambda=2 \pi \lambda_{c}(1-\cos \theta)=4 \pi \lambda_{c} \sin ^{2} \frac{\theta}{2}
\]

Hence

\[
\tan \varphi=\frac{2 \sin \frac{\theta}{2} \cos \frac{\theta}{2}}{\frac{\Delta \lambda}{\lambda}+\frac{\Delta \lambda}{2 \pi \lambda_{c}}}
\]

But

\[
\sin \theta=2 \sqrt{\frac{\Delta \lambda}{4 \pi \lambda_{c}}} \sqrt{1-\frac{\Delta \lambda}{4 \pi \lambda_{c}}}=\frac{\Delta \lambda}{2 \pi \lambda_{c}} \sqrt{\frac{4 \pi \lambda_{c}}{\Delta \lambda}-1}
\]

Thus

\[
\tan \varphi=\frac{\sqrt{\frac{4 \pi \hbar}{m c \Delta \lambda}-1}}{1+\frac{2 \pi \hbar}{m c \lambda}}=\frac{\sqrt{\frac{4 \pi \hbar}{m c \Delta \lambda}-1}}{1+\frac{\hbar \omega}{m c^{2}}}=31.3^{\circ}
\]

5.291 By head on collision we understand that the electron moves on in the direction of the incident photon after the collision and the photon is scaltered backwards. Then, let us write

\[
\begin{gathered}
\hbar \omega=\eta m c^{2} \\
\hbar \omega^{\prime}=\sigma m c^{2} \\
(E, p)=\left(\varepsilon m c^{2}, \mu m c\right) \text { of the electron. }
\end{gathered}
\]

Then by energy momentum conservation (cancelling factors of \(\boldsymbol{m} \boldsymbol{c}^{\mathbf{2}}\) and \(\boldsymbol{m} \boldsymbol{c}\) )

\[
\begin{aligned}
1+\eta & =\sigma+\varepsilon \\
\eta & =\mu-\sigma \\
\varepsilon^{2} & =1+\mu^{2}
\end{aligned}
\]

So eliminating \(\sigma \& \varepsilon\)

\[
1+\eta=-\eta+\mu+\sqrt{\mu^{2}+1}
\]

or

\[
(1+2 \eta-\mu)=\sqrt{\mu^{2}+1}
\]

Squaring

\[
\begin{gathered}
(1+2 \eta)^{2}-2 \mu(1+2 \eta)=1 \\
4 \eta+4 \eta^{2}=2 \mu(1+2 \eta)
\end{gathered}
\]

or

\[
\mu=\frac{2 \eta(1+\eta)}{1+2 \eta}
\]

Thus the momentum of the Compton electron is

Now in a magnetic field

\[
p=\mu m c=\frac{2 \eta(1+\eta) m c}{1+2 \eta} .
\]

Thus

\[
p=B e \rho
\]

Substituting the values

\[
\rho=2 \eta(1+\eta) /(1+2 \eta) \frac{m c}{B e} .
\]

\[
\rho=3.412 \mathrm{~cm}
\]

5.292 This is the inverse of usual compton scattering.

When we write down the energy-momentum conservation equation for this process we find that they are the same for the inverse process as they are for the usual process. If follows that the formula for compton shift is applicable except that the energy (frequency) of the photon is increased on scattering and the wavelength is shifted downward. With this understanding, we write

\[
\begin{aligned}
\Delta \lambda & =2 \pi \frac{\hbar}{m c}(1-\cos \theta) \\
& =4 \pi\left(\frac{\hbar}{m c}\right) \sin ^{2} \frac{\theta}{2}=1.21 \mathrm{pm}
\end{aligned}
\]

\begin{center}
\includegraphics[max width=\textwidth, alt={}]{da0156b6-4133-434e-90d4-ec01b2070afe-266_446_550_1427_846}
\end{center}

\section*{PART SIX}
\section*{ATOMIC AND NUCLEAR PHYSICS}
In this chapter the formulas in the book are given in the CGS units. Since most students are familar only with MKS units, we shall do the problems in MKS units. However, where needed, we shall also write the formulas in the Gaussian units.

\subsection*{6.1 SCATTERING OF PARTICLES. RUTHERFORD-BOHR ATOM}
6.1 The Thomson model consists of a uniformly charged nucleus in which the electrons are at rest at certain equilibrium points (the plum in the pudding model). For the hydrogen nucleus the charge on the nucleus is \(+e\) while the charge on the electron is - \(e\). The electron by symmetry must be at the centre of the nuclear charge where the potential (from problem (3.38a)) is

\[
\varphi_{0}=\left(\frac{1}{4 \pi \varepsilon_{0}}\right) \frac{3 e}{2 R}
\]

where \(R\) is the radius of the nucleur charge distribution. The potential energy of the electron is \(-e \varphi_{0}\) and since the electron is at rest, this is also the total energy. To ionize such an electron will require an energy of \(E=e \varphi_{0}\)

From this we find

\[
R=\left(\frac{1}{4 \pi \varepsilon_{0}}\right) \frac{3 e^{2}}{2 E}
\]

In Gaussian system the factor \(\frac{1}{4 \pi \varepsilon_{0}}\) is missing.\\
Putting the values we get \(\quad R=0.159 \mathrm{~nm}\).\\
Light is emitted when the electron vibrates. If we displace the electron slightly inside the nucleus by giving it a push \(r\) in some radial direction and an energy \(\delta E\) of oscillation then since the potential at a distance \(r\) in the nucleus is

\[
\varphi(r)=\left(\frac{1}{4 \pi \varepsilon_{0}}\right) \frac{e}{R}\left(\frac{3}{2}-\frac{r^{2}}{2 R^{2}}\right) .
\]

the total energy of the nucleus becomes\\
or

\[
\begin{gathered}
\frac{1}{2} m r^{2}-\left(\frac{1}{4 \pi \varepsilon_{0}}\right) \frac{e^{2}}{R}\left(\frac{3}{2}-\frac{r^{2}}{2 R^{2}}\right)=-e \varphi_{0}+\delta E \\
\delta E=\frac{1}{2} m \dot{r}^{2}+\left(\frac{1}{4 \pi \varepsilon_{0}}\right) \frac{e^{2}}{2 R^{3}} r^{2}
\end{gathered}
\]

This is the energy of a harmonic oscillator whose frequency is :

\[
\omega^{2}=\left(\frac{1}{4 \pi \varepsilon_{0}}\right) \frac{e^{2}}{m R^{3}}
\]

The vibrating electron emits radiation of frequency \(\omega\) whose wavelength is

\[
\lambda=\frac{2 \pi c}{\omega}=\frac{2 \pi c}{e} \sqrt{m R^{3}}\left(4 \pi \varepsilon_{0}\right)^{1 / 2}
\]

In Gaussian units the factor \(\left(4 \pi \varepsilon_{0}\right)^{1 / 2}\) is missing.\\
Putting the values we get \(\lambda=0.237 \mu \mathrm{~m}\).\\
6.2 Equation (6.1a) of the book reads in MKS units

\[
\tan \theta / 2=\left(\frac{q_{1} q_{2}}{4 \pi \varepsilon_{0}}\right) / 2 b T
\]

Thus

\[
b=\left(\frac{q_{1} q_{2}}{4 \pi \varepsilon_{0}}\right) \frac{\cot \theta / 2}{2 T}
\]

For \(\boldsymbol{\alpha}\) particl

\[
q_{1}=2 e, \text { for gold } q_{2}=79 e
\]

(In Gaussian units there is no factor \(\left(\frac{1}{4 \pi \varepsilon_{0}}\right)\).)\\
Substituting we get

\[
b=0.731 \mathrm{pm} .
\]

6.3 (a) In the Pb case we shall ignore the recoil of the nucleus both because Pb is quite heavy \(\left(A_{\mathrm{Pb}}=208=52 \times A_{H e}\right)\) as well as because Pb in not free. Then for a head on collision, at the distance of closest approach, the K.E. of the \(\alpha\)-particle must become zero (because \(\alpha\) - particle will turn back at this point). Then

\[
\frac{2 Z e^{2}}{\left(4 \pi \varepsilon_{0}\right) r_{\min }}=T
\]

(No ( \(4 \pi \varepsilon_{0}\) ) in Gaussian units.). Thus putting the values

\[
r_{\min }=0.591 \mathrm{pm} .
\]

(b) Here we have to take account of the fact that part of the energy is spent in the recoil of Li nucleus. Suppose \(x_{1}=\) coordinate of the \(\alpha\)-particle from some arbitrary point on the line joining it to the Li nucleus, \(x_{2}=\) coordinate of the Li nucleus with respect to the same point. Then we have the energy momentum equations

\[
\begin{gathered}
\frac{1}{2} m_{1} \dot{x}_{1}^{2}+\frac{1}{2} m_{2} \dot{x}_{2}^{2}+\frac{2 \times 3 e^{2}}{\left(4 \pi \varepsilon_{0}\right)\left|x_{1}-x_{2}\right|}=T \\
m_{1} \dot{x}_{1}+m_{2} \dot{x}_{2}=\sqrt{2 m_{1} T}
\end{gathered}
\]

Here \(m_{1}=\) mass of \(\mathrm{He}^{++}\)nucleus, \(m_{2}=\) mass of Li nucleus. Eliminating \(\dot{x}_{2}\)

\[
T=\frac{1}{2} m_{1} \dot{x}_{1}^{2}+\frac{1}{2 m_{2}}\left(\sqrt{2 m_{1} T}-m_{1} \dot{x}_{1}\right)^{2}+\frac{6 e^{2}}{\left(4 \pi \varepsilon_{0}\right)\left(x_{1}-x_{2}\right)}
\]

We complete the square on the right hand side and rewrite the above equation as

\[
\begin{gathered}
\frac{m_{2}}{m_{1}+m_{2}} T=\frac{1}{2 m_{2}}\left[\sqrt{m_{1}\left(m_{1}+m_{2}\right)} \dot{x}_{1}-\sqrt{\frac{m_{1}}{m_{1}+m_{2}}} \sqrt{2 m_{1} T}\right]^{2} \\
+\frac{6 e^{2}}{\left(4 \pi \varepsilon_{0}\right)\left|x_{1}-x_{2}\right|}
\end{gathered}
\]

For the least distance of apporach, the second term on the right must be greatest which implies that the first term must vanish.

Thus

\[
\left|x_{1}-x_{2}\right|_{\min }=\frac{6 e^{2}}{\left(4 \pi \varepsilon_{0}\right) T}\left(1+\frac{m_{1}}{m_{2}}\right)
\]

Using \(\frac{m_{1}}{m_{2}}=\frac{4}{7}\) and other values we get

\[
\left|x_{1}-x_{2}\right|_{\min }=.034 \mathrm{pm}
\]

(In Gaussian units the factor \(4 \pi \varepsilon_{0}\) is absent).\\
6.4 We shall ignore the recoil of Hg nucleus.\\
(a) Let \(A\) be the point of closest approach to the centre \(C, A C=r_{\text {min }}\). At \(A\) the motion is instantaneously circular because the radial velocity vanishes. Then if \(v_{0}\) is the speed of the particle at \(A\), the following equations hold


\begin{align*}
& \Gamma=\frac{1}{2} m v_{0}^{2}+\frac{Z_{1} Z_{2} e^{2}}{\left(4 \pi \varepsilon_{0}\right) r_{\min }}  \tag{1}\\
& m v_{0} r_{\min }=\sqrt{2 m T} b  \tag{2}\\
& \frac{m v_{0}^{2}}{\rho_{\min }}=\frac{Z_{1} Z_{2} e^{2}}{\left(4 \pi \varepsilon_{0}\right) r_{\min }^{2}} \tag{3}
\end{align*}


\includegraphics[max width=\textwidth, alt={}, center]{da0156b6-4133-434e-90d4-ec01b2070afe-269_462_561_958_872}\\
(This is Newton's law. Here \(\rho=\rho_{\min }\) is the radius of curvature of the path at \(A\) and \(\rho\) is minimum at \(A\) by symmetry.) Finally we have Eqn. (6.1 a) in the form


\begin{align*}
& \text { From (2) and (3) } \\
& \text { or } \\
& \begin{aligned}
\frac{2 T b^{2}}{\rho_{\min }} & =\frac{Z_{1} Z_{2} e^{2}}{\left(4 \pi \varepsilon_{0}\right)} \\
\text { with } & \rho_{\min }=\frac{Z_{1} Z_{2} e^{2}}{\left(4 \pi \varepsilon_{0}\right) 2 . T} \cot ^{2} \frac{\theta}{2}, \\
z_{1} & =2, z_{2}=80 \mathrm{we} \text { get } \\
\rho_{\min } & =0.231 \mathrm{pm} .
\end{aligned}
\end{align*}


(b) From (2) and (4) we write

\[
\begin{aligned}
& r_{\min }=\frac{Z_{1} Z_{2} e^{2}}{\left(4 \pi \varepsilon_{0}\right) \sqrt{2 m T}} \frac{\cot \theta / 2}{v_{0}} \\
& T=\frac{1}{2} m v_{0}^{2}+\sqrt{2 m T} v_{0} \tan \theta / 2
\end{aligned}
\]

Solving for \(v_{0}\) we get

Then

\[
\begin{aligned}
v_{0} & =\sqrt{\frac{2 T}{m}}\left(\sec \frac{\theta}{2}-\tan \frac{\theta}{2}\right) \\
r_{\min } & =\frac{Z_{1} Z_{2} e^{2}}{\left(4 \pi \varepsilon_{0}\right) 2 T} \frac{\cot \frac{\theta}{2}}{\sec \frac{\theta}{2}-\tan \frac{\theta}{2}} \\
& =\frac{Z_{1} Z_{2} e^{2}}{\left(4 \pi \varepsilon_{0}\right) 2 T} \cot \frac{\theta}{2}\left(\sec \frac{\theta}{2}+\tan \frac{\theta}{2}\right) \\
& =\frac{Z_{1} Z_{2} e^{2}}{\left(4 \pi \varepsilon_{0}\right) 2 T}\left(1+\operatorname{cosec} \frac{\theta}{2}\right)=0.557 \mathrm{pm}
\end{aligned}
\]

6.5 By momentum conservation

\[
\vec{P}+\overrightarrow{P_{i}}=\vec{P}+\vec{P}_{f}
\]

(proton) (Au) (proton) (Au)\\
Thus the momentum transfered to the gold nucleus is clearly

\[
\Delta \vec{P}=\vec{P}_{f}-\vec{P}_{i}=\vec{p}-\vec{p}
\]

Although the momentum transfered to the Au nucleus is not small, the energy associated with this recoil is quite small and its effect back on the motion of the proton can\\
\includegraphics[max width=\textwidth, alt={}, center]{da0156b6-4133-434e-90d4-ec01b2070afe-270_379_515_885_954}\\
be neglected to a first approximation. Then

\[
\Delta \vec{P} \Longleftarrow \sqrt{2 m T}(1-\cos \theta) \hat{i} \sqrt{2 m T} \sin \theta \hat{j}
\]

Here \(\hat{i}\) is the unit vector in the direction of the incident proton and \(\hat{j}\) is normal to it on the side on which it is scattered. Thus

\[
|\Delta \vec{P}| \sim 2 \sqrt{2 m T} \sin \frac{\theta}{2}
\]

Or using \(\tan \theta / 2=\frac{Z e^{2}}{\left(4 \pi \varepsilon_{0}\right) 2 b T}\) for the proton we get

\[
|\Delta \vec{P}| \sim 2 \sqrt{2 m T /\left(1+\left(\frac{2 b T\left(4 \pi \varepsilon_{0}\right)}{z e^{2}}\right)^{2}\right)}
\]

6.6 The proton moving by the electron first accelerates and then decelerates and it not easy to calculate the energy lost by the proton so energy conservation does not do the trick. Rather\\
we must directly calculate the momentum acquire by the electron. By symmetry that momentum is along \(\overrightarrow{O A}\) and its magnitude is

\[
P_{d} \propto \int F d t
\]

where \(F_{1}\) is the component along \(O A\) of the force on electron. Thus

\[
\begin{aligned}
P_{d} & =\int_{-\infty}^{\infty} \frac{e^{2}}{4 \pi \varepsilon_{0}} \frac{b}{\sqrt{b^{2}+v^{2} t^{2}}} \cdot \frac{1}{b^{2}+v^{2} t^{2}} d t \\
& =\frac{e^{2} b}{4 \pi \varepsilon_{0} v} \cdot \int_{-\infty}^{\infty} \frac{d x}{\left(b^{2}+x^{2}\right)^{3 / 2}}
\end{aligned}
\]

Evaluate the integral by substituting

Then

\[
\begin{gathered}
x=b \tan \theta \\
P_{e}=\frac{2 e^{2}}{\left(4 \pi \varepsilon_{0}\right) \nu b} \\
T_{e}=\frac{P_{e}^{2}}{2 m_{e}}=\frac{m_{p} e^{4}}{\left(4 \pi \varepsilon_{0}\right)^{2} T b^{2} m_{e}}
\end{gathered}
\]

\begin{center}
\includegraphics[max width=\textwidth, alt={}]{da0156b6-4133-434e-90d4-ec01b2070afe-271_292_451_501_924}
\end{center}

Then

In Gaussian units there is no factor \(\left(4 \pi \varepsilon_{0}\right)^{2}\). Substituting the values we get

\[
T_{e}=3.82 \mathrm{eV}
\]

6.7 See the diagram on the next page. In the region where potential is nonzero, the kinetic energy of the particle is, by energy conservation,\\
\(T+U_{0}\) and the momentum of the particle has the magnitude \(\sqrt{2 m\left(T+U_{0}\right)}\). On the boundary the force is radial, so the tangential component of the momentum does not change :\\
\(\sqrt{2 m T} \sin \alpha=\sqrt{2 m\left(T+U_{0}\right)} \sin \varphi\)\\
so \(\sin \varphi=\sqrt{\frac{T}{T+U_{0}}} \sin \alpha=\frac{\sin \alpha}{n}\)\\
where \(n=\sqrt{1+\frac{U_{0}}{T}}\). We also have

\[
\theta=2(\alpha-\varphi)
\]

Therefore\\
\includegraphics[max width=\textwidth, alt={}, center]{da0156b6-4133-434e-90d4-ec01b2070afe-271_373_597_1370_700}

\[
\sin \frac{\theta}{2}=\sin (\alpha-\varphi)=\sin \alpha \cos \varphi-\cos \alpha \sin \varphi
\]

or

\[
=\sin \alpha\left(\cos \varphi-\frac{\cos \alpha}{n}\right)
\]

or

\[
\frac{n \sin \theta / 2}{\sin \alpha}=\sqrt{n^{2}-\sin ^{2} \alpha}-\cos \alpha
\]

or\\
or

\[
\begin{gathered}
\left(\frac{n \sin \theta / 2}{\sin \alpha}+\cos \alpha\right)^{2}=n^{2}-\sin ^{2} \alpha \\
n^{2} \sin ^{2} \frac{\theta}{2} \cot ^{2} \alpha+2 n \sin \frac{\theta}{2} \cot \alpha+1=n^{2} \cos ^{2} \frac{\theta}{2} \\
\cot \alpha=\frac{n \cos \frac{\theta}{2}-1}{n \sin \frac{\theta}{2}} \\
\sin \alpha=\frac{n \sin \frac{\theta}{2}}{\sqrt{1+n^{2}-2 n \cos \frac{\theta}{2}}}
\end{gathered}
\]

Hence

Finally, the impact parameter is

\[
b=R \sin \alpha=\frac{n R \sin \frac{\theta}{2}}{\sqrt{1+n^{2}-2 n \cos \frac{\theta}{2}}} .
\]

6.8 It is implied that the ball is too heavy to recoil.\\
(a) The trajectory of the particle is symmetrical about the radius vector through the point of impact. It is clear from the diagram that

Also

\[
\begin{aligned}
& \theta=\pi-2 \varphi \text { or } \varphi=\frac{\pi}{2}-\frac{\theta}{2} . \\
& b=(R+r) \sin \varphi=(R+r) \cos \frac{\theta}{2} .
\end{aligned}
\]

\includegraphics[max width=\textwidth, alt={}, center]{da0156b6-4133-434e-90d4-ec01b2070afe-272_500_507_1023_974}\\
(b) With \(b\) defined above, the fraction of particles scattered between \(\theta\) and \(\theta+d \theta\) (or the probability of the same) is

\[
d P=\frac{|2 \pi b d b|}{\pi(R+r)^{2}}=\frac{1}{2} \sin \theta d \theta
\]

(c) This is

\[
P=\int_{0}^{\pi / 2} \frac{1}{2} \sin \theta d \theta=\frac{1}{2} \int_{-1}^{0} d(-\cos \theta)=\frac{1}{2}
\]

6.9 From the formula (6.1 b) of the book

\[
\frac{d N}{N}=n\left(\frac{Z e^{2}}{\left(4 \pi \varepsilon_{0}\right) 2 T}\right)^{2} \frac{d \Omega}{\sin ^{4} \frac{\theta}{2}}
\]

We have put \(q_{1}=2 e, q_{2}=Z e\) here. Also \(n=\) no. of Pt nuclei in the foil per unit area

\[
=\begin{array}{cc}
(A t \rho) & \frac{N_{A}}{A_{P t}} \cdot \frac{1}{A}=\frac{N_{A} \rho t}{A_{P t}} \\
\downarrow & \downarrow \\
\begin{array}{c}
\text { mass of } \\
\text { the foil }
\end{array} & \begin{array}{c}
\downarrow \\
\text { no. of } \\
\text { nuclei per } \\
\text { unit mass }
\end{array}
\end{array}
\]

Using the values \(A_{P t}=195, \rho=21.5 \times 10^{3} \mathrm{~kg} / \mathrm{m}^{3}\)

\[
\begin{aligned}
N_{A} & =6.023 \times 10^{26} / \text { kilo mole } \\
n & =6.641 \times 10^{22} \text { per } m^{2}
\end{aligned}
\]

we get\\
Also

\[
d \Omega=\frac{d S_{n}}{r^{2}}=10^{-2} \mathrm{Sr}
\]

Substituting we get

\[
\frac{d N}{N}=3.36 \times 10^{-5}
\]

6.10 A scattered flux density of \(J\) (perticles per unit area per second) equals \(J / \frac{1}{r^{2}}=r^{2} J\) particles scattered per unit time per steradian in the given direction. Let \(n=\) concentration of the gold nuclei in the foil. Then

\[
n=\frac{N_{A} \rho}{A_{A u}}
\]

and the number of Au nuclei per unit area of the foil is \(n d\) where \(d=\) thickness of the foil Then from Eqn. (6.1 b) (note that \(n \rightarrow n d\) here)

\[
r^{2} J=d N=n d I\left(\frac{Z e^{2}}{\left(4 \pi \varepsilon_{0}\right) 2 T}\right)^{2} \operatorname{cosec}^{4} \frac{\theta}{2}
\]

Here \(I\) is the number of \(\alpha\)-particles falling on the foil per second\\
Hence

\[
d=\frac{4 T^{2} r^{2} J}{n I\left(\frac{Z e^{2}}{4 \pi \varepsilon_{0}}\right)^{2}} \sin ^{4} \theta / 2
\]

using \(Z=79, A_{A u}=197, \rho=19.3 \times 10^{3} \mathrm{~kg} / \mathrm{m}^{3} N_{A}=6.023 \times 10^{26} /\) kilo mole and other data from the problem we get

\[
d=1.47 \mu \mathrm{~m}
\]

6.11 From the formula ( 6.1 b) of the book, we find

\[
\frac{d N_{P t}}{d N_{A g}}=\frac{n_{P t}}{n_{A g}} \cdot \frac{Z_{P t}^{2}}{Z_{A g}^{2}}=\eta
\]

But since the foils have the same mass thickness ( \(-\rho d\) ) , we have

\[
\frac{n_{P t}}{n_{A g}}=\frac{A_{A g}}{A_{P t}}
\]

see the problem (6.9). Hence

\[
Z_{P t}=Z_{A g} \cdot \sqrt{\frac{\eta A_{A g}}{A_{P t}}}
\]

Substituting \(Z_{A g}=47, A_{A g}=108, A_{P t}=195\) and \(\eta=1.52\) we get

\[
Z_{P t}=77.86 \approx 78
\]

6.12 (a) From Eqn. ( 6.1 b) we get

\[
d N=I_{0} \tau \frac{\rho \cdot d \cdot N_{A}}{A_{A u}}\left(\frac{Z e^{2}}{\left(4 \pi \varepsilon_{0}\right) 2 T}\right)^{2} \frac{2 \pi \sin \theta d \theta}{\sin ^{4} \theta / 2}
\]

(we have used \(d \Omega=2 \pi \sin \theta d \theta\) and \(N=I_{0} I\) )\\
From the data

\[
d \theta=2^{\circ}=\frac{2}{57 \cdot 3} \text { radian }
\]

Also \(Z_{A u}=79, A_{A u}=197\). Putting the values we get

\[
d N=1.63 \times 10^{6}
\]

(b) This number is

\[
N\left(\theta_{0}\right)=I_{0} \tau\left(\frac{\rho d N_{A}}{A_{A u}}\right)\left(\frac{Z e^{2}}{\left(4 \pi \varepsilon_{0}\right) 2 T}\right)^{2} 4 \pi \int_{\theta_{0}}^{\pi} \frac{\cos \frac{\theta}{2} d \theta}{\sin ^{3} \frac{\theta}{2}}
\]

The integral is

\[
2 \int_{\sin \frac{\theta_{0}}{2}}^{1} \frac{d x}{x^{2}}=\frac{1}{2}\left[\frac{-1}{2 x^{2}}\right]_{\sin \frac{\theta_{0}}{2}}^{1}=\cot ^{2} \frac{\theta_{0}}{2}
\]

Thus

\[
N\left(\theta_{0}\right)=\pi n d\left(\frac{Z e^{2}}{\left(4 \pi \varepsilon_{0}\right) T}\right)^{2} I_{0} \tau \cot ^{2} \frac{\theta_{0}}{2}
\]

where \(n\) is the concentration of nuclei in the foil. ( \(n=\rho N_{1} / A_{A u}\) ) Substitution gives

\[
N\left(\theta_{0}\right)=2.02 \times 10^{7}
\]

6.13 The requisite probability can be written easily by analogy with (b) of the previous problem. It is

\[
P=\frac{N(\pi / 2)}{I_{0} \tau}=n d\left(\frac{Z e^{2}}{\left(4 \pi \varepsilon_{0}\right) 2 m v^{2}}\right)^{2} 4 \pi \int_{\pi / 2}^{\pi} \frac{\cos \theta / 2 d \theta}{\sin ^{3} \frac{\theta}{2}}
\]

The integral is unity. Thus

\[
P=\pi n d\left(\frac{Z e^{2}}{\left(4 \pi \varepsilon_{0}\right) m v^{2}}\right)^{2}
\]

Substitution gives using

\[
n=\frac{\rho_{A g} N_{A}}{A_{A g}}=\frac{10.5 \times 10^{3} \times 6.023 \times 10^{26}}{108}, P=0.06
\]

6.14 Because of the \(\operatorname{cosec}^{4} \frac{\theta}{2}\) dependence of the scattering, the number of particles (or fraction) scattered through \(\theta<\theta_{0}\) cannot be calculated directly. But we can write this fraction as

\[
P\left(\theta_{0}\right)=1-Q\left(\theta_{0}\right)
\]

where \(Q\left(\theta_{0}\right)\) is the fraction of particles scattered through \(\theta \geq \theta_{0}\). This fraction has been calculated before and is (see the results of 6.12 (b))

\[
Q\left(\theta_{0}\right)=\pi n\left(\frac{Z e^{2}}{\left(4 \pi \varepsilon_{0}\right) T}\right)^{2} \cot ^{2} \frac{\theta_{0}}{2}
\]

where \(n\) here is number of nuclei \(/ \mathrm{cm}^{2}\). Using the data we get

\[
Q=0.4
\]

Thus

\[
P\left(\theta_{0}\right)=0.6
\]

6.15 The relevant fraction can be immediately written down (see 6.12 (b)) (Note that the projectiles are protons)

\[
\frac{\Delta N}{N}=\left(\frac{e^{2}}{\left(4 \pi \varepsilon_{0}\right) 2 T}\right)^{2} \pi \cot ^{2} \frac{\theta_{0}}{2} \cdot\left(n_{1} Z_{1}^{2}+n_{2} Z_{2}^{2}\right)
\]

Here \(n_{1}\left(n_{2}\right)\) is the number of \(Z_{n}(C u)\) nuclei per \(\mathrm{cm}^{2}\) of the foil and \(Z_{1}\left(Z_{2}\right)\) is the atomic number of \(Z_{n}(C u)\). Now

\[
n_{1}=\frac{\rho d N_{A}}{M_{1}}=0.7, n_{2}=\frac{\rho d N_{A}}{M_{2}}=0.3
\]

Here \(M_{1}, M_{2}\) are the mass numbers of \(Z_{n}\) and \(C u\).\\
Then, substituting the values \(Z_{1}=30, Z_{2}=29, M_{1}=65 \cdot 4, M_{2}=63 \cdot 5\), we get

\[
\frac{\Delta N}{N}=1.43 \times 10^{-3}
\]

6.16 From the Rutherford scattering formula\\
or

\[
\begin{aligned}
\frac{d \sigma}{d \Omega} & =\left(\frac{Z e^{2}}{\left(4 \pi \varepsilon_{0}\right) 2 T}\right)^{2} \frac{1}{\sin ^{4} \frac{\theta}{2}} \\
d \sigma & =\left(\frac{Z e^{2}}{\left(4 \pi \varepsilon_{0}\right) 2 T}\right)^{2} \frac{2 \pi \sin \theta d \theta}{\sin ^{4} \frac{\theta}{2}} \\
& =\left(\frac{Z e^{2}}{\left(4 \pi \varepsilon_{0}\right) T}\right)^{2} \pi \frac{\cos \theta / 2 d \theta}{\sin ^{3} \theta / 2}
\end{aligned}
\]

Then integrating from \(\theta=\theta_{0}\) to \(\theta=\pi\) we get the required cross section

\[
\begin{aligned}
\Delta \sigma & =\left(\frac{Z e^{2}}{\left(4 \pi \varepsilon_{0}\right) T}\right)^{2} \pi \int_{\theta_{0}}^{\pi} \frac{\cos \theta / 2 d \theta}{\sin ^{3} \frac{\theta}{2}} \\
& =\left(\frac{Z e^{2}}{\left(4 \pi \varepsilon_{0}\right) 2 T}\right)^{2} \cot ^{2} \frac{\theta_{0}}{2}
\end{aligned}
\]

For U nucleus \(Z=92\) and we get on putting the values

\[
\Delta \sigma=737 \mathrm{~b}=0.737 \mathrm{~kb}
\]

( \(1 b=1\) barn \(=10^{-28} \mathrm{~m}^{2}\) ).\\
6.17 (a) From the previous formula\\
or

\[
\begin{aligned}
\Delta \sigma & =\left(\frac{Z e^{2}}{\left(4 \pi \varepsilon_{0}\right) 2 T}\right)^{2} \pi \cot ^{2} \frac{\theta_{0}}{2} \\
T & =\frac{Z e^{2}}{4 \pi \varepsilon_{0}} \cot \frac{\theta_{0}}{2} \sqrt{\frac{\pi}{\Delta \sigma}}
\end{aligned}
\]

Substituting the values with \(Z=79\) we get ( \(\theta_{0}=90^{\circ}\) )

\[
T=0.903 \mathrm{MeV}
\]

(b) The differential scattering cross section is

\[
\frac{d \sigma}{d \Omega}=C \operatorname{cosec}^{4} \frac{\theta}{2}
\]

where

\[
\Delta \sigma\left(\theta>\theta_{0}\right)=4 \pi C \cot ^{2} \frac{\theta_{0}}{2}
\]

Thus from the given data

\[
C=\frac{500}{4 \pi} b=39 \cdot 79 \mathrm{~b} / \mathrm{sr}
\]

So

\[
\frac{d \sigma}{d \Omega}\left(\theta=60^{\circ}\right)=39.79 \times 16 \mathrm{~b} / \mathrm{sr}=0.637 \mathrm{~kb} / \mathrm{sr} .
\]

6.18 The formula in MKS units is

\[
\frac{d E}{d t}=-\frac{\mu_{0} e^{2}}{6 \pi c} \overrightarrow{w^{2}}
\]

For an electron performing (linear) harmonic vibrations \(\vec{w}\) is in some definite directions with

\[
\begin{aligned}
w_{x} & =-\omega^{2} x \text { say. } \\
\frac{d E}{d t} & =-\frac{\mu_{0} e^{2} \omega^{4}}{6 \pi c} x^{2}
\end{aligned}
\]

Thus\\
If the radiation loss is small (i.e. if \(\omega\) is not too large), then the motion of the electron is always close to simple harmonic with slowly decreasing amplitude. Then we can write\\
and

\[
\begin{aligned}
E & =\frac{1}{2} m \omega^{2} a^{2} \\
x & =a \cos \omega t
\end{aligned}
\]

and average the above equation ignoring the variation of \(a\) in any cycle. Thus we get the equation, on using \(\left\langle x^{2}\right\rangle=\frac{1}{2} a^{2}\)

\[
\frac{d E}{d t}=-\frac{\mu_{0} e^{2} \omega^{4}}{6 \pi c} \cdot \frac{1}{2} a^{2}=-\frac{\mu_{0} e^{2} \omega^{2}}{6 \pi m c} E
\]

since \(E=\frac{1}{2} m \omega^{2} a^{2}\) for a harmonic oscillator.\\
This equation integrates to\\
where

\[
\begin{gathered}
E=E_{o} e^{-t / T} \\
T=6 \pi m c / e^{2} \omega^{2} \mu_{0}
\end{gathered}
\]

It is then seen that energy decreases \(\eta\) times in

\[
t_{0}=T \ln \eta=\frac{6 \pi m c}{e^{2} \omega^{2} \mu_{0}} \ln \eta=14.7 \mathrm{~ns}
\]

6.19 Moving around the nucleus, the electron radiates and its energy decreases. This means that the electron gets nearer the nucleus. By the statement of the problem we can assume that the electron is always moving in a circular orbit and the radial acceleration by Newton's law is

\[
w=\frac{e^{2}}{\left(4 \pi \varepsilon_{0}\right) m r^{2}}
\]

directed inwards. Thus

\[
\frac{d E}{d t}=-\frac{\mu_{0} e^{6}}{6 \pi c} \frac{1}{\left(4 \pi \varepsilon_{0}\right)^{2} m^{2} r^{4}}
\]

On the other hand in a circular orbit

\[
E=-\frac{e^{2}}{\left(4 \pi \varepsilon_{0}\right) 2 r}
\]

so

\[
\frac{e^{2}}{\left(4 \pi \varepsilon_{0}\right) 2 r^{2}} \frac{d r}{d t}=-\frac{\mu_{0} e^{6}}{\left(4 \pi \varepsilon_{0}\right)^{2} 6 \pi c m^{2} r^{4}}
\]

or

\[
\frac{d r}{d t}=-\frac{\mu_{0} e^{4}}{\left(4 \pi \varepsilon_{0}\right) 3 \pi c m^{2} r^{2}}
\]

Integrating

\[
r^{3}=r_{0}^{3}-\frac{\mu_{0} e^{4}}{4 \pi^{2} \varepsilon_{0} \mathrm{~cm}^{2}} t
\]

and the radius falls to zero in

\[
t_{0}=\frac{4 \pi^{2} \varepsilon_{0} c m^{2} r_{0}^{3}}{\mu_{0} e^{4}} \mathrm{sec} .=13 \cdot 1 \mathrm{ps}
\]

6.20 In a circular orbit we have the following formula

Then

\[
\begin{gathered}
\frac{m v^{2}}{r}=\frac{Z \dot{e}^{2}}{\left(4 \pi \varepsilon_{0}\right) r^{2}} \\
m v r=n \hbar \\
v=\frac{Z e^{2}}{\left(4 \pi \varepsilon_{0}\right) n \hbar} \\
r=\frac{n^{2}}{Z} \frac{\hbar\left(4 \pi \varepsilon_{0}\right)}{m e^{2}}
\end{gathered}
\]

The energy \(E\) is \(\quad E_{n}=\frac{1}{2} m \mathrm{v}^{2}-\frac{Z e^{2}}{\left(4 \pi \varepsilon_{0}\right) r}\)

\[
=\left(\frac{Z e^{2}}{4 \pi \varepsilon_{0}}\right)^{2} \frac{m}{2 \hbar^{2} n^{2}}-\left(\frac{Z e^{2}}{4 \pi \varepsilon_{0}}\right)^{2} \frac{m}{\hbar^{2} n^{2}}=m\left(\frac{Z e^{2}}{4 \pi \varepsilon_{0}}\right)^{2} / 2 \hbar^{2} n^{2}
\]

and the circular frequency of this orbit is

\[
\omega_{n}=\frac{v}{r}=\left(\frac{Z e^{2}}{4 \pi \varepsilon_{0}}\right)^{2} \cdot m / \hbar^{3} n^{3}
\]

On the other hand the frequency \(\omega\) of the light emitted when the electron makes a transition \(n+1 \rightarrow n\) is

\[
\omega=\left(\frac{Z e^{2}}{4 \pi \varepsilon_{0}}\right)^{2} \frac{m}{2 \hbar^{2}}\left(\frac{1}{n^{2}}-\frac{1}{(n+1)^{2}}\right)
\]

Thus the inequality

\[
\omega_{n}>\omega>\omega_{n+1}
\]

will result if

\[
\frac{1}{n^{3}}>\frac{1}{2}\left(\frac{1}{n^{2}}-\frac{1}{(n+1)^{2}}\right)>\frac{1}{(n+1)^{3}}
\]

Or multiplying by \(n^{2}(n+1)^{2}\) we have to prove

\[
\frac{(n+1)^{2}}{n}>\frac{1}{2}(2 n+1)>\frac{n^{2}}{n+1}
\]

This can be written as

\[
n+2+\frac{1}{n}>n+\frac{1}{2}>n+1-2+\frac{1}{n+1}
\]

This is obvious because \(-1+\frac{1}{n+1}<-\frac{1}{2}\) since \(n \geq 1\)\\
For large \(n\)

\[
\frac{\omega_{n}}{\omega_{n+1}}=\left(\frac{n+1}{n}\right)^{3}=1+\frac{3}{n}
\]

so

\[
\frac{\omega_{n}}{\omega_{n+1}} \rightarrow 1 \text { and we may say } \frac{\omega}{\omega_{n}} \rightarrow 1
\]

6.21 We have the following equation (we ignore reduced mass effects)

\[
\begin{aligned}
\frac{m v^{2}}{r} & =k r \\
m v r & =n \hbar \\
m \cdot v & =\sqrt{m k} r
\end{aligned}
\]

so

\[
r=\sqrt{\frac{n \hbar}{\sqrt{m k}}},
\]

and

\[
\mathrm{v}=\sqrt{n \hbar \sqrt{m k}} / \mathrm{m}
\]

The energy levels are

\[
\begin{aligned}
E_{n} & =\frac{1}{2} m \mathrm{v}^{2}+\frac{1}{2} k r^{2} \\
& =\frac{1}{2} \frac{n \hbar \sqrt{m k}}{m}+\frac{1}{2} k \frac{n \hbar}{\sqrt{m k}} \\
& =n \hbar \sqrt{\frac{k}{m}}
\end{aligned}
\]

6.22 The basic equations have been derived in the problem (6.20). We rewrite them here and determine the the required values.\\
(a) \(\quad r_{1}=\frac{\hbar^{2}}{m\left(Z e^{2} / 4 \pi \varepsilon_{0}\right)}, Z=1\) for \(H, Z=2\) for \(H e^{+}\)

Thus

\[
\begin{gathered}
r_{1}=52.8 \mathrm{pm}, \text { for } \mathrm{H} \text { atom } \\
r_{1}=26.4 \mathrm{pm}, \text { for } \mathrm{He}^{+} \text {ion } \\
v_{1}=\frac{Z e^{2}}{\left(4 \pi \varepsilon_{0}\right) h} \\
v_{1}=2.191 \times 10^{6} \mathrm{~m} / \mathrm{s} \text { for } \mathrm{H} \text { atom } \\
=4.382 \times 10^{6} \mathrm{~m} / \mathrm{s} \text { for } \mathrm{He}^{+} \text {ion }
\end{gathered}
\]

(b) \(T=\frac{1}{2} m v_{1}^{2}=\frac{m\left(Z e^{2}\right)^{2}}{\left(4 \pi \varepsilon_{0}\right)^{2} 2 \hbar^{2}}\)

\[
\begin{aligned}
& T=13.65 \mathrm{eV} \text { for } \mathrm{H} \text { atom } \\
& T=54.6 \mathrm{eV} \text { for } \mathrm{He}^{+} \text {ion }
\end{aligned}
\]

In both cases \(E_{b}=T\) because \(E_{b}=-E\) and \(E=-T\) (Recall that for coulomb force \(V=-2 T\) )\\
(c) The ionization potential \(\varphi_{i}\) is given by

\[
e \varphi_{i}=E_{b}
\]

so \(\varphi_{i}=13.65\) volts for H atom \(\varphi_{i}=54.6\) volts for \(\mathrm{He}^{+}\)ion\\
The energy levels are \(E_{n}=-\frac{13 \cdot 65}{n^{2}} \mathrm{eV}\) for H atom and

\[
E_{n}=-\frac{54 \cdot 6}{n^{2}} \mathrm{eV} \text { for } \mathrm{He}^{+} \text {ion }
\]

Thus

\[
\begin{aligned}
& \varphi_{1}=13.65\left(1-\frac{1}{4}\right) \text { volts }=10.23 \text { volts for } \mathrm{H} \text { atom } \\
& \varphi_{1}=4 \times 10.23=40.9 \text { volts for } \mathrm{He}^{+} \text {ion }
\end{aligned}
\]

The wavelength of the resonance line

\[
\begin{gathered}
\left(n^{\prime}=2 \rightarrow n=1\right) \text { is given by } \\
\frac{2 \pi \hbar c}{\lambda}=-\frac{13.6}{4}+\frac{13.6}{1}=10.23 \mathrm{eV} \text { for } \mathrm{H} \text { atom }
\end{gathered}
\]

so

\[
\lambda=121 \cdot 2 \mathrm{~nm} \text { for } \mathrm{H} \text { atom }
\]

For \(\mathrm{He}^{+}\)ion

\[
\lambda=\frac{121 \cdot 2}{4}=30 \cdot 3 \mathrm{~nm} .
\]

6.23 This has been calculated before in problem (6.20). It is

\[
\omega=\frac{m\left(Z e^{2} / 4 \pi \varepsilon_{0}\right)^{2}}{\hbar^{3} n^{3}}=2.08 \times 10^{16} \mathrm{rad} / \mathrm{sec}
\]

6.24 An electron moving in a circle with a time period \(T\) constitutes a current

\[
I=\frac{e}{T}
\]

and forms a current loop of area \(\pi r^{2}\). This is equivalent to magnetic moment,

\[
\mu=I \pi r^{2}=\frac{e \pi r^{2}}{T}=\frac{e v r}{2}
\]

on using \(\mathrm{v}=2 \pi r / T\). Thus

\[
\mu_{n}=\frac{e m v r}{2 m}=\frac{n e \hbar}{2 m}
\]

for the \(n^{\text {th }}\) orbit. (In Gaussian units

\[
\left.\mu_{n}=n e \hbar / 2 m c\right)
\]

We see that

\[
\mu_{n}=\frac{e}{2 m} M_{n}
\]

where \(M_{n}=n \hbar=m v r\) is the angular momentum

Thus

\[
\begin{gathered}
\frac{\mu_{n}}{M_{n}}=\frac{e}{2 m} \\
\mu_{1}=\frac{e \hbar}{2 m}=\mu_{B}=9.27 \times 10^{-24} \mathrm{~A} \cdot \mathrm{~m}^{2} \\
\text { (In CGS units } \mu_{1}=\mu_{B}=9.27 \times 10^{-21} \mathrm{erg} / \text { gauss) }
\end{gathered}
\]

6.25 The revolving electron is equivalent to a circular current

\[
I=\frac{e}{T}=\frac{e}{2 \pi r / v}=\frac{e v}{2 \pi r}
\]

The magnetic induction

\[
\begin{aligned}
B & =\frac{\mu_{0} I}{2 r}=\frac{\mu_{0} e v}{4 \pi r^{2}}=\frac{\mu_{0}}{4 \pi} \cdot e \cdot \frac{e^{2}}{\left(4 \pi \varepsilon_{0}\right) \hbar} \cdot\left[\frac{m e^{2}}{\hbar^{2}\left(4 \pi \varepsilon_{0}\right)}\right]^{2} \\
& =\frac{\mu_{0} m^{2} e^{7}}{256 \pi^{4} \varepsilon_{0}^{3} \hbar^{5}}
\end{aligned}
\]

Substitution gives \(B=12 \cdot 56 \mathrm{~T}\) at the centre.\\
(In Gaussian units

\[
B=\frac{m^{2} e^{7}}{c \hbar^{5}}=125 \cdot 6 \mathrm{kG}
\]

6.26 From the general formula for the transition \(n_{2} \rightarrow n_{1}\)\\
where

\[
\begin{aligned}
\hbar \omega & =E_{H}\left(\frac{1}{n_{1}^{2}}-\frac{1}{n_{2}^{2}}\right) \\
E_{H} & =13.65 \mathrm{eV} . \text { Then }
\end{aligned}
\]

(1) Lyman, \(n_{1}=1, n_{2}=2,3\). Thus

\[
\hbar \omega \geq \frac{3}{4} E_{H}=10.238 \mathrm{eV}
\]

This corresponds to

\[
\lambda=\frac{2 \pi c \hbar}{\hbar \omega}=0.121 \mu \mathrm{~m}
\]

and Lyman lines have \(\lambda \leq 0.121 \mu \mathrm{~m}\) with the series limit at \(\cdot 0909 \mu \mathrm{~m}\)\\
(2) Balmer : \(n_{2}=2, n_{3}=3,4\),

\[
\hbar \omega \geq E_{H}\left(\frac{1}{4}-\frac{1}{9}\right)=\frac{5}{36} E_{H}=1.876 \mathrm{eV}
\]

This corresponds to

\[
\lambda=0.65 \mu \mathrm{~m}
\]

and Balmer series has \(\lambda \leq 0.65 \mu \mathrm{~m}\) with the series limit at \(\lambda=0.363 \mu \mathrm{~m}\).\\
(3) Paschen : \(n_{2}=3, n_{1}^{i}=4,5, \ldots\)

This corresponds to

\[
h \omega \geq E_{H}\left(\frac{1}{9}-\frac{1}{16}\right)=\frac{7}{144} E_{H}=0.6635 \mathrm{eV}
\]

with the series limit at \(\lambda=0.818 \mu \mathrm{~m}\)\\
\includegraphics[max width=\textwidth, alt={}, center]{da0156b6-4133-434e-90d4-ec01b2070afe-282_234_867_1314_310}\\
6.27 The Balmer line of wavelength 486.1 nm is due to the transition \(4 \rightarrow 2\) while the Balmer line of wavelingth \(410 \cdot 2 \mathrm{~nm}\) is due to the transition \(6 \rightarrow 2\). The line whose wave number corresponds to the difference in wave numbers of these two lines is due to the transition \(6 \rightarrow 4\). That line belongs to the Brackett series. The wavelength of this line is

\[
\lambda=\frac{1}{\frac{1}{\lambda_{2}}-\frac{1}{\lambda_{1}}}=\frac{\lambda_{1} \lambda_{2}}{\lambda_{1}-\lambda_{2}}=2.627 \mu \mathrm{~m}
\]

6.28 The energies are

\[
\begin{gathered}
E_{H}\left(\frac{1}{4}-\frac{1}{9}\right)=\frac{5}{36} E_{H}, E_{H}\left(\frac{1}{4}-\frac{1}{16}\right)=\frac{3}{16} E_{H} \\
E_{H}\left(\frac{1}{4}-\frac{1}{25}\right)=\frac{21}{100} E_{H}
\end{gathered}
\]

They correspond to wavelengths

\[
654 \cdot 2 \mathrm{~nm}, 484 \cdot 6 \mathrm{~nm} \text { and } 433 \mathrm{~nm}
\]

The \(\boldsymbol{n}^{\text {th }}\) line of the Balmer series has the energy

\[
E_{H}\left(\frac{1}{4}-\frac{1}{(n+2)^{2}}\right)
\]

For \(n=19\), we get the wavelength \(366 \cdot 7450 \mathrm{~nm}\)\\
For \(n=20\) we get the wavelength \(366 \cdot 4470 \mathrm{~nm}\)\\
To resolve these lines we require a resolving power of

\[
R \approx \frac{\lambda}{\delta \lambda}=\frac{366.6}{0.298}=1.23 \times 10^{3}
\]

6.29 For the Balmer series

\[
\hbar \omega_{n}=\hbar R\left(\frac{1}{4}-\frac{2}{n^{2}}\right), n \geq 3,
\]

where \(\hbar R=E_{H}=13.65 \mathrm{eV}\). Thus\\
or

\[
\begin{gathered}
\frac{2 \pi \hbar c}{\lambda_{n}}=\hbar R\left(\frac{1}{4}-\frac{1}{n^{2}}\right) \\
\frac{2 \pi \hbar c}{\lambda_{n+1}}-\frac{2 \pi \hbar c}{\lambda_{n}}=\hbar R\left(\frac{1}{n^{2}}-\frac{1}{(n+1)^{2}}\right) \\
=\hbar R\left(\frac{2 n+1}{n^{2}(n+1)^{2}}\right) \approx \frac{2 R}{n^{3}} \text { for } n \gg 1 \\
\frac{2 \pi \hbar c}{\lambda_{n}^{2}} \delta \lambda \approx \frac{2 R \hbar}{n^{3}} \\
\frac{\lambda_{n}}{\delta \lambda} \approx \frac{\pi \hbar c n^{3}}{\lambda_{n} R \hbar}=\frac{\pi c n^{3}}{\lambda_{n} R}
\end{gathered}
\]

or\\
Thus

On the other hand for just resolution in a diffraction grating

Hence

\[
\begin{gathered}
\frac{\lambda}{\delta \lambda}=k N=k \frac{l}{d}=\frac{l}{\lambda d} k \lambda=\frac{l}{\lambda d} d \sin \theta=\frac{l}{\lambda} \sin \theta \\
\sin \theta=\frac{\pi c n^{3}}{l R}
\end{gathered}
\]

Substitution gives \(\theta \approx 59 \cdot 4^{\circ}\).\\
6.30 If all wavelengths are four times shorter but otherwise similar to the hydrogen atom spectrum then the energy levels of the given atom must be four times greater.

This means

\[
E_{n}=-\frac{4 E_{H}}{n^{2}}
\]

compared to \(E_{n}=-\frac{E_{H}}{n^{2}}\) for hydrogen atom. Therefore the spectrum is that of \(\mathrm{He}^{+}\)ion ( \(Z=2\) ).\\
6.31 Because of cascading all possible transitions are seen. Thus we look for the number of ways in which we can select upper and lower levels. The number of ways we can do this is

\[
\frac{1}{2} n(n-1)
\]

where the factor \(\frac{1}{2}\) takes account of the fact that the photon emission always arises from upper → lower transition.\\
6.32 These are the Lyman lines

\[
\hbar \omega=E_{H}\left(\frac{1}{1}-\frac{1}{n^{2}}\right) n=2,3,4, \ldots
\]

For

\[
n=2 \text { we get } \lambda=121 \cdot 1 \mathrm{~nm}
\]

For

\[
n=3 \text { we get } \lambda=102 \cdot 2 \mathrm{~nm}
\]

For

\[
n=4 \text { we get } \lambda=96.9 \mathrm{~nm}
\]

For

\[
n=5 \text { we get } \lambda=94.64 \mathrm{~nm}
\]

For

\[
n=6 \text { we get } \lambda=93.45 \mathrm{~nm}
\]

Thus at the level of accuracy of our calculation, there are four lines

\[
121 \cdot 1 \mathrm{~nm}, 102 \cdot 2 \mathrm{~nm}, 96 \cdot 9 \mathrm{~nm} \text { and } 94 \cdot 64 \mathrm{~nm} .
\]

6.33 If the wavelengths are \(\lambda_{1}, \lambda_{2}\) then the total energy of the excited start must be

\[
E_{n}=E_{1}+\frac{2 \pi c h}{\lambda_{1}}+\frac{2 \pi c h}{\lambda_{2}}
\]

But \(E_{1}=-4 E_{H}\) and \(E_{n}=-\frac{4 E_{H}}{n^{2}}\) where we are ignoring reduced mass effects.\\
Then

\[
4 E_{H}=\frac{4 E_{H}}{n^{2}}+\frac{2 \pi c \hbar}{\lambda_{1}}+\frac{2 \pi c \hbar}{\lambda_{2}}
\]

Substituting the values we get

\[
n^{2}=23
\]

which we take to mean \(n=5\). (The result is sensitive to the values of the various quantities and small differences get multiplied because difference of two large quantities is involved :

\[
n^{2}=\frac{E_{H}}{E_{H}-\frac{\pi c h}{2}\left(\frac{1}{\lambda_{1}}+\frac{1}{\lambda_{2}}\right)} .
\]

6.34 For the longest wavelength (first) line of the Balmer series we have on using the generalized Balmer formula\\
the result\\
the result

Then\\
so

\[
\begin{gathered}
\omega=Z^{2} R\left(\frac{1}{n^{2}}-\frac{1}{m^{2}}\right) \\
\lambda_{1 \text { Lyman }}=\frac{2 \pi c}{Z^{2} R\left(1-\frac{1}{4}\right)}=\frac{8 \pi c}{3 Z^{2} R} \\
\Delta \lambda=\lambda_{1 \text { Balmar }}-\lambda_{1 \text { Lyman }}=\frac{176 \pi c}{15 Z^{2} R} \\
R=\frac{176 \pi c}{15 Z^{2} \Delta \lambda}=2.07 \times 10^{16} \mathrm{sec}^{-1}
\end{gathered}
\]

Then\\
6.35 From the formula of the previous problem\\
or

\[
\begin{aligned}
& \text { problem } \\
& \Delta \lambda=\frac{176 \pi c}{15 Z^{2} R} \\
& Z=\sqrt{\frac{176 \pi c}{15 R \Delta \lambda}}
\end{aligned}
\]

Substitution of \(\Delta \lambda=59.3 \mathrm{~nm}\) and \(R\) and the previous problem gives \(Z=3\)\\
This identifies the ion as \(\boldsymbol{L} \boldsymbol{i}^{++}\)\\
6.36 We start from the generalized Balmer formula

Here

\[
\begin{aligned}
\omega & =R Z^{2}\left(\frac{1}{n^{2}}-\frac{1}{m^{2}}\right) \\
m & =n+1, n+2, \ldots \infty
\end{aligned}
\]

The interval between extreme lines of this series (series \(n\) ) is

\[
\Delta \omega=R Z^{2}\left(\frac{1}{n^{2}}-\frac{1}{(\infty)^{2}}\right)-R Z^{2}\left(\frac{1}{n^{2}}-\frac{1}{(n+1)^{2}}\right)=R Z^{2} /(n+1)^{2}
\]

Hence

\[
n=Z \sqrt{\frac{R}{\Delta \omega}}-1
\]

Then the angular frequency of the first line of this series (series \(n\) ) is

\[
\begin{aligned}
\omega_{1} & =R Z^{2}\left(\frac{1}{n^{2}}-\frac{1}{(n+1)^{2}}\right)=\Delta \omega\left(\left(\frac{n+1}{n}\right)^{2}-1\right) \\
& \left.=\Delta \omega\left[\left\{\frac{z \sqrt{\frac{R}{\Delta \omega}}}{z \sqrt{\frac{R}{\Delta \omega}}-1}\right\}^{2}\right]-1\right]=\Delta \omega \frac{2 z \sqrt{\frac{R}{\Delta \omega}}-1}{\left(z \sqrt{\frac{R}{\Delta \omega}}-1\right)^{2}}
\end{aligned}
\]

Then the wavelength will be

\[
\lambda_{1}=\frac{2 \pi c}{\omega_{1}}=\frac{2 \pi c}{\Delta \omega} \frac{\left(Z \sqrt{\frac{R}{\Delta \omega}}-1\right)^{2}}{2 Z \sqrt{\frac{R}{\Delta \omega}}-1}
\]

Substitution (with the value of \(R\) from problem 6.34 which is also the correct value determined directly) gives

\[
\lambda_{1}=0.468 \mu \mathrm{~m}
\]

6.37 For the third line of of Balmer series

Hence

\[
\begin{gathered}
\omega=R Z^{2}\left(\frac{1}{2^{2}}-\frac{1}{5^{2}}\right)=\frac{21}{100} R Z^{2} \\
\lambda=\frac{2 \pi c}{\omega}=\frac{200 \pi c}{21 R Z^{2}} \\
Z=\sqrt{\frac{200 \pi c}{21 R \lambda}}
\end{gathered}
\]

or

\[
E_{b}=4 E_{H}=4 \times 13.65=54.6 \mathrm{eV}
\]

The ion is \(\mathrm{He}^{+}\).\\
6.38 To remove one electron requires 24.6 eV .

The ion that is left is \(\mathrm{He}^{+}\)which in its ground start has a binding energy of \(4 E_{H}=4 \hbar R\).\\
The complete binding energy of both electrons is then

Substitution gives

\[
\begin{aligned}
E & =E_{0}+4 \hbar R \\
E & =79 \cdot 1 \mathrm{eV}
\end{aligned}
\]

6.39 By conservation of energy

\[
\frac{1}{2} m v^{2}=\frac{2 \pi \hbar c}{\lambda}-E_{b}
\]

where \(E_{b}=4 \hbar R\) is the binding energy of the electron in the ground state of \(\mathrm{He}^{+}\). (Recoil of \(\mathrm{He}^{++}\)nucleus is neglected). Then

\[
v=\sqrt{\frac{2}{m}\left(\frac{2 \pi h c}{\lambda}-E_{b}\right)}
\]

Substitution gives

\[
v=2.25 \times 10^{6} \mathrm{~m} / \mathrm{s}
\]

6.40 Photon can be emitted in \(H-H\) collision only if one of the \(H\) is excited to an \(n=2\) state which then dexcites to \(n=1\) state by emitting a photon. Let \(v_{1}\) and \(v_{2}\) be the velocities of the two Hydrogen atoms after the collision and \(M\) their masses. Then, energy momentum conservation

\[
M v_{1}+M v_{2}=\sqrt{2 M T}
\]

(in the frame of the stationary \(H\) atom)

\[
\frac{1}{2} M v_{1}^{2}+\frac{1}{2} M v_{2}^{2}+\frac{3}{4} \hbar R=T
\]

\(\frac{3}{4} \hbar R=\hbar R\left(1-\frac{1}{4}\right)\) is the excitation energy of the \(n=2\) state from the ground state.

Eliminating \(\mathbf{v}_{\mathbf{2}}\)\\
or

\[
\begin{gathered}
\frac{1}{2} M\left\{\mathrm{v}_{1}^{2}+\left(\sqrt{\frac{2 T}{M}}-\mathrm{v}_{1}\right)^{2}\right\}+\frac{3}{4} \hbar R=T \\
\frac{1}{2} M\left\{2 \mathrm{v}_{1}^{2}-2 \sqrt{\frac{2 T}{M}} \mathrm{v}_{1}+\frac{2 T}{M}\right\}+\frac{3}{4} \hbar R=T \\
M\left\{\left(\mathrm{v}_{1}-\frac{1}{2} \sqrt{\frac{2 T}{M}}\right)^{2}\right\}+\frac{1}{2} T+\frac{3}{4} \hbar R=T \\
M\left\{\left(\mathrm{v}_{1}-\frac{1}{2} \sqrt{\frac{2 T}{M}}\right)^{2}\right\}+\frac{3}{4} \hbar R=\frac{1}{2} T
\end{gathered}
\]

For minimum \(T\), the square on the left should vanish. Thus \(T=\frac{3}{2} \hbar R=20.4 \mathrm{eV}\)\\
6.41 In the rest frame of the original excited nucleus we have the equations \(O=\overrightarrow{p_{\gamma}}+\overrightarrow{p_{H}}\)

\[
\frac{3}{4} \hbar R=c\left|\vec{p}_{y}\right|+p_{H}^{2} / 2 M
\]

( \(\frac{3}{4} \hbar R\) is the energy available in \(n=2 \rightarrow=1\) transition corresponding to the first Lyman line.)\\
Then

\[
p_{H}^{2}+2 M c p_{H}-\frac{3 \hbar R M}{2}=0
\]

or

\[
\begin{gathered}
\left(p_{H}+M c\right)^{2}=M^{2} c^{2}+\frac{3}{2} \hbar R M \\
p_{H}=-M c+\sqrt{M^{2} c^{2}+\frac{3}{2} \hbar R M}=-M c+M c\left(1+\frac{3 \hbar R}{2 M c^{2}}\right)^{1 / 2} \cong \frac{3 \hbar R}{4 c}
\end{gathered}
\]

(We could have written this directly by noting that \(p_{H}^{2} / 2 M \ll c p_{\gamma}\).) Then

\[
\mathrm{v}_{\mathrm{H}}=\frac{3 \hbar R}{4 M c}=3.3 \mathrm{~m} / \mathrm{s}
\]

6.42 We have

Then

\[
\begin{aligned}
\varepsilon & =\frac{3}{4} \hbar R \quad \text { and } \quad \varepsilon^{\prime} \propto \frac{3}{4} \hbar R-\frac{1}{2 M}\left(\frac{3}{4} \hbar R / c\right)^{2} \\
\frac{\varepsilon-\varepsilon^{\prime}}{\varepsilon} & =\frac{3 \hbar R}{8 M c^{2}}=\frac{v_{H}}{2 c}=5.5 \times 10^{-9}=0.55 \times 10^{-6} \%
\end{aligned}
\]

6.43 We neglect recoil effects. The energy of the first Lyman line photon emitted by \(\mathrm{He}^{+}\)is

\[
4 \hbar R\left(1-\frac{1}{4}\right)=3 \hbar R
\]

The velocity \(v\) of the photoelectron that this photon liberates is given by

\[
3 \hbar R=\frac{1}{2} m \mathrm{v}^{2}+\hbar R
\]

where \(h R\) on the right is the binding energy of the \(n=1\) electron in H atom. Thus

\[
\mathrm{v}=\sqrt{\frac{4 \hbar R}{m}}=2 \sqrt{\frac{\hbar R}{m}}=3.1 \times 10^{6} \mathrm{~m} / \mathrm{s}
\]

Here \(m\) is the mass of the electron.\\
6.44 Since \(\Delta \lambda(=0.20 \mathrm{~nm}) \ll \lambda(=121 \mathrm{~nm})\) of the first Lyman line of H atom, we need not worry about \(\mathrm{v}^{2} / c^{2}\) effects. Then

\[
\omega^{\prime}=\frac{\omega}{1-\beta \cos \theta}, \beta=\frac{v}{c}
\]

Hence

\[
1-\beta \cos \theta=\frac{\omega}{\omega^{\prime}}=\frac{\lambda^{\prime}}{\lambda}
\]

or

\[
\beta \cos \theta=1-\frac{\lambda^{\prime}}{\lambda}=\frac{\Delta \lambda}{\lambda}
\]

But

\[
\omega=\frac{3}{4} R \text { so } \lambda=\frac{2 \pi c}{\omega}=\frac{8 \pi c}{3 R}
\]

Hence

\[
v=c \beta=\frac{3 R \Delta \lambda}{8 \pi \cos \theta}
\]

Substitution gives \(\left(\cos \theta=\frac{1}{\sqrt{2}}\right)\)

\[
\mathrm{v}=7.0 \times 10^{5} \mathrm{~m} / \mathrm{s}
\]

6.45 (a) If we measure energy from the bottom of the well, then \(V(x)=0\) inside the walls. Then the quantization conditon reads \(\oint p d x=2 l p=2 \pi n \hbar\)\\
or \(p=\pi \boldsymbol{n} \boldsymbol{h} / \boldsymbol{l}\)\\
Hence

\[
E_{n}=\frac{p^{2}}{2 m}=\frac{\pi^{2} n^{2} \hbar}{2 m l^{2}}
\]

\(\oint p d x=2 l p\) because we have to consider the integral form \(-\frac{l}{2}\) to \(\frac{l}{2}\) and then back to \(-\frac{l}{2}\).)\\
(b) Here \(\oint p d x=2 \pi r p=2 \pi n \hbar\)\\
or

\[
p=\frac{n \hbar}{r}
\]

Hence

\[
E_{n}=\frac{n^{2} \hbar^{2}}{2 m r^{2}}
\]

(c) By energy conservation \(\quad \frac{p^{2}}{2 m}+\frac{1}{2} \alpha x^{2}=E\)\\
so

\[
p=\sqrt{2 m E-m \alpha x^{2}}
\]

Then \(\oint p d x=\oint \sqrt{2 m E-m \alpha x^{2}} d x\)

\[
\begin{gathered}
\sqrt{\frac{2 E}{\alpha}} \\
=2 \sqrt{m \alpha} \int^{\frac{2 E}{\alpha}-x^{2}} d x \\
-\sqrt{\frac{2 E}{\alpha}}
\end{gathered}
\]

The integral is \(\quad \int_{-a}^{a} \sqrt{a^{2}-x^{2}} d x=a^{2} \int_{-\pi / 2}^{\pi / 2} \cos ^{2} \theta d \theta\)

\[
=\frac{a^{2}}{2} \int_{-\pi / 2}^{\pi / 2}(1+\cos 2 \theta) d \theta=a^{2} \frac{\pi}{2}
\]

Thus \(\quad \oint p d x=\pi \sqrt{m \alpha} \cdot \frac{2 E}{\alpha}=E \cdot 2 \pi \cdot \sqrt{\frac{m}{\alpha}}=2 \pi n \hbar\)

Hence

\[
E_{n}=n \hbar \sqrt{\frac{\alpha}{m}}
\]

(b) It is requined to find the energy levels of the circular orbait for the rotential \(U\left(r_{i}\right)=-\frac{\alpha}{r}\)

In a circular orbit, the particle only has tangential velocity and the quantization condition reads \(\oint p d x=m v \cdot 2 \pi r=2 \pi n \hbar\).\\
so

\[
m \vee r=M=n \hbar
\]

The energy of the particle is

\[
E=\frac{n^{2} \hbar^{2}}{2 m r^{2}}-\frac{\alpha}{r}
\]

Equilibrium requires that the energy as a function of \(r\) be minimum. Thus

Hence

\[
\begin{gathered}
\frac{n^{2} \hbar^{2}}{m r^{3}}=\frac{\alpha}{r^{2}} \text { or } r=\frac{n^{2} \hbar^{2}}{m \alpha} \\
E_{n}=-\frac{m \alpha^{2}}{2 n^{2} h^{2}}
\end{gathered}
\]

6.46 The total energy of the H -atom in an arbitrary frame is

\[
E=\frac{1}{2} m \overrightarrow{V_{1}}+\frac{1}{2} M \overrightarrow{V_{2}}-\frac{e^{2}}{\left(4 \pi \varepsilon_{0}\right)\left|r_{1}-\overrightarrow{r_{2}}\right|}
\]

Here \(\overrightarrow{V_{1}}=\dot{\overrightarrow{r_{1}}}, \overrightarrow{V_{2}}=\dot{\overrightarrow{r_{2}}}, \vec{r}\) \& \(\overrightarrow{r_{2}}\) are the coordinates of the electron and protons.\\
We define

Then

\[
\begin{aligned}
\vec{R} & =\frac{m \overrightarrow{r_{1}}+M \overrightarrow{r_{2}}}{M+m} \\
\vec{r} & =\overrightarrow{r_{1}}-\overrightarrow{r_{2}} \\
\vec{V} & =\frac{m \overrightarrow{V_{1}}+M \overrightarrow{V_{2}}}{m+M} \\
\overrightarrow{\mathrm{v}} & =\vec{V}_{1}-\vec{V}_{2}
\end{aligned}
\]

or

\[
\begin{aligned}
& \vec{V}_{1}=\vec{V}+\frac{M}{m+M} \overrightarrow{\mathrm{v}} \\
& \vec{V}_{2}=\vec{V}-\frac{m}{m+M} \overrightarrow{\mathrm{v}}
\end{aligned}
\]

and we get

\[
E=\frac{1}{2}(m+M) \overrightarrow{V^{2}}+\frac{1}{2} \frac{m M}{m+M} v^{2}-\frac{e^{2}}{4 \pi \varepsilon_{0} r}
\]

In the frame \(\vec{V}=0\), this reduces to the energy of a particle of mass

\[
\mu=\frac{m M}{m+M}
\]

\(\mu\) is called the reduced mass.\\
Then

\[
E_{b}=\frac{\mu e^{4}}{2 \hbar^{2}} \quad \text { and } \quad R=\frac{\mu e^{4}}{2 \hbar^{3}}
\]

Since

\[
\mu=\frac{m}{1+\frac{m}{M}} \approx m\left(1-\frac{m}{M}\right)
\]

these values differ by \(\frac{m}{M}(\approx 0.54 \%)\) from the values obtained without considering nuclear motion. ( \(M=1837 \mathrm{~m}\) )\\
6.47 The difference between the binding energies is

\[
\begin{aligned}
\Delta E_{b} & =E_{b}(D)-E_{b}(H) \\
& =\Delta \frac{m}{1+\frac{m}{M}} \frac{e^{4}}{2 \hbar^{2}}+\frac{m}{1+\frac{m}{2 M}} \frac{e^{4}}{2 \hbar^{2}} \\
& =\frac{m e^{4}}{2 \hbar^{2}}\left(\frac{m}{2 M}\right)
\end{aligned}
\]

Substitution gives \(\Delta E_{b}=3.7 \mathrm{meV}\).\\
For the first line of the Lyman series\\
or

\[
\begin{aligned}
\frac{2 \pi \hbar c}{\lambda} & =\hbar R\left(\frac{1}{1}-\frac{1}{4}\right)=\frac{3}{4} \hbar R \\
\lambda & =\frac{8 \pi c}{3 R}=\frac{8 \pi \hbar c}{3 E_{b}}
\end{aligned}
\]

Hence

\[
\begin{aligned}
\lambda_{H}-\lambda_{D} & =\frac{8 \pi \hbar c}{3}\left(\frac{1}{E_{b}(H)}-\frac{1}{E_{b}(D)}\right) \\
& =\frac{8 \pi \hbar c}{3} \cdot\left(\frac{m e^{4}}{2 \hbar^{2}}\right)^{-1}\left(1+\frac{m}{M}-1-\frac{m}{2 M}\right) \\
& =\frac{8 \pi \hbar c}{3\left(\frac{m e^{4}}{2 \hbar^{2}}\right)} \cdot \frac{m}{2 M} \\
& =\frac{m}{2 M} \times \lambda_{1}
\end{aligned}
\]

(where \(\lambda_{1}\) is the wavelength of the first line of Lyman series without considering nuclear motion).\\
Substitution gives (see 6.21 for \(\lambda_{1}\) ) using \(\lambda_{1}=121 \mathrm{~nm}\)

\[
\Delta \lambda=33 p m
\]

6.48 (a) In the mesonic system, the reduced mass of the system is related to the masses of the meson ( \(m_{\mu}\) ) and proton ( \(m_{p}\) ) by

\[
\mu=\frac{m_{\mu} m_{p}}{m_{\mu}+M_{p}}=186.04 m_{e}
\]

Then,

\[
\begin{aligned}
& \text { separation between the particles in the ground state }=\frac{\hbar^{2}}{\mu e^{2}} \\
&=\frac{1}{186} \frac{\hbar^{2}}{m e^{2}} \\
&=0.284 \mathrm{pm} \\
& E_{b}(\text { meson })=\frac{\mu e^{4}}{2 \hbar^{2}}=186 \times 13.65 \mathrm{eV} \\
&=2.54 \mathrm{keV} \\
& \lambda_{1}=\frac{8 \pi \hbar c}{3 E_{b}(\text { meson })}=\frac{\lambda_{1}(\text { Hydrogen })}{186}=0.65 \mathrm{~nm} \\
&\text { (on using } \left.\lambda_{1}(\text { Hydrogen })=121 \mathrm{~nm}\right) .
\end{aligned}
\]

(b) In the positronium

\[
\mu=\frac{m_{e}^{2}}{2 m_{e}}=\frac{m_{e}}{2}
\]

Thus separation between the particles is the ground state

\[
\begin{gathered}
=2 \frac{\hbar^{2}}{m_{e} e^{2}}=105.8 \mathrm{pm} \\
E_{b}(\text { positronium })=\frac{m_{e}}{2} \cdot \frac{e^{4}}{2 \hbar^{2}}=\frac{1}{2} E_{b}(H)=6.8 \mathrm{eV} \\
\lambda_{1}(\text { positronium })=2 \lambda_{1}(\text { Hydrogen })=0.243 \mathrm{~nm}
\end{gathered}
\]

\subsection*{6.2 WAVE PROPERTIES OF PARTICLES. SCHRODINGER EQUATION}
6.49 The kinetic energy is nonrelativistic in all three cases. Now

\[
\lambda=\frac{2 \pi \hbar}{p}=\frac{2 \pi \hbar}{\sqrt{2 m T}}
\]

using

\[
\begin{gathered}
T=1.602 \times 10^{-17} \text { Joules, we get } \\
\lambda_{e}=122.6 \mathrm{pm} \\
\lambda_{P}=2.86 \mathrm{pm} \\
\lambda_{U}=\frac{\lambda_{P}}{\sqrt{238}}=0.185 \mathrm{pm} .
\end{gathered}
\]

(where we have used a mass number of 238 for the \(\boldsymbol{U}\) nucleus).\\
6.50 From \(\lambda=\frac{2 \pi \hbar}{p}=\frac{2 \pi \hbar}{\sqrt{2 m T}}\)\\
we find

\[
T=\frac{4 \pi^{2} \hbar^{2}}{2 m \lambda^{2}}=\frac{2 \pi^{2} \hbar^{2}}{m \lambda^{2}}
\]

Thus

\[
T_{2}-T_{1}=\frac{2 \pi^{2} \hbar^{2}}{m}\left(\frac{1}{\lambda_{2}^{2}}-\frac{1}{\lambda_{1}^{2}}\right)
\]

Substitution gives \(\Delta T=451 \mathrm{eV}=0.451 \mathrm{keV}\).\\
6.51 We shall use \(M_{0} \approx 2 M_{n}\). The CM is moving with velocity

\[
V=\frac{\sqrt{2 M_{n} T}}{3 M_{n}}=\sqrt{\frac{2 T}{9 M_{n}}}
\]

with respect to the Lab frame. In the CM frame the velocity of neutron is\\
and

\[
\begin{gathered}
\mathrm{v}_{n}^{\prime}=\mathrm{v}_{\mathrm{n}}-V=\sqrt{\frac{2 T}{M_{n}}}-\sqrt{\frac{2 T}{9 M_{n}}}=\sqrt{\frac{2 T}{M_{n}}} \cdot \frac{2}{3} \\
\lambda_{n}^{\prime}=\frac{2 \pi \hbar}{M_{n} \mathrm{v}_{n}^{\prime}}=\frac{3 \pi \hbar}{\sqrt{2 M_{n} T}}
\end{gathered}
\]

Substitution gives \(\lambda_{n}^{\prime}=8.6 \mathrm{pm}\)\\
Since the momenta are equal in the CM frame the de Broglie wavelengths will also be equal.\\
If we do not assume \(M_{d} \approx 2 M_{n}\) we shall get

\[
\lambda_{n}^{\prime}=\frac{2 \pi \hbar\left(1+M_{n} / M_{d}\right)}{\sqrt{2 M_{n} T}}
\]

6.52 If \(\overrightarrow{p_{1}}, \overrightarrow{p_{2}}\) are the momenta of the two particles then their momenta in the CM frame will be \(\pm\left(\overrightarrow{p_{1}}-\overrightarrow{p_{2}}\right) / 2\) as the particle are identical.\\
Hence their de Broglie wavelength will be

\[
\begin{aligned}
\tilde{\lambda} & =\frac{2 \pi \hbar}{\frac{1}{2}\left|\overrightarrow{p_{1}}-\overrightarrow{p_{2}}\right|}=\frac{4 \pi \hbar}{\sqrt{p_{1}^{2}+p_{2}^{2}}}\left(\text { because } \overrightarrow{p_{1}} \perp \overrightarrow{p_{2}}\right) \\
& =\frac{2}{\sqrt{\frac{1}{\lambda_{1}^{2}}+\frac{1}{\lambda_{2}^{2}}}}=\frac{2 \lambda_{1} \lambda_{2}}{\sqrt{\lambda_{1}^{2}+\lambda_{2}^{2}}}
\end{aligned}
\]

6.53 In thermodynamic equilibrium, Maxwell's velocity distribution law holds :

\[
d N(\mathrm{v})=\Phi(\mathrm{v}) d \mathrm{v}=A \mathrm{v}^{2} e^{-\mathrm{m} \mathrm{v}^{2} / 2 \mathrm{kT}} d \mathrm{v}
\]

\(\Phi(v)\) is maximum when

\[
\Phi^{\prime}(v)=\Phi(v)\left[\frac{2}{v}-\frac{m v}{k T}\right]=0
\]

The difines the most probable velocity.

\[
\mathrm{v}_{\mathrm{pr}}=\sqrt{\frac{2 k T}{m}}
\]

The de Broglie wavelength of H molecules with the most probable velocity is

\[
\lambda=\frac{2 \pi \hbar}{m v_{p r}}=\frac{2 \pi \hbar}{\sqrt{2 m k T}}
\]

Substituting the appropriate value especially

\[
\begin{gathered}
m=m_{H_{2}}=2 m_{H}, T=300 \mathrm{~K}, \text { we get } \\
\lambda=126 \mathrm{pm}
\end{gathered}
\]

6.54 To find the most probable de Broglie wavelength of a gas in thermodynamic equilibrium we determine the distribution is \(\lambda\) corresponding to Maxwellian velocity distribution.\\
It is given by

\[
\psi(\lambda) d \lambda=-\Phi(\mathrm{v}) d \mathrm{v}
\]

(where - sign takes account of the fact that \(\lambda\) decreaes as \(v\) increases). Now

\[
\begin{aligned}
\lambda & =\frac{2 \pi \hbar}{m \mathrm{v}} \text { or } \mathrm{v}=\frac{2 \pi \hbar}{m \lambda} \\
d \mathrm{v} & =-\frac{2 \pi \hbar}{m \lambda^{2}} d \lambda
\end{aligned}
\]

Thus

\[
\Psi(\lambda)=+A \mathrm{v}^{2} e^{-m \mathrm{v}^{2} / 2 k T}\left(-\frac{d \mathrm{v}}{d \lambda}\right)
\]

\[
\begin{aligned}
& =A\left(\frac{2 \pi \hbar}{m \lambda}\right)^{2}\left(\frac{2 \pi \hbar}{m \lambda^{2}}\right) e^{-\frac{m}{2 k T} \cdot\left(\frac{2 \pi h}{m \lambda}\right)^{2}} \\
& =\text { Const } \cdot \lambda^{-4} e^{-a / \lambda^{2}} \\
& a=\frac{2 \pi^{2} \hbar^{2}}{m k T}
\end{aligned}
\]

where\\
This is maximum when\\
or

\[
\begin{aligned}
\psi^{\prime}(\lambda) & =0=\psi(\lambda)\left[\frac{-4}{\lambda}+\frac{2 a}{\lambda^{3}}\right] \\
\lambda_{p r} & =\sqrt{a / 2}=\pi \hbar / \sqrt{m k T}
\end{aligned}
\]

Using the result of the previous problem it is

\[
\lambda_{p r}=\frac{126}{\sqrt{2}} \mathrm{pm}=89 \cdot 1 \mathrm{pm}
\]

6.55 For a relativistic particle

\[
T+m c^{2}=\text { total energy }=\sqrt{c^{2} p^{2}+m^{2} c^{4}}
\]

Squaring

Hence

\[
\begin{aligned}
\sqrt{T\left(T+2 m c^{2}\right)}=c p & \\
\lambda & =\frac{2 \pi \hbar c}{\sqrt{T\left(T+2 m c^{2}\right)}} \\
& =\frac{2 \pi \hbar}{\sqrt{2 m T\left(1+\frac{T}{2 m c^{2}}\right)}}
\end{aligned}
\]

If we use nonrelativistic formula,

\[
\begin{gathered}
\lambda_{N R}=\frac{2 \pi \hbar}{\sqrt{2 m T}} \\
\frac{\Delta \lambda}{\lambda}=\frac{\lambda_{N R}-\lambda}{\lambda_{N R}} \equiv \frac{T}{4 m c^{2}} \\
\left(\text { If } T / 2 m c^{2} \ll 1, \text { we can write }\left(1+\frac{T}{2 m c^{2}}\right)^{-1 / 2} \cong 1-\frac{T}{4 m c^{2}}\right)
\end{gathered}
\]

so

Thus \(T \leq \frac{4 m c^{2} \Delta \lambda}{\lambda}\) if the error is less than \(\Delta \lambda\)\\
For electron the error is not more than \(1 \%\) if

\[
\begin{aligned}
T \leq & 4 \times 0.511 \times .01 \mathrm{MeV} \\
& \leq 20.4 \mathrm{keV}
\end{aligned}
\]

For a proton, the error is not more than \(1 \%\) if\\
i.e.

\[
\begin{gathered}
T \leq 4 \times 938 \times 0.01 \mathrm{MeV} \\
T \leq 37.5 \mathrm{MeV}
\end{gathered}
\]

6.56 The de Broglie wavelength is

\[
\lambda_{d B}=\frac{\frac{2 \pi \hbar}{m_{0} v}}{\sqrt{1-v^{2} / c^{2}}}=\frac{2 \pi \hbar}{m_{0} v} \sqrt{1-v^{2} / c^{2}}
\]

and the Compton wavelength is

\[
\lambda_{c}=\frac{2 \pi \hbar}{m_{0} c}
\]

The two are equal if \(\beta=\sqrt{1-\beta^{2}}\), where \(\beta=\frac{v}{c}\)\\
or

\[
\beta=\frac{1}{\sqrt{2}}
\]

The corresponding kinetic energy is

\[
T=\frac{m_{0} c^{2}}{\sqrt{1-\beta^{2}}}-m_{0} c^{2}=(\sqrt{2}-1) m_{0} c^{2}
\]

Here \(m_{0}\) is th rest mass of the particle (here an electron).\\
6.57 For relativistic electrons, the formula for the short wavelength limit of \(X\)-rays will be\\
or

\[
\frac{2 \pi \hbar c}{\lambda_{s h}}=m_{0} c^{2}\left(\frac{1}{\sqrt{1-\beta^{2}}}-1\right)=c \sqrt{p^{2}+m^{2} c^{2}}-m c^{2}
\]

or

\[
\left(\frac{2 \pi \hbar}{\lambda_{s h}}+m c\right)^{2}=p^{2}+m^{2} c^{2}
\]

\[
\left(\frac{2 \pi \hbar}{\lambda_{s h}}\right)\left(\frac{2 \pi \hbar}{\lambda_{s h}}+2 m c\right)=p^{2}
\]

or

\[
p=\frac{2 \pi \hbar}{\lambda_{s h}} \sqrt{1+\frac{m c \lambda_{s h}}{\pi \hbar}}
\]

Hence

\[
\lambda_{d B}=\lambda_{s h} / \sqrt{1+\frac{m c \lambda_{s h}}{\pi \hbar}}=3.29 \mathrm{pm}
\]

6.58 The first minimum in a Fraunhofer diffraction is given by ( \(b\) is the width of the slit) \(b \sin \theta=\lambda\)

Here

\[
\begin{gathered}
\sin \theta=\frac{\Delta x / 2}{\sqrt{l^{2}+\left(\frac{\Delta x}{2}\right)^{2}}} \approx \frac{\Delta x}{2 l} \\
\quad \lambda=\frac{b \Delta x}{2 l}=\frac{2 \pi \hbar}{m \mathrm{v}} \\
v=\frac{4 \pi \hbar l}{m b \Delta x}=2.02 \times 10^{6} \mathrm{~m} / \mathrm{s}
\end{gathered}
\]

Thus\\
so\\
6.59 From the Young slit formula

\[
\Delta x=\frac{l \lambda}{d}=\frac{l}{d} \cdot \frac{2 \pi \hbar}{\sqrt{2 m e V}}
\]

Substitution gives

\[
\Delta x=4.90 \mu \mathrm{~m}
\]

6.60 From Bragg's law, for the first case

\[
2 d \sin \theta=n_{0} \lambda=n_{0} \frac{2 \pi \hbar}{\sqrt{2 m e V_{0}}}
\]

where \(n_{0}\) is an unknown integer. For the next higher voltage

\[
2 d \sin \theta=\left(n_{0}+1\right) \frac{2 \pi \hbar}{\sqrt{2 m e \eta V_{0}}}
\]

Thus

\[
n_{0}=\frac{n_{0}+1}{\sqrt{\eta}}
\]

or

\[
n_{0}\left(1-\frac{1}{\sqrt{\eta}}\right)=\frac{1}{\sqrt{\eta}} \quad \text { or } \quad n_{0}=\frac{1}{\sqrt{\eta}-1}
\]

Going back we get

\[
V_{0}=\frac{\pi^{2} \hbar^{2}}{2 m e d^{2} \sin ^{2} \theta} \frac{1}{(\sqrt{\eta}-1)^{2}}=0.150 \mathrm{keV}
\]

Note :- In the Bragg's formula, \(\theta\) is the glancing angle and not the angle of incidence. We have obtained correct result by taking \(\theta\) to be the glancing angle. If \(\theta\) is the angle of incidence, then the glancing angle will be \(90-\theta\). Then the final answer will be smaller by a factor \(\tan ^{2} \theta=\frac{1}{3}\).\\
6.61 Path difference is\\
\(d+d \cos \theta=2 d \cos ^{2} \frac{\theta}{2}\).\\
Thus for reflection maximum of the \(\boldsymbol{k}^{\text {th }}\) order

\[
2 d \cos ^{2} \frac{\theta}{2}=k \lambda=k \frac{2 \pi \hbar}{\sqrt{2 m T}}
\]

Hence \(\quad d=\frac{k \pi \hbar}{\sqrt{2 m T}} \operatorname{Sec}^{2} \frac{\theta}{2}\).\\
Substitution with \(k=4\) gives\\
\includegraphics[max width=\textwidth, alt={}, center]{da0156b6-4133-434e-90d4-ec01b2070afe-298_494_483_73_828}\\
6.62 See the analogous problem with \(X\) - rays (5.156)

The glancing angle is obtained from

\[
\tan 2 \theta=\frac{D}{2 l}
\]

where \(D=\) diameter of the ring, \(l=\) distance from the foil to the screen.\\
Then for the third order Bragg reflection

\[
2 d \sin \theta=k \lambda=k \frac{2 \pi \hbar}{\sqrt{2 m T}},(k=3)
\]

Thus

\[
d=\frac{\pi \hbar k}{\sqrt{2 m T} \sin \theta}=0.232 \mathrm{~nm}
\]

6.63 Inside the metal, there is a negative potential energy of \(-e V_{i}\). (This potential energy prevents electrons from leaking out and can be measured in photoelectric effect etc.) An electron whose K.E. is eV outside the metal will find its K.E. increased to \(e\left(V+V_{i}\right)\) in the metal. Then\\
(a) de Broglie wavelength in the metal

\[
=\lambda_{m}=\frac{2 \pi \hbar}{\sqrt{2 m e\left(V+V_{i}\right)}}
\]

Also de Broglie wavelength in vacuum

\[
=\lambda_{0}=\frac{2 \pi \hbar}{\sqrt{2 m V e}}
\]

Hence refractive index \(n=\frac{\lambda_{0}}{\lambda_{m}}=\sqrt{1+\frac{V_{i}}{V}}\)\\
Substituting we get

\[
n=\sqrt{1+\frac{1}{10}} \simeq 1.05
\]

(b) \(n-1=\sqrt{1+\frac{V_{i}}{V}}-1 \leq \eta\)\\
then

\[
1+\frac{V_{i}}{V} \leq(1+\eta)^{2}
\]

or

\[
V_{i} \leq \eta(2+\eta) V
\]

or

\[
\frac{V}{V_{i}} \geq \frac{1}{\eta(2+\eta)}
\]

For

\[
\begin{aligned}
\eta=1 \% & =0 \\
\frac{V}{V_{i}} & \geq 50
\end{aligned}
\]

we get\\
6.64 The energy inside the well is all kinetic if energy is measured from the value inside. We require\\
or

\[
\begin{aligned}
l & =n \lambda / 2=n \frac{\pi \hbar}{\sqrt{2 m E}} \\
E_{n} & =\frac{n^{2} \pi^{2} \hbar^{2}}{2 m l^{2}}, n=1,2, \ldots
\end{aligned}
\]

6.65 The Bohr condition

\[
\oint p d x=\oint \frac{2 \pi \hbar}{\lambda} d x=2 \pi n \hbar
\]

For the case when \(\lambda\) is constant (for example in circular orbits) this means

\[
2 \pi r=n \lambda
\]

Here \(r\) is the radius of the circular orbit.\\
6.66 From the uncertainty principle (Eqn. (6.2b))

\[
\Delta x \Delta p_{x} \geq \hbar
\]

Thus

\[
\Delta p_{x}=m \Delta \mathrm{v}_{\mathrm{x}} \geq \frac{\hbar}{\Delta x}
\]

or

\[
\Delta v_{x} \geq \frac{\hbar}{m \Delta x}
\]

For an electron this means an uncertainty in velocity of \(116 \mathrm{~m} / \mathrm{s}\) if \(\Delta x=10^{-6} \mathrm{~m}=1 \mu \mathrm{~m}\) For a proton

\[
\Delta v_{x}=6.3 \mathrm{~cm} / \mathrm{s}
\]

For a ball

\[
\Delta v_{x}=1 \times 10^{-20} \mathrm{~cm} / \mathrm{s}
\]

6.67 As in the previous problem

\[
\Delta \mathrm{v} \sim \frac{\hbar}{m l}=1.16 \times 10^{6} \mathrm{~m} / \mathrm{s}
\]

The actual velocity \(v_{1}\) has been calculated in problem 6.21. It is

\[
v_{1}=2.21 \times 10^{6} \mathrm{~m} / \mathrm{s}
\]

Thus \(\Delta v v_{1}\) (They are of the same order of magnitude)\\
6.68 If \(\Delta x=\lambda / 2 \pi=\frac{2 \pi \hbar}{p} \cdot \frac{1}{2 \pi}=\frac{\hbar}{p}=\frac{\hbar}{m v}\)

Thus

\[
\Delta v \geq \frac{\hbar}{m \Delta x}=v
\]

Thus \(\Delta v\) is of the same order as \(v\).\\
6.69 Initial uncertainty \(\Delta v>\frac{\hbar}{m l}\). With this incertainty the wave train will spread out to a distance \(\eta l\) long in time

\[
t_{0} \approx \eta l / \frac{\hbar}{m l} \approx \frac{\eta m l^{2}}{\hbar} \mathrm{sec} .=8.6 \times 10^{16} \mathrm{sec} . \sim 10^{-15} \mathrm{sec}
\]

6.70 Clearly \(\Delta x \leq l\) so \(\Delta p_{x} ; \geq \frac{\hbar}{l}\)

Now \(p_{x} \geq \Delta p_{x}\) and so

Thus

\[
\begin{aligned}
T & =\frac{p_{x}^{2}}{2 m} \geq \frac{\hbar^{2}}{2 m l^{2}} \\
T_{\min } & =\frac{\hbar^{2}}{2 m l^{2}} \simeq 0.95 \mathrm{eV}
\end{aligned}
\]

6.71 The momentum the electron is \(\Delta p_{x}=\sqrt{2 m T}\)

Uncertainty in its momentum is

\[
\Delta p_{x} \geq \hbar / \Delta x=h / l
\]

Hence relative uncertainty

\[
\frac{\Delta p_{x}}{p_{x}}=\frac{\hbar}{l \sqrt{2 m T}}=\sqrt{\frac{\hbar^{2}}{2 m l^{2}} / T}=\frac{\Delta \mathrm{v}}{\mathrm{v}}
\]

Substitution gives

\[
\frac{\Delta v}{v}=\frac{\Delta p}{p}=9.75 \times 10^{-5} \approx 10^{-4}
\]

6.72 By uncertainty principle, the uncertainty in momentum

\[
\Delta p \geqslant \frac{\hbar}{l}
\]

For the ground state, we expect \(\Delta p_{\sim} p\) so

\[
E \sim \frac{\hbar^{2}}{2 m l^{2}}
\]

The force excerted on the wall can be obtained most simply from

\[
F=-\frac{\partial U}{\partial l}=\frac{\hbar^{2}}{m l^{3}} .
\]

6.73 We write

\[
p \sim \Delta p \sim \frac{\hbar}{\Delta x} \sim \frac{\hbar}{x}
\]

i.e. all four quantities are of the same order of magnitude. Then

\[
E \approx \frac{\hbar^{2}}{2 m x^{2}}+\frac{1}{2} k x^{2}=\frac{1}{2 m}\left(\frac{\hbar}{x}-\sqrt{m k} x\right)^{2}+\hbar \sqrt{\frac{k}{m}}
\]

Thus we get an equilibrium situation \((E=\) minimum \()\) when\\
and then

\[
\begin{gathered}
x=x_{0}=\sqrt{\frac{\hbar}{\sqrt{m k}}} \\
E=E_{0} \sim \hbar \sqrt{\frac{k}{m}}=\hbar \omega
\end{gathered}
\]

Quantum mechanics gives

\[
E_{0}=\hbar \omega / 2
\]

6.74 Hence we write

Then

\[
\begin{aligned}
r \sim \Delta r & , p \sim \Delta p \sim \hbar / \Delta r \\
E & =\frac{\hbar^{2}}{2 m r^{2}}-\frac{e^{2}}{r} \\
& =\frac{1}{2 m}\left(\frac{\hbar}{r}-\frac{m e^{2}}{\hbar}\right)^{2}-\frac{m e^{4}}{2 \hbar^{2}}
\end{aligned}
\]

Hence \(r_{\text {eff }}=\frac{\hbar^{2}}{m e^{2}}=53 \mathrm{pm}\) for the equilibrium state.\\
and then

\[
E=-\frac{m e^{4}}{2 \hbar^{2}}=-13.6 \mathrm{eV}
\]

6.75 Suppose the width of the slit (its extension along the \(y\)-axis) is \(\delta\). Then each electron has an uncertainty \(\Delta y_{\sim} \delta\). This translates to an uncertainty \(\Delta p_{y} \sim \hbar / \delta\). We must therefore have \(p_{y} \geq \hbar / \delta\).

For the image, brodening has two sources.\\
We write

\[
\Delta(\delta)=\delta+\Delta^{\prime}(\delta)
\]

where \(\Delta^{\prime}\) is the width caused by the spreading of electrons due to their transverse momentum. We have

\[
\Delta^{\prime}=v_{y} \frac{l}{v_{x}} \approx p_{y} \frac{l}{p}=\frac{l \hbar}{m v \delta}
\]

Thus

\[
\Delta(\delta)=\delta+\frac{l h}{m v \delta}
\]

For large \(\delta, \Delta(\delta) \sim \delta\) and quantum effect is unimportant. For small \(\delta\), quantum effects are large. But \(\Delta(\delta)\) is minimum when

\[
\delta=\sqrt{\frac{l h}{m v}}
\]

as we see by completing the square. Substitution gives

\[
\delta=1.025 \times 10^{-5} \mathrm{~m} \approx 0.01 \mathrm{~mm}
\]

6.76 The Schrodinger equation in one dimension for a free particle is

\[
i \hbar \frac{\partial \psi}{\partial t}=-\frac{\hbar^{2}}{2 m} \frac{\partial^{2} \psi}{\partial x^{2}}
\]

we write \(\psi(x, t)=\varphi(x) \chi(t)\). Then

Then

\[
\begin{gathered}
\frac{i \hbar}{\chi} \frac{d \chi}{d t}=-\frac{\hbar^{2}}{2 m} \frac{1}{\varphi} \frac{d^{2} \varphi}{d x^{2}}=E, \text { say } \\
\chi(t) \sim \exp \left(-\frac{i E t}{h}\right) \\
\varphi(x) \stackrel{\sim}{\sim} \exp \left(i \frac{\sqrt{2 m E}}{\hbar} x\right)
\end{gathered}
\]

E must be real and positive if \(\varphi(x)\) is to be bounded everywhere. Then

\[
\psi(x, t)=\text { Const } \exp \left(\frac{i}{\hbar}(\sqrt{2 m E} x-E t)\right)
\]

This particular solution describes plane waves.\\
6.77 We look for the solution of Schrodinger eqn. with


\begin{equation*}
-\frac{\hbar^{2}}{2 m} \frac{d^{2} \psi}{d x^{2}}=E \psi, 0 \leq x \leq l \tag{1}
\end{equation*}


The boundary condition of impenetrable walls means

\[
\begin{gathered}
\psi(x)=0 \text { for } x=0 \text { and } x=l \\
\text { (as } \psi(x)=0 \text { for } x<0 \text { and } x>l,)
\end{gathered}
\]

The solution of (1) is

Then

\[
\begin{aligned}
& \psi(x)= A \sin \frac{\sqrt{2 m E}}{\hbar} x+B \cos \frac{\sqrt{2 m E}}{\hbar} \cdot x \\
& \psi(0)=0 \Rightarrow B=0 \\
& \psi(l)=0 \Rightarrow A \sin \frac{\sqrt{2 m E}}{\hbar} l=0
\end{aligned}
\]

\(A \approx 0\) so

Hence

\[
\begin{gathered}
\frac{\sqrt{2 m E}}{n} l=n \pi \\
E_{n}=\frac{n^{2} \pi^{2} \hbar^{2}}{2 m l^{2}}, n=1,2,3, \ldots
\end{gathered}
\]

Thus the ground state wave function is

\[
\psi(x)=A \sin \frac{\pi x}{l}
\]

We evaluate \(\boldsymbol{A}\) by nomalization

Thus

\[
\begin{aligned}
1=A^{2} \int_{0}^{l} \sin ^{2} \frac{\pi x}{l} d x & =A^{2} \frac{l}{\pi} \int_{0}^{\pi} \sin ^{2} \theta d \theta=A^{2} \frac{l}{\pi} \cdot \frac{\pi}{2} \\
A & =\sqrt{\frac{2}{l}}
\end{aligned}
\]

Finally, the probability \(P\) for the particle to lie in \(\frac{l}{2} \leq x \leq \frac{2 l}{3}\) is

\[
\begin{aligned}
& P=P\left(\frac{l}{3} \leq x \leq \frac{2 l}{3}\right)=\frac{2}{l} \int_{\frac{l}{3}}^{2 L / 3} \sin ^{2} \frac{\pi x}{l} d x \\
& =\frac{2}{\pi} \int_{\pi / 3}^{2 \pi / 3} \sin ^{2} \theta d \theta=\frac{1}{\pi} \int_{\pi / 3}^{2 \pi / 3}(1-\cos 2 \theta) d \theta
\end{aligned}
\]

\[
\begin{aligned}
& =\frac{1}{\pi}\left(\theta-\frac{1}{2} \sin 2 \theta\right)_{\pi / 3}^{2 \pi / 3}=\frac{1}{\pi}\left(\frac{2 \pi}{3}-\frac{\pi}{3}-\frac{1}{2} \sin \frac{4 \pi}{3}+\frac{1}{2} \sin \frac{2 \pi}{3}\right) \\
& =\frac{1}{\pi}\left(\frac{\pi}{3}+\frac{1}{2} \frac{\sqrt{3}}{2}+\frac{1}{2} \frac{\sqrt{3}}{2}\right)=\frac{1}{3}+\frac{\sqrt{3}}{2 \pi}=0.609
\end{aligned}
\]

6.78 Here \(-\frac{l}{2} \leq x \leq \frac{l}{2}\). Again we have

\[
\psi(x)=B \cos \frac{\sqrt{2 m E} x}{\hbar}+A \sin \frac{\sqrt{2 m E} x}{\hbar}
\]

Then the boundary condition \(\psi\left( \pm \frac{l}{2}\right)=0\)\\
gives

\[
B \cos \frac{\sqrt{2 m E} l}{2 \hbar} \pm A \sin \frac{\sqrt{2 m E} l}{2 \hbar}=0
\]

There are two cases.\\
(1) \(A=0, \frac{\sqrt{2 m E} l}{2 \hbar}=n \pi+\frac{\pi}{2}\)\\
gives even solution. Here\\
and

\[
\begin{gathered}
\sqrt{2 m E}=(2 n+1) \frac{\pi \hbar}{l} \\
E_{n}=(2 n+1)^{2} \frac{\pi^{2} \hbar^{2}}{2 m l^{2}} \\
\psi_{n}^{e}(x)=\sqrt{\frac{2}{l}} \cos (2 n+1) \frac{\pi x}{l} \\
n=0,1,2,3, \ldots
\end{gathered}
\]

This solution is even under \(x \rightarrow-x\).\\
(2) \(B=0\),

\[
\begin{gathered}
\frac{\sqrt{2 m E} l}{2 h}=n \pi, n=1,2, \ldots \\
E_{n}=(2 n \pi)^{2} \frac{\hbar^{2}}{2 m l^{2}} \\
\psi_{n}^{0}=\sqrt{\frac{2}{l}} \sin \frac{2 n \pi x}{l}, n=1,2, \ldots \text { This solution is odd. }
\end{gathered}
\]

6.79 The wave function is given in 6.77. We see that

\[
\int_{0}^{l} \psi_{n}(x) \psi_{n^{\prime}}(x) d x=\frac{2}{l} \int_{0}^{l} \sin \frac{n \pi x}{l} \sin \frac{n^{\prime} \pi x}{l} d x
\]

\[
\begin{aligned}
& =\frac{1}{l} \int_{0}^{l}\left[\cos \left(n-n^{\prime}\right) \frac{\pi x}{l}-\cos \left(n+n^{\prime}\right) \frac{\pi x}{l}\right] d x \\
& =\frac{1}{l}\left[\frac{\sin \left(n-n^{\prime}\right) \pi x / l}{\left(n-n^{\prime}\right) \frac{\pi}{l}}-\frac{\sin \left(n+n^{\prime}\right) \frac{\pi x}{l}}{\left(n+n^{\prime}\right) \frac{\pi x}{l}}\right]_{0}^{l} e
\end{aligned}
\]

If \(n \neq n^{\prime}\), this is zero as \(n\) and \(n^{\prime}\) are integers.\\
6.80 We have found that

\[
E_{n}=\frac{n^{2} \pi^{2} \hbar^{2}}{2 m l^{2}}
\]

Let \(N(E)=\) number of states upto \(E\). This number is \(n\). The number of states upto \(E+d E\) is \(N(E+d E)=N(E)+d N(E)\). Then \(d N(E)=1\) and

\[
\frac{d N(E)}{d E}=\frac{1}{\Delta E}
\]

where \(\Delta E=\) difference in energies between the \(n^{\text {th }} \&(n+1)^{\text {th }}\) level

\[
\begin{aligned}
& =\frac{(n+1)^{2}-n^{2}}{2 m l^{2}} \pi^{2} \hbar^{2}=\frac{2 n+1}{2 m l^{2}} \pi^{2} \hbar^{2} \\
& \equiv \frac{\pi^{2} \hbar^{2}}{2 m l^{2}} 2^{*} n, \quad(\text { neglecting } 1 \ll n) \\
& =\frac{\pi^{2} \hbar^{2}}{2 m l^{2}} \times \sqrt{\frac{2 m l^{2}}{\pi^{2} \hbar^{2}}} \sqrt{E} \times 2 \\
& =\frac{\pi \hbar}{l} \sqrt{\frac{2}{m}} \sqrt{E} \\
& \quad \frac{d N(E)}{d E}=\frac{l}{\pi \hbar} \sqrt{\frac{m}{2 E}} .
\end{aligned}
\]

Thus\\
For the given case this gives \(\frac{d N(E)}{d E}=0.816 \times 10^{7}\) levels per eV\\
6.81 (a) Here the schrodinger equation is

\[
-\frac{\hbar^{2}}{2 m}\left(\frac{\partial^{2}}{\partial x^{2}}+\frac{\partial^{2}}{\partial y^{2}}\right) \psi=E \psi
\]

we take the origin at one of the corners of the rectangle where the particle can lie. Then the wave function must vanish for

\[
x=0 \quad \text { or } \quad x=l_{1}
\]

or

\[
y=0 \text { or } y=l_{2} .
\]

we look for a solution in the form

\[
\psi=A \sin k_{1} x \sin k_{2} y
\]

cosines are not permitted by the boundary condition. Then\\
and

\[
\begin{gathered}
k_{1}=\frac{n_{i} \pi}{l_{1}}, k_{2}=n_{2} \frac{\pi}{l_{2}} \\
E=\frac{k_{1}^{2}+k_{2}^{2}}{2 m} \hbar^{2}=\frac{\pi^{2} \hbar^{2}}{2 m}\left(\frac{n_{1}^{2}}{l_{1}^{2}}+\frac{n_{2}^{2}}{l_{2}^{2}}\right)
\end{gathered}
\]

Here \(n_{1}, n_{2}\) are nonzero integers.\\
(b) If \(l_{1}=l_{2}=l\) then

\[
\begin{array}{lr} 
& \frac{E}{\hbar^{2} / m l^{2}}=\frac{n_{1}^{2}+n_{2}^{2}}{2} \pi^{2} \\
1^{\text {st }} \text { level : } & n_{1}=n_{2}=1 \rightarrow \pi^{2}=9.87 \\
2^{\text {nd }} \text { level : } & \left.\begin{array}{l}
n_{1}=1, n_{2}=2 \\
n_{1}=2, n_{2}=1
\end{array}\right\} \rightarrow \frac{5}{2} \pi^{2}=24.7 \\
& \text { or } \quad \begin{array}{l}
n_{1}=2,=n_{2}=2 \rightarrow 4 \pi^{2}=39.5 \\
3^{\text {st }} \text { level : }
\end{array} \\
4^{\text {nd }} \text { level : } & \left.\begin{array}{l}
n_{1}=1, n_{2}=3 \\
n_{1}=3, n_{2}=1
\end{array}\right\} \rightarrow 5 \pi^{2}=49.3
\end{array}
\]

6.82 The wave function for the ground state is

\[
\psi_{11}(x, y)=A \sin \frac{\pi x}{a} \sin \frac{\pi y}{b}
\]

we find \(A\) by normalization

\[
1=A^{2} \int_{0}^{a} d x \int_{0}^{d} d y \sin ^{2} \frac{\pi x}{a} \sin ^{2} \frac{\pi y}{b}=A^{2} \frac{a b}{4}
\]

Thus

\[
A=\frac{2}{\sqrt{a b}}
\]

Then the requisite probability is

\[
\begin{aligned}
P & =\int_{0}^{a / 3} d x \int_{0}^{b} d y \frac{4}{a b} \sin ^{2} \frac{\pi x}{a} \sin ^{2} \frac{\pi y}{b} \\
& =\frac{2}{a} \int_{0}^{a / 3} d x \sin ^{2} \frac{\pi x}{a} \text { on doing the } y \text { integral }
\end{aligned}
\]

\[
\begin{aligned}
& =\frac{1}{a} \int_{0}^{a / 3} d\left(1-\cos \frac{2 \pi x}{a}\right)=\frac{1}{a}\left(\frac{a}{3}-\frac{\sin \frac{2 \pi}{3}}{2 \pi / a}\right) \\
& =\frac{1}{3}-\frac{\sqrt{3}}{4 \pi}=0.196=19.6 \%
\end{aligned}
\]

6.83 We proceed axactly as in (6.81). The wave function is chosen in the form \(\psi(x, y, z)=A \sin k_{1} x \sin k_{2} y \sin k_{3} z\).\\
(The origin is at one corner of the box and the axes of coordinates are along the edges.) The boundary conditions are that \(\psi=0\) for

\[
x=0, x=a, y=0, y=a, z=0, z=a
\]

This gives

\[
k_{1}=\frac{n_{1} \pi}{a}, k_{2}=\frac{n_{2} \pi}{a}, k_{3}=\frac{n_{3} \pi}{a}
\]

The energy eigenvalues are

\[
E\left(n_{1}, n_{2}, n_{3}\right)=\frac{\pi_{2} \hbar^{2}}{2 m a^{2}}\left(n_{1}^{2}+n_{2}^{2}+n_{3}\right)
\]

The first level is \((1,1,1)\). The second has \((1,1,2),(1,2,1) \&(2,1,1)\). The third level is \((1,2,2)\) or \((2,1,2)\) or \((2,2,1)\). Its energy is

\[
\frac{9 \pi^{2} \hbar^{2}}{2 m a^{2}}
\]

The fourth energy level is \((1,1,3)\) or \((1,3,1)\) or \((3,1,1)\)

Its energy is

\[
E=\frac{11 \pi^{2} \hbar^{2}}{2 m a^{2}}
\]

(b) Thus

\[
\Delta=E_{4}-E_{3}=\frac{\hbar^{2} \pi^{2}}{m a^{2}}
\]

(c) The fifth level is \((2,2,2)\). The sixth level is \((1,2,3),(1,3,2),(2,1,3),(2,3,1),(3\), \(1,2),(3,2,1)\)\\
Its energy is

\[
\frac{7 \hbar^{2} \pi^{2}}{m a^{2}}
\]

and its degree of degeneracy is 6 (six).\\
6.84 We can for definiteness assume that the discontinuity occurs at the point \(x=0\). Now the schrodinger equation is

\[
-\frac{\hbar^{2}}{2 m} \frac{d^{2} \psi}{d x^{2}}+U(x) \psi(x)=E \psi(x)
\]

We integrate this equation around \(x=0\) i.e., from \(x=-\varepsilon_{1}\) to \(x=+\varepsilon_{2}\) where \(\varepsilon_{1}, \varepsilon_{2}\) are small positive numbers. Then\\
or

\[
\begin{gathered}
-\frac{\hbar^{2}}{2 m} \int_{-\varepsilon_{1}}^{+\varepsilon_{2}} \frac{d^{2} \psi}{d x^{2}} d x=\int_{-\varepsilon_{1}}^{+\varepsilon_{2}}(E-U(x) \psi(x) d x \\
\left(\frac{d \psi}{d x}\right)_{+\varepsilon_{2}}-\left(\frac{d \psi}{d x}\right)_{-\varepsilon_{1}}=-\frac{2 m}{\hbar^{2}} \int_{-\varepsilon_{1}}^{\varepsilon_{2}}(E-U(x))_{d x} \psi(x)
\end{gathered}
\]

Since the potential and the energy \(E\) are finite and \(\psi(x)\) is bounded by assumption, the integral on the right exists and \(\rightarrow 0\) as \(\varepsilon_{1}, \varepsilon_{2} \rightarrow 0\)

Thus

\[
\left(\frac{d \psi}{d x}\right)_{+\varepsilon_{2}}=\left(\frac{d \psi}{d x}\right)_{-\varepsilon_{1}} \text { as } \varepsilon_{1}, \varepsilon_{2} \rightarrow 0
\]

So \(\left(\frac{d \psi}{d x}\right)\) is continuous at \(x=0\) (the point where \(U(x)\) has a finite jump discontinuity.)\\
6.85\\
\includegraphics[max width=\textwidth, alt={}, center]{da0156b6-4133-434e-90d4-ec01b2070afe-308_371_697_1083_415}\\
(a) Starting from the Schrodinger equation in the regions \(I \& I I\)


\begin{gather*}
\frac{d^{2} \psi}{d x^{2}}+\frac{2 m E}{\hbar^{2}} \psi=0 \quad x \text { in } I  \tag{1}\\
\frac{d^{2} \psi}{d x^{2}}-\frac{2 m E\left(U_{0}-E\right)}{\hbar^{2}} \psi=0 \quad x \text { in } I I \tag{2}
\end{gather*}


where \(U_{0}>E>0\), we easily derive the solutions in \(I \& I I\)


\begin{gather*}
\Psi_{I}(x)=A \sin k x+B \cos k x  \tag{3}\\
\Psi_{I I}(x)=C e^{\alpha x}+D e^{-\alpha x} \tag{4}
\end{gather*}


where

\[
k^{2}=\frac{2 m E}{\hbar^{2}}, \alpha^{2}=\frac{2 m\left(U_{0}-E\right)}{\hbar^{2}}
\]

The boundary conditions are


\begin{equation*}
\Psi(o)=0 \tag{5}
\end{equation*}


and \(\Psi \&\left(\frac{d \Psi}{d x}\right)\) are continuous at \(x=l\), and \(\Psi\) must vanish at \(x=+\infty\).\\
Then

\[
\psi_{I}=A \sin k x
\]

and

\[
\Psi_{I I}=D e^{-\alpha x}
\]

so

\[
\begin{aligned}
A \sin k l & =D e^{-\alpha l} \\
k A \cos k l & =-\alpha D e^{-\alpha l}
\end{aligned}
\]

From this we get

\[
\tan k l=-\frac{k}{\alpha}
\]

or


\begin{align*}
\sin k l & = \pm k l / \sqrt{k^{2} l^{2}+\alpha^{2} l^{2}} \\
& = \pm k l / \sqrt{\frac{2 m U_{0} l^{2}}{h^{2}}} \\
: & = \pm k l \sqrt{\hbar^{2} / 2 m U_{0} l^{2}} \tag{6}
\end{align*}


\begin{center}
\includegraphics[max width=\textwidth, alt={}]{da0156b6-4133-434e-90d4-ec01b2070afe-309_552_715_1039_312}
\end{center}

Plotting the left and right sides of this equation we can find the points at which the straight lines cross the sine curve. The roots of the equation corresponding to the eigen values of encrgy \(E_{i}\) and found from the inter section points \((k l)_{i}\), for which tan \((k l)_{i}<0\) (i.e. \(2^{\text {nd }}\) \& \(4^{\text {th }}\) and other even quadrants). It is seen that bound states do not always exist. For the first bound state to appear (refer to the line (b) above)

\[
(k l)_{1, \min }=\frac{\pi}{2}
\]

(b) Substituting, we get \(\left(l^{2} U_{0}\right)_{1, \min }=\frac{\pi^{2} h^{2}}{8 m}\)\\
as the condition for the appearance of the first bound state. The second bound state will appear when \(k l\) is in the fourth quadrant. The magnitude of the slope of the straight line must then be less than

\[
\frac{1}{3 \pi / 2}
\]

Corresponding to \((k l)_{2, \min }=\frac{3 \pi}{2}=(3) \frac{\pi}{2}=(2 \times 2-1) \frac{\pi}{2}\)\\
For \(n\) bound states, it is easy to convince one self that the slope of the appropriate straight line (upper or lower) must be less than

Then

\[
\begin{aligned}
(k l)_{n, \min } & =(2 n-1) \frac{\pi}{2} \\
\left(l^{2} U_{0}\right)_{n, \min } & =\frac{(2 n-1)^{2} \pi^{2} \hbar^{2}}{8 m}
\end{aligned}
\]

Do not forget to note that for large \(n\) both + and - signs in the Eq. (6) contribute to solutions.\\
\(6.86 U_{0} l^{2}=\left(\frac{3}{4} \pi\right)^{2} \frac{\hbar^{2}}{m}\)\\
and

\[
\begin{gathered}
E l^{2}=\left(\frac{3}{4} \pi\right)^{2} \frac{\hbar^{2}}{2 m} \\
k l=\frac{3}{4} \pi
\end{gathered}
\]

It is easy to check that the condition of the boud state is satisfied. Also

\[
\alpha l=\sqrt{\frac{2 m}{\hbar^{2}}\left(U_{0}-E\right) l^{2}}=\sqrt{\frac{m U_{0}}{\hbar^{2}} l^{2}}=\frac{3}{4} \pi
\]

Then from the previous problem

\[
D=A e^{\alpha l} \sin k l=A \frac{e^{3 \pi / 4}}{\sqrt{2}}
\]

By normalization

\[
I=A^{2}\left[\int_{0}^{l} \sin ^{2} k x d x+\int_{l}^{\infty} \frac{e^{3 \pi / 2}}{2} e^{-(3 \pi / 2) x / l} d x\right]
\]

\[
\begin{aligned}
&= A^{2}\left[\frac{1}{2} \int_{0}^{l}(1-\cos 2 k x) d x+l \int_{0}^{\infty} \frac{1}{2} e^{-\frac{3 \pi}{2} y} d y\right] \\
&= A^{2}\left[\frac{1}{2}\left[-\frac{\sin 2 k l}{2 k}\right]+\frac{1}{2} \cdot \frac{\frac{l}{3 \pi}}{2}\right]=A^{2} l\left[\frac{1}{2}\left[1+\frac{\frac{1}{3 \pi}}{2}\right]+\frac{1}{2} \frac{\frac{1}{3 \pi}}{2}\right] \\
&=A^{2} l\left[\frac{1}{2}+\frac{2}{3 \pi}\right]=A^{2} \frac{l}{2}\left(1+\frac{4}{3 \pi}\right) \text { or } A=\sqrt{\frac{2}{l}}\left(1+\frac{4}{3 \pi}\right)^{-1 / 2}
\end{aligned}
\]

The probability of the particle to be located in the region \(x>l\) is

\[
\begin{aligned}
P & =\int_{l}^{\infty} \psi^{2} d x=\frac{2}{l}\left(1+\frac{4}{3 \pi}\right)^{-1} \int_{l}^{\infty} \frac{e^{3 \pi / 2}}{2} e^{-\frac{3 \pi}{2} \frac{x}{l}} d x \\
& =\left(1+\frac{4}{3 \pi}\right)^{-1} \int_{l}^{\infty} e^{3 \pi / 2} e^{-(3 \pi / 2) y} d y=\frac{2}{3 \pi} \times \frac{3 \pi}{3 \pi+4}=14.9 \%
\end{aligned}
\]

6.87 The Schrodinger equation is

\[
\nabla^{2} \psi+\frac{2 m}{\hbar^{2}}(E-U(r)) \psi=0
\]

when \(\psi\) depends on \(\mathbf{r}\) only,

\[
\nabla^{2} \psi=\frac{1}{r^{2}} \frac{d}{d r}\left(r^{2} \frac{d \psi}{d r}\right)
\]

If we put

\[
\psi=\frac{\chi(r)}{r}, \frac{d \psi}{d r}=\frac{\chi^{\prime}}{r}-\frac{\chi}{r^{2}}
\]

and

\[
\nabla^{2} \psi=\frac{\chi^{\prime \prime}}{r} . \text { Thus we get }
\]

\[
\frac{d^{2} \chi}{d r^{2}}+\frac{2 m}{\hbar^{2}}(E-U(r)) \chi=0
\]

The solution is

\[
\chi=A \sin k r, r<r_{0}
\]

\[
\begin{array}{r}
k^{2}=\frac{2 m E}{\hbar^{2}} \\
\chi=0 \quad r>r_{0}
\end{array}
\]

and\\
(For \(r<r_{0}\) we have rejected a term \(B \cos k r\) as it does not vanish at \(r=0\) ). Continuity of the wavefunction at \(r=r_{0}\) requires

Hence

\[
\begin{gathered}
k r_{0}=n \pi \\
E_{n}=\frac{n^{2} \pi^{2} \hbar^{2}}{2 m r_{0}^{2}}
\end{gathered}
\]

6.88 (a) The nomalized wave functions are obtained from the normalization

\[
\begin{aligned}
1 & =\int|\psi|^{2} d V=\int|\psi|^{2} 4 \pi r^{2} d r \\
& =\int_{0}^{r_{0}} A^{2} 4 \pi \chi^{2} d r=4 \pi A^{2} \int_{0}^{r_{0}} \sin ^{2} \frac{n \pi r}{r_{0}} d r \\
& =4 \pi A^{2} \frac{r_{0}}{n \pi} \int_{0}^{\pi} \sin ^{2} x d r=4 \pi A^{2} \frac{r_{0}}{n \pi} \cdot \frac{n \pi}{2}=r_{0} \cdot 2 \pi A^{2} \\
A & =\frac{1}{\sqrt{2 \pi r_{0}}} \text { and } \psi=\frac{1}{\sqrt{2 \pi \cdot r_{0}}} \frac{\sin \frac{n \pi r}{r_{0}}}{r}
\end{aligned}
\]

(b) The radial probability distribution function is

\[
P_{n}(r)=4 \pi r^{2}(\psi)^{2}=\frac{2}{r_{0}} \sin ^{2} \frac{n \pi r}{r_{0}}
\]

For the ground state

\[
n=1
\]

so

\[
P_{1}(r)=\frac{2}{r_{0}} \sin ^{2} \frac{\pi r}{r_{0}}
\]

By inspection this is maximum for \(r=\frac{r_{0}}{2}\). Thus \(r_{p r}=\frac{r_{0}}{2}\)\\
The probability for the particle to be found in the region \(r<r_{p r}\) is clearly \(50 \%\) as one can immediately see from a graph of \(\sin ^{2} x\).\\
6.89 If we put \(\psi=\frac{\chi(r)}{r}\)\\
the equation for \(\chi(r)\) has the form\\
which can be written as

\[
\chi^{\prime \prime}+\frac{2 m}{\hbar^{2}}[E-U(r)] \chi(r)=0
\]

and

\[
\chi^{\prime \prime}+k^{2} \chi=0,0 \leq r<r_{0}
\]

where

\[
k^{2}=\frac{2 m E}{\hbar^{2}}, \alpha^{2}=\frac{2 m\left(U_{0}-E\right)}{\hbar^{2}} .
\]

The boundary condition is

\[
\left.\begin{array}{c}
\chi(0)=0 \\
\text { and } \chi, \chi^{\prime} \text { are continuous at } r=r_{0}
\end{array}\right\}
\]

These are exactly same as in the one dimensional problem in problem (6.85) We therefore omit further details\\
6.90 The Schrodinger equation is \(\frac{d^{2} \Psi}{d x^{2}}+\frac{2 m}{h^{2}}\left(E-\frac{1}{2} k x^{2}\right) \Psi=0\)

We are given

\[
\Psi=A e^{-\alpha x^{2} / 2}
\]

Then

\[
\Psi^{\prime}=-\propto x A e^{-\propto x^{2} / 2}
\]

\[
\Psi^{\prime \prime}=-\propto A e^{-\propto x^{2} / 2}+\propto^{2} x^{2} A e^{-\propto x^{2} / 2}
\]

Substituting we find that following equation must hold

\[
\left[\left(\propto^{2} x^{2}-\propto\right)+\frac{2 m}{h^{2}}\left(E-\frac{1}{2} k x^{2}\right)\right] \Psi=0
\]

since \(\Psi \neq 0\), the bracket must vanish identicall. This means that the coefficient of \(x^{2}\) as well the term independent of \(x\) must vanish. We get

\[
\propto^{2}=\frac{m k}{h^{2}} \text { and } \propto=\frac{2 m E}{h^{2}}
\]

Putting \(k / m=\omega^{2}\), this leads to \(\propto=\frac{m \omega}{h^{2}}\) and \(E=\frac{h^{2} \omega}{2}\)\\
6.91 The Schrödinger equation for the problem in Gaussian units

\[
\nabla^{2} \psi+\frac{2 m}{\hbar^{2}}\left[E+\frac{e^{2}}{r}\right] \psi=0
\]

In MKS units we should read \(\left(e^{2} / 4 \pi \varepsilon_{0}\right)\) for \(e^{2}\).\\
we put


\begin{equation*}
\psi=\frac{\chi(r)}{r} . \text { Then } \chi^{\prime \prime}+\frac{2 m}{\hbar^{2}}\left[E+\frac{e^{2}}{r}\right] \chi=0 \tag{1}
\end{equation*}


We are given that

\[
\chi=r \psi=A r(1+a r) e^{-\alpha r}
\]

so

\[
\begin{gathered}
\chi^{\prime}=A(1+2 a r) e^{-\alpha r}-\alpha A r(1+a r) e^{-\alpha r} \\
\chi^{\prime \prime}=\alpha^{2} A r(1+a r) e^{-\alpha r}-2 \alpha A(1+2 a r) e^{-\alpha r}+2 a A e^{-\alpha r}
\end{gathered}
\]

Substitution in (1) gives the condition

\[
\alpha^{2}\left(r+a r^{2}\right)-2 \alpha(1+2 a r)+2 a+\frac{2 m}{\hbar^{2}}\left(E r+e^{2}\right) \times(1+a r)=0
\]

Equating the coefficients of \(r^{2}, r\), and constant term to zero we get


\begin{gather*}
2 a-2 \alpha+\frac{2 m e^{2}}{\hbar^{2}}=0  \tag{2}\\
a \alpha^{2}+\frac{2 m}{\hbar^{2}} E a=0  \tag{3}\\
\alpha^{2}-4 a \alpha+\frac{2 m}{\hbar^{2}}\left(E+e^{2} a\right)=0 \tag{4}
\end{gather*}


From (3) either \(a=0, \quad\) or \(\quad E=-\frac{\hbar^{2} \alpha^{2}}{2 m}\)

In the first case

\[
\alpha=\frac{m e^{2}}{\hbar^{2}}, E=-\frac{\hbar^{2}}{2 m} \alpha^{2}=-\frac{m e^{4}}{2 \hbar^{2}}
\]

This state is the ground state.\\
In the second case

\[
\begin{gathered}
a=\alpha-\frac{m e^{2}}{\hbar^{2}}, \alpha=\frac{1}{2} \frac{m e^{2}}{\hbar^{2}} \\
E=-\frac{m e^{4}}{8 \hbar^{2}} \quad \text { and } a=-\frac{1}{2} \frac{m e^{2}}{\hbar^{2}}
\end{gathered}
\]

This state is one with \(n=2(2 \mathrm{~s})\).\\
6.92 We first find A by normalization

\[
1=\int_{0}^{\infty} 4 \pi A^{2} e^{-2 r / r_{1}} r^{2} d r=\frac{\pi A^{2}}{2} r_{1}^{3} \int_{0}^{\infty} e^{-x} x^{2} d x=\pi A^{2} r_{1}^{3}
\]

since the integral has the value 2 .\\
Thus

\[
A^{2}=\frac{1}{\pi r_{1}^{3}} \text { or } A=\frac{1}{\sqrt{r_{1}^{3} \pi}}
\]

(a) The most probable distance \(r_{p r}\) is that value of \(r\) for which

\[
P(r)=4 \pi r^{2}|\psi(r)|^{2}=\frac{4}{r_{1}^{3}} r^{2} e^{-2 r / r_{1}}
\]

is maximum. This requires\\
or

\[
P^{\prime}(r)=\frac{4}{r_{1}^{3}}\left[2 r-\frac{2 r^{2}}{r_{1}}\right] e^{-2 r / r_{1}}=0
\]

\[
r=r_{1}=r_{p r}
\]

(b) The coulomb force being given by \(-e^{2} / r^{2}\), the mean value of its modulus is

\[
\begin{aligned}
<F> & =\int_{0}^{\infty} 4 \pi r^{2} \frac{1}{\pi r_{1}^{3}} e^{-2 r / r_{1}} \frac{e^{2}}{r^{2}} d r \\
& =\int_{0}^{\infty} \frac{4 e^{2}}{r_{1}^{3}} e^{-2 r / r_{1}} d r=\frac{2 e^{2}}{r_{1}^{2}} \int_{0}^{\infty} e^{-x} d x=\frac{2 e^{2}}{r_{1}^{2}}
\end{aligned}
\]

In MKS units we should read \(\left(e^{2} / 4 \pi \varepsilon_{0}\right)\) for \(e^{2}\)\\
(c)

\[
<U>=\int_{0}^{\infty} 4 \pi r^{2} \frac{1}{\pi r_{1}^{3}} e^{-2 r / r_{3}} \frac{-e^{2}}{r} d r=-\frac{e^{2}}{r_{1}} \int_{0}^{\infty} x e^{-x} d x=-\frac{e^{2}}{r_{1}}
\]

In MKS units we should read \(\left(e^{2} / 4 \pi \varepsilon_{0}\right)\) for \(e^{2}\).\\
6.93 We find \(A\) by normalization as above. We get

\[
A=\frac{1}{\sqrt{\pi r_{1}^{3}}}
\]

Then the electronic charge density is

\[
\rho=-e|\psi|^{2}=-e \frac{e^{-2 r / r_{1}}}{\pi r_{1}^{3}}=\rho(\vec{r})
\]

The potential \(\psi(\vec{r})\) due to this charge density is

\[
\varphi(\vec{r})=\frac{1}{4 \pi \varepsilon_{0}} \int \frac{\rho\left(\overrightarrow{r^{\prime}}\right)}{|\vec{r}-\vec{r}|} d^{3} r^{\prime}
\]

so at the origin \(\varphi(0)=\frac{1}{4 \pi \varepsilon_{0}} \int_{0}^{\infty} \frac{\rho\left(r^{\prime}\right)}{r^{\prime}} 4 \pi r^{\prime 2} d r^{\prime}=\frac{-e}{4 \pi \varepsilon_{0}} \int_{0}^{\infty} \frac{4 r^{\prime}}{r_{1}^{3}} e^{-2 r^{\prime} / r_{1}} d r\)

\[
=-\frac{e}{4 \pi \varepsilon_{0} r_{1}} \int_{0}^{\infty} x e^{-x} d x=-\frac{e}{\left(4 \pi \varepsilon_{0}\right) r_{1}}
\]

6.94 (a) We start from the Schrodinger equation \(\frac{d^{2} \psi}{d x^{2}}+\frac{2 m}{\hbar^{2}}(E-U(x)) \psi=0\) which we write as \(\Psi_{I}^{\prime \prime}+k^{2} \Psi_{I}=0, x<0\)\\
and

\[
\begin{gathered}
k^{2}=\frac{2 m E}{\hbar^{2}} \\
\Psi_{I I}^{\prime \prime}+\alpha^{2} \Psi_{I I}=0 \quad x>0 \\
\alpha^{2}=\frac{2 m}{\hbar^{2}}\left(E-U_{0}\right)>0
\end{gathered}
\]

It is convenient to look for solutions in the form

\[
\begin{gathered}
\psi_{I}=e^{i k x}+R e^{-i k x} x<0 \\
\Psi_{\Pi}=A e^{i \alpha x}+B e^{-i \alpha x} x>0
\end{gathered}
\]

In region \(I(x<0)\), the amplitude of \(e^{i k x}\) is written as unity by convention. In \(I I\) we expect only a transmitted wave to the right, \(B=0\) then. So

\[
\Psi_{I I}=A e^{i k x} x>0
\]

The boundary conditions follow from the continuity of \(\Psi \& \frac{d \psi}{d x}\) at \(x=0\).

Then

\[
\begin{aligned}
& 1+R=A \\
& i K(1-R)=i \alpha A \\
& \frac{1-R}{1+R}=\frac{\alpha}{k} \quad \text { or } \quad R=\frac{k-\alpha}{k+\alpha}
\end{aligned}
\]

The reflection coefficient is the absolute square of \(R\) :

\[
r=|R|^{2}=\left|\frac{k-\alpha}{k+\alpha}\right|^{2}
\]

(b) In this case \(E<U_{0}, \alpha^{2}=-\beta^{2}<0\). Then \(\Psi_{I}\) is unchanged in form but

\[
\Psi_{I I}=A e^{-\beta x}+B e^{+\beta x}
\]

we must have \(B=0\) since otherwise \(\Psi(x)\) will become unbounded as \(x \rightarrow \infty\). Finally

\[
\Psi_{\Pi}=A e^{-\beta x}
\]

Inside the barrier, the particle then has a probability density equal to

\[
\left|\psi_{I I}\right|^{2}=|A|^{2} e^{-2 \beta x}
\]

This decreases to \(\frac{1}{e}\) of its value in

\[
x_{\text {eff }}=\frac{1}{2 \beta}=\frac{\hbar}{2 \sqrt{2 m\left(U_{0}-E\right)}}
\]

6.95 The formula is

\[
D \approx \exp \left[-\frac{2}{\hbar} \int_{x_{1}}^{x_{2}} \sqrt{2 m(V(x)-E)} d x\right]
\]

Here \(V\left(x_{2}\right)=V\left(x_{1}\right)=E\) and \(V(x)>E\) in the region \(x_{2}>x>x_{1}\).\\
(a) For the problem, the integral is trivial

\[
D \approx \exp \left[\begin{array}{ll}
-\frac{2 l}{\hbar} & \sqrt{2 m\left(U_{0}-E\right)}
\end{array}\right]
\]

(b) We can without loss of generality take \(\boldsymbol{x} \boldsymbol{=} \mathbf{0}\) at the point the potential begins to climb. Then

Then

\[
\begin{aligned}
& U(x)=\left\{\begin{array}{cl}
0 & x<0 \\
U_{0} \frac{x}{l} & 0<x<l \\
0 & x>l
\end{array}\right. \\
& D \approx \exp \left[-\frac{2}{\hbar} \int_{l \frac{\varepsilon}{U_{0}}}^{l} \sqrt{2 m\left(U_{0} \frac{x}{l}-E\right)} d x\right] \\
&=\exp \left[-\frac{2}{\hbar} \sqrt{\frac{2 m U_{0}}{l}} \int_{x_{0}}^{l} \sqrt{x-x_{0}} d x\right] x_{0}=l \frac{E}{U_{0}}
\end{aligned}
\]

\[
\begin{aligned}
& =\exp \left[-\left.\frac{2}{\hbar} \sqrt{\frac{2 m U_{0}}{l}} \frac{2}{3}\left(x-x_{0}\right)^{3 / 2}\right|_{x_{0}} ^{l}\right] \\
& =\exp \left[-\frac{4}{3 \hbar} \sqrt{\frac{2 m U_{0}}{l}}\left(l-l \frac{E}{U_{0}}\right)^{3 / 2}\right] \\
& =\exp \left[-\frac{4 l}{3 \hbar U_{0}}\left(U_{0}-E\right)^{3 / 2} \sqrt{2 m}\right]
\end{aligned}
\]

6.96 The potential is \(U(x)=U_{0}\left(1-\frac{x^{2}}{l^{2}}\right)\). The turning points are

Then

\[
\begin{aligned}
& \frac{E}{U_{0}}=1-\frac{x^{2}}{l^{2}} \text { or } x= \pm l \sqrt{1-\frac{E}{U_{0}}} . \\
& D \approx \exp \left[-\frac{4}{\hbar} \int_{0}^{l \sqrt{1-\left(E / U_{0}\right)}} \sqrt{2 m\left\{U_{0}\left(1-\frac{x^{2}}{l^{2}}\right)-E\right\}} d x\right] \\
&=\exp \left[-\frac{4}{\hbar} \int_{0}^{l \sqrt{1-\left(E / U_{0}\right)}} \sqrt{2 m U_{0}} \sqrt{1-\frac{E}{U_{0}}-\frac{x^{2}}{l^{2}}} d x\right. \\
&=\exp \left[-\frac{4 l}{\hbar} \sqrt{2 m V_{0}} \int_{0}^{x_{0}} \sqrt{x_{0}^{2}-x^{2}} d x\right], x_{0}=\sqrt{1-E / V_{0}}
\end{aligned}
\]

The integral is

Thus

\[
\begin{gathered}
\int_{0}^{x_{0}} \sqrt{x_{0}^{2}-x^{2}} d x=x_{0}^{2} \int_{0}^{\pi / 2} \cos ^{2} \theta d \theta=\frac{\pi}{4} x_{0}^{2} \\
D \approx \exp \left[-\frac{\pi l}{\hbar} \sqrt{2 m U_{0}}\left(1-\frac{E}{U_{0}}\right)\right] \\
\quad=\exp \left[-\frac{\pi l}{\hbar} \sqrt{\frac{2 m}{U_{0}}}\left(U_{0}-E\right)\right]
\end{gathered}
\]

\subsection*{6.3 PROPERTIES OF ATOMS. SPECTRA}
6.97 From the Rydberg formula we write

\[
E_{n}=-\frac{\hbar R}{\left(n+\alpha_{l}\right)^{2}}
\]

we use \(\hbar R=13.6 \mathrm{eV}\). Then for \(n=2\) state

\[
\begin{gathered}
5 \cdot 39=-\frac{13 \cdot 6}{\left(2+\alpha_{0}\right)^{2}}, l=0(S) \text { state } \\
\alpha_{0} \approx-0 \cdot 41
\end{gathered}
\]

for \(p\) state

\[
\begin{gathered}
3.54=-\frac{13.6}{\left(2+\alpha_{1}\right)^{2}} \\
\alpha_{1}=-0.039
\end{gathered}
\]

6.98 The energy of the \(3 p\) state must be \(-\left(E_{0}-e \varphi\right)\) where \(-E_{0}\) is the energy of the \(3 S\) state. Then\\
so

\[
\begin{gathered}
E_{0}-e \varphi_{1}=\frac{\hbar R}{\left(3+\alpha_{1}\right)^{2}} \\
\alpha_{1}=\sqrt{\frac{\hbar R}{E_{0}-e \varphi_{1}}}-3=-0.885
\end{gathered}
\]

6.99 For the first line of the sharp series \((3 S \rightarrow 2 P)\) in a Li atom

\[
\frac{2 \pi \hbar c}{\lambda_{1}}=-\frac{\hbar R}{\left(3+\alpha_{0}\right)^{2}}+\frac{\hbar R}{\left(2+\alpha_{1}\right)^{2}}
\]

For the short wave cut-off wave-length of the same series

\[
\frac{2 \pi \hbar c}{\lambda_{2}}=\frac{\hbar R}{\left(2+\alpha_{1}\right)^{2}}
\]

From these two equations we get on subtraction

\[
\begin{aligned}
3+\alpha_{0} & =\sqrt{\hbar R / \frac{2 \pi \hbar c\left(\lambda_{1}-\lambda_{2}\right)}{\lambda_{1} \lambda_{2}}} \\
& =\sqrt{\frac{R \lambda_{1} \lambda_{2}}{2 \pi c \Delta \lambda}}, \Delta \lambda=\lambda_{1}-\lambda_{2}
\end{aligned}
\]

Thus in the ground state, the binding energy of the electron is

\[
E_{b}=\frac{\hbar R}{\left(2+\alpha_{0}\right)^{2}}=\hbar R /\left(\sqrt{\frac{R \lambda_{1} \lambda_{2}}{2 \pi c \Delta \lambda}}-1\right)^{2}=5.32 \mathrm{eV}
\]

6.100 The energy of the \(3 S\) state is

\[
E(3 S)=-\frac{\hbar R}{(3-0.41)^{2}}=-2.03 \mathrm{eV}
\]

The energy of a \(2 S\) state is

\[
E(2 S)=-\frac{\hbar R}{(2-0.41)^{2}}=-5.39 \mathrm{eV}
\]

The energy of a \(2 P\) state is

\[
E(2 P)=-\frac{\hbar R}{(2-.04)^{2}}=-3.55 \mathrm{eV}
\]

We see that

\[
E(2 S)<E(2 P)<E(3 S)
\]

The transitions are \(3 S \rightarrow 2 P\) and \(2 P \rightarrow 2 S\).\\
Direct \(3 S \rightarrow 2 S\) transition is forbidden by selection rules. The wavelengths are determined by

\[
E_{2}-E_{1}=\Delta E=\frac{2 \pi \hbar c}{\lambda}
\]

Substitution gives\\
and

\[
\begin{aligned}
& \lambda=0.816 \mu \mathrm{~m}(3 S \rightarrow 2 P) \\
& \lambda=0.674 \mu \mathrm{~m}(2 P \rightarrow 2 S)
\end{aligned}
\]

6.101 The splitting of the \(N a\) lines is due to the fine structure splitting of \(3 p\) lines (The \(3 s\) state is nearly single except for possible hyperfine effects.) The splitting of the \(3 p\) level then equals the energy difference

\[
\Delta E=\frac{2 \pi \hbar c}{\lambda_{1}}-\frac{2 \pi \hbar c}{\lambda_{2}}=\frac{2 \pi \hbar c\left(\lambda_{2}-\lambda_{1}\right)}{\lambda_{1} \lambda_{2}} \approx \frac{2 \pi \hbar c \Delta \lambda}{\lambda^{2}}
\]

Here \(\Delta \lambda=\) wavelength difference \(\& \lambda=\) average wavelength. Substitution gives

\[
\Delta E=2.0 \mathrm{meV}
\]

6.102 The sharp series arise from the transitions \(n s \rightarrow m p\). The \(s\) lines are unsplit so the splitting is due entirely to the \(p\) level. The frequency difference between sequent lines is \(\frac{\Delta E}{\hbar}\) and is the same for all lines of the sharp series. It is

\[
\frac{1}{\hbar}\left(\frac{2 \pi \hbar c}{\lambda_{1}}-\frac{2 \pi \hbar c}{\lambda_{2}}\right)=\frac{2 \pi c \Delta \lambda}{\lambda_{1} \lambda_{2}}
\]

Evaluation gives

\[
1.645 \times 10^{14} \mathrm{rad} / \mathrm{s}
\]

6.103 We shall ignore hyperfine interaction. The state with principal quantum number \(n=3\) has orbital angular momentum quantum number

\[
l=0,1,2
\]

The levels with these terms are \(3 S, 3 P, 3 D\). The total angular momentum is obtained by combining spin and angular momentum. For a single electron this leads to

\[
\begin{aligned}
& J=\frac{1}{2}, \text { if } L=0 \\
& J=L-\frac{1}{2} \text { and } L+\frac{1}{2} \text { if } L \neq 0
\end{aligned}
\]

We then get the final designations

\[
3 S_{\frac{1}{2}}, 3 P_{\frac{1}{2}}, 3 P_{3 / 2}, 3 D_{3 / 2}, 3 D_{5 / 2}
\]

6.104 The rule is that if \(\vec{J}=\vec{L}+\vec{S}\) then \(J\) takes the values

\[
|L-S| \text { to } L+S
\]

in step of 1 . Thus :\\
(a) The values are \(1,2,3,4,5\)\\
(b) The values are \(0,1,2,3,4,5,6\)\\
(c) The values are \(\frac{1}{2}, \frac{3}{2}, \frac{5}{2}, \frac{7}{2}, \frac{9}{2}\).\\
6.105 For the state \(4 p, L=1, S=\frac{3}{2}\) (since \(2 s+1=4\) ). For the state \(5 d, L=2, s=2\). The possible values of \(J\) are

\[
\begin{aligned}
& J: \frac{5}{2}, \frac{3}{2}, \frac{1}{2} \text { for } 4 p \\
& J: 4,3,2,1,0 \text { for } 5 d
\end{aligned}
\]

The value of the magnitude of angular momentum is \(\hbar \sqrt{J(J+1)}\). Substitution gives the values\\
\(4 P\) :

\[
\begin{gathered}
\hbar \sqrt{\frac{1}{2} \cdot \frac{3}{2}}=\frac{\hbar \sqrt{3}}{2}, \hbar \sqrt{\frac{3}{2} \cdot \frac{5}{2}}=\frac{\hbar \sqrt{15}}{2} \\
\hbar \sqrt{\frac{5}{2} \cdot \frac{7}{2}}=\frac{\hbar \sqrt{35}}{2} \\
5 D: 0, \hbar \sqrt{2}, \hbar \sqrt{6}, \hbar \sqrt{12}, \hbar \sqrt{20}
\end{gathered}
\]

and\\
6.106 (a) For the Na atoms the valence electron has principal quantem number \(n=4\), and the possible values of orbital angular momentum are \(l=0,1,2,3\) so \(l_{\max }=3\). The state is \({ }^{2} F\), maximum value of \(J\) is \(\frac{7}{2}\).\\
Thus the state with maximum angular momentum will be

For this state

\[
M_{\max }=\frac{{ }^{2} F_{7 / 2}}{\hbar} \sqrt{\frac{7}{2} \cdot \frac{9}{2}}=\frac{\hbar \sqrt{63}}{2}
\]

(b) For the atom with electronic configuration \(1 s^{2} 2 p 3 d\). There are two inequivalent valence electrons. The total orbital angular moments will be \(1,2,3\) so we pick \(l=3\). The total spin angular momentum will be \(s=0,1\) so we pick up \(s=1\). Finally \(J\) will be \(2,3,4\) so we pick up 4 . Thus maximum angular momentum state is

For this state

\[
M_{\max }=\hbar \sqrt{4 \times 5}=2 \hbar \sqrt{5} .
\]

6.107 For the \(f\) state \(L=3\), For the \(d\) state \(L=2\). Now if the state has spin \(s\) the possible angular momentum are

\[
|L-S| \text { to } L+S
\]

The number of \(J\) angular momentum values is \(2 S+1\) if \(L \geq S\) and \(2 L+1\) if \(L<S\). Since the number of states is 5 , we must have \(S \geq L=2\) for \(D\) state while \(S \leq 3\) and \(2 S+1=5\) in ply \(S=2\) for \(F\) state. Thus for the \(F\) state total spin angular momentum\\
while for \(D\) state

\[
\begin{gathered}
M_{s}=\hbar \sqrt{2 \cdot 3}=\hbar \sqrt{6} \\
M_{s} \geq \hbar \sqrt{6}
\end{gathered}
\]

6.108 Multiplicity is \(2 S+1\) so \(S=1\).

Total angular momentum is \(\hbar \sqrt{J(J+1)}\) so \(J=4\). Then

\[
L \text { must equal } 3,4,5
\]

in order that \(J=4\) may be included in

\[
|L-S| \text { to } L+S
\]

6.109 (a) Here \(J=2, L=2\). Then \(S=0,1,2,3,4\) and the multiplicities( \(2 S-1\) ) are

\[
1,3,5,7,9 .
\]

(b) Here \(J=3 / 2, L=1\) Then

\[
S=\frac{5}{2}, \frac{3}{2}, \frac{1}{2}
\]

and the mulitiplicities are

\[
6,4,2
\]

(c) Here \(J=1, L=3\). Then

\[
S=2,3,4
\]

and the multiplicities are

\[
5,7,9
\]

6.110 The total angular momentum is greatest when \(L, S\) are both greatest and add to form \(J\). Now for a triplet of \(s, p, d\) electrons\\
Maximum spin \(\rightarrow S=\frac{3}{2}\) corresponding to

\[
M_{s}=\hbar \sqrt{\frac{3}{2} \cdot \frac{5}{2}}=\frac{\hbar \sqrt{15}}{2}
\]

Maximum orbital angular momentum →

\[
L=3
\]

corresponding to

\[
M_{L}=\hbar \sqrt{\frac{3}{2} \frac{5}{2}}=\frac{\hbar \sqrt{15}}{2}
\]

Maximum total angular momentum \(\quad J=\frac{9}{2}\)\\
corresponding to

\[
M=\frac{h}{2} \sqrt{99}
\]

In vector model

\[
\vec{L}=\vec{J}-\vec{S}
\]

or in magnitude squared

Thus

\[
L(L+1) \hbar^{2}=J(J+1) \hbar^{2}+S(S+1) \hbar^{2}-2 \vec{J} \cdot \vec{S}
\]

Substitution gives

\[
\cos (\prec \vec{J}, \vec{S})=\frac{J(J+1)+S(S+1)-L(L+1)}{2 \sqrt{J(J+1)} \sqrt{S(S+1)}}
\]

\[
<(\vec{J}, \vec{S})=31 \cdot 1^{\circ} .
\]

6.111 Total angular momentum \(\hbar \sqrt{6}\) means \(J=2\). It is gives that \(S=1\).

This means that \(L=1,2\), or 3 . From vector model relation

\[
\begin{aligned}
L(L+1) \hbar^{2} & =6 \hbar^{2}+2 \hbar^{2}-2 \hbar^{2} \sqrt{6} \sqrt{2} \cos 73 \cdot 2^{\circ} \\
& =5 \cdot 998 \hbar^{2} \approx 6 \hbar^{2}
\end{aligned}
\]

Thus \(L=2\) and the spectral symbol of the state is

\[
{ }^{3} D_{2}
\]

6.112 In a system containing a \(p\) electron and a \(d\) electron

\[
\begin{gathered}
S=0,1 \\
L=1,2,3
\end{gathered}
\]

For \(S=0\) we have the terms

\[
{ }^{1} P_{1},{ }^{1} D_{2},{ }^{1} F_{3}
\]

For \(S=1\) we have the terms

\[
{ }^{3} P_{0},{ }^{3} P_{1},{ }^{3} P_{2},{ }^{3} D_{1},{ }^{3} D_{2},{ }^{3} D_{3},{ }^{3} F_{2},{ }^{3} F_{3},{ }^{3} F_{4}
\]

6.113 The atom has \(\$_{1}=1 / 2, l_{1}=1, j,=\frac{3}{2}\)

The electron has \(\$_{2}=\frac{1}{2}, l_{2}=2\) so the total angular momentum quantum number must be

\[
j_{2}=\frac{3}{2} \quad \text { or } \quad \frac{5}{2}
\]

In \(L-S\) compling we get \(S=0,1 . L=1,2,3\) and the terms that can be formed are the same as written in the problem above. The possible values of angular momentum are consistant with the addition \(j_{1}=\frac{3}{2}\) to \(j_{2}=\frac{3}{2}\) or \(\frac{5}{2}\).\\
The latter gives us \(\quad J=0,1,2,3 ; 1,2,3,4\)\\
All these values are reached above.\\
6.114 Selection rules are \(\quad \Delta S=0\)

\[
\begin{gathered}
\Delta L= \pm 1 \\
\Delta J=0, \pm 1(\text { no } 0 \rightarrow 0) \\
{ }^{2} D_{3 / 2} \rightarrow{ }^{2} P_{1 / 2} \text { is allowed } \\
{ }^{3} P_{1} \rightarrow{ }^{2} S_{1 / 2} \text { not allowed } \\
\rightarrow{ }^{3} P_{2} \text { is not allowed }(\Delta L=2) \\
\rightarrow{ }^{4} D_{5 / 2} \text { is allowed }
\end{gathered}
\]

6.115 For a \(3 d\) state of a \(L i\) atom, \(S=\frac{1}{2}\) because there is only one electron and \(L=2\).

The total degeneracy is

\[
g=(2 L+1)(2 S+1)=5 \times 2=10 .
\]

The states are \({ }^{2} D_{\frac{3}{2}}\) and \({ }^{2} D_{5 / 2}\) and we check that

\[
g=4+6=\left(2 \times \frac{3}{2}+1\right)+\left(2 \times \frac{5}{2}+1\right)
\]

6.116 The state with greatest possible total angular momentum are

For a \({ }^{2} P\) state

\[
J=\frac{1}{2}+1=\frac{3}{2} \text { i.e. }{ }^{2} P_{3 / 2}
\]

Its degeneracy is 4 .\\
For a \({ }^{3} D\) state

\[
J=1+2=3 \text { i.e. }{ }^{3} D_{3}
\]

Its degeneracy is \(2 \times 3+1=7\)\\
For a \({ }^{\mathbf{4}} \boldsymbol{F}\) state

\[
J=\frac{3}{2}+3=\frac{9}{2} \text { i.e. }{ }^{4} F_{\frac{4}{2}}
\]

Its degeneracy is

\[
2 \times \frac{9}{2}+1=10 .
\]

6.117 The degeneracy is \(2 J+1\). So we must have \(J=3\). From \(L=3 S\), we see that \(S\) must be an integer since \(L\) is integral and \(S\) can be either integral or half integral. If \(S=0\) then \(L=0\) but this is consistent with \(J=3\). For \(S \geq 2, L \geq 6\) and then \(J=3\). Thus the state is

\[
{ }^{3} F_{3}
\]

6.118 The order of filling is\\
\(K, L, M\) shells, then \(4 s^{2}, 3 d^{10}\) then \(4 p^{3}\). The electronic configuration of the element will be

\[
1 s^{2} 2 s^{2} 2 p^{6} 3 s^{2} 3 p^{6} 4 s^{2} 3 d^{10} 4 p^{3}
\]

(There must be three \(4 p\) electrons)\\
The number of electrons is \(Z=33\) and the element is \(A s\). (The \(3 d\) subshell must be filled before 4p fills up.)\\
6.119 (a) when the partially filled shell contains three \(p\) electrons, the total spin \(S\) must equal \(S=\frac{1}{2}\) or \(\frac{3}{2}\). The state \(S=\frac{3}{2}\) has maximum spin and is totally symmetric under exchange of spin lables. By Pauli's exclusion principle this implies that the angular part of the wavefunction must be totally anti symmetric. Since the angular part of the wave function a \(p\) electron is vector \(\vec{r}\), the total wavefunction of three \(p\) electrons is the totally antisymmetric combination of \(\overrightarrow{r_{1}}, \overrightarrow{r_{2}}\), and \(\overrightarrow{r_{3}}\). The only such combination is

\[
\overrightarrow{r_{1}} \cdot\left(\overrightarrow{r_{2}} \times \overrightarrow{r_{3}}\right)=\left|\begin{array}{lll}
x_{1} & x_{2} & x_{3} \\
y_{1} & y_{2} & y_{3} \\
z_{1} & z_{2} & z_{3}
\end{array}\right|
\]

This combination is a scalar and hence has \(L=0\). The spectral term of the ground state is then

\[
{ }^{4} S_{\frac{3}{2}} \text { since } J=\frac{3}{2}
\]

(b) We can think of four \(p\) electrons as consisting of a full \(p\) shell with two \(p\) holes. The state of maximum spin \(S\) is then \(S=1\). By Pauli's principle the orbital angular momentum part must be antisymmetric and can only have the form

\[
\overrightarrow{r_{1}} \times \overrightarrow{r_{2}}
\]

where \(\overrightarrow{r_{1}}, \overrightarrow{r_{2}}\) are the coordinates of holes. The result is harder to see if we do not use the concept of holes. Four \(p\) electrons can have \(S=0,1,2\) but the \(S=2\) state is totally symmetric. The corresponding angular wavefunction must be totally antisymmetric. But this is impossible : there is no quantity which is amtisymmetric in four vectors. Thus the maximum allowed \(S\) is \(S=1\). We can construct such a state by coupling the spins of electrons \(1 \& 2\) to \(S=1\) and of electrons \(3 \& 4\) to \(S=1\) and then coupling the resultant spin states to \(S=1\). Such a state is symmetric under the exchange of spins of 1 \& 2nd 3 and 4 but antisymmetric under the simultaneous exchange of \((1,2) \&(3,4)\). the con-\\
jugate angular wavefunction must be antisymmetric under the exchange of ( 1,2 ) and under the exchange of \((3,4)\) by Pauli principle. It must also be antisymmetric under the simultaneous exchange of \((1,2)\) and \((3,4)\). (This is because two exchanges of electrons are involved.) The required angular wavefunction then has the form

\[
\left(\overrightarrow{r_{1}} \times \overrightarrow{r_{2}}\right) \times\left(\overrightarrow{r_{3}} \times \overrightarrow{r_{4}}\right)
\]

and is a vector, \(L=1\). Thus, using also the fact that the shell is more than half full, we find the spectral term \({ }^{3} P_{2}\)

\[
(J=L+S) .
\]

6.120 (a) The maximum spin angular momentum of three electrons can be \(S=\frac{3}{2}\). This state is totally symmetric and hence the conjugate angular wavefunction must be antisymmetric By Pauli's exclusion principle the totally antisymmetric state must have different magnetic quantum numbers. It is easy to see that for d electrons the maximum value of the magnetic quantum number for orbital angular momentum \(\left|M_{L-z}\right|=3\) (from \(2+1+0\) ). Higher values violate Pauli's principle. Thus the state of highest orbital angular momentum consistent with Pauli's principle is \(L=3\).

The state of the atom is then \({ }^{4} F_{J}\) where \(J=L-S\) by Hund's rule. Thus we get

\[
{ }^{4} F_{3 / 2}
\]

The magnitude of the angular momentum is

\[
\hbar \sqrt{\frac{3}{2} \cdot \frac{5}{2}}=\frac{\hbar}{2} \sqrt{15}
\]

(b) Seven \(d\) electrons mean three holes. Then \(S=\frac{3}{2}\) and \(L=3\) as before. But \(J=L+S=\frac{9}{2}\) by Hund's rule for more than half filled shell. Thus the state is

\[
{ }^{4} F_{9 / 2}
\]

Total angular momentum has the magnitude

\[
\hbar \sqrt{\frac{9}{2} \cdot \frac{11}{2}}=\frac{3 \hbar}{2} \sqrt{11} .
\]

6.121 (a) \({ }^{3} F_{2}\) : The maximum value of spin is \(S=1\) here. This means there are 2 electrons. \(L=3\) so \(s\) and \(p\) electrons are ruled out. Thus the simplest possibility is \(d\) electrons. This is the correct choice for if we were considering \(f\) electrons, the maximum value of \(L\) allowed by Pauli principle will be \(L=5\) (maximum value of the magnitude of magnetic quantum number will be \(3+2=5\).)\\
Thus the atom has two \(d\) electrons in the unfilled shell.\\
(b) \({ }^{2} P_{3 / 2}\) Here \(L=1, S=\frac{1}{2}\) and \(J=\frac{3}{2}\)

Since \(J=L+S\), Hund's rule implies the shell is more than half full. This means one electron less than a full shell. On the basis of hole picture it is easy to see that we have \(p\) electrons. Thus the atom has \(5 p\) electrons.\\
(c) \({ }^{6} S_{5 / 2}\) Here \(S=\frac{5}{2}, L=0\). We either have five electrons or five holes. The angular part is antisymmetric. For five \(d\) electrons, the maximum value of the quantum number consistent with Pauli exclusion principle is \((2+1+0-1-2)=0\) so \(L=0\). For \(f\) or g electrons \(L>0\) whether the shell has five electrons or five holes. Thus the atom has five \(d\) electrons.\\
6.122 (a) If \(S=1\) is the maximum spin then there must be two electrons (If there are two holes then the shell will be more than half full.). This means that there are 6 electrons in the full shell so it is a \(p\) shell. By Paul's principle the only antisymmetric combination of two electrons has \(L=1\) Also \(J=L-S\) as the shell is less than half full. Thus the term is \({ }^{3} P_{0}\)\\
(b) \(S=\frac{3}{2}\) means either 3 electrons or 3 holes. As the shell is more than half full the former possibility is ruled out. Thus we must have seven d electrons. Then as in problem 6.120 we get the term \({ }^{\mathbf{4}} \boldsymbol{F}_{\mathbf{9} \boldsymbol{/} \mathbf{2}}\)\\
6.123 With three electrons \(S=\frac{3}{2}\) and the spin part is totally symmetric. It is given that the basic term has \(L=3\) so \(L=3\) is the state of highest orbital angular momentum. This is not possible with \(p\) electron so we must have \(d\) electrons for which \(L=3\) for 3 electrons. For three \(f, g\) electrons \(L>3\). Thus we have \(3 d\) electrons. Then as in (6.120) the ground state is

\[
{ }^{4} F_{\frac{3}{2}}
\]

6.124 We have \(5 d\) electrons in the only unfilled shell. Then \(S=\frac{5}{2}\). Maximum value of \(L\) consistent with Pauli's principle is \(L=0\). Then \(J=\frac{5}{2}\).\\
So by Lande's formula

Thus

\[
\begin{gathered}
g=1+\frac{\frac{5}{2}\left(\frac{7}{2}\right)+\frac{5}{2}\left(\frac{7}{2}\right)-0}{2 \frac{5}{2}\left(\frac{7}{2}\right)}=2 \\
\mu=g \sqrt{J(J+1)} \mu_{B}=2 \frac{\sqrt{35}}{2} \mu_{B}=2 \sqrt{35} \mu_{B} .
\end{gathered}
\]

The ground state is \({ }^{6} S_{5 / 2}\).\\
6.125 By Boltzmann formula

\[
\frac{N_{2}}{N_{1}}=\frac{g_{2}}{g_{1}} e^{-\Delta E / K T}
\]

Here \(\Delta E=\) energy difference between \(n=1\) and \(n=2\) states

\[
\begin{gathered}
=13.6\left(1-\frac{1}{4}\right) \mathrm{eV}=10.22 \mathrm{eV} \\
g_{1}=2 \text { and } g_{2}=8 \text { (counting } 2 \mathrm{~s} \& 2 P \text { states.) Thus } \\
\frac{N_{2}}{N_{1}}=4 e^{-10.22 \times 1.602 \times 10^{-19} / 1.38 \times 10^{-23} \times 3000}=2.7 \times 10^{-17} \\
\eta=\frac{N_{2}}{N_{1}}=n^{2} e^{-\Delta E_{n} / K T}, \Delta E_{n}=\hbar R\left(1-\frac{1}{n^{2}}\right)
\end{gathered}
\]

for the nth excited state because the degeneracy of the state with principal quantum number \(n\) is \(2 n^{2}\).\\
6.126 We have

\[
\frac{N}{N_{0}}=\frac{g}{g_{0}} e^{-h \omega / k T}=\frac{g}{g_{0}} e^{-2 \pi h c / \lambda k T}
\]

Here \(\mathrm{g}=\) degeneracy ofthe \(3 P\) sate \(=6, g_{0}=\) degeneracy of the \(3 S\) state \(=2\) and \(\lambda=\) wavelength of the \(3 P \rightarrow 3 S\) line \(\left(\frac{2 \pi \hbar c}{\lambda}=\right.\) energy difference between \(3 P \& 3 S\) levels.

Substitution gives

\[
\frac{N}{N_{0}}=1.13 \times 10^{-4}
\]

6.127 Let \(T=\) mean life time of the excited atoms. Then the number of excited atoms will decrease with time as \(e^{-t / T}\). In time \(t\) the atom travels a distance \(\mathrm{v} t\) so \(t=\frac{l}{\mathrm{v}}\). Thus the number of excited atoms in a beam that has traversed a distance \(l\) has decreased by

\[
e^{-L / v T}
\]

The intensity of the line is proportional to the number of excited atoms in the beam. Thus

\[
e^{-l / v \tau}=\frac{1}{\eta} \text { or } \tau=\frac{l}{v \ln \eta}=1.29 \times 10^{-6} \text { second. }
\]

6.128 As a result of the lighting by the mercury lamp a number of atoms are pumped to the excited state. In equilibrium the number of such atoms is \(N\). Since the mean life time of the atom is \(T\), the number decaying per unit time is \(\frac{N}{\tau}\). Since a photon of energy \(\frac{2 \pi \hbar c}{\lambda}\) results from each decay, the total radiated power will be \(\frac{2 \pi \hbar c}{\lambda} \cdot \frac{N}{\tau}\). This must equal \(P\). Thus

\[
N=P \tau / \frac{2 \pi \hbar c}{\lambda}=\frac{P \tau \lambda}{2 \pi h}=6.7 \times 10^{9}
\]

6.129 The number of excited atoms per unit volume of the gas in \(2 P\) state is

\[
N=n \frac{g_{p}}{g_{s}} e^{-2 \pi h c / \lambda k T}
\]

Here \(g_{p}=\) degeneracy of the \(2 p\) state \(=6, g_{s}=\) degeneacy of the \(2 s\) state \(=2\) and \(\lambda=\) wavelength of the resonant line \(2 p \rightarrow 2 s\). The rate of decay of these atoms is \(\frac{N}{\tau}\) per sec. per unit volume. Since each such atom emits light of wavelength \(\lambda\), we must have

\[
\frac{1}{\tau} \frac{2 \pi \hbar c}{\lambda} n \frac{g_{p}}{g_{s}} e^{-2 \pi h c / \lambda k T}=P
\]

Thus

\[
\tau=\frac{1}{P} \frac{2 \pi \hbar c}{\lambda} n \frac{g_{p}}{g_{s}} e^{-2 \pi h c / \lambda k T}=65.4 \times 10^{-9} \mathrm{~s}=65.4 \mathrm{~ns}
\]

6.130 (a) We know that

\[
\begin{aligned}
& P_{21}^{s p}=A_{21} \\
& P_{21}^{i n d}=B_{21} u_{a s} \\
&=\frac{\pi^{2} c^{3}}{\hbar \omega^{3}} A_{21} \cdot \frac{\hbar \omega^{3}}{\pi^{2} c^{3}} \frac{1}{e^{-h / k T}-1} \quad=\frac{A_{21}}{e^{-h \omega / k T}-1} \\
& \frac{P_{21}^{i n d}}{P_{21}^{s p}}=\frac{1}{e^{-h \omega / k T}-1}
\end{aligned}
\]

Thus

For the transition \(2 P \rightarrow 1 S \hbar \omega=\frac{3}{4} \hbar R\) and\\
we get

\[
\frac{P_{21}^{1 i n}}{P_{21}^{s p}} \approx e^{-h \omega k T}
\]

substitution gives \(7 \times 10^{-18}\)\\
(b) The two rates become equal when \(e^{-h \omega / k T}=2\)\\
or

\[
T=(\hbar \omega / k \ln 2)=1.71 \times 10^{5} \mathrm{~K}
\]

6.131 Because of the resonant nature of the processes we can ignore nonresonant processes. We also ignore spontaneous emission since it does not contribute to the absorption coefficient and is a small term if the beam is intense enough.\\
Suppose \(I\) is the intensity of the beam at some point. The decrease in the value of this intensity on passing through the layer of the substance of thickness \(d x\) is equal to

\[
-d I=X I d x=\left(N_{1} B_{12}-N_{2} B_{21}\right)\left(\frac{I}{c}\right) \hbar \omega d x
\]

Here \(N_{1}=\) No. of atoms in lower level\\
\(N_{2}=\) No of atoms in the upper level per unit volume.\\
\(B_{12}, B_{21}\) are Einstein coefficients and \(I_{c}=\) energy density in the beam, \(c=\) velocity of light.

A factor \(\hbar \omega\) arises because each transition result in a loss or gain of energy \(\hbar \omega\)

Hence

\[
x=\frac{\hbar \omega}{c} N_{1} B_{12}\left(1-\frac{N_{2} B_{21}}{N_{1} B_{12}}\right)
\]

But

\[
g_{1} B_{12}=g_{2} B_{21} \text { so }
\]

By Boltzman factor

\[
x=\frac{\hbar \omega}{c} N_{1} B_{12}\left(1-\frac{g_{1}}{g_{2}} \frac{N_{2}}{N_{1}}\right)
\]

\[
\frac{N_{2}}{N_{1}}=\frac{g_{2}}{g_{1}} e^{-h \omega / k T}
\]

When \(\hbar \omega \gg k T\) we can put \(N_{1}=N_{0}\) the total number of atoms per unit volume.\\
Then

\[
x=x_{0}\left(1-e^{-h \omega / k T}\right)
\]

where \(x_{0}=\frac{\hbar \omega}{c} N_{0} B_{12}\) is the absorption coefficient for \(T \rightarrow 0\).\\
6.132 A short lived state of mean life \(T\) has an uncertainty in energy of \(\Delta E \sim \frac{\hbar}{T}\) which is transmitted to the photon it emits as natural broadening. Then

\[
\Delta \omega_{n a t}=\frac{1}{T} \quad \text { so } \quad \Delta \lambda_{n a t}=\frac{\lambda^{2}}{2 \pi c \tau}
\]

The Döppler broadening on the other hand arises from the thermal motion of radiating atoms. The effect is non-relativistic and the maximum broadening can be written as

Thus

\[
\begin{aligned}
& \frac{\Delta \lambda_{\text {Dop }}}{\lambda}=2 \beta=\frac{2 v_{p r}}{c} \\
& \frac{\Delta \lambda_{\text {Dopp }}}{\Delta \lambda_{\text {nat }}}=\frac{4 \pi v_{p r} \tau}{\lambda}
\end{aligned}
\]

Substitution gives using \(v_{p r}=\sqrt{\frac{2 R T}{M}}=157 \mathrm{~m} / \mathrm{s}\),

\[
\frac{\Delta \lambda_{D o p p}}{\Delta \lambda_{n a t}} \approx 1.2 \times 10^{3}
\]

Note :- Our formula is an order of magnitude estimate.\\
6.133 From Moseley's law\\
or

\[
\omega_{K_{\mathrm{a}}}=\frac{3}{4} R(Z-1)^{2}
\]

\[
\lambda_{K_{\mathrm{a}}}=\frac{4}{3 R} \frac{1}{(Z-1)^{2}}
\]

Thus

\[
\frac{\lambda_{K_{\mathrm{a}}}(C u)}{\lambda_{K_{\mathrm{a}}}(F e)}=\left(\frac{25}{28}\right)^{2}=\left(\frac{Z_{F e}-1}{Z_{c u}-1}\right)^{2}
\]

Substitution gives

\[
\lambda_{K_{a}}(C u)=153.9 \mathrm{pm}
\]

6.134 (a) From Moseley's law\\
or

\[
\begin{gathered}
\omega_{k a}=\frac{3}{4} R(Z-\sigma)^{2} \\
\lambda_{K_{a}}=\frac{2 \pi c}{\omega_{K_{a}}}=\frac{8 \pi c}{3 R} \frac{1}{(Z-\sigma)^{2}}
\end{gathered}
\]

We shall take \(\sigma=1\). For Aluminium ( \(Z=13\) )

\[
\lambda_{K_{\mathrm{a}}}(A l)^{\prime}=843.2 \mathrm{pm}
\]

and for cobalt ( \(Z=27\) )

\[
\lambda_{K_{\mathrm{a}}}(C o)=179.6 \mathrm{pm}
\]

(b) This difference is nearly equal to the energy of the \(K_{\alpha}\) line which by Moseley's law is equal to ( \(Z=23\) for vanadium)

\[
\Delta E=\hbar \omega_{K_{\mathrm{a}}}=\frac{3}{4} \times 13.62 \times 22 \times 22=4.94 \mathrm{keV}
\]

6.135 We calculate the \(Z\) values corresponding to the given wavelengths using Moseley's law. Sec problem (134).\\
Substitution gives that\\
and

\[
\begin{aligned}
& Z=23 \text { corresponding to } \lambda=250 \mathrm{pm} \\
& Z=27 \text { corresponding to } \lambda=179 \mathrm{pm}
\end{aligned}
\]

There are thus three elements in a row between those whose wavelengths of \(K_{\alpha}\) lines are equal to \(\mathbf{2 5 0 ~ p m}\) and \(\mathbf{1 7 9 ~ p m}\).\\
6.136 From Moseley's law

\[
\lambda_{K_{\mathrm{a}}}(N i)=\frac{8 \pi c}{3 R} \frac{1}{(Z-1)^{2}}
\]

where \(Z=28\) for \(N i\). Substitution gives

\[
\lambda_{K_{\alpha}}(N i)=166.5 \mathrm{pm}
\]

Now the short wave cut of off the continuous spectrum must be more energetic (smaller wavelength) otherwise \(K_{\alpha}\) lines will not emerge. Then since \(\Delta \lambda=\lambda_{K_{\alpha}}-\lambda_{0}=84 \mathrm{pm}\) we get

\[
\lambda_{0}=82.5 \mathrm{pm}
\]

This corresponds to a voltage of

\[
V=\frac{2 \pi \hbar c}{e \lambda_{0}}
\]

Substitution gives \(V=15 \cdot 0 k \mathrm{~V}\)\\
6.137 Since the short wavelength cut off of the continuous spectrum is\\
the voltage applied must be

\[
\begin{gathered}
\lambda_{0}=0.50 \mathrm{~nm} \\
V=\frac{2 \pi \hbar c}{e \lambda_{0}}=2.48 \mathrm{kV}
\end{gathered}
\]

since this is greater than the excitation potential of the \(K\) series of the characteristic spectrum (which is only \(1.56 k \mathrm{~V}\) ) the latter will be observed.\\
6.138 Suppose \(\lambda_{0}=\) wavelength of the characteristic \(X\)-ray line. Then using the formula for short wavelength limit of continuous radiation

Hence

\[
\begin{gathered}
\frac{\lambda_{0}-\frac{2 \pi \hbar c}{e V_{1}}}{\lambda_{0}-\frac{2 \pi \hbar c}{e V_{2}}}=\frac{1}{n} \\
\lambda_{0}=\frac{2 \pi \hbar c}{e V_{1}} \frac{\left(n-\frac{V_{1}}{V_{2}}\right)}{n-1}
\end{gathered}
\]

Using also Moseley's law, we get

\[
Z=1+\sqrt{\frac{8 \pi c}{3 R \lambda}}=1+2 \sqrt{\frac{n-1}{3 \hbar R} \frac{e V_{1}}{n-\frac{V_{1}}{V_{2}}}}=29
\]

6.139 The difference in frequencies of the \(K\) and \(L\)\\
\includegraphics[max width=\textwidth, alt={}, center]{da0156b6-4133-434e-90d4-ec01b2070afe-331_436_630_1064_800}\\
absorption edges is equal, according to the Bohr picture, to the frequency of the \(K_{\alpha}\) line (see the diagram below). Thus by Moseley's formule

\[
\begin{array}{r}
\Delta \omega=\frac{3}{4} R(Z-1)^{2} \\
\text { or } \quad Z=1+\sqrt{\frac{4 \Delta \omega}{3 R}}=22
\end{array}
\]

The metal is titanium.\\
6.140 From the diagram above we see that the binding energy \(E_{b}\) of a \(K\) electron is the sum of the energy of a \(K_{\alpha}\) line and the energy corresponding to the \(L\) edge of absorption spectrum

\[
E_{b}=\frac{2 \pi \hbar c}{\lambda_{L}}+\frac{3}{4} \hbar R(Z 1)^{2}
\]

For vanadium \(Z=23\) and the energy of \(K_{\alpha}\) line of vanadium has been calculated in problem 134 (b). Using

\[
\frac{2 \pi \hbar c}{\lambda_{L}}=0.51 \mathrm{keV} \text { for } \lambda_{L}=2.4 \mathrm{~nm}
\]

we get

\[
E_{b}=5.46 \mathrm{keV}
\]

6.141 By Moseley's law

\[
\hbar \omega=\frac{2 \pi \hbar c}{\lambda} E_{K}-E_{L}=\frac{3}{4} \hbar R(Z-1)^{2}
\]

where \(-E_{K}\) is the energy of the \(K\) electron and \(-E_{L}\) of the \(L\) electron. Also the energy of the line corresponding to the short wave cut off of the \(K\) series is

Hence

\[
\begin{aligned}
E_{K} & =\frac{2 \pi \hbar c}{\lambda-\Delta \lambda}=\frac{2 \pi \hbar c}{\frac{2 \pi c}{\omega}-\Delta \lambda} \\
& =\frac{\hbar}{\frac{1}{\omega}-\frac{\Delta \lambda}{2 \pi c}}=\frac{\hbar \omega}{1-\frac{\omega \Delta \lambda}{2 \pi c}} \\
E_{L} & =\frac{\hbar \omega}{1-\frac{\omega \Delta \lambda}{2 \pi c}}-\hbar \omega=\frac{\hbar \omega}{\frac{2 \pi c}{\omega \Delta \lambda}-1}
\end{aligned}
\]

Substitution gives for titanium ( \(Z=22\) )

\[
\omega=6.85 \times 10^{18} s^{-1}
\]

and hence \(E_{L}=0.47 \mathrm{keV}\)\\
6.142 The energy of the \(K_{\alpha}\) radiation of \(Z n\) is

\[
\hbar \omega=\frac{3}{4} \hbar R(Z-1)^{2}
\]

where \(Z=\) atomic number of Zinc \(=30\). The binding energy of the \(K\) electrons in iron is obtained from the wavelength of \(K\) absorption edge as \(E_{K}=2 \pi \hbar c / \lambda_{K}\)\\
Hence by Einstein equation

\[
T=\frac{3}{4} \hbar R(Z-1)^{2}-\frac{2 \pi \hbar c}{\lambda_{K}}
\]

Substitution gives

\[
T=1.463 \mathrm{keV}
\]

This corresponds to a velocity of the photo electrons of

\[
v=2.27 \times 10^{6} \mathrm{~m} / \mathrm{s}
\]

6.143 From the Lande formula

\[
g=1+\frac{J(J+1)+S(S+1)-L(L+1)}{2 J(J+1)}
\]

(a) For \(S\) states \(L=0\). This implies \(J=S\). Then, if \(S \neq 0\)

\[
g=2
\]

(For singlet states g is not defined if \(L=0\) )\\
(b) For singlet states, \(J=L\)

\[
g=1+\frac{J(J+1)-L(L+1)}{2 J(J+1)}=1
\]

6.144 (a) \({ }^{6} F_{\frac{1}{2}}\) Here \(S=\frac{5}{2}, L=3, J=\frac{1}{2}\)

\[
g=1+\frac{\frac{3}{4}+\frac{35}{4}-12}{2 \times \frac{3}{4}}=1+\frac{38-48}{6}=-\frac{2}{3}
\]

(b) \({ }^{4} D_{1 / 2}\) : Here \(S=\frac{3}{2}, L=2, J=\frac{1}{2}\)

\[
g=1+\frac{\frac{3}{4}+\frac{15}{4}-6}{2 \times \frac{3}{4}}=1+\frac{18-24}{6}=0
\]

(c) \({ }^{5} F_{2}\) Here \(S=2, L=3, J=2\)

\[
g=1+\frac{6+6-12}{2 \times 6}=1
\]

(d) \({ }^{5} P_{1}\) Here \(S=2, L=1, J=1\)

\[
g=1+\frac{2+6-2}{2 \times 2}=\frac{5}{2}
\]

(e) \({ }^{3} P_{0}\). For states with \(J=0, L=S\) the \(g\) gactor is indeterminate.\\
6.145 (a) For the \({ }^{1} F\) state \(S=0, L=3=J\)

Hence

\[
\begin{aligned}
g & =1+\frac{3 \times 4-3 \times 4}{2 \times 3 \times 4}=1 \\
\mu & =\sqrt{3 \times 4} \mu_{B}=2 \sqrt{3} \mu_{B}
\end{aligned}
\]

(b) For the \({ }^{2} D_{3 / 2}\) state \(S=\frac{1}{2}, L=2, J=\frac{3}{2}\)

\[
g=1+\frac{\frac{15}{4}+\frac{3}{4}-6}{2 \times \frac{15}{4}}=1+\frac{18-24}{30}=\frac{4}{5}
\]

Hence

\[
\mu=\frac{4}{5} \sqrt{15 / 4} \mu_{B}=\frac{2}{5} \sqrt{15} \mu_{B}=2 \sqrt{\frac{3}{5}} \mu_{B} .
\]

(c) We have

\[
\begin{array}{ll} 
& \frac{4}{3}=1+\frac{J(J+1)+2-6}{2 J(J+1)} \\
\text { or } & \frac{4}{3} J(J+1)=J(J+1)-4 \\
\text { or } & J(J+1)=12 \Rightarrow J=3 \\
\text { Hence } & \mu=\frac{4}{3} \sqrt{12} \mu_{B}=\frac{8}{\sqrt{3}} \mu_{B}
\end{array}
\]

6.146 The expression for the projection of the magnetic moment is

\[
\mu_{Z}=g m_{J} \mu_{B}
\]

where \(m_{J}\) is the projection of \(\vec{J}\) on the \(Z\)-axis.\\
Maximum value of the \(\boldsymbol{m}_{\boldsymbol{J}}\) is \(\boldsymbol{J}\). Thus

\[
g J=4
\]

Since \(J=2\), we get \(g=2\). Now

\[
\begin{aligned}
2 & =1+\frac{J(J+1)+S(S+1)-L(L+1)}{2 J(J+1)} \\
& =1+\frac{6+S(S+1)-6}{2 \times 6}, \text { as } L=2 \\
& =1+\frac{S(S+1)}{12}
\end{aligned}
\]

Hence

\[
S(S+1)=12 \text { or } S=3
\]

Thus

\[
M_{S}=\hbar \sqrt{3 \times 4}=2 \sqrt{3} \hbar
\]

6.147 The angle between the angular momentum vector and the field direction is the least when the angular mommentum projection is maximum i.e. \(J \hbar\).\\
Thus

\[
J \hbar=\sqrt{J(J+1)} \hbar \cos 30^{\circ}
\]

or

\[
\sqrt{\frac{J}{J+1}}=\frac{\sqrt{3}}{2}
\]

Hence

\[
J=3
\]

Then

\[
g=1+\frac{3 \times 4+1 \times 2-2 \times 3}{2 \times 3 \times 4}=1+\frac{8}{24}=\frac{4}{3}
\]

and

\[
\mu=\frac{4}{3} \sqrt{3 \times 4} \mu_{B}=\frac{8}{\sqrt{3}} \mu_{B}
\]

6.148 For a state with \(n=3, l=2\). Thus the state with maximum angular momentum is

\[
{ }^{2} D_{5 / 2}
\]

Then

Hence

\[
\begin{aligned}
g & =1+\frac{\frac{5}{2} \times \frac{7}{2}+\frac{1}{2} \times \frac{3}{2}-2 \times 3}{2 \times \frac{5}{2} \times \frac{7}{2}} \\
& =1+\frac{35+3-24}{70}=1+\frac{1}{5}=\frac{6}{5} \\
\mu & =\frac{6}{5} \sqrt{\frac{5}{2} \times \frac{7}{2}} \mu_{B}=3 \sqrt{\frac{7}{5}} \mu_{B}
\end{aligned}
\]

6.149 To get the greatest possible angular momentum we must have \(S=S_{\max }=1\)

\[
L=L_{\max }=1+2=3 \text { and } J=L+S=4
\]

Then

\[
g=1+\frac{4 \times 5+1 \times 2-3 \times 4}{2 \times 4 \times 5}=1+\frac{10}{40}=\frac{5}{4}
\]

and

\[
\mu=\frac{5}{4} \sqrt{4 \times 5} \mu_{B}=\frac{5 \sqrt{5}}{2} \mu_{B}
\]

6.150 Since \(\mu=0\) we must have either \(J=0\) or \(g=0\). But \(J=0\) is incompatible with \(L=2\) and \(S=\frac{3}{2}\). Hence \(g=0\). Thus\\
or

\[
\begin{aligned}
& 0=1+\frac{J(J+1)+\frac{3}{2} \times \frac{5}{2}-2 \times 3}{2 J(J+1)} \\
& -3 J(J+1)=\frac{15}{4}-6=-\frac{9}{4}
\end{aligned}
\]

Hence

\[
J=\frac{1}{2}
\]

Thus

\[
M=\hbar \sqrt{\frac{1}{2} \times \frac{3}{2}}=\frac{\hbar \sqrt{3}}{2}
\]

6.151 From \(M=\hbar \sqrt{J+1}=\sqrt{2} \hbar\)\\
we find \(J=1\). From the zero value of the magnetic moment we find\\
or

\[
\begin{gathered}
1+\frac{1 \times 2 L(L+1)+2 \times 3}{2 \times 1 \times 2}=0 \\
1+\frac{-L(L+1)+8}{4}=0
\end{gathered}
\]

or

\[
12=L(L+1)
\]

Hence \(L=3\). The state is

\[
{ }^{5} F_{1}
\]

6.152 If \(\vec{M}\) is the total angular momentum vector of the atom then there is a magnetic moment

\[
\vec{\mu}_{m}=g \mu_{B} \vec{M} / \hbar
\]

associated with it; here \(g\) is the Lande factor. In a magnetic field of induction \(\vec{B}\), an energy

\[
H^{\prime}=-g \mu_{B} \vec{M} \cdot \vec{B} / \hbar
\]

is associated with it. This interaction term corresponds to a presession of the angular momentum vector because if leads to an equation of motion of the angular momentum vector of the form\\
where

\[
\begin{array}{r}
\frac{d \vec{M}}{d t}=\vec{\Omega} \times M \\
\vec{\Omega}=\frac{g \mu_{B} \vec{B}}{h}
\end{array}
\]

Using Gaussian unit expression of \(\mu_{B} \mu_{B}=0.927 \times 10^{-20}\) erg/gauss, \(B=10^{3}\) gauss \(h=1.054 \times 10^{-27} \mathrm{erg}\) sec and for the \({ }^{2} P_{3 / 2}\) state

\[
\begin{gathered}
g=1+\frac{\frac{3}{2} \times \frac{5}{2}+\frac{1}{2} \times \frac{3}{2}-1 \times 2}{2 \times \frac{3}{2} \times \frac{5}{2}}=1+\frac{1}{3}=\frac{4}{3} \\
\Omega=1.17 \times 10^{10} \mathrm{rad} / \mathrm{s}
\end{gathered}
\]

and\\
The same formula is valid in MKS units also But \(\mu_{B}=0.927 \times 10^{-23} \mathrm{~A} \cdot \mathrm{~m}^{2}, B=10^{-1} T\) and \(\hbar=1.054 \times 10^{-34}\) Joule sec. The answer is the same.\\
6.153 The force on an atom with magnetic moment \(\vec{\mu}\) in a magnetic field of induction \(\vec{B}\) is given by

\[
\vec{F}=(\vec{\mu} \cdot \vec{\nabla}) \vec{B}
\]

In the present case, the maximum force arise when \(\vec{\mu}\) is along the axis or close to it.\\
Then

\[
F_{Z}=\left(\mu_{Z}\right)_{\max } \frac{\partial B}{\partial Z}
\]

Here \(\left(\mu_{Z}\right)_{\max }=g \mu_{B} J\). The Lande factor \(g\) is for \({ }^{2} P_{1 / 2}\)

\[
g=1+\frac{\frac{1}{2} \times \frac{3}{2}+\frac{1}{2} \times \frac{3}{2}-1 \times 2}{2 \frac{1}{2} \times \frac{3}{2}}=1-\frac{1 / 2}{3 / 2}=\frac{2}{3} .
\]

and

\[
J=\frac{1}{2} \text { so }\left(\mu_{z}\right)_{\max }=\frac{1}{3} \mu_{B} \text {. }
\]

The magnetic field is given by

\[
B_{Z}=\frac{\mu_{0}}{4 \pi} \cdot \frac{2 I \pi r^{2}}{\left(r^{2}+Z^{2}\right)^{3 / 2}}
\]

or

\[
\begin{gathered}
\frac{\partial B_{Z}}{\partial Z}=-\frac{\mu_{0}}{4 \pi} 6 I \pi r^{2} \frac{Z}{\left(r^{2}+Z^{2}\right)^{5 / 2}} . \\
\left(\frac{\partial B_{Z}}{\partial Z}\right)_{Z=r}=\frac{\mu_{0}}{4 \pi} \frac{3 I \pi}{\sqrt{8} r^{2}} .
\end{gathered}
\]

Thus

\[
F=\frac{1}{3} \mu_{B} \frac{\mu_{0}}{4 \pi} \frac{3 \pi}{\sqrt{8}} \frac{I}{r^{2}}
\]

Substitution gives (using data in MKS units)

\[
F=4.1 \times 10^{-27} \mathrm{~N}
\]

6.154 The magnetic field at a distance \(r\) from a long current carrying wire is mostly tangential and given by

\[
B_{\varphi}=\frac{\mu_{0} I}{2 \pi r}=\frac{\mu_{0}}{4 \pi} \frac{2 I}{r} .
\]

The force on a magnetic dipole of moment \(\vec{\mu}\) due to this magnetic field is also tangential and has a magnitude

\[
\left(\vec{\mu} \cdot \nabla_{0}\right) B_{\varphi}
\]

This force is nonvanishing only when the component of \(\vec{\mu}\) along \(\vec{r}\) non zero. Then

\[
F=\mu_{r} \frac{\partial}{\partial r} B_{\varphi}=-\mu_{r} \frac{\mu_{0}}{4 \pi} \frac{2 I}{r^{2}}
\]

Now the maximum value of \(\mu_{r}= \pm \mu_{B}\). Thus the force is

\[
F_{\max }=\mu_{B} \frac{\mu_{0}}{4 \pi} \frac{2 I}{r^{2}}=2.97 \times 10^{-26} \mathrm{~N}
\]

6.155 In the homogeneous magnetic field the atom experinces a force

\[
F=g J \mu_{B} \frac{\partial B}{\partial Z}
\]

Depending on the sign of \(J\), this can be either upward or downward. Suppose the latter is true. The atom then traverses first along a parabola inside the field and, once outside, in a straight line. The total distance between extreme lines on the screen will be

\[
\delta=2 g J \mu_{B} \frac{\partial B}{\partial Z}\left\{\frac{1}{2}\left(\frac{l_{1}}{v}\right)^{2}+\frac{l_{1}}{v} \cdot \frac{l_{2}}{v}\right\} / m_{V}
\]

Here \(m_{V}\) is the mass of the vanadium atom. (The first term is the displacement within the field and the second term is the displacement due to the transverse velocity acquired in the magnetic field).\\
Thus using \(\frac{1}{2} m_{V} \mathrm{v}^{2}=T\)\\
\includegraphics[max width=\textwidth, alt={}, center]{da0156b6-4133-434e-90d4-ec01b2070afe-337_288_610_1806_544}\\
we get

\[
\frac{\partial B}{\partial Z}=\frac{2 T \delta}{g \mu_{B} J l_{1}\left(l_{1}+2 l_{2}\right)}
\]

For vanadium atom in the ground state \({ }^{\mathbf{4}} \boldsymbol{F}_{\mathbf{3} \boldsymbol{/} \mathbf{2}}\).

\[
g=1+\frac{\frac{3 \times 5}{4}+\frac{3 \times 5}{4}-3 \times 4}{2 \times \frac{3 \times 5}{4}}=1+\frac{30-48}{30}=1-\frac{18}{30}=\frac{2}{5}
\]

\(J=\frac{3}{2}\), using other data, and substituting\\
we get

\[
\frac{\partial B}{\partial Z}=1.45 \times 10^{13} \mathrm{G} / \mathrm{cm}
\]

This value differs from the answer given in the book by almost a factor of \(10^{9}\). For neutral atoms in stern Gerlach experiments, the value \(T=22 \mathrm{MeV}\) is much too large. A more appropriate value will be \(T=22 \mathrm{meV}\) i.e. \(10^{9}\) times smaller. Then one gets the right answer.\\
6.156 (a) The term \(3 P_{0}\) does not split in weak magnetic field as it has zero total angular momentum.\\
(b) The term \({ }^{2} F_{5 / 2}\) will split into \(2 \times \frac{5}{2}+1=6\) sublevels. The shift in each sublevel is given by

\[
\Delta E=-g \mu_{B} M_{Z} B
\]

where \(M_{j}=-J(J-1), \ldots, J\) and \(g\) is the Landi factor

\[
g=1+\frac{\frac{5 \times 7}{4}+\frac{1 \times 3}{4}-3 \times 4}{2 \times \frac{5 \times 7}{4}}=1+\frac{38-48}{70}=\frac{6}{7}
\]

(c) In this case for the \({ }^{4} D_{1 / 2}\) term

\[
g=1+\frac{\frac{1 \times 3}{4}+\frac{3 \times 5}{4}-2 \times 3}{2 \times \frac{1 \times 3}{4}}=1+\frac{3+15-24}{6}=1-1=0
\]

Thus the energy differences vanish and the level does not split.\\
6.157 (a) For the \({ }^{1} D_{2}\) term

\[
g=1+\frac{2 \times 3+0-2 \times 3}{2 \times 2 \times 3}=1
\]

and

\[
\Delta E=-\mu_{B} M_{J} B
\]

\(M_{J}=-2,-1,0,+1,+2\). Thus the splitting is

\[
\delta E=4 \mu_{B} B
\]

Substitution gives \(\delta E=57.9 \mu \mathrm{eV}\)\\
(b) For the \({ }^{3} F_{4}\) term \(g=1+\frac{4 \times 5+1 \times 2-3 \times 4}{2 \times 4 \times 5}=1+\frac{10}{40}=\frac{5}{4}\).\\
and

\[
\Delta=-\frac{5}{4} \mu_{B} B M_{J}
\]

where

\[
M_{J}=-4 \text { to }+4 . \text { Thus }
\]

\[
\delta E=\frac{5}{4} \mu_{B} B \times 8=10 \mu_{B} B\left(=2 g J \mu_{B}\right)
\]

Substitution gives \(\delta E=144.7 \mu \mathrm{eV}\)\\
6.158 (a) The term \({ }^{1} P_{1}\) splits into 3 lines with \(M_{Z}= \pm 1,0\) in accordance with the formula\\
where

\[
\begin{gathered}
\Delta E=-g \mu_{B} B M_{Z} \\
g=1+\frac{1 \times 2+0-1 \times 2}{2 \times 1 \times 2}=1
\end{gathered}
\]

The term \({ }^{1} S_{0}\) does not split in weak magnetic field. Thus the transitions between \({ }^{1} P_{1} \&{ }^{1} S_{0}\) will result in 3 lines i.e. a normal Zeeman triplet.\\
(b) The term \({ }^{2} D_{5 / 2}\) will split ot in 6 terms in accordance with the formula

\[
\begin{aligned}
& \Delta E=-g \mu_{B} M_{Z} \\
& M_{Z}= \pm \frac{5}{2}, \pm \frac{3}{2}, \pm \frac{1}{2}, \text { and } \\
& g=1+\frac{5 \times 7+1 \times 3-4 \times 2 \times 3}{2 \times 5 \times 7}=\frac{6}{5}
\end{aligned}
\]

\begin{center}
\includegraphics[max width=\textwidth, alt={}]{da0156b6-4133-434e-90d4-ec01b2070afe-339_252_307_818_1039}
\end{center}

Ther term \({ }^{2} P_{3 / 2}\) will also split into 4 lines in accordance with the above formula with

\[
M_{Z}= \pm \frac{3}{2}, \pm \frac{1}{2} \text { and } g=1+\frac{3 \times 5+1 \times 3-4 \times 1 \times 2}{2 \times 3 \times 5}=\frac{4}{3}
\]

It is seen that the \(Z\) eeman splitting is auomalous as \(g\) factors are different.\\
(c) \({ }^{3} D_{1} \rightarrow{ }^{3} P_{0}\)

The term \({ }^{3} D_{1}\) splits into 3 levels \((g=5 / 2)\)\\
The term \({ }^{3} P_{0}\) does not split. Thus the \(Z\) eeman spectrum is normal.\\
(d) For the \(5 I_{5}\) term

\[
\begin{aligned}
g & =1+\frac{5 \times 6+2 \times 3-6 \times 7}{2 \times 5 \times 6} \\
& =1+\frac{36-42}{60}=1-\frac{1}{10}=\frac{9}{10}
\end{aligned}
\]

For the \({ }^{\mathbf{5}} \boldsymbol{H}_{\mathbf{4}}\) term

\[
g=1+\frac{4 \times 5+2 \times 3-5 \times 6}{2 \times 4 \times 5}=1+\frac{26-30}{40}=\frac{9}{10}
\]

We see that the splitting in the two levels given by \(\Delta E=-g \mu_{B} B M_{Z}\) is the same though the number of levels is different (11 and 9). It is then easy to see that only the lines with following energies occur

\[
\hbar \omega_{0}, \hbar \omega_{0} \pm g \mu_{B} B .
\]

The \(Z\) eeman pattern is normal\\
6.159 For a singlet term \(S=0, L=J, g=1\)

Then the total splitting is

\[
\delta E=2 J \mu_{B} B
\]

Substitution gives

\[
J=3\left(=\delta E / 2 \mu_{B} B\right)
\]

The term is \({ }^{1} F_{3}\).\\
6.160 As the spectral line is caused by transition between singlet terms, the \(Z\) eeman effect will be normal (since \(g=1\) for both terms). The energy difference between extreme components of the line will be \(2 \mu_{B} B\).This must equal

\[
\begin{aligned}
& -\Delta\left(\frac{2 \pi \hbar c}{\lambda}\right)=\frac{2 \pi \hbar c \Delta \lambda}{\lambda^{2}} \\
& \Delta \lambda=\frac{\mu_{B} B \lambda^{2}}{\pi \hbar c}=35 \mathrm{pm} .
\end{aligned}
\]

Thus\\
6.161 From the previous problem, if the components are \(\lambda, \lambda \pm \Delta \lambda\), then

\[
\frac{\lambda}{\Delta \lambda}=\frac{2 \pi \hbar c}{\mu_{B} B \lambda}
\]

For resolution \(\frac{\lambda}{\Delta \lambda} \leq R=\frac{\lambda}{\delta \lambda}\) of the instrument.\\
Thus \(\quad \frac{2 \pi \hbar c}{\mu_{B} B \lambda} \leq R\) or \(B \geq \frac{2 \pi \hbar c}{\mu_{B} \lambda R}\)\\
Hence the minimum megnetic induction is

\[
B_{\min }=\frac{2 \pi \hbar c}{\mu_{B} \lambda R}=4 k G=0.4 \mathrm{~T}
\]

6.162 The \({ }^{3} P_{0}\) term does not split. The \({ }^{3} D_{1}\) term splits into 3 lines corresponding to the shift.

\[
\Delta E=-g \mu_{B} B M_{Z}
\]

with \(M_{Z}= \pm 1,0\). The interval between neighbouring components is then given by

\[
\hbar \Delta \omega=g \mu_{B} B
\]

Hence

\[
B=\frac{\hbar \Delta \omega}{g \mu_{B}}
\]

Now for the \({ }^{3} D_{1}\) term

\[
g=1+\frac{1 \times 2+1 \times 2-2 \times 3}{2 \times 1 \times 2}=1+\frac{4-6}{4}=\frac{1}{2} .
\]

Substitution gives \(B=3.00 \mathrm{kG} .=0.3 \mathrm{~T}\).\\
6.163 (a) For the \({ }^{2} P_{3 / 2}\) term

\[
g=1+\frac{\frac{3}{2} \times \frac{5}{2}+\frac{1}{2} \times \frac{3}{2}-1 \times 2}{2 \times \frac{3}{2} \times \frac{5}{2}}=1+\frac{10}{30}=\frac{4}{3}
\]

and the energy of the \({ }^{2} P_{3 / 2}\) sublevels will be

\[
E\left(M_{Z}\right)=E_{0}-\frac{4}{3} \mu_{B} B M_{Z}
\]

where \(M_{Z}= \pm \frac{3}{2}, \pm \frac{1}{2}\). Thus, between neighbouring sublevels.

\[
\delta E\left({ }^{2} P_{3 / 2}\right)=\frac{4}{3} \mu_{B} B
\]

For the \({ }^{2} P_{1 / 2}\) terms

\[
\begin{aligned}
g & =1+\frac{\frac{1}{2} \times \frac{3}{2}+\frac{1}{2} \times \frac{3}{2}-1 \times 2}{2 \times \frac{1}{2} \times \frac{3}{2}} \\
& =1+\frac{6-8}{6}=1-\frac{1}{3}=\frac{2}{3}
\end{aligned}
\]

and the separation between the two sublevels into which the \({ }^{2} P_{1 / 2}\) term will split is

\[
\delta E\left({ }^{2} P_{1 / 2}\right)=\frac{2}{3} \mu_{B} B
\]

The ratio of the two splittings is \(2: 1\).\\
(b) The interval between neighbouring Zeeman sublevels of the \({ }^{2} P_{3 / 2}\) term is \(\frac{4}{3} \mu_{B} B\). The energy separation between \(D_{1}\) and \(D_{2}\) lines is \(\frac{2 \pi \hbar c}{\lambda^{2}} \Delta \lambda\) (this is the natural separation of the \({ }^{\mathbf{2}} \boldsymbol{P}\) them)

Thus

\[
\frac{4}{3} \mu_{B} B=\frac{2 \pi \hbar c \Delta \lambda}{\lambda^{2} \eta}
\]

or

\[
B=\frac{3 \pi \hbar c \Delta \lambda}{2 \mu_{B} \lambda^{2} \eta}
\]

Substitution gives

\[
B=5.46 \mathrm{k} \mathrm{G}
\]

6.164 For the \({ }^{2} P_{3 / 2}\) level \(g=4 / 3\) (see above) and the energies of sublevels are

\[
E^{\prime}=E_{0}^{\prime}-\frac{4}{3} \mu_{B} B M_{-z}^{\prime}
\]

where \(M_{z}^{\prime}= \pm \frac{3}{2}, \pm \frac{1}{2}\) for the four sublevels\\
For the \({ }^{2} S_{\frac{1}{2}}\) level, \(g=2(\) since \(L=0)\) and

\[
E=E_{0}-2 \mu_{B} B M_{Z}
\]

where

\[
M_{Z}= \pm \frac{1}{2}
\]

Permitted transitions must have

\[
\Delta M_{Z}=0, \pm 1
\]

Thus only the following transitions occur

\[
\begin{aligned}
\begin{array}{c}
\frac{3}{2} \\
\rightarrow-3 / 2
\end{array} \rightarrow-\frac{1}{2} \\
\left.\begin{array}{l}
\frac{1}{2} \rightarrow \frac{1}{2} \\
-\frac{1}{2} \rightarrow-\frac{1}{2}
\end{array}\right\} \Delta \omega= \pm \mu_{B} B / \hbar=3.96 \times 10^{10} \mathrm{rad} / \mathrm{s} \\
\left.\begin{array}{r}
\frac{1}{2} \rightarrow-\frac{1}{2} \\
-\frac{1}{2} \rightarrow \frac{1}{2}
\end{array}\right\} \Delta \omega= \pm \frac{1}{3} \mu_{B} B / \hbar=1.32 \times 10^{10} \mathrm{rad} / \mathrm{s}
\end{aligned}
\]

These six lines are shown below\\
\includegraphics[max width=\textwidth, alt={}, center]{da0156b6-4133-434e-90d4-ec01b2070afe-342_630_652_1297_375}\\
6.165 The difference arises because of different selection rules in the two cases. In (1) the line is emitted perpendicular to the field. The selection rules are then

\[
\Delta M_{Z}=0, \pm 1
\]

In (2) the light is emitted along the direction of the field. Then the selection rules are

\[
\begin{gathered}
\Delta M_{Z}= \pm 1 \\
\Delta M_{Z}=0 \text { is forbidden. }
\end{gathered}
\]

(a) In the transition \({ }^{2} P_{3 / 2} \rightarrow{ }^{2} S_{1 / 2}\)

This has been considered above. In (1) we get all the six lines shown in the problem above\\
In (2) the line corresponding to \(\frac{1}{2} \rightarrow \frac{1}{2}\) and \(-\frac{1}{2} \rightarrow-\frac{1}{2}\) is forbidden.\\
Then we get four lines\\
(b) \({ }^{\mathbf{3}} P_{\mathbf{2}} \rightarrow{ }^{\mathbf{3}} S_{1}\)

For the \({ }^{3} P_{2}\) level, \(g=1+\frac{2 \times 3+1 \times 2-1 \times 2}{2 \times 2 \times 3}=\frac{3}{2}\)\\
so the energies of the sublevels are\\
where

\[
\begin{gathered}
E^{\prime}\left(M_{z}^{\prime}\right)=E_{0}^{\prime}-\frac{3}{2} \mu_{B} B M_{z}^{\prime} \\
M_{z}^{\prime}= \pm 2, \pm 1,0
\end{gathered}
\]

For the \({ }^{3} S_{1}\) line, \(g=2\) and the energies of the sublevels are

\[
E\left(M_{Z}\right)=E_{0}-2 \mu_{B} B M_{Z}
\]

where

\[
M_{z}= \pm 1,0 . \text { The lines are }
\]

\[
\begin{gathered}
\Delta M_{Z}=M_{Z}-M_{Z}=+1:-2 \rightarrow-1,-1 \rightarrow 0 \text { and } 0 \rightarrow 1 \\
\Delta M_{Z}=0-1 \rightarrow-1,0 \rightarrow 0,1 \rightarrow 1 \\
\Delta M_{Z}=1,2 \rightarrow 1,1 \rightarrow 0,0 \rightarrow-1
\end{gathered}
\]

All energy differences are unequal because the two \(g\) values are unequal. There are then nine lines if viewed along (1) and Six lines if viewed along (2).\\
6.166 For the two levels

\[
\begin{aligned}
E_{0}^{\prime} & =E_{0}-g^{\prime} \mu_{B}^{\prime} M_{Z}^{\prime} B \\
E_{0} & =E_{0}-g \mu_{B} M_{Z} B
\end{aligned}
\]

and hence the shift of the component is the value of

\[
\Delta \omega=\frac{\mu_{B} B}{\hbar}\left[g^{\prime} M_{Z}^{\prime}-g M_{Z}\right]
\]

subject to the selection rule \(\Delta M_{Z}=0, \pm 1\). For \({ }^{3} D_{3}\)

\[
g^{\prime}=1+\frac{3 \times 4+1 \times 2-2 \times 3}{2 \times 3 \times 4}=1+\frac{8}{24}=\frac{4}{3}
\]

For \({ }^{\mathbf{3}} \boldsymbol{P}_{\mathbf{2}}\),

\[
g=1+\frac{2 \times 3+1 \times 2-1 \times 2}{2 \times 2 \times 3}=\frac{3}{2}
\]

Thus

\[
\Delta \omega=\frac{\mu_{B} B}{\hbar}\left|\frac{4}{3} M_{z}^{\prime}-\frac{3}{2} M_{z}\right|
\]

For the different transition we have the following table

\[
M_{z}^{\prime} g^{\prime}-M_{Z} g
\]

\begin{center}
\begin{tabular}{|l|l|l|l|}
\hline
\(3 \rightarrow 2\) & \(\mu_{B} B\) & \(\mathbf{0} \rightarrow \mathbf{1}\) & \(-\frac{3}{2} \mu_{B} B\) \\
\hline
\(2 \rightarrow 2\) & \(-\frac{1}{3} \mu_{B} B\) & \(0 \rightarrow 0\) & 0 \\
\hline
\(2 \rightarrow 1\) & \(7 / 6 \mu_{B} B\) & \(0 \rightarrow-1\) & \(3 / 2 \mu_{B} B\) \\
\hline
\(1 \rightarrow 2\) & \(-5 / 3 \mu_{B} B\) & \(-1 \rightarrow 0\) & \(-4 / 3 \mu_{B} B\) \\
\hline
\(1 \rightarrow 1\) & \(-\mathbf{1} / \mathbf{6} \mu_{B} B\) & \(-1 \rightarrow-1\) & \(1 / 6 \mu_{B} B\) \\
\hline
\(1 \rightarrow 0\) & \(4 / 3 \mu_{B} B\) & \(-1 \rightarrow-2 \rightarrow\) & \(5 / 3 \mu_{B} B\) \\
\hline
 &  & \(-2 \rightarrow-1 \rightarrow\) & \(-7 / 6 \mu_{B} B\) \\
\hline
 &  & \(-2 \rightarrow-2 \rightarrow\) & \(1 / 3 \mu_{B} B\) \\
\hline
 &  & \(-3 \rightarrow-2 \rightarrow\) & \(-\mu_{B} B\) \\
\hline
\end{tabular}
\end{center}

There are 15 lines in all.\\
The lines farthest out are \(1 \rightarrow 2\) and \(-1 \rightarrow-2\).\\
The splitting between them is the total splitting. It is

\[
\Delta \omega=\frac{10}{3} \mu_{B} B / \hbar
\]

Substitution gives \(\Delta \omega=7.8 \times 10^{10} \mathrm{rad} / \mathrm{sec}\).

\subsection*{6.4 MOLECULES AND CRYSTALS}
6.167 In the first excited rotational level \(J=1\)\\
so

\[
\begin{gathered}
E_{J}=1 \times 2 \frac{\hbar^{2}}{2 I}=\frac{1}{2} I \omega^{2} \text { classically } \\
\omega=\sqrt{2} \frac{\hbar}{I} \\
I=\Sigma m_{i} r_{i}^{2}=\frac{m}{2} \frac{d^{2}}{4}+\frac{m}{2} \frac{d^{2}}{4}=m \frac{d^{2}}{4}
\end{gathered}
\]

Now\\
Thus\\
where \(m\) is the mass of the mole cub and \(r_{i}\) is the distance of the atom from the axis.\\
Thus

\[
\omega=\frac{4 \sqrt{2} \hbar}{m d^{2}}=1.56 \times 10^{11} \mathrm{rad} / \mathrm{s}
\]

6.168 The axis of rotation passes through the centre of mass of the HCl molecule. The distances of the two atoms from the centre of mass are

\[
d_{H}=\frac{m_{C l}}{m_{H C l}} d, d_{C l}=\frac{m_{H}}{m_{H C l}} d
\]

Thus \(I=\) moment of inertia about the axis

\[
=\frac{4}{2} m_{H} d_{H}^{2}+m_{C l} d_{C l}^{2}=\frac{m_{H} m_{C l}}{m_{H}+m_{C l}} d^{2}
\]

The energy difference between two neighbouring levels whose quantum numbers are \(J \& J-1\) is

\[
\frac{\hbar^{2}}{2 I} \cdot 2 J=\frac{J \hbar^{2}}{I}=7.86 \mathrm{meV}
\]

Hence \(J=3\) and the levels have quantum numbers \(2 \& 3\).\\
6.169 The angular momentum is \(\sqrt{2 I E}=M\)

Now

\[
I=\frac{m d^{2}}{4}\left(m=\text { mass of } O_{2} \text { molecule }\right)=1.9584 \times 10^{-39} \mathrm{gm} \mathrm{c}^{2}
\]

So

\[
M=3.68 \times 10^{-27} \mathrm{erg} \mathrm{sec} .=3.49 \hbar
\]

(This corresponds to \(J=3\) )\\
6.170 From \(E_{J}=\frac{\hbar^{2}}{2 I} J(J+1)\)\\
and the selection rule \(\Delta J=1\) or \(J \rightarrow J-1\) for a pure rotational spectrum we get

\[
\omega(J, J-1)=\frac{\hbar J}{I}
\]

Thus transition lines are equispaced in frequency \(\Delta \omega=\frac{\hbar}{I}\).

In the case of CH molecule

\[
I=\frac{\hbar}{\Delta \omega}=1.93 \times 10^{-40} \mathrm{gm} \mathrm{~cm}^{2}
\]

Also

\[
I=\frac{m_{c} m_{H}}{m_{c}+m_{H}} d^{2}
\]

so

\[
d=1.117 \times 10^{-8} \mathrm{~cm}=111.7 \mathrm{pm}
\]

6.171 If the vibrational frequency is \(\omega_{0}\) the excitation energy of the first vibrational level will be \(\hbar \omega_{0}\). Thus if there are \(J\) rotational levels contained in the band between the ground state and the first vibrational excitation, then

\[
\hbar \omega_{0}=\frac{J(J+1) \hbar^{2}}{2 I}
\]

where as stated in the problem we have ignored any coupling between the two. For \(H F\) molecule

Then

\[
I=\frac{m_{H} m_{F}}{m_{H}+m_{F}} d^{2}=1.336 \times 10^{-4} \mathrm{gm} \mathrm{~cm}^{2}
\]

\[
J(J+1)=\frac{2 I \omega_{0}}{\hbar}=197.4
\]

For \(J=14, J(J+1)=210\). For \(J=13, J(J+1)=182\). Thus there lie 13 levels between the ground state and the first vibrational excitation.\\
6.172 We proceed as above. Calculating \(\frac{2 I \omega_{0}}{\hbar}\) we get

\[
\frac{2 I \omega_{0}}{\hbar} \cong 1118
\]

Now this must equal \(J(J+1) \approx\left(J+\frac{1}{2}\right)^{2}\)\\
Taking the square root we get \(J \approx 33\).\\
6.173 From the formula\\
or

\[
J(J+1) \frac{\hbar^{2}}{2 I}=E \text { we get } J(J+1)=2 I E / \hbar^{2}
\]

or

\[
\left(J+\frac{1}{2}\right)^{2}-\frac{1}{4}=\frac{2 I E}{\hbar^{2}}
\]

Hence

\[
J=-\frac{1}{2}+\sqrt{\frac{1}{4}+\frac{2 I E}{\hbar^{2}}}
\]

writing

\[
J+1=-\frac{1}{2}+\sqrt{\frac{1}{4}+\frac{2 I}{\hbar^{2}}(E+\Delta E)}
\]

we find

\[
\begin{aligned}
1 & =\sqrt{\frac{1}{4}+\frac{2 I}{\hbar^{2}} E+\frac{2 I}{\hbar^{2}} \Delta E}-\sqrt{\frac{1}{4}+\frac{2 I E}{\hbar^{2}}} \\
& =\sqrt{\frac{1}{4}+\frac{2 I}{\hbar^{2}} E}\left[\left(1+\frac{\Delta E}{E+\frac{\hbar^{2}}{8 I}}\right)^{1 / 2}-1\right] \\
& \approx \sqrt{\frac{1}{4}+\frac{2 I}{\hbar^{2}} E} \cdot \frac{\Delta E}{2\left(E+\frac{\hbar^{2}}{8 I}\right)} \\
& =\sqrt{\frac{2 I}{\hbar^{2}}} \frac{\Delta E}{2 \sqrt{E+\frac{\hbar^{2}}{8 I}}}
\end{aligned}
\]

The quantity \(\frac{d N}{d E}\) is \(\frac{1}{\Delta E}\). For large \(E\) it is

\[
\frac{d N}{d E}=\sqrt{\frac{I}{2 \hbar^{2} E}}
\]

For an iodine molecule

\[
I=m_{I} d^{2} / 2=7.57 \times 10^{-38} \mathrm{gm} \mathrm{~cm}^{2}
\]

Thus for \(J=10\)

\[
\frac{d N}{d E}=\sqrt{\frac{I}{2 \hbar^{2} \cdot \frac{\hbar^{2}}{2 I} J(J+1)}}=\frac{I}{\sqrt{J(J+1)} \hbar^{2}}
\]

Substitution gives

\[
\frac{d N}{d E}=1.04 \times 10^{4} \text { levels per } \mathrm{eV}
\]

6.174 For the first rotational level

\[
E_{r o t}=2 \frac{\hbar^{2}}{2 I}=\frac{\hbar^{2}}{I} \text { and }
\]

for the first vibrational level \(E_{v i b}=\hbar \omega\)\\
Thus

\[
\xi=\frac{E_{v i b}}{E_{r o f}}=\frac{I \omega}{h}
\]

Here \(\omega=\) frequency of vibration. Now

\[
I=\mu d^{2}=\frac{m_{1} m_{2}}{m_{1}+m_{2}} d^{2}
\]

(a) For \(H_{2}\) molecule \(I=4.58 \times 10^{-41} \mathrm{gm} \mathrm{cm}^{2}\) and \(\xi=36\)\\
(b) For HI molecule,

\[
I=4.247 \times 10^{-40} \mathrm{gm} \mathrm{~cm}^{2} \text { and } \xi=175
\]

(c) For \(I_{2}\) molecule

\[
I=7.57 \times 10^{-38} \mathrm{gm} \mathrm{~cm}^{2} \text { and } \xi=2872
\]

6.175 The energy of the molecule in the first rotational level will be \(\frac{\hbar^{2}}{I}\).The ratio of the number of molecules at the first excited vibrational level to the number of molecules at the first excited rotational level is

\[
\begin{gathered}
\frac{e^{-\hbar \omega / k T}}{(2 J+1) e^{-\hbar^{2} J(J+1) / 2 I k T}} \\
=\frac{1}{3} e^{-\hbar \omega / k T} \times e^{-\hbar^{2} / 1 k T}=\frac{1}{3} e^{-\hbar(\omega-2 B / k T} \\
B=\hbar / 2 I
\end{gathered}
\]

where\\
For the hydrogen molecule \(I=\frac{1}{2} m_{H} d^{2}\)

\[
=4.58 \times 10^{-41} \mathrm{gm} \mathrm{~cm}^{2}
\]

Substitution gives \(3.04 \times 10^{-4}\)\\
6.176 By definition

\[
\begin{aligned}
<E> & =\frac{\sum E_{v} e^{-E_{v} / k T}}{\sum \exp \left(-E_{v} / k T\right)}=\frac{\frac{\partial}{\partial \beta} \sum_{v=0}^{\infty} e^{-\beta E_{v}}}{\sum_{v=0}^{\infty} e^{-\beta E_{v}}} \\
& =-\frac{\partial}{\partial \beta} \ln \sum_{v=0}^{\infty} e^{-\beta(v+1 / 2) \hbar \omega}, \beta=\frac{1}{k T} \\
& =-\frac{\partial}{\partial \beta} \ln e^{-1 / 2 \beta \hbar \omega} \frac{1}{1-e^{-\beta \hbar \omega}} \\
& =-\frac{\partial}{\partial \beta}\left[-\frac{1}{2} \hbar \omega \beta-\ln \left(1-e^{-\beta \hbar \omega}\right)\right] \\
& =\frac{1}{2} \hbar \omega+\frac{\hbar \omega}{e^{\hbar \omega / k T}-1} .
\end{aligned}
\]

Thus for one gm mole of diatomic gas

\[
C_{V_{1 v \omega}}=N \frac{\partial\langle E\rangle}{\partial T}=\frac{R\left(\frac{\hbar \omega}{k T}\right)^{2} e^{\hbar \omega / k T}}{\left(e^{\hbar \omega / k T}-1\right)^{2}}
\]

where \(R=N k\) is the gas constant.\\
In the present case

\[
\frac{\hbar \omega}{K T}=2.7088
\]

and

\[
C_{V_{1 \text { wib }}}=0.56 R
\]

6.177 In the rotation vibration band the main transition is due to change in vibrational quantum number \(v \rightarrow v-1\). Together with this rotational quantum number may change. The "Zeroeth line" \(0 \rightarrow 0\) is forbidden in this case so the neighbouring lines arise due to \(1 \rightarrow 0\) or \(0 \rightarrow 1\) in the rotational quantum number. Now

Thus

\[
E=E_{v}+\frac{\hbar^{2}}{2 I} J(J+1)
\]

Thus

\[
\hbar \omega=\hbar \omega_{0}+\frac{\hbar^{2}}{2 I}( \pm 2)
\]

Hence

\[
\Delta \omega=\frac{2 \hbar}{I}=\frac{2 \hbar}{\mu d^{2}}
\]

so

\[
d=\sqrt{\frac{2 \hbar}{\mu \Delta \omega}}
\]

Substitution gives \(\boldsymbol{d} \boldsymbol{=} \mathbf{0 . 1 2 8} \mathbf{n m}\).\\
6.178 If \(\lambda_{R}=\) wavelength of the red satellite\\
and \(\lambda_{V}=\) wavelength of the violet satellite\\
then

\[
\begin{aligned}
& \frac{2 \pi \hbar c}{\lambda_{R}}=\frac{2 \pi \hbar c}{\lambda_{0}}-\hbar \omega \\
& \frac{2 \pi \hbar c}{\lambda_{V}}=\frac{2 \pi \hbar c}{\lambda_{0}}+\hbar \omega
\end{aligned}
\]

and

\[
\begin{aligned}
& \lambda_{R}=424.3 \mathrm{~nm} \\
& \lambda_{V}=386.8 \mathrm{~nm}
\end{aligned}
\]

The two formulas can be combined to give

\[
\lambda=\frac{2 \pi c}{\frac{2 \pi c}{\lambda_{0}} \mp \omega}=\frac{\lambda_{0}}{1 \mp \frac{\lambda_{0} \omega}{2 \pi c}}
\]

6.179 As in the previous problem

\[
\omega=\pi c\left(\frac{1}{\lambda_{V}}-\frac{1}{\lambda_{R}}\right)=\frac{\pi c\left(\lambda_{R}-\lambda_{V}\right)}{\lambda_{R} \lambda_{V}}=1.368 \times 10^{14} \mathrm{rad} / \mathrm{s}
\]

The force constant \(x\) is defined by

\[
x=\mu \omega^{2}
\]

where \(\mu=\) reduced mass of the \(S_{2}\) molecule.\\
Substitution gives

\[
x=5.01 \mathrm{~N} / \mathrm{cm}
\]

6.180 The violet satellite arises from the transition \(1 \rightarrow 0\) in the vibrational state of the scattering molecule while the red satellite arises from the transition \(0 \rightarrow 1\). The intensities of these two transitions are in the ratio of initial populations of the two states i.e. in the ratio

Thus

\[
\begin{gathered}
e^{-\hbar \omega / k T} \\
\frac{I_{v}}{I_{r}}=e^{-\hbar \omega / k T}=0.067
\end{gathered}
\]

If the temperature is doubled, the rato increases to 0.259 , an increase of 3.9 times.\\
6.181 (a) \(\mathrm{CO}_{2}(\mathrm{O}-\mathrm{C}-\mathrm{O})\)

The molecule has 9 degrees of freedom 3 for each atom. This means that it can have up to nine frequencies. 3 degrees of freedom correspond to rigid translation, the frequency associated with this is zero as the potential energy of the system can not change under rigid translation. The P.E. will not change under rotations about axes passing through the \(C\)-atom and perpendicular to the \(O-C-O\) line. Thus there can be at most four non zero frequencies. We must look for modes different from the above.

One mode is\\
\includegraphics[max width=\textwidth, alt={}, center]{da0156b6-4133-434e-90d4-ec01b2070afe-350_93_541_1211_512}

Another mode is\\
\includegraphics[max width=\textwidth, alt={}, center]{da0156b6-4133-434e-90d4-ec01b2070afe-350_98_547_1355_514}

These are the only collinear modes.

A third mode is doubly degenerate :\\
\includegraphics[alt={}]{smile-17d05052a18a46d5d4183332f804ed8cb635cb73}\\
(vibration in \& \(\perp\) to the plane of paper).\\
(b) \(\mathrm{C}_{2} \mathrm{H}_{2}(\mathrm{H}-\mathrm{C}-\mathrm{C}-\mathrm{H})\)

There are \(4 \times 3-3-2=7\) different vibrations. There are three collinear modes.\\
\(\omega_{1}\)\\
\includegraphics[max width=\textwidth, alt={}, center]{da0156b6-4133-434e-90d4-ec01b2070afe-351_95_534_133_464}\\
\includegraphics[max width=\textwidth, alt={}, center]{da0156b6-4133-434e-90d4-ec01b2070afe-351_100_615_273_383}\\
\includegraphics[max width=\textwidth, alt={}, center]{da0156b6-4133-434e-90d4-ec01b2070afe-351_96_649_423_374}

Two other doubly degenerate frequencies are\\
\includegraphics[alt={}]{smile-a6545726c638733b9a5d748663c1ffac3e1a6a4f}\\
\includegraphics[alt={}]{smile-242f9abe93d5a1e9efa4e0438060aeb9179a4652}\\
together with their counterparts in the plane \(\perp^{r}\) to the paper.\\
6.182 Suppose the string is stretched along the \(x\) axis from \(x=0\) to \(x=l\) with the end points fixed. Suppose \(y(x, t)\) is the transverse displacement of the element at \(x\) at time \(t\). Then \(y(x, t)\) obeys

\[
\frac{\partial^{2} y}{\partial t^{2}}=v^{2} \frac{\partial^{2} y}{\partial x^{2}}
\]

We look for a stationary wave solution of this equation

\[
y(x, t)=A \sin \frac{\omega}{v} x \sin (\omega t+\delta)
\]

where \(A \& \delta\) are constants.. In this from \(y=0\) at \(x=0\). The further condition\\
implies

\[
\begin{gathered}
y=0 \text { at } x=l \\
\frac{\omega l}{\mathrm{v}}=N \pi, N>0 \\
N=\frac{l}{\pi \mathrm{v}} \omega
\end{gathered}
\]

or

\[
l N=\frac{l}{\pi v} d \omega
\]

Thus\\
6.183 Let \(\xi(x, y, t)\) be the displacement of the element at \((x, y)\) at time \(t\). Then it obeys the equation\\
where

\[
\begin{gathered}
\frac{\partial^{2} \xi}{\partial t^{2}}=\mathrm{v}^{2}\left(\frac{\partial^{2} \xi}{\partial x^{2}}+\frac{\partial^{2} \xi}{\partial y^{2}}\right) \\
\xi=0 \text { at } x=0, x=l, y=0 \text { and } y=l
\end{gathered}
\]

We look for a solution in the form

Then

\[
\begin{gathered}
\xi=A \sin k_{1} x \sin k_{2} y \sin (\omega t+\delta) \\
\omega^{2}=v^{2}\left(k_{1}^{2}+k_{2}^{2}\right) \\
k_{1}=\frac{n \pi}{l}, k_{2}=\frac{m \pi}{l}
\end{gathered}
\]

we write this as

\[
n^{2}+m^{2}=\left(\frac{l \omega}{\pi \mathrm{v}}\right)^{2}
\]

Here \(n, m>0\). Each pair ( \(n, m\) ) determines a mode. The total mumber of modes whose frequency is \(\leq \omega\) is the area of the quadrant of a circle of radius \(\frac{l \omega}{\pi v}\) i.e.

Then

\[
\begin{gathered}
N=\frac{\pi}{4}\left(\frac{l \omega}{\pi \mathrm{v}}\right)^{2} \\
d N=\frac{l^{2}}{2 \pi \mathrm{v}^{2}} \omega d \omega=\frac{S}{2 \pi \mathrm{v}^{2}} \omega d \omega
\end{gathered}
\]

where \(S=l^{2}\) is the area of the membrane.\\
6.184 For transverse vibrations of a 3-dimensional continuum (in the form of a cube say) we have the equation

\[
\frac{\partial^{2} \vec{\xi}}{\partial t^{2}}=v^{2} \vec{\nabla}^{2} \vec{\xi}, \operatorname{div} \vec{\xi}=0
\]

Here \(\vec{\xi}=\vec{\xi}(x, y, z, t)\). We look for solutions in the form

\[
\vec{\xi}=\vec{A} \sin k_{1} x, \sin k_{2} y, \sin k_{3} z, \sin (\omega t+\delta)
\]

This requires \(\omega^{2}=v^{2}\left(k_{1}^{2}+k_{2}^{2}+k_{3}^{2}\right)\)\\
From the boundary condition that \(\vec{\xi}=0\) for \(x=0, x=l, y=0, y=l, z=0, z=l\), we get

\[
k_{1}=\frac{n_{1} \pi}{l}, k_{2}=\frac{n_{2} \pi}{l}, k_{3}=\frac{n_{3} \pi}{l}
\]

where \(n_{1} ; n_{2}, n_{3}\) are nonzero positive integers.\\
We then get

\[
n_{1}^{2}+n_{2}^{2}+n_{3}^{2}=\left(\frac{l \omega}{\pi \mathrm{v}}\right)^{2}
\]

Each triplet ( \(n_{1}, n_{2}, n_{3}\) ) determines a possible mode and the number of such modes whose frequency \(\leq \omega\) is the volume of the all positive octant of a sphere of radius \(\frac{l \omega}{\pi v}\). Considering also the fact that the subsidiary condition \(\operatorname{div} \vec{\xi}=0\) implies two independent values of \(\vec{A}\) for each choice of the wave vector \(\left(k_{1}, k_{2}, k_{3}\right)\)\\
we find

Thus

\[
\begin{gathered}
N(\omega)=\frac{1}{8} \cdot \frac{4 \pi}{3}\left(\frac{l \omega}{\pi \mathrm{v}}\right)^{3} \cdot 2=\frac{V \omega^{3}}{3 \pi^{2} \mathrm{v}^{3}} \\
d N=\frac{V \omega^{2}}{\pi^{2} \mathrm{v}^{3}} d \omega
\end{gathered}
\]

6.185 To determine the Debye temperature we cut off the high frequency modes in such a way as to get the total number of modes correctly.\\
(a) In a linear crystal with \(n_{0} l\) atoms, the number of modes of transverse vibrations in any given plane cannot exceed \(n_{0} l\). Then

\[
n_{0} l=\frac{l}{\pi \mathrm{v}} \int_{0}^{\omega_{0}} d \omega=\frac{l}{\pi \mathrm{v}} \omega_{0}
\]

The cut off frequency \(\omega_{0}\) is related to the Debye temperature (1) by

Thus

\[
\Theta=\left(\frac{h}{k}\right) \pi n_{0} \mathrm{v}
\]

(b) In a square lattice, the number of modes of transverse oscillations cannot exceed \(n_{0} S\). Thus\\
or

\[
\begin{aligned}
n_{0} S & =\frac{S}{2 \pi \mathrm{v}^{2}} \int_{0}^{\omega_{0}} \omega d \omega=\frac{S}{4 \pi \mathrm{v}^{2}} \omega_{0}^{2} \\
\Theta & =\frac{\hbar}{k} \omega_{0}=\left(\frac{\hbar}{k}\right)\left(\sqrt{4 \pi n_{0}}\right) \mathrm{v}
\end{aligned}
\]

(c) In a cubic crystal, the maximum number of transverse waves must be \(2 n_{0} V\) (two for each atom). Thus

\[
2 n_{0} V=\frac{V}{\pi^{2} \mathrm{v}^{3}} \int_{0}^{\omega_{0}} \omega^{2} d \omega=\frac{V \omega_{0}^{3}}{3 \pi^{2} \mathrm{v}^{3}}
\]

Thus

\[
\Theta=\left(\frac{\hbar}{k}\right) \mathrm{v}\left(6 \pi^{2} n_{0}\right)^{1 / 3}
\]

6.186 We proceed as in the previous example. The total number of modes must be \(3 n_{0} v\) (total transverse and one longitudinal per atom). On the other hand the number of transverse modes per unit frequency interval is given by

\[
d N^{\perp}=\frac{V \omega^{2}}{\pi^{2} v_{\perp}^{3}} d \omega
\]

while the number of longitudinal modes ner unit frequency tnterval is given by

\[
d N^{\prime}=\frac{V \omega^{2}}{2 \pi^{2} v_{1}^{3}} d \omega
\]

The total number per unit frequency interval is

\[
d N=\frac{V \omega^{2}}{2 \pi^{2}}\left(\frac{2}{v_{\perp}^{3}}+\frac{1}{v_{\|}^{3}}\right) d \omega
\]

If the high frequency cut off is at \(\omega_{0}=\frac{k \Theta}{\hbar}\), the total number of modes will be

\[
3 n_{0} V=\frac{V}{6 \pi^{2}}\left(\frac{2}{v_{\perp}^{3}}+\frac{1}{v_{\|}^{3}}\right)\left(\frac{k \Theta}{\hbar}\right)^{3}
\]

Here \(n_{0}\) is the number of iron atoms per unit volume. Thus

\[
: \Theta=\frac{\hbar}{k}\left[18 \pi^{2} n_{0} /\left(\frac{2}{v_{\perp}^{3}}+\frac{1}{v_{\|}^{3}}\right)\right]^{1 / 3}
\]

For iron

\[
n_{0}=N_{A} / \frac{M}{\rho}=\frac{\rho N_{A}}{M}
\]

( \(\rho=\) density, \(M=\) atomic weight of iron \(N_{A}=\) Avogadro number).

\[
n_{0}=8.389 \times 10^{22} \text { per cc }
\]

Substituting the data we get

\[
\Theta=469 \cdot 1 \mathrm{~K}
\]

6.187 We apply the same formula but assume \(v_{\|} \approx v_{\perp}\). Then\\
or

\[
\begin{gathered}
\Theta=\frac{\hbar}{k} v\left(6 \pi^{2} n_{0}\right)^{1 / 3} \\
v=k \Theta /\left[\hbar\left(6 \pi^{2} n_{0}\right)^{1 / 3}\right]
\end{gathered}
\]

For \(A l\)

\[
n_{0}=\frac{\rho N_{A}}{M}=6.023 \times 10^{22} \text { per c.c }
\]

Thus

\[
\mathrm{v}=3.39 \mathrm{k} \mathrm{~m} / \mathrm{s} .
\]

The tabulated values are \(v_{\|}=6.3 \mathrm{~km} / \mathrm{s}\)\\
and

\[
v_{\perp}=3.1 \mathrm{~km} / \mathrm{s} .
\]

6.188 In the Debye approximation the number of modes per unit frequency interval is given by

\[
d N=\frac{l}{\pi \mathrm{v}} d \omega 0 \leq \omega \leq \frac{k \Theta}{\hbar}
\]

But

\[
\frac{k \Theta}{\hbar}=\pi n_{0} v
\]

Thus

\[
d N=\frac{l}{\pi \mathrm{v}} d \omega, 0 \leq \omega \leq \pi n_{0} \mathrm{v}
\]

The energy per mode is \(\quad<E>=\frac{1}{2} \hbar \omega+\frac{\hbar \omega}{e^{\hbar \omega / k T}-1}\)\\
Then the total interval energy of the chain is

\[
\begin{gathered}
U=\frac{l}{\pi \mathrm{v}} \int_{0}^{\pi n_{0} \mathrm{v}} \frac{1}{2} \hbar \omega d \omega \\
+\frac{l}{\pi \mathrm{v}} \int_{0}^{\pi n_{0} \mathrm{v}} \frac{\hbar \omega}{e^{\hbar \omega / K T}-1} d \omega=\frac{l \hbar}{4 \pi \mathrm{v}}\left(\pi n_{0} v\right)^{2}+\frac{l}{\pi v \hbar}(k T)^{2} \int_{0}^{T} \frac{x d x}{e^{x}-1}=l n_{0} k \cdot \frac{\hbar}{k}\left(\pi n_{0} v\right) \cdot \frac{1}{4} \\
+l n_{0} k \frac{T^{2}}{\left(\pi n_{0} v \hbar / k\right)} \int_{0}^{\frac{x}{e^{x}-1}}
\end{gathered}
\]

We put \(l n_{0} k=R\) for 1 mole of the chain.

Then

\[
U=R \Theta\left\{\frac{1}{4}+\left(\frac{I}{\Theta}\right)^{2} \int_{0}^{\partial / T} \frac{x d x}{e^{x}-1}\right\}
\]

Hence the molar heat capacity is by differentiation

\[
C_{V}=\left(\frac{\partial U}{\partial T}\right)_{\Theta}=R\left[2\left(\frac{T}{\Theta}\right) \int_{0}^{\Theta / T} \frac{x d x}{e^{x}-1}-\frac{\Theta / T}{e^{\Theta / T}-1}\right]
\]

when

\[
T \gg \Theta, C_{v} \approx R .
\]

6.189 If the chain has \(N\) atoms, we can assume atom number 0 and \(N+1\) held ficed. Then the displacement of the \(\boldsymbol{n}^{\text {th }}\) atom has the form

\[
u_{n}=A\left(\sin \frac{m \pi}{L} \cdot n a\right) \sin \omega t
\]

Here \(k=\frac{m \pi}{L}\). Allowed frequencies then have the form

\[
\omega=\omega_{\max } \sin \frac{k a}{2}
\]

In our form only + ve \(k\) values are allowed.\\
The number of modes in a wave number range \(d k\) is

But

\[
d N=\frac{L d k}{\pi}=\frac{L}{\pi} \frac{d k}{d \omega} d \omega
\]

Hence

\[
d \omega=\frac{a}{2} \omega_{\max } \cos \frac{k a}{2} d k
\]

So

\[
\frac{d \omega}{d k}=\frac{a}{2} \sqrt{\omega_{\max }^{2}-\omega^{2}}
\]

\[
d N=\frac{2 L}{\pi a} \frac{d \omega}{\sqrt{\omega_{\max }^{2}-\omega^{2}}}
\]

(b) The total number of modes is

\[
N=\int_{0}^{\omega_{\max }} \frac{2 L}{\pi a} \frac{d \omega}{\sqrt{\omega_{\max }^{2}-\omega^{2}}}=\frac{2 L}{\pi a} \cdot \frac{\pi}{2}=\frac{L}{a}
\]

i.e. the number of atoms in the chain.\\
6.190 Molar zero point energy is \(\frac{9}{8} R \Theta\). The zero point energy per gm of copper is \(\frac{9 R \Theta}{8 M_{C u}}, M_{C u}\) is the atomic weight of the copper.\\
Substitution gives \(48.6 \mathrm{~J} / \mathrm{gm}\).\\
6.191 (a) By Dulong and Petit's law, the classical heat capacity is \(3 R=24.94 \mathrm{~J} / \mathrm{K}\) - mole. Thus

\[
\frac{C}{C_{C l}}=0.6014
\]

From the graph we see that this\\
value of \(\frac{C}{C_{C l}}\) corresponds to \(\frac{T}{\Theta}=0.29\)\\
Hence

\[
\Theta=\frac{65}{0.29} \approx 224 \mathrm{~K}
\]

(b) \(22.4 \mathrm{~J} /\) mole- K corresponds to \(\frac{22.4}{3 \times 8.314}=0.898\). From the graph this corresponds to

\[
\frac{T}{\Theta} \approx 0.65 . \text { This gives } \Theta=\frac{250}{0.65} \approx 385 \mathrm{~K}
\]

Then 80 K corresponds to \(\frac{T}{\Theta}=0.208\)\\
The corresponding value of \(\frac{C}{C_{C l}}\) is 0.42 . Hence \(C=10.5 \mathrm{~J} / \mathrm{mole}-\mathrm{K}\).\\
(c) We calculate \(\Theta\) from the datum that \(\frac{C}{C_{C l}}=0.75\) at \(T=125 \mathrm{~K}\).

The \(x\)-coordinate corresponding to 0.75 is 0.40 . Hence

\[
\Theta=\frac{125}{0 \cdot 4}=3125 \mathrm{~K}
\]

Now

\[
k \Theta=\hbar \omega_{\max }
\]

So

\[
\omega_{\max }=4.09 \times 10^{13} \mathrm{rad} / \mathrm{sec}
\]

6.192 We use the formula (6.4d)

\[
\begin{aligned}
U & =9 R \Theta\left[\frac{1}{8}+\left(\frac{T}{\Theta}\right)^{4} \int_{0}^{\Theta / T} \frac{x^{3} d x}{e^{x}-1}\right] \\
& =9 R \Theta\left[\frac{1}{8}+\left\{\int_{0}^{\infty} \frac{x^{3} d x}{e^{x}-1}\right\}\left(\frac{T}{\Theta}\right)^{4}-\left(\frac{T}{\Theta}\right)^{4} \int_{\Theta / T}^{\infty} \frac{x^{3} d x}{e^{x}-1}\right]
\end{aligned}
\]

In the limit \(T \ll \Theta\), the third term in the bracket is exponentially small together with its derivatives.\\
Then we can drop the last term

Thus

\[
\begin{gathered}
U=\text { Const }+\frac{9 R}{\Theta^{3}} T^{4} \int_{0}^{\infty} \frac{x^{3} d x}{e^{x}-1} \\
C_{V}=\left(\frac{\partial U}{\partial T}\right)_{V}^{\infty}=\left(\frac{\partial U}{\partial T}\right)_{\Theta}=36 R\left(\frac{T}{\Theta}\right)^{3} \int_{0}^{\infty} \frac{x^{3} d x}{e^{x}-1}
\end{gathered}
\]

Now from the table in the book

\[
\int_{0}^{\infty} \frac{x^{3} d x}{e^{x}-1}=\frac{\pi^{4}}{15}
\]

Thus

\[
C_{V}=\frac{12 \pi^{4}}{5}\left(\frac{T}{\Theta}\right)^{3}
\]

Note :- Call the \(3^{\text {rd }}\) term in the bracket above \(-U_{3}\). Then

\[
U_{3}=\left(\frac{T}{\Theta}\right)^{4} \int_{\Theta / T}^{\infty} \frac{x^{3}}{2 \sin h(x / 2)} \cdot e^{-x / 2} d x
\]

The maximum value of \(\frac{x^{3}}{2 \sin h \frac{x}{2}}\) is a finite +ve quantity \(C_{0}\) for \(0 \leq x<\infty\). Thus

\[
U_{3} \leq 2 C_{0}\left(\frac{T}{\Theta}\right)^{4} e^{-\Theta / 2 T}
\]

we see that \(U_{3}\) is exponentially small as \(T \rightarrow 0\). So is \(\frac{d U_{3}}{d T}\).\\
6.193 At low temperatures \(C \propto T^{3}\). This is also a test of the "lowness" of the temperature We see that

\[
\left(\frac{C_{1}}{C_{2}}\right)^{1 / 3}=1.4982 \approx 1.5=\frac{T_{1}}{T_{2}}=\frac{30}{20}
\]

Thus \(T^{3}\) law is obeyed and \(T_{1}, T_{2}\) can we regarded low.\\
6.194

The total zero point energy of 1 mole of the solid is \(\frac{9}{8} R \Theta\). Dividing this by the number of modes 3 N we get the average zero point energy per mode. It is

\[
\frac{3}{8} k \Theta .
\]

6.195 In the Debye model

\[
d N_{\omega}=A \omega^{2}, 0 \leq \omega \leq \omega_{m}
\]

Then

\[
3 N=\int_{0}^{\omega_{m}} d N_{\omega}=\frac{A \omega_{m}^{3}}{3} .(\text { Total no. of modes is } 3 \mathrm{~N})
\]

Thus

\[
A=\frac{9 N}{\omega_{m}^{3}} .
\]

we get

\[
\begin{aligned}
U & =\frac{9 N}{\omega_{m}^{3}} \int_{0}^{\omega_{m}} \frac{\omega^{2} \cdot \hbar \omega}{e^{-h} \omega / k T-1} d \omega \quad \text { ignoring } \mathrm{z} \\
& =9 N \hbar \omega_{m} \int_{0}^{1} \frac{x^{3} d x}{e^{\hbar \omega_{m} x / k T}-1}, x=\frac{\omega}{\omega_{m}} \\
& =9 R \Theta \int_{0}^{1} \frac{x^{3} d x}{e^{x \Theta / T}-1}, \Theta=\hbar \omega_{m} / k
\end{aligned}
\]

Thus

\[
\frac{1}{9 R \Theta} \frac{d U(x)}{d x}=\frac{x^{3}}{e^{x \Theta / T}-1} \text { for } 0 \leq x \leq 1
\]

For

\[
T=\Theta / 2, \text { this is } \frac{x^{3}}{e^{2 x}-1} ; \text { for }
\]

\(T=\frac{\Theta}{4}\), it is \(\frac{x^{3}}{e^{4 x}-1}\). Plotting then we get the figures given in the answer.\\
6.196 The maximum energy of the phonon is

\[
\hbar \omega_{m}=k \Theta=28.4 \mathrm{meV}
\]

On substituting \(\Theta=330 \mathrm{~K}\).\\
To get the corresponding value of the maximum momentum we must know the dispersion relation \(\omega=\omega(\vec{k})\). For small \((\vec{k})\) we know \(\omega=v|\vec{k}|\) where \(v\) is velocity of sound in the crystal. For an order of magnitude estimate we continue to use this result for high \(|\vec{k}|\). Then we estimate \(v\) from the values of the modulus of elasticity and density

\[
\mathrm{v} \sim \sqrt{\frac{E}{\rho}}
\]

We write

\[
E \sim 100 \mathrm{GPa}, \rho=8.9 \times 10^{3} \mathrm{~kg} / \mathrm{m}^{3}
\]

Then

\[
v \sim^{3 \times 10^{3} \mathrm{~m} / \mathrm{s}}
\]

Hence

\[
\hbar|\vec{k}|_{\max } \sim \frac{\hbar \omega_{m}}{\mathrm{v}} \sim 1.5 \times 10^{-19} \mathrm{gm} \mathrm{~cm} \mathrm{~s}^{-1}
\]

6.197 (a) From the formula

\[
d n=\frac{\sqrt{2} m^{3 / 2}}{\pi^{2} \hbar^{3}} E^{1 / 2} d E
\]

the maximum value \(E_{\text {max }}\) of \(E\) is determined in terms of \(n\) by

Or

\[
\begin{aligned}
n & =\frac{\sqrt{2} m^{3 / 2}}{\pi^{2} \hbar^{3}} \int_{0}^{E_{\max }} E^{1 / 2} d E \\
& =\frac{\sqrt{2} m^{3 / 2}}{\pi^{2} \hbar^{3}} \frac{2}{3} E_{\max }^{3 / 2} \\
E_{\max }^{3 / 2} & =\left(\frac{\hbar^{2}}{2 m}\right)^{3 / 2}\left(3 \pi^{2} n\right) \\
E_{\max } & =\frac{\hbar^{2}}{2 m}\left(3 \pi^{2} n\right)^{2 / 3}
\end{aligned}
\]

(b) Mean K.E. < \(E\) >is

\[
\begin{aligned}
<E> & =\int_{0}^{E_{\max }} E d n / \int_{0}^{E_{\max }} d n \\
& =\int_{0}^{E_{\max }} E^{3 / 2} d E / \int_{0}^{1 / 2} d E=\frac{2}{5} E_{\max }^{5 / 2} / \frac{2}{3} E_{\max }^{3 / 2}=\frac{3}{5} E_{\max }
\end{aligned}
\]

6.198 The fraction is

\[
\eta=\int_{\frac{1}{2} E_{\max }}^{E_{\max }} E^{1 / 2} d E / \int_{0}^{E_{\max }} E^{1 / 2} d E=1-2^{-3 / 2}=0.646 \text { or } 64.6 \%
\]

6.199 We calculate the concentration \(n\) of electron in the Na metal from\\
we get from

\[
\begin{gathered}
E_{\max }=E_{F}=\frac{\hbar^{2}}{2 m}\left(3 \pi^{2} n\right)^{2 / 3} \\
E_{F}=3.07 \mathrm{eV} \\
n=2.447 \times 10^{22} \text { per c.c. }
\end{gathered}
\]

From this we get the number of electrons per one \(N a\) atom as

\[
\frac{n}{\rho} \cdot \frac{M}{N_{A}}
\]

where \(\rho=\) density of \(N a, M=\) molar weight in gm of \(N a, N_{A}=\) Avogadro number we get\\
0.963 elecrons per one Na atom.\\
6.200 The mean K.E. of electrons in a Fermi gas is \(\frac{3}{5} E_{F}\). This must equal \(\frac{3}{2} k T\). Thus

\[
T=\frac{2 E_{F}}{5 k}
\]

We calculate \(E_{F}\) first. For Cu

\[
n=\frac{N_{A}}{M / \rho}=\frac{\rho N_{A}}{M}=" 8.442 \times 10^{22} \text { per c.c. }
\]

Then

\[
E_{F}=7.01 \mathrm{eV}
\]

and

\[
T=3.25 \times 10^{4} \mathrm{~K}
\]

6.201 We write the expression for the number of electrons as

\[
d N=\frac{V \sqrt{2} m^{3 / 2}}{\pi^{2} \hbar^{3}} E^{1 / 2} d E
\]

Hence if \(\Delta E\) is the spacing between neighbouring levels near the Fermi level we must have

\[
2=\frac{V \sqrt{2} m^{3 / 2}}{\pi^{2} \hbar^{3}} E_{F}^{1 / 2} \Delta E
\]

( 2 on the RHS is to take care of both spins \(f\) electrons). Thus

But

\[
\Delta E=\frac{\sqrt{2} \pi^{2} \hbar^{3}}{V m^{3 / 2} E_{F}^{1 / 2}}
\]

So

\[
E_{F}^{1 / 2}=\frac{\hbar}{\sqrt{2} m^{1 / 2}}\left(3 \pi^{2} n\right)^{1 / 3}
\]

So \(\quad \Delta E=\frac{2 \pi^{2} \hbar^{2}}{m V\left(3 \pi^{2} n\right)^{1 / 3}}\)\\
Substituting the data we get

\[
\Delta E=1.79 \times 10^{-22} \mathrm{eV}
\]

6.202 (a) From

\[
d n(E)=\frac{\sqrt{2} m^{3 / 2}}{\pi^{2} \hbar^{3}} E^{1 / 2} d E
\]

we get on using \(E=\frac{1}{2} m v^{2}, d n(E)=d n(v)\)

\[
d n(\mathrm{v})=\frac{\sqrt{2} m^{3 / 2}}{\pi^{2} \hbar^{3}} \frac{1}{\sqrt{2}} m^{1 / 2} \mathrm{v} m \mathrm{v} d \mathrm{v}=\frac{m^{3}}{\pi^{2} \hbar^{3}} \mathrm{v}^{2} d \mathrm{v}
\]

This holds for \(0<v<v_{\mathrm{F}}\) where \(\frac{1}{2} m v_{\mathrm{F}}^{2}=E_{F}\)\\
and

\[
d n(v)=0 \text { for } v>v_{F} .
\]

(b) Mean velocity is

\[
\begin{gathered}
\langle v\rangle=\int_{0}^{v_{F}} v^{3} d v / \int_{0}^{v_{F}} v^{2} d v=\frac{3}{4} v_{F} \\
\therefore \frac{\langle v\rangle}{v_{F}}=\frac{3}{4} .
\end{gathered}
\]

6.203 Using the formula of the previous section

\[
d n(v)=\frac{m^{3}}{\pi^{2} h^{3}} v^{2} d v
\]

We put \(m \mathrm{v}=\frac{2 \pi \hbar}{\lambda}\), where \(\lambda=\) de Broglie wavelength

\[
m d v=-\frac{2 \pi \hbar}{\lambda^{2}} d \lambda
\]

Taking account of the fact that \(\lambda\) decreases when \(v\) increases we write

\[
d n(\lambda)=-d n(v)=\frac{(2 \pi)^{3} d \lambda}{\pi^{2} \lambda^{4}} \cdot=\frac{8 \pi}{\lambda^{4}} d \lambda
\]

6.204 From the kinetic theory of gasses we know

\[
p=\frac{2}{3} \frac{U}{V}
\]

Here \(U\) is the total interval energy of the gas. This result is applicable to Fermi gas also Now at \(T=0\)\\
so

\[
\begin{aligned}
U=U_{0}= & N<E>=n V<E> \\
p & =\frac{2}{3} n<E> \\
& =\frac{2}{3} n \times \frac{3}{5} E_{F}=\frac{2}{5} n E_{F} \\
& =\frac{\hbar^{2}}{5 m}\left(3 \pi^{2}\right)^{2 / 3} n^{5 / 3}
\end{aligned}
\]

Substituting the values we get

\[
p=4.92 \times 10^{4} \text { atmos }
\]

6.205 From Richardson's equation

\[
I=a T^{2} e^{-A / k T}
\]

where \(A\) is the work function in eV . When \(T\) increases by \(\Delta T, I\) increases to \((1+\eta) I\). Then

\[
1+\eta=\left(\frac{T+\Delta T}{T}\right)^{2} e^{-\frac{A}{k T}\left(\frac{T}{T+\Delta T}-1\right)}=\left(1+\frac{\Delta T}{T}\right)^{2} e^{+\frac{A}{k T} \cdot \frac{\Delta T}{T+\Delta T}}
\]

Expanding and neglecting higher powers of \(\frac{\Delta T}{T}\) we get

\[
\eta=2 \frac{\Delta T}{T}+\frac{A}{k T^{2}} \Delta T
\]

Thus

\[
\begin{gathered}
A=k T\left(\frac{\eta T}{\Delta T}-2\right) \\
A=4.48 \mathrm{eV}
\end{gathered}
\]

Substituting we get\\
6.206\\
\includegraphics[max width=\textwidth, alt={}, center]{da0156b6-4133-434e-90d4-ec01b2070afe-363_435_665_639_165}\\
\includegraphics[max width=\textwidth, alt={}, center]{da0156b6-4133-434e-90d4-ec01b2070afe-363_419_484_639_930}

The potential energy inside the metal is \(-U_{0}\) for the electron and it related to the work function A by

\[
U_{0}=E_{F}+A
\]

If \(T\) is the K.E. of electrons outside the metal, its K.E. inside the metal will be \(\left(E+U_{0}\right)\). On entering the metal electron connot experience any tangential force so the tangential component of momentum is unchanged. Then

\[
\sqrt{2 m T} \sin \alpha=\sqrt{2 m\left(T+U_{0}\right)} \sin \beta
\]

Hence

\[
\frac{\sin \alpha}{\sin \beta}=\sqrt{1+\frac{U_{0}}{T}}=n \text { by definition of refractive index. }
\]

In sodium with one free electron per Na atom

\[
\begin{aligned}
n=2.54 & \times 10^{22} \text { per c.c. } \\
E_{F} & =3.15 \mathrm{eV} \\
A & =2.27 \mathrm{eV} \text { (from table) } \\
U_{0} & =5.42 \mathrm{eV} \\
n & =1.02
\end{aligned}
\]

6.207 In a pure (intrinsic) semiconductor the conductivity is related to the temperature by the following formula very closely :

\[
\sigma=\sigma_{0} e^{-\Delta a / 2 k T}
\]

where \(\Delta \varepsilon\) is the energy gap between the top of valence band and the bottom of conduction band; it is also the minimum energy required for the formation of electron-hole pair. The conductivity increases with temperature and we have\\
or

\[
\eta=e^{+\frac{\Delta \varepsilon}{2 k}\left(\frac{1}{T_{1}}-\frac{1}{T_{2}}\right)}
\]

\[
\ln \eta=\frac{\Delta \varepsilon}{2 k} \frac{T_{2}-T_{1}}{T_{1} T_{2}}
\]

Hence

\[
\Delta \varepsilon=\frac{2 k T_{1} T_{2}}{T_{2}-T_{1}} \ln \eta
\]

Substitution gives

\[
\Delta \varepsilon=0.333 \mathrm{eV}=E_{\mathrm{min}}
\]

6.208 The photoelectric threshold determines the band gap \(\Delta \varepsilon\) by

\[
\Delta \varepsilon=\frac{2 \pi \hbar c}{. \lambda_{r h}}
\]

On the other hand the temperature coefficient of resistance is defined by ( \(\rho\) is resistivity)

\[
\alpha=\frac{1}{\rho} \frac{d \rho}{d T}=\frac{d}{d T} \ln \rho=-\frac{d}{d T} \ln \sigma
\]

where \(\sigma\) is the conductivity. But

\[
\ln \sigma=\ln \sigma_{0}-\frac{\Delta \varepsilon}{2 k T}
\]

Then

\[
\alpha=-\frac{\Delta \varepsilon}{2 k T^{2}}=-\frac{\pi \hbar c}{k T^{2} \lambda_{t h}}=-0.047 \mathrm{~K}^{-1}
\]

6.209 At high temperatures (small values of \(T^{-1}\) ) most of the conductivity is infrinsic i.e. it is due to the transition of electrons from the upper levels of the valance band into the lower levels of conduction vands.\\
For this we can apply approximately the formula\\
or

\[
\begin{aligned}
& \sigma=\sigma_{0} \exp \left(-\frac{E_{g}}{2 k T}\right) \\
& \ln \sigma=\ln \sigma_{0}-\frac{E_{g}}{2 k t}
\end{aligned}
\]

From this we get the band gap

\[
E_{g}=-2 k \frac{\Delta \ln \sigma}{\Delta(1 / T)}
\]

The slope must be calculated at small \(\frac{1}{T}\). Evaluation gives \(-\frac{\Delta \ln \sigma}{\Delta\left(\frac{1}{T}\right)}=7000 \mathrm{~K}\)\\
Hence

\[
E_{g}=1.21 \mathrm{eV}
\]

At low temperatures (high values of \(\frac{1}{T}\) ) the conductance is mostly due to impurities. If \(E_{0}\) is the ionization energy of donor levels then we can write the approximate formula (valid at low temperature)

So

\[
\begin{aligned}
& \sigma^{\prime}=\sigma_{0}^{\prime} \exp \left(-\frac{E_{0}}{2 k T}\right) \\
& E_{0}=-2 k T \frac{\Delta \ln \sigma^{\prime}}{\Delta\left(\frac{1}{T}\right)}
\end{aligned}
\]

The slope must be calculated at low temperatures. Evaluation gives the slope

Then

\[
\begin{gathered}
-\frac{\Delta \ln \sigma^{\prime}}{\Delta\left(\frac{1}{T}\right)}=\frac{1}{3} \times 1000 \mathrm{~K} \\
E_{0} \sim 0.057 \mathrm{eV}
\end{gathered}
\]

6.210 We write the conductivity of the sample as \(\sigma=\sigma_{i}+\sigma_{\gamma}\)\\
where \(\sigma_{i}=\) intrinsic conductivity and \(\sigma_{\gamma}\) is the photo conductivity. At \(t=0\), assuming saturation we have

\[
\frac{1}{\rho_{1}}=\frac{1}{\rho}+\sigma_{\gamma_{0}} \quad \text { or } \quad \sigma_{\gamma_{0}}=\frac{1}{\rho_{1}}-\frac{1}{\rho}
\]

Time \(t\) after light source is switched off we have because of recombination of electron and holes in the sample

\[
\sigma=\sigma_{i}+\sigma_{\gamma_{0}} e^{-t / T}
\]

where \(T=\) mean lifetime of electrons and holes.\\
Thus

\[
\frac{1}{\rho_{2}}=\frac{1}{\rho}+\left(\frac{1}{\rho_{1}}-\frac{1}{\rho}\right) e^{-t / T}
\]

or

\[
\frac{1}{\rho_{2}}-\frac{1}{\rho}=\left(\frac{1}{\rho_{1}}-\frac{1}{\rho}\right) e^{-t / T}
\]

or

\[
e^{t / T}=\frac{\frac{1}{\rho_{1}}-\frac{1}{\rho}}{\frac{1}{\rho_{2}}-\frac{1}{\rho}}=\frac{\rho_{2}\left(\rho-\rho_{1}\right)}{\rho_{1}\left(\rho-\rho_{2}\right)}
\]

Hence

\[
T=t / \ln \left\{\frac{\rho_{2}\left(\rho-\rho_{1}\right)}{\rho_{1}\left(\rho-\rho_{2}\right)}\right\}
\]

Substitution gives \(T=9.87 \mathrm{~ms} \sim 0.01 \mathrm{sec}\)\\
\includegraphics[max width=\textwidth, alt={}, center]{da0156b6-4133-434e-90d4-ec01b2070afe-366_279_788_69_322}

We shall ignore minority carriers.\\
Drifting holes experience a sideways force in the magnetic field and react by setting up a Hall electric field \(E_{y}\) to counterbalance it. Thus

\[
v_{x} B=E_{y}=\frac{V_{H}}{h}
\]

If the concentration of carriers is \(n\) then

Hence

\[
\begin{gathered}
j_{x}=n e v_{x} \\
n=\frac{J_{x}}{e v_{x}}=\frac{\frac{j_{x}}{e V_{H}}}{h B}=\frac{j_{x} h B}{e V_{H}}
\end{gathered}
\]

Also using

\[
j_{x}=\sigma E_{x}=E_{x} / \rho=\frac{V}{\rho l}
\]

we get

\[
n=\frac{V h B}{e \rho l V_{H}}
\]

Substituting the data (note that in MKS units \(B=5.0 \mathrm{kG}=0.5 \mathrm{~T}\) )

\[
\rho=2.5 \times 10^{-2} \text { ohm }-\mathrm{m}
\]

we get

\[
\begin{aligned}
n & =4.99 \times 10^{21} \mathrm{~m}^{-3} \\
& =4.99 \times 10^{15} \text { per cm}
\end{aligned}
\]

Also the mobility is

\[
u_{0}=\frac{v_{x}}{E_{x}}=\frac{V_{H}}{n B} \times \frac{l}{V}=\frac{V_{H} l}{h B V}
\]

Substitution gives

\[
u_{0}=0.05 \mathrm{~m}^{2} / \mathrm{V}-\mathrm{s}
\]

6.212\\
\includegraphics[max width=\textwidth, alt={}, center]{da0156b6-4133-434e-90d4-ec01b2070afe-366_167_520_1576_426}

If an electric field \(E_{x}\) is present in a sample containing equal amounts of both electrons and holes, the two drift in opposite directions.\\
In the presence of a magnetic field \(B_{z}=B\) they set up Hall voltages in opposite directions.

The net Hall electric field is given by

But

\[
\begin{aligned}
\frac{E_{y}}{E_{x}} & =\frac{1}{\eta} \text { Hence } \\
\left|u_{0}^{+}-u_{0}^{-}\right| & =\frac{1}{\eta B}
\end{aligned}
\]

Substitution gives \(\left|u_{0}^{+}-u_{0}^{-}\right|=0.2 \mathrm{~m}^{2} /\) volt - sec\\
\includegraphics[max width=\textwidth, alt={}, center]{da0156b6-4133-434e-90d4-ec01b2070afe-367_171_661_587_406}

When the sample contains unequal number of carriers of both types whose mobilities are different, static equilibrium (i.e. no transverse movement of either electron or holes) is impossible in a magnetic field. The transverse electric field acts differently on electrons and holes. If the \(E_{y}\) that is set up is as shown, the net Lorentz force per unit charge (effective transverse electric field) on electrons is

\[
E_{y}-v_{x}^{-} B
\]

and on holes

\[
E_{y}+v_{x}^{+} B
\]

(we are assuming \(B=B_{z}\) ). There is then a transverse drift of electrons and holes and the net transverse current must vanish in equilibrium. Using mobility\\
or

\[
\begin{gathered}
u_{0}^{-} N_{e} e\left(E_{y}-u_{0}^{-} E_{x} B\right)+N_{h} e u_{0}^{+}\left(E_{y}+u_{0}^{+} E_{x} B\right)=0 \\
E_{y}=\frac{N_{e} u_{0}^{-2}-N u_{0}^{+2}}{N_{e} u_{0}^{-}+N_{h} u_{0}^{+}} E_{x} B
\end{gathered}
\]

On the other hand

\[
J_{x}=\left(N_{e} u_{0}^{-}+N_{h} u_{0}^{+}\right) e E_{x}
\]

Thus, the Hall coefficient is

\[
R_{H}=\frac{E_{y}}{j_{x} B}=\frac{1}{e} \frac{N_{e} u_{0}^{-2}-N_{h} u_{0}^{+2}}{\left(N_{e} u_{0}^{-}+N_{h} u_{0}^{+}\right)^{2}}
\]

We see that \(R_{H}=0\) when

\[
\frac{N_{e}}{N_{h}}=\left(\frac{u_{0}^{+}}{u_{0}^{-}}\right)^{2}=\frac{1}{\eta^{2}}=\frac{1}{4}
\]

\subsection*{6.5 RADIOACTIVITY}
6.214 (a) The probability of survival (i.e. not decaying) in time \(t\) is \(e^{-\lambda t}\). Hence the probability of decay is \(1-e^{-\lambda t}\)\\
(b) The probability that the particle decays in time \(d t\) around time \(t\) is the difference

\[
e^{-\lambda t}-e^{-\lambda(t+d t)}=e^{-\lambda t}\left[1-e^{-e \lambda d t}\right]=\lambda e^{-\lambda t} d t
\]

Therefore the mean life time is

\[
T=\int_{0}^{\infty} t \lambda e^{-\lambda t} d t / \int_{0}^{\infty} \lambda e^{-\lambda t} d t=\frac{1}{\lambda} \int_{0}^{\infty} x e^{-x} d x / \int_{0}^{\infty} e^{-x} d x=\frac{1}{\lambda}
\]

6.215 We calculate \(\lambda\) first

\[
\lambda=\frac{\ln 2}{T_{1 / 2}}=9.722 \times 10^{-3} \text { per day }
\]

Hence\\
fraction decaying in a month \(=1-e^{-\lambda t}=0.253\)\\
6.216 Here \(N_{0}=\frac{1 \mu g}{24 g} \times 6.023 \times 10^{23}=2.51 \times 10^{16}\)

Also

\[
\lambda=\frac{\ln 2}{T_{1 / 2}}=0.04621 \text { per hour }
\]

So the number of \(\beta\) rays emitted in one hour is

\[
N_{0}\left(1-e^{-\lambda}\right)=1.13 \times 10^{15}
\]

6.217 If \(N_{0}\) is the number of radionuclei present initially, then

\[
\begin{gathered}
N_{1}=N_{0}\left(1-e^{-t_{1} / \tau}\right) \\
\eta N_{1}=N_{0}\left(1-e^{-t_{2} / \tau}\right)
\end{gathered}
\]

where

\[
\eta=2.66 \text { and } t_{2}=3 t_{1} \text {. Then }
\]

or

\[
\begin{gathered}
\eta=\frac{1-e^{-t_{2} / \tau}}{1-e^{-t_{1} / \tau}} \\
\eta-\eta e^{-t_{1} / \tau}=1-e^{-t_{2} / \tau}
\end{gathered}
\]

Substituting the values

\[
1.66=2.66 e^{-2 / \tau}-e^{-6 / \tau}
\]

Put \(e^{-2 / \tau}=x\). Then

\[
\begin{gathered}
x^{3}-2.66 x+1.66=0 \\
\left(x^{2}-1\right) x-1.66(x-1)=0
\end{gathered}
\]

or

\[
(x-1)\left(x^{2}+x-1 \cdot 66\right)=0
\]

Now

\[
\begin{aligned}
& x=1 \text { so } x^{2}+x-1.66=0 \\
& x=\frac{-1 \pm \sqrt{1+4 \times 1.66}}{2}
\end{aligned}
\]

Negative sign has to be rejected as \(x>0\).\\
Thus

\[
x=0.882
\]

This gives

\[
\tau=\frac{-2}{\ln 0.882}=15.9 \mathrm{sec}
\]

6.218 If the half-life is \(T\) days

\[
(2)^{-7 / T}=\frac{1}{2 \cdot 5}
\]

Hence

\[
\frac{7}{T}=\frac{\ln 2.5}{\ln 2}
\]

or

\[
T=\frac{7 \ln 2}{\ln 2.5}=5.30 \text { days. }
\]

6.219 The activity is proportional to the number of parent nuclei (assuming that the daughter is not radioactive). In half its half-life period, the number of parent nucli decreases by a factor

\[
(2)^{-1 / 2}=\frac{1}{\sqrt{2}}
\]

So activity decreases to \(\frac{650}{\sqrt{2}}=460\) particles per minute.\\
6.220 If the decay constant (in (hour) \({ }^{-1}\) ) is \(\lambda\), then the activity after one hour will decrease by a factor \(e^{-\lambda}\). Hence\\
or

\[
\begin{gathered}
0.96=e^{-\lambda} \\
\lambda=1.11 \times 10^{-5} \mathrm{~s}^{-1}=0.0408 \text { per hour }
\end{gathered}
\]

The mean life time is \(24 \cdot 5\) hour\\
6.221 Here \(N_{0}=\frac{1}{238} \times 6.023 \times 10^{23}\).

\[
=2.531 \times 10^{21}
\]

The activity is

\[
A=1.24 \times 10^{4} \mathrm{dis} / \mathrm{sec}
\]

Then

\[
\lambda=\frac{A}{N_{0}}=4.90 \times 10^{-18} \text { per sec }
\]

Hence the half life is

\[
T_{\frac{1}{2}}=\frac{\ln 2}{\lambda}=4.49 \times 10^{9} \text { years }
\]

6.222 In old wooden atoms the number of \(C^{14}\) nuclei steadily decreases because of radioactive decay. (In live trees biological processes keep replenishing \(C^{14}\) nuclei maintaining a balance. This balance starts getting disrupted as soon as the tree is felled.)\\
If \(T_{1 / 2}\) is the half life of \(C^{14}\) then \(e^{-t \times \frac{\ln 2}{T_{\nu 2}}}=\frac{3}{5}\)\\
Hence

\[
t=T_{1 / 2} \frac{\ln 5 / 3}{\ln 2}=4105 \text { years } \approx 4.1 \times 10^{3} \text { years }
\]

6.223 What this implies is that in the time since the ore was formed, \(\frac{\eta}{1+\eta} U^{238}\) nuclei have remained undecayed. Thus\\
or

\[
\begin{aligned}
& \frac{\eta}{1+\eta}=e^{-t \times \frac{\ln 2}{T_{1 / 2}}} \\
& t=T_{\frac{1}{2}} \frac{\ln \frac{1+\eta}{\eta}}{\ln 2}
\end{aligned}
\]

Substituting \(T_{\frac{1}{2}}=4.5 \times 10^{9}\) years, \(\eta=2.8\)\\
we get

\[
t=1.98 \times 10^{9} \text { years. }
\]

6.224 The specific activity of \(\mathrm{Na}^{24}\) is

\[
\lambda \frac{N_{A}}{M}=\frac{N_{A} \ln 2}{M T_{\frac{1}{2}}}=3.22 \times 10^{17} \mathrm{dis} /(\mathrm{gm} . \mathrm{sec})
\]

Here \(M=\) molar weight of \(N_{a}^{24}=24 \mathrm{gm}, N_{A}\) is Avogadro number \(\& T_{\frac{1}{2}}\) is the half-life of \(\mathrm{Na}{ }^{24}\).

Similarly the specific activity of \(U^{235}\) is

\[
\begin{aligned}
& \frac{6.023 \times 10^{23} \times \ln 2}{235 \times 10^{8} \times 365 \times 86400} \\
& =0.793 \times 10^{5} \mathrm{dis} /(\mathrm{gm}-\mathrm{s})
\end{aligned}
\]

6.225 Let \(V=\) volume of blood in the body of the human being. Then the total activity of the blood is \(A^{\prime} V\). Assuming all this activity is due to the injected \(\mathrm{Na}^{24}\) and taking account of the decay of this radionuclide, we get

\[
V \dot{A}^{\prime}=A e^{-\lambda t}
\]

Now

\[
\lambda=\frac{\ln 2}{15} \text { per hour }, t=5 \text { hour }
\]

Thus

\[
V=\frac{A}{A^{\prime}} e^{-\ln 2 / 3}=\frac{2.0 \times 10^{3}}{(16 / 60)} e^{-\ln 2 / 3} \mathrm{cc}=5.95 \text { litre }
\]

\subsection*{6.226 We see that}
Specific activity of the sample

\[
=\frac{1}{M+M^{\prime}}\left\{\text { Activity of } M \mathrm{gm} \text { of } C o^{58} \text { in the sample }\right\}
\]

Here \(M\) and \(M^{\prime}\) are the masses of \(C o^{58}\) and \(C o^{59}\) in the sample. Now activity of \(M \mathrm{gm}\) of Co \({ }^{58}\)

\[
\begin{gathered}
=\frac{M}{58} \times 6.023 \times 10^{23} \times \frac{\ln 2}{71.3 \times 86400} \mathrm{dis} / \mathrm{sec} \\
=1.168 \times 10^{15} \mathrm{M}
\end{gathered}
\]

Thus from the problem\\
or

\[
\begin{aligned}
& 1.168 \times 10^{15} \frac{M}{M+M^{\prime}}=2.2 \times 10^{12} \\
& \frac{M}{M+M^{\prime}}=1.88 \times 10^{-3} \text { i.e. } 0.188 \%
\end{aligned}
\]

6.227 Suppose \(N_{1}, N_{2}\) are the initial number of component nuclei whose decay constants are \(\lambda_{1}, \lambda_{2}\left(\right.\) in \(\left.(\text { hour })^{-1}\right)\)\\
Then the activity at any instant is

\[
A=\lambda_{1} N_{1} e^{-\lambda_{1} t}+\lambda_{2} N_{2} e^{-\lambda_{2} t}
\]

The activity so defined is in units dis/hour. We assume that data \(\ln A\) given is of its natural logarithm. The daughter nuclei are assumed nonradioactive.\\
We see from the data that at large \(t\) the change in \(\ln A\) per hour of elapsed time is constant and equal to -0.07 . Thus

\[
\lambda_{2}=0.07 \text { per hour }
\]

We can then see that the best fit to data is obtained by

\[
A(t)=51 \cdot 1 e^{-0.66 t}+10 \cdot 0 e^{-0.07 t}
\]

[To get the fit we calculate \(A(t) e^{0.07 t}\). We see that it reaches the constant value 10.0 at \(t=7,10,14,20\) very nearly. This fixes the second term. The first term is then obtained by subtracting out the constant value 10.0 from each value of \(A(t) e^{0.07 t}\) in the data for small t.]

Thus we get \(\lambda_{1}=0.66\) per hour

\[
\left.\begin{array}{l}
T_{1}=1.05 \text { hour } \\
T_{2}=9.9 \text { hours }
\end{array}\right\} \text { half-lives }
\]

Ratio

\[
\frac{N_{1}}{N_{2}}=\frac{51.1}{10.0} \times \frac{\lambda_{2}}{\lambda_{1}}=0.54
\]

The answer given in the book is misleading.\\
6.228 Production of the nucleus is governed by the equation\\
\includegraphics[max width=\textwidth, alt={}, center]{da0156b6-4133-434e-90d4-ec01b2070afe-372_181_357_107_568}

We see that \(N\) will approach a constant value \(\frac{g}{\lambda}\). This can also be proved directly. Multiply by \(e^{\lambda t}\) and write

Then

\[
\begin{gathered}
\frac{d N}{d t} e^{\lambda t}+\lambda e^{\lambda t} N=g e^{\lambda t} \\
\frac{d}{d t}\left(N e^{\lambda t}\right)=g e^{\lambda t} \\
N e^{\lambda t}=\frac{g}{\lambda} e^{\lambda t}+\mathrm{const}
\end{gathered}
\]

or\\
At \(t=0\) when the production is starteed, \(N=0\)

\[
\begin{aligned}
0 & =\frac{g}{\lambda}+\text { constant } \\
N & =\frac{g}{\lambda}\left(1-e^{-\lambda t}\right)
\end{aligned}
\]

Hence

\[
A=\lambda N=g\left(1-e^{-\lambda t}\right)
\]

From the problem

\[
\frac{1}{2 \cdot 7}=1-e^{-\lambda t}
\]

This gives \(\lambda t=0.463\)\\
so

\[
t=\frac{0.463}{\lambda}=\frac{0.463 \times T}{0.693}=9.5 \text { days } .
\]

Algebraically

\[
t=-\frac{T}{\ln 2} \ln \left(1-\frac{A}{g}\right)
\]

6.229 (a) Suppose \(N_{1}\) and \(N_{2}\) are the number of two radionuclides \(A_{1}, A_{2}\) at time \(t\). Then


\begin{gather*}
\frac{d N_{1}}{d t}=-\lambda_{1} N_{1}  \tag{1}\\
\frac{d N_{2}}{d t}=\lambda_{1} N_{1}-\lambda_{2} N_{2} \tag{2}
\end{gather*}


From (1)

\[
N_{1}=N_{10} e^{-\lambda_{1} t}
\]

where \(N_{10}\) is the initial number of nuclides \(A_{1}\) at time \(t=0\)\\
From (2)\\
or

\[
\left(\frac{d N_{2}}{d t}+\lambda_{2} N_{2}\right) e^{\lambda_{2} t}=\lambda_{1} N_{10} e^{-\left(\lambda_{1}-\lambda_{2}\right) t}
\]

since

\[
\left(N_{2} e^{\lambda_{2} t}\right)=\text { const } \frac{\lambda_{1} N_{10}}{\lambda_{1}-\lambda_{2}} e^{-\left(\lambda_{1}-\lambda_{2}\right) t}
\]

\[
N_{2}=0 \text { at } t=0
\]

Constant

\[
N_{2}=\frac{\lambda_{1} N_{10}}{\lambda_{1}-\lambda_{2}}
\]

Thus

\[
=\frac{\lambda_{1} N_{10}}{\lambda_{1}-\lambda_{2}}\left(e^{-\lambda_{2} t}-e^{-\lambda_{1} t}\right)
\]

(b) The activity of nuclide \(A_{2}\) is \(\lambda_{2} N_{2}\). This is maximum when \(N_{2}\) is maximum. That happens when \(\frac{d N_{2}}{d t}=0\)\\
This requires

\[
\lambda_{2} e^{-\lambda_{2} t_{m}}=\lambda_{1} e^{-\lambda_{1} t_{m}}
\]

or

\[
t_{m}=\frac{\ln \left(\lambda_{1} / \lambda_{2}\right)}{\lambda_{1}-\lambda_{2}}
\]

6.230 (a) This case can be obtained from the previous one on putting

\[
\lambda_{2}=\lambda_{1}-\varepsilon
\]

where \(\varepsilon\) is very small and letting \(\varepsilon \rightarrow 0\) at the end. Then

\[
N_{2}=\frac{\lambda_{1} N_{10}}{\varepsilon}\left(e^{\varepsilon t}-1\right) e^{-\lambda_{1} t}=\lambda_{1} t e^{-\lambda_{1} t} N_{10}
\]

or dropping the subscript 1 as the two values are equal

\[
N_{2}=N_{10} \lambda t e^{-\lambda t}
\]

(b) This is maximum when

\[
\frac{d N_{2}}{d t}=0 \quad \text { or } \quad t=\frac{1}{\lambda}
\]

6.231 Here we have the equations

\[
\begin{gathered}
\frac{d N_{1}}{d t}=-\lambda_{1} N_{1} \\
\frac{d N_{2}}{d t}=\lambda_{1} N_{1}-\lambda_{2} N_{2} \text { and } \frac{d N_{3}}{d t}=\lambda_{2} N_{2}
\end{gathered}
\]

From problem 229

Then

\[
\begin{gathered}
N_{1}=N_{10} e^{-\lambda_{1} t} \\
N_{2}=\frac{\lambda_{1} N_{10}}{\lambda_{1}-\lambda_{2}}\left(e^{-\lambda_{2} t}-e^{-\lambda_{1} t}\right) \\
\frac{d N_{3}}{d t}=\frac{\lambda_{1} \lambda_{2}}{\lambda_{1}-\lambda_{2}} N_{10}\left(e^{-\lambda_{2} t}-e^{-\lambda_{1} t}\right)
\end{gathered}
\]

or

\[
N_{3}=\text { Const }-\frac{\lambda_{1} \lambda_{2}}{\lambda_{1}-\lambda_{2}}\left(\frac{e^{-\lambda_{2} t}}{\lambda_{2}}-\frac{e^{-\lambda_{1} t}}{\lambda_{1}}\right) N_{10}
\]

since \(N_{3}=0\) initially

So

\[
\begin{gathered}
\text { Const }=\frac{\lambda_{1} \lambda_{2}}{\lambda_{1}-\lambda_{2}} N_{10}\left(\frac{1}{\lambda_{2}}-\frac{1}{\lambda_{1}}\right) \\
N_{3}=\frac{\lambda_{1} \lambda_{2} N_{10}}{\lambda_{1}-\lambda_{2}}\left[\frac{1}{\lambda_{2}}\left(1-e^{-\lambda_{2} t}\right)-\frac{1}{\lambda_{1}}\left(1-e^{-\lambda_{1} t}\right)\right] \\
=N_{10}\left[1+\frac{\lambda_{1} e^{-\lambda_{2} t}-\lambda_{2} e^{-\lambda_{1} t}}{\lambda_{2}-\lambda_{1}}\right]
\end{gathered}
\]

6.232 We have the chain

\[
\begin{gathered}
B_{i}^{210} \xrightarrow[\lambda_{1}]{\beta} \longrightarrow P_{0}^{210} \xrightarrow[\lambda_{2}]{\alpha} \longrightarrow P_{b}^{206} \\
A_{1} \\
A_{2}
\end{gathered}
\]

of the previous problem initially

\[
N_{10}=\frac{10^{-3}}{210} \times 6.023 \times 10^{23}=2.87 \times 10^{18}
\]

A month after preparation

\[
\begin{aligned}
& N_{1}=4.54 \times 10^{16} \\
& N_{2}=2.52 \times 10^{18}
\end{aligned}
\]

using the results of the previous problem.\\
Then

\[
\begin{aligned}
& A_{\beta}=\lambda_{1} N_{1}=0.725 \times 10^{11} \mathrm{dis} / \mathrm{sec} \\
& A_{\alpha}=\lambda_{2} N_{2}=1.46 \times 10^{11} \mathrm{dis} / \mathrm{sec}
\end{aligned}
\]

6.233 (a) Ra has \(Z=88, A=226\)

After \(5 \alpha\) emission and \(4 \beta\) (electron) emission

\[
A=206
\]

\[
Z=88+4-5 \times 2=82
\]

The product is \({ }^{\mathbf{8 2}} \mathbf{P b}^{\mathbf{2 0 6}}\)\\
(b) We require

\[
\begin{array}{r}
-\Delta Z=10=2 n-m \\
-\Delta A=32=n \times 4
\end{array}
\]

Here

\[
n=\text { no. of } \alpha \text { emissions }
\]

Thus

\[
n=8, m=6
\]

6.234 The momentum of the \(\alpha\)-particle is \(\sqrt{2 M_{\alpha} T}\). This is also the recoil momentum of the daughter nuclear in opposite direction. The recoil velocity of the daughter nucleus is

\[
\frac{\sqrt{2 M_{\alpha} T}}{M_{d}}=\frac{2}{196} \sqrt{\frac{2 T}{M_{p}}}=3.39 \times 10^{5} \mathrm{~m} / \mathrm{s}
\]

The energy of the daughter nucleus is \(\frac{M_{\alpha}}{M_{d}} T\) and this represents a fraction

\[
\frac{M_{\alpha} / M_{d}}{1+\frac{M_{\alpha}}{M_{d}}}=\frac{M_{\alpha}}{M_{\alpha}+M_{d}}=\frac{4}{200}=\frac{1}{50}=0.02
\]

of total energy. Here \(\boldsymbol{M}_{\boldsymbol{d}}\) is the mass of the daughter nucleus.\\
6.235 The number of nuclei initially present is

\[
\frac{10^{-3}}{210} \times 6.023 \times 10^{23}=2.87 \times 10^{18}
\]

In thé mean life time of these nuclei the number decaying is the fraction \(1-\frac{1}{e}=0.632\). Thus the energy released is

\[
2.87 \times 10^{18} \times 0.632 \times 5.3 \times 1.602 \times 10^{-13} \mathrm{~J} .=1.54 \mathrm{M} \mathrm{~J}
\]

6.236 We neglect all recoil effects. Then the following diagram gives the energy of the gamma ray

\[
P_{o} 210
\]

\includegraphics[max width=\textwidth, alt={}, center]{da0156b6-4133-434e-90d4-ec01b2070afe-375_324_706_1235_502}\\
6.237 (a) For an alpha particle with initial K.E. 7.0 MeV , the initial velocity is

Thus

\[
\begin{aligned}
v_{0} & =\sqrt{\frac{2 T}{M_{\alpha}}} \\
& =\sqrt{\frac{2 \times 7 \times 1.602 \times 10^{-6}}{4 \times 1.672 \times 10^{-24}}} \\
& =1.83 \times 10^{9} \mathrm{~cm} / \mathrm{sec} \\
R & =6.02 \mathrm{~cm}
\end{aligned}
\]

(b) Over the whole path the number of ion pairs is

\[
\frac{7 \times 10^{6}}{34}=2.06 \times 10^{5}
\]

Over the first half of the path :- We write the formula for the mean path as \(R \alpha E^{3 / 2}\) where \(E\) is the initial energy. Thus if the energy of the \(\alpha\)-particle after traversing the first half of the path is \(E_{1}\) then

\[
R_{0} E_{1}^{3 / 2}=\frac{1}{2} R_{0} E_{0}^{3 / 2} \quad \text { or } \quad E_{1}=2^{-2 / 2} / E_{0}
\]

Hence number of ion pairs formed in the first half of the path length is

\[
\frac{E_{0}-E_{1}}{34 e V}=\left(1-2^{-2 / 3}\right) \times 2.06 \times 10^{5}=0.76 \times 10^{5}
\]

6.238 In \(\boldsymbol{\beta}^{-}\)decay

\[
\begin{aligned}
Z^{X^{A}} & \rightarrow_{Z+1} Y^{A}+e^{-}+Q \\
Q & =\left(M_{x}-M_{y}-m_{e}\right) c^{2} \\
& =\left[\left(M_{x}+Z m_{e}\right)-\left(M_{y}+Z m_{e}+m_{e}\right)\right] c^{2} \\
& =\left(M_{p}-M_{d}\right) c^{2}
\end{aligned}
\]

since \(M_{p}, M_{d}\) are the masses of the atoms. The binding energy of the electrons in ignored. In \(K\) capture

In \(\boldsymbol{\beta}^{\boldsymbol{+}}\) decay

\[
\begin{aligned}
e_{K}^{-} & +{ }_{z} X^{A} \rightarrow{ }_{z-1} Y^{A}+Q \\
Q & =\left(M_{X}-M_{Y}\right) c^{2}+m_{e} c^{2} \\
& =\left(M_{x}^{c^{2}}+Z m_{e} c^{2}\right)-\left(M_{Y} c^{2}+(Z-1) m_{e} c^{2}\right) \\
& =c^{2}\left(M_{p}-M_{d}\right)
\end{aligned}
\]

Then

\[
\begin{aligned}
{ }_{z} X^{A} & \rightarrow z_{-1} Y^{A}+e^{+}+Q \\
Q & =\left(M_{x}-M_{y}-m_{e}\right) c^{2} \\
& =\left[M_{x}+Z m_{e}\right] c^{2}-\left[M_{y}+(Z-1) m_{e}\right] c^{2}-2 m_{e} c^{2} \\
& =\left(M_{p}-M_{d}-2 m_{e}\right) c^{2}
\end{aligned}
\]

6.239 The reaction is

\[
B e^{10} \rightarrow B^{10}+e^{-}+\bar{v}_{e}
\]

For maximum K.E. of electrons we can put the energy of \(\bar{v}_{e}\) to be zero. The atomic masses are

\[
\begin{aligned}
B e^{10} & =10.016711 \mathrm{amu} \\
B^{10} & =10.016114 \mathrm{amu}
\end{aligned}
\]

So the K.E. of electrons is (see previous problem)

\[
597 \times 10^{-6} \mathrm{amu} \times c^{2}=0.56 \mathrm{MeV}
\]

The momentum of electrons with this K.E. is \(0.941 \frac{M e V}{c}\)\\
and the recoil energy of the daughter is

\[
\frac{(0.941)^{2}}{2 \times M_{d} c^{2}}=\frac{(0.941)^{2}}{2 \times 10 \times 938} \mathrm{MeV}=47.2 \mathrm{eV}
\]

6.240 The masses are

\[
\mathrm{Na}^{24}=24-0.00903 \mathrm{amu} \text { and } \mathrm{Mg}^{24}=24-0.01496 \mathrm{amu}
\]

The reaction is

\[
\mathrm{Na}^{24} \rightarrow \mathrm{Mg}^{24}+e^{-}+\bar{v}_{e}
\]

The maximum K.E. of electrons is

\[
0.00593 \times 93 \mathrm{MeV}=5.52 \mathrm{MeV}
\]

Average K.E. according to the problem is then \(\frac{5 \cdot 52}{3}=1.84 \mathrm{MeV}\)\\
The initial number of \(\mathrm{Na}^{24}\) is

\[
\frac{10^{-3} \times 6.023 \times 10^{23}}{24}=2.51 \times 10^{19}
\]

The fraction decaying in a day is

\[
1-(2)^{-24 / 15}=0.67
\]

Hence the heat produced in a day is\\
\(0.67 \times 2.51 \times 10^{19} \times 1.84 \times 1.602 \times 10^{-13}\) Joul \(=4.95 \mathrm{MJ}\)\\
6.241 We assume that the parent nucleus is at rest. Then since the daughter nucleus does not recoil, we have

\[
\vec{p}=-\overrightarrow{p_{v}}
\]

i.e. positron \& \(v\) mometum are equal and opposite. On the other hand\\
\(\sqrt{c^{2} p^{2}+m_{e}^{2}} c^{4}+c p=Q=\) total energy released. (Here we have used the fact that energy of the neutrino is \(c\left|\overrightarrow{p_{v}}\right|=c p\) )

Now

\[
\begin{aligned}
Q & =\left[\left(\text { Mass of } C^{\|} \text {nucleus }\right)-\left(\text { Mass of } B \|_{\text {nucleus }}\right)\right] c^{2} \\
& =\left[\text { Mass of } C^{\|} \text {atom }- \text { Mass of } B^{\|} \text {atom }-m_{e}\right] c^{2} \\
& =0.00213 \mathrm{amu} \times c^{2}-m_{e} c^{2} \\
& =(0.00213 \times 931-0.511) \mathrm{MeV}=1.47 \mathrm{MeV}
\end{aligned}
\]

Then

\[
c^{2} p^{2}+(0.511)^{2}=(1.47-c p)^{2}=(1.47)^{2}-2.94 c p+c^{2} p^{2}
\]

Thus \(c p=0.646 \mathrm{MeV}=\) energy of neutrino\\
Also K.E. of electron \(=1.47-0.646-0.511=0.313 \mathrm{MeV}\)\\
6.242 The K.E. of the positron is maximum when the energy of neutrino is zero. Since the recoil energy of the nucleus is quite small, it can be calculated by successive approximation. The reaction is

\[
N^{13} \rightarrow C^{13}+e^{+}+v_{e}
\]

The maximum energy available to the positron (including its rest energy) is

\[
\begin{aligned}
& c^{2}\left(\text { Mass of } N^{13} \text { nucleus - Mass of } C^{13} \text { nucleus }\right) \\
= & c^{2}\left(\text { Mass of } N^{13} \text { atom - Mass of } C^{13} \text { atom - } m_{e}\right) \\
= & 0.00239 c^{2}-m_{e} c^{2} \\
= & (0.00239 \times 931-0.511) \mathrm{MeV} \\
= & 1.71 \mathrm{MeV}
\end{aligned}
\]

The momentum corresponding to this energy is \(1.636 \mathrm{MeV} / \mathrm{c}\).\\
The recoil energy of the nucleus is then

\[
E=\frac{p^{2}}{2 M}=\frac{(1.636)^{2}}{2 \times 13 \times 931}=111 \mathrm{eV}=0.111 \mathrm{keV}
\]

on using \(\boldsymbol{M} c^{\mathbf{2}} \boldsymbol{=} \mathbf{1 3} \boldsymbol{\times} \mathbf{9 3 1 ~ M e V}\)\\
6.243 The process is

\[
e_{k}^{-}+B e^{7} \rightarrow L i^{7}+v
\]

The energy available in the process is

\[
\begin{aligned}
Q & =c^{2}\left(\text { Mass of } B e^{7} \text { atom }- \text { Mass of } L i^{7} \text { atom }\right) \\
& =0.00092 \times 931 \mathrm{MeV}=0.86 \mathrm{MeV}
\end{aligned}
\]

The momentum of a \(K\) electron is negligible. So in the rest frame of the \(B e^{7}\) atom, most of the energy is taken by neutrino whose momentum is very nearly \(0.86 \mathrm{MeV} / \mathrm{c}\)\\
The momentum of the recoiling nucleus is equal and opposite. The velocity of recoil is

\[
\frac{0.86 \mathrm{MeV} / \mathrm{c}}{M_{L i}}=c \times \frac{0.86}{7 \times 931}=3.96 \times 10^{6} \mathrm{~cm} / \mathrm{s}
\]

6.244 In internal conversion, the total energy is used to knock out \(K\) electrons. The K.E. of these electrons is energy available-B.E. of \(K\) electrons

\[
=(87-26)=61 \mathrm{keV}
\]

The total energy including rest mass of electrons is \(0.511+0.061=0.572 \mathrm{MeV}\)\\
The momentum corresponding to this total energy is \(\sqrt{(0.572)^{2}-(0.511)^{2}} / c=0257 \mathrm{MeV} / \mathrm{c}\).

The velocity is then

\[
\frac{c^{2} p}{E}=c \times \frac{0.257}{0.572}=0.449 c
\]

6.245 With recoil neglected, the \(\gamma\)-quantrum will have 129 keV energy. To a first approximation, its momentrum will be \(129 \mathrm{keV} / \mathrm{c}\) and the energy of recoil will be

\[
\frac{(0.129)^{2}}{2 \times 191 \times 931} \mathrm{MeV}=4.18 \times 10^{-8} \mathrm{MeV}
\]

In the next approximation we therefore write \(E_{\gamma} \simeq \cdot 129-8 \cdot 2 \times 10^{-8} \mathrm{MeV}\)\\
Therefore

\[
\frac{\delta E_{\gamma}}{E_{\gamma}}=3.63 \times 10^{-7}
\]

6.246 For maximum (resonant) absorption, the absorbing nucleus must be moving with enough speed to cancel the momentum of the oncoming photon and have just right energy ( \(\varepsilon=129 \mathrm{keV}\) ) available for transition to the excited state.\\
\includegraphics[max width=\textwidth, alt={}, center]{da0156b6-4133-434e-90d4-ec01b2070afe-379_190_738_660_385}

Since \(\delta E_{\gamma} \sim \frac{\varepsilon^{2}}{2 M c^{2}}\) and momentum of the photon is \(\frac{\varepsilon}{c}\), these condition can be satisfied if the velocity of the nucleus is

\[
\frac{\varepsilon}{M c}=c \frac{\varepsilon}{M c^{2}}=218 \mathrm{~m} / \mathrm{s}=0.218 \mathrm{k} \mathrm{~m} / \mathrm{s}
\]

6.247 Because of the gravitational shift the frequency of the gamma ray at the location of the absorber is increased by

\[
\frac{\delta \omega}{\omega}=\frac{g h}{c^{2}}
\]

For this to be compensated by the Doppler shift (assuming that resonant absorption is possible in the absence of gravitational field) we must have

\[
\frac{g h}{c^{2}}=\frac{v}{c} \quad \text { or } \quad v=\frac{g h}{c}=0.65 \mu \mathrm{~m} / \mathrm{s}
\]

6.248 The natural life time is

\[
\Gamma=\frac{\hbar}{\tau}=4.7 \times 10^{-10} \mathrm{eV}
\]

Thus the condition \(\delta E_{\gamma} \geq \Gamma\) implies \(\frac{g h}{c^{2}} \geq \frac{\Gamma}{\varepsilon}=\frac{\hbar}{\tau \varepsilon}\)\\
or

\[
h \geq \frac{c^{2} \hbar}{\tau \varepsilon g}=4.64 \text { metre }
\]

( \(h\) here is height of the place, not planck's constant.)

\subsection*{6.6 NUCLEAR REACTIONS}
6.249 Initial momentum of the \(\alpha\) particle is \(\sqrt{2 m T_{\alpha}} \hat{i}\) (where \(\hat{i}\) is a unit vector in the incident direction). Final momenta are respectively \(\vec{p}_{\alpha}\) and \(\vec{p}_{L i}\). Conservation of momentum reads

Squaring


\begin{gather*}
\overrightarrow{p_{\alpha}}+\overrightarrow{p_{L i}}=\sqrt{2 m T_{\alpha}} \hat{i} \\
p_{\alpha}^{2}+p_{L i}^{2}+2 p_{\alpha} p_{L i} \cos \Theta=2 m T_{\alpha} \tag{1}
\end{gather*}


where \(\Theta\) is the angle between \(\overrightarrow{p_{\alpha}}\) and \(\overrightarrow{p_{L i}}\).\\
Also by energy conservation

\[
\frac{p_{\alpha}^{2}}{2 m}+\frac{p_{L_{t}}^{2}}{2 M}=T_{\alpha}
\]

( \(m \& M\) are respectively the masses of \(\alpha\) particle and \(L i^{6}\).) So


\begin{equation*}
p_{\alpha}^{2}+\frac{m}{M} p_{L i}^{2}=2 m T_{\alpha} \tag{2}
\end{equation*}


Substracting (2) from (1) we see that

\[
p_{L i}\left[\left(1-\frac{m}{M}\right) p_{L i}+2 p_{\alpha} \cos \Theta\right]=0
\]

Thus if

\[
p_{L i} \neq 0
\]

\[
p_{\alpha}=-\frac{1}{2}\left(1-\frac{m}{M}\right) p_{L i} \sec \Theta
\]

Since \(p_{\alpha}, p_{L i}\) are both positive number (being magnitudes of vectors) we must have

\[
-1 \leq \cos \Theta<0 \text { if } m<M
\]

This being understood, we write

\[
\frac{p_{L i}^{2}}{2 M}\left[1+\frac{M}{4 m}\left(1-\frac{m}{M}\right)^{2} \sec ^{2} \Theta\right]=T_{\alpha}
\]

Hence the recoil energy of the \(L_{i}\) nucleus is

\[
\frac{p_{L i}^{2}}{2 M}=\frac{T_{\alpha}}{1+\frac{(M-m)^{2}}{4 m M} \sec ^{2} \Theta}
\]

As we pointed out above \(\Theta \neq 60^{\circ}\). If we take \(\Theta=120^{\circ}\), we get recoil energy of \(L i=6 \mathrm{MeV}\)\\
6.250 (a) In a head on collision

\[
\begin{gathered}
\sqrt{2 m T}=p_{d}+p_{n} \\
T=\frac{p_{d}^{2}}{2 M}+\frac{p_{n}^{2}}{2 m}
\end{gathered}
\]

Where \(p_{d}\) and \(p_{n}\) are the momenta of deuteron and neutron after the collision. Squaring

\[
p_{d}^{2}+p_{n}^{2}+2 p_{d} p_{n}=2 m T
\]

\[
p_{n}^{2}+\frac{m}{M} p_{d}^{2}=2 m T
\]

or since \(p_{d} \neq 0\) in a head on collisions

\[
p_{n}=-\frac{1}{2}\left(1-\frac{m}{M}\right) p_{d}
\]

Going back to energy conservation

So

\[
\begin{gathered}
\frac{p_{d}^{2}}{2 M}\left[1+\frac{M}{4 m}\left(1-\frac{m}{M}\right)^{2}\right]=T \\
\frac{p_{d}^{2}}{2 M}=\frac{4 m M}{(m+M)^{2}} T
\end{gathered}
\]

This is the energy lost by neutron. So the fraction of energy lost is

\[
\eta=\frac{4 m M}{(m+M)^{2}}=\frac{8}{9}
\]

(b) In this case neutron is scattered by \(90^{\circ}\). Then\\
we have from the diagram\\
\(\overrightarrow{p_{d}}=p_{n} \hat{j}+\sqrt{2 m T} \hat{i}\)\\
Then by energy conservation\\
\(\frac{p_{n}^{2}+2 m T}{2 M}+\frac{p_{n}^{2}}{2 m}=T\)\\
or \(\quad \frac{p_{n}^{2}}{2 m}\left(1+\frac{m}{M}\right)=T\left(1-\frac{m}{M}\right)\)\\
or \(\quad \frac{p_{n}^{2}}{2 m}=\frac{M-m}{M+m} \cdot T\)\\
\includegraphics[max width=\textwidth, alt={}, center]{da0156b6-4133-434e-90d4-ec01b2070afe-381_428_562_895_759}

The energy lost by neutron in then

\[
T-\frac{p_{n}^{2}}{2 m}=\frac{2 m}{M+m} T
\]

or fraction of energy lost is \(\eta=\frac{2 m}{M+m}=\frac{2}{3}\)\\
6.251 From conservation of momentum

\[
\sqrt{2 M T} \hat{i}=\overrightarrow{p_{d}}+\overrightarrow{p_{p}}
\]

or \(p_{p}^{2}=2 M T+p_{d}^{2}-2 \sqrt{2 M T} p_{d} \cos \theta\)\\
From energy conservation

\[
T=\frac{p_{d}^{2}}{2 M}+\frac{p_{p}^{2}}{2 m}
\]

\includegraphics[max width=\textwidth, alt={}, center]{da0156b6-4133-434e-90d4-ec01b2070afe-381_374_663_1587_616}\\
( \(M=\) mass of denteron, \(m=\) mass of proton)\\
So

\[
p_{p}^{2}=2 m T-\frac{m}{M} p_{d}^{2}
\]

Hence

\[
p_{d}^{2}\left(1+\frac{m}{M}\right)-2 \sqrt{2 M T} p_{d} \cos \theta+2(M-m) T=0
\]

For real roots

\[
\begin{gathered}
4(2 M T) \cos ^{2} \theta-4 \times 2(M-m) T\left(1+\frac{m}{M}\right) \geq 0 \\
\cos ^{2} \theta \geq\left(1-\frac{m^{2}}{M^{2}}\right)
\end{gathered}
\]

Hence

\[
\sin ^{2} \theta \leq \frac{m^{2}}{M^{2}}
\]

i.e.

\[
\theta \leq \sin ^{-1} \frac{m}{M}
\]

For deuteron-proton scaltering \(\theta_{\text {max }}=30^{\circ}\).\\
6.252 This problem has a misprint. Actually the radius \(R\) of a nucleus is given by

\[
R=1.3 \sqrt[3]{A} \mathrm{fm}
\]

where

\[
f m=10^{-15} \mathrm{~m}
\]

Then the number of nucleous per unit volume is

\[
\frac{A}{\frac{4 \pi}{3} R^{3}}=\frac{3}{4 \pi} \times(1.3)^{-3} \times 10^{+39} \mathrm{~cm}^{-3}=1.09 \times 10^{38} \text { per cc }
\]

The corresponding mass density is

\[
\left(1.09 \times 10^{38} \times \text { mass of a nucleon }\right) \text { per } \mathrm{cc}=1.82 \times 10^{11} \mathrm{~kg} / \mathrm{cc}
\]

6.253 (a) The particle \(x\) must carry two nucleons and a unit of positive charge.

The reaction is

\[
B^{10}(d, \alpha) B_{e}^{8}
\]

(b) The particle \(x\) must contain a proton in addition to the constituents of \(O^{17}\). Thus the reaction is

\[
O^{17}(d, n) F^{18}
\]

(c) The particle \(x\) must carry nucleon number 4 and two units of + ve charge. Thus the particle must be \(x=\alpha\) and the reaction is

\[
N a^{23}(p, \alpha) N e^{20}
\]

(d) The particle \(x\) must carry mass number 37 and have one unit less of positive charge. Thus \(x=\mathrm{Cl}^{37}\) and the reaction is

\[
C l^{37}(p, n) A r^{37}
\]

6.254 From the basic formula

We define

\[
\begin{aligned}
& E_{b}=Z m_{H}+(A-Z) m_{n}-M \\
& \Delta_{H}=m_{H}-1 \mathrm{amu} \\
& \Delta_{n}=m_{n}-1 \mathrm{amu} \\
& \Delta=M-A \mathrm{amu}
\end{aligned}
\]

Then clearly \(E_{b}=Z \Delta_{H}+(A-Z) \Delta_{n}-\Delta\)\\
6.255 The mass number of the given nucleus must be

\[
27 /\left(\frac{3}{2}\right)^{3}=8
\]

Thus the nucleus is \(B e^{8}\). Then the binding energy is

\[
\begin{gathered}
E_{b}=4 \times 0.00867+4+\times 0.00783-0.00531 \mathrm{amu} \\
=0.06069 \mathrm{amu}=56.5 \mathrm{MeV}
\end{gathered}
\]

On using \(1 \mathrm{amu}=931 \mathrm{MeV}\).\\
6.256 (a) Total binding energy of the \(O^{16}\) nucleus is

\[
\begin{gathered}
E_{b}=8 \times .00867+8 \times .00783+0.00509 \mathrm{amu} \\
=0.13709 \mathrm{amu}=127.6 \mathrm{MeV}
\end{gathered}
\]

So B.E. per nucleon is \(\mathbf{7 . 9 8 ~ M e v}\) / nucleon\\
(b) B.E. of neutron in \(B^{11}\) nucleus

\[
\text { = B.E. of } B^{11} \text { - B.E. of } B^{10}
\]

(since on removing a neutron from \(B^{11}\) we get \(B^{10}\) )

\[
\begin{gathered}
=\Delta_{n}-\Delta_{B_{11}}+\Delta_{B_{10}}=0.0867-0.0930+.01294 \\
=0.01231 \mathrm{amu}=11.46 \mathrm{MeV}
\end{gathered}
\]

B.E. of ( an \(\boldsymbol{\alpha}\)-particle in \(\boldsymbol{B}^{\mathbf{1 1}}\) )

\[
\text { = B.E. of } B^{1} \text { - B.E. of } L i^{7} \text { - B.E. of } \alpha
\]

(since on removing an \(\alpha\) from \(B^{11}\) we get \(L i^{7}\) )

\[
\begin{gathered}
=-\Delta_{B_{11}}+\Delta_{L \varphi_{1}}+\Delta_{\alpha} \\
=-0.00930+0.01601+0.00260 \\
=0.00931 \mathrm{amu}=8.67 \mathrm{MeV}
\end{gathered}
\]

(c) This energy is\\[0pt]
[B.E. of \(O^{16}+4\) (B.E. of \(\alpha\) particles)]

\[
\begin{aligned}
& =-\Delta_{O^{16}}+4 \Delta_{\alpha} \\
& =4 \times 0.00260+0.00509 \\
& =0.01549 \mathrm{amu}=14.42 \mathrm{MeV}
\end{aligned}
\]

6.257 B.E. of a neutron in \(B^{11}\) - B.E. of a proton in \(B^{11}\)

\[
\begin{gathered}
=\left(\Delta_{n}-\Delta_{B}^{11}+\Delta_{B}^{10}\right)-\left(\Delta_{p}-\Delta_{B}^{11}+\Delta_{B_{e}^{10}}\right) \\
=\Delta_{n}-\Delta_{p}+\Delta_{B}^{10}-\Delta_{B_{e}^{10}}=0.00867-0.00783 \\
+0.01294-0.01354=0.00024 \mathrm{amu}=0.223 \mathrm{MeV}
\end{gathered}
\]

The difference in binding energy is essentially due to the coulomb repulsion between the proton and the residual nucleus \(B e^{10}\) which together constitute \(B^{11}\).\\
6.258 Required energy is simply the difference in total binding energies-

\[
\begin{gathered}
=\text { B.E. of } N e^{20}-2\left(\text { B.E. of } H e^{4}\right)-\text { B.E. of } C^{12} \\
=20 \varepsilon_{N e}-8 \varepsilon_{\alpha}-12 \varepsilon_{C}
\end{gathered}
\]

( \(\varepsilon\) is binding energy per unit nucleon.)\\
Substitution gives

\[
11.88 \mathrm{MeV} \text {. }
\]

6.259 (a) We have for \(L i^{8}\)

\[
41.3 \mathrm{MeV}=0.044361 \mathrm{amu}=3 \Delta_{H}+5 \Delta_{n}-\Delta
\]

Hence

\[
\Delta=3 \times 0.00783+5 \times 0.00867-0.09436=0.02248 \mathrm{amu}
\]

(b) For \(C^{10}\)

\[
\begin{gathered}
10 \times 6.04=60.4 \mathrm{MeV} \\
=0.06488 \mathrm{amu} \\
=6 \Delta_{H}+4 \Delta_{n}-\Delta
\end{gathered}
\]

Hence \(\quad \Delta=6 \times 0.00783+4 \times 0.00867-0.06488=0.01678 \mathrm{amu}\)\\
Hence the mass of \(C^{10}\) is\\
10.01678 amu\\
6.260 Suppose \(M_{1}, M_{2}, M_{3}, M_{4}\) are the rest masses of the nuclei \(A_{1}, A_{2}, A_{3}\) and \(A_{4}\) perticipating in the reaction

\[
A_{1}+A_{2} \rightarrow A_{3}+A_{4}+Q
\]

Here \(Q\) is the energy released. Then by conservation of energy.

\[
Q=c^{2}\left(M_{1}+M_{2}-M_{3}-M_{4}\right)
\]

Now

\[
M_{1} c^{2}=c^{2}\left(Z_{1} m_{H}+\left(A_{1}-Z\right) m_{n}\right)-E_{1} \text { etc. and }
\]

\[
Z_{1}+Z_{2}=Z_{3}+Z_{4}(\text { conservation of change })
\]

\[
A_{1}+A_{2}=A_{3}+A_{4} \text { (conservation of heavy particles) }
\]

Hence

\[
Q=\left(E_{3}+E_{4}\right)-\left(E_{1}+E_{2}\right)
\]

6.261 (a) the energy liberated in the fission of 1 kg of \(U^{235}\) is

\[
\frac{1000}{235} \times 6.023 \times 10^{23} \times 200 \mathrm{MeV}=8.21 \times 10^{10} \mathrm{~kJ}
\]

The mass of coal with equivalent calorific value is

\[
\frac{8.21 \times 10^{10}}{30000} \mathrm{~kg}=2.74 \times 10^{6} \mathrm{~kg}
\]

(b) The required mass' is

\[
\frac{30 \times 10^{9} \times 4.1 \times 10^{3}}{200 \times 1.602 \times 10^{-13} \times 6.023 \times 10^{23}} \times \frac{235}{1000} \mathrm{~kg}=1.49 \mathrm{~kg}
\]

6.262 The reaction is (in effect).

Then

\[
\begin{aligned}
H^{2} & +H^{2} \rightarrow H e^{4}+Q \\
Q & =2 \Delta_{H^{2}}-\Delta_{H e}{ }^{4}+Q \\
& =0.02820-0.00260 \\
& =0.02560 \mathrm{amu}=23.8 \mathrm{MeV}
\end{aligned}
\]

Hence the energy released in 1 gm of \(\mathrm{He}^{4}\) is

\[
\frac{6.023 \times 10^{23}}{4} \times 23.8 \times 16.02 \times 10^{-13} \mathrm{Joule}=5.75 \times 10^{8} \mathrm{~kJ}
\]

This energy can be derived from

\[
\frac{5.75 \times 10^{8}}{30000} \mathrm{~kg}=1.9 \times 10^{4} \mathrm{~kg} \text { of Coal. }
\]

6.263 The energy released in the reaction\\
is

\[
\begin{gathered}
L i^{6}+H^{2} \rightarrow 2 H e^{4} \\
\Delta_{L i}^{6}+\Delta_{H}^{2}-2 \Delta_{H e}^{4} \\
=0.01513+0.01410-2 \times 0.00260 \mathrm{amu} \\
=0.02403 \mathrm{amu}=22.37 \mathrm{MeV}
\end{gathered}
\]

(This result for change in B.E. is correct because the contribution of \(\Delta_{n}\) \& \(\Delta_{H}\) cancels out by conservation law for protons \& neutrons.)-\\
Energy per nucleon is then

\[
\frac{22.37}{8}=2.796 \mathrm{MeV} / \text { nucleon } .
\]

This should be compared with the value \(\frac{200}{235}=0.85 \mathrm{MeV} /\) nucleon\\
6.264 The energy of reaction

\[
\mathrm{Li}^{7}+p \rightarrow 2 \mathrm{He}^{4}
\]

is,

\[
2 \times \text { B.E. of } \mathrm{He}^{4}-\text { B.E. of } \mathrm{Li}^{7}
\]

\[
=8 \varepsilon_{\alpha}-7 \varepsilon_{L i}=8 \times 7.06-7 \times 5.60=17.3 \mathrm{MeV}
\]

6.265 The reaction is \(N^{14}(\alpha, p) O^{17}\).

It is given that (in the Lab frame where \(N^{14}\) is at rest) \(T_{\alpha}=4.0 \mathrm{MeV}\).\\
The momentum of incident \(\alpha\) particle is

\[
\sqrt{2 m_{\alpha} T_{\alpha}} \hat{i}=\sqrt{2 \eta_{\alpha} m_{0} T_{\alpha}} \hat{i}
\]

The momentum of outgoing proton is

\[
\begin{aligned}
& \sqrt{2 m_{p} T_{p}}(\cos \theta \hat{i}+\sin \theta \hat{j}) \\
& =\sqrt{2 \eta_{p} m_{0} T_{p}}(\cos \theta \hat{i}+\sin \theta \hat{j})
\end{aligned}
\]

where \(\eta_{p}=\frac{m_{p}}{m_{0}}, \eta_{\alpha}=\frac{m_{\alpha}}{m_{0}}\),\\
and \(m_{0}\) is the mass of \(O^{17}\).\\
The momentum of \(O^{17}\) is\\
\includegraphics[max width=\textwidth, alt={}, center]{da0156b6-4133-434e-90d4-ec01b2070afe-386_356_667_166_731}

\[
\begin{aligned}
\left(\sqrt{2 \eta_{\alpha} m_{0} T_{\alpha}}\right. & \left.-\sqrt{2 \eta_{p} m_{0} T_{p}} \cos \theta\right) \hat{i} \\
& -\sqrt{2 m_{0} \eta_{p} T_{p}} \sin \theta \hat{j}
\end{aligned}
\]

By energy conservation (conservation of energy including rest mass energy and kinetic energy)

\[
\begin{aligned}
M_{14} c^{2}+ & M_{\alpha} c^{2}+T_{\alpha} \\
=M_{p} c^{2}+ & T_{p}+M_{17} c^{2} \\
& +\left[\left(\sqrt{\eta_{\alpha} T_{\alpha}}-\sqrt{\eta_{p} T_{p}} \cos \theta\right)^{2}+\eta_{p} T_{p} \sin ^{2} \theta+\eta_{p} T_{p} \sin ^{2} \theta\right]
\end{aligned}
\]

Hence by definition of the \(Q\) of reaction

\[
\begin{aligned}
Q= & M_{14} c^{2}+M_{\alpha} c^{2}-M_{p} c^{2}-M_{17} c^{2} \\
= & T_{p}+\eta_{\alpha} T_{\alpha}+\eta_{p} T_{p}-2 \sqrt{\eta_{p} \eta_{\alpha} T_{\alpha} T_{p}} \times \cos \theta-T_{\alpha} \\
= & \left(1+\eta_{p}\right) T_{p}+T_{\alpha}\left(1-\eta_{\alpha}\right) \\
& -2 \sqrt{\eta_{p} \eta_{\alpha} T_{\alpha} T_{p}} \cos \theta=-1.19 \mathrm{MeV}
\end{aligned}
\]

6.266 (a) The reaction is \(L i^{7}(p, n) B e^{7}\) and the energy of reaction is

\[
\begin{aligned}
Q & =\left(M_{B e}^{7}+M_{L i}^{7}\right) c^{2}+\left(M_{p}-M_{n}\right) c^{2} \\
& =\left(\Delta_{L i_{7}}-\Delta_{B e}^{7}\right) c^{2}+\Delta_{p}-\Delta_{n} \\
& =[0.01601+0.00783-0.01693-0.00867] \text { amu } \times c^{2} \\
& =-1.64 \mathrm{MeV}
\end{aligned}
\]

(b) The reaction is \(B e^{9}(n, \gamma) B e^{10}\).

Mass of \(\gamma\) is taken zero. Then

\[
\begin{aligned}
Q & =\left(M_{B e^{9}}+M_{n}-M_{B e^{10}}\right) c^{2} \\
& =\left(\Delta_{B e^{9}}+\Delta_{n}-\Delta_{B e^{10}}\right) c^{2} \\
& =(0.01219+0.00867-0.01354), \text { amu } \times c^{2} \\
& =6.81 \mathrm{MeV}
\end{aligned}
\]

(c) The reaction is \(L i^{7}(\alpha, n) B^{10}\). The energy is

\[
\begin{aligned}
Q & =\left(\Delta_{L i}{ }^{7}+\Delta_{\alpha}-\Delta_{n}-\Delta_{B}{ }^{10}\right) c^{2} \\
& =(0.01601+0.00260-0.00867-.01294) \text { amu } \times c^{2} \\
& =-2.79 \mathrm{MeV}
\end{aligned}
\]

(d) The reaction is \(O^{16}(d, \alpha) N^{14}\). The energy of reaction is

\[
\begin{aligned}
Q & =\left(\Delta_{O}^{16}+\Delta_{d}-\Delta_{\alpha}-\Delta_{N}^{14}\right) c^{2} \\
& =(-0.00509+0.01410-0.00260-0.00307) \mathrm{amu} \times c^{2} \\
& =3.11 \mathrm{MeV}
\end{aligned}
\]

6.267 The reaction is \(B^{10}(n, \alpha)\) Li \(^{7}\). The energy of the reaction is

\[
\begin{aligned}
Q & =\left(\Delta_{B}^{10}+\Delta_{n}-\Delta_{\alpha}-\Delta_{L i}^{7}\right) c^{2} \\
& =(0.01294+0.00867-0.00260-0.01601) \mathrm{amu} \times c^{2} \\
& =2.79 \mathrm{MeV}
\end{aligned}
\]

Since the incident neutron is very slow and \(B^{10}\) is stationary, the final total momentum must also be zero. So the reaction products must emerge in opposite directions. If their speeds are, repectively, \(v_{\alpha}\) and \(v_{L i}\)\\
then

\[
4 v_{\alpha}=7 v_{L i}
\]

and

\[
\frac{1}{2}\left(4 v_{\alpha}^{2}+7 v_{L i}^{2}\right) \times 1.672 \times 10^{-24}=2.79 \times 1.602 \times 10^{-6}
\]

So

\[
\frac{1}{2} \times 4 v_{\alpha}^{2}\left(1+\frac{4}{7}\right)=2.70 \times 10^{18} \mathrm{~cm}^{2} / \mathrm{s}^{2}
\]

or

\[
v_{\alpha}=9.27 \times 10^{6} \mathrm{~m} / \mathrm{s}
\]

Then

\[
v_{L i}=5.3 \times 10^{6} \mathrm{~m} / \mathrm{s}
\]

\(6.268 Q\) of this reaction \(\left(\mathrm{Li}^{7}(p, n) \mathrm{Be}^{7}\right)\) was calculated in problem \(266(\mathrm{a})\). If is -1.64 MeV .\\
We have by conservation of momentum and energy \(p_{p}=p_{B e}\) (since initial \(L i\) and final neutron are both at rest)

\[
\frac{p_{p}^{2}}{2 m_{p}}=\frac{p_{B e}^{2}}{2 m_{L i}}+1.64
\]

Then

\[
\frac{P_{p}^{2}}{2 m_{p}}\left(1-\frac{m_{p}}{m_{B e}}\right)=1.64
\]

Hence

\[
T_{p}=\frac{p_{p}^{2}}{2 m_{p}}=\frac{7}{6} \times 1.64 \mathrm{MeV}=1.91 \mathrm{MeV}
\]

6.269 It is understood that \(B e^{9}\) is initially at rest. The moment of the outgoing neutron is \(\sqrt{2 m_{n} T_{n}} \hat{j}\). The momenturm of \(C^{12}\) is\\
\(\sqrt{2 m_{\alpha} T} \hat{i}-\sqrt{2 m_{n} T_{n}} \hat{j}\)\\
Then by energy conservation

\[
T+Q=T_{n}+\frac{2 m_{\alpha} T+2 m_{n} T_{n}}{2 m_{c i}}
\]

( \(m_{c}\) is the mass of \(C^{12}\) )\\
\includegraphics[max width=\textwidth, alt={}, center]{da0156b6-4133-434e-90d4-ec01b2070afe-388_367_703_98_684}

Thus \(\quad T_{n}=\frac{m_{c}(T+Q)-m_{\alpha} T}{m_{c}+m_{n}}\)

\[
=\frac{\left(m_{c}-m_{\alpha}\right) T+m_{c} Q}{m_{c}+m_{n}}=\frac{Q+\left(1-\frac{m_{\alpha}}{m_{c}}\right) T}{1+\frac{m_{n}}{m_{c}}}=8.52 \mathrm{MeV}
\]

6.270 The \(Q\) value of the reaction \(\operatorname{Li}^{7}(p, \alpha) H e^{4}\) is

\[
\begin{aligned}
Q & =\left(\Delta_{L i}{ }^{7}+\Delta_{H}-2 \Delta_{H e}{ }^{4}\right) c^{2} \\
& =(0.01601+0.00783-0.00520) \mathrm{amu} \times c^{2} \\
& =0.01864 \mathrm{amu} \times c^{2}=17.35 \mathrm{MeV}
\end{aligned}
\]

Since the direction of \(\mathrm{He}^{4}\) nuclei is symmetrical, their momenta must also be equal. Let \(T\) be the K.E. of each \(H e^{4}\). Then

\[
p_{p}=2 \sqrt{2 m_{H e} T} \cos \frac{\theta}{2}
\]

( \(\boldsymbol{p}_{\boldsymbol{p}}\) is the momentum of proton). Also

\[
\frac{p_{p}^{2}}{2 m_{p}}+Q=2 T=T_{p}+Q
\]

\begin{center}
\includegraphics[max width=\textwidth, alt={}]{da0156b6-4133-434e-90d4-ec01b2070afe-388_351_616_1127_731}
\end{center}

Hence

\[
\begin{aligned}
T_{p}+Q & =2 \frac{p_{p}^{2} \sec ^{2} \frac{\theta}{2}}{8 m_{H e}} \\
& =T_{p} \frac{m_{p}}{2 m_{H e}} \sec ^{2} \frac{\theta}{2}
\end{aligned}
\]

Hence

\[
\cos \frac{\theta}{2}=\sqrt{\frac{m_{p}}{2 m_{H e}} \frac{T_{p}}{T_{p}+Q}}
\]

Substitution gives

Also

\[
\begin{gathered}
\theta=170.53^{\circ} \\
T=\frac{1}{2}\left(T_{p}+Q\right)=9.18 \mathrm{MeV}
\end{gathered}
\]

6.271 Energy required is minimum when the reaction products all move in the direction of the incident particle with the same velocity (so that the combination is at rest in the centre of mass frame). We then have

\[
\sqrt{2 m T_{t h}}=(m+M) \mathrm{v}
\]

(Total mass is constant in the nonrelativistic limit).

\[
\begin{gathered}
T_{t h}-|Q|=\frac{1}{2}(m+M) \mathrm{v}^{2}=\frac{m T_{t h}}{m+M} \\
T_{t h} \frac{M}{m+M}=|Q| \\
T_{t h}=\left(1+\frac{m}{M}\right)|Q|
\end{gathered}
\]

Hence\\
6.272 The result of the previous problem applies and we find that energy required to split a deuteron is

\[
T \geq\left(1+\frac{M_{p}}{M_{d}}\right) E_{b}=3.3 \mathrm{MeV}
\]

6.273 Since the reaction \(\mathrm{Li}^{7}(p, n) \mathrm{Be}^{7}(Q=-1.65 \mathrm{MeV})\) is initiated, the incident proton energy must be

\[
2\left(1+\frac{M_{p}}{M_{L i}}\right) \times 1.65=1.89 \mathrm{MeV}
\]

since the reaction \(B e^{9}(p, n) B^{9}(Q=-1.85 \mathrm{MeV})\) is not initiated,

\[
T \leq\left(1+\frac{M_{p}}{M_{V e}}\right) \times 1.85=2.06 \mathrm{MeV} \text { Thus } 1.89 \mathrm{MeV} \leq T_{p} \leq 2.06 \mathrm{MeV}
\]

6.274 We have \(4.0=\left(1+\frac{m_{n}}{M_{B}{ }^{11}}\right)|Q|\)\\
or

\[
Q=-\frac{11}{12} \times 4 \mathrm{MeV}=-3.67 \mathrm{MeV}
\]

6.275 The \(Q\) of the reaction \(\mathrm{Li}^{7}(p, n) \mathrm{Be}^{7}\) was calculated in problem 266 (a). It is -1.64 MeV Hence, the threshold K.E. of protons for initiating this reaction is

\[
T_{t h}=\left(1+\frac{m_{p}}{m_{L i}}\right)|Q|=\frac{8}{7} \times 1.64=1.87 \mathrm{MeV}
\]

For the reaction \(L i^{7}(p, d) L i^{6}\)\\
we find

\[
\begin{aligned}
Q & =\left(\Delta_{L i}{ }^{7}+\Delta_{M}-\Delta_{d}-\Delta_{L i}{ }^{6}\right) c^{2} \\
& =(0.01601+0.00783-0.01410-0.01513) \mathrm{amu} \times c^{2} \\
& =-5.02 \mathrm{MeV}
\end{aligned}
\]

The threshold proton energy for initiating this reaction is

\[
T_{t h}=\left(1+\frac{m_{p}}{m_{L i}{ }^{7}}\right) \times|Q|=5.73 \mathrm{MeV}
\]

6.276 The \(Q\) of \(\operatorname{Li}^{7}(\alpha, n) B^{10}\) was calculated in problem 266 (c). It is \(Q=2.79 \mathrm{MeV}\) Then the threshold energy of \(\alpha\)-particle is

\[
T_{t h}=\left(1+\frac{m_{\alpha}}{m_{L i}}\right)|Q|=\left(1+\frac{4}{7}\right) 2.79=4.38 \mathrm{MeV}
\]

The velocity of \(B^{10}\) in this case is simply the volocity of centre of mass :-

\[
v=\frac{\sqrt{2 m_{\alpha} T_{t h}}}{m_{\alpha}+m_{L i}}=\frac{1}{1+\frac{m_{L i}}{m_{\alpha}}} \sqrt{\frac{2 T_{t h}}{m_{\alpha}}}
\]

This is because both \(B^{10}\) and \(n\) are at rest in the \(C M\) frame at theshold.\\
Substituting the values of various quantities\\
we get

\[
v=5.27 \times 10^{6} \mathrm{~m} / \mathrm{s}
\]

6.277 The momentum of incident neutron is \(\sqrt{2 m_{n} T} \hat{i}\), that of \(\alpha\) particle is \(\sqrt{2 m_{\alpha} T_{\alpha}} \hat{j}\) and of \(B e^{9}\) is\\
\(-\sqrt{2 m_{\alpha} T_{\alpha}} \hat{j}+\sqrt{2 m_{n} T} \hat{i}\)\\
By conservation of energy\\
\(T=T_{\alpha}+\frac{m_{\alpha} T_{\alpha}+m_{n} T}{M}+|Q|\)\\
( \(M\) is the mass of \(B e^{9}\) ). Thus\\
\(T_{\alpha}=\left[T\left(1-\frac{m_{n}}{M}\right)-|Q|\right] \frac{M}{M+m_{\alpha}}\).\\
Using\\
\includegraphics[max width=\textwidth, alt={}, center]{da0156b6-4133-434e-90d4-ec01b2070afe-390_239_233_826_1200}\\
we get

\[
\begin{gathered}
T_{t h}=\left(1+\frac{m_{n}}{M}\right)|Q| \\
T_{\alpha}=\frac{M}{M+m_{\alpha}}\left[\left(1-\frac{m_{n}}{M}\right) T-\frac{T_{t h}}{1+\frac{m_{n}}{M}}\right]
\end{gathered}
\]

\(M^{\prime}\) is the mass of \(C^{12}\) nucleus.\\
or

\[
T_{\alpha}=\frac{1}{M+m_{\alpha}}\left[\left(M-m_{n}\right) T-\frac{M M^{\prime}}{M^{\prime}+m_{n}} T_{t h}\right]=2.21 \mathrm{MeV}
\]

6.278 The formula of problem 6.271 does not apply here because the photon is always reletivistic. At threshold, the energy of the photon \(E_{\gamma}\) implies a momentum \(\frac{E_{\gamma}}{c}\). The velocity of centre of mass with respect to the rest frame of initial \(H^{2}\) is

\[
\frac{E_{\gamma}}{\left(m_{n}+m_{p}\right) c}
\]

Since both \(n \& p\) are at rest in CM frame at threshold, we write

\[
E_{\gamma}=\frac{E_{\gamma}^{2}}{2\left(m_{n}+m_{p}\right) c^{2}}+E_{b}
\]

by conservation of energy. Since the first term is a small correction, we have

Thus

\[
\begin{gathered}
E_{\gamma} \cong E_{b}+\frac{E_{b}^{2}}{2\left(m_{n}+m_{p}\right) c^{2}} \\
\frac{\delta E}{E_{b}}=\frac{E_{b}}{2\left(m_{n}+m_{p}\right) c^{2}}=\frac{2.2}{2 \times 2 \times 938}=5.9 \times 10^{-4}
\end{gathered}
\]

or nearly \(0.06 \%\).\\
6.279 The reaction is

\[
p+d \rightarrow H e^{3}
\]

Excitation energy of \(\mathrm{He}^{3}\) is just the energy available in centre of mass. The velocity of the centre of mass is

\[
\frac{\sqrt{2 m_{p} T_{p}}}{m_{p}+m_{d}} \approx \frac{1}{3} \sqrt{\frac{2 T_{p}}{m_{p}}}
\]

In the CM frame, the kinetic energy available is ( \(m_{d} \approx 2 m_{p}\) )

\[
\frac{1}{2} m_{p}\left(\frac{2}{3} \sqrt{\frac{2 T_{p}}{m_{p}}}\right)^{2}+\frac{1}{2} 2 m_{p}\left(\frac{1}{3} \sqrt{\frac{2 T}{m_{p}}}\right)^{2}=\frac{2 T}{3}
\]

The total energy available is then \(Q+\frac{2 T}{3}\)\\
where

\[
\begin{aligned}
Q & =c^{2}\left(\Delta_{n}+\Delta_{d}-\Delta_{H e}^{3}\right) \\
& =c^{2} \times(0.00783+0.01410-0.01603) \mathrm{amu} \\
& =5.49 \mathrm{MeV}
\end{aligned}
\]

Finally

\[
E=6.49 \mathrm{MeV}
\]

6.280 The reaction is

\[
d+C^{13} \rightarrow N^{15 *} \rightarrow n+N^{14}
\]

Maxima of yields determine the energy levels of \(N^{15}\). As in the previous problem the excitation energy is

\[
E_{e x c}=Q+E_{K}
\]

where \(E_{K}=\) available kinetic energy. This is found is as in the previous problem. The velocity of the centre of mass is

\[
\frac{\sqrt{2 m_{d} T_{i}}}{m_{d}+m_{c}}=\frac{m_{d}}{m_{d}+m_{c}} \sqrt{\frac{2 T_{i}}{m_{d}}}
\]

So

\[
E_{K}=\frac{1}{2} m_{d}\left(1-\frac{m_{d}}{m_{d}+m_{c}}\right)^{2} \frac{2 T_{i}}{m_{d}}+\frac{1}{2} m_{c}\left(\frac{m_{d}}{m_{d}+m_{c}}\right)^{2} \frac{2 T_{i}}{m_{d}}=\frac{m_{c}}{m_{d}+m_{c}} T_{i}
\]

\(Q\) is the \(Q\) value for the ground state of \(N^{15}\) : We have

\[
\begin{aligned}
Q & =c^{2} \times\left(\Delta_{d}+\Delta_{C}^{13}-\Delta_{N}^{15}\right) \\
& =c^{2} \times(0.01410+0.00335-0.00011) \mathrm{amu} \\
& =16.14 \mathrm{MeV}
\end{aligned}
\]

The excitation energies then are

\[
\begin{gathered}
16 \cdot 66 \mathrm{MeV}, 16 \cdot 92 \mathrm{MeV} \\
17 \cdot 49 \mathrm{MeV} \text { and } 17 \cdot 70 \mathrm{MeV} .
\end{gathered}
\]

6.281 We have the relation

\[
\frac{1}{\eta}=e^{-n \sigma d}
\]

Here \(\frac{1}{\eta}=\) attenuation factor\\
\(n=\) no. of Cd nuclei per unit volume\\
\(\boldsymbol{\sigma}=\) effective cross section\\
\(d=\) thickness of the plate\\
Now

\[
n=\frac{\rho N_{A}}{M}
\]

( \(\rho=\) density, \(M=\) Molar weight of \(\mathrm{Cd}, N_{A}=\) Avogadro number.)\\
Thus

\[
\sigma=\frac{M}{\rho N_{A} d} \ln \eta=2.53 \mathrm{~kb}
\]

6.282 Here

\[
\frac{1}{\eta}=e^{-\left(n_{2} \sigma_{2}+n_{1} \sigma_{1}\right) d}
\]

where 1 refers to \(O^{1}\) and 2 to \(D\) nuclei\\
Using \(n_{2}=2 n, n_{1}=n=\) concentration of \(O\) nuclei in heavy water we get

\[
\frac{1}{\eta}=e^{-\left(2 \sigma_{2}+\sigma_{1}\right) n d}
\]

Now using the data for heavy water

\[
n=\frac{1.1 \times 6.023 \times 10^{23}}{20}=3.313 \times 10^{22} \text { per cc }
\]

Thus substituting the values

\[
\eta=20 \cdot 4=\frac{I_{0}}{I} .
\]

6.283 In traversing a distance \(d\) the fraction which is either scattered or absorbed is clearly

\[
1-e^{-n\left(\sigma_{s}+\sigma_{a}\right) d}
\]

by the usual definition of the attenuation factor. Of this, the fraction scattered is (by definition of scattering and absorption cross section)

\[
w=\left\{1-e^{-n\left(\sigma_{s}+\sigma_{a}\right) d}\right\} \frac{\sigma_{s}}{\sigma_{s}+\sigma_{a}}
\]

In iron

\[
n=\frac{\rho \times N_{A}}{M}=8.39 \times 10^{22} \text { per cc }
\]

Substitution gives

\[
w=0.352
\]

6.284 (a) Assuming of course, that each reaction produces a radio nuclide of the same type, the decay constant \(\alpha\) of the radionuclide is \(k / w\). Hence \(T=\frac{\ln 2}{\lambda}=\frac{w}{k} \ln 2\)\\
(b) number of bombarding particles is : \(\frac{I t}{e}\)\\
\(\left(e=\right.\) charge on proton). Then the number of \(B e^{7}\) produced is : \(\frac{I t}{e} w\)\\
If \(\lambda=\) decay constant of \(B e^{7}=\frac{\ln 2}{T}\), then the activity is \(A=\frac{I t}{e} w \cdot \frac{\ln 2}{T}\)\\
Hence

\[
w=\frac{e A T}{I t \ln 2}=1.98 \times 10^{-3}
\]

6.285 (a) Suppose \(N_{0}=\) no. of \(A u^{197}\) nuclei in the foil. Then the number of \(A u^{197}\) nuclei transformed in time \(t\) is

\[
N_{0} \cdot J \cdot \sigma \cdot t
\]

For this to equal \(\eta N_{0}\), we must have

\[
t=\eta /(J \cdot \sigma)=323 \text { years }
\]

(b) Rate of formation of the \(A u^{198}\) nuclei is \(N_{0} \cdot J \cdot \sigma\) per sec\\
and rate of decay is \(\lambda n\), where \(n\) is the number of \(A u^{198}\) at any instant.\\
Thus

\[
\frac{d n}{d t}=n_{0} \cdot J \cdot \sigma-\lambda n
\]

The maximum number of \(A u^{198}\) is clearly

\[
n_{\max }=\frac{N_{0} \cdot J \cdot \sigma}{\lambda}=\frac{N_{0} \cdot J \cdot \sigma \cdot T}{\ln 2}
\]

because if \(n\) is smaller, \(\frac{d n}{d t}>0\) and \(n\) will increase further and if \(n\) is larger \(\frac{d n}{d t}<0\) and \(n\) will decrease. (Actually \(n_{\text {max }}\) is approached steadily as \(t \rightarrow \infty\) )\\
Substitution gives using \(N_{0}=3.057 \times 10^{19}, n_{\max }=1.01 \times 10^{13}\)\\
6.286 Rate of formation of the radionuclide is \(n \cdot J \cdot \sigma\) per unit area per -sec. Rate of decay is \(\lambda N\).

Thus

\[
\frac{d N}{d t}=n \cdot J \cdot \sigma-\lambda N \text { per unit area per second }
\]

Then

\[
\left(\frac{d N}{d t}+\lambda N\right) e^{\lambda t}=n \cdot J \cdot \sigma e^{\lambda t} \text { or } \frac{d}{d t}\left(N e^{\lambda t}\right)=n \cdot J \cdot \sigma \cdot e^{\lambda t}
\]

Hence

\[
N e^{\lambda t}=\text { Const }+\frac{n \cdot J \cdot \sigma}{\lambda} e^{\lambda t}
\]

The number of radionuclide at \(t=0\) when the process starts is zero. So constant \(=-\frac{n \cdot J \cdot \sigma}{\lambda}\)\\
Then

\[
N=\frac{n \cdot J \cdot \sigma}{\lambda}\left(1-e^{-\lambda t}\right)
\]

6.287 We apply the formula of the previous problem except that have\\
\(N=\) no. of radio nuclide and no. of host nuclei originally.\\
Here

\[
n=\frac{0.2}{197} \times 6.023 \times 10^{23}=6.115 \times 10^{20}
\]

Then after time \(\boldsymbol{t}\)

\[
N=\frac{n \cdot J \cdot \sigma \cdot T}{\ln 2}\left(1-e^{-\frac{t \ln 2}{T}}\right)
\]

\(T=\) half life of the radionuclide.\\
After the source of neutrons is cut off the activity after time \(T\) will be

Thus

\[
A=\frac{n \cdot J \cdot \sigma \cdot T}{\ln 2}\left(1-e^{-t \ln 2 / T}\right) e^{-\tau \ln 2 / T \times \frac{\ln 2}{T}}=n \cdot J \cdot \sigma\left(1-e^{-t \ln 2 / T}\right) e^{-\tau \ln 2 / T}
\]

\[
J=A e^{\tau \ln 2 / T} / n \sigma\left(1-e^{-t \ln 2 / T}\right)=5.92 \times 10^{9} \mathrm{part} / \mathrm{cm}^{2} \cdot \mathrm{~s}
\]

6.288 No. of nuclei in the first generation \(=\) No. of nuclei initially \(=N_{0}\)\\
\(N_{0}\) in the second generation \(=N_{0} \times\) multiplication factor \(=N_{0} \cdot k\)\\
\(N_{0}\) in the the 3 rd generation \(=N_{0} \cdot k \cdot k=N_{0} k^{2}\)\\
\(N_{0}\) in the nth generation \(=N_{0} k^{n-1}\)\\
Substitution gives \(1.25 \times 10^{5}\) neutrons\\
\(6.289 N_{0}\). of fissions per unit time is clearly \(P / E\). Hence no. of neutrons produced per unit time to \(\frac{v P}{E}\). Substitution gives \(7.80 \times 10^{18}\) neutrons \(/ \mathrm{sec}\)\\
6.290 (a) This number is \(k^{n-1}\) where \(n=\) no. of generations in time \(t=t / T\) Substitution gives 388.\\
(b) We write

\[
k^{n-1}=e^{\left(\frac{T}{\tau}-1\right) \ln k}=e
\]

or

\[
\frac{T}{\tau}-1=\frac{1}{\ln k} \quad \text { and } \quad T=\tau\left(1+\frac{1}{\ln k}\right)=10.15 \mathrm{sec}
\]

\subsection*{6.7 ELEMENTARY PARTICLES}
6.291 The formula is

\[
T=\sqrt{c^{2} p^{2}+m_{0}^{2} c^{4}}-m_{0} c^{2}
\]

Thus

\[
\begin{aligned}
T=5.3 \mathrm{MeV} & \text { for } p=0.10 \frac{\mathrm{GeV}}{\mathrm{c}}=5.3 \times 10^{-3} \mathrm{GeV} \\
T & =0.433 \mathrm{GeV} \text { for } p=1.0 \frac{\mathrm{GeV}}{\mathrm{c}} \\
T & =9.106 \mathrm{GeV} \text { for } p=10 \frac{\mathrm{GeV}}{\mathrm{c}}
\end{aligned}
\]

Here we have used \(m_{0} c^{2}=0.938 \mathrm{GeV}\)\\
6.292 Energy of pions is \((1+\eta) m_{0} c^{2}\) so

Hence

\[
\begin{gathered}
(1+\eta) m_{0} c^{2}=\frac{m_{0} c^{2}}{\sqrt{1-\beta^{2}}} \\
\frac{1}{\sqrt{1-\beta^{2}}}=1+\eta \text { or } \beta=\frac{\sqrt{\eta(2+\eta)}}{1+\eta}
\end{gathered}
\]

Here \(\beta=\frac{v}{c}\) of pion. Hence time dilation factor is \(1+\eta\) and the distance traversed by the pion in its lifetime will be

\[
\frac{c \beta \tau_{0}}{\sqrt{1-\beta^{2}}}=c \tau_{0} \sqrt{\eta(2+\eta)}=15.0 \text { metres }
\]

on substituting the values of various quantities. (Note. The factor \(\frac{1}{\sqrt{1-\beta^{2}}}\) can be looked at as a time dilation effect in the laboratory frame or as length contraction factor brought to the other side in the proper frame of the pion).\\
6.293 From the previous problem \(l=c \tau_{0} \sqrt{\eta(\eta+2)}\)\\
where \(\eta=\frac{T}{m_{\pi} c^{2}}, m_{\pi}\) is the rest mass of pions.\\
substitution gives

\[
\begin{aligned}
\tau_{0} & =\frac{l}{c \sqrt{\eta(2+\eta)}}=2.63 \mathrm{~ns} \\
& =\frac{l m_{\pi} c}{\sqrt{T\left(T+2 m_{\pi} c^{2}\right)}}
\end{aligned}
\]

where we have used \(\eta=\frac{100}{139.6}=0.716\)\\
6.294 Here \(\eta=\frac{T}{m c^{2}}=1\) so the life time of the pion in the laboratory frame is

\[
\eta=(1+\eta) \tau_{0}=2 \tau_{0}
\]

The law of radioactive decay implies that the flux decrease by the factor.

\[
\begin{aligned}
\frac{J}{J_{0}} & =e^{-t / \tau}=e^{-l / v \tau}=e^{-l / c \tau_{0} \sqrt{\eta(2+\eta)}} \\
& =\exp \left(-\frac{m c l}{\tau_{0} \sqrt{T\left(T+2 m c^{2}\right)}}\right)=0.221
\end{aligned}
\]

6.295 Energy-momentum conservation implies

\[
O=\overrightarrow{p_{\mu}}+\overrightarrow{p_{v}}
\]

\[
m_{\pi} c^{2}=E_{\mu}+E_{v} \quad \text { or } \quad m_{\pi} c^{2}-E_{v}=E_{\mu}
\]

But

\[
E_{v}=c\left|\overrightarrow{p_{v}}\right|=c\left|p_{\mu}\right| . \text { Thus }
\]

\[
m_{\pi}^{2} c^{4}-2 m_{\pi} c^{2} \cdot c\left|\overrightarrow{p_{\mu}}\right|+c^{2} p_{\mu}^{2}=E_{\mu}^{2}=c^{2} p_{\mu}^{2}+m_{\mu}^{2} c^{4}
\]

Hence

\[
c\left|\overrightarrow{p_{\mu}}\right|=\frac{m_{\pi}^{2}-m_{\mu}^{2}}{2 m_{\pi}} \cdot c^{2}
\]

So

\[
T_{\mu}=\sqrt{c^{2} p_{\mu}^{2}+m_{\mu}^{2} c^{4}}-m_{\mu} c^{2}=\sqrt{\frac{\left(m_{\pi}^{2}-m_{\mu}^{2}\right)^{2}}{4 m_{\pi}^{2}}+m_{\mu}^{2}} \cdot c^{2}-m_{\mu} c^{2}
\]

\[
=\frac{m_{\pi}^{2}+m_{\mu}^{2}}{2 m_{\pi}} c^{2}-m_{\mu} c^{2}=\frac{\left(m_{\pi}-m_{\mu}\right)^{2}}{2 m_{\pi}} \cdot c^{2}
\]

Substituting

\[
\begin{gathered}
m_{\pi} c^{2}=139.6 \mathrm{MeV} \\
m_{\mu} c^{2}=105.7 \mathrm{MeV} \text { we get } \\
T_{\mu}=4.12 \mathrm{MeV} \\
E_{v}=\frac{m_{\pi}^{2}-m_{\mu}^{2}}{2 m_{\pi}} c^{2}=29.8 \mathrm{MeV}
\end{gathered}
\]

Also\\
6.296 We have\\
or

\[
\left(m_{\Sigma} c^{2}-E_{n}\right)^{2}=E_{\pi}^{2}
\]

or

\[
m_{\Sigma}^{2} c^{4}-2 m_{\Sigma} c^{2} E_{n}=E_{\pi}^{2}-E_{n}^{2}=c^{4} m_{\pi}^{2}-c^{4} m_{n}^{2}
\]

because (1) implies

\[
E_{\pi}^{2}-E_{n}^{2}=m_{\pi}^{2} c^{4}-m_{n}^{2} c^{4}
\]

Hence

\[
E_{n}=\frac{m_{\Sigma}^{2}+m_{n}^{2}-m_{\pi}^{2}}{2 m_{\Sigma}} c^{2}
\]

and

\[
\begin{gathered}
T_{n}=\left(\frac{m_{\Sigma}^{2}+m_{n}^{\dot{2}}-m_{\pi}^{2}}{2 m_{\Sigma}}-m_{n}\right) c^{2}=\frac{\left(m_{\Sigma}-m_{n}\right)^{2}-m_{\pi}^{2}}{2 m_{\Sigma}} c^{2} . \\
T_{n}=19.55 \mathrm{MeV}
\end{gathered}
\]

Substitution gives

\[
\mu^{+} \rightarrow e^{+}+\bar{v}_{e}+\bar{v}_{\mu}
\]

The neutrinoes are massless. The positron will carry largest momentum if both neutriones ( \(v_{e} \& \bar{v}_{\mu}\) ) move in the same direction in the rest frame of the nuon. Then the final product is effectively a two body system and we get from problem (295)\\
6.297 The reaction is

\[
\left(T_{e}^{+}\right)_{\max }=\frac{\left(m_{\mu}-m_{e}\right)^{2}}{2 m_{\mu}} c^{2}
\]

Substitution gives

\[
\left(T_{e}^{+}\right)_{\max }=52.35 \mathrm{MeV}
\]

6.298 By conservation of energy-momentum

Then

\[
\begin{aligned}
M c^{2} & =E_{p}+E_{\pi} \\
O & =\overrightarrow{p_{p}}+\overrightarrow{p_{\pi}} \\
m_{\pi}^{2} c^{4} & =E_{\pi}^{2}-\overrightarrow{p_{\pi}^{2}} c^{2}=\left(M c^{2}-E_{p}\right)^{2}-c^{2} \overrightarrow{p_{p}^{2}} \\
& =M^{2} c^{4}-2 M c^{2} E_{p}+m_{p}^{2} c^{4}
\end{aligned}
\]

This is a quadratic equation in M

\[
M^{2}-2 \frac{E_{p}}{c^{2}} M+m_{p}^{2}-m_{\pi}^{2}=0
\]

or using \(E_{p}=m_{p} c^{2}+T\) and solving

\[
\left(M-\frac{E_{p}}{c^{2}}\right)^{2}=\frac{E_{p}^{2}}{c^{4}}-m_{p}^{2}+m_{\pi}^{2}
\]

Hence,

\[
M=\frac{E_{p}}{c^{2}}+\sqrt{\frac{E_{p}^{2}}{c^{4}}-m_{p}^{2}+m_{\pi}^{2}}
\]

taking the positive sign. Thus

\[
M=m_{p}+\frac{T}{c^{2}}+\sqrt{m_{\pi}^{2}+\frac{T}{c^{2}}\left(2 m_{p}+\frac{T}{c^{2}}\right)}
\]

Substitution gives

\[
M=1115 \cdot 4 \frac{\mathrm{MeV}}{\mathrm{c}^{2}}
\]

From the table of masses we identify the particle as a \(\wedge\) particle\\
\includegraphics[max width=\textwidth, alt={}, center]{da0156b6-4133-434e-90d4-ec01b2070afe-398_365_708_118_351}

See the diagram. By conservation of energy\\
or\\
or

\[
\begin{gathered}
\sqrt{m_{\pi}^{2} c^{4}+c^{2} p_{\pi}^{2}}=c p_{v}+\sqrt{m_{\mu}^{2} c^{4}+\dot{p}_{\pi}^{2} c^{2}+c^{2} p_{v}^{2}} \\
\left(\sqrt{m_{\pi}^{2} c^{4}+c^{2} p_{\pi}^{2}}-c p_{v}\right)^{2}=m_{\mu}^{2} c^{4}+c^{2} p_{\pi}^{2}+c^{2} p_{v}^{2} \\
m_{\pi}^{2} c^{4}-2 c p_{v} \sqrt{m_{\pi}^{2} c^{4}+c^{2} p_{\pi}^{2}}=m_{\mu}^{2} c^{4}
\end{gathered}
\]

Hence the energy of the neutrino is\\
on writing

\[
\begin{gathered}
E_{\mathrm{v}}=c p_{\mathrm{v}}=\frac{m_{\pi}^{2} c^{4}-m_{\mu}^{2} c^{4}}{2\left(m_{\pi} c^{2}+T\right)} \\
\sqrt{m_{\pi}^{2} c^{4}+c^{2} p_{\pi}^{2}}=m_{\pi} c^{2}+T \\
E_{\mathrm{v}}=21.93 \mathrm{MeV}
\end{gathered}
\]

Substitution gives\\
\includegraphics[max width=\textwidth, alt={}, center]{da0156b6-4133-434e-90d4-ec01b2070afe-398_383_755_1160_371}

By energy conservation\\
or

\[
\begin{gathered}
\sqrt{m_{\Sigma}^{2} c^{4}+c^{2} p_{\Sigma}^{2}}=\sqrt{m_{\pi}^{2} c^{4}+c^{2} p_{\pi}^{2}}+\sqrt{m_{n}^{2} c^{4}+c^{2} p_{\pi}^{2}+c^{2} p_{\Sigma}^{2}} \\
\left(\sqrt{m_{\Sigma}^{2} c^{4}+c^{2} p_{\Sigma}^{2}}-\sqrt{m_{\pi}^{2} c^{4}+c^{2} p_{\pi}^{2}}\right)^{2}=m_{n}^{2} c^{4}+c^{2} p_{\pi}^{2}+c^{2} p_{\Sigma}^{2} \\
m_{\Sigma}^{2} c^{4}+c^{2} p_{\Sigma}^{2}+m_{\pi}^{2} c^{4}+c^{2} p_{\pi}^{2}-2 \sqrt{m_{\pi}^{2} c^{4}+c^{2} p_{\Sigma}^{2}} \sqrt{m_{\pi}^{2} c^{2}+c^{2} p_{\pi}^{2}} \\
=m_{n}^{2} c^{2}+c^{2} p_{\pi}^{2}+c^{2} p_{\Sigma}^{2}
\end{gathered}
\]

or\\
or using the K.E. of \(\Sigma \& \pi\)\\
and

\[
\begin{gathered}
m_{n}^{2}=m_{\Sigma}^{2}+m_{\pi}^{2}-2\left(m_{\Sigma}+\frac{T_{\Sigma}}{c^{2}}\right)\left(m_{\pi}+\frac{T_{\pi}}{c^{2}}\right) \\
m_{n}=\sqrt{m_{\Sigma}^{2}+m_{\pi}^{2}-2\left(m_{\Sigma}+\frac{T_{\Sigma}}{c^{2}}\right)\left(m_{\pi}+\frac{T_{\pi}}{c^{2}}\right)}=0.949 \frac{\mathrm{GeV}}{c^{2}}
\end{gathered}
\]

6.301 Here by conservation of momentum\\
\(p_{\pi}=2 \times \frac{E_{\pi}}{2 c} \times \cos \frac{\theta}{2}\)\\
or \(\quad c p_{\pi}=E_{\pi} \cos \frac{\theta}{2}\)\\
Thus \(\quad E_{\pi}^{2} \cos ^{2} \frac{\theta}{2}=E_{n}^{2}-m_{\pi}^{2} c^{4}\)\\
or\\
\includegraphics[max width=\textwidth, alt={}, center]{da0156b6-4133-434e-90d4-ec01b2070afe-399_319_552_416_684}\\
and

\[
T_{\pi}=m_{\pi} c^{2}\left(\operatorname{cosec} \frac{\theta}{2}-1\right)
\]

substitution gives \(T_{\pi}=m_{\pi} c^{2}=135 \mathrm{MeV}\) for \(\theta=60^{\circ}\).

Also

\[
\begin{aligned}
E_{y} & =\frac{m_{\pi} c^{2}+T_{\pi}}{2}=\frac{m_{\pi} c^{2}}{2} \operatorname{cosec} \frac{\theta}{2} \\
& =m_{\pi} c^{2} \text { in this case }\left(\theta=60^{\circ}\right)
\end{aligned}
\]

6.302 With particle masses standing for the names of the particles, the reaction is

\[
m+M \rightarrow m_{1}+m_{2}+\ldots
\]

On R.H.S. let the energy momenta be \(\left(E_{1}, c \overrightarrow{p_{1}}\right),\left(E_{2}, c \overrightarrow{p_{2}}\right)\) etc. On the left the energy momentum of the particle \(m\) is \((E, c \vec{p})\) and that of the other particle is \(\left(M c^{2}, \vec{O}\right)\), where, ofcourse, the usual relations

\[
E^{2}-c^{2} \vec{p}=m^{2} c^{4} \text { etc }
\]

hold. From the conservation of energy momentum we see that

\[
\left(E+M c^{2}\right)^{2}-c^{2} \vec{p}^{2}=\left(\Sigma E_{i}\right)^{2}-\left(\Sigma c \overrightarrow{p_{i}}\right)^{2}
\]

Left hand side is

\[
m^{2} c^{4}+M^{2} c^{4}+2 M c^{2} E
\]

We evaluate the R.H.S. in the frame where \(\Sigma \overrightarrow{p_{i}}=0\) (CM frame of the decay product).\\
Then

\[
\text { R.H.S. }=\left(\Sigma E_{i}\right)^{2} \geq\left(\Sigma m_{i} c^{2}\right)^{2}
\]

because all energies are +ve. Therefore we have the result

\[
E \geq \frac{\left(\Sigma m_{i}\right)^{2}-m^{2}-M^{2}}{2 M} c^{2}
\]

or since \(E=m c^{2}+T\), we see that \(T \geq T_{t h}\) where

\[
T_{t h}=\frac{\left(\Sigma m_{i}\right)^{2}-(m+M)^{2}}{2 M} c^{2}
\]

6.303 By momentum conservation

\[
\begin{aligned}
& \sqrt{E^{2}-m_{e}^{2} c^{4}}=2 \frac{E+m_{e} c^{2}}{2} \cos \frac{\theta}{2} \\
& \text { or } \cos _{2}^{\theta}=\sqrt{\frac{E-m_{e} c^{2}}{E+m_{e} c^{2}}}=\sqrt{\frac{T}{T+2 m_{e} c^{2}}} \cdots \frac{E+M c^{2}}{2} \\
& \text { Substitution gives } \\
& \theta=98.8^{\circ}
\end{aligned}
\]

6.304 The formula of problem 3.02 gives

\[
E_{t h}=\frac{\left(\Sigma m_{i}\right)^{2}-M^{2}}{2 M} c^{2}
\]

when the projectile is a photon\\
(a) For \(\gamma+e^{-} \rightarrow e^{-}+e^{-}+e^{+}\)

\[
E_{t h}=\frac{9 m_{e}^{2}-m_{e}^{2}}{2 m_{e}} c^{2}=4 m_{e} c^{2}=2.04 \mathrm{MeV}
\]

(b) For

\[
\gamma+p \rightarrow p+\pi^{+} \pi^{-}
\]

\[
E_{t h}=\frac{\left(M_{p}+2 m_{\pi}\right)^{2}-M_{p}^{2}}{2 M_{p}} c^{2}=\frac{4 m_{\pi} M_{p}+4 m_{\pi}^{2}}{2 M_{p}} c^{2}=2\left(m_{\pi}+\frac{m_{\pi}^{2}}{M_{p}}\right) c^{2}=320.8 \mathrm{MeV}
\]

6.305 (a) For \(p+p \rightarrow p+p+p+\bar{p}\)

\[
T \geq T_{t h}=\frac{16 m_{p}^{2}-4 m_{p}^{2}}{2 m_{p}} c^{2}=6 m_{p} c^{2}=5.63 \mathrm{GeV}
\]

(b) For \(p+p \rightarrow p+p+\pi^{\circ}\)

\[
\begin{aligned}
T \geq T_{t h} & =\frac{\left(2 m_{p}+m_{\pi^{\circ}}\right)^{2}-4 m_{p}^{2}}{2 m_{p}} c^{2} \\
& =\left(2 m_{\pi}+\frac{m_{\pi^{\circ}}^{2}}{2 m_{p}}\right) c^{2}=0.280 \mathrm{GeV}
\end{aligned}
\]

6.306 (a) Here

\[
T_{t h}=\frac{\left(m_{K}+m_{\Sigma}\right)^{2}-\left(m_{\pi}+m_{p}\right)^{2}}{2 m_{p}} c^{2}
\]

Substitution gives \(\boldsymbol{T}_{\text {th }}=0.904 \mathrm{GeV}\)\\
(b) \(T_{t h}=\frac{\left(m_{K^{+}}+m_{\lambda}\right)^{2}-\left(m_{\pi^{0}}+m_{p}\right)^{2}}{2 m_{p}} c^{2}\)

Substitution gives \(\boldsymbol{T}_{\text {th }}=0.77 \mathrm{GeV}\).\\
6.307 From the Gell-Mann Nishijima formula

\[
Q=T_{Z}+\frac{Y}{2}
\]

we get

\[
O=\frac{1}{2}+\frac{Y}{2} \text { or } Y=-1
\]

Also

\[
Y=B+S \Rightarrow S=-2 \text {. Thus the particle is } m^{\circ} 0 \text {. }
\]

6.308 (1) The process \(n \rightarrow p+e^{-}+v_{e}\) cannot occur as there are 2 more leptons ( \(e^{-}, v_{e}\) ) on the right comopared to zero on the left.\\
(2) The process \(\pi^{+} \rightarrow \mu^{+}+e^{-}+e^{+}\)is forbidden because this corresponds to a change of lepton number by, ( 0 on the left -1 on the right)\\
(3) The process \(\pi^{-} \rightarrow \mu^{-}+v_{\mu}\) is forbidden because \(\mu^{-}, v_{\mu}\) being both leptons \(\Delta L=2\) hre.\\
(4), (5), (6) are allowed (except that one must distinguish between muon neutrinoes and electron neutrinoes). The correct names would be\\
(4) \(p+e^{-} \rightarrow n+v_{e}\)\\
(5) \(\mu^{+} \rightarrow e^{+}+v_{e}+\tilde{v}_{\mu}\)\\
(6) \(K^{-} \rightarrow \mu^{-}+\tilde{v}_{\mu}\).\\
6.309 (1) \(\pi^{-}+p \rightarrow \Sigma^{-}+K^{+}\)\\
so

\[
\begin{array}{lccr}
0 \quad 0 & -1 & 1 \\
\Delta S & =0 . \text { allowed }
\end{array}
\]

(2) \(\pi^{-}+p \rightarrow \Sigma^{+}+K^{-}\)

\[
0.00-11
\]

so

\[
\Delta S=-2 . \text { forbidden }
\]

(3) \(\pi^{-}+p \rightarrow K^{-}+K^{+}+n\)

\[
0 \quad 0 \rightarrow-1 \quad 1 \quad 0
\]

\(S=0\), allowed.\\
(4) \(n+p \rightarrow \Lambda^{\circ}+\Sigma^{+}\)\\
so

\[
\begin{array}{llll}
0 & 0 & -1 & -1
\end{array}
\]

\[
\Delta S=-2 . \text { forbidden }
\]

(5) \(\pi^{-}+n \rightarrow m^{-}+K^{+}+K^{-}\)

\[
0 \quad 0 \rightarrow-2 \quad 1 \quad-1
\]

so

\[
\Delta S=-2 . \text { forbidden } .
\]

(6) \(K^{-}+p \rightarrow \Omega^{-}+K^{+} K^{\circ}\)

\[
\begin{array}{lll}
-1 & 0 & -3+1+1
\end{array}
\]

so

\[
\Delta S=0 . \text { allowed. }
\]

(1) \(\Sigma^{-} \rightarrow \Lambda^{\circ}+\pi^{-}\)\\
is forbidden by energy conservation. The mass difference

\[
M_{\Sigma^{-}}-M_{\Lambda^{\circ}}=82 \frac{\mathrm{MeV}}{c^{2}}<m_{\pi}-
\]

(The process \(1 \rightarrow 2+3\) will be allowed only if \(m_{1}>m_{2}+m_{3}\).)\\
(2) \(\pi^{-}+p \rightarrow K^{+}+K^{-}\)\\
is disallowed by conservation of baryon number.\\
(3) \(K^{-}+n \rightarrow \Omega^{-}+K^{+}+K^{\circ}\)\\
is forbidden by conservation of charge\\
(4) \(n+p \rightarrow \Sigma^{+}+\Lambda^{\circ}\)\\
is forbidden by strangeness conservation.\\
(5) \(\pi^{-} \rightarrow \mu^{-}+e^{-}+e^{+}\)\\
is forbidden by conservation of muon number (or lepton number).\\
(6) \(\mu^{-} \rightarrow e^{-}+v_{e}+\tilde{v}_{\mu}\)\\
is forbidden by the separate conservation of muon number as well as lepton number.


\end{document}